\documentclass[11pt,oneside]{book}
\usepackage[margin=1in]{geometry}
\usepackage{amsmath,amssymb,amsthm,mathtools}
\usepackage{booktabs,longtable,array}
\usepackage{enumitem}
\usepackage[colorlinks=true,linkcolor=blue,urlcolor=blue,citecolor=blue]{hyperref}
\usepackage[nameinlink,noabbrev]{cleveref}

\newtheorem{theorem}{Theorem}[chapter]
\newtheorem{lemma}[theorem]{Lemma}
\newtheorem{proposition}[theorem]{Proposition}
\newtheorem{corollary}[theorem]{Corollary}
\newtheorem{assessment}[theorem]{Assessment}
\newtheorem{firewall}[theorem]{Claim firewall}
\newtheorem{openproblem}[theorem]{Open problem}
\newtheorem{record}[theorem]{Record}
\newtheorem{programme}[theorem]{Programme}
\newtheorem{audit}[theorem]{Audit}
\newtheorem{correction}[theorem]{Correction}
\newtheorem{warning}[theorem]{Warning}
\newtheorem{decision}[theorem]{Decision}
\newtheorem{scope}[theorem]{Scope}
\theoremstyle{definition}
\newtheorem{definition}[theorem]{Definition}
\theoremstyle{remark}
\newtheorem{remark}[theorem]{Remark}

\newcommand{\MGclosed}{\textsc{closed}}
\newcommand{\MGopen}{\textsc{open}}
\newcommand{\MGpartial}{\textsc{partial}}
\newcommand{\Tr}{\operatorname{Tr}}
\newcommand{\dist}{\operatorname{dist}}

\title{Renewal Yang--Mills Mass-Gap Research Notes\\
Seventh Series, Version V1}
\author{Compact proof-indexed continuation after archives f1--f6}
\date{7 September 2026}

\begin{document}
\maketitle
\tableofcontents

\chapter{Context, cumulative results, and the live proof frontier}
\label{ch:s7v1-context}

\section{Purpose and status}
\label{sec:s7v1-purpose}

The ultimate objective is a rigorous construction of four-dimensional
quantum Yang--Mills theory for a compact nonabelian gauge group, followed
by Osterwalder--Schrader reconstruction and a proof that the physical
Hamiltonian has a strictly positive gap above its vacuum.  The present
programme works upstream from that objective.  It develops exact
finite-cutoff identities, conditional spectral-gap compilers, local
Wilson-action geometry, and cofinal estimates that would be needed in a
continuum construction.

Six long source cycles have now been frozen.  The immediately preceding
cycle reached V100 and was archived as \(\texttt{\_f6}\).  This compact V1
does not restart the mathematics.  It gives a proof-indexed summary of the
major established results and corrections, states the smallest live
interface, and then continues the attack at the complete-action boundary
problem identified in f6-V100.

\begin{firewall}[No mass-gap claim]
\label{fire:s7v1-no-claim}
The full Yang--Mills mass gap has not been proved.  In particular, the
notes do not yet provide the complete-action conditional smoothing
estimate, response-ground stability, block tensorization and physical
form incidence, cutoff-uniform Osterwalder--Schrader measure, or continuum
Hamiltonian gap.  Every implication using one of those rows remains
conditional.
\end{firewall}

\begin{record}[Immutable archive ledger]
\label{rec:s7v1-archives}
\begin{center}
\small
\begin{tabular}{p{0.49\textwidth}rr}
\toprule
archive & source lines & bytes\\
\midrule
\texttt{Renewal\_Yang\_Mills\_Mass\_Gap\_Research\_Notes\_V135\_f1.tex}
& \(88\,517\) & \(3\,083\,467\)\\
\texttt{Renewal\_Yang\_Mills\_Mass\_Gap\_Research\_Notes\_V100\_f2.tex}
& \(49\,581\) & \(1\,706\,928\)\\
\texttt{Renewal\_Yang\_Mills\_Mass\_Gap\_Research\_Notes\_V100\_f3.tex}
& \(48\,381\) & \(1\,720\,945\)\\
\texttt{Renewal\_Yang\_Mills\_Mass\_Gap\_Research\_Notes\_V100\_f4.tex}
& \(44\,712\) & \(1\,483\,037\)\\
\texttt{Renewal\_Yang\_Mills\_Mass\_Gap\_Research\_Notes\_V100\_f5.tex}
& \(45\,189\) & \(1\,467\,995\)\\
\texttt{Renewal\_Yang\_Mills\_Mass\_Gap\_Research\_Notes\_V100\_f6.tex}
& \(47\,291\) & \(1\,647\,389\)\\
\bottomrule
\end{tabular}
\end{center}
Their SHA--256 digests, in the same order, are
\begin{align*}
&\texttt{82f6bbffeecdad3c5c63551660b19a0bdc7f3d6dffe4c68b4e9bba3a783a2b71},\\
&\texttt{6444a1044e74dfcb1a72ed9c7248bc0f234efc4a0e4248c736cd7db212dfca9a},\\
&\texttt{a5e05a144c753426d8a461c8290e9538a4cd19d8672bab8b08b509d82faf892f},\\
&\texttt{069cdb5806423cb3a871776eb571b63919feb56d4dd638be694c1991ccb0ae16},\\
&\texttt{8b9965394878a038518c180907033555ceb2494c17f2bcb0be7e46ea7e736857},\\
&\texttt{63a6ee472d2127ff2d928791f1c720acff1dbb5193047c2d3cea14f7db759a4a}.
\end{align*}
Archived labels mentioned below are descriptive proof indices, not local
\LaTeX\ cross-references.
\end{record}

\begin{record}[Version protocol]
\label{rec:s7v1-protocol}
Each successor version copies the complete active source, appends one
proof note, passes exact-prefix, case-sensitive label/reference,
environment-stack, and terminator checks, and then replaces its active
predecessor.  Every version divisible by five audits the newest active
Standard-Model--Einstein and Navier--Stokes files read-only.  On reaching
V100, the active file is frozen under the next suffix
\(\texttt{\_fX}\), and a new contextual V1 begins.
\end{record}

\section{Major mathematical results already obtained}
\label{sec:s7v1-summary}

\begin{record}[Compact-group and one-link foundation]
\label{rec:s7v1-foundation}
Archives f1--f5 establish the following exact finite-cutoff tools:
\begin{enumerate}[label=\textup{(\roman*)},leftmargin=3.5em]
\item Peter--Weyl output pinching and representation-normalized compact
      group heat-kernel identities;
\item dimension-free heat-increment moments and exact
      Duhamel/interpolation formulae;
\item endpoint-factorizing BKAR forest representations;
\item reserve-weighted Cayley--Pr\"ufer, rooted-forest, and
      maximum-spanning-tree identities;
\item connected hypergraph and bounded-excess primitive-core estimates
      with explicit activity bookkeeping; and
\item strong one-link connected estimates on the parameter ranges stated
      in the archived ledgers.
\end{enumerate}
These results are combinatorial and analytic inputs.  They do not by
themselves construct a continuum measure or a Hamiltonian gap.
\end{record}

\begin{record}[Exact conditional and ground-transform identities]
\label{rec:s7v1-conditional}
The later archives prove the following operator statements.
\begin{enumerate}[label=\textup{(\roman*)},leftmargin=3.5em]
\item If \(E_i\) is conditional expectation and \(D_i=I-E_i\), then for
      a differentiating field \(X_a\),
      \[
       X_aD_if
       =D_i(X_af)-\operatorname{Cov}_i
          \bigl(f,X_a\log\rho_i\bigr).
      \]
      The score covariance is mandatory when the conditional density
      moves.
\item For the positive ground transform \({\cal U}\) and the associated
      projection \(Q\), the physical projection identities are
      \[
       {\cal U}P{\cal U}^{-1}=s_1^{-2}Q,\qquad
       {\cal U}(PM_BP){\cal U}^{-1}=s_1^{-4}QM_BQ.
      \]
\item A direct integral of fibrewise positive self-adjoint Markov
      transfers on the same conditional block obeys
      \(EP=PE=E\), hence \([I-E,P]=0\).
\item This commutation gives the projected smoothing estimate
      \[
       \|M_BD F\|_2^2
       \le (1+\theta)\eta_B\|DF\|_2^2
       +(1+\theta^{-1})\langle F,(I-P)F\rangle .
      \]
\end{enumerate}
The proofs occur in f6-V91--V94, with the independently derived
Hilbert-valued form audited in SM V35.  Their scope is a fixed positive
carrier and a block containing the declared event.
\end{record}

\begin{record}[Corrections and no-go conclusions that govern the route]
\label{rec:s7v1-no-go}
The archive chain contains several negative results which must remain
active.
\begin{enumerate}[label=\textup{(\roman*)},leftmargin=3.5em]
\item Small unconditional bad-event probability does not imply a uniform
      conditional probability or a restricted operator norm.
\item A conditional Fisher-information integral does not imply an
      essential-supremum influence bound.
\item One-link Wilson fibres can have exterior-induced competing centre
      wells, so a uniform one-link contraction cannot be assumed at weak
      coupling.
\item Tree gauge on an isolated star changes links in a pointwise-fixed
      exterior; it is a coordinate transport between fibres, not an
      internal symmetry of one fibre.
\item Reflection positivity, static correlation estimates, or a
      Euclidean area estimate alone do not identify the required
      finite-volume Hamiltonian gap.
\item A pure unitary tangent may be cancelled exactly, but an actual
      physical deformation cannot be relabelled as gauge transport.
\end{enumerate}
These statements explain why the active route uses enlarged conditional
blocks, explicit bad-set smoothing, and separately typed physical
response rows.
\end{record}

\begin{record}[Local Wilson concentration and convexity]
\label{rec:s7v1-local-Wilson}
For the six staples \(S_1,\ldots,S_6\), let
\[
 \mathfrak F
 :=\min_{V\in SU(3)}
   \sum_{j=1}^6\bigl(3-\Re\Tr(VS_j)\bigr).
\]
The alignment analysis in f6-V95 proves that, below an explicit
dimension-only threshold, the matrix-Fisher phase has a unique global
maximizer and a quadratic well.  If
\(\epsilon=\sqrt{12\mathfrak F}\), then after alignment
\[
 \left\|\sum_{j=1}^6S_j-6I\right\|_{\mathrm F}\le\epsilon,
\]
and the loss from the maximizer dominates a fixed multiple of squared
group distance.

For the isolated \(19\)-link star, f6-V96 gives a \(14\)-vertex spanning
tree with six chord holonomies and an exact common/relative variance
decomposition.  Its later boundary-gauge correction is essential: this
isolated quotient is not the full pointwise-fixed physical fibre.

For the star defect
\[
 \mathfrak D_e
 =\sum_{j=1}^6(3-\Re\Tr S_j)
 ={1\over2}\sum_{j=1}^6\|S_j-I\|_{\mathrm F}^2,
\]
f6-V97 proves \(\mathfrak F_e\le\mathfrak D_e\), a bare Wilson action
margin \(s a\) on \(\{\mathfrak D_e\ge a\}\), strong convexity of a
small exponential-chart core, and a core/bad-set gluing inequality.  The
resulting conditional Poincar\'e compiler is valid once its restricted
smoothing and physical incidence hypotheses are supplied.
\end{record}

\begin{record}[Exact collar topology and reconstruction]
\label{rec:s7v1-collar}
The principal one-layer collar of f6-V98--V100 has
\[
 |P_0|=6,\quad |E_0|=19,\quad |V_0|=14,\qquad
 |P^+|=72,\quad |E^+|=123,\quad |V^+|=64.
\]
Its scalar curvature matrix has exact rational rank \(60\), and its
\(63\)-dimensional kernel is exactly the vertex-gauge image.  After
fixing the \(13\)-edge inner tree, the remaining \(50\) null directions
are precisely boundary frames.

The \(123\) links split into a \(63\)-edge spanning tree and \(60\)
curvature chords.  A displayed \(60\times60\) integer minor has determinant
\(-1\), \(102\) nonzero entries, and gives the explicit squared
singular-value floor
\[
 \sigma_{\min}^2\ge102^{-59}.
\]
The V100 nonabelian presentation has \(60\) chord generators and \(72\)
plaquette relators.  A sixty-step singleton-relator Tietze certificate
eliminates every generator, proving
\[
 \pi_1(K_e^+)=1.
\]
Consequently every fixed full-tree frame has exactly one flat chord
completion.  Compactness and the local inverse theorem then give constants
\(a_{\rm glob},C_{\rm glob}>0\) such that
\[
 \mathfrak D_{K_e^+}(U)<a_{\rm glob}
 \quad\Longrightarrow\quad
 \dist(U,U^{\rm fl})^2
 \le C_{\rm glob}\mathfrak D_{K_e^+}(U),
\]
uniformly in the finite volume, collar location, and conditioned tree
frames.  Thus the census, curvature, nonlinear flat-sector, and global
reconstruction rows of the chord-conditioned Wilson collar card are
closed.
\end{record}

\section{The live implication interface}
\label{sec:s7v1-interface}

\begin{definition}[Remaining chord-conditioned rows]
\label{def:s7v1-remaining}
The unresolved part of the chord-conditioned Wilson collar card consists
of:
\begin{enumerate}[label=\textup{(C\arabic*)},start=5,leftmargin=4em]
\item a complete-action restricted-smoothing estimate outside the
      reconstructed core;
\item response-ground stability inside the good core; and
\item a translated family of chord updates which covers every physical
      fluctuation with cutoff-uniform approximate tensorization and
      Dirichlet-form incidence.
\end{enumerate}
\end{definition}

\begin{record}[What these rows would buy]
\label{rec:s7v1-compiler}
The proved f6-V97--V100 compiler has the following conditional form.
If rows (C5)--(C7) hold with strict constants, then the block
core/bad-set inequality gives a cutoff-uniform conditional Poincar\'e
constant.  Approximate tensorization and physical form incidence then
give the corresponding finite-volume gap row.  Separate continuum,
reflection-positivity, Ward-quotient, vacuum-sector, and scaling
identifications are still required before this can become a physical
Yang--Mills mass gap.
\end{record}

\begin{assessment}[Current bottleneck]
\label{ass:s7v1-bottleneck}
The finite collar topology is complete.  The first open question is now
whether literal complete-action internality can mean containment of every
plaquette touching an updated link in one uniformly bounded block.  The
next chapter proves that it cannot.
\end{assessment}

\chapter{The finite action-closure obstruction}
\label{ch:s7v1-closure}

\section{Wilson incidence closure}
\label{sec:s7v1-closure}

Work on the nearest-neighbour hypercubic lattice \(\mathbb Z^d\),
\(d\ge2\).  A positive oriented link is denoted \((x;\mu)\), and a
positive plaquette by \((x;\mu,\nu)\), \(\mu<\nu\), with boundary
\begin{equation}
 (x;\mu)+(x+e_\mu;\nu)
 -(x+e_\nu;\mu)-(x;\nu).                                 \label{eq:s7v1-boundary}
\end{equation}

\begin{definition}[One-step action closure]
\label{def:s7v1-action-closure}
For a link set \(A\), define
\begin{align}
 {\cal P}(A)
 &:=\{p:E(p)\cap A\ne\varnothing\},                       \label{eq:s7v1-PA}\\
 {\cal L}(A)
 &:=\bigcup_{p\in{\cal P}(A)}E(p).                        \label{eq:s7v1-LA}
\end{align}
Thus \({\cal L}(A)\) is the complete link support of all Wilson
plaquettes whose action depends on at least one variable in \(A\).
Call \(A\) action closed if \({\cal L}(A)\subseteq A\).
\end{definition}

\begin{theorem}[No nonempty finite Wilson-action-closed block]
\label{thm:s7v1-no-finite-closure}
On \(\mathbb Z^d\), \(d\ge2\), no nonempty finite link set is action
closed.
\end{theorem}

\begin{proof}
Let \((x;\mu)\in A\), and choose \(\nu\ne\mu\).  The positive plaquette
containing this link on its \(+\nu\) side also contains the parallel link
\((x+e_\nu;\mu)\).  Hence action closure implies
\[
 (x+e_\nu;\mu)\in A.
\]
Applying the same argument inductively gives
\[
 (x+ne_\nu;\mu)\in A\qquad(n=0,1,2,\ldots).
\]
These are distinct links, contradicting finiteness.  Therefore a finite
action-closed set must be empty.
\end{proof}

\begin{corollary}[Quantitative growth of iterated closure]
\label{cor:s7v1-growth}
Let \(A_0\) contain a link \((x;\mu)\), and define
\(A_{n+1}={\cal L}(A_n)\).  For any fixed \(\nu\ne\mu\),
\begin{equation}
 \{(x+ke_\nu;\mu):-n\le k\le n\}\subseteq A_n,
 \qquad |A_n|\ge2n+1.                                    \label{eq:s7v1-growth}
\end{equation}
In particular, repeated finite-collar enlargement cannot stabilize with
a cutoff-independent number of links.
\end{corollary}

\begin{proof}
The assertion is true for \(n=0\).  If all translates with
\(|k|\le n\) lie in \(A_n\), the two plaquettes on the
\(\pm\nu\) sides of the endpoint links add the translates with
\(k=n+1\) and \(k=-n-1\) to \(A_{n+1}\).  Induction proves the inclusion
and the count.
\end{proof}

\begin{corollary}[Periodic regulators do not give a cofinal local block]
\label{cor:s7v1-periodic}
On a periodic lattice, action closure of one link contains its complete
parallel translate orbit in every transverse direction.  Hence the size
of an action-closed block grows at least as a transverse circumference
and cannot remain uniformly bounded along a growing-volume regulator
sequence.
\end{corollary}

\begin{proof}
Repeat the translation argument in
\Cref{thm:s7v1-no-finite-closure} until the orbit closes periodically.
The orbit length is the corresponding transverse circumference.
\end{proof}

\section{Exact boundary cost of the V100 chord update}
\label{sec:s7v1-count}

Let \(C\) be the \(60\)-link chord set in the full-tree decomposition of
the V100 collar, and retain \(P^+\) and \(E^+\) from the archive.

\begin{proposition}[Exact first crossing shell]
\label{prop:s7v1-crossing-count}
The complete set of plaquettes incident to the \(60\) chord links has
\[
 |{\cal P}(C)|=250.
\]
Exactly \(72\) of these are the principal plaquettes \(P^+\), while
\[
 |{\cal P}(C)\setminus P^+|=178.
\]
Their total link support has
\[
 |{\cal L}(C)|=382,\qquad
 |{\cal L}(C)\setminus E^+|=259.                         \label{eq:s7v1-crossing-count}
\]
\end{proposition}

\begin{proof}
Use the duplicate-free convention \eqref{eq:s7v1-boundary}.  Each
positive link \((x;\mu)\) belongs, for every \(\nu\ne\mu\), to the
plaquette based at \(x\) and the plaquette based at \(x-e_\nu\).
Enumerate these six plaquettes for each of the sixty lexicographically
specified V100 chords, remove duplicate tuples, and then enumerate their
four boundary links.  Set subtraction against the archived \(72\)
plaquettes and \(123\) links gives the displayed counts.

The calculation is reproduced by
\texttt{scripts/verify\_v1\_boundary.py}, which uses only integer tuples,
set insertion, and free-word reduction.  At creation it had \(233\)
source lines, \(8\,094\) bytes, and SHA--256
\[
 \texttt{932FCDAC0EC6AB6440C9F2F3F7FAE3CDCDBC99FB4E876C1E054ECC79310A3233}.
\]
\end{proof}

\begin{assessment}[Meaning of the count]
\label{ass:s7v1-count}
The V100 \(72\)-plaquette action is the correct principal curvature
action and is sufficient for topology and reconstruction.  It is not the
complete conditional Wilson action for a \(60\)-chord update: \(178\)
additional plaquettes depend on those variables.  Enlarging once to their
\(259\) new links merely moves the same issue outward by
\Cref{thm:s7v1-no-finite-closure}.
\end{assessment}

\section{Why a bounded-oscillation repair has the wrong scale}
\label{sec:s7v1-scale}

\begin{proposition}[Core-size incompatibility for termwise Holley repair]
\label{prop:s7v1-Holley-no-go}
Suppose a principal collar core has defect threshold \(a_s\), its
reconstruction gives chord radius \(O(\sqrt{a_s})\), and every crossing
Wilson term is treated as a bounded perturbation using its Lipschitz
oscillation on that core.  Then:
\begin{enumerate}[label=\textup{(\roman*)},leftmargin=3.5em]
\item a coupling-uniform perturbation constant requires
      \(s\sqrt{a_s}=O(1)\), hence \(a_s=O(s^{-2})\);
\item the bare action-margin exponent then satisfies
      \(s a_s=O(s^{-1})\) and cannot tend to infinity.
\end{enumerate}
Therefore this termwise bounded-perturbation strategy cannot
simultaneously give a uniform core comparison and a vanishing action-tail
debit as \(s\to\infty\).
\end{proposition}

\begin{proof}
A crossing plaquette is a product of a bounded number of group variables.
With the exterior fixed, its trace is uniformly Lipschitz in the chord
product metric.  The Wilson coefficient multiplies this oscillation by
\(s\).  Reconstruction bounds the chord diameter of the principal core
by \(C\sqrt{a_s}\), so the standard termwise Holley exponent is bounded by
\[
 C_\partial s\sqrt{a_s}.
\]
Keeping it bounded forces \(a_s\le C's^{-2}\).  The principal Wilson
action gap outside the core is \(s a_s\), which is then \(O(s^{-1})\).
Thus it cannot supply the required growing suppression exponent.  This
is precisely the scaling conflict; no assertion is made about a
non-Holley cancellation mechanism.
\end{proof}

\begin{firewall}[Scope of the no-go theorem]
\label{fire:s7v1-Holley-scope}
\Cref{prop:s7v1-Holley-no-go} rules out a proof architecture, not the
conditional gap.  It does not preclude direct convexity of the complete
conditional action on a stronger good event, an averaged boundary-field
estimate, a multiscale shell contraction, or a reflection-positive
comparison.
\end{firewall}

\chapter{A viable replacement: fully good boundary curvature}
\label{ch:s7v1-boundary}

\section{The complete incident action}
\label{sec:s7v1-complete}

For fixed values \(\eta\) of all non-chord links, write
\begin{equation}
 {\cal V}^{\eta}_{\rm inc}(U)
 :=
 \sum_{p\in{\cal P}(C)}
 \bigl(3-\Re\Tr W_p(U,\eta)\bigr),                        \label{eq:s7v1-Vinc}
\end{equation}
where \(U\in SU(3)^{60}\) are the chord variables.  Split
\[
 {\cal V}^{\eta}_{\rm inc}
 ={\cal V}_{+}+{\cal V}^{\eta}_{\partial}
\]
over \(P^+\) and \({\cal P}(C)\setminus P^+\), respectively.

\begin{definition}[Fully good incident set]
\label{def:s7v1-good}
For \(a,b>0\), let
\begin{equation}
 {\cal G}_{a,b}^{\eta}
 :=
 \left\{
 U:\ {\cal V}_{+}(U)<a,\quad
 \max_{p\in{\cal P}(C)\setminus P^+}
 \bigl(3-\Re\Tr W_p(U,\eta)\bigr)<b
 \right\}.                                               \label{eq:s7v1-good}
\end{equation}
\end{definition}

\begin{theorem}[Uniform complete-action Hessian on the fully good set]
\label{thm:s7v1-complete-Hessian}
Use tree-parallel-transport exponential coordinates and the product
bi-invariant metric.  There exist dimension-only constants
\(a_*,b_*,\lambda_*>0\) such that, for every exterior assignment \(\eta\),
every \(U\in{\cal G}_{a_*,b_*}^{\eta}\), and every chord tangent \(X\),
\begin{equation}
 \operatorname{Hess}{\cal V}^{\eta}_{\rm inc}(U)[X,X]
 \ge\lambda_*\|X\|^2.                                   \label{eq:s7v1-complete-Hessian}
\end{equation}
One may take \(\lambda_*\) to be a fixed positive fraction of the V99
linear floor \(102^{-59}\).  For Wilson coupling \(s\), the Hessian of
\(s{\cal V}^{\eta}_{\rm inc}\) is at least \(s\lambda_*\).
\end{theorem}

\begin{proof}
At the unique principal flat completion \(U^{\rm fl}\), the Hessian of
\({\cal V}_{+}\) is the squared linearized curvature map.  The selected
minor from f6-V99 gives
\[
 \operatorname{Hess}{\cal V}_{+}(U^{\rm fl})[X,X]
 \ge102^{-59}\|X\|^2
\]
with the normalization used there.  The f6-V100 global reconstruction
theorem implies that \({\cal V}_{+}(U)<a\) places \(U\) uniformly close
to \(U^{\rm fl}\) when \(a\) is sufficiently small.  Smoothness of the
finite principal action then gives, for some \(a_*>0\),
\begin{equation}
 \operatorname{Hess}{\cal V}_{+}(U)[X,X]
 \ge {1\over2}102^{-59}\|X\|^2.                          \label{eq:s7v1-principal-Hessian}
\end{equation}

For one crossing plaquette, the function
\(\phi(W)=3-\Re\Tr W\) has zero first derivative and nonnegative Hessian
at \(W=I\).  The plaquette multiplication map and all its first two
derivatives are uniformly bounded on the compact finite product of
\(SU(3)\)'s.  Hence uniform continuity gives a constant
\(\varepsilon(b)\downarrow0\) as \(b\downarrow0\) such that
\[
 \operatorname{Hess}\bigl(\phi\circ W_p\bigr)(U)[X,X]
 \ge-\varepsilon(b)\|X\|^2
\]
whenever \(\phi(W_p)<b\).  There are exactly \(178\) crossing
plaquettes.  Choose \(b_*>0\) so that
\[
 178\,\varepsilon(b_*)\le {1\over4}102^{-59}.
\]
Summing this estimate with
\eqref{eq:s7v1-principal-Hessian} gives
\eqref{eq:s7v1-complete-Hessian} with
\(\lambda_*={1\over4}102^{-59}\).  All constants are uniform in
\(\eta\), because the exterior variables range over a compact group and
enter only through left and right translations and finite products.
\end{proof}

\begin{remark}[What was gained]
\label{rem:s7v1-gain}
The \(178\) boundary plaquettes need not be a bounded perturbation of the
principal action.  When they themselves have small Wilson defect, their
Hessians are almost nonnegative, so the principal curvature floor absorbs
their entire negative part.  This proves the good-sector coercivity half
of the complete-action row without shrinking the principal threshold as
\(s^{-2}\).
\end{remark}

\begin{firewall}[The remaining probability problem]
\label{fire:s7v1-boundary-probability}
For an arbitrary pointwise-fixed exterior \(\eta\), the set
\({\cal G}_{a_*,b_*}^{\eta}\) may be empty or have very small conditional
mass: the exterior can frustrate a crossing plaquette.  Therefore
\Cref{thm:s7v1-complete-Hessian} does not give an
essential-supremum conditional Poincar\'e estimate.  The missing theorem
must control bad boundary fields in distribution or pass their cost to a
larger-scale update; it cannot assume that every exterior is good.
\end{firewall}

\section{Replacement card and next target}
\label{sec:s7v1-card}

\begin{definition}[Boundary Wilson influence card, BWIC]
\label{def:s7v1-BWIC}
For a translated family of chord blocks, BWIC consists of:
\begin{enumerate}[label=\textup{(B\arabic*)},leftmargin=4em]
\item the principal \(72\)-plaquette reconstruction and curvature floor
      proved in f6-V99--V100;
\item the exact incident-action split
      \eqref{eq:s7v1-Vinc} and the fully good Hessian theorem
      \Cref{thm:s7v1-complete-Hessian};
\item a restricted smoothing or capacity estimate for the event that at
      least one of the \(178\) crossing plaquettes is bad;
\item an assignment of that bad boundary event to the next shell with
      an influence factor strictly below one after overlap;
\item summability of the iterated shell debits without an
      essential-supremum-over-all-exteriors shortcut;
\item response-ground stability on the resulting good components; and
\item approximate tensorization and physical Dirichlet-form incidence
      for the translated update family.
\end{enumerate}
\end{definition}

\begin{theorem}[BWIC is the non-circular replacement for literal
complete-action internality]
\label{thm:s7v1-BWIC-reduction}
Literal containment of all action contacts in a uniformly bounded block
is impossible by \Cref{thm:s7v1-no-finite-closure}.  If BWIC holds, then
its rows (B1)--(B5) supply CCWCC row (C5), row (B6) supplies (C6), and
row (B7) supplies (C7).  The proved f6 conditional-gap compiler then
applies with the BWIC constants.
\end{theorem}

\begin{proof}
The first assertion is exactly
\Cref{thm:s7v1-no-finite-closure,cor:s7v1-growth}.  On the fully good
sector, (B1)--(B2) give uniform complete-action coercivity.  Rows
(B3)--(B5) control its complement by the required restricted
smoothing/capacity mechanism and sum the boundary debits; together these
are the complete-action estimate demanded by (C5).  The identifications
of (B6) and (B7) with (C6) and (C7) are their definitions.  Substitution
into the f6-V97--V100 implication chain yields the final conditional
statement.
\end{proof}

\begin{decision}[Seventh-series V1 bottleneck]
\label{dec:s7v1-bottleneck}
The attempt to obtain (C5) by finite action closure is terminated: such a
block cannot exist.  The next note will attack BWIC (B3)--(B5).  The first
candidate is an oriented shell ownership rule for the \(178\) crossing
plaquettes, followed by a martingale estimate that charges each bad
boundary plaquette once and tests whether the resulting shell influence
is strictly contractive after exact overlap counting.
\end{decision}

\begin{programme}[Seventh-series V2]
\label{prog:s7v1-V2}
V2 will:
\begin{enumerate}[leftmargin=3em]
\item construct a translation-covariant owner map for crossing
      plaquettes;
\item enumerate the exact overlap of owned boundary shells;
\item derive the corresponding conditional variance recursion;
\item prove a strict influence criterion or exhibit the first
      noncontracting boundary mode; and
\item go upstream to a larger-scale or reflection-positive estimate if
      the one-shell recursion is noncontractive.
\end{enumerate}
\end{programme}

\begin{firewall}[Claim boundary after seventh-series V1]
\label{fire:s7v1-boundary}
This note summarizes the six frozen archives; proves that no nonempty
finite link block can contain every Wilson plaquette touching its updated
variables; computes the V100 chord update's exact first crossing shell;
rules out the termwise bounded-oscillation repair at the required weak
coupling scale; and proves uniform complete-action Hessian coercivity on
the fully good incident set.

It does not prove the boundary bad-event capacity, shell contraction,
response-ground stability, block tensorization, cutoff removal,
Osterwalder--Schrader reconstruction, or the Yang--Mills mass gap.
\end{firewall}

\begin{record}[Rollover parent]
\label{rec:s7v1-parent}
This compact series begins after
\texttt{Renewal\_Yang\_Mills\_Mass\_Gap\_Research\_Notes\_V100\_f6.tex},
which has \(47\,291\) logical source lines, \(1\,647\,389\) bytes, and
SHA--256
\[
 \texttt{63a6ee472d2127ff2d928791f1c720acff1dbb5193047c2d3cea14f7db759a4a}.
\]
\end{record}

% =============================================================================
% BEGIN APPENDED SEVENTH-SERIES V2 MATERIAL
% =============================================================================

\chapter{Oriented shell ownership and the exact influence threshold}
\label{ch:s7v2}

V1 proved that finite action closure is impossible and replaced it by the
boundary Wilson influence card.  The first task is to prevent the
\(178\) crossing plaquettes of each chord block from being counted
arbitrarily many times when their cost is passed to neighbouring blocks.
This chapter constructs an explicit translation-covariant owner,
computes its optimal orientation load, and derives the exact
\(\ell^2\) contraction threshold which any analytic boundary transfer
must beat.

\section{The translated oriented block family}
\label{sec:s7v2-family}

For each positive link \(e=(x;\mu)\), let \(C_e\) be the translated
V100 chord set after the cyclic coordinate permutation
\[
 \rho_\mu(j)=j+\mu\pmod4
\]
which sends the reference direction \(0\) to \(\mu\).  Let \(P_e^+\)
be its translated principal collar and put
\begin{equation}
 \partial P_e:={\cal P}(C_e)\setminus P_e^+.             \label{eq:s7v2-boundary-shell}
\end{equation}
The construction is translation covariant and uses the same sixty chord
indices in every orientation.

\begin{lemma}[Uniform row count]
\label{lem:s7v2-row-count}
For every positive link \(e\),
\begin{equation}
 |C_e|=60,\qquad |P_e^+|=72,\qquad
 |\partial P_e|=178.                                    \label{eq:s7v2-row-count}
\end{equation}
\end{lemma}

\begin{proof}
Coordinate permutations and translations are bijections of links and
plaquettes and preserve incidence.  Apply the exact reference count in
\Cref{prop:s7v1-crossing-count}.
\end{proof}

\section{How many translated blocks cross one plaquette?}
\label{sec:s7v2-multiplicity}

For an oriented coordinate pair \(0\le a<b\le3\), let \(m_{ab}\) be the
number of block centres \(e=(x;\mu)\) for which the reference plaquette
\((0;a,b)\) belongs to \(\partial P_e\).  Translation covariance makes
this independent of the plaquette basepoint.

\begin{theorem}[Exact crossing multiplicities]
\label{thm:s7v2-multiplicity}
The six multiplicities are
\begin{equation}
\begin{array}{c|cccccc}
(a,b)&(0,1)&(0,2)&(0,3)&(1,2)&(1,3)&(2,3)\\
\hline
m_{ab}&118&120&118&118&120&118.
\end{array}                                               \label{eq:s7v2-multiplicity}
\end{equation}
\end{theorem}

\begin{proof}
For each \(\mu\), cyclically permute the \(178\) reference crossing
plaquettes.  For a permuted crossing plaquette
\((y;a,b)\), there is exactly one translation of its block centre which
makes it equal to \((0;a,b)\), namely translation by \(-y\).  Thus
\(m_{ab}\) is obtained by counting the entries with direction pair
\((a,b)\) in the four permuted finite lists.  Duplicate block centres are
removed.  The resulting six counts are exactly
\eqref{eq:s7v2-multiplicity}.

The enumeration is reproduced by
\texttt{scripts/verify\_v2\_shell.py}.  The script has \(166\) lines,
\(5\,577\) bytes, and SHA--256
\[
 \texttt{AA3062D188FB2CC3C6285458C91938F42F94A52827B2660CB972759CCD12475C}.
\]
It constructs the reference collar from
\eqref{eq:s7v1-boundary}, checks all four rotated counts, and obtains the
six numbers by exact integer-tuple set operations.
\end{proof}

\section{An optimal axis-owner rule}
\label{sec:s7v2-owner}

Choose, for each direction pair, one of its two axes:
\begin{equation}
\begin{array}{c|cccccc}
(a,b)&(0,1)&(0,2)&(0,3)&(1,2)&(1,3)&(2,3)\\
\hline
\alpha(a,b)&0&2&3&2&1&3.
\end{array}                                               \label{eq:s7v2-owner-axis}
\end{equation}
For \(p=(x;a,b)\), define its owner to be the positive link
\begin{equation}
 o(p):=(x;\alpha(a,b)).                                  \label{eq:s7v2-owner}
\end{equation}
This is one of the four boundary links of \(p\), even when it occurs with
negative sign in the oriented boundary.

\begin{lemma}[Ownership converts boundary badness to principal badness]
\label{lem:s7v2-owner-principal}
Let
\begin{align}
 \phi_p(U)&:=3-\Re\Tr W_p(U),                            \label{eq:s7v2-phi}\\
 A_o(b)&:=\left\{\sum_{q\in P_o^+}\phi_q\ge b\right\},   \label{eq:s7v2-A}\\
 B_e^\partial(b)&:=
 \left\{\max_{p\in\partial P_e}\phi_p\ge b\right\}.      \label{eq:s7v2-B}
\end{align}
Then
\begin{equation}
 {\bf1}_{B_e^\partial(b)}
 \le
 \sum_{p\in\partial P_e}{\bf1}_{A_{o(p)}(b)}.            \label{eq:s7v2-owner-indicator}
\end{equation}
\end{lemma}

\begin{proof}
If \(\phi_p\ge b\), then \(p\) is incident to its owner link \(o(p)\).
It is therefore one of the six central plaquettes of that owner's star,
hence belongs to \(P_{o(p)}^+\).  Nonnegativity of every Wilson defect
gives
\[
 \sum_{q\in P_{o(p)}^+}\phi_q\ge\phi_p\ge b.
\]
Take the union bound in indicator form over
\(p\in\partial P_e\).
\end{proof}

\begin{theorem}[Optimal translation-covariant orientation load]
\label{thm:s7v2-optimal-owner}
For the owner rule \eqref{eq:s7v2-owner-axis}, the incoming occurrence
loads at a fixed owner link of directions \(0,1,2,3\), respectively, are
\begin{equation}
 (L_0,L_1,L_2,L_3)=(118,120,238,236).                    \label{eq:s7v2-owner-loads}
\end{equation}
In particular,
\[
 L_*:=\max_\mu L_\mu=238.
\]
Among all rules which assign each direction pair \((a,b)\) to one of
its two axes and use the based boundary link \((x;\alpha(a,b))\), the
value \(238\) is optimal.
\end{theorem}

\begin{proof}
A fixed owner link of direction \(\mu\) owns exactly one based plaquette
for every pair assigned to \(\mu\).  Each such plaquette is a crossing
plaquette of \(m_{ab}\) translated source blocks.  Therefore
\[
 L_\mu=\sum_{\substack{a<b\\\alpha(a,b)=\mu}}m_{ab}.
\]
Substitution of
\eqref{eq:s7v2-multiplicity,eq:s7v2-owner-axis} gives
\eqref{eq:s7v2-owner-loads}.

For optimality, suppose the maximum load were at most \(237\).  The two
weight-\(120\) pairs are \((0,2)\) and \((1,3)\); they are disjoint.
Whichever endpoint owns either such pair cannot own any second pair,
because every other weight is at least \(118\) and
\(120+118=238\).  Let \(u\in\{0,2\}\) and
\(v\in\{1,3\}\) be the chosen owners of the two weight-\(120\) pairs.
The pair \((u,v)\) is one of the four weight-\(118\) pairs and must be
owned by either \(u\) or \(v\), contradicting the preceding restriction.
Thus every rule has maximum load at least \(238\), and the displayed rule
attains it.
\end{proof}

\begin{remark}[Why lexicographic ownership was rejected]
\label{rem:s7v2-lex}
Choosing the lexicographically first boundary link gives directional
occurrence loads \((356,238,118,0)\).  The balanced rule
\eqref{eq:s7v2-owner-axis} reduces the maximum from \(356\) to its optimal
value \(238\).  This is a genuine constant improvement, although the
remaining amplification is still large.
\end{remark}

\section{The shell incidence operator}
\label{sec:s7v2-operator}

Let \({\cal I}\) be a finite periodic family of oriented block centres
large enough that the reference shells do not wrap around themselves.
Define the nonnegative matrix
\begin{equation}
 R_{eo}
 :=
 \#\{p\in\partial P_e:o(p)=o\}.                          \label{eq:s7v2-R}
\end{equation}
The same inequalities below hold with no larger constants for a
finite-volume subfamily with boundary.

\begin{theorem}[Exact row and column Schur data]
\label{thm:s7v2-Schur-data}
The shell incidence matrix satisfies
\begin{equation}
 \|R\|_{\infty}=178,\qquad
 \|R\|_1=238,\qquad
 \|R\|_{2\to2}\le\sqrt{178\cdot238}=\sqrt{42364}.        \label{eq:s7v2-Schur}
\end{equation}
Moreover
\begin{equation}
 R{\bf1}=178\,{\bf1},                                    \label{eq:s7v2-constant-mode}
\end{equation}
so \(\|R\|_{2\to2}\ge178\).
\end{theorem}

\begin{proof}
Every row sum is \(|\partial P_e|=178\) by
\Cref{lem:s7v2-row-count}.  The column sum at a direction-\(\mu\) owner
is \(L_\mu\), so the maximum is \(238\) by
\Cref{thm:s7v2-optimal-owner}.  The Schur bound
\[
 \|R\|_{2\to2}^2\le\|R\|_1\|R\|_\infty
\]
gives \eqref{eq:s7v2-Schur}.  Equal row sums give
\eqref{eq:s7v2-constant-mode}; taking the norm of the constant vector
gives the lower bound.
\end{proof}

\begin{corollary}[Pointwise global owner inequality]
\label{cor:s7v2-global-owner}
For every configuration,
\begin{equation}
 \bigl({\bf1}_{B_e^\partial(b)}\bigr)_{e\in{\cal I}}
 \le
 R\bigl({\bf1}_{A_o(b)}\bigr)_{o\in{\cal I}}
                                                               \label{eq:s7v2-vector-owner}
\end{equation}
componentwise.
\end{corollary}

\begin{proof}
Regroup the sum in
\eqref{eq:s7v2-owner-indicator} according to \(o(p)\).  Its multiplicity
is exactly \eqref{eq:s7v2-R}.
\end{proof}

\section{The exact analytic contraction window}
\label{sec:s7v2-window}

\begin{theorem}[Unnormalized shell-transfer criterion]
\label{thm:s7v2-window}
Suppose an analytic smoothing step produces nonnegative Hilbert-load
vectors \(x,z,r\in\ell^2({\cal I})\) satisfying
\begin{equation}
 z\le q R x+r                                             \label{eq:s7v2-transfer}
\end{equation}
componentwise, where \(q\ge0\) is the per-owned-occurrence transfer
coefficient.  Then
\begin{equation}
 \|z\|_2
 \le q\sqrt{42364}\,\|x\|_2+\|r\|_2.                    \label{eq:s7v2-transfer-bound}
\end{equation}
Consequently the shell recursion is strictly contractive whenever
\begin{equation}
 \boxed{q<42364^{-1/2}.}                                 \label{eq:s7v2-window}
\end{equation}
On periodic regulators, \(q<1/178\) is necessary for a scalar multiple
of this unnormalized incidence map to contract the constant mode.
\end{theorem}

\begin{proof}
For nonnegative vectors, componentwise domination implies norm
domination.  Apply \Cref{thm:s7v2-Schur-data} and the triangle inequality
to \eqref{eq:s7v2-transfer}.  If
\eqref{eq:s7v2-window} holds, the linear coefficient is below one and
the Neumann series for repeated shells converges.  Conversely,
\eqref{eq:s7v2-constant-mode} shows that \(qR\) has eigenvalue \(178q\)
on the constant vector, proving the necessary condition.
\end{proof}

\begin{corollary}[Normalized shell-transfer criterion]
\label{cor:s7v2-normalized}
If the analytic step instead produces
\[
 z\le {q\over178}Rx+r,
\]
then a sufficient strict window is
\begin{equation}
 \boxed{q<\sqrt{178/238}=\sqrt{89/119}.}                 \label{eq:s7v2-normalized}
\end{equation}
The constant mode requires \(q<1\).
\end{corollary}

\begin{proof}
Divide \eqref{eq:s7v2-Schur} by \(178\) and apply
\Cref{thm:s7v2-window}.  The constant-mode statement follows from
\eqref{eq:s7v2-constant-mode}.
\end{proof}

\begin{firewall}[The indicator owner map is not yet an operator transfer]
\label{fire:s7v2-operator}
\Cref{cor:s7v2-global-owner} is a pointwise event inequality.  It does
not imply \eqref{eq:s7v2-transfer} for conditional fluctuation operators:
the source and owner blocks overlap, their conditional projections do
not generally commute, and the physical ground transform may move.
The coefficient \(q\) must come from a proved restricted smoothing or
capacity estimate.  It cannot be replaced by the unconditional
probability of \(A_o(b)\).
\end{firewall}

\begin{assessment}[What the exact threshold says]
\label{ass:s7v2-threshold}
The ownership geometry is finite and completely controlled, but a raw
union transfer pays between \(178\) and \(\sqrt{42364}\) in its constant
or Schur mode.  A normalized average has a much milder sufficient
threshold \(\sqrt{89/119}\), but the event union in
\eqref{eq:s7v2-owner-indicator} supplies no factor \(1/178\).  The next
analytic task is therefore not more counting: it is to obtain an
energy-weighted or martingale average which creates a normalization
without discarding rare large boundary defects.
\end{assessment}

\begin{decision}[V2 bottleneck decision]
\label{dec:s7v2-bottleneck}
BWIC's owner and overlap rows are now explicit.  The next note must test
whether the Wilson sum itself supplies the missing normalization.  It
will replace the maximum event by the boundary defect
\[
 D_e^\partial:=\sum_{p\in\partial P_e}\phi_p
\]
and seek a conditional form inequality for
\(D_e^\partial/178\).  If the conversion from the maximum bad event to
this average loses the factor \(178\) again, the one-shell route is
noncontractive and the programme will go upstream to a multi-shell
energy or reflection-positive estimate.
\end{decision}

\begin{programme}[Seventh-series V3]
\label{prog:s7v2-V3}
V3 will derive the exact double-counting identity for
\(\sum_eD_e^\partial\), compare maximum and averaged boundary events,
and determine whether an energy-weighted Efron--Stein inequality yields
the normalized transfer in \Cref{cor:s7v2-normalized}.  The distinction
between probability, multiplication-operator norm, and restricted
smoothing norm will remain explicit.
\end{programme}

\begin{firewall}[Claim boundary after seventh-series V2]
\label{fire:s7v2-boundary}
V2 proves the exact orientation multiplicities, an optimal
translation-covariant plaquette owner, the row/column shell loads, and
the unnormalized and normalized \(\ell^2\) contraction windows.

V2 does not prove the analytic transfer coefficient \(q\), boundary
capacity, response-ground stability, tensorization, cutoff removal,
Osterwalder--Schrader reconstruction, or the Yang--Mills mass gap.
\end{firewall}

\paragraph{Parent-version record.}
V2 was formed from
\texttt{Renewal\_Yang\_Mills\_Mass\_Gap\_Research\_Notes\_V1.tex}.
The V1 parent had \(703\) logical source lines, \(28\,473\) bytes, and
SHA--256
\[
 \texttt{73E8194BC4415A65ED8B5A2F9C50FBE1EFBC4D916AD62A28EF72B1CD07DEEBAF}.
\]

% =============================================================================
% END APPENDED SEVENTH-SERIES V2 MATERIAL
% =============================================================================

% =============================================================================
% BEGIN APPENDED SEVENTH-SERIES V3 MATERIAL
% =============================================================================

\chapter{Boundary-energy double counting and the normalization fork}
\label{ch:s7v3}

V2 found the exact shell incidence matrix.  A raw union transfer has row
sum \(178\) and cannot acquire a normalization from ownership alone.
There is nevertheless a strict global factor hidden in the Wilson
\emph{sum}: after division by \(178\), every plaquette defect is counted
only \(118/178\) or \(120/178\) times.  This chapter proves the resulting
\(60/89\) contraction and, just as importantly, proves that it cannot be
obtained by a deterministic maximum-to-average conversion.

\section{Exact global boundary-energy identity}
\label{sec:s7v3-double}

For a block centre \(e\), set
\begin{equation}
 D_e^\partial
 :=\sum_{p\in\partial P_e}\phi_p,
 \qquad
 \overline D_e^\partial
 :={1\over178}D_e^\partial.                              \label{eq:s7v3-Dbar}
\end{equation}
Let \({\cal P}_{ab}\) be the plaquettes with direction pair
\((a,b)\), and let
\[
 D_{\rm tot}:=\sum_p\phi_p
\]
be the total Wilson plaquette defect.  Work first on a periodic regulator
larger than the fixed shell diameter, so that no reference shell wraps
around itself.

\begin{theorem}[Exact shell double counting]
\label{thm:s7v3-double}
\begin{equation}
 \sum_e D_e^\partial
 =
 118\!\!\sum_{\substack{a<b\\(a,b)\ne(0,2),(1,3)}}
       \ \sum_{p\in{\cal P}_{ab}}\phi_p
 +120\!\!\sum_{(a,b)=(0,2),(1,3)}
       \ \sum_{p\in{\cal P}_{ab}}\phi_p.                 \label{eq:s7v3-double}
\end{equation}
Consequently,
\begin{equation}
 {59\over89}D_{\rm tot}
 \le
 \sum_e\overline D_e^\partial
 \le
 {60\over89}D_{\rm tot}.                                \label{eq:s7v3-strict}
\end{equation}
The upper coefficient is strictly below one.
\end{theorem}

\begin{proof}
Interchange the two finite sums:
\[
 \sum_eD_e^\partial
 =\sum_p\phi_p\,\#\{e:p\in\partial P_e\}.
\]
The multiplicity in braces is \(m_{ab}\) for a plaquette of direction
pair \((a,b)\).  Substitute the six exact values from
\eqref{eq:s7v2-multiplicity} to obtain
\eqref{eq:s7v3-double}.  Divide by \(178=2\cdot89\);
the two possible coefficients are
\[
 {118\over178}={59\over89},
 \qquad
 {120\over178}={60\over89}.
\]
Nonnegativity of the defects proves
\eqref{eq:s7v3-strict}.
\end{proof}

\begin{corollary}[Finite-boundary version]
\label{cor:s7v3-finite}
For a finite collection of source blocks, with plaquettes outside the
regulator omitted rather than periodically identified,
\begin{equation}
 \sum_e\overline D_e^\partial
 \le {60\over89}D_{\rm tot}.                             \label{eq:s7v3-finite}
\end{equation}
\end{corollary}

\begin{proof}
Removing source blocks or crossing plaquettes can only decrease each
nonnegative plaquette multiplicity from its full periodic value.
\end{proof}

\section{A physical isolated-shell-defect witness}
\label{sec:s7v3-witness}

\begin{proposition}[Boundary support degree census]
\label{prop:s7v3-degree}
For the reference block, the \(178\) crossing plaquettes use \(343\)
links.  Of these, \(259\) lie outside the principal support \(E^+\).
Their crossing-shell degree histogram is
\begin{equation}
\begin{array}{c|ccc}
\text{number of incident crossing plaquettes}&1&2&3\\
\hline
\text{number of exterior links}&130&108&21.
\end{array}                                               \label{eq:s7v3-degree}
\end{equation}
In particular, \(130\) exterior links occur in exactly one plaquette of
the reference crossing shell.
\end{proposition}

\begin{proof}
Enumerate the four boundary links of each of the \(178\) crossing
plaquettes and count incidences.  Subtract the \(123\)-link principal
support.  The exact tuple census gives
\eqref{eq:s7v3-degree}.  A lexicographically first degree-one witness is
\[
 \ell_*=\bigl((-2,0,0,0);1\bigr),
\]
and its unique crossing plaquette is
\[
 p_*=\bigl((-2,0,0,0);0,1\bigr).
\]
The calculation is reproduced by
\texttt{scripts/verify\_v3\_boundary\_energy.py}, with \(47\) lines,
\(1\,700\) bytes, and SHA--256
\[
 \texttt{21403126B85A979110B483E9C93228B36268E4CE49F7AC3480FAC981F0364967}.
\]
\end{proof}

\begin{theorem}[No deterministic \(1/178\) conversion of the maximum
event]
\label{thm:s7v3-no-max-average}
Let
\[
 B_e^\partial(b)
 =\{\max_{p\in\partial P_e}\phi_p\ge b\}.
\]
Any pointwise inequality of the form
\begin{equation}
 {\bf1}_{B_e^\partial(b)}
 \le c\,{\overline D_e^\partial\over b}                  \label{eq:s7v3-max-average}
\end{equation}
valid for all lattice gauge configurations and all
\(0<b\le9/2\) must have
\begin{equation}
 c\ge178.                                                \label{eq:s7v3-c-lower}
\end{equation}
\end{theorem}

\begin{proof}
Use the degree-one exterior link \(\ell_*\) from
\Cref{prop:s7v3-degree}.  Set every link to the identity except
\[
 U_{\ell_*}=zI,\qquad z=e^{2\pi i/3}.
\]
Because \(\ell_*\notin E^+\), every principal plaquette remains flat.
Inside the reference crossing shell, only \(p_*\) contains \(\ell_*\).
Its holonomy is \(zI\) or \(z^{-1}I\), so
\[
 \phi_{p_*}
 =3-\Re\Tr(zI)
 =3-3\cos(2\pi/3)
 ={9\over2}.
\]
Thus for \(b=9/2\),
\[
 {\bf1}_{B_e^\partial(b)}=1,\qquad
 {\overline D_e^\partial\over b}={1\over178}.
\]
Substitution in \eqref{eq:s7v3-max-average} gives
\eqref{eq:s7v3-c-lower}.  The other five plaquettes incident to
\(\ell_*\) lie beyond this source shell; their nontriviality is
consistent with the lattice Bianchi constraints and is exactly the
outward propagation that BWIC must control.
\end{proof}

\begin{corollary}[Union bounds erase the strict global factor]
\label{cor:s7v3-union-loss}
The strict coefficient \(60/89\) in
\eqref{eq:s7v3-strict} cannot be combined with
\eqref{eq:s7v3-max-average} to control the maximum bad event with a
coefficient below one: the sharp universal conversion cost is at least
\(178\).
\end{corollary}

\begin{proof}
The product of any valid universal conversion coefficient with
\(60/89\) is at least
\[
 178\cdot{60\over89}=120>1.
\]
\end{proof}

\section{Source-independent normalized form transfer}
\label{sec:s7v3-form}

The double-counting gain survives for analytic debits only if the debit
attached to a plaquette is independent of which source block sees that
plaquette.

\begin{theorem}[Normalized plaquette-transfer compiler]
\label{thm:s7v3-compiler}
Let \(Z_e(f)\ge0\) be block boundary debits and
\(Q_p(f)\ge0\) plaquette debits such that, for every \(e\),
\begin{equation}
 Z_e(f)
 \le {1\over178}\sum_{p\in\partial P_e}Q_p(f)+r_e(f),
 \qquad r_e(f)\ge0.                                     \label{eq:s7v3-local-transfer}
\end{equation}
Crucially, \(Q_p(f)\) depends on \(p\) and \(f\), but not on the source
block \(e\).  Then
\begin{equation}
 \sum_e Z_e(f)
 \le {60\over89}\sum_pQ_p(f)+\sum_e r_e(f).              \label{eq:s7v3-global-transfer}
\end{equation}
If in addition
\begin{equation}
 \sum_pQ_p(f)
 \le\kappa\,{\cal E}(f)+r_0(f),                          \label{eq:s7v3-incidence}
\end{equation}
then the relative form coefficient is strictly below one whenever
\begin{equation}
 \boxed{\kappa<{89\over60}.}                             \label{eq:s7v3-kappa-window}
\end{equation}
\end{theorem}

\begin{proof}
Sum \eqref{eq:s7v3-local-transfer} in \(e\) and interchange the finite
sums.  A fixed \(Q_p(f)\) appears \(m_{ab}\) times when \(p\) has
direction pair \((a,b)\).  Therefore the same calculation as
\Cref{thm:s7v3-double} gives
\[
 {1\over178}\sum_e\sum_{p\in\partial P_e}Q_p(f)
 \le {60\over89}\sum_pQ_p(f).
\]
This proves \eqref{eq:s7v3-global-transfer}.  Insert
\eqref{eq:s7v3-incidence}; the coefficient of \({\cal E}(f)\) is
\((60/89)\kappa\), which is below one exactly under
\eqref{eq:s7v3-kappa-window}.
\end{proof}

\begin{remark}[Comparison with V2]
\label{rem:s7v3-comparison}
The V2 matrix estimate treats every owned occurrence separately and
pays the shell operator \(R\).  The present theorem identifies equal
plaquette debits before summation and divides at the block level by the
row count \(178\).  This changes the geometric factor from at least
\(178\) to at most \(60/89\).  The gain is valid only if the analytic
construction really produces the source-independent
\(Q_p(f)\) and the normalized local inequality
\eqref{eq:s7v3-local-transfer}.
\end{remark}

\begin{firewall}[Source dependence destroys double counting]
\label{fire:s7v3-source}
If the local estimate produces \(Q_{e,p}(f)\), with different conditional
projections or ground transforms for different source blocks \(e\), then
there is no identity reducing it to one \(Q_p(f)\).  Replacing
\(Q_{e,p}\) by a supremum recreates an essential-supremum debit and may
lose the strict factor.  Equality of the plaquette debit across all
incident source blocks is therefore a substantive covariance and
commutation theorem, not a relabelling.
\end{firewall}

\begin{definition}[Plaquette-normalized capacity card, PNCC]
\label{def:s7v3-PNCC}
PNCC consists of:
\begin{enumerate}[label=\textup{(P\arabic*)},leftmargin=4em]
\item a positive self-adjoint plaquette transfer and debit \(Q_p(f)\)
      defined in the owner block of \(p\);
\item transport or commutation identities making this same \(Q_p(f)\)
      valid for every source block with \(p\in\partial P_e\);
\item the normalized local boundary estimate
      \eqref{eq:s7v3-local-transfer}, obtained analytically rather than
      from the maximum-event indicator;
\item physical form incidence
      \eqref{eq:s7v3-incidence} with
      \(\kappa<89/60\);
\item summable remainder terms and compatibility with the conditional
      ground transform; and
\item the response-ground and translated-block rows already isolated in
      BWIC.
\end{enumerate}
\end{definition}

\begin{corollary}[PNCC closes the scalar shell-summation rows]
\label{cor:s7v3-PNCC}
PNCC rows (P1)--(P5) imply BWIC rows (B3)--(B5), with strict geometric
factor \(60/89\).  Row (P6) retains the separate response and
tensorization obligations.
\end{corollary}

\begin{proof}
Rows (P1)--(P3) give
\eqref{eq:s7v3-local-transfer}; row (P4) and
\Cref{thm:s7v3-compiler} give strict absorption.  Row (P5) controls the
remainders and ground-transform contacts.  These are precisely the
boundary-capacity and summation content of BWIC (B3)--(B5).
\end{proof}

\begin{decision}[V3 bottleneck decision]
\label{dec:s7v3-bottleneck}
Boundary-energy double counting is favorable: a genuinely normalized,
source-independent plaquette debit has a strict \(60/89\) global factor.
The maximum-event route is closed by
\Cref{thm:s7v3-no-max-average}.  The next note must therefore build a
randomized plaquette transfer whose Dirichlet form contains the
\(1/178\) normalization at the operator level and test whether its
conditional projections can be transported to one owner debit \(Q_p\).
\end{decision}

\begin{programme}[Seventh-series V4]
\label{prog:s7v3-V4}
For each source block, V4 will average the \(178\) owner-plaquette
transfers, prove positivity and its exact Dirichlet-form identity, and
compute the projection/ground commutators which obstruct
source-independent \(Q_p\).  If exact transport fails, the note will
separate the gauge-coordinate part from the physical conditional-density
part and state the smallest remaining covariance estimate.
\end{programme}

\begin{firewall}[Claim boundary after seventh-series V3]
\label{fire:s7v3-boundary}
V3 proves the exact boundary-energy identity and strict \(60/89\)
normalized factor, an explicit physical configuration showing the sharp
loss \(178\) in maximum-to-average conversion, and the
source-independent normalized form compiler with window
\(\kappa<89/60\).

V3 does not construct PNCC's plaquette transfer, prove its projection or
ground covariance, establish response stability or tensorization, remove
the cutoff, reconstruct the continuum theory, or prove the Yang--Mills
mass gap.
\end{firewall}

\paragraph{Parent-version record.}
V3 was formed from
\texttt{Renewal\_Yang\_Mills\_Mass\_Gap\_Research\_Notes\_V2.tex}.
The V2 parent had \(1\,071\) logical source lines, \(42\,336\) bytes,
and SHA--256
\[
 \texttt{B9DA47F32B5DAD2C083C6B4550996CC9BAA54F81816DA35B6C21CC61D393A5F4}.
\]

% =============================================================================
% END APPENDED SEVENTH-SERIES V3 MATERIAL
% =============================================================================

% =============================================================================
% BEGIN APPENDED SEVENTH-SERIES V4 MATERIAL
% =============================================================================

\chapter{Randomized plaquette transfer and the exact leakage debit}
\label{ch:s7v4}

The normalized compiler in V3 requires one source-independent debit
\(Q_p(f)\) for every plaquette and a factor \(1/178\) before shell
summation.  Convex averaging of positive self-adjoint plaquette transfers
produces exactly this structure.  Those transfers generally update
variables in the source block's exterior, so they need not commute with
the source conditional projection.  This chapter retains the resulting
failure as one explicit positive leakage operator and proves the corrected
projected-smoothing inequality.

\section{The averaged positive transfer}
\label{sec:s7v4-average}

Let \((\Omega,\mu)\) be a finite-volume positive Yang--Mills carrier.
For every plaquette \(p\), let \(K_p\) be a positive self-adjoint Markov
contraction on \(L^2(\mu)\):
\begin{equation}
 K_p1=1,\qquad K_p=K_p^*,\qquad 0\le K_p\le I.           \label{eq:s7v4-K}
\end{equation}
The choice of \(K_p\) is made in the owner block of \(p\) and is
independent of which source block sees \(p\) in its boundary shell.
For a source block \(e\), define
\begin{equation}
 P_e:={1\over178}\sum_{p\in\partial P_e}K_p,
 \qquad A_e:=I-P_e.                                     \label{eq:s7v4-P}
\end{equation}

\begin{theorem}[Exact randomized-transfer identity]
\label{thm:s7v4-average}
For every source block \(e\),
\[
 P_e1=1,\qquad P_e=P_e^*,\qquad0\le P_e\le I.
\]
Moreover, with
\begin{equation}
 Q_p(f):=\langle f,(I-K_p)f\rangle,                     \label{eq:s7v4-Qp}
\end{equation}
one has the exact normalized form identity
\begin{equation}
 \boxed{
 \langle f,A_ef\rangle
 ={1\over178}\sum_{p\in\partial P_e}Q_p(f).}            \label{eq:s7v4-form-identity}
\end{equation}
Consequently,
\begin{equation}
 \sum_e\langle f,A_ef\rangle
 \le {60\over89}\sum_pQ_p(f).                           \label{eq:s7v4-global-form}
\end{equation}
\end{theorem}

\begin{proof}
The cone of positive self-adjoint contractions is convex, as is the
affine set of constant-preserving operators.  Since the shell has exactly
\(178\) plaquettes, their average has all properties in
\eqref{eq:s7v4-K}.  Expanding \(I-P_e\) proves
\eqref{eq:s7v4-form-identity}.  The debit \(Q_p(f)\) is
source independent by construction, so
\Cref{thm:s7v3-compiler} gives
\eqref{eq:s7v4-global-form} with zero remainder.
\end{proof}

\begin{remark}[A concrete but not automatically non-circular choice]
\label{rem:s7v4-choice}
One may take \(K_p\) to be a conditional heat-bath or short-time
reversible transfer in the owner block \(o(p)\).  This makes
\eqref{eq:s7v4-K} automatic.  Taking a full owner-block heat-bath is not
by itself a proof of physical form incidence: bounding its conditional
variance by the local Yang--Mills Dirichlet form may already require the
gap being sought.  The choice of transfer and its incidence estimate must
therefore be audited together.
\end{remark}

\section{Source conditional projection and leakage}
\label{sec:s7v4-leakage}

Let \(E_e\) be conditional expectation over the source chord variables,
with the source exterior held fixed, and put
\[
 D_e:=I-E_e.
\]
The owner transfers in \(P_e\) can update source-exterior variables.
Define their positive source-mean leakage by
\begin{equation}
 L_e:=E_eA_eE_e
 ={1\over178}\sum_{p\in\partial P_e}E_e(I-K_p)E_e,
 \qquad
 \ell_e:=\|L_e\|_{2\to2}.                               \label{eq:s7v4-leakage}
\end{equation}

\begin{lemma}[Leakage properties]
\label{lem:s7v4-leakage}
\begin{equation}
 0\le L_e\le E_e,\qquad 0\le\ell_e\le1.                 \label{eq:s7v4-leakage-bounds}
\end{equation}
Moreover, \(L_e=0\) if and only if \(P_eE_e=E_e\).  This condition holds
when \(P_e\) is a fibrewise Markov transfer for the same disintegration
as \(E_e\); commutation by itself is not sufficient.
\end{lemma}

\begin{proof}
Because \(0\le A_e\le I\), compression by the orthogonal projection
\(E_e\) gives \(0\le E_eA_eE_e\le E_e\), proving the bounds.  If
\(P_eE_e=E_e\), then \(A_eE_e=0\), hence \(L_e=0\).  Conversely,
\(L_e=(A_e^{1/2}E_e)^*(A_e^{1/2}E_e)=0\) implies \(A_eE_e=0\), or
\(P_eE_e=E_e\).  A fibrewise Markov transfer on the same conditional
fibres fixes every exterior-measurable function, which gives this
identity.  An operator may commute with \(E_e\) while acting
nontrivially on its range, so commutation alone does not imply it.
\end{proof}

\begin{warning}[Commutation is not automatic here]
\label{warn:s7v4-commutation}
The last sentence of \Cref{lem:s7v4-leakage} applies when the transfer and
conditional expectation use the same disintegration.  A boundary-owner
transfer generally changes variables in the source exterior and does not
have that property.  Thus the exact V94 commuting-block theorem cannot be
invoked for \(P_e\) without a new nested-disintegration argument.
\end{warning}

\section{Corrected projected smoothing without commutation}
\label{sec:s7v4-projected}

Let \(B_e\) be any source-block bad event and define the restricted
averaged-transfer coefficient
\begin{equation}
 \eta_e
 :=\|M_{B_e}P_eD_e\|_{2\to2}^2.                         \label{eq:s7v4-eta}
\end{equation}

\begin{lemma}[Compression of a noncommuting form]
\label{lem:s7v4-compression}
For every \(f\in L^2(\mu)\) and every \(\alpha>0\),
\begin{equation}
 \langle D_ef,A_eD_ef\rangle
 \le
 (1+\alpha)\langle f,A_ef\rangle
 +(1+\alpha^{-1})\ell_e\|E_ef\|_2^2.                   \label{eq:s7v4-compression}
\end{equation}
\end{lemma}

\begin{proof}
Since \(A_e\ge0\),
\[
 A_e^{1/2}D_ef
 =A_e^{1/2}f-A_e^{1/2}E_ef.
\]
The tunable square inequality gives
\begin{align*}
 \|A_e^{1/2}D_ef\|_2^2
 &\le(1+\alpha)\|A_e^{1/2}f\|_2^2\\
 &\quad +(1+\alpha^{-1})
        \|A_e^{1/2}E_ef\|_2^2.
\end{align*}
The last square equals
\(\langle E_ef,L_eE_ef\rangle\) and is at most
\(\ell_e\|E_ef\|_2^2\).  This proves
\eqref{eq:s7v4-compression}.
\end{proof}

\begin{theorem}[Randomized projected-smoothing inequality]
\label{thm:s7v4-projected}
For every \(\theta,\alpha>0\),
\begin{align}
 \|M_{B_e}D_ef\|_2^2
 \le{}&
 (1+\theta)\eta_e\|D_ef\|_2^2 \notag\\
 &+(1+\theta^{-1})(1+\alpha)
   {1\over178}\sum_{p\in\partial P_e}Q_p(f) \notag\\
 &+(1+\theta^{-1})(1+\alpha^{-1})
   \ell_e\|E_ef\|_2^2.                                  \label{eq:s7v4-projected}
\end{align}
\end{theorem}

\begin{proof}
Split
\[
 M_{B_e}D_ef
 =M_{B_e}P_eD_ef+M_{B_e}A_eD_ef
\]
and use the tunable square inequality.  The first square is at most
\(\eta_e\|D_ef\|_2^2\).  Since \(0\le A_e\le I\),
\[
 \|M_{B_e}A_eD_ef\|_2^2
 \le\|A_eD_ef\|_2^2
 =\langle D_ef,A_e^2D_ef\rangle
 \le\langle D_ef,A_eD_ef\rangle.
\]
Apply \Cref{lem:s7v4-compression}, followed by the exact identity
\eqref{eq:s7v4-form-identity}.
\end{proof}

\begin{corollary}[Global normalized form debit]
\label{cor:s7v4-global}
Summing the middle line of \eqref{eq:s7v4-projected} over source blocks
gives
\begin{equation}
 (1+\theta^{-1})(1+\alpha){60\over89}
 \sum_pQ_p(f).                                          \label{eq:s7v4-global-debit}
\end{equation}
If
\[
 \sum_pQ_p(f)\le\kappa{\cal E}(f)+r_0(f),
\]
its relative coefficient is strictly below one whenever
\begin{equation}
 (1+\theta^{-1})(1+\alpha){60\over89}\kappa<1.          \label{eq:s7v4-global-window}
\end{equation}
\end{corollary}

\begin{proof}
Apply \eqref{eq:s7v4-global-form} and then the assumed incidence bound.
\end{proof}

\begin{assessment}[The precise new debit]
\label{ass:s7v4-debit}
The desired \(1/178\) normalization and source-independent plaquette
forms now hold exactly.  Noncommutation does not destroy them; it creates
the last line of \eqref{eq:s7v4-projected}.  The remaining projection
problem is therefore the positive operator \(L_e=E_e(I-P_e)E_e\), not an
unspecified commutator.
\end{assessment}

\section{A two-variable diagnostic}
\label{sec:s7v4-binary}

\begin{proposition}[Leakage can be order one]
\label{prop:s7v4-binary}
Let \(x,y\in\{-1,+1\}\) have probability
\[
 \mu_\beta(x,y)
 ={e^{\beta xy}\over4\cosh\beta}.
\]
Let \(E_y\) be conditional expectation onto functions of \(y\), and
\(K_x\) conditional expectation onto functions of \(x\).  Put
\(t=\tanh\beta\).  Then on the centered linear span of \(x,y\),
\begin{align}
 \|[E_y,K_x]\|_{2\to2}
 &=|t|\sqrt{1-t^2},                                     \label{eq:s7v4-binary-comm}\\
 \|E_y(I-K_x)E_y\|_{2\to2}
 &=1-t^2=\operatorname{sech}^2\beta.                   \label{eq:s7v4-binary-leak}
\end{align}
Thus owner/source projections need not commute, and the leakage is not
uniformly small over all coupling strengths.
\end{proposition}

\begin{proof}
The conditional means are
\[
 E_yx=ty,\qquad K_xy=tx,
\]
while \(E_yy=y\) and \(K_xx=x\).  On centered functions,
\(E_y\) and \(K_x\) are rank-one orthogonal projections onto the unit
vectors \(y\) and \(x\), whose inner product is \(t\).  The commutator of
two rank-one projections at angle cosine \(t\) has norm
\(|t|\sqrt{1-t^2}\).  Also
\[
 E_y(I-K_x)E_y\,y
 =E_y(y-tx)
 =(1-t^2)y,
\]
and the operator vanishes on the orthogonal complement of \(y\) in the
centered space.  This proves both identities.
\end{proof}

\begin{remark}[Concentration may help, but must be proved]
\label{rem:s7v4-binary}
In the binary model the leakage decays as \(4e^{-2|\beta|}\) at strong
alignment.  This suggests that the V95 alignment margin could control
\(\ell_e\) on a good Wilson sector.  The formula is only a diagnostic:
the Yang--Mills owner/source overlap has many group variables and a
physical ground transform, so no binary estimate may be imported without
a comparison theorem.
\end{remark}

\section{The randomized transfer card}
\label{sec:s7v4-card}

\begin{definition}[Randomized plaquette transfer card, RPTC]
\label{def:s7v4-RPTC}
RPTC consists of:
\begin{enumerate}[label=\textup{(R\arabic*)},leftmargin=4em]
\item source-independent positive self-adjoint owner transfers \(K_p\);
\item the restricted averaged norm
      \(\eta_e=\|M_{B_e}P_eD_e\|^2\), uniformly small on the declared
      bad boundary event;
\item physical form incidence
      \(\sum_pQ_p\le\kappa{\cal E}+r_0\) satisfying
      \eqref{eq:s7v4-global-window};
\item a leakage estimate for \(L_e=E_e(I-P_e)E_e\), either relative to
      the source variance/form or summable after block overlap;
\item response-ground stability for the same \(K_p\); and
\item translated-block approximate tensorization with all remainder
      terms summable.
\end{enumerate}
\end{definition}

\begin{theorem}[RPTC populates PNCC without a hidden union bound]
\label{thm:s7v4-RPTC}
Rows (R1)--(R4) give the PNCC normalized capacity and shell-summation
rows through \eqref{eq:s7v4-projected}.  If (R5)--(R6) also hold, RPTC
supplies the remaining BWIC rows and the f6 conditional-gap compiler
applies.
\end{theorem}

\begin{proof}
Row (R1) gives
\Cref{thm:s7v4-average}; (R2) and (R4) control the first and last lines
of \eqref{eq:s7v4-projected}; (R3) and
\Cref{cor:s7v4-global} give strict normalized form absorption.  No
maximum-event indicator is converted to an average.  Rows (R5)--(R6)
are exactly the remaining response and tensorization entries.
\end{proof}

\begin{decision}[V4 bottleneck decision]
\label{dec:s7v4-bottleneck}
The operator-level \(1/178\) normalization is now exact.  The next
bottleneck is the source-mean leakage \(L_e\), together with the
restricted norm \(\eta_e\).  V5 is a compulsory companion audit.  It will
also test whether alignment of the complete good incident action yields a
uniform estimate
\[
 L_e\le\ell_{\rm good}E_e+\text{bad-boundary debit},
\qquad \ell_{\rm good}<1,
\]
or whether nested-block martingales are required.
\end{decision}

\begin{programme}[Seventh-series V5]
\label{prog:s7v4-V5}
V5 will read the current active SM and NS notes, compare their moving
carrier and conditional transfer identities with \(L_e\), and derive the
best non-circular leakage decomposition available from the V1 complete
good-sector Hessian.
\end{programme}

\begin{firewall}[Claim boundary after seventh-series V4]
\label{fire:s7v4-boundary}
V4 constructs the normalized randomized transfer, proves its exact
source-independent Dirichlet identity, derives projected smoothing with
an explicit positive leakage debit, and identifies the precise strict
form window.

V4 does not bound the Yang--Mills leakage or restricted norm, prove
response-ground stability or tensorization, remove the cutoff, construct
the Osterwalder--Schrader continuum, or prove the Yang--Mills mass gap.
\end{firewall}

\paragraph{Parent-version record.}
V4 was formed from
\texttt{Renewal\_Yang\_Mills\_Mass\_Gap\_Research\_Notes\_V3.tex}.
The V3 parent had \(1\,420\) logical source lines, \(54\,792\) bytes,
and SHA--256
\[
 \texttt{6E0152C18AC0303761317949073877BAB7217F8E6C108E7EE9C0A272414F34C3}.
\]

% =============================================================================
% END APPENDED SEVENTH-SERIES V4 MATERIAL
% =============================================================================

% =============================================================================
% BEGIN APPENDED SEVENTH-SERIES V5 MATERIAL
% =============================================================================

\chapter{V5 companion audit and martingale control of leakage}
\label{ch:s7v5}

V4 reduced the source/owner mismatch to the positive leakage form
\(\langle E_ef,A_eE_ef\rangle\).  Bounding its full \(L^2\) operator norm
is too strong: the source exterior contains arbitrary functions of
variables which an owner transfer may resample.  The compulsory companion
audit points to the correct replacement.  SM V45--V46 uses an
owner/influence resolvent with operator Schur bounds.  Applying the same
valid functional analysis to an orthogonal martingale resolution turns
the collection of leakage forms into one auditable influence matrix.

\section{Compulsory SM/NS audit}
\label{sec:s7v5-audit}

\begin{audit}[Current companion files read at V5]
\label{audit:s7v5-files}
The active files are:
\begin{enumerate}[label=\textup{(\roman*)},leftmargin=3.8em]
\item
\texttt{Renewal\_SM\_Einstein\_Continuum\_Research\_Notes\_V46.tex},
with \(18\,989\) logical source lines, \(685\,323\) bytes, one live
terminator, and SHA--256
\[
 \texttt{F99DD3A359FE2AF2E43A3E94F018A7A66A4E2763AC995D009178EF45C4738698};
\]
\item
\texttt{Renewal\_NS\_Stability\_Research\_Notes\_V26.tex},
with \(5\,319\) logical source lines, \(278\,283\) bytes, one live
terminator, and SHA--256
\[
 \texttt{82C887F8C2C69E4F28FAD7C3DC8DE5B9099CB5A4BF431EBA1FFF3E91935978AD}.
\]
\end{enumerate}
The SM file has advanced from V39 at the preceding YM audit to V46.  The
NS file is unchanged.
\end{audit}

\begin{theorem}[Transferable conclusions of the V5 audit]
\label{thm:s7v5-audit}
The audit yields the following precise guidance.
\begin{enumerate}[leftmargin=3em]
\item SM V40 proves that a finite curvature floor controls a transverse
      configuration tangent but does not imply positive Hamiltonian
      spectral order.  Thus the V100 collar floor cannot be used as a
      response-ground estimate.
\item SM V41--V44 replace divergent global vacuum-vector response, where
      appropriate, by local or bounded--trace observable response.  This
      is relevant to response rows, but not a substitute for controlling
      all physical martingale fluctuations in a gap proof.
\item SM V45 proves the Neumann boundary resolvent and the
      operator-valued Schur estimate
      \(\|K\|\le\sqrt{r_*c_*}\).  It explicitly confirms that finite
      shell ownership does not prove strictness.
\item SM V46 gives Gaussian and polynomial-growth sufficient conditions
      for strict boundary influence, conditional on a decaying restricted
      bad-event norm.  It does not prove that norm; it warns that an
      unconditional probability cannot replace the essential-supremum
      conditional capacity.
\item NS V26 supplies no new boundary-capacity, martingale, or physical
      ground-response estimate.
\end{enumerate}
\end{theorem}

\begin{proof}
The statements are the terminal result ledgers and firewalls of SM
V40--V46 fixed by \Cref{audit:s7v5-files}.  In particular, the SM V45
resolvent assumes \(\sqrt{r_*c_*}<1\), and the V46 one-step kernel assumes
the restricted norm \(\eta_*\).  Neither assumption is supplied by the
companion audit for Yang--Mills.  The NS conclusion follows from its
unchanged digest and V26 ledger.
\end{proof}

\begin{assessment}[Audit-guided change]
\label{ass:s7v5-audit}
The Schur/resolvent mechanism is reusable; the SM physical capacity is
not.  We therefore apply the Schur mechanism to the exact V4 leakage
operators and leave their Yang--Mills decay as a separately named row.
\end{assessment}

\section{Why the full leakage norm is generally the wrong target}
\label{sec:s7v5-no-go}

\begin{theorem}[Decoupled-spectator leakage obstruction]
\label{thm:s7v5-spectator}
Positivity, self-adjointness, and the Markov property do not imply
\[
 \|E(I-K)E\|<1
\]
when \(E\) and \(K\) use different update fibres.  In fact, there is a
product probability carrier and two conditional expectations for which
\begin{equation}
 \|E(I-K)E\|=1.                                         \label{eq:s7v5-spectator}
\end{equation}
\end{theorem}

\begin{proof}
Let
\[
 (\Omega,\mu)
 =(\Omega_x,\mu_x)\times(\Omega_y,\mu_y)\times
   (\Omega_z,\mu_z),
\]
where \(\Omega_y\) supports a nonzero centered
\(h\in L^2(\mu_y)\).  Let \(E\) average only \(x\), so its range contains
all functions of \(y,z\).  Let \(K\) average only \(y\).  Both are
positive self-adjoint Markov projections.  Since \(h\) is independent of
\(x\) and centered in \(y\),
\[
 Eh=h,\qquad Kh=0.
\]
After normalizing \(h\),
\[
 \langle h,E(I-K)Eh\rangle=1.
\]
The operator is a positive contraction, so its norm is both at least and
at most one, proving \eqref{eq:s7v5-spectator}.
\end{proof}

\begin{corollary}[No abstract full-space strictness]
\label{cor:s7v5-no-abstract}
The V4 leakage coefficient \(\ell_e\) cannot be made strictly below one
from the abstract transfer axioms \eqref{eq:s7v4-K}, finite shell
incidence, or ownership alone.
\end{corollary}

\begin{proof}
The construction in \Cref{thm:s7v5-spectator} satisfies all those
operator axioms and is a one-contact finite-incidence example.
\end{proof}

\begin{firewall}[What the spectator example does not show]
\label{fire:s7v5-spectator}
The example does not prove that the physical Yang--Mills leakage equals
one.  Wilson alignment may suppress the action of owner updates on the
specific fluctuations generated by a source block.  It proves that such
suppression must use dynamics, locality, a restricted fluctuation space,
or a martingale structure; it is not a formal consequence of Markov
positivity.
\end{firewall}

\section{Orthogonal martingale resolution}
\label{sec:s7v5-martingale}

Let \(\Pi_0 f=\int f\,d\mu\).  Choose a finite filtration or orthogonal
block resolution with mutually orthogonal projections
\(\Delta_j\), \(j\in{\cal J}\), satisfying
\begin{equation}
 I-\Pi_0=\sum_{j\in{\cal J}}\Delta_j,\qquad
 \Delta_i\Delta_j=0\quad(i\ne j).                       \label{eq:s7v5-resolution}
\end{equation}
The sum is exact at a finite regulator.  For source block \(e\), define
\begin{equation}
 T_{ej}
 :=A_e^{1/2}E_e\Delta_j:
 \operatorname{Ran}\Delta_j\longrightarrow L^2(\mu).    \label{eq:s7v5-Tej}
\end{equation}

\begin{lemma}[Exact martingale representation of leakage]
\label{lem:s7v5-representation}
For every \(f\in L^2(\mu)\),
\begin{equation}
 A_e^{1/2}E_ef
 =\sum_{j\in{\cal J}}T_{ej}\Delta_jf,                   \label{eq:s7v5-leak-representation}
\end{equation}
and
\begin{equation}
 \sum_j\|\Delta_jf\|_2^2
 =\|f-\Pi_0f\|_2^2=\operatorname{Var}_\mu(f).           \label{eq:s7v5-orthogonality}
\end{equation}
\end{lemma}

\begin{proof}
The Markov property gives \(A_e1=0\), hence
\(A_e^{1/2}E_e\Pi_0=0\).  Insert
\eqref{eq:s7v5-resolution} into \(f-\Pi_0f\) to obtain
\eqref{eq:s7v5-leak-representation}.  Orthogonality of the ranges of
\(\Delta_j\) gives Parseval's identity
\eqref{eq:s7v5-orthogonality}.
\end{proof}

Define the operator-valued Schur constants
\begin{equation}
 r_T:=\sup_e\sum_j\|T_{ej}\|,
 \qquad
 c_T:=\sup_j\sum_e\|T_{ej}\|.                           \label{eq:s7v5-Schur-data}
\end{equation}

\begin{theorem}[Martingale leakage compiler]
\label{thm:s7v5-compiler}
If \(r_T,c_T<\infty\), then
\begin{equation}
 \boxed{
 \sum_e\langle E_ef,A_eE_ef\rangle
 \le r_Tc_T\,\operatorname{Var}_\mu(f).}                \label{eq:s7v5-leakage-bound}
\end{equation}
Thus a collective leakage estimate can be strict even when some
individual full-space norms
\(\|E_eA_eE_e\|\) equal one.
\end{theorem}

\begin{proof}
Let
\[
 {\cal T}:\bigoplus_{j\in{\cal J}}\operatorname{Ran}\Delta_j
 \longrightarrow\bigoplus_eL^2(\mu),
 \qquad
 ({\cal T}x)_e=\sum_jT_{ej}x_j.
\]
The operator-valued Schur argument gives
\[
 \|{\cal T}\|\le\sqrt{r_Tc_T}.
\]
Apply this to \(x_j=\Delta_jf\).  By
\Cref{lem:s7v5-representation},
\[
 \|{\cal T}x\|^2
 =\sum_e\|A_e^{1/2}E_ef\|_2^2
 =\sum_e\langle E_ef,A_eE_ef\rangle.
\]
Use \eqref{eq:s7v5-orthogonality} to obtain
\eqref{eq:s7v5-leakage-bound}.
\end{proof}

\begin{corollary}[Projected smoothing with collective leakage]
\label{cor:s7v5-projected}
Summing the last line of the V4 inequality
\eqref{eq:s7v4-projected} over source blocks gives the relative variance
debit
\begin{equation}
 (1+\theta^{-1})(1+\alpha^{-1})r_Tc_T\,
 \operatorname{Var}_\mu(f).                             \label{eq:s7v5-projective-debit}
\end{equation}
The middle line retains the normalized \(60/89\) form factor from
\eqref{eq:s7v4-global-debit}.
\end{corollary}

\begin{proof}
Apply \Cref{thm:s7v5-compiler} to the leakage terms in the sum of
\eqref{eq:s7v4-projected}.  The plaquette-form term is already summed in
\Cref{cor:s7v4-global}.
\end{proof}

\begin{remark}[Budget interaction]
\label{rem:s7v5-budget}
The coefficient in \eqref{eq:s7v5-projective-debit} is not required to
be below one in isolation; it must fit, together with the restricted
averaged norm, normalized plaquette form, tensorization, and all tunable
square factors, inside the final strict variance budget.  Writing it
explicitly prevents it from being counted twice as both a commutator and
a boundary response.
\end{remark}

\section{Decay criteria and the audited SM comparison}
\label{sec:s7v5-decay}

\begin{theorem}[Gaussian martingale leakage criterion]
\label{thm:s7v5-Gaussian}
Suppose block and martingale indices carry a graph distance and, for
some \(a,\tau>0\),
\begin{equation}
 \|T_{ej}\|\le\tau e^{-a d(e,j)^2}.                      \label{eq:s7v5-Gaussian}
\end{equation}
If both index families have sphere growth
\[
 N(r)\le C_{\rm g}(1+r)^{q-1},
\]
then
\begin{equation}
 r_T,c_T
 \le
 \tau C_{\rm g}C_q(1+a^{-q/2}),                         \label{eq:s7v5-Gaussian-Schur}
\end{equation}
and the collective leakage coefficient in
\eqref{eq:s7v5-leakage-bound} is at most the square of the right side.
\end{theorem}

\begin{proof}
For fixed \(e\), sum \eqref{eq:s7v5-Gaussian} over distance spheres and
compare
\[
 \sum_{r\ge0}(1+r)^{q-1}e^{-ar^2}
\]
with the Gaussian integral after the change of variables
\(u=\sqrt a\,r\).  This gives the row bound.  The same hypothesis with
the roles of the two index families interchanged gives the column bound.
Apply \Cref{thm:s7v5-compiler}.
\end{proof}

\begin{assessment}[What SM V45--V46 contributes]
\label{ass:s7v5-SM}
SM V45 supplies exactly the Schur/Neumann functional analysis used in
\Cref{thm:s7v5-compiler}; SM V46 supplies the analogous polynomial-growth
Gaussian summation.  Neither companion result proves
\eqref{eq:s7v5-Gaussian} for
\[
 T_{ej}=A_e^{1/2}E_e\Delta_j
\]
on the Yang--Mills carrier.  In particular, its assumed factor
\(\sqrt{\eta_*}\) is the same kind of physical strictness input which
remains open here.
\end{assessment}

\begin{definition}[Martingale leakage card, MLC]
\label{def:s7v5-MLC}
MLC consists of:
\begin{enumerate}[label=\textup{(M\arabic*)},leftmargin=4em]
\item an explicit gauge-compatible filtration and the orthogonal
      resolution \eqref{eq:s7v5-resolution};
\item localization of \(A_e^{1/2}E_e\Delta_j\) in the source/owner block
      graph;
\item row and column sums satisfying a strict final variance budget;
\item the V4 restricted averaged norm and normalized plaquette-form
      incidence;
\item compatibility with the physical gauge quotient and the
      response-ground transfer; and
\item a cutoff-uniform construction with summable boundary and regulator
      remainders.
\end{enumerate}
\end{definition}

\begin{theorem}[MLC replaces a full-space leakage norm]
\label{thm:s7v5-MLC}
Rows (M1)--(M4) replace the V4 pointwise leakage coefficient
\(\ell_e\) by the collective coefficient \(r_Tc_T\) in the global
projected-smoothing budget.  If (M5)--(M6) and the remaining RPTC rows
hold, the BWIC conditional-gap handoff applies.
\end{theorem}

\begin{proof}
Rows (M1)--(M3) give
\Cref{thm:s7v5-compiler}; row (M4) supplies the other two terms in
\eqref{eq:s7v4-projected}.  Rows (M5)--(M6) are precisely the physical
and cofinal qualifications not contained in the Schur estimate.
\end{proof}

\begin{decision}[V5 bottleneck decision]
\label{dec:s7v5-bottleneck}
The full leakage norm is abandoned as structurally too strong.
The live analytic target is now the off-diagonal martingale estimate
\[
 \|A_e^{1/2}E_e\Delta_j\|
 \le\tau_s e^{-a_s d(e,j)^2}
\]
with row/column budget compatible with V4.  V6 will first compute this
operator in a finite Gaussian surrogate.  If its zero-frequency mode does
not decay, the filtration must be changed before attempting the
nonlinear Yang--Mills carrier.
\end{decision}

\begin{programme}[Seventh-series V6]
\label{prog:s7v5-V6}
V6 will diagonalize the martingale leakage operator for a translation
invariant Gaussian nearest-neighbour field, distinguish massless and
massive symbols, and determine whether local owner averaging produces
the required spatial decay or retains a noncontracting infrared mode.
\end{programme}

\begin{firewall}[Claim boundary after seventh-series V5]
\label{fire:s7v5-boundary}
V5 performs the compulsory SM/NS audit; proves that a full-space leakage
norm can equal one; gives an exact orthogonal-martingale representation
of total leakage; and proves its operator-valued Schur and Gaussian
compilers.

V5 does not establish the Yang--Mills martingale decay, restricted norm,
response-ground stability, tensorization, continuum construction, or
mass gap.
\end{firewall}

\paragraph{Parent-version record.}
V5 was formed from
\texttt{Renewal\_Yang\_Mills\_Mass\_Gap\_Research\_Notes\_V4.tex}.
The V4 parent had \(1\,786\) logical source lines, \(67\,950\) bytes,
and SHA--256
\[
 \texttt{0C01D81E7518468E0D6F76A028EB942CD9281212D1EAA05F7A9A73A5CDA11CED}.
\]

% =============================================================================
% END APPENDED SEVENTH-SERIES V5 MATERIAL
% =============================================================================

% =============================================================================
% BEGIN APPENDED SEVENTH-SERIES V6 MATERIAL
% =============================================================================

\chapter{Gaussian infrared test of martingale leakage}
\label{ch:s7v6}

V5 reduced the leakage problem to decay of
\(T_{ej}=A_e^{1/2}E_e\Delta_j\).  Before seeking such a bound in the
nonlinear Wilson carrier, this chapter tests the generic infrared
mechanism in a translation-invariant Gaussian surrogate.  The result is
sharp: locality of the source stencil is not sufficient in a massless
field.  A square-root response requires a first-order zero at momentum
zero, while a full inverse requires a second-order zero.

\section{Fourier model}
\label{sec:s7v6-model}

Let
\[
 {\mathbb T}_L^d=(\mathbb Z/L\mathbb Z)^d
\]
with normalized counting measure, and let \(Q_0\) project away the
constant mode.  The nearest-neighbour nonnegative Laplacian has Fourier
symbol
\begin{equation}
 \lambda(k)
 =4\sum_{\nu=1}^d\sin^2(k_\nu/2),
 \qquad
 k\in{2\pi\over L}\{0,\ldots,L-1\}^d.                  \label{eq:s7v6-lambda}
\end{equation}
Set
\[
 H_m=m^2-\Delta,\qquad m\ge0.
\]
Let \(S\) be a finite-range translation-invariant scalar convolution
with trigonometric-polynomial symbol \(s(k)\).  For
\(\gamma>0\), define
\begin{equation}
 T_{m,\gamma}
 :=S H_m^{-\gamma}Q_0.                                  \label{eq:s7v6-T}
\end{equation}

\begin{theorem}[Exact Gaussian response norm]
\label{thm:s7v6-norm}
\begin{equation}
 \boxed{
 \|T_{m,\gamma}\|_{2\to2}
 =
 \max_{k\ne0}
 {|s(k)|\over(m^2+\lambda(k))^\gamma}.}                 \label{eq:s7v6-norm}
\end{equation}
For \(m>0\),
\begin{equation}
 \|T_{m,\gamma}\|
 \le m^{-2\gamma}\|S\|_{2\to2}.                         \label{eq:s7v6-massive}
\end{equation}
\end{theorem}

\begin{proof}
The discrete Fourier transform is unitary.  On each nonzero Fourier mode,
\(S\), \(H_m\), and \(Q_0\) are simultaneous scalar multipliers with
values \(s(k)\), \(m^2+\lambda(k)\), and \(1\), respectively.  The norm
of a diagonal multiplier is the maximum modulus of its multiplier,
proving \eqref{eq:s7v6-norm}.  Since
\(m^2+\lambda(k)\ge m^2\), the massive bound follows.
\end{proof}

\section{Necessary cancellation order}
\label{sec:s7v6-order}

\begin{theorem}[Cofinal infrared order test]
\label{thm:s7v6-order}
Suppose the first nonzero Taylor term of \(s\) at zero is a homogeneous
polynomial \(P_q\) of order \(q\ge0\).  Choose a nonzero integer vector
\(n\in\mathbb Z^d\) with \(P_q(n)\ne0\); such a vector exists because a
nonzero polynomial cannot vanish on all of \(\mathbb Z^d\).  Then
\begin{equation}
 s(tn)=t^qP_q(n)+O(t^{q+1}).                            \label{eq:s7v6-Taylor}
\end{equation}
Then on the massless tori,
\begin{equation}
 \|T_{0,\gamma}\|
 \ge c\,L^{2\gamma-q}                                  \label{eq:s7v6-growth}
\end{equation}
for all sufficiently large \(L\).  In particular:
\begin{enumerate}[label=\textup{(\roman*)},leftmargin=3.5em]
\item if \(q<2\gamma\), the norms are not cutoff uniform;
\item a square-root response \(\gamma=1/2\) requires at least a
      first-order zero;
\item a full inverse \(\gamma=1\) requires at least a second-order zero.
\end{enumerate}
\end{theorem}

\begin{proof}
For \(k\to0\),
\[
 \lambda(k)=|k|^2+O(|k|^4).
\]
Choose the lattice momenta \(k_L=2\pi n/L\), represented in the
principal Brillouin zone for all sufficiently large \(L\).  Then
\(|k_L|\asymp L^{-1}\).  By
\eqref{eq:s7v6-Taylor},
\[
 |s(k_L)|\ge c_1L^{-q},
 \qquad
 \lambda(k_L)^\gamma\le c_2L^{-2\gamma}
\]
for all large \(L\).  Insert these estimates in
\eqref{eq:s7v6-norm}.  The three conclusions follow.
\end{proof}

\begin{corollary}[A local contact can still be infrared divergent]
\label{cor:s7v6-local}
Finite range of \(S\) does not imply a uniform massless response.  If
\(s(0)\ne0\), then
\[
 \|T_{0,\gamma}\|\gtrsim L^{2\gamma}.
\]
If \(s(0)=0\) but \(Ds(0)\ne0\), then the full inverse
\(\gamma=1\) still has norm at least of order \(L\).
\end{corollary}

\begin{proof}
Apply \Cref{thm:s7v6-order} with \(q=0\) and \(q=1\),
respectively.
\end{proof}

\begin{firewall}[Spatial decay is not automatic at zero mass]
\label{fire:s7v6-decay}
A finite-range source followed by a massless inverse is generally
long-ranged.  Therefore the Gaussian Schur hypothesis
\eqref{eq:s7v5-Gaussian} cannot be inferred from locality of the
unpropagated owner contact.  A Ward, divergence, curvature, or martingale
cancellation must remove the low-momentum singularity.
\end{firewall}

\section{Sufficient factorizations}
\label{sec:s7v6-factor}

\begin{theorem}[Symbol factorization criterion]
\label{thm:s7v6-factor}
If, uniformly in \(L\),
\begin{equation}
 |s(k)|\le\kappa\,\lambda(k)^\gamma
 \qquad(k\ne0),                                         \label{eq:s7v6-factor}
\end{equation}
then
\begin{equation}
 \|T_{0,\gamma}\|_{2\to2}\le\kappa.                     \label{eq:s7v6-factor-bound}
\end{equation}
The response is strictly contractive whenever \(\kappa<1\).
\end{theorem}

\begin{proof}
Substitute \eqref{eq:s7v6-factor} into the exact multiplier formula
\eqref{eq:s7v6-norm}.
\end{proof}

\begin{proposition}[Exact first- and second-order examples]
\label{prop:s7v6-examples}
Let \(\nabla\) be the forward lattice gradient.  Then
\begin{align}
 \|\nabla(-\Delta)^{-1/2}Q_0\|_{\ell^2\to\ell^2(\mathbb C^d)}
 &=1,                                                    \label{eq:s7v6-gradient}\\
 \|(-\Delta)(-\Delta)^{-1}Q_0\|_{2\to2}
 &=1.                                                    \label{eq:s7v6-Laplacian}
\end{align}
Multiplication of the sources by \(\kappa\) changes the two norms to
\(|\kappa|\).
\end{proposition}

\begin{proof}
The \(\nu\)-th gradient multiplier is \(e^{ik_\nu}-1\), and
\[
 \sum_{\nu=1}^d|e^{ik_\nu}-1|^2
 =4\sum_{\nu=1}^d\sin^2(k_\nu/2)
 =\lambda(k).
\]
Thus the squared vector multiplier in
\eqref{eq:s7v6-gradient} is one on every nonzero mode.
The multiplier in \eqref{eq:s7v6-Laplacian} is also one on every
nonzero mode.  The scaled statements are immediate.
\end{proof}

\begin{remark}[Massive versus massless mechanisms]
\label{rem:s7v6-mass}
At fixed \(m>0\), \eqref{eq:s7v6-massive} controls every bounded local
source, with exponential position-space decay.  That estimate degenerates
as \(m\downarrow0\).  The factorizations in
\Cref{thm:s7v6-factor,prop:s7v6-examples} remain uniform because the
source cancels the same low-frequency power that the response would
otherwise invert.
\end{remark}

\section{Implication for the Yang--Mills leakage target}
\label{sec:s7v6-YM}

\begin{definition}[Gaussian leakage factorization card, GLFC]
\label{def:s7v6-GLFC}
For a linearized, translation-invariant surrogate of
\(A_e^{1/2}E_e\Delta_j\), GLFC consists of:
\begin{enumerate}[label=\textup{(G\arabic*)},leftmargin=4em]
\item identification of the massless response power
      \(\gamma\in\{1/2,1\}\);
\item a source symbol with a zero of order at least \(2\gamma\);
\item a quantitative factorization
      \(|s(k)|\le\kappa\lambda(k)^\gamma\);
\item a strict post-overlap constant compatible with the V4--V5 budget;
\item a position-space row/column estimate for the residual multiplier;
      and
\item control of nonlinear, gauge, boundary-frame, and ground-transform
      remainders.
\end{enumerate}
\end{definition}

\begin{theorem}[Gaussian decision for MLC]
\label{thm:s7v6-decision}
In a massless Gaussian surrogate, MLC cannot follow from finite range of
the owner contact alone.  GLFC rows (G1)--(G5) are sufficient for the
linearized martingale leakage Schur row.  Row (G6) is necessary before
the conclusion can be transferred to the nonlinear physical
Yang--Mills carrier.
\end{theorem}

\begin{proof}
The failure of locality alone is
\Cref{cor:s7v6-local,fire:s7v6-decay}.  Rows (G1)--(G3) give the uniform
Fourier bound by \Cref{thm:s7v6-factor}; (G4)--(G5) give the strict Schur
constants required in \Cref{thm:s7v5-compiler}.  The Gaussian
calculation contains none of the remainders listed in (G6), so that row
cannot be omitted in the nonlinear handoff.
\end{proof}

\begin{assessment}[Connection with the companion audit]
\label{ass:s7v6-SM}
The cancellation order is the same structural issue isolated by SM
V40--V43: a configuration-space curvature floor does not imply
Hamiltonian infrared order, while a genuine factorization by a positive
Hamiltonian power does.  Here the statement is used only as a Gaussian
test of leakage.  No SM physical-source factorization is imported into
Yang--Mills.
\end{assessment}

\begin{decision}[V6 bottleneck decision]
\label{dec:s7v6-bottleneck}
The Gaussian surrogate rejects a purely local leakage estimate and
selects a derivative factorization.  The next note will compute the
Fourier symbol of the lattice gauge differential.  On the transverse
one-form sector, the target identity is
\[
 d_1^*d_1=\lambda(k)I,
\]
which would give the exact square-root cancellation analogous to
\eqref{eq:s7v6-gradient}.  Gauge-longitudinal and boundary-frame modes
must be projected explicitly rather than discarded.
\end{decision}

\begin{programme}[Seventh-series V7]
\label{prog:s7v6-V7}
V7 will prove the periodic lattice Hodge-symbol identity for one-forms,
construct the transverse projector without dividing at the zero mode,
and determine precisely which part of the linearized collar/owner
leakage carries a curvature derivative.
\end{programme}

\begin{firewall}[Claim boundary after seventh-series V6]
\label{fire:s7v6-boundary}
V6 proves the exact Gaussian multiplier norm, the necessary infrared
cancellation order, and sufficient first- and second-order
factorizations.  It shows that local support alone cannot yield the V5
massless Schur estimate.

V6 does not prove that the Yang--Mills leakage has the required
curvature factor, control nonlinear remainders, establish the restricted
norm or response-ground rows, construct the continuum, or prove the
mass gap.
\end{firewall}

\paragraph{Parent-version record.}
V6 was formed from
\texttt{Renewal\_Yang\_Mills\_Mass\_Gap\_Research\_Notes\_V5.tex}.
The V5 parent had \(2\,173\) logical source lines, \(81\,984\) bytes,
and SHA--256
\[
 \texttt{A8787684E8BD964166A90F4C5EA61963C0638A63B11EEE48E469ADFEFE0E0C74}.
\]

% =============================================================================
% END APPENDED SEVENTH-SERIES V6 MATERIAL
% =============================================================================

% =============================================================================
% BEGIN APPENDED SEVENTH-SERIES V7 MATERIAL
% =============================================================================

\chapter{Lattice Hodge symbol and transverse square-root cancellation}
\label{ch:s7v7}

The V6 Gaussian test asks whether the owner/source leakage carries a
first-order derivative before a square-root massless response.  For the
linearized lattice gauge complex the answer is exact on every nonzero
transverse Fourier mode.  The curvature differential has modulus
\(\sqrt{\lambda(k)}\), while its kernel is precisely the longitudinal
gauge line.  This chapter proves the identity and records the harmonic
zero-mode qualification which a finite collar calculation cannot remove.

\section{Periodic lattice gauge complex}
\label{sec:s7v7-complex}

On \({\mathbb T}_L^d\), let \(C^0,C^1,C^2\) be complex-valued lattice
zero-, one-, and oriented two-cochains.  The color-valued result follows
by tensoring with the Lie algebra.  Put
\begin{equation}
 q_\mu(k):=e^{ik_\mu}-1,\qquad
 |q(k)|^2=\sum_{\mu=1}^d|q_\mu(k)|^2=\lambda(k).         \label{eq:s7v7-q}
\end{equation}
The forward coboundaries have Fourier symbols
\begin{align}
 (d_0\varphi)_\mu(k)
 &=q_\mu(k)\varphi(k),                                  \label{eq:s7v7-d0}\\
 (d_1a)_{\mu\nu}(k)
 &=q_\mu(k)a_\nu(k)-q_\nu(k)a_\mu(k),
 \qquad\mu<\nu.                                         \label{eq:s7v7-d1}
\end{align}

\begin{lemma}[Fourier Lagrange identity]
\label{lem:s7v7-Lagrange}
For \(q,a\in\mathbb C^d\),
\begin{equation}
 \sum_{\mu<\nu}|q_\mu a_\nu-q_\nu a_\mu|^2
 =|q|^2|a|^2-|q^*a|^2.                                 \label{eq:s7v7-Lagrange}
\end{equation}
\end{lemma}

\begin{proof}
Expand the left side.  The diagonal terms give
\[
 \sum_{\mu\ne\nu}|q_\mu|^2|a_\nu|^2
 =|q|^2|a|^2-\sum_\mu|q_\mu|^2|a_\mu|^2.
\]
The cross terms are
\[
 -\sum_{\mu\ne\nu}
   q_\mu\overline{q_\nu}\,
   a_\nu\overline{a_\mu}
 =
 -|q^*a|^2+\sum_\mu|q_\mu|^2|a_\mu|^2.
\]
Their sum is \eqref{eq:s7v7-Lagrange}.
\end{proof}

\begin{theorem}[Exact one-form Hodge symbol]
\label{thm:s7v7-Hodge}
At every momentum,
\begin{equation}
 d_1(k)^*d_1(k)+d_0(k)d_0(k)^*
 =\lambda(k)I_{C^1}.                                   \label{eq:s7v7-Hodge}
\end{equation}
Equivalently,
\begin{align}
 d_0(k)d_0(k)^*&=q(k)q(k)^*,                            \label{eq:s7v7-longitudinal}\\
 d_1(k)^*d_1(k)&=\lambda(k)I-q(k)q(k)^*.                \label{eq:s7v7-curvature}
\end{align}
\end{theorem}

\begin{proof}
The quadratic form of \(d_1^*d_1\) at \(a\) is the left side of
\eqref{eq:s7v7-Lagrange}.  The quadratic form of \(d_0d_0^*=qq^*\) is
\(|q^*a|^2\).  Their sum is
\(\lambda(k)|a|^2\).  Equality of the Hermitian quadratic forms gives the
operator identities.
\end{proof}

\section{Transverse and longitudinal modes}
\label{sec:s7v7-projectors}

For \(k\ne0\), define
\begin{equation}
 P_{\rm L}(k):={q(k)q(k)^*\over\lambda(k)},
 \qquad
 P_{\rm T}(k):=I-P_{\rm L}(k).                          \label{eq:s7v7-projectors}
\end{equation}

\begin{corollary}[Exact transverse curvature]
\label{cor:s7v7-transverse}
For every \(k\ne0\),
\begin{align}
 \operatorname{Ran}P_{\rm L}(k)
 &=\operatorname{Ran}d_0(k)=\ker d_1(k),                \label{eq:s7v7-long-space}\\
 \operatorname{Ran}P_{\rm T}(k)
 &=\ker d_0(k)^*,                                       \label{eq:s7v7-trans-space}\\
 d_1(k)^*d_1(k)
 &=\lambda(k)P_{\rm T}(k).                              \label{eq:s7v7-trans-curvature}
\end{align}
Thus the nonzero singular values of \(d_1(k)\) on transverse one-forms
are all \(\sqrt{\lambda(k)}\).
\end{corollary}

\begin{proof}
The matrices in \eqref{eq:s7v7-projectors} are the orthogonal projections
onto the line spanned by \(q(k)\) and its orthogonal complement.
Equation \eqref{eq:s7v7-d1} gives \(d_1q=0\), while
\eqref{eq:s7v7-Lagrange} shows that equality \(d_1a=0\) forces \(a\) to
be parallel to \(q\) when \(q\ne0\).  The remaining identities follow
from \eqref{eq:s7v7-curvature}.
\end{proof}

\begin{theorem}[Transverse square-root isometry]
\label{thm:s7v7-isometry}
On the nonzero Fourier modes,
\begin{equation}
 {\cal C}
 :=d_1(-\Delta)^{-1/2}P_{\rm T}                         \label{eq:s7v7-C}
\end{equation}
is an isometry from transverse one-forms into two-forms:
\begin{equation}
 \|{\cal C}a\|_2=\|P_{\rm T}Q_0a\|_2.                  \label{eq:s7v7-isometry}
\end{equation}
In particular,
\[
 \|{\cal C}\|_{2\to2}=1.
\]
\end{theorem}

\begin{proof}
At \(k\ne0\), use
\eqref{eq:s7v7-trans-curvature}:
\begin{align*}
 \|d_1(k)\lambda(k)^{-1/2}P_{\rm T}(k)a(k)\|^2
 &=
 \lambda(k)^{-1}
 \langle P_{\rm T}a,d_1^*d_1P_{\rm T}a\rangle\\
 &=\|P_{\rm T}(k)a(k)\|^2.
\end{align*}
Sum over Fourier modes.  The zero mode is removed by \(Q_0\).
\end{proof}

\begin{corollary}[Exact GLFC cancellation on the transverse Gaussian
sector]
\label{cor:s7v7-GLFC}
For a square-root response \(\gamma=1/2\), the curvature source
\(S=d_1P_{\rm T}\) satisfies the V6 factorization with coefficient one.
For a source \(\kappa d_1P_{\rm T}\), the response norm is
\(|\kappa|\).
\end{corollary}

\begin{proof}
This is \Cref{thm:s7v7-isometry}, or equivalently
\[
 \|d_1(k)P_{\rm T}(k)a\|^2
 =\lambda(k)\|P_{\rm T}(k)a\|^2.
\]
\end{proof}

\section{The harmonic and gauge qualifications}
\label{sec:s7v7-zero}

\begin{firewall}[The zero mode is not longitudinal gauge]
\label{fire:s7v7-zero}
At \(k=0\), one has \(q(0)=0\), so
\[
 d_0(0)=0,\qquad d_1(0)=0.
\]
All constant one-forms are harmonic.  The quotient by periodic
infinitesimal gauge transformations does not remove them, because
constant gauge parameters also have zero derivative.  They represent
linearized toron/global-holonomy directions and must be fixed by a
sector condition, boundary condition, or separate global argument.
Defining \(P_{\rm T}(0)\) by division by \(\lambda(0)\) would hide this
kernel.
\end{firewall}

\begin{firewall}[Finite collar rank is not the periodic Hodge gap]
\label{fire:s7v7-collar}
The f6 collar becomes injective after its full tree frame is conditioned,
with a dimension-only finite-complex floor.  The periodic global complex
has momenta approaching zero and a harmonic zero mode.  The collar floor
therefore cannot replace \(\lambda(k)\) in
\eqref{eq:s7v7-trans-curvature}; it controls a local Dirichlet problem,
not the global massless Hodge spectrum.
\end{firewall}

\begin{proposition}[Color multiplicity does not change the symbol]
\label{prop:s7v7-color}
At a flat \(SU(3)\) background, tensoring the complex with
\(\mathfrak{su}(3)\) repeats every scalar singular value eight times.
Equations \eqref{eq:s7v7-Hodge}--\eqref{eq:s7v7-isometry} remain valid
with an additional identity factor \(I_8\).
\end{proposition}

\begin{proof}
The linearized flat covariant derivative acts componentwise in a
parallel color frame.  Tensor products with \(I_8\) preserve products,
adjoints, projections, and operator norms.
\end{proof}

\section{The strict geometric slack and its limit}
\label{sec:s7v7-slack}

\begin{theorem}[Curvature isometry fits the normalized shell form
window]
\label{thm:s7v7-slack}
Suppose the source-independent plaquette debit in
\eqref{eq:s7v4-Qp} is identified, on the flat transverse Gaussian sector,
with the curvature square in
\Cref{thm:s7v7-isometry}.  Then its physical form-incidence coefficient
is \(\kappa=1\), while shell double counting contributes
\begin{equation}
 {60\over89}\kappa={60\over89}<1.                       \label{eq:s7v7-slack}
\end{equation}
The available arithmetic margin is \(29/89\).
\end{theorem}

\begin{proof}
The incidence coefficient one is the isometry
\eqref{eq:s7v7-isometry}.  Substitute \(\kappa=1\) in
\eqref{eq:s7v3-kappa-window}; since \(1<89/60\), the condition holds.
The resulting coefficient and margin are
\[
 {60\over89},\qquad1-{60\over89}={29\over89}.
\]
\end{proof}

\begin{warning}[Tunable squares and leakage spend the margin]
\label{warn:s7v7-margin}
The raw margin \(29/89\) is not yet a gap.  In the V4 projected inequality
the form factor is multiplied by
\((1+\theta^{-1})(1+\alpha)\), while the restricted averaged norm and
martingale leakage add separate debits.  Those quantities must be
optimized jointly and fit strictly inside the final compiler.  The
identity \eqref{eq:s7v7-slack} only proves that the geometric incidence
coefficient is not itself supercritical.
\end{warning}

\section{Hodge factorization card}
\label{sec:s7v7-card}

\begin{definition}[Transverse Hodge leakage card, THLC]
\label{def:s7v7-THLC}
THLC consists of:
\begin{enumerate}[label=\textup{(H\arabic*)},leftmargin=4em]
\item an explicit removal or conditioning of harmonic and boundary-frame
      modes;
\item the physical gauge quotient and transverse projector;
\item a factorization of the linearized martingale leakage source through
      \(d_1P_{\rm T}\);
\item identification of its square-root response with
      \eqref{eq:s7v7-C};
\item nonlinear stability of the covariant Hodge identity on the fully
      good incident set;
\item row/column localization of the remaining multiplier; and
\item response-ground and regulator compatibility.
\end{enumerate}
\end{definition}

\begin{theorem}[THLC would populate the Gaussian part of MLC]
\label{thm:s7v7-THLC}
Rows (H1)--(H4) give exact first-order infrared cancellation with
coefficient one.  Rows (H5)--(H6) would transfer this to the V5
martingale Schur estimate with perturbative constants.  Row (H7) remains
necessary for the physical conditional-gap handoff.
\end{theorem}

\begin{proof}
The exact cancellation is
\Cref{cor:s7v7-transverse,thm:s7v7-isometry}.  Harmonic and longitudinal
qualifications are
\Cref{fire:s7v7-zero,fire:s7v7-collar}.  Stability and localization are
exactly the nonlinear and position-space rows absent from the Fourier
identity.
\end{proof}

\begin{decision}[V7 bottleneck decision]
\label{dec:s7v7-bottleneck}
The transverse Gaussian symbol has the required square-root cancellation
and fits the \(60/89\) normalized incidence window.  The next bottleneck
is nonlinear: compare the covariant lattice differential at a
small-defect background with the flat Hodge complex while retaining the
conditioned tree frames.  The target is a quantitative perturbation of
\eqref{eq:s7v7-trans-curvature} whose error is absorbed by the V100
reconstruction radius.
\end{decision}

\begin{programme}[Seventh-series V8]
\label{prog:s7v7-V8}
V8 will derive the exact covariant coboundary, estimate its difference
from the flat coboundary after tree parallel transport, and prove a
near-flat transverse curvature lower bound or identify a boundary-frame
commutator that prevents it.
\end{programme}

\begin{firewall}[Claim boundary after seventh-series V7]
\label{fire:s7v7-boundary}
V7 proves the periodic lattice Hodge identity, the exact transverse
curvature symbol, and square-root isometry on nonzero modes.  It identifies
the harmonic zero mode and shows that curvature incidence coefficient one
fits inside the normalized \(60/89\) shell window.

V7 does not identify the physical leakage with this curvature source,
prove nonlinear covariant stability, control harmonic sectors,
response-ground motion or tensorization, construct the continuum, or
prove the Yang--Mills mass gap.
\end{firewall}

\paragraph{Parent-version record.}
V7 was formed from
\texttt{Renewal\_Yang\_Mills\_Mass\_Gap\_Research\_Notes\_V6.tex}.
The V6 parent had \(2\,465\) logical source lines, \(92\,452\) bytes,
and SHA--256
\[
 \texttt{49885FB13F3E0E1F1157623F14C4CB78CD08CEF895E5942ACFE52D95DCE5DB7D}.
\]

% =============================================================================
% END APPENDED SEVENTH-SERIES V7 MATERIAL
% =============================================================================

% =============================================================================
% BEGIN APPENDED SEVENTH-SERIES V8 MATERIAL
% =============================================================================

\chapter{Nonlinear collar polar factorization}
\label{ch:s7v8}

V7 proved the flat periodic identity
\(d_1^*d_1=\lambda P_{\rm T}\).  Its cancellation is more robust than
the Fourier representation suggests.  For every injective curvature
linearization \(B\), the polar factor
\(B(B^*B)^{-1/2}\) is an exact isometry.  This remains true at a
nonlinear near-flat collar background.  The small-defect geometry is
needed to preserve injectivity; no uniform global Hodge gap is needed.
The substantive remaining condition is that the physical leakage source
have a controlled curvature-adjoint factorization.

\section{Selected nonlinear curvature map}
\label{sec:s7v8-map}

Fix the V100 full-tree frame \(\tau\) and use its sixty chord variables.
Let
\[
 {\cal F}_\tau:SU(3)^{60}\longrightarrow\mathfrak{su}(3)^{60}
\]
be the selected plaquette-log map obtained from the sixty plaquette rows
in the unimodular V99 minor.  It is defined on the uniform neighbourhood
where the selected plaquette holonomies avoid the logarithm cut.  Write
\begin{equation}
 B_{\tau,U}:=D{\cal F}_\tau(U):
 \mathfrak{su}(3)^{60}\longrightarrow
 \mathfrak{su}(3)^{60}.                                 \label{eq:s7v8-B}
\end{equation}
At the unique flat completion \(U^{\rm fl}(\tau)\), tree parallel
transport identifies
\begin{equation}
 B_{\tau,U^{\rm fl}}=A\otimes I_8,                       \label{eq:s7v8-flat-B}
\end{equation}
where \(A\) is the integer minor with \(\det A=-1\).

\begin{lemma}[Uniform derivative perturbation]
\label{lem:s7v8-perturbation}
There are collar constants \(r_B,C_B>0\), independent of \(\tau\), such
that
\begin{equation}
 \|B_{\tau,U}-B_{\tau,U^{\rm fl}}\|
 \le C_B\dist(U,U^{\rm fl})                             \label{eq:s7v8-perturbation}
\end{equation}
whenever \(\dist(U,U^{\rm fl})<r_B\).
\end{lemma}

\begin{proof}
In exponential chord coordinates, every selected plaquette holonomy is
a product of a fixed number of group variables, inverses, and tree-frame
translations.  On a compact neighbourhood avoiding the logarithm cut,
this map has uniformly bounded second derivative.  The mean-value
theorem therefore bounds the difference of its first derivatives by
\(C_B\) times the chord distance.  Left and right tree transports are
isometries, and the tree-frame space \(SU(3)^{63}\) is compact, so one
finite subcover makes \(r_B,C_B\) uniform in \(\tau\).
\end{proof}

\begin{theorem}[Uniform nonlinear injectivity on a principal
small-defect core]
\label{thm:s7v8-injectivity}
Let
\[
 \sigma_0:=102^{-59/2}.
\]
There is \(a_8>0\), depending only on the collar, such that
\[
 {\cal V}_+(U)<a_8
\]
implies
\begin{equation}
 \sigma_{\min}(B_{\tau,U})\ge{\sigma_0\over2}.           \label{eq:s7v8-sigma}
\end{equation}
This is uniform in the finite volume, collar location, and fixed tree
frame.
\end{theorem}

\begin{proof}
The V99 determinant/Frobenius certificate gives
\[
 \sigma_{\min}(B_{\tau,U^{\rm fl}})^2\ge102^{-59},
\]
hence the flat singular value is at least \(\sigma_0\).  The V100 global
reconstruction theorem gives
\[
 \dist(U,U^{\rm fl})^2
 \le C_{\rm glob}{\cal V}_+(U)
\]
below a fixed principal threshold.  Choose \(a_8\) smaller than that
threshold and so that
\[
 C_B\sqrt{C_{\rm glob}a_8}\le{\sigma_0\over2},
\qquad
 \sqrt{C_{\rm glob}a_8}<r_B.
\]
Then \Cref{lem:s7v8-perturbation} and the singular-value perturbation
inequality
\[
 |\sigma_{\min}(B)-\sigma_{\min}(C)|\le\|B-C\|
\]
give \eqref{eq:s7v8-sigma}.
\end{proof}

\section{The polar curvature isometry}
\label{sec:s7v8-polar}

\begin{theorem}[Injective polar factorization]
\label{thm:s7v8-polar}
Let \(B:X\to Y\) be an injective operator between finite-dimensional
Hilbert spaces and put \(H=B^*B\).  Then
\begin{equation}
 V:=BH^{-1/2}                                             \label{eq:s7v8-V}
\end{equation}
is an isometry \(X\to Y\):
\begin{equation}
 V^*V=I_X,\qquad \|Vx\|=\|x\|.                          \label{eq:s7v8-isometry}
\end{equation}
Moreover,
\begin{equation}
 H^{-1/2}B^*=V^*,\qquad
 VV^*=P_{\operatorname{Ran}B}.                          \label{eq:s7v8-adjoint-polar}
\end{equation}
\end{theorem}

\begin{proof}
Since \(B\) is injective, \(H>0\).  Direct multiplication gives
\[
 V^*V=H^{-1/2}B^*BH^{-1/2}=I.
\]
Thus \(V\) is an isometry and \(VV^*\) is the orthogonal projection onto
its range, which equals \(\operatorname{Ran}B\).  The polar decomposition
\(B=VH^{1/2}\) gives \(B^*=H^{1/2}V^*\), proving the first identity in
\eqref{eq:s7v8-adjoint-polar}.
\end{proof}

\begin{corollary}[Nonlinear collar square-root cancellation]
\label{cor:s7v8-cancellation}
For every \(U\) in the core of \Cref{thm:s7v8-injectivity}, put
\[
 H_{\tau,U}:=B_{\tau,U}^*B_{\tau,U}.
\]
Then
\begin{equation}
 \|B_{\tau,U}H_{\tau,U}^{-1/2}x\|=\|x\|,                \label{eq:s7v8-forward-cancel}
\end{equation}
and, for every \(y\),
\begin{equation}
 \|H_{\tau,U}^{-1/2}B_{\tau,U}^*y\|
 =\|P_{\operatorname{Ran}B_{\tau,U}}y\|
 \le\|y\|.                                              \label{eq:s7v8-adjoint-cancel}
\end{equation}
The constant is exactly one and does not pay the very small certified
floor \(\sigma_0\).
\end{corollary}

\begin{proof}
Apply \Cref{thm:s7v8-polar} to \(B_{\tau,U}\), which is injective by
\Cref{thm:s7v8-injectivity}.
\end{proof}

\begin{remark}[Recovery of the flat Hodge identity]
\label{rem:s7v8-Hodge}
For \(B=d_1P_{\rm T}\) on a nonzero periodic Fourier mode,
\(B^*B=\lambda P_{\rm T}\).  The forward polar isometry in
\eqref{eq:s7v8-forward-cancel} is exactly the V7 operator
\(d_1(-\Delta)^{-1/2}P_{\rm T}\).  Thus the collar theorem is the
finite nonlinear form of the same square-root cancellation.
\end{remark}

\section{Controlled curvature-range sources}
\label{sec:s7v8-source}

\begin{theorem}[Curvature-adjoint source estimate]
\label{thm:s7v8-source}
Let a chord-cotangent source have the decomposition
\begin{equation}
 S=B_{\tau,U}^*Y+R.                                     \label{eq:s7v8-source}
\end{equation}
On the V8 core,
\begin{equation}
 \|H_{\tau,U}^{-1/2}S\|
 \le
 \|Y\|+{2\over\sigma_0}\|R\|.                           \label{eq:s7v8-source-bound}
\end{equation}
If \(R=0\), the estimate has constant one and is independent of the
curvature singular floor.
\end{theorem}

\begin{proof}
Use
\[
 H^{-1/2}S
 =H^{-1/2}B^*Y+H^{-1/2}R.
\]
The first term is bounded by \(\|Y\|\) using
\eqref{eq:s7v8-adjoint-cancel}.  By
\eqref{eq:s7v8-sigma},
\[
 \|H^{-1/2}\|
 ={1\over\sigma_{\min}(B_{\tau,U})}
 \le {2\over\sigma_0}.
\]
The triangle inequality proves
\eqref{eq:s7v8-source-bound}.
\end{proof}

\begin{firewall}[Formal range membership is circular]
\label{fire:s7v8-range}
Because \(B_{\tau,U}\) is injective and square, \(B_{\tau,U}^*\) is
surjective.  Hence every source can be written formally as
\[
 S=B_{\tau,U}^*Y,\qquad
 Y=(B_{\tau,U}^*)^{-1}S.
\]
This observation supplies no cancellation: the norm of this \(Y\) can
pay \(1/\sigma_{\min}(B_{\tau,U})\), exactly the inverse loss one sought
to avoid.  A useful factorization must produce \(Y\) directly from a
local Ward, curvature, or martingale identity with a cutoff-uniform norm
and locality bound.
\end{firewall}

\begin{definition}[Curvature-range leakage card, CRLC]
\label{def:s7v8-CRLC}
For every martingale leakage contact, CRLC consists of:
\begin{enumerate}[label=\textup{(C\arabic*)},leftmargin=4em]
\item the small-defect injectivity domain of
      \Cref{thm:s7v8-injectivity};
\item a decomposition
      \(S_{ej}=B_{\tau,U}^*Y_{ej}+R_{ej}\);
\item a direct local construction of \(Y_{ej}\), not a pseudoinverse
      definition;
\item row/column decay of \(Y_{ej}\) compatible with the V5 Schur
      budget;
\item residual bounds for \(R_{ej}\) strong enough to absorb the factor
      \(2/\sigma_0\);
\item transport between collar coordinates, the physical gauge quotient,
      and the conditional ground representation; and
\item control of the complement of the small-defect core by the V4
      randomized transfer.
\end{enumerate}
\end{definition}

\begin{theorem}[CRLC supplies the nonlinear square-root row]
\label{thm:s7v8-CRLC}
Rows (C1)--(C5) give a cutoff-uniform nonlinear square-root response
estimate for the martingale leakage sources through
\eqref{eq:s7v8-source-bound}.  With (C6)--(C7), the estimate enters MLC
and RPTC without paying the principal curvature inverse on the
curvature-adjoint component.
\end{theorem}

\begin{proof}
Rows (C1)--(C2) allow
\Cref{thm:s7v8-source}; (C3)--(C4) control its first term in the V5
operator-valued Schur sums; (C5) controls the inverse-weighted residual.
Rows (C6)--(C7) are exactly the coordinate, physical, and bad-set
contacts required to insert that estimate in the earlier compilers.
\end{proof}

\section{Gauge-orbit qualification away from flatness}
\label{sec:s7v8-gauge}

\begin{lemma}[Curvature changes by conjugation along a gauge orbit]
\label{lem:s7v8-gauge}
Let \(W_p(U)\) be a plaquette holonomy based at a vertex \(v\), and let
\(g_v(t)=\exp(t\xi_v)\) be a gauge path.  Then
\begin{equation}
 {d\over dt}\bigg|_{t=0}W_p(U^{g(t)})
 =[\xi_v,W_p(U)].                                       \label{eq:s7v8-gauge}
\end{equation}
Thus an infinitesimal gauge tangent lies in the kernel of the raw
holonomy derivative at a flat plaquette, but not necessarily at a
nonflat plaquette.
\end{lemma}

\begin{proof}
Gauge covariance of based holonomy gives
\[
 W_p(U^{g(t)})=g_v(t)W_p(U)g_v(t)^{-1}.
\]
Differentiate at zero.
\end{proof}

\begin{warning}[Do not identify the off-flat kernel with gauge]
\label{warn:s7v8-gauge}
The V99 identity
\(\ker B=\operatorname{im}d_0\) is a flat-background statement.
\Cref{lem:s7v8-gauge} shows that the kernel of a raw nonlinear curvature
derivative need not contain the gauge tangent away from flatness.  The
tree-conditioned chord map avoids this issue by fixing all frames and
being injective.  A physical gauge-invariant formulation must instead
transport the full action/observable or use a covariant quotient; it may
not define the physical space as \((\ker B_{\tau,U})^\perp\) without
proof.
\end{warning}

\begin{decision}[V8 bottleneck decision]
\label{dec:s7v8-bottleneck}
Nonlinear near-flat square-root cancellation is now exact once a
controlled curvature-adjoint source is present.  The remaining issue is
not the small singular floor but CRLC (C3): derive the leakage source
factorization directly.  V9 will differentiate a reversible
owner-plaquette transfer with respect to a source chord and determine
whether its score is a curvature-adjoint divergence plus a ground-score
residual.
\end{decision}

\begin{programme}[Seventh-series V9]
\label{prog:s7v8-V9}
V9 will prove the derivative formula for a conditional reversible
transfer, isolate the plaquette score, apply discrete integration by
parts to its curvature component, and retain every derivative of the
conditional normalization and ground transform as an explicit residual.
\end{programme}

\begin{firewall}[Claim boundary after seventh-series V8]
\label{fire:s7v8-boundary}
V8 proves uniform injectivity of the selected nonlinear collar curvature
map on a fixed small-defect core, its exact polar isometry, and the
curvature-adjoint source estimate.  It also corrects the off-flat
gauge-kernel interpretation.

V8 does not prove curvature-range factorization of the physical leakage
source, bound its ground-score residual, establish the remaining
restricted norm or tensorization, construct the continuum, or prove the
Yang--Mills mass gap.
\end{firewall}

\paragraph{Parent-version record.}
V8 was formed from
\texttt{Renewal\_Yang\_Mills\_Mass\_Gap\_Research\_Notes\_V7.tex}.
The V7 parent had \(2\,791\) logical source lines, \(104\,225\) bytes,
and SHA--256
\[
 \texttt{A0B367FEAF66B621487CCBD32920C4221A88DB3E16BE0A99260E9F440DEE0706}.
\]

% =============================================================================
% END APPENDED SEVENTH-SERIES V8 MATERIAL
% =============================================================================

% =============================================================================
% BEGIN APPENDED SEVENTH-SERIES V9 MATERIAL
% =============================================================================

\chapter{Conditional score differentiation and curvature-adjoint action
sources}
\label{ch:s7v9}

V8 shows that a controlled source of the form \(B_U^*Y\) receives exact
square-root cancellation.  This chapter derives the source generated by
a Wilson conditional density.  The raw Wilson action derivative is
exactly curvature adjoint and its polar preimage is bounded by the
plaquette defect.  Conditional normalization and the physical ground
density add separate terms.  An integration-by-parts identity cancels
those terms only for directions internal to the updated fibre; it does
not apply to a source chord held in the exterior.

\section{Derivative of a normalized conditional heat bath}
\label{sec:s7v9-conditional}

Let \(z\) be a compact owner-fibre variable with reference measure
\(dm(z)\), let \(t\) be an exterior parameter, and let
\begin{equation}
 d\mu_t(z)=\rho_t(z)\,dm(z),\qquad
 \rho_t={\widetilde\rho_t\over Z_t}.                    \label{eq:s7v9-density}
\end{equation}
Write
\[
 K_tf:=\int f(z)\,d\mu_t(z).
\]
All statements apply fibrewise in any additional far-exterior variable.

\begin{theorem}[Exact conditional score derivative]
\label{thm:s7v9-score}
Assume \(t\mapsto\widetilde\rho_t\) and \(t\mapsto f_t\) are
differentiable under the integral.  Put
\begin{align}
 h_t&:=\partial_t\log\widetilde\rho_t,                  \label{eq:s7v9-raw-score}\\
 \sigma_t&:=h_t-K_th_t.                                 \label{eq:s7v9-score}
\end{align}
Then
\begin{align}
 K_t\sigma_t&=0,                                        \label{eq:s7v9-score-centered}\\
 \partial_tK_tf_t
 &=K_t(\partial_tf_t)+K_t(f_t\sigma_t)                  \label{eq:s7v9-derivative}\\
 &=K_t(\partial_tf_t)+\operatorname{Cov}_{\mu_t}(f_t,h_t).
                                                                  \label{eq:s7v9-covariance}
\end{align}
\end{theorem}

\begin{proof}
Differentiate
\[
 \log\rho_t=\log\widetilde\rho_t-\log Z_t.
\]
Since
\[
 \partial_t\log Z_t
 ={1\over Z_t}\int\partial_t\widetilde\rho_t\,dm
 =K_th_t,
\]
one has \(\partial_t\log\rho_t=\sigma_t\), whose mean is zero.
Differentiating the normalized integral gives
\[
 \partial_tK_tf_t
 =K_t(\partial_tf_t)+K_t(f_t\partial_t\log\rho_t).
\]
This is \eqref{eq:s7v9-derivative}; centering the second factor gives
\eqref{eq:s7v9-covariance}.
\end{proof}

\begin{corollary}[Moving conditional projection]
\label{cor:s7v9-moving}
For \(D_t=I-K_t\),
\begin{equation}
 \partial_t(D_tf_t)
 =D_t(\partial_tf_t)-K_t(f_t\sigma_t).                  \label{eq:s7v9-moving}
\end{equation}
Thus differentiating a conditional fluctuation necessarily produces the
centered score covariance.
\end{corollary}

\begin{proof}
Subtract \eqref{eq:s7v9-derivative} from
\(\partial_tf_t\).
\end{proof}

\section{The full incident curvature map}
\label{sec:s7v9-curvature}

On the fully good incident set of V1, define
\[
 \widetilde{\cal F}_\eta(U)
 :=\bigl(\log W_p(U,\eta)\bigr)_{p\in{\cal P}(C)}
 \in\mathfrak{su}(3)^{250}
\]
and
\begin{equation}
 \widetilde B_{\eta,U}:=
 D\widetilde{\cal F}_\eta(U):
 \mathfrak{su}(3)^{60}\longrightarrow
 \mathfrak{su}(3)^{250}.                               \label{eq:s7v9-Btilde}
\end{equation}
Its selected sixty rows are \(B_{\tau,U}\), so
\(\widetilde B_{\eta,U}\) is injective on the V8 core.
Let
\[
 \Psi(\xi):=3-\Re\Tr(e^\xi)
\]
and put
\begin{equation}
 G_{\eta,U}:=
 \bigl(\nabla\Psi(\log W_p(U,\eta))\bigr)_{p\in{\cal P}(C)}.
                                                                  \label{eq:s7v9-G}
\end{equation}

\begin{theorem}[Wilson gradient is exactly curvature adjoint]
\label{thm:s7v9-action-source}
For the complete incident potential
\[
 {\cal V}_{\rm inc}^\eta(U)
 =\sum_{p\in{\cal P}(C)}\Psi(\log W_p(U,\eta)),
\]
one has
\begin{equation}
 \boxed{
 \nabla_U{\cal V}_{\rm inc}^\eta
 =\widetilde B_{\eta,U}^*G_{\eta,U}.}                  \label{eq:s7v9-action-source}
\end{equation}
There is a dimension-only constant \(C_G\) such that on a sufficiently
small fully good set,
\begin{equation}
 \|G_{\eta,U}\|^2
 \le C_G{\cal V}_{\rm inc}^\eta(U).                     \label{eq:s7v9-G-bound}
\end{equation}
\end{theorem}

\begin{proof}
For a chord tangent \(X\), the chain rule gives
\begin{align*}
 D{\cal V}_{\rm inc}^\eta(U)[X]
 &=
 \sum_p
 \left\langle
 \nabla\Psi(\log W_p),
 D(\log W_p)[X]\right\rangle\\
 &=
 \langle G_{\eta,U},\widetilde B_{\eta,U}X\rangle
 =
 \langle\widetilde B_{\eta,U}^*G_{\eta,U},X\rangle,
\end{align*}
which proves \eqref{eq:s7v9-action-source}.

The analytic function \(\Psi\) has a nondegenerate quadratic minimum at
zero.  Hence on a fixed small ball there are constants \(c_1,c_2>0\)
such that
\[
 \|\nabla\Psi(\xi)\|^2\le c_1\|\xi\|^2
 \le c_2\Psi(\xi).
\]
Apply this to all \(250\) plaquette logarithms and sum.
\end{proof}

\begin{corollary}[Exact polar control of the Wilson action source]
\label{cor:s7v9-polar}
Let
\[
 \widetilde H_{\eta,U}
 :=\widetilde B_{\eta,U}^*\widetilde B_{\eta,U}.
\]
On the V8 core,
\begin{equation}
 \left\|
 \widetilde H_{\eta,U}^{-1/2}
 \nabla_U{\cal V}_{\rm inc}^\eta
 \right\|
 \le
 \|G_{\eta,U}\|
 \le
 \sqrt{C_G{\cal V}_{\rm inc}^\eta(U)}.                  \label{eq:s7v9-polar}
\end{equation}
\end{corollary}

\begin{proof}
The selected rows make \(\widetilde B_{\eta,U}\) injective.  Apply
\Cref{thm:s7v8-polar} to \(\widetilde B_{\eta,U}\), then use
\eqref{eq:s7v9-action-source,eq:s7v9-G-bound}.
\end{proof}

\begin{assessment}[CRLC action row]
\label{ass:s7v9-action}
The raw Wilson action derivative now satisfies the non-circular part of
CRLC: its preimage \(G_{\eta,U}\) is constructed directly from plaquette
holonomies, is local, and is controlled by the incident defect.  No
pseudoinverse or curvature singular floor appears.
\end{assessment}

\section{Physical conditional score ledger}
\label{sec:s7v9-ledger}

Suppose the unnormalized owner-fibre density has the physical form
\begin{equation}
 \widetilde\rho_t
 =
 \exp\{-s{\cal V}_t\}\,u_{0,t}^{\,2}\,J_t,              \label{eq:s7v9-physical-density}
\end{equation}
where \(J_t\) is any declared reference/Jacobian density.  For a
source-chord tangent \(X\), denote the corresponding vector derivatives
by
\begin{align}
 S_{\rm W}&:=\nabla_X{\cal V}_t
 =\widetilde B^*G,                                      \label{eq:s7v9-SW}\\
 S_0&:=\nabla_X\log u_{0,t},                            \label{eq:s7v9-S0}\\
 S_J&:=\nabla_X\log J_t.                                \label{eq:s7v9-SJ}
\end{align}
The conditional expectation \(K_t\) acts componentwise on these
finite-dimensional source vectors.

\begin{theorem}[Exact centered physical score decomposition]
\label{thm:s7v9-score-ledger}
The vector score whose pairing with \(X\) appears in
\eqref{eq:s7v9-score} is
\begin{equation}
 \Sigma_t
 =
 -sS_{\rm W}+2S_0+S_J
 -K_t(-sS_{\rm W}+2S_0+S_J).                            \label{eq:s7v9-Sigma}
\end{equation}
Equivalently,
\begin{equation}
 \boxed{
 \Sigma_t
 =-s\widetilde B^*G+R_{\rm gn},}                        \label{eq:s7v9-score-split}
\end{equation}
where the ground/normalization residual is
\begin{equation}
 R_{\rm gn}
 :=
 2S_0+S_J+sK_tS_{\rm W}-2K_tS_0-K_tS_J.                \label{eq:s7v9-Rgn}
\end{equation}
\end{theorem}

\begin{proof}
Differentiate \eqref{eq:s7v9-physical-density}:
\[
 h_t=-s\langle S_{\rm W},X\rangle
     +2\langle S_0,X\rangle+\langle S_J,X\rangle.
\]
Subtract its conditional mean as in
\eqref{eq:s7v9-score}.  Since the identity holds for every \(X\), the
vector score is \eqref{eq:s7v9-Sigma}.  Insert
\eqref{eq:s7v9-SW} and collect the remaining terms.
\end{proof}

\begin{corollary}[What V8 now controls]
\label{cor:s7v9-score-bound}
On the fully good core,
\begin{equation}
 \|\widetilde H^{-1/2}\Sigma_t\|
 \le
 s\sqrt{C_G{\cal V}_{\rm inc}^\eta(U)}
 +{2\over\widetilde\sigma_0}\|R_{\rm gn}\|,             \label{eq:s7v9-score-bound}
\end{equation}
where \(\widetilde\sigma_0\) is any uniform lower bound for
\(\sigma_{\min}(\widetilde B)\), and one may use the selected-map bound
\(\widetilde\sigma_0\ge\sigma_0/2\).
\end{corollary}

\begin{proof}
Apply \Cref{thm:s7v8-source} with
\(Y=-sG\), \(R=R_{\rm gn}\), and
\(B=\widetilde B\), then use
\eqref{eq:s7v9-G-bound}.  Adding plaquette rows cannot decrease the least
singular value, so the selected-map bound applies.
\end{proof}

\begin{firewall}[Conditional centering is a real residual]
\label{fire:s7v9-centering}
Even though \(S_{\rm W}(z)=\widetilde B_z^*G_z\) pointwise in the owner
variable, its conditional mean \(K_tS_{\rm W}\) need not equal
\(\widetilde B_z^*Y(z)\) with a uniformly bounded local \(Y\).
The derivative \(\widetilde B_z\) varies across the owner fibre.  Moving
\(K_tS_{\rm W}\) into the curvature-adjoint term by a pseudoinverse would
be the circular construction forbidden in
\Cref{fire:s7v8-range}.
\end{firewall}

\section{Internal integration by parts}
\label{sec:s7v9-ibp}

\begin{theorem}[Exact cancellation for an internal Haar direction]
\label{thm:s7v9-ibp}
Let \(X\) be a left-invariant divergence-free vector field acting only on
a compact group variable integrated by \(K_t\), with all exterior
parameters fixed.  For the normalized density
\eqref{eq:s7v9-physical-density}, let
\[
 h_X=X\log\rho_t
 =-sX{\cal V}_t+2X\log u_{0,t}+X\log J_t.
\]
Then
\begin{equation}
 K_th_X=0,\qquad
 K_t(fh_X)=-K_t(Xf).                                    \label{eq:s7v9-ibp}
\end{equation}
\end{theorem}

\begin{proof}
Haar invariance and compactness give
\[
 0=\int X(f\rho_t)\,dm
 =K_t(Xf)+K_t(fX\log\rho_t).
\]
This proves the second identity.  Set \(f=1\) for the first.
\end{proof}

\begin{corollary}[Internal score is a generator derivative]
\label{cor:s7v9-internal}
For an internal direction, the score covariance in
\eqref{eq:s7v9-derivative} is exactly \(-K_t(Xf)\).  It need not be
estimated by a Fisher-information or essential-supremum score bound.
\end{corollary}

\begin{proof}
Use \Cref{thm:s7v9-ibp} and \(K_th_X=0\).
\end{proof}

\begin{warning}[A source-exterior tangent cannot be integrated by parts]
\label{warn:s7v9-external}
If \(X\) differentiates a source chord which is held fixed in the
owner-fibre integral, then \(X\) is not a vector field on the integrated
variable.  The identity
\eqref{eq:s7v9-ibp} is unavailable.  This is precisely the crossing
contact which produces \(R_{\rm gn}\) and the V4--V5 source/owner
leakage.  Treating it as internal would change the conditional fibre.
\end{warning}

\begin{definition}[Conditional curvature-score card, CCSC]
\label{def:s7v9-CCSC}
CCSC consists of:
\begin{enumerate}[label=\textup{(S\arabic*)},leftmargin=4em]
\item the exact raw action factorization
      \eqref{eq:s7v9-action-source};
\item an internal/external tangent split for each owner transfer;
\item use of \eqref{eq:s7v9-ibp} only on the internal part;
\item a local curvature-adjoint representation or summable estimate for
      the external normalization mean \(K_tS_{\rm W}\);
\item a corresponding bound for the physical ground and Jacobian score
      \(S_0,S_J\);
\item insertion of the residual in the V8 polar and V5 martingale Schur
      budgets; and
\item response-ground, gauge-quotient, and regulator compatibility.
\end{enumerate}
\end{definition}

\begin{theorem}[CCSC identifies the exact remaining CRLC source row]
\label{thm:s7v9-CCSC}
Rows (S1)--(S3) close the raw Wilson and internal-score parts of CRLC.
Rows (S4)--(S5) are exactly the missing ground/normalization residual
estimate in \eqref{eq:s7v9-score-bound}.  If (S4)--(S7) hold with the
V4--V5 strict budgets, the conditional score enters RPTC without an
uncentered action-score loss.
\end{theorem}

\begin{proof}
The raw factorization and polar bound are
\Cref{thm:s7v9-action-source,cor:s7v9-polar}.  The internal component is
\Cref{thm:s7v9-ibp}.  The exact residual left after these steps is
\eqref{eq:s7v9-Rgn}; rows (S4)--(S7) state precisely the estimates and
physical identifications needed by
\Cref{thm:s7v8-CRLC,thm:s7v5-MLC,thm:s7v4-RPTC}.
\end{proof}

\begin{decision}[V9 bottleneck decision]
\label{dec:s7v9-bottleneck}
The Wilson action score is no longer the obstruction: it is exactly
curvature adjoint, and internal score terms are exact generator
derivatives.  The remaining source is the external conditional
normalization plus physical ground/Jacobian score
\(R_{\rm gn}\).  V10 is the compulsory companion audit and will test
whether the newest SM or NS notes contain a non-circular moving-ground or
nested-conditioning identity for this residual.
\end{decision}

\begin{programme}[Seventh-series V10]
\label{prog:s7v9-V10}
V10 will perform the SM/NS audit, derive a two-level nested-conditional
formula for \(K_tS_{\rm W}\), and decide whether its covariance can be
charged to the owner martingale kernel without an essential-supremum
score assumption.
\end{programme}

\begin{firewall}[Claim boundary after seventh-series V9]
\label{fire:s7v9-boundary}
V9 proves the exact derivative of a normalized conditional heat bath,
the curvature-adjoint factorization and defect bound for the complete
Wilson action gradient, its polar control, the centered physical score
ledger, and internal Haar integration by parts.

V9 does not control the external normalization or ground-score residual,
prove the remaining restricted norm, response stability or
tensorization, construct the continuum, or prove the Yang--Mills mass
gap.
\end{firewall}

\paragraph{Parent-version record.}
V9 was formed from
\texttt{Renewal\_Yang\_Mills\_Mass\_Gap\_Research\_Notes\_V8.tex}.
The V8 parent had \(3\,136\) logical source lines, \(116\,402\) bytes,
and SHA--256
\[
 \texttt{576D70B7E5E70D0424A72E0BE7DF7142BEAD73354A015AE9AE2A565EFF57D6B4}.
\]

% =============================================================================
% END APPENDED SEVENTH-SERIES V9 MATERIAL
% =============================================================================

% =============================================================================
% BEGIN APPENDED SEVENTH-SERIES V10 MATERIAL
% =============================================================================

\chapter{Companion audit and the exact two-conditioning score split}
\label{ch:s7v10}

This is the compulsory five-version companion audit.  Its purpose is not
to import an analogy as a proof.  It tests the precise V9
ground/normalization residual against the newest Standard--Model and
Navier--Stokes calculations, records what survives that test, and then
proves the corresponding conditional-projection statement in the
Yang--Mills probability space.

\section{Compulsory V10 companion audit}
\label{sec:s7v10-audit}

\begin{audit}[Files inspected]
\label{aud:s7v10-files}
The current companion files were:
\begin{enumerate}[label=\textup{(\alph*)},leftmargin=3em]
\item
\texttt{Renewal\_SM\_Einstein\_Continuum\_Research\_Notes\_V49.tex},
with \(20\,274\) logical source lines, \(730\,129\) bytes, and SHA--256
\[
 \texttt{345B10AA597C15891384565DE31DF0ABCC3E942D856F8C20B7F4A5CCFD91CCFC};
\]
\item
\texttt{Renewal\_NS\_Stability\_Research\_Notes\_V26.tex},
with \(5\,319\) logical source lines, \(278\,283\) bytes, and SHA--256
\[
 \texttt{82C887F8C2C69E4F28FAD7C3DC8DE5B9099CB5A4BF431EBA1FFF3E91935978AD}.
\]
\end{enumerate}
The SM file advanced from V46, inspected at YM V5, through V47--V49.
The NS file is unchanged since that audit.
\end{audit}

\begin{assessment}[SM V47--V49]
\label{ass:s7v10-SM}
SM V47 proves that a fixed one-cell absolute chart has worst conditional
bad probability one under translated neighbours, and replaces it by a
common/relative block split.  V48 proves the exact finite-block
positivity threshold for its relative quadratic form and shows that the
cosmological mass does not by itself give a cofinal capacity estimate.
V49 performs the relevant upstream step: after the linearized ADM
constraint is solved, a generic matter source acquires an inverse
Laplacian, whereas a source of the form
\(\rho(k)=ik\cdot J(k)\) has gradient response
\[
 {ik\over |k|^2}\rho(k)
 =-{k\otimes k\over |k|^2}J(k)
\]
and therefore an \(L^2\) bound with constant one.

The transferable statement is structural: an inverse elliptic weight is
uniform only after the needed derivative factor has been proved for the
actual source.  This agrees with the YM V6--V8 curvature-adjoint route.
The gravitational constraint, conformal mode, and cosmological
coefficients themselves do not transfer to the compact gauge fibre.
\end{assessment}

\begin{assessment}[NS V26]
\label{ass:s7v10-NS}
NS V26 computes exact rank-one norms of the first two Oseen-kernel
derivatives.  The first-derivative refinement gives no gain; the
second-derivative refinement gives only a small local improvement.  This
confirms the value of source-specific operator norms, but supplies no
moving-ground identity and no estimate for the YM external conditional
score.  No NS regularity result is imported.
\end{assessment}

\begin{decision}[Audit decision]
\label{dec:s7v10-audit}
The next YM step should retain the V8 curvature square root and factor
the \emph{actual conditional mean} before applying it.  The exact
analogue of the SM divergence cancellation is an orthogonal
source/owner conditioning split.  A full-space norm is again the wrong
quantity because V5 already exhibits a spectator on which it equals
one.
\end{decision}

\section{Three conditional projections}
\label{sec:s7v10-projections}

Let \((\Omega,\mathcal F,\mu)\) be a probability space and let
\[
 \mathcal C\subset\mathcal A\cap\mathcal B\subset\mathcal F
\]
be sigma fields.  In the application, \(\mathcal A\) retains the source
exterior, \(\mathcal B\) retains the owner exterior, and \(\mathcal C\)
retains their declared common exterior.  On
\(L^2(\mu;\mathfrak h)\), where \(\mathfrak h\) is a finite-dimensional
real Hilbert space, write
\begin{equation}
 E=\mathbb E[\cdot\mid\mathcal A],\qquad
 K=\mathbb E[\cdot\mid\mathcal B],\qquad
 R=\mathbb E[\cdot\mid\mathcal C].                       \label{eq:s7v10-EKR}
\end{equation}
This \(K\) is one owner conditional projection, not the randomized
average \(P_e\) of V4.

\begin{lemma}[Nested projection algebra]
\label{lem:s7v10-nested}
The operators in \eqref{eq:s7v10-EKR} are orthogonal projections and
\begin{equation}
 ER=RE=KR=RK=R.                                          \label{eq:s7v10-nesting}
\end{equation}
Consequently \(E-R\) and \(K-R\) are the orthogonal projections onto
\[
 L^2(\mathcal A;\mathfrak h)\ominus L^2(\mathcal C;\mathfrak h),
 \qquad
 L^2(\mathcal B;\mathfrak h)\ominus L^2(\mathcal C;\mathfrak h),
\]
respectively.
\end{lemma}

\begin{proof}
Conditional expectation is the orthogonal projection onto the
corresponding measurable subspace.  Since
\(L^2(\mathcal C)\) is a closed subspace of both
\(L^2(\mathcal A)\) and \(L^2(\mathcal B)\), projecting first onto either
larger subspace and then onto \(L^2(\mathcal C)\), in either admissible
order, gives \(R\).  Thus, for example,
\[
 (E-R)^2=E-ER-RE+R=E-R
\]
and \(E-R\) is self-adjoint.  Its range and kernel give the stated
orthogonal difference.  The proof for \(K-R\) is identical.
\end{proof}

The curvature weight below may depend on the common exterior.  Let
\(M:\Omega\to\mathcal L(\mathfrak h)\) be bounded and
\(\mathcal C\)-measurable.  Multiplication by \(M\) commutes with
\(E,K,R\); this follows first for simple \(M\), then by bounded
approximation.

\begin{theorem}[Exact two-conditioning Pythagoras identity]
\label{thm:s7v10-Pythagoras}
For every \(S\in L^2(\mu;\mathfrak h)\),
\begin{align}
 (I-E)KS
  &=(I-E)(K-R)S,                                         \label{eq:s7v10-coarse-cancel}\\
 \|M(I-E)KS\|_2^2
  &=\|M(K-R)S\|_2^2
    -\|M(E-R)(K-R)S\|_2^2.                               \label{eq:s7v10-Pythagoras}
\end{align}
Thus the common-exterior part \(RS\) cancels exactly; no norm estimate is
spent on it.
\end{theorem}

\begin{proof}
Since \((I-E)R=0\), the first identity is immediate.  Put
\[
 h=M(K-R)S.
\]
The \(\mathcal C\)-measurability of \(M\) and
\(R(K-R)=0\) give \(Rh=0\).  Hence
\[
 Eh=(E-R)h=M(E-R)(K-R)S.
\]
The orthogonal decomposition \(h=Eh+(I-E)h\) yields
\[
 \|(I-E)h\|_2^2=\|h\|_2^2-\|Eh\|_2^2.
\]
Commute \(M\) through \(I-E\) and substitute the displayed
expressions.
\end{proof}

\begin{corollary}[Only the source fluctuation sees the residual]
\label{cor:s7v10-pairing}
For scalar \(f\in L^2(\mu)\), \(X\in\mathfrak h\), and
\(D=I-E\),
\begin{equation}
 \big\langle Df,\langle KS,X\rangle_{\mathfrak h}\big\rangle_{L^2}
 =
 \big\langle Df,
   \langle (I-E)(K-R)S,X\rangle_{\mathfrak h}
 \big\rangle_{L^2}.                                     \label{eq:s7v10-pairing}
\end{equation}
\end{corollary}

\begin{proof}
The function \(Df\) is orthogonal to every
\(\mathcal A\)-measurable scalar function.  Subtract \(EKS\) inside the
second factor and use \eqref{eq:s7v10-coarse-cancel}.
\end{proof}

\section{The correct lower-angle quantity}
\label{sec:s7v10-angle}

Let \(\mathscr S\subset L^2(\mu;\mathfrak h)\) be the closed linear
source space actually generated by the incident action scores and the
martingale contacts under consideration.  Define
\begin{equation}
 \gamma_{M}(\mathscr S)
 :=
 \inf_{\substack{S\in\mathscr S\\M(K-R)S\ne0}}
 {\|M(E-R)(K-R)S\|_2\over
  \|M(K-R)S\|_2},                                       \label{eq:s7v10-gamma}
\end{equation}
with value one if the set is empty.

\begin{theorem}[Restricted lower-angle leakage bound]
\label{thm:s7v10-angle}
One has \(0\le\gamma_M(\mathscr S)\le1\), and every
\(S\in\mathscr S\) satisfies
\begin{equation}
 \|M(I-E)KS\|_2
 \le
 \sqrt{1-\gamma_M(\mathscr S)^2}\,
 \|M(K-R)S\|_2.                                         \label{eq:s7v10-angle}
\end{equation}
The constant is sharp for the declared source space.
\end{theorem}

\begin{proof}
Because \(E-R\) is an orthogonal projection and commutes with \(M\), the
quotient in \eqref{eq:s7v10-gamma} lies in \([0,1]\).
Insert its lower bound in
\eqref{eq:s7v10-Pythagoras}.  Sharpness follows by taking an infimizing
sequence in the definition of \(\gamma_M(\mathscr S)\).
\end{proof}

\begin{warning}[This is not the maximal-correlation angle]
\label{warn:s7v10-angle}
An upper bound on
\(\|(E-R)(K-R)\|\) does not prove
\eqref{eq:s7v10-angle}; the needed information is a \emph{lower} bound
on the source-generated range.  On the full \(K-R\) range that lower
bound can be zero.  The spectator obstruction of V5 is exactly such a
zero-angle direction.  The correlated binary computation of V4 has
\(\gamma=|t|\) on its one-dimensional active range and leakage
\(\sqrt{1-t^2}\), consistent with
\eqref{eq:s7v10-angle}.  Neither example proves a YM value.
\end{warning}

\section{Fixed-reference curvature transport}
\label{sec:s7v10-reference}

The nonlinear square root in V9 is owner dependent and therefore cannot
simply be commuted through \(K\).  A fixed flat reference removes that
commutator, provided the comparison is made in the relative curvature
metric.

Let \(B_\star:\mathfrak x\to\mathfrak y\) be the injective flat
incident-curvature derivative on the fixed collar, and let
\[
 H_\star=B_\star^*B_\star.
\]
It may depend on the common exterior frames, but is assumed
\(\mathcal C\)-measurable.  For the nonlinear derivative \(B=B(U)\), put
\(H=B^*B\).

\begin{lemma}[Relative collar comparison]
\label{lem:s7v10-relative}
Suppose
\begin{equation}
 \varepsilon:=
 \|(B-B_\star)H_\star^{-1/2}\|_{\mathfrak x\to\mathfrak y}<1.
                                                               \label{eq:s7v10-epsilon}
\end{equation}
Then
\begin{align}
 H&\succeq(1-\varepsilon)^2H_\star,                    \label{eq:s7v10-Loewner}\\
 \|H^{-1/2}v\|
 &\le {1\over1-\varepsilon}\|H_\star^{-1/2}v\|,
 \qquad v\in\mathfrak x,                               \label{eq:s7v10-inverse-compare}\\
 \|H_\star^{-1/2}B^*Y\|
 &\le(1+\varepsilon)\|Y\|,
 \qquad Y\in\mathfrak y.                               \label{eq:s7v10-adjoint-compare}
\end{align}
\end{lemma}

\begin{proof}
For \(x\in\mathfrak x\), set \(z=H_\star^{1/2}x\).
The map \(B_\star H_\star^{-1/2}\) is an isometry, so
\[
 \|(B-B_\star)x\|\le\varepsilon\|z\|
 =\varepsilon\|B_\star x\|.
\]
The reverse triangle inequality gives
\(\|Bx\|\ge(1-\varepsilon)\|B_\star x\|\), proving
\eqref{eq:s7v10-Loewner}.  Inversion reverses the Loewner order and
gives \eqref{eq:s7v10-inverse-compare}.  Finally,
\[
 \|B H_\star^{-1/2}\|
 \le\|B_\star H_\star^{-1/2}\|
    +\|(B-B_\star)H_\star^{-1/2}\|
 \le1+\varepsilon.
\]
Take adjoints.
\end{proof}

\begin{remark}[Why the relative core exists]
\label{rem:s7v10-relative-core}
The collar is finite and \(B(U)\) is analytic in its link variables on
the chosen logarithm chart.  Since \(H_\star\) is invertible, continuity
gives a fixed neighbourhood on which
\eqref{eq:s7v10-epsilon} holds for any prescribed
\(\varepsilon\in(0,1)\).  This is an existence statement for a fixed
core, not yet a useful probability estimate for its complement.  In
particular it does not conceal the very small V100 singular floor; the
core may have to be correspondingly small.
\end{remark}

\begin{theorem}[Conditional normalization compiler on the active score space]
\label{thm:s7v10-compiler}
Assume \eqref{eq:s7v10-epsilon} uniformly on the declared good core and
let the incident Wilson score have its V9 form
\[
 S_{\rm W}=B^*G.
\]
Let
\(\gamma_\star=\gamma_{H_\star^{-1/2}}(\mathscr S_{\rm W})\)
for the closed span of these scores.  Then
\begin{equation}
 \boxed{
 \|H^{-1/2}(I-E)KS_{\rm W}\|_2
 \le
 {1+\varepsilon\over1-\varepsilon}
 \sqrt{1-\gamma_\star^2}\,\|G\|_2.}                     \label{eq:s7v10-compiler}
\end{equation}
The estimate is cutoff independent if
\(\varepsilon<1\) and \(\gamma_\star\) have cutoff-independent bounds.
\end{theorem}

\begin{proof}
Use \eqref{eq:s7v10-inverse-compare} pointwise, then
\Cref{thm:s7v10-angle} with \(M=H_\star^{-1/2}\):
\[
 \|H^{-1/2}(I-E)KS_{\rm W}\|_2
 \le { \sqrt{1-\gamma_\star^2}\over1-\varepsilon}
       \|H_\star^{-1/2}(K-R)S_{\rm W}\|_2.
\]
Because \(H_\star^{-1/2}\) is \(\mathcal C\)-measurable, it commutes
with \(K-R\), which is a contraction.  Thus the last norm is at most
\(\|H_\star^{-1/2}S_{\rm W}\|_2\).  Apply
\eqref{eq:s7v10-adjoint-compare} to \(S_{\rm W}=B^*G\).
\end{proof}

\begin{corollary}[Defect form]
\label{cor:s7v10-defect}
On the small Wilson chart of V9,
\begin{equation}
 \|H^{-1/2}(I-E)KS_{\rm W}\|_2
 \le
 {1+\varepsilon\over1-\varepsilon}
 \sqrt{1-\gamma_\star^2}\,
 \sqrt{C_G\,\mathbb E{\cal V}_{\rm inc}}.                \label{eq:s7v10-defect}
\end{equation}
\end{corollary}

\begin{proof}
Square and integrate \eqref{eq:s7v9-G-bound}, then use
\Cref{thm:s7v10-compiler}.
\end{proof}

\section{Exact residual after source projection}
\label{sec:s7v10-residual}

Set
\[
 S_{\rm ph}:=2S_0+S_J.
\]
The V9 ledger becomes
\begin{equation}
 R_{\rm gn}=(I-K)S_{\rm ph}+sKS_{\rm W}.                 \label{eq:s7v10-Rgn}
\end{equation}

\begin{theorem}[Projected ground/normalization ledger]
\label{thm:s7v10-ledger}
For \(D=I-E\),
\begin{equation}
 \boxed{
 DR_{\rm gn}
 =D(I-K)S_{\rm ph}
  +sD(K-R)S_{\rm W}.}                                   \label{eq:s7v10-ledger}
\end{equation}
The second term is controlled by
\eqref{eq:s7v10-compiler}.  No part of the first term is controlled by
the conditional-angle calculation alone.
\end{theorem}

\begin{proof}
Apply \(D\) to \eqref{eq:s7v10-Rgn}.  Since \(DR=0\),
\(DKS_{\rm W}=D(K-R)S_{\rm W}\).  This proves the identity.  The stated
division of estimates follows from
\Cref{thm:s7v10-compiler}; that theorem contains no hypothesis on
\(S_{\rm ph}\).
\end{proof}

\begin{firewall}[What an angle bound would and would not prove]
\label{fire:s7v10-angle}
Even a strict \(\gamma_\star>0\) would only control the external
\emph{Wilson normalization mean} after source projection.  It would not
prove a corresponding bound for
\(D(I-K)(2S_0+S_J)\), the good-core probability, the restricted
martingale Schur sums, continuum construction, or the spectral mass
gap.  Conversely, estimating the full-space lower angle is unnecessary
and generally impossible by the V5 spectator theorem; only the active
score space enters \eqref{eq:s7v10-compiler}.
\end{firewall}

\begin{definition}[Conditional source-angle card, CSAC]
\label{def:s7v10-CSAC}
CSAC consists of:
\begin{enumerate}[label=\textup{(A\arabic*)},leftmargin=4em]
\item a global physical measure for which the source, owner, and common
      fibres are the conditional expectations \(E,K,R\);
\item a common-exterior flat curvature metric \(H_\star\);
\item a relative nonlinear collar bound
      \eqref{eq:s7v10-epsilon} with \(\varepsilon<1\);
\item a positive lower angle \(\gamma_\star\) on the incident
      action-score space, not on the full Hilbert space;
\item a separate factorization or resolvent estimate for
      \(D(I-K)(2S_0+S_J)\);
\item insertion of both terms of \eqref{eq:s7v10-ledger} into the V5
      martingale row/column sums with the strict V3 shell budget;
\item good-core, gauge, regulator, response-ground, and observable
      compatibility.
\end{enumerate}
\end{definition}

\begin{theorem}[What CSAC would close]
\label{thm:s7v10-CSAC}
Rows (A1)--(A4) place the external Wilson normalization mean in the
curvature square-root norm without an essential-supremum score loss.
Rows (A5)--(A7), if proved with the existing strict budgets, would close
the remaining CCSC source row and feed the result into MLC and RPTC.
\end{theorem}

\begin{proof}
Rows (A1)--(A4) are precisely
\Cref{thm:s7v10-Pythagoras,thm:s7v10-compiler}.  The exact unresolved
term is identified by \Cref{thm:s7v10-ledger}.  The downstream
interfaces are
\Cref{thm:s7v9-CCSC,thm:s7v5-MLC,thm:s7v4-RPTC}.
\end{proof}

\begin{assessment}[Progress at seventh-series V10]
\label{ass:s7v10-status}
The conditional normalization mean is no longer an undifferentiated
residual.  Its common-exterior component cancels exactly, its remaining
source fluctuation obeys a sharp Pythagoras identity, and a fixed
curvature reference converts that identity into
\eqref{eq:s7v10-compiler}.  The new numerical datum is a restricted
\emph{lower} conditional angle \(\gamma_\star\).  No positive value has
yet been proved for the Wilson measure.  The other unresolved term is
the physical ground/Jacobian response
\(D(I-K)(2S_0+S_J)\).
\end{assessment}

\begin{decision}[V10 bottleneck decision]
\label{dec:s7v10-bottleneck}
The ground score is now the upstream target.  Bounding it as an
arbitrary derivative would merely reintroduce the inverse-curvature
loss.  V11 should instead differentiate the finite-volume fibre ground
equation, derive the exact reduced-resolvent/Poisson equation for
\(\partial\log u_0\), and determine whether the V9
curvature-adjoint Wilson source survives the resolvent.  The same
calculation will state the spectral input needed for a positive
source-specific angle.
\end{decision}

\begin{programme}[Seventh-series V11]
\label{prog:s7v10-V11}
V11 will derive the moving-ground equation for a parameter-dependent
positive fibre ground state, separate the normalization gauge, express
the ground score through the reduced generator inverse, and test that
formula in the curvature square-root norm.  Any required fibre gap or
commutator estimate will remain an explicit hypothesis rather than be
inferred from finite dimensionality.
\end{programme}

\begin{firewall}[Claim boundary after seventh-series V10]
\label{fire:s7v10-boundary}
V10 performs the required SM/NS audit and proves the exact
three-projection cancellation, the sharp restricted-angle identity, the
relative fixed-curvature comparison, and the resulting Wilson
normalization compiler.

V10 does not prove a positive Wilson source angle, control the
ground/Jacobian response, prove the remaining martingale budgets,
construct the continuum, or prove the Yang--Mills mass gap.
\end{firewall}

\paragraph{Parent-version record.}
V10 was formed from
\texttt{Renewal\_Yang\_Mills\_Mass\_Gap\_Research\_Notes\_V9.tex}.
The V9 parent had \(3\,554\) logical source lines, \(130\,457\) bytes,
and SHA--256
\[
 \texttt{DEA91B46CE0AECF3B4966E448F674C10EE44AA8FA6643A9877EE2B89C3E39F67}.
\]

% =============================================================================
% END APPENDED SEVENTH-SERIES V10 MATERIAL
% =============================================================================

% =============================================================================
% BEGIN APPENDED SEVENTH-SERIES V11 MATERIAL
% =============================================================================

\chapter{The moving fibre ground and the exact tree-gauge measure}
\label{ch:s7v11}

V10 isolates two genuinely different pieces of the V9 residual: the
owner normalization of the Wilson score and the derivative of the
positive response ground.  This chapter goes upstream to the eigenvalue
equation defining that ground.  The logarithmic derivative is an exact
reduced-generator inverse of a centered Hamiltonian source.  The
inverse is harmless in the required curvature metric only if a
cutoff-uniform fibre gap and a curvature-adjoint Hamiltonian derivative
are proved.  Separately, maximal-tree gauge fixing preserves product
Haar measure exactly, so the lattice tree change of variables creates no
source-dependent Jacobian score.

\section{A parameter-dependent positive ground}
\label{sec:s7v11-ground}

Let \((Z,m)\) be a compact measured space and let
\(\mathsf H_t\) be a differentiable family of self-adjoint operators in
\(L^2(m)\) with common form domain.  Assume that near \(t=0\) it has a
simple normalized ground pair
\begin{equation}
 \mathsf H_tu_t=\lambda_tu_t,\qquad
 u_t>0,\qquad \int_Zu_t^2\,dm=1.                         \label{eq:s7v11-ground}
\end{equation}
All derivatives below are at \(t=0\).  Write
\[
 u=u_0,\qquad \lambda=\lambda_0,\qquad
 d\nu=u^2\,dm,\qquad K_0f=\int_Zf\,d\nu.
\]
Define the ground-transformed operator
\begin{equation}
 \mathcal Lf:=u^{-1}(\mathsf H_0-\lambda)(uf).            \label{eq:s7v11-L}
\end{equation}
It is understood on the transported form domain in \(L^2(\nu)\).

\begin{lemma}[Ground-transform spectral facts]
\label{lem:s7v11-ground-transform}
The operator \(\mathcal L\) is nonnegative and self-adjoint in
\(L^2(\nu)\), \(\mathcal L1=0\), and its nonzero spectrum is the spectrum
of \(\mathsf H_0-\lambda\) off the ground state.
\end{lemma}

\begin{proof}
Multiplication \(Uf=uf\) is unitary from \(L^2(\nu)\) to \(L^2(m)\),
because
\(\|Uf\|_{L^2(m)}^2=\int|f|^2u^2\,dm\).
Equation \eqref{eq:s7v11-L} says
\(\mathcal L=U^{-1}(\mathsf H_0-\lambda)U\).  Unitary conjugacy gives
self-adjointness, nonnegativity, and the spectrum.  Finally
\(\mathcal L1=u^{-1}(\mathsf H_0-\lambda)u=0\).
\end{proof}

For a parameter direction \(X\), put
\begin{equation}
 q_X:=\partial_X\log u_t\big|_{t=0},\qquad
 W_X:=u^{-1}\bigl((\partial_X\mathsf H_t)u\bigr)\big|_{t=0}.
                                                               \label{eq:s7v11-qW}
\end{equation}
If only the potential moves, then \(W_X\) is its ordinary derivative.

\begin{theorem}[Exact moving-ground Poisson equation]
\label{thm:s7v11-Poisson}
Under the preceding differentiability assumptions,
\begin{align}
 K_0q_X&=0,                                               \label{eq:s7v11-q-centered}\\
 \partial_X\lambda_t\big|_{t=0}&=K_0W_X,                 \label{eq:s7v11-HF}\\
 \mathcal Lq_X&=-(I-K_0)W_X.                             \label{eq:s7v11-Poisson}
\end{align}
If the spectral gap of \(\mathcal L\) on \(1^\perp\) is
\(g>0\), then
\begin{equation}
 q_X=-\mathcal L_\perp^{-1}(I-K_0)W_X.                   \label{eq:s7v11-resolvent}
\end{equation}
\end{theorem}

\begin{proof}
Differentiate the normalization in \eqref{eq:s7v11-ground}:
\[
 0=2\int_Zu\,\partial_Xu\,dm=2K_0q_X,
\]
which proves \eqref{eq:s7v11-q-centered}.  Differentiating the
eigenvalue equation gives
\[
 (\mathsf H_0-\lambda)\partial_Xu
 =-\bigl(\partial_X\mathsf H_t-\partial_X\lambda_t\bigr)u.
\]
Take the \(L^2(m)\) inner product with \(u\).  The left side vanishes,
and the right side gives
\(\partial_X\lambda_t=\langle u,(\partial_X\mathsf H_t)u\rangle
=K_0W_X\), proving \eqref{eq:s7v11-HF}.  Since
\(\partial_Xu=uq_X\), division by \(u\) and
\eqref{eq:s7v11-L} give \eqref{eq:s7v11-Poisson}.  The right side and
\(q_X\) are in \(1^\perp\); invert \(\mathcal L\) there when \(g>0\).
\end{proof}

\begin{corollary}[Sharp gap estimates]
\label{cor:s7v11-gap}
If \(g>0\), then
\begin{align}
 \|q_X\|_{L^2(\nu)}
 &\le g^{-1}\|(I-K_0)W_X\|_{L^2(\nu)},                  \label{eq:s7v11-L2-gap}\\
 \|\mathcal L^{1/2}q_X\|_{L^2(\nu)}
 &\le g^{-1/2}\|(I-K_0)W_X\|_{L^2(\nu)}.                \label{eq:s7v11-energy-gap}
\end{align}
The powers of \(g\) cannot be improved for a general source.
\end{corollary}

\begin{proof}
Use the spectral resolution of \(\mathcal L\) on \(1^\perp\), where
its spectral parameter is at least \(g\).  A source in the
\(g\)-eigenspace gives equality in both estimates.
\end{proof}

\begin{warning}[Finite fibre does not give a uniform gap]
\label{warn:s7v11-finite}
Compactness gives a positive gap for each fixed elliptic fibre problem;
it gives no lower bound uniform in parameters or cutoffs.  Indeed, on
any fixed connected compact Riemannian \(Z\), the family
\(\mathsf H_a=-a\Delta\), \(a>0\), has positive constant ground and gap
\(a\lambda_1(-\Delta)\to0\).  Thus a worst-conditional lower bound for
\(g\) is a new quantitative row, not a consequence of finite
dimensionality.
\end{warning}

\section{Vector source form and curvature transport}
\label{sec:s7v11-vector}

Let the exterior parameter space be the finite-dimensional chord
tangent space \(\mathfrak x\).  Define the vector-valued Hamiltonian and
ground scores by
\begin{equation}
 W_X=\langle S_{\mathsf H},X\rangle_{\mathfrak x},
 \qquad
 q_X=\langle S_0,X\rangle_{\mathfrak x}.                 \label{eq:s7v11-vector-scores}
\end{equation}
The operator \(\mathcal L_\perp^{-1}\) acts componentwise.

\begin{corollary}[Vector moving-ground identity]
\label{cor:s7v11-vector}
One has
\begin{equation}
 K_0S_0=0,\qquad
 S_0=-\mathcal L_\perp^{-1}(I-K_0)S_{\mathsf H}.         \label{eq:s7v11-vector}
\end{equation}
\end{corollary}

\begin{proof}
Pair the asserted identities with every
\(X\in\mathfrak x\) and apply
\Cref{thm:s7v11-Poisson}.
\end{proof}

Retain the V10 flat reference
\(H_\star=B_\star^*B_\star\), assumed measurable with respect to the
common exterior.  It therefore commutes with the owner-fibre operators
\(K_0\) and \(\mathcal L\).

\begin{theorem}[Curvature-adjoint source survives the ground resolvent]
\label{thm:s7v11-curvature-ground}
Assume the relative collar estimate
\[
 \|(B-B_\star)H_\star^{-1/2}\|\le\varepsilon<1
\]
and suppose the derivative of the fibre Hamiltonian has the explicit
source form
\begin{equation}
 S_{\mathsf H}=B^*Y_{\mathsf H}.                         \label{eq:s7v11-H-source}
\end{equation}
If the ground-transformed fibre gap is at least \(g>0\), then
\begin{align}
 \|H_\star^{-1/2}S_0\|_{L^2(\nu)}
 &\le {1+\varepsilon\over g}
       \|Y_{\mathsf H}\|_{L^2(\nu)},                     \label{eq:s7v11-ground-curvature}\\
 \|\mathcal L^{1/2}H_\star^{-1/2}S_0\|_{L^2(\nu)}
 &\le {1+\varepsilon\over\sqrt g}
       \|Y_{\mathsf H}\|_{L^2(\nu)}.                     \label{eq:s7v11-ground-energy}
\end{align}
\end{theorem}

\begin{proof}
Since \(H_\star^{-1/2}\) depends only on the common exterior, commute it
through \(K_0\) and the ground resolvent in
\eqref{eq:s7v11-vector}.  The projection \(I-K_0\) is a contraction.
Pointwise,
\[
 \|H_\star^{-1/2}B^*Y_{\mathsf H}\|
 \le(1+\varepsilon)\|Y_{\mathsf H}\|
\]
by \eqref{eq:s7v10-adjoint-compare}.  Apply the two spectral estimates
of \Cref{cor:s7v11-gap} componentwise and integrate.
\end{proof}

\begin{firewall}[The source identity is substantive]
\label{fire:s7v11-H-source}
Equation \eqref{eq:s7v11-H-source} holds if the response Hamiltonian
depends on the exterior chord only through a Wilson potential whose
derivative is the V9 curvature adjoint.  It need not hold if the kinetic
metric, boundary condition, gauge slice, or regulator also moves.
Those derivatives must be computed and either factored or retained as
an explicit residual.  Defining \(Y_{\mathsf H}\) by a pseudoinverse
would repeat the circular step excluded in V8.
\end{firewall}

\section{Ground measure versus physical owner measure}
\label{sec:s7v11-measure-change}

The expectation \(K_0\) in the eigenvalue equation is generally not the
physical owner expectation \(K\) of V9.  On each exterior fibre write
\begin{equation}
 d\mu=w\,d\nu,\qquad w>0,\qquad K_0w=1,\qquad
 Kf=K_0(wf).                                             \label{eq:s7v11-w}
\end{equation}

\begin{lemma}[Exact measure-change discrepancy]
\label{lem:s7v11-measure-change}
If \(K_0F=0\) for a vector-valued \(F\), then
\begin{align}
 KF&=K_0((w-1)F),                                        \label{eq:s7v11-K-discrepancy}\\
 \|KF\|^2
 &\le K_0((w-1)^2)\,K_0(\|F\|^2),                       \label{eq:s7v11-chi2}\\
 \|F\|_{L^2(\mu)}
 &\le\|w\|_\infty^{1/2}\|F\|_{L^2(\nu)}.                \label{eq:s7v11-density-ratio}
\end{align}
\end{lemma}

\begin{proof}
Since \(K_0F=0\),
\[
 KF=K_0(wF)=K_0((w-1)F).
\]
Cauchy--Schwarz proves \eqref{eq:s7v11-chi2}.  The last estimate follows
by inserting \(d\mu=w\,d\nu\).
\end{proof}

\begin{theorem}[Physical centred ground-score bound]
\label{thm:s7v11-physical-ground}
In addition to the hypotheses of
\Cref{thm:s7v11-curvature-ground}, assume
\begin{equation}
 \|w\|_\infty\le W_\infty<\infty.                        \label{eq:s7v11-W}
\end{equation}
Let \(D=I-E\) be the source fluctuation of V10 and
\(H=B^*B\).  Then
\begin{equation}
 \boxed{
 \|H^{-1/2}D(I-K)S_0\|_{L^2(\mu)}
 \le
 {\sqrt{W_\infty}(1+\varepsilon)\over
  (1-\varepsilon)g}
 \|Y_{\mathsf H}\|_{L^2(\nu)}.}                          \label{eq:s7v11-physical-ground}
\end{equation}
\end{theorem}

\begin{proof}
The Loewner comparison \eqref{eq:s7v10-inverse-compare} gives
\[
 \|H^{-1/2}D(I-K)S_0\|_{L^2(\mu)}
 \le {1\over1-\varepsilon}
 \|H_\star^{-1/2}D(I-K)S_0\|_{L^2(\mu)}.
\]
Because \(H_\star\) is common-exterior measurable, its inverse square
root commutes with \(D\) and \(K\).  Both \(D\) and \(I-K\) are
orthogonal contractions in the physical \(L^2(\mu)\) space.  The last
norm is therefore at most
\(\|H_\star^{-1/2}S_0\|_{L^2(\mu)}\).  Use
\eqref{eq:s7v11-density-ratio}, followed by
\eqref{eq:s7v11-ground-curvature}.
\end{proof}

\begin{warning}[A density ratio is not free]
\label{warn:s7v11-density}
For the physical weight
\[
 w={e^{-s{\cal V}}J\over K_0(e^{-s{\cal V}}J)},
\]
the constant \(W_\infty\) can grow with the coupling, boundary
condition, or chart.  A good-core oscillation estimate may control it,
but that estimate and the complementary bad-set debit must be placed in
the cofinal budget.  V11 does not infer
\eqref{eq:s7v11-W} uniformly from positivity.
\end{warning}

\section{Maximal-tree gauge has unit Haar Jacobian}
\label{sec:s7v11-tree}

Let \(\Gamma=(V,E)\) be a finite connected oriented graph, \(r\in V\) a
root, and \(T\subset E\) a spanning tree.  Let \(G\) be a compact group
with normalized Haar measure.  A reversal of an oriented edge replaces
its variable by its inverse, which also preserves Haar measure.

\begin{theorem}[Exact product-Haar tree factorization]
\label{thm:s7v11-tree-Haar}
There is a bijective change of variables, up to the fixed root frame,
\begin{equation}
 (U_e)_{e\in E}
 \longleftrightarrow
 \bigl((g_v)_{v\ne r},(W_c)_{c\in E\setminus T}\bigr),   \label{eq:s7v11-tree-map}
\end{equation}
such that every tree link is gauge transformed to the identity and
every \(W_c\) is the corresponding based chord holonomy.  Moreover,
\begin{equation}
 \prod_{e\in E}dU_e
 =
 \left(\prod_{v\ne r}dg_v\right)
 \left(\prod_{c\in E\setminus T}dW_c\right).             \label{eq:s7v11-Haar-factor}
\end{equation}
Thus maximal-tree gauge fixing has constant Faddeev--Popov density one
in group variables.
\end{theorem}

\begin{proof}
Orient each tree edge temporarily from its parent toward its child.
Set \(g_r=1\) and determine \(g_v\) recursively so that
\[
 g_{\operatorname{par}(v)}U_{e_v}g_v^{-1}=1.
\]
Hence \(U_{e_v}=g_{\operatorname{par}(v)}^{-1}g_v\);
this gives a bijection between tree-link variables and
\((g_v)_{v\ne r}\).  Order vertices so parents precede children.
Conditioned on all earlier variables, the substitution
\(U_{e_v}\mapsto g_v\) is a left translation, and therefore
\(dU_{e_v}=dg_v\).  Recursion yields
\[
 \prod_{e\in T}dU_e=\prod_{v\ne r}dg_v.
\]
For a chord \(c=(a,b)\), set
\[
 W_c=g_aU_cg_b^{-1}.
\]
Conditioned on the tree frames this is a left and right translation, so
\(dU_c=dW_c\).  Multiplying over chords proves
\eqref{eq:s7v11-Haar-factor}.  Restoring any original tree-edge
orientation only inserts inversion, which preserves Haar measure.
\end{proof}

\begin{corollary}[No tree-gauge Jacobian score]
\label{cor:s7v11-no-J}
For the finite lattice with product Haar reference measure, if the
owner and source chords are kept as group variables in the maximal-tree
coordinates of \Cref{thm:s7v11-tree-Haar}, then
\begin{equation}
 J\equiv1,\qquad S_J=\nabla_X\log J=0.                   \label{eq:s7v11-no-J}
\end{equation}
The V9 physical residual reduces to
\begin{equation}
 R_{\rm gn}=2(I-K)S_0+sKS_{\rm W}.                       \label{eq:s7v11-Rgn-reduced}
\end{equation}
\end{corollary}

\begin{proof}
Equation \eqref{eq:s7v11-Haar-factor} is the reference measure
factorization, so its Radon--Nikodym density in the displayed group
coordinates is one and has zero exterior derivative.  Substitute
\(S_J=0\) in \eqref{eq:s7v10-Rgn}.
\end{proof}

\begin{remark}[Logarithm charts]
\label{rem:s7v11-log-chart}
Writing a chord locally as \(W=\exp\xi\) introduces the intrinsic Haar
density \(j(\xi)\,d\xi\).  If the chart is fixed independently of the
exterior source, this density has no source derivative and still
contributes no \(S_J\).  A moving chart would create a score and must be
listed separately; it is not part of maximal-tree gauge fixing.
\end{remark}

\section{Combined residual compiler}
\label{sec:s7v11-combined}

\begin{theorem}[Conditional ground/normalization estimate]
\label{thm:s7v11-combined}
Assume:
\begin{enumerate}[label=\textup{(\roman*)},leftmargin=3.5em]
\item the hypotheses of the V10 normalization compiler, with lower
      angle \(\gamma_\star\);
\item the hypotheses of
      \Cref{thm:s7v11-physical-ground}, including the fibre gap \(g\)
      and density bound \(W_\infty\);
\item the response Hamiltonian source factorization
      \(S_{\mathsf H}=B^*Y_{\mathsf H}\); and
\item fixed maximal-tree group coordinates, so \(S_J=0\).
\end{enumerate}
Then on the declared good core,
\begin{align}
 \|H^{-1/2}DR_{\rm gn}\|_{L^2(\mu)}
 &\le
 {2\sqrt{W_\infty}(1+\varepsilon)\over
  (1-\varepsilon)g}
 \|Y_{\mathsf H}\|_{L^2(\nu)}
 \notag\\
 &\quad+
 {s(1+\varepsilon)\over1-\varepsilon}
 \sqrt{1-\gamma_\star^2}\,
 \|G\|_{L^2(\mu)}.                                      \label{eq:s7v11-combined}
\end{align}
\end{theorem}

\begin{proof}
Apply \(D\) to \eqref{eq:s7v11-Rgn-reduced}.  Bound the first term by
twice \eqref{eq:s7v11-physical-ground}.  Bound the second by
\eqref{eq:s7v10-compiler} and multiply by \(s\).  Use the triangle
inequality.
\end{proof}

\begin{definition}[Moving-ground curvature card, MGCC]
\label{def:s7v11-MGCC}
MGCC consists of:
\begin{enumerate}[label=\textup{(G\arabic*)},leftmargin=4em]
\item a precise regulated response Hamiltonian and positive simple
      ground;
\item a worst-conditional lower bound \(g>0\) for its transformed fibre
      gap;
\item direct differentiation of every potential, kinetic, boundary,
      gauge-slice, and regulator coefficient;
\item a curvature-adjoint representation for the resulting
      \(S_{\mathsf H}\), with any unfactored term explicit;
\item a cofinal physical/ground density-ratio or weaker transfer
      estimate;
\item the V10 positive active-score lower angle;
\item insertion of \eqref{eq:s7v11-combined} into the MLC row/column
      sums and V3 strict shell budget; and
\item good-core, bad-set, gauge, reflection, observable, continuum, and
      spectral compatibility.
\end{enumerate}
\end{definition}

\begin{theorem}[What MGCC would close]
\label{thm:s7v11-MGCC}
Rows (G1)--(G6) give a non-circular curvature-square-root estimate for
the complete source-projected V9 physical score.  Rows (G7)--(G8), if
proved with strict cutoff-uniform constants, close the remaining local
CCSC source interface.  They do not by themselves construct the
continuum or its Hamiltonian spectrum.
\end{theorem}

\begin{proof}
The moving-ground identity is
\Cref{thm:s7v11-Poisson}; its curvature estimate is
\Cref{thm:s7v11-curvature-ground}; the Haar Jacobian is removed by
\Cref{cor:s7v11-no-J}; and the combined conditional estimate is
\Cref{thm:s7v11-combined}.  The downstream requirements are those of
CSAC, MLC, and RPTC.
\end{proof}

\begin{assessment}[Progress at seventh-series V11]
\label{ass:s7v11-status}
The ground score is no longer a formal derivative: it is the reduced
fibre-generator inverse in \eqref{eq:s7v11-vector}.  The
curvature-adjoint source survives that inverse with cost \(g^{-1}\), and
maximal-tree gauge contributes no Jacobian score in product-Haar group
coordinates.  The exact remaining quantitative inputs are now visible:
a uniform fibre gap, a physical/ground measure transfer, direct
factorization of every response-Hamiltonian derivative, and the V10
active-score lower angle.
\end{assessment}

\begin{decision}[V11 bottleneck decision]
\label{dec:s7v11-bottleneck}
The next upstream question is whether the fibre gap required above can
be obtained without assuming the desired field-theory mass gap.  V12
will distinguish the two notions rigorously.  It will derive a
finite-dimensional conditional Poincar\'e estimate from convexity only
where the Wilson Hessian actually has a positive lower bound, and will
test compact-group/Doeblin minorization on the full fibre.  Any constant
that degenerates with \(s\), the boundary, or the regulator will be
recorded rather than promoted to a cofinal gap.
\end{decision}

\begin{programme}[Seventh-series V12]
\label{prog:s7v11-V12}
V12 will compare Bakry--\'Emery, Holley--Stroock, and compact-kernel
minorization bounds for the owner fibre; compute their dependence on
the Wilson coupling and boundary oscillation; and decide which, if any,
can supply the \(g\) and \(W_\infty\) rows of MGCC without circular use
of the continuum mass gap.
\end{programme}

\begin{firewall}[Claim boundary after seventh-series V11]
\label{fire:s7v11-boundary}
V11 proves the exact moving-ground Poisson/resolvent equation, its sharp
gap estimates, curvature transport under explicit hypotheses, the
ground/physical measure-change ledger, and the unit-Jacobian
maximal-tree Haar factorization.

V11 does not prove a cutoff-uniform fibre gap, density ratio,
Hamiltonian-source factorization for a specified regulated response
operator, positive conditional angle, remaining martingale budgets,
continuum construction, or the Yang--Mills mass gap.
\end{firewall}

\paragraph{Parent-version record.}
V11 was formed from
\texttt{Renewal\_Yang\_Mills\_Mass\_Gap\_Research\_Notes\_V10.tex}.
The V10 parent had \(4\,054\) logical source lines, \(148\,874\) bytes,
and SHA--256
\[
 \texttt{2FF5799930960D4CCFA5D1FD6A6A93F9786F9CC1927F1A0996407C64F477C56E}.
\]

% =============================================================================
% END APPENDED SEVENTH-SERIES V11 MATERIAL
% =============================================================================

% =============================================================================
% BEGIN APPENDED SEVENTH-SERIES V12 MATERIAL
% =============================================================================

\chapter{Owner-fibre gap mechanisms and their exact losses}
\label{ch:s7v12}

V11 requires a worst-conditional spectral gap for a fixed
sixty-chord owner fibre.  Such a gap is local and finite dimensional; it
is not the desired infinite-volume Yang--Mills mass gap.  Nevertheless,
using the latter to justify the former would be circular.  This chapter
therefore derives three independent finite-fibre estimates and records
their parameter dependence.  The outcome is a sharp routing decision:
global oscillation methods are valid but lose exponentially in the
Wilson coupling, while global convexity is topologically unavailable.
The remaining viable route is a good-core local coercivity estimate
combined with a separately paid escape probability.

\section{The ground transform is a weighted Poincar\'e problem}
\label{sec:s7v12-ground}

Let \(Z=G^N\) with normalized product Haar measure \(m\), where in the
collar application \(G=\mathrm{SU}(3)\) and \(N=60\).  Fix the
bi-invariant product metric and let \(-\Delta\) be its nonnegative
Laplacian.  Consider
\begin{equation}
 \mathsf H=-a\Delta+V,\qquad a>0,                         \label{eq:s7v12-H}
\end{equation}
with positive normalized ground \(u\), ground energy \(\lambda\), and
\(d\nu=u^2dm\).  Let \(\mathcal L\) be its V11 ground transform.

\begin{theorem}[Exact transformed Dirichlet form]
\label{thm:s7v12-Dirichlet}
For every smooth real \(f\),
\begin{equation}
 \langle f,\mathcal Lf\rangle_{L^2(\nu)}
 =a\int_Z|\nabla f|^2\,d\nu.                             \label{eq:s7v12-Dirichlet}
\end{equation}
Consequently the fibre gap \(g\) in V11 is exactly the best constant in
\begin{equation}
 g\,\operatorname{Var}_\nu(f)
 \le a\int_Z|\nabla f|^2\,d\nu.                          \label{eq:s7v12-Poincare}
\end{equation}
\end{theorem}

\begin{proof}
By definition and \(d\nu=u^2dm\),
\[
 \langle f,\mathcal Lf\rangle_\nu
 =\langle uf,(\mathsf H-\lambda)(uf)\rangle_m.
\]
Expanding the kinetic form gives
\[
 a\int|u\nabla f+f\nabla u|^2dm
 +\int(V-\lambda)u^2f^2dm.
\]
The terms containing \(f^2\) and \(f\nabla f\) cancel after testing the
weak ground equation
\((-a\Delta+V-\lambda)u=0\) against \(uf^2\).  The remainder is
\(a\int u^2|\nabla f|^2dm\), proving
\eqref{eq:s7v12-Dirichlet}.  The variational characterization of the
first nonzero eigenvalue gives \eqref{eq:s7v12-Poincare}.
\end{proof}

\begin{lemma}[Product-Haar reference gap]
\label{lem:s7v12-product-gap}
If the one-factor Haar Poincar\'e gap is \(\lambda_G>0\), then the
product \(G^N\) has gap \(\lambda_G\), independently of \(N\).
\end{lemma}

\begin{proof}
Tensorization of conditional variance gives
\[
 \operatorname{Var}_m(f)
 \le {1\over\lambda_G}
 \sum_{j=1}^N\int|\nabla_jf|^2\,dm.
\]
This proves a product gap at least \(\lambda_G\).  A function depending
only on one factor and approaching a first-factor minimizer proves the
reverse inequality.
\end{proof}

\section{Oscillation comparison}
\label{sec:s7v12-oscillation}

\begin{theorem}[Two-sided density comparison]
\label{thm:s7v12-density}
Let \(d\nu=r\,dm\), \(\int r\,dm=1\), and suppose
\begin{equation}
 0<r_-\le r\le r_+<\infty.                               \label{eq:s7v12-ratio}
\end{equation}
Then
\begin{equation}
 \operatorname{Var}_\nu(f)
 \le {r_+\over r_-\lambda_G}
       \int|\nabla f|^2\,d\nu.                           \label{eq:s7v12-HS}
\end{equation}
For \(\mathsf H\) in \eqref{eq:s7v12-H},
\begin{equation}
 g\ge a\lambda_G\,{r_-\over r_+}.                        \label{eq:s7v12-gap-ratio}
\end{equation}
\end{theorem}

\begin{proof}
Using the variational form of variance,
\[
 \operatorname{Var}_\nu(f)
 =\inf_c\int|f-c|^2r\,dm
 \le r_+\operatorname{Var}_m(f).
\]
Apply \Cref{lem:s7v12-product-gap}, then use
\(\int|\nabla f|^2dm\le r_-^{-1}\int|\nabla f|^2d\nu\).
Multiply by \(a\) and invoke
\Cref{thm:s7v12-Dirichlet}.
\end{proof}

\begin{corollary}[Ground-oscillation bound]
\label{cor:s7v12-u-osc}
For \(r=u^2\),
\begin{equation}
 g\ge
 a\lambda_G
 \exp\{-2\operatorname{osc}_Z(\log u)\}.                 \label{eq:s7v12-u-osc}
\end{equation}
\end{corollary}

\begin{proof}
The normalizing constant cancels in
\(r_+/r_-\), which equals
\(\exp\{2\operatorname{osc}(\log u)\}\).  Apply
\eqref{eq:s7v12-gap-ratio}.
\end{proof}

\begin{assessment}[What oscillation comparison costs]
\label{ass:s7v12-oscillation}
The number of collar chords does not degrade the Haar reference gap,
because it tensorizes.  All loss is concentrated in the ground-density
oscillation.  Any Harnack estimate used to bound
\(\operatorname{osc}\log u\) must therefore be tracked in the Wilson
coupling, the kinetic coefficient, and the worst boundary.  Merely
noting that \(u\) is continuous and positive proves a fixed-fibre
constant, not a cofinal one.
\end{assessment}

\section{A direct min--max estimate}
\label{sec:s7v12-minmax}

Let \(0=\lambda_0^\Delta<\lambda_1^\Delta=\lambda_G\le\cdots\)
be the spectrum of \(-\Delta\) on \(Z\), and let
\(\lambda_0^{\mathsf H}<\lambda_1^{\mathsf H}\le\cdots\)
be the spectrum of \(\mathsf H\).

\begin{theorem}[Bounded-potential fibre gap]
\label{thm:s7v12-minmax}
For bounded real \(V\),
\begin{equation}
 \lambda_1^{\mathsf H}-\lambda_0^{\mathsf H}
 \ge a\lambda_G-\operatorname{osc}_ZV.                   \label{eq:s7v12-minmax}
\end{equation}
\end{theorem}

\begin{proof}
The form inequalities
\[
 -a\Delta+\inf V\le\mathsf H\le-a\Delta+\sup V
\]
and the min--max principle give
\[
 \lambda_1^{\mathsf H}\ge a\lambda_G+\inf V.
\]
Testing the ground Rayleigh quotient with the constant function gives
\(\lambda_0^{\mathsf H}\le\int V\,dm\le\sup V\).
Subtract.
\end{proof}

\begin{corollary}[Why the crude Wilson oscillation route fails]
\label{cor:s7v12-Wilson-osc}
For an incident Wilson potential
\[
 V=s\sum_{p\in{\cal P}(C)}\bigl(3-\Re\Tr W_p\bigr),
\]
any bound obtained only from the range of each plaquette term has the
form
\begin{equation}
 g\ge a\lambda_G-s\,|{\cal P}(C)|\,c_{\mathrm{osc}},
 \qquad |{\cal P}(C)|=250,                               \label{eq:s7v12-Wilson-osc}
\end{equation}
for a fixed group constant \(c_{\mathrm{osc}}>0\).  It becomes vacuous
when the Wilson coefficient grows unless \(a\) grows at least
proportionally.
\end{corollary}

\begin{proof}
The oscillation of a sum is at most the sum of oscillations, so
\(\operatorname{osc}V\le
s|{\cal P}(C)|c_{\mathrm{osc}}\).  Insert this in
\eqref{eq:s7v12-minmax}.  Once the right side is nonpositive it supplies
no spectral lower bound.
\end{proof}

\begin{warning}[A negative lower bound is not gap collapse]
\label{warn:s7v12-vacuous}
The failure of \eqref{eq:s7v12-Wilson-osc} at large \(s\) does not show
that the true fibre gap closes.  It shows only that a global
bounded-perturbation estimate discards the local well geometry that can
make a strongly coupled finite-dimensional oscillator more coercive.
\end{warning}

\section{Global convexity is unavailable on the compact fibre}
\label{sec:s7v12-convexity}

\begin{theorem}[No smooth strictly convex potential on a closed fibre]
\label{thm:s7v12-no-convex}
Let \(Z\) be a closed Riemannian manifold of positive dimension.  There
is no smooth function \(V\) satisfying
\begin{equation}
 \operatorname{Hess}V\succeq\kappa\,g_Z
 \quad\hbox{everywhere for any }\kappa>0.                 \label{eq:s7v12-global-convex}
\end{equation}
\end{theorem}

\begin{proof}
Taking the metric trace of \eqref{eq:s7v12-global-convex} gives
\(\Delta V\ge\kappa\dim Z\).  Integration over the closed manifold
would give
\[
 0=\int_Z\Delta V\,d\operatorname{vol}
 \ge\kappa\dim Z\,\operatorname{vol}(Z)>0,
\]
a contradiction.
\end{proof}

\begin{corollary}[Scope of the V1 Hessian coercivity]
\label{cor:s7v12-local-only}
The positive incident Wilson Hessian proved in V1 can hold only on its
declared fully good chart.  A Bakry--\'Emery argument based on that
Hessian must stop at the chart boundary and pay an escape/bad-set term;
it cannot be extended to the whole compact owner fibre with the same
strict Hessian constant.
\end{corollary}

\begin{proof}
Otherwise the Wilson potential restricted to the compact fibre would
satisfy \eqref{eq:s7v12-global-convex}, contradicting
\Cref{thm:s7v12-no-convex}.
\end{proof}

\section{Compact-kernel minorization}
\label{sec:s7v12-Doeblin}

\begin{theorem}[Exact Doeblin-to-gap implication]
\label{thm:s7v12-Doeblin}
Let \(\mathcal L\) be a nonnegative self-adjoint Markov generator in
\(L^2(\nu)\), and suppose \(T_\tau=e^{-\tau\mathcal L}\) has a transition
density \(p_\tau(z,z')\) with respect to the invariant probability
\(\nu\).  If
\begin{equation}
 p_\tau(z,z')\ge\delta>0
 \quad\hbox{for all }z,z',                                \label{eq:s7v12-minorization}
\end{equation}
then
\begin{equation}
 \operatorname{gap}(\mathcal L)
 \ge-{1\over\tau}\log(1-\delta).                         \label{eq:s7v12-Doeblin-gap}
\end{equation}
\end{theorem}

\begin{proof}
Let \(\Pi f=\int f\,d\nu\).  The kernel
\[
 S={T_\tau-\delta\Pi\over1-\delta}
\]
is positive, preserves constants, and preserves \(\nu\); hence Jensen's
inequality makes it an \(L^2(\nu)\) contraction.  On \(1^\perp\),
\(T_\tau=(1-\delta)S\), so
\(\|T_\tau\|_{1^\perp}\le1-\delta\).  Spectral calculus gives
\(\|T_\tau\|_{1^\perp}=e^{-\tau g}\), which proves
\eqref{eq:s7v12-Doeblin-gap}.
\end{proof}

\begin{assessment}[Minorization is pointwise, not automatically uniform]
\label{ass:s7v12-Doeblin}
For each fixed uniformly elliptic operator on a compact connected
fibre, a positive continuous heat kernel has a positive minimum at
fixed \(\tau>0\).  The resulting \(\delta\) may nevertheless tend to
zero when the kinetic coefficient, potential, boundary, or regulator
moves.  Compactness therefore proves fixed-parameter positivity but not
the uniform row required by MGCC.
\end{assessment}

\section{The local-core alternative}
\label{sec:s7v12-local}

\begin{definition}[Local fibre-gap card, LFGC]
\label{def:s7v12-LFGC}
LFGC consists of:
\begin{enumerate}[label=\textup{(F\arabic*)},leftmargin=4em]
\item a specified response Hamiltonian
      \(-a\Delta+V\) and its precise cutoff scaling;
\item the fully good collar domain on which the V1/V8 curvature Hessian
      is positive;
\item a reflected, stopped, or IMS-localized Poincar\'e estimate inside
      that domain with an explicit boundary convention;
\item a capacity or escape estimate for reaching its complement under
      both the ground and physical fibre measures;
\item a gluing inequality that combines the local gap and escape debit
      without using global convexity;
\item a uniform physical/ground density transfer;
\item insertion of the resulting \(g\) into the V11 combined residual
      estimate; and
\item the V10 lower-angle, MLC, shell, bad-set, and continuum contacts.
\end{enumerate}
\end{definition}

\begin{theorem}[Decision supplied by the three tests]
\label{thm:s7v12-decision}
The product-Haar estimate supplies an \(N\)-independent reference gap.
The density-comparison and Doeblin routes supply rigorous fibre gaps
only with, respectively, a density oscillation or minorization constant.
The min--max route is useful only in the weak-oscillation regime.
Global Wilson convexity cannot supply the missing row.  Hence outside a
separately proved uniform oscillation/minorization estimate, LFGC is the
remaining non-circular route based on the existing local curvature
coercivity.
\end{theorem}

\begin{proof}
The four assertions are
\Cref{lem:s7v12-product-gap,thm:s7v12-density,
thm:s7v12-Doeblin,thm:s7v12-minmax,thm:s7v12-no-convex}.
The V1 Hessian is local by
\Cref{cor:s7v12-local-only}, so its use requires precisely the
localization and escape rows listed in LFGC.
\end{proof}

\begin{assessment}[Progress at seventh-series V12]
\label{ass:s7v12-status}
The V11 gap hypothesis has been separated from the ultimate mass gap
and reduced to explicit finite-fibre constants.  Haar tensorization is
uniform in the sixty owner chords, but neither positivity of the ground
nor compactness makes its perturbation uniform.  A global strictly
convex Wilson proof is impossible on the compact fibre.  The next
calculation must therefore quantify confinement to the already
constructed good collar and glue its local Poincar\'e estimate to the
complement.
\end{assessment}

\begin{decision}[V12 bottleneck decision]
\label{dec:s7v12-bottleneck}
V13 will formulate and prove a two-region Poincar\'e gluing lemma with
all constants exposed.  It will then identify the exact collar
probability, overlap, and local Neumann-gap data needed from the Wilson
measure.  If the lemma demands an uncontrolled inverse rare-event
probability, the route will be rejected and the programme will return
upstream to a path-coupling or canonical-path fibre estimate.
\end{decision}

\begin{programme}[Seventh-series V13]
\label{prog:s7v12-V13}
V13 will decompose variance into good-core, bad-region, and mean-jump
terms; prove a quantitative gluing estimate through an overlap collar;
and test its constants against the V1--V3 good/bad shell budgets.
\end{programme}

\begin{firewall}[Claim boundary after seventh-series V12]
\label{fire:s7v12-boundary}
V12 proves the exact transformed Dirichlet form, the tensorized Haar
gap, two-sided density comparison, a min--max gap bound, the
compact-fibre no-global-convexity theorem, and the
Doeblin-to-gap implication.

V12 does not prove uniform ground oscillation, minorization, a local
core gluing estimate, the MGCC gap row, the active-score angle,
remaining martingale budgets, continuum construction, or the
Yang--Mills mass gap.
\end{firewall}

\paragraph{Parent-version record.}
V12 was formed from
\texttt{Renewal\_Yang\_Mills\_Mass\_Gap\_Research\_Notes\_V11.tex}.
The V11 parent had \(4\,561\) logical source lines, \(168\,189\) bytes,
and SHA--256
\[
 \texttt{3D3749B3C9A19E112522418C040F900820381D018ECC7EC03C90D43C7B2B35A8}.
\]

% =============================================================================
% END APPENDED SEVENTH-SERIES V12 MATERIAL
% =============================================================================

% =============================================================================
% BEGIN APPENDED SEVENTH-SERIES V13 MATERIAL
% =============================================================================

\chapter{Buffered Poincar\'e gluing and the rare-bad-set obstruction}
\label{ch:s7v13}

V12 leaves one possible use of the local collar Hessian: prove a
Poincar\'e inequality on a good core and glue it across an overlap to
the rest of the owner fibre.  This chapter proves the exact elementary
gluing lemma and then tests the data supplied by the present programme.
The test rules out a tempting shortcut.  Small bad-set probability and
a good-core gap do not control the global fibre gap without an
interface conductance or a second local inequality.

\section{An overlap-cover gluing lemma}
\label{sec:s7v13-gluing}

Let \((\Omega,\mu)\) be a probability space equipped with a nonnegative
local energy density \(\Gamma(f)\).  For a measurable set \(A\) of
positive measure, write
\[
 f_A={1\over\mu(A)}\int_Af\,d\mu.
\]
The boundary convention implicit in \(\Gamma\) is part of any local
Poincar\'e hypothesis below.

\begin{theorem}[Two-region overlap gluing]
\label{thm:s7v13-gluing}
Let \(A,B\subset\Omega\) satisfy
\[
 A\cup B=\Omega,\qquad C:=A\cap B,\qquad\mu(C)>0.
\]
Suppose that for every function in a common form core,
\begin{align}
 \int_A|f-f_A|^2d\mu
 &\le C_A\int_A\Gamma(f)d\mu,                            \label{eq:s7v13-local-A}\\
 \int_B|f-f_B|^2d\mu
 &\le C_B\int_B\Gamma(f)d\mu.                            \label{eq:s7v13-local-B}
\end{align}
Then
\begin{align}
 \operatorname{Var}_\mu(f)
 &\le
 2\left(1+{\mu(A)\over\mu(C)}\right)C_A
       \int_A\Gamma(f)d\mu \notag\\
 &\quad+
 2\left(1+{\mu(B)\over\mu(C)}\right)C_B
       \int_B\Gamma(f)d\mu                               \label{eq:s7v13-gluing}\\
 &\le
 4\max\left\{
 \left(1+{\mu(A)\over\mu(C)}\right)C_A,\,
 \left(1+{\mu(B)\over\mu(C)}\right)C_B
 \right\}
 \int_\Omega\Gamma(f)d\mu.                               \label{eq:s7v13-gluing-max}
\end{align}
\end{theorem}

\begin{proof}
Variance is the least squared distance to a constant, so
\[
 \operatorname{Var}_\mu(f)
 \le\int_\Omega|f-f_C|^2d\mu
 \le\int_A|f-f_C|^2d\mu+\int_B|f-f_C|^2d\mu.
\]
For the \(A\) term,
\[
 \int_A|f-f_C|^2d\mu
 \le2\int_A|f-f_A|^2d\mu
   +2\mu(A)|f_A-f_C|^2.
\]
Moreover,
\[
 |f_A-f_C|^2
 =\left|{1\over\mu(C)}\int_C(f_A-f)d\mu\right|^2
 \le {1\over\mu(C)}\int_C|f-f_A|^2d\mu
 \le {1\over\mu(C)}\int_A|f-f_A|^2d\mu.
\]
Apply \eqref{eq:s7v13-local-A}.  The identical calculation on \(B\)
proves \eqref{eq:s7v13-gluing}.  Finally
\[
 \int_A\Gamma(f)d\mu+\int_B\Gamma(f)d\mu
 \le2\int_\Omega\Gamma(f)d\mu,
\]
which gives \eqref{eq:s7v13-gluing-max}.
\end{proof}

\begin{corollary}[Quantitative buffered-core form]
\label{cor:s7v13-buffer}
Let \(A\) be the good collar, let \(B\) contain the bad region and a
buffer inside the collar, and let \(C=A\cap B\) be that buffer.  If
\[
 \mu(C)\ge\theta_A\mu(A),\qquad
 \mu(C)\ge\theta_B\mu(B)
\]
with \(\theta_A,\theta_B>0\), then the global Poincar\'e constant is at
most
\begin{equation}
 4\max\{(1+\theta_A^{-1})C_A,\,
        (1+\theta_B^{-1})C_B\}.                          \label{eq:s7v13-buffer}
\end{equation}
\end{corollary}

\begin{proof}
Insert the two overlap-ratio bounds in
\eqref{eq:s7v13-gluing-max}.
\end{proof}

\begin{assessment}[What the lemma actually asks for]
\label{ass:s7v13-gluing}
The local collar Hessian can in principle supply \(C_A\).  It supplies
neither \(C_B\) on the buffer-plus-bad region nor a lower bound on the
overlap ratios.  If \(\mu(C)\) is very small, the explicit inverse
overlap factor in \eqref{eq:s7v13-gluing-max} destroys the estimate.
The lemma is therefore a compiler for interface data, not a method for
creating those data.
\end{assessment}

\section{Rarity alone cannot replace interface conductance}
\label{sec:s7v13-no-go}

\begin{theorem}[Rare bad set plus a good local gap is insufficient]
\label{thm:s7v13-rare-no-go}
Fix any \(\epsilon\in(0,1)\).  There is a family of reversible
three-state Dirichlet forms with:
\begin{enumerate}[label=\textup{(\roman*)},leftmargin=3.5em]
\item a bad state of invariant probability \(\epsilon\);
\item a good two-state restriction whose Poincar\'e constant is fixed
      independently of the family parameter; and
\item global spectral gap tending to zero.
\end{enumerate}
\end{theorem}

\begin{proof}
Take
\(\Omega=\{g_1,g_2,b\}\) and
\[
 \mu(g_1)=\mu(g_2)={1-\epsilon\over2},\qquad\mu(b)=\epsilon.
\]
For \(\delta>0\), define
\begin{equation}
 {\cal E}_\delta(f,f)
 =(f(g_1)-f(g_2))^2
 +\delta(f(g_2)-f(b))^2.                                \label{eq:s7v13-three-state}
\end{equation}
This is a reversible Dirichlet form with invariant measure \(\mu\)
after assigning the corresponding symmetric edge conductances.
On \(A=\{g_1,g_2\}\), the conditional measure is uniform and
\[
 \operatorname{Var}_{\mu(\cdot\mid A)}(f)
 ={1\over4}(f(g_1)-f(g_2))^2,
\]
so the restricted Poincar\'e constant relative to the first term of
\eqref{eq:s7v13-three-state} is \(1/4\), independently of \(\delta\).
For the test function \(f=\mathbf1_{\{b\}}\),
\[
 \operatorname{Var}_\mu(f)=\epsilon(1-\epsilon),
 \qquad
 {\cal E}_\delta(f,f)=\delta.
\]
The global variational gap is therefore at most
\[
 {\delta\over\epsilon(1-\epsilon)}\longrightarrow0.
\]
\end{proof}

\begin{corollary}[The V1--V3 bad probability is not a fibre-gap proof]
\label{cor:s7v13-probability-not-gap}
Even a cutoff-uniform upper bound on the probability of leaving the V1
good collar, together with its local Hessian coercivity, cannot by
itself supply the V11 fibre gap.  A conductance/capacity estimate across
the collar boundary or a second local Poincar\'e inequality is
logically necessary.
\end{corollary}

\begin{proof}
\Cref{thm:s7v13-rare-no-go} keeps both the bad probability and good
local gap fixed while the interface conductance tends to zero.
Therefore no lower bound depending only on those two data can exist.
\end{proof}

\section{Application ledger for the Wilson collar}
\label{sec:s7v13-Wilson}

Choose defect thresholds
\[
 0<\delta_{\rm in}<\delta_{\rm out}
\]
and define an inner good core \(A\) by defect below
\(\delta_{\rm out}\).  Let \(B\) contain the complement of the
\(\delta_{\rm in}\) core, so the annulus between the two thresholds is
\(C=A\cap B\).  This construction creates a genuine overlap but proves
no estimate on it.

\begin{definition}[Buffered Wilson fibre card, BWFC]
\label{def:s7v13-BWFC}
BWFC consists of:
\begin{enumerate}[label=\textup{(B\arabic*)},leftmargin=4em]
\item explicit inner and outer defect thresholds inside the V8
      logarithm and relative-curvature core;
\item a local Neumann or reflected Poincar\'e constant \(C_A\) on the
      outer core, derived from the incident Hessian with its boundary
      convention stated;
\item a Poincar\'e constant \(C_B\) on the annulus-plus-bad region, or a
      canonical-path/capacity substitute controlling passage to the
      core;
\item lower bounds for
      \(\mu(C)/\mu(A)\) and \(\mu(C)/\mu(B)\), not merely an upper bound
      for \(\mu(B\setminus A)\);
\item uniformity of these data under every allowed exterior boundary,
      Wilson coupling, and regulator;
\item transfer of the resulting gap through the V11 ground/physical
      measure comparison;
\item the V10 active-score angle and V5 martingale Schur budgets; and
\item bad-set, gauge, observable, continuum, and spectral contacts.
\end{enumerate}
\end{definition}

\begin{theorem}[BWFC to the moving-ground gap row]
\label{thm:s7v13-BWFC}
If rows (B1)--(B5) hold for the ground measure \(\nu\), then
\Cref{thm:s7v13-gluing} gives a worst-conditional fibre gap
\begin{equation}
 g\ge
 {a\over
 4\max\{(1+\theta_A^{-1})C_A,\,
        (1+\theta_B^{-1})C_B\}},                         \label{eq:s7v13-BWFC-gap}
\end{equation}
where \(a\Gamma\) is the V12 transformed Dirichlet density.
Rows (B6)--(B8) would then insert it into MGCC.
\end{theorem}

\begin{proof}
Apply \Cref{cor:s7v13-buffer} to the unscaled gradient form.  Multiplying
that form by \(a\) makes the spectral gap the reciprocal of the
resulting Poincar\'e constant, giving
\eqref{eq:s7v13-BWFC-gap}.  The downstream statement is exactly the
V11 MGCC interface.
\end{proof}

\begin{firewall}[No hidden use of a rare-event inverse]
\label{fire:s7v13-rare}
Replacing the overlap lower bounds in BWFC by
\(\mu(B\setminus A)\le p\) is invalid.  The latter gives no positive
lower bound on \(\mu(C)\), and
\Cref{thm:s7v13-rare-no-go} shows that an arbitrarily thin interface can
destroy the global gap while \(p\) stays fixed and small.
\end{firewall}

\begin{assessment}[Progress at seventh-series V13]
\label{ass:s7v13-status}
The local-to-global fibre step now has an exact theorem and an exact
obstruction.  The theorem exposes two data absent from the present
collar ledger: mixing on the buffer-plus-bad region and a lower overlap
mass.  The existing V1--V3 defect and bad-event estimates are
one-sided and cannot manufacture either datum.  Thus the pure
local-Hessian-plus-rarity route is incomplete.
\end{assessment}

\begin{decision}[V13 bottleneck decision]
\label{dec:s7v13-bottleneck}
Go upstream to a dynamics that crosses the interface.  The archived
programme already contains strong one-link conditional estimates and
exact conditional projections.  V14 will derive a block-dynamics
Poincar\'e comparison: a uniform one-link conditional gap plus a strict
interdependence matrix yields a fibre gap.  This avoids inverse overlap
mass if the influence norm is genuinely below one; if the exact
250-plaquette incidence makes that norm non-strict, the failure will be
recorded.
\end{decision}

\begin{programme}[Seventh-series V14]
\label{prog:s7v13-V14}
V14 will prove the conditional-variance block-dynamics comparison,
define the exact owner-chord influence matrix, and reduce its strictness
to computable source derivatives of one-link conditional laws.  It will
not substitute a Dobrushin constant from a different measure or a
different ground transform.
\end{programme}

\begin{firewall}[Claim boundary after seventh-series V13]
\label{fire:s7v13-boundary}
V13 proves a quantitative overlap-cover Poincar\'e gluing lemma and a
three-state no-go theorem showing that rare bad mass plus a good local
gap is insufficient.  It identifies the complete buffered data that
would supply the V11 fibre-gap row.

V13 does not prove those interface data, a uniform owner-fibre gap, the
active-score angle, remaining martingale budgets, continuum
construction, or the Yang--Mills mass gap.
\end{firewall}

\paragraph{Parent-version record.}
V13 was formed from
\texttt{Renewal\_Yang\_Mills\_Mass\_Gap\_Research\_Notes\_V12.tex}.
The V12 parent had \(4\,951\) logical source lines, \(182\,456\) bytes,
and SHA--256
\[
 \texttt{41C9DC129ECA8276AE0CF206213D3B1A9C59CAEEF0E58E7FBFB4BC0DC6D6916A}.
\]

% =============================================================================
% END APPENDED SEVENTH-SERIES V13 MATERIAL
% =============================================================================

% =============================================================================
% BEGIN APPENDED SEVENTH-SERIES V14 MATERIAL
% =============================================================================

\chapter{Block heat-bath dynamics and the exact influence Gram matrix}
\label{ch:s7v14}

V13 shows that a local collar gap cannot be glued through a rare bad
region without interface information.  A heat-bath dynamics supplies
such information only if its conditional blocks collectively see every
nonconstant mode.  This chapter gives an exact Hilbert-space comparison
which separates two questions: coercivity inside each conditional
block, and mixing among the blocks.  The second question is encoded by
a positive operator-valued Gram matrix; a scalar diagonal-dominance
test is a sufficient, auditable influence criterion.

\section{Conditional block generator}
\label{sec:s7v14-generator}

Let \((\Omega,\mu)\) be the finite-cutoff physical probability space and
let \(\Pi f=\int f\,d\mu\).  For conditional blocks
\(\Lambda_1,\ldots,\Lambda_M\), let
\begin{equation}
 K_i=\mathbb E_\mu[\cdot\mid\mathcal F_{\Lambda_i^c}],
 \qquad D_i=I-K_i.                                       \label{eq:s7v14-KD}
\end{equation}
Thus \(K_i\) is an orthogonal projection and
\[
 \|D_if\|_2^2
 =\mathbb E_\mu\operatorname{Var}
   (f\mid\mathcal F_{\Lambda_i^c}).
\]
Define the unnormalized heat-bath generator
\begin{equation}
 {\cal A}_{\rm hb}:=\sum_{i=1}^M D_i.                    \label{eq:s7v14-Ahb}
\end{equation}

\begin{lemma}[Exact heat-bath Dirichlet form]
\label{lem:s7v14-form}
For every \(f\in L^2(\mu)\),
\begin{equation}
 \langle f,{\cal A}_{\rm hb}f\rangle
 =\sum_{i=1}^M\|D_if\|_2^2.                              \label{eq:s7v14-form}
\end{equation}
Its kernel is
\[
 \ker{\cal A}_{\rm hb}=\bigcap_{i=1}^M\operatorname{Ran}K_i.
\]
\end{lemma}

\begin{proof}
Each \(D_i\) is an orthogonal projection, so
\(\langle f,D_if\rangle=\|D_if\|^2\).  Summation proves
\eqref{eq:s7v14-form}.  A sum of nonnegative terms vanishes exactly
when every \(D_if=0\), equivalently when
\(f\in\operatorname{Ran}K_i\) for every \(i\).
\end{proof}

\begin{definition}[Block-dynamics gap]
\label{def:s7v14-lambda}
Assuming the common invariant subspace is the constants, define
\begin{equation}
 \lambda_{\rm hb}
 :=
 \inf_{\substack{f\perp1\\f\ne0}}
 {\sum_i\|D_if\|_2^2\over\|f\|_2^2}.                    \label{eq:s7v14-lambda}
\end{equation}
\end{definition}

\section{From conditional fibre gaps to the owner-fibre gap}
\label{sec:s7v14-local-global}

Suppose a local diffusion form \({\cal E}_i\) acts only in block
\(\Lambda_i\), and set \({\cal E}=\sum_i{\cal E}_i\).

\begin{theorem}[Exact block-dynamics comparison]
\label{thm:s7v14-comparison}
Assume that, for almost every exterior configuration,
\begin{equation}
 {\cal E}_i(f,f)
 \ge g_i\,
 \mathbb E_\mu\!\left[
   \operatorname{Var}(f\mid\mathcal F_{\Lambda_i^c})
 \right],
 \qquad g_i\ge g_{\rm loc}>0.                            \label{eq:s7v14-local-gap}
\end{equation}
Then
\begin{equation}
 {\cal E}(f,f)
 \ge g_{\rm loc}\lambda_{\rm hb}
       \operatorname{Var}_\mu(f).                        \label{eq:s7v14-comparison}
\end{equation}
\end{theorem}

\begin{proof}
Sum \eqref{eq:s7v14-local-gap} and use
\eqref{eq:s7v14-form}:
\[
 {\cal E}(f,f)\ge g_{\rm loc}\sum_i\|D_if\|_2^2.
\]
Apply \eqref{eq:s7v14-lambda} to \(f-\Pi f\); every \(D_i\) annihilates
constants.
\end{proof}

\begin{warning}[Two independent strictness requirements]
\label{warn:s7v14-two-gaps}
Neither factor in \(g_{\rm loc}\lambda_{\rm hb}\) can be omitted.
A conditional block can have an excellent internal gap while the
collection leaves a common nonconstant mode invariant, making
\(\lambda_{\rm hb}=0\).  Conversely, block geometry can mix all modes
while a worst-conditional internal gap tends to zero.  The one-link
competing-centre-well obstruction recorded in V1 concerns
\(g_{\rm loc}\), not merely \(\lambda_{\rm hb}\).
\end{warning}

\section{An exact operator-valued influence matrix}
\label{sec:s7v14-Gram}

Choose any orthogonal resolution of the nonconstant subspace,
\begin{equation}
 I-\Pi=\sum_{j\in{\cal J}}\Delta_j,\qquad
 \Delta_j\Delta_k=0\quad(j\ne k),                        \label{eq:s7v14-resolution}
\end{equation}
with strong convergence if \({\cal J}\) is infinite.  For
\(h_j\in\operatorname{Ran}\Delta_j\), define the block operators
\begin{equation}
 {\mathbb G}_{jk}
 :=
 \sum_{i=1}^M
 \Delta_jD_i\Delta_k:
 \operatorname{Ran}\Delta_k\longrightarrow
 \operatorname{Ran}\Delta_j.                            \label{eq:s7v14-Gjk}
\end{equation}

\begin{theorem}[Influence Gram identity]
\label{thm:s7v14-Gram}
The block matrix
\(\mathbb G=(\mathbb G_{jk})_{j,k\in{\cal J}}\) is positive
self-adjoint on
\(\bigoplus_j\operatorname{Ran}\Delta_j\), and for
\(f=\sum_jh_j\perp1\),
\begin{equation}
 \sum_{i=1}^M\|D_if\|_2^2
 =
 \sum_{j,k}
 \langle h_j,\mathbb G_{jk}h_k\rangle.                  \label{eq:s7v14-Gram-identity}
\end{equation}
Consequently
\begin{equation}
 \lambda_{\rm hb}=\inf\sigma(\mathbb G).                 \label{eq:s7v14-Gram-gap}
\end{equation}
\end{theorem}

\begin{proof}
Since \(D_i=D_i^*=D_i^2\),
\[
 \sum_i\|D_if\|^2
 =\sum_i\langle f,D_if\rangle.
\]
Insert \(f=\sum_jh_j\) on both sides of each inner product and use
\(\Delta_jh_j=h_j\), obtaining
\eqref{eq:s7v14-Gram-identity}.  It also shows that \(\mathbb G\) is
the compression of the nonnegative self-adjoint operator
\(\sum_iD_i\) under the unitary orthogonal decomposition
\eqref{eq:s7v14-resolution}.  Hence it is positive self-adjoint.
The Rayleigh characterization is exactly
\eqref{eq:s7v14-lambda}.
\end{proof}

\begin{theorem}[Scalar diagonal-dominance certificate]
\label{thm:s7v14-Gershgorin}
Suppose numbers \(a_j\ge0\) and symmetric
\(c_{jk}=c_{kj}\ge0\) satisfy
\begin{align}
 \langle h,\mathbb G_{jj}h\rangle
 &\ge a_j\|h\|^2,
 &&h\in\operatorname{Ran}\Delta_j,                       \label{eq:s7v14-diagonal}\\
 \|\mathbb G_{jk}\|&\le c_{jk},
 &&j\ne k.                                               \label{eq:s7v14-offdiagonal}
\end{align}
Then
\begin{equation}
 \lambda_{\rm hb}
 \ge
 \inf_j\left(a_j-\sum_{k\ne j}c_{jk}\right).             \label{eq:s7v14-Gershgorin}
\end{equation}
\end{theorem}

\begin{proof}
Let \(x_j=\|h_j\|\).  From
\Cref{thm:s7v14-Gram},
\[
 \langle h,\mathbb Gh\rangle
 \ge\sum_ja_jx_j^2-\sum_{j\ne k}c_{jk}x_jx_k.
\]
Use \(2x_jx_k\le x_j^2+x_k^2\) and symmetry of \(c\).  The coefficient
left on \(x_j^2\) is at least
\(a_j-\sum_{k\ne j}c_{jk}\).  Sum and apply
\eqref{eq:s7v14-Gram-gap}.
\end{proof}

\begin{remark}[Why this matrix is not a guessed Dobrushin matrix]
\label{rem:s7v14-exact}
The blocks \(\mathbb G_{jk}\) are defined in the actual physical
\(L^2(\mu)\) space and use its actual conditional projections.  Their
smallest spectral value is the heat-bath gap itself.  Replacing them by
scalar \(c_{jk}\) is optional and loses information; any such bound must
be proved for the same ground transform and boundary family.
\end{remark}

\section{Differentiating a one-block conditional law}
\label{sec:s7v14-score}

Let \(K_i\) integrate the variable in block \(\Lambda_i\), and let
\(X_j\) differentiate an exterior variable belonging to another block.
Write the conditional density as \(\rho_i\) and set
\begin{equation}
 \sigma_{ij}:=
 X_j\log\rho_i-K_i(X_j\log\rho_i).                       \label{eq:s7v14-sigma}
\end{equation}

\begin{theorem}[Exact source-to-block commutator]
\label{thm:s7v14-commutator}
For every differentiable \(f\),
\begin{equation}
 X_jK_if-K_iX_jf=K_i(f\sigma_{ij}).                      \label{eq:s7v14-commutator}
\end{equation}
Fibrewise,
\begin{equation}
 |X_jK_if-K_iX_jf|^2
 \le
 \operatorname{Var}_i(f)\,
 K_i(\sigma_{ij}^2).                                    \label{eq:s7v14-Fisher}
\end{equation}
\end{theorem}

\begin{proof}
The first identity is the V9 normalized conditional-score formula.
Since \(K_i\sigma_{ij}=0\),
\[
 K_i(f\sigma_{ij})
 =K_i((f-K_if)\sigma_{ij}).
\]
Conditional Cauchy--Schwarz proves
\eqref{eq:s7v14-Fisher}.
\end{proof}

\begin{definition}[Physical score-influence coefficient]
\label{def:s7v14-chi}
For unit exterior tangents in block \(j\), define
\begin{equation}
 \chi_{ij}^2
 :=
 \operatorname*{ess\,sup}_{\rm exterior}
 \sup_{\substack{X_j\in T\Lambda_j\\|X_j|=1}}
 K_i(\sigma_{ij,X_j}^2).                                \label{eq:s7v14-chi}
\end{equation}
Set \(\chi_{ij}=0\) when the conditional density of block \(i\) is
independent of block \(j\).
\end{definition}

\begin{corollary}[Incidence sparsity]
\label{cor:s7v14-sparsity}
For a Wilson action in fixed group coordinates,
\(\chi_{ij}=0\) whenever no plaquette contains links from both blocks
\(\Lambda_i\) and \(\Lambda_j\).  For an incident pair, the score is the
centered derivative of exactly the shared plaquette terms, together
with any explicitly moving ground or regulator score.
\end{corollary}

\begin{proof}
Terms in the conditional log density which do not contain block \(j\)
have zero \(X_j\)-derivative.  If no action or ground/regulator term
couples the two blocks, \(X_j\log\rho_i=0\), hence
\(\sigma_{ij}=0\).  Otherwise the derivative and conditional centering
in \eqref{eq:s7v14-sigma} leave precisely the shared terms.
\end{proof}

\begin{firewall}[Fisher influence is not yet Gram strictness]
\label{fire:s7v14-Fisher}
Equation \eqref{eq:s7v14-Fisher} controls a derivative commutator.  To
deduce \eqref{eq:s7v14-offdiagonal}, one must relate the selected
martingale levels \(\Delta_j\) to block derivatives by a proved local
Poincar\'e or Riesz estimate.  V1 already records that an integrated
conditional Fisher bound cannot be promoted to an
essential-supremum influence coefficient.  Neither step may be skipped.
\end{firewall}

\section{The one-link partition fails the declared uniformity test}
\label{sec:s7v14-one-link}

\begin{proposition}[One-link no-go propagates to the comparison]
\label{prop:s7v14-one-link}
Consider the partition of the collar into its sixty individual chord
variables.  The existing competing-centre-well construction implies
that the hypothesis
\[
 \inf_{\text{allowed exterior}}\min_i g_i>0
\]
has not been established and cannot be inferred from the archived
one-link contraction bounds outside their stated parameter range.
Therefore \Cref{thm:s7v14-comparison} cannot presently yield a uniform
owner-fibre gap for this partition, even if a strict scalar influence
matrix were obtained.
\end{proposition}

\begin{proof}
The comparison theorem requires
\eqref{eq:s7v14-local-gap} for every allowed exterior fibre.  The
competing-centre-well result supplies allowed exterior data for which a
uniform one-link contraction was explicitly ruled out at weak
coupling.  The archived positive estimates apply only on their declared
parameter ranges and hence do not prove the displayed infimum.  Since
the product in \eqref{eq:s7v14-comparison} contains \(g_{\rm loc}\), an
unproved or vanishing first factor cannot be repaired by the second.
\end{proof}

\begin{decision}[Required block enlargement]
\label{dec:s7v14-enlarge}
The block must contain the modes responsible for the competing
one-link wells.  The sixty-chord V100 collar is one valid maximal
candidate because its fixed tree frames and sixty selected curvature
rows give a unique flat completion on the good core.  Smaller
overlapping blocks may be preferable, but their uniqueness and
conditional gaps must be checked rather than inferred link by link.
\end{decision}

\begin{definition}[Physical block-dynamics card, PBDC]
\label{def:s7v14-PBDC}
PBDC consists of:
\begin{enumerate}[label=\textup{(D\arabic*)},leftmargin=4em]
\item an explicit overlapping block cover of the sixty collar chords;
\item a worst-conditional internal gap \(g_i\) for each enlarged block;
\item proof that the common invariant subspace of all \(K_i\) is exactly
      the constants in the physical sector;
\item an orthogonal martingale resolution and the exact Gram blocks
      \(\mathbb G_{jk}\);
\item either a direct lower spectral bound for \(\mathbb G\) or the
      strict diagonal-dominance data
      \eqref{eq:s7v14-diagonal}--\eqref{eq:s7v14-offdiagonal};
\item derivation of every off-diagonal estimate from the physical
      score \(\sigma_{ij}\), including ground and regulator motion;
\item insertion of \(g_{\rm loc}\lambda_{\rm hb}\) into the V11
      moving-ground bound and the V5 martingale sums; and
\item good/bad, gauge, reflection, continuum, and spectral contacts.
\end{enumerate}
\end{definition}

\begin{theorem}[What PBDC would provide]
\label{thm:s7v14-PBDC}
Rows (D1)--(D6) imply a positive owner-fibre gap
\[
 g_{\rm owner}\ge g_{\rm loc}\lambda_{\rm hb}
\]
through \Cref{thm:s7v14-comparison}.  Rows (D7)--(D8), if proved with
cutoff-uniform strict constants, insert that gap into MGCC.  No
infinite-volume spectral conclusion follows from the finite-fibre
comparison alone.
\end{theorem}

\begin{proof}
Use \Cref{thm:s7v14-Gram,thm:s7v14-Gershgorin} to bound
\(\lambda_{\rm hb}\), then
\Cref{thm:s7v14-comparison}.  The remaining interfaces are precisely
those listed in MGCC.
\end{proof}

\begin{assessment}[Progress at seventh-series V14]
\label{ass:s7v14-status}
The interface problem has been converted into an exact positive
operator matrix.  Its diagonal is block visibility and its
off-diagonal part is the true physical influence.  The conditional
score formula identifies the quantities from which scalar influence
bounds must be proved.  The individual-link partition does not pass the
already established worst-boundary test, so the next construction must
use enlarged overlapping blocks.
\end{assessment}

\begin{programme}[Seventh-series V15]
\label{prog:s7v14-V15}
V15 is the compulsory companion audit.  After that audit it will build
the smallest explicit overlapping block family justified by the V100
collar incidence data, compute its coverage multiplicities and common
invariant sigma field, and state the first exact Gram diagonal that
remains to be bounded.
\end{programme}

\begin{firewall}[Claim boundary after seventh-series V14]
\label{fire:s7v14-boundary}
V14 proves the block-dynamics comparison, the exact influence Gram
representation, a scalar diagonal-dominance certificate, and the
physical conditional-score commutator.  It also shows why the
individual-link partition cannot use the archived estimates as a
uniform proof.

V14 does not construct a successful enlarged block cover, prove its
conditional gaps or Gram strictness, establish the active-score angle,
close the martingale budgets, construct the continuum, or prove the
Yang--Mills mass gap.
\end{firewall}

\paragraph{Parent-version record.}
V14 was formed from
\texttt{Renewal\_Yang\_Mills\_Mass\_Gap\_Research\_Notes\_V13.tex}.
The V13 parent had \(5\,256\) logical source lines, \(193\,564\) bytes,
and SHA--256
\[
 \texttt{4FAE700A2316CADE772844637EA4A006134668DCE78F46E7A628C08009078BB9}.
\]

% =============================================================================
% END APPENDED SEVENTH-SERIES V14 MATERIAL
% =============================================================================

% =============================================================================
% BEGIN APPENDED SEVENTH-SERIES V15 MATERIAL
% =============================================================================

\chapter{Companion audit and a certified overlapping collar cover}
\label{ch:s7v15}

This is the compulsory five-version companion audit.  The audit agrees
with V14 that interface mixing must be proved in the actual regulated
measure and adds no shortcut.  The constructive part of this chapter
returns to the V100 nonabelian certificate.  Its elimination order is
stronger than mere triviality of the collar presentation: the original
pivot relators are triangular.  This produces five explicit
twenty-chord blocks, each with an invertible selected curvature
Jacobian at the flat point, while their union sees every collar chord.

\section{Compulsory V15 companion audit}
\label{sec:s7v15-audit}

\begin{audit}[Files inspected]
\label{aud:s7v15-files}
The current companion files are:
\begin{enumerate}[label=\textup{(\alph*)},leftmargin=3em]
\item
\texttt{Renewal\_SM\_Einstein\_Continuum\_Research\_Notes\_V55.tex},
with \(22\,440\) logical lines, \(804\,549\) bytes, and SHA--256
\[
 \texttt{93D56FE924930BE60C9DE08A8320D3D86B57DBE29D7894AC6E4668B29924D93A};
\]
\item
\texttt{Renewal\_NS\_Stability\_Research\_Notes\_V26.tex},
with \(5\,319\) logical lines, \(278\,283\) bytes, and SHA--256
\[
 \texttt{82C887F8C2C69E4F28FAD7C3DC8DE5B9099CB5A4BF431EBA1FFF3E91935978AD}.
\]
\end{enumerate}
The SM programme advanced from V49 to V55.  The NS file is unchanged
since the V10 audit.
\end{audit}

\begin{assessment}[SM V50--V55]
\label{ass:s7v15-SM}
SM V50 extends the polar cancellation to an operator with a kernel by
working on \((\ker B)^\perp\); exact adjoint sources again enter with
coefficient one, while non-adjoint residuals pay the least nonzero
singular value.  V51--V54 derive time-window, reference-current,
diffeomorphism, and alignment-defect formulae.  Their transferable
lesson is the same source ledger used here: factor an actual adjoint or
divergence before inversion and retain equation, regulator, charge, and
boundary defects separately.

SM V55 audits YM V9--V13 and independently proves that good-chart
rarity plus local coercivity cannot create a global carrier gap without
interface mixing.  It selects an exact finite-regulator Ward score and
an active-score projection rather than a full-space norm.  This
supports the V14 physical Gram route.  Its gravitational currents,
ADM constraint, and diffeomorphism score do not provide numerical
Wilson block-influence constants.
\end{assessment}

\begin{assessment}[NS V26]
\label{ass:s7v15-NS}
NS V26 remains the exact rank-one Oseen-kernel calculation already
audited at V10.  It gives no new conditional block gap, no Wilson
influence estimate, and no reason to change the present direction.
\end{assessment}

\begin{decision}[Audit decision]
\label{dec:s7v15-audit}
Continue with the physical block Gram matrix.  The useful SM transfer is
the kernel-aware polar formulation: if an enlarged block retains a
gauge or harmonic kernel, project it explicitly and estimate only the
orthogonal physical sector.  Do not infer Gram strictness from score
centering, local good-core probability, or finite dimensionality.
\end{decision}

\section{Certified triangular data from V100}
\label{sec:s7v15-certificate}

The reproducibility script
\texttt{scripts/analyze\_v15\_collar\_blocks.py} has \(126\) logical
lines, \(4\,521\) bytes, and SHA--256
\[
 \texttt{787E44E797F419EAA301C33DD5E02AB4F37C8862D99353C297CBD015395F9CA4}.
\]
It imports the original V100 collar constructor, rebuilds all link and
plaquette sets, reruns the deterministic Tietze elimination, and then
checks the supports of the \emph{original} pivot relators.

Let \(q_1,\ldots,q_{60}\) denote chord generators in elimination order.
In terms of the zero-based V100 chord indices, that order is
\begin{align}
 &(3,0,1,2,4,5,6,7,8,36,9,10,37,39,11,12,13,41,14,40,
 \notag\\
 &\quad15,16,17,42,45,18,19,20,46,21,50,22,23,52,53,54,
 \notag\\
 &\quad26,27,49,29,30,25,57,31,44,32,48,33,34,28,59,35,
 \notag\\
 &\quad24,56,51,38,43,47,58,55).                        \label{eq:s7v15-order}
\end{align}

\begin{record}[Machine-checked triangular census]
\label{rec:s7v15-census}
For the pivot relator \(r_k\) used at elimination step \(k\):
\begin{enumerate}[label=\textup{(\roman*)},leftmargin=3.5em]
\item every generator in the original word \(r_k\) has elimination
      position at most \(k\);
\item \(q_k\) occurs exactly once in \(r_k\);
\item the support-size histogram is
      \(\{1:34,\ 2:12,\ 3:14\}\);
\item the backward-span histogram is
\[
 \{0:34,1:7,2:6,3:1,4:1,5:2,6:1,15:1,19:1,
   21:1,24:2,27:1,32:1,42:1\};
\]
\item the largest backward span is \(42\); and
\item each chord belongs to between one and six of the sixty original
      pivot-relator supports.
\end{enumerate}
The script terminates with
\texttt{triangular=True} and
\texttt{pivot\_multiplicity\_one=True}.
\end{record}

\begin{theorem}[Triangular flat curvature matrices]
\label{thm:s7v15-triangular}
Fix integers \(1\le a\le b\le60\).  Hold
\(q_1,\ldots,q_{a-1}\) fixed and consider the selected plaquette map
\begin{equation}
 {\cal F}_{[a,b]}:
 (q_a,\ldots,q_b)\longmapsto(r_a,\ldots,r_b).            \label{eq:s7v15-block-map}
\end{equation}
At the flat configuration, in left-trivialized Lie-algebra coordinates,
its derivative is
\[
 L_{[a,b]}\otimes I_{\mathfrak{su}(3)},
\]
where \(L_{[a,b]}\) is lower triangular with diagonal entries
\(\pm1\).  Hence the derivative is invertible.
\end{theorem}

\begin{proof}
By \Cref{rec:s7v15-census}(i), the original word \(r_k\) contains no
generator with elimination position greater than \(k\).  Therefore row
\(k\) of the derivative of \eqref{eq:s7v15-block-map} has no column
greater than \(k\).  Generators before \(a\) are fixed exterior data and
do not appear among the derivative columns.  By
\Cref{rec:s7v15-census}(ii), \(q_k\) occurs exactly once in \(r_k\).
At the flat configuration, differentiating a group word gives the
signed sum of its letter tangents, so the coefficient of the
\(q_k\) tangent is \(+I\) or \(-I\).  Thus the displayed matrix is
lower triangular with invertible diagonal.
\end{proof}

\begin{corollary}[Explicit crude singular floor for a twenty-chord window]
\label{cor:s7v15-floor}
For \(b-a+1\le20\), let
\[
 C_{20}:=\sum_{\ell=0}^{19}12^{\ell/2}.
\]
Then the flat selected derivative satisfies
\begin{equation}
 \sigma_{\min}(D{\cal F}_{[a,b]})\ge C_{20}^{-1}.         \label{eq:s7v15-floor}
\end{equation}
\end{corollary}

\begin{proof}
After multiplying rows by their diagonal signs, write
\(L_{[a,b]}=I+N\), with \(N\) strictly lower triangular and
\(N^{20}=0\).  Each pivot word has at most three distinct generators,
so every row of \(N\) has at most two nonzero entries of magnitude one.
Each chord occurs in at most six pivot supports, so every column has at
most six such entries.  Hence
\[
 \|N\|_2\le\sqrt{\|N\|_1\|N\|_\infty}\le\sqrt{12}.
\]
The finite Neumann identity gives
\[
 (I+N)^{-1}=\sum_{\ell=0}^{19}(-N)^\ell,
 \qquad
 \|(I+N)^{-1}\|_2\le C_{20}.
\]
Tensoring with the identity on \(\mathfrak{su}(3)\) does not change
singular values, proving \eqref{eq:s7v15-floor}.
\end{proof}

\begin{remark}[The floor is deliberately crude]
\label{rem:s7v15-floor}
The bound \(C_{20}^{-1}\) is an exact positive certificate, not a useful
numerical estimate.  Direct singular-value computation of the five
integer matrices can sharpen it.  More importantly, a local curvature
floor does not by itself imply a conditional Poincar\'e gap outside the
flat chart.
\end{remark}

\section{Five half-overlapping blocks}
\label{sec:s7v15-cover}

Define blocks by positions in \eqref{eq:s7v15-order}:
\begin{equation}
 \Lambda_1=[1,20],\quad
 \Lambda_2=[11,30],\quad
 \Lambda_3=[21,40],\quad
 \Lambda_4=[31,50],\quad
 \Lambda_5=[41,60].                                     \label{eq:s7v15-blocks}
\end{equation}
Here \([a,b]\) means the chord set
\(\{q_a,\ldots,q_b\}\), not a lattice interval.

\begin{proposition}[Coverage and selected local coordinates]
\label{prop:s7v15-cover}
The five blocks cover all sixty chords.  Twenty endpoint-position
chords belong to one block and forty interior-position chords belong to
two blocks.  For each block, its twenty pivot plaquettes define a local
analytic coordinate map at flatness with singular floor at least
\(C_{20}^{-1}\).
\end{proposition}

\begin{proof}
The interval census in \eqref{eq:s7v15-blocks} is immediate: positions
\(1\)--\(10\) and \(51\)--\(60\) occur once, and positions \(11\)--\(50\)
occur twice.  Apply
\Cref{thm:s7v15-triangular,cor:s7v15-floor} to each interval.  Analytic
local coordinates follow from the inverse function theorem.
\end{proof}

\begin{theorem}[Common invariant sigma field of the cover]
\label{thm:s7v15-common-kernel}
Assume the physical collar measure has a strictly positive density with
respect to product Haar measure on the sixty chord variables.  Let
\[
 K_\alpha
 =\mathbb E[\cdot\mid\mathcal F_{\Lambda_\alpha^c}],
 \qquad \alpha=1,\ldots,5.
\]
Then
\begin{equation}
 \bigcap_{\alpha=1}^5\operatorname{Ran}K_\alpha
 =\{\hbox{constants}\}.                                 \label{eq:s7v15-common-kernel}
\end{equation}
Thus row (D3) of PBDC holds for this finite-cutoff cover.
\end{theorem}

\begin{proof}
A function in \(\operatorname{Ran}K_\alpha\) is measurable, modulo
null sets, with respect to the coordinates outside
\(\Lambda_\alpha\).  Strict positivity makes the physical and
product-Haar null sets identical.  For product coordinate sigma fields,
the intersection over \(\alpha\) is the sigma field generated by
coordinates outside \(\bigcup_\alpha\Lambda_\alpha\); this follows
successively from Fubini's theorem by varying one covered coordinate at
a time.  The union is all sixty coordinates by
\Cref{prop:s7v15-cover}, so the intersection sigma field is trivial.
\end{proof}

\begin{corollary}[Fixed-parameter heat-bath kernel]
\label{cor:s7v15-A-kernel}
For the five-block heat-bath generator
\[
 {\cal A}_{20}:=\sum_{\alpha=1}^5(I-K_\alpha),
\]
one has
\[
 \ker{\cal A}_{20}=\{\hbox{constants}\}.
\]
\end{corollary}

\begin{proof}
Combine \Cref{lem:s7v14-form} with
\Cref{thm:s7v15-common-kernel}.
\end{proof}

\begin{warning}[Trivial kernel is not a uniform spectral floor]
\label{warn:s7v15-kernel}
The physical \(L^2\) space is infinite dimensional even at fixed
lattice cutoff.  A nonnegative sum of projections can have trivial
nonconstant kernel while its spectrum accumulates at zero.  Hence
\Cref{cor:s7v15-A-kernel} does not prove a positive, much less uniform,
value of \(\lambda_{\rm hb}\).  The exact Gram lower bound remains
necessary.
\end{warning}

\section{The first Gram diagonal for this cover}
\label{sec:s7v15-Gram}

Let
\[
 \mathcal F_k=\sigma(q_1,\ldots,q_k),\qquad
 \Delta_k=\mathbb E[\cdot\mid\mathcal F_k]
           -\mathbb E[\cdot\mid\mathcal F_{k-1}],
 \quad k=1,\ldots,60,
\]
with all far-exterior variables held fixed and with the corresponding
conditional physical measure.  These are orthogonal martingale
differences.

\begin{definition}[Certified-cover Gram data]
\label{def:s7v15-Gram-data}
For \(j,k\in\{1,\ldots,60\}\), put
\begin{align}
 \mathbb G^{(20)}_{jk}
 &:=
 \sum_{\alpha=1}^5
 \Delta_j(I-K_\alpha)\Delta_k,                           \label{eq:s7v15-Gjk}\\
 a_j^{(20)}
 &:=
 \inf_{\substack{h\in\operatorname{Ran}\Delta_j\\h\ne0}}
 {\sum_{\alpha=1}^5\|(I-K_\alpha)h\|_2^2\over\|h\|_2^2},
                                                               \label{eq:s7v15-aj}\\
 c_{jk}^{(20)}
 &:=\|\mathbb G^{(20)}_{jk}\|,
 \qquad j\ne k.                                         \label{eq:s7v15-cjk}
\end{align}
\end{definition}

\begin{theorem}[Exact strictness target for the certified cover]
\label{thm:s7v15-strictness}
If
\begin{equation}
 \delta_{20}
 :=
 \min_{1\le j\le60}
 \left(a_j^{(20)}-\sum_{k\ne j}c_{jk}^{(20)}\right)>0,   \label{eq:s7v15-delta}
\end{equation}
then
\begin{equation}
 \lambda_{\rm hb}^{(20)}\ge\delta_{20}.                  \label{eq:s7v15-hb-gap}
\end{equation}
If, in addition, every twenty-chord conditional diffusion has
worst-exterior gap at least \(g_{20}>0\), then the whole owner-fibre
diffusion has gap at least
\begin{equation}
 g_{\rm owner}\ge g_{20}\delta_{20}.                     \label{eq:s7v15-owner-gap}
\end{equation}
\end{theorem}

\begin{proof}
The definitions
\eqref{eq:s7v15-Gjk}--\eqref{eq:s7v15-cjk} are the diagonal and
off-diagonal hypotheses of
\Cref{thm:s7v14-Gershgorin}.  This proves
\eqref{eq:s7v15-hb-gap}.  Apply
\Cref{thm:s7v14-comparison} with \(g_{\rm loc}=g_{20}\) for
\eqref{eq:s7v15-owner-gap}.
\end{proof}

\begin{firewall}[What triangular curvature does not prove]
\label{fire:s7v15-curvature-gap}
The invertibility of \(D{\cal F}_{\Lambda_\alpha}\) proves local
coordinates and a local Wilson Hessian floor near flatness.  It does
not prove \(g_{20}>0\) uniformly over the full twenty-chord conditional
fibre, nor does it bound \(a_j^{(20)}\) or
\(c_{jk}^{(20)}\) in the physical ground measure.  Promoting the
integer-matrix floor directly to \eqref{eq:s7v15-owner-gap} would repeat
the V12 global-convexity error.
\end{firewall}

\begin{assessment}[Progress at seventh-series V15]
\label{ass:s7v15-status}
The mandatory audit supports the active-score/block-dynamics route.
The V100 certificate has now been mined for an explicit five-block
cover.  Every block has twenty chord variables, twenty triangular pivot
plaquettes, an explicit flat singular floor, and half-overlap with its
neighbours; together the blocks have only constants as common
invariants.  The remaining local quantitative data are exactly
\(g_{20}\) and the physical Gram margin \(\delta_{20}\).
\end{assessment}

\begin{decision}[V15 bottleneck decision]
\label{dec:s7v15-bottleneck}
The long backward spans in
\Cref{rec:s7v15-census} show that elimination order is not geometrically
local, so scalar Gershgorin may be wasteful.  V16 will compute the five
flat integer block matrices and the exact block-overlap incidence.
It will compare the scalar margin with the full operator Gram symbol in
the Gaussian flat surrogate.  If that surrogate already has zero
margin, the cover must be changed before any nonlinear score estimate
is attempted.
\end{decision}

\begin{programme}[Seventh-series V16]
\label{prog:s7v15-V16}
V16 will extend the reproducibility script to construct
\(L_{\Lambda_\alpha}\), compute exact determinants, numerical singular
values with certified rational lower bounds where feasible, and form
the Gaussian conditional-projection Gram matrix for the five-block
cover.  The calculation is a necessary linearized test, not a proof for
the nonlinear physical measure.
\end{programme}

\begin{firewall}[Claim boundary after seventh-series V15]
\label{fire:s7v15-boundary}
V15 performs the required SM/NS audit, certifies triangular support of
the sixty original pivot relators, constructs a five-block
half-overlapping cover, proves its flat local-coordinate property and
common invariant sigma field, and states its exact Gram strictness
target.

V15 does not prove the twenty-chord conditional gaps, the physical Gram
margin, the active-score angle, remaining martingale budgets, continuum
construction, or the Yang--Mills mass gap.
\end{firewall}

\paragraph{Parent-version record.}
V15 was formed from
\texttt{Renewal\_Yang\_Mills\_Mass\_Gap\_Research\_Notes\_V14.tex}.
The V14 parent had \(5\,670\) logical source lines, \(208\,728\) bytes,
and SHA--256
\[
 \texttt{F584BDB8047FF818954C1AA2962A53C37564EE5A99F027D6F95125EFB1B23C2C}.
\]

% =============================================================================
% END APPENDED SEVENTH-SERIES V15 MATERIAL
% =============================================================================

% =============================================================================
% BEGIN APPENDED SEVENTH-SERIES V16 MATERIAL
% =============================================================================

\chapter{Sharp flat collar linear algebra and the Gaussian block gap}
\label{ch:s7v16}

V15 used only a generic nilpotent-series estimate for its five block
matrices.  Exact integer inversion reveals that this was extremely
wasteful.  The complete selected flat curvature matrix has a
cutoff-independent singular floor \(1/\sqrt{28}\), and every
twenty-chord block has a substantially larger certified floor.  These
data also prove a positive spectral gap for the entire selected
correlated Gaussian block heat bath, on all Wiener chaoses rather than
only on linear observables.

\section{Reproducibility record}
\label{sec:s7v16-record}

The extended script
\texttt{scripts/analyze\_v15\_collar\_blocks.py} now has \(227\)
logical lines, \(8\,341\) bytes, and SHA--256
\[
 \texttt{2C8C7B8831BF376C885F691F458D138DBF3B84DA28274BD0A4335B3CF1E4673E}.
\]
It rebuilds the V100 collar, verifies the triangular support and pivot
multiplicity, constructs the exact integer matrices, inverts each by
integer forward substitution, checks the products against the identity,
and computes the overlap and Gaussian projection matrices.  Floating
singular values are recorded only as diagnostics; all lower bounds used
in the theorems below come from exact integer norms.

\section{The complete integer curvature matrix}
\label{sec:s7v16-full}

Let \(L\in M_{60}(\mathbb Z)\) be the scalar derivative at the identity
of the sixty pivot relators in the order
\eqref{eq:s7v15-order}, with columns in the same chord order.  Thus
\[
 D{\cal F}_\star=L\otimes I_{\mathfrak{su}(3)}.
\]
V15 proves that \(L\) is lower triangular with diagonal entries
\(\pm1\).

\begin{record}[Exact inverse norms]
\label{rec:s7v16-full}
Integer forward substitution and exact multiplication give
\begin{equation}
 \|L^{-1}\|_1=7,\qquad
 \|L^{-1}\|_\infty=4.                                   \label{eq:s7v16-inverse-norms}
\end{equation}
\end{record}

\begin{theorem}[Sharpened full-collar singular floor]
\label{thm:s7v16-full-floor}
The selected flat curvature derivative satisfies
\begin{equation}
 \boxed{
 \sigma_{\min}(D{\cal F}_\star)
 =\sigma_{\min}(L)
 \ge {1\over\sqrt{28}}.}                                \label{eq:s7v16-full-floor}
\end{equation}
Equivalently,
\begin{equation}
 (D{\cal F}_\star)^*D{\cal F}_\star
 \succeq {1\over28}I.                                   \label{eq:s7v16-full-Hessian}
\end{equation}
\end{theorem}

\begin{proof}
For every matrix,
\[
 \|L^{-1}\|_2
 \le\sqrt{\|L^{-1}\|_1\|L^{-1}\|_\infty}.
\]
Insert \eqref{eq:s7v16-inverse-norms} to obtain
\(\|L^{-1}\|_2\le\sqrt{28}\).  Since
\(\sigma_{\min}(L)=\|L^{-1}\|_2^{-1}\), the scalar bound follows.
Tensoring with the identity on \(\mathfrak{su}(3)\) repeats singular
values and proves both displayed conclusions.
\end{proof}

\begin{corollary}[Improved nonlinear collar core]
\label{cor:s7v16-nonlinear}
Let \(B(U)\) be the nonlinear selected curvature derivative in the same
tree-gauge coordinates.  If
\begin{equation}
 \|B(U)-D{\cal F}_\star\|\le {1\over2\sqrt{28}},          \label{eq:s7v16-perturbation}
\end{equation}
then
\begin{equation}
 \sigma_{\min}(B(U))\ge {1\over2\sqrt{28}}.              \label{eq:s7v16-nonlinear-floor}
\end{equation}
\end{corollary}

\begin{proof}
The least singular value is \(1\)-Lipschitz in operator norm:
\[
 \sigma_{\min}(B(U))
 \ge\sigma_{\min}(D{\cal F}_\star)
   -\|B(U)-D{\cal F}_\star\|.
\]
Use \eqref{eq:s7v16-full-floor,eq:s7v16-perturbation}.
\end{proof}

\begin{assessment}[Relation to the V100 floor]
\label{ass:s7v16-improvement}
The earlier \(102^{-59/2}\) singular floor came from a generic
step-by-step elimination estimate.  Equation
\eqref{eq:s7v16-full-floor} uses the exact certified matrix and
supersedes that crude flat bound for this selected row set.  It does not
enlarge the nonlinear good core until
\eqref{eq:s7v16-perturbation} is converted into an explicit plaquette
defect radius.
\end{assessment}

\section{Exact twenty-chord block floors}
\label{sec:s7v16-blocks}

Let \(L_\alpha\) be the principal row/column matrix for the block
\(\Lambda_\alpha\) of \eqref{eq:s7v15-blocks}.

\begin{record}[Exact block data]
\label{rec:s7v16-blocks}
The verifier gives:
\begin{center}
\begin{tabular}{c|c|c|c|c|c}
\(\alpha\)&positions&\(\det L_\alpha\)&
\(\|L_\alpha^{-1}\|_1\)&
\(\|L_\alpha^{-1}\|_\infty\)&certified floor\\ \hline
1&\(1\)--\(20\)&\(1\)&6&3&\(1/\sqrt{18}\)\\
2&\(11\)--\(30\)&\(-1\)&2&3&\(1/\sqrt6\)\\
3&\(21\)--\(40\)&\(-1\)&2&3&\(1/\sqrt6\)\\
4&\(31\)--\(50\)&\(-1\)&3&3&\(1/3\)\\
5&\(41\)--\(60\)&\(-1\)&6&3&\(1/\sqrt{18}\)
\end{tabular}
\end{center}
The corresponding floating diagnostics for the least singular values
are \(0.381966011250\), \(0.517638090205\),
\(0.517638090205\), \(0.445041867913\), and
\(0.381966011250\).
\end{record}

\begin{theorem}[Certified block curvature floors]
\label{thm:s7v16-block-floors}
For each block,
\begin{equation}
 \sigma_{\min}(L_\alpha)
 \ge
 \bigl(
 \|L_\alpha^{-1}\|_1
 \|L_\alpha^{-1}\|_\infty
 \bigr)^{-1/2},                                         \label{eq:s7v16-block-floor}
\end{equation}
giving exactly the five certified floors in
\Cref{rec:s7v16-blocks}.
\end{theorem}

\begin{proof}
Apply
\(\|A\|_2\le\sqrt{\|A\|_1\|A\|_\infty}\)
to \(A=L_\alpha^{-1}\), then invert.  All displayed matrix norms were
computed by exact integer arithmetic and verified by
\(L_\alpha L_\alpha^{-1}=I\).
\end{proof}

The exact overlap matrix is
\begin{equation}
 \bigl(|\Lambda_\alpha\cap\Lambda_\beta|\bigr)_{\alpha,\beta=1}^5
 =
 \begin{pmatrix}
 20&10&0&0&0\\
 10&20&10&0&0\\
 0&10&20&10&0\\
 0&0&10&20&10\\
 0&0&0&10&20
 \end{pmatrix}.                                         \label{eq:s7v16-overlap}
\end{equation}

\section{The selected correlated Gaussian surrogate}
\label{sec:s7v16-Gaussian}

It suffices to treat one Lie-algebra component.  On
\(\mathbb R^{60}\), define
\begin{equation}
 d\gamma_L(x)
 =Z^{-1}\exp\!\left(-{1\over2}\|Lx\|^2\right)dx.         \label{eq:s7v16-gaussian}
\end{equation}
The change of variables \(y=Lx\) sends \(\gamma_L\) to standard
Gaussian measure.  For a block \(\Lambda_\alpha\), let \(K_\alpha^G\)
be conditional expectation given \(x_{\Lambda_\alpha^c}\) and
\(D_\alpha^G=I-K_\alpha^G\).

\begin{lemma}[First-chaos block subspace]
\label{lem:s7v16-first-chaos}
In standard \(y\)-coordinates, the restriction of
\(D_\alpha^G\) to the first Wiener chaos is the orthogonal projection
\(P_\alpha\) onto
\begin{equation}
 U_\alpha
 :=\operatorname{Ran}L_{:\!,\Lambda_\alpha}.             \label{eq:s7v16-U}
\end{equation}
\end{lemma}

\begin{proof}
Since \(x=L^{-1}y\), the exterior coordinate linear forms span
\[
 V_\alpha
 =\operatorname{span}\{L^{-T}e_j:j\notin\Lambda_\alpha\}.
\]
Conditional expectation of a standard-Gaussian linear observable is
orthogonal projection onto its measurable linear subspace
\(V_\alpha\).  A vector \(z\) is orthogonal to \(V_\alpha\) exactly when
\[
 (L^{-1}z)_j=0\quad(j\notin\Lambda_\alpha),
\]
which is equivalent to
\(z\in\operatorname{Ran}L_{:\!,\Lambda_\alpha}=U_\alpha\).
Thus \(D_\alpha^G\) is projection onto \(V_\alpha^\perp=U_\alpha\).
\end{proof}

Put
\begin{equation}
 {\mathbb G}_1^G:=\sum_{\alpha=1}^5P_\alpha.             \label{eq:s7v16-G1}
\end{equation}

\begin{theorem}[Rigorous first-chaos frame bound]
\label{thm:s7v16-frame}
The Gaussian first-chaos Gram operator satisfies
\begin{equation}
 {\mathbb G}_1^G\succeq {1\over952}I.                   \label{eq:s7v16-frame}
\end{equation}
\end{theorem}

\begin{proof}
Write \(A_\alpha=L_{:\!,\Lambda_\alpha}\).  Since \(A_\alpha\) has full
column rank,
\[
 P_\alpha
 =A_\alpha(A_\alpha^TA_\alpha)^{-1}A_\alpha^T
 \succeq {1\over\|A_\alpha\|_2^2}A_\alpha A_\alpha^T.
\]
The exact squared Frobenius norms of the five \(A_\alpha\) are
\(34,33,33,33,33\), so
\[
 P_\alpha\succeq {1\over34}A_\alpha A_\alpha^T.
\]
If \(D_{\rm cov}\) is the diagonal matrix of block coverage
multiplicities, then
\[
 \sum_\alpha A_\alpha A_\alpha^T
 =LD_{\rm cov}L^T\succeq LL^T
 \succeq {1\over28}I
\]
by \eqref{eq:s7v16-full-Hessian}.  Therefore
\[
 {\mathbb G}_1^G
 \succeq {1\over34\cdot28}I={1\over952}I.
\]
\end{proof}

\begin{theorem}[All-chaos Gaussian block gap]
\label{thm:s7v16-all-chaos}
The complete Gaussian heat-bath generator
\begin{equation}
 {\cal A}_{20}^G
 :=\sum_{\alpha=1}^5(I-K_\alpha^G)                       \label{eq:s7v16-AG}
\end{equation}
has spectral gap
\begin{equation}
 \operatorname{gap}({\cal A}_{20}^G)
 =\lambda_{\min}({\mathbb G}_1^G)
 \ge {1\over952}.                                       \label{eq:s7v16-gaussian-gap}
\end{equation}
\end{theorem}

\begin{proof}
Under the standard-Gaussian identification, \(K_\alpha^G\) is the
second quantization of the first-chaos exterior projection
\(Q_\alpha=I-P_\alpha\).  On the \(n\)-th Wiener chaos it acts as the
restriction of \(Q_\alpha^{\otimes n}\) to symmetric tensors.

For any orthogonal projection \(Q\),
\begin{equation}
 I-Q^{\otimes n}
 \succeq {1\over n}\sum_{r=1}^n
 I^{\otimes(r-1)}\otimes(I-Q)\otimes I^{\otimes(n-r)}.   \label{eq:s7v16-tensor-projection}
\end{equation}
Indeed, the commuting projections on the right can be diagonalized
simultaneously; on a joint eigenvector the left side is one when at
least one factor eigenvalue is zero, whereas the right side is the
fraction of zero factors.

Sum \eqref{eq:s7v16-tensor-projection} over \(\alpha\).
On every tensor position,
\[
 \sum_\alpha(I-Q_\alpha)=\sum_\alpha P_\alpha
 ={\mathbb G}_1^G
 \succeq\lambda_{\min}({\mathbb G}_1^G)I.
\]
Hence every nonzero chaos has gap at least the first-chaos minimum.
The first chaos is itself an invariant subspace and attains that
minimum, proving equality.  Apply
\Cref{thm:s7v16-frame}.
\end{proof}

\begin{record}[Numerical diagnostic]
\label{rec:s7v16-numerical}
Double-precision orthogonal projection gives
\[
 \lambda_{\min}({\mathbb G}_1^G)
 \approx0.381966011250,\qquad
 \lambda_{\max}({\mathbb G}_1^G)
 \approx3.41421356237.
\]
These values are not used to certify
\eqref{eq:s7v16-gaussian-gap}; the rigorous lower bound is \(1/952\).
\end{record}

\begin{assessment}[What the Gaussian test decides]
\label{ass:s7v16-Gaussian}
The five-block cover passes the complete selected Gaussian test with a
strict, explicit margin.  Thus its overlap geometry and the
long-range triangular dependencies do not create a linearized zero
mode.  Any failure of the nonlinear physical Gram margin must come from
measure deformation, additional action/ground/regulator terms, a bad
boundary fibre, or leaving the curvature core.
\end{assessment}

\begin{definition}[Gaussian-to-physical block transport card, GPBTC]
\label{def:s7v16-GPBTC}
GPBTC consists of:
\begin{enumerate}[label=\textup{(T\arabic*)},leftmargin=4em]
\item the exact flat frame bound \eqref{eq:s7v16-frame};
\item a common probability space or unitary identification comparing
      physical and selected Gaussian conditional projections;
\item a bound
      \(\|K_\alpha-K_\alpha^G\|\le\eta_\alpha\) on the active
      martingale subspace;
\item a perturbation estimate preserving a positive fraction of the
      \(1/952\) Gram margin;
\item direct inclusion of the moving ground, complete incident Wilson
      rows, regulator, and allowed boundary;
\item uniform twenty-chord internal gaps;
\item insertion into PBDC, MGCC, MLC, and the strict shell budget; and
\item good/bad, gauge, reflection, continuum, and spectral contacts.
\end{enumerate}
\end{definition}

\begin{decision}[V16 bottleneck decision]
\label{dec:s7v16-bottleneck}
The flat geometry is now quantitatively adequate.  The next target is
not a different cover but stability of conditional projections under
measure deformation.  V17 will prove an operator perturbation theorem
for conditional expectations whose fibre density ratios are close in
\(L^\infty\), derive the resulting Gram-gap loss, and calculate the
maximum admissible deformation relative to the certified \(1/952\)
margin.  The estimate will then be tested against the Wilson
good-core oscillation.
\end{decision}

\begin{firewall}[Claim boundary after seventh-series V16]
\label{fire:s7v16-boundary}
V16 proves the sharpened \(1/\sqrt{28}\) full flat curvature floor, five
exact block floors, the first-chaos frame bound, and the all-chaos
selected Gaussian block gap \(1/952\).

V16 does not prove stability of that gap for the nonlinear physical
measure, uniform block conditional gaps, the active-score angle,
remaining martingale budgets, continuum construction, or the
Yang--Mills mass gap.
\end{firewall}

\paragraph{Parent-version record.}
V16 was formed from
\texttt{Renewal\_Yang\_Mills\_Mass\_Gap\_Research\_Notes\_V15.tex}.
The V15 parent had \(6\,086\) logical source lines, \(224\,600\) bytes,
and SHA--256
\[
 \texttt{16AB6BE306C5FB85750C57A6E551F226E993A40C1E7D1BFB15CEA568D74B77F4}.
\]

% =============================================================================
% END APPENDED SEVENTH-SERIES V16 MATERIAL
% =============================================================================

% =============================================================================
% BEGIN APPENDED SEVENTH-SERIES V17 MATERIAL
% =============================================================================

\chapter{Stability of conditional projections under density deformation}
\label{ch:s7v17}

V16 proves a strict block heat-bath gap for the selected correlated
Gaussian measure.  This chapter gives an exact perturbation theorem for
conditional expectations on a common probability space.  It quantifies
two different losses: conditional density oscillation moves each
block projection, while global density oscillation compares the two
variance norms.  The resulting theorem is rigorous but exposes a new
obstruction: a local logarithm chart can make the first loss small, yet
does not control the second over the full compact fibre.

\section{Conditional density ratios}
\label{sec:s7v17-ratios}

Let \(\mu_0\) be a probability measure and let
\begin{equation}
 d\mu=w\,d\mu_0,\qquad w>0,\qquad\int w\,d\mu_0=1.        \label{eq:s7v17-w}
\end{equation}
For the five blocks of V15, write
\[
 K_\alpha^0
 =\mathbb E_{\mu_0}[\cdot\mid\mathcal F_{\Lambda_\alpha^c}],
 \qquad
 K_\alpha
 =\mathbb E_{\mu}[\cdot\mid\mathcal F_{\Lambda_\alpha^c}].
\]
The conditional Radon--Nikodym density is
\begin{equation}
 r_\alpha
 :={w\over K_\alpha^0w},\qquad
 K_\alpha f=K_\alpha^0(r_\alpha f).                     \label{eq:s7v17-r}
\end{equation}

\begin{lemma}[Conditional oscillation controls the normalized ratio]
\label{lem:s7v17-ratio}
Suppose that on almost every exterior fibre
\begin{equation}
 \operatorname*{ess\,osc}_{\Lambda_\alpha}\log w
 \le\omega_\alpha.                                      \label{eq:s7v17-omega}
\end{equation}
Then
\begin{equation}
 e^{-\omega_\alpha}\le r_\alpha\le e^{\omega_\alpha},
 \qquad
 |r_\alpha-1|\le\delta_\alpha,
 \quad
 \delta_\alpha:=e^{\omega_\alpha}-1.                    \label{eq:s7v17-delta}
\end{equation}
\end{lemma}

\begin{proof}
On a fixed exterior fibre let \(w_-\) and \(w_+\) be the essential
infimum and supremum.  The conditional average satisfies
\[
 w_-\le K_\alpha^0w\le w_+,
\]
while \(w_+/w_-\le e^{\omega_\alpha}\).  Divide \(w\) by its conditional
average to obtain the two-sided bound.  Since
\(e^{\omega_\alpha}-1\ge1-e^{-\omega_\alpha}\), the final estimate
follows.
\end{proof}

\begin{theorem}[Conditional-projection perturbation]
\label{thm:s7v17-projection}
Under \eqref{eq:s7v17-omega},
\begin{equation}
 \|K_\alpha-K_\alpha^0\|_{L^2(\mu_0)\to L^2(\mu_0)}
 \le\delta_\alpha.                                      \label{eq:s7v17-projection}
\end{equation}
The same bound holds for
\((I-K_\alpha)-(I-K_\alpha^0)\).
\end{theorem}

\begin{proof}
Using \eqref{eq:s7v17-r} and conditional Cauchy--Schwarz,
\begin{align*}
 |(K_\alpha-K_\alpha^0)f|^2
 &=|K_\alpha^0((r_\alpha-1)f)|^2\\
 &\le K_\alpha^0((r_\alpha-1)^2)
       K_\alpha^0(f^2)
 \le\delta_\alpha^2K_\alpha^0(f^2).
\end{align*}
Integrate with respect to \(\mu_0\) and use invariance of conditional
expectation.  Subtracting each projection from the identity gives the
last assertion.
\end{proof}

\section{A quantitative heat-bath gap perturbation}
\label{sec:s7v17-gap}

Let
\[
 D_\alpha^0=I-K_\alpha^0,\qquad
 D_\alpha=I-K_\alpha,\qquad
 \eta:=\left(\sum_{\alpha=1}^5\delta_\alpha^2\right)^{1/2}.
\]
Assume that the reference heat bath has gap
\begin{equation}
 \sum_{\alpha=1}^5\|D_\alpha^0f\|_{L^2(\mu_0)}^2
 \ge\lambda_0\operatorname{Var}_{\mu_0}(f).              \label{eq:s7v17-reference-gap}
\end{equation}

\begin{theorem}[Five-block density-perturbation compiler]
\label{thm:s7v17-gap}
Suppose
\begin{equation}
 0<m\le w\le M<\infty                                    \label{eq:s7v17-global-ratio}
\end{equation}
and \(\eta<\sqrt{\lambda_0}\).  Then the physical conditional
projections satisfy
\begin{equation}
 \boxed{
 \sum_{\alpha=1}^5\|D_\alpha f\|_{L^2(\mu)}^2
 \ge
 {m\over M}
 \bigl(\sqrt{\lambda_0}-\eta\bigr)^2
 \operatorname{Var}_{\mu}(f).}                          \label{eq:s7v17-gap}
\end{equation}
\end{theorem}

\begin{proof}
Replace \(f\) by
\(g=f-\int f\,d\mu_0\); all \(D_\alpha\) and \(D_\alpha^0\) annihilate
constants.  In the Hilbert direct sum of five copies of
\(L^2(\mu_0)\), the reverse triangle inequality and
\Cref{thm:s7v17-projection} give
\begin{align*}
 \left(\sum_\alpha\|D_\alpha g\|_0^2\right)^{1/2}
 &\ge
 \left(\sum_\alpha\|D_\alpha^0g\|_0^2\right)^{1/2}
 -
 \left(\sum_\alpha
 \|(D_\alpha-D_\alpha^0)g\|_0^2\right)^{1/2}\\
 &\ge
 \bigl(\sqrt{\lambda_0}-\eta\bigr)\|g\|_0.
\end{align*}
Here \(\|g\|_0^2=\operatorname{Var}_{\mu_0}(f)\).
The lower density bound yields
\[
 \sum_\alpha\|D_\alpha f\|_\mu^2
 \ge m\sum_\alpha\|D_\alpha g\|_0^2.
\]
The upper density bound and the variational definition of variance give
\[
 \operatorname{Var}_\mu(f)
 \le M\operatorname{Var}_{\mu_0}(f).
\]
Combine the three inequalities.
\end{proof}

\begin{corollary}[Oscillation form]
\label{cor:s7v17-oscillation}
If
\[
 \operatorname*{ess\,osc}_\Omega\log w\le\Omega,
\]
then \(m/M\ge e^{-\Omega}\), and
\begin{equation}
 \lambda_{\rm hb}(\mu)
 \ge
 e^{-\Omega}
 \left(
 \sqrt{\lambda_0}
 -\left[\sum_{\alpha=1}^5
 (e^{\omega_\alpha}-1)^2\right]^{1/2}
 \right)^2                                             \label{eq:s7v17-oscillation-gap}
\end{equation}
whenever the parenthesis is positive.
\end{corollary}

\begin{proof}
The ratio \(M/m\) is at most
\(\exp(\operatorname{ess\,osc}\log w)\).  Apply
\Cref{thm:s7v17-gap}.
\end{proof}

\section{Numerical window supplied by the certified margin}
\label{sec:s7v17-window}

\begin{corollary}[Uniform per-block oscillation threshold]
\label{cor:s7v17-window}
Use the rigorous V16 value \(\lambda_0=1/952\).  If all five conditional
oscillations are at most \(\omega\), then the strictness condition is
\begin{equation}
 \sqrt5(e^\omega-1)<{1\over\sqrt{952}},
 \quad\hbox{or equivalently}\quad
 \omega<
 \log\!\left(1+{1\over\sqrt{4760}}\right)
 \approx0.0143902379724.                                \label{eq:s7v17-window}
\end{equation}
If the left side is at most \(1/(2\sqrt{952})\), then
\begin{equation}
 \lambda_{\rm hb}(\mu)\ge {m\over3808M}.                 \label{eq:s7v17-half-margin}
\end{equation}
\end{corollary}

\begin{proof}
Here
\(\eta\le\sqrt5(e^\omega-1)\).  Solve
\(\eta<1/\sqrt{952}\) for \(\omega\).  Under the half-margin
hypothesis,
\[
 (\sqrt{\lambda_0}-\eta)^2
 \ge {1\over4}\lambda_0={1\over3808}.
\]
Use \eqref{eq:s7v17-gap}.
\end{proof}

\begin{remark}[The exact numerical spectrum is much better]
\label{rem:s7v17-numerical}
The V16 floating diagnostic gives
\(\lambda_{\min}\approx0.381966011250\), whereas \(1/952\) is about
\(0.00105042\).  The rigorous perturbation window can therefore be
greatly enlarged if that algebraic-looking spectrum is certified by
exact rational characteristic-polynomial or interval arithmetic.
V17 does not use the floating value as proof.
\end{remark}

\section{Testing the Wilson good core}
\label{sec:s7v17-Wilson}

On a common logarithm chart, let the selected Gaussian reference be the
quadratic part of the physical fibre density and write
\begin{equation}
 \log w=R_{\rm W}+R_0+R_{\rm reg},                       \label{eq:s7v17-remainders}
\end{equation}
where the terms are respectively the nonlinear Wilson remainder, the
moving-ground contribution, and all remaining regulator/reference
contributions.

\begin{lemma}[Core Taylor criterion]
\label{lem:s7v17-Taylor}
Suppose on a core of radius \(\rho\),
\begin{equation}
 |R_{\rm W}|\le C_{\rm W}s\rho^3,\qquad
 |R_0|+|R_{\rm reg}|\le A_{\rm gr}.                      \label{eq:s7v17-Taylor-assumption}
\end{equation}
Then on every conditional block fibre restricted to that core,
\begin{equation}
 \omega_\alpha
 \le2(C_{\rm W}s\rho^3+A_{\rm gr}).                      \label{eq:s7v17-Taylor}
\end{equation}
\end{lemma}

\begin{proof}
The oscillation of a function on a set is at most twice its supremum
absolute value.  Apply this to the sum
\eqref{eq:s7v17-remainders}.
\end{proof}

\begin{corollary}[Strong-coupling local scale]
\label{cor:s7v17-scale}
Ignoring \(A_{\rm gr}\) only for the purpose of reading the Wilson
scaling, the certified perturbation window permits
\[
 \rho=O(s^{-1/3}).
\]
This is larger than the \(O(s^{-1/2})\) fluctuation scale of a
nondegenerate strong-coupling Gaussian well.
\end{corollary}

\begin{proof}
Requiring \(2C_{\rm W}s\rho^3\) to stay below a fixed positive constant
gives \(\rho=O(s^{-1/3})\).  The Gaussian scale follows by rescaling a
quadratic density \(e^{-cs|x|^2}\).
\end{proof}

\begin{warning}[The local chart does not supply the global ratio]
\label{warn:s7v17-global}
\Cref{lem:s7v17-Taylor} controls conditional oscillation only after
restriction to the logarithm core.  The V16 Gaussian measure lives on
the tangent space, whereas the physical chord measure lives on the
compact group.  There is no global Radon--Nikodym ratio
\eqref{eq:s7v17-w} between those two measures until a common truncated
chart, wrapped reference, or other explicit identification is chosen.
Even after such a choice, the full-fibre oscillation \(\Omega\) can be
of order \(s\), making \(m/M\) exponentially small.  Local Taylor
control cannot be inserted into \eqref{eq:s7v17-gap} as if it were a
global comparison.
\end{warning}

\begin{firewall}[Why rarity still does not repair the theorem]
\label{fire:s7v17-rarity}
One cannot apply \Cref{thm:s7v17-gap} on the good core and then discard
its complement using only a small probability.  V13 proves that this
loses interface conductance.  A localized perturbation theorem must
control the spectral flow or capacity across the chart boundary, not
only the conditional density ratio inside it.
\end{firewall}

\begin{definition}[Conditional density transport card, CDTC]
\label{def:s7v17-CDTC}
CDTC consists of:
\begin{enumerate}[label=\textup{(P\arabic*)},leftmargin=4em]
\item a common global reference space for physical and comparison
      conditional projections;
\item conditional oscillations \(\omega_\alpha\) satisfying the strict
      margin in \eqref{eq:s7v17-oscillation-gap};
\item a nondegenerate global norm comparison, or a replacement avoiding
      the factor \(m/M\);
\item exact inclusion of the moving ground and regulator in
      \(\log w\);
\item a localized spectral-flow/capacity estimate if only core control
      is available;
\item uniform twenty-chord internal gaps;
\item insertion into PBDC, MGCC, MLC, and the shell budget; and
\item gauge, reflection, continuum, and spectral contacts.
\end{enumerate}
\end{definition}

\begin{assessment}[Progress at seventh-series V17]
\label{ass:s7v17-status}
The Gaussian-to-physical projection perturbation now has an exact
operator theorem and a numerical admissible window based solely on the
certified \(1/952\) margin.  The conditional part is compatible with a
shrinking \(s^{-1/3}\) Wilson core.  The global norm comparison is not:
the tangent Gaussian and compact-group measures are not yet on a common
global chart, and a crude full-fibre density ratio can decay
exponentially.  The remaining issue is therefore spectral transport
across the core boundary, not flat block geometry.
\end{assessment}

\begin{decision}[V17 bottleneck decision]
\label{dec:s7v17-bottleneck}
V18 will place a smooth interpolation \(\mu_t\) on the compact group
itself and differentiate its conditional projections after a unitary
transport to one fixed \(L^2\) space.  The goal is a spectral-flow
bound driven by centered conditional scores, which can exploit the
active-source estimates of V9--V11 and may avoid the global
essential-infimum \(m\).  All score norms and gap denominators will
remain explicit.
\end{decision}

\begin{firewall}[Claim boundary after seventh-series V17]
\label{fire:s7v17-boundary}
V17 proves conditional-projection stability under fibre density
oscillation and the resulting five-block heat-bath gap perturbation.
It computes the certified equal-block window and the local Wilson
Taylor scale.

V17 does not provide a legal global Gaussian/group identification,
avoid the global density-ratio loss, prove nonlinear physical Gram
strictness, uniform block internal gaps, remaining martingale budgets,
continuum construction, or the Yang--Mills mass gap.
\end{firewall}

\paragraph{Parent-version record.}
V17 was formed from
\texttt{Renewal\_Yang\_Mills\_Mass\_Gap\_Research\_Notes\_V16.tex}.
The V16 parent had \(6\,472\) logical source lines, \(237\,803\) bytes,
and SHA--256
\[
 \texttt{2C3FF58735842599CC3A406CC0B1B61EFC855889E4D7D21F6B3897CCC2B82483}.
\]

% =============================================================================
% END APPENDED SEVENTH-SERIES V17 MATERIAL
% =============================================================================

% =============================================================================
% BEGIN APPENDED SEVENTH-SERIES V18 MATERIAL
% =============================================================================

\chapter{Intrinsic conditional-projection spectral flow}
\label{ch:s7v18}

V17 compares two endpoint measures and therefore pays a global density
ratio.  This chapter keeps every interpolating measure on the compact
group and transports its \(L^2\) space unitarily to a fixed Hilbert
space.  The derivative of a conditional projection is then exactly a
symmetrized multiplication by the \emph{centered conditional score}.
This removes the global essential-infimum loss.  A first spectral-flow
bound still uses an essential-supremum score; its failure identifies the
precise place where the V9--V11 active curvature-source estimates must
enter next.

\section{A differentiable family of conditional expectations}
\label{sec:s7v18-family}

Let \(m\) be a fixed probability measure on the compact collar fibre
and let
\[
 d\mu_t=\rho_t\,dm,\qquad \rho_t>0,\qquad
 \int\rho_t\,dm=1,\qquad 0\le t\le1.
\]
Assume enough differentiability and boundedness to justify all operator
derivatives below, and put
\begin{equation}
 h_t:=\partial_t\log\rho_t,\qquad
 \int h_t\,d\mu_t=0.                                    \label{eq:s7v18-global-score}
\end{equation}
For block \(\Lambda_\alpha\), let
\[
 K_{\alpha,t}
 =\mathbb E_{\mu_t}[\cdot\mid\mathcal F_{\Lambda_\alpha^c}],
 \qquad
 D_{\alpha,t}=I-K_{\alpha,t},
\]
and define the centered conditional score
\begin{equation}
 s_{\alpha,t}:=
 h_t-K_{\alpha,t}h_t.                                   \label{eq:s7v18-score}
\end{equation}

\begin{lemma}[Derivative before Hilbert-space transport]
\label{lem:s7v18-raw-derivative}
On a fixed test function \(f\),
\begin{equation}
 \partial_tK_{\alpha,t}f
 =
 K_{\alpha,t}(fh_t)
 -(K_{\alpha,t}f)(K_{\alpha,t}h_t)
 =
 K_{\alpha,t}(f s_{\alpha,t}).                          \label{eq:s7v18-raw-derivative}
\end{equation}
\end{lemma}

\begin{proof}
Disintegrate \(\mu_t\) over the block and differentiate its normalized
conditional density.  Its logarithmic derivative is
\(h_t-K_{\alpha,t}h_t=s_{\alpha,t}\), exactly as in V9.  This gives the
last expression; expanding it gives the middle one.
\end{proof}

\section{Unitary transport to one Hilbert space}
\label{sec:s7v18-unitary}

Fix \(\mu_0\) as the Hilbert-space reference and define
\begin{equation}
 U_t:L^2(\mu_t)\longrightarrow L^2(\mu_0),
 \qquad
 U_tf=\left({\rho_t\over\rho_0}\right)^{1/2}f.           \label{eq:s7v18-U}
\end{equation}
Set
\[
 \widehat K_{\alpha,t}=U_tK_{\alpha,t}U_t^{-1},
 \qquad
 \widehat D_{\alpha,t}=I-\widehat K_{\alpha,t}.
\]
These are orthogonal projections in the fixed space \(L^2(\mu_0)\).

\begin{theorem}[Exact transported projection derivative]
\label{thm:s7v18-Kdot}
Let \(M_s\) denote multiplication by \(s\).  Then
\begin{equation}
 \boxed{
 \partial_t\widehat K_{\alpha,t}
 ={1\over2}U_t
 \bigl(
 K_{\alpha,t}M_{s_{\alpha,t}}
 +M_{s_{\alpha,t}}K_{\alpha,t}
 \bigr)U_t^{-1}.}                                      \label{eq:s7v18-Kdot}
\end{equation}
In particular the derivative is self-adjoint and off-diagonal with
respect to
\(\operatorname{Ran}\widehat K_{\alpha,t}\oplus
\operatorname{Ran}\widehat D_{\alpha,t}\).
\end{theorem}

\begin{proof}
Differentiating \eqref{eq:s7v18-U} gives
\[
 U_t^{-1}\dot U_t={1\over2}M_{h_t}.
\]
Thus, after conjugating back to \(L^2(\mu_t)\),
\begin{align*}
 U_t^{-1}\dot{\widehat K}_{\alpha,t}U_t
 &=
 \dot K_{\alpha,t}
 +{1\over2}(M_{h_t}K_{\alpha,t}
            -K_{\alpha,t}M_{h_t}).
\end{align*}
By \Cref{lem:s7v18-raw-derivative},
\[
 \dot K_{\alpha,t}
 =K_{\alpha,t}M_{h_t}
  -M_{K_{\alpha,t}h_t}K_{\alpha,t}.
\]
Since \(K_{\alpha,t}h_t\) is exterior measurable, its multiplication
operator commutes with \(K_{\alpha,t}\).  Substitution and
\(s_{\alpha,t}=h_t-K_{\alpha,t}h_t\) give
\eqref{eq:s7v18-Kdot}.

The right side is self-adjoint.  Moreover
\[
 K_{\alpha,t}M_{s_{\alpha,t}}K_{\alpha,t}=0,
\]
because \(K_{\alpha,t}s_{\alpha,t}=0\); hence both diagonal
compressions of the derivative vanish, proving the off-diagonal claim.
\end{proof}

\begin{corollary}[Exact quadratic-form derivative]
\label{cor:s7v18-form-derivative}
For \(F=U_tf\),
\begin{equation}
 \left\langle F,
 \partial_t\widehat D_{\alpha,t}F
 \right\rangle_{L^2(\mu_0)}
 =
 -\Re\left\langle
 K_{\alpha,t}f,\,
 s_{\alpha,t}D_{\alpha,t}f
 \right\rangle_{L^2(\mu_t)}.                            \label{eq:s7v18-form-derivative}
\end{equation}
\end{corollary}

\begin{proof}
Use \(\dot{\widehat D}=-\dot{\widehat K}\) and
\eqref{eq:s7v18-Kdot}.  The two symmetrized terms are complex
conjugates.  The component
\(\langle Kf,sKf\rangle\) vanishes because \(|Kf|^2\) is exterior
measurable and \(Ks=0\), leaving the displayed cross term.
\end{proof}

\section{Additive and eigenbranch bounds}
\label{sec:s7v18-flow}

Define the transported five-block generator
\begin{equation}
 \widehat{\cal A}_t
 :=\sum_{\alpha=1}^5\widehat D_{\alpha,t}.               \label{eq:s7v18-A}
\end{equation}

\begin{corollary}[Full operator-norm flow]
\label{cor:s7v18-operator-flow}
If
\[
 b_{\alpha}(t):=\|s_{\alpha,t}\|_{L^\infty(\mu_t)},
\]
then
\begin{align}
 \|\partial_t\widehat K_{\alpha,t}\|
 &\le b_\alpha(t),                                      \label{eq:s7v18-Kdot-bound}\\
 \|\widehat{\cal A}_t-\widehat{\cal A}_0\|
 &\le\int_0^t\sum_{\alpha=1}^5b_\alpha(r)\,dr.           \label{eq:s7v18-A-bound}
\end{align}
\end{corollary}

\begin{proof}
Conditional expectation and multiplication have norms at most one and
\(b_\alpha(t)\), respectively.  Apply
\eqref{eq:s7v18-Kdot}, then integrate the sum.
\end{proof}

\begin{warning}[The additive norm bound is too expensive]
\label{warn:s7v18-additive}
To preserve a reference gap \(\lambda_0\) directly through
\eqref{eq:s7v18-A-bound}, the integrated sum of five
essential-supremum scores must be smaller than \(\lambda_0\).  For a
Haar-to-Wilson interpolation the action score is of order \(s\) in
full-fibre supremum norm.  Splitting the path into more steps does not
reduce its integral.  The global density-ratio loss of V17 has been
removed, but a full-space score loss remains.
\end{warning}

There is a sharper bound along a positive eigenbranch.

\begin{theorem}[Square-root spectral-flow inequality]
\label{thm:s7v18-eigenflow}
Assume that on an interval \(I\), the bottom positive spectral value of
\(\widehat{\cal A}_t\) is an isolated differentiable eigenbranch
\(\lambda(t)>0\) with normalized eigenvector \(F_t=U_tf_t\).
Put
\begin{equation}
 \beta(t):=
 \left(\sum_{\alpha=1}^5b_\alpha(t)^2\right)^{1/2}.      \label{eq:s7v18-beta}
\end{equation}
Then almost everywhere on \(I\),
\begin{align}
 |\dot\lambda(t)|
 &\le\beta(t)\sqrt{\lambda(t)},                          \label{eq:s7v18-lambda-dot}\\
 \sqrt{\lambda(t)}
 &\ge
 \sqrt{\lambda(0)}
 -{1\over2}\int_0^t\beta(r)\,dr.                        \label{eq:s7v18-sqrt-flow}
\end{align}
\end{theorem}

\begin{proof}
The Hellmann--Feynman formula and
\Cref{cor:s7v18-form-derivative} give
\[
 \dot\lambda(t)
 =-\sum_{\alpha=1}^5
 \Re\langle K_{\alpha,t}f_t,
 s_{\alpha,t}D_{\alpha,t}f_t\rangle_t.
\]
Since \(\|K_{\alpha,t}f_t\|_t\le1\), Cauchy--Schwarz over the five
blocks yields
\begin{align*}
 |\dot\lambda(t)|
 &\le
 \left(\sum_\alpha
 b_\alpha(t)^2\|K_{\alpha,t}f_t\|_t^2\right)^{1/2}
 \left(\sum_\alpha\|D_{\alpha,t}f_t\|_t^2\right)^{1/2}\\
 &\le\beta(t)\sqrt{\lambda(t)}.
\end{align*}
Here the last sum equals
\(\langle F_t,\widehat{\cal A}_tF_t\rangle=\lambda(t)\).
Where \(\lambda(t)>0\),
\[
 {d\over dt}\sqrt{\lambda(t)}
 ={\dot\lambda(t)\over2\sqrt{\lambda(t)}}
 \ge-{1\over2}\beta(t).
\]
Integrate.
\end{proof}

\begin{firewall}[Spectral-edge regularity is a hypothesis]
\label{fire:s7v18-isolation}
The sum of conditional projections acts on an infinite-dimensional
group \(L^2\) space.  Its positive spectral edge need not be an isolated
eigenvalue merely because the lattice has finitely many links.
\Cref{thm:s7v18-eigenflow} is rigorous on isolated eigenbranches and on
finite Peter--Weyl truncations.  Extending it to a non-isolated edge
requires a lower-form or spectral-projection argument and is a separate
row.
\end{firewall}

\section{A legal compact reference measure}
\label{sec:s7v18-Haar}

Unlike the tangent Gaussian, product Haar measure lives on the same
compact chord space as the physical measure.

\begin{theorem}[Five-block product-Haar gap]
\label{thm:s7v18-Haar-gap}
For product Haar measure on the sixty chord variables and the five
blocks \eqref{eq:s7v15-blocks}, the heat-bath generator has exact gap
\begin{equation}
 \lambda_{\rm hb}^{\rm Haar}=1.                          \label{eq:s7v18-Haar-gap}
\end{equation}
\end{theorem}

\begin{proof}
Decompose \(L^2\) as the tensor product of constants and mean-zero
one-coordinate subspaces.  On a pure tensor whose nonconstant
coordinate support is the nonempty set \(S\), the block fluctuation
\(D_\alpha\) is the identity exactly when
\(\Lambda_\alpha\cap S\ne\varnothing\), and is zero otherwise.
The generator eigenvalue is therefore the number of blocks meeting
\(S\).  Every coordinate belongs to at least one block, so this number
is at least one.  A nonconstant function of a coordinate in positions
\(1\)--\(10\) or \(51\)--\(60\), which has block coverage one, attains
eigenvalue one.
\end{proof}

\begin{corollary}[Haar-to-physical sufficient condition]
\label{cor:s7v18-Haar-flow}
Under the isolated-edge hypotheses of
\Cref{thm:s7v18-eigenflow}, a Haar-to-physical interpolation preserves
a positive block gap if
\begin{equation}
 \int_0^1\beta(t)\,dt<2.                                 \label{eq:s7v18-Haar-flow}
\end{equation}
It then gives
\[
 \lambda_{\rm hb}(1)
 \ge
 \left(1-\frac12\int_0^1\beta(t)\,dt\right)^2.
\]
\end{corollary}

\begin{proof}
Use \(\lambda(0)=1\) from
\Cref{thm:s7v18-Haar-gap} in
\eqref{eq:s7v18-sqrt-flow}.
\end{proof}

\section{The Wilson interpolation score}
\label{sec:s7v18-Wilson}

For a path
\begin{equation}
 \rho_t
 =Z_t^{-1}
 \exp\{-t s{\cal V}\}\,
 u_t^2\,J_t,                                             \label{eq:s7v18-physical-path}
\end{equation}
the global scalar score is
\begin{equation}
 h_t
 =-s{\cal V}
 +2\partial_t\log u_t
 +\partial_t\log J_t
 -\partial_t\log Z_t.                                   \label{eq:s7v18-h-physical}
\end{equation}
The final scalar normalization disappears after conditional centering.

\begin{proposition}[Exact conditional path-score ledger]
\label{prop:s7v18-score-ledger}
For every block,
\begin{equation}
 s_{\alpha,t}
 =
 -sD_{\alpha,t}{\cal V}
 +2D_{\alpha,t}(\partial_t\log u_t)
 +D_{\alpha,t}(\partial_t\log J_t).                     \label{eq:s7v18-score-ledger}
\end{equation}
In fixed maximal-tree group coordinates,
\(\partial_t\log J_t=0\) for the Haar reference Jacobian.
\end{proposition}

\begin{proof}
Apply \(D_{\alpha,t}=I-K_{\alpha,t}\) to
\eqref{eq:s7v18-h-physical}.  It annihilates the scalar
\(\partial_t\log Z_t\).  The tree-Haar statement is
\Cref{cor:s7v11-no-J}.
\end{proof}

\begin{assessment}[Why the \(L^\infty\) condition is not the target]
\label{ass:s7v18-sup}
Even though \({\cal V}\) is bounded on the compact fibre,
\(\|sD_\alpha{\cal V}\|_\infty\) generally grows like \(s\).  The
moving-ground term can have its own inverse-gap cost.  Therefore
\eqref{eq:s7v18-Haar-flow} is not presently uniform in the continuum
scaling.  The exact form derivative
\eqref{eq:s7v18-form-derivative}, however, pairs the score only between
\(K_\alpha f\) and \(D_\alpha f\).  This is an active off-diagonal
matrix element, precisely the setting of the V9 curvature-adjoint and
V10 lower-angle estimates.
\end{assessment}

\begin{definition}[Active spectral-flow card, ASFC]
\label{def:s7v18-ASFC}
ASFC consists of:
\begin{enumerate}[label=\textup{(S\arabic*)},leftmargin=4em]
\item a compact-group interpolation \eqref{eq:s7v18-physical-path};
\item the exact transported derivative \eqref{eq:s7v18-Kdot};
\item a fixed or controlled positive spectral-edge subspace;
\item a bound on the off-diagonal matrix elements in
      \eqref{eq:s7v18-form-derivative} using the active
      curvature-source space rather than \(\|s_{\alpha,t}\|_\infty\);
\item the V10 conditional lower angle and V11 moving-ground
      factorization along the entire path;
\item control of bad fibres and any loss of spectral-edge isolation;
\item uniform block internal gaps and insertion into PBDC/MGCC/MLC;
      and
\item gauge, reflection, continuum, and Hamiltonian spectral contacts.
\end{enumerate}
\end{definition}

\begin{assessment}[Progress at seventh-series V18]
\label{ass:s7v18-status}
The global density-ratio obstruction of V17 has been removed at the
operator level.  Conditional projection motion is exactly the
off-diagonal centered-score operator
\eqref{eq:s7v18-Kdot}; product Haar supplies a legal compact reference
with gap one.  A full-supremum flow estimate is still nonuniform, but
the exact quadratic form shows that only an active score matrix element
must be controlled.
\end{assessment}

\begin{decision}[V18 bottleneck decision]
\label{dec:s7v18-bottleneck}
V19 will combine
\eqref{eq:s7v18-form-derivative} with the V9 Wilson
curvature-adjoint factorization and the V10 fixed-reference polar
metric.  It will derive a relative form bound on the spectral
subspace, keeping the moving-ground residual separate.  If the action
coefficient \(s\) survives without compensation from the block
Dirichlet form, the interpolation parameterization or reference
generator must be rescaled rather than hidden in a score norm.
\end{decision}

\begin{firewall}[Claim boundary after seventh-series V18]
\label{fire:s7v18-boundary}
V18 proves the exact transported conditional-projection derivative,
its quadratic-form identity, additive and isolated-eigenbranch flow
bounds, the exact product-Haar block gap, and the physical path-score
ledger.

V18 does not prove an active matrix-element bound, spectral-edge
isolation in the untruncated space, a uniform physical block gap,
twenty-chord internal gaps, remaining martingale budgets, continuum
construction, or the Yang--Mills mass gap.
\end{firewall}

\paragraph{Parent-version record.}
V18 was formed from
\texttt{Renewal\_Yang\_Mills\_Mass\_Gap\_Research\_Notes\_V17.tex}.
The V17 parent had \(6\,836\) logical source lines, \(250\,571\) bytes,
and SHA--256
\[
 \texttt{71B7F707094873D98BCE96ED21F55B156F0968CCE03F1DD4504E35CB506D9BCE}.
\]

% =============================================================================
% END APPENDED SEVENTH-SERIES V18 MATERIAL
% =============================================================================

% =============================================================================
% BEGIN APPENDED SEVENTH-SERIES V19 MATERIAL
% =============================================================================

\chapter{Configuration transport, divergence scores, and a type correction}
\label{ch:s7v19}

V18 proposed inserting the V9 curvature-adjoint action source into the
Haar-to-physical spectral flow.  A type check shows that this cannot be
done for the coupling interpolation: its score is an action
\emph{value} fluctuation, whereas V9 controls a coordinate
\emph{derivative}.  This chapter records that correction and derives
the setting in which the polar cancellation does apply.  If the measure
path is generated by a configuration-space velocity, its score is a
weighted divergence and its action term is the directional derivative
\(\langle B^*G,X\rangle\).  Constructing a generic such velocity by a
Poisson equation, however, costs precisely the fibre gap one is trying
to prove.

\section{The coupling-score type mismatch}
\label{sec:s7v19-type}

For the direct Gibbs interpolation
\[
 d\mu_t=Z_t^{-1}e^{-ts{\cal V}}\,d\mu_{\rm ref},
\]
the score is
\begin{equation}
 h_t=-s({\cal V}-\mathbb E_{\mu_t}{\cal V}),             \label{eq:s7v19-coupling-score}
\end{equation}
and its block-centered version is
\[
 s_{\alpha,t}=-sD_{\alpha,t}{\cal V}.
\]
By contrast, V9 proves
\begin{equation}
 \nabla{\cal V}=B^*G,                                   \label{eq:s7v19-gradient-source}
\end{equation}
an identity in the chord tangent space.

\begin{proposition}[Action fluctuation is not the curvature-adjoint gradient]
\label{prop:s7v19-type}
The scalar \(D_\alpha{\cal V}\) in
\eqref{eq:s7v19-coupling-score} cannot be replaced by the vector
\(B^*G\) in \eqref{eq:s7v19-gradient-source}.  Near a nondegenerate
flat point they also have different vanishing order: in local
coordinates,
\begin{align}
 {\cal V}(x)&={1\over2}\langle x,Hx\rangle+O(|x|^3),
                                                               \label{eq:s7v19-V-Taylor}\\
 \nabla{\cal V}(x)&=Hx+O(|x|^2).                         \label{eq:s7v19-grad-Taylor}
\end{align}
Conditional centering changes the quadratic scalar by another scalar
quadratic moment; it does not turn it into the linear tangent vector.
\end{proposition}

\begin{proof}
The two expressions lie in different spaces:
\(D_\alpha{\cal V}\in L^2(\mu_t;\mathbb R)\), while
\(\nabla{\cal V}\in L^2(\mu_t;T\Omega)\).
The Taylor formulae follow from
\({\cal V}(0)=0\), \(\nabla{\cal V}(0)=0\), and Hessian \(H\).
Conditional expectation preserves scalar type and maps a quadratic
Taylor term to its conditional scalar average.  Thus neither type nor
order agrees.
\end{proof}

\begin{firewall}[Correction to the V18 programme]
\label{fire:s7v19-correction}
The V9 polar identity may enter a measure spectral flow only after the
path score contains a directional derivative \(X{\cal V}\).  Applying
it directly to the coupling score
\({\cal V}-K_\alpha{\cal V}\) is invalid.  All V18 identities remain
correct; only that proposed next substitution is rejected.
\end{firewall}

\section{Scores generated by a configuration flow}
\label{sec:s7v19-flow}

Let \(M\) be the compact collar group manifold with product Haar volume
\(dm\).  For a positive density \(\rho\), define the weighted divergence
\begin{equation}
 \operatorname{div}_{\mu}X
 :={1\over\rho}\operatorname{div}_m(\rho X)
 =\operatorname{div}_mX+X\log\rho,
 \qquad d\mu=\rho\,dm.                                  \label{eq:s7v19-weighted-div}
\end{equation}

\begin{theorem}[Continuity-equation score]
\label{thm:s7v19-continuity}
Suppose a smooth measure path satisfies
\begin{equation}
 \partial_t\rho_t+\operatorname{div}_m(\rho_tX_t)=0.     \label{eq:s7v19-continuity}
\end{equation}
Then its logarithmic score is
\begin{equation}
 h_t:=\partial_t\log\rho_t
 =-\operatorname{div}_{\mu_t}X_t.                       \label{eq:s7v19-div-score}
\end{equation}
For every smooth \(f\),
\begin{equation}
 \int_M f h_t\,d\mu_t
 =\int_M X_tf\,d\mu_t.                                  \label{eq:s7v19-transport-ibp}
\end{equation}
\end{theorem}

\begin{proof}
Divide \eqref{eq:s7v19-continuity} by \(\rho_t\) and use
\eqref{eq:s7v19-weighted-div}.  On the closed manifold,
\[
 \int f h_t\rho_t\,dm
 =-\int f\,\operatorname{div}_m(\rho_tX_t)\,dm
 =\int X_tf\,\rho_t\,dm,
\]
which proves \eqref{eq:s7v19-transport-ibp}.
\end{proof}

\begin{corollary}[Exact block score conversion]
\label{cor:s7v19-block-conversion}
Let \(K_{\alpha,t}\) and \(D_{\alpha,t}\) be the V18 block projections
and \(s_{\alpha,t}=D_{\alpha,t}h_t\).  Then
\begin{equation}
 \left\langle K_{\alpha,t}f,\,
 s_{\alpha,t}D_{\alpha,t}f\right\rangle_t
 =
 \int_M
 X_t\!\left(
 \overline{K_{\alpha,t}f}\,D_{\alpha,t}f
 \right)d\mu_t.                                        \label{eq:s7v19-block-conversion}
\end{equation}
\end{corollary}

\begin{proof}
The term containing \(K_{\alpha,t}h_t\) vanishes:
\[
 \int\overline{K_\alpha f}\,(K_\alpha h_t)D_\alpha f\,d\mu_t=0,
\]
because the first two factors are exterior measurable and
\(K_\alpha D_\alpha f=0\).  The remaining integral is
\(\int\overline{K_\alpha f}\,h_tD_\alpha f\,d\mu_t\).
Apply \eqref{eq:s7v19-transport-ibp} to the product.
\end{proof}

\begin{corollary}[Block-internal velocity]
\label{cor:s7v19-internal}
If \(X_t\) differentiates only coordinates in \(\Lambda_\alpha\), then
\begin{equation}
 \left\langle K_{\alpha,t}f,\,
 s_{\alpha,t}D_{\alpha,t}f\right\rangle_t
 =
 \int_M\overline{K_{\alpha,t}f}\,X_tf\,d\mu_t.           \label{eq:s7v19-internal}
\end{equation}
\end{corollary}

\begin{proof}
An exterior-measurable representative of \(K_{\alpha,t}f\) is constant
along block-internal directions, so
\(X_tK_{\alpha,t}f=0\).  Consequently
\(X_tD_{\alpha,t}f=X_tf\).  Insert these facts in
\eqref{eq:s7v19-block-conversion}.
\end{proof}

\section{The directional Wilson term}
\label{sec:s7v19-Wilson}

For the physical density
\[
 \rho=Z^{-1}e^{-s{\cal V}}u_0^2J,
\]
\eqref{eq:s7v19-weighted-div} gives
\begin{equation}
 -\operatorname{div}_{\mu}X
 =
 -\operatorname{div}_mX
 +sX{\cal V}
 -2X\log u_0
 -X\log J.                                               \label{eq:s7v19-physical-div}
\end{equation}

\begin{theorem}[Polar control of the transport action score]
\label{thm:s7v19-polar-transport}
On a core where the selected curvature derivative \(B\) is injective,
let \(H=B^*B\).  If a transport velocity has the form
\begin{equation}
 X=H^{-1/2}Z,                                            \label{eq:s7v19-X}
\end{equation}
then
\begin{equation}
 |X{\cal V}|
 =|\langle G,BH^{-1/2}Z\rangle|
 \le\|P_{\operatorname{Ran}B}G\|\,\|Z\|
 \le\|G\|\,\|Z\|.                                       \label{eq:s7v19-polar-transport}
\end{equation}
No inverse singular value appears.
\end{theorem}

\begin{proof}
By V9,
\[
 X{\cal V}
 =\langle\nabla{\cal V},X\rangle
 =\langle B^*G,H^{-1/2}Z\rangle
 =\langle G,BH^{-1/2}Z\rangle.
\]
The polar map \(BH^{-1/2}\) is an isometry onto
\operatorname{Ran}B\); its adjoint annihilates the orthogonal
complement.  Cauchy--Schwarz proves the estimate.
\end{proof}

\begin{corollary}[Defect-weighted transport score]
\label{cor:s7v19-defect}
On the V9 small Wilson chart,
\begin{equation}
 |X{\cal V}|
 \le\sqrt{C_G{\cal V}_{\rm inc}}\,\|Z\|.                \label{eq:s7v19-defect}
\end{equation}
\end{corollary}

\begin{proof}
Use \eqref{eq:s7v9-G-bound} in
\eqref{eq:s7v19-polar-transport}.
\end{proof}

\begin{warning}[The velocity must be constructed]
\label{warn:s7v19-velocity}
Equation \eqref{eq:s7v19-X} is useful only if the continuity equation
produces \(Z=H^{1/2}X\) with a controlled norm.  Defining \(X\) by an
uncontrolled inverse divergence and then invoking the polar identity
would simply move the missing gap into the transport velocity.
\end{warning}

\section{Generic Moser transport is gap-equivalent}
\label{sec:s7v19-Moser}

For a fixed positive measure \(\mu\), define the weighted Laplacian
\begin{equation}
 {\cal L}_\mu
 :=-\operatorname{div}_\mu\nabla.                       \label{eq:s7v19-Lmu}
\end{equation}
It is nonnegative in \(L^2(\mu)\) and has Dirichlet form
\(\int|\nabla f|^2d\mu\).

\begin{theorem}[Minimum-norm divergence solver]
\label{thm:s7v19-divergence-solver}
Assume \({\cal L}_\mu\) has gap \(g_\mu>0\).  For centered
\(h\in L^2(\mu)\), the gradient field
\begin{equation}
 X_h:=\nabla{\cal L}_\mu^{-1}h                           \label{eq:s7v19-Xh}
\end{equation}
satisfies
\begin{equation}
 -\operatorname{div}_\mu X_h=h,\qquad
 \|X_h\|_{L^2(\mu;TM)}^2
 =\langle h,{\cal L}_\mu^{-1}h\rangle
 \le {1\over g_\mu}\|h\|_2^2.                           \label{eq:s7v19-Xh-bound}
\end{equation}
Among \(L^2\) vector fields solving the divergence equation, \(X_h\)
has minimum norm.
\end{theorem}

\begin{proof}
The first identity follows from
\eqref{eq:s7v19-Lmu}.  Integration by parts gives
\[
 \|\nabla{\cal L}_\mu^{-1}h\|_2^2
 =\langle{\cal L}_\mu^{-1}h,
 {\cal L}_\mu{\cal L}_\mu^{-1}h\rangle
 =\langle h,{\cal L}_\mu^{-1}h\rangle.
\]
The spectral theorem on the centered subspace gives the upper bound.
If another solution is \(X_h+Y\), then
\(\operatorname{div}_\mu Y=0\), so
\[
 \langle X_h,Y\rangle
 =\langle\nabla{\cal L}_\mu^{-1}h,Y\rangle
 =-\langle{\cal L}_\mu^{-1}h,\operatorname{div}_\mu Y\rangle=0.
\]
Pythagoras proves minimality.
\end{proof}

\begin{theorem}[Uniform right inverse is equivalent to a Poincar\'e gap]
\label{thm:s7v19-equivalence}
The optimal operator norm of the minimum-norm solution map
\[
 h\longmapsto X_h:
 L_0^2(\mu)\longrightarrow L^2(\mu;TM)
\]
is \(g_\mu^{-1/2}\).  Consequently a cutoff-uniform \(L^2\) Moser
velocity bound for arbitrary centered scores is equivalent to a
cutoff-uniform weighted Poincar\'e gap.
\end{theorem}

\begin{proof}
By spectral calculus,
\[
 \|X_h\|_2^2
 =\langle h,{\cal L}_\mu^{-1}h\rangle.
\]
The operator norm squared is
\(\|{\cal L}_\mu^{-1}\|_{1^\perp}=g_\mu^{-1}\).  Approximate the bottom
of the positive spectrum to see that this bound is sharp.
\end{proof}

\begin{firewall}[Generic transport is circular]
\label{fire:s7v19-Moser}
Moser's continuity-equation construction proves existence of a
transport between smooth positive compact-fibre densities.  If its
velocity is estimated through
\({\cal L}_{\mu_t}^{-1}\), then
\Cref{thm:s7v19-equivalence} assumes exactly a uniform fibre Poincar\'e
gap.  Such a generic construction cannot prove the V11 gap row.
The transport must instead be explicit on the active Wilson source or
use a weaker norm already controlled by the programme.
\end{firewall}

\section{A block-local transport target}
\label{sec:s7v19-block-target}

\begin{definition}[Curvature-adapted transport card, CATC]
\label{def:s7v19-CATC}
CATC consists of:
\begin{enumerate}[label=\textup{(C\arabic*)},leftmargin=4em]
\item a compact-group measure path and a decomposition
      \(X_t=\sum_{\alpha=1}^5X_{\alpha,t}\) into block-supported
      velocities;
\item the exact continuity equation
      \(\partial_t\rho_t+\operatorname{div}_m(\rho_tX_t)=0\);
\item representations
      \(X_{\alpha,t}=H_{\alpha,t}^{-1/2}Z_{\alpha,t}\)
      with active \(L^2\) or form bounds on \(Z_{\alpha,t}\);
\item the polar action estimate
      \eqref{eq:s7v19-polar-transport};
\item direct divergence, moving-ground, and regulator estimates in
      \eqref{eq:s7v19-physical-div};
\item control of cross-block terms in
      \eqref{eq:s7v19-block-conversion};
\item a spectral-flow theorem on the positive edge without a full-space
      essential-supremum score; and
\item PBDC, MGCC, MLC, good/bad, gauge, reflection, continuum, and
      spectral contacts.
\end{enumerate}
\end{definition}

\begin{theorem}[What CATC would repair]
\label{thm:s7v19-CATC}
Rows (C1)--(C6) convert the V18 score matrix element into derivatives of
the test and curvature-weighted transport fields, where the V9 polar
identity is type-correct.  Rows (C7)--(C8), if proved with strict
cofinal constants, would transport the known Haar block gap toward the
physical measure without the global density ratio of V17 or the
generic inverse-gap construction of
\Cref{thm:s7v19-divergence-solver}.
\end{theorem}

\begin{proof}
The score conversion is
\Cref{thm:s7v19-continuity,cor:s7v19-block-conversion}; its action term
is \Cref{thm:s7v19-polar-transport}.  The circular generic alternative
is excluded by \Cref{thm:s7v19-equivalence}.  The remaining rows state
the downstream interfaces explicitly.
\end{proof}

\begin{assessment}[Progress at seventh-series V19]
\label{ass:s7v19-status}
The coupling-score/curvature-gradient mismatch has been corrected.
Configuration transport supplies the right object: a directional
Wilson derivative receiving exact polar cancellation.  The full score
matrix element becomes a derivative of the test by
\eqref{eq:s7v19-block-conversion}.  Generic Moser transport is not a
solution, because its optimal \(L^2\) velocity norm is exactly the
inverse square root of the missing fibre gap.
\end{assessment}

\begin{decision}[V19 bottleneck decision]
\label{dec:s7v19-bottleneck}
V20 is the compulsory companion audit.  After the audit it will test
whether the triangular curvature coordinates of V15--V16 generate an
explicit block-local transport.  The first target is the Jacobian
formula for pushing a density along the selected curvature map and the
exact residual created by plaquette rows outside the selected sixty.
If the transport again requires solving a global divergence equation,
that branch will be stopped.
\end{decision}

\begin{firewall}[Claim boundary after seventh-series V19]
\label{fire:s7v19-boundary}
V19 proves the coupling-score type correction, the
continuity-equation score and block conversion identities, polar
control of a curvature-adapted transport action term, and the exact
equivalence between a generic \(L^2\) divergence solver and the missing
Poincar\'e gap.

V19 does not construct the required block-local transport, prove a
physical heat-bath or fibre gap, establish active spectral flow,
remaining martingale budgets, continuum construction, or the
Yang--Mills mass gap.
\end{firewall}

\paragraph{Parent-version record.}
V19 was formed from
\texttt{Renewal\_Yang\_Mills\_Mass\_Gap\_Research\_Notes\_V18.tex}.
The V18 parent had \(7\,268\) logical source lines, \(265\,260\) bytes,
and SHA--256
\[
 \texttt{9BF7F4ADB8EC92B8B33F0176CFC933647C3A8087E5309CD887D7F63ECAF7BD57}.
\]

% =============================================================================
% END APPENDED SEVENTH-SERIES V19 MATERIAL
% =============================================================================

% =============================================================================
% BEGIN APPENDED SEVENTH-SERIES V20 MATERIAL
% =============================================================================

\chapter{Companion audit and curvature-coordinate dilation transport}
\label{ch:s7v20}

This is the compulsory five-version companion audit.  The newest
Standard--Model notes independently develop the Gaussian block-frame
route and then exhibit a volume obstruction for fixed local
charge-preserving blocks.  That result changes the Yang--Mills
programme at the global level: the five collar blocks are a local
compiler, not yet a volume-uniform dynamics.  The constructive part of
this chapter remains local and resolves the V19 transport question.
Scaling the selected curvature coordinates gives an explicit
non-Poisson velocity.  It transports the quadratic Gaussian exactly;
for the physical measure the entire failure is the centered material
derivative of one declared nonlinear/ground/Jacobian remainder.

\section{Compulsory V20 companion audit}
\label{sec:s7v20-audit}

\begin{audit}[Files inspected]
\label{aud:s7v20-files}
At the audit time the active companion files were:
\begin{enumerate}[label=\textup{(\alph*)},leftmargin=3em]
\item
\texttt{Renewal\_SM\_Einstein\_Continuum\_Research\_Notes\_V63.tex},
with \(25\,305\) logical lines, \(899\,300\) bytes, and SHA--256
\[
 \texttt{666801AFFF6EC6D0D155B016B52FA06F825BF85B4A5A4A98343C4C2009042CD2};
\]
\item
\texttt{Renewal\_NS\_Stability\_Research\_Notes\_V26.tex},
with \(5\,319\) logical lines, \(278\,283\) bytes, and SHA--256
\[
 \texttt{82C887F8C2C69E4F28FAD7C3DC8DE5B9099CB5A4BF431EBA1FFF3E91935978AD}.
\]
\end{enumerate}
The SM programme advanced from V55 through V63.  The NS file remains
unchanged.
\end{audit}

\begin{assessment}[SM V56--V59]
\label{ass:s7v20-SM-score}
SM V56 derives the exact finite-regulator Ward score and its
observable-active dual norm.  V57 corrects a type error between a
diffeomorphism Ward derivative and a lapse-multiplier constraint
derivative.  This directly supports the YM V19 correction: centered
scores with different parameter directions cannot be identified.
V58--V59 derive positive Gaussian constraint penalties, their coarea
limit, and the normal/surface Hessian split.  Their normal stiffness
does not by itself control surface mixing.
\end{assessment}

\begin{assessment}[SM V60--V63]
\label{ass:s7v20-SM-frame}
SM V60 audits YM V14--V16 and transports the all-chaos frame theorem to
a linearized constraint-normal Gaussian.  V61 proves that nearest
neighbour charge-preserving pair updates have exact gap
\[
 { \bigl(4\sin^2(\pi/L)\bigr)^3\over4d(2d+3)}
 \asymp L^{-6}.
\]
V62 proves that the \(L^{-6}\) collapse holds for every fixed finite
family of bounded-range rank-one charge-preserving translated shapes.
Fourier-shell updates have gap one but are nonlocal.  V63 constructs a
real-space dyadic frame with margin \(1/3\), but the required total
update intensity grows like \(L^2/12\); normalizing the physical clock
restores an \(O(L^{-2})\) slow mode.

This does not prove a corresponding Yang--Mills obstruction, because
the constraint symbol and physical sectors differ.  It does prove that
a fixed local collar gap cannot simply be declared volume uniform.
The Yang--Mills global frame and physical clock must be computed.
\end{assessment}

\begin{assessment}[NS V26]
\label{ass:s7v20-NS}
No new NS note appeared.  Its rank-one kernel refinement does not
alter the transport or multiscale decisions.
\end{assessment}

\begin{decision}[Audit-adjusted direction]
\label{dec:s7v20-audit}
Retain the V15--V16 collar blocks as local conditional tools.  In
parallel, add a mandatory global first-chaos/multiscale frame test
before treating their local gap as a field-theory gap.  For the present
local transport problem, use the explicit selected curvature map and
measure every departure from quadratic scaling.
\end{decision}

\section{Selected curvature coordinates}
\label{sec:s7v20-coordinates}

Let \({\cal U}\) be a star-shaped tree-gauge neighbourhood of the flat
collar on which the selected curvature map
\begin{equation}
 F:{\cal U}\longrightarrow{\cal Y}\subset
 \mathfrak{su}(3)^{60},\qquad
 F(U)=(\log W_{r_k}(U))_{k=1}^{60}                       \label{eq:s7v20-F}
\end{equation}
is a diffeomorphism onto its image.  Write
\[
 B_U=DF_U,\qquad H_U=B_U^*B_U.
\]
By V16, at flatness
\[
 \sigma_{\min}(B_\star)\ge1/\sqrt{28},
\]
and on the V16 perturbation core
\(\sigma_{\min}(B_U)\ge1/(2\sqrt{28})\).
The real coordinate dimension is
\begin{equation}
 n=60\dim\mathfrak{su}(3)=480.                           \label{eq:s7v20-n}
\end{equation}

\section{Explicit curvature dilation}
\label{sec:s7v20-dilation}

Let \(s(t)>0\) be differentiable and define in curvature coordinates
\begin{equation}
 Y_t(y):=-{\dot s(t)\over2s(t)}\,y.                      \label{eq:s7v20-Y}
\end{equation}
Pull it back to the link chart:
\begin{equation}
 X_t(U):=B_U^{-1}Y_t(F(U)).                              \label{eq:s7v20-X}
\end{equation}

\begin{theorem}[Exact non-Poisson curvature velocity]
\label{thm:s7v20-velocity}
The field \eqref{eq:s7v20-X} is explicitly defined on \({\cal U}\) and
satisfies
\begin{align}
 B_UX_t(U)
 &=-{\dot s\over2s}F(U),                                 \label{eq:s7v20-BX}\\
 \|H_U^{1/2}X_t(U)\|
 &={|\dot s|\over2s}\|F(U)\|.                            \label{eq:s7v20-Xnorm}
\end{align}
No divergence equation or fibre spectral gap is used.
\end{theorem}

\begin{proof}
The first identity is the definition.  For the second,
\[
 \|H_U^{1/2}X_t\|^2
 =\langle X_t,B_U^*B_UX_t\rangle
 =\|B_UX_t\|^2.
\]
Insert \eqref{eq:s7v20-BX}.
\end{proof}

\begin{corollary}[Flow in curvature coordinates]
\label{cor:s7v20-flow}
As long as the trajectory stays in \({\cal Y}\), the flow obeys
\begin{equation}
 F(U_t)
 =
 \left({s(0)\over s(t)}\right)^{1/2}F(U_0).              \label{eq:s7v20-flow}
\end{equation}
\end{corollary}

\begin{proof}
Equation \eqref{eq:s7v20-BX} is
\[
 {d\over dt}F(U_t)=-{\dot s\over2s}F(U_t).
\]
Integrate the scalar linear equation.
\end{proof}

\begin{corollary}[Polar action estimate for the explicit velocity]
\label{cor:s7v20-action}
For the complete incident Wilson potential of V9,
\begin{equation}
 |X_t{\cal V}_{\rm inc}|
 \le {|\dot s|\over2s}
 \|G_U\|\,\|F(U)\|
 \le {|\dot s|\over2s}
 \sqrt{C_G{\cal V}_{\rm inc}(U)}\,\|F(U)\|.              \label{eq:s7v20-action}
\end{equation}
\end{corollary}

\begin{proof}
Apply \Cref{thm:s7v19-polar-transport} with
\(Z=H_U^{1/2}X_t\), then use
\eqref{eq:s7v20-Xnorm} and the V9 defect bound.
\end{proof}

\section{Exact density defect of curvature dilation}
\label{sec:s7v20-defect}

Work first in a curvature-coordinate domain on which the flow has no
boundary flux.  Let
\begin{equation}
 Q(y):={1\over2}\langle y,A_{\rm inc}y\rangle            \label{eq:s7v20-Q}
\end{equation}
be the complete quadratic Taylor polynomial of
\({\cal V}_{\rm inc}(F^{-1}(y))\).  The matrix
\(A_{\rm inc}\) is positive definite because it contains the positive
selected-row form as a summand.  Consider a normalized density
\begin{equation}
 d\nu_t(y)
 =
 Z_t^{-1}
 \exp\!\left\{-s(t)Q(y)+{\cal R}_t(y)\right\}dy.
                                                               \label{eq:s7v20-density}
\end{equation}
Let \(h_t=\partial_t\log(d\nu_t/dy)\), and define the continuity defect
\begin{equation}
 {\mathfrak r}_t
 :=h_t+\operatorname{div}_{\nu_t}Y_t.                   \label{eq:s7v20-r}
\end{equation}

\begin{theorem}[Centered material-remainder identity]
\label{thm:s7v20-remainder}
For the dilation field \eqref{eq:s7v20-Y},
\begin{equation}
 \boxed{
 {\mathfrak r}_t
 =
 (I-\mathbb E_{\nu_t})
 \bigl[(\partial_t+Y_t){\cal R}_t\bigr].}                \label{eq:s7v20-remainder}
\end{equation}
In particular, when \({\cal R}_t\equiv0\), curvature dilation solves the
continuity equation exactly.
\end{theorem}

\begin{proof}
Differentiate the normalized log density:
\[
 h_t
 =-\dot s\,Q(y)+\partial_t{\cal R}_t
 -\mathbb E_{\nu_t}
   \left[-\dot s\,Q(y)+\partial_t{\cal R}_t\right].
\]
Since
\(\operatorname{div}Y_t=-n\dot s/(2s)\),
\[
 Y_t(-sQ)=\dot s\,Q,
\]
because \(Q\) is homogeneous of degree two.
Thus the nonconstant quadratic terms cancel in
\(h_t+\operatorname{div}_{\nu_t}Y_t\).  The remaining nonconstant
part is
\((\partial_t+Y_t){\cal R}_t\).  Finally,
\[
 \mathbb E_{\nu_t}\operatorname{div}_{\nu_t}Y_t=0
\]
by the no-flux assumption, while \(\mathbb E_{\nu_t}h_t=0\);
therefore the total defect has mean zero.  Subtracting the mean of the
remaining part gives \eqref{eq:s7v20-remainder}.
\end{proof}

\begin{corollary}[Exact quadratic transport]
\label{cor:s7v20-Gaussian}
The family
\[
 {\det(sA_{\rm inc})^{1/2}\over(2\pi)^{n/2}}
 e^{-s\langle y,A_{\rm inc}y\rangle/2}dy
\]
is transported by
\(Y_t=-(\dot s/2s)y\) with zero residual.
\end{corollary}

\begin{proof}
Set \({\cal R}_t=0\) in
\Cref{thm:s7v20-remainder}.
\end{proof}

\section{The physical remainder ledger}
\label{sec:s7v20-ledger}

Let \(U=F^{-1}(y)\).  The product-Haar measure in curvature coordinates
has a positive local density \(j_F(y)\,dy\).  Write the complete
physical density on the chart as in
\eqref{eq:s7v20-density}, with
\begin{equation}
 \begin{split}
 {\cal R}_t(y)
 &:=
 -s(t)\left[
 {\cal V}_{\rm inc}(F^{-1}(y))-Q(y)
 \right]\\
 &\quad
 +2\log u_t(F^{-1}(y))
 +\log j_F(y)+\log J_{{\rm reg},t}(y).
 \end{split}                                             \label{eq:s7v20-physical-R}
\end{equation}
Here \(Q\) includes the quadratic contributions of all \(250\) incident
plaquette rows.  The first line therefore contains only their nonlinear
Taylor remainder, including the \(190\) rows not used as coordinates.

\begin{theorem}[Exact physical curvature-dilation defect]
\label{thm:s7v20-physical-defect}
On the declared chart and under the no-flux condition, the difference
between the target physical score and the score generated by the
explicit velocity \(X_t\) is exactly
\begin{equation}
 {\mathfrak r}_t
 =
 (I-\mathbb E_{\nu_t})
 (\partial_t+Y_t)
 \left\{
 -s\left({\cal V}_{\rm inc}-Q(F)\right)
 +2\log u_t+\log j_F+\log J_{{\rm reg},t}
 \right\}.                                               \label{eq:s7v20-physical-defect}
\end{equation}
\end{theorem}

\begin{proof}
Substitute \eqref{eq:s7v20-physical-R} into
\eqref{eq:s7v20-remainder}.  The pullback relation
\eqref{eq:s7v20-X} identifies \(Y_t\) acting on coordinate functions
with \(X_t\) acting on link functions.
\end{proof}

\begin{lemma}[Cubic Wilson material bound]
\label{lem:s7v20-cubic}
Suppose on \(\|y\|\le\rho\),
\begin{align}
 \left|{\cal V}_{\rm inc}(F^{-1}(y))
       -Q(y)\right|
 &\le C_3\|y\|^3,                                       \label{eq:s7v20-cubic-a}\\
 \left|y\cdot\nabla_y
 \left({\cal V}_{\rm inc}(F^{-1}(y))
       -Q(y)\right)\right|
 &\le C_3'\|y\|^3.                                      \label{eq:s7v20-cubic-b}
\end{align}
Then the Wilson part of the uncentered material remainder is bounded by
\begin{equation}
 |\dot s|\left(C_3+{C_3'\over2}\right)\rho^3.            \label{eq:s7v20-cubic}
\end{equation}
\end{lemma}

\begin{proof}
Put
\[
 R_3(y)={\cal V}_{\rm inc}(F^{-1}(y))-Q(y).
\]
Since \(R_3\) has no explicit \(t\)-dependence,
\[
 (\partial_t+Y_t)(-sR_3)
 =-\dot s\,R_3+{\dot s\over2}y\cdot\nabla R_3.
\]
Apply \eqref{eq:s7v20-cubic-a,eq:s7v20-cubic-b}.
\end{proof}

\begin{warning}[Chart boundary and omitted sectors are real residuals]
\label{warn:s7v20-boundary}
The radial field is not tangent to the boundary of an arbitrary
logarithm chart.  A smooth cutoff that makes it vanish near the chart
boundary creates an additional annular divergence and action residual.
Extending it by zero leaves the target score unchanged outside the
chart.  Neither contribution is controlled by
\Cref{lem:s7v20-cubic}.  Likewise, the unselected plaquette,
moving-ground, and regulator terms in
\eqref{eq:s7v20-physical-defect} must be estimated rather than absorbed
into the selected quadratic form.
\end{warning}

\section{Local versus global block frames}
\label{sec:s7v20-multiscale}

\begin{definition}[Yang--Mills multiscale frame test, YMFT]
\label{def:s7v20-YMFT}
Before a local collar gap is used in a volume-uniform conclusion, YMFT
requires:
\begin{enumerate}[label=\textup{(M\arabic*)},leftmargin=4em]
\item the linearized physical Yang--Mills configuration space after
      gauge and harmonic sectors are declared;
\item translated local curvature-block update subspaces on a periodic
      box;
\item the exact Fourier multiplier of their first-chaos frame;
\item its least nonzero value as volume grows;
\item if it collapses, an explicit multiscale/coarse block family and
      its update intensities;
\item normalization by the intended physical time and a new least-mode
      calculation;
\item nonlinear conditional-projection transport for the successful
      frame; and
\item reflection, continuum, Hamiltonian, and spectral contacts.
\end{enumerate}
\end{definition}

\begin{firewall}[The Einstein \(L^{-6}\) exponent is not imported]
\label{fire:s7v20-exponent}
The SM obstruction uses a scalar Laplacian constraint and
charge-preserving update shapes.  Yang--Mills has a one-derivative
curvature symbol on the transverse gauge sector and separate toron
modes.  Its low-frequency exponent must be computed from its own
symbol.  What transfers is the requirement to perform the frame and
clock calculation, not the value \(6\).
\end{firewall}

\begin{assessment}[Progress at seventh-series V20]
\label{ass:s7v20-status}
The required audit has added a global multiscale gate to the programme.
Locally, the selected curvature chart supplies an explicit dilation
velocity with exact curvature norm and no inverse divergence.  Its
quadratic Gaussian part is transported exactly.  Every physical
failure is consolidated in the centered material remainder
\eqref{eq:s7v20-physical-defect}, whose Wilson nonlinear part is cubic
on the core.  Boundary, unselected-row, moving-ground, and regulator
terms remain.
\end{assessment}

\begin{decision}[V20 bottleneck decision]
\label{dec:s7v20-bottleneck}
The global frame test is logically prior to further optimization of the
local transport residual.  V21 will linearize the translated
Yang--Mills collar updates on a periodic four-dimensional lattice,
project to the transverse one-form sector, and compute the Fourier
symbol of the local block frame.  Harmonic torons will remain explicit.
The result will decide whether fixed collar blocks have a
volume-uniform Gaussian heat-bath margin or require a dyadic hierarchy.
\end{decision}

\begin{firewall}[Claim boundary after seventh-series V20]
\label{fire:s7v20-boundary}
V20 performs the required SM/NS audit, constructs the explicit
curvature-coordinate dilation, proves its exact quadratic transport and
centered material-remainder identity, bounds the local cubic Wilson
term, and introduces the mandatory YM global frame test.

V20 does not control the complete material remainder, chart boundary,
global physical block frame, local conditional gaps, active spectral
flow, remaining martingale budgets, continuum construction, or the
Yang--Mills mass gap.
\end{firewall}

\paragraph{Parent-version record.}
V20 was formed from
\texttt{Renewal\_Yang\_Mills\_Mass\_Gap\_Research\_Notes\_V19.tex}.
The V19 parent had \(7\,679\) logical source lines, \(279\,890\) bytes,
and SHA--256
\[
 \texttt{FB5903451D8A22637A60651ED87D8DD047A0B8E7CCD3D2C8745AB8D2DCCA4F31}.
\]

% =============================================================================
% END APPENDED SEVENTH-SERIES V20 MATERIAL
% =============================================================================

% =============================================================================
% BEGIN APPENDED SEVENTH-SERIES V21 MATERIAL
% =============================================================================

\chapter{The global Yang--Mills block frame and its infrared obstruction}
\label{ch:s7v21}

V20 makes the global first-chaos frame a mandatory gate.  This chapter
computes that gate for the linearized Yang--Mills curvature Gaussian.
The simplest translated one-link updates have exact all-chaos gap
proportional to \(L^{-2}\).  More generally, every fixed finite family
of bounded-range curvature-update templates has a frame mode of order
at most \(L^{-2}\).  Thus the five twenty-chord collar blocks can remain
valuable local estimates, but their fixed translates cannot by
themselves supply a volume-uniform heat-bath compiler.  The exponent is
two, not the scalar-constraint exponent six found in the SM audit.

\section{Periodic linearized curvature complex}
\label{sec:s7v21-complex}

Let \(\mathbb T_L^d=(\mathbb Z/L\mathbb Z)^d\), \(d\ge2\), and treat one
Lie-algebra colour.  A link one-form is
\(A=(A_\mu(x))\), and its plaquette curvature is
\begin{equation}
 (d_1A)_{\mu\nu}(x)
 =
 A_\nu(x+e_\mu)-A_\nu(x)
 -A_\mu(x+e_\nu)+A_\mu(x),
 \qquad\mu<\nu.                                         \label{eq:s7v21-d1}
\end{equation}
At momentum
\(k\in(2\pi/L)\mathbb Z^d\), put
\begin{equation}
 q_\mu(k)=e^{ik_\mu}-1,\qquad
 \lambda(k)=\sum_{\mu=1}^d|q_\mu(k)|^2
 =4\sum_{\mu=1}^d\sin^2(k_\mu/2).                       \label{eq:s7v21-lambda}
\end{equation}
Then
\begin{equation}
 \widehat{d_1A}_{\mu\nu}(k)
 =q_\mu(k)\widehat A_\nu(k)
  -q_\nu(k)\widehat A_\mu(k).                            \label{eq:s7v21-symbol}
\end{equation}

\begin{lemma}[Transverse curvature identity]
\label{lem:s7v21-transverse}
For \(k\ne0\), let
\[
 X_T(k)=\{a\in\mathbb C^d:q(k)^*a=0\}.
\]
Then
\begin{equation}
 d_1(k)^*d_1(k)=\lambda(k)I
 \quad\hbox{on }X_T(k).                                 \label{eq:s7v21-transverse}
\end{equation}
The nonzero spectrum of \(d_1(k)d_1(k)^*\) on
\operatorname{Ran}d_1(k)\) is also \(\lambda(k)\).
\end{lemma}

\begin{proof}
Expanding \eqref{eq:s7v21-symbol} gives
\[
 d_1(k)^*d_1(k)
 =\lambda(k)I-q(k)q(k)^*.
\]
The second term vanishes on \(X_T(k)\).  Nonzero spectra of
\(D^*D\) and \(DD^*\) agree, proving the last assertion.
\end{proof}

\begin{firewall}[Harmonic torons]
\label{fire:s7v21-torons}
At \(k=0\), \(d_1(k)=0\).  Constant link one-forms are harmonic toron
directions and are invisible to the curvature Gaussian.  They are not
included in the transverse nonzero-momentum estimates below and must be
retained as a separate compact/global sector, as already required in
V7.
\end{firewall}

\section{Exact one-link curvature frame}
\label{sec:s7v21-link}

Let \(e_{x,\mu}\) be the unit link one-form and
\[
 u_{x,\mu}:=d_1e_{x,\mu}.
\]
Every link belongs to two oriented plaquettes for each transverse
direction, so
\begin{equation}
 \|u_{x,\mu}\|^2=2(d-1).                                \label{eq:s7v21-link-norm}
\end{equation}
Let \(P_{x,\mu}\) be orthogonal projection in the curvature range onto
\(\operatorname{span}\{u_{x,\mu}\}\).

\begin{theorem}[Exact translated one-link frame]
\label{thm:s7v21-link-frame}
On \(\operatorname{Ran}d_1\),
\begin{equation}
 {\mathbb F}_{\rm link}
 :=\sum_{x,\mu}P_{x,\mu}
 ={1\over2(d-1)}d_1d_1^*.                              \label{eq:s7v21-link-frame}
\end{equation}
Its least nonzero eigenvalue on the periodic box is
\begin{equation}
 \delta_{\rm link}(L,d)
 ={4\sin^2(\pi/L)\over2(d-1)}
 \sim {2\pi^2\over(d-1)L^2}.                            \label{eq:s7v21-link-margin}
\end{equation}
\end{theorem}

\begin{proof}
For a rank-one vector \(u\), its orthogonal projection is
\(|u\rangle\langle u|/\|u\|^2\).  Therefore
\[
 \sum_{x,\mu}P_{x,\mu}
 ={1\over2(d-1)}
 \sum_{x,\mu}|d_1e_{x,\mu}\rangle
              \langle d_1e_{x,\mu}|
 ={1\over2(d-1)}d_1d_1^*.
\]
By \Cref{lem:s7v21-transverse}, its nonzero eigenvalue at momentum
\(k\) is \(\lambda(k)/(2(d-1))\).  The least nonzero
\(\lambda(k)\) is attained at
\((\pm2\pi/L,0,\ldots,0)\) and equals
\(4\sin^2(\pi/L)\).  The asymptotic follows from
\(\sin(\pi/L)\sim\pi/L\).
\end{proof}

\begin{theorem}[Exact all-chaos one-link heat-bath gap]
\label{thm:s7v21-link-gap}
For the centered curvature Gaussian on
\(\operatorname{Ran}d_1\), let \(K_{x,\mu}\) condition on the
orthogonal complement of the one-dimensional update direction
\operatorname{span}\{u_{x,\mu}\}\).  Then
\begin{equation}
 \operatorname{gap}
 \left(\sum_{x,\mu}(I-K_{x,\mu})\right)
 =\delta_{\rm link}(L,d).                               \label{eq:s7v21-link-gap}
\end{equation}
\end{theorem}

\begin{proof}
The first-chaos difference projections are precisely
\(P_{x,\mu}\), so their frame minimum is
\eqref{eq:s7v21-link-margin}.  The all-chaos projection inequality
\eqref{eq:s7v16-tensor-projection} shows that every higher Wiener chaos
has at least the same lower bound.  The first chaos attains it.
\end{proof}

\begin{assessment}[Meaning of the exact \(L^{-2}\) gap]
\label{ass:s7v21-link}
This is a gap of an auxiliary local heat-bath dynamics, not the
Yang--Mills Hamiltonian mass gap.  Its collapse does not disprove the
latter.  It proves that this particular local conditional-variance
compiler cannot pass to infinite volume with a uniform constant.
\end{assessment}

\section{Universal obstruction for fixed local curvature blocks}
\label{sec:s7v21-local}

Let \({\cal V}_1,\ldots,{\cal V}_R\) be a fixed finite family of
finite-dimensional link one-form subspaces on \(\mathbb Z^d\), each
supported in a ball of fixed radius.  Let
\[
 {\cal U}_a=d_1{\cal V}_a
\]
and choose an orthonormal basis
\(\{u_{a,r}\}_{r=1}^{m_a}\) of its nonzero curvature image.  Translate
these shapes over \(\mathbb T_L^d\), with a fixed number of orientations
and fixed update intensities \(c_a>0\).  Their first-chaos frame is
\begin{equation}
 {\mathbb F}_{\rm loc}
 =\sum_{a=1}^R c_a\sum_zP_{\tau_zu_{a,r}},
                                                               \label{eq:s7v21-local-frame}
\end{equation}
where the basis index \(r\) is also summed.

\begin{lemma}[Every local curvature shape vanishes at zero momentum]
\label{lem:s7v21-zero}
For each fixed template,
\begin{equation}
 \widehat u_{a,r}(k)=d_1(k)\widehat v_{a,r}(k)
 \quad\hbox{for some finitely supported }v_{a,r},
                                                               \label{eq:s7v21-factor}
\end{equation}
and there is a constant \(C_{a,r}\), independent of \(L\), such that
\begin{equation}
 \|\widehat u_{a,r}(k)\|\le C_{a,r}|k|
 \quad\hbox{for }|k|\le1.                               \label{eq:s7v21-zero}
\end{equation}
\end{lemma}

\begin{proof}
Membership in \(d_1{\cal V}_a\) gives a finitely supported preimage
\(v_{a,r}\).  Its Fourier polynomial is uniformly bounded.  From
\(|e^{ik_\mu}-1|\le|k_\mu|\), the operator norm of
\(d_1(k)\) is at most a dimension-dependent multiple of \(|k|\).
Multiply the two bounds.
\end{proof}

\begin{theorem}[Fixed-range Yang--Mills frame obstruction]
\label{thm:s7v21-obstruction}
For every family \eqref{eq:s7v21-local-frame}, there is
\(C_{\cal V}<\infty\), independent of \(L\), such that
\begin{equation}
 \lambda_{\min}
 \left(
 {\mathbb F}_{\rm loc}
 \big|_{\operatorname{Ran}d_1}
 \right)
 \le {C_{\cal V}\over L^2}.                              \label{eq:s7v21-obstruction}
\end{equation}
Consequently no fixed finite family of bounded-range curvature-image
blocks with bounded translation intensities has a volume-uniform
Gaussian heat-bath gap.
\end{theorem}

\begin{proof}
With unitary Fourier normalization, summing all translations of a
rank-one template gives at momentum \(k\) the outer product of its
unnormalized Fourier polynomial with itself.  Therefore the frame
symbol is
\[
 {\mathbb F}_{\rm loc}(k)
 =\sum_{a,r}c_a\,
   \widehat u_{a,r}(k)\widehat u_{a,r}(k)^*.
\]
By \Cref{lem:s7v21-zero},
\[
 \|{\mathbb F}_{\rm loc}(k)\|
 \le\left(\sum_{a,r}c_aC_{a,r}^2\right)|k|^2.
\]
At a least nonzero momentum \(k_L=(2\pi/L,0,\ldots,0)\), the curvature
range is nontrivial.  Its least frame eigenvalue is no greater than its
operator norm, which is at most \(C_{\cal V}L^{-2}\).
The first Wiener chaos is invariant, so the full Gaussian heat-bath gap
is no larger than this first-chaos value.
\end{proof}

\begin{corollary}[Application to translated V15 collar blocks]
\label{cor:s7v21-collar}
Any cofinal dynamics formed from finitely many orientations of the
fixed five V15 twenty-chord block shapes, translated with uniformly
bounded per-translate intensity, has a selected linearized curvature
frame margin at most \(C L^{-2}\).
\end{corollary}

\begin{proof}
Each V15 block has fixed finite support, and its linearized curvature
image is of the form \(d_1{\cal V}_a\).  Apply
\Cref{thm:s7v21-obstruction}.
\end{proof}

\begin{warning}[Local \(1/952\) and global \(CL^{-2}\) are compatible]
\label{warn:s7v21-compatible}
V16 conditions a single sixty-chord collar and proves that five
overlapping subblocks frame that fixed collar.  V21 translates fixed
collars through a growing periodic lattice and tests the long
wavelength curvature modes.  The local lower bound and global upper
bound refer to different frame operators; there is no contradiction.
\end{warning}

\section{The ideal multiscale target}
\label{sec:s7v21-multiscale}

Let
\[
 \operatorname{Ran}d_1
 =\bigoplus_{j=1}^J{\cal S}_j
\]
be any orthogonal decomposition into nonzero Fourier shells, and let
\(K_j^{\rm shell}\) condition on all curvature coefficients outside
\({\cal S}_j\).

\begin{theorem}[Unit Yang--Mills Fourier-shell frame]
\label{thm:s7v21-shell}
For the linearized curvature Gaussian,
\begin{equation}
 \sum_{j=1}^J(I-K_j^{\rm shell})
\quad\hbox{has all-chaos gap }1.                         \label{eq:s7v21-shell}
\end{equation}
\end{theorem}

\begin{proof}
On the first chaos, the difference projections are the mutually
orthogonal shell projections and sum to the identity on
\(\operatorname{Ran}d_1\).  The tensor projection argument of V16 gives
the same lower bound on every higher chaos, and the first chaos attains
one.
\end{proof}

\begin{firewall}[Fourier shells are not a local construction]
\label{fire:s7v21-shell}
The shell update is spatially nonlocal and ignores the compact
nonabelian group geometry, large-field sectors, reflection interfaces,
and the intended physical clock.  It is an ideal frame proving that the
infrared obstruction is caused by fixed update range, not by a
curvature zero mode at \(k\ne0\).
\end{firewall}

\begin{theorem}[Necessary change for a uniform translated frame]
\label{thm:s7v21-necessary}
A volume-uniform linearized curvature frame cannot be obtained from a
fixed finite family of bounded-support templates with bounded
intensities.  At least one of the following must occur:
\begin{enumerate}[label=\textup{(\roman*)},leftmargin=3.5em]
\item update support or first moment grows with scale;
\item low-frequency update intensity grows with scale; or
\item nonlocal or auxiliary coarse variables directly update the
      long modes.
\end{enumerate}
\end{theorem}

\begin{proof}
If none occurs, the hypotheses and constant in
\Cref{thm:s7v21-obstruction} remain uniform, forcing the margin to zero.
\end{proof}

\begin{definition}[Multiscale Yang--Mills curvature frame card, MYCFC]
\label{def:s7v21-MYCFC}
MYCFC consists of:
\begin{enumerate}[label=\textup{(Y\arabic*)},leftmargin=4em]
\item dyadic or wavelet link one-form templates whose curvature images
      cover every nonzero transverse momentum;
\item an exact first-chaos frame multiplier and a positive
      volume-uniform lower bound;
\item all-chaos lifting as in V16;
\item total update intensity and physical-clock normalization;
\item separate treatment of gauge, harmonic toron, electric, and
      boundary sectors;
\item insertion of the V15--V16 collar estimates at the finest scale;
\item nonlinear physical projection transport and bad-set control at
      every scale; and
\item continuum, reflection, observable, and Hamiltonian spectral
      contacts.
\end{enumerate}
\end{definition}

\begin{assessment}[Progress at seventh-series V21]
\label{ass:s7v21-status}
The global local-block question is decided at linear order.  A
translated one-link curvature heat bath has exact gap
\(4\sin^2(\pi/L)/(2(d-1))\), and every fixed bounded-range curvature
block family has a mode at most \(CL^{-2}\).  Hence fixed collar
translation cannot yield the uniform constant needed by the current
compiler.  Fourier shells restore gap one and identify multiscale range
as the missing ingredient.
\end{assessment}

\begin{decision}[V21 bottleneck decision]
\label{dec:s7v21-bottleneck}
V22 will construct real-space dyadic link templates on the transverse
sector and compute their curvature-frame multiplier and total
intensity.  The desired result is a volume-uniform frame with a clock
cost no worse than the one-derivative scaling predicts.  Gauge
projection and torons will remain explicit; no scalar-constraint
multiplier will be imported from the SM calculation.
\end{decision}

\begin{firewall}[Claim boundary after seventh-series V21]
\label{fire:s7v21-boundary}
V21 proves the exact one-link curvature frame and all-chaos gap, the
universal \(L^{-2}\) obstruction for fixed bounded-range curvature
blocks, and the ideal Fourier-shell gap.

V21 does not construct a local multiscale Yang--Mills frame, reconcile
its physical clock, transport it to the nonlinear measure, control
torons and bad sectors, close the martingale/continuum interfaces, or
prove the Yang--Mills mass gap.
\end{firewall}

\paragraph{Parent-version record.}
V21 was formed from
\texttt{Renewal\_Yang\_Mills\_Mass\_Gap\_Research\_Notes\_V20.tex}.
The V20 parent had \(8\,122\) logical source lines, \(295\,605\) bytes,
and SHA--256
\[
 \texttt{9D3BB208C541A0638E7D1340F58A42F317B0D50A4A6E6FD8F1300F2133979EEA}.
\]

% =============================================================================
% END APPENDED SEVENTH-SERIES V21 MATERIAL
% =============================================================================

% =============================================================================
% BEGIN APPENDED SEVENTH-SERIES V22 MATERIAL
% =============================================================================

\chapter{A bounded-clock dyadic curvature frame}
\label{ch:s7v22}

V21 proves that fixed-range curvature updates have an \(L^{-2}\) slow
mode.  The one-derivative structure also suggests its repair.  A link
field constant on a cube of side \(s\) has curvature only on the cube
faces: its curvature norm grows like \(s^{(d-1)/2}\), while its
low-frequency Fourier amplitude grows like \(s^d\).  Translating these
templates over dyadic scales and assigning rate \(s^{-(d-1)}\) gives a
uniform first-chaos curvature frame.  Unlike the scalar-constraint
construction audited at V20, the total rate per lattice translate stays
bounded for every \(d>1\).

\section{Dyadic cube link templates}
\label{sec:s7v22-templates}

Assume \(L\ge8\) is a power of two and let
\begin{equation}
 {\cal S}_L=\{1,2,4,\ldots,L/4\}.                        \label{eq:s7v22-scales}
\end{equation}
For \(s\in{\cal S}_L\), let
\[
 b_s(x)=\mathbf1_{\{0,\ldots,s-1\}^d}(x)
\]
on \(\mathbb T_L^d\).  For a translation \(z\) and link orientation
\(\mu\), define the link one-form
\begin{equation}
 v_{s,z,\mu}(x,\nu)
 :=b_s(x-z)\mathbf1_{\{\nu=\mu\}},\qquad
 u_{s,z,\mu}:=d_1v_{s,z,\mu}.                           \label{eq:s7v22-vu}
\end{equation}

\begin{lemma}[Exact face norm]
\label{lem:s7v22-face}
For \(s\le L/4\),
\begin{equation}
 \|u_{s,z,\mu}\|_2^2=2(d-1)s^{d-1}.                     \label{eq:s7v22-face}
\end{equation}
\end{lemma}

\begin{proof}
Only the \(\mu\nu\) curvature components with \(\nu\ne\mu\) are
nonzero.  For each such \(\nu\), the discrete difference of \(b_s\) in
the \(\nu\)-direction is supported on the two opposite faces normal to
\(\nu\).  Each face has \(s^{d-1}\) sites and the values are \(\pm1\).
The faces do not overlap because \(s\le L/4\).  Summing their squared
values over the \(d-1\) transverse orientations gives
\eqref{eq:s7v22-face}.
\end{proof}

The Fourier polynomial of the cube is
\begin{equation}
 B_s(k)
 :=\sum_xb_s(x)e^{-ik\cdot x}
 =\prod_{\ell=1}^d
   {1-e^{-is k_\ell}\over1-e^{-ik_\ell}},               \label{eq:s7v22-B}
\end{equation}
with the quotient interpreted as \(s\) at \(k_\ell=0\).

\begin{lemma}[Small-frequency cube lower bound]
\label{lem:s7v22-B-lower}
If \(s\|k\|_\infty\le\pi/2\), then
\begin{equation}
 |B_s(k)|\ge\left({2s\over\pi}\right)^d.                 \label{eq:s7v22-B-lower}
\end{equation}
\end{lemma}

\begin{proof}
For \(|s\theta|/2\le\pi/4\),
\[
 \left|\sum_{r=0}^{s-1}e^{-ir\theta}\right|
 ={|\sin(s\theta/2)|\over|\sin(\theta/2)|}.
\]
Using \(\sin x\ge2x/\pi\) on \(0\le x\le\pi/2\) in the numerator and
\(\sin x\le x\) in the denominator gives at least \(2s/\pi\).
The value at \(\theta=0\) is \(s\), which obeys the same lower bound.
Multiply over coordinates.
\end{proof}

\section{Exact dyadic frame multiplier}
\label{sec:s7v22-frame}

Let \(P_{s,z,\mu}\) be orthogonal projection in
\(\operatorname{Ran}d_1\) onto
\(\operatorname{span}\{u_{s,z,\mu}\}\).  Give scale \(s\) the update
intensity
\begin{equation}
 a_s:=s^{-(d-1)}                                        \label{eq:s7v22-rate}
\end{equation}
and define
\begin{equation}
 {\mathbb F}_{\rm dyad}
 :=
 \sum_{s\in{\cal S}_L}a_s
 \sum_{z\in\mathbb T_L^d}\sum_{\mu=1}^d
 P_{s,z,\mu}.                                            \label{eq:s7v22-frame}
\end{equation}

\begin{theorem}[Exact Fourier multiplier]
\label{thm:s7v22-multiplier}
On the curvature range at nonzero momentum \(k\),
\({\mathbb F}_{\rm dyad}\) is scalar multiplication by
\begin{equation}
 f_{\rm dyad}(k)
 =
 {\lambda(k)\over2(d-1)}
 \sum_{s\in{\cal S}_L}
 {|B_s(k)|^2\over s^{2(d-1)}}.                          \label{eq:s7v22-multiplier}
\end{equation}
\end{theorem}

\begin{proof}
For fixed \(s,\mu\),
\[
 \widehat u_{s,0,\mu}(k)
 =B_s(k)d_1(k)e_\mu.
\]
Summing all translations of a normalized rank-one shape gives its
Fourier outer product.  By \Cref{lem:s7v22-face}, the contribution at
scale \(s\), after summing orientations and inserting \(a_s\), is
\[
 {a_s|B_s(k)|^2\over2(d-1)s^{d-1}}
 d_1(k)d_1(k)^*
 =
 {|B_s(k)|^2\over2(d-1)s^{2(d-1)}}
 d_1(k)d_1(k)^*.
\]
On \(\operatorname{Ran}d_1(k)\),
\(d_1(k)d_1(k)^*=\lambda(k)I\) by
\Cref{lem:s7v21-transverse}.  Sum over scales.
\end{proof}

\section{Uniform frame floor}
\label{sec:s7v22-floor}

\begin{theorem}[Bounded-clock dyadic Yang--Mills frame]
\label{thm:s7v22-floor}
For every nonzero lattice momentum,
\begin{equation}
 f_{\rm dyad}(k)\ge\delta_d,                             \label{eq:s7v22-floor}
\end{equation}
where
\begin{equation}
 \delta_d
 :=
 \min\left\{
 {2\sin^2(\pi/8)\over d-1},\,
 {2^{\,2d-3}\over(d-1)\pi^{2d}}
 \right\}>0.                                             \label{eq:s7v22-delta}
\end{equation}
For \(d=4\), the second term is the smaller one and
\begin{equation}
 \delta_4={32\over3\pi^8}
 \approx0.00112416417764.                               \label{eq:s7v22-delta4}
\end{equation}
\end{theorem}

\begin{proof}
Choose representatives with \(k_\ell\in[-\pi,\pi]\) and put
\(r=\|k\|_\infty\).

If \(r\ge\pi/4\), use the \(s=1\) term.  Since \(B_1=1\) and at least
one component has magnitude \(r\),
\[
 \lambda(k)\ge4\sin^2(r/2)\ge4\sin^2(\pi/8).
\]
Equation \eqref{eq:s7v22-multiplier} gives the first constant in
\eqref{eq:s7v22-delta}.

If \(0<r<\pi/4\), choose the largest dyadic \(s\) satisfying
\(sr\le\pi/2\).  Then \(sr>\pi/4\).  The least nonzero lattice
frequency is \(2\pi/L\), so this \(s\) belongs to
\({\cal S}_L\).  By \Cref{lem:s7v22-B-lower},
\[
 |B_s(k)|^2\ge(2s/\pi)^{2d}.
\]
Also, for \(r/2\le\pi/2\),
\[
 \lambda(k)\ge4\sin^2(r/2)\ge {4r^2\over\pi^2}.
\]
The single chosen scale in \eqref{eq:s7v22-multiplier} is therefore at
least
\begin{align*}
 {1\over2(d-1)}
 {4r^2\over\pi^2}
 {(2s/\pi)^{2d}\over s^{2d-2}}
 &=
 {2^{2d+1}\over(d-1)\pi^{2d+2}}s^2r^2\\
 &>
 {2^{2d-3}\over(d-1)\pi^{2d}},
\end{align*}
using \(sr>\pi/4\).  This is the second constant.
Substitution of \(d=4\) gives \eqref{eq:s7v22-delta4}.
\end{proof}

\begin{theorem}[All-chaos weighted heat-bath gap]
\label{thm:s7v22-all-chaos}
For the linearized curvature Gaussian, let \(K_{s,z,\mu}\) condition on
the orthogonal complement of
\(\operatorname{span}\{u_{s,z,\mu}\}\).  Then the weighted generator
\begin{equation}
 {\cal A}_{\rm dyad}
 =
 \sum_{s\in{\cal S}_L}a_s
 \sum_{z,\mu}(I-K_{s,z,\mu})                             \label{eq:s7v22-A}
\end{equation}
has
\begin{equation}
 \operatorname{gap}({\cal A}_{\rm dyad})
 =\min_{k\ne0}f_{\rm dyad}(k)
 \ge\delta_d.                                            \label{eq:s7v22-gap}
\end{equation}
\end{theorem}

\begin{proof}
The first-chaos frame is
\eqref{eq:s7v22-frame} with multiplier
\eqref{eq:s7v22-multiplier}.  Apply the weighted version of
\eqref{eq:s7v16-tensor-projection}: multiply each projection inequality
by \(a_s\) before summing.  Every higher chaos has at least the
first-chaos frame floor, and the first chaos attains its minimum.
\end{proof}

\section{The physical-clock cost}
\label{sec:s7v22-clock}

\begin{theorem}[Uniform total scale intensity]
\label{thm:s7v22-clock}
The total scale intensity per translation and link orientation obeys
\begin{equation}
 W_{L,d}:=\sum_{s\in{\cal S}_L}a_s
 \le {1\over1-2^{-(d-1)}}.                              \label{eq:s7v22-clock}
\end{equation}
Including all \(d\) link orientations, the per-site intensity is at
most
\[
 {d\over1-2^{-(d-1)}}.
\]
For \(d=4\), this is \(32/7\approx4.57143\), independently of \(L\).
\end{theorem}

\begin{proof}
The dyadic sum is a truncated geometric series:
\[
 \sum_{j\ge0}2^{-j(d-1)}
 ={1\over1-2^{-(d-1)}}.
\]
\end{proof}

\begin{assessment}[Why the clock differs from the SM scalar constraint]
\label{ass:s7v22-clock}
The Yang--Mills curvature symbol has one derivative.  A cube link
average has a codimension-one curvature boundary and does not need a
charge-preserving difference.  The rate \(s^{-(d-1)}\) is summable.
The SM scalar constraint has two derivatives and a charge condition;
its audited dyadic construction required growing total intensity.
Thus the SM obstruction prompted the correct test but its clock cost
does not transfer.
\end{assessment}

\section{Gauge and nonabelian qualification}
\label{sec:s7v22-qualification}

\begin{firewall}[The theorem is on the linearized curvature range]
\label{fire:s7v22-linear}
The construction works on \(\operatorname{Ran}d_1\) and repeats
independently for each Lie-algebra colour.  It does not update harmonic
torons, prove a gauge-fixed link-space gap, or define a nonlinear
simultaneous group update.  The cube link template has a harmonic
component before applying \(d_1\); only its curvature image enters the
Gaussian theorem.
\end{firewall}

\begin{firewall}[Coarse blocks grow in physical diameter]
\label{fire:s7v22-range}
Although total update intensity is bounded, the largest cube has side
\(L/4\).  A nonlinear implementation is genuinely multiscale and
long-range.  Its ownership, reflection-plane crossings, boundary
conditions, and compatibility with a local physical Hamiltonian must
be proved.  Bounded clock does not mean bounded spatial range.
\end{firewall}

\begin{definition}[Nonlinear dyadic curvature update card, NDCUC]
\label{def:s7v22-NDCUC}
NDCUC consists of:
\begin{enumerate}[label=\textup{(D\arabic*)},leftmargin=4em]
\item a compact-group realization of the collective cube link
      direction \(v_{s,z,\mu}\);
\item a block conditional measure and internal gap at every scale;
\item nonlinear comparison of its curvature image with
      \(u_{s,z,\mu}\) on a declared good set;
\item a scale-uniform conditional score and moving-ground ledger;
\item bad-set and interface estimates uniform over all scales;
\item ownership and row/column sums using the summable rates
      \(s^{-(d-1)}\);
\item explicit toron, gauge, reflection, and physical-clock treatment;
      and
\item insertion into PBDC, MGCC, MLC, continuum construction, and the
      Hamiltonian spectral argument.
\end{enumerate}
\end{definition}

\begin{theorem}[What V22 resolves]
\label{thm:s7v22-result}
Rows corresponding only to the linearized frame and total clock in
NDCUC are now proved with
\[
 \operatorname{gap}({\cal A}_{\rm dyad})\ge\delta_d,
 \qquad
 \sum_sa_s\le(1-2^{-(d-1)})^{-1}.
\]
Thus the V21 fixed-range infrared obstruction has an explicit
bounded-clock multiscale repair at Gaussian curvature level.
\end{theorem}

\begin{proof}
Use
\Cref{thm:s7v22-floor,thm:s7v22-all-chaos,thm:s7v22-clock}.
\end{proof}

\begin{assessment}[Progress at seventh-series V22]
\label{ass:s7v22-status}
The global linearized frame bottleneck is closed without a divergent
per-site clock.  Translated dyadic cube link fields, weighted by
\(s^{-(d-1)}\), give an explicit uniform all-chaos curvature-Gaussian
gap; in four dimensions the certified floor is
\(32/(3\pi^8)\).  The remaining difficulty is nonlinear realization of
these growing collective blocks and their physical conditional scores.
\end{assessment}

\begin{decision}[V22 bottleneck decision]
\label{dec:s7v22-bottleneck}
V23 will define the compact-group collective cube update by simultaneous
left multiplication of the selected links, compute the exact
first and second derivatives of every affected plaquette, and separate
interior cancellations from the cube-face curvature.  This will test
whether the nonlinear score grows with volume \(s^d\) or only with the
boundary \(s^{d-1}\), which is decisive for summability with the rates
\eqref{eq:s7v22-rate}.
\end{decision}

\begin{firewall}[Claim boundary after seventh-series V22]
\label{fire:s7v22-boundary}
V22 constructs a real-space dyadic curvature frame, proves its exact
Fourier multiplier, a uniform all-chaos Gaussian gap, and bounded total
per-site update intensity.

V22 does not construct the nonlinear compact-group updates, control
their conditional measures or scores, treat torons and reflection
interfaces, close the local ground/martingale budgets, construct the
continuum, or prove the Yang--Mills mass gap.
\end{firewall}

\paragraph{Parent-version record.}
V22 was formed from
\texttt{Renewal\_Yang\_Mills\_Mass\_Gap\_Research\_Notes\_V21.tex}.
The V21 parent had \(8\,505\) logical source lines, \(309\,759\) bytes,
and SHA--256
\[
 \texttt{DC1404D21C88CB7939CBCE33F663AEA1A12EAD8736A5C2AB094A7C4611294FDC}.
\]

% =============================================================================
% END APPENDED SEVENTH-SERIES V22 MATERIAL
% =============================================================================

\newpage

% =============================================================================
% BEGIN APPENDED SEVENTH-SERIES V23 MATERIAL
% =============================================================================

\chapter{Nonabelian cube updates: the bulk defect, a cut-gauge repair,
and the true local clock}
\label{ch:s7v23}

V22 solves the infrared frame problem for the linearized curvature
Gaussian by using cube one-forms on all dyadic scales.  This note tests
the first nonlinear implementation suggested there.  The test has
three outcomes.
\begin{enumerate}[label=\textup{(\roman*)},leftmargin=3.5em]
\item Simultaneous left multiplication of the selected parallel links
      does \emph{not} retain the codimension-one support of its
      linearization: a noncommutative bulk term occurs on every
      interior plaquette.
\item Replacing that update by a compact cut-gauge block action restores
      exact boundary localization at every field, and its flat tangent
      space contains the V22 cube one-form exactly.
\item The rate in V22 has bounded event density but linear-in-\(L\)
      link-participation load.  Thus the claim of a bounded physical
      clock was too strong: the present comparison with the local
      electric form still loses a factor of order \(L\).
\end{enumerate}
The second point is a genuine nonlinear repair; the third identifies
the next upstream computation rather than hiding a clock
renormalization.

\section{The naive parallel-link update}
\label{sec:s7v23-naive}

Let \(G\subset U(n)\) be a compact matrix group with nonabelian
semisimple Lie algebra \(\mathfrak g\), and use the Hilbert--Schmidt
norm.  For a positively oriented plaquette \(p=(x;\mu,\nu)\), put
\begin{equation}
 A=U_\mu(x),\quad B=U_\nu(x+\widehat\mu),\quad
 C=U_\mu(x+\widehat\nu),\quad D=U_\nu(x),\quad
 P=ABC^{-1},\quad W=PD^{-1}.                            \label{eq:s7v23-ABCD}
\end{equation}
For the cube indicator \(b_s\) of V22, write
\[
 \alpha=b_s(x-z),\qquad
 \gamma=b_s(x+\widehat\nu-z).
\]
The proposed update with generator \(X\in\mathfrak g\) is
\[
 U_\mu(y;t)=e^{t b_s(y-z)X}U_\mu(y),
 \qquad U_\rho(y;t)=U_\rho(y)\quad(\rho\ne\mu).
\]

\begin{lemma}[Exact first and second plaquette derivatives]
\label{lem:s7v23-derivatives}
Under the naive update,
\begin{align}
 W(t)&=e^{t\alpha X}P e^{-t\gamma X}D^{-1},
                                                               \label{eq:s7v23-Wt}\\
 \dot W(0)&=\alpha XW-\gamma PXD^{-1},                         \label{eq:s7v23-Wdot}\\
 \ddot W(0)&=\alpha^2X^2W-2\alpha\gamma XPXD^{-1}
                 +\gamma^2PX^2D^{-1}.                         \label{eq:s7v23-Wddot}
\end{align}
For the normalized Wilson plaquette potential
\[
 q(W)=1-\frac1n\operatorname{Re}\operatorname{Tr}W,
\]
its first and second derivatives are
\(-n^{-1}\operatorname{Re}\operatorname{Tr}\dot W(0)\) and
\(-n^{-1}\operatorname{Re}\operatorname{Tr}\ddot W(0)\).
\end{lemma}

\begin{proof}
Only \(A\) and \(C\) can be selected.  Left multiplication of \(C\)
becomes right multiplication of \(C^{-1}\) by \(e^{-t\gamma X}\),
which gives \eqref{eq:s7v23-Wt}.  Two differentiations give
\eqref{eq:s7v23-Wdot}--\eqref{eq:s7v23-Wddot}; differentiating the
definition of \(q\) finishes the proof.
\end{proof}

\begin{lemma}[The interior term is a curvature commutator]
\label{lem:s7v23-interior}
If both parallel links are selected, so that
\(\alpha=\gamma=1\), then
\begin{equation}
 {d\over dt}\operatorname{Re}\operatorname{Tr}W(t)\big|_{t=0}
 =
 \operatorname{Re}\operatorname{Tr}
 X\bigl(W-D^{-1}WD\bigr).                              \label{eq:s7v23-commutator}
\end{equation}
Consequently
\begin{equation}
 \left|
 {d\over dt}q(W(t))\big|_{t=0}
 \right|
 \le {2\over n}\|X\|_{\rm HS}\|W-I\|_{\rm HS}.          \label{eq:s7v23-interior-bound}
\end{equation}
The right side vanishes at the flat field, but it is not a boundary
term at a general noncommuting field.
\end{lemma}

\begin{proof}
Cyclicity of the trace and \(P=WD\) turn the second term of
\eqref{eq:s7v23-Wdot} into
\(\operatorname{Tr}(XD^{-1}WD)\).  This proves
\eqref{eq:s7v23-commutator}.  Since unitary conjugation is an isometry,
\[
 \|W-D^{-1}WD\|_{\rm HS}
 \le \|W-I\|_{\rm HS}+\|D^{-1}(W-I)D\|_{\rm HS}
 =2\|W-I\|_{\rm HS}.
\]
Cauchy--Schwarz for the trace gives
\eqref{eq:s7v23-interior-bound}.
\end{proof}

\begin{theorem}[Genuine volume-order obstruction for the naive update]
\label{thm:s7v23-volume}
There are constant link configurations, a generator
\(X\in\mathfrak g\), and constants \(c,C>0\), independent of \(s\),
for which the derivative of the Wilson potential under the naive cube
update satisfies
\begin{equation}
 {d\over dt}S_{\rm W}(U(t))\big|_{t=0}
 \ge c(s-1)s^{d-1}-Cs^{d-1}.                           \label{eq:s7v23-volume}
\end{equation}
Thus the nonlinear score of this update can be of order \(s^d\), even
though its derivative at the flat field has curvature on only the
cube faces.
\end{theorem}

\begin{proof}
Choose \(R,S\in\mathfrak g\) with
\([S,[R,S]]\ne0\), which is possible in every nonabelian compact
semisimple Lie algebra.  Put
\[
 A=e^{\varepsilon R},\qquad B=e^{\varepsilon S}
\]
for sufficiently small fixed \(\varepsilon>0\).  The group commutator
\(W=ABA^{-1}B^{-1}\) obeys
\[
 W-B^{-1}WB
 =\varepsilon^3[S,[R,S]]+O(\varepsilon^4).
\]
Taking \(X=[S,[R,S]]\), and reversing its sign if necessary, makes the
derivative of \(q\) in \eqref{eq:s7v23-commutator} a fixed positive
number \(c\) for all sufficiently small fixed \(\varepsilon\).

Now set every \(\mu\)-link equal to \(A\), every \(\nu\)-link equal to
\(B\), and all remaining oriented links equal to the identity.  There
are exactly
\[
 (s-1)s^{d-1}
\]
canonical \(\mu\nu\)-plaquettes whose two \(\mu\)-links lie in the
selected cube.  They all give the same positive interior derivative.
Every other affected plaquette lies in a one-plaquette-thick boundary
of the cube; its number is at most \(C_1s^{d-1}\), and compactness of
\(G\) bounds each of its derivatives by a constant \(C_2\).
The remaining interior orientations have trivial plaquette and zero
derivative.  Combining the interior contribution and the boundary
bound proves \eqref{eq:s7v23-volume}.
\end{proof}

\begin{corollary}[The V22 nonlinear proposal must be changed]
\label{cor:s7v23-naive-no}
No estimate depending only on the number of cube-face plaquettes can
uniformly bound the naive nonlinear score.  With
\(a_s=s^{-(d-1)}\), the worst allowed bulk contribution is not summable
over the dyadic scales.
\end{corollary}

\begin{proof}
The lower bound \eqref{eq:s7v23-volume} divided by \(s^{d-1}\) grows
linearly in \(s\).
\end{proof}

\section{An exact compact cut-gauge block}
\label{sec:s7v23-cut}

The obstruction in \Cref{thm:s7v23-volume} comes from asking a
nonabelian parallel translation to behave like an abelian gradient.
There is a compact group action which realizes the required tangent
without making that request.

For \(s\le L/4\), a translation \(z\), and an orientation \(\mu\),
define the anisotropic vertex box
\begin{equation}
 Q_{s,z}^{\mu}
 :=
 \left\{x:
 0\le x_\mu-z_\mu\le s,\quad
 0\le x_\rho-z_\rho\le s-1\ \ (\rho\ne\mu)
 \right\}.                                               \label{eq:s7v23-Q}
\end{equation}
Let \(E(Q)\) be the positive links whose two endpoints lie in \(Q\).
The compact block group
\[
 {\cal G}_Q:=G^Q
\]
acts on link configurations by
\begin{equation}
 (T_gU)_{xy}
 :=
 \begin{cases}
 g_xU_{xy}g_y^{-1},& (x,y)\in E(Q),\\
 U_{xy},& (x,y)\notin E(Q).
 \end{cases}                                             \label{eq:s7v23-cut-action}
\end{equation}
Boundary-crossing links are deliberately frozen.

\begin{lemma}[Compact action and exact boundary support]
\label{lem:s7v23-boundary-action}
Equation \eqref{eq:s7v23-cut-action} defines a measure-preserving
action of the compact group \({\cal G}_Q\) on product Haar link
measure.  A Wilson plaquette wholly contained in \(Q\) changes only by
conjugation, and a plaquette having no internal edge is unchanged.
Therefore the Wilson action varies only on the mixed plaquette set
\begin{equation}
 \partial_2Q
 :=
 \{p:0<|E(p)\cap E(Q)|<4\}.                              \label{eq:s7v23-partial2}
\end{equation}
For \(s\ge2\),
\begin{equation}
 |\partial_2Q|
 \le 8d^2(d-1)(s+3)^{d-1}.                              \label{eq:s7v23-boundary-count}
\end{equation}
\end{lemma}

\begin{proof}
Composition in \({\cal G}_Q\) gives composition of the transformations,
and left and right multiplication preserve Haar measure on every
internal link.  On a plaquette with all vertices in \(Q\), the endpoint
factors cancel around the loop, leaving \(g_xW_pg_x^{-1}\).  The
normalized real trace is unchanged.  The assertion for plaquettes
without internal edges is immediate.

Every mixed plaquette meets the one-step vertex boundary of \(Q\).
The number of vertices in that layer is at most
\(4d(s+3)^{d-1}\), by expanding and shrinking the rectangular box and
telescoping one coordinate at a time.  A vertex belongs to
\(2d(d-1)\) unoriented plaquettes.  Counting incidences, with
overcounting allowed, gives \eqref{eq:s7v23-boundary-count}.
\end{proof}

\begin{theorem}[The V22 cube direction is an exact flat tangent]
\label{thm:s7v23-tangent}
Let \(X\in\mathfrak g\) and, for \(x\in Q_{s,z}^{\mu}\), set
\begin{equation}
 \phi_x:=-(x_\mu-z_\mu)X,\qquad g_x(t):=e^{t\phi_x}.
                                                               \label{eq:s7v23-affine}
\end{equation}
At the flat configuration, the tangent of
\(T_{g(t)}\) is exactly
\begin{equation}
 \dot U_\rho(x;0)
 =b_s(x-z)\mathbf1_{\{\rho=\mu\}}X
 =v_{s,z,\mu}(x,\rho)X.                                 \label{eq:s7v23-exact-tangent}
\end{equation}
In particular its linearized curvature is exactly
\(u_{s,z,\mu}X\) from V22.
\end{theorem}

\begin{proof}
For an internal link \(e=(x,x+\widehat\rho)\), differentiation at the
flat field gives
\[
 \dot U_e(0)=\phi_x-\phi_{x+\widehat\rho}.
\]
This equals \(X\) if \(\rho=\mu\), and zero otherwise.  The internal
\(\mu\)-links have precisely the base points
\(x-z\in\{0,\ldots,s-1\}^d\).  Links outside \(E(Q)\) are frozen.
This proves \eqref{eq:s7v23-exact-tangent}; applying \(d_1\) gives the
last statement.
\end{proof}

\begin{theorem}[Gaussian frame domination by cut-gauge blocks]
\label{thm:s7v23-frame-domination}
At the flat field, let \({\cal T}_{s,z,\mu}\) be the link tangent
space of the action \({\cal G}_{Q_{s,z}^{\mu}}\), and let
\(\Pi_{s,z,\mu}\) be orthogonal projection onto
\(d_1{\cal T}_{s,z,\mu}\) in the curvature range.  Then
\begin{equation}
 \Pi_{s,z,\mu}\succeq P_{s,z,\mu},                       \label{eq:s7v23-proj-domination}
\end{equation}
where \(P_{s,z,\mu}\) is the rank-one V22 projection, independently in
each colour.  Hence
\begin{equation}
 \sum_{s,z,\mu}a_s\Pi_{s,z,\mu}
 \succeq{\mathbb F}_{\rm dyad}
 \succeq\delta_d I                                      \label{eq:s7v23-frame-domination}
\end{equation}
on \(\operatorname{Ran}d_1\).  The corresponding Gaussian
cut-gauge heat-bath generator has all-chaos gap at least \(\delta_d\).
\end{theorem}

\begin{proof}
\Cref{thm:s7v23-tangent} says
\(u_{s,z,\mu}X\in d_1{\cal T}_{s,z,\mu}\).  Orthogonal projection onto
a subspace dominates projection onto any line in that subspace, which
proves \eqref{eq:s7v23-proj-domination}.  Sum with the positive weights
\(a_s\), and invoke \Cref{thm:s7v22-floor}.  The tensor-projection
argument of V16 and V22 lifts the first-chaos lower bound to every
Gaussian chaos.
\end{proof}

\begin{theorem}[Exact nonlinear boundary-score reduction]
\label{thm:s7v23-score-support}
Along every cut-gauge orbit,
\begin{equation}
 S_{\rm W}(T_gU)
 =
 C(U)+\sum_{p\in\partial_2Q}q\bigl(W_p(T_gU)\bigr),       \label{eq:s7v23-score-support}
\end{equation}
where \(C(U)\) is independent of \(g\).  Thus the conditional density
of the block variable relative to Haar measure on \({\cal G}_Q\)
depends on at most \(C_d s^{d-1}\) plaquettes, not \(s^d\) plaquettes.
\end{theorem}

\begin{proof}
Split the plaquettes into those wholly internal, wholly unaffected,
and mixed.  The first two contributions are independent of \(g\) by
\Cref{lem:s7v23-boundary-action}; the last set obeys
\eqref{eq:s7v23-boundary-count}.  Product Haar invariance gives the
stated reference measure on the compact orbit, with the usual quotient
by its stabilizer if the action is not free.
\end{proof}

\begin{warning}[Boundary support does not imply a uniform amplitude]
\label{warn:s7v23-amplitude}
For the affine tangent \eqref{eq:s7v23-affine}, the left-trivialized
velocity on an internal link \(e=(x,y)\) is
\begin{equation}
 \xi_e=\phi_x-\operatorname{Ad}_{U_e}\phi_y.             \label{eq:s7v23-velocity}
\end{equation}
If \(\|U_e-I\|_{\rm op}\le\epsilon\), then
\begin{equation}
 \|\xi_e-(\phi_x-\phi_y)\|_{\rm HS}
 \le2s\epsilon\|X\|_{\rm HS}.                           \label{eq:s7v23-velocity-error}
\end{equation}
Thus exact support localization removes the bulk \emph{count}, but a
scale-growing parallel-transport error remains unless the good-field
control is covariant and scale adapted.
\end{warning}

\begin{proof}
Subtract the flat velocity from \eqref{eq:s7v23-velocity}:
\[
 \xi_e-(\phi_x-\phi_y)
 =(I-\operatorname{Ad}_{U_e})\phi_y.
\]
For unitary \(U_e\),
\(\|(I-\operatorname{Ad}_{U_e})Y\|_{\rm HS}
\le2\|U_e-I\|_{\rm op}\|Y\|_{\rm HS}\), while
\(\|\phi_y\|_{\rm HS}\le s\|X\|_{\rm HS}\).
\end{proof}

\section{Event rate versus local-electric participation}
\label{sec:s7v23-clock}

The V22 quantity \(d\sum_sa_s\) counts block events per lattice site.
A comparison with the local electric Hamiltonian instead counts how
many block moves touch a fixed link.

\begin{theorem}[Exact dyadic link-participation load]
\label{thm:s7v23-load}
Fix a positive link \(e\) of orientation \(\rho\).  At scale \(s\), the
number of translated, oriented boxes \(Q_{s,z}^{\mu}\) whose internal
edge set contains \(e\) is
\begin{equation}
 N_s
 =s^d+(d-1)(s^2-1)s^{d-2}
 =ds^d-(d-1)s^{d-2}.                                   \label{eq:s7v23-Ns}
\end{equation}
With \(a_s=s^{-(d-1)}\), its weighted participation load is
\begin{align}
 \Lambda_{L,d}
 &:=\sum_{s\in{\cal S}_L}a_sN_s                         \label{eq:s7v23-Lambda}\\
 &=\sum_{s\in{\cal S}_L}
       \left(ds-\frac{d-1}{s}\right)                    \label{eq:s7v23-Lambda-sum}\\
 &=d\left(\frac L2-1\right)
   -(d-1)\left(2-\frac4L\right).                        \label{eq:s7v23-Lambda-exact}
\end{align}
In four dimensions,
\begin{equation}
 \Lambda_{L,4}=2L-10+\frac{12}{L}.                      \label{eq:s7v23-Lambda4}
\end{equation}
Therefore the block-event density is bounded as in V22, but the
per-link participation is \(\Theta(L)\).
\end{theorem}

\begin{proof}
If the block orientation \(\mu\) equals \(\rho\), there are \(s\)
allowed translations in the link direction and \(s\) in every
transverse direction, giving \(s^d\).  If \(\mu\ne\rho\), there are
\(s+1\) translations in the \(\mu\)-direction, \(s-1\) in the
\(\rho\)-direction, and \(s\) in the remaining \(d-2\) directions.
This gives \((s^2-1)s^{d-2}\) for each of the \(d-1\) choices of
\(\mu\ne\rho\), proving \eqref{eq:s7v23-Ns}.

Division by \(s^{d-1}\) gives
\eqref{eq:s7v23-Lambda-sum}.  If \(L=2^J\), then
\[
 \sum_{s\in{\cal S}_L}s=\frac L2-1,\qquad
 \sum_{s\in{\cal S}_L}s^{-1}=2-\frac4L.
\]
Substitution proves \eqref{eq:s7v23-Lambda-exact} and
\eqref{eq:s7v23-Lambda4}.
\end{proof}

\begin{proposition}[What the standard local-form comparison yields]
\label{prop:s7v23-form-comparison}
Suppose, optimistically, that every scale-\(s\) cut-gauge block
projection satisfies a scale-independent one-block estimate
\begin{equation}
 \langle f,(I-K_Q)f\rangle
 \le C_0\sum_{e\in E(Q)}
       \mathbb E\|\nabla_e f\|^2.                       \label{eq:s7v23-one-block}
\end{equation}
Then the weighted sum of all block forms satisfies
\begin{equation}
 {\cal E}_{\rm cut}(f,f)
 \le C_0\Lambda_{L,d}{\cal E}_{\rm electric}(f,f).       \label{eq:s7v23-global-comparison}
\end{equation}
Combining this comparison with the Gaussian event-clock gap
\(\delta_d\) gives only the local-electric lower bound
\begin{equation}
 \operatorname{gap}_{\rm electric}
 \ge{\delta_d\over C_0\Lambda_{L,d}}
 =O(L^{-1})                                             \label{eq:s7v23-gap-loss}
\end{equation}
by this route.  Equation \eqref{eq:s7v23-gap-loss} is a limitation of
the current comparison, not an upper bound on the Yang--Mills gap.
\end{proposition}

\begin{proof}
Multiply \eqref{eq:s7v23-one-block} by \(a_s\) and sum over
\((s,z,\mu)\).  The coefficient of a fixed link gradient is precisely
the load \(\Lambda_{L,d}\) computed in
\Cref{thm:s7v23-load}.  This proves
\eqref{eq:s7v23-global-comparison}.  If
\({\cal E}_{\rm cut}\ge\delta_d\operatorname{Var}\), then
\eqref{eq:s7v23-global-comparison} implies
\({\cal E}_{\rm electric}\ge
\delta_d(C_0\Lambda_{L,d})^{-1}\operatorname{Var}\).
\end{proof}

\begin{correction}[Precise replacement for the V22 clock claim]
\label{corr:s7v23-clock}
The theorem \Cref{thm:s7v22-clock} is correct as a statement about
total scale rate per translation, orientation, or block-event density.
It does not establish a bounded comparison with the local electric
Hamiltonian.  The missing column ledger is
\eqref{eq:s7v23-Lambda-exact}, and it grows linearly in \(L\).
\end{correction}

\section{Revised nonlinear multiscale programme}
\label{sec:s7v23-programme}

\begin{theorem}[Rows now proved in the nonlinear update card]
\label{thm:s7v23-card}
The following parts of NDCUC are now rigorous.
\begin{enumerate}[label=\textup{(\arabic*)},leftmargin=3.5em]
\item The naive simultaneous-left-multiplication realization fails by
      the volume obstruction \eqref{eq:s7v23-volume}.
\item The compact cut-gauge block
      \eqref{eq:s7v23-cut-action} gives an exact nonlinear
      realization whose flat tangent contains every V22 cube
      direction.
\item The Wilson action and its conditional block density depend on
      only \(O(s^{d-1})\) mixed plaquettes along a cut-gauge orbit.
\item The Gaussian cut-gauge first-chaos frame, and hence its
      all-chaos heat-bath gap, is at least \(\delta_d\).
\item The exact physical column ledger for the V22 rates is
      \(\Lambda_{L,d}=\Theta(L)\).
\end{enumerate}
\end{theorem}

\begin{proof}
These are respectively
\Cref{thm:s7v23-volume,thm:s7v23-tangent,
thm:s7v23-score-support,thm:s7v23-frame-domination,
thm:s7v23-load}.
\end{proof}

\begin{assessment}[Progress at seventh-series V23]
\label{ass:s7v23-status}
The first nonlinear multiscale obstruction is no longer ambiguous.
Noncommutativity destroys the boundary cancellation for the naive
parallel-link update, and the failure can be volume order.  A compact
cut-gauge group action repairs this exactly: it changes only
boundary-straddling Wilson terms and contains the desired cube
curvature direction at the flat field.  However, the local-electric
column load of the V22 rates is linear in the lattice diameter, so a
uniform event-clock gap is not yet a physical Hamiltonian gap.
\end{assessment}

\begin{decision}[V23 bottleneck decision]
\label{dec:s7v23-bottleneck}
V24 will compute the full flat tangent projection of the cut-gauge
block, rather than retain only its single affine direction.  The
decisive test is whether its \(O(s^{d-1})\)-dimensional boundary
curvature range supplies the missing factor allowing rates of order
\(s^{-d}\), which reduce the per-link participation from \(O(L)\) to
\(O(\log L)\).  The same calculation must decide whether sparse
ownership or a further scale weight can remove that remaining
logarithm without destroying the frame.  If it cannot, a different
coarse variable is required before returning to nonlinear conditional
gaps.
\end{decision}

\begin{firewall}[Claim boundary after seventh-series V23]
\label{fire:s7v23-boundary}
V23 proves the volume-order failure of the naive cube update, constructs
an exact compact cut-gauge replacement, proves boundary localization
of its Wilson score, embeds the V22 Gaussian frame in its tangent
frame, and computes the exact linear growth of the local participation
load.

V23 does not prove a scale-uniform nonlinear conditional gap, a bounded
comparison with the local electric Hamiltonian, covariant large-scale
good-field control, toron and reflection estimates, continuum
construction, or the Yang--Mills mass gap.
\end{firewall}

\paragraph{Parent-version record.}
V23 was formed from
\texttt{Renewal\_Yang\_Mills\_Mass\_Gap\_Research\_Notes\_V22.tex}.
The V22 parent had \(8\,875\) logical source lines, \(322\,409\) bytes,
and SHA--256
\[
 \texttt{EB8E96B0378FBA49A52C6F3CC13C67D48EE4859145E6FC8B9826363FF533EE04}.
\]

% =============================================================================
% END APPENDED SEVENTH-SERIES V23 MATERIAL
% =============================================================================

\newpage

% =============================================================================
% BEGIN APPENDED SEVENTH-SERIES V24 MATERIAL
% =============================================================================

\chapter{A clock--frame incompatibility theorem for cut-gauge blocks}
\label{ch:s7v24}

V23 left open whether the full cut-gauge tangent space, which is much
larger than the single affine cube direction used in V22, could supply
enough additional curvature directions to lower the rates from
\(s^{-(d-1)}\) to \(s^{-d}\).  It cannot.  The reason is geometric and
does not require diagonalizing the block projection: every cut-gauge
curvature tangent is supported on the block boundary and is orthogonal
to constant curvature.  A lowest Fourier mode is almost constant on a
box, so its projection onto that tangent space has an additional
\((|k|s)^2\) suppression.

This yields two conclusions.  First, rates \(s^{-d}\) give at most an
\(O(L^{-1})\) lowest-mode frame, even before the residual
\(O(\log L)\) participation is normalized.  Second, allowing sparse
translations does not help: any family of these blocks with uniformly
bounded link-participation load has frame margin at most \(C/L\).

\section{Boundary support and harmonic annihilation}
\label{sec:s7v24-harmonic}

Work in one colour at the flat field; the argument repeats for every
Lie-algebra component.  For a vertex box \(Q=Q_{s,z}^{\mu}\), let
\[
 R_Q:C^1(\mathbb T_L^d)\longrightarrow C^1(\mathbb T_L^d)
\]
retain the internal links \(E(Q)\) and set all other links to zero.
The flat link tangent space and curvature tangent space of the
cut-gauge action are
\begin{equation}
 {\cal T}_Q=\{R_Qd_0\phi:\phi\in C^0(Q)\},\qquad
 {\cal C}_Q=d_1{\cal T}_Q.                              \label{eq:s7v24-TC}
\end{equation}
Write \(\Pi_Q\) for orthogonal projection onto \({\cal C}_Q\).

\begin{lemma}[Support of the linear cut-gauge curvature]
\label{lem:s7v24-support}
Every \(w\in{\cal C}_Q\) is supported on the mixed plaquette set
\(\partial_2Q\) of \eqref{eq:s7v23-partial2}.
\end{lemma}

\begin{proof}
Let \(w=d_1R_Qd_0\phi\).  On a plaquette wholly inside \(Q\), \(R_Q\)
is the identity on all four edges, so
\[
 w=d_1d_0\phi=0.
\]
On a plaquette with no internal edge, \(R_Qd_0\phi=0\).  Only mixed
plaquettes remain.
\end{proof}

\begin{lemma}[Constant-curvature annihilation]
\label{lem:s7v24-constant}
If \(h\) is any constant two-cochain on the periodic lattice, then
\begin{equation}
 \Pi_Qh=0.                                               \label{eq:s7v24-constant}
\end{equation}
\end{lemma}

\begin{proof}
Constant two-cochains obey \(d_1^*h=0\).  For
\(w=d_1R_Qd_0\phi\),
\[
 \langle h,w\rangle
 =\langle d_1^*h,R_Qd_0\phi\rangle=0.
\]
Thus \(h\perp{\cal C}_Q\), which is
\eqref{eq:s7v24-constant}.
\end{proof}

\section{The almost-constant Fourier estimate}
\label{sec:s7v24-Fourier}

For a nonzero lattice momentum \(k\), choose a unit polarization
\(\eta(k)\in\operatorname{Ran}d_1(k)\), and define the normalized
curvature plane wave
\begin{equation}
 y_k(p)=L^{-d/2}e^{ik\cdot x(p)}\eta(k),\qquad
 \|y_k\|_2=1.                                           \label{eq:s7v24-yk}
\end{equation}
Here \(x(p)\) is the base vertex of the positively oriented plaquette;
changing this convention only changes the constants below.

\begin{theorem}[One-block low-frequency projection bound]
\label{thm:s7v24-one-block}
There is a dimension-dependent constant \(C_d\) such that, for every
cut-gauge box of scale \(s\),
\begin{equation}
 \|\Pi_Qy_k\|_2^2
 \le {C_d\over L^d}|k|^2s^{d+1}.                        \label{eq:s7v24-one-block}
\end{equation}
\end{theorem}

\begin{proof}
Let
\[
 h_{k,z}(p)=L^{-d/2}e^{ik\cdot z}\eta(k)
\]
be the constant two-cochain with the phase frozen at the box
translation \(z\).  By \Cref{lem:s7v24-constant},
\(\Pi_Qh_{k,z}=0\), while
\Cref{lem:s7v24-support} and contractivity of orthogonal projection
give
\begin{align*}
 \|\Pi_Qy_k\|_2^2
 &=\|\Pi_Q(y_k-h_{k,z})\|_2^2\\
 &\le
 \|\mathbf1_{\partial_2Q}(y_k-h_{k,z})\|_2^2.
\end{align*}
For \(p\in\partial_2Q\), the distance from \(x(p)\) to \(z\) is at
most \(C'_ds\).  The elementary bound
\[
 |e^{ia}-e^{ib}|\le|a-b|
\]
therefore gives
\[
 |y_k(p)-h_{k,z}(p)|
 \le L^{-d/2}C'_d|k|s|\eta(k)|.
\]
Finally
\(|\partial_2Q|\le C''_ds^{d-1}\) by
\eqref{eq:s7v23-boundary-count}.  Multiplication proves
\eqref{eq:s7v24-one-block}.
\end{proof}

\begin{corollary}[Weighted translated-frame upper bound]
\label{cor:s7v24-translated}
For nonnegative scale rates \(a_s\), define the full cut-gauge
first-chaos frame
\begin{equation}
 {\mathbb F}_{\rm cut}(a)
 :=
 \sum_{s\in{\cal S}_L}a_s
 \sum_{z\in\mathbb T_L^d}\sum_{\mu=1}^d
 \Pi_{Q_{s,z}^{\mu}}.                                  \label{eq:s7v24-F}
\end{equation}
Then
\begin{equation}
 \langle y_k,{\mathbb F}_{\rm cut}(a)y_k\rangle
 \le C_d|k|^2
       \sum_{s\in{\cal S}_L}a_ss^{d+1}.                 \label{eq:s7v24-F-upper}
\end{equation}
\end{corollary}

\begin{proof}
Apply \Cref{thm:s7v24-one-block} to the \(dL^d\) oriented translates
at each scale and absorb \(d\) into \(C_d\).
\end{proof}

\begin{corollary}[The proposed \(s^{-d}\) rates fail]
\label{cor:s7v24-sminusd}
If \(a_s=s^{-d}\), then at a lowest nonzero momentum
\(|k_{\min}|=2\pi/L\),
\begin{equation}
 \langle y_{k_{\min}},
 {\mathbb F}_{\rm cut}(a)y_{k_{\min}}\rangle
 \le {C_d\over L}.                                      \label{eq:s7v24-sminusd}
\end{equation}
Consequently the full tangent projection does not compensate for the
extra factor \(s^{-1}\) in the rate.
\end{corollary}

\begin{proof}
Equation \eqref{eq:s7v24-F-upper} gives
\[
 C_d|k_{\min}|^2\sum_{s\in{\cal S}_L}s.
\]
The dyadic sum is \(L/2-1\), so this is \(O(L^{-1})\).
\end{proof}

\begin{correction}[The residual clock at rate \(s^{-d}\)]
\label{corr:s7v24-log}
By the exact incidence count \eqref{eq:s7v23-Ns},
\[
 \sum_{s\in{\cal S}_L}s^{-d}N_s
 =
 d|{\cal S}_L|
 -(d-1)\sum_{s\in{\cal S}_L}s^{-2}
 =\Theta(\log L).
\]
Thus \(s^{-d}\) has logarithmic, not bounded, link participation; even
before division by this logarithm its frame already vanishes by
\eqref{eq:s7v24-sminusd}.
\end{correction}

\section{A sparse-family no-go theorem}
\label{sec:s7v24-sparse}

The preceding argument is not an artefact of summing every
translation.  Let \({\cal Q}\) be any finite collection of anisotropic
boxes of the form \eqref{eq:s7v23-Q}, possibly using only selected
translations and orientations, with nonnegative weights \(a_Q\).
Define
\begin{equation}
 {\mathbb F}_{\cal Q}:=\sum_{Q\in{\cal Q}}a_Q\Pi_Q,\qquad
 \Lambda_{\cal Q}
 :=\sup_e\sum_{\substack{Q\in{\cal Q}\\e\in E(Q)}}a_Q.   \label{eq:s7v24-sparse-def}
\end{equation}

\begin{theorem}[Bounded local load forces a vanishing curvature frame]
\label{thm:s7v24-no-go}
There is a constant \(C_d\), independent of \(L\), the selected boxes,
and their weights, such that
\begin{equation}
 \lambda_{\min}
 \left(
 {\mathbb F}_{\cal Q}\big|_{\operatorname{Ran}d_1}
 \right)
 \le {C_d\Lambda_{\cal Q}\over L}.                       \label{eq:s7v24-no-go}
\end{equation}
In particular, no cut-gauge box family with
\(\sup_L\Lambda_{\cal Q}<\infty\) has a volume-uniform linearized
curvature frame.
\end{theorem}

\begin{proof}
Use the lowest plane wave \(y_{k_{\min}}\) in
\eqref{eq:s7v24-yk}.  The proof of
\Cref{thm:s7v24-one-block} gives the slightly more geometric estimate
\begin{equation}
 \|\Pi_Qy_{k_{\min}}\|_2^2
 \le {C_d|k_{\min}|^2\over L^d}
       \operatorname{diam}(Q)^2|\partial_2Q|.            \label{eq:s7v24-geometric}
\end{equation}
For a scale-\(s\) box,
\[
 \operatorname{diam}(Q)^2|\partial_2Q|
 \le C_ds^{d+1}
 \le C_dL\,|E(Q)|,
\]
because \(s\le L/4\) and \(|E(Q)|\ge s^d\).  Hence
\begin{align*}
 \langle y_{k_{\min}},
 {\mathbb F}_{\cal Q}y_{k_{\min}}\rangle
 &\le {C_d|k_{\min}|^2L\over L^d}
       \sum_{Q\in{\cal Q}}a_Q|E(Q)|\\
 &= {C_d|k_{\min}|^2L\over L^d}
       \sum_e\sum_{\substack{Q\in{\cal Q}\\e\in E(Q)}}a_Q\\
 &\le C_d|k_{\min}|^2L\Lambda_{\cal Q}.
\end{align*}
There are \(dL^d\) positive links; their factor \(d\) is absorbed into
\(C_d\).  Since \(|k_{\min}|=2\pi/L\), the last expression is at most
\(C_d\Lambda_{\cal Q}/L\).  The smallest eigenvalue is no larger than
this Rayleigh quotient.
\end{proof}

\begin{corollary}[A linear load is necessary for a uniform cut frame]
\label{cor:s7v24-linear-load}
If
\[
 {\mathbb F}_{\cal Q}\succeq\delta I
 \quad\hbox{on }\operatorname{Ran}d_1
\]
with \(\delta>0\) independent of \(L\), then
\begin{equation}
 \Lambda_{\cal Q}\ge c_d\delta L.                        \label{eq:s7v24-linear-load}
\end{equation}
Thus the \(\Theta(L)\) load found for the V22 rates in
\eqref{eq:s7v23-Lambda-exact} has the necessary order of magnitude
within the entire cut-gauge boundary-block class.
\end{corollary}

\begin{proof}
Rearrange \eqref{eq:s7v24-no-go}.
\end{proof}

\begin{firewall}[What the no-go theorem does and does not exclude]
\label{fire:s7v24-scope}
The theorem concerns comparison schemes whose flat update spaces are
cut-gauge curvature images, whose Wilson-action variation is confined
to codimension-one box boundaries, and whose physical cost is measured
by local link participation.  It permits arbitrary box selection,
orientation, and nonnegative weights.

It does not rule out a Yang--Mills mass gap.  It also does not rule out
a direct renormalization proof for the local Hamiltonian, an enlarged
system with auxiliary coarse fields, a comparison exploiting a
different local form, or collective variables whose flat curvature
tangents are not boundary-supported.
\end{firewall}

\section{Upstream consequence}
\label{sec:s7v24-upstream}

\begin{theorem}[Closure of the cut-gauge clock branch]
\label{thm:s7v24-closure}
The following three requirements are mutually incompatible for
cut-gauge box blocks:
\begin{enumerate}[label=\textup{(\roman*)},leftmargin=3.5em]
\item a volume-uniform first-chaos curvature frame;
\item a volume-uniform local link-participation bound; and
\item use only the boundary-supported flat tangent spaces
      \({\cal C}_Q\).
\end{enumerate}
\end{theorem}

\begin{proof}
Requirements (ii) and (iii) imply
\(\lambda_{\min}\le C_d/L\) by
\Cref{thm:s7v24-no-go}, contradicting (i).
\end{proof}

\begin{assessment}[Progress at seventh-series V24]
\label{ass:s7v24-status}
The nonlinear cut-gauge construction of V23 remains useful because it
gives exact boundary localization and a compact conditional orbit.
It cannot, however, be the sole bridge from a uniform multiscale
heat-bath frame to the local electric Hamiltonian.  The lowest Fourier
mode proves a quantitative tradeoff:
\[
 \hbox{frame margin}\le C_d
 {\hbox{maximum link participation}\over L}.
\]
This closes the proposed full-tangent and sparse-ownership repairs
within that block class.
\end{assessment}

\begin{decision}[V24 bottleneck decision]
\label{dec:s7v24-bottleneck}
V25 is a mandatory cross-programme audit.  It will inspect the current
SM and NS notes specifically for mechanisms that avoid a local
column-load comparison: auxiliary coarse fields, exact multilevel
factorizations, or a direct coercive local Hamiltonian estimate.  The
Yang--Mills branch will import only a type-correct mechanism; if none
exists, it will formulate the required bulk-curvature or extended-state
replacement explicitly.
\end{decision}

\begin{firewall}[Claim boundary after seventh-series V24]
\label{fire:s7v24-boundary}
V24 proves the one-block almost-constant Fourier estimate, the failure
of \(s^{-d}\) rates, and a sparse-family no-go theorem showing that a
uniform cut-gauge curvature frame forces at least linear local
participation.

V24 does not construct the required replacement coarse variables,
prove a nonlinear conditional gap, compare uniformly with the local
electric Hamiltonian, control torons and interfaces, construct the
continuum theory, or prove the Yang--Mills mass gap.
\end{firewall}

\paragraph{Parent-version record.}
V24 was formed from
\texttt{Renewal\_Yang\_Mills\_Mass\_Gap\_Research\_Notes\_V23.tex}.
The V23 parent had \(9\,413\) logical source lines, \(343\,328\) bytes,
and SHA--256
\[
 \texttt{A3083D787387F471C15C8951B90064719C5B882A14297260E187526D81326864}.
\]

% =============================================================================
% END APPENDED SEVENTH-SERIES V24 MATERIAL
% =============================================================================

\newpage

% =============================================================================
% BEGIN APPENDED SEVENTH-SERIES V25 MATERIAL
% =============================================================================

\chapter{Mandatory companion audit and the direct spectral route}
\label{ch:s7v25}

V24 proves that a uniform multiscale heat-bath frame built only from
cut-gauge boundary blocks necessarily has local link-participation of
order at least \(L\).  This mandatory audit asks whether the current
Standard Model--Einstein or Navier--Stokes notes contain a
type-correct mechanism which avoids that load comparison.  They do
not.  The SM programme instead makes a useful strategic move: after
proving its own clock obstruction, it restricts to physical local
pairings and explicitly declines to call those estimates a full gap.

For Yang--Mills the analogous honest pivot is not to weaken the final
claim, but to change the sufficient criterion.  A mass gap of the
Osterwalder--Schrader Hamiltonian follows directly from a
regulator-uniform exponential connected-correlation bound on a dense
gauge-invariant core.  Such a proof does not require a spectral gap for
an auxiliary sampling generator.  This note proves that abstract
criterion and installs it as the next route.

\section{Companion files inspected}
\label{sec:s7v25-audit-record}

The audit inspected the current files
\begin{align}
 &\texttt{Renewal\_SM\_Einstein\_Continuum\_Research\_Notes\_V72.tex},
 \notag\\
 &\hspace{2em}28\,437\ \text{logical lines},\quad
 999\,419\ \text{bytes},\quad
 \operatorname{SHA256}
 =\texttt{1AED43FFCDA4D5CF70C62AA715BA34B25392C1357D16FBA8C737EA2A20DDE5DC},
                                                               \label{eq:s7v25-SM-file}\\
 &\texttt{Renewal\_NS\_Stability\_Research\_Notes\_V26.tex},
 \notag\\
 &\hspace{2em}5\,319\ \text{logical lines},\quad
 278\,283\ \text{bytes},\quad
 \operatorname{SHA256}
 =\texttt{82C887F8C2C69E4F28FAD7C3DC8DE5B9099CB5A4BF431EBA1FFF3E91935978AD}.
                                                               \label{eq:s7v25-NS-file}
\end{align}
Each has one live document terminator.  The NS file is byte-for-byte
unchanged from the V20 audit; the SM file has advanced from V63 to
V72.

\begin{audit}[SM V64: a matching clock obstruction]
\label{audit:s7v25-SM64}
SM V64 replaces all translated dyadic scalar blocks by randomized
nonoverlapping tilings.  It proves that recovering the uniform
Gaussian frame requires per-site clock of order \(L^3\), and that every
rate choice in that tiling family obeys a matching slow-mode upper
bound.  Randomization reduces overlap but does not create
low-frequency strength.  This is the scalar, higher-derivative analogue
of the Yang--Mills tradeoff in V24, not an escape from it.
\end{audit}

\begin{audit}[SM V65--V68: observable-restricted response]
\label{audit:s7v25-SM65-68}
After the clock obstruction, SM V65 rules out generic Poisson--Moser
transport because the minimum-norm inverse carries the reciprocal
missing gap.  V66--V67 instead factor declared scalar sources through
direct divergence or balance-law primitives and move a derivative to a
fixed, spacetime-smoothed test.  V68 proves weak--strong cofinal
convergence of those local pairings.  The notes explicitly state that
this supplies neither full mixing nor a Hamiltonian mass gap.
\end{audit}

\begin{audit}[SM V69--V72: fixed-patch form convergence]
\label{audit:s7v25-SM69-72}
SM V69--V71 prove local Lax--Milgram solvability, adjoint
test-resolvent identities, and stability under uniformly coercive
variable-metric and lower-order forms.  V72 derives the conformal ADM
principal/potential/source ledger.  These are fixed-patch results:
their Poincaré and coercivity constants are not asserted uniformly on
growing domains, and metric-chart probability, curvature convergence,
boundary influence, and positivity remain open.
\end{audit}

\begin{audit}[NS V26]
\label{audit:s7v25-NS26}
The unchanged NS note refines rank-one norms of Oseen-kernel
derivatives.  It neither constructs a gauge-field coarse variable nor
provides a spectral or clock mechanism relevant to the V24
cut-gauge obstruction.
\end{audit}

\begin{theorem}[Audit non-transfer]
\label{thm:s7v25-no-transfer}
None of the audited results supplies a bounded-local-clock,
volume-uniform Yang--Mills block gap.  The following transfers would
be invalid:
\begin{enumerate}[label=\textup{(\roman*)},leftmargin=3.5em]
\item importing the SM scalar derivative exponent into the
      one-derivative curvature complex;
\item treating a fixed-test inverse-elliptic pairing as a full
      Hamiltonian spectral gap;
\item applying fixed-patch Poincaré coercivity on a growing spatial
      volume without its diameter dependence; or
\item importing the NS rank-one kernel improvement into a compact
      gauge-field frame.
\end{enumerate}
\end{theorem}

\begin{proof}
Items (i)--(iv) have, respectively, different Fourier symbols,
different operator claims, nonuniform domain constants, and different
state spaces.  The audit statements above record each mismatch.  In
particular, SM V65--V72 preserve their restricted-observable and
fixed-patch firewalls, so they cannot imply the full Yang--Mills gap.
\end{proof}

\section{The dense-core spectral criterion}
\label{sec:s7v25-spectral}

The audit nevertheless suggests avoiding the obstructed auxiliary
sampling comparison.  The following elementary theorem states exactly
what direct Euclidean decay must prove.

\begin{theorem}[Dense-core exponential criterion for a mass gap]
\label{thm:s7v25-dense-core}
Let \({\cal H}\) be a Hilbert space, let \(H\ge0\) be self-adjoint, and
let \(\Omega\in{\cal H}\) be a normalized vector with \(H\Omega=0\).
Suppose \({\cal D}\subset\Omega^\perp\) is dense in
\(\Omega^\perp\).  The following are equivalent for \(m>0\).
\begin{enumerate}[label=\textup{(\roman*)},leftmargin=3.5em]
\item
\[
 \operatorname{spec}\bigl(H|_{\Omega^\perp}\bigr)
 \subset[m,\infty).
                                                               \label{eq:s7v25-gap}
\]
\item For every \(\psi\in{\cal D}\), there are
\(C_\psi<\infty\) and \(t_\psi<\infty\) such that
\begin{equation}
 0\le\langle\psi,e^{-tH}\psi\rangle
 \le C_\psi e^{-mt}
 \qquad(t\ge t_\psi).                                   \label{eq:s7v25-core-decay}
\end{equation}
\end{enumerate}
\end{theorem}

\begin{proof}
If (i) holds, the spectral theorem gives
\[
 0\le\langle\psi,e^{-tH}\psi\rangle
 \le e^{-mt}\|\psi\|^2
\]
for every \(\psi\perp\Omega\), so (ii) follows.

Conversely, fix \(\epsilon\in(0,m)\) and
\(\psi\in{\cal D}\).  Let \(\mu_\psi\) be the spectral measure of
\(\psi\).  If
\(\mu_\psi([0,m-\epsilon])>0\), then
\[
 e^{mt}\langle\psi,e^{-tH}\psi\rangle
 =
 \int e^{t(m-\lambda)}\,d\mu_\psi(\lambda)
 \ge
 e^{\epsilon t}\mu_\psi([0,m-\epsilon]),
\]
contradicting \eqref{eq:s7v25-core-decay}.  Hence the spectral
projection \(E_H([0,m-\epsilon])\) annihilates every vector in
\({\cal D}\).  It is bounded and \({\cal D}\) is dense in
\(\Omega^\perp\), so it annihilates all of \(\Omega^\perp\).  Let
\(\epsilon\downarrow0\) to obtain (i).
\end{proof}

\begin{remark}[Observable-dependent prefactors are harmless]
\label{rem:s7v25-prefactor}
The decay rate \(m\) must be common to the dense core, but
\(C_\psi\) and \(t_\psi\) may depend on the observable.  No uniform
norm bound on the whole unit sphere is needed for the spectral-support
argument.
\end{remark}

\section{Osterwalder--Schrader form of the criterion}
\label{sec:s7v25-OS}

Let \(\omega\) be a reflection-positive Euclidean state with time
reflection \(\Theta\), time translation \(\tau_t\), reconstructed
vacuum \(\Omega\), and Hamiltonian \(H\).  For a positive-time
gauge-invariant observable \(A\), put
\[
 A^\circ:=A-\omega(A)\mathbf1,\qquad
 \psi_A:=[A^\circ].
\]
The reconstruction identity is
\begin{equation}
 \omega\!\left((\Theta A^\circ)\tau_tA^\circ\right)
 =
 \langle\psi_A,e^{-tH}\psi_A\rangle                    \label{eq:s7v25-OS-identity}
\end{equation}
once the translated supports are separated in the declared convention.

\begin{theorem}[Gauge-invariant Euclidean mass-gap criterion]
\label{thm:s7v25-OS-gap}
Assume:
\begin{enumerate}[label=\textup{(\roman*)},leftmargin=3.5em]
\item \(\omega\) satisfies the Osterwalder--Schrader hypotheses needed
      for reconstruction and has a unique vacuum;
\item a gauge-invariant local algebra \({\cal A}_{\rm loc}^{\rm inv}\)
      has
\[
 \{\psi_A:A\in{\cal A}_{\rm loc}^{\rm inv}\}
\quad\hbox{dense in }\Omega^\perp;
\]
\item there is \(m_*>0\) such that, for every
      \(A\in{\cal A}_{\rm loc}^{\rm inv}\), one has
\begin{equation}
 0\le
 \omega\!\left((\Theta A^\circ)\tau_tA^\circ\right)
 \le C_Ae^{-m_*t}
 \qquad(t\ge t_A).                                      \label{eq:s7v25-OS-decay}
\end{equation}
\end{enumerate}
Then the reconstructed physical Hamiltonian has mass gap at least
\(m_*\).
\end{theorem}

\begin{proof}
By \eqref{eq:s7v25-OS-identity}, hypothesis (iii) is
\eqref{eq:s7v25-core-decay} on the dense subspace in (ii).  Apply
\Cref{thm:s7v25-dense-core}.  Vacuum uniqueness identifies its
orthogonal complement as the nonvacuum physical sector.
\end{proof}

\begin{corollary}[Regulator-uniform handoff]
\label{cor:s7v25-regulator}
Suppose a cofinal family of reflection-positive regulated
Yang--Mills states converges on a fixed gauge-invariant local algebra,
the limiting state obeys the reconstruction and density hypotheses of
\Cref{thm:s7v25-OS-gap}, and the regulated connected correlators obey
\eqref{eq:s7v25-OS-decay} with one \(m_*>0\) and constants \(C_A,t_A\)
independent of the regulator.  Then the limiting Hamiltonian has mass
gap at least \(m_*\).
\end{corollary}

\begin{proof}
Pass to the limit in each fixed local correlator.  The same exponential
bound holds for the limiting state.  Apply
\Cref{thm:s7v25-OS-gap}.
\end{proof}

\begin{firewall}[A selected finite family is not a dense core]
\label{fire:s7v25-density}
Decay for one Wilson loop, finitely many plaquettes, or a
non-dense test class does not prove a full mass gap.  Density in the
reconstructed nonvacuum Hilbert space, or a theorem extending the
bound from a core to that space, is an independent row.
\end{firewall}

\begin{firewall}[Polynomial response bounds are not exponential time
decay]
\label{fire:s7v25-exponential}
The SM V65--V72 estimates control inverse-elliptic pairings on fixed
local tests.  They do not give
\eqref{eq:s7v25-OS-decay}.  Yang--Mills still needs a common positive
decay exponent in physical Euclidean time, uniform through volume and
continuum limits.
\end{firewall}

\section{The direct gauge-invariant decay card}
\label{sec:s7v25-card}

\begin{definition}[Direct Osterwalder--Schrader gap card, DOSGC]
\label{def:s7v25-DOSGC}
DOSGC consists of:
\begin{enumerate}[label=\textup{(G\arabic*)},leftmargin=4em]
\item a reflection-positive regulated Yang--Mills measure and transfer
      operator in the gauge-invariant sector;
\item a fixed algebra of local spin-network or Wilson observables and
      proof that its centered reconstructed vectors form a dense core;
\item a connected two-slab representation with all gauge projections,
      boundary data, and normalizations explicit;
\item one positive physical-time decay rate \(m_*\), uniform in spatial
      volume and ultraviolet regulator, for every observable in the
      core;
\item observable-dependent prefactors controlled well enough to pass
      each fixed local observable to the limit;
\item uniqueness of the limiting vacuum and exclusion of residual
      toron or topological zero-energy sectors;
\item convergence of the reflected local correlators and preservation
      of positivity; and
\item application of \Cref{thm:s7v25-OS-gap} to the reconstructed
      continuum Hamiltonian.
\end{enumerate}
\end{definition}

\begin{theorem}[What the audit changes]
\label{thm:s7v25-route}
The V24 cut-gauge no-go blocks only the route
\[
 \hbox{uniform nonlocal block gap}
 \Longrightarrow
 \hbox{bounded local-form comparison}
 \Longrightarrow
 \hbox{physical mass gap}.
\]
It does not block the direct implication
\[
 \hbox{uniform connected Euclidean decay on a dense physical core}
 \Longrightarrow
 \hbox{physical mass gap},
\]
which is rigorous by \Cref{thm:s7v25-OS-gap}.  DOSGC is therefore the
selected upstream branch.
\end{theorem}

\begin{proof}
The first route fails within the cut-gauge class by
\Cref{thm:s7v24-no-go}.  The second implication is
\Cref{thm:s7v25-OS-gap} and contains no auxiliary sampling generator
or link-participation comparison.
\end{proof}

\begin{assessment}[Progress at seventh-series V25]
\label{ass:s7v25-status}
The mandatory audit finds no type-correct companion mechanism that
repairs the V24 clock--frame incompatibility.  It does identify the
proper upstream escape: prove exponential physical-time clustering
directly on a dense gauge-invariant OS core.  The exact spectral
theorem needed for that handoff is now proved, including
observable-dependent prefactors and the cofinal limiting statement.
The difficult new row is the regulator-uniform decay estimate itself.
\end{assessment}

\begin{decision}[V25 bottleneck decision]
\label{dec:s7v25-bottleneck}
V26 will work at fixed lattice regulator and derive the exact
gauge-projected transfer-kernel representation of
\eqref{eq:s7v25-OS-identity}.  It will identify a concrete dense
spin-network core and express the desired common decay exponent as a
uniform contraction on its centered two-slab kernel.  Strong-coupling
character expansion may verify this at fixed bare coupling, but any
failure to remain uniform on the continuum scaling trajectory will be
recorded rather than promoted to a continuum gap.
\end{decision}

\begin{firewall}[Claim boundary after seventh-series V25]
\label{fire:s7v25-boundary}
V25 records the mandatory SM V72/NS V26 audit and proves a dense-core
Osterwalder--Schrader exponential criterion sufficient for the full
physical Hamiltonian mass gap.

V25 does not prove the required uniform exponential correlator bound,
density and uniqueness in the continuum limit, the transfer-kernel
contraction, regulator removal, or the Yang--Mills mass gap.
\end{firewall}

\paragraph{Parent-version record.}
V25 was formed from
\texttt{Renewal\_Yang\_Mills\_Mass\_Gap\_Research\_Notes\_V24.tex}.
The V24 parent had \(9\,774\) logical source lines, \(355\,735\) bytes,
and SHA--256
\[
 \texttt{D75E00B82E02C9A84EC336F84CFE9B06340D443DBABE6CBC8CBB134C9E3A3561}.
\]

% =============================================================================
% END APPENDED SEVENTH-SERIES V25 MATERIAL
% =============================================================================

\newpage

% =============================================================================
% BEGIN APPENDED SEVENTH-SERIES V26 MATERIAL
% =============================================================================

\chapter{The gauge-projected transfer kernel and a dense spin-network
core}
\label{ch:s7v26}

V25 replaces the obstructed auxiliary block-gap route by a direct
physical criterion: exponential Euclidean-time decay on a dense
gauge-invariant Osterwalder--Schrader core.  The present note supplies
the exact fixed-cutoff transfer object and proves that finite
spin-network observables give such a core.  It also separates two
facts which are often conflated:
\begin{enumerate}[label=\textup{(\roman*)},leftmargin=3.5em]
\item at every fixed finite cutoff a strictly positive heat-kernel
      transfer has a unique vacuum and a positive eigenvalue gap; but
\item neither finiteness nor a global operator-norm perturbation makes
      that gap uniform in volume and continuum scaling.
\end{enumerate}

\section{Finite spatial lattice and gauge projection}
\label{sec:s7v26-space}

Let \(\Gamma=(V,E)\) be a finite connected oriented spatial graph and
let
\[
 X_\Gamma:=G^E,\qquad {\cal H}_\Gamma:=L^2(X_\Gamma,dU),
\]
where \(dU\) is product normalized Haar measure.  The vertex gauge
group \({\cal G}_\Gamma=G^V\) acts by
\begin{equation}
 (g\cdot U)_e
 =g_{s(e)}U_eg_{t(e)}^{-1}.                             \label{eq:s7v26-gauge}
\end{equation}
Let
\begin{equation}
 (P_\Gamma f)(U)
 :=
 \int_{{\cal G}_\Gamma}f(g\cdot U)\,dg                  \label{eq:s7v26-P}
\end{equation}
and
\[
 {\cal H}^{\rm phys}_\Gamma
 :=P_\Gamma{\cal H}_\Gamma.
\]

Fix \(\tau>0\).  Let \(k_\tau\) be the normalized central heat kernel
on \(G\), in the convention
\begin{equation}
 k_\tau(h)
 =
 \sum_{\pi\in\widehat G}
 d_\pi e^{-\tau C_2(\pi)}\chi_\pi(h).                   \label{eq:s7v26-heat}
\end{equation}
The archived compact-group calculations establish this
representation-normalized expansion; here it is used to define the
one-step electric convolution
\begin{equation}
 (C_\tau f)(U)
 :=
 \int_{X_\Gamma}
 \prod_{e\in E}k_\tau(U_eV_e^{-1})f(V)\,dV.             \label{eq:s7v26-C}
\end{equation}

\begin{lemma}[Gauge covariance and positivity]
\label{lem:s7v26-C}
The operator \(C_\tau\) is self-adjoint, injective, positive, and a
contraction on \({\cal H}_\Gamma\).  It commutes with
\(P_\Gamma\).
\end{lemma}

\begin{proof}
On a tensor product of Peter--Weyl matrix coefficients,
\(C_\tau\) multiplies the edge representation \(\pi_e\) by
\(e^{-\tau C_2(\pi_e)}\).  These multipliers lie in \((0,1]\), proving
self-adjointness, injectivity, positivity, and contraction.

Under simultaneous gauge transformation of \(U\) and \(V\),
\(U_eV_e^{-1}\) is conjugated at \(s(e)\).  Centrality of
\(k_\tau\) and Haar invariance prove gauge covariance, hence
\([C_\tau,P_\Gamma]=0\).
\end{proof}

\section{The symmetrized physical transfer operator}
\label{sec:s7v26-transfer}

Let \(S_{\rm sp}:X_\Gamma\to[0,\infty)\) be a continuous,
gauge-invariant spatial Wilson potential and put
\begin{equation}
 M_{\rm sp}(U):=e^{-S_{\rm sp}(U)/2}.                   \label{eq:s7v26-M}
\end{equation}
Define
\begin{equation}
 T_\Gamma
 :=
 M_{\rm sp}P_\Gamma C_\tau P_\Gamma M_{\rm sp}
 \quad\hbox{on }{\cal H}^{\rm phys}_\Gamma.             \label{eq:s7v26-T}
\end{equation}
Since \(M_{\rm sp}\) is gauge invariant, it commutes with
\(P_\Gamma\).

\begin{theorem}[Exact gauge-projected transfer kernel]
\label{thm:s7v26-kernel}
The operator \(T_\Gamma\) is a positive self-adjoint compact
contraction.  Its integral kernel on the physical space is
\begin{align}
 K_\Gamma(U,V)
 &=
 e^{-S_{\rm sp}(U)/2}e^{-S_{\rm sp}(V)/2}
 \notag\\
 &\quad\times
 \int_{{\cal G}_\Gamma}
 \prod_{e\in E}
 k_\tau\!\left(
 U_e\bigl(g_{s(e)}V_eg_{t(e)}^{-1}\bigr)^{-1}
 \right)\,dg.                                           \label{eq:s7v26-kernel}
\end{align}
Moreover \(K_\Gamma(U,V)>0\) for every \(U,V\).
\end{theorem}

\begin{proof}
On the physical subspace,
\[
 T_\Gamma=M_{\rm sp}C_\tau M_{\rm sp}.
\]
Positivity and self-adjointness follow from
\[
 T_\Gamma
 =(C_\tau^{1/2}M_{\rm sp})^*
  (C_\tau^{1/2}M_{\rm sp}).
\]
Both multiplication factors have norm at most one, so \(T_\Gamma\)
is a contraction.  Inserting the gauge average
\eqref{eq:s7v26-P} into \eqref{eq:s7v26-C} gives
\eqref{eq:s7v26-kernel}.  The heat kernel is continuous and strictly
positive for \(\tau>0\), while the spatial factors are strictly
positive.  Thus the kernel is continuous and strictly positive on the
compact space \(X_\Gamma\times X_\Gamma\).  A continuous kernel on a
compact space is Hilbert--Schmidt and hence compact.
\end{proof}

\begin{theorem}[Unique finite-cutoff vacuum]
\label{thm:s7v26-vacuum}
The spectral radius \(\lambda_{0,\Gamma}>0\) of \(T_\Gamma\) is a
simple eigenvalue.  Its normalized eigenfunction
\(\Omega_\Gamma\) can be chosen continuous, gauge invariant, and
strictly positive.  Every other eigenvalue is strictly smaller than
\(\lambda_{0,\Gamma}\).
\end{theorem}

\begin{proof}
\Cref{thm:s7v26-kernel} makes \(T_\Gamma\) positivity improving:
if \(f\ge0\) is nonzero, then \(T_\Gamma f>0\) everywhere.  The compact
positivity-improving Perron--Frobenius--Jentzsch theorem gives a simple
positive top eigenvalue and a strictly positive eigenfunction.  The
kernel maps \(L^2\) into continuous functions, so the eigenfunction is
continuous.  Gauge invariance follows because the operator acts on
\({\cal H}^{\rm phys}_\Gamma\).  Simplicity and positivity improvement
exclude any other eigenvalue of the same modulus.
\end{proof}

\begin{remark}[Finiteness is not a uniform estimate]
\label{rem:s7v26-finite}
\Cref{thm:s7v26-vacuum} proves a nonzero spectral separation for each
fixed finite graph because the compact spectrum can accumulate only at
zero.  It gives no lower bound uniform in \(|E|\), \(\tau\), spatial
coupling, or lattice spacing.
\end{remark}

\section{Spin-network density and the exact correlator}
\label{sec:s7v26-core}

Let \({\cal S}_\Gamma\) be the algebraic span of finite
gauge-invariant spin-network functions: tensor products of edge matrix
coefficients with invariant intertwiners at the vertices.

\begin{theorem}[Spin networks are dense in the physical space]
\label{thm:s7v26-spin-density}
The space \({\cal S}_\Gamma\) is dense in
\({\cal H}^{\rm phys}_\Gamma\).  In addition,
\begin{equation}
 {\cal D}_\Gamma
 :=
 \left\{
 \bigl(F-\langle\Omega_\Gamma,F\Omega_\Gamma\rangle\bigr)
 \Omega_\Gamma:
 F\in{\cal S}_\Gamma
 \right\}                                               \label{eq:s7v26-D}
\end{equation}
is dense in \(\Omega_\Gamma^\perp\).
\end{theorem}

\begin{proof}
Finite tensor products of Peter--Weyl matrix coefficients are dense
in \(L^2(G^E)\).  Applying the bounded orthogonal projection
\(P_\Gamma\) preserves density in its range.  Gauge averaging contracts
the representation indices with invariant vertex tensors, so the
resulting algebraic range is precisely the spin-network span.  This
proves the first assertion.

For cyclic density, suppose \(\psi\in{\cal H}^{\rm phys}_\Gamma\) is
orthogonal to \(F\Omega_\Gamma\) for every
\(F\in{\cal S}_\Gamma\).  Since spin-network polynomials are uniformly
dense in the continuous gauge-invariant functions on the compact
quotient, the finite measure
\[
 d\nu_\psi=\overline\psi\,\Omega_\Gamma\,dU
\]
annihilates every continuous gauge-invariant test.  Hence
\(\overline\psi\,\Omega_\Gamma=0\) almost everywhere.  Strict
positivity of \(\Omega_\Gamma\) gives \(\psi=0\).  Thus
\(\{F\Omega_\Gamma\}\) is dense.  Subtracting its component along
\(\Omega_\Gamma\) proves density of \eqref{eq:s7v26-D} in
\(\Omega_\Gamma^\perp\).
\end{proof}

\begin{definition}[Normalized transfer and Hamiltonian]
\label{def:s7v26-H}
For time spacing \(a_t>0\), set
\begin{equation}
 \widehat T_\Gamma:={T_\Gamma\over\lambda_{0,\Gamma}},
 \qquad
 H_\Gamma:=-a_t^{-1}\log\widehat T_\Gamma.              \label{eq:s7v26-H}
\end{equation}
The logarithm is defined by spectral calculus.  Injectivity of
\(C_\tau\) and strict positivity of \(M_{\rm sp}\) imply that
\(\widehat T_\Gamma\) is injective.
\end{definition}

\begin{theorem}[Exact two-slab connected transfer identity]
\label{thm:s7v26-correlator}
For \(F\in{\cal S}_\Gamma\), put
\[
 F^\circ
 :=F-\langle\Omega_\Gamma,F\Omega_\Gamma\rangle,\qquad
 \psi_F:=F^\circ\Omega_\Gamma.
\]
Then the normalized connected correlation at integer time separation
\(n\) is
\begin{equation}
 C_{\Gamma,F}(n)
 :=
 \langle\psi_F,\widehat T_\Gamma^{\,n}\psi_F\rangle
 =
 \langle\psi_F,e^{-na_tH_\Gamma}\psi_F\rangle
 \ge0.                                                   \label{eq:s7v26-correlator}
\end{equation}
Equivalently, expanding each transfer factor with
\eqref{eq:s7v26-kernel} gives the gauge-projected Euclidean path
integral on the \(n\)-slab cylinder with insertions \(F^\circ\) at its
two ends.
\end{theorem}

\begin{proof}
The first equality is the transfer definition of the two-end
correlator after division by the vacuum eigenvalue at each step.  The
second is \eqref{eq:s7v26-H}.  Positivity follows from
\(\widehat T_\Gamma^{\,n}=e^{-na_tH_\Gamma}\ge0\).
Fubini's theorem and repeated insertion of
\eqref{eq:s7v26-kernel} give the slab integral, including one gauge
average at each interface.
\end{proof}

\begin{corollary}[Exact finite-cutoff mass ratio]
\label{cor:s7v26-ratio}
Let \(\lambda_{1,\Gamma}\) be the largest eigenvalue of \(T_\Gamma\)
on \(\Omega_\Gamma^\perp\).  Then
\begin{equation}
 m_\Gamma
 :=
 \inf\operatorname{spec}
 (H_\Gamma|_{\Omega_\Gamma^\perp})
 =
 {1\over a_t}
 \log{\lambda_{0,\Gamma}\over\lambda_{1,\Gamma}}.        \label{eq:s7v26-ratio}
\end{equation}
Moreover \(m_\Gamma\ge m_*\) if and only if every
\(F\in{\cal S}_\Gamma\) obeys
\begin{equation}
 C_{\Gamma,F}(n)
 \le\|\psi_F\|^2e^{-m_*na_t}
 \qquad(n\ge0).                                         \label{eq:s7v26-uniform-decay}
\end{equation}
\end{corollary}

\begin{proof}
The nonvacuum norm of \(\widehat T_\Gamma\) is
\(\lambda_{1,\Gamma}/\lambda_{0,\Gamma}\).  Spectral calculus gives
\eqref{eq:s7v26-ratio}.  The gap implies
\eqref{eq:s7v26-uniform-decay}.  Conversely, the latter holds on the
dense subspace \({\cal D}_\Gamma\) from
\Cref{thm:s7v26-spin-density}; apply the discrete-time version of
\Cref{thm:s7v25-dense-core}, or take \(t=na_t\) in its spectral-measure
proof.
\end{proof}

\section{The solvable electric point and the crude perturbation}
\label{sec:s7v26-electric}

At \(S_{\rm sp}=0\), one has \(T_{\Gamma,0}=C_\tau\) on the physical
space.  Define the least nontrivial admissible spin-network Casimir
\begin{equation}
 \kappa_\Gamma
 :=
 \min\left\{
 \sum_{e\in E}C_2(\pi_e):
 \boldsymbol\pi\ \hbox{admits invariant vertex intertwiners and is
 not identically trivial}
 \right\}.                                               \label{eq:s7v26-kappa}
\end{equation}
If the physical space contains no nonconstant function, put
\(\kappa_\Gamma=\infty\).

\begin{theorem}[Exact electric transfer gap]
\label{thm:s7v26-electric}
At \(S_{\rm sp}=0\),
\begin{equation}
 \lambda_{0,\Gamma}=1,\qquad
 \lambda_{1,\Gamma}=e^{-\tau\kappa_\Gamma},\qquad
 m_{\Gamma,0}={\tau\over a_t}\kappa_\Gamma.              \label{eq:s7v26-electric}
\end{equation}
\end{theorem}

\begin{proof}
The constant spin network has multiplier one.  Every nonconstant
admissible spin network with edge labels \(\boldsymbol\pi\) has
multiplier
\[
 \prod_e e^{-\tau C_2(\pi_e)}
 =
 e^{-\tau\sum_eC_2(\pi_e)}.
\]
Maximizing this nontrivial multiplier is equivalent to minimizing the
Casimir sum, giving \eqref{eq:s7v26-electric}.
\end{proof}

\begin{proposition}[Finite-volume operator-norm perturbation]
\label{prop:s7v26-perturbation}
Let
\[
 R_\Gamma:=\|S_{\rm sp}\|_\infty,\qquad
 \eta_\Gamma:=1-e^{-R_\Gamma/2}.
\]
Then
\begin{equation}
 \|T_\Gamma-C_\tau\|\le2\eta_\Gamma.                    \label{eq:s7v26-perturbation}
\end{equation}
If
\begin{equation}
 4\eta_\Gamma<1-e^{-\tau\kappa_\Gamma},                 \label{eq:s7v26-perturb-condition}
\end{equation}
then
\begin{equation}
 {\lambda_{1,\Gamma}\over\lambda_{0,\Gamma}}
 \le
 {e^{-\tau\kappa_\Gamma}+2\eta_\Gamma
  \over1-2\eta_\Gamma}
 <1.                                                     \label{eq:s7v26-perturb-ratio}
\end{equation}
\end{proposition}

\begin{proof}
Since \(0<M_{\rm sp}\le1\),
\(\|M_{\rm sp}-I\|=\eta_\Gamma\).  Therefore
\[
 M_{\rm sp}C_\tau M_{\rm sp}-C_\tau
 =(M_{\rm sp}-I)C_\tau M_{\rm sp}
 +C_\tau(M_{\rm sp}-I),
\]
whose norm is at most \(2\eta_\Gamma\).  The min--max principle moves
each ordered eigenvalue by at most that amount, so
\[
 \lambda_{0,\Gamma}\ge1-2\eta_\Gamma,\qquad
 \lambda_{1,\Gamma}\le
 e^{-\tau\kappa_\Gamma}+2\eta_\Gamma.
\]
Condition \eqref{eq:s7v26-perturb-condition} makes the numerator smaller
than the denominator and gives \eqref{eq:s7v26-perturb-ratio}.
\end{proof}

\begin{firewall}[The crude perturbation is extensive]
\label{fire:s7v26-extensive}
For the normalized Wilson potential
\[
 S_{\rm sp}(U)
 =\beta_{\rm sp}\sum_{p\in P_\Gamma}
 \left(1-\frac1n\operatorname{Re}\operatorname{Tr}U_p\right),
\]
one has
\begin{equation}
 R_\Gamma\le2\beta_{\rm sp}|P_\Gamma|.                  \label{eq:s7v26-extensive}
\end{equation}
Thus \eqref{eq:s7v26-perturb-condition} forces a coupling window
shrinking at least inversely with volume.  It proves only a
finite-volume neighbourhood of the electric point.  A
volume-uniform strong-coupling result requires a connected polymer or
cluster estimate, not global operator norm.
\end{firewall}

\begin{firewall}[Strong coupling is not the continuum trajectory]
\label{fire:s7v26-continuum}
Even a volume-uniform cluster expansion at sufficiently small
\(\beta_{\rm sp}\) would prove a lattice strong-coupling gap.  In
four-dimensional asymptotically free Yang--Mills theory, the isotropic
continuum trajectory has bare Wilson parameter tending to weak
coupling.  Reaching that trajectory requires a controlled
renormalization flow or a universality bridge; analytic continuation
from the strong-coupling region cannot be assumed.
\end{firewall}

\section{Transfer-contraction target}
\label{sec:s7v26-target}

\begin{definition}[Gauge-projected two-slab contraction card, GPTCC]
\label{def:s7v26-GPTCC}
GPTCC consists of:
\begin{enumerate}[label=\textup{(T\arabic*)},leftmargin=4em]
\item the kernel \eqref{eq:s7v26-kernel} with a declared temporal
      action and normalization;
\item positivity improvement, vacuum uniqueness, and the spin-network
      dense core;
\item a connected expansion for
      \(C_{\Gamma,F}(n)\), not for the extensive partition function
      alone;
\item a common bound
      \(C_{\Gamma,F}(n)\le C_Fe^{-m_*na_t}\) with \(m_*>0\)
      uniform in spatial volume and cutoff for each fixed local
      spin network \(F\);
\item boundary and winding polymers separated from local connected
      clusters;
\item a renormalization or universality bridge from any convergent
      strong-coupling expansion to the continuum scaling trajectory;
\item preservation of reflection positivity and convergence of the
      local observable algebra; and
\item the DOSGC handoff through
      \Cref{thm:s7v25-OS-gap}.
\end{enumerate}
\end{definition}

\begin{assessment}[Progress at seventh-series V26]
\label{ass:s7v26-status}
At fixed cutoff the physical transfer kernel, gauge projection,
strictly positive vacuum, dense spin-network core, connected
correlator, and exact mass eigenvalue ratio are now explicit.  The
electric point has a computable spin-network Casimir gap.  The
elementary spatial-potential perturbation is volume-extensive, exposing
the need for a connected expansion, and the strong-coupling region
still lies away from the four-dimensional continuum trajectory.
\end{assessment}

\begin{decision}[V26 bottleneck decision]
\label{dec:s7v26-bottleneck}
V27 will formulate a connected two-slab polymer expansion directly for
\(C_{\Gamma,F}(n)\), using the archived BKAR and connected-hypergraph
tools without reproducing their proofs.  The first target is a
volume-uniform strong-coupling decay bound for fixed local spin
networks, with temporal length and spatial volume counted separately.
The continuum weak-coupling bridge will remain an explicit later row.
\end{decision}

\begin{firewall}[Claim boundary after seventh-series V26]
\label{fire:s7v26-boundary}
V26 constructs and analyzes the fixed-cutoff gauge-projected transfer
operator, proves unique-vacuum and spin-network-core statements,
identifies the exact finite-cutoff mass ratio and electric gap, and
proves the limitation of the crude perturbation.

V26 does not prove a volume-uniform connected expansion, a
regulator-uniform decay exponent, a weak-coupling renormalization
bridge, continuum reconstruction, or the Yang--Mills mass gap.
\end{firewall}

\paragraph{Parent-version record.}
V26 was formed from
\texttt{Renewal\_Yang\_Mills\_Mass\_Gap\_Research\_Notes\_V25.tex}.
The V25 parent had \(10\,141\) logical source lines, \(370\,452\) bytes,
and SHA--256
\[
 \texttt{9004DB5B23280FB2C3FF80AE2F74A37EBC30F4C82A54F08904047AD92233FB10}.
\]

% =============================================================================
% END APPENDED SEVENTH-SERIES V26 MATERIAL
% =============================================================================

\newpage

% =============================================================================
% BEGIN APPENDED SEVENTH-SERIES V27 MATERIAL
% =============================================================================

\chapter{A volume-uniform strong-coupling correlator expansion}
\label{ch:s7v27}

V26 identifies the exact physical target at fixed cutoff:
volume-uniform exponential decay of connected spin-network
correlators.  A global operator-norm perturbation of the electric
transfer fails because the Wilson potential is extensive.  The proper
small-coupling object is instead the connected plaquette polymer.  This
note gives a self-contained, deliberately coarse convergence criterion
for that expansion and obtains an explicit lattice strong-coupling
mass gap.

The regulator in this chapter is the isotropic Wilson plaquette
regulator.  Its transfer kernel differs from the heat-kernel temporal
choice used for the exact electric calculation in V26.  Both are
reflection-positive lattice regulators, but relating their continuum
limits is a universality problem and is not assumed here.

\section{Normalized plaquette activities}
\label{sec:s7v27-activities}

Let \(\Lambda\) be a finite four-dimensional hypercubic box, with
arbitrary spatial size and either open temporal boundary or temporal
extent larger than every separation under consideration.  For
\(G=SU(n)\), define
\begin{equation}
 w_\beta(h)
 :=\exp\left\{
 {\beta\over n}\operatorname{Re}\operatorname{Tr}h
 \right\},\qquad
 c_0(\beta):=\int_Gw_\beta(h)\,dh.                      \label{eq:s7v27-w}
\end{equation}
Put
\begin{equation}
 \widehat w_\beta:=c_0(\beta)^{-1}w_\beta,\qquad
 r_\beta:=\widehat w_\beta-1,\qquad
 q_\beta:=e^{2\beta}-1.                                \label{eq:s7v27-r}
\end{equation}
The harmless factor \(c_0(\beta)^{|P_\Lambda|}\) cancels from all
normalized expectations.

\begin{lemma}[Uniform single-plaquette activity bound]
\label{lem:s7v27-q}
For every \(\beta\ge0\),
\begin{equation}
 \|r_\beta\|_\infty\le q_\beta.                         \label{eq:s7v27-q}
\end{equation}
\end{lemma}

\begin{proof}
Since
\(-1\le n^{-1}\operatorname{Re}\operatorname{Tr}h\le1\),
\[
 e^{-\beta}\le c_0(\beta)\le e^\beta.
\]
Hence
\[
 e^{-2\beta}\le\widehat w_\beta(h)\le e^{2\beta}.
\]
The larger of \(e^{2\beta}-1\) and \(1-e^{-2\beta}\) is
\(q_\beta\).
\end{proof}

\begin{lemma}[Nonnegative character coefficients]
\label{lem:s7v27-positive-type}
The central function \(w_\beta\) has a character expansion with
nonnegative coefficients.  Consequently the Wilson plaquette
regulator is link-reflection positive and admits a positive transfer
operator on the gauge-invariant Hilbert space.
\end{lemma}

\begin{proof}
If \(\chi_{\rm f}\) is the fundamental character, then
\[
 w_\beta
 =
 \exp\left\{
 {\beta\over2n}(\chi_{\rm f}+\chi_{\overline{\rm f}})
 \right\}.
\]
Expand both exponentials in power series.  Every coefficient is
nonnegative, and every product of characters is the character of a
tensor product, which decomposes into irreducibles with nonnegative
integer multiplicities.  Thus every irreducible character coefficient
is nonnegative.  For a plaquette crossing a reflection plane, split
its representation matrix across the plane; the character contraction
is a sum of a function times its reflected conjugate with the same
nonnegative coefficient.  Multiplication over plaquettes proves
reflection positivity and hence positivity of the transfer operator.
\end{proof}

\section{The plaquette polymer gas}
\label{sec:s7v27-polymers}

Two plaquettes are adjacent if they share a link.  Let \(\Delta_D\)
be the maximum degree of this adjacency graph in \(D\) Euclidean
dimensions.  The elementary incidence bound is
\begin{equation}
 \Delta_D\le4\bigl(2(D-1)-1\bigr)=8D-12,
 \qquad \Delta_4\le20.                                  \label{eq:s7v27-Delta}
\end{equation}
A polymer is a nonempty connected plaquette set \(\gamma\).  Its link
support is denoted \(E(\gamma)\), and its activity is
\begin{equation}
 \rho(\gamma)
 :=
 \int_{G^{E(\gamma)}}
 \prod_{p\in\gamma}r_\beta(U_p)\,dU_{E(\gamma)}.         \label{eq:s7v27-rho}
\end{equation}

\begin{lemma}[Polymer representation and activity majorant]
\label{lem:s7v27-polymer}
The normalized partition function is the hard-core polymer gas
\begin{equation}
 {\cal Z}_\Lambda
 =
 \sum_{\{\gamma_1,\ldots,\gamma_m\}\ {\rm compatible}}
 \prod_{j=1}^m\rho(\gamma_j),                           \label{eq:s7v27-polymer-Z}
\end{equation}
where compatible polymers have disjoint, nonadjacent plaquette
supports.  Moreover
\begin{equation}
 |\rho(\gamma)|\le q_\beta^{|\gamma|}.                  \label{eq:s7v27-rho-bound}
\end{equation}
\end{lemma}

\begin{proof}
Expand
\[
 \prod_{p\in P_\Lambda}(1+r_\beta(U_p))
\]
as a sum over plaquette subsets.  Every subset splits uniquely into
connected components in the plaquette adjacency graph.  Components
with no shared link depend on disjoint Haar variables, so the integral
factorizes into \eqref{eq:s7v27-rho}; this gives
\eqref{eq:s7v27-polymer-Z}.  The absolute-value bound follows from
\Cref{lem:s7v27-q} and normalized Haar measure.
\end{proof}

\begin{lemma}[Coarse connected-set count]
\label{lem:s7v27-count}
In a graph of maximum degree \(\Delta\ge2\), the number \(N_m(p)\) of
connected \(m\)-vertex sets containing a fixed vertex \(p\) obeys
\begin{equation}
 N_m(p)\le\Delta^{2(m-1)}.                              \label{eq:s7v27-count}
\end{equation}
\end{lemma}

\begin{proof}
Give the graph and every finite connected set a fixed deterministic
ordering.  Choose the corresponding canonical spanning tree rooted at
\(p\) and traverse it depth first.  The traversal is a graph walk of
length \(2(m-1)\), and its visited vertices recover the set.  There are
at most \(\Delta^{2(m-1)}\) such walks.
\end{proof}

\section{A tilted Koteck\'y--Preiss bound}
\label{sec:s7v27-KP}

Let \(\operatorname{diam}_t(\gamma)\) be the temporal diameter of a
polymer in lattice units.  Set \(\Delta=\Delta_4\).

\begin{theorem}[Explicit strong-coupling convergence condition]
\label{thm:s7v27-KP}
Assume
\begin{equation}
 0<q_\beta<{1\over2e\Delta^2}                           \label{eq:s7v27-small}
\end{equation}
and define
\begin{equation}
 b_\beta
 :=
 \log{1\over2e\Delta^2q_\beta}>0.                       \label{eq:s7v27-b}
\end{equation}
Then, uniformly in the finite box and its spatial volume,
\begin{equation}
 \sup_{p}
 \sum_{\gamma\ni p}
 |\rho(\gamma)|
 \exp\{(1+b_\beta)|\gamma|\}
 \le {1\over\Delta^2}
 \le {1\over\Delta+1}.                                  \label{eq:s7v27-KP}
\end{equation}
In particular the hard-core polymer logarithm and all fixed local
source derivatives converge absolutely.
\end{theorem}

\begin{proof}
The left side is at most
\begin{align*}
 \sum_{m\ge1}
 N_m(p)q_\beta^m e^{(1+b_\beta)m}
 &\le
 q_\beta e^{1+b_\beta}
 \sum_{m\ge1}
 \bigl(
 \Delta^2q_\beta e^{1+b_\beta}
 \bigr)^{m-1}.
\end{align*}
By \eqref{eq:s7v27-b},
\[
 \Delta^2q_\beta e^{1+b_\beta}={1\over2},
 \qquad
 q_\beta e^{1+b_\beta}={1\over2\Delta^2}.
\]
The geometric sum is therefore \(1/\Delta^2\).  A polymer incompatible
with a fixed \(\gamma_0\) contains a plaquette equal or adjacent to one
of the plaquettes of \(\gamma_0\).  There are at most
\((\Delta+1)|\gamma_0|\) choices.  Thus
\eqref{eq:s7v27-KP} implies the standard Koteck\'y--Preiss inequality
with size function \(|\gamma_0|\), including the temporal tilt.  The
usual inductive proof of the polymer logarithm then converges
absolutely; differentiating a finite number of local source parameters
only marks finitely many polymers and preserves convergence.
\end{proof}

\begin{remark}[A fully explicit but deliberately small interval]
\label{rem:s7v27-small}
Using \(\Delta_4=20\), condition \eqref{eq:s7v27-small} is
\[
 e^{2\beta}-1<{1\over800e}.
\]
The constant is not optimized.  Its role is to give a certified
nonempty strong-coupling interval independent of lattice volume.
\end{remark}

\section{Connected correlations must span the time separation}
\label{sec:s7v27-decay}

For bounded observables \(A,B\) depending on finite link sets, let
\(\operatorname{supp}A,\operatorname{supp}B\) denote the plaquettes
incident to those links.  Write
\[
 \langle A;B\rangle_\beta
 :=\langle AB\rangle_\beta
   -\langle A\rangle_\beta\langle B\rangle_\beta.
\]

\begin{theorem}[Volume-uniform strong-coupling clustering]
\label{thm:s7v27-clustering}
Under \eqref{eq:s7v27-small}, for every pair of fixed bounded local
observables there is a finite constant \(C_{A,B}(\beta)\), independent
of the spatial volume, such that
\begin{equation}
 |\langle A;B\rangle_\beta|
 \le
 C_{A,B}(\beta)
 e^{-b_\beta d_t(A,B)},                                 \label{eq:s7v27-clustering}
\end{equation}
where \(d_t(A,B)\) is their temporal support separation in lattice
units.
\end{theorem}

\begin{proof}
Introduce source parameters \(s,t\) by multiplying the measure by
\(e^{sA+tB}\).  The connected correlation is
\[
 \partial_s\partial_t\log{\cal Z}_\Lambda(s,t)
 \big|_{s=t=0}.
\]
The absolutely convergent sourced polymer logarithm from
\Cref{thm:s7v27-KP} expresses this derivative as a sum of connected
clusters containing both marked supports.  All clusters meeting only
one marked support cancel in the logarithmic mixed derivative.

Every surviving connected Ursell cluster has a union joining the two
marked supports.  The total number of plaquettes in its polymers is at
least \(d_t(A,B)\).  Multiply the activity of every polymer
\(\gamma\) by \(e^{b_\beta|\gamma|}\), sum using the tilted
Koteck\'y--Preiss majorant \eqref{eq:s7v27-KP}, and then remove the
tilt.  The product of the removed factors contributes at most
 \(e^{-b_\beta d_t(A,B)}\).  The remaining rooted sums meet only the
fixed finite marked supports and are bounded by an absolutely
convergent constant depending on \(A,B,\beta\), but not on the box.
This proves \eqref{eq:s7v27-clustering}.
\end{proof}

\begin{corollary}[Infinite-volume uniqueness and local convergence]
\label{cor:s7v27-limit}
Under \eqref{eq:s7v27-small}, local expectations have a unique
infinite-volume limit, independent of boundary conditions whose
distance from the observable tends to infinity.  The limiting
connected correlations obey \eqref{eq:s7v27-clustering}.
\end{corollary}

\begin{proof}
The difference between two boundary conditions has a connected
polymer representation in which a cluster joins the fixed observable
support to the boundary.  The tilted bound makes this difference
exponentially small in their distance.  It is therefore Cauchy and
boundary independent.  Pass to the limit in
\eqref{eq:s7v27-clustering}.
\end{proof}

\section{The lattice strong-coupling mass gap}
\label{sec:s7v27-gap}

\begin{theorem}[Strong-coupling Yang--Mills lattice gap]
\label{thm:s7v27-gap}
Assume \eqref{eq:s7v27-small}.  The infinite-volume
reflection-positive Wilson lattice state has a unique vacuum and, on
the gauge-invariant local spin-network core, satisfies
\begin{equation}
 0\le
 \omega_\beta\!\left(
 (\Theta A^\circ)\tau_nA^\circ
 \right)
 \le C_A(\beta)e^{-b_\beta n}.                          \label{eq:s7v27-OS-decay}
\end{equation}
Consequently its reconstructed lattice Hamiltonian has
\begin{equation}
 m_{\rm latt}\ge b_\beta.                               \label{eq:s7v27-lattice-gap}
\end{equation}
For isotropic lattice spacing \(a\), this is the physical-unit bound
\begin{equation}
 m_{\rm phys}\ge {b_\beta\over a}.                      \label{eq:s7v27-physical-gap}
\end{equation}
\end{theorem}

\begin{proof}
Reflection positivity follows from
\Cref{lem:s7v27-positive-type}.  The polymer uniqueness statement in
\Cref{cor:s7v27-limit} gives the unique translation-invariant vacuum
state in this strong-coupling regime.  Apply
\Cref{thm:s7v27-clustering} with
\(B=\tau_nA^\circ\) and the reflected first insertion.  The finite
diameter of the fixed observable changes \(d_t\) from \(n\) by only an
observable-dependent constant, which is absorbed into \(C_A(\beta)\).
This proves \eqref{eq:s7v27-OS-decay}.  The spin-network core is dense
by the same Peter--Weyl and OS cyclicity argument as
\Cref{thm:s7v26-spin-density}.  Apply
\Cref{thm:s7v25-OS-gap}; conversion from lattice steps to physical time
gives \eqref{eq:s7v27-physical-gap}.
\end{proof}

\begin{firewall}[This is a lattice strong-coupling theorem]
\label{fire:s7v27-regime}
The hypothesis is \(q_\beta=e^{2\beta}-1\ll1\), hence
\(\beta\ll1\).  Four-dimensional continuum Yang--Mills requires an
asymptotically free weak-bare-coupling trajectory, conventionally
\(\beta(a)\to\infty\) as \(a\downarrow0\).  The bound
\eqref{eq:s7v27-physical-gap} at fixed small \(\beta\) therefore does
not define the desired continuum mass.  It closes the
volume-uniform strong-coupling row only.
\end{firewall}

\begin{firewall}[No analytic continuation across a phase diagram]
\label{fire:s7v27-analytic}
Analyticity of local expectations in the small-\(\beta\) polymer disk
does not prove that the gap remains open to \(\beta=\infty\).
Zeros, loss of convergence, or a vanishing spectral edge may occur
outside that disk.  A continuum conclusion requires a controlled
renormalization or universality bridge with a quantitative physical
decay rate.
\end{firewall}

\section{Strong-coupling row of the direct programme}
\label{sec:s7v27-programme}

\begin{theorem}[GPTCC strong-coupling insertion]
\label{thm:s7v27-card}
For the Wilson regulator under \eqref{eq:s7v27-small}, the following
GPTCC rows are closed uniformly in volume:
\begin{enumerate}[label=\textup{(\arabic*)},leftmargin=3.5em]
\item reflection positivity and a positive gauge-invariant transfer;
\item a connected, rather than extensive, plaquette expansion;
\item absolute polymer convergence and boundary-condition uniqueness;
\item exponential connected decay for every fixed local spin network;
      and
\item the dense-core lattice Hamiltonian gap
      \(m_{\rm latt}\ge b_\beta\).
\end{enumerate}
The cutoff-uniform continuum row is not closed.
\end{theorem}

\begin{proof}
Use
\Cref{lem:s7v27-positive-type,thm:s7v27-KP,
cor:s7v27-limit,thm:s7v27-clustering,thm:s7v27-gap}.
\end{proof}

\begin{assessment}[Progress at seventh-series V27]
\label{ass:s7v27-status}
The global-volume loss in V26 has been removed in a certified
strong-coupling interval.  A tilted polymer expansion gives a common
lattice decay exponent
\[
 b_\beta=\log\!\left({1\over2e\Delta_4^2(e^{2\beta}-1)}\right)>0
\]
for the full local gauge-invariant spin-network core, and hence a
genuine infinite-volume lattice Hamiltonian gap.  The central
bottleneck is now sharply separated: propagate a positive
\emph{physical} decay scale from the asymptotically free ultraviolet
trajectory to a nonperturbative coarse scale.
\end{assessment}

\begin{decision}[V27 bottleneck decision]
\label{dec:s7v27-bottleneck}
V28 will formulate an exact block-spin/transfer matching identity for
local reflected correlators and a quantitative renormalization bridge
card.  It will compute the number of blocking steps required by the
one-loop asymptotically free running and state the precise stability
estimate needed to hand the ultraviolet theory to a convergent
strong-coupling or other massive coarse measure.  Perturbative running
alone will not be counted as the nonperturbative handoff.
\end{decision}

\begin{firewall}[Claim boundary after seventh-series V27]
\label{fire:s7v27-boundary}
V27 proves an explicit, volume-uniform strong-coupling polymer
convergence criterion, exponential connected spin-network clustering,
infinite-volume uniqueness, and a positive Wilson-lattice Hamiltonian
gap in that regime.

V27 does not connect the strong-coupling region to the asymptotically
free continuum trajectory, prove a cutoff-uniform finite physical
mass, construct the continuum measure, or prove the full continuum
Yang--Mills mass gap.
\end{firewall}

\paragraph{Parent-version record.}
V27 was formed from
\texttt{Renewal\_Yang\_Mills\_Mass\_Gap\_Research\_Notes\_V26.tex}.
The V26 parent had \(10\,629\) logical source lines, \(387\,781\) bytes,
and SHA--256
\[
 \texttt{99D7E7C896B0C54977499EE33D5401502495417BB879C2741D17DEAD4EFC7A42}.
\]

% =============================================================================
% END APPENDED SEVENTH-SERIES V27 MATERIAL
% =============================================================================

\newpage

% =============================================================================
% BEGIN APPENDED SEVENTH-SERIES V28 MATERIAL
% =============================================================================

\chapter{Exact blocking and the quantitative continuum bridge}
\label{ch:s7v28}

V27 proves a genuine infinite-volume lattice gap in a certified
strong-coupling interval.  The remaining task is not analytic
continuation in the bare Wilson parameter.  It is to transport
physical correlators, reflection positivity, and a spectral
contraction through a sequence of gauge-covariant blocking steps from
the asymptotically free ultraviolet theory to a nonperturbative coarse
theory.  This note separates the exact kinematic part of that statement
from the missing dynamical estimate.

\section{A gauge-covariant coarse-link map}
\label{sec:s7v28-block-map}

Let the fine lattice have spacing \(a\), and let \(b\ge2\) divide every
period.  The coarse vertices are the sublattice \(b\mathbb Z^4\).  For
a coarse oriented edge \(E=(x,x+b\widehat\mu)\), define
\begin{equation}
 ({\cal B}_bU)_E
 :=
 U_{x,\mu}U_{x+\widehat\mu,\mu}\cdots
 U_{x+(b-1)\widehat\mu,\mu}.                            \label{eq:s7v28-B}
\end{equation}

\begin{lemma}[Gauge covariance of path-product blocking]
\label{lem:s7v28-covariance}
Under the fine gauge action,
\begin{equation}
 {\cal B}_b(g\cdot U)_E
 =
 g_x({\cal B}_bU)_Eg_{x+b\widehat\mu}^{-1}.             \label{eq:s7v28-covariance}
\end{equation}
Consequently the pushforward of any fine gauge-invariant measure by
\({\cal B}_b\) is a coarse gauge-invariant measure.
\end{lemma}

\begin{proof}
Insert \eqref{eq:s7v26-gauge} in every factor of
\eqref{eq:s7v28-B}.  The gauge matrices at all \(b-1\) interior
vertices cancel pairwise, leaving only the two endpoint matrices.
Gauge invariance of the pushforward follows by change of variables.
\end{proof}

\begin{lemma}[Reflection and translation compatibility]
\label{lem:s7v28-reflection}
If the blocking grid is aligned with a lattice reflection plane and
with translations by \(b\) fine steps, then
\begin{equation}
 {\cal B}_b\Theta=\Theta{\cal B}_b,\qquad
 {\cal B}_b\tau_{b\ell}=\tau_\ell{\cal B}_b.             \label{eq:s7v28-reflection}
\end{equation}
\end{lemma}

\begin{proof}
Reflection reverses every crossed path and inverts its ordered
product, exactly as it reverses the corresponding coarse link.
Translation by \(b\ell\) maps each blocking path to the path defining
the translated coarse link.
\end{proof}

Let \(\mu_a\) be a fine regulated Euclidean measure and define
\begin{equation}
 \mu_{ba}:=({\cal B}_b)_*\mu_a.                          \label{eq:s7v28-push}
\end{equation}

\begin{theorem}[Exact preservation of coarse correlators]
\label{thm:s7v28-correlator}
For bounded coarse observables \(F_1,\ldots,F_m\),
\begin{equation}
 \int\prod_{j=1}^mF_j(V)\,d\mu_{ba}(V)
 =
 \int\prod_{j=1}^mF_j({\cal B}_bU)\,d\mu_a(U).           \label{eq:s7v28-correlator}
\end{equation}
In particular, for reflected observables and time separations divisible
by \(b\), their Osterwalder--Schrader correlators agree exactly after
rescaling the time coordinate by \(b\).
\end{theorem}

\begin{proof}
Equation \eqref{eq:s7v28-correlator} is the definition of pushforward
measure.  Apply \Cref{lem:s7v28-reflection} to the reflected and
translated insertions.
\end{proof}

\section{The effective action is not a one-coupling Wilson action}
\label{sec:s7v28-effective}

Disintegrate \(\mu_a\) over \({\cal B}_b\).  Relative to a declared
coarse reference measure \(dV\), define the exact effective action by
\begin{equation}
 e^{-S_{ba}^{\rm eff}(V)}
 :=
 {d({\cal B}_b)_*\mu_a\over dV}(V).                     \label{eq:s7v28-effective}
\end{equation}

\begin{proposition}[Gauge-invariant effective-action ledger]
\label{prop:s7v28-action-ledger}
The action \(S_{ba}^{\rm eff}\) is coarse gauge invariant and
reflection invariant whenever \(\mu_a\) has those symmetries.  Its
complete Peter--Weyl basis includes every allowed coarse spin-network
interaction, and the coefficients generated in general are not
confined to fundamental plaquettes: higher representations, larger
Wilson loops, multi-loop couplings, and nonlocal terms are allowed.
Exact blocking gives no
identity of the form
\[
 S_{ba}^{\rm eff}=S_{\rm Wilson}(\beta')
\]
with a single renormalized coupling \(\beta'\).
\end{proposition}

\begin{proof}
The symmetry statement follows from
\Cref{lem:s7v28-covariance,lem:s7v28-reflection} and uniqueness of the
Radon--Nikodym derivative.  Gauge-invariant Peter--Weyl functions form
a complete basis on the coarse quotient by
\Cref{thm:s7v26-spin-density}, so the logarithm of the pushforward
density has an expansion in all such interactions.  No symmetry
annihilates the larger-loop or multi-loop coefficients.  Integrating
fine links coupled to more than one plaquette produces them already in
the ordinary cumulant expansion.  Therefore closure in the
one-parameter Wilson family would be an additional theorem, not a
blocking identity.
\end{proof}

\begin{firewall}[Coupling recursion alone is not exact RG]
\label{fire:s7v28-one-coupling}
A scalar recursion \(\beta\mapsto\beta'\) obtained by projecting
\(S^{\rm eff}\) onto the fundamental plaquette character discards the
remaining interaction coordinates.  To use it in a mass-gap proof one
must bound the discarded action in a norm which controls reflected
connected correlators and transfer spectra at every step.
\end{firewall}

\section{Coarse observables and the unresolved fine complement}
\label{sec:s7v28-complement}

Let \({\cal F}_b=\sigma({\cal B}_b)\) and let
\[
 E_bA:=\mathbb E_{\mu_a}[A\mid{\cal F}_b],\qquad
 R_bA:=A-E_bA.
\]
There is a coarse function \(A_b\) with
\(E_bA=A_b\circ{\cal B}_b\).

\begin{lemma}[Exact one-step correlation decomposition]
\label{lem:s7v28-decomposition}
For every pair \(A,B\in L^2(\mu_a)\),
\begin{align}
 \langle A;B\rangle_{\mu_a}
 &=
 \langle A_b;B_b\rangle_{\mu_{ba}}
 +\langle E_bA;R_bB\rangle_{\mu_a}
 \notag\\
 &\quad
 +\langle R_bA;E_bB\rangle_{\mu_a}
 +\langle R_bA;R_bB\rangle_{\mu_a}.                    \label{eq:s7v28-decomposition}
\end{align}
\end{lemma}

\begin{proof}
Insert \(A=E_bA+R_bA\) and \(B=E_bB+R_bB\) in the covariance and
expand.  The first term is the coarse covariance by
\Cref{thm:s7v28-correlator}.  The other three terms are retained
because conditional orthogonality at one time does not automatically
annihilate a translated residual at another time.
\end{proof}

\begin{warning}[Exact coarse matching controls only a subspace]
\label{warn:s7v28-subspace}
Equation \eqref{eq:s7v28-correlator} transfers decay for observables
measurable with respect to \({\cal F}_b\).  A full mass gap requires a
dense fine algebra, so the three residual terms in
\eqref{eq:s7v28-decomposition} need their own contraction.  Ignoring
them proves at most a coarse-observable gap.
\end{warning}

\section{An abstract multiscale transfer compiler}
\label{sec:s7v28-compiler}

The exact operator requirement can be stated without choosing an RG
scheme.  Let \(a_{j+1}=b_ja_j\), with integers \(b_j\ge2\).  On Hilbert
spaces \({\cal H}_j\), let \(T_j\) be positive self-adjoint
contractions with simple vacuum vectors \(\Omega_j\) and
\(T_j\Omega_j=\Omega_j\).  Let
\[
 Q_j:{\cal H}_{j+1}\longrightarrow{\cal H}_j
\]
be an isometry with \(Q_j\Omega_{j+1}=\Omega_j\), and put
\(P_j=Q_jQ_j^*\).

\begin{theorem}[Approximate block-transfer gap recursion]
\label{thm:s7v28-recursion}
Assume for \(A_j:=T_j^{b_j}\) that
\begin{align}
 \|P_jA_jP_j-Q_jT_{j+1}Q_j^*\|
 &\le\epsilon_j,                                       \label{eq:s7v28-coarse-error}\\
 \|P_j^\perp A_jP_j^\perp\|
 &\le r_j,                                              \label{eq:s7v28-residual}\\
 \|P_jA_jP_j^\perp\|
 &\le\epsilon_j.                                       \label{eq:s7v28-offdiag}
\end{align}
Let
\[
 q_J:=\|T_J|_{\Omega_J^\perp}\|
\]
and recursively define
\begin{equation}
 q_j
 :=
 \left[
 \max\{q_{j+1}+2\epsilon_j,\ r_j+\epsilon_j\}
 \right]^{1/b_j}.                                      \label{eq:s7v28-q-recursion}
\end{equation}
If every bracket in \eqref{eq:s7v28-q-recursion} is smaller than one,
then
\begin{equation}
 \|T_0|_{\Omega_0^\perp}\|\le q_0,\qquad
 \operatorname{gap}(H_0)
 \ge-{1\over a_0}\log q_0.                             \label{eq:s7v28-gap-recursion}
\end{equation}
\end{theorem}

\begin{proof}
Decompose
\(\Omega_j^\perp\) into
\[
 \bigl(P_j{\cal H}_j\cap\Omega_j^\perp\bigr)
 \oplus P_j^\perp{\cal H}_j.
\]
On the first summand, \eqref{eq:s7v28-coarse-error} and the inductive
coarse bound give diagonal norm at most
\(q_{j+1}+\epsilon_j\).  On the second, the diagonal norm is at most
\(r_j\).  The off-diagonal block has norm at most \(\epsilon_j\).
For a unit vector \(x\oplus y\), self-adjointness and
\(2\|x\|\|y\|\le\|x\|^2+\|y\|^2\) give
\[
 \langle x\oplus y,A_j(x\oplus y)\rangle
 \le
 \max\{q_{j+1}+2\epsilon_j,r_j+\epsilon_j\}.
\]
Since \(A_j\ge0\), this bounds its norm on the nonvacuum space.
Spectral calculus for \(A_j=T_j^{b_j}\) takes the \(b_j\)-th root and
gives \eqref{eq:s7v28-q-recursion}.  Iterate to \(j=0\), then use
\(H_0=-a_0^{-1}\log T_0\).
\end{proof}

\begin{corollary}[Exact uniform-mass handoff]
\label{cor:s7v28-exact}
If \(\epsilon_j=0\) and
\begin{equation}
 q_J\le e^{-m_*a_J},\qquad
 r_j\le e^{-m_*a_{j+1}}                                \label{eq:s7v28-exact-hyp}
\end{equation}
for one \(m_*>0\), then
\begin{equation}
 \operatorname{gap}(H_0)\ge m_*                         \label{eq:s7v28-exact-gap}
\end{equation}
independently of the number of blocking steps.
\end{corollary}

\begin{proof}
Backward induction in \eqref{eq:s7v28-q-recursion} gives
\[
 q_j\le
 (e^{-m_*a_{j+1}})^{1/b_j}
 =e^{-m_*a_j}.
\]
Use \eqref{eq:s7v28-gap-recursion}.
\end{proof}

\begin{corollary}[Error budget needed for an approximate handoff]
\label{cor:s7v28-error-budget}
Fix \(m_*>\widetilde m>0\).  A sufficient stepwise condition for
\(\operatorname{gap}(H_0)\ge\widetilde m\) is
\begin{align}
 q_J&\le e^{-\widetilde m a_J},                         \label{eq:s7v28-final}\\
 q_{j+1}+2\epsilon_j
 &\le e^{-\widetilde m a_{j+1}},\qquad
 r_j+\epsilon_j
 \le e^{-\widetilde m a_{j+1}}.                        \label{eq:s7v28-step-budget}
\end{align}
Thus the RG error must fit inside the physical exponential margin at
every scale; convergence of couplings without
\eqref{eq:s7v28-step-budget} is insufficient.
\end{corollary}

\begin{proof}
Insert \eqref{eq:s7v28-step-budget} in
\eqref{eq:s7v28-q-recursion} and argue backward as in
\Cref{cor:s7v28-exact}.
\end{proof}

\section{How many ultraviolet steps must be controlled}
\label{sec:s7v28-running}

For pure \(SU(N)\) Yang--Mills, write the one-loop ultraviolet beta
function in the convention
\begin{equation}
 \mu{dg\over d\mu}
 =
 -b_0g^3+O(g^5),\qquad
 b_0={11N\over48\pi^2}.                                \label{eq:s7v28-beta}
\end{equation}
Ignoring only the displayed higher-order remainder gives
\begin{equation}
 {1\over g(\mu)^2}
 =
 2b_0\log{\mu\over\Lambda}.                             \label{eq:s7v28-running}
\end{equation}

\begin{proposition}[One-loop blocking count]
\label{prop:s7v28-step-count}
Start at ultraviolet spacing \(a\) and block by a fixed factor \(b\).
The scale \(\mu_J=(ab^J)^{-1}\) at which the one-loop coupling reaches
a declared value \(g_*\) obeys
\begin{equation}
 J(a;g_*)
 =
 {\log(1/(a\Lambda))\over\log b}
 -{1\over2b_0g_*^2\log b}
 +O(1),                                                 \label{eq:s7v28-step-count}
\end{equation}
where the \(O(1)\) records integer rounding.  The corresponding coarse
physical spacing is
\begin{equation}
 a_J
 \simeq
 \Lambda^{-1}
 \exp\left\{-{1\over2b_0g_*^2}\right\}.                 \label{eq:s7v28-coarse-scale}
\end{equation}
Thus a continuum construction must control
\(\Theta(\log(1/a))\) RG steps while preserving a fixed physical
spectral margin.
\end{proposition}

\begin{proof}
Set \(g(\mu_J)=g_*\) in \eqref{eq:s7v28-running}:
\[
 \log{\mu_J\over\Lambda}
 ={1\over2b_0g_*^2}.
\]
Use \(\mu_J=(ab^J)^{-1}\), solve for \(J\), and invert \(\mu_J\) to
obtain \eqref{eq:s7v28-coarse-scale}.
\end{proof}

\begin{firewall}[One-loop running stops before the handoff]
\label{fire:s7v28-perturbative}
Equation \eqref{eq:s7v28-running} is controlled only while \(g\) is
small.  The V27 polymer condition is a large-\(g\), small-Wilson-\(\beta\)
condition.  Extrapolating the one-loop formula to that region does not
prove \eqref{eq:s7v28-coarse-error}--\eqref{eq:s7v28-residual}, does
not control the generated effective interactions, and is not the
nonperturbative RG bridge.
\end{firewall}

\begin{firewall}[Dimensional transmutation supplies a scale, not a
theorem]
\label{fire:s7v28-transmutation}
The appearance of \(\Lambda\) and the fixed coarse length in
\eqref{eq:s7v28-coarse-scale} explains how a physical mass scale can
survive \(a\downarrow0\).  It does not prove positivity of that mass.
The missing input is a uniform transfer contraction at or before the
nonperturbative scale, with the accumulated errors controlled by
\eqref{eq:s7v28-step-budget}.
\end{firewall}

\section{Renormalization bridge card}
\label{sec:s7v28-card}

\begin{definition}[Nonperturbative transfer RG card, NTRGC]
\label{def:s7v28-NTRGC}
NTRGC consists of:
\begin{enumerate}[label=\textup{(R\arabic*)},leftmargin=4em]
\item gauge-covariant, reflection-compatible blocking maps such as
      \eqref{eq:s7v28-B};
\item exact pushforward effective actions retaining every generated
      spin-network interaction;
\item Hilbert-space isometries \(Q_j\) matching the coarse OS spaces
      to fine coarse-observable subspaces;
\item the coarse transfer error \(\epsilon_j\), off-diagonal error, and
      fine-complement contraction \(r_j\);
\item the stepwise physical-margin inequalities
      \eqref{eq:s7v28-step-budget} through
      \(J(a)=\Theta(\log(1/a))\) steps;
\item a terminal coarse theory with a proven physical transfer gap,
      whether by V27 polymers or another nonperturbative method;
\item uniform control of local observable representatives and their
      residuals as in \eqref{eq:s7v28-decomposition}; and
\item convergence, vacuum uniqueness, reflection positivity, and the
      DOSGC handoff.
\end{enumerate}
\end{definition}

\begin{theorem}[Exact statement of the continuum bridge bottleneck]
\label{thm:s7v28-bottleneck}
Given NTRGC rows (R1)--(R6) with a common
\(\widetilde m>0\), the ultraviolet regulated Hamiltonians have gap at
least \(\widetilde m\) by
\Cref{thm:s7v28-recursion,cor:s7v28-error-budget}.  With rows
(R7)--(R8), \Cref{cor:s7v25-regulator} transfers that gap to the
continuum Hamiltonian.  The currently unproved mathematical content is
precisely rows (R2)--(R6) along the asymptotically free trajectory, not
the kinematic block identity or the final spectral implication.
\end{theorem}

\begin{proof}
Rows (R3)--(R6) supply the hypotheses of
\Cref{thm:s7v28-recursion} at every scale, and
\Cref{cor:s7v28-error-budget} gives the uniform regulated gap.
Rows (R7)--(R8) are the convergence and density hypotheses of
\Cref{cor:s7v25-regulator}.
\end{proof}

\begin{assessment}[Progress at seventh-series V28]
\label{ass:s7v28-status}
The strong-coupling result is now connected to the ultraviolet problem
by an exact, falsifiable interface.  Path-product blocking preserves
gauge covariance, reflection, and coarse correlators, but generates a
full spin-network effective action.  The transfer recursion shows that
every scale requires both coarse matching and contraction of the
discarded fine complement.  One-loop asymptotic freedom predicts
\(\Theta(\log(1/a))\) steps to the scale \(\Lambda^{-1}\), but supplies
neither of those nonperturbative operator estimates.
\end{assessment}

\begin{decision}[V28 bottleneck decision]
\label{dec:s7v28-bottleneck}
V29 will attack the first NTRGC step at weak coupling in a finite
small-field chart.  It will derive the exact Gaussian block
decomposition into coarse and wavelet curvature modes and compute the
fine-complement transfer contraction.  The nonlinear Jacobian,
large-field probability, and accumulation over
\(\Theta(\log(1/a))\) steps will remain separate rather than being
absorbed into perturbative notation.
\end{decision}

\begin{firewall}[Claim boundary after seventh-series V28]
\label{fire:s7v28-boundary}
V28 proves exact gauge/reflection-compatible coarse-link matching,
exposes the complete effective-action ledger, proves an approximate
multiscale transfer-gap recursion, and computes the one-loop number
and physical scale of the required blocking steps.

V28 does not construct the nonperturbative effective-action flow,
prove fine-complement contraction and summable RG errors, reach the V27
strong-coupling region from the continuum trajectory, construct the
continuum state, or prove the full continuum Yang--Mills mass gap.
\end{firewall}

\paragraph{Parent-version record.}
V28 was formed from
\texttt{Renewal\_Yang\_Mills\_Mass\_Gap\_Research\_Notes\_V27.tex}.
The V27 parent had \(11\,068\) logical source lines, \(404\,167\) bytes,
and SHA--256
\[
 \texttt{42EB3B0FAFFEE7F8437C3854B85F0C7B7233169656BD6BF15A232D568D92A29F}.
\]

% =============================================================================
% END APPENDED SEVENTH-SERIES V28 MATERIAL
% =============================================================================

\newpage

% =============================================================================
% BEGIN APPENDED SEVENTH-SERIES V29 MATERIAL
% =============================================================================

\chapter{Exact Gaussian RG splitting and fine-complement contraction}
\label{ch:s7v29}

V28 isolates two operator estimates at every blocking step: matching
on the retained coarse subspace and contraction on its orthogonal fine
complement.  This note computes both exactly for the quadratic
transverse gauge theory.  The result is diagnostically important.  A
scale-independent contraction of discarded ultraviolet modes is
available; what remains massless is the retained infrared theory.
Thus perturbative high-mode elimination alone cannot create the
terminal mass required by NTRGC.

\section{Quadratic transverse lattice Hamiltonian}
\label{sec:s7v29-free}

Work on a periodic three-dimensional spatial lattice of spacing \(a\).
After linear gauge reduction, omit the harmonic \(k=0\) sector and let
\({\cal K}_a^\times\) be the nonzero spatial momenta in the Brillouin
zone.  For every transverse polarization and Lie-algebra colour, the
normal-ordered quadratic Hamiltonian has oscillator frequency
\begin{equation}
 \omega_a(k)
 :=
 {2\over a}
 \left(
 \sum_{i=1}^3\sin^2{ak_i\over2}
 \right)^{1/2}.                                        \label{eq:s7v29-omega}
\end{equation}
Thus, on the bosonic Fock space,
\begin{equation}
 H_a^{(0)}
 =
 \sum_{\substack{k\in{\cal K}_a^\times\\
                  \lambda,\mathsf c}}
 \omega_a(k)
 N_{k,\lambda,\mathsf c}.                              \label{eq:s7v29-H}
\end{equation}

\begin{firewall}[The zero mode is not discarded silently]
\label{fire:s7v29-zero}
Constant harmonic link modes have zero quadratic curvature and are not
contained in \eqref{eq:s7v29-H}.  Their compact toron, gauge, and
topological treatment remains a separate NTRGC row.  V29 concerns the
nonzero transverse oscillator sector only.
\end{firewall}

\section{Exact low/high factorization}
\label{sec:s7v29-split}

Fix an integer block factor \(b\ge2\).  Define
\begin{align}
 {\cal K}_{\rm lo}
 &:=
 \{k\in{\cal K}_a^\times:
   |ak_i|\le\pi/b\ \hbox{for all }i\},                  \label{eq:s7v29-low}\\
 {\cal K}_{\rm hi}
 &:={\cal K}_a^\times\setminus{\cal K}_{\rm lo}.        \label{eq:s7v29-high}
\end{align}
Let \({\cal F}_{\rm lo}\) and \({\cal F}_{\rm hi}\) be the Fock spaces
over the corresponding one-particle subspaces.

\begin{theorem}[Exact Gaussian tensor factorization]
\label{thm:s7v29-factor}
There are canonical identifications
\begin{equation}
 {\cal F}_a
 \cong{\cal F}_{\rm lo}\otimes{\cal F}_{\rm hi},\qquad
 H_a^{(0)}
 =
 H_{\rm lo}^{(0)}\otimes I+I\otimes H_{\rm hi}^{(0)}.    \label{eq:s7v29-factor}
\end{equation}
The vacuum factorizes as
\(\Omega_a=\Omega_{\rm lo}\otimes\Omega_{\rm hi}\).
\end{theorem}

\begin{proof}
The one-particle Hilbert space is the orthogonal direct sum of its
low- and high-momentum coordinate subspaces.  Bosonic second
quantization sends an orthogonal direct sum to the tensor product of
Fock spaces and sends the diagonal one-particle frequency operator to
the sum in \eqref{eq:s7v29-factor}.  The zero-particle vector
factorizes.
\end{proof}

Let \(a'=ba\).  The retained low-mode covariance defines a
\emph{perfect Gaussian coarse action}: it is the exact pushforward
under Fourier restriction, without replacing its dispersion by a
nearest-neighbour ansatz.  Let
\[
 Q_b:{\cal F}_{a'}^{\rm perf}\longrightarrow{\cal F}_a
\]
embed the coarse Fock state as that low-mode state tensored with
\(\Omega_{\rm hi}\).

\begin{theorem}[Exact coarse transfer matching]
\label{thm:s7v29-matching}
The map \(Q_b\) is an isometry, maps vacuum to vacuum, and obeys
\begin{equation}
 e^{-baH_a^{(0)}}Q_b
 =
 Q_b e^{-a'H_{a'}^{\rm perf}}.                          \label{eq:s7v29-matching}
\end{equation}
Hence the Gaussian coarse matching and off-diagonal errors in
\eqref{eq:s7v28-coarse-error} and
\eqref{eq:s7v28-offdiag} are exactly zero.
\end{theorem}

\begin{proof}
By definition the perfect coarse Hamiltonian is the restriction
\(H_{\rm lo}^{(0)}\) written in coarse variables.  Under the
factorization \eqref{eq:s7v29-factor},
\[
 Q_b\psi=\psi\otimes\Omega_{\rm hi}.
\]
Since \(H_{\rm hi}^{(0)}\Omega_{\rm hi}=0\), exponentiating
\eqref{eq:s7v29-factor} for time \(ba=a'\) gives
\eqref{eq:s7v29-matching}.  Both the retained range and its orthogonal
complement are invariant, so the off-diagonal block vanishes.
\end{proof}

\begin{warning}[The perfect action is spatially nonlocal]
\label{warn:s7v29-perfect}
Sharp Fourier restriction gives an exact benchmark, not a local block
map.  The perfect coarse covariance has long-range spatial couplings.
A local gauge-covariant RG must approximate this split and prove the
errors required in \eqref{eq:s7v28-step-budget}.
\end{warning}

\section{Uniform contraction of the discarded sector}
\label{sec:s7v29-contraction}

\begin{lemma}[High-mode frequency floor]
\label{lem:s7v29-floor}
Every \(k\in{\cal K}_{\rm hi}\) obeys
\begin{equation}
 \omega_a(k)
 \ge {2\over a}\sin{\pi\over2b}.                        \label{eq:s7v29-floor}
\end{equation}
\end{lemma}

\begin{proof}
For a high momentum, some coordinate satisfies
\(|ak_i|>\pi/b\).  On the Brillouin interval
\(|ak_i|\le\pi\), the sine in \eqref{eq:s7v29-omega} is
nonnegative and increasing up to the needed lower endpoint.  Keeping
only that coordinate gives \eqref{eq:s7v29-floor}.
\end{proof}

\begin{theorem}[Scale-independent fine-complement contraction]
\label{thm:s7v29-contraction}
Let \(P_b=Q_bQ_b^*\).  On \(P_b^\perp{\cal F}_a\), every state has at
least one high-mode quantum, and
\begin{equation}
 \left\|
 P_b^\perp e^{-baH_a^{(0)}}P_b^\perp
 \right\|
 \le
 r_b^{(0)}
 :=
 \exp\left\{
 -2b\sin{\pi\over2b}
 \right\}<1.                                            \label{eq:s7v29-contraction}
\end{equation}
For \(b=2\),
\begin{equation}
 r_2^{(0)}=e^{-2\sqrt2}.                                \label{eq:s7v29-b2}
\end{equation}
The bound is independent of \(a\), spatial volume, colour
multiplicity, and the number of prior blocking steps.
\end{theorem}

\begin{proof}
The complement of
\({\cal F}_{\rm lo}\otimes\Omega_{\rm hi}\) contains at least one
high oscillator excitation.  Its energy is at least the one-particle
floor in \eqref{eq:s7v29-floor}.  Evolution for time \(ba\) therefore
has norm at most
\[
 \exp\left\{
 -ba\,{2\over a}\sin{\pi\over2b}
 \right\},
\]
which is \eqref{eq:s7v29-contraction}.  Substitution of \(b=2\) gives
\eqref{eq:s7v29-b2}.
\end{proof}

\begin{corollary}[The Gaussian step closes the V28 operator rows]
\label{cor:s7v29-compiler}
In the nonzero transverse Gaussian sector, one RG step satisfies
\begin{equation}
 \epsilon_j=0,\qquad
 r_j\le r_{b_j}^{(0)}<1.                                \label{eq:s7v29-rows}
\end{equation}
For every fixed proposed physical mass
\(\widetilde m\), the residual inequality
\[
 r_j\le e^{-\widetilde m a_{j+1}}
\]
holds automatically at all sufficiently ultraviolet scales.
\end{corollary}

\begin{proof}
The first statement combines
\Cref{thm:s7v29-matching,thm:s7v29-contraction}.  The left side of the
second inequality is a fixed number below one, whereas
\(e^{-\widetilde m a_{j+1}}\to1\) as
\(a_{j+1}\to0\).
\end{proof}

\section{Why the Gaussian recursion still has no mass}
\label{sec:s7v29-massless}

\begin{theorem}[The retained Gaussian sector remains gapless]
\label{thm:s7v29-gapless}
At every blocking scale, the lowest retained nonzero frequency is
\begin{equation}
 \omega_{\min}(L)
 =
 {2\over a}
 \sin{\pi a\over L}
 \longrightarrow0
 \qquad(L\to\infty),                                    \label{eq:s7v29-gapless}
\end{equation}
where \(L\) is the physical spatial period.  Exact Gaussian blocking
therefore preserves a massless coarse transfer:
\begin{equation}
 \|T_{\rm coarse}^{(0)}|_{\Omega^\perp}\|
 \longrightarrow1
 \qquad(L\to\infty).                                    \label{eq:s7v29-coarse-norm}
\end{equation}
\end{theorem}

\begin{proof}
The least nonzero momentum has magnitude \(2\pi/L\) in one coordinate.
Insert it in \eqref{eq:s7v29-omega}.  Exact blocking retains this mode
at every scale until the box itself becomes one coarse cell, and
\Cref{thm:s7v29-matching} preserves its physical frequency.  Its
one-particle transfer eigenvalue tends to one, proving
\eqref{eq:s7v29-coarse-norm}.
\end{proof}

\begin{corollary}[Ultraviolet complement control is not mass
generation]
\label{cor:s7v29-no-mass}
Even perfect matching, zero RG error, and uniform fine-complement
contraction at every perturbative Gaussian step do not supply the
terminal hypothesis \(q_J<1\) uniformly in volume in
\Cref{thm:s7v28-recursion}.  A positive terminal mass must arise from
the interacting coarse dynamics.
\end{corollary}

\begin{proof}
Use \eqref{eq:s7v29-coarse-norm} for the terminal retained theory.
\end{proof}

\section{Nonlinear error ledger}
\label{sec:s7v29-errors}

\begin{proposition}[First interacting deviations from the benchmark]
\label{prop:s7v29-errors}
For a local gauge-covariant replacement of the sharp Gaussian split,
the quantities which must replace \eqref{eq:s7v29-rows} are:
\begin{enumerate}[label=\textup{(\roman*)},leftmargin=3.5em]
\item localization error between the local wavelet covariance and the
      perfect low covariance;
\item nonlinear curvature and Haar-Jacobian terms in the small-field
      chart;
\item low--high interaction blocks contributing to
      \(\epsilon_j\);
\item large-field and chart-exit contributions;
\item generated irrelevant and marginal spin-network interactions in
      the coarse action; and
\item toron, gauge, reflection-interface, and boundary terms.
\end{enumerate}
Each term must be bounded in the transfer-operator norms of
\eqref{eq:s7v28-coarse-error}--\eqref{eq:s7v28-offdiag}, not merely in
a formal action expansion.
\end{proposition}

\begin{proof}
Items (i)--(vi) are exactly the ways in which the interacting local
block differs from the orthogonal tensor factorization used in
\Cref{thm:s7v29-factor,thm:s7v29-matching}.  The V28 recursion accepts
only operator errors and a complement norm, so action coefficients
must be converted to those quantities before insertion.
\end{proof}

\begin{assessment}[Progress at seventh-series V29]
\label{ass:s7v29-status}
The weak-coupling benchmark is now exact.  A sharp low/high curvature
split has zero matching and off-diagonal error, and the discarded
sector contracts by
\(\exp\{-2b\sin(\pi/(2b))\}\) per blocking step, uniformly in cutoff
and volume.  The retained sector nevertheless has the original
massless lowest mode.  Hence the central continuum problem is neither
kinematic blocking nor ultraviolet complement damping: it is
nonperturbative mass generation in the coarse interacting theory,
together with operator control of low--high mixing.
\end{assessment}

\begin{decision}[V29 bottleneck decision]
\label{dec:s7v29-bottleneck}
V30 will perform the mandatory SM/NS audit, then compare their newest
source--test and form methods with the nonlinear error ledger in
\Cref{prop:s7v29-errors}.  If no physical transfer estimate is
available, the next Yang--Mills step will derive a local wavelet
covariance split and quantify its error relative to the perfect
Gaussian benchmark before adding interactions.
\end{decision}

\begin{firewall}[Claim boundary after seventh-series V29]
\label{fire:s7v29-boundary}
V29 proves exact Gaussian coarse matching, a scale-independent
fine-complement transfer contraction, and persistence of the massless
retained mode.

V29 does not control a local nonlinear RG map, low--high interaction
errors, chart exits and torons, accumulation through the
nonperturbative scale, terminal mass generation, continuum
reconstruction, or the full Yang--Mills mass gap.
\end{firewall}

\paragraph{Parent-version record.}
V29 was formed from
\texttt{Renewal\_Yang\_Mills\_Mass\_Gap\_Research\_Notes\_V28.tex}.
The V28 parent had \(11\,538\) logical source lines, \(421\,616\) bytes,
and SHA--256
\[
 \texttt{819BA8653A966C26E1629D6DDE8F1DED1BB892A6FA27C0100BD5B82B283C67B2}.
\]

% =============================================================================
% END APPENDED SEVENTH-SERIES V29 MATERIAL
% =============================================================================

\newpage

% =============================================================================
% BEGIN APPENDED SEVENTH-SERIES V30 MATERIAL
% =============================================================================

\chapter{Mandatory V30 companion audit}
\label{ch:s7v30}

This audit follows the exact Gaussian RG benchmark of V29.  Its
question is narrow: do the newest SM--Einstein or Navier--Stokes notes
supply a local wavelet construction, a low--high transfer estimate, or
a nonperturbative terminal mass for NTRGC?  The answer is no.  The SM
notes do supply a compatible bookkeeping lesson---global weak
densities should be formed before artificial patch boundaries are
charged---but this does not alter the Yang--Mills transfer bottleneck.

\section{Files and integrity}
\label{sec:s7v30-files}

The current companion files are
\begin{align}
 &\texttt{Renewal\_SM\_Einstein\_Continuum\_Research\_Notes\_V75.tex},
 \notag\\
 &\hspace{2em}29\,665\ \text{logical lines},\quad
 1\,043\,350\ \text{bytes},\quad
 \operatorname{SHA256}
 =\texttt{51661B6F7DE0DD5E6FB04A20FB63AE443E71DE6395B9ED0B8B0DDB856B33C46D},
                                                               \label{eq:s7v30-SM-file}\\
 &\texttt{Renewal\_NS\_Stability\_Research\_Notes\_V26.tex},
 \notag\\
 &\hspace{2em}5\,319\ \text{logical lines},\quad
 278\,283\ \text{bytes},\quad
 \operatorname{SHA256}
 =\texttt{82C887F8C2C69E4F28FAD7C3DC8DE5B9099CB5A4BF431EBA1FFF3E91935978AD}.
                                                               \label{eq:s7v30-NS-file}
\end{align}
Each file has one live \(\backslash\texttt{end\{document\}}\) command.
The two additional text matches in SM V75 are literal
\(\backslash\texttt{verb}\) records, not terminators.  NS V26 is
byte-for-byte unchanged from V25.

\section{New SM material}
\label{sec:s7v30-SM}

\begin{audit}[SM V73: weak curvature as a first-derivative form]
\label{audit:s7v30-SM73}
SM V73 rewrites the scalar-curvature density as a quadratic connection
term plus a divergence.  On a fixed three-dimensional patch,
uniformly elliptic \(W^{1,3}\cap L^\infty\) metrics produce a bounded
\(H^1_0\to H^{-1}\) curvature form, and strong
\(W^{1,3}\cap L^\infty\) convergence gives form-norm convergence.
This is a deterministic fixed-patch statement, not a gauge-transfer
or spectral-gap estimate.
\end{audit}

\begin{audit}[SM V74: genuine boundary trace]
\label{audit:s7v30-SM74}
SM V74 proves by an explicit conformal example that the V73 topology
does not control curvature flux at a boundary.  The stronger
\(W^{2,3/2}\cap W^{1,3}\cap L^\infty\) topology supplies an
\(L^2\) normal trace and a bounded full \(H^1\) form.  The result
concerns actual geometric boundaries and does not contract the
discarded modes of an RG step.
\end{audit}

\begin{audit}[SM V75: global-before-local and connected normalization]
\label{audit:s7v30-SM75}
After auditing YM V23--V27, SM V75 decides to form one global
reference-connection weak curvature density before using a partition
of unity, so artificial chart faces cancel.  It also records the
connected-cluster principle that source derivatives of the logarithm
remove disconnected extensive vacuum pieces.  This agrees with the
V27 correlator expansion, but SM explicitly has no positive Einstein
polymer majorant and no transfer contraction to import.
\end{audit}

\section{Type-correct audit conclusion}
\label{sec:s7v30-conclusion}

\begin{theorem}[V30 companion non-transfer]
\label{thm:s7v30-no-transfer}
None of SM V73--V75 or NS V26 supplies any of the nonlinear
Yang--Mills quantities
\[
 \epsilon_j,\qquad r_j,\qquad q_J
\]
in the transfer recursion \eqref{eq:s7v28-q-recursion}.  In
particular:
\begin{enumerate}[label=\textup{(\roman*)},leftmargin=3.5em]
\item convergence of a fixed-patch elliptic form does not imply
      spectral stability of a global gauge transfer;
\item a boundary normal-trace estimate is not a low--high oscillator
      contraction;
\item cancellation of artificial coordinate interfaces is not mass
      generation in the retained infrared sector; and
\item the NS Oseen rank-one constants act on a different kernel,
      field space, and evolution.
\end{enumerate}
\end{theorem}

\begin{proof}
The audited SM inputs are forms on scalar \(H^1\) patch spaces and
their boundary traces.  The NTRGC inputs are operator norms between
gauge-invariant OS Hilbert spaces and a terminal transfer spectral
edge.  No intertwining identity identifies these objects.  The NS
input is a pointwise tensor norm for the Oseen kernel, so it has the
same type mismatch.  Therefore none can be inserted into
\eqref{eq:s7v28-coarse-error}--\eqref{eq:s7v28-residual}.
\end{proof}

\begin{proposition}[Reusable global-before-interface rule]
\label{prop:s7v30-global-first}
In a local Yang--Mills wavelet RG, exact gauge constraints and Bianchi
identities must be imposed on the global fine complex before wavelet
patch boundaries are assigned residual errors.  Otherwise a
coordinate or ownership interface can be mistaken for physical
low--high mixing.
\end{proposition}

\begin{proof}
The lattice identities \(d_1d_0=0\) and the gauge covariance of
\eqref{eq:s7v28-B} hold globally.  Restricting first creates boundary
commutators of the restriction with \(d_0,d_1\), exactly as the
cut-gauge curvature in V23--V24.  Those commutators are bookkeeping
interfaces, not bulk interacting mass terms.  Forming the global
complex first preserves the identities; only the subsequently proved
localization error belongs to \(\epsilon_j\).
\end{proof}

\begin{assessment}[Progress at seventh-series V30]
\label{ass:s7v30-status}
The required companion audit is complete.  It confirms the V29
direction and adds one discipline: construct the global covariant
low/high complex before charging localization boundaries.  No
companion result changes the main frontier.  The exact Gaussian
discarded sector contracts uniformly; the unproved rows remain local
nonlinear low--high operator control and terminal nonperturbative mass
generation.
\end{assessment}

\begin{decision}[V30 bottleneck decision]
\label{dec:s7v30-bottleneck}
V31 will build a global gauge-compatible smooth Fourier partition and
then localize it by finite-range wavelet filters.  It will quantify the
commutator with the lattice curvature complex and translate that
commutator into the coarse matching, off-diagonal, and residual norms
of V28.  Interaction and large-field terms will be added only after
the free localization debit is explicit.
\end{decision}

\begin{firewall}[Claim boundary after seventh-series V30]
\label{fire:s7v30-boundary}
V30 records the current SM V75/NS V26 audit, proves the type
non-transfer, and imports only the global-before-interface bookkeeping
rule.

V30 does not construct the local wavelet RG, prove nonlinear transfer
errors, generate a terminal mass, remove the cutoff, or prove the full
Yang--Mills mass gap.
\end{firewall}

\paragraph{Parent-version record.}
V30 was formed from
\texttt{Renewal\_Yang\_Mills\_Mass\_Gap\_Research\_Notes\_V29.tex}.
The V29 parent had \(11\,879\) logical source lines, \(433\,971\) bytes,
and SHA--256
\[
 \texttt{F7C7AC00F6FE5950433B0406ADF1AF962A41011B233B095551D557F4F04FD067}.
\]

% =============================================================================
% END APPENDED SEVENTH-SERIES V30 MATERIAL
% =============================================================================

\newpage

% =============================================================================
% BEGIN APPENDED SEVENTH-SERIES V31 MATERIAL
% =============================================================================

\chapter{Local commuting filters and the single-channel projection
obstruction}
\label{ch:s7v31}

V29 uses a sharp Fourier projection to split retained and discarded
Gaussian modes.  V30 proposes localizing that split without first
breaking the global gauge complex.  The first calculation gives a
useful correction: a scalar convolution applied componentwise to
lattice cochains commutes \emph{exactly} with the exterior derivative,
regardless of its range.  The obstruction is elsewhere.  No
nonconstant finite-range scalar convolution can be an orthogonal
projection, and a positive contraction which is uniformly almost
idempotent is necessarily close to either zero or the identity.

Thus a nontrivial local low/high decomposition requires downsampling,
multiple filter channels, or loss of exact projection.  This note
proves that no-go and quantifies the last option.

\section{Scalar cochain convolutions are exact chain maps}
\label{sec:s7v31-chain}

On the periodic lattice, let \(C^q\) be the Hilbert space of
real or Lie-algebra-valued oriented \(q\)-cochains, and let
\[
 d_q:C^q\longrightarrow C^{q+1}
\]
be the translation-invariant lattice exterior derivative.  For a
scalar kernel \(h\), define componentwise convolution
\begin{equation}
 (S_h^{(q)}f)_I(x)
 :=
 \sum_yh(y)f_I(x-y),                                    \label{eq:s7v31-convolution}
\end{equation}
using the same kernel for every oriented \(q\)-cell component \(I\).

\begin{theorem}[Exact convolution--complex commutation]
\label{thm:s7v31-commutation}
For every kernel \(h\), finite range or not,
\begin{equation}
 d_qS_h^{(q)}
 =
 S_h^{(q+1)}d_q.                                        \label{eq:s7v31-commutation}
\end{equation}
In particular,
\[
 S_h^{(2)}d_1d_0=d_1d_0S_h^{(0)}=0.
\]
\end{theorem}

\begin{proof}
In Fourier variables, \(S_h^{(q)}\) has scalar symbol
\(\widehat h(k)I\), while \(d_q\) has its incidence symbol \(D_q(k)\).
Scalar multiplication commutes with every matrix:
\[
 D_q(k)\widehat h(k)I
 =
 \widehat h(k)I D_q(k).
\]
Inverse Fourier transformation gives
\eqref{eq:s7v31-commutation}.
\end{proof}

\begin{corollary}[No localization curvature commutator]
\label{cor:s7v31-no-commutator}
Finite-range truncation of a scalar translation-invariant filter
creates no commutator debit with \(d_0\) or \(d_1\).  Any free
low--high matching error comes from non-idempotence, downsampling, or
the choice of retained band, not from the lattice Bianchi identity.
\end{corollary}

\begin{proof}
Finite-range truncation changes \(h\) but not the algebra in
\Cref{thm:s7v31-commutation}.
\end{proof}

\section{No nontrivial finite-range scalar projection}
\label{sec:s7v31-no-go}

\begin{theorem}[Finite-range scalar projection no-go]
\label{thm:s7v31-projection-no-go}
Let \(S_h\) be a translation-invariant finite-range scalar convolution
on a periodic lattice family of arbitrarily large periods.  If
\[
 S_h=S_h^*=S_h^2,
\]
then \(S_h=0\) or \(S_h=I\).
\end{theorem}

\begin{proof}
Finite range makes the multiplier
\[
 p(k)=\sum_{|x|\le R}h(x)e^{-ik\cdot x}
\]
a continuous trigonometric polynomial on the connected Brillouin
torus.  Self-adjointness makes \(p\) real, and idempotence gives
\[
 p(k)^2=p(k).
\]
Thus \(p(k)\in\{0,1\}\) for every \(k\).  A continuous map from a
connected space to the discrete set \(\{0,1\}\) is constant.  Hence
the convolution is zero or the identity.
\end{proof}

\begin{theorem}[Quantitative near-projection no-go]
\label{thm:s7v31-near-projection}
Let \(S_h\) be a self-adjoint positive-contraction scalar convolution
on \(\ell^2(\mathbb Z^d)\), or a fixed convolution kernel with the
stated bounds uniformly on periodic lattices of arbitrarily large
period.  Assume
\begin{equation}
 \|S_h^2-S_h\|\le\epsilon<\frac14.                       \label{eq:s7v31-near}
\end{equation}
Set
\begin{equation}
 \delta(\epsilon)
 :=
 {1-\sqrt{1-4\epsilon}\over2}.                          \label{eq:s7v31-delta}
\end{equation}
Then either
\begin{equation}
 \|S_h\|\le\delta(\epsilon)
 \quad\hbox{or}\quad
 \|I-S_h\|\le\delta(\epsilon).                          \label{eq:s7v31-trivial}
\end{equation}
Therefore a scalar filter cannot be uniformly close to a nontrivial
orthogonal band projection while remaining continuous and
translation invariant.
\end{theorem}

\begin{proof}
The scalar multiplier satisfies \(0\le p(k)\le1\) and
\[
 p(k)(1-p(k))\le\epsilon.
\]
Solving the quadratic inequality gives
\[
 p(k)\in[0,\delta(\epsilon)]
 \cup[1-\delta(\epsilon),1].
\]
Because \(\epsilon<1/4\), these intervals are disjoint.  Continuity on
the connected torus forces the image of \(p\) into one interval.
Taking the essential supremum gives
\eqref{eq:s7v31-trivial}.
\end{proof}

\begin{firewall}[Scope of the obstruction]
\label{fire:s7v31-scope}
Theorems \ref{thm:s7v31-projection-no-go} and
\ref{thm:s7v31-near-projection} concern one scalar,
translation-invariant channel.  Matrix-valued multichannel filter
banks, critically sampled wavelets, and block-periodic projections are
not excluded.  They are exactly the remaining local possibilities.
\end{firewall}

\section{Controlled finite-range approximation of a smooth split}
\label{sec:s7v31-smooth}

Let \(K_{\rm lo}\) and \(K_{\rm hi}\) be disjoint closed momentum sets
with positive separation.  Choose a real even
\(\ell\in C^m(\mathbb T^3)\), \(m>3\), such that
\begin{equation}
 0\le\ell\le1,\qquad
 \ell=1\ \hbox{on }K_{\rm lo},\qquad
 \ell=0\ \hbox{on }K_{\rm hi}.                          \label{eq:s7v31-cutoff}
\end{equation}
Let
\[
 \ell(k)=\sum_{x\in\mathbb Z^3}\widehat\ell(x)e^{-ik\cdot x}
\]
and truncate at range \(R\):
\begin{equation}
 \ell_R(k)
 :=
 \sum_{\|x\|_\infty\le R}
 \widehat\ell(x)e^{-ik\cdot x}.                         \label{eq:s7v31-truncation}
\end{equation}

\begin{theorem}[Finite-range smooth-filter error]
\label{thm:s7v31-filter-error}
There is \(C_{\ell,m}<\infty\), independent of lattice period, such
that
\begin{equation}
 \|S_{\ell_R}-S_\ell\|_{\ell^2\to\ell^2}
 \le
 \varepsilon_R
 :=
 C_{\ell,m}R^{3-m}.                                     \label{eq:s7v31-filter-error}
\end{equation}
Both filters commute exactly with the complete lattice de Rham
complex.  Moreover,
\begin{align}
 \|(S_{\ell_R}-I)P_{\rm lo}\|&\le\varepsilon_R,
                                                               \label{eq:s7v31-low-error}\\
 \|S_{\ell_R}P_{\rm hi}\|&\le\varepsilon_R,              \label{eq:s7v31-high-error}
\end{align}
where \(P_{\rm lo}\) and \(P_{\rm hi}\) are the sharp projections onto
the two declared bands.
\end{theorem}

\begin{proof}
Repeated integration by parts in the Fourier coefficient gives
\[
 |\widehat\ell(x)|
 \le C_{\ell,m}(1+\|x\|_\infty)^{-m}.
\]
The number of lattice points in a radius-\(r\) shell is \(O(r^2)\);
hence
\[
 \sum_{\|x\|_\infty>R}|\widehat\ell(x)|
 \le C_{\ell,m}\sum_{r>R}r^{2-m}
 \le C_{\ell,m}R^{3-m}.
\]
The convolution operator norm is the supremum norm of its multiplier,
which is bounded by this Fourier tail.  Exact commutation is
\Cref{thm:s7v31-commutation}.  On \(K_{\rm lo}\), \(S_\ell=I\), and on
\(K_{\rm hi}\), \(S_\ell=0\); restricting the operator estimate proves
\eqref{eq:s7v31-low-error}--\eqref{eq:s7v31-high-error}.
\end{proof}

\begin{corollary}[Free transfer leakage away from the transition band]
\label{cor:s7v31-transfer}
For the Gaussian transfer of V29, \(S_{\ell_R}\) commutes exactly with
\(e^{-tH_a^{(0)}}\).  If \(K_{\rm hi}\) begins at
\(\|ak\|_\infty\ge\kappa>0\), then every vector spectrally supported
there contracts for time \(ba\) by at most
\begin{equation}
 \exp\{-2b\sin(\kappa/2)\}.                              \label{eq:s7v31-high-contraction}
\end{equation}
The filter leakage into or out of the separated low/high bands is at
most \(\varepsilon_R\).
\end{corollary}

\begin{proof}
Both the free Hamiltonian and the filter are scalar Fourier
multipliers on each transverse polarization, so they commute.  The
frequency bound follows from \eqref{eq:s7v29-omega} exactly as in
\Cref{lem:s7v29-floor,thm:s7v29-contraction}.  The leakage bounds are
\eqref{eq:s7v31-low-error}--\eqref{eq:s7v31-high-error}.
\end{proof}

\begin{warning}[The transition band is mathematically substantive]
\label{warn:s7v31-transition}
On the region where \(0<\ell<1\), the exact multiplier has
\(\ell(1-\ell)>0\); if it crosses \(1/2\), its idempotence defect is
\(1/4\).  Increasing \(R\) cannot remove this defect.  The smooth
single-channel filter supplies localized analysis and separated-band
leakage bounds, but not the isometry \(Q_j\) and orthogonal
decomposition assumed by \Cref{thm:s7v28-recursion}.
\end{warning}

\section{Revised free localization card}
\label{sec:s7v31-card}

\begin{definition}[Cochain wavelet filter-bank card, CWFBC]
\label{def:s7v31-CWFBC}
CWFBC consists of:
\begin{enumerate}[label=\textup{(W\arabic*)},leftmargin=4em]
\item a finite-range multichannel analysis/synthesis bank on every
      cochain degree;
\item exact or quantitatively approximate intertwining with
      \(d_0,d_1\);
\item an isometric perfect-reconstruction identity after downsampling;
\item separation of the retained low channel from all high channels,
      including the transition frequencies;
\item a scale-independent free transfer contraction on the high
      synthesis range;
\item gauge, toron, reflection, and boundary compatibility;
\item conversion of filter errors to
      \(\epsilon_j,r_j\) in the V28 operator norms; and
\item nonlinear interaction and chart-exit perturbations added only
      after rows (W1)--(W7).
\end{enumerate}
\end{definition}

\begin{theorem}[What V31 decides]
\label{thm:s7v31-result}
The free localization problem cannot be solved by a nontrivial
single-channel finite-range scalar projection.  Smooth finite-range
scalar filters preserve the gauge complex exactly and have
superalgebraically small separated-band leakage, but leave a
non-idempotent transition sector.  A multichannel downsampled cochain
filter bank is therefore necessary for the local version of the V29
split.
\end{theorem}

\begin{proof}
Exact complex preservation is
\Cref{thm:s7v31-commutation}; the single-channel obstruction is
\Cref{thm:s7v31-projection-no-go,thm:s7v31-near-projection}; and the
available approximation is \Cref{thm:s7v31-filter-error}.
\end{proof}

\begin{assessment}[Progress at seventh-series V31]
\label{ass:s7v31-status}
The anticipated curvature-commutator debit vanishes for scalar
convolution filters.  The true free obstruction is topological:
continuity of a scalar Fourier multiplier forbids a nontrivial local
projection.  A controlled smooth finite-range filter exists on
separated bands, but its transition band prevents direct insertion as
the V28 coarse isometry.  This narrows the next construction to
multichannel cochain wavelets.
\end{assessment}

\begin{decision}[V31 bottleneck decision]
\label{dec:s7v31-bottleneck}
V32 will construct a critically sampled two-scale filter bank on the
one-dimensional lattice, derive the exact polyphase perfect-
reconstruction identities, and then tensorize it across spatial
coordinates.  Cochain-degree filters will be chosen to satisfy the
discrete derivative intertwining identity.  The first test is whether
the resulting high-channel space retains the V29 frequency floor
without a volume-dependent constant.
\end{decision}

\begin{firewall}[Claim boundary after seventh-series V31]
\label{fire:s7v31-boundary}
V31 proves exact convolution--complex commutation, exact and
quantitative single-channel projection no-go theorems, and a
finite-range smooth-filter leakage bound.

V31 does not construct the required multichannel cochain wavelet bank,
prove its transfer complement contraction, control nonlinear RG
errors, generate the terminal mass, remove the cutoff, or prove the
full Yang--Mills mass gap.
\end{firewall}

\paragraph{Parent-version record.}
V31 was formed from
\texttt{Renewal\_Yang\_Mills\_Mass\_Gap\_Research\_Notes\_V30.tex}.
The V30 parent had \(12\,054\) logical source lines, \(441\,384\) bytes,
and SHA--256
\[
 \texttt{A4BA314FFE5E3FEC58336F7475679B07DC10E427A7797C5A0D6B95899C6A8D85}.
\]

% =============================================================================
% END APPENDED SEVENTH-SERIES V31 MATERIAL
% =============================================================================

\newpage

% =============================================================================
% BEGIN APPENDED SEVENTH-SERIES V32 MATERIAL
% =============================================================================

\chapter{A finite-range cochain lifting split and its orthogonalization
debit}
\label{ch:s7v32}

V31 proves that a single scalar local convolution cannot be a
nontrivial projection.  The simplest multichannel replacement is a
lifting split.  In one dimension it can be made an exact cochain map:
the fine derivative becomes one coarse difference channel plus one
detail channel, with local perfect reconstruction and volume-uniform
condition numbers.  The detail channel also has exactly the ultraviolet
energy floor needed in V29.

There is, however, a new distinction.  The lifting coordinates form a
Riesz decomposition, not an orthogonal one.  Orthogonalizing the coarse
embedding produces a quasi-local perfect coarse derivative and an
explicit coarse--detail derivative block of norm \(2-\sqrt2\).  Hence
the exact local chain split cannot be inserted into the orthogonal V28
transfer compiler with zero error.

\section{One-dimensional lifting analysis}
\label{sec:s7v32-lifting}

Let the fine periodic vertex lattice have even length \(2M\).  For a
vertex cochain \(f\), define
\begin{align}
 c_n&:=f_{2n},                                           \label{eq:s7v32-c}\\
 d_n&:=f_{2n+1}-{f_{2n}+f_{2n+2}\over2}.                \label{eq:s7v32-d}
\end{align}
For an edge cochain \(a\), with \(a_j\) oriented from \(j\) to \(j+1\),
define
\begin{align}
 b_n&:=a_{2n}+a_{2n+1},                                 \label{eq:s7v32-b}\\
 e_n&:={a_{2n}-a_{2n+1}\over2}.                         \label{eq:s7v32-e}
\end{align}
All indices are periodic.

\begin{theorem}[Local perfect reconstruction]
\label{thm:s7v32-reconstruction}
The vertex and edge analysis maps are bijective, with local synthesis
\begin{align}
 f_{2n}&=c_n,&
 f_{2n+1}&=d_n+{c_n+c_{n+1}\over2},                    \label{eq:s7v32-f-synth}\\
 a_{2n}&={b_n\over2}+e_n,&
 a_{2n+1}&={b_n\over2}-e_n.                            \label{eq:s7v32-a-synth}
\end{align}
Their norms and inverse norms are bounded independently of \(M\):
\begin{align}
 {1\over2}\|f\|_2
 &\le\|(c,d)\|_2\le2\|f\|_2,                            \label{eq:s7v32-vertex-Riesz}\\
 {1\over\sqrt2}\|a\|_2
 &\le\|(b,e)\|_2\le\sqrt2\|a\|_2.                       \label{eq:s7v32-edge-Riesz}
\end{align}
\end{theorem}

\begin{proof}
Equations \eqref{eq:s7v32-f-synth}--\eqref{eq:s7v32-a-synth} follow
by solving the definitions.  In coarse Fourier variable \(z=e^{i\theta}\),
the vertex analysis matrix on the even/odd polyphase vector is
\[
 A_0(\theta)
 =
 \begin{pmatrix}
 1&0\\
 -q(\theta)&1
 \end{pmatrix},
 \qquad q(\theta)={1+z\over2},\qquad |q|\le1.
\]
Both \(A_0\) and \(A_0^{-1}\) have norm at most \(2\), proving
\eqref{eq:s7v32-vertex-Riesz}.  The edge matrix is
\[
 A_1=
 \begin{pmatrix}
 1&1\\
 1/2&-1/2
 \end{pmatrix}.
\]
Its rows are orthogonal with squared norms \(2\) and \(1/2\), so its
singular values are \(\sqrt2\) and \(1/\sqrt2\).  This proves
\eqref{eq:s7v32-edge-Riesz}.
\end{proof}

\section{Exact derivative intertwining}
\label{sec:s7v32-chain}

Let
\[
 (D f)_j:=f_{j+1}-f_j,\qquad
 (D_cf)_n:=f_{n+1}-f_n.
\]

\begin{theorem}[Exact lifting cochain identity]
\label{thm:s7v32-chain}
If \(a=Df\), then its analyzed channels satisfy
\begin{equation}
 b=D_cc,\qquad e=d.                                     \label{eq:s7v32-chain}
\end{equation}
Equivalently, in lifting coordinates the fine differential is
\begin{equation}
 \widetilde D(c,d)=(D_cc,d).                            \label{eq:s7v32-Dtilde}
\end{equation}
\end{theorem}

\begin{proof}
Telescoping gives
\[
 b_n=(f_{2n+1}-f_{2n})+(f_{2n+2}-f_{2n+1})
 =c_{n+1}-c_n.
\]
Also
\begin{align*}
 e_n
 &={1\over2}
 \bigl[
 (f_{2n+1}-f_{2n})-(f_{2n+2}-f_{2n+1})
 \bigr]\\
 &=f_{2n+1}-{f_{2n}+f_{2n+2}\over2}
 =d_n.
\end{align*}
This is \eqref{eq:s7v32-chain}.
\end{proof}

\begin{corollary}[Exact high-channel energy]
\label{cor:s7v32-high-energy}
On the pure lifting-detail subspace \(c=0\),
\begin{equation}
 \|Df\|_2^2=2\|f\|_2^2.                                \label{eq:s7v32-high-energy}
\end{equation}
Thus a spatial lattice derivative \(a^{-1}D\) has detail frequency
\(\sqrt2/a\), and evolution for a doubled time step \(2a\) contracts
that pure detail channel by \(e^{-2\sqrt2}\).
\end{corollary}

\begin{proof}
If \(c=0\), synthesis gives \(f_{2n}=0\) and \(f_{2n+1}=d_n\).
The two adjacent edge differences are \(d_n\) and \(-d_n\), so
\(\|Df\|_2^2=2\sum_n|d_n|^2=2\|f\|_2^2\).
\end{proof}

\section{Coarse embeddings and the perfect derivative}
\label{sec:s7v32-orthogonal}

Let \(J_0\) and \(J_1\) be the lifting synthesis maps with the detail
channel set to zero:
\begin{align}
 (J_0c)_{2n}&=c_n,&
 (J_0c)_{2n+1}&={c_n+c_{n+1}\over2},                   \label{eq:s7v32-J0}\\
 (J_1b)_{2n}&={b_n\over2},&
 (J_1b)_{2n+1}&={b_n\over2}.                           \label{eq:s7v32-J1}
\end{align}

\begin{lemma}[Exact unnormalized coarse chain map]
\label{lem:s7v32-J-chain}
One has
\begin{equation}
 DJ_0=J_1D_c.                                           \label{eq:s7v32-J-chain}
\end{equation}
The Gram multipliers are
\begin{equation}
 G_0(\theta):=J_0^*J_0
 =1+\cos^2{\theta\over2},\qquad
 G_1:=J_1^*J_1={1\over2}.                              \label{eq:s7v32-Gram}
\end{equation}
\end{lemma}

\begin{proof}
Both fine differences of the linear interpolant
\eqref{eq:s7v32-J0} equal
\((c_{n+1}-c_n)/2\), which is
\eqref{eq:s7v32-J1} applied to \(D_cc\).  In Fourier variables the
vertex synthesis column is
\[
 j_0(\theta)=
 \begin{pmatrix}1\\q(\theta)\end{pmatrix},
 \qquad |q(\theta)|^2=\cos^2(\theta/2),
\]
while the edge synthesis column is
\((1/2,1/2)^{\mathsf T}\).  Their squared norms give
\eqref{eq:s7v32-Gram}.
\end{proof}

Define the isometric coarse embeddings
\begin{equation}
 Q_0:=J_0G_0^{-1/2},\qquad
 Q_1:=J_1G_1^{-1/2},                                   \label{eq:s7v32-Q}
\end{equation}
and the perfect coarse derivative
\begin{equation}
 D_c^{\rm perf}
 :=
 G_1^{1/2}D_cG_0^{-1/2}.                               \label{eq:s7v32-Dperf}
\end{equation}

\begin{theorem}[Orthogonal coarse chain map]
\label{thm:s7v32-perfect-chain}
The isometries \eqref{eq:s7v32-Q} obey
\begin{equation}
 DQ_0=Q_1D_c^{\rm perf}.                               \label{eq:s7v32-perfect-chain}
\end{equation}
The multiplier \(G_0^{-1/2}\) is analytic in a complex strip around
the real torus, so \(Q_0\) and \(D_c^{\rm perf}\) are exponentially
quasi-local, but not finite range.
\end{theorem}

\begin{proof}
Use \eqref{eq:s7v32-J-chain}:
\[
 DQ_0
 =J_1D_cG_0^{-1/2}
 =Q_1G_1^{1/2}D_cG_0^{-1/2}.
\]
This is \eqref{eq:s7v32-perfect-chain}.  The symbol
\(G_0=1+\cos^2(\theta/2)\) lies in \([1,2]\); its analytic continuation
has no zero in a sufficiently narrow strip.  Holomorphic functional
calculus makes \(G_0^{-1/2}\) analytic there, and its Fourier
coefficients therefore decay exponentially.
\end{proof}

\section{The orthogonal coarse--detail coupling}
\label{sec:s7v32-debit}

Let \(H_0(\theta)\) be the normalized column spanning the orthogonal
complement of \(Q_0\):
\begin{equation}
 H_0(\theta)
 =
 {1\over\sqrt{1+|q|^2}}
 \begin{pmatrix}-\overline q\\1\end{pmatrix}.           \label{eq:s7v32-H0}
\end{equation}

\begin{theorem}[Exact orthogonal high energy and mixing]
\label{thm:s7v32-mixing}
For every coarse momentum,
\begin{align}
 \|DH_0(\theta)\|^2
 &=
 {2(1+3\cos^2(\theta/2))
  \over1+\cos^2(\theta/2)}
 \ge2,                                                  \label{eq:s7v32-orth-energy}\\
 \|Q_1^*DH_0\|_{\rm op}
 &=
 2-\sqrt2.                                              \label{eq:s7v32-mixing}
\end{align}
Thus the orthogonal detail space retains the frequency floor
\(\sqrt2/a\), but its derivative has a nonzero coarse-edge component
of order one.
\end{theorem}

\begin{proof}
Ignoring the common normalization in \eqref{eq:s7v32-H0}, the
even/odd vertex components are
\((-\overline q,1)\).  Their fine edge differences are
\[
 1+\overline q,\qquad -(1+q).
\]
Therefore
\[
 \|DH_0\|^2
 ={2|1+q|^2\over1+|q|^2}.
\]
With \(q=(1+e^{i\theta})/2\),
\[
 |q|^2=\cos^2(\theta/2),\qquad
 |1+q|^2=1+3\cos^2(\theta/2),
\]
which proves \eqref{eq:s7v32-orth-energy}.

The normalized coarse-edge column is
\((1,1)^{\mathsf T}/\sqrt2\).  Its inner product with the displayed
two edge differences has magnitude
\[
 {|\overline q-q|
  \over\sqrt{2(1+|q|^2)}}
 =
 {|\sin\theta|
  \over\sqrt{2(1+\cos^2(\theta/2))}}.
\]
Writing \(x=\cos^2(\theta/2)\), its square is
\[
 {2x(1-x)\over1+x}.
\]
The maximum occurs at \(x=\sqrt2-1\) and equals
\(6-4\sqrt2=(2-\sqrt2)^2\), proving
\eqref{eq:s7v32-mixing}.
\end{proof}

\begin{corollary}[Why orthogonalization is not a zero-error insertion]
\label{cor:s7v32-not-zero}
Although the lifting differential \eqref{eq:s7v32-Dtilde} is exactly
channel diagonal, the orthogonal coarse/detail decomposition has a
coarse--detail derivative block of norm \(2-\sqrt2\).  Hence an
orthogonal transfer Hamiltonian built from \(D^*D\) requires an
explicit off-diagonal analysis; the local lifting identity alone does
not set \(\epsilon_j=0\) in \eqref{eq:s7v28-offdiag}.
\end{corollary}

\begin{proof}
Use \eqref{eq:s7v32-mixing}.  A nonzero off-diagonal block of \(D\)
generically contributes cross blocks to \(D^*D\) and its transfer
semigroup.
\end{proof}

\section{Tensor-product extension}
\label{sec:s7v32-tensor}

\begin{theorem}[Multidimensional Riesz filter bank]
\label{thm:s7v32-tensor}
Tensoring the one-dimensional lifting analysis in \(d_s\) spatial
coordinates gives \(2^{d_s}\) vertex channels and the corresponding
tensor-product cochain channels.  It has:
\begin{enumerate}[label=\textup{(\roman*)},leftmargin=3.5em]
\item finite-range perfect reconstruction;
\item an exact tensor-product de Rham chain identity before
      orthogonalization;
\item Riesz bounds depending only on \(d_s\), not volume or scale; and
\item a gradient-energy floor at least \(2/a^2\) on every pure channel
      with at least one detail factor.
\end{enumerate}
\end{theorem}

\begin{proof}
Tensor products preserve invertibility and perfect reconstruction.
The de Rham complex on a rectangular lattice is the graded tensor
product of the one-dimensional two-term complexes; apply
\eqref{eq:s7v32-chain} in each factor with the standard Koszul signs.
The Riesz constants multiply a fixed number \(d_s\) of times.  In a
channel containing a detail factor, the derivative in that coordinate
alone contributes the lower bound \eqref{eq:s7v32-high-energy}; all
other coordinate energies are nonnegative.
\end{proof}

\begin{warning}[Tensorization does not remove the orthogonal debit]
\label{warn:s7v32-tensor}
At least one coordinate factor in a multidimensional orthogonal
high channel carries the one-dimensional coupling
\eqref{eq:s7v32-mixing}.  Tensorization keeps it volume independent
but does not make it perturbatively small.
\end{warning}

\section{Filter-bank status}
\label{sec:s7v32-status}

\begin{theorem}[CWFBC rows closed at V32]
\label{thm:s7v32-card}
The Riesz analogues of rows (W1)--(W5) of CWFBC are closed:
finite-range multichannel analysis/synthesis, exact chain
intertwining, stable perfect reconstruction, uniform channel
separation, and a scale-independent high-channel energy floor.  The
isometric requirement in row (W3), as needed literally by the V28
orthogonal compiler, is not closed by the finite-range lifting map.
After orthogonalization the perfect coarse derivative is quasi-local
and the coarse--detail debit is the explicit constant \(2-\sqrt2\).
\end{theorem}

\begin{proof}
Use
\Cref{thm:s7v32-reconstruction,thm:s7v32-chain,
cor:s7v32-high-energy,thm:s7v32-perfect-chain,
thm:s7v32-mixing,thm:s7v32-tensor}.
\end{proof}

\begin{assessment}[Progress at seventh-series V32]
\label{ass:s7v32-status}
A concrete local cochain wavelet now exists.  In lifting coordinates it
is an exact finite-range chain map with uniformly stable reconstruction
and the same \(e^{-2\sqrt2}\) ultraviolet scale seen in V29.  Passing
to an orthogonal coarse Hilbert subspace preserves the high energy
floor but creates an order-one derivative cross block.  The next issue
is therefore not existence of local wavelets; it is whether the free
semigroup suppresses that cross block enough over one block time, or
whether the V28 compiler should be generalized from orthogonal to
uniform Riesz decompositions.
\end{assessment}

\begin{decision}[V32 bottleneck decision]
\label{dec:s7v32-bottleneck}
V33 will compute the exact \(2\times2\) polyphase matrix of the free
Laplacian and its blocked semigroup in the orthogonal
\((Q_0,H_0)\) basis.  It will obtain the true off-diagonal transfer
norm after time \(2a\), rather than estimate it from
\(\|Q_1^*DH_0\|\), and test the V28 step budget quantitatively.
\end{decision}

\begin{firewall}[Claim boundary after seventh-series V32]
\label{fire:s7v32-boundary}
V32 constructs an exact finite-range multichannel cochain lifting
split, proves uniform Riesz bounds and high-channel energy, constructs
the orthogonal quasi-local coarse derivative, and computes the exact
coarse--detail derivative coupling.

V32 does not yet bound the off-diagonal free transfer semigroup in the
orthogonal split, incorporate gauge projection and nonlinear
interactions, generate a terminal mass, remove the cutoff, or prove the
full Yang--Mills mass gap.
\end{firewall}

\paragraph{Parent-version record.}
V32 was formed from
\texttt{Renewal\_Yang\_Mills\_Mass\_Gap\_Research\_Notes\_V31.tex}.
The V31 parent had \(12\,399\) logical source lines, \(453\,919\) bytes,
and SHA--256
\[
 \texttt{76B105570DE440B122FA130B6C47D67B9F0F9C95665D685CA2BBDA1B80A5FAEE}.
\]

% =============================================================================
% END APPENDED SEVENTH-SERIES V32 MATERIAL
% =============================================================================

\newpage

% =============================================================================
% BEGIN APPENDED SEVENTH-SERIES V33 MATERIAL
% =============================================================================

\chapter{The exact Haar-block free semigroup}
\label{ch:s7v33}

V32 computes an order-one derivative cross block after orthogonalizing
the local lifting split.  The V28 compiler, however, needs the
off-diagonal block of the transfer over a complete block time, not the
derivative itself.  This note diagonalizes the \(2\times2\) polyphase
Laplacian and evaluates that block exactly.  The semigroup strongly
suppresses it, but only to a fixed constant.  Since the admissible
physical-margin error at an ultraviolet RG step is \(O(a)\), a fixed
Haar filter still fails the continuum step budget.

\section{Polyphase Laplacian}
\label{sec:s7v33-L}

Use the V32 even/odd vertex vector
\[
 F(\theta)=
 \begin{pmatrix}f_{\rm ev}(\theta)\\f_{\rm odd}(\theta)\end{pmatrix},
 \qquad z=e^{i\theta}.
\]
The fine forward derivative and Laplacian have symbols
\begin{align}
 D(\theta)
 &=
 \begin{pmatrix}
 -1&1\\ z&-1
 \end{pmatrix},                                        \label{eq:s7v33-D}\\
 L(\theta):=D(\theta)^*D(\theta)
 &=
 \begin{pmatrix}
 2&-1-\overline z\\
 -1-z&2
 \end{pmatrix}.                                        \label{eq:s7v33-L}
\end{align}

\begin{lemma}[Aliased fine eigenvalues]
\label{lem:s7v33-eigenvalues}
For \(\theta\in[-\pi,\pi]\), put
\[
 c:=\cos(\theta/2)\in[0,1].
\]
The eigenvalues of \(L(\theta)\) are
\begin{equation}
 \lambda_-(c)=2-2c,\qquad
 \lambda_+(c)=2+2c.                                    \label{eq:s7v33-eigenvalues}
\end{equation}
\end{lemma}

\begin{proof}
The trace is \(4\), and the off-diagonal magnitude is
\(|1+z|=2c\).  The eigenvalues are therefore \(2\pm2c\).
\end{proof}

Let
\[
 q={1+z\over2},\qquad G=1+|q|^2=1+c^2.
\]
The orthonormal coarse/detail basis of V32 is
\begin{equation}
 U(\theta)
 :=
 \begin{pmatrix}Q_0&H_0\end{pmatrix}
 =
 {1\over\sqrt G}
 \begin{pmatrix}
 1&-\overline q\\
 q&1
 \end{pmatrix}.                                        \label{eq:s7v33-U}
\end{equation}

\begin{lemma}[Laplacian in the orthogonal lifting basis]
\label{lem:s7v33-block-L}
Writing
\[
 U^*LU=
 \begin{pmatrix}A&B\\\overline B&C\end{pmatrix},
\]
one has
\begin{align}
 A(c)&={2(1-c^2)\over1+c^2},                            \label{eq:s7v33-A}\\
 C(c)&={2(1+3c^2)\over1+c^2},                           \label{eq:s7v33-C}\\
 |B(c)|&={2c(1-c^2)\over1+c^2}.                        \label{eq:s7v33-B}
\end{align}
\end{lemma}

\begin{proof}
The diagonal entries are the squared derivative norms of \(Q_0\) and
\(H_0\).  Equation \eqref{eq:s7v32-perfect-chain} gives
\[
 A={|z-1|^2\over2G}
 ={2(1-c^2)\over1+c^2},
\]
while \eqref{eq:s7v32-orth-energy} gives \(C\).  Since
\(\det(U^*LU)=\lambda_-\lambda_+=4(1-c^2)\), the identity
\(|B|^2=AC-4(1-c^2)\) yields
\eqref{eq:s7v33-B}.
\end{proof}

\section{Exact doubled-step semigroup blocks}
\label{sec:s7v33-semigroup}

In lattice units the continuous-time free transfer over the doubled
step \(2a\) is
\begin{equation}
 E(\theta):=\exp\{-2\sqrt{L(\theta)}\}.                 \label{eq:s7v33-E}
\end{equation}

\begin{theorem}[Exact coarse--detail transfer block]
\label{thm:s7v33-offdiag}
The off-diagonal block in the basis \eqref{eq:s7v33-U} has magnitude
\begin{align}
 \varepsilon_{\rm H}(c)
 &:=
 |Q_0^*EH_0|                                            \label{eq:s7v33-epsilon}\\
 &=
 {1-c^2\over2(1+c^2)}
 \left[
 e^{-2\sqrt{\,2-2c\,}}
 -
 e^{-2\sqrt{\,2+2c\,}}
 \right].                                               \label{eq:s7v33-epsilon-formula}
\end{align}
Uniformly in momentum,
\begin{equation}
 \varepsilon_{\rm H}(c)
 \le
 {1+\sqrt2\over8e^2}
 <0.041.                                                \label{eq:s7v33-epsilon-bound}
\end{equation}
\end{theorem}

\begin{proof}
For any scalar function \(f\) on the two eigenvalues,
\[
 f(L)
 =
 \alpha I+\beta L,\qquad
 \beta=
 {f(\lambda_+)-f(\lambda_-)\over\lambda_+-\lambda_-}.
\]
Only the \(\beta L\) term contributes off diagonal.  Insert
\eqref{eq:s7v33-B}, \(\lambda_+-\lambda_-=4c\), and
\(f(\lambda)=e^{-2\sqrt\lambda}\); continuity supplies the endpoint
values.  This proves \eqref{eq:s7v33-epsilon-formula}.

For the bound, put \(u=\sqrt{2-2c}\).  Discarding the second positive
exponential gives
\[
 \varepsilon_{\rm H}(c)
 \le
 {1-c^2\over2(1+c^2)}e^{-2u}.
\]
The elementary maximum
\[
 \max_{0\le c\le1}{1+c\over2(1+c^2)}
 ={1+\sqrt2\over4}
\]
and \(1-c=u^2/2\) imply
\[
 {1-c^2\over2(1+c^2)}
 \le {1+\sqrt2\over8}u^2.
\]
Finally \(\max_{u\ge0}u^2e^{-2u}=e^{-2}\), proving
\eqref{eq:s7v33-epsilon-bound}.
\end{proof}

\begin{theorem}[Exact compressed high transfer]
\label{thm:s7v33-high}
The high-channel diagonal block is
\begin{align}
 r_{\rm H}(c)
 &:=
 H_0^*EH_0                                               \label{eq:s7v33-rH}\\
 &=
 {(1-c)^2\over2(1+c^2)}
 e^{-2\sqrt{\,2-2c\,}}
 +
 {(1+c)^2\over2(1+c^2)}
 e^{-2\sqrt{\,2+2c\,}}.                                 \label{eq:s7v33-rH-formula}
\end{align}
It obeys the volume- and cutoff-independent estimate
\begin{equation}
 r_{\rm H}(c)
 \le {3\over2}e^{-2\sqrt2}
 <0.089.                                                \label{eq:s7v33-rH-bound}
\end{equation}
\end{theorem}

\begin{proof}
The spectral weight of \(H_0\) on the
\(\lambda_-\)-eigenspace is
\[
 w_-={\lambda_+-C\over\lambda_+-\lambda_-}
 ={(1-c)^2\over2(1+c^2)},
\]
and \(w_+=1-w_-=(1+c)^2/[2(1+c^2)]\).  Spectral calculus gives
\eqref{eq:s7v33-rH-formula}.

For the second term, \(\lambda_+\ge2\), so it is at most
\(e^{-2\sqrt2}\).  For the first, with
\(u=\sqrt{2-2c}\),
\[
 w_-\le{(1-c)^2\over2}={u^4\over8}.
\]
On \(0\le u\le\sqrt2\), the function \(u^4e^{-2u}\) is increasing
because its unconstrained maximum is at \(u=2\).  Hence the first term
is at most
\[
 {(\sqrt2)^4\over8}e^{-2\sqrt2}
 ={1\over2}e^{-2\sqrt2}.
\]
Adding the two bounds proves \eqref{eq:s7v33-rH-bound}.
\end{proof}

\begin{remark}[Diagnostic sharp values]
\label{rem:s7v33-numerics}
Direct evaluation of the exact formulas gives
\[
 \sup_c\varepsilon_{\rm H}(c)\approx0.0328111,\qquad
 \sup_cr_{\rm H}(c)=e^{-2\sqrt2}\approx0.0591057.
\]
These values are not used in any proof; the rigorous bounds are
\eqref{eq:s7v33-epsilon-bound} and
\eqref{eq:s7v33-rH-bound}.
\end{remark}

\section{Insertion into the V28 step budget}
\label{sec:s7v33-budget}

Define the exact compressed coarse transfer
\begin{equation}
 T_c(\theta):=Q_0(\theta)^*E(\theta)Q_0(\theta).         \label{eq:s7v33-Tc}
\end{equation}
With this definition the coarse diagonal matching error is zero, while
the off-diagonal and residual estimates are
\begin{equation}
 \epsilon_j\le
 {1+\sqrt2\over8e^2},\qquad
 r_j\le{3\over2}e^{-2\sqrt2}.                          \label{eq:s7v33-step-values}
\end{equation}

\begin{theorem}[Fixed Haar lifting fails the ultraviolet physical
margin]
\label{thm:s7v33-budget-failure}
Let \(\widetilde m>0\) be a fixed proposed physical mass.  The V28
step condition requires
\begin{equation}
 q_{j+1}+2\epsilon_j
 \le e^{-2\widetilde m a_j}.                            \label{eq:s7v33-required}
\end{equation}
Assume the retained coarse dynamics has a finite physical spectral
scale, in the minimal quantitative sense
\begin{equation}
 q_{j+1}\ge1-Ca_j                                      \label{eq:s7v33-finite-scale}
\end{equation}
for some \(C<\infty\) independent of the ultraviolet cutoff.  Then the
fixed Haar off-diagonal block cannot satisfy
\eqref{eq:s7v33-required} uniformly as \(a_j\downarrow0\).
Therefore it is not a continuum-uniform approximate intertwiner for
\Cref{thm:s7v28-recursion} in a retained sector with finite physical
energies.
\end{theorem}

\begin{proof}
The available physical margin per doubled ultraviolet step is
\[
 1-e^{-2\widetilde m a_j}
 =2\widetilde m a_j+O(a_j^2)\longrightarrow0.
\]
The exact off-diagonal norm is bounded below, independently of \(a_j\),
by its value at \(c=1/2\):
\begin{equation}
 \epsilon_j\ge\epsilon_0
 :={3\over10}\left(e^{-2}-e^{-2\sqrt3}\right)>0.         \label{eq:s7v33-positive}
\end{equation}
Combining \eqref{eq:s7v33-finite-scale} and
\eqref{eq:s7v33-positive}, the left side of
\eqref{eq:s7v33-required} is at least
\[
 1-Ca_j+2\epsilon_0,
\]
which exceeds one for all sufficiently small \(a_j\), while the right
side is strictly below one.  This is the contradiction.
\end{proof}

\begin{firewall}[Small numerical mixing is not small in physical
units]
\label{fire:s7v33-units}
The constant \(0.033\) is small on one lattice step but does not scale
to zero with the ultraviolet time spacing.  Treating it as a harmless
perturbation at every one of \(O(\log(1/a))\) steps would violate the
physical-margin recursion before nonlinear errors are added.
\end{firewall}

\section{Upstream alternatives}
\label{sec:s7v33-alternatives}

\begin{theorem}[What a successful free local split must change]
\label{thm:s7v33-alternatives}
At least one of the following changes is required before adding the
Yang--Mills interaction:
\begin{enumerate}[label=\textup{(\roman*)},leftmargin=3.5em]
\item use a family of increasingly accurate filters whose
      coarse--detail semigroup block is \(O(a_j)\);
\item use an exact spectral or perfect-action intertwiner and prove
      sufficient quasi-locality instead of finite range;
\item generalize the transfer compiler to the uniformly conditioned
      nonorthogonal lifting decomposition and prove that its exact
      triangular chain identity controls physical correlations; or
\item group several ultraviolet steps into a single map whose total
      error is measured against a fixed physical time.
\end{enumerate}
\end{theorem}

\begin{proof}
The fixed finite-range orthogonal Haar choice fails by
\Cref{thm:s7v33-budget-failure}.  The four alternatives remove,
respectively, the fixed error, the localization approximation, the
orthogonalization debit, or the vanishing one-step physical margin.
\end{proof}

\begin{assessment}[Progress at seventh-series V33]
\label{ass:s7v33-status}
The actual free transfer debit of the local Haar split is now exact.
The doubled-step high compression is uniformly strict and the
coarse--detail transfer is much smaller than the raw derivative block.
Nevertheless its nonzero constant size is fatal to the per-step
continuum margin.  The next construction must make the filter accuracy
scale with \(a\), retain an exact quasi-local perfect split, or replace
the orthogonal one-step compiler.
\end{assessment}

\begin{decision}[V33 bottleneck decision]
\label{dec:s7v33-bottleneck}
V34 will pursue alternative (ii).  Starting from the exact Gaussian
spectral projection of V29, it will replace the sharp band edge by an
analytic two-band contour separated by a fixed annulus and prove
exponential kernel decay for the associated spectral flattening.  It
will determine whether the localization length can grow only
logarithmically in \(1/a\) while forcing the transfer off-diagonal
below \(O(a)\).
\end{decision}

\begin{firewall}[Claim boundary after seventh-series V33]
\label{fire:s7v33-boundary}
V33 computes the exact orthogonal Haar polyphase Laplacian, the exact
doubled-step transfer blocks, rigorous uniform high and off-diagonal
bounds, and the failure of the fixed filter to meet the ultraviolet
physical-margin budget.

V33 does not construct a scale-accurate quasi-local filter, control
nonlinear RG errors, generate terminal nonperturbative mass, remove the
cutoff, or prove the full Yang--Mills mass gap.
\end{firewall}

\paragraph{Parent-version record.}
V33 was formed from
\texttt{Renewal\_Yang\_Mills\_Mass\_Gap\_Research\_Notes\_V32.tex}.
The V32 parent had \(12\,813\) logical source lines, \(467\,894\) bytes,
and SHA--256
\[
 \texttt{E4B3DCE1B7710B905A3220B380B67A4B12E91617F5258F86DC239997CD231AD5}.
\]

% =============================================================================
% END APPENDED SEVENTH-SERIES V33 MATERIAL
% =============================================================================

\newpage

% =============================================================================
% BEGIN APPENDED SEVENTH-SERIES V34 MATERIAL
% =============================================================================

\chapter{Buffered analytic spectral flattening}
\label{ch:s7v34}

V33 shows that a fixed local Haar split has a nonzero transfer
off-diagonal and therefore cannot fit an \(O(a)\) ultraviolet
physical-margin budget.  This note asks whether filter accuracy can
scale with the cutoff while the filter remains quasi-local.  With a
fixed spectral buffer the answer is yes: a logistic spectral
flattening has leakage \(O(a^p)\) and exponential localization length
\(O(\log(1/a))\).  A literal finite-range polynomial approximation
requires range \(O((\log(1/a))^2)\).

There is also an unavoidable qualification.  The free lattice
Laplacian has spectrum throughout the buffer.  No continuous
multiplier can approximate the sharp low-band projection in global
operator norm below \(1/2\).  The buffer must therefore be a third
transition channel, not silently assigned by an approximate
two-channel projection.

\section{Finite-propagation functional calculus}
\label{sec:s7v34-functional}

Let \({\cal L}\) be a bounded nonnegative self-adjoint operator on
\(\ell^2(\mathbb Z^{d_s})\otimes\mathbb C^m\), with
\begin{equation}
 \operatorname{spec}{\cal L}\subset[0,\Lambda],\qquad
 \langle\delta_x,{\cal L}\delta_y\rangle=0
 \quad\hbox{if }\operatorname{dist}(x,y)>r_0.            \label{eq:s7v34-propagation}
\end{equation}
The dimensionless spatial lattice Laplacian in V29 is the principal
example.

\begin{lemma}[Analytic functional calculus is exponentially
quasi-local]
\label{lem:s7v34-locality}
Suppose \(f\) is analytic and bounded by \(M\) in the complex
\(\rho\)-neighbourhood of \([0,\Lambda]\), with
\(0<\rho\le\Lambda\).  Then there are numerical constants
\(C,c>0\) such that
\begin{equation}
 |\langle\delta_x,f({\cal L})\delta_y\rangle|
 \le
 CM\exp\left\{
 -{c\rho\over\Lambda r_0}\operatorname{dist}(x,y)
 \right\}.                                               \label{eq:s7v34-locality}
\end{equation}
The same statement holds uniformly on periodic lattices larger than
the displayed distance.
\end{lemma}

\begin{proof}
Rescale \([0,\Lambda]\) to \([-1,1]\).  Analyticity in a
\(\rho\)-neighbourhood supplies a Bernstein ellipse whose parameter is
at least \(\exp(c\rho/\Lambda)\).  Truncating the Chebyshev expansion
at degree \(N\) gives a polynomial \(p_N\) with
\[
 \|f-p_N\|_{[0,\Lambda]}
 \le CM e^{-c\rho N/\Lambda}.
\]
Because \({\cal L}\) has propagation \(r_0\), the polynomial
\(p_N({\cal L})\) has propagation \(Nr_0\).  If
\(\operatorname{dist}(x,y)>Nr_0\), its \((x,y)\) matrix element
vanishes, while the matrix element of the difference is bounded by
the operator norm.  Choose
\(N=\lfloor\operatorname{dist}(x,y)/r_0\rfloor-1\) and absorb bounded
short distances into \(C\).
\end{proof}

\section{Logistic low-pass with a spectral buffer}
\label{sec:s7v34-logistic}

Fix
\begin{equation}
 0<\alpha<\beta<\Lambda,\qquad
 m_0:={\alpha+\beta\over2},\qquad
 g:={\beta-\alpha\over2}.                               \label{eq:s7v34-buffer}
\end{equation}
Define, for \(\kappa>0\),
\begin{equation}
 f_\kappa(\lambda)
 :=
 {1\over1+e^{\kappa(\lambda-m_0)}},\qquad
 F_\kappa:=f_\kappa({\cal L}).                          \label{eq:s7v34-logistic}
\end{equation}
Let
\[
 P_{\rm lo}:=\mathbf1_{[0,\alpha]}({\cal L}),\quad
 P_{\rm tr}:=\mathbf1_{(\alpha,\beta)}({\cal L}),\quad
 P_{\rm hi}:=\mathbf1_{[\beta,\Lambda]}({\cal L}).
\]

\begin{theorem}[Buffered leakage and localization]
\label{thm:s7v34-buffered}
The operator \(F_\kappa\) is a self-adjoint positive contraction and
\begin{align}
 \|(I-F_\kappa)P_{\rm lo}\|
 &\le e^{-\kappa g},                                    \label{eq:s7v34-low-leak}\\
 \|F_\kappa P_{\rm hi}\|
 &\le e^{-\kappa g}.                                    \label{eq:s7v34-high-leak}
\end{align}
Its kernel obeys
\begin{equation}
 |\langle\delta_x,F_\kappa\delta_y\rangle|
 \le
 C\exp\left\{
 -{c\over\kappa\Lambda r_0}
 \operatorname{dist}(x,y)
 \right\}.                                               \label{eq:s7v34-logistic-locality}
\end{equation}
\end{theorem}

\begin{proof}
For \(\lambda\le\alpha=m_0-g\),
\[
 1-f_\kappa(\lambda)
 ={e^{\kappa(\lambda-m_0)}
   \over1+e^{\kappa(\lambda-m_0)}}
 \le e^{-\kappa g}.
\]
For \(\lambda\ge\beta=m_0+g\),
\[
 f_\kappa(\lambda)\le e^{-\kappa g}.
\]
Spectral calculus proves
\eqref{eq:s7v34-low-leak}--\eqref{eq:s7v34-high-leak}.

The poles of \(f_\kappa\) are
\[
 m_0+{(2j+1)\pi i\over\kappa}.
\]
In the strip \(|\operatorname{Im}\lambda|\le\pi/(2\kappa)\), the
phase of \(e^{\kappa(\lambda-m_0)}\) lies between
\(-\pi/2\) and \(\pi/2\), so the denominator has modulus at least one.
Thus \(f_\kappa\) is analytic and bounded there.  Apply
\Cref{lem:s7v34-locality} with
\(\rho=\min\{\Lambda,\pi/(2\kappa)\}\); enlarging \(C\) handles bounded
\(\kappa\) and gives \eqref{eq:s7v34-logistic-locality}.
\end{proof}

\begin{corollary}[Cutoff-scaled accuracy]
\label{cor:s7v34-scaled}
Let \(0<a<e^{-1}\), \(p>0\), and choose
\begin{equation}
 \kappa(a):={p\over g}\log{1\over a}.                   \label{eq:s7v34-kappa}
\end{equation}
Then
\begin{equation}
 \|(I-F_{\kappa(a)})P_{\rm lo}\|
 +\|F_{\kappa(a)}P_{\rm hi}\|
 \le2a^p,                                               \label{eq:s7v34-ap}
\end{equation}
while the exponential localization length is
\begin{equation}
 \xi(a)=O(\log(1/a)).                                   \label{eq:s7v34-xi}
\end{equation}
\end{corollary}

\begin{proof}
Insert \eqref{eq:s7v34-kappa} in
\eqref{eq:s7v34-low-leak}--\eqref{eq:s7v34-high-leak}.  The exponent
in \eqref{eq:s7v34-logistic-locality} has inverse rate proportional to
\(\kappa(a)\).
\end{proof}

\section{Literal finite range costs a second logarithm}
\label{sec:s7v34-finite}

\begin{theorem}[Finite-propagation approximation]
\label{thm:s7v34-polynomial}
For every \(N\ge1\), there is a polynomial \(p_{\kappa,N}\) of degree
\(N\) such that
\begin{equation}
 \|F_\kappa-p_{\kappa,N}({\cal L})\|
 \le C e^{-cN/\kappa},                                  \label{eq:s7v34-polynomial}
\end{equation}
and \(p_{\kappa,N}({\cal L})\) has propagation at most \(Nr_0\).
With \(\kappa(a)\) from \eqref{eq:s7v34-kappa}, achieving truncation
error at most \(a^p\) is possible with
\begin{equation}
 N(a)=O((\log(1/a))^2).                                 \label{eq:s7v34-log2}
\end{equation}
\end{theorem}

\begin{proof}
Use the Chebyshev approximation in the proof of
\Cref{lem:s7v34-locality} with analytic width
\(\rho\asymp\kappa^{-1}\).  This gives
\eqref{eq:s7v34-polynomial}.  A degree-\(N\) polynomial in a
range-\(r_0\) operator has range \(Nr_0\).  To make the right side at
most \(a^p\), it suffices that
\[
 N\ge C\kappa(a)\log(1/a),
\]
which is \eqref{eq:s7v34-log2}.
\end{proof}

\begin{assessment}[Quasi-local versus finite-range cost]
\label{ass:s7v34-locality}
An exponentially decaying exact filter needs characteristic length
only \(O(\log(1/a))\).  Replacing it by a strictly finite-range
operator with the same \(O(a^p)\) norm accuracy multiplies that length
by a second logarithm.  Either growth is subpower in \(a^{-1}\), but
its incidence and nonlinear action cost must still be entered in the
NTRGC ledger.
\end{assessment}

\section{Why a transition channel is unavoidable}
\label{sec:s7v34-transition}

\begin{theorem}[No global continuous approximation below one half]
\label{thm:s7v34-half}
Let the spectrum of \({\cal L}\) contain points arbitrarily close to
\(\lambda_*\) from both sides.  Put
\[
 P_*:=\mathbf1_{[0,\lambda_*]}({\cal L}).
\]
For every continuous scalar function \(h\) on
\(\operatorname{spec}{\cal L}\),
\begin{equation}
 \|h({\cal L})-P_*\|\ge{1\over2}.                       \label{eq:s7v34-half}
\end{equation}
\end{theorem}

\begin{proof}
If the norm were less than \(1/2\), then spectral calculus would give
\[
 |h(\lambda)-1|<1/2\quad(\lambda\le\lambda_*),
 \qquad
 |h(\lambda)|<1/2\quad(\lambda>\lambda_*).
\]
Along spectral sequences approaching \(\lambda_*\) from the two sides,
continuity would force \(h(\lambda_*)\) to be simultaneously strictly
above and at most \(1/2\), a contradiction.  The non-strict inequality
follows by closure.
\end{proof}

\begin{corollary}[The logistic estimate is necessarily buffered]
\label{cor:s7v34-buffer-necessary}
No choice of \(\kappa\) makes \(F_\kappa\) an
\(o(1)\) operator-norm approximation to a sharp projection whose edge
lies in continuous spectrum.  Equations
\eqref{eq:s7v34-low-leak}--\eqref{eq:s7v34-high-leak} are possible
only because the transition interval \((\alpha,\beta)\) is excluded
from both estimates.
\end{corollary}

\begin{proof}
Apply \Cref{thm:s7v34-half}; the logistic multiplier is continuous.
\end{proof}

\section{Free transfer treatment of the transition band}
\label{sec:s7v34-transfer}

Let the physical free Hamiltonian be
\[
 H_a^{(0)}=a^{-1}\sqrt{\cal L}
\]
on the one-particle sector, and evolve for block time \(ba\).

\begin{theorem}[Uniform transition/high contraction]
\label{thm:s7v34-transition-contraction}
The entire non-low sector
\(P_{\rm tr}+P_{\rm hi}=\mathbf1_{(\alpha,\Lambda]}({\cal L})\)
obeys
\begin{equation}
 \left\|
 e^{-baH_a^{(0)}}(P_{\rm tr}+P_{\rm hi})
 \right\|
 \le e^{-b\sqrt\alpha}<1.                               \label{eq:s7v34-transition-contraction}
\end{equation}
This bound is independent of \(a\) and volume.
\end{theorem}

\begin{proof}
Every spectral value in the declared sector is at least \(\alpha\).
The transfer multiplier is \(e^{-b\sqrt\lambda}\), so spectral
calculus gives \eqref{eq:s7v34-transition-contraction}.
\end{proof}

\begin{theorem}[Exact soft two-channel isometry]
\label{thm:s7v34-soft-isometry}
The analysis map
\begin{equation}
 {\cal W}_\kappa\psi
 :=
 \left(
 F_\kappa^{1/2}\psi,\,
 (I-F_\kappa)^{1/2}\psi
 \right)                                                \label{eq:s7v34-W}
\end{equation}
is an isometry into \({\cal H}\oplus{\cal H}\) and exactly intertwines
every bounded Borel function of \({\cal L}\), including the free
transfer.  Its two output channels are not orthogonal independent
subspaces: their overlap multiplier is
\begin{equation}
 \sqrt{f_\kappa(1-f_\kappa)},                           \label{eq:s7v34-overlap}
\end{equation}
whose norm is \(1/2\) whenever the spectrum meets \(m_0\).
\end{theorem}

\begin{proof}
\[
 {\cal W}_\kappa^*{\cal W}_\kappa
 =F_\kappa+(I-F_\kappa)=I.
\]
All three operators in \eqref{eq:s7v34-W} are functions of
\({\cal L}\), so they commute with every other spectral multiplier.
The cross-channel synthesis product is
\(F_\kappa^{1/2}(I-F_\kappa)^{1/2}\), giving
\eqref{eq:s7v34-overlap}; at \(f_\kappa=1/2\) it equals \(1/2\).
\end{proof}

\begin{firewall}[A tight analysis frame is not the V28 decomposition]
\label{fire:s7v34-frame}
The isometry \({\cal W}_\kappa\) embeds one field into a doubled channel
space.  It does not decompose the original Hilbert space into an
orthogonal coarse range plus residual range, because the channel
overlap is order one on the transition band.  To use the buffered
\(O(a^p)\) leakage in NTRGC, the transition channel must be retained
explicitly in a three-band or redundant-frame version of the transfer
compiler.
\end{firewall}

\section{Analytic-filter status}
\label{sec:s7v34-status}

\begin{theorem}[What V34 proves]
\label{thm:s7v34-result}
For every fixed spectral buffer and every \(p>0\), there is an exact
quasi-local filter with:
\begin{enumerate}[label=\textup{(\roman*)},leftmargin=3.5em]
\item low/high leakage at most \(O(a^p)\);
\item exponential localization length \(O(\log(1/a))\);
\item a finite-range approximation of range
      \(O((\log(1/a))^2)\) and error \(O(a^p)\);
\item exact commutation with the free transfer; and
\item a uniformly contracting transition-plus-high sector.
\end{enumerate}
It does not give a two-channel orthogonal coarse/residual
decomposition in global operator norm.
\end{theorem}

\begin{proof}
Use
\Cref{thm:s7v34-buffered,cor:s7v34-scaled,
thm:s7v34-polynomial,thm:s7v34-transition-contraction,
thm:s7v34-half,thm:s7v34-soft-isometry}.
\end{proof}

\begin{assessment}[Progress at seventh-series V34]
\label{ass:s7v34-status}
The V33 fixed-error problem can be reduced to \(O(a^p)\) on separated
spectral bands with only logarithmic quasi-locality.  Strict finite
range costs a second logarithm.  The obstruction is now isolated to
the populated transition annulus: continuity forbids approximating
its sharp assignment, while a soft tight-frame split has channel
overlap \(1/2\) there.  Fortunately all transition modes have a
uniform free transfer contraction.  The next compiler must treat that
channel as residual instead of forcing a two-band projection.
\end{assessment}

\begin{decision}[V34 bottleneck decision]
\label{dec:s7v34-bottleneck}
V35 is a mandatory SM/NS audit.  After that audit it will formulate a
three-channel transfer recursion with a retained low subspace, a
buffer transition subspace, and a high subspace.  The target is to
insert the \(O(a^p)\) buffered leakage and the exact
\(e^{-b\sqrt\alpha}\) residual contraction without an order-one
channel-overlap penalty.
\end{decision}

\begin{firewall}[Claim boundary after seventh-series V34]
\label{fire:s7v34-boundary}
V34 proves the logarithmic quasi-locality and squared-log finite-range
cost of cutoff-accurate buffered spectral flattening, the global
one-half obstruction at a populated band edge, and uniform free
contraction of the transition sector.

V34 does not construct the required three-channel transfer compiler,
control nonlinear interactions and quasi-local incidence, generate
the terminal nonperturbative mass, remove the cutoff, or prove the full
Yang--Mills mass gap.
\end{firewall}

\paragraph{Parent-version record.}
V34 was formed from
\texttt{Renewal\_Yang\_Mills\_Mass\_Gap\_Research\_Notes\_V33.tex}.
The V33 parent had \(13\,188\) logical source lines, \(479\,953\) bytes,
and SHA--256
\[
 \texttt{7DE055211744088F4E556474960952BBA22BB40E58CCF2F6A8EEB70D9E57164F}.
\]

% =============================================================================
% END APPENDED SEVENTH-SERIES V34 MATERIAL
% =============================================================================

% =============================================================================
% BEGIN APPENDED SEVENTH-SERIES V35 MATERIAL
% =============================================================================

\chapter{Three-channel transfer recursion and the physical-time pivot}
\label{ch:s7v35}

V34 proved that a fixed spectral buffer permits \(O(a^p)\) separated-band
leakage with logarithmic quasi-locality, but that a continuous two-channel
filter has order-one overlap on the populated transition annulus.  The first
task of this note is the compulsory companion audit.  The second is to write
the exact three-channel transfer inequality and determine precisely which
errors must scale with the ultraviolet physical-time margin.

The result corrects an over-optimistic reading of V34.  Making the transition
annulus an explicit orthogonal channel removes the artificial \(1/2\)
tight-frame overlap.  It does not, by itself, control the true
low--transition block of the interacting transfer operator.  A one-lattice-
step recursion still charges every block incident on the retained low sector
against an \(O(a)\) margin.  This sends the programme upstream to a fixed
physical-time slab.  On such a slab the transition and high sectors have
energy \(O(a^{-1})\), and a direct Duhamel argument suppresses their transfer
blocks by \(a\|V_aP_{\rm r}\|\).  Thus an asymptotically free interaction of
size \(O(g(a)/a)\) costs \(O(g(a))\) against a fixed, rather than vanishing,
physical margin.

\section{Compulsory V35 companion audit}
\label{sec:s7v35-audit}

\subsection{Files inspected}
\label{sec:s7v35-files}

The newest active companion files in \texttt{latex\_notes} were inspected
read-only:
\begin{align}
 &\texttt{Renewal\_SM\_Einstein\_Continuum\_Research\_Notes\_V79.tex},
 \notag\\[-2pt]
 &\qquad 31\,495\ \hbox{logical source lines},\quad
 1\,113\,043\ \hbox{bytes},                             \notag\\[-2pt]
 &\qquad
 \texttt{3EEE24069FC5CA08FE2CF21DEE74707D2EE557F063EA4BACDD400AF4E9636FA4},
 \label{eq:s7v35-sm-checkpoint}\\
 &\texttt{Renewal\_NS\_Stability\_Research\_Notes\_V26.tex},
 \notag\\[-2pt]
 &\qquad 5\,319\ \hbox{logical source lines},\quad
 278\,283\ \hbox{bytes},                                \notag\\[-2pt]
 &\qquad
 \texttt{82C887F8C2C69E4F28FAD7C3DC8DE5B9099CB5A4BF431EBA1FFF3E91935978AD}.
 \label{eq:s7v35-ns-checkpoint}
\end{align}
The next compulsory companion audit is V40.

\subsection{SM V76--V79}
\label{sec:s7v35-sm}

Since the V30 audit, the SM programme has replaced patchwise scalar-curvature
fluxes by a global reference-connection weak form, split the constant mode
from the mean-zero sector, derived the exact scalar Schur complement, and
constructed a strong fixed-volume implicit-function chart with an explicit
multiplier.  The reusable lesson is structural:
\begin{equation}
 \hbox{diagonal block}
 \quad+\quad
 \hbox{mixing debit}
 \quad+\quad
 \hbox{separate kernel or multiplier row}.             \label{eq:s7v35-sm-lesson}
\end{equation}
In particular, invertibility of a restricted block does not erase its
coupling to a distinguished low mode.  This is directly analogous to the
transition-channel issue in V34.  No Einstein or ADM estimate is imported
into Yang--Mills theory.

\subsection{NS V26}
\label{sec:s7v35-ns}

The NS programme remains at V26.  Its newest calculation replaces unrestricted
tensor norms by rank-one norms for derivatives of the Oseen kernel.  The first
derivative gains nothing and the second gains only a modest amount on a
bounded radial interval.  The transferable lesson is that a restricted norm
must be computed only after the admissible input channel has been identified;
numerical smallness in a larger norm does not repair an incorrectly typed
decomposition.  No Navier--Stokes kernel constant enters the present transfer
operator.

\begin{theorem}[V35 companion-audit decision]
\label{thm:s7v35-audit}
The companion audit changes the order of attack but imports no mathematical
estimate.  The Yang--Mills transfer is first decomposed into exact orthogonal
low, transition, and high blocks; every off-diagonal block is retained in the
ledger.  Only after this type-correct decomposition may restricted or
quasi-local estimates be used.
\end{theorem}

\begin{proof}
The SM Schur formula shows that a contracting or coercive complementary block
still contributes through its coupling to the distinguished mode.  The NS
rank-one calculation shows that restriction can sharpen a constant only
after the restriction is exact.  Applying both lessons to V34 gives the
stated order.  The field equations, state spaces, and norms in the companion
programmes differ from the Yang--Mills transfer problem, so their numerical
constants cannot be imported.
\end{proof}

\section{The exact three-channel transfer bound}
\label{sec:s7v35-three}

Let \(T\) be a positive self-adjoint contraction on a Hilbert space
\({\cal H}\), let \(T\Omega=\Omega\), and suppose the vacuum is unique.
On \({\cal H}_\circ:=\Omega^\perp\), fix an orthogonal decomposition
\begin{equation}
 {\cal H}_\circ
 =
 {\cal H}_{\rm L}\oplus{\cal H}_{\rm T}\oplus{\cal H}_{\rm H},
 \qquad
 P_{\rm L}+P_{\rm T}+P_{\rm H}=I_{{\cal H}_\circ}.      \label{eq:s7v35-decomposition}
\end{equation}
For an integer blocking factor \(b\ge2\), put
\begin{equation}
 A:=T^b\big|_{{\cal H}_\circ},\qquad
 A_{ij}:=P_iAP_j
 \quad(i,j\in\{{\rm L,T,H}\}).                          \label{eq:s7v35-blocks}
\end{equation}
Define
\begin{align}
 q_i&:=\|A_{ii}\|,                                      \label{eq:s7v35-qi}\\
 \epsilon_{ij}&:=\|A_{ij}\|=\|A_{ji}\|
 \quad(i\ne j).                                        \label{eq:s7v35-eps}
\end{align}

\begin{theorem}[Three-channel row-sum compiler]
\label{thm:s7v35-row}
With the preceding notation,
\begin{equation}
 \|A\|
 \le
 \max\left\{
 \begin{aligned}
  &q_{\rm L}+\epsilon_{\rm LT}+\epsilon_{\rm LH},\\
  &q_{\rm T}+\epsilon_{\rm LT}+\epsilon_{\rm TH},\\
  &q_{\rm H}+\epsilon_{\rm LH}+\epsilon_{\rm TH}
 \end{aligned}
 \right\}.                                              \label{eq:s7v35-row}
\end{equation}
No independence, commutation, or probabilistic assumption is used.
\end{theorem}

\begin{proof}
Write \(x=x_{\rm L}+x_{\rm T}+x_{\rm H}\).  Positivity of \(A\) gives
\(\|A\|=\sup_{\|x\|=1}\langle Ax,x\rangle\).  Self-adjointness and
Cauchy--Schwarz give
\begin{align*}
 \langle Ax,x\rangle
 &\le
 \sum_iq_i\|x_i\|^2
 +2\sum_{i<j}\epsilon_{ij}\|x_i\|\|x_j\|\\
 &\le
 \sum_i
 \left(q_i+\sum_{j\ne i}\epsilon_{ij}\right)\|x_i\|^2,
\end{align*}
where \(2uv\le u^2+v^2\) was used in the second line.  Divide by
\(\sum_i\|x_i\|^2\) and take the supremum.
\end{proof}

\begin{remark}[Why the transition channel removes the frame penalty]
\label{rem:s7v35-frame}
For the exact spectral projections of V34 the three spaces in
\eqref{eq:s7v35-decomposition} are orthogonal.  Hence the \(1/2\) overlap of
the redundant soft analysis map is absent from
\eqref{eq:s7v35-row}.  Any nonzero \(\epsilon_{ij}\) is now a genuine transfer
block, not a synthesis overlap.
\end{remark}

\subsection{Coarse matching}
\label{sec:s7v35-matching}

Let \(T_c\) be the next-scale transfer on its vacuum complement, and let
\(Q:{\cal H}_{c,\circ}\to{\cal H}_{\rm L}\) be an isometry.  Suppose
\begin{equation}
 \|A_{\rm LL}-QT_cQ^*\|\le\eta_{\rm L},\qquad
 \|T_c\|=q_c,                                           \label{eq:s7v35-match}
\end{equation}
and
\begin{equation}
 q_{\rm T}\le r_{\rm T},\qquad
 q_{\rm H}\le r_{\rm H}.                               \label{eq:s7v35-residual}
\end{equation}

\begin{corollary}[Three-channel RG recursion]
\label{cor:s7v35-recursion}
Define
\begin{equation}
 R:=
 \max\left\{
 \begin{aligned}
  &q_c+\eta_{\rm L}+\epsilon_{\rm LT}+\epsilon_{\rm LH},\\
  &r_{\rm T}+\epsilon_{\rm LT}+\epsilon_{\rm TH},\\
  &r_{\rm H}+\epsilon_{\rm LH}+\epsilon_{\rm TH}
 \end{aligned}
 \right\}.                                              \label{eq:s7v35-R}
\end{equation}
Then
\begin{equation}
 \|T|_{{\cal H}_\circ}\|\le R^{1/b}.                   \label{eq:s7v35-recursion}
\end{equation}
\end{corollary}

\begin{proof}
Equation \eqref{eq:s7v35-match} gives
\(q_{\rm L}\le q_c+\eta_{\rm L}\).  Insert this and
\eqref{eq:s7v35-residual} into \Cref{thm:s7v35-row}.  Since \(T\) is
positive,
\(\|T|_{{\cal H}_\circ}\|^b=\|T^b|_{{\cal H}_\circ}\|\).
\end{proof}

\begin{firewall}[Every low-incident block consumes low-mode margin]
\label{fire:s7v35-low-row}
Uniform contraction of the transition and high diagonal blocks does not
remove \(\epsilon_{\rm LT}\) or \(\epsilon_{\rm LH}\) from the first row of
\eqref{eq:s7v35-R}.  In a one-step ultraviolet recursion that row is close
to one, so both low-incident transfer blocks, as well as the coarse matching
error, must fit the physical mass margin.
\end{firewall}

\section{Physical mass bookkeeping and ultraviolet summability}
\label{sec:s7v35-mass}

Let \(a_{j+1}=ba_j\), let \(q_j\) denote the fine one-step norm on the vacuum
complement, and define its effective physical mass by
\begin{equation}
 m_j:=-{1\over a_j}\log q_j.                            \label{eq:s7v35-mj}
\end{equation}

\begin{theorem}[Exact mass update]
\label{thm:s7v35-mass-update}
If the \(j\)-th three-channel compiler has row maximum \(R_j<1\), then
\begin{equation}
 m_j\ge-{1\over a_{j+1}}\log R_j.                       \label{eq:s7v35-mass-update}
\end{equation}
If, in addition,
\begin{equation}
 R_j\le e^{-m_{j+1}a_{j+1}}+\delta_j,                  \label{eq:s7v35-delta}
\end{equation}
then
\begin{equation}
 m_j
 \ge
 m_{j+1}
 -
 {1\over a_{j+1}}
 \log\left(1+\delta_je^{m_{j+1}a_{j+1}}\right)
 \ge
 m_{j+1}
 -
 {\delta_je^{m_{j+1}a_{j+1}}\over a_{j+1}}.            \label{eq:s7v35-mass-loss}
\end{equation}
\end{theorem}

\begin{proof}
\Cref{cor:s7v35-recursion} gives
\(q_j\le R_j^{1/b}\).  Taking \(-a_j^{-1}\log\) and using
\(ba_j=a_{j+1}\) proves \eqref{eq:s7v35-mass-update}.  Factor
\(e^{-m_{j+1}a_{j+1}}\) from the right side of
\eqref{eq:s7v35-delta}; monotonicity of \(-\log\) gives the first inequality
in \eqref{eq:s7v35-mass-loss}.  The second is
\(\log(1+x)\le x\).
\end{proof}

\begin{corollary}[Power-small one-step errors are ultraviolet summable]
\label{cor:s7v35-summable}
Suppose that on the levels with \(a_j\le a_*\),
\begin{equation}
 \delta_j\le Ca_j^p,\qquad p>1,\qquad
 m_{j+1}a_{j+1}\le M.                                  \label{eq:s7v35-power}
\end{equation}
Then the total mass loss over those levels is at most
\begin{equation}
 {Ce^M\over b}\,
 {a_*^{p-1}\over1-b^{-(p-1)}}.                         \label{eq:s7v35-sum}
\end{equation}
In particular, the ultraviolet tail of the debit can be made arbitrarily
small by choosing \(a_*\) small.
\end{corollary}

\begin{proof}
Sum the last term of \eqref{eq:s7v35-mass-loss}.  Since
\(a_{j+1}=ba_j\), each summand is at most
\((Ce^M/b)a_j^{p-1}\).  The scales form a geometric sequence, and summing
backward from the largest scale \(a_*\) gives
\[
 \sum_{a_j\le a_*}a_j^{p-1}
 \le {a_*^{p-1}\over1-b^{-(p-1)}}.
\]
\end{proof}

\begin{assessment}[What the V34 \(O(a^p)\) target would buy]
\label{ass:s7v35-summability}
If all three terms incident on the low row in
\eqref{eq:s7v35-R} were \(O(a^p)\) with \(p>1\), they would cause only a
summable ultraviolet mass loss.  The exponent \(p>1\), not merely \(p>0\),
is the type-correct requirement after conversion from dimensionless transfer
error to physical mass.
\end{assessment}

\section{What the buffered filter does and does not control}
\label{sec:s7v35-buffer}

Return to the V34 operator \({\cal L}\) and its exact spectral projections
\[
 P_{\rm L}=\mathbf1_{[0,\alpha]}({\cal L}),\qquad
 P_{\rm T}=\mathbf1_{(\alpha,\beta)}({\cal L}),\qquad
 P_{\rm H}=\mathbf1_{[\beta,\Lambda]}({\cal L}).
\]
For the free blocked transfer
\begin{equation}
 A_0:=e^{-b\sqrt{\cal L}},                              \label{eq:s7v35-A0}
\end{equation}
all three projections commute with \(A_0\).

\begin{proposition}[Exact free three-band ledger]
\label{prop:s7v35-free}
For \(A_0\),
\begin{equation}
 \epsilon_{\rm LT}^{(0)}
 =\epsilon_{\rm LH}^{(0)}
 =\epsilon_{\rm TH}^{(0)}=0,                           \label{eq:s7v35-free-off}
\end{equation}
and
\begin{equation}
 r_{\rm T}^{(0)}\le e^{-b\sqrt\alpha},\qquad
 r_{\rm H}^{(0)}\le e^{-b\sqrt\beta}.                  \label{eq:s7v35-free-r}
\end{equation}
\end{proposition}

\begin{proof}
All operators are scalar Borel functions of \({\cal L}\), proving
\eqref{eq:s7v35-free-off}.  The transfer multiplier is
\(e^{-b\sqrt\lambda}\), and the lower spectral edges of the transition and
high sectors are \(\alpha\) and \(\beta\), respectively.
\end{proof}

\subsection{Two separated-band estimates}
\label{sec:s7v35-separated}

\begin{lemma}[Logistic commutator estimate]
\label{lem:s7v35-logistic-commutator}
Let \(A\) be bounded and let \(F_\kappa\) be the V34 logistic filter.  Put
\(\delta_\kappa=e^{-\kappa(\beta-\alpha)/2}\).  Then
\begin{equation}
 \|P_{\rm L}AP_{\rm H}\|
 \le
 2\delta_\kappa\|A\|+\|[F_\kappa,A]\|.                 \label{eq:s7v35-logistic-commutator}
\end{equation}
\end{lemma}

\begin{proof}
Because \(F_\kappa\) commutes with the spectral projections,
\begin{align*}
 P_{\rm L}AP_{\rm H}
 &=
 P_{\rm L}(I-F_\kappa)AP_{\rm H}
 +P_{\rm L}F_\kappa AP_{\rm H}\\
 &=
 P_{\rm L}(I-F_\kappa)AP_{\rm H}
 +P_{\rm L}AF_\kappa P_{\rm H}
 +P_{\rm L}[F_\kappa,A]P_{\rm H}.
\end{align*}
V34 bounds the first and second terms by
\(\delta_\kappa\|A\|\); the third is bounded by the commutator norm.
\end{proof}

\begin{lemma}[Spectral-separation commutator estimate]
\label{lem:s7v35-gap-commutator}
If \([{\cal L},A]\) is bounded, then
\begin{equation}
 \|P_{\rm L}AP_{\rm H}\|
 \le
 {1\over\beta-\alpha}\|P_{\rm L}[{\cal L},A]P_{\rm H}\|.
 \label{eq:s7v35-gap-commutator}
\end{equation}
\end{lemma}

\begin{proof}
Let \(L_{\rm L}\) and \(L_{\rm H}\) be the restrictions of \({\cal L}\) to
the low and high spaces, and put
\[
 X=P_{\rm L}AP_{\rm H},\qquad
 C=P_{\rm L}[{\cal L},A]P_{\rm H}.
\]
Then \(L_{\rm L}X-XL_{\rm H}=C\).  Since
\(\operatorname{spec}L_{\rm L}\subset[0,\alpha]\) and
\(\operatorname{spec}L_{\rm H}\subset[\beta,\Lambda]\),
\begin{equation}
 X=-\int_0^\infty e^{tL_{\rm L}}Ce^{-tL_{\rm H}}\,dt.  \label{eq:s7v35-Sylvester}
\end{equation}
Indeed, differentiating the integrand and integrating from zero to infinity
recovers the Sylvester equation; the endpoint at infinity vanishes by the
spectral separation.  Taking norms in \eqref{eq:s7v35-Sylvester} gives
\[
 \|X\|
 \le\|C\|\int_0^\infty e^{-t(\beta-\alpha)}\,dt,
\]
which is \eqref{eq:s7v35-gap-commutator}.
\end{proof}

\begin{theorem}[The adjacent transition block is an independent debit]
\label{thm:s7v35-adjacent}
The V34 low/high leakage estimates and either of
\Cref{lem:s7v35-logistic-commutator,lem:s7v35-gap-commutator} do not imply
any bound tending to zero for \(P_{\rm L}AP_{\rm T}\).  This failure persists
even for positive self-adjoint contractions.
\end{theorem}

\begin{proof}
On \(\mathbb C^2={\cal H}_{\rm L}\oplus{\cal H}_{\rm T}\), take
\begin{equation}
 A_\gamma=
 \begin{pmatrix}
 1/2&\gamma\\
 \gamma&1/2
 \end{pmatrix},
 \qquad 0<\gamma\le1/2.                                \label{eq:s7v35-counter}
\end{equation}
Its eigenvalues are \(1/2\pm\gamma\), so it is a positive contraction, while
\(\|P_{\rm L}A_\gamma P_{\rm T}\|=\gamma\).  Adjoin any high space on which
\(A_\gamma\) is zero.  The low/high block is then exactly zero, irrespective
of the separated-band filter accuracy, but the low/transition block remains
\(\gamma\).  Thus no estimate for the latter follows from the former.
\end{proof}

\begin{firewall}[Corrected V34 handoff]
\label{fire:s7v35-v34}
V34 supplies a quasi-local separated-band extractor and a strict free
transition contraction.  It does not supply the
\(\epsilon_{\rm LT}=O(a^p)\) hypothesis required by
\Cref{cor:s7v35-summable}.  Claiming the one-step RG bridge from V34 alone
would therefore be circular.
\end{firewall}

\section{Fixed physical-time slab}
\label{sec:s7v35-slab}

The vanishing one-step margin is not compulsory.  Fix a physical Euclidean
time \(\tau>0\).  Let
\begin{equation}
 H_a=H_{0,a}+V_a,\qquad
 S_a(\tau):=e^{-\tau H_a},                             \label{eq:s7v35-H}
\end{equation}
on a finite-volume vacuum complement.  Assume \(H_{0,a}\) and \(H_a\) are
nonnegative self-adjoint, \(V_a\) is bounded and self-adjoint, and
\(P_{\rm L}\) commutes with \(H_{0,a}\).  Put
\begin{equation}
 P_{\rm r}:=I-P_{\rm L}=P_{\rm T}+P_{\rm H}.            \label{eq:s7v35-Pr}
\end{equation}
Assume the free residual energy bound
\begin{equation}
 P_{\rm r}H_{0,a}P_{\rm r}
 \ge {\Delta\over a}P_{\rm r}                           \label{eq:s7v35-free-energy}
\end{equation}
for some \(\Delta>0\).  For the V34 transition-plus-high sector one may take
\(\Delta=\sqrt\alpha\).

\begin{theorem}[Duhamel suppression of the ultraviolet sector]
\label{thm:s7v35-Duhamel}
Let
\begin{equation}
 v_a:=\|V_aP_{\rm r}\|.                                \label{eq:s7v35-va}
\end{equation}
Then
\begin{align}
 \|P_{\rm L}S_a(\tau)P_{\rm r}\|
 &\le {a v_a\over\Delta},                              \label{eq:s7v35-cross}\\
 \|P_{\rm r}S_a(\tau)P_{\rm L}\|
 &\le {a v_a\over\Delta},                              \label{eq:s7v35-cross-adjoint}\\
 \|P_{\rm r}S_a(\tau)P_{\rm r}\|
 &\le
 e^{-\tau\Delta/a}+{a v_a\over\Delta}.                 \label{eq:s7v35-residual-slab}
\end{align}
The bounds are uniform in \(\tau\) from above; no factor proportional to the
number of microscopic time steps occurs.
\end{theorem}

\begin{proof}
The bounded-perturbation Duhamel formula is
\begin{equation}
 e^{-\tau H_a}-e^{-\tau H_{0,a}}
 =
 -\int_0^\tau
 e^{-(\tau-s)H_a}V_ae^{-sH_{0,a}}\,ds.                \label{eq:s7v35-Duhamel}
\end{equation}
Both semigroups are contractions.  Since \(P_{\rm r}\) commutes with
\(H_{0,a}\), \eqref{eq:s7v35-free-energy} gives
\[
 \|e^{-sH_{0,a}}P_{\rm r}\|
 \le e^{-s\Delta/a}.
\]
The free low/residual cross block vanishes.  Multiply
\eqref{eq:s7v35-Duhamel} by \(P_{\rm L}\) on the left and \(P_{\rm r}\) on
the right, use
\(V_ae^{-sH_{0,a}}P_{\rm r}
=V_aP_{\rm r}e^{-sH_{0,a}}P_{\rm r}\), and integrate:
\[
 \|P_{\rm L}S_a(\tau)P_{\rm r}\|
 \le
 v_a\int_0^\tau e^{-s\Delta/a}\,ds
 \le {av_a\over\Delta}.
\]
Self-adjointness of \(S_a(\tau)\) proves
\eqref{eq:s7v35-cross-adjoint}.  Multiplying
\eqref{eq:s7v35-Duhamel} by \(P_{\rm r}\) on both sides gives the free term
\(e^{-\tau\Delta/a}\) plus the same integral, proving
\eqref{eq:s7v35-residual-slab}.
\end{proof}

\begin{corollary}[Fixed-slab transfer compiler]
\label{cor:s7v35-slab-compiler}
Suppose an isometry \(Q\) into the low sector and a coarse fixed-slab
transfer \(S_c(\tau)\) satisfy
\begin{equation}
 \|P_{\rm L}S_a(\tau)P_{\rm L}-QS_c(\tau)Q^*\|
 \le\eta_a,\qquad
 \|S_c(\tau)\|_{{\cal H}_{c,\circ}}\le q_c.            \label{eq:s7v35-slab-match}
\end{equation}
Then
\begin{equation}
 \|S_a(\tau)\|_{{\cal H}_\circ}
 \le
 \max\left\{
 q_c+\eta_a+{av_a\over\Delta},\
 e^{-\tau\Delta/a}+{2av_a\over\Delta}
 \right\}.                                             \label{eq:s7v35-slab-compiler}
\end{equation}
\end{corollary}

\begin{proof}
Apply the two-channel instance of \Cref{thm:s7v35-row} to
\({\cal H}_{\rm L}\oplus{\cal H}_{\rm r}\).
The low diagonal is at most \(q_c+\eta_a\), the cross block is bounded by
\eqref{eq:s7v35-cross}, and the residual diagonal by
\eqref{eq:s7v35-residual-slab}.
\end{proof}

\begin{corollary}[Asymptotic-freedom-sized perturbations fit a fixed margin]
\label{cor:s7v35-AF}
Assume
\begin{equation}
 v_a\le {C g(a)\over a},\qquad g(a)\longrightarrow0,\qquad
 \eta_a\longrightarrow0,                              \label{eq:s7v35-AF-size}
\end{equation}
uniformly in the volume in the declared state-space norm.  If
\begin{equation}
 q_c\le e^{-m\tau}                                     \label{eq:s7v35-coarse-mass}
\end{equation}
for some \(m>0\), then for every \(0<m'<m\), all sufficiently small \(a\)
obey
\begin{equation}
 \|S_a(\tau)\|_{{\cal H}_\circ}\le e^{-m'\tau}.         \label{eq:s7v35-mprime}
\end{equation}
\end{corollary}

\begin{proof}
Under \eqref{eq:s7v35-AF-size}, the two perturbative debits in
\eqref{eq:s7v35-slab-compiler} tend to zero, as does
\(e^{-\tau\Delta/a}\).  The fixed positive number
\[
 e^{-m'\tau}-e^{-m\tau}
\]
is therefore larger than the first-row debit for sufficiently small \(a\);
the second row also eventually lies below \(e^{-m'\tau}\).
\end{proof}

\begin{firewall}[The volume-uniform interaction bound is not yet proved]
\label{fire:s7v35-interaction}
For a many-link Hamiltonian, the global operator norm of \(V_a\) is normally
volume extensive.  Equation \eqref{eq:s7v35-AF-size} is therefore a target,
not an established Yang--Mills estimate.  It must be replaced by a
volume-uniform relative-form, connected-kernel, or local-observable version
of \Cref{thm:s7v35-Duhamel}.  Using the extensive finite-volume norm and then
passing to infinite volume would not prove the required statement.
\end{firewall}

\section{Comparison of the two recursion clocks}
\label{sec:s7v35-clocks}

\begin{theorem}[Why the fixed physical slab is the correct upstream branch]
\label{thm:s7v35-clock-decision}
The established estimates distinguish the two clocks as follows.
\begin{enumerate}[label=\textup{(\roman*)},leftmargin=3.5em]
\item A one-step compiler preserves physical mass through the ultraviolet
      tail if its low-row debit is \(O(a^p)\) with \(p>1\), by
      \Cref{cor:s7v35-summable}.
\item V34 controls separated low/high filter leakage at that order but does
      not control the adjacent low/transition transfer block, by
      \Cref{thm:s7v35-adjacent}.
\item On a fixed physical-time slab, the residual energy denominator produces
      the exact factor \(a/\Delta\), by
      \Cref{thm:s7v35-Duhamel}.
\item Consequently a perturbation of natural asymptotically free size
      \(g(a)/a\) gives a vanishing \(O(g(a))\) slab debit, which can fit the
      fixed mass margin between \(m\) and any \(m'<m\).
\end{enumerate}
Thus the programme should not try to manufacture an \(O(a^p)\)
low/transition estimate from the V34 scalar filter.  It should prove a
volume-uniform connected version of the physical-slab Duhamel bound and match
its low block to the next RG scale.
\end{theorem}

\begin{proof}
Items (i)--(iv) are respectively
\Cref{cor:s7v35-summable,thm:s7v35-adjacent,
thm:s7v35-Duhamel,cor:s7v35-AF}.  The conclusion follows because the
fixed-slab route uses the actual ultraviolet energy denominator and does not
discard the transition mixing term.
\end{proof}

\begin{assessment}[Progress at seventh-series V35]
\label{ass:s7v35-status}
The three-channel compiler is now exact, its physical mass loss is explicit,
and the buffered-filter handoff has been type-corrected.  The transition
channel removes the soft-frame overlap but leaves a genuine adjacent transfer
block.  Moving to fixed physical time converts the free transition energy
\(\Delta/a\) into the rigorous suppression factor \(a/\Delta\).  This is the
first recursion in the present series whose perturbative margin is compatible
with a logarithmically vanishing asymptotically free coupling without
requiring \(g(a)=O(a)\).
\end{assessment}

\begin{decision}[V36 acquisition order]
\label{dec:s7v35-V36}
V36 will localize \Cref{thm:s7v35-Duhamel} to the gauge-invariant
spin-network core used by the direct Osterwalder--Schrader criterion.  It will:
\begin{enumerate}[label=\textup{(\roman*)},leftmargin=3.5em]
\item replace the extensive norm \(v_a\) by a connected local interaction
      form;
\item derive the spatial support enlargement over a fixed Euclidean-time
      slab;
\item identify the precise power of \(g(a)\) and the quasi-local
      \(O(\log(1/a))\) incidence debit; and
\item state the low-block matching estimate needed to hand the slab to the
      next scale and ultimately to the V27 strong-coupling terminal gap.
\end{enumerate}
If a volume-uniform connected bound cannot be proved at fixed \(\tau\), the
next upstream move will be to a polymer/Dyson expansion of the slab kernel,
not back to the one-step soft filter.
\end{decision}

\begin{firewall}[Claim boundary after seventh-series V35]
\label{fire:s7v35-boundary}
V35 proves the exact three-channel row-sum recursion, its physical mass-loss
ledger, ultraviolet summability for \(O(a^p)\), \(p>1\), two separated-band
commutator estimates, the independence of the adjacent transition debit, and
the fixed-physical-time Duhamel suppression theorem.

V35 does not prove the volume-uniform interacting estimate
\eqref{eq:s7v35-AF-size}, the low-block RG matching
\eqref{eq:s7v35-slab-match}, nonlinear gauge-sector control, convergence to
the terminal strong-coupling regime, cutoff removal, the Wightman axioms, or
the full Yang--Mills mass gap.
\end{firewall}

\paragraph{Parent-version record.}
V35 was formed from
\texttt{Renewal\_Yang\_Mills\_Mass\_Gap\_Research\_Notes\_V34.tex}.
The V34 parent had \(13\,592\) logical source lines, \(494\,269\) bytes,
and SHA--256
\[
 \texttt{A6762C08333B443F1C93A44F33CC64E709A59ABF613B2E01F55E0EAFEC12261B}.
\]
The next compulsory companion audit is V40.

% =============================================================================
% END APPENDED SEVENTH-SERIES V35 MATERIAL
% =============================================================================

% =============================================================================
% BEGIN APPENDED SEVENTH-SERIES V36 MATERIAL
% =============================================================================

\chapter{Cutoff-compatible quasi-local incidence on a physical slab}
\label{ch:s7v36}

V35 replaces the vanishing one-step transfer margin by a fixed
physical-time slab and proves that a residual energy floor
\(\Delta/a\) converts an off-diagonal interaction norm \(v_a\) into the
dimensionless debit \(av_a/\Delta\).  The remaining problem is locality:
the global norm of a sum of local interactions is volume extensive.

This note separates two issues.  First, it proves a volume-uniform row-kernel
calculus for the linearized quasi-local layer.  Second, it corrects the V34
filter schedule.  The \(O(a^p)\) accuracy used for a one-step recursion is
unnecessarily expensive on a fixed physical slab.  It is enough that filter
leakage and interaction incidence tend to zero.  Under the one-loop
asymptotically free scale \(g_{\rm YM}(a)=O((\log(1/a))^{-1/2})\), there are
explicit sublogarithmic filter schedules for which both quantities vanish,
even after literal finite-range truncation.

\section{Volume-uniform row-kernel calculus}
\label{sec:s7v36-row}

Let \(\Gamma\) be a graph of polynomial growth dimension \(D\): there is
\(C_{\rm vol}\) such that
\begin{equation}
 \#B(x,R)\le C_{\rm vol}(1+R)^D
 \quad(x\in\Gamma,\ R\ge0).                             \label{eq:s7v36-growth}
\end{equation}
For an operator-valued matrix kernel \(K=(K_{xy})_{x,y\in\Gamma}\), define
\begin{equation}
 {\cal I}(K)
 :=
 \max\left\{
 \sup_x\sum_y\|K_{xy}\|,\
 \sup_y\sum_x\|K_{xy}\|
 \right\}.                                              \label{eq:s7v36-I}
\end{equation}
This is a volume-uniform incidence norm; it does not sum over all lattice
sites.

\begin{lemma}[Schur and product bounds]
\label{lem:s7v36-Schur}
If \({\cal I}(K)<\infty\), then \(K\) is bounded on
\(\ell^2(\Gamma)\) and
\begin{equation}
 \|K\|_{2\to2}\le{\cal I}(K).                           \label{eq:s7v36-Schur}
\end{equation}
If \(K,L\) have finite incidence norm, then
\begin{equation}
 {\cal I}(KL)\le{\cal I}(K){\cal I}(L).                \label{eq:s7v36-product}
\end{equation}
Both bounds are uniform under restriction to finite periodic volumes.
\end{lemma}

\begin{proof}
The first estimate is the Schur test, using the two row and column sums in
\eqref{eq:s7v36-I}.  For the row sum of the product,
\begin{align*}
 \sup_x\sum_z\|(KL)_{xz}\|
 &\le
 \sup_x\sum_y\|K_{xy}\|\sum_z\|L_{yz}\|\\
 &\le{\cal I}(K){\cal I}(L).
\end{align*}
The same calculation with rows and columns exchanged gives the column
bound.  No volume occurs in either estimate.
\end{proof}

\begin{lemma}[Incidence of an exponentially local kernel]
\label{lem:s7v36-exp}
Suppose
\begin{equation}
 \|K_{xy}\|
 \le C_0\exp\{-c_0d(x,y)/\kappa\},
 \qquad \kappa\ge1.                                    \label{eq:s7v36-exp}
\end{equation}
Then
\begin{equation}
 {\cal I}(K)\le C_1\kappa^D,                            \label{eq:s7v36-exp-I}
\end{equation}
where \(C_1\) depends only on
\(C_0,c_0,C_{\rm vol},D\), not on volume or \(\kappa\).
\end{lemma}

\begin{proof}
Equation \eqref{eq:s7v36-growth} implies that the number of vertices in the
shell \(n\le d(x,y)<n+1\) is at most
\(C(1+n)^{D-1}\).  Therefore
\[
 \sum_y\|K_{xy}\|
 \le
 C\sum_{n=0}^\infty(1+n)^{D-1}e^{-c_0n/\kappa}.
\]
Comparison with the integral of
\((1+t)^{D-1}e^{-c_0t/\kappa}\) gives \(C_1\kappa^D\).
The same argument gives the column bound.
\end{proof}

\begin{corollary}[The V34 quasi-local incidence]
\label{cor:s7v36-filter-I}
Let \(F_\kappa=f_\kappa({\cal L})\) be the V34 logistic filter on a
\(D\)-dimensional finite-range lattice.  Its exponentially decaying part has
\begin{equation}
 {\cal I}(F_\kappa)\le C\kappa^D.                       \label{eq:s7v36-filter-I}
\end{equation}
The same estimate holds for \(I-F_\kappa\) after separating its on-site
identity kernel.
\end{corollary}

\begin{proof}
Apply \Cref{lem:s7v36-exp} to the V34 kernel bound.  Subtracting the on-site
identity changes \({\cal I}\) by at most one.
\end{proof}

\section{Local interaction incidence}
\label{sec:s7v36-interaction}

Let \(W_a\) denote the off-diagonal interaction kernel at the linearized
one-particle or Gaussian-comparison layer.  Assume the volume-uniform local
bound
\begin{equation}
 {\cal I}(W_a)\le {C_Wg_{\rm YM}(a)\over a}.            \label{eq:s7v36-W}
\end{equation}
The factor \(a^{-1}\) is the ultraviolet energy scale, while the
dimensionless interaction coefficient is the running coupling.

\begin{theorem}[Quasi-local interaction debit]
\label{thm:s7v36-debit}
If \(K_\kappa\) obeys
\({\cal I}(K_\kappa)\le C_K\kappa^D\), then
\begin{equation}
 a\|W_aK_\kappa\|_{2\to2}
 \le
 C_WC_Kg_{\rm YM}(a)\kappa^D.                          \label{eq:s7v36-debit}
\end{equation}
This bound is uniform in finite periodic volume.
\end{theorem}

\begin{proof}
Use \Cref{lem:s7v36-Schur}:
\[
 a\|W_aK_\kappa\|_{2\to2}
 \le
 a{\cal I}(W_aK_\kappa)
 \le
 a{\cal I}(W_a){\cal I}(K_\kappa).
\]
Insert \eqref{eq:s7v36-W} and the assumed bound on \(K_\kappa\).
\end{proof}

\begin{remark}[What has become volume uniform]
\label{rem:s7v36-volume}
The global sum \(\sum_xW_{a,x}\) may have norm proportional to the volume.
The incidence norm instead measures the amount of interaction entering or
leaving one anchor.  The algebra in \Cref{lem:s7v36-Schur} is exact for
matrix kernels.  Extending the same estimate to nonlinear gauge-invariant
many-body connected kernels remains a separate step.
\end{remark}

\section{A filter schedule compatible with asymptotic freedom}
\label{sec:s7v36-schedule}

To avoid confusing the running coupling with the V34 buffer half-width, put
\begin{equation}
 h:={\beta-\alpha\over2}>0,\qquad
 L(a):=\log{1\over a\Lambda_0},                         \label{eq:s7v36-L}
\end{equation}
where \(a\Lambda_0<e^{-1}\).  Assume
\begin{equation}
 g_{\rm YM}(a)\le C_gL(a)^{-\nu}                       \label{eq:s7v36-running}
\end{equation}
for some \(\nu>0\).  One-loop four-dimensional asymptotic freedom has
\(\nu=1/2\), up to lower-order logarithms.

\begin{theorem}[Quasi-local schedule]
\label{thm:s7v36-quasi-schedule}
Choose
\begin{equation}
 \kappa(a)=L(a)^s,\qquad 0<s<{\nu\over D}.              \label{eq:s7v36-kappa}
\end{equation}
Then
\begin{align}
 e^{-h\kappa(a)}&\longrightarrow0,                     \label{eq:s7v36-leak-zero}\\
 g_{\rm YM}(a)\kappa(a)^D
 &\le C_gL(a)^{-\nu+sD}\longrightarrow0.               \label{eq:s7v36-inc-zero}
\end{align}
Thus the V34 separated-band leakage and the V36 quasi-local interaction
incidence vanish simultaneously.
\end{theorem}

\begin{proof}
Since \(s>0\), \(\kappa(a)\to\infty\), proving
\eqref{eq:s7v36-leak-zero}.  Substitution of
\eqref{eq:s7v36-kappa} into \eqref{eq:s7v36-running} gives
\eqref{eq:s7v36-inc-zero}; its exponent is negative by the declared range of
\(s\).
\end{proof}

\begin{corollary}[Explicit four-dimensional Yang--Mills choice]
\label{cor:s7v36-explicit}
For a spatial lattice of growth dimension \(D=3\) and the one-loop bound
\(\nu=1/2\), the choice
\begin{equation}
 \kappa(a)=L(a)^{1/16}                                 \label{eq:s7v36-explicit-kappa}
\end{equation}
gives
\begin{equation}
 e^{-h\kappa(a)}=e^{-hL(a)^{1/16}},\qquad
 g_{\rm YM}(a)\kappa(a)^3=O(L(a)^{-5/16}).             \label{eq:s7v36-explicit-rates}
\end{equation}
\end{corollary}

\begin{proof}
Here \(s=1/16<1/6=\nu/D\), and
\(-1/2+3/16=-5/16\).
\end{proof}

\subsection{Literal finite range}
\label{sec:s7v36-finite}

V34 gives a degree-\(N\) polynomial approximation with error
\(Ce^{-cN/\kappa}\).  Matching that truncation error to the logistic leakage
\(e^{-h\kappa}\) requires only \(N=O(\kappa^2)\).

\begin{theorem}[Finite-range schedule]
\label{thm:s7v36-finite-schedule}
Choose
\begin{equation}
 N(a)=\left\lceil C_N\kappa(a)^2\right\rceil,\qquad
 C_N>{h\over c},                                      \label{eq:s7v36-N}
\end{equation}
where \(c\) is the V34 polynomial-approximation constant.  Then
\begin{equation}
 \|F_{\kappa(a)}-p_{\kappa(a),N(a)}({\cal L})\|
 \le Ce^{-h\kappa(a)},                                 \label{eq:s7v36-finite-error}
\end{equation}
after increasing \(C_N\) to absorb the fixed prefactor.  If a strictly
range-\(N\) kernel is charged by the ball-incidence bound \(O(N^D)\), its
interaction debit vanishes whenever
\begin{equation}
 0<s<{\nu\over2D}.                                     \label{eq:s7v36-finite-window}
\end{equation}
\end{theorem}

\begin{proof}
The V34 bound gives
\[
 Ce^{-cN(a)/\kappa(a)}
 \le
 Ce^{-cC_N\kappa(a)}.
\]
Choose \(C_N\) so that the exponent and the fixed prefactor are bounded by
the right side of \eqref{eq:s7v36-finite-error} for small \(a\).  The
range-\(N\) incidence is at most \(C(1+N)^D=O(\kappa^{2D})\).  Hence its
dimensionless interaction debit is bounded by
\[
 Cg_{\rm YM}(a)\kappa(a)^{2D}
 \le C L(a)^{-\nu+2sD},
\]
which tends to zero exactly in the range
\eqref{eq:s7v36-finite-window}.
\end{proof}

\begin{corollary}[The explicit choice also survives finite-range truncation]
\label{cor:s7v36-explicit-finite}
For \(D=3\), \(\nu=1/2\), and \(s=1/16\),
\begin{equation}
 N(a)=O(L(a)^{1/8}),\qquad
 g_{\rm YM}(a)N(a)^3=O(L(a)^{-1/8}).                  \label{eq:s7v36-explicit-finite}
\end{equation}
\end{corollary}

\begin{proof}
One has \(2s=1/8\) and
\(-1/2+3/8=-1/8\).
\end{proof}

\begin{firewall}[The old \(O(a^p)\) schedule is too costly for this norm]
\label{fire:s7v36-old-schedule}
The V34 one-step choice \(\kappa(a)\asymp L(a)\) has incidence
\(\kappa(a)^D\).  Under
\(g_{\rm YM}(a)=O(L(a)^{-1/2})\), the product
\(g_{\rm YM}(a)\kappa(a)^D\) does not tend to zero for \(D\ge1\).
This does not contradict V34, which proved operator-norm leakage and
localization.  It shows that the filter accuracy must be matched to the
clock and the incidence norm in which the interacting estimate is made.
\end{firewall}

\section{Fixed-slab error card}
\label{sec:s7v36-card}

\begin{definition}[Quasi-local physical-slab card, QLPSC]
\label{def:s7v36-card}
At cutoff \(a\), QLPSC records:
\begin{enumerate}[label=\textup{(Q\arabic*)},leftmargin=4em]
\item a fixed physical time \(\tau>0\) and a residual free-energy floor
      \(\Delta/a\);
\item a filter parameter \(\kappa(a)\to\infty\);
\item separated-band leakage
      \(\ell_a\le C e^{-h\kappa(a)}\);
\item a volume-uniform local interaction incidence
      \(j_a\le Cg_{\rm YM}(a)\kappa(a)^D\), or
      \(Cg_{\rm YM}(a)N(a)^D\) after finite-range truncation;
\item a low-block coarse matching error \(\eta_a\);
\item a connected nonlinear remainder \(n_a\); and
\item the total slab debit
\begin{equation}
 d_a:=\eta_a+\ell_a+j_a+n_a+e^{-\tau\Delta/a}.         \label{eq:s7v36-da}
\end{equation}
\end{enumerate}
\end{definition}

\begin{theorem}[Conditional fixed-slab gap propagation]
\label{thm:s7v36-card}
Assume the exact gauge-invariant slab transfer admits the V35 low/residual
decomposition and that its row bound is at most
\begin{equation}
 \max\{q_c+C d_a,\ C d_a\}.                            \label{eq:s7v36-row-card}
\end{equation}
If
\begin{equation}
 q_c\le e^{-m\tau},\qquad d_a\longrightarrow0,         \label{eq:s7v36-card-hyp}
\end{equation}
then for every \(0<m'<m\),
\begin{equation}
 \|S_a(\tau)\|_{\Omega^\perp}\le e^{-m'\tau}           \label{eq:s7v36-card-gap}
\end{equation}
for all sufficiently small \(a\).
\end{theorem}

\begin{proof}
The positive margins
\[
 e^{-m'\tau}-e^{-m\tau}>0,\qquad e^{-m'\tau}>0
\]
are fixed independently of \(a\).  Since \(d_a\to0\), both rows in
\eqref{eq:s7v36-row-card} eventually lie below \(e^{-m'\tau}\).
\end{proof}

\begin{corollary}[Filter and linearized incidence rows close]
\label{cor:s7v36-linear-close}
Under \eqref{eq:s7v36-running}, either schedule
\eqref{eq:s7v36-kappa} or
\eqref{eq:s7v36-finite-window} makes
\begin{equation}
 \ell_a+j_a+e^{-\tau\Delta/a}\longrightarrow0.         \label{eq:s7v36-three-close}
\end{equation}
Hence the only unpaid QLPSC rows are the interacting low-block match
\(\eta_a\) and the connected nonlinear remainder \(n_a\).
\end{corollary}

\begin{proof}
Use
\Cref{thm:s7v36-quasi-schedule,thm:s7v36-finite-schedule}; the last
exponential tends to zero faster than every power of \(a\).
\end{proof}

\section{Boundary between the linearized and nonlinear problems}
\label{sec:s7v36-boundary}

\begin{theorem}[What the incidence calculation establishes]
\label{thm:s7v36-result}
At the linearized matrix-kernel layer, the V34 filter can be made
simultaneously accurate, quasi-local, volume uniform, and compatible with
the asymptotically free coupling.  In spatial dimension three, the explicit
choice \(\kappa=L^{1/16}\) gives leakage
\(e^{-hL^{1/16}}\), quasi-local interaction debit
\(O(L^{-5/16})\), and, after truncation at range \(O(L^{1/8})\), a worst-case
finite-range incidence debit \(O(L^{-1/8})\).
\end{theorem}

\begin{proof}
Combine
\Cref{lem:s7v36-Schur,lem:s7v36-exp,thm:s7v36-debit,
cor:s7v36-explicit,cor:s7v36-explicit-finite}.
\end{proof}

\begin{firewall}[No nonlinear connected-kernel theorem yet]
\label{fire:s7v36-nonlinear}
The Wilson interaction acts on a compact-group many-link Hilbert space, not
merely on a one-particle matrix kernel.  V36 has not proved that its
gauge-invariant connected slab kernels obey
\eqref{eq:s7v36-W}, nor that vacuum normalization deletes every disconnected
volume factor with the same constants.  Those statements are exactly the
unpaid nonlinear remainder \(n_a\).  Substituting the linearized incidence
bound for \(n_a\) would be circular.
\end{firewall}

\begin{assessment}[Progress at seventh-series V36]
\label{ass:s7v36-status}
The logarithmic localization cost is no longer a conflict with asymptotic
freedom once the transfer clock is fixed in physical units.  The filter need
not achieve \(a^p\) accuracy: a sublogarithmic
\(\kappa(a)\to\infty\) makes leakage vanish while keeping the local incidence
smaller than the running-coupling gain.  The exact remaining bridge is now
twofold: a connected nonlinear slab-kernel estimate and a low-block RG
matching theorem.
\end{assessment}

\begin{decision}[V37 acquisition order]
\label{dec:s7v36-V37}
V37 will construct the connected-kernel algebra for normalized
gauge-invariant slab correlations.  It will first prove the cancellation of
vacuum-disconnected interaction components in the ratio defining a connected
two-point function.  It will then derive an anchored polymer norm and test
whether the Wilson plaquette Dyson coefficients are summable on a fixed
physical slab with the schedule \eqref{eq:s7v36-explicit-kappa}.  If the
Dyson activity is not small, the programme will move upstream to a
block-Gaussian interpolation rather than hide the volume factor.
\end{decision}

\begin{firewall}[Claim boundary after seventh-series V36]
\label{fire:s7v36-claim}
V36 proves the volume-uniform row-kernel calculus, the quasi-local and
finite-range scheduling windows, the explicit \(D=3\) one-loop rates, and
the conditional fixed-slab error card.

V36 does not prove the nonlinear connected Wilson-kernel bound, the low-block
RG match, a nonperturbative terminal handoff, cutoff removal, or the full
Yang--Mills mass gap.
\end{firewall}

\paragraph{Parent-version record.}
V36 was formed from
\texttt{Renewal\_Yang\_Mills\_Mass\_Gap\_Research\_Notes\_V35.tex}.
The V35 parent had \(14\,265\) logical source lines, \(519\,659\) bytes,
and SHA--256
\[
 \texttt{B5B8391D62A6D4ADCCE5EAC8BD6015770916E896A86DFA2B0781E9D4AD3677DC}.
\]
The next compulsory companion audit is V40.

% =============================================================================
% END APPENDED SEVENTH-SERIES V36 MATERIAL
% =============================================================================

% =============================================================================
% BEGIN APPENDED SEVENTH-SERIES V37 MATERIAL
% =============================================================================

\chapter{High-energy Feshbach reduction and the self-energy ledger}
\label{ch:s7v37}

V36 proposed an anchored connected Dyson expansion for the fixed physical
slab.  Before building that expansion, one must distinguish two kinds of
high-sector contribution.  An open excursion ending in the residual sector
has the \(O(g_{\rm YM})\) mixing size found in V35.  A closed excursion that
returns to the low sector carries two interaction vertices and one
high-energy denominator.  Its natural size is
\[
 {g_{\rm YM}(a)^2\over a}.
\]
This is not a vanishing transfer error.  It is a low-energy self-energy of
the same physical dimension as the Hamiltonian and must be absorbed into the
renormalized coarse block.

This note makes that distinction exact by a Feshbach--Schur reduction.  It
proves the operator factorization, the mixing and self-energy scales, a
volume-uniform incidence version, preservation of gauge symmetry, and a
spectral-gap handoff from the energy-dependent effective low Hamiltonian to
the full Hamiltonian.

\section{Why the naive Dyson remainder does not close}
\label{sec:s7v37-Dyson}

Let
\begin{equation}
 {\cal H}_\circ={\cal H}_{\rm L}\oplus{\cal H}_{\rm r},
 \qquad
 H_d=H_{\rm L}\oplus H_{\rm r},                        \label{eq:s7v37-Hd}
\end{equation}
where \(H_{\rm L}\ge0\) and
\begin{equation}
 H_{\rm r}\ge{\Delta\over a}I_{\rm r}.                 \label{eq:s7v37-Hr-gap}
\end{equation}
Let the interaction be purely off diagonal,
\begin{equation}
 W=
 \begin{pmatrix}0&K\\K^*&0\end{pmatrix}.               \label{eq:s7v37-W}
\end{equation}
The diagonal parts of an interaction can always be included in
\(H_{\rm L}\) and \(H_{\rm r}\) before this split.

\begin{lemma}[One high sojourn]
\label{lem:s7v37-sojourn}
For every \(\tau>0\),
\begin{equation}
 \int_0^\tau ds_2\int_0^{s_2}ds_1\,
 e^{-(s_2-s_1)\Delta/a}
 \le{\tau a\over\Delta}.                               \label{eq:s7v37-sojourn}
\end{equation}
Consequently the norm estimate for the second-order low-to-low Dyson return
is
\begin{equation}
 \|{\cal D}^{(2)}_{\rm LL}(\tau)\|
 \le{\tau a\over\Delta}\|K\|^2.                        \label{eq:s7v37-D2}
\end{equation}
\end{lemma}

\begin{proof}
For fixed \(s_2\), set \(u=s_2-s_1\).  The inner integral is at most
\[
 \int_0^\infty e^{-u\Delta/a}\,du={a\over\Delta}.
\]
Integrating \(s_2\) over \([0,\tau]\) proves
\eqref{eq:s7v37-sojourn}.  In the Dyson term, the low semigroups are
contractions, the residual semigroup is bounded by the displayed
exponential, and each off-diagonal vertex contributes \(\|K\|\).
\end{proof}

\begin{corollary}[The return has Hamiltonian size]
\label{cor:s7v37-return-size}
If
\begin{equation}
 \|K\|\le {C_Kg_{\rm YM}(a)\over a},                   \label{eq:s7v37-Kscale}
\end{equation}
then
\begin{equation}
 \|{\cal D}^{(2)}_{\rm LL}(\tau)\|
 \le
 {C_K^2\tau\over\Delta}
 {g_{\rm YM}(a)^2\over a}.                             \label{eq:s7v37-return-size}
\end{equation}
Thus the high return cannot be assigned to the vanishing slab debit merely
because \(g_{\rm YM}(a)\to0\).
\end{corollary}

\begin{proof}
Insert \eqref{eq:s7v37-Kscale} into \eqref{eq:s7v37-D2}.
\end{proof}

\section{Exact Feshbach factorization}
\label{sec:s7v37-Feshbach}

Let
\begin{equation}
 H=
 \begin{pmatrix}
 H_{\rm L}&K\\
 K^*&H_{\rm r}
 \end{pmatrix}                                         \label{eq:s7v37-H}
\end{equation}
be self-adjoint, with \(K\) bounded.  Fix a real \(z<\Delta/a\), and define
\begin{align}
 R_{\rm r}(z)&:=(H_{\rm r}-z)^{-1},                    \label{eq:s7v37-Rr}\\
 \Sigma_a(z)&:=KR_{\rm r}(z)K^*,                       \label{eq:s7v37-Sigma}\\
 F_a(z)&:=H_{\rm L}-z-\Sigma_a(z).                     \label{eq:s7v37-F}
\end{align}

\begin{theorem}[Feshbach--Schur factorization]
\label{thm:s7v37-factor}
On the natural block domain,
\begin{equation}
 H-z
 =
 \begin{pmatrix}I&KR_{\rm r}(z)\\0&I\end{pmatrix}
 \begin{pmatrix}F_a(z)&0\\0&H_{\rm r}-z\end{pmatrix}
 \begin{pmatrix}I&0\\R_{\rm r}(z)K^*&I\end{pmatrix}.
 \label{eq:s7v37-factor}
\end{equation}
Hence \(H-z\) is invertible if and only if \(F_a(z)\) is invertible.  If
\(\psi=(\psi_{\rm L},\psi_{\rm r})\) obeys
\((H-z)\psi=0\), then
\begin{align}
 F_a(z)\psi_{\rm L}&=0,                                \label{eq:s7v37-eigen-low}\\
 \psi_{\rm r}&=-R_{\rm r}(z)K^*\psi_{\rm L}.           \label{eq:s7v37-eigen-high}
\end{align}
\end{theorem}

\begin{proof}
Multiplying the last two factors in \eqref{eq:s7v37-factor} gives
\[
 \begin{pmatrix}
 F_a(z)&0\\
 K^*&H_{\rm r}-z
 \end{pmatrix}.
\]
Multiplication by the first factor gives the off-diagonal entries \(K,K^*\);
the upper-left entry is
\[
 F_a(z)+KR_{\rm r}(z)K^*=H_{\rm L}-z.
\]
This proves the factorization.  The triangular factors are boundedly
invertible, and \(H_{\rm r}-z\) is invertible by
\eqref{eq:s7v37-Hr-gap}; the invertibility equivalence follows.  The two
kernel equations are the block equations with the residual variable
eliminated.
\end{proof}

\begin{corollary}[Exact resolvent blocks]
\label{cor:s7v37-resolvent}
Whenever \(F_a(z)\) is invertible,
\begin{equation}
 (H-z)^{-1}
 =
 \begin{pmatrix}
 F_a(z)^{-1}
 &
 -F_a(z)^{-1}KR_{\rm r}(z)
 \\
 -R_{\rm r}(z)K^*F_a(z)^{-1}
 &
 R_{\rm r}(z)+
 R_{\rm r}(z)K^*F_a(z)^{-1}KR_{\rm r}(z)
 \end{pmatrix}.
 \label{eq:s7v37-resolvent}
\end{equation}
\end{corollary}

\begin{proof}
Invert the three factors in \eqref{eq:s7v37-factor} in reverse order and
multiply them.
\end{proof}

\section{Mixing and self-energy scales}
\label{sec:s7v37-scales}

Put
\begin{equation}
 \delta_z:=\Delta-az>0.                                \label{eq:s7v37-deltaz}
\end{equation}

\begin{theorem}[Sharp dimensional ledger]
\label{thm:s7v37-scales}
For \(z<\Delta/a\),
\begin{align}
 \|R_{\rm r}(z)\|
 &\le {a\over\delta_z},                                \label{eq:s7v37-Rbound}\\
 \|R_{\rm r}(z)K^*\|
 &\le {a\|K\|\over\delta_z},                           \label{eq:s7v37-mixing}\\
 \|\Sigma_a(z)\|
 &\le {a\|K\|^2\over\delta_z},                         \label{eq:s7v37-Sigma-bound}\\
 \|\Sigma_a'(z)\|
 &\le {a^2\|K\|^2\over\delta_z^2}.                    \label{eq:s7v37-Sigma-prime}
\end{align}
Under \eqref{eq:s7v37-Kscale}, these become
\begin{align}
 \|R_{\rm r}(z)K^*\|
 &\le {C_Kg_{\rm YM}(a)\over\delta_z},                 \label{eq:s7v37-mixing-g}\\
 \|\Sigma_a(z)\|
 &\le {C_K^2g_{\rm YM}(a)^2\over a\delta_z},           \label{eq:s7v37-Sigma-g}\\
 \|\Sigma_a'(z)\|
 &\le {C_K^2g_{\rm YM}(a)^2\over\delta_z^2}.           \label{eq:s7v37-Sigma-prime-g}
\end{align}
\end{theorem}

\begin{proof}
The residual spectral floor gives
\eqref{eq:s7v37-Rbound}.  Submultiplicativity gives the next two estimates.
The resolvent derivative is
\[
 R_{\rm r}'(z)=R_{\rm r}(z)^2,
\]
so
\(\Sigma_a'(z)=KR_{\rm r}(z)^2K^*\), proving
\eqref{eq:s7v37-Sigma-prime}.  Substitute
\eqref{eq:s7v37-Kscale}.
\end{proof}

\begin{example}[The self-energy scale is attained by a positive matrix]
\label{ex:s7v37-sharp}
On \(\mathbb C\oplus\mathbb C\), set
\begin{equation}
 H_{\rm r}={\Delta\over a},\qquad
 K={g\over a},\qquad
 H_{\rm L}={g^2\over a\Delta}.                         \label{eq:s7v37-model}
\end{equation}
Then \(H\ge0\), \(\det H=0\), and
\begin{equation}
 \Sigma_a(0)=KH_{\rm r}^{-1}K^*
 ={g^2\over a\Delta}=H_{\rm L}.                        \label{eq:s7v37-model-Sigma}
\end{equation}
Thus the \(g^2/a\) term can be the exact counterterm required for a
nonnegative low mode; it is not an artifact of an upper bound.
\end{example}

\begin{proof}
Both diagonal entries in \eqref{eq:s7v37-model} are nonnegative and the
determinant is zero, so the \(2\times2\) Hermitian matrix is positive
semidefinite.  Direct substitution gives
\eqref{eq:s7v37-model-Sigma}.
\end{proof}

\begin{firewall}[Closed high excursions belong to the coarse action]
\label{fire:s7v37-self-energy}
The open high component of a low-energy state is \(O(g_{\rm YM})\) by
\eqref{eq:s7v37-mixing-g}.  The closed high return is
\(O(g_{\rm YM}^2/a)\) by \eqref{eq:s7v37-Sigma-g}.  Only the first is a
vanishing mixing debit.  The second must be included in the low effective
Hamiltonian
\begin{equation}
 H_{{\rm eff},a}:=F_a(0)=H_{\rm L}-KH_{\rm r}^{-1}K^*.
 \label{eq:s7v37-Heff}
\end{equation}
\end{firewall}

\section{Volume-uniform incidence form of the reduction}
\label{sec:s7v37-incidence}

Use the V36 incidence norm \({\cal I}\).  Assume that the block kernels obey
\begin{align}
 {\cal I}(K)&\le {C_Kg_{\rm YM}(a)\over a},             \label{eq:s7v37-I-K}\\
 {\cal I}(R_{\rm r}(z))&\le {C_Ra\over\delta_z}.        \label{eq:s7v37-I-R}
\end{align}

\begin{theorem}[Anchored Feshbach ledger]
\label{thm:s7v37-incidence}
Under \eqref{eq:s7v37-I-K}--\eqref{eq:s7v37-I-R},
\begin{align}
 {\cal I}(R_{\rm r}(z)K^*)
 &\le {C_RC_Kg_{\rm YM}(a)\over\delta_z},              \label{eq:s7v37-I-mixing}\\
 {\cal I}(\Sigma_a(z))
 &\le
 {C_RC_K^2g_{\rm YM}(a)^2\over a\delta_z}.             \label{eq:s7v37-I-Sigma}
\end{align}
The constants are independent of finite periodic volume.
\end{theorem}

\begin{proof}
The V36 product inequality gives
\[
 {\cal I}(R_{\rm r}K^*)
 \le{\cal I}(R_{\rm r}){\cal I}(K)
\]
and
\[
 {\cal I}(KR_{\rm r}K^*)
 \le{\cal I}(K){\cal I}(R_{\rm r}){\cal I}(K^*).
\]
The incidence norm is invariant under adjoint.  Substitute
\eqref{eq:s7v37-I-K}--\eqref{eq:s7v37-I-R}.
\end{proof}

\begin{lemma}[A sufficient resolvent-incidence condition]
\label{lem:s7v37-Neumann}
Suppose
\begin{equation}
 H_{\rm r}-z={\delta_z\over a}(I+B_z),\qquad
 {\cal I}(B_z)\le\vartheta<1.                          \label{eq:s7v37-Neumann-hyp}
\end{equation}
Then
\begin{equation}
 {\cal I}(R_{\rm r}(z))
 \le {a\over\delta_z(1-\vartheta)}.                    \label{eq:s7v37-Neumann}
\end{equation}
If \(B_z\) has range at most \(R_0\), the inverse kernel also obeys an
exponential tail with decay rate at least
\(|\log\vartheta|/R_0\).
\end{lemma}

\begin{proof}
The incidence algebra of V36 is submultiplicative, so the Neumann series
\[
 R_{\rm r}(z)
 ={a\over\delta_z}\sum_{n=0}^\infty(-B_z)^n
\]
converges in incidence norm and gives
\eqref{eq:s7v37-Neumann}.  If
\(d(x,y)>nR_0\), then \((B_z^n)_{xy}=0\).  Hence only
\(n\ge\lceil d(x,y)/R_0\rceil\) contribute to an off-diagonal matrix
element, and their geometric sum is bounded by a constant times
\(\vartheta^{d(x,y)/R_0}\).
\end{proof}

\begin{firewall}[Spectral gap does not automatically give incidence decay]
\label{fire:s7v37-resolvent-incidence}
The operator bound \(\|R_{\rm r}(z)\|\le a/\delta_z\) alone does not imply
\eqref{eq:s7v37-I-R}.  A Combes--Thomas estimate, a convergent Neumann
representation such as \Cref{lem:s7v37-Neumann}, or a direct connected
kernel bound is still required for the Wilson residual Hamiltonian.
\end{firewall}

\section{Gauge covariance of the effective Hamiltonian}
\label{sec:s7v37-gauge}

\begin{theorem}[Feshbach reduction preserves exact gauge symmetry]
\label{thm:s7v37-gauge}
Let \(U(\gamma)\) be a unitary representation of the lattice gauge group.
Assume
\begin{equation}
 [H,U(\gamma)]=0,\qquad
 [P_{\rm L},U(\gamma)]=[P_{\rm r},U(\gamma)]=0         \label{eq:s7v37-gauge-hyp}
\end{equation}
for every \(\gamma\).  Then
\begin{equation}
 [F_a(z),U_{\rm L}(\gamma)]=0,\qquad
 [H_{{\rm eff},a},U_{\rm L}(\gamma)]=0,                \label{eq:s7v37-gauge-eff}
\end{equation}
where \(U_{\rm L}\) is the restriction to \({\cal H}_{\rm L}\).
\end{theorem}

\begin{proof}
The hypotheses imply that every block of \(H\) intertwines the corresponding
restricted representations.  Functional calculus gives
\([R_{\rm r}(z),U_{\rm r}(\gamma)]=0\).  Therefore
\(KR_{\rm r}(z)K^*\) commutes with \(U_{\rm L}(\gamma)\), and so do
\(F_a(z)\) and its value at zero.
\end{proof}

\begin{corollary}[No gauge-breaking counterterm is generated]
\label{cor:s7v37-gauge-counterterm}
If the low/residual projections are gauge invariant, the self-energy
\(\Sigma_a(z)\) belongs to the gauge-invariant low operator algebra.  Its
local expansion may require counterterms, but none may violate the exact
lattice gauge symmetry.
\end{corollary}

\begin{proof}
This is the commutation statement in
\Cref{thm:s7v37-gauge}.
\end{proof}

\section{Spectral-gap handoff}
\label{sec:s7v37-gap}

Work on the vacuum complement, so that a positive lower bound means a mass
gap.  Suppose
\begin{equation}
 H_{{\rm eff},a}\ge m_{\rm eff}I_{\rm L}               \label{eq:s7v37-Heff-gap}
\end{equation}
for some \(m_{\rm eff}>0\).

\begin{theorem}[Effective-to-full gap theorem]
\label{thm:s7v37-gap}
Assume \eqref{eq:s7v37-Hr-gap} and
\eqref{eq:s7v37-Heff-gap}.  Let \(m_0<\Delta/a\) and suppose
\begin{equation}
 \sup_{0\le z\le m_0}\|\Sigma_a'(z)\|\le\theta.         \label{eq:s7v37-theta}
\end{equation}
Then the full Hamiltonian has no spectrum in
\begin{equation}
 0<z<
 \min\left\{
 m_0,\ {m_{\rm eff}\over1+\theta}
 \right\}.                                             \label{eq:s7v37-full-gap}
\end{equation}
In particular, its gap on the vacuum complement is at least the right-hand
minimum.
\end{theorem}

\begin{proof}
For \(0\le z\le m_0\), the fundamental theorem of calculus and
\eqref{eq:s7v37-theta} give
\[
 \|\Sigma_a(z)-\Sigma_a(0)\|\le\theta z.
\]
Therefore
\begin{align*}
 F_a(z)
 &=H_{{\rm eff},a}-z-\bigl(\Sigma_a(z)-\Sigma_a(0)\bigr)\\
 &\ge
 \bigl(m_{\rm eff}-(1+\theta)z\bigr)I_{\rm L}.
\end{align*}
For \(z\) in \eqref{eq:s7v37-full-gap}, both \(F_a(z)\) and
\(H_{\rm r}-z\) are strictly positive and hence invertible.  The exact
factorization \eqref{eq:s7v37-factor} makes \(H-z\) invertible.  Thus that
interval is disjoint from the spectrum.
\end{proof}

\begin{corollary}[Asymptotically full recovery of the effective gap]
\label{cor:s7v37-gap-AF}
Assume \eqref{eq:s7v37-Kscale}, fix \(m_0<\Delta/(2a)\), and suppose
\(m_0=O(1)\) as \(a\downarrow0\).  Then
\begin{equation}
 \theta
 \le {4C_K^2\over\Delta^2}g_{\rm YM}(a)^2             \label{eq:s7v37-theta-AF}
\end{equation}
for sufficiently small \(a\).  If
\(m_{\rm eff}\ge m_*>0\) uniformly, the full gap is at least
\begin{equation}
 {m_*\over1+O(g_{\rm YM}(a)^2)}                        \label{eq:s7v37-gap-AF}
\end{equation}
up to the irrelevant residual threshold \(\Delta/a\).
\end{corollary}

\begin{proof}
On \(0\le z\le m_0<\Delta/(2a)\), one has
\(\delta_z\ge\Delta/2\).  Apply
\eqref{eq:s7v37-Sigma-prime-g} and then
\Cref{thm:s7v37-gap}.
\end{proof}

\begin{assessment}[Meaning of the gap theorem]
\label{ass:s7v37-gap}
High-sector elimination does not have to preserve the bare low Hamiltonian.
It has to preserve the gap of the renormalized effective operator
\(H_{{\rm eff},a}\).  Once that operator has a uniform gap and the
off-diagonal mixing is asymptotically free, the energy dependence of the
Feshbach map loses only a relative \(O(g_{\rm YM}^2)\) factor.
\end{assessment}

\section{Revised connected-kernel target}
\label{sec:s7v37-target}

\begin{definition}[Gauge-invariant Feshbach matching card, GIFMC]
\label{def:s7v37-card}
At cutoff \(a\), GIFMC records:
\begin{enumerate}[label=\textup{(F\arabic*)},leftmargin=4em]
\item gauge-invariant projections \(P_{\rm L},P_{\rm r}\);
\item a residual lower bound \(H_{\rm r}\ge\Delta/a\);
\item an anchored off-diagonal bound
      \({\cal I}(K)\le Cg_{\rm YM}(a)/a\);
\item an anchored resolvent bound
      \({\cal I}(R_{\rm r}(z))\le C a/\delta_z\);
\item the complete self-energy
      \(\Sigma_a(z)=KR_{\rm r}(z)K^*\);
\item a local gauge-invariant expansion and counterterm ledger for
      \(H_{{\rm eff},a}=H_{\rm L}-\Sigma_a(0)\);
\item a coarse matching bound for \(H_{{\rm eff},a}\), not for the bare
      \(H_{\rm L}\); and
\item a uniform effective gap plus the derivative bound
      \eqref{eq:s7v37-theta}.
\end{enumerate}
\end{definition}

\begin{theorem}[What V37 closes]
\label{thm:s7v37-result}
Under the declared block and incidence hypotheses, high-energy elimination
has a volume-uniform, gauge-invariant, type-correct ledger:
\begin{equation}
 \hbox{open mixing}=O(g_{\rm YM}),\qquad
 \hbox{closed self-energy}=O(g_{\rm YM}^2/a).           \label{eq:s7v37-ledger}
\end{equation}
The second term is incorporated into the effective low Hamiltonian.  A
uniform gap of that effective Hamiltonian passes to the full Hamiltonian with
only relative \(O(g_{\rm YM}^2)\) loss.
\end{theorem}

\begin{proof}
Use
\Cref{thm:s7v37-scales,thm:s7v37-incidence,
thm:s7v37-gauge,thm:s7v37-gap,cor:s7v37-gap-AF}.
\end{proof}

\begin{assessment}[Progress at seventh-series V37]
\label{ass:s7v37-status}
The volume problem and the renormalization problem are now separated.
V36 gives a compatible quasi-local incidence schedule for open mixing.
V37 proves that closed high excursions are the gauge-invariant self-energy
that must define the matched low block.  The central continuum bridge has
therefore been reduced to a concrete statement: construct the Wilson
Hamiltonian GIFMC uniformly in volume and prove a gap for its renormalized
effective low operator through the RG flow to the V27 strong-coupling
terminal scale.
\end{assessment}

\begin{decision}[V38 acquisition order]
\label{dec:s7v37-V38}
V38 will test GIFMC on the compact-group lattice Hamiltonian.  It will derive
the exact electric-energy change produced by multiplication by a matrix
coefficient or a plaquette character, identify the residual energy
denominators in Peter--Weyl variables, and compute the support and gauge type
of the second-order self-energy.  The purpose is to determine whether the
anchored bounds (F2)--(F4) can be proved before any perturbative continuum
claim is made.
\end{decision}

\begin{firewall}[Claim boundary after seventh-series V37]
\label{fire:s7v37-claim}
V37 proves the exact high-energy Feshbach factorization, the sharp
mixing/self-energy dimensional ledger, its incidence and gauge-covariance
forms, and an effective-to-full spectral-gap theorem.

V37 does not prove the GIFMC hypotheses for the Wilson Hamiltonian, the
counterterm expansion or coarse match of \(H_{{\rm eff},a}\), a uniform
effective gap, the terminal strong-coupling handoff, cutoff removal, or the
full Yang--Mills mass gap.
\end{firewall}

\paragraph{Parent-version record.}
V37 was formed from
\texttt{Renewal\_Yang\_Mills\_Mass\_Gap\_Research\_Notes\_V36.tex}.
The V36 parent had \(14\,703\) logical source lines, \(535\,505\) bytes,
and SHA--256
\[
 \texttt{6056C624519C0A498EC280E022D3F38AC82293FE9C581BCA566DF9F32F4465BF}.
\]
The next compulsory companion audit is V40.

% =============================================================================
% END APPENDED SEVENTH-SERIES V37 MATERIAL
% =============================================================================

% =============================================================================
% BEGIN APPENDED SEVENTH-SERIES V38 MATERIAL
% =============================================================================

\chapter{The compact-link split and the Gaussian oscillator correction}
\label{ch:s7v38}

V37 formulates the gauge-invariant Feshbach matching card with a residual
energy floor \(\Delta/a\) and off-diagonal size \(O(g_{\rm YM}/a)\).  This
note tests those scalings against the compact-group Hamiltonian before using
them.  The test fails for a split based on the electric Peter--Weyl operator
alone: a fundamental plaquette loop has electric energy
\(O(g_{\rm YM}^2/a)\), while an individual Wilson magnetic multiplier has
norm \(O(1/(ag_{\rm YM}^2))\).

The failure is not a failure of the weak-coupling programme.  It identifies
the correct unperturbed operator.  Electric and quadratic magnetic energies
must be combined and the link coordinate rescaled by \(g_{\rm YM}\).  The
result is a harmonic oscillator with coupling-independent frequencies
\(O(1/a)\); cubic nonabelian terms then have the desired
\(O(g_{\rm YM}/a)\) size.  The calculation below is exact at the quadratic
level and keeps the compact-chart and gauge-zero-mode debits explicit.

\section{Electric Peter--Weyl spectrum}
\label{sec:s7v38-PW}

Let \(G\) be a compact connected matrix Lie group with normalized Haar
measure.  Choose an invariant inner product on its Lie algebra and let
\(-\Delta_G\ge0\) be the corresponding Laplace--Beltrami operator.  For an
irreducible unitary representation \(\lambda\), write \(C_2(\lambda)\) for
its quadratic Casimir eigenvalue.

\begin{lemma}[Peter--Weyl electric eigenvalues]
\label{lem:s7v38-PW}
Every matrix coefficient \(D^\lambda_{mn}\) obeys
\begin{equation}
 -\Delta_GD^\lambda_{mn}
 =C_2(\lambda)D^\lambda_{mn}.                           \label{eq:s7v38-PW}
\end{equation}
On a finite link set \(E\), the product basis
\(\prod_{\ell\in E}D^{\lambda_\ell}_{m_\ell n_\ell}(U_\ell)\) is an
eigenbasis of
\begin{equation}
 H_E={g_0^2\over2a}\sum_{\ell\in E}(-\Delta_\ell)       \label{eq:s7v38-HE}
\end{equation}
with eigenvalue
\begin{equation}
 {g_0^2\over2a}\sum_{\ell\in E}C_2(\lambda_\ell).
 \label{eq:s7v38-HE-eigen}
\end{equation}
\end{lemma}

\begin{proof}
The Casimir element is central in the universal enveloping algebra and acts
as the scalar \(C_2(\lambda)\) in an irreducible representation.  Its left
or right regular action on a matrix coefficient is
\(-\Delta_G\), proving \eqref{eq:s7v38-PW}.  The link Laplacians act on
different tensor factors and their eigenvalues add, giving
\eqref{eq:s7v38-HE-eigen}.
\end{proof}

\begin{lemma}[Representation change under a link multiplier]
\label{lem:s7v38-product}
Multiplication by a matrix coefficient \(D^F_{ab}\) maps the
\(\lambda\)-Peter--Weyl sector into the sum of sectors
\begin{equation}
 \bigoplus_{\mu\subset F\otimes\lambda}\mu.             \label{eq:s7v38-tensor-product}
\end{equation}
The corresponding electric-energy changes are
\begin{equation}
 {g_0^2\over2a}
 \bigl(C_2(\mu)-C_2(\lambda)\bigr),\qquad
 \mu\subset F\otimes\lambda.                           \label{eq:s7v38-energy-change}
\end{equation}
They need not be positive.
\end{lemma}

\begin{proof}
The product of the two matrix coefficients is a matrix coefficient of the
tensor-product representation \(F\otimes\lambda\).  Decompose that
representation into irreducibles and apply
\Cref{lem:s7v38-PW}.  Some summands can have smaller Casimir than
\(\lambda\), so no sign assertion follows.
\end{proof}

\section{The plaquette-loop obstruction}
\label{sec:s7v38-plaquette}

Specialize to \(G=SU(N)\), with
\begin{equation}
 C_F={N^2-1\over2N}                                    \label{eq:s7v38-CF}
\end{equation}
for the fundamental representation.  Let \(p\) be a plaquette with four
distinct oriented links, and put
\begin{equation}
 \Phi_p(U):=\operatorname{Tr}
 \left(U_{\ell_1}U_{\ell_2}U_{\ell_3}^{-1}U_{\ell_4}^{-1}\right).
 \label{eq:s7v38-Phi}
\end{equation}

\begin{theorem}[Exact electric energy of a fundamental plaquette]
\label{thm:s7v38-plaquette-energy}
The function \(\Phi_p\) is gauge invariant, orthogonal to the constant
vacuum, and
\begin{equation}
 H_E\Phi_p
 ={2g_0^2C_F\over a}\Phi_p.                            \label{eq:s7v38-plaquette-energy}
\end{equation}
Consequently the first nonzero gauge-invariant electric energy is at most
\begin{equation}
 {2g_0^2C_F\over a}.                                   \label{eq:s7v38-electric-upper}
\end{equation}
\end{theorem}

\begin{proof}
Gauge transformations conjugate the based plaquette holonomy, so its trace
is invariant.  Integrating over any one of the four distinct links gives
zero because every nontrivial irreducible matrix coefficient has zero Haar
mean; hence \(\Phi_p\) is orthogonal to constants.  On each plaquette link,
\(\Phi_p\) is a matrix coefficient of the fundamental or dual fundamental
representation, both with Casimir \(C_F\).  The four link Laplacians
therefore contribute \(4C_F\).  Multiplication by the coefficient
\(g_0^2/(2a)\) in \eqref{eq:s7v38-HE} gives
\eqref{eq:s7v38-plaquette-energy}.  The Rayleigh quotient of this nonzero
gauge-invariant vector proves \eqref{eq:s7v38-electric-upper}.
\end{proof}

\begin{corollary}[No coupling-independent residual floor from \(H_E\)]
\label{cor:s7v38-no-electric-floor}
As \(g_0\to0\), no residual subspace containing the fundamental plaquette
state can satisfy
\begin{equation}
 H_E\ge{\Delta\over a}                                 \label{eq:s7v38-false-floor}
\end{equation}
with \(\Delta>0\) independent of \(g_0\).
\end{corollary}

\begin{proof}
The ratio between the eigenvalue in
\eqref{eq:s7v38-plaquette-energy} and \(a^{-1}\) is
\(2C_Fg_0^2\to0\).
\end{proof}

\subsection{Magnetic multiplier size}
\label{sec:s7v38-magnetic}

Use the standard nonnegative plaquette function
\begin{equation}
 w_p(U):=
 1-{1\over N}\operatorname{Re}\operatorname{Tr}U_p.
 \label{eq:s7v38-wp}
\end{equation}

\begin{lemma}[Exact uniform multiplier bound]
\label{lem:s7v38-magnetic}
For every \(U_p\in SU(N)\),
\begin{equation}
 0\le w_p(U)\le2.                                      \label{eq:s7v38-wp-bound}
\end{equation}
Thus an individual term
\begin{equation}
 H_{B,p}:={1\over ag_0^2}w_p                           \label{eq:s7v38-HBp}
\end{equation}
obeys
\begin{equation}
 {4\over Nag_0^2}
 \le\|H_{B,p}\|
 \le {2\over ag_0^2}.                                  \label{eq:s7v38-HBp-bound}
\end{equation}
\end{lemma}

\begin{proof}
Every eigenvalue of \(U_p\) has real part in \([-1,1]\), so
\(-N\le\operatorname{Re}\operatorname{Tr}U_p\le N\).
Equation \eqref{eq:s7v38-wp-bound} follows.  The matrix
\[
 U_*=\operatorname{diag}(-1,-1,1,\ldots,1)\in SU(N)
\]
has
\(\operatorname{Re}\operatorname{Tr}U_*=N-4\), hence
\(w_p(U_*)=4/N\).  A plaquette holonomy can equal \(U_*\), and continuity
gives sets of positive Haar measure with values arbitrarily close to it.
The multiplication-operator norm is the essential supremum, proving both
sides of \eqref{eq:s7v38-HBp-bound}.
\end{proof}

\begin{theorem}[Raw electric/magnetic Feshbach split has the wrong powers]
\label{thm:s7v38-raw-no-go}
A Feshbach split in which \(H_E\) alone supplies the residual operator and
the full \(H_B\) is treated as its perturbation cannot verify the V37 scale
card
\[
 H_{\rm r}\ge{\Delta\over a},\qquad
 \|K\|_{\rm anchor}=O(g_0/a)
\]
uniformly as \(g_0\to0\).  The electric test vector gives only an
\(O(g_0^2/a)\) residual scale, while one magnetic plaquette multiplier has
the upper scale \(O(1/(ag_0^2))\).
\end{theorem}

\begin{proof}
The first assertion is \Cref{cor:s7v38-no-electric-floor}.  The second is
\Cref{lem:s7v38-magnetic}: if the full magnetic multiplier is assigned to
the perturbation, even one plaquette has norm at least
\(4/(Nag_0^2)\).  These powers are incompatible with the stated V37
hypotheses.  The conclusion concerns this raw split only; cancellations
between electric and magnetic quadratic terms have not yet been used.
\end{proof}

\section{Exact quadratic oscillator rescaling}
\label{sec:s7v38-oscillator}

Let \(q\in\mathbb R^n\) be local Lie-algebra coordinates after removal of
gauge zero modes, and let \(L\) be a positive real symmetric matrix.  Consider
\begin{equation}
 H_{\rm quad}(g)
 :=
 {g^2\over2a}(-\Delta_q)
 +{1\over2g^2a}q^{\mathsf T}Lq.                        \label{eq:s7v38-Hquad}
\end{equation}
Define the unitary dilation
\begin{equation}
 ({\cal D}_g\psi)(y):=g^{n/2}\psi(gy),
 \qquad y={q\over g}.                                  \label{eq:s7v38-Dg}
\end{equation}

\begin{theorem}[Coupling cancels from the Gaussian Hamiltonian]
\label{thm:s7v38-dilation}
The dilation gives the exact unitary equivalence
\begin{equation}
 {\cal D}_gH_{\rm quad}(g){\cal D}_g^{-1}
 =
 {1\over2a}
 \left(-\Delta_y+y^{\mathsf T}Ly\right).               \label{eq:s7v38-dilation}
\end{equation}
If \(\lambda_1,\ldots,\lambda_n>0\) are the eigenvalues of \(L\), then,
after subtraction of the ground energy, the spectral gap is
\begin{equation}
 {1\over a}\min_j\sqrt{\lambda_j}.                     \label{eq:s7v38-osc-gap}
\end{equation}
\end{theorem}

\begin{proof}
The change \(q=gy\) gives
\(\partial_{q_j}=g^{-1}\partial_{y_j}\), so the kinetic factor \(g^2\)
cancels.  The quadratic potential contributes
\((g^2y^{\mathsf T}Ly)/(2g^2a)\).  This proves
\eqref{eq:s7v38-dilation}.  Orthogonally diagonalize \(L\).  Each coordinate
is a harmonic oscillator
\[
 {1\over2a}\left(-{d^2\over dy_j^2}+\lambda_jy_j^2\right)
\]
whose consecutive energy levels differ by
\(\sqrt{\lambda_j}/a\).  The smallest excitation proves
\eqref{eq:s7v38-osc-gap}.
\end{proof}

\begin{corollary}[The V37 residual scale is a joint quadratic scale]
\label{cor:s7v38-joint}
If the transverse quadratic magnetic symbol has a dimensionless lower edge
\(\lambda_*>0\) on the residual band, then the joint electric--magnetic
Gaussian Hamiltonian has
\begin{equation}
 H_{{\rm quad},\rm r}-E_0
 \ge{\sqrt{\lambda_*}\over a}P_{\rm r}.                \label{eq:s7v38-joint-floor}
\end{equation}
Thus the V37 constant is
\(\Delta=\sqrt{\lambda_*}\), independent of \(g\).
\end{corollary}

\begin{proof}
Apply \eqref{eq:s7v38-osc-gap} to every residual normal mode.
\end{proof}

\begin{firewall}[Gauge zero modes must be removed before the oscillator gap]
\label{fire:s7v38-zero}
The magnetic Hessian \(L\) vanishes on infinitesimal gauge directions and,
on a torus, on the flat toron sector identified in V7.  The positivity
hypothesis in \Cref{thm:s7v38-dilation} applies only after exact gauge
projection and separate treatment of global flat modes.  Applying
\eqref{eq:s7v38-osc-gap} to the unreduced Hessian would be false.
\end{firewall}

\section{Power counting after dilation}
\label{sec:s7v38-power}

\begin{lemma}[Exact monomial scaling rule]
\label{lem:s7v38-monomial}
Let
\[
 M_{r,s}(q,\partial_q)
\]
be a differential monomial homogeneous of degree \(r\) in \(q\) and degree
\(s\) in derivatives.  Then
\begin{equation}
 {\cal D}_g
 \left[
 {g^\eta\over a}M_{r,s}(q,\partial_q)
 \right]
 {\cal D}_g^{-1}
 =
 {g^{\eta+r-s}\over a}
 M_{r,s}(y,\partial_y).                                \label{eq:s7v38-monomial}
\end{equation}
\end{lemma}

\begin{proof}
Every coordinate contributes \(q=gy\), and every derivative contributes
\(\partial_q=g^{-1}\partial_y\).  Multiply the powers.
\end{proof}

\begin{corollary}[Target weak-coupling expansion]
\label{cor:s7v38-target-expansion}
Suppose the gauge-fixed compact Hamiltonian in exponential coordinates has,
after inclusion of the quadratic electric and magnetic parts, the local
analytic expansion
\begin{equation}
 {\cal D}_g(H-E_0){\cal D}_g^{-1}
 =
 {1\over a}
 \left(
 H_2+gH_3+g^2H_4+\cdots
 \right).                                              \label{eq:s7v38-expansion}
\end{equation}
Then its leading off-diagonal interaction between Gaussian spectral bands is
\begin{equation}
 K={g\over a}P_{\rm L}H_3P_{\rm r}
 +O(g^2/a),                                            \label{eq:s7v38-Ktarget}
\end{equation}
which has exactly the V37 dimension.
\end{corollary}

\begin{proof}
Project \eqref{eq:s7v38-expansion} between the two spectral sectors.  The
quadratic operator \(H_2\) is block diagonal by construction.  The first
remaining term has coefficient \(g/a\).
\end{proof}

\begin{firewall}[The analytic expansion is local, not yet global]
\label{fire:s7v38-chart}
Equation \eqref{eq:s7v38-expansion} is the target form of the weak-field
chart.  V38 has proved its scaling consequences, not a global unitary
equivalence between \(L^2(G^E)\) and a Euclidean oscillator space.  The
exponential map is not globally injective, Haar density is nonconstant, and
large-field configurations lie outside a single chart.
\end{firewall}

\section{Gaussian control of the chart boundary}
\label{sec:s7v38-tail}

The compact-chart obstruction is nonperturbative in the Gaussian reference
state.  Let \(\psi_0\) be the normalized ground state of the rescaled
oscillator in \eqref{eq:s7v38-dilation}.  Its squared modulus has the form
\begin{equation}
 |\psi_0(y)|^2=Z^{-1}e^{-y^{\mathsf T}L^{1/2}y}.        \label{eq:s7v38-ground}
\end{equation}

\begin{lemma}[Gaussian chart-exit bound]
\label{lem:s7v38-tail}
Let \(\lambda_*=\min\operatorname{spec}L>0\).  For every fixed chart radius
\(\rho>0\), there are \(C,c>0\), depending on \(n,L,\rho\) but not on \(g\),
such that
\begin{equation}
 \int_{|q|\ge\rho}|\psi_{0,g}(q)|^2\,dq
 =
 \int_{|y|\ge\rho/g}|\psi_0(y)|^2\,dy
 \le
 C e^{-c\rho^2/g^2}.                                   \label{eq:s7v38-tail}
\end{equation}
\end{lemma}

\begin{proof}
Since \(L^{1/2}\ge\sqrt{\lambda_*}I\),
\[
 |\psi_0(y)|^2\le C_0e^{-\sqrt{\lambda_*}|y|^2}.
\]
For \(R\ge1\), radial integration and one integration by parts bound the
Gaussian tail by \(C_1R^{n-2}e^{-\sqrt{\lambda_*}R^2}\), which is at most
\(Ce^{-cR^2}\) after decreasing \(c\).  Set \(R=\rho/g\); bounded remaining
values of \(g\) are absorbed into \(C\).
\end{proof}

\begin{corollary}[Chart exit is smaller than every coupling power]
\label{cor:s7v38-tail-power}
For every \(M>0\),
\begin{equation}
 e^{-c\rho^2/g^2}=o(g^M)\qquad(g\downarrow0).          \label{eq:s7v38-tail-power}
\end{equation}
\end{corollary}

\begin{proof}
Take logarithms:
\(-c\rho^2/g^2-M\log g\to-\infty\).
\end{proof}

\begin{firewall}[Reference-state tails are not interacting tails]
\label{fire:s7v38-tail}
\Cref{lem:s7v38-tail} controls the Gaussian ground state only.  A
cutoff-uniform interacting large-field bound must still show that the Wilson
measure or transfer ground state remains dominated by a comparable
small-field tail, including compact copies and flat directions.  That bound
cannot be inferred solely from perturbation theory inside the chart.
\end{firewall}

\section{Corrected Feshbach input}
\label{sec:s7v38-card}

\begin{theorem}[What the compact-link test establishes]
\label{thm:s7v38-result}
The V37 Feshbach card cannot be based on the electric Hamiltonian alone.
It can have the required dimensions after the following exact order:
\begin{enumerate}[label=\textup{(\roman*)},leftmargin=3.5em]
\item remove gauge and toron zero modes;
\item combine electric energy with the quadratic magnetic Hessian;
\item dilate the weak link coordinate \(q=gy\);
\item define low and residual bands spectrally for the resulting
      coupling-independent oscillator; and
\item include the \(O(g^2/a)\) closed high return in the effective low
      Hamiltonian.
\end{enumerate}
At the quadratic level the residual floor is \(\Delta/a\); under the local
analytic expansion \eqref{eq:s7v38-expansion}, the leading open mixing is
\(O(g/a)\), exactly as required by V37.
\end{theorem}

\begin{proof}
The failure of the electric split is
\Cref{thm:s7v38-plaquette-energy,thm:s7v38-raw-no-go}.  The corrected
quadratic floor is
\Cref{thm:s7v38-dilation,cor:s7v38-joint}; the interaction power is
\Cref{lem:s7v38-monomial,cor:s7v38-target-expansion}; the self-energy
assignment is the V37 Feshbach identity.
\end{proof}

\begin{assessment}[Progress at seventh-series V38]
\label{ass:s7v38-status}
The residual energy denominator used in V35--V37 now has the correct
microscopic origin.  It is the frequency of the transverse joint
electric--magnetic oscillator, not the compact electric Casimir by itself.
The raw Peter--Weyl split is ruled out by an explicit gauge-invariant
plaquette eigenfunction.  The weak-field oscillator split has the desired
\(\Delta/a\), \(g/a\), and \(g^2/a\) hierarchy, but its global compact-group
and large-field justification remains open.
\end{assessment}

\begin{decision}[V39 acquisition order]
\label{dec:s7v38-V39}
V39 will construct a cutoff weak-field localization for the compact link
space.  It will use a smooth gauge-invariant small-plaquette cutoff, derive
an IMS localization formula for the electric Laplacian, and compare the
inside operator with the rescaled oscillator.  The target is a dichotomy:
oscillator Feshbach control inside the weak-field region and a positive
large-field energy barrier outside.  If no volume-uniform barrier exists
because of flat compact connections, the toron and pure-gauge strata will be
split off before continuing.
\end{decision}

\begin{firewall}[Claim boundary after seventh-series V38]
\label{fire:s7v38-claim}
V38 proves the exact Peter--Weyl electric spectrum, the plaquette-loop
obstruction to an electric-only residual gap, the magnetic multiplier bound,
the oscillator dilation and gap, monomial power counting, and the Gaussian
chart-exit estimate.

V38 does not prove the global compact weak-field chart, an interacting
large-field barrier, the incidence-resolvent hypotheses, a uniform effective
low gap, the terminal strong-coupling handoff, cutoff removal, or the full
Yang--Mills mass gap.
\end{firewall}

\paragraph{Parent-version record.}
V38 was formed from
\texttt{Renewal\_Yang\_Mills\_Mass\_Gap\_Research\_Notes\_V37.tex}.
The V37 parent had \(15\,250\) logical source lines, \(554\,321\) bytes,
and SHA--256
\[
 \texttt{0A8957D4164409C75D19BC74BE63B1674B9596CBF862587C2D2CE6369112D57D}.
\]
The next compulsory companion audit is V40.

% =============================================================================
% END APPENDED SEVENTH-SERIES V38 MATERIAL
% =============================================================================

% =============================================================================
% BEGIN APPENDED SEVENTH-SERIES V39 MATERIAL
% =============================================================================

\chapter{Gauge-invariant weak-field IMS localization}
\label{ch:s7v39}

V38 identifies the correct ultraviolet reference operator as the joint
electric--magnetic oscillator in a gauge-reduced weak-field chart.  A fixed
link-coordinate chart is not gauge invariant and misses compact flat
connections.  This note instead localizes with the gauge-invariant
plaquette action.  The key estimate is a self-bounding gradient:
\[
 |\nabla S_B|^2\le C S_B.
\]
It removes a potential volume factor from the IMS derivative.  With a
shrinking threshold
\(\varepsilon=g^2R^2\), the exterior magnetic barrier is \(R^2/a\), the
IMS debit is \(1/(aR^2)\), and, on a fixed contractible block after tree
gauge, the interior oscillator-form error is \(O(gR)\).  Thus
\[
 R\longrightarrow\infty,\qquad gR\longrightarrow0
\]
separates the weak oscillator sector from large curvature.

\section{The plaquette action controls its gradient}
\label{sec:s7v39-gradient}

Let \(G=SU(N)\) with a fixed bi-invariant metric.  Retain
\[
 w(U)=1-{1\over N}\operatorname{Re}\operatorname{Tr}U.
\]
Let \({\mathscr B}\) be a fixed finite spatial lattice block, with link set
\(E({\mathscr B})\) and plaquette set \(P({\mathscr B})\).  Define
\begin{equation}
 S_{\mathscr B}(U):=\sum_{p\in P({\mathscr B})}w(U_p).
 \label{eq:s7v39-S}
\end{equation}

\begin{lemma}[Single-holonomy gradient bound]
\label{lem:s7v39-single}
There is \(C_G<\infty\) such that
\begin{equation}
 |\nabla_Gw(U)|^2\le C_Gw(U)
 \qquad(U\in G).                                       \label{eq:s7v39-single}
\end{equation}
\end{lemma}

\begin{proof}
The function \(w\) is smooth, nonnegative, and vanishes only at the identity.
In exponential coordinates \(U=e^X\) near the identity,
\[
 w(e^X)=Q(X)+O(|X|^3),\qquad
 |\nabla w(e^X)|^2=O(|X|^2),
\]
where \(Q\) is a positive quadratic form for the chosen invariant metric.
Hence \(|\nabla w|^2/w\) is bounded in a punctured neighborhood of the
identity and extends there with a finite upper limit.  On the compact
complement of a smaller neighborhood, \(w\) has a positive minimum while
\(|\nabla w|\) is bounded.  Combining the two regions proves
\eqref{eq:s7v39-single}.
\end{proof}

\begin{theorem}[Plaquette-action self-bound]
\label{thm:s7v39-self-bound}
Let
\begin{equation}
 \nabla S_{\mathscr B}
 :=
 \bigl(\nabla_\ell S_{\mathscr B}\bigr)_{\ell\in E({\mathscr B})}.
\end{equation}
If at most \(d_*\) plaquettes meet a link and every plaquette has four links,
then
\begin{equation}
 \sum_{\ell\in E({\mathscr B})}
 |\nabla_\ell S_{\mathscr B}|^2
 \le 4d_*C_G S_{\mathscr B}.                           \label{eq:s7v39-self-bound}
\end{equation}
The constant is independent of the number of plaquettes in the block.
\end{theorem}

\begin{proof}
For a fixed link, Cauchy--Schwarz gives
\[
 |\nabla_\ell S_{\mathscr B}|^2
 \le
 d_*\sum_{p\ni\ell}|\nabla_\ell w(U_p)|^2.
\]
Left or right multiplication by the other link variables in \(U_p\) is an
isometry for the bi-invariant metric.  Therefore
\(\lvert\nabla_\ell w(U_p)\rvert^2\le C_Gw(U_p)\) by
\Cref{lem:s7v39-single}.  Sum over links.  Each plaquette appears four times,
which proves \eqref{eq:s7v39-self-bound}.
\end{proof}

\section{An exact gauge-invariant IMS partition}
\label{sec:s7v39-IMS}

Let \(\vartheta\in C^\infty([0,\infty);[0,\pi/2])\) satisfy
\[
 \vartheta(t)=0\quad(t\le1),\qquad
 \vartheta(t)={\pi\over2}\quad(t\ge2),\qquad
 |\vartheta'|\le C_\vartheta.
\]
For \(\varepsilon>0\), define
\begin{equation}
 \chi_{\rm in}
 :=
 \cos\vartheta(S_{\mathscr B}/\varepsilon),\qquad
 \chi_{\rm out}
 :=
 \sin\vartheta(S_{\mathscr B}/\varepsilon).            \label{eq:s7v39-cutoffs}
\end{equation}
Then
\begin{equation}
 \chi_{\rm in}^2+\chi_{\rm out}^2=1.                   \label{eq:s7v39-partition}
\end{equation}
Both multipliers are gauge invariant.

On \(L^2(G^{E({\mathscr B})})\), consider the block Hamiltonian form
\begin{equation}
 q_{{\mathscr B},g}(\psi)
 :=
 {g^2\over2a}\sum_{\ell}\|\nabla_\ell\psi\|_2^2
 +{1\over ag^2}
 \int S_{\mathscr B}|\psi|^2\,dU.                      \label{eq:s7v39-q}
\end{equation}

\begin{theorem}[Gauge-invariant IMS identity]
\label{thm:s7v39-IMS}
For every form-domain vector \(\psi\),
\begin{align}
 q_{{\mathscr B},g}(\psi)
 ={}&
 q_{{\mathscr B},g}(\chi_{\rm in}\psi)
 +q_{{\mathscr B},g}(\chi_{\rm out}\psi)               \notag\\
 &-
 {g^2\over2a}
 \int
 \left(
 |\nabla\chi_{\rm in}|^2+
 |\nabla\chi_{\rm out}|^2
 \right)|\psi|^2\,dU.                                  \label{eq:s7v39-IMS}
\end{align}
Moreover,
\begin{equation}
 |\nabla\chi_{\rm in}|^2+
 |\nabla\chi_{\rm out}|^2
 \le {C_{\rm IMS}\over\varepsilon},                    \label{eq:s7v39-IMS-gradient}
\end{equation}
where
\begin{equation}
 C_{\rm IMS}:=8d_*C_GC_\vartheta^2                    \label{eq:s7v39-CIMS}
\end{equation}
is independent of block volume.
\end{theorem}

\begin{proof}
The usual quadratic-form IMS expansion follows from the product rule.
The cross terms cancel because
\(\chi_{\rm in}\nabla\chi_{\rm in}
+\chi_{\rm out}\nabla\chi_{\rm out}
=\frac12\nabla(\chi_{\rm in}^2+\chi_{\rm out}^2)=0\).
The magnetic multiplier commutes with the cutoffs and splits exactly by
\eqref{eq:s7v39-partition}.

The chain rule gives
\[
 |\nabla\chi_{\rm in}|^2+
 |\nabla\chi_{\rm out}|^2
 =
 {|\vartheta'(S_{\mathscr B}/\varepsilon)|^2\over\varepsilon^2}
 |\nabla S_{\mathscr B}|^2.
\]
The derivative is supported where
\(\varepsilon\le S_{\mathscr B}\le2\varepsilon\).  Apply
\Cref{thm:s7v39-self-bound} and the bound on \(\vartheta'\) to obtain
\eqref{eq:s7v39-IMS-gradient}.
\end{proof}

\begin{corollary}[Interior/exterior energy dichotomy]
\label{cor:s7v39-dichotomy}
For every \(\psi\),
\begin{align}
 q_{{\mathscr B},g}(\psi)
 \ge{}&
 q_{{\mathscr B},g}(\chi_{\rm in}\psi)
 +{\varepsilon\over ag^2}
 \|\chi_{\rm out}\psi\|_2^2
 -{C_{\rm IMS}g^2\over2a\varepsilon}\|\psi\|_2^2.
 \label{eq:s7v39-dichotomy}
\end{align}
\end{corollary}

\begin{proof}
On the support of \(\chi_{\rm out}\), one has
\(S_{\mathscr B}\ge\varepsilon\).  Drop its nonnegative electric energy in
\eqref{eq:s7v39-IMS} and use
\eqref{eq:s7v39-IMS-gradient}.
\end{proof}

\section{The shrinking weak-field threshold}
\label{sec:s7v39-threshold}

Let \(R=R(g)\ge1\) and choose
\begin{equation}
 \varepsilon(g):=g^2R(g)^2.                            \label{eq:s7v39-epsilon}
\end{equation}

\begin{theorem}[Barrier, localization debit, and scale separation]
\label{thm:s7v39-scales}
With \eqref{eq:s7v39-epsilon},
\begin{equation}
 q_{{\mathscr B},g}(\psi)
 \ge
 q_{{\mathscr B},g}(\chi_{\rm in}\psi)
 +{R(g)^2\over a}\|\chi_{\rm out}\psi\|_2^2
 -{C_{\rm IMS}\over2aR(g)^2}\|\psi\|_2^2.              \label{eq:s7v39-scales}
\end{equation}
If
\begin{equation}
 R(g)\longrightarrow\infty,\qquad
 gR(g)\longrightarrow0,                               \label{eq:s7v39-window}
\end{equation}
then the exterior barrier diverges relative to \(a^{-1}\), the IMS debit is
\(o(a^{-1})\), and the plaquette holonomies on the support of
\(\chi_{\rm in}\) converge uniformly to the identity.
\end{theorem}

\begin{proof}
Substitute \eqref{eq:s7v39-epsilon} in
\eqref{eq:s7v39-dichotomy}.  On the support of
\(\chi_{\rm in}\), one has
\[
 0\le w(U_p)\le S_{\mathscr B}\le2g^2R(g)^2
\]
for every plaquette \(p\).  Since \(w(U)\) has a nondegenerate quadratic
minimum at \(I\), there are \(c,\rho>0\) such that
\[
 w(U)\ge c\,d_G(U,I)^2
\]
whenever \(w(U)\le\rho\).  The last condition in
\eqref{eq:s7v39-window} makes the displayed plaquette actions smaller than
\(\rho\), and hence
\[
 \max_p d_G(U_p,I)\le CgR(g)\longrightarrow0.
\]
The two energy-scale assertions are immediate.
\end{proof}

\begin{corollary}[Explicit asymptotically free schedule]
\label{cor:s7v39-explicit}
Let
\[
 L(a)=\log{1\over a\Lambda_0},\qquad
 g(a)\le C L(a)^{-1/2},
\]
and choose
\begin{equation}
 R(a)=L(a)^{1/16}.                                     \label{eq:s7v39-R}
\end{equation}
Then
\begin{align}
 g(a)R(a)&=O(L(a)^{-7/16}),                            \label{eq:s7v39-chart-rate}\\
 {R(a)^2\over a}&={L(a)^{1/8}\over a},                 \label{eq:s7v39-barrier-rate}\\
 {1\over aR(a)^2}&={1\over aL(a)^{1/8}}.               \label{eq:s7v39-IMS-rate}
\end{align}
The IMS debit is smaller than the exterior barrier by
\(O(L(a)^{-1/4})\).
\end{corollary}

\begin{proof}
These are direct exponent calculations:
\(-1/2+1/16=-7/16\), and
\(R^2=L^{1/8}\).
\end{proof}

\begin{remark}[The IMS debit is a residual-scale debit]
\label{rem:s7v39-relative}
The quantity \(1/(aR^2)\) need not tend to zero in physical energy units.
It is small relative to the ultraviolet oscillator and exterior scales
\(1/a\) and \(R^2/a\).  It therefore belongs in the residual lower-bound
ledger.  It may not be charged directly against an \(O(1)\) continuum mass
without an additional low-state boundary-tail estimate.
\end{remark}

\section{From small plaquettes to a link chart}
\label{sec:s7v39-tree}

Let \({\mathscr B}\) now be a fixed finite contractible cubical block with a
rooted maximal tree \({\cal T}\) in its one-skeleton.  Gauge fixing along
\({\cal T}\) sets every tree link to the identity.  For every chord
\(\ell\notin{\cal T}\), fix a plaquette filling of its fundamental cycle and
let \(A_\ell\) be the number of plaquettes in that filling.  Put
\begin{equation}
 A_{\mathscr B}:=\max_{\ell\notin{\cal T}}A_\ell<\infty.
 \label{eq:s7v39-area}
\end{equation}

\begin{lemma}[Finite-block nonabelian reconstruction]
\label{lem:s7v39-reconstruction}
There are \(\rho_{\mathscr B},C_{\mathscr B}>0\) such that, if every
plaquette holonomy obeys
\[
 d_G(U_p,I)\le\rho_{\mathscr B},
\]
then in maximal-tree gauge every chord link obeys
\begin{equation}
 d_G(U_\ell,I)
 \le
 C_{\mathscr B}\max_p d_G(U_p,I).                      \label{eq:s7v39-reconstruction}
\end{equation}
One may take \(C_{\mathscr B}\) proportional to \(A_{\mathscr B}\).
\end{lemma}

\begin{proof}
Discrete nonabelian elimination across the chosen plaquette filling writes
the fundamental-cycle holonomy as an ordered product of at most
\(A_\ell\) plaquette holonomies and their conjugates.  In tree gauge the
fundamental-cycle holonomy is the chord link itself.  A bi-invariant distance
obeys
\[
 d_G(V_1\cdots V_k,I)
 \le\sum_{j=1}^kd_G(V_j,I)
\]
and is unchanged by conjugation.  Hence
\[
 d_G(U_\ell,I)
 \le A_\ell\max_p d_G(U_p,I).
\]
Choosing the plaquette neighborhood inside the injectivity radius supplies
the stated constants.
\end{proof}

\begin{corollary}[Gauge-fixed chart on the IMS interior]
\label{cor:s7v39-chart}
Under \eqref{eq:s7v39-window}, the support of \(\chi_{\rm in}\) on a fixed
contractible block admits maximal-tree gauge coordinates
\begin{equation}
 U_\ell=\exp(q_\ell),\qquad
 \max_{\ell\notin{\cal T}}|q_\ell|
 \le C_{\mathscr B}gR(g).                              \label{eq:s7v39-q-bound}
\end{equation}
\end{corollary}

\begin{proof}
Combine the plaquette estimate in
\Cref{thm:s7v39-scales} with
\Cref{lem:s7v39-reconstruction}.  For sufficiently small \(g\), all chord
links lie in the exponential-map injectivity neighborhood.
\end{proof}

\begin{firewall}[Contractibility is essential]
\label{fire:s7v39-toron}
On a periodic torus, small or zero plaquette action does not force the
noncontractible holonomies near the identity.  Flat torons survive tree
gauge.  Therefore \Cref{lem:s7v39-reconstruction} is a fixed-block statement,
not a global periodic-lattice chart.  Toron variables must remain separate
compact low channels, as already required by V7.
\end{firewall}

\section{Interior oscillator-form comparison}
\label{sec:s7v39-comparison}

After tree gauge and removal of the remaining infinitesimal gauge directions,
write \(q=gy\).  Let \(q_{\rm osc}\) be the quadratic form of the transverse
joint oscillator from V38, including its Haar-flat kinetic term and quadratic
magnetic Hessian.

\begin{theorem}[Uniform fixed-block weak-field comparison]
\label{thm:s7v39-comparison}
There are constants \(c_0,C_{\mathscr B}>0\) such that, whenever
\[
 R\ge1,\qquad gR\le c_0,
\]
the gauge-fixed Hamiltonian form on functions supported in
\(\max|y_\ell|\le C_{\mathscr B}R\) obeys
\begin{equation}
 \left|
 q_{{\rm in},g}(\phi)-q_{\rm osc}(\phi)
 \right|
 \le
 C_{\mathscr B}gR
 \left(
 q_{\rm osc}(\phi)+{1\over a}\|\phi\|_2^2
 \right).
 \label{eq:s7v39-comparison}
\end{equation}
The comparison is uniform in \(a\); its constants depend on the fixed block
and the declared gauge chart.
\end{theorem}

\begin{proof}
In exponential coordinates, the inverse Haar metric, Haar density, and
plaquette holonomies are analytic.  Their quadratic parts are exactly the
flat electric metric and the magnetic Hessian used in \(q_{\rm osc}\).
On the coordinate ball \(|q|\le C_{\mathscr B}gR\), the coefficient
differences from their values at zero are bounded by
\(C_{\mathscr B}gR\).  The magnetic Taylor remainder obeys
\[
 |S_{\mathscr B}(q)-Q_2(q)|
 \le C_{\mathscr B}|q|\,Q_2(q)
\]
after restriction to the transverse finite-dimensional space.  With
\(q=gy\), division by \(ag^2\) bounds it by
\(C_{\mathscr B}gR\) times the quadratic magnetic form.

For the electric term, the analytic metric coefficients give the same
relative bound on the principal gradient form.  Conjugating the smooth Haar
density to flat measure produces first- and zeroth-order terms; Cauchy--
Schwarz absorbs the former into the principal form, while bounded analytic
coefficients bound the latter by
\(C_{\mathscr B}gR\,a^{-1}\|\phi\|_2^2\) for \(R\ge1\).
Adding the estimates proves \eqref{eq:s7v39-comparison}.
\end{proof}

\begin{corollary}[Residual oscillator floor survives]
\label{cor:s7v39-residual}
Let the oscillator residual sector have
\[
 q_{\rm osc}(\phi)\ge{\Delta\over a}\|\phi\|_2^2.
\]
For sufficiently small \(gR\), the interior comparison gives
\begin{equation}
 q_{{\rm in},g}(\phi)
 \ge{\Delta'\over a}\|\phi\|_2^2                       \label{eq:s7v39-residual}
\end{equation}
for every fixed \(0<\Delta'<\Delta\), after the common vacuum-energy
normalization is included in both forms.
\end{corollary}

\begin{proof}
Use the lower side of \eqref{eq:s7v39-comparison}.  The error is at most a
constant times \(gR\) times \(q_{\rm osc}+a^{-1}\|\phi\|^2\).  Since the
oscillator band under consideration has a fixed dimensionless lower edge,
choosing \(gR\) sufficiently small leaves any prescribed smaller edge
\(\Delta'\).  The common scalar ground-energy shift does not change a gap.
\end{proof}

\begin{firewall}[Upper spectral control is needed in the relative estimate]
\label{fire:s7v39-upper}
The last absorption is immediate on a fixed oscillator band or spectral
window.  On the entire unbounded oscillator residual space, a two-sided
relative-form argument must also control the high-energy coefficient error;
V39 does not claim a uniform comparison of all arbitrarily high oscillator
levels inside a finite coordinate chart.
\end{firewall}

\section{Weak-field localization card}
\label{sec:s7v39-card}

\begin{definition}[Gauge-invariant weak-field localization card, GIWLC]
\label{def:s7v39-card}
For each fixed contractible RG block, GIWLC records:
\begin{enumerate}[label=\textup{(W\arabic*)},leftmargin=4em]
\item the action \(S_{\mathscr B}\) and its self-bound
      \eqref{eq:s7v39-self-bound};
\item a scale \(R(a)\to\infty\) with \(g(a)R(a)\to0\);
\item the cutoff threshold \(\varepsilon(a)=g(a)^2R(a)^2\);
\item the exterior barrier \(R(a)^2/a\);
\item the IMS debit \(C/(aR(a)^2)\);
\item a maximal-tree chart of radius \(O(g(a)R(a))\);
\item separate toron and boundary-holonomy channels; and
\item the oscillator comparison error \(O(g(a)R(a))\) on the declared
      residual spectral window.
\end{enumerate}
\end{definition}

\begin{theorem}[What V39 proves]
\label{thm:s7v39-result}
On every fixed contractible lattice block, GIWLC items (W1)--(W6) and (W8)
hold.  For the one-loop schedule \(R=L^{1/16}\), the chart radius tends to
zero as \(L^{-7/16}\), the exterior/IMS scale ratio grows as \(L^{1/4}\),
and the joint oscillator retains a coupling-independent residual energy
floor.  Item (W7) is absent on a contractible block and becomes compulsory
under periodic gluing.
\end{theorem}

\begin{proof}
Use
\Cref{thm:s7v39-self-bound,thm:s7v39-IMS,
thm:s7v39-scales,cor:s7v39-explicit,
lem:s7v39-reconstruction,thm:s7v39-comparison,
cor:s7v39-residual}.
\end{proof}

\begin{assessment}[Progress at seventh-series V39]
\label{ass:s7v39-status}
The compact weak-field chart is no longer assumed globally.  A
gauge-invariant plaquette cutoff separates large curvature with a strong
magnetic barrier, and maximal-tree gauge converts the interior of each fixed
contractible block into the shrinking oscillator chart required by V38.
The remaining global difficulty is not local compactness; it is gluing the
block localizations without an extensive IMS debit while retaining torons,
boundary holonomies, and the connected low effective action.
\end{assessment}

\begin{decision}[V40 acquisition order]
\label{dec:s7v39-V40}
V40 is the compulsory SM/NS audit.  After that audit it will formulate an
owner assignment for overlapping GIWLC blocks, test whether the IMS
derivative can be charged once per plaquette, and isolate the boundary
holonomy/toron Schur complement.  The next construction must remain
gauge-invariant and must not sum a fixed-block localization debit over the
entire spatial volume.
\end{decision}

\begin{firewall}[Claim boundary after seventh-series V39]
\label{fire:s7v39-claim}
V39 proves the action-gradient self-bound, the exact gauge-invariant IMS
identity, the shrinking-threshold barrier/debit ledger, fixed-block
nonabelian reconstruction, and the interior oscillator-form comparison.

V39 does not prove volume-uniform block gluing, periodic toron control, the
nonlinear connected Feshbach incidence bounds, a coarse effective gap, the
terminal strong-coupling handoff, cutoff removal, or the full Yang--Mills
mass gap.
\end{firewall}

\paragraph{Parent-version record.}
V39 was formed from
\texttt{Renewal\_Yang\_Mills\_Mass\_Gap\_Research\_Notes\_V38.tex}.
The V38 parent had \(15\,734\) logical source lines, \(572\,703\) bytes,
and SHA--256
\[
 \texttt{98A2A2980410912A1FF4E6165BD6FEC788EA6B5CAA8D31C98D21A519C8FA2AB8}.
\]
The next version is the compulsory V40 companion audit.

% =============================================================================
% END APPENDED SEVENTH-SERIES V39 MATERIAL
% =============================================================================

% =============================================================================
% BEGIN APPENDED SEVENTH-SERIES V40 MATERIAL
% =============================================================================

\chapter{Companion audit and block-owned residual elimination}
\label{ch:s7v40}

V39 proves a gauge-invariant weak-field localization on every fixed
contractible block, but a global sum of fixed-block IMS debits would be
volume extensive.  This compulsory-audit note first examines the current SM
and NS programmes.  It then changes the organization of the Yang--Mills
elimination: microscopic links are assigned to disjoint real-space blocks,
the complete slow space of every block is retained, and only the
tensor-product residual is eliminated.  Boundary plaquettes receive unique
owners.  This gives a volume-independent residual floor and an anchored
interaction norm.

At second Feshbach order, an exact overlap rule appears: two boundary
interactions with disjoint block supports cannot make a low-to-low return.
The self-energy is therefore a sum over overlapping block pairs and has
anchored size \(O(g_{\rm YM}^2/a)\), with no factor equal to the total number
of blocks.

\section{Compulsory V40 companion audit}
\label{sec:s7v40-audit}

\subsection{Files inspected}
\label{sec:s7v40-files}

The current active companion files are
\begin{align}
 &\texttt{Renewal\_SM\_Einstein\_Continuum\_Research\_Notes\_V82.tex},
 \notag\\[-2pt]
 &\qquad 32\,934\ \hbox{logical source lines},\quad
 1\,170\,293\ \hbox{bytes},                             \notag\\[-2pt]
 &\qquad
 \texttt{1DFB2BAD454C171100F2D8AB134770F529F918DA65D138D0DB9526A78C3D62E5},
 \label{eq:s7v40-sm-checkpoint}\\
 &\texttt{Renewal\_NS\_Stability\_Research\_Notes\_V26.tex},
 \notag\\[-2pt]
 &\qquad 5\,319\ \hbox{logical source lines},\quad
 278\,283\ \hbox{bytes},                                \notag\\[-2pt]
 &\qquad
 \texttt{82C887F8C2C69E4F28FAD7C3DC8DE5B9099CB5A4BF431EBA1FFF3E91935978AD}.
 \label{eq:s7v40-ns-checkpoint}
\end{align}
The next compulsory companion audit is V45.

\subsection{SM V80--V82}
\label{sec:s7v40-sm}

The SM programme has added a finite-range Helmholtz-projector no-go, an exact
quotient by conformal Killing fields, a coupled scalar--vector Schur margin,
and a nonlinear augmented constraint chart.  Its V82 slow-sector theorem is
the relevant new structural input: under a changing background, exact kernel
vectors and lifted near-kernel vectors must be transported together by a
buffered spectral projection.  A lifted mode is not gauge merely because it
was once in a kernel.

The Yang--Mills analogue is immediate at the level of programme design.
Block boundary holonomies, torons, and oscillator modes below the declared
buffer must all remain in the block slow space.  Only the uniformly separated
complement may be assigned the residual floor.  No Einstein constraint
constant is imported.

\subsection{NS V26}
\label{sec:s7v40-ns}

The NS programme remains byte-for-byte at V26.  Its rank-one Oseen-kernel
refinement contains no block gauge quotient, compact-link IMS estimate, or
Feshbach self-energy bound.  No numerical or analytic NS input transfers to
the present construction.

\begin{theorem}[V40 audit decision]
\label{thm:s7v40-audit}
The audit changes the Yang--Mills block definition as follows:
\begin{enumerate}[label=\textup{(\roman*)},leftmargin=3.5em]
\item the block slow projector is a buffered spectral projector, not merely
      the exact block kernel;
\item gauge, toron, boundary-holonomy, and lifted near-kernel directions are
      retained;
\item the residual estimate is proved only on the complementary spectral
      band; and
\item no finite-range exact transverse projector is postulated.
\end{enumerate}
No companion theorem supplies a Yang--Mills gap estimate.
\end{theorem}

\begin{proof}
Items (i)--(iii) are the type-correct transfer of the SM slow-sector
principle, while item (iv) follows from the projector no-go reviewed in SM
V80 and independently from YM V31.  The companion operators and state spaces
are not Yang--Mills transfer operators, so only the structural ordering
transfers.
\end{proof}

\section{Disjoint link blocks and the physical gauge subspace}
\label{sec:s7v40-blocks}

Let a finite spatial lattice have link set \(E\), vertex set \(V\), and
plaquette set \({\cal P}\).  Partition the links into nonempty disjoint sets
\begin{equation}
 E=\bigsqcup_{C\in{\cal C}}E_C.                        \label{eq:s7v40-link-partition}
\end{equation}
The kinematic Hilbert space then factorizes exactly:
\begin{equation}
 {\cal H}_{\rm kin}
 =L^2(G^E)
 =\bigotimes_{C\in{\cal C}}{\cal H}_C,
 \qquad
 {\cal H}_C:=L^2(G^{E_C}).                             \label{eq:s7v40-tensor}
\end{equation}
Blocks may share vertices but not link variables.

Let \({\cal G}=G^V\) act by the usual left/right endpoint transformations,
and let
\begin{equation}
 \Pi_{\rm phys}
 :=
 \int_{\cal G}U(\gamma)\,d\gamma                       \label{eq:s7v40-Gauss}
\end{equation}
be the orthogonal projection onto the gauge-invariant Hilbert space.

\begin{lemma}[Tensorization is compatible with Gauss projection]
\label{lem:s7v40-gauge}
The representation \(U(\gamma)\) acts as a tensor product over the link
blocks in \eqref{eq:s7v40-tensor}.  If a block operator \(h_C\) is invariant
under every endpoint gauge action on its links, then
\begin{equation}
 \left[\sum_Ch_C,\Pi_{\rm phys}\right]=0.               \label{eq:s7v40-gauge-commute}
\end{equation}
Every spectral projection of \(h_C\) has the same covariance.
\end{lemma}

\begin{proof}
Each link variable belongs to exactly one tensor factor, and a vertex gauge
transformation acts on all incident link factors by their endpoint
representations.  Thus the total unitary is the tensor product of those
actions.  Gauge invariance of every \(h_C\) implies commutation with
\(U(\gamma)\), hence with its Haar average
\(\Pi_{\rm phys}\).  Functional calculus preserves commutation.
\end{proof}

\begin{firewall}[The block factorization is kinematic]
\label{fire:s7v40-kinematic}
The physical gauge projection does not factor over blocks at shared vertices.
Equation \eqref{eq:s7v40-tensor} is an exact factorization of link variables,
not an assertion that the physical Hilbert space is a tensor product.
All later projectors commute with \(\Pi_{\rm phys}\) before restriction to the
physical subspace.
\end{firewall}

\section{Buffered block slow spaces}
\label{sec:s7v40-slow}

Let \(h_C\ge E_C\) be a self-adjoint block Hamiltonian containing the electric
terms of links in \(E_C\), every oscillator counterterm assigned internally,
and those plaquette terms whose links all belong to \(E_C\).  Let \(P_C\) be
a gauge-covariant spectral projection containing the complete block spectrum
below a buffer.  Put
\begin{equation}
 Q_C:=I-P_C.                                           \label{eq:s7v40-QC}
\end{equation}
Assume the uniform block residual bound
\begin{equation}
 h_C-E_C
 \ge{\Delta\over a}Q_C                                \label{eq:s7v40-block-gap}
\end{equation}
as a quadratic-form inequality, with \(\Delta>0\) independent of \(C\),
volume, and cutoff.

Define
\begin{align}
 H_0&:=\sum_{C\in{\cal C}}h_C,                         \label{eq:s7v40-H0}\\
 E_0&:=\sum_{C\in{\cal C}}E_C,                         \label{eq:s7v40-E0}\\
 P&:=\bigotimes_{C\in{\cal C}}P_C,\qquad Q:=I-P.       \label{eq:s7v40-PQ}
\end{align}

\begin{theorem}[Tensor residual floor]
\label{thm:s7v40-tensor-gap}
The noninteracting block sum satisfies
\begin{equation}
 Q(H_0-E_0)Q
 \ge{\Delta\over a}Q.                                 \label{eq:s7v40-tensor-gap}
\end{equation}
The same inequality holds after restriction to
\(\Pi_{\rm phys}{\cal H}_{\rm kin}\).
\end{theorem}

\begin{proof}
The operators \(Q_C\) act on distinct tensor factors and commute.  Summing
\eqref{eq:s7v40-block-gap} gives
\begin{equation}
 H_0-E_0
 \ge{\Delta\over a}\sum_CQ_C.                         \label{eq:s7v40-number}
\end{equation}
On the complement of
\(P=\prod_C(I-Q_C)\), at least one block lies in its \(Q_C\) sector, so
\[
 Q\left(\sum_CQ_C\right)Q\ge Q.
\]
This proves \eqref{eq:s7v40-tensor-gap}.  By
\Cref{lem:s7v40-gauge}, all operators commute with
\(\Pi_{\rm phys}\); restricting a form inequality to an invariant subspace
preserves it.
\end{proof}

\begin{remark}[Why all slow modes are retained]
\label{rem:s7v40-slow}
The rank of \(P_C\) may exceed the block ground-state multiplicity.  It may
contain torons, boundary electric fluxes, and lifted modes below the chosen
buffer.  The proof of \Cref{thm:s7v40-tensor-gap} uses only the complementary
floor and is unchanged by enlarging \(P_C\).
\end{remark}

\section{Unique owners and anchored boundary incidence}
\label{sec:s7v40-owner}

Call a plaquette internal if all its links lie in one \(E_C\); all other
plaquettes are boundary plaquettes.  For every boundary plaquette \(p\), let
\begin{equation}
 {\rm supp}_{\cal C}(p)
 :=
 \{C:E_C\cap E(p)\ne\varnothing\}.                     \label{eq:s7v40-block-support}
\end{equation}
Choose once and for all an owner
\begin{equation}
 o(p)\in{\rm supp}_{\cal C}(p).                        \label{eq:s7v40-owner}
\end{equation}
The owner is bookkeeping; the operator still acts on every block in its
support.

Let \(v_p\) be the oscillator-recentered boundary interaction and write
\begin{equation}
 V:=\sum_{p\ {\rm boundary}}v_p.                       \label{eq:s7v40-V}
\end{equation}
Define the anchored block-incidence norm
\begin{equation}
 {\cal J}(V)
 :=
 \sup_{C\in{\cal C}}
 \sum_{p:\,C\in{\rm supp}_{\cal C}(p)}\|v_p\|.          \label{eq:s7v40-J}
\end{equation}

\begin{lemma}[No volume factor in boundary incidence]
\label{lem:s7v40-incidence}
For cubic blocks of fixed side \(b\) in spatial dimension \(D\), there is a
geometric constant \(c_D\) such that
\begin{equation}
 \#\{p:C\in{\rm supp}_{\cal C}(p),\ p\ {\rm boundary}\}
 \le c_Db^{D-1}.                                       \label{eq:s7v40-boundary-count}
\end{equation}
If
\begin{equation}
 \|v_p\|\le {C_vg_{\rm YM}(a)\over a},                 \label{eq:s7v40-vp}
\end{equation}
then
\begin{equation}
 {\cal J}(V)
 \le {c_Db^{D-1}C_vg_{\rm YM}(a)\over a}.              \label{eq:s7v40-J-bound}
\end{equation}
\end{lemma}

\begin{proof}
A boundary plaquette meeting a block lies within one lattice unit of a
codimension-one block face.  The number of such plaquettes is bounded by a
dimension-dependent constant times the total face area \(b^{D-1}\), with
edges and corners absorbed into the same constant for \(b\ge1\).
Sum \eqref{eq:s7v40-vp} over this set.
\end{proof}

\begin{firewall}[The local interaction estimate remains a target]
\label{fire:s7v40-vp}
Equation \eqref{eq:s7v40-vp} can apply only to the nonlinear remainder after
the complete quadratic magnetic form has been included in the reference
operator.  Oscillator recentering by itself is insufficient: a plaquette
crossing a block face still has an \(O(1/a)\) quadratic coupling between the
adjacent block oscillators.  It is also false for the raw Wilson multiplier,
whose norm is \(O(1/(ag_{\rm YM}^2))\).  Therefore the pure block sum
\eqref{eq:s7v40-H0} and the estimate \eqref{eq:s7v40-vp} are not yet known to
hold simultaneously.  V40 proves the conditional block combinatorics; it
does not claim that the compact Hamiltonian estimate has been established.
\end{firewall}

\section{Second-order Feshbach connectedness}
\label{sec:s7v40-second}

Because \(H_0\) and every \(P_C\) are obtained by block functional calculus,
\begin{equation}
 [H_0,P_C]=0\qquad(C\in{\cal C}).                       \label{eq:s7v40-H0-PC}
\end{equation}
Let
\begin{equation}
 R_0
 :=
 Q\bigl(Q(H_0-E_0)Q\bigr)^{-1}Q.                      \label{eq:s7v40-R0}
\end{equation}
By \Cref{thm:s7v40-tensor-gap},
\begin{equation}
 \|R_0\|\le {a\over\Delta}.                            \label{eq:s7v40-R0-bound}
\end{equation}
The second-order low self-energy is
\begin{equation}
 \Sigma^{(2)}
 :=
 PVR_0VP
 =
 \sum_{p,q}Pv_pR_0v_qP.                               \label{eq:s7v40-Sigma2}
\end{equation}

\begin{theorem}[Exact overlap rule]
\label{thm:s7v40-overlap}
If
\begin{equation}
 {\rm supp}_{\cal C}(p)
 \cap
 {\rm supp}_{\cal C}(q)=\varnothing,                   \label{eq:s7v40-disjoint}
\end{equation}
then
\begin{equation}
 Pv_pQR_0Qv_qP=0                                      \label{eq:s7v40-overlap-zero}
\end{equation}
after replacing each \(v_p\) by its low--residual component in the Feshbach
off-diagonal \(K=QVP\) and its adjoint.  Equivalently,
\begin{equation}
 \Sigma^{(2)}
 =
 \sum_{\substack{p,q:\\
 {\rm supp}_{\cal C}(p)\cap{\rm supp}_{\cal C}(q)\ne\varnothing}}
 Pv_pQR_0Qv_qP.                                       \label{eq:s7v40-overlap-sum}
\end{equation}
\end{theorem}

\begin{proof}
Starting from the global slow space \(P{\cal H}_{\rm kin}\), the vector
\(Qv_qP\psi\) has at least one excited block, and every such block belongs to
\({\rm supp}_{\cal C}(q)\), because \(v_q\) acts trivially elsewhere.
Equation \eqref{eq:s7v40-H0-PC} implies that \(R_0\) preserves the
slow/residual status of every block.  If the supports in
\eqref{eq:s7v40-disjoint} are disjoint, \(v_p\) acts trivially on every
excited block created by \(v_q\).  At least one of those excitations
therefore remains, and the left projection \(P\) annihilates the result.
This proves \eqref{eq:s7v40-overlap-zero} and hence the restricted sum.
\end{proof}

\subsection{Anchored self-energy bound}
\label{sec:s7v40-Sigma-bound}

For a family of terms \(A_{pq}\) labelled by overlapping plaquette supports,
define
\begin{equation}
 {\cal J}_2(A)
 :=
 \sup_C
 \sum_{\substack{p,q:\,
 {\rm supp}(p)\cap{\rm supp}(q)\ne\varnothing;\\
 C\in{\rm supp}(p)\cup{\rm supp}(q)}}
 \|A_{pq}\|.                                           \label{eq:s7v40-J2}
\end{equation}
Let
\begin{equation}
 s_*:=\sup_p|{\rm supp}_{\cal C}(p)|\le4.              \label{eq:s7v40-sstar}
\end{equation}

\begin{theorem}[Volume-uniform second-order self-energy]
\label{thm:s7v40-Sigma-bound}
The terms
\[
 A_{pq}:=Pv_pQR_0Qv_qP
\]
obey
\begin{equation}
 {\cal J}_2(\Sigma^{(2)})
 \le {2s_*a\over\Delta}{\cal J}(V)^2.                 \label{eq:s7v40-Sigma-bound}
\end{equation}
Under \eqref{eq:s7v40-J-bound},
\begin{equation}
 {\cal J}_2(\Sigma^{(2)})
 \le
 {2s_*c_D^2b^{2D-2}C_v^2\over\Delta}
 {g_{\rm YM}(a)^2\over a}.                            \label{eq:s7v40-Sigma-g}
\end{equation}
\end{theorem}

\begin{proof}
Equation \eqref{eq:s7v40-R0-bound} gives
\[
 \|A_{pq}\|
 \le {a\over\Delta}\|v_p\|\|v_q\|.
\]
Fix an anchor block \(C\).  Split the sum in
\eqref{eq:s7v40-J2} according to whether \(C\in{\rm supp}(p)\) or
\(C\in{\rm supp}(q)\), overcounting pairs that satisfy both.  In the first
case,
\begin{align*}
 \sum_{p:\,C\in{\rm supp}(p)}
 \|v_p\|
 \sum_{q:\,{\rm supp}(q)\cap{\rm supp}(p)\ne\varnothing}
 \|v_q\|
 &\le
 \sum_{p:\,C\in{\rm supp}(p)}
 \|v_p\|
 \sum_{C'\in{\rm supp}(p)}
 \sum_{q:\,C'\in{\rm supp}(q)}\|v_q\|\\
 &\le s_*{\cal J}(V)^2.
\end{align*}
The second case gives the same bound.  Multiply by \(a/\Delta\), proving
\eqref{eq:s7v40-Sigma-bound}.  Insert
\eqref{eq:s7v40-J-bound} for the last assertion.
\end{proof}

\begin{corollary}[The owner rule has the correct renormalization size]
\label{cor:s7v40-owner-size}
At fixed block side \(b\), the anchored open-mixing scale is
\[
 a{\cal J}(V)=O(g_{\rm YM}(a)),
\]
and the anchored closed-return scale is
\[
 {\cal J}_2(\Sigma^{(2)})=O(g_{\rm YM}(a)^2/a),
\]
with constants independent of the number of blocks.
\end{corollary}

\begin{proof}
Use \eqref{eq:s7v40-J-bound} and
\eqref{eq:s7v40-Sigma-g}.
\end{proof}

\section{What the overlap rule does not yet prove}
\label{sec:s7v40-limits}

\begin{firewall}[Labelled incidence is not yet a local effective operator]
\label{fire:s7v40-labelled}
The resolvent \(R_0\) preserves the excitation set, which proves the exact
overlap zero.  If the retained \(P_C\) band has nonconstant internal
energies, \(R_0\) can still depend on all retained block energies through its
denominator.  Therefore the label \( {\rm supp}(p)\cup{\rm supp}(q)\) does
not by itself prove that \(A_{pq}\) acts as the identity outside that union.
A local effective interaction requires a time-kernel representation,
operator-valued energy denominators, or a connected resolvent expansion.
\end{firewall}

\begin{firewall}[Second order is not the full linked-cluster theorem]
\label{fire:s7v40-higher}
At higher order, disconnected collections can make separate low-to-low
returns and contribute products.  They disappear from a normalized
connected generator only after taking a logarithm or cumulant.  V40 proves
the exact second-order overlap rule and anchored bound; it does not infer the
all-orders linked-cluster expansion from that rule.
\end{firewall}

\begin{definition}[Block-owned Feshbach card, BOFC]
\label{def:s7v40-card}
BOFC consists of:
\begin{enumerate}[label=\textup{(B\arabic*)},leftmargin=4em]
\item a disjoint link partition and exact Gauss projection;
\item buffered gauge-covariant slow projectors \(P_C\);
\item a uniform block residual floor \(\Delta/a\);
\item unique owners for every cross-block plaquette;
\item the anchored open-interaction bound
      \({\cal J}(V)=O(g_{\rm YM}/a)\);
\item a connected time-kernel or cumulant expansion extending
      \Cref{thm:s7v40-overlap} to all orders;
\item local self-energy counterterms of anchored size
      \(O(g_{\rm YM}^2/a)\);
\item a quasi-local effective transfer on the tensor slow space; and
\item physical Gauss gluing of the retained boundary and toron channels.
\end{enumerate}
\end{definition}

\begin{theorem}[What V40 proves]
\label{thm:s7v40-result}
BOFC items (B1)--(B4) imply an exact volume-independent residual floor.
Under the local oscillator interaction estimate (B5), the open mixing is
\(O(g_{\rm YM})\) per physical slab and the complete second-order closed
return is a sum over overlapping block supports with anchored size
\(O(g_{\rm YM}^2/a)\).  Gauge restriction preserves all these inequalities.
\end{theorem}

\begin{proof}
Use
\Cref{lem:s7v40-gauge,thm:s7v40-tensor-gap,
lem:s7v40-incidence,thm:s7v40-overlap,
thm:s7v40-Sigma-bound,cor:s7v40-owner-size}.
\end{proof}

\begin{assessment}[Progress at seventh-series V40]
\label{ass:s7v40-status}
The fixed-block IMS construction is now embedded in a real-space
renormalization architecture that does not sum a global operator norm over
the volume.  Every block retains its exact and near-zero slow modes; the
tensor complement has the same \(\Delta/a\) gap as one block.  Boundary
plaquettes have bounded anchored incidence, and their second-order
self-energy is automatically overlap-connected.  The next bottleneck is the
order-one quadratic coupling across block faces.  It must be incorporated
into a globally matched Gaussian reference, or controlled by a stable
overlapping-domain decomposition, before the all-orders normalized connected
expansion and the locality of the energy-dependent denominator can be
addressed.
\end{assessment}

\begin{decision}[V41 acquisition order]
\label{dec:s7v40-V41}
V41 will first compute the cross-block quadratic magnetic form and prove that
it is \(O(1/a)\), not \(O(g_{\rm YM}/a)\), after the V38 dilation.  It will
then test two upstream repairs: a globally exact Gaussian spectral split and
an overlapping block-domain decomposition with a stable frame bound.  Only
after the full quadratic form is in the reference operator may the
interaction picture use \(g_{\rm YM}(a)\) as its connected expansion
parameter.
\end{decision}

\begin{firewall}[Claim boundary after seventh-series V40]
\label{fire:s7v40-claim}
V40 performs the compulsory companion audit and proves the exact link-block
factorization, gauge-compatible tensor residual floor, unique-owner
incidence bound, second-order overlap rule, and volume-uniform self-energy
estimate.

V40 does not prove the local oscillator interaction hypothesis, the
all-orders linked-cluster expansion, locality of the full Feshbach map,
coarse slow-transfer matching, a terminal nonperturbative gap, cutoff
removal, or the full Yang--Mills mass gap.
\end{firewall}

\paragraph{Parent-version record.}
V40 was formed from
\texttt{Renewal\_Yang\_Mills\_Mass\_Gap\_Research\_Notes\_V39.tex}.
The V39 parent had \(16\,267\) logical source lines, \(591\,378\) bytes,
and SHA--256
\[
 \texttt{B48BF6ABFC5797A14EC060B674BA0AB852C9BED83BEAA0E1A5693939D2007867}.
\]
The next compulsory companion audit is V45.

% =============================================================================
% END APPENDED SEVENTH-SERIES V40 MATERIAL
% =============================================================================

% =============================================================================
% BEGIN APPENDED SEVENTH-SERIES V41 MATERIAL
% =============================================================================

\chapter{The cross-block Gaussian obstruction and the exact reference split}
\label{ch:s7v41}

V40 gives a conditional block-owned Feshbach estimate, but its final audit
identifies an unpaid term: a plaquette crossing a block face has a quadratic
magnetic coupling of order \(1/a\) after weak-coupling dilation.  This note
computes that term exactly.  It proves that a disjoint tensor block sum cannot
simultaneously be the exact Gaussian reference and have only
\(O(g_{\rm YM}/a)\) boundary interactions.

Two repairs are then separated.  An overlapping curvature frame represents
the full quadratic magnetic form locally and exactly, but loses tensor
factorization.  A global Gaussian spectral split keeps the full quadratic
operator exact and gives the required \(\Delta/a\) residual gap, but its
spectral projections are nonlocal and must use the V34--V36 buffered
quasi-local calculus.  The programme selects the second route for the
Feshbach split and retains real-space owners only for the nonlinear Wilson
remainder.

\section{Exact cross-block quadratic form}
\label{sec:s7v41-cross}

Let \(y=(y_\ell)_{\ell\in E}\) be a Lie-algebra-valued one-cochain, and let
\(d_1y\) be its oriented plaquette coboundary:
\begin{equation}
 (d_1y)_p
 :=
 \sum_{\ell\in\partial p}\sigma_{p\ell}y_\ell.
 \label{eq:s7v41-d}
\end{equation}
The weak-coupling quadratic magnetic form after \(q=gy\) is
\begin{equation}
 {\cal Q}_2(y)
 :=
 {1\over2a}\sum_p\|(d_1y)_p\|_{\mathfrak g}^2.         \label{eq:s7v41-Q2}
\end{equation}

For the disjoint link partition of V40, define
\begin{equation}
 b_{p,C}(y)
 :=
 \sum_{\ell\in\partial p\cap E_C}
 \sigma_{p\ell}y_\ell.                                 \label{eq:s7v41-bpC}
\end{equation}
Then
\begin{equation}
 (d_1y)_p=\sum_{C\in{\rm supp}_{\cal C}(p)}b_{p,C}(y).
 \label{eq:s7v41-sum-b}
\end{equation}

\begin{theorem}[Exact boundary-plaquette decomposition]
\label{thm:s7v41-decomposition}
For every plaquette,
\begin{align}
 {1\over2a}\|(d_1y)_p\|^2
 ={}&
 {1\over2a}\sum_C\|b_{p,C}(y)\|^2                      \notag\\
 &+
 {1\over a}\sum_{C<D}
 \langle b_{p,C}(y),b_{p,D}(y)\rangle.                 \label{eq:s7v41-decomposition}
\end{align}
If \(p\) crosses a block face, the second line is a nonzero quadratic
coupling of order \(1/a\), independent of \(g\).
\end{theorem}

\begin{proof}
Expand the squared norm of the sum in
\eqref{eq:s7v41-sum-b}.  The dilation \(q=gy\) has already cancelled the
factor \(g^{-2}\) in the magnetic Hamiltonian, so no running-coupling factor
multiplies either line.
\end{proof}

\begin{example}[The cross term is not relatively small]
\label{ex:s7v41-sharp}
Suppose a plaquette meets exactly two link blocks and choose a cochain for
which
\[
 b_{p,C}=v,\qquad b_{p,D}=v\ne0.
\]
Then the diagonal contribution in
\eqref{eq:s7v41-decomposition} is \(\|v\|^2/a\), and the cross contribution
is also \(\|v\|^2/a\).  If instead \(b_{p,D}=-v\), the full plaquette
curvature vanishes although the diagonal block contribution remains
\(\|v\|^2/a\).
\end{example}

\begin{proof}
Substitute the two choices into
\eqref{eq:s7v41-decomposition}.
\end{proof}

\begin{corollary}[No small perturbation of a disjoint quadratic block sum]
\label{cor:s7v41-no-small}
There is no constant \(\eta(g)\to0\) such that every cross-block quadratic
form obeys
\begin{equation}
 |{\cal Q}_{\partial,{\rm cross}}(y)|
 \le
 \eta(g){\cal Q}_{\partial,{\rm diag}}(y)              \label{eq:s7v41-false-relative}
\end{equation}
uniformly in \(y\).  Hence the V40 estimate
\(\|v_p\|=O(g/a)\) cannot include the quadratic part of a boundary
plaquette.
\end{corollary}

\begin{proof}
The first configuration in \Cref{ex:s7v41-sharp} makes the ratio of the two
sides before \(\eta\) equal to one for every \(g\).
\end{proof}

\begin{theorem}[Tensor locality versus exact Gaussian locality]
\label{thm:s7v41-tensor-no-go}
Assume a link partition has at least one plaquette meeting two blocks.  A
reference Hamiltonian of the disjoint tensor form
\[
 H_0=\sum_Ch_C
\]
cannot contain the exact quadratic plaquette form
\eqref{eq:s7v41-Q2} while every \(h_C\) acts only on \(E_C\).
The omitted boundary operator has coefficient \(O(1/a)\), not
\(O(g/a)\).
\end{theorem}

\begin{proof}
A sum of operators acting on disjoint link factors has a quadratic form with
no bilinear monomial involving a link in \(E_C\) and a link in \(E_D\),
\(C\ne D\).  The second line of
\eqref{eq:s7v41-decomposition} contains precisely such monomials for a
crossing plaquette.  Their coefficient is \(1/a\).
\end{proof}

\begin{firewall}[Correction to the block-owned card]
\label{fire:s7v41-correction}
The V40 tensor residual theorem remains correct for a declared disjoint block
sum.  Its application to the weak Yang--Mills Hamiltonian is incomplete
because the exact Gaussian boundary curvature is missing from that sum.
Real-space owner combinatorics may be applied to the nonlinear remainder
only after the complete quadratic boundary form has been put into the
reference dynamics.
\end{firewall}

\section{Repair A: an exact overlapping curvature frame}
\label{sec:s7v41-frame}

Let \(\{{\mathscr O}_\alpha\}\) be any finite-overlap family of plaquette
patches.  Choose weights \(\omega_{\alpha p}\) such that
\begin{equation}
 \omega_{\alpha p}=0\quad(p\notin{\mathscr O}_\alpha),
 \qquad
 \sum_\alpha\omega_{\alpha p}^2=1
 \quad(p\in{\cal P}).                                  \label{eq:s7v41-weights}
\end{equation}
Define the local curvature analysis maps
\begin{equation}
 (B_\alpha y)_p:=\omega_{\alpha p}(d_1y)_p.            \label{eq:s7v41-Balpha}
\end{equation}

\begin{theorem}[Exact local curvature tight frame]
\label{thm:s7v41-frame}
For every one-cochain \(y\),
\begin{equation}
 \sum_\alpha\|B_\alpha y\|_{\ell^2({\cal P})}^2
 =
 \|d_1y\|_{\ell^2({\cal P})}^2.                        \label{eq:s7v41-frame}
\end{equation}
Thus the complete quadratic magnetic energy is an exact sum of overlapping
local curvature energies, with neither a boundary cross error nor a volume
dependent frame constant.
\end{theorem}

\begin{proof}
Interchange the finite sums and use
\eqref{eq:s7v41-weights}:
\[
 \sum_\alpha\sum_p\omega_{\alpha p}^2\|(d_1y)_p\|^2
 =
 \sum_p\|(d_1y)_p\|^2.
\]
\end{proof}

\begin{firewall}[The tight curvature frame is not a tensor decomposition]
\label{fire:s7v41-frame}
Different \(B_\alpha\) use the same link coordinates.  Their outputs obey
consistency constraints and do not define independent Hilbert factors.
Consequently \Cref{thm:s7v41-frame} repairs the quadratic energy ledger but
does not supply the tensor-product residual projection used in V40.
\end{firewall}

\section{Repair B: keep the global Gaussian operator exact}
\label{sec:s7v41-global}

Let \({\cal A}_{\rm T}\) be the finite-volume transverse one-cochain space
after exact gauge reduction, with harmonic torons removed to a separate slow
factor.  Define
\begin{equation}
 H_{2,a}
 :=
 {1\over2a}
 \left(
 -\Delta_{{\cal A}_{\rm T}}
 +\langle y,d_1^*d_1y\rangle
 \right).                                              \label{eq:s7v41-H2}
\end{equation}
This operator contains every quadratic plaquette term, including those
crossing any real-space block boundary.

Let
\[
 P_{\rm L},P_{\rm T},P_{\rm H}
\]
be the exact Gaussian spectral bands of V29 and V34, now defined for the
complete operator \eqref{eq:s7v41-H2}.

\begin{theorem}[Exact global Gaussian residual floor]
\label{thm:s7v41-global-gap}
On the Fock sector containing at least one transition or high one-particle
quantum with dimensionless curvature eigenvalue at least \(\alpha>0\),
\begin{equation}
 H_{2,a}-E_{2,0}
 \ge{\sqrt\alpha\over a}
 (P_{\rm T}+P_{\rm H}).                                \label{eq:s7v41-global-gap}
\end{equation}
The Gaussian transfer commutes with all three spectral projections, so all
free inter-band blocks vanish exactly.
\end{theorem}

\begin{proof}
Diagonalize the positive transverse matrix \(d_1^*d_1\).  V38 gives one
oscillator frequency \(\sqrt\lambda/a\) for a mode of curvature eigenvalue
\(\lambda\).  Any state in the declared residual sector carries at least one
quantum with \(\lambda\ge\alpha\), proving the lower bound.  The transfer and
the spectral projections are Borel functions of the same diagonal number
operators, so they commute.
\end{proof}

\begin{corollary}[Correct division of labour]
\label{cor:s7v41-division}
The exact Gaussian operator \eqref{eq:s7v41-H2} supplies the residual energy
denominator and free three-band split.  The V34 analytic filter and V36
incidence schedule supply quasi-local approximations for separated-band
observables.  Real-space owners are applied only to
\begin{equation}
 W_{\rm nl}(g)
 :=
 H_{\rm Wilson}(g)-H_{2,a}-E_{\rm ct}(g),              \label{eq:s7v41-Wnl}
\end{equation}
where \(E_{\rm ct}\) includes the declared scalar and quadratic
counterterms.
\end{corollary}

\begin{proof}
\Cref{thm:s7v41-global-gap} gives the exact free floor.  V34--V36 apply to
functions of the same global quadratic operator.  Subtracting the complete
quadratic Taylor part from the Wilson Hamiltonian leaves only nonlinear and
counterterm remainders for the owner ledger.
\end{proof}

\begin{firewall}[Global exactness reintroduces quasi-locality]
\label{fire:s7v41-global}
The exact spectral projections of \(H_{2,a}\) are not finite range.  The
choice in \Cref{cor:s7v41-division} abandons the false tensor factorization,
not the locality requirement.  The quasi-local filters, transition channel,
and toron slow factor remain compulsory.
\end{firewall}

\section{Local Taylor order of the nonlinear remainder}
\label{sec:s7v41-Taylor}

For one plaquette, put
\begin{equation}
 U_p(g,y)
 :=
 \prod_{\ell\in\partial p}^{\longrightarrow}
 \exp\bigl(\sigma_{p\ell}g y_\ell\bigr).               \label{eq:s7v41-Up}
\end{equation}
Let
\[
 w_p(g,y)=1-{1\over N}\operatorname{Re}\operatorname{Tr}U_p(g,y).
\]

\begin{lemma}[Uniform plaquette Taylor remainder]
\label{lem:s7v41-Taylor}
For every fixed \(\rho>0\), there are \(C_\rho<\infty\) and a positive
quadratic form \(Q_p(y)\), equal to the abelianized curvature form up to the
chosen Lie-algebra normalization, such that
\begin{equation}
 \left|
 {1\over g^2}w_p(g,y)-Q_p(y)
 \right|
 \le
 C_\rho g(1+\|y\|)^3                                  \label{eq:s7v41-Taylor}
\end{equation}
whenever \(0<g\le1\) and \(g\max_{\ell\in\partial p}|y_\ell|\le\rho\).
\end{lemma}

\begin{proof}
The finite product of matrix exponentials is analytic in the variables
\(g y_\ell\).  The Wilson function has zero constant and linear terms at the
identity.  Its Hessian is the positive quadratic form of the oriented sum
\((d_1y)_p\).  Taylor's theorem with integral remainder bounds the cubic and
higher terms by \(C_\rho g^3(1+\|y\|)^3\) on the declared compact coordinate
set.  Divide by \(g^2\).
\end{proof}

\begin{corollary}[The nonlinear owner coefficient starts at \(g/a\)]
\label{cor:s7v41-owner-coefficient}
On the V39 weak-field region, the nonlinear magnetic remainder of one
plaquette satisfies the multiplication bound
\begin{equation}
 \left|
 {1\over ag^2}w_p(g,y)-{1\over a}Q_p(y)
 \right|
 \le
 {C_\rho g\over a}(1+\|y\|)^3.                         \label{eq:s7v41-owner-coefficient}
\end{equation}
Thus the coefficient, but not the unbounded polynomial field norm, has the
V40 scale \(g/a\).
\end{corollary}

\begin{proof}
Divide \eqref{eq:s7v41-Taylor} by \(a\).
\end{proof}

\begin{firewall}[Coefficient smallness is not operator-norm smallness]
\label{fire:s7v41-unbounded}
The cubic polynomial in \eqref{eq:s7v41-owner-coefficient} is unbounded on
oscillator Hilbert space.  Therefore that equation does not prove
\(\|v_p\|\le Cg/a\) globally.  One needs an oscillator-number weighted norm,
a spectral occupation cutoff with a tail estimate, or a relative-form
bound.
\end{firewall}

\section{Oscillator-number weighted bound}
\label{sec:s7v41-number}

On a fixed finite oscillator block, let
\[
 {\cal N}=\sum_{j=1}^n a_j^*a_j
\]
be the number operator and let
\begin{equation}
 P_M:=\mathbf1_{[0,M]}({\cal N}).                      \label{eq:s7v41-PM}
\end{equation}

\begin{lemma}[Polynomial fields on a number window]
\label{lem:s7v41-number}
For every polynomial \(P_r(y)\) of total degree at most \(r\), there is
\(C_{r,n}\) such that
\begin{equation}
 \|P_r(y)P_M\|
 \le C_{r,n}(M+1)^{r/2}.                               \label{eq:s7v41-number}
\end{equation}
\end{lemma}

\begin{proof}
Each coordinate is a fixed linear combination of \(a_j+a_j^*\).  The
standard relations give
\[
 \|a_j({\cal N}+1)^{-1/2}\|\le1,\qquad
 \|a_j^*({\cal N}+1)^{-1/2}\|\le1.
\]
Move one factor \(({\cal N}+1)^{1/2}\) to the right for each coordinate,
using
\(({\cal N}+1)^{s}a_j^*=a_j^*({\cal N}+2)^{s}\) and the analogous lowering
identity.  A monomial of degree \(r\) is bounded by
\(C_{r,n}({\cal N}+r+1)^{r/2}\).  On
\(\operatorname{Ran}P_M\), this is at most a constant times
\((M+1)^{r/2}\).  Sum the finitely many monomials.
\end{proof}

\begin{theorem}[Nonlinear interaction on a growing occupation window]
\label{thm:s7v41-number-bound}
Suppose the weak-field Hamiltonian remainder through quartic order is
\[
 W_{\rm nl}(g)
 ={1\over a}\bigl(gH_3+g^2H_4+{\cal R}_{\ge5}\bigr)
\]
on a fixed block, with polynomial coefficients of the indicated degrees and
an analytic higher remainder on the V39 chart.  Then
\begin{equation}
 \|P_MW_{\rm nl}(g)P_M\|
 \le
 {C\over a}
 \left[
 g(M+1)^{3/2}
 +g^2(M+1)^2
 +{\cal E}_{\ge5}(g,M)
 \right],                                              \label{eq:s7v41-number-bound}
\end{equation}
where
\({\cal E}_{\ge5}(g,M)=O(g^3(M+1)^{5/2})\) while
\(g\sqrt{M+1}\) remains in a fixed compact interval.
\end{theorem}

\begin{proof}
Apply \Cref{lem:s7v41-number} term by term.  Analyticity makes the
degree-\(r\) coefficient of the post-quadratic expansion \(O(g^{r-2})\).
The resulting majorant series is bounded by a constant times its first
degree-five term when \(g\sqrt{M+1}\) lies in a sufficiently small fixed
interval.
\end{proof}

\begin{corollary}[Explicit one-loop window]
\label{cor:s7v41-window}
Assume
\[
 g(a)=O(L(a)^{-1/2}),\qquad
 M(a)=L(a)^{1/8}.
\]
Then
\begin{align}
 g(a)M(a)^{3/2}&=O(L(a)^{-5/16}),                     \label{eq:s7v41-cubic-rate}\\
 g(a)^2M(a)^2&=O(L(a)^{-3/4}),                        \label{eq:s7v41-quartic-rate}\\
 g(a)\sqrt{M(a)}&=O(L(a)^{-7/16}).                    \label{eq:s7v41-chart-rate}
\end{align}
Thus the dimensionless nonlinear interaction on the declared number window
tends to zero.
\end{corollary}

\begin{proof}
Use
\(-1/2+3/16=-5/16\),
\(-1+1/4=-3/4\), and
\(-1/2+1/16=-7/16\).
\end{proof}

\section{Corrected global-to-local card}
\label{sec:s7v41-card}

\begin{definition}[Gaussian-first Feshbach card, GFFC]
\label{def:s7v41-card}
GFFC records:
\begin{enumerate}[label=\textup{(G\arabic*)},leftmargin=4em]
\item exact gauge and toron separation;
\item the complete global quadratic operator \(H_{2,a}\), including every
      cross-block plaquette;
\item its exact low/transition/high spectral projections and residual floor
      \(\Delta/a\);
\item the V34--V36 quasi-local filter schedule;
\item a V39 gauge-invariant weak-field cutoff;
\item a growing oscillator-number window \(P_{M(a)}\);
\item local owners only for the nonlinear remainder \(W_{\rm nl}\);
\item an anchored weighted bound with dimensionless activity
      \(O(gM^{3/2}+g^2M^2)\);
\item a high-occupation tail or Feshbach bound; and
\item coarse matching of the complete self-energy-corrected low operator.
\end{enumerate}
\end{definition}

\begin{theorem}[What V41 establishes]
\label{thm:s7v41-result}
The disjoint-block Gaussian reference is obstructed by an exact
coupling-independent cross form.  GFFC items (G1)--(G3) remove that
obstruction by retaining the complete global Gaussian operator.  On the
weak-field number window \(M=L^{1/8}\), the owned nonlinear coefficient has
dimensionless size \(O(L^{-5/16})\) at cubic order and
\(O(L^{-3/4})\) at quartic order.
\end{theorem}

\begin{proof}
The obstruction is
\Cref{thm:s7v41-decomposition,cor:s7v41-no-small,
thm:s7v41-tensor-no-go}.  The corrected residual floor is
\Cref{thm:s7v41-global-gap}.  The nonlinear coefficient and number-window
bounds are
\Cref{cor:s7v41-owner-coefficient,thm:s7v41-number-bound,
cor:s7v41-window}.
\end{proof}

\begin{assessment}[Progress at seventh-series V41]
\label{ass:s7v41-status}
The last false small parameter has been removed.  Cross-block quadratic
curvature is part of the exact Gaussian reference and is never charged to the
asymptotically free interaction.  After that correction, the nonlinear
Wilson remainder really begins at \(g/a\); on an explicit growing oscillator
window its dimensionless anchored size tends to zero.  The new bottleneck is
the complement of that number window and the construction of a connected,
quasi-local low effective transfer for the global Gaussian split.
\end{assessment}

\begin{decision}[V42 acquisition order]
\label{dec:s7v41-V42}
V42 will control the oscillator-number complement.  It will combine the
number-operator gap \(M/a\), the weak-field IMS barrier, and a Feshbach
estimate for polynomial creation/annihilation vertices.  The target is a
two-cutoff theorem in which both high momentum and high occupation are
residual, while the nonlinear low-window activity is the vanishing quantity
in \Cref{cor:s7v41-window}.
\end{decision}

\begin{firewall}[Claim boundary after seventh-series V41]
\label{fire:s7v41-claim}
V41 proves the exact cross-block quadratic obstruction, the overlapping
curvature-frame identity, the corrected global Gaussian residual floor, the
local Wilson Taylor order, and the oscillator-number window bounds.

V41 does not prove the high-occupation Feshbach estimate, nonlinear
connected-kernel convergence, self-energy coarse matching, a terminal
nonperturbative gap, cutoff removal, or the full Yang--Mills mass gap.
\end{firewall}

\paragraph{Parent-version record.}
V41 was formed from
\texttt{Renewal\_Yang\_Mills\_Mass_Gap\_Research\_Notes\_V40.tex}.
The V40 parent had \(16\,807\) logical source lines, \(612\,111\) bytes,
and SHA--256
\[
 \texttt{AB9EFE6CE182A03F4104B26D343E413ED18D96112A445D6DA0FC3AB451C2EEC8}.
\]
The next compulsory companion audit is V45.

% =============================================================================
% END APPENDED SEVENTH-SERIES V41 MATERIAL
% =============================================================================

% =============================================================================
% BEGIN APPENDED SEVENTH-SERIES V42 MATERIAL
% =============================================================================

\chapter{High-occupation elimination and nested Feshbach consistency}
\label{ch:s7v42}

V41 puts the complete quadratic lattice gauge operator into the exact
Gaussian reference and proves that the nonlinear Wilson remainder is small
on a growing oscillator-number window.  This note controls the complement of
that window under an explicit relative-form hypothesis.  The key gain is an
additional occupation-energy denominator.  A cubic vertex has cutoff norm
\(O(gM^{3/2}/a)\), but the sector \(N>M\) begins at energy \(O(M/a)\).
Therefore
\[
 \hbox{high-occupation mixing}=O(g\sqrt M),\qquad
 \hbox{self-energy}=O(g^2M^2/a).
\]
For \(g=O(L^{-1/2})\) and \(M=L^{1/8}\), the first is
\(O(L^{-7/16})\), the dimensionless self-energy coefficient is
\(O(L^{-3/4})\), and the energy dependence costs only
\(O(L^{-7/8})\).

The final section proves that high momentum and high occupation may be
eliminated successively: nested Schur complements agree with the direct
Feshbach map whenever all displayed inverses exist.

\section{Oscillator number and polynomial vertices}
\label{sec:s7v42-number}

Let
\begin{equation}
 H_{0,a}
 :=
 {1\over a}\sum_{j=1}^n\omega_jN_j,\qquad
 \omega_j\ge\omega_*>0,                               \label{eq:s7v42-H0}
\end{equation}
where the Gaussian ground energy has been subtracted.  Put
\begin{equation}
 {\cal N}:=\sum_{j=1}^nN_j,\qquad
 P_M:=\mathbf1_{[0,M]}({\cal N}),\qquad
 Q_M:=I-P_M.                                           \label{eq:s7v42-PQ}
\end{equation}
Then
\begin{equation}
 Q_MH_{0,a}Q_M
 \ge{\omega_*(M+1)\over a}Q_M.                        \label{eq:s7v42-free-floor}
\end{equation}

Let \(H_r\) be a polynomial of total degree at most \(r\) in the creation and
annihilation operators of this fixed oscillator family.

\begin{lemma}[Window-to-complement vertex bound]
\label{lem:s7v42-vertex}
There is \(C_{r,n,\omega}<\infty\) such that
\begin{equation}
 \|Q_MH_rP_M\|
 \le C_{r,n,\omega}(M+r+1)^{r/2}.                     \label{eq:s7v42-vertex}
\end{equation}
Moreover \(Q_MH_rP_M\) has range in the number shell
\begin{equation}
 M<{\cal N}\le M+r.                                    \label{eq:s7v42-shell}
\end{equation}
\end{lemma}

\begin{proof}
Every monomial of degree \(r\) changes total occupation by at most \(r\).
This proves \eqref{eq:s7v42-shell}.  The creation/annihilation estimate used
in V41 bounds a degree-\(r\) monomial on
\(\operatorname{Ran}P_M\) by a constant times
\((M+r+1)^{r/2}\).  Sum the finitely many monomials.
\end{proof}

Consider
\begin{equation}
 H_a
 =
 H_{0,a}+W_a,\qquad
 W_a={1\over a}\bigl(gH_3+g^2H_4+{\cal R}_a\bigr).
 \label{eq:s7v42-H}
\end{equation}
Define the dimensionless window vertex
\begin{equation}
 k_M(g)
 :=
 C_3g(M+4)^{3/2}
 +C_4g^2(M+5)^2+\rho_a(M),                            \label{eq:s7v42-k}
\end{equation}
where
\begin{equation}
 \|Q_M{\cal R}_aP_M\|\le\rho_a(M).                    \label{eq:s7v42-rho}
\end{equation}

\begin{corollary}[Off-diagonal occupation block]
\label{cor:s7v42-K}
The block
\begin{equation}
 K_M:=P_MH_aQ_M=P_MW_aQ_M                             \label{eq:s7v42-K}
\end{equation}
obeys
\begin{equation}
 \|K_M\|\le{k_M(g)\over a}.                            \label{eq:s7v42-K-bound}
\end{equation}
\end{corollary}

\begin{proof}
The quadratic number operator commutes with \(P_M\).  Apply
\Cref{lem:s7v42-vertex} to \(H_3,H_4\), use
\eqref{eq:s7v42-rho}, and take adjoints.
\end{proof}

\section{A full residual floor from a relative-form estimate}
\label{sec:s7v42-floor}

The nonlinear Hamiltonian inside the weak-field chart is unbounded, so the
free floor \eqref{eq:s7v42-free-floor} must be transferred by a form
estimate, not an operator-norm estimate.

\begin{assumption}[Weak-field relative-form card]
\label{ass:s7v42-form}
On the gauge-fixed V39 interior form domain,
\begin{equation}
 |\langle\psi,W_a\psi\rangle|
 \le
 \eta_a\langle\psi,H_{0,a}\psi\rangle
 +{c_a\over a}\|\psi\|^2,                              \label{eq:s7v42-form}
\end{equation}
where
\begin{equation}
 0\le\eta_a<1.                                         \label{eq:s7v42-eta}
\end{equation}
\end{assumption}

\begin{theorem}[Interacting high-occupation floor]
\label{thm:s7v42-floor}
Under \Cref{ass:s7v42-form},
\begin{equation}
 Q_MH_aQ_M
 \ge{\Delta_M(a)\over a}Q_M,                           \label{eq:s7v42-floor}
\end{equation}
where
\begin{equation}
 \Delta_M(a)
 :=
 (1-\eta_a)\omega_*(M+1)-c_a.                         \label{eq:s7v42-DeltaM}
\end{equation}
In particular, if \(\eta_a\to0\), \(c_a=o(M)\), and \(M\to\infty\), then
\begin{equation}
 \Delta_M(a)\ge{\omega_*M\over2}                       \label{eq:s7v42-half-floor}
\end{equation}
for all sufficiently small \(a\).
\end{theorem}

\begin{proof}
For \(\psi\in\operatorname{Ran}Q_M\), use the lower side of
\eqref{eq:s7v42-form} and
\eqref{eq:s7v42-free-floor}:
\begin{align*}
 \langle\psi,H_a\psi\rangle
 &\ge
 (1-\eta_a)\langle\psi,H_{0,a}\psi\rangle
 -{c_a\over a}\|\psi\|^2\\
 &\ge
 {[(1-\eta_a)\omega_*(M+1)-c_a]\over a}\|\psi\|^2.
\end{align*}
This is \eqref{eq:s7v42-floor}.  The asymptotic hypotheses make the bracket
at least \(\omega_*M/2\).
\end{proof}

\begin{remark}[Input supplied and not supplied by V39]
\label{rem:s7v42-V39}
V39 proves analytic coefficient control on a fixed block and identifies the
relative small parameter \(gR\), while its IMS formula controls the chart
exterior relative to the ultraviolet scale.  A complete proof of
\eqref{eq:s7v42-form} for the gauge-fixed Wilson Hamiltonian must still
include the Haar kinetic coefficients, the shrinking cutoff derivatives,
torons, and the common vacuum-energy normalization.  It is therefore stated
as an assumption here and retained in the final claim boundary.
\end{remark}

\section{High-occupation Feshbach estimates}
\label{sec:s7v42-Feshbach}

Fix a real spectral parameter \(z\) with
\begin{equation}
 az<\Delta_M(a),\qquad
 \delta_M(z):=\Delta_M(a)-az>0.                        \label{eq:s7v42-delta}
\end{equation}
Let
\begin{align}
 R_M(z)&:=
 \bigl(Q_M(H_a-z)Q_M\bigr)^{-1},                       \label{eq:s7v42-R}\\
 \Sigma_M(z)&:=K_MR_M(z)K_M^*.                         \label{eq:s7v42-Sigma}
\end{align}

\begin{theorem}[Occupation mixing and self-energy ledger]
\label{thm:s7v42-ledger}
Under the hypotheses above,
\begin{align}
 \|R_M(z)\|
 &\le{a\over\delta_M(z)},                              \label{eq:s7v42-R-bound}\\
 \|R_M(z)K_M^*\|
 &\le{k_M(g)\over\delta_M(z)},                         \label{eq:s7v42-mixing}\\
 \|\Sigma_M(z)\|
 &\le{k_M(g)^2\over a\delta_M(z)},                     \label{eq:s7v42-Sigma-bound}\\
 \|\Sigma_M'(z)\|
 &\le{k_M(g)^2\over\delta_M(z)^2}.                     \label{eq:s7v42-Sigma-prime}
\end{align}
\end{theorem}

\begin{proof}
\Cref{thm:s7v42-floor} gives the resolvent estimate
\eqref{eq:s7v42-R-bound}.  Combine it with
\eqref{eq:s7v42-K-bound} to obtain
\eqref{eq:s7v42-mixing} and
\eqref{eq:s7v42-Sigma-bound}.  Finally
\[
 R_M'(z)=R_M(z)^2,\qquad
 \Sigma_M'(z)=K_MR_M(z)^2K_M^*,
\]
which gives \eqref{eq:s7v42-Sigma-prime}.
\end{proof}

\begin{corollary}[Leading dimensional orders]
\label{cor:s7v42-orders}
Assume \eqref{eq:s7v42-half-floor}, and suppose the cubic term dominates
\(k_M(g)\).  Uniformly for \(az=o(M)\),
\begin{align}
 \|R_M(z)K_M^*\|
 &=O(gM^{1/2}),                                        \label{eq:s7v42-mixing-order}\\
 \|\Sigma_M(z)\|
 &=O(g^2M^2/a),                                        \label{eq:s7v42-Sigma-order}\\
 \|\Sigma_M'(z)\|
 &=O(g^2M).                                             \label{eq:s7v42-derivative-order}
\end{align}
\end{corollary}

\begin{proof}
Use \(k_M(g)=O(gM^{3/2})\) and
\(\delta_M(z)\asymp M\) in
\Cref{thm:s7v42-ledger}.
\end{proof}

\begin{corollary}[Explicit one-loop rates]
\label{cor:s7v42-rates}
Let
\begin{equation}
 g(a)=O(L(a)^{-1/2}),\qquad M(a)=L(a)^{1/8},           \label{eq:s7v42-schedule}
\end{equation}
and assume
\begin{equation}
 \rho_a(M)=o(gM^{3/2}),\quad
 \eta_a\to0,\quad c_a=o(M).                            \label{eq:s7v42-remainders}
\end{equation}
Then, for bounded physical \(z\),
\begin{align}
 \|R_M(z)K_M^*\|
 &=O(L(a)^{-7/16}),                                    \label{eq:s7v42-rate-mixing}\\
 \|\Sigma_M(z)\|
 &=O\!\left({L(a)^{-3/4}\over a}\right),               \label{eq:s7v42-rate-Sigma}\\
 \|\Sigma_M'(z)\|
 &=O(L(a)^{-7/8}).                                     \label{eq:s7v42-rate-prime}
\end{align}
\end{corollary}

\begin{proof}
The exponents are
\[
 -\frac12+\frac1{16}=-\frac7{16},\qquad
 -1+\frac14=-\frac34,\qquad
 -1+\frac18=-\frac78.
\]
Apply \Cref{cor:s7v42-orders}; quartic and higher terms are smaller under
\eqref{eq:s7v42-remainders}.
\end{proof}

\begin{firewall}[The self-energy is not discarded]
\label{fire:s7v42-Sigma}
Although the dimensionless coefficient in
\eqref{eq:s7v42-rate-Sigma} tends to zero, the physical operator scale
contains \(a^{-1}\).  The high-occupation self-energy must be included in
the matched low Hamiltonian.  Only the mixing norm
\eqref{eq:s7v42-rate-mixing} and the energy dependence
\eqref{eq:s7v42-rate-prime} are vanishing gap-transfer debits.
\end{firewall}

\section{Eigenvector tail and effective gap}
\label{sec:s7v42-gap}

Define the number-window effective operator
\begin{equation}
 F_M(z)
 :=
 P_M(H_a-z)P_M-\Sigma_M(z).                            \label{eq:s7v42-F}
\end{equation}

\begin{theorem}[High-occupation tail of a low-energy eigenvector]
\label{thm:s7v42-tail}
If \(H_a\psi=z\psi\) with \(az<\Delta_M(a)\), then
\begin{equation}
 Q_M\psi=-R_M(z)K_M^*P_M\psi                          \label{eq:s7v42-tail}
\end{equation}
and therefore
\begin{equation}
 \|Q_M\psi\|
 \le{k_M(g)\over\delta_M(z)}\|P_M\psi\|.              \label{eq:s7v42-tail-bound}
\end{equation}
Under \eqref{eq:s7v42-schedule}--\eqref{eq:s7v42-remainders}, the
high-occupation norm is \(O(L^{-7/16})\) relative to the low component.
\end{theorem}

\begin{proof}
The \(Q_M\) block of the eigenvalue equation is
\[
 Q_M(H_a-z)Q_MQ_M\psi=-K_M^*P_M\psi.
\]
Apply \(R_M(z)\) and then
\eqref{eq:s7v42-mixing}.  The rate is
\eqref{eq:s7v42-rate-mixing}.
\end{proof}

\begin{corollary}[Occupation-effective gap handoff]
\label{cor:s7v42-gap}
Suppose
\begin{equation}
 F_M(0)\ge m_MI_{P_M{\cal H}_\circ},\qquad m_M>0,      \label{eq:s7v42-effective-gap}
\end{equation}
and
\begin{equation}
 \sup_{0\le z\le z_0}\|\Sigma_M'(z)\|\le\theta_M,
 \qquad z_0<\Delta_M(a)/a.                             \label{eq:s7v42-theta}
\end{equation}
Then \(H_a\) has no spectrum in
\begin{equation}
 0<z<
 \min\left\{
 z_0,\ {m_M\over1+\theta_M}
 \right\}.                                             \label{eq:s7v42-gap}
\end{equation}
\end{corollary}

\begin{proof}
This is the V37 effective-to-full gap proof with
\({\cal H}_{\rm L}=\operatorname{Ran}P_M\) and
\({\cal H}_{\rm r}=\operatorname{Ran}Q_M\).
Explicitly,
\[
 F_M(z)
 =
 F_M(0)-zP_M-\bigl(\Sigma_M(z)-\Sigma_M(0)\bigr)
 \ge
 [m_M-(1+\theta_M)z]P_M.
\]
The high block and the Feshbach operator are invertible throughout the
displayed interval, so the full block factorization excludes spectrum there.
\end{proof}

\section{Nested elimination}
\label{sec:s7v42-nested}

Let a Hilbert space split orthogonally as
\begin{equation}
 {\cal H}={\cal H}_{\rm L}\oplus{\cal H}_{\rm M}
 \oplus{\cal H}_{\rm H},                               \label{eq:s7v42-three}
\end{equation}
where \({\cal H}_{\rm H}\) may represent high momentum and
\({\cal H}_{\rm M}\) high occupation inside the retained momentum band.

\begin{theorem}[Associativity of Schur elimination]
\label{thm:s7v42-associativity}
Let \(A\) be a block operator on \eqref{eq:s7v42-three}.  Assume the high
block \(A_{\rm HH}\), its first Schur complement on
\({\cal H}_{\rm L}\oplus{\cal H}_{\rm M}\), and every inverse below exist.
Then eliminating \({\cal H}_{\rm H}\) first and
\({\cal H}_{\rm M}\) second gives the same operator on
\({\cal H}_{\rm L}\) as eliminating
\({\cal H}_{\rm M}\oplus{\cal H}_{\rm H}\) in one step.
\end{theorem}

\begin{proof}
Write \(A\) as a \(2\times2\) block matrix relative to
\(({\cal H}_{\rm L}\oplus{\cal H}_{\rm M})\oplus{\cal H}_{\rm H}\) and
apply the exact triangular factorization of V37.  This eliminates the high
block and leaves its Schur complement \(S_{\rm LM}\).  Apply the same
factorization to \(S_{\rm LM}\) relative to
\({\cal H}_{\rm L}\oplus{\cal H}_{\rm M}\).

The product of the two lower and two upper triangular elimination matrices
is itself a triangular elimination of the combined
\({\cal H}_{\rm M}\oplus{\cal H}_{\rm H}\) block.  The remaining diagonal
low block is unique because multiplication reconstructs \(A\).  Equivalently,
both constructions equal the upper-left block of \(A^{-1}\) inverted:
\[
 S_{\rm L}=\bigl(P_{\rm L}A^{-1}P_{\rm L}\bigr)^{-1}.
\]
Hence the two Schur complements coincide.
\end{proof}

\begin{corollary}[High momentum and high occupation may be treated in order]
\label{cor:s7v42-two-cutoff}
Provided their residual inverses and buffers remain uniform, the global
Gaussian high-momentum sector of V41 may be eliminated first and the
high-occupation sector \(Q_M\) second, or conversely.  No additional
order-of-cutoffs counterterm is generated; all self-energies are those of
the common direct Feshbach map.
\end{corollary}

\begin{proof}
Apply \Cref{thm:s7v42-associativity} to the two residual sectors.
\end{proof}

\section{Two-cutoff status}
\label{sec:s7v42-status}

\begin{theorem}[What V42 proves]
\label{thm:s7v42-result}
Under the weak-field relative-form card
\eqref{eq:s7v42-form} and the analytic remainder bounds
\eqref{eq:s7v42-remainders}, the occupation complement has a residual floor
\(\asymp M/a\).  For \(g=O(L^{-1/2})\), \(M=L^{1/8}\), its mixing with the
number window is \(O(L^{-7/16})\), its self-energy has size
\(O(L^{-3/4}/a)\), and the derivative debit is \(O(L^{-7/8})\).  Nested
high-momentum and high-occupation elimination is exactly consistent.
\end{theorem}

\begin{proof}
Use
\Cref{thm:s7v42-floor,thm:s7v42-ledger,
cor:s7v42-rates,thm:s7v42-tail,
thm:s7v42-associativity}.
\end{proof}

\begin{assessment}[Progress at seventh-series V42]
\label{ass:s7v42-status}
The growing number cutoff introduced in V41 no longer leaves an untyped
complement.  Its energy denominator improves the cubic mixing from
\(gM^{3/2}\) to \(gM^{1/2}\), and the energy dependence of the resulting
self-energy vanishes rapidly.  What remains is to prove the relative-form
card for the compact gauge-fixed Wilson operator and to organize the
self-energy-corrected low-window interactions into a connected,
volume-uniform RG map.
\end{assessment}

\begin{decision}[V43 acquisition order]
\label{dec:s7v42-V43}
V43 will attack \Cref{ass:s7v42-form}.  It will derive a coordinate-invariant
relative bound for the Haar Laplacian and Wilson Taylor remainder under the
V39 plaquette cutoff, retain the IMS derivative as a residual-scale term,
and state separately the toron and chart-boundary contributions.  No
connected expansion will be attempted until this analytic hypothesis is
either proved or reduced to an explicit finite-block constant.
\end{decision}

\begin{firewall}[Claim boundary after seventh-series V42]
\label{fire:s7v42-claim}
V42 proves the polynomial window-to-complement bound, the interacting
occupation floor conditional on a relative-form estimate, the complete
occupation Feshbach ledger, the eigenvector tail, the effective-gap handoff,
and associativity of nested eliminations.

V42 does not prove the Wilson relative-form card, the analytic remainder
rate on the full gauge-fixed space, connected-kernel convergence, coarse
self-energy matching, a terminal nonperturbative gap, cutoff removal, or the
full Yang--Mills mass gap.
\end{firewall}

\paragraph{Parent-version record.}
V42 was formed from
\texttt{Renewal\_Yang\_Mills\_Mass_Gap\_Research\_Notes\_V41.tex}.
The V41 parent had \(17\,326\) logical source lines, \(630\,757\) bytes,
and SHA--256
\[
 \texttt{72E8B61385DA792CE7196CD875CC1B173D86A82D4B2614EF60F09252E67E1E93}.
\]
The next compulsory companion audit is V45.

% =============================================================================
% END APPENDED SEVENTH-SERIES V42 MATERIAL
% =============================================================================

% =============================================================================
% BEGIN APPENDED SEVENTH-SERIES V43 MATERIAL
% =============================================================================

\chapter{The fixed-block Wilson relative-form estimate}
\label{ch:s7v43}

V42 reduces high-occupation control to the weak-field relative-form card
\[
 |\langle\psi,W_a\psi\rangle|
 \le\eta_a\langle\psi,H_{0,a}\psi\rangle
 +{c_a\over a}\|\psi\|^2.
\]
This note proves that card on a fixed contractible, gauge-reduced block.  The
analytic link metric and Wilson potential give
\(\eta_a=O(gR)\) on the V39 interior chart, while the gauge-invariant IMS
localization contributes \(c_a=O(R^{-2})\).  With the common schedule
\[
 R=L^{1/16},\qquad M=R^2=L^{1/8},
\]
one has \(\eta_a=O(L^{-7/16})\) and \(c_a=O(L^{-1/8})=o(M)\), precisely the
hypotheses used in V42.

The proof is deliberately local.  It does not turn the toron sector or
overlapping block boundary variables into small perturbations.  Those remain
global slow channels.

\section{Analytic gauge-reduced coordinates}
\label{sec:s7v43-coordinates}

Fix a finite contractible lattice block \({\mathscr B}\), impose maximal-tree
gauge, and quotient the remaining infinitesimal gauge kernel.  The resulting
coordinate space is a finite-dimensional real vector space
\({\cal A}_{\mathscr B}^{\rm T}\) with coordinate \(q\).  Choose the
Lie-algebra inner product so that at \(q=0\) the electric metric is the
identity.

On a ball \(|q|\le\rho_0\) inside the exponential injectivity radius, write
the Haar measure and electric inverse metric as
\begin{equation}
 d\mu(q)=J(q)\,dq,\qquad
 A(q)=\bigl(A^{ij}(q)\bigr),                            \label{eq:s7v43-JA}
\end{equation}
with
\begin{equation}
 J(0)=1,\qquad A(0)=I.                                 \label{eq:s7v43-origin}
\end{equation}

\begin{lemma}[Uniform analytic coefficient bounds]
\label{lem:s7v43-coefficients}
After decreasing \(\rho_0\), there is \(C_{\mathscr B}\) such that
\begin{align}
 \|A(q)-I\|+|J(q)-1|
 &\le C_{\mathscr B}|q|,                               \label{eq:s7v43-A-bound}\\
 |\nabla_q\log J(q)|
 &\le C_{\mathscr B}|q|                                \label{eq:s7v43-J-derivative}
\end{align}
for \(|q|\le\rho_0\).
\end{lemma}

\begin{proof}
The exponential-coordinate metric and Haar density are analytic.  The first
estimate is Taylor's theorem at \eqref{eq:s7v43-origin}.  In exponential
normal coordinates for a compact unimodular group, inversion \(q\mapsto-q\)
preserves Haar measure, so \(J(q)=J(-q)\).  Hence
\(\nabla J(0)=0\), and Taylor's theorem gives
\eqref{eq:s7v43-J-derivative}.  Passing to the finite gauge-reduced product
chart changes only the block-dependent constant.
\end{proof}

\begin{remark}[A weaker density bound would also suffice]
\label{rem:s7v43-density}
If a chosen gauge slice produces only
\(|\nabla\log J|\le C_{\mathscr B}\), the proof below gains a harmless
\(O(g^2/a)\) zeroth-order term instead of the sharper
\(O(g^2R^2/a)\) term.  Both are \(o(M/a)\) for the declared schedule.
\end{remark}

\section{Electric-form comparison after dilation}
\label{sec:s7v43-electric}

Set \(q=gy\) and let
\begin{equation}
 ({\cal U}_g\psi)(y)
 :=
 g^{n/2}J(gy)^{1/2}\psi(gy)                            \label{eq:s7v43-U}
\end{equation}
map the coordinate Haar space unitarily to flat \(L^2(dy)\).  Let
\begin{equation}
 t_0(\phi)
 :=
 {1\over2a}\int|\nabla_y\phi|^2\,dy                    \label{eq:s7v43-t0}
\end{equation}
be the flat electric form.

\begin{theorem}[Electric relative-form bound]
\label{thm:s7v43-electric}
There are \(c_0,C_{\mathscr B}>0\) such that, for \(R\ge1\),
\(gR\le c_0\), and \(\phi\) supported in \(|y|\le R\),
\begin{equation}
 |t_g(\phi)-t_0(\phi)|
 \le
 C_{\mathscr B}gR\,t_0(\phi)
 +{C_{\mathscr B}g^2R^2\over a}\|\phi\|_2^2,           \label{eq:s7v43-electric}
\end{equation}
where \(t_g\) is the conjugated Haar electric form.
\end{theorem}

\begin{proof}
Before conjugation, dilation cancels the factor \(g^2\) multiplying the link
Laplacian and gives
\[
 {1\over2a}
 \int
 \langle\nabla_y\widetilde\phi,A(gy)\nabla_y\widetilde\phi\rangle
 J(gy)\,dy.
\]
Conjugation by \(J(gy)^{1/2}\) replaces
\(\nabla\widetilde\phi\) by
\[
 J(gy)^{-1/2}
 \left[
 \nabla\phi-\frac12g\nabla_q\log J(gy)\,\phi
 \right].
\]
On \(|y|\le R\), \Cref{lem:s7v43-coefficients} gives
\[
 \|A(gy)-I\|\le C_{\mathscr B}gR,\qquad
 g|\nabla_q\log J(gy)|
 \le C_{\mathscr B}g^2R.
\]
Expand the square.  The principal coefficient error is bounded by
\(C_{\mathscr B}gR\,t_0\).  Cauchy--Schwarz with a small fraction of the
principal form absorbs the cross term; the remaining square is at most
\(C_{\mathscr B}g^4R^2a^{-1}\|\phi\|^2\), which is bounded by the displayed
zeroth-order term for \(0<g\le1\).  This proves
\eqref{eq:s7v43-electric}.
\end{proof}

\section{Magnetic-form comparison}
\label{sec:s7v43-magnetic}

Let
\begin{equation}
 S_{\mathscr B}(q)
 :=
 \sum_{p\subset{\mathscr B}}
 \left(
 1-{1\over N}\operatorname{Re}\operatorname{Tr}U_p(q)
 \right).                                              \label{eq:s7v43-S}
\end{equation}
On the transverse coordinate space its Hessian at zero is a positive
quadratic form \(Q_2(q)\).  Normalize \(Q_2\) so that the quadratic magnetic
form after dilation is
\begin{equation}
 v_0(\phi)
 :=
 {1\over a}\int Q_2(y)|\phi(y)|^2\,dy.                 \label{eq:s7v43-v0}
\end{equation}

\begin{lemma}[Relative Taylor remainder]
\label{lem:s7v43-magnetic}
There are \(c_0,C_{\mathscr B}>0\) such that
\begin{equation}
 \left|
 {1\over g^2}S_{\mathscr B}(gy)-Q_2(y)
 \right|
 \le C_{\mathscr B}g|y|Q_2(y)                         \label{eq:s7v43-magnetic}
\end{equation}
whenever \(g|y|\le c_0\).
\end{lemma}

\begin{proof}
Analyticity and the vanishing of the constant and linear terms give
\[
 S_{\mathscr B}(q)=Q_2(q)+R_3(q),\qquad
 |R_3(q)|\le C_{\mathscr B}|q|^3.
\]
On the finite-dimensional transverse quotient,
\(Q_2(q)\ge c_{\mathscr B}|q|^2\).  Hence
\[
 |R_3(q)|
 \le C_{\mathscr B}|q|Q_2(q).
\]
Set \(q=gy\), use \(Q_2(gy)=g^2Q_2(y)\), and divide by \(g^2\).
\end{proof}

\begin{theorem}[Magnetic relative-form bound]
\label{thm:s7v43-magnetic}
For \(\phi\) supported in \(|y|\le R\) with \(gR\le c_0\),
\begin{equation}
 |v_g(\phi)-v_0(\phi)|
 \le C_{\mathscr B}gR\,v_0(\phi),                      \label{eq:s7v43-magnetic-form}
\end{equation}
where
\[
 v_g(\phi)
 ={1\over ag^2}
 \int S_{\mathscr B}(gy)|\phi(y)|^2\,dy.
\]
\end{theorem}

\begin{proof}
Multiply \eqref{eq:s7v43-magnetic} by
\(a^{-1}|\phi|^2\), use \(|y|\le R\), and integrate.
\end{proof}

\section{Insertion of the gauge-invariant IMS debit}
\label{sec:s7v43-IMS}

Let
\begin{equation}
 h_{0,a}:=t_0+v_0                                    \label{eq:s7v43-h0}
\end{equation}
be the fixed-block transverse oscillator form.  Use the V39 cutoff threshold
\begin{equation}
 \varepsilon=g^2R^2.                                   \label{eq:s7v43-epsilon}
\end{equation}

\begin{theorem}[Fixed-block Wilson relative-form card]
\label{thm:s7v43-card}
After the common scalar ground-energy normalization, the weak-field
localized Wilson form \(h_{g,R}\) satisfies
\begin{equation}
 |h_{g,R}(\phi)-h_{0,a}(\phi)|
 \le
 \eta_{g,R}h_{0,a}(\phi)
 +{c_R\over a}\|\phi\|^2,                              \label{eq:s7v43-card}
\end{equation}
where
\begin{align}
 \eta_{g,R}&\le C_{\mathscr B}gR,                      \label{eq:s7v43-eta}\\
 c_R&\le C_{\mathscr B}\left(R^{-2}+g^2R^2\right).
 \label{eq:s7v43-c}
\end{align}
The constants are independent of \(a\) and \(g\) for the fixed block.
\end{theorem}

\begin{proof}
Add
\Cref{thm:s7v43-electric,thm:s7v43-magnetic}.  The relative terms are
bounded by \(C_{\mathscr B}gR\,h_{0,a}\), while the electric density
conjugation contributes \(C_{\mathscr B}g^2R^2/a\).

The exact V39 IMS identity subtracts at most
\[
 {C_{\mathscr B}\over aR^2}\|\phi\|^2
\]
when \(\varepsilon=g^2R^2\).  Insert this as a two-sided localization debit;
enlarging the fixed constant gives
\eqref{eq:s7v43-card}--\eqref{eq:s7v43-c}.  A common scalar
ground-energy shift changes neither side's excitation gap and is included
before comparison.
\end{proof}

\begin{corollary}[The V42 hypotheses close on a fixed block]
\label{cor:s7v43-V42}
Assume
\begin{equation}
 g(a)=O(L(a)^{-1/2}),\qquad
 R(a)=L(a)^{1/16},\qquad
 M(a)=R(a)^2=L(a)^{1/8}.                               \label{eq:s7v43-schedule}
\end{equation}
Then
\begin{align}
 \eta_{g,R}
 &=O(L(a)^{-7/16})\longrightarrow0,                   \label{eq:s7v43-eta-rate}\\
 c_R
 &=O(L(a)^{-1/8}),                                    \label{eq:s7v43-c-rate}\\
 {c_R\over M(a)}
 &=O(L(a)^{-1/4})\longrightarrow0.                    \label{eq:s7v43-c-over-M}
\end{align}
Therefore the assumptions
\(\eta_a\to0\) and \(c_a=o(M)\) in V42 hold for the fixed-block transverse
weak-field form.
\end{corollary}

\begin{proof}
The first exponent is
\(-1/2+1/16=-7/16\).  In \eqref{eq:s7v43-c},
\[
 R^{-2}=L^{-1/8},\qquad
 g^2R^2=O(L^{-7/8}),
\]
so the first term dominates.  Divide by \(M=L^{1/8}\).
\end{proof}

\begin{corollary}[Explicit occupation floor]
\label{cor:s7v43-floor}
Let \(\omega_*>0\) be the lowest fixed-block transverse oscillator
frequency.  The V42 residual coefficient satisfies
\begin{equation}
 \Delta_M(a)
 =
 (1-O(L^{-7/16}))\omega_*(M+1)-O(L^{-1/8})
 \ge{\omega_*M\over2}                                 \label{eq:s7v43-floor}
\end{equation}
for all sufficiently small \(a\).
\end{corollary}

\begin{proof}
Insert \eqref{eq:s7v43-eta-rate}--\eqref{eq:s7v43-c-rate} into the V42
formula for \(\Delta_M(a)\).  Since \(M\to\infty\), the displayed lower
bound follows.
\end{proof}

\section{Toron and boundary rows}
\label{sec:s7v43-slow}

\begin{firewall}[Torons are outside the transverse coercive form]
\label{fire:s7v43-toron}
The proof of \Cref{lem:s7v43-magnetic} uses positivity of \(Q_2\) on the
finite transverse quotient.  Periodic torons and exact gauge directions lie
in its kernel.  Gauge directions are quotiented; torons are physical compact
slow variables and must remain in the low operator.  The relative bound does
not assign them an oscillator frequency.
\end{firewall}

\begin{firewall}[Block boundary holonomies are slow gluing data]
\label{fire:s7v43-boundary}
Maximal-tree gauge on one contractible block leaves boundary transformations
that are paired with neighboring blocks in the global Gauss law.  V43 proves
a form estimate after fixing one block chart.  It does not prove that the
product of those charts has a volume-uniform Jacobian or that boundary
holonomy integration preserves the same constants.
\end{firewall}

\begin{firewall}[The global Gaussian split is still nonlocal]
\label{fire:s7v43-global}
The fixed-block relative estimate controls local coefficients of the
nonlinear remainder.  The momentum-band projection selected in V41 is a
global Gaussian spectral projection.  Converting the local form estimate
into the global anchored Feshbach norm still requires the V34--V36
quasi-local filter and a connected support expansion.
\end{firewall}

\section{Relative-form status}
\label{sec:s7v43-status}

\begin{theorem}[What V43 proves]
\label{thm:s7v43-result}
On every fixed contractible gauge-reduced block, the weak-field localized
Wilson Hamiltonian is relatively form-close to its joint
electric--magnetic oscillator:
\[
 \eta_{g,R}=O(gR),\qquad
 c_R=O(R^{-2}+g^2R^2).
\]
For the schedules \(g=O(L^{-1/2})\), \(R=L^{1/16}\), and \(M=R^2\), this
proves the local analytic hypotheses used in the V42 high-occupation
Feshbach theorem.
\end{theorem}

\begin{proof}
Use
\Cref{lem:s7v43-coefficients,thm:s7v43-electric,
lem:s7v43-magnetic,thm:s7v43-card,
cor:s7v43-V42,cor:s7v43-floor}.
\end{proof}

\begin{assessment}[Progress at seventh-series V43]
\label{ass:s7v43-status}
The high-occupation theorem is no longer conditional on an unspecified
local analytic estimate.  Haar kinetic coefficients, Wilson curvature, and
the gauge-invariant IMS derivative fit one explicit relative-form card, and
the common polylogarithmic schedule closes its numerical inequalities.
The remaining obstruction has moved to the correct global level: toron and
boundary slow variables, quasi-local spectral filtering, and connected
volume-uniform assembly of the self-energy-corrected low transfer.
\end{assessment}

\begin{decision}[V44 acquisition order]
\label{dec:s7v43-V44}
V44 will formulate the connected assembly theorem for the global Gaussian
split using scalar normalized slab correlations on the gauge-invariant OS
core.  It will prove the exact cumulant cancellation of block-disconnected
vacuum factors and state a tree-summability criterion in the anchored
activity produced by V41--V43.  Torons and boundary holonomies will be
retained as external slow labels rather than integrated into the residual
polymer activity.
\end{decision}

\begin{firewall}[Claim boundary after seventh-series V43]
\label{fire:s7v43-claim}
V43 proves the fixed-block analytic coefficient bounds, electric and
magnetic relative-form comparisons, the explicit IMS-inclusive card, and
closure of the V42 local occupation-floor hypotheses.

V43 does not prove global slow-variable gluing, the connected cumulant
expansion, coarse self-energy matching, a terminal nonperturbative gap,
cutoff removal, or the full Yang--Mills mass gap.
\end{firewall}

\paragraph{Parent-version record.}
V43 was formed from
\texttt{Renewal\_Yang\_Mills\_Mass_Gap\_Research\_Notes\_V42.tex}.
The V42 parent had \(17\,806\) logical source lines, \(647\,224\) bytes,
and SHA--256
\[
 \texttt{DAF5C2B3151112F93EA5C00123F3B4D71AC2122067BB330B251959A6B3026FA6}.
\]
The next compulsory companion audit is V45.

% =============================================================================
% END APPENDED SEVENTH-SERIES V43 MATERIAL
% =============================================================================

% =============================================================================
% BEGIN APPENDED SEVENTH-SERIES V44 MATERIAL
% =============================================================================

\chapter{Normalized connected slab assembly}
\label{ch:s7v44}

V41--V43 isolate a vanishing dimensionless nonlinear activity on a
weak-field oscillator window, but a sum of local interactions still has an
extensive global norm.  The mass-gap criterion concerns normalized connected
correlations, where vacuum-disconnected terms must cancel before any norm is
taken.  This note proves that cancellation by differentiating the logarithm
of the normalized generating functional.  It then gives a tree-summability
theorem for the surviving connected families.

The reference is treated in two layers.  A product block reference gives
literal vanishing of disconnected cumulants.  The actual eliminated Gaussian
sector is not a product in position space; its replacement is a uniform
tree-cumulant bound whose edges carry the exponentially decaying residual
covariance.  Slow momentum modes, torons, and boundary holonomies are fixed
external labels throughout.  No mass is attributed to them by the residual
cluster expansion.

\section{Normalized connected correlation as a logarithmic derivative}
\label{sec:s7v44-log}

Let \((\Omega_0,\mu_{0,\xi})\) be a finite-volume reference probability space
depending on an external slow label \(\xi\).  Let
\begin{equation}
 V=\sum_{X\in{\cal P}_\Lambda}v_X                         \label{eq:s7v44-V}
\end{equation}
be a finite sum of bounded local slab interactions.  For bounded observables
\(A,B\), define
\begin{equation}
 Z_\xi(s,t,u)
 :=
 \mathbb E_{0,\xi}
 \exp\{sA+tB-uV\}.                                      \label{eq:s7v44-Z}
\end{equation}
The interacting normalized expectation is
\begin{equation}
 \mathbb E_{V,\xi}[F]
 :=
 {\mathbb E_{0,\xi}[Fe^{-V}]
  \over
  \mathbb E_{0,\xi}[e^{-V}]}.                           \label{eq:s7v44-EV}
\end{equation}

\begin{lemma}[Exact logarithmic identity]
\label{lem:s7v44-log}
Whenever \(Z_\xi\) is nonzero near \((0,0,1)\),
\begin{equation}
 \mathbb E_{V,\xi}[A;B]
 :=
 \mathbb E_{V,\xi}[AB]
 -\mathbb E_{V,\xi}[A]\mathbb E_{V,\xi}[B]
 =
 \partial_s\partial_t
 \log Z_\xi(s,t,1)\big|_{s=t=0}.                        \label{eq:s7v44-log}
\end{equation}
\end{lemma}

\begin{proof}
Differentiate \(\log Z\) once:
\[
 \partial_s\log Z={\mathbb E_{0,\xi}[Ae^{sA+tB-V}]
 \over \mathbb E_{0,\xi}[e^{sA+tB-V}]}.
\]
Differentiating in \(t\) and setting \(s=t=0\) gives the covariance under
\eqref{eq:s7v44-EV}.
\end{proof}

Let
\(\kappa_{0,\xi}(F_1,\ldots,F_m)\) denote the joint cumulant under
\(\mu_{0,\xi}\), defined by
\begin{equation}
 \kappa_{0,\xi}(F_1,\ldots,F_m)
 :=
 \partial_{z_1}\cdots\partial_{z_m}
 \log\mathbb E_{0,\xi}
 e^{\sum_jz_jF_j}\bigg|_{z=0}.                          \label{eq:s7v44-cumulant}
\end{equation}

\begin{theorem}[Exact connected perturbation series]
\label{thm:s7v44-series}
If the series below is absolutely convergent, then
\begin{align}
 \mathbb E_{V,\xi}[A;B]
 =
 \sum_{n=0}^\infty{(-1)^n\over n!}
 \sum_{X_1,\ldots,X_n\in{\cal P}_\Lambda}
 \kappa_{0,\xi}
 \bigl(A,B,v_{X_1},\ldots,v_{X_n}\bigr).
 \label{eq:s7v44-series}
\end{align}
The denominator in \eqref{eq:s7v44-EV} has disappeared into cumulants; no
vacuum partition-function factor remains.
\end{theorem}

\begin{proof}
Expand \(\log Z_\xi(s,t,u)\) at \(u=0\) using
\eqref{eq:s7v44-cumulant}.  The coefficient of
\(st\,u^n\) is
\[
 {(-1)^n\over n!}
 \kappa_{0,\xi}(A,B,V,\ldots,V).
\]
Multilinearity expands each copy of \(V\) according to
\eqref{eq:s7v44-V}.  Absolute convergence justifies evaluation at \(u=1\)
and interchange of sums and derivatives, yielding
\eqref{eq:s7v44-series}.
\end{proof}

\section{Exact cancellation for a product reference}
\label{sec:s7v44-product}

Suppose the reference variables factor over a set of blocks \({\cal C}\):
\begin{equation}
 \mu_{0,\xi}
 =
 \bigotimes_{C\in{\cal C}}\mu_{C,\xi}.                  \label{eq:s7v44-product}
\end{equation}
For a random variable \(F\), let \({\rm supp}(F)\subset{\cal C}\) be a block
support on which it is measurable.

\begin{lemma}[Cumulants vanish across an independence split]
\label{lem:s7v44-independent}
Suppose \(\{1,\ldots,m\}=I\sqcup J\), with \(I,J\ne\varnothing\), and
\begin{equation}
 \left(\bigcup_{i\in I}{\rm supp}(F_i)\right)
 \cap
 \left(\bigcup_{j\in J}{\rm supp}(F_j)\right)
 =\varnothing.                                         \label{eq:s7v44-independent}
\end{equation}
Then
\begin{equation}
 \kappa_{0,\xi}(F_1,\ldots,F_m)=0.                     \label{eq:s7v44-cumulant-zero}
\end{equation}
\end{lemma}

\begin{proof}
The two families are independent by
\eqref{eq:s7v44-product}, so their joint moment generating function
factorizes:
\[
 \mathbb E_{0,\xi}e^{\sum_i z_iF_i}
 =
 \mathbb E_{0,\xi}e^{\sum_{i\in I}z_iF_i}
 \mathbb E_{0,\xi}e^{\sum_{j\in J}z_jF_j}.
\]
Its logarithm is a sum of a function of the \(I\)-variables and a function
of the \(J\)-variables.  The mixed derivative defining the full joint
cumulant is therefore zero.
\end{proof}

\begin{corollary}[Only support-connected families survive]
\label{cor:s7v44-connected}
Under \eqref{eq:s7v44-product}, a term in
\eqref{eq:s7v44-series} can be nonzero only if the hypergraph with vertices
\[
 A,\ B,\ X_1,\ldots,X_n
\]
and edges given by support intersection is connected.  In particular its
union contains a chain from \({\rm supp}(A)\) to
\({\rm supp}(B)\).
\end{corollary}

\begin{proof}
If the intersection hypergraph is disconnected, its vertices split into
two nonempty families with disjoint unions of block supports.  Apply
\Cref{lem:s7v44-independent}.
\end{proof}

\begin{firewall}[The global Gaussian is not a product reference]
\label{fire:s7v44-not-product}
The exact Gaussian operator selected in V41 contains cross-block quadratic
curvature.  Its position-space measure does not satisfy
\eqref{eq:s7v44-product}.  Therefore
\Cref{cor:s7v44-connected} cannot be applied to it by merely declaring
real-space blocks independent.  The gapped residual covariance replaces
literal independence by the tree bound below.
\end{firewall}

\section{Tree-cumulant hypothesis for the residual Gaussian}
\label{sec:s7v44-tree}

Let \(d(S,T)\) be the graph distance between finite block supports.  Give
each local random variable \(F\) a seminorm \(\|F\|_*\) that controls the
derivatives required by Gaussian integration by parts.  Assume the following
uniform bound for fixed slow labels \(\xi\).

\begin{assumption}[Residual tree-cumulant bound]
\label{ass:s7v44-tree}
There are \(C_0,\mu>0\), independent of volume and \(\xi\), such that for
every \(m\ge2\),
\begin{align}
 |\kappa_{0,\xi}(F_1,\ldots,F_m)|
 \le
 C_0^m\prod_{i=1}^m\|F_i\|_*
 \sum_{T\in{\cal T}_m}
 \prod_{\{i,j\}\in E(T)}
 e^{-\mu d({\rm supp}F_i,{\rm supp}F_j)}.              \label{eq:s7v44-tree}
\end{align}
Here \({\cal T}_m\) is the set of labelled spanning trees on
\{1,\ldots,m\}.
\end{assumption}

\begin{lemma}[Origin of the tree edge]
\label{lem:s7v44-Gaussian-edge}
Let \(\zeta\) be a centered finite-dimensional Gaussian with covariance
\({\cal C}\).  For differentiable functions \(F,G\) supported on coordinate
sets \(S,T\),
\begin{equation}
 |\operatorname{Cov}(F,G)|
 \le
 \sum_{x\in S,y\in T}
 |{\cal C}_{xy}|\,
 \|\partial_xF\|_\infty\|\partial_yG\|_\infty.          \label{eq:s7v44-Gaussian-edge}
\end{equation}
If
\begin{equation}
 \sup_x\sum_y e^{\mu d(x,y)}|{\cal C}_{xy}|
 \le C_{\cal C},                                       \label{eq:s7v44-Cdecay}
\end{equation}
the right side is bounded by a constant times
\(e^{-\mu d(S,T)}\) times the corresponding local derivative norms.
\end{lemma}

\begin{proof}
Let \((\zeta,\zeta')\) be independent copies and interpolate
\[
 \zeta_t=\sqrt t\,\zeta+\sqrt{1-t}\,\zeta'.
\]
Differentiating the standard Gaussian interpolation between the joint law
and the product law, then integrating by parts, gives
\[
 \operatorname{Cov}(F,G)
 =
 \int_0^1
 \sum_{x,y}{\cal C}_{xy}
 \mathbb E[
 \partial_xF(\zeta)\,
 \partial_yG(\zeta_t)]\,dt
\]
in a choice of equivalent interpolation; taking absolute values yields
\eqref{eq:s7v44-Gaussian-edge}.  If \(x\in S,y\in T\), then
\(d(x,y)\ge d(S,T)\).  Pull out
\(e^{-\mu d(S,T)}\) and apply \eqref{eq:s7v44-Cdecay}.
\end{proof}

\begin{remark}[Higher cumulants]
\label{rem:s7v44-higher}
Iterated Gaussian interpolation and integration by parts produce one
covariance edge at each step needed to connect the \(m\) insertions.  Summing
the possible connection histories gives a spanning-tree bound of the form
\eqref{eq:s7v44-tree}.  V44 uses that bound as a declared hypothesis because
the derivative seminorms for the cutoff Wilson polynomials and their
uniformity in external toron/boundary labels must still be checked.
\end{remark}

\section{Anchored activity and tree summation}
\label{sec:s7v44-activity}

Assume every interaction support has at most \(s_*\) blocks.  Define
\begin{equation}
 \rho_\mu
 :=
 \sup_{C\in{\cal C}}
 \sum_{X\ni C}
 e^{\mu\,{\rm diam}(X)}
 C_0\|v_X\|_* .                                        \label{eq:s7v44-rho}
\end{equation}
Let
\[
 r=d({\rm supp}A,{\rm supp}B).
\]

\begin{theorem}[Tree-summability of normalized connected correlations]
\label{thm:s7v44-summability}
Under \Cref{ass:s7v44-tree}, there is a numerical constant
\(C_{\rm tr}\), depending only on \(s_*\) and the bounded local geometry,
such that if
\begin{equation}
 C_{\rm tr}\rho_\mu<1,                                 \label{eq:s7v44-small}
\end{equation}
the series \eqref{eq:s7v44-series} converges absolutely and uniformly in
volume and \(\xi\).  Moreover, for every \(0<\mu'<\mu\),
\begin{equation}
 |\mathbb E_{V,\xi}[A;B]|
 \le
 C_{A,B,\mu'}e^{-\mu' r},                              \label{eq:s7v44-decay}
\end{equation}
where the constant is independent of volume and the external slow label.
\end{theorem}

\begin{proof}
Insert \eqref{eq:s7v44-tree} into the absolute value of
\eqref{eq:s7v44-series}.  Root each spanning tree at the vertex carrying
\(A\).  Sum interaction supports successively from the leaves toward the
root.  If a child support is attached to a parent support \(S\), choose one
of at most \(|S|\le s_*\) anchor blocks and use
\eqref{eq:s7v44-rho}; the factor
\(e^{-\mu d(S,X)}\) is absorbed by the exponential diameter weight and the
bounded number of block positions at each distance.  Each interaction
vertex therefore costs at most \(C_{\rm tr}\rho_\mu\).

There are \((n+2)^n\) labelled trees on \(n+2\) vertices by Cayley's formula.
The factor \(1/n!\) in \eqref{eq:s7v44-series} and Stirling's bound give
\begin{equation}
 {(n+2)^n\over n!}
 \le C_1(C_2)^n                                       \label{eq:s7v44-Cayley}
\end{equation}
for numerical \(C_1,C_2\).  Increase \(C_{\rm tr}\) to include \(C_2\),
the two observable vertices, and the finite support choices.  The resulting
geometric series converges under \eqref{eq:s7v44-small}, uniformly in the
number of available blocks.

Every spanning tree contains a path from the \(A\)-vertex to the
\(B\)-vertex.  By the triangle inequality, the sum of support distances and
diameters along that path is at least \(r\).  Reserve
\(\mu' r\) from the exponential edge/diameter factors and perform the same
summation with the remaining rate \(\mu-\mu'>0\).  This gives
\eqref{eq:s7v44-decay}.
\end{proof}

\begin{corollary}[Exact volume cancellation]
\label{cor:s7v44-volume}
Under \eqref{eq:s7v44-small}, the connected correlation bound contains no
factor proportional to \(|{\cal C}|\).  Interactions not joined to the
observable cluster contribute only to
\(\mathbb E_{0,\xi}[e^{-V}]\) and cancel through the logarithm in
\eqref{eq:s7v44-log}.
\end{corollary}

\begin{proof}
The rooted tree summation anchors every surviving interaction to
\({\rm supp}A\cup{\rm supp}B\); it never chooses an unrestricted first block.
The exact logarithmic identity already incorporates the vacuum
normalization.
\end{proof}

\section{Insertion of the V41--V43 activity}
\label{sec:s7v44-insertion}

Let
\begin{equation}
 L(a)=\log{1\over a\Lambda_0},\qquad
 M(a)=L(a)^{1/8},\qquad R(a)=L(a)^{1/16}.              \label{eq:s7v44-schedule}
\end{equation}
The local cubic, quartic, filter, and form debits found earlier have
dimensionless sizes
\begin{align}
 \rho_3(a)&=O(L^{-5/16}),                              \label{eq:s7v44-rho3}\\
 \rho_4(a)&=O(L^{-3/4}),                               \label{eq:s7v44-rho4}\\
 \rho_{\rm mix}(a)&=O(L^{-7/16}),                     \label{eq:s7v44-rhomix}\\
 \rho_{\rm form}(a)&=O(L^{-7/16}+L^{-1/8}).           \label{eq:s7v44-rhoform}
\end{align}

\begin{assumption}[Renormalized slab activity map]
\label{ass:s7v44-map}
After the \(O(g^2M^2/a)\) self-energy is inserted into the low effective
Hamiltonian, the remaining normalized fixed-slab kernels admit local
activities with
\begin{equation}
 \rho_\mu(a)
 \le
 C\bigl(
 \rho_3(a)+\rho_4(a)+\rho_{\rm mix}(a)
 +\rho_{\rm form}(a)+\rho_{\rm tail}(a)
 \bigr),                                               \label{eq:s7v44-map}
\end{equation}
uniformly in volume and external slow labels, where
\(\rho_{\rm tail}(a)\to0\).
\end{assumption}

\begin{theorem}[Conditional ultraviolet connected closure]
\label{thm:s7v44-UV}
Under \Cref{ass:s7v44-tree,ass:s7v44-map}, if
\begin{equation}
 \rho_{\rm tail}(a)\longrightarrow0,                   \label{eq:s7v44-tail}
\end{equation}
then the tree-summability condition
\eqref{eq:s7v44-small} holds for all sufficiently small \(a\).  The
normalized residual correction to every pair of local OS-core observables
is volume uniform and exponentially clustered at every rate
\(\mu'<\mu\).
\end{theorem}

\begin{proof}
All four explicit powers in
\eqref{eq:s7v44-rho3}--\eqref{eq:s7v44-rhoform} tend to zero.  Together
with \eqref{eq:s7v44-tail}, the right side of
\eqref{eq:s7v44-map} tends to zero.  Apply
\Cref{thm:s7v44-summability}.
\end{proof}

\begin{firewall}[The activity map is the remaining nonlinear theorem]
\label{fire:s7v44-map}
V44 proves what follows from an anchored residual activity; it does not
derive \eqref{eq:s7v44-map} from the compact Wilson transfer.  That derivation
must show that high-sojourn integration supplies the energy denominators of
V37 and V42, that every \(g^2/a\) self-energy term is moved into the coarse
Hamiltonian, and that the external toron/boundary dependence is uniform.
Using the raw time integral \(\tau\|W_a\|\) would give the wrong parameter.
\end{firewall}

\section{Interface with the mass-gap recursion}
\label{sec:s7v44-gap}

\begin{theorem}[Connected residuals are sufficient only after low matching]
\label{thm:s7v44-gap-interface}
Suppose the hypotheses of \Cref{thm:s7v44-UV} hold and, in addition, the
self-energy-corrected low slab transfer matches a next-scale transfer with
error \(\eta_a\to0\).  If that next-scale transfer has vacuum-complement norm
at most \(e^{-m\tau}\), then the V35 fixed-slab compiler gives, for every
\(0<m'<m\),
\begin{equation}
 \|S_a(\tau)\|_{\Omega^\perp}\le e^{-m'\tau}            \label{eq:s7v44-gap}
\end{equation}
for sufficiently small \(a\), provided the connected activity bound controls
the operator row required by that compiler.
\end{theorem}

\begin{proof}
The low-row matching debit is \(\eta_a\).  The high-momentum and
high-occupation mixing debits tend to zero by V36 and V42, while
\Cref{thm:s7v44-UV} removes the extensive vacuum factor from the nonlinear
residual.  Under the stated operator-row conversion, the total fixed-slab
debit tends to zero.  Apply the V35 fixed-margin argument.
\end{proof}

\begin{firewall}[Scalar clustering is not automatically an operator norm]
\label{fire:s7v44-scalar}
\Cref{thm:s7v44-UV} controls scalar connected correlations on a local
observable core.  The last hypothesis of
\Cref{thm:s7v44-gap-interface} is a genuine conversion step.  One may avoid
it by using the V25 theorem that uniform exponential time decay on a dense
OS core is itself sufficient for a mass gap, but then the connected
expansion must be proved in Euclidean time uniformly for that entire core.
\end{firewall}

\section{Connected-assembly status}
\label{sec:s7v44-status}

\begin{theorem}[What V44 proves]
\label{thm:s7v44-result}
Normalized interacting covariances have the exact cumulant expansion
\eqref{eq:s7v44-series}; product-reference disconnected families vanish
identically; and a uniform residual tree-cumulant bound plus small anchored
activity yields volume-uniform exponential clustering.  The explicit
V41--V43 ultraviolet powers satisfy the resulting smallness criterion once
the renormalized slab activity map and tail row are proved.
\end{theorem}

\begin{proof}
Use
\Cref{lem:s7v44-log,thm:s7v44-series,
lem:s7v44-independent,cor:s7v44-connected,
thm:s7v44-summability,thm:s7v44-UV}.
\end{proof}

\begin{assessment}[Progress at seventh-series V44]
\label{ass:s7v44-status}
The volume-extensive interaction norm has been replaced by the correct
normalized connected object.  Vacuum-disconnected components cancel
exactly, while the gapped Gaussian residual can be handled by covariance
trees rather than false block independence.  The remaining ultraviolet
theorem is now sharply stated: construct the renormalized slab activity map
with uniform slow-label derivative norms and convert its time-connected
decay to the dense OS core.  Independently, the low self-energy-corrected
transfer still has to be matched along the RG trajectory to the V27
strong-coupling terminal gap.
\end{assessment}

\begin{decision}[V45 acquisition order]
\label{dec:s7v44-V45}
V45 is the compulsory SM/NS audit.  It will then isolate the tail term
\(\rho_{\rm tail}\) in \eqref{eq:s7v44-map}: chart exit, high occupation,
transition momentum, toron variation, and boundary holonomy will receive
separate rows.  The first target is to prove which of these rows are already
controlled by V34--V43 and which still obstruct uniformity in the external
slow labels.
\end{decision}

\begin{firewall}[Claim boundary after seventh-series V44]
\label{fire:s7v44-claim}
V44 proves the logarithmic connected-correlation identity, exact
product-reference cancellation, the abstract residual tree-summability
theorem, and conditional ultraviolet connected closure.

V44 does not prove the residual Gaussian tree bound in the required Wilson
seminorms, the renormalized activity map, uniform toron/boundary-label
control, low coarse matching, a terminal gap handoff, cutoff removal, or the
full Yang--Mills mass gap.
\end{firewall}

\paragraph{Parent-version record.}
V44 was formed from
\texttt{Renewal\_Yang\_Mills\_Mass_Gap\_Research\_Notes\_V43.tex}.
The V43 parent had \(18\,213\) logical source lines, \(661\,105\) bytes,
and SHA--256
\[
 \texttt{BC261D02DCB45DBDADC968B718B45776147239CABBC914AFC003985C511CB1AC}.
\]
The next compulsory companion audit is V45.

% =============================================================================
% END APPENDED SEVENTH-SERIES V44 MATERIAL
% =============================================================================

% =============================================================================
% BEGIN APPENDED SEVENTH-SERIES V45 MATERIAL
% =============================================================================

\chapter{Companion audit and the residual-tail ledger}
\label{ch:s7v45}

V44 reduces volume-uniform connected assembly to a residual tree-cumulant
bound, a renormalized slab activity map, and a tail
\(\rho_{\rm tail}(a)\to0\).  This compulsory-audit note identifies what is
already available for those rows.  The newest SM note supplies a
two-sided exponentially weighted Schur algebra and a source-marked
vacuum-loop cancellation.  Only the weighted inverse mechanism transfers:
the determinant and Einstein constraint factors do not.

The tail then splits into high momentum, filter truncation, high occupation,
chart exit, IMS localization, toron dependence, and boundary-holonomy
dependence.  The first five have explicit vanishing schedules or residual
energy floors.  Torons and boundary holonomies are not tails at all; they
are external slow variables.  The remaining uniformity problem is to cover
their compact parameter space by buffered spectral patches that retain every
emerging near-zero mode.

\section{Compulsory V45 companion audit}
\label{sec:s7v45-audit}

\subsection{Files inspected}
\label{sec:s7v45-files}

The current active companion files are
\begin{align}
 &\texttt{Renewal\_SM\_Einstein\_Continuum\_Research\_Notes\_V86.tex},
 \notag\\[-2pt]
 &\qquad 34\,838\ \hbox{logical source lines},\quad
 1\,241\,661\ \hbox{bytes},                             \notag\\[-2pt]
 &\qquad
 \texttt{846B27B83DF6F5F91B339050E20797755E0DF23537FAF2AE69770E1859B6BC06},
 \label{eq:s7v45-sm-checkpoint}\\
 &\texttt{Renewal\_NS\_Stability\_Research\_Notes\_V26.tex},
 \notag\\[-2pt]
 &\qquad 5\,319\ \hbox{logical source lines},\quad
 278\,283\ \hbox{bytes},                                \notag\\[-2pt]
 &\qquad
 \texttt{82C887F8C2C69E4F28FAD7C3DC8DE5B9099CB5A4BF431EBA1FFF3E91935978AD}.
 \label{eq:s7v45-ns-checkpoint}
\end{align}
The next compulsory companion audit is V50.

\subsection{SM V83--V86}
\label{sec:s7v45-sm}

SM V83--V84 derive finite-regulator constrained-measure Jacobians, chart
cocycles, and adjacent-regulator density criteria.  V85 audits YM through
V41 and selects an anchored rather than extensive determinant response.
V86 proves:
\begin{enumerate}[label=\textup{(\roman*)},leftmargin=3.5em]
\item a two-sided exponentially weighted Schur algebra;
\item Neumann inversion with volume-uniform exponential kernel decay;
\item a remote anchored sandwich estimate with \(e^{-2\mu D}\);
\item exact cancellation of unmarked determinant loops in a relative
      source-marked logarithm.
\end{enumerate}

Items (i)--(iii) are type-correct for a Yang--Mills residual covariance or
resolvent once its block kernel is identified.  Item (iv) is determinant
specific, although its organizational principle agrees with the normalized
cumulant cancellation proved directly in V44.  None of the SM constraint
determinants or coarea factors enters the Yang--Mills measure.

\subsection{NS V26}
\label{sec:s7v45-ns}

The NS note remains unchanged.  Its rank-one heat-kernel derivative norms
do not provide a residual Gaussian inverse, a toron chart, or a Wilson slab
activity estimate.

\begin{theorem}[V45 audit decision]
\label{thm:s7v45-audit}
The V44 residual tree hypothesis will be tested through a two-sided weighted
Schur bound
\begin{equation}
 \|C_{\rm r}\|_{{\cal S}_\mu}
 :=
 \max\left\{
 \sup_x\sum_ye^{\mu d(x,y)}\|(C_{\rm r})_{xy}\|,\
 \sup_y\sum_xe^{\mu d(x,y)}\|(C_{\rm r})_{xy}\|
 \right\}<\infty.                                      \label{eq:s7v45-Schur}
\end{equation}
This imports an abstract norm mechanism, not an SM estimate.  The
Yang--Mills residual covariance must still be proved to satisfy
\eqref{eq:s7v45-Schur}, uniformly in the slow labels and cutoff.
\end{theorem}

\begin{proof}
The weighted product and Neumann-series arguments depend only on the metric
kernel algebra and hence transfer without changing field type.  The actual
kernel, fibre, and constants are theory specific, so no numerical or
operator conclusion transfers from SM.
\end{proof}

\section{Weighted residual covariance gives the Gaussian tree edge}
\label{sec:s7v45-covariance}

\begin{lemma}[Weighted Schur inversion]
\label{lem:s7v45-inverse}
Let \(T\) be a block kernel with
\begin{equation}
 \|T\|_{{\cal S}_\mu}=q<1.                             \label{eq:s7v45-q}
\end{equation}
Then
\begin{equation}
 (I+T)^{-1}=\sum_{n=0}^\infty(-T)^n,\qquad
 \|(I+T)^{-1}\|_{{\cal S}_\mu}\le{1\over1-q},          \label{eq:s7v45-inverse}
\end{equation}
and every block satisfies
\begin{equation}
 \|[(I+T)^{-1}]_{xy}\|
 \le{e^{-\mu d(x,y)}\over1-q}.                         \label{eq:s7v45-inverse-decay}
\end{equation}
\end{lemma}

\begin{proof}
The triangle inequality for \(d\) makes the two-sided weighted Schur norm
submultiplicative.  The geometric series converges in that Banach algebra
and has norm at most \((1-q)^{-1}\).  A single term in a weighted row or
column sum gives \eqref{eq:s7v45-inverse-decay}.
\end{proof}

\begin{corollary}[Sufficient covariance row for V44]
\label{cor:s7v45-tree-edge}
Suppose the conditional residual Gaussian covariance has the representation
\begin{equation}
 C_{{\rm r},\xi}
 =
 C_0^{1/2}(I+T_\xi)^{-1}C_0^{1/2},                    \label{eq:s7v45-Crep}
\end{equation}
where \(C_0^{1/2}\in{\cal S}_\mu\) uniformly and
\(\sup_\xi\|T_\xi\|_{{\cal S}_\mu}\le q<1\).  Then
\begin{equation}
 \sup_\xi\|C_{{\rm r},\xi}\|_{{\cal S}_\mu}
 \le
 {\|C_0^{1/2}\|_{{\cal S}_\mu}^2\over1-q}.             \label{eq:s7v45-Cbound}
\end{equation}
Consequently the two-point Gaussian interpolation edge in V44 is uniform in
volume and \(\xi\).
\end{corollary}

\begin{proof}
Apply \Cref{lem:s7v45-inverse} and the weighted product inequality to
\eqref{eq:s7v45-Crep}.  The V44 Gaussian covariance lemma then supplies the
edge \(e^{-\mu d(S,T)}\).
\end{proof}

\begin{firewall}[A spectral gap is weaker than a weighted inverse]
\label{fire:s7v45-gap}
The residual energy floor \(\Delta/a\) bounds the unweighted operator norm
of the inverse.  It does not alone imply
\eqref{eq:s7v45-Schur}.  A finite-range Neumann comparison, a
Combes--Thomas estimate, or a direct heat-kernel bound must prove the
weighted spatial decay.
\end{firewall}

\section{The seven residual-tail rows}
\label{sec:s7v45-tail}

Use the common schedules
\begin{equation}
 L=\log{1\over a\Lambda_0},\qquad
 \kappa=L^{1/16},\qquad
 R=L^{1/16},\qquad
 M=L^{1/8}.                                            \label{eq:s7v45-schedule}
\end{equation}

\subsection{Separated momentum and filter rows}
\label{sec:s7v45-momentum}

V34 gives logistic leakage
\begin{equation}
 \rho_{\rm filt}(a)
 \le C e^{-hL^{1/16}},                                 \label{eq:s7v45-filter}
\end{equation}
and V36 gives quasi-local and finite-range incidence debits
\begin{align}
 \rho_{\rm qloc}(a)&=O(L^{-5/16}),                     \label{eq:s7v45-qloc}\\
 \rho_{\rm fr}(a)&=O(L^{-1/8}).                        \label{eq:s7v45-fr}
\end{align}
For fixed slab time \(\tau\), the free transition/high transfer obeys
\begin{equation}
 \rho_{\rm mom}(a)
 \le e^{-\tau\sqrt\alpha/a}.                           \label{eq:s7v45-momentum}
\end{equation}

\begin{theorem}[Momentum/filter tail closes]
\label{thm:s7v45-momentum}
Every quantity in
\eqref{eq:s7v45-filter}--\eqref{eq:s7v45-momentum} tends to zero, and their
sum does so independently of spatial volume.
\end{theorem}

\begin{proof}
All powers of \(L\) have negative exponents, while the two displayed
exponentials have arguments tending to \(-\infty\).  The underlying V34 and
V36 bounds are weighted row or operator bounds, not global site sums.
\end{proof}

\subsection{High-occupation row}
\label{sec:s7v45-occupation}

V42--V43 give
\begin{align}
 \rho_{\rm occ,mix}(a)&=O(L^{-7/16}),                  \label{eq:s7v45-occ-mix}\\
 \rho_{\rm occ,der}(a)&=O(L^{-7/8}),                   \label{eq:s7v45-occ-der}\\
 \Sigma_{\rm occ}(a)&=O(L^{-3/4}/a).                  \label{eq:s7v45-occ-Sigma}
\end{align}

\begin{theorem}[Occupation tail closes after renormalization]
\label{thm:s7v45-occupation}
The mixing and energy-dependent occupation rows tend to zero.  The
self-energy row \eqref{eq:s7v45-occ-Sigma} is not a tail; after insertion
into the low effective Hamiltonian, its remaining Feshbach debit is
\eqref{eq:s7v45-occ-der}.
\end{theorem}

\begin{proof}
The first two powers tend to zero.  V37 and V42 prove that a closed return
belongs to the effective low Hamiltonian, whereas the derivative of its
energy-dependent part controls the gap-transfer loss.
\end{proof}

\subsection{Chart-exit and IMS rows}
\label{sec:s7v45-chart}

For the Gaussian reference state, V38 gives
\begin{equation}
 \rho_{\rm Gexit}(a)
 \le C e^{-cR^2}
 =
 C e^{-cL^{1/8}}.                                      \label{eq:s7v45-Gexit}
\end{equation}
The fixed-block relative-form and IMS rows are
\begin{align}
 \eta_{\rm chart}(a)&=O(gR)=O(L^{-7/16}),              \label{eq:s7v45-chart}\\
 c_{\rm IMS}(a)&=O(R^{-2})=O(L^{-1/8}).                \label{eq:s7v45-IMS}
\end{align}

\begin{theorem}[Reference chart tail closes locally]
\label{thm:s7v45-chart}
On every fixed contractible gauge-reduced block, the Gaussian chart-exit,
relative coefficient, and IMS residual ratios tend to zero with the rates
\eqref{eq:s7v45-Gexit}--\eqref{eq:s7v45-IMS}.
\end{theorem}

\begin{proof}
Insert \(R=L^{1/16}\) into the V38 Gaussian tail and the V43 relative-form
card.
\end{proof}

\begin{firewall}[Interacting chart exit still needs the activity map]
\label{fire:s7v45-chart}
Gaussian tail decay and local relative-form control do not by themselves
give a normalized interacting chart-exit probability uniformly over all
blocks.  That conclusion belongs to the V44 renormalized activity map or an
anchored IMS/ground-state comparison.  The local rates are proved; their
global connected assembly remains open.
\end{firewall}

\subsection{Toron and boundary-holonomy rows}
\label{sec:s7v45-slow}

Let \(\Xi_{\mathscr B}\) denote the compact space of toron and boundary
holonomy labels retained by a fixed block.  Let
\begin{equation}
 \xi\longmapsto L_{\mathscr B}(\xi)                    \label{eq:s7v45-Lxi}
\end{equation}
be the finite-dimensional nonnegative quadratic oscillator matrix after
gauge reduction.

\begin{lemma}[Continuous slow-label spectrum]
\label{lem:s7v45-continuity}
If the coefficients of \(L_{\mathscr B}(\xi)\) are continuous in \(\xi\),
then its ordered eigenvalues are continuous, and every spectral projection
defined by a contour in a common resolvent set depends continuously on
\(\xi\).
\end{lemma}

\begin{proof}
For finite Hermitian matrices, the min--max formula bounds each ordered
eigenvalue difference by
\(\|L_{\mathscr B}(\xi)-L_{\mathscr B}(\xi')\|\).
Resolvent continuity follows from the resolvent identity, and contour
integration gives continuity of the Riesz projection.
\end{proof}

\begin{theorem}[Finite buffered cover of the slow-label space]
\label{thm:s7v45-cover}
For every \(\xi_0\in\Xi_{\mathscr B}\), choose numbers
\begin{equation}
 0<\alpha_{\xi_0}<\beta_{\xi_0}                        \label{eq:s7v45-buffer}
\end{equation}
whose closed interval contains no eigenvalue of
\(L_{\mathscr B}(\xi_0)\).  Then there is a neighborhood
\({\cal U}_{\xi_0}\) on which the same interval remains in the resolvent set,
the slow Riesz projection below it has constant rank, and the complement
has spectral floor at least \(\beta_{\xi_0}\).  Compactness of
\(\Xi_{\mathscr B}\) gives a finite collection of such patches.
\end{theorem}

\begin{proof}
The distance from the finite spectrum at \(\xi_0\) to
\([\alpha_{\xi_0},\beta_{\xi_0}]\) is positive.  Operator-norm continuity
preserves half that distance on a neighborhood.  The Riesz projection then
varies continuously and hence has locally constant integer rank.  Its
complement contains only eigenvalues above the upper contour edge.  A finite
subcover follows from compactness.
\end{proof}

\begin{firewall}[No single fixed-rank chart need cover every toron]
\label{fire:s7v45-rank}
Eigenvalues can cross a proposed band edge as \(\xi\) varies.  At such a
crossing the slow rank changes, and a fixed-rank spectral projection cannot
remain continuous.  The transition mode must be retained on the overlap of
buffered patches.  Declaring it gauge or assigning it to the residual would
destroy the uniform complement floor.
\end{firewall}

\begin{definition}[Residual-tail ledger, RTL]
\label{def:s7v45-RTL}
For every cutoff and fixed local observable support, RTL records:
\begin{enumerate}[label=\textup{(T\arabic*)},leftmargin=4em]
\item separated filter leakage \(\rho_{\rm filt}\);
\item high-momentum transfer \(\rho_{\rm mom}\);
\item high-occupation mixing and energy dependence;
\item Gaussian and interacting chart-exit rows;
\item the IMS relative residual row;
\item a finite buffered cover of toron/boundary slow labels;
\item weighted covariance constants uniform on each patch;
\item transition-mode rules on patch overlaps; and
\item the sum of genuinely residual rows
\begin{equation}
 \rho_{\rm tail}
 :=
 \rho_{\rm filt}+\rho_{\rm mom}
 +\rho_{\rm occ,mix}+\rho_{\rm occ,der}
 +\rho_{\rm exit}+\rho_{\rm IMS}.                      \label{eq:s7v45-tail-sum}
\end{equation}
\end{enumerate}
\end{definition}

\begin{theorem}[Tail classification at V45]
\label{thm:s7v45-tail}
The filter, free momentum, occupation-mixing, occupation-derivative,
Gaussian chart-exit, and local IMS ratios in RTL tend to zero with explicit
rates.  Torons, boundary holonomies, and self-energies are not residual-tail
terms.  The unpaid tail rows are:
\begin{enumerate}[label=\textup{(\roman*)},leftmargin=3.5em]
\item normalized interacting chart exit;
\item weighted residual-covariance bounds uniform on every slow-label patch;
\item compatible transition-sector estimates on patch overlaps; and
\item conversion of these local rows into the V44 renormalized activity map.
\end{enumerate}
\end{theorem}

\begin{proof}
Use
\Cref{thm:s7v45-momentum,thm:s7v45-occupation,
thm:s7v45-chart,thm:s7v45-cover}.  The exclusions follow from the Feshbach
ledger: a self-energy modifies the low Hamiltonian, while toron and boundary
variables parametrize that Hamiltonian.
\end{proof}

\section{Audit-guided next step}
\label{sec:s7v45-next}

\begin{assessment}[Progress at seventh-series V45]
\label{ass:s7v45-status}
The compulsory audit has converted the abstract Gaussian tree edge into a
concrete weighted-Schur target and has removed four quantities from the
undifferentiated tail symbol.  High momentum, high occupation, Gaussian
chart exit, and local IMS ratios have compatible schedules.  The genuine
new bottleneck is uniform slow-label locality: quadratic ranks change over
the toron/boundary parameter space, so the residual covariance must be
constructed patchwise with transition modes retained.
\end{assessment}

\begin{decision}[V46 acquisition order]
\label{dec:s7v45-V46}
V46 will prove a parameter-uniform Combes--Thomas estimate on one buffered
slow-label patch.  It will start from a finite-range positive quadratic
kernel with complement floor \(\beta>0\), conjugate by an exponential
spatial weight, and derive an explicit admissible decay rate
\(\mu<\mu(\beta,\|L\|_{\rm hop})\).  The estimate will then be tested under
the smooth analytic filters and patch transitions of
\Cref{thm:s7v45-cover}.
\end{decision}

\begin{firewall}[Claim boundary after seventh-series V45]
\label{fire:s7v45-claim}
V45 performs the compulsory companion audit, proves the weighted inverse
criterion, records the explicit residual-tail rates, and constructs a finite
buffered cover of the compact slow-label space.

V45 does not prove the parameter-uniform Yang--Mills Combes--Thomas bound,
the interacting chart-exit activity, patch-overlap connected assembly,
coarse low matching, a terminal nonperturbative gap, cutoff removal, or the
full Yang--Mills mass gap.
\end{firewall}

\paragraph{Parent-version record.}
V45 was formed from
\texttt{Renewal\_Yang\_Mills\_Mass_Gap\_Research\_Notes\_V44.tex}.
The V44 parent had \(18\,723\) logical source lines, \(680\,217\) bytes,
and SHA--256
\[
 \texttt{6B73C076256E95882EF353596A84988B2E28489222F4478879B89C59CE6C69E5}.
\]
The next compulsory companion audit is V50.

% =============================================================================
% END APPENDED SEVENTH-SERIES V45 MATERIAL
% =============================================================================

% =============================================================================
% BEGIN APPENDED SEVENTH-SERIES V46 MATERIAL
% =============================================================================

\chapter{Uniform buffered Combes--Thomas locality}
\label{ch:s7v46}

V45 identifies a two-sided weighted residual covariance as the missing
Gaussian input to the V44 tree expansion.  A spectral gap gives an
unweighted inverse but not spatial decay.  This note proves the required
decay for a finite-range quadratic kernel, uniformly on one compact
slow-label patch.  The correct object is a buffered high-band analytic
functional calculus.  A raw formula that subtracts a changing low
projection can inherit nonlocal slow eigenfunctions; a contour enclosing
only the high spectrum avoids that mistake.

\section{Exponential conjugation of a finite-range kernel}
\label{sec:s7v46-conjugation}

Let \((X,d)\) be a finite or countable metric graph, with finite-dimensional
fibres \(E_x\).  Let \(L=(L_{xy})\) be bounded and self-adjoint on
\(\ell^2\bigoplus_xE_x\).  Assume
\begin{equation}
 L_{xy}=0\quad\hbox{if }d(x,y)>R_0                    \label{eq:s7v46-range}
\end{equation}
and define the hopping row bound
\begin{equation}
 J:=
 \max\left\{
 \sup_x\sum_{y\ne x}\|L_{xy}\|,
 \sup_y\sum_{x\ne y}\|L_{xy}\|
 \right\}.                                             \label{eq:s7v46-J}
\end{equation}

Fix \(x_0\in X\), set
\[
 \varphi_{x_0}(x):=d(x,x_0),
\]
and let \(W_{\mu,x_0}\) multiply the \(x\)-fibre by
\(e^{\mu\varphi_{x_0}(x)}\).

\begin{lemma}[Conjugation debit]
\label{lem:s7v46-conjugation}
For \(\mu\ge0\),
\begin{equation}
 B_{\mu,x_0}:=
 W_{\mu,x_0}LW_{\mu,x_0}^{-1}-L                       \label{eq:s7v46-B}
\end{equation}
extends to a bounded operator and
\begin{equation}
 \|B_{\mu,x_0}\|
 \le J\bigl(e^{\mu R_0}-1\bigr).                       \label{eq:s7v46-B-bound}
\end{equation}
The bound is independent of \(x_0\) and volume.
\end{lemma}

\begin{proof}
The block of \(B_{\mu,x_0}\) is
\[
 (B_{\mu,x_0})_{xy}
 =
 \left(
 e^{\mu(\varphi_{x_0}(x)-\varphi_{x_0}(y))}-1
 \right)L_{xy}.
\]
It vanishes on the diagonal.  If \(L_{xy}\ne0\), the triangle inequality and
\eqref{eq:s7v46-range} give
\[
 |\varphi_{x_0}(x)-\varphi_{x_0}(y)|
 \le d(x,y)\le R_0.
\]
Thus the absolute multiplier is at most
\(e^{\mu R_0}-1\).  The row and column sums are bounded by the right side of
\eqref{eq:s7v46-B-bound}; the Schur test proves the operator estimate.
\end{proof}

\section{Combes--Thomas resolvent estimate}
\label{sec:s7v46-resolvent}

\begin{theorem}[Uniform finite-range Combes--Thomas bound]
\label{thm:s7v46-CT}
Let \(z\in\mathbb C\) satisfy
\begin{equation}
 \operatorname{dist}(z,\operatorname{spec}L)\ge\delta>0.
 \label{eq:s7v46-distance}
\end{equation}
Choose \(\mu>0\) so that
\begin{equation}
 \varepsilon_\mu
 :=
 J(e^{\mu R_0}-1)<\delta.                              \label{eq:s7v46-mu}
\end{equation}
Then
\begin{equation}
 \|(L-z)^{-1}_{xy}\|
 \le
 {e^{-\mu d(x,y)}\over\delta-\varepsilon_\mu}.         \label{eq:s7v46-CT}
\end{equation}
\end{theorem}

\begin{proof}
The unweighted spectral theorem gives
\(\|(L-z)^{-1}\|\le\delta^{-1}\).  By
\Cref{lem:s7v46-conjugation},
\[
 W_{\mu,x_0}(L-z)W_{\mu,x_0}^{-1}
 =
 L-z+B_{\mu,x_0}.
\]
The resolvent identity and Neumann series show that the inverse exists and
\begin{equation}
 \|W_{\mu,x_0}(L-z)^{-1}W_{\mu,x_0}^{-1}\|
 \le {1\over\delta-\varepsilon_\mu}.                   \label{eq:s7v46-weighted-resolvent}
\end{equation}
On an infinite graph this argument is first applied to the bounded weight
\[
 \exp\{\mu\min(d(x,x_0),N)\}.
\]
Its logarithm is still one-Lipschitz, so all bounds are independent of
\(N\).  Fixed matrix elements stabilize to the untruncated weighted formula
as \(N\to\infty\).
Set \(x_0=y\), sandwich between the \(x\)- and \(y\)-fibre inclusions, and
move the weights to the scalar side.  Since
\(\varphi_y(y)=0\) and \(\varphi_y(x)=d(x,y)\),
\eqref{eq:s7v46-CT} follows.
\end{proof}

\begin{corollary}[Explicit admissible rate]
\label{cor:s7v46-rate}
If \(J>0\), every
\begin{equation}
 0<\mu<
 {1\over R_0}\log\left(1+{\delta\over J}\right)        \label{eq:s7v46-rate}
\end{equation}
is admissible.  Choosing
\[
 \mu_*={1\over R_0}\log\left(1+{\delta\over2J}\right)
\]
gives \(\varepsilon_{\mu_*}=\delta/2\) and prefactor \(2/\delta\).
\end{corollary}

\begin{proof}
Solve \eqref{eq:s7v46-mu} for \(\mu\).
\end{proof}

\section{From pointwise decay to a weighted Schur bound}
\label{sec:s7v46-Schur}

Assume polynomial volume growth
\begin{equation}
 \#B(x,r)\le C_{\rm vol}(1+r)^D.                       \label{eq:s7v46-growth}
\end{equation}

\begin{lemma}[Loss of decay rate in the Schur sum]
\label{lem:s7v46-Schur}
If
\[
 \|K_{xy}\|\le C_Ke^{-\mu_*d(x,y)},
\]
then, for every \(0\le\mu<\mu_*\),
\begin{equation}
 \|K\|_{{\cal S}_\mu}
 \le
 C_KC(D,C_{\rm vol},\mu_*-\mu).                        \label{eq:s7v46-Schur}
\end{equation}
The constant is independent of volume.
\end{lemma}

\begin{proof}
The weighted row sum is bounded by
\[
 C_K\sum_y e^{-(\mu_*-\mu)d(x,y)}.
\]
Decompose into distance shells and apply
\eqref{eq:s7v46-growth}; the resulting polynomial times exponential series
converges uniformly in \(x\).  The column argument is identical.
\end{proof}

\begin{corollary}[Weighted resolvent]
\label{cor:s7v46-weighted-resolvent}
Under \Cref{thm:s7v46-CT} and
\eqref{eq:s7v46-growth}, for every \(0\le\mu<\mu_*\),
\begin{equation}
 \|(L-z)^{-1}\|_{{\cal S}_\mu}
 \le
 {2C(D,C_{\rm vol},\mu_*-\mu)\over\delta}.             \label{eq:s7v46-resolvent-Schur}
\end{equation}
\end{corollary}

\begin{proof}
Use \Cref{thm:s7v46-CT} at rate \(\mu_*\) and then
\Cref{lem:s7v46-Schur}.
\end{proof}

\section{Buffered high-band functional calculus}
\label{sec:s7v46-functional}

Let \(\xi\) range over one slow-label patch \({\cal U}\).  Suppose
\(L(\xi)\) is a norm-continuous family satisfying
\eqref{eq:s7v46-range}--\eqref{eq:s7v46-J} with common \(R_0,J\), and
\begin{equation}
 \operatorname{spec}L(\xi)
 \subset[0,\alpha]\cup[\beta,\Lambda],
 \qquad 0<\alpha<\beta<\Lambda                         \label{eq:s7v46-buffer}
\end{equation}
for every \(\xi\in{\cal U}\).

Choose a positively oriented contour \(\Gamma_{\rm H}\) enclosing
\([\beta,\Lambda]\), excluding \([0,\alpha]\), and remaining at distance at
least \(\delta_\Gamma>0\) from both spectral sets.  Define
\begin{align}
 P_{\rm H}(\xi)
 &:=
 {1\over2\pi i}
 \int_{\Gamma_{\rm H}}
 (z-L(\xi))^{-1}\,dz,                                  \label{eq:s7v46-PH}\\
 C_{\rm H}(\xi)
 &:=
 {1\over2\pi i}
 \int_{\Gamma_{\rm H}}
 z^{-1/2}(z-L(\xi))^{-1}\,dz,                          \label{eq:s7v46-CH}
\end{align}
where one branch of \(z^{-1/2}\) is fixed near the high contour.

\begin{theorem}[Uniform high-band locality]
\label{thm:s7v46-high}
There is \(\mu_{\cal U}>0\), depending only on
\(R_0,J,\delta_\Gamma\), such that for every
\(0\le\mu<\mu_{\cal U}\),
\begin{equation}
 \sup_{\xi\in{\cal U}}
 \left(
 \|P_{\rm H}(\xi)\|_{{\cal S}_\mu}
 +
 \|C_{\rm H}(\xi)\|_{{\cal S}_\mu}
 \right)<\infty.                                      \label{eq:s7v46-high}
\end{equation}
The constants are independent of finite periodic volume when the range,
hopping, buffer, contour length, and volume-growth constants are uniform.
\end{theorem}

\begin{proof}
Apply \Cref{cor:s7v46-weighted-resolvent} uniformly to every
\(z\in\Gamma_{\rm H}\), with
\(\delta=\delta_\Gamma\).  The contour has finite length, and
\(z^{-1/2}\) is bounded on it.  Integrating in the complete weighted Schur
space gives the two bounds.
\end{proof}

\begin{corollary}[The V45 Gaussian edge closes on one patch]
\label{cor:s7v46-V45}
If the eliminated Gaussian covariance is the high-band operator
\(C_{\rm H}(\xi)\), then the V45 weighted covariance target and the V44
two-point Gaussian tree edge hold uniformly for
\(\xi\in{\cal U}\).
\end{corollary}

\begin{proof}
\Cref{thm:s7v46-high} is the required two-sided weighted Schur bound.
Insert it into the Gaussian interpolation estimate of V44.
\end{proof}

\begin{firewall}[Why the raw reduced inverse is not used]
\label{fire:s7v46-reduced}
Writing
\[
 (I-P_{\rm L})L^{-1}(I-P_{\rm L})
\]
as a difference between a full inverse and a low-mode contribution is
unsafe when \(L\) has zero modes or extended slow eigenvectors.  The two
terms may be separately singular or nonlocal.  The contour
\eqref{eq:s7v46-CH} defines the high inverse directly and never estimates a
subtracted slow kernel.
\end{firewall}

\section{Patch overlaps and common refinement}
\label{sec:s7v46-overlap}

Suppose two slow-label patches \({\cal U}_1,{\cal U}_2\) retain spectral
subspaces \(P_{{\rm L},1}\) and \(P_{{\rm L},2}\).  On an overlap, let
\(P_{\rm L}^{\rm ref}\) be the spectral projection retaining every
eigenvalue retained by either patch, including all modes in either
transition buffer.

\begin{theorem}[Overlap refinement preserves the residual floor]
\label{thm:s7v46-refinement}
If the projections are spectral functions of the same \(L(\xi)\), then they
commute and
\begin{equation}
 P_{\rm L}^{\rm ref}
 =
 P_{{\rm L},1}\vee P_{{\rm L},2}
 =
 P_{{\rm L},1}+P_{{\rm L},2}
 -P_{{\rm L},1}P_{{\rm L},2}.                         \label{eq:s7v46-refinement}
\end{equation}
The common residual
\[
 P_{\rm H}^{\rm ref}:=I-P_{\rm L}^{\rm ref}
\]
is contained in both original residual spaces and hence inherits the larger
of the two declared lower spectral thresholds.  If a common high contour
exists, \Cref{thm:s7v46-high} applies on the overlap.
\end{theorem}

\begin{proof}
Spectral projections of one self-adjoint matrix commute, and the displayed
formula is the orthogonal projection onto the sum of their ranges.  Its
orthogonal complement is the intersection of the two residual ranges.
Every vector there satisfies both residual spectral lower bounds.  The last
assertion is the uniform contour theorem.
\end{proof}

\begin{firewall}[Refinement may enlarge the slow rank]
\label{fire:s7v46-rank}
The overlap construction keeps locality of the residual covariance by
moving ambiguous transition modes into the slow sector.  It does not keep a
fixed slow rank.  Those extra modes require transition maps in the coarse
effective action; they cannot be integrated out merely to simplify chart
bookkeeping.
\end{firewall}

\section{Parameter derivatives}
\label{sec:s7v46-derivative}

\begin{theorem}[Uniform derivative of the high covariance]
\label{thm:s7v46-derivative}
Assume \(L(\xi)\) is differentiable and
\begin{equation}
 \sup_{\xi\in{\cal U}}
 \|\partial_\xi L(\xi)\|_{{\cal S}_\mu}
 \le K_\xi.                                            \label{eq:s7v46-dL}
\end{equation}
Then
\begin{align}
 \partial_\xi C_{\rm H}(\xi)
 =
 {1\over2\pi i}
 \int_{\Gamma_{\rm H}}
 z^{-1/2}(z-L)^{-1}
 (\partial_\xi L)
 (z-L)^{-1}\,dz,                                      \label{eq:s7v46-dC}
\end{align}
and its \({\cal S}_\mu\) norm is uniformly bounded on the patch.
\end{theorem}

\begin{proof}
Differentiate the resolvent identity; the contour stays in a common
resolvent set.  The weighted Schur algebra is submultiplicative, and each
resolvent factor has the uniform bound from
\Cref{thm:s7v46-high}.  Integrate the resulting product norm.
\end{proof}

\begin{corollary}[Uniform slow-label Gaussian seminorms]
\label{cor:s7v46-slow-uniform}
For local polynomial observables whose coefficient derivatives in \(\xi\)
are uniformly bounded, the Gaussian interpolation constants and their first
slow-label derivatives are uniform on \({\cal U}\).
\end{corollary}

\begin{proof}
Wick contractions and Gaussian integration by parts use only finitely many
local coefficient norms and covariance factors.  Apply
\Cref{thm:s7v46-high,thm:s7v46-derivative}.
\end{proof}

\section{Combes--Thomas status}
\label{sec:s7v46-status}

\begin{theorem}[What V46 proves]
\label{thm:s7v46-result}
On every buffered compact slow-label patch, a finite-range quadratic kernel
with uniform hopping bound has a volume-uniform exponentially localized
high-band projection and high-band inverse square root.  The decay rate may
be chosen explicitly below
\[
 {1\over R_0}
 \log\left(1+{\delta_\Gamma\over2J}\right).
\]
Common refinement on patch overlaps retains ambiguous modes and preserves
the residual floor; differentiable slow-label dependence preserves the
weighted covariance bounds.
\end{theorem}

\begin{proof}
Use
\Cref{thm:s7v46-CT,cor:s7v46-weighted-resolvent,
thm:s7v46-high,thm:s7v46-refinement,
thm:s7v46-derivative}.
\end{proof}

\begin{assessment}[Progress at seventh-series V46]
\label{ass:s7v46-status}
The residual Gaussian tree edge assumed in V44 is now proved at the
quadratic level on each buffered toron/boundary patch, with explicit
volume-independent decay.  Patch overlaps are handled by enlarging the slow
sector, not by subtracting nonlocal eigenfunctions.  The remaining
connected-assembly inputs are nonlinear derivative seminorms on the
occupation/chart window, interacting chart exit, and compatibility of the
patchwise slow effective actions.
\end{assessment}

\begin{decision}[V47 acquisition order]
\label{dec:s7v46-V47}
V47 will combine the V43 polynomial number-window estimates with the V46
weighted covariance to prove a finite-order Gaussian cumulant tree bound
uniform in slow labels.  It will identify how the constants grow with
cumulant order and whether the resulting factorial is compatible with the
\(1/n!\) in the normalized V44 expansion.
\end{decision}

\begin{firewall}[Claim boundary after seventh-series V46]
\label{fire:s7v46-claim}
V46 proves finite-range Combes--Thomas decay, weighted Schur resolvent
bounds, uniform high-band functional calculus, patch-overlap refinement, and
slow-label derivative bounds.

V46 does not prove all-order Wilson cumulant seminorms, interacting chart
exit, patchwise coarse-action cocycles, terminal strong-coupling matching,
cutoff removal, or the full Yang--Mills mass gap.
\end{firewall}

\paragraph{Parent-version record.}
V46 was formed from
\texttt{Renewal\_Yang\_Mills\_Mass_Gap\_Research\_Notes\_V45.tex}.
The V45 parent had \(19\,144\) logical source lines, \(696\,851\) bytes,
and SHA--256
\[
 \texttt{DAD2A182C2C29CC49387067997453E391E77B2E22458032069A94DF3B40516CD}.
\]
The next compulsory companion audit is V50.

% =============================================================================
% END APPENDED SEVENTH-SERIES V46 MATERIAL
% =============================================================================

% =============================================================================
% BEGIN APPENDED SEVENTH-SERIES V47 MATERIAL
% =============================================================================

\chapter{The zero-radius obstruction and the bounded-window repair}
\label{ch:s7v47}

\section{Why the acquisition order must change}
\label{sec:s7v47-context}

V46 ended by proposing an all-order Gaussian cumulant estimate for the
untruncated Wilson polynomial.  Before attempting the spatial tree estimate,
one must check that the perturbative series exists as a convergent series at
all.  It does not, already for a one-dimensional positive quartic
perturbation.  This is not a technical defect of a particular tree formula:
it is an analytic obstruction at the origin.

The correction made in this note is therefore upstream.  The interaction is
first compressed to the growing occupation window already used in
V42--V43.  On that window it is a bounded operator whose norm tends to zero
under the established schedule.  Its complement is returned by the
Feshbach/IMS mechanism, rather than being expanded into unbounded Gaussian
moments.  V46 supplies the spatially decaying covariance needed for the
remaining bounded-window connected estimate.

\section{A zero-radius theorem}
\label{sec:s7v47-zero}

Let \(X\) be a standard real Gaussian and, for \(\lambda\geq0\), put
\begin{equation}
 Z_4(\lambda):=
 {\mathbb E}\exp(-\lambda X^4).
 \label{eq:s7v47-Z4}
\end{equation}

\begin{theorem}[The positive quartic Gaussian integral is not analytic at
the origin]
\label{thm:s7v47-zero}
The formal Taylor series of \(Z_4\) at \(0\) is
\begin{equation}
 \sum_{n\geq0}{(-\lambda)^n\over n!}\,
 {\mathbb E}X^{4n},
 \label{eq:s7v47-formal}
\end{equation}
and it has radius of convergence zero.  Consequently neither \(Z_4\) nor
\(\log Z_4\) admits a holomorphic extension to a complex neighbourhood of
\(\lambda=0\).
\end{theorem}

\begin{proof}
For every fixed \(n\), differentiation from the positive real half-line is
permitted at \(\lambda=0\): the difference quotients are dominated after
one uses the integral remainder for the exponential, or equivalently one
first differentiates for \(\lambda>0\) and applies monotone domination to
the nonnegative absolute value.  Thus the formal coefficient in absolute
value is
\[
 c_n={{\mathbb E}X^{4n}\over n!}.
\]
Gaussian moments give
\[
 {\mathbb E}X^{4n}=(4n-1)!!,
\]
and hence
\begin{equation}
 {c_{n+1}\over c_n}
 ={(4n+1)(4n+3)\over n+1}
 \longrightarrow\infty.
 \label{eq:s7v47-ratio}
\end{equation}
The ratio test proves that \eqref{eq:s7v47-formal} has radius zero.

If \(Z_4\) had a holomorphic extension through \(0\), its Taylor
coefficients would be the derivatives just computed, contradicting the
zero radius.  If \(\log Z_4\) had such an extension, its exponential would
give a holomorphic extension of \(Z_4\), because \(Z_4(0)=1\).  This gives
the second assertion.
\end{proof}

\begin{corollary}[The untruncated V44 expansion is not an all-order
argument]
\label{cor:s7v47-V44}
An absolute all-order justification of the V44 cumulant expansion cannot be
deduced merely from small polynomial Wilson coefficients and Gaussian
moments on the full field space.
\end{corollary}

\begin{proof}
The one-site choice \(v=\lambda X^4\) is a special case of precisely that
proposed mechanism.  Its partition function already fails to be analytic
at zero.
\end{proof}

\begin{firewall}[What the obstruction does and does not say]
\label{fire:s7v47-zero}
\Cref{thm:s7v47-zero} invalidates a convergent untruncated weak-coupling
power series.  It does not say that the positive quartic measure fails to
exist, that its correlations fail to decay, or that every resummation must
fail.  In particular it leaves constructive expansions with large-field
control open.  The present programme uses the previously constructed
occupation/chart decomposition instead.
\end{firewall}

\section{A universal bounded-cumulant budget}
\label{sec:s7v47-bounded}

For bounded random variables \(F_1,\ldots,F_m\), write
\(\kappa(F_1,\ldots,F_m)\) for their joint cumulant.  Let
\({m\brace k}\) denote a Stirling number of the second kind and set
\begin{equation}
 S_m:=\sum_{k=1}^m {m\brace k}(k-1)!.
 \label{eq:s7v47-Sm}
\end{equation}

\begin{lemma}[Factorial bound for bounded cumulants]
\label{lem:s7v47-cumulant}
For \(m\geq1\),
\begin{equation}
 |\kappa(F_1,\ldots,F_m)|
 \leq S_m\prod_{j=1}^m\|F_j\|_\infty
 \leq e^m m!\prod_{j=1}^m\|F_j\|_\infty .
 \label{eq:s7v47-cumulant}
\end{equation}
\end{lemma}

\begin{proof}
The moment--cumulant formula is
\begin{equation}
 \kappa(F_1,\ldots,F_m)
 =
 \sum_{\pi\in\Pi_m}
 (|\pi|-1)!(-1)^{|\pi|-1}
 \prod_{B\in\pi}{\mathbb E}\prod_{j\in B}F_j.
 \label{eq:s7v47-moment-cumulant}
\end{equation}
Taking absolute values and sorting partitions by their number \(k\) of
blocks gives the first inequality.  The number
\(k!{m\brace k}\) counts surjections from an \(m\)-element labelled set to
a \(k\)-element labelled set and is at most \(k^m\).  Therefore
\[
 S_m\leq\sum_{k=1}^m k^{m-1}\leq m^m.
\]
Finally \(m!\geq(m/e)^m\), which proves \(m^m\leq e^m m!\).
\end{proof}

\begin{corollary}[Compatibility with the normalized \(1/n!\)]
\label{cor:s7v47-factorial}
If \(A,B\) are bounded and \(v_1,\ldots,v_n\) satisfy
\(\|v_j\|_\infty\leq\rho\), then
\begin{equation}
 {1\over n!}
 |\kappa(A,B,v_1,\ldots,v_n)|
 \leq
 e^{n+2}(n+2)(n+1)
 \|A\|_\infty\|B\|_\infty\rho^n.
 \label{eq:s7v47-factorial}
\end{equation}
In particular the purely order-dependent majorant is summable whenever
\(e\rho<1\).
\end{corollary}

\begin{proof}
Apply \Cref{lem:s7v47-cumulant} with \(m=n+2\) and divide by \(n!\).
The remaining factor is
\((n+2)!/n!=(n+2)(n+1)\).
\end{proof}

\begin{firewall}[Order control is not spatial control]
\label{fire:s7v47-order}
The estimate \eqref{eq:s7v47-factorial} controls the growth with cumulant
order.  Applied to the extensive sum \(\sum_Xv_X\), its sup norm would grow
with volume and the estimate would be useless.  A volume-uniform proof must
retain the individual supports and recover a connected tree joining them.
That separate spatial step is the target left after this note.
\end{firewall}

\section{Operator bounds on the occupation window}
\label{sec:s7v47-window}

Let \({\cal N}\) be the total oscillator number operator of one fixed
finite block, and let
\begin{equation}
 P_M:=\mathbf 1_{[0,M]}({\cal N}),\qquad Q_M:=I-P_M.
 \label{eq:s7v47-PM}
\end{equation}
The number of oscillator components per block is fixed independently of
the ultraviolet and volume parameters.

\begin{lemma}[Words in creation and annihilation operators]
\label{lem:s7v47-word}
If \(W_r=b_1\cdots b_r\), where every \(b_j\) is a creation or annihilation
operator of one of the fixed block modes, then
\begin{equation}
 \|P_M W_r P_M\|\leq(M+r)^{r/2}.
 \label{eq:s7v47-word}
\end{equation}
The same estimate, with a fixed coefficient depending only on the number
of block modes, holds for a fixed finite sum of such words.
\end{lemma}

\begin{proof}
On the \(n\)-particle subspace an annihilation operator has norm at most
\(\sqrt n\), while a creation operator has norm at most \(\sqrt{n+1}\)
into the next subspace.  Starting from a vector in \(P_M{\cal H}\), every
intermediate occupation is at most \(M+r\).  Multiplication of the \(r\)
one-step bounds gives \((M+r)^{r/2}\).  The triangle inequality handles a
fixed finite sum.
\end{proof}

Write the block-local nonlinear Wilson remainder, after the V41 global
Gaussian part has been extracted, as
\begin{equation}
 V_{\rm nl}(g)=gV_3+g^2V_4+V_{\geq5}(g),
 \label{eq:s7v47-Vnl}
\end{equation}
where the coefficients of the finitely many cubic and quartic words are
uniform on a buffered slow-label patch.

\begin{proposition}[Bounded-window Wilson activity]
\label{prop:s7v47-activity}
There are patch-uniform constants \(C_3,C_4\) such that
\begin{equation}
 \|P_MV_{\rm nl}(g)P_M\|
 \leq
 C_3g(M+3)^{3/2}
 +C_4g^2(M+4)^2
 +\rho_{\geq5}(g,M),
 \label{eq:s7v47-activity}
\end{equation}
where
\(\rho_{\geq5}(g,M):=\|P_MV_{\geq5}(g)P_M\|\).
For the schedule
\begin{equation}
 g=L^{-1/2},\qquad M=L^{1/8},
 \label{eq:s7v47-schedule}
\end{equation}
the displayed cubic and quartic contributions are respectively
\begin{equation}
 O(L^{-5/16}),\qquad O(L^{-3/4}).
 \label{eq:s7v47-rates}
\end{equation}
\end{proposition}

\begin{proof}
Apply \Cref{lem:s7v47-word} term by term to the cubic and quartic
polynomials.  The fixed number of terms and their uniformly bounded
coefficients are absorbed in \(C_3,C_4\).  Substitution of
\eqref{eq:s7v47-schedule} gives
\[
 L^{-1/2}(L^{1/8})^{3/2}=L^{-5/16},
 \qquad
 L^{-1}(L^{1/8})^2=L^{-3/4}.
\]
Adding fixed constants to \(M\) does not change these powers.
\end{proof}

\begin{corollary}[A genuine order-smallness condition]
\label{cor:s7v47-small}
If
\begin{equation}
 \rho_{\geq5}(L^{-1/2},L^{1/8})=o(1),
 \label{eq:s7v47-tail}
\end{equation}
then, for all sufficiently large \(L\), the compressed interaction norm
\(\rho_M:=\|P_MV_{\rm nl}P_M\|\) satisfies \(e\rho_M<1\).
Consequently the bounded cumulant order majorant in
\Cref{cor:s7v47-factorial} is summable.
\end{corollary}

\begin{proof}
\Cref{prop:s7v47-activity} and \eqref{eq:s7v47-tail} imply
\(\rho_M\to0\).
\end{proof}

\begin{firewall}[The higher Wilson tail remains a named estimate]
\label{fire:s7v47-tail}
Equation \eqref{eq:s7v47-tail} is not inferred from formal analyticity.
It must follow from a uniform remainder theorem for the Wilson chart on
the chosen block window.  V43 supplies the finite-degree relative-form
mechanism but not yet an all-degree, block-uniform proof of
\eqref{eq:s7v47-tail}.
\end{firewall}

\section{Returning the discarded occupation complement}
\label{sec:s7v47-Feshbach}

The use of \(P_M\) is not a replacement of the theory by a finite
dimensional model.  It is one half of the nested decomposition in V42.
For clarity, the elementary operator estimate behind that return is
recorded here.

\begin{lemma}[Norm control of the occupation Feshbach correction]
\label{lem:s7v47-Feshbach}
Let \(H\) be self-adjoint on
\({\cal H}=P{\cal H}\oplus Q{\cal H}\).  If, for a real \(z\),
\begin{equation}
 Q(H-z)Q\geq\Delta Q,\qquad \Delta>0,
 \label{eq:s7v47-Qgap}
\end{equation}
and \(\|PHQ\|\leq b\), then the Feshbach correction
\begin{equation}
 \Sigma_P(z):=
 PHQ\,[Q(H-z)Q]^{-1}QHP
 \label{eq:s7v47-Sigma}
\end{equation}
satisfies
\begin{equation}
 0\leq\Sigma_P(z)\leq {b^2\over\Delta}P.
 \label{eq:s7v47-Sigma-bound}
\end{equation}
\end{lemma}

\begin{proof}
The assumption gives
\(0\leq[Q(H-z)Q]^{-1}\leq\Delta^{-1}Q\).
Sandwiching by \(PHQ\) and \(QHP\) proves positivity and the stated norm
bound.
\end{proof}

\begin{proposition}[Logical separation of the two estimates]
\label{prop:s7v47-separation}
Suppose the V42 high-occupation hypothesis provides
\(\Delta_M>0\) and \(b_M^2/\Delta_M=o(1)\), uniformly in finite volume and
slow label.  Then:
\begin{enumerate}
\item the \(P_M\)-sector interaction can be treated by the genuinely
bounded small parameter of \Cref{cor:s7v47-small};
\item the \(Q_M\)-sector changes the compressed operator by at most
\(b_M^2/\Delta_M\);
\item no untruncated polynomial moment expansion is used in either step.
\end{enumerate}
\end{proposition}

\begin{proof}
The first item is \Cref{cor:s7v47-small}; the second is
\Cref{lem:s7v47-Feshbach} with \(P=P_M\); the third follows by inspection:
the first estimate is in operator norm after compression and the second is
a resolvent inequality.
\end{proof}

\begin{firewall}[Hard spectral cutoffs and interpolation derivatives]
\label{fire:s7v47-cutoff}
The spectral projection \(P_M\) is used inside the operator/Feshbach
decomposition.  It should not be silently replaced by a smooth classical
large-field cutoff and then differentiated arbitrarily many times in a
BKAR formula: high cutoff derivatives can recreate factorial growth.
Any spatial interpolation in the next note must either keep \(P_M\) as an
operator compression or state and prove a uniform derivative budget for
the chosen cutoff.
\end{firewall}

\section{Corrected connected-assembly interface}
\label{sec:s7v47-interface}

\begin{theorem}[What V47 proves]
\label{thm:s7v47-result}
The following statements now hold.
\begin{enumerate}
\item The naive untruncated weak-coupling expansion has a genuine
zero-radius obstruction.
\item On the occupation window, cumulant coefficients have at most the
explicit factorial growth in \eqref{eq:s7v47-factorial}, which is compatible
with the normalized \(1/n!\) series when the compressed local activity is
small.
\item The cubic and quartic window activities tend to zero at the explicit
rates \(L^{-5/16}\) and \(L^{-3/4}\).
\item The high-occupation complement is assigned to a resolvent estimate,
with correction bounded by \(b_M^2/\Delta_M\), rather than to polynomial
Gaussian moments.
\end{enumerate}
\end{theorem}

\begin{proof}
Use
\Cref{thm:s7v47-zero,cor:s7v47-factorial,
prop:s7v47-activity,lem:s7v47-Feshbach}.
\end{proof}

\begin{assessment}[Progress at seventh-series V47]
\label{ass:s7v47-status}
The order-growth part of the V46 acquisition problem is settled, but only
after correcting the domain on which the expansion is performed.  The
remaining central step is spatial: combine the V46 exponentially weighted
Gaussian covariance with bounded local window activities to obtain a
volume-uniform connected tree bound.  That proof must keep the support
incidence constants explicit and must not differentiate an uncontrolled
cutoff.
\end{assessment}

\begin{decision}[V48 acquisition order]
\label{dec:s7v47-V48}
V48 will formulate and prove the bounded-variable spatial cumulant tree
estimate from a covariance-decoupling interpolation.  It will separate the
abstract interpolation identity from the Yang--Mills verification and will
record every volume-independent incidence constant.  If the required
higher derivative seminorm cannot be derived from the operator window,
the programme will move upstream to a Duhamel/operator connected
interpolation rather than assuming it.
\end{decision}

\begin{firewall}[Claim boundary after seventh-series V47]
\label{fire:s7v47-claim}
V47 does not prove the all-support connected tree bound, the higher Wilson
tail \eqref{eq:s7v47-tail}, the interacting chart-exit estimate, patchwise
coarse-action cocycles, terminal strong-coupling matching, cutoff removal,
or the full Yang--Mills mass gap.
\end{firewall}

\paragraph{Parent-version record.}
V47 was formed from
\texttt{Renewal\_Yang\_Mills\_Mass\_Gap\_Research\_Notes\_V46.tex}.
The V46 parent had \(19\,580\) logical source lines, \(711\,342\) bytes,
and SHA--256
\[
 \texttt{AC109F612DED0B6262329E2412952CCBD422491661E51096A597139F3F9C808B}.
\]
The next compulsory companion audit is V50.

% =============================================================================
% END APPENDED SEVENTH-SERIES V47 MATERIAL
% =============================================================================

% =============================================================================
% BEGIN APPENDED SEVENTH-SERIES V48 MATERIAL
% =============================================================================

\chapter{Normalized Gaussian mixing and the missing tree lift}
\label{ch:s7v48}

\section{The correct bounded-observable question}
\label{sec:s7v48-context}

V47 removed the order-growth obstruction by compressing the nonlinear
interaction to a bounded occupation window.  Spatial summability still
requires a statement about the reference Gaussian field.  V46 proves an
exponentially weighted bound on its covariance, but a raw covariance
estimate is not by itself a mixing estimate for arbitrary bounded
functions.  The correct dimensionless quantity is the normalized
cross-covariance, also called the canonical-correlation operator.

This note proves the exact two-cluster theorem, derives its uniform
exponential consequence under a local marginal floor, and then shows why
this result cannot silently be promoted to an all-order tree theorem.  The
result both repairs the V46--V47 interface and isolates the next upstream
geometric estimate.

\section{Canonical correlation of Gaussian blocks}
\label{sec:s7v48-canonical}

Let \((U,V)\) be a centred jointly Gaussian vector in finite-dimensional
real Hilbert spaces \(E_U\oplus E_V\), with covariance
\begin{equation}
 {\cal C}=
 \begin{pmatrix}
 A&D\\ D^*&B
 \end{pmatrix}\geq0.
 \label{eq:s7v48-blockcov}
\end{equation}
Zero-variance directions are deleted, or equivalently all inverse square
roots below are Moore--Penrose inverses on the covariance ranges.  Define
\begin{equation}
 K_{UV}:=A^{-1/2}DB^{-1/2},
 \qquad
 \vartheta(U,V):=\|K_{UV}\|.
 \label{eq:s7v48-theta}
\end{equation}
Positivity of \({\cal C}\) implies \(0\leq\vartheta(U,V)\leq1\).

\begin{theorem}[Gaussian maximal-correlation theorem]
\label{thm:s7v48-maxcorr}
For all square-integrable \(f(U)\) and \(g(V)\),
\begin{equation}
 |\operatorname{Cov}(f(U),g(V))|
 \leq
 \vartheta(U,V)
 \sqrt{\operatorname{Var}f(U)}
 \sqrt{\operatorname{Var}g(V)}.
 \label{eq:s7v48-maxcorr}
\end{equation}
The constant \(\vartheta(U,V)\) is optimal.
\end{theorem}

\begin{proof}
Whiten the nondegenerate covariance ranges, so that \(A=B=I\) and the
cross-covariance is \(K:=K_{UV}\).  The Gaussian \(L^2\) spaces have the
orthogonal Wiener-chaos decompositions
\[
 L^2(U)=\bigoplus_{r\geq0}{\cal H}^{U}_r,
 \qquad
 L^2(V)=\bigoplus_{r\geq0}{\cal H}^{V}_r.
\]
Conditional expectation from the \(r\)-th \(V\)-chaos to the \(r\)-th
\(U\)-chaos is the symmetric tensor power \(K^{\otimes_s r}\); different
chaos degrees are orthogonal.  This follows first for Wick monomials from
Wick's formula, where every contraction crossing from \(V\) to \(U\)
contributes one copy of \(K\), and then by density of polynomials.
Consequently its norm on the nonconstant subspace is
\[
 \sup_{r\geq1}\|K\|^r=\|K\|,
\]
because \(\|K\|\leq1\).  If \(f_0=f-{\mathbb E}f\) and
\(g_0=g-{\mathbb E}g\), then
\[
 |\operatorname{Cov}(f,g)|
 =
 |\langle f_0,{\mathbb E}(g_0\mid U)\rangle|
 \leq\|K\|\,\|f_0\|_2\|g_0\|_2.
\]
This is \eqref{eq:s7v48-maxcorr}.  Linear functions in singular vectors
corresponding to the largest singular value of \(K\) attain the constant.
\end{proof}

\begin{corollary}[Bounded two-cluster mixing]
\label{cor:s7v48-bounded}
For bounded \(f,g\),
\begin{equation}
 |\operatorname{Cov}(f(U),g(V))|
 \leq
 \vartheta(U,V)\|f\|_\infty\|g\|_\infty.
 \label{eq:s7v48-bounded}
\end{equation}
\end{corollary}

\begin{proof}
Use
\(\operatorname{Var}f\leq{\mathbb E}|f|^2\leq\|f\|_\infty^2\)
and similarly for \(g\) in \Cref{thm:s7v48-maxcorr}.
\end{proof}

\section{From V46 weighted covariance to canonical mixing}
\label{sec:s7v48-weighted}

Let \(C=(C_{xy})\) be the covariance on a finite graph
\((\Lambda,d)\), with a fixed finite fibre at each \(x\).  Recall the
two-sided weighted Schur norm
\begin{equation}
 \|C\|_{{\cal S}_\mu}
 :=
 \max\left\{
 \sup_x\sum_y e^{\mu d(x,y)}\|C_{xy}\|,
 \sup_y\sum_x e^{\mu d(x,y)}\|C_{xy}\|
 \right\}.
 \label{eq:s7v48-Schur}
\end{equation}
For \(S\subset\Lambda\), \(C_{SS}\) denotes the covariance compressed to
the coordinate fibres over \(S\).

\begin{definition}[Local stochastic floor]
\label{def:s7v48-floor}
For a class \({\mathfrak S}\) of allowed observable supports, a number
\(c_->0\) is a local stochastic floor if
\begin{equation}
 C_{SS}\geq c_- I_S
 \quad\hbox{for every }S\in{\mathfrak S},
 \label{eq:s7v48-floor}
\end{equation}
after deterministic zero-variance directions have been quotiented out.
\end{definition}

\begin{theorem}[Weighted covariance implies bounded local mixing, with a
floor]
\label{thm:s7v48-localmix}
Assume \(\|C\|_{{\cal S}_\mu}\leq K_\mu\) and
\eqref{eq:s7v48-floor}.  If \(S,T\in{\mathfrak S}\) are disjoint and
\(F,G\) are bounded functions of the Gaussian coordinates in \(S,T\),
respectively, then
\begin{equation}
 |\operatorname{Cov}(F,G)|
 \leq
 \min\left\{1,{K_\mu\over c_-}e^{-\mu d(S,T)}\right\}
 \|F\|_\infty\|G\|_\infty.
 \label{eq:s7v48-localmix}
\end{equation}
\end{theorem}

\begin{proof}
The cross block \(D=C_{ST}\) has
\[
 \sup_{x\in S}\sum_{y\in T}\|C_{xy}\|
 \leq K_\mu e^{-\mu d(S,T)}
\]
and the analogous column bound.  The Schur test therefore gives
\[
 \|C_{ST}\|\leq K_\mu e^{-\mu d(S,T)}.
\]
The floor implies
\(\|C_{SS}^{-1/2}\|,\|C_{TT}^{-1/2}\|\leq c_-^{-1/2}\), and hence
\[
 \vartheta(S,T)
 \leq {K_\mu\over c_-}e^{-\mu d(S,T)}.
\]
Combine this with \Cref{cor:s7v48-bounded}; the universal bound
\(\vartheta\leq1\) gives the minimum in
\eqref{eq:s7v48-localmix}.
\end{proof}

\begin{corollary}[Volume-uniform two-cluster interface]
\label{cor:s7v48-uniform}
On any slow-label patch on which the V46 weighted covariance constant and
the local stochastic floor are uniform, bounded observables with allowed
supports obey a volume- and slow-label-uniform exponential covariance
bound.
\end{corollary}

\begin{proof}
All constants in \Cref{thm:s7v48-localmix} are then uniform.
\end{proof}

\section{Why the floor cannot be omitted}
\label{sec:s7v48-necessity}

\begin{proposition}[Small raw covariance need not mix bounded functions]
\label{prop:s7v48-counter}
For every fixed \(r\in(0,1)\) there are jointly Gaussian scalar variables
\((U_\epsilon,V_\epsilon)\) such that
\[
 |\operatorname{Cov}(U_\epsilon,V_\epsilon)|\longrightarrow0
\]
as \(\epsilon\downarrow0\), while the covariance of two bounded functions
of those variables stays strictly positive.
\end{proposition}

\begin{proof}
Let \(X,Y\) be independent standard Gaussians and set
\[
 U_\epsilon=\epsilon X,\qquad
 V_\epsilon=\epsilon(rX+\sqrt{1-r^2}\,Y).
\]
Then
\(\operatorname{Cov}(U_\epsilon,V_\epsilon)=r\epsilon^2\).
On the other hand,
\[
 \operatorname{Cov}(\operatorname{sgn}U_\epsilon,
                     \operatorname{sgn}V_\epsilon)
 ={2\over\pi}\arcsin r,
\]
independently of \(\epsilon\).  The last identity follows by evaluating the
angular fraction of the rotationally invariant \((X,Y)\)-plane on which
the two linear forms have the same sign.
\end{proof}

\begin{assessment}[Correction to the V46 interface]
\label{ass:s7v48-correction}
V46 proves exponentially decaying raw covariance blocks.  For polynomial
Gaussian observables with controlled derivative seminorms that is already
a usable edge estimate.  For the arbitrary bounded observables suggested
by the V47 occupation compression, \Cref{prop:s7v48-counter} shows that a
uniform local stochastic floor, or an equivalent normalized-correlation
estimate, must also be proved.
\end{assessment}

\section{Two-cluster mixing is not yet a tree theorem}
\label{sec:s7v48-treegap}

For \(m\) random variables, the joint cumulant is zero if their joint law
factors across any nontrivial partition.  This connectedness property is
necessary for a tree estimate, but pairwise covariance bounds alone do not
quantify the repeated decouplings needed to obtain \(m-1\) spatial edges.

\begin{proposition}[Pairwise independence does not control a third
cumulant]
\label{prop:s7v48-pairwise}
There exist bounded variables \(R_1,R_2,R_3\) which are pairwise
independent but satisfy
\(\kappa(R_1,R_2,R_3)=1\).
\end{proposition}

\begin{proof}
Let \(R_1,R_2\) be independent symmetric Rademacher variables and set
\(R_3=R_1R_2\).  Every pair is independent and every variable has mean
zero.  Hence
\[
 \kappa(R_1,R_2,R_3)
 ={\mathbb E}R_1R_2R_3=1.
\]
\end{proof}

\begin{firewall}[The example is logical, not Gaussian]
\label{fire:s7v48-pairwise}
The variables in \Cref{prop:s7v48-pairwise} are not Gaussian.  The example
does not refute an all-order Gaussian tree theorem.  It proves that such a
theorem must use more of the Gaussian structure than the two-function
conclusion \eqref{eq:s7v48-localmix}.  In particular, repeatedly citing
two-cluster mixing would leave a gap.
\end{firewall}

\begin{lemma}[What a single cut can prove]
\label{lem:s7v48-onecut}
Let \(\nu\) and \(\nu_0\) be two laws for
\((F_1,\ldots,F_m)\), with \(|F_i|\leq a_i\), and suppose
\(\nu_0\) factors across a nontrivial partition of the indices.  Then
\begin{equation}
 |\kappa_\nu(F_1,\ldots,F_m)|
 \leq
 2\left(\sum_{k=1}^m {m\brace k}k!\right)
 \|\nu-\nu_0\|_{\rm TV}\prod_{i=1}^m a_i
 \leq
 2m^{m+1}\|\nu-\nu_0\|_{\rm TV}\prod_{i=1}^m a_i,
 \label{eq:s7v48-onecut}
\end{equation}
where the total-variation convention is
\(\|\nu-\nu_0\|_{\rm TV}:=\sup_A|\nu(A)-\nu_0(A)|\).
\end{lemma}

\begin{proof}
The cumulant under \(\nu_0\) vanishes by factorization.  Subtract the two
moment--cumulant formulae.  For each partition \(\pi\), telescope the
difference between the products of its block moments.  Every moment has
absolute value at most the product of the corresponding \(a_i\), and for
any bounded \(H\),
\[
 |{\mathbb E}_\nu H-{\mathbb E}_{\nu_0}H|
 \leq2\|\nu-\nu_0\|_{\rm TV}\|H\|_\infty.
\]
The telescoping bound for a partition with \(k\) blocks is therefore at
most \(2k\|\nu-\nu_0\|_{\rm TV}\prod_i a_i\).  Multiplication by its
Möbius coefficient \((k-1)!\), followed by summation over the
\({m\brace k}\) such partitions, gives the first bound.  Finally
\({m\brace k}k!\leq k^m\), so
\[
 \sum_{k=1}^m {m\brace k}k!
 \leq\sum_{k=1}^m k^m\leq m^{m+1}.
\]
\end{proof}

\begin{firewall}[A single cut remains insufficient]
\label{fire:s7v48-onecut}
The conservative bound \eqref{eq:s7v48-onecut} is superfactorial and
contains only one decoupling factor.  It is recorded as a diagnostic, not
as an input to the normalized connected expansion.  Neither defect may be
removed by silently iterating the two-cluster estimate.
\end{firewall}

\section{The precise next target}
\label{sec:s7v48-target}

\begin{definition}[Normalized Gaussian tree lift]
\label{def:s7v48-treelift}
For allowed local supports \(S_1,\ldots,S_m\), a normalized tree lift is an
estimate
\begin{equation}
 |\kappa(F_1,\ldots,F_m)|
 \leq
 C_0^m\prod_{i=1}^m\|F_i\|_\infty
 \sum_{T\in{\mathfrak T}_m}
 \prod_{\{i,j\}\in E(T)}
 e^{-\mu_0 d(S_i,S_j)}
 \label{eq:s7v48-treelift}
\end{equation}
with \(C_0,\mu_0>0\) independent of volume, slow label, and \(m\).
\end{definition}

\begin{proposition}[Why \eqref{eq:s7v48-treelift} is the right interface]
\label{prop:s7v48-interface}
If \eqref{eq:s7v48-treelift} holds for the bounded-window reference
observables and each local interaction insertion has norm at most
\(\rho_M\), then the \(1/n!\) normalized connected expansion is dominated
by an anchored labelled-tree sum with per-added-vertex weight
\(C\rho_M\), where \(C\) depends only on \(C_0,\mu_0\) and the lattice
volume-growth constant.
\end{proposition}

\begin{proof}
Apply \eqref{eq:s7v48-treelift} to \(A,B\) and the \(n\) interaction
insertions.  Root a spanning tree at the support of \(A\), sum successively
over the position of every nonroot vertex, and use
\[
 \sup_x\sum_y e^{-\mu_0d(x,y)}<\infty.
\]
The support-incidence multiplicity is fixed because every interaction has
uniformly bounded range.  There are
\((n+2)^n\) labelled trees on the \(n+2\) vertices.  Since
\[
 n!\geq(n/e)^n,\qquad
 {(n+2)^n\over n!}
 \leq e^n\left(1+{2\over n}\right)^n
 \leq e^{n+2},
\]
the normalized \(1/n!\) absorbs the labelled-tree count up to a constant
to the power \(n\).  Thus every added interaction carries \(C\rho_M\).
\end{proof}

\begin{firewall}[Why no extra factorial is allowed]
\label{fire:s7v48-interface}
The bound \eqref{eq:s7v48-treelift} deliberately has no prefactor \(m!\).
If one multiplied its labelled-tree sum by \(m!\), the single
normalization \(1/n!\) would leave Cayley's superexponential tree count
uncancelled.  Thus the universal bounded-cumulant estimate of V47 settles
only the nonspatial order diagnostic; the spatial forest identity must
produce the sharper normalized tree structure stated here.
\end{firewall}

\begin{theorem}[What V48 proves]
\label{thm:s7v48-result}
The V46 weighted covariance yields exponential mixing of arbitrary bounded
local Gaussian observables provided a uniform local stochastic floor is
available.  The dependence on that floor is explicit.  Neither raw
covariance decay nor the resulting two-cluster estimate alone proves the
normalized tree lift \eqref{eq:s7v48-treelift}.
\end{theorem}

\begin{proof}
Use
\Cref{thm:s7v48-maxcorr,thm:s7v48-localmix,
prop:s7v48-counter,prop:s7v48-pairwise}.
\end{proof}

\begin{assessment}[Progress at seventh-series V48]
\label{ass:s7v48-status}
The bounded-observable spatial interface is now dimensionless and exact:
the relevant edge is canonical correlation, not raw covariance.  Two
previously hidden obligations are exposed: a local marginal floor for the
projected high Gaussian field, and an all-order Gaussian forest identity
with the normalized tree budget \eqref{eq:s7v48-treelift}.  This prevents the
programme from mistaking pair clustering for the required connected
assembly.
\end{assessment}

\begin{decision}[V49 acquisition order]
\label{dec:s7v48-V49}
V49 will go upstream to the first obligation.  It will express the local
stochastic floor in terms of the angle between a coordinate support and
the slow spectral subspace, prove the floor for uniformly delocalized
finite-rank toron modes, and isolate localized boundary or transition modes
that must remain explicit coarse variables.  Only after this geometric
floor is secured will the all-order forest identity be attempted.
\end{decision}

\begin{firewall}[Claim boundary after seventh-series V48]
\label{fire:s7v48-claim}
V48 does not prove the local stochastic floor for every Yang--Mills patch,
the normalized Gaussian tree lift, the exact normalized forest
combinatorics, the higher Wilson tail, interacting chart exit, patchwise
coarse-action cocycles, terminal strong-coupling matching, cutoff removal,
or the full Yang--Mills mass gap.
\end{firewall}

\paragraph{Parent-version record.}
V48 was formed from
\texttt{Renewal\_Yang\_Mills\_Mass\_Gap\_Research\_Notes\_V47.tex}.
The V47 parent had \(20\,001\) logical source lines, \(726\,337\) bytes,
and SHA--256
\[
 \texttt{79216C07B0364D36CFC04CE255AC33F6B69590FEE75E11AB4C0EB52549E5B0B3}.
\]
The next compulsory companion audit is V50.

% =============================================================================
% END APPENDED SEVENTH-SERIES V48 MATERIAL
% =============================================================================

% =============================================================================
% BEGIN APPENDED SEVENTH-SERIES V49 MATERIAL
% =============================================================================

\chapter{Local principal angles and the high-Gaussian marginal floor}
\label{ch:s7v49}

\section{The geometric content of nondegeneracy}
\label{sec:s7v49-context}

The stochastic floor required in V48 is not an independent probabilistic
assumption.  For the spectrally projected high Gaussian it is controlled
by one principal angle: a local coordinate subspace must not lie almost
inside the retained slow spectral subspace.  This note proves the exact
operator inequality, verifies it for uniformly delocalized finite-rank
toron modes, and states the obstruction created by localized boundary and
transition modes.

Let \({\cal H}_\Lambda=\bigoplus_{x\in\Lambda}E_x\), with fixed finite
fibre dimension, and let
\[
 P_{\rm L}+P_{\rm H}=I
\]
be orthogonal spectral projections.  For \(S\subset\Lambda\), denote the
coordinate projection by \(E_S\).

\section{The principal-angle identity}
\label{sec:s7v49-angle}

\begin{definition}[Local slow leakage]
\label{def:s7v49-leak}
The slow leakage into \(S\) is
\begin{equation}
 \ell(S):=\|P_{\rm L}E_S\|^2
 =\|E_SP_{\rm L}E_S\|.
 \label{eq:s7v49-leak}
\end{equation}
Thus \(1-\ell(S)\) is the squared sine of the smallest principal angle
between \(E_S{\cal H}_\Lambda\) and \(P_{\rm L}{\cal H}_\Lambda\).
\end{definition}

\begin{lemma}[Compression of the high projection]
\label{lem:s7v49-compression}
For every \(u\in E_S{\cal H}_\Lambda\),
\begin{equation}
 \langle u,E_SP_{\rm H}E_Su\rangle
 =\|P_{\rm H}u\|^2
 \geq(1-\ell(S))\|u\|^2.
 \label{eq:s7v49-compression}
\end{equation}
The constant is optimal.
\end{lemma}

\begin{proof}
Since \(P_{\rm H}=I-P_{\rm L}\) and the two projections are orthogonal,
\[
 \|P_{\rm H}u\|^2
 =\|u\|^2-\|P_{\rm L}u\|^2.
\]
The definition of \(\ell(S)\) gives the inequality.  The norm of the
finite-dimensional positive operator \(E_SP_{\rm L}E_S\) is an eigenvalue,
so an associated unit eigenvector gives equality.
\end{proof}

Let the high covariance be defined by functional calculus,
\begin{equation}
 C_{\rm H}=P_{\rm H}f(L)P_{\rm H},
 \label{eq:s7v49-CH}
\end{equation}
where \(P_{\rm H}\) is a spectral projection of the positive self-adjoint
quadratic kernel \(L\), and suppose
\begin{equation}
 f(\lambda)\geq c_{\rm H}>0
 \quad\hbox{on the high spectral band}.
 \label{eq:s7v49-fmin}
\end{equation}
For the V46 inverse-square-root covariance,
\(f(\lambda)=\lambda^{-1/2}\), one may take
\(c_{\rm H}=\Lambda_{\rm H}^{-1/2}\) if
\(\lambda\leq\Lambda_{\rm H}\) on that band.

\begin{theorem}[Principal-angle marginal floor]
\label{thm:s7v49-floor}
Under \eqref{eq:s7v49-fmin},
\begin{equation}
 E_SC_{\rm H}E_S
 \geq
 c_{\rm H}(1-\ell(S))E_S.
 \label{eq:s7v49-floor}
\end{equation}
Consequently, if \(\ell(S)\leq1-\eta\) for every allowed support, then the
V48 stochastic floor holds with
\begin{equation}
 c_-=c_{\rm H}\eta.
 \label{eq:s7v49-cminus}
\end{equation}
\end{theorem}

\begin{proof}
The spectral functional calculus and \eqref{eq:s7v49-fmin} give
\[
 C_{\rm H}\geq c_{\rm H}P_{\rm H}.
\]
Compress by \(E_S\) and apply \Cref{lem:s7v49-compression}.
\end{proof}

\begin{corollary}[Explicit normalized mixing constant]
\label{cor:s7v49-mixing}
If the V46 covariance satisfies
\(\|C_{\rm H}\|_{{\cal S}_\mu}\leq K_\mu\) and
\(\ell(S),\ell(T)\leq1-\eta\), then bounded observables on disjoint
supports \(S,T\) obey
\begin{equation}
 |\operatorname{Cov}(F,G)|
 \leq
 \min\left\{1,{K_\mu\over c_{\rm H}\eta}
 e^{-\mu d(S,T)}\right\}
 \|F\|_\infty\|G\|_\infty.
 \label{eq:s7v49-mixing}
\end{equation}
\end{corollary}

\begin{proof}
Insert \eqref{eq:s7v49-cminus} into
\Cref{thm:s7v48-localmix}.
\end{proof}

\section{Finite-rank delocalized slow modes}
\label{sec:s7v49-toron}

\begin{lemma}[Local leakage from an orthonormal slow basis]
\label{lem:s7v49-basis}
Let \(\psi_1,\ldots,\psi_r\) be an orthonormal basis of
\(\operatorname{Ran}P_{\rm L}\).  Then
\begin{equation}
 \ell(S)
 \leq
 \sum_{j=1}^r\|E_S\psi_j\|^2
 =
 \operatorname{Tr}(E_SP_{\rm L}E_S).
 \label{eq:s7v49-trace}
\end{equation}
\end{lemma}

\begin{proof}
The positive operator \(E_SP_{\rm L}E_S\) has norm at most its trace.
Expanding the trace in the orthonormal basis of the slow range gives the
identity.
\end{proof}

\begin{theorem}[Uniform floor for delocalized toron modes]
\label{thm:s7v49-toron}
Suppose \(r\) is independent of volume and
\begin{equation}
 \sup_{x\in\Lambda}\|\psi_j(x)\|
 \leq {K_{\rm del}\over\sqrt{|\Lambda|}},
 \qquad 1\leq j\leq r.
 \label{eq:s7v49-deloc}
\end{equation}
Then
\begin{equation}
 \ell(S)\leq
 {rK_{\rm del}^2|S|\over|\Lambda|}.
 \label{eq:s7v49-toron-leak}
\end{equation}
In particular, for supports with \(|S|\leq s_0\) and volumes satisfying
\begin{equation}
 |\Lambda|\geq2rK_{\rm del}^2s_0,
 \label{eq:s7v49-volume}
\end{equation}
one has the uniform stochastic floor
\begin{equation}
 E_SC_{\rm H}E_S\geq {c_{\rm H}\over2}E_S.
 \label{eq:s7v49-toron-floor}
\end{equation}
\end{theorem}

\begin{proof}
For each \(j\),
\[
 \|E_S\psi_j\|^2
 =\sum_{x\in S}\|\psi_j(x)\|^2
 \leq {|S|K_{\rm del}^2\over|\Lambda|}.
\]
Sum over \(j\) in \Cref{lem:s7v49-basis}.  Under
\eqref{eq:s7v49-volume}, the result is at most \(1/2\), and
\Cref{thm:s7v49-floor} gives \eqref{eq:s7v49-toron-floor}.
\end{proof}

\begin{corollary}[Constant torons satisfy the hypothesis]
\label{cor:s7v49-constant}
If the slow toron sector consists of a fixed number of normalized constant
fibre vectors, then \eqref{eq:s7v49-deloc} holds with a
volume-independent \(K_{\rm del}\), so every fixed-size local support has a
uniform high-Gaussian floor for all sufficiently large volumes.
\end{corollary}

\begin{proof}
A normalized constant vector has pointwise norm proportional to
\(|\Lambda|^{-1/2}\).  Apply \Cref{thm:s7v49-toron}.
\end{proof}

\section{A basis-free density criterion}
\label{sec:s7v49-density}

Define the diagonal slow density
\begin{equation}
 \varrho_{\rm L}:=
 \sup_{x\in\Lambda}\operatorname{tr}_{E_x}
 (E_xP_{\rm L}E_x).
 \label{eq:s7v49-density}
\end{equation}

\begin{proposition}[Density-to-angle estimate]
\label{prop:s7v49-density}
For every finite support \(S\),
\begin{equation}
 \ell(S)\leq |S|\varrho_{\rm L}.
 \label{eq:s7v49-density-angle}
\end{equation}
Hence any family with
\(\varrho_{\rm L}\leq(1-\eta)/s_0\) has the angle gap
\(\ell(S)\leq1-\eta\) on all supports of size at most \(s_0\).
\end{proposition}

\begin{proof}
As in \Cref{lem:s7v49-basis}, norm is at most trace, and
\[
 \operatorname{Tr}(E_SP_{\rm L}E_S)
 =\sum_{x\in S}\operatorname{tr}_{E_x}(E_xP_{\rm L}E_x)
 \leq|S|\varrho_{\rm L}.
\]
\end{proof}

\begin{assessment}[What must be uniform on patch overlaps]
\label{ass:s7v49-overlap}
The V46 common refinement enlarges \(P_{\rm L}\), so it can only increase
\(\ell(S)\).  Spectral buffering therefore preserves the high-band
spectral floor but does not automatically preserve the local marginal
floor.  The additional modes placed in the overlap slow sector must obey a
uniform density estimate such as
\eqref{eq:s7v49-density-angle}, or they must be represented explicitly in
the local coarse coordinate system.
\end{assessment}

\section{Localized modes and the quotient formulation}
\label{sec:s7v49-localized}

If a slow vector is supported entirely in \(S\), then
\(\ell(S)=1\), and \eqref{eq:s7v49-floor} gives no positive floor on the
full coordinate space.  This is correct: that local direction has been
removed from the high Gaussian and is deterministic after conditioning on
its slow amplitude.

Let
\begin{equation}
 {\cal D}_S:=\ker(E_SP_{\rm H}E_S),\qquad
 {\cal R}_S:={\cal D}_S^\perp\cap E_S{\cal H}_\Lambda.
 \label{eq:s7v49-quotient}
\end{equation}
The first space consists of exact deterministic local directions.

\begin{proposition}[Floor on the stochastic quotient]
\label{prop:s7v49-quotient}
Let
\begin{equation}
 \eta_S^+:=
 \min\bigl(\sigma(E_SP_{\rm H}E_S)\setminus\{0\}\bigr).
 \label{eq:s7v49-etaplus}
\end{equation}
Then
\begin{equation}
 E_SC_{\rm H}E_S\big|_{{\cal R}_S}
 \geq c_{\rm H}\eta_S^+ I_{{\cal R}_S}.
 \label{eq:s7v49-quotient-floor}
\end{equation}
Thus exact local slow directions may be quotiented out, but a
volume-uniform bounded-observable mixing theorem still requires
\(\inf_S\eta_S^+>0\).
\end{proposition}

\begin{proof}
Compress \(C_{\rm H}\geq c_{\rm H}P_{\rm H}\) to \(E_S\), then restrict to
the positive spectral subspace of \(E_SP_{\rm H}E_S\).  Its lowest
eigenvalue is \(\eta_S^+\).
\end{proof}

\begin{firewall}[Almost localized modes]
\label{fire:s7v49-localized}
Quotienting removes eigenvalue \(0\), not a sequence of positive
eigenvalues tending to zero.  An almost localized transition mode can make
\(\eta_S^+\downarrow0\), causing the canonical-correlation constant to
blow up.  Such a mode must be promoted to an explicit coarse coordinate or
controlled by a quantitative principal-angle theorem; calling it
\emph{already slow} is not sufficient.
\end{firewall}

\section{Uniform upper edge of the high band}
\label{sec:s7v49-upper}

\begin{lemma}[Finite-range row bounds give a volume-independent upper
spectral edge]
\label{lem:s7v49-upper}
If the self-adjoint block kernel \(L=(L_{xy})\) satisfies
\begin{equation}
 \sup_x\sum_y\|L_{xy}\|\leq J,
 \qquad
 \sup_y\sum_x\|L_{xy}\|\leq J,
 \label{eq:s7v49-LSchur}
\end{equation}
then \(\|L\|\leq J\).  Therefore the inverse-square-root high covariance
has \(c_{\rm H}\geq J^{-1/2}\), provided its high band lies in
\((0,J]\).
\end{lemma}

\begin{proof}
The block Schur test gives \(\|L\|\leq J\).  Since
\(\lambda^{-1/2}\geq J^{-1/2}\) for \(0<\lambda\leq J\), the functional
calculus gives the stated lower value on the high band.
\end{proof}

\begin{corollary}[Closed two-cluster theorem for delocalized finite-rank
slow sectors]
\label{cor:s7v49-closed}
Assume the V46 weighted covariance bound, the row bound
\eqref{eq:s7v49-LSchur}, a fixed support-size bound \(s_0\), and the
delocalization hypothesis \eqref{eq:s7v49-deloc}.  Then, for all volumes
satisfying \eqref{eq:s7v49-volume}, the V48 bounded two-cluster estimate
holds uniformly with
\[
 c_-\geq {1\over2\sqrt J}.
\]
\end{corollary}

\begin{proof}
Combine
\Cref{lem:s7v49-upper,thm:s7v49-toron,
thm:s7v48-localmix}.
\end{proof}

\section{Principal-angle status}
\label{sec:s7v49-status}

\begin{theorem}[What V49 proves]
\label{thm:s7v49-result}
The local stochastic floor is exactly controlled by the principal angle
to the slow spectral subspace.  It is uniform for fixed-size supports when
the slow sector has fixed rank and uniformly delocalized eigenvectors.
Exact localized slow directions can be quotiented, but almost localized
directions require a separate uniform angle estimate or explicit
promotion to the coarse sector.
\end{theorem}

\begin{proof}
Use
\Cref{thm:s7v49-floor,thm:s7v49-toron,
prop:s7v49-density,prop:s7v49-quotient}.
\end{proof}

\begin{assessment}[Progress at seventh-series V49]
\label{ass:s7v49-status}
For the toron part of the slow sector, the missing V48 normalized
covariance floor is now proved with explicit volume dependence and becomes
uniform in the thermodynamic regime.  The remaining geometric danger is
concentrated in boundary and buffered transition modes.  This is a sharper
target than a generic covariance lower bound and is stable under the
finite-range row estimates used in V46.
\end{assessment}

\begin{decision}[V50 audit and acquisition order]
\label{dec:s7v49-V50}
V50 is a compulsory companion audit.  It will inspect the current Standard
Model and Navier--Stokes notes for a principal-angle, Schur-complement, or
connected-forest mechanism relevant to localized boundary modes and
all-order assembly.  It will then either import a proved compatible
estimate or continue directly to the normalized Gaussian forest identity,
without copying determinant- or PDE-specific conclusions.
\end{decision}

\begin{firewall}[Claim boundary after seventh-series V49]
\label{fire:s7v49-claim}
V49 does not prove a uniform principal-angle gap for all boundary and
transition modes, the normalized Gaussian tree lift, the exact forest
identity, the higher Wilson tail, interacting chart exit, patchwise
coarse-action cocycles, terminal strong-coupling matching, cutoff removal,
or the full Yang--Mills mass gap.
\end{firewall}

\paragraph{Parent-version record.}
V49 was formed from
\texttt{Renewal\_Yang\_Mills\_Mass\_Gap\_Research\_Notes\_V48.tex}.
The V48 parent had \(20\,430\) logical source lines, \(741\,627\) bytes,
and SHA--256
\[
 \texttt{B4685EAAB9E496ACCEA85753934DAD334815D0E2C8C2606F477F1119F38DA582}.
\]
The next compulsory companion audit is V50.

% =============================================================================
% END APPENDED SEVENTH-SERIES V49 MATERIAL
% =============================================================================

% =============================================================================
% BEGIN APPENDED SEVENTH-SERIES V50 MATERIAL
% =============================================================================

\chapter{Companion audit and the finite-rank locality correction}
\label{ch:s7v50}

\section{Compulsory companion audit}
\label{sec:s7v50-audit}

The current companion files in \texttt{latex\_notes} were inspected
read-only:
\begin{align}
 &\texttt{Renewal\_SM\_Einstein\_Continuum\_Research\_Notes\_V92.tex},
 \label{eq:s7v50-SM-file}\\
 &\texttt{Renewal\_NS\_Stability\_Research\_Notes\_V26.tex}.
 \label{eq:s7v50-NS-file}
\end{align}
Their exact checkpoints are
\begin{align}
 {\rm SM\ V92:}\quad&
 37\,596\ {\rm lines},\quad 1\,338\,393\ {\rm bytes},\notag\\[-2pt]
 &\texttt{313088C76E3DEB0A805025D31F95D44D964544F40F0ACD0A96926BE5C30B979E},
 \label{eq:s7v50-SM-hash}\\
 {\rm NS\ V26:}\quad&
 5\,319\ {\rm lines},\quad 278\,283\ {\rm bytes},\notag\\[-2pt]
 &\texttt{82C887F8C2C69E4F28FAD7C3DC8DE5B9099CB5A4BF431EBA1FFF3E91935978AD}.
 \label{eq:s7v50-NS-hash}
\end{align}

\subsection{Standard Model import}
\label{sec:s7v50-SM}

The relevant SM result is the V87 separation of a finite-dimensional slow
determinant from the normalized fast quotient.  It also proves directly
that the constant-mode projector has no volume-uniform exponentially
weighted Schur norm on expanding ordinary boxes.  Two warnings transfer
without using any Einstein-specific determinant conclusion:
\begin{enumerate}
\item a delocalized slow projector belongs in an exact finite-dimensional
coordinate ledger, not in the exponentially weighted local kernel;
\item under ultraviolet refinement, a fixed physical anchor need not have
fixed cell cardinality, so source incidence must be normalized in physical
units.
\end{enumerate}
SM V88--V92 concern weighted elliptic inverses and localized or
finite-part determinants.  Their determinant and heat-kernel conclusions
are not Yang--Mills correlation estimates and are not imported.

\subsection{Navier--Stokes audit}
\label{sec:s7v50-NS}

NS V26 sharpens rank-one Oseen-kernel derivative norms.  Its distinction
between full tensor norms and norms on the actual rank-one input class is a
useful methodological reminder, but the parabolic kernel, Gaussian weight,
and nonlinear input class differ from the present spectral Gaussian
problem.  No NS estimate supplies a Yang--Mills principal-angle or
connected-forest bound, so no mathematical result is imported from NS
V26.

\begin{theorem}[V50 companion-audit decision]
\label{thm:s7v50-audit}
The SM fast-quotient/slow-ledger separation and its nonlocal-projector
obstruction are type-correct and are imported.  No determinant, heat
coefficient, Oseen-kernel constant, or Navier--Stokes regularity statement
is imported.
\end{theorem}

\begin{proof}
The imported statements are Hilbert-space identities about finite-rank
orthogonal projections and weighted row sums; they do not depend on the
SM field equations.  The rejected statements involve different operators,
measures, or nonlinearities and therefore cannot close a Yang--Mills
estimate.
\end{proof}

\section{The obstruction to pure weighted locality}
\label{sec:s7v50-obstruction}

Let \(\Lambda_N\) be a connected graph with \(N\) vertices and graph
distance \(d\).  Assume there are constants \(a_0,a_1>0\), independent of
\(N\), and a vertex \(x_N\) such that at least \(a_0N\) vertices lie at
distance at least \(a_1\operatorname{diam}(\Lambda_N)\) from \(x_N\).
Ordinary expanding boxes have this property.  Let
\begin{equation}
 e_0=N^{-1/2}(1,\ldots,1),\qquad P_0=e_0e_0^*.
 \label{eq:s7v50-P0}
\end{equation}

\begin{theorem}[Delocalized constant projectors are not exponentially
Schur local]
\label{thm:s7v50-projector}
For every fixed \(\mu>0\),
\begin{equation}
 \|P_0\|_{{\cal S}_\mu}
 \geq
 a_0\exp\!\left(
 \mu a_1\operatorname{diam}(\Lambda_N)\right).
 \label{eq:s7v50-projector}
\end{equation}
In particular the norm diverges when the diameter tends to infinity.
Moreover, if \(P_{\rm H}=I-P_0\), then
\begin{equation}
 \|P_{\rm H}\|_{{\cal S}_\mu}
 \geq\|P_0\|_{{\cal S}_\mu}-1,
 \label{eq:s7v50-PH}
\end{equation}
so the complementary high projection is not uniformly weighted local
either.
\end{theorem}

\begin{proof}
Every matrix entry of \(P_0\) equals \(N^{-1}\).  The weighted row based at
\(x_N\) is at least
\[
 {1\over N}(a_0N)
 e^{\mu a_1\operatorname{diam}(\Lambda_N)},
\]
which proves \eqref{eq:s7v50-projector}.  Since
\(\|I\|_{{\cal S}_\mu}=1\) and
\(P_0=I-P_{\rm H}\), the triangle inequality gives
\eqref{eq:s7v50-PH}.
\end{proof}

\begin{corollary}[Compatibility restriction on V46]
\label{cor:s7v50-V46}
The uniform hypotheses and conclusion of the V46 contour-locality theorem
cannot all hold with a fixed decay exponent on expanding ordinary boxes
when its slow space is only the constant toron and
\(P_{\rm H}=I-P_0\).  At least one spectral-buffer, range, hopping, decay,
or slow-rank parameter must deteriorate or change with volume.
\end{corollary}

\begin{proof}
V46 concludes a uniform \({\cal S}_\mu\) bound for \(P_{\rm H}\).
\Cref{thm:s7v50-projector} contradicts that conclusion in the stated
constant-toron family.
\end{proof}

\begin{firewall}[Conditional theorem versus Yang--Mills verification]
\label{fire:s7v50-V46}
The abstract finite-volume Combes--Thomas theorem of V46 remains correct.
The audit changes its application: one may not simultaneously assert a
volume-independent contour buffer, a purely finite-rank delocalized slow
space, and a uniformly exponentially local sharp complementary
projection on expanding local graphs.
\end{firewall}

\section{Exponential locality plus a finite-rank tail}
\label{sec:s7v50-twoscale}

The correct finite-volume class keeps delocalized slow contributions
separate.  Let
\begin{equation}
 C=C^{\rm loc}+F,\qquad
 F=\sum_{j=1}^r a_j\,\psi_j\psi_j^*,
 \label{eq:s7v50-decomp}
\end{equation}
where
\begin{align}
 \|C^{\rm loc}\|_{{\cal S}_\mu}&\leq K_\mu,
 \label{eq:s7v50-local}\\
 |a_j|&\leq A,\qquad
 \|\psi_j\|_2=1,\qquad
 \sup_x\|\psi_j(x)\|\leq {K_{\rm del}\over\sqrt N}.
 \label{eq:s7v50-tail}
\end{align}
The sign of \(a_j\) is unrestricted; subtracting a conditioned slow mode
produces a negative finite-rank term.

\begin{lemma}[Pointwise and unweighted bounds for the finite-rank tail]
\label{lem:s7v50-tail}
Under \eqref{eq:s7v50-tail},
\begin{align}
 \|F_{xy}\|
 &\leq {rAK_{\rm del}^2\over N},
 \label{eq:s7v50-tail-point}\\
 \|F\|_{{\cal S}_0}
 &\leq rAK_{\rm del}.
 \label{eq:s7v50-tail-row}
\end{align}
\end{lemma}

\begin{proof}
For the pointwise bound, sum
\[
 |a_j|\,\|\psi_j(x)\|\,\|\psi_j(y)\|
 \leq {AK_{\rm del}^2\over N}.
\]
For a row sum, Cauchy--Schwarz gives
\(\sum_y\|\psi_j(y)\|\leq\sqrt N\|\psi_j\|_2=\sqrt N\).
Thus
\[
 \sup_x\sum_y
 |a_j|\|\psi_j(x)\|\|\psi_j(y)\|
 \leq AK_{\rm del}.
\]
The column estimate is identical; sum over \(j\).
\end{proof}

\begin{theorem}[Two-scale canonical-correlation bound]
\label{thm:s7v50-twoscale}
Assume the local marginal covariances of \(C\) on supports of size at most
\(s_0\) have stochastic floor \(c_->0\).  If \(S,T\) are disjoint,
\(|S|,|T|\leq s_0\), and \(G_S,G_T\) are bounded functions of the
corresponding Gaussian coordinates, then
\begin{align}
 |\operatorname{Cov}(G_S,G_T)|
 \leq
 {1\over c_-}\left(
 K_\mu e^{-\mu d(S,T)}
 +{rAK_{\rm del}^2s_0\over N}
 \right)
 \|G_S\|_\infty\|G_T\|_\infty,
 \label{eq:s7v50-twoscale}
\end{align}
with the right side truncated at the universal coefficient \(1\).
\end{theorem}

\begin{proof}
The local cross block obeys
\[
 \|E_SC^{\rm loc}E_T\|
 \leq K_\mu e^{-\mu d(S,T)}
\]
by the weighted Schur test, as in V48.  For each rank-one term,
\[
 \|E_S(\psi_j\psi_j^*)E_T\|
 =
 \|E_S\psi_j\|\,\|E_T\psi_j\|
 \leq {K_{\rm del}^2\sqrt{|S||T|}\over N}
 \leq {K_{\rm del}^2s_0\over N}.
\]
Sum over \(j\), normalize on both sides by the marginal inverse square
roots, whose norms are at most \(c_-^{-1/2}\), and apply the Gaussian
maximal-correlation theorem \Cref{thm:s7v48-maxcorr}.
\end{proof}

\begin{corollary}[Thermodynamic disappearance at fixed local supports]
\label{cor:s7v50-thermo}
For a thermodynamic sequence with uniform constants in
\eqref{eq:s7v50-twoscale}, the finite-rank term tends to zero for every two
fixed-size local supports.  Every limiting two-cluster covariance therefore
inherits the exponential term.
\end{corollary}

\begin{proof}
The second term in \eqref{eq:s7v50-twoscale} is \(O(N^{-1})\).
\end{proof}

\section{Row summability of the corrected edge}
\label{sec:s7v50-rowsum}

Define the scalar majorant
\begin{equation}
 q_N(x,y):=
 K_\mu e^{-\mu d(x,y)}+{B\over N},
 \qquad B:=rAK_{\rm del}^2.
 \label{eq:s7v50-q}
\end{equation}
Assume the graphs have uniform exponential volume growth below \(\mu\), so
that
\begin{equation}
 Q_\mu:=\sup_{N,x}\sum_{y\in\Lambda_N}
 e^{-\mu d(x,y)}<\infty.
 \label{eq:s7v50-growth}
\end{equation}

\begin{lemma}[The finite-rank correction does not destroy tree
summability]
\label{lem:s7v50-rowsum}
The corrected edge satisfies
\begin{equation}
 \sup_{N,x}\sum_y q_N(x,y)
 \leq K_\mu Q_\mu+B.
 \label{eq:s7v50-rowsum}
\end{equation}
\end{lemma}

\begin{proof}
Sum the two terms in \eqref{eq:s7v50-q}; the constant tail contributes
\(N(B/N)=B\).
\end{proof}

\begin{proposition}[Fixed-tree volume bound]
\label{prop:s7v50-fixedtree}
Let \(T\) be a labelled tree with one vertex fixed at \(x_0\).  Summing the
remaining vertex positions gives
\begin{equation}
 \sum_{x_2,\ldots,x_m}
 \prod_{\{i,j\}\in E(T)}q_N(x_i,x_j)
 \leq
 (K_\mu Q_\mu+B)^{m-1}.
 \label{eq:s7v50-fixedtree}
\end{equation}
\end{proposition}

\begin{proof}
A finite tree has a leaf distinct from the fixed root.  Sum its position
using \eqref{eq:s7v50-rowsum}, remove that leaf, and iterate.
\end{proof}

\begin{firewall}[Row summability is not decay between two anchors]
\label{fire:s7v50-anchors}
\Cref{prop:s7v50-fixedtree} proves a volume bound.  To prove exponential
clustering between two fixed anchors, one must retain the unique path
between them.  Terms using only local edges on that path decay
exponentially; terms using at least one \(B/N\) edge should vanish in the
thermodynamic limit.  The latter statement requires a path-resolved sum,
not merely \eqref{eq:s7v50-fixedtree}, and is assigned to V51.
\end{firewall}

\section{Order of limits and physical supports}
\label{sec:s7v50-limits}

\begin{firewall}[Fixed cells versus fixed physical anchors]
\label{fire:s7v50-physical}
The bound \(|S|\leq s_0\) is appropriate for a lattice-local observable at
a fixed regulator.  Under ultraviolet refinement, a fixed physical region
may contain a growing number of cells.  Then both the V49 leakage estimate
and the \(s_0/N\) term in \eqref{eq:s7v50-twoscale} must be rewritten using
cell weights and physical-volume normalization.  The thermodynamic and
ultraviolet limits may not be interchanged until that weighted incidence
card is proved.
\end{firewall}

\begin{decision}[Corrected slow/fast ledger]
\label{dec:s7v50-ledger}
From V50 onward:
\begin{enumerate}
\item delocalized toron and global gauge coordinates remain in a separate
finite-dimensional slow ledger;
\item the fast local kernel is estimated in the weighted Schur algebra;
\item any residual delocalized correction is recorded explicitly as a
finite-rank \(O(N^{-1})\) kernel with a bounded unweighted row sum;
\item sharp high projections are not declared uniformly exponentially
local unless their slow-rank and spectral-buffer scaling has been checked;
\item fixed physical sources carry cell weights before any refinement
claim.
\end{enumerate}
\end{decision}

\section{Audit status}
\label{sec:s7v50-status}

\begin{theorem}[What V50 proves]
\label{thm:s7v50-result}
The finite-rank delocalized slow projector obstructs pure weighted locality
of the sharp high projection.  An exponential local kernel plus a
delocalized finite-rank tail instead gives an exponential-plus-\(N^{-1}\)
bounded-observable mixing estimate and retains a uniform tree row sum.
The exact slow ledger and source normalization are indispensable.
\end{theorem}

\begin{proof}
Use
\Cref{thm:s7v50-projector,thm:s7v50-twoscale,
lem:s7v50-rowsum,prop:s7v50-fixedtree}.
\end{proof}

\begin{assessment}[Progress at seventh-series V50]
\label{ass:s7v50-status}
The compulsory audit prevented an inconsistent combination of V46 and
V49.  Pure exponential locality is retained for the genuinely fast local
kernel, while global toron effects are moved to an explicit finite-rank
tail.  The corrected edge remains summable and disappears pointwise in the
thermodynamic limit.  This repair must precede the all-order forest
identity.
\end{assessment}

\begin{decision}[V51 acquisition order]
\label{dec:s7v50-V51}
V51 will prove the path-resolved tree convolution for
\(q_N=q_{\rm loc}+B/N\).  It will show that a tree joining two fixed
anchors is the sum of a uniformly exponentially decaying all-local term
and an \(O(N^{-1})\) term after all off-path vertices are summed.  It will
then state exactly how this result interfaces with the still-unproved
normalized Gaussian tree lift.
\end{decision}

\begin{firewall}[Claim boundary after seventh-series V50]
\label{fire:s7v50-claim}
V50 does not prove the normalized Gaussian tree lift, its exact forest
identity, weighted physical-source incidence under ultraviolet refinement,
the higher Wilson tail, interacting chart exit, patchwise coarse-action
cocycles, terminal strong-coupling matching, cutoff removal, or the full
Yang--Mills mass gap.
\end{firewall}

\paragraph{Parent-version record.}
V50 was formed from
\texttt{Renewal\_Yang\_Mills\_Mass\_Gap\_Research\_Notes\_V49.tex}.
The V49 parent had \(20\,832\) logical source lines, \(754\,563\) bytes,
and SHA--256
\[
 \texttt{CE68CCC66A2D6B0B1BD42FE6DD30CC3E9E66E3A216A852E464625460B81D8E2D}.
\]
The next compulsory companion audit is V55.

% =============================================================================
% END APPENDED SEVENTH-SERIES V50 MATERIAL
% =============================================================================

% =============================================================================
% BEGIN APPENDED SEVENTH-SERIES V51 MATERIAL
% =============================================================================

\chapter{Path-resolved trees for exponential plus finite-rank edges}
\label{ch:s7v51}

\section{Purpose}
\label{sec:s7v51-context}

V50 replaced the inconsistent pure weighted-locality claim by the
two-scale edge
\[
 q_N=p+h_N,
\]
where \(p\) is exponentially local and \(h_N\) is a delocalized
finite-rank majorant of size \(B/N\).  An unweighted row sum shows that
trees remain volume-uniform, but clustering between two anchors requires
more: the unique anchor-to-anchor path must be retained.  This note gives
that path-resolved calculation exactly.

The result is deliberately conditional on the normalized Gaussian tree
lift of V48.  It proves that the finite-rank repair is compatible with that
lift; it does not pretend to prove the lift itself.

\section{Kernel hypotheses}
\label{sec:s7v51-kernel}

Let \(\Lambda_N\) have \(N\) vertices.  Let \(p_N(x,y)\geq0\) be symmetric
and suppose, for some \(\nu>0\),
\begin{align}
 P_0&:=\sup_{N,x}\sum_y p_N(x,y)<\infty,
 \label{eq:s7v51-P0}\\
 P_\nu&:=\sup_{N,x}\sum_y
 e^{\nu d(x,y)}p_N(x,y)<\infty.
 \label{eq:s7v51-Pnu}
\end{align}
Put
\begin{equation}
 h_N(x,y):={B\over N},\qquad
 q_N:=p_N+h_N,
 \qquad
 R:=\max\{P_\nu,P_0+B,1\}.
 \label{eq:s7v51-q}
\end{equation}
Symmetry gives the same row and column constants.

\begin{lemma}[Unweighted convolution masses]
\label{lem:s7v51-mass}
For every \(k\geq0\),
\begin{align}
 \sup_x\sum_y q_N^{*k}(x,y)&\leq(P_0+B)^k,
 \label{eq:s7v51-qmass}\\
 \sup_x\sum_y p_N^{*k}(x,y)&\leq P_0^k,
 \label{eq:s7v51-pmass}
\end{align}
where the zeroth convolution power is the identity kernel.
\end{lemma}

\begin{proof}
The row sum of \(h_N\) is \(B\), so the row sum of \(q_N\) is at most
\(P_0+B\).  Induction and Tonelli's theorem for nonnegative kernels prove
the two estimates.
\end{proof}

\begin{lemma}[Weighted local convolution]
\label{lem:s7v51-weighted}
For \(k\geq1\),
\begin{equation}
 p_N^{*k}(x,y)
 \leq P_\nu^k e^{-\nu d(x,y)}.
 \label{eq:s7v51-weighted}
\end{equation}
\end{lemma}

\begin{proof}
For a chain \(x=x_0,x_1,\ldots,x_k=y\), the triangle inequality gives
\[
 e^{\nu d(x,y)}
 \leq\prod_{j=1}^k e^{\nu d(x_{j-1},x_j)}.
\]
Multiply the convolution summand by the left side, use this inequality,
and sum successively with \eqref{eq:s7v51-Pnu}.
\end{proof}

\begin{lemma}[A path containing a mean-field edge]
\label{lem:s7v51-path-tail}
For \(k\geq1\),
\begin{equation}
 0\leq q_N^{*k}(x,y)-p_N^{*k}(x,y)
 \leq {kB\over N}R^{k-1}.
 \label{eq:s7v51-path-tail}
\end{equation}
\end{lemma}

\begin{proof}
The noncommutative telescoping identity for convolution kernels is
\begin{equation}
 q_N^{*k}-p_N^{*k}
 =
 \sum_{j=1}^k
 q_N^{*(j-1)}*h_N*p_N^{*(k-j)}.
 \label{eq:s7v51-telescope}
\end{equation}
For one summand, use that \(h_N(u,v)=B/N\):
\begin{align*}
 &(q_N^{*(j-1)}*h_N*p_N^{*(k-j)})(x,y)\\
 &\quad={B\over N}
 \left(\sum_u q_N^{*(j-1)}(x,u)\right)
 \left(\sum_v p_N^{*(k-j)}(v,y)\right)
 \leq {B\over N}R^{k-1}.
\end{align*}
The second factor is a column sum, equal to a row sum by symmetry.
Sum the \(k\) terms.
\end{proof}

\section{A fixed tree with two anchors}
\label{sec:s7v51-fixedtree}

Let \(T\) be a labelled tree on vertices
\(\{a,b,1,\ldots,n\}\).  Fix
\(x_a=x\), \(x_b=y\), and sum the other vertex positions.  Denote
\begin{equation}
 {\cal W}_{T,N}(x,y):=
 \sum_{x_1,\ldots,x_n\in\Lambda_N}
 \prod_{\{i,j\}\in E(T)}q_N(x_i,x_j).
 \label{eq:s7v51-WT}
\end{equation}

\begin{theorem}[Path-resolved fixed-tree bound]
\label{thm:s7v51-fixedtree}
For every such tree,
\begin{equation}
 {\cal W}_{T,N}(x,y)
 \leq
 R^{n+1}e^{-\nu d(x,y)}
 +{(n+1)B\over N}R^n.
 \label{eq:s7v51-fixedtree}
\end{equation}
\end{theorem}

\begin{proof}
Let \({\cal P}_{ab}\) be the unique path from \(a\) to \(b\), with \(k\)
edges, where \(1\leq k\leq n+1\).  Repeatedly remove and sum every leaf
which is not on \({\cal P}_{ab}\).  Each removed \(q_N\)-edge costs at
most \(P_0+B\leq R\).  If \(n+1-k\) edges are off the path, this gives
\begin{equation}
 {\cal W}_{T,N}(x,y)
 \leq R^{n+1-k}q_N^{*k}(x,y).
 \label{eq:s7v51-prune}
\end{equation}
Split the last kernel into \(p_N^{*k}\) and its difference from
\(q_N^{*k}\).  By
\Cref{lem:s7v51-weighted,lem:s7v51-path-tail},
\[
 q_N^{*k}(x,y)
 \leq R^k e^{-\nu d(x,y)}
 +{kB\over N}R^{k-1}.
\]
Insert this into \eqref{eq:s7v51-prune} and use \(k\leq n+1\).
\end{proof}

\begin{assessment}[Meaning of the two terms]
\label{ass:s7v51-two}
Branches may use arbitrary finite-rank edges without spoiling exponential
decay, because the anchor path can remain entirely local.  Loss of
exponential decay occurs only when the unique anchor path contains a
finite-rank edge; \Cref{lem:s7v51-path-tail} then preserves its explicit
\(N^{-1}\) factor after all path vertices are summed.
\end{assessment}

\section{Normalized labelled-tree summation}
\label{sec:s7v51-labelled}

Cayley's formula gives \((n+2)^n\) labelled trees on the \(n+2\) vertices.

\begin{corollary}[The \(1/n!\) tree sum]
\label{cor:s7v51-treesum}
For every \(n\geq0\),
\begin{align}
 {1\over n!}\sum_{T\in{\mathfrak T}_{n+2}}
 {\cal W}_{T,N}(x,y)
 \leq&
 e^2R\,(eR)^n e^{-\nu d(x,y)}
 \notag\\
 &+{e^2B\over N}(n+1)(eR)^n.
 \label{eq:s7v51-treesum}
\end{align}
\end{corollary}

\begin{proof}
For \(n\geq1\), the elementary Stirling lower bound gives
\[
 {(n+2)^n\over n!}
 \leq e^n\left(1+{2\over n}\right)^n
 \leq e^{n+2}.
\]
The same final inequality is true at \(n=0\).  Multiply
\eqref{eq:s7v51-fixedtree} by the tree count and \(1/n!\).
\end{proof}

\section{Conditional connected-series closure}
\label{sec:s7v51-series}

\begin{hypothesis}[Normalized two-scale Gaussian tree lift]
\label{hyp:s7v51-lift}
For local bounded observables \(A_x,B_y\) and \(n\) bounded-window
interaction insertions \(v_{x_1},\ldots,v_{x_n}\), assume
\begin{align}
 &|\kappa_0(A_x,B_y,v_{x_1},\ldots,v_{x_n})|
 \notag\\
 &\quad\leq
 C_*^{n+2}\|A_x\|_\infty\|B_y\|_\infty
 \prod_{j=1}^n\|v_{x_j}\|_\infty
 \sum_{T\in{\mathfrak T}_{n+2}}
 \prod_{\{i,j\}\in E(T)}q_N(x_i,x_j).
 \label{eq:s7v51-lift}
\end{align}
Finite interaction range and support incidence are absorbed into \(C_*\).
\end{hypothesis}

\begin{theorem}[Exponential plus vanishing finite-volume remainder]
\label{thm:s7v51-series}
Suppose \(\|v_z\|_\infty\leq\rho\) and
\begin{equation}
 \alpha:=eC_*\rho R<1.
 \label{eq:s7v51-alpha}
\end{equation}
Then the absolute normalized connected series is bounded by
\begin{align}
 &\sum_{n\geq0}{1\over n!}
 \sum_{x_1,\ldots,x_n}
 |\kappa_0(A_x,B_y,v_{x_1},\ldots,v_{x_n})|
 \notag\\
 &\quad\leq
 e^2C_*^2\|A_x\|_\infty\|B_y\|_\infty
 \left[
 {R\over1-\alpha}e^{-\nu d(x,y)}
 +{B\over N(1-\alpha)^2}
 \right].
 \label{eq:s7v51-series}
\end{align}
\end{theorem}

\begin{proof}
Insert \eqref{eq:s7v51-lift}, the activity bound, and
\Cref{cor:s7v51-treesum}.  The local terms form
\[
 e^2C_*^2R e^{-\nu d(x,y)}
 \sum_{n\geq0}\alpha^n
 ={e^2C_*^2R\over1-\alpha}e^{-\nu d(x,y)}.
\]
The finite-rank terms form
\[
 {e^2C_*^2B\over N}
 \sum_{n\geq0}(n+1)\alpha^n
 ={e^2C_*^2B\over N(1-\alpha)^2}.
\]
Restore the two observable norms.
\end{proof}

\begin{corollary}[Thermodynamic exponential clustering, conditional on
the tree lift]
\label{cor:s7v51-thermo}
If the finite-volume connected correlations converge for fixed local
observables and fixed separation, with the constants above uniform, then
every limiting correlation satisfies
\begin{equation}
 |\langle A_x;B_y\rangle_\infty|
 \leq
 {e^2C_*^2R\over1-\alpha}
 \|A_x\|_\infty\|B_y\|_\infty
 e^{-\nu d(x,y)}.
 \label{eq:s7v51-thermo}
\end{equation}
\end{corollary}

\begin{proof}
Take \(N\to\infty\) in \eqref{eq:s7v51-series}.
\end{proof}

\begin{firewall}[The hypothesis is the remaining analytic core]
\label{fire:s7v51-lift}
\Cref{thm:s7v51-series} proves all subsequent tree counting and
thermodynamic removal of the finite-rank edge.  It does not prove
\eqref{eq:s7v51-lift}.  V47's universal bounded-cumulant bound is too
large, and V48's two-cluster maximal-correlation theorem has too few
edges.  The normalized forest identity remains a genuine bottleneck.
\end{firewall}

\section{Averaging the explicit slow label}
\label{sec:s7v51-slow}

Even after conditioning on toron or boundary labels, the full covariance
has the exact decomposition
\begin{equation}
 \operatorname{Cov}(A,B)
 =
 {\mathbb E}_{\rm s}\operatorname{Cov}(A,B\mid s)
 +
 \operatorname{Cov}_{\rm s}
 \bigl({\mathbb E}(A\mid s),{\mathbb E}(B\mid s)\bigr).
 \label{eq:s7v51-totalcov}
\end{equation}
The second term is not controlled by the conditional tree calculation.

\begin{lemma}[Compact slow-label sensitivity]
\label{lem:s7v51-slow}
Let the slow-label space have diameter \(D_{\rm s}\).  If
\[
 f_A(s):={\mathbb E}(A\mid s),\qquad
 f_B(s):={\mathbb E}(B\mid s)
\]
are Lipschitz with constants \(L_A/\sqrt N\) and \(L_B/\sqrt N\), then
\begin{equation}
 |\operatorname{Cov}_{\rm s}(f_A,f_B)|
 \leq {D_{\rm s}^2L_AL_B\over4N}.
 \label{eq:s7v51-slow}
\end{equation}
\end{lemma}

\begin{proof}
Every real random variable taking values in an interval of length
\(\omega\) has variance at most \(\omega^2/4\).  The oscillations of
\(f_A,f_B\) are at most
\(D_{\rm s}L_A/\sqrt N\) and
\(D_{\rm s}L_B/\sqrt N\).  Apply Cauchy--Schwarz to their covariance.
For complex observables, apply the estimate to real and imaginary parts,
with an inessential factor two.
\end{proof}

\begin{firewall}[Slow-label sensitivity is not automatic]
\label{fire:s7v51-slow}
Compactness of the slow-label space does not prove the
\(N^{-1/2}\) Lipschitz constants.  They must come from the normalization
of delocalized modes together with the V46-type parameter derivative
bound for the genuinely fast covariance and the nonlinear effective
action.  Localized boundary labels generally do not have this volume
suppression and must be handled by distance-to-boundary or explicit
boundary-sector estimates.
\end{firewall}

\section{Path-resolved status}
\label{sec:s7v51-status}

\begin{theorem}[What V51 proves]
\label{thm:s7v51-result}
For the corrected edge \(q_N=p_N+B/N\), every tree connecting two anchors
is bounded by an exponentially decaying all-local path plus an explicit
\(O(N^{-1})\) path remainder.  Cayley counting and the normalized
\(1/n!\) give the convergent bound \eqref{eq:s7v51-series} whenever the
normalized tree lift holds and \(eC_*\rho R<1\).
\end{theorem}

\begin{proof}
Use
\Cref{thm:s7v51-fixedtree,cor:s7v51-treesum,
thm:s7v51-series}.
\end{proof}

\begin{assessment}[Progress at seventh-series V51]
\label{ass:s7v51-status}
The V50 finite-rank repair is now fully compatible with anchored tree
summation and thermodynamic exponential clustering.  No delocalized
projector is hidden in a weighted norm.  The remaining central analytic
problem is sharply isolated as the normalized Gaussian tree lift
\eqref{eq:s7v51-lift}; the separate slow-label average additionally
requires a normalized sensitivity estimate.
\end{assessment}

\begin{decision}[V52 acquisition order]
\label{dec:s7v51-V52}
V52 will test whether a Gaussian covariance-interpolation forest formula
can prove \eqref{eq:s7v51-lift} for the actual bounded-window observable
class.  It will write the exact positive interpolation covariance and
derivative identity.  If repeated derivatives require uncontrolled
classical cutoff seminorms, the note will stop that route and formulate the
operator-Duhamel replacement rather than concealing the mismatch.
\end{decision}

\begin{firewall}[Claim boundary after seventh-series V51]
\label{fire:s7v51-claim}
V51 does not prove the normalized Gaussian tree lift, slow-label
sensitivity, weighted physical-source incidence under ultraviolet
refinement, the higher Wilson tail, interacting chart exit, patchwise
coarse-action cocycles, terminal strong-coupling matching, cutoff removal,
or the full Yang--Mills mass gap.
\end{firewall}

\paragraph{Parent-version record.}
V51 was formed from
\texttt{Renewal\_Yang\_Mills\_Mass\_Gap\_Research\_Notes\_V50.tex}.
The V50 parent had \(21\,235\) logical source lines, \(768\,659\) bytes,
and SHA--256
\[
 \texttt{278CF60C7AB393D92EEEE300356E28C7A2C026C43BFA309A07463BFF2EA366BF}.
\]
The next compulsory companion audit is V55.

% =============================================================================
% END APPENDED SEVENTH-SERIES V51 MATERIAL
% =============================================================================

% =============================================================================
% BEGIN APPENDED SEVENTH-SERIES V52 MATERIAL
% =============================================================================

\chapter{The Gaussian forest identity and its cutoff-domain obstruction}
\label{ch:s7v52}

\section{Purpose}
\label{sec:s7v52-context}

V51 reduced connected clustering to the normalized tree lift
\eqref{eq:s7v51-lift}.  A Gaussian covariance interpolation is the most
direct candidate because the BKAR forest formula produces exactly one
covariance derivative per tree edge and no external factorial.  This note
proves that identity and its sharp derivative-seminorm bound.

The calculation also decides whether this route applies to the V47
bounded occupation window.  It does not: a spectral occupation projection
is an operator compression, not a smooth function of the classical
Gaussian coordinate, while a nonzero smooth compact field cutoff cannot
obey the all-order derivative budget needed for the normalized tree bound.
This is a domain mismatch, not a missing constant.  The next route is
therefore operator Duhamel expansion, where the occupation projector is
never differentiated.

\section{Positive forest interpolation}
\label{sec:s7v52-positive}

Let \(X=(X_1,\ldots,X_m)\) be a centred Gaussian vector, where
\(X_i\in{\mathbb R}^{d_i}\), with covariance blocks \(C_{ij}\).  For a
symmetric scalar matrix \(s=(s_{ij})\) with \(s_{ii}=1\), define
\begin{equation}
 C[s]_{ij}:=s_{ij}C_{ij}.
 \label{eq:s7v52-Cs}
\end{equation}

For a forest \(F\) on \(\{1,\ldots,m\}\) and parameters
\(t=(t_e)_{e\in F}\in[0,1]^F\), put
\begin{equation}
 s^F_{ii}(t)=1,\qquad
 s^F_{ij}(t)=
 \begin{cases}
 \displaystyle\min_{e\in{\cal P}_{ij}^F}t_e,
 &i,j\hbox{ in the same component of }F,\\
 0,&i,j\hbox{ in different components}.
 \end{cases}
 \label{eq:s7v52-sF}
\end{equation}

\begin{lemma}[Positivity of the forest weakening]
\label{lem:s7v52-positive}
The scalar matrix \(s^F(t)\) is positive semidefinite.  Consequently
\(C[s^F(t)]\geq0\).
\end{lemma}

\begin{proof}
For \(u\in[0,1]\), retain from \(F\) only edges with \(t_e\geq u\), and
let \(R_u(i,j)\) be \(1\) when \(i,j\) lie in the same resulting component
and \(0\) otherwise.  Each matrix \(R_u\) is positive semidefinite because
it is the Gram matrix of the component-indicator vectors.  Moreover,
\[
 s^F_{ij}(t)=\int_0^1 R_u(i,j)\,du.
\]
Thus \(s^F(t)\geq0\).  To prove the block statement, choose Gram
representations
\[
 C_{i\alpha,j\beta}
 =\langle u_{i\alpha},u_{j\beta}\rangle,\qquad
 s^F_{ij}=\langle v_i,v_j\rangle.
\]
Then
\[
 C[s^F]_{i\alpha,j\beta}
 =\langle u_{i\alpha}\otimes v_i,
          u_{j\beta}\otimes v_j\rangle,
\]
which is positive semidefinite.
\end{proof}

\section{Covariance differentiation}
\label{sec:s7v52-derivative}

For \(i\ne j\), introduce the cross-derivative operator
\begin{equation}
 {\cal D}_{ij}:=
 \sum_{\alpha=1}^{d_i}\sum_{\beta=1}^{d_j}
 (C_{ij})_{\alpha\beta}
 {\partial^2\over\partial x_{i,\alpha}\partial x_{j,\beta}}.
 \label{eq:s7v52-Dij}
\end{equation}

\begin{lemma}[Gaussian covariance derivative]
\label{lem:s7v52-derivative}
Let \(\Phi\) be smooth with bounded derivatives through order two.  At a
positive definite \(C[s]\),
\begin{equation}
 {\partial\over\partial s_{ij}}
 {\mathbb E}_{C[s]}\Phi(X)
 =
 {\mathbb E}_{C[s]}{\cal D}_{ij}\Phi(X),
 \qquad i<j.
 \label{eq:s7v52-derivative}
\end{equation}
The identity extends to positive semidefinite \(C[s]\) by adding
\(\epsilon I\) and taking \(\epsilon\downarrow0\).
\end{lemma}

\begin{proof}
For a differentiable covariance path \(C(u)\), Gaussian integration by
parts gives
\[
 {d\over du}{\mathbb E}_{C(u)}\Phi
 ={1\over2}\sum_{a,b}\dot C_{ab}(u)
 {\mathbb E}_{C(u)}\partial_a\partial_b\Phi.
\]
Changing \(s_{ij}\) changes both the \(ij\) block by \(C_{ij}\) and the
\(ji\) block by \(C_{ji}=C_{ij}^*\); their equal contributions cancel the
factor \(1/2\) and give \eqref{eq:s7v52-derivative}.  Bounded derivatives
permit dominated convergence in the semidefinite regularization.
\end{proof}

\section{The exact connected tree formula}
\label{sec:s7v52-tree}

\begin{theorem}[Gaussian BKAR cumulant identity]
\label{thm:s7v52-BKAR}
Let \(f_i\in C^\infty({\mathbb R}^{d_i})\) have bounded derivatives of all
orders.  Then
\begin{align}
 \kappa_C(f_1(X_1),\ldots,f_m(X_m))
 =
 \sum_{T\in{\mathfrak T}_m}
 \int_{[0,1]^{E(T)}}&
 {\mathbb E}_{C[s^T(t)]}
 \left[
 \left(\prod_{\{i,j\}\in E(T)}{\cal D}_{ij}\right)
 \prod_{k=1}^m f_k(X_k)
 \right]\,dt .
 \label{eq:s7v52-BKAR}
\end{align}
\end{theorem}

\begin{proof}
Apply the fundamental theorem of calculus successively in the
\(\binom m2\) off-diagonal covariance parameters and organize the result
by the first edges that connect previously distinct components.  The
standard forest interpolation identity obtained in this way is
\begin{align}
 {\mathbb E}_{C}\prod_{i=1}^m f_i(X_i)
 =
 \sum_{F\ {\rm forest}}
 \int_{[0,1]^{E(F)}}
 {\mathbb E}_{C[s^F(t)]}
 \left[
 \left(\prod_{\{i,j\}\in E(F)}{\cal D}_{ij}\right)
 \prod_{k=1}^m f_k(X_k)
 \right]dt .
 \label{eq:s7v52-forest}
\end{align}
For completeness, the minimum in \eqref{eq:s7v52-sF} is the parameter
value at which the endpoints first become connected in this ordered
fundamental-theorem construction; induction on the number of unordered
pairs proves \eqref{eq:s7v52-forest}.  Positivity at every interpolation
point is \Cref{lem:s7v52-positive}, and every parameter derivative is
\Cref{lem:s7v52-derivative}.

For a fixed forest \(F\), the covariance
\(C[s^F(t)]\) has zero cross blocks between distinct components of \(F\).
The Gaussian law and the differential operator in
\eqref{eq:s7v52-forest} therefore factor over those components.  Thus the
right side is the set-partition product of its connected forest weights.
Möbius inversion between moments and cumulants selects the weights with
one component.  A connected forest on \(m\) vertices is a tree, giving
\eqref{eq:s7v52-BKAR}.
\end{proof}

\begin{firewall}[No hidden factorial]
\label{fire:s7v52-factorial}
Equation \eqref{eq:s7v52-BKAR} is a sum over labelled trees with one unit
cube integral per tree.  There is no additional \(m!\).  This is exactly
the combinatorial normalization required by V48 and V51.
\end{firewall}

\section{The derivative-seminorm tree bound}
\label{sec:s7v52-seminorm}

For \(k\geq0\), define the coordinate projective derivative seminorm
\begin{equation}
 M_k(f):=
 \sup_x
 \sum_{\alpha_1,\ldots,\alpha_k}
 |\partial_{\alpha_1}\cdots\partial_{\alpha_k}f(x)|,
 \qquad M_0(f)=\|f\|_\infty.
 \label{eq:s7v52-Mk}
\end{equation}
For covariance blocks put
\begin{equation}
 c_{ij}:=\max_{\alpha,\beta}|(C_{ij})_{\alpha\beta}|.
 \label{eq:s7v52-cij}
\end{equation}

\begin{theorem}[Sharp degree-resolved tree estimate]
\label{thm:s7v52-degree}
Under the hypotheses of \Cref{thm:s7v52-BKAR},
\begin{equation}
 |\kappa_C(f_1,\ldots,f_m)|
 \leq
 \sum_{T\in{\mathfrak T}_m}
 \left(\prod_{\{i,j\}\in E(T)}c_{ij}\right)
 \prod_{i=1}^m M_{\deg_T(i)}(f_i).
 \label{eq:s7v52-degree}
\end{equation}
\end{theorem}

\begin{proof}
For a fixed tree, expand every \({\cal D}_{ij}\) in coordinates.  Vertex
\(i\) receives exactly \(\deg_T(i)\) derivatives.  Take the absolute value
inside the expectation and use the supremum in
\eqref{eq:s7v52-Mk}.  Every edge coefficient is at most \(c_{ij}\), and
the coordinate sums are bounded by the corresponding projective
seminorms.  The interpolation cube has volume one.
\end{proof}

\begin{corollary}[A sufficient analytic-type budget]
\label{cor:s7v52-analytic}
If
\begin{equation}
 M_k(f_i)\leq A_iR_i^k
 \quad\hbox{for every }k\geq0,
 \label{eq:s7v52-analytic}
\end{equation}
then
\begin{equation}
 |\kappa_C(f_1,\ldots,f_m)|
 \leq
 \left(\prod_{i=1}^m A_i\right)
 \sum_{T\in{\mathfrak T}_m}
 \prod_{\{i,j\}\in E(T)}
 \bigl(R_iR_jc_{ij}\bigr).
 \label{eq:s7v52-normalized}
\end{equation}
Thus \eqref{eq:s7v52-analytic} produces the normalized tree lift with no
extra factorial.
\end{corollary}

\begin{proof}
In \eqref{eq:s7v52-degree},
\[
 \prod_iR_i^{\deg_T(i)}
 =\prod_{\{i,j\}\in E(T)}R_iR_j.
\]
\end{proof}

\section{Why the current cutoff is not in this domain}
\label{sec:s7v52-cutoff}

\begin{proposition}[No nonzero compact cutoff has the analytic-type
budget]
\label{prop:s7v52-no-cutoff}
Let \(\chi\in C^\infty_c({\mathbb R}^d)\).  If there are \(A,R<\infty\)
such that
\begin{equation}
 M_k(\chi)\leq AR^k\qquad(k\geq0),
 \label{eq:s7v52-cutoff-budget}
\end{equation}
then \(\chi=0\).
\end{proposition}

\begin{proof}
Restrict \(\chi\) to any affine line.  Taylor's theorem and
\eqref{eq:s7v52-cutoff-budget} bound the remainder after order \(n\) on a
segment of length \(h\) by
\[
 {A R^{n+1}|h|^{n+1}\over(n+1)!},
\]
which tends to zero.  Hence every line restriction is real analytic.
Because \(\chi\) vanishes on the nonempty open complement of its compact
support, analytic continuation along lines gives \(\chi=0\) everywhere.
\end{proof}

\begin{proposition}[The three candidate activity classes]
\label{prop:s7v52-classes}
None of the following currently proves the V51 tree-lift hypothesis:
\begin{enumerate}
\item the untruncated cubic/quartic Wilson polynomial has
\(M_k=\infty\) at the low derivative orders and has the V47 zero-radius
moment obstruction;
\item multiplying it by a nonzero smooth compact field cutoff introduces
derivatives of every order and cannot satisfy
\eqref{eq:s7v52-analytic} by
\Cref{prop:s7v52-no-cutoff};
\item the occupation projection \(P_MV_{\rm nl}P_M\) is a bounded operator
on a Fock subspace, not a classical smooth function \(f_i(X_i)\), so
\eqref{eq:s7v52-BKAR} does not apply to it.
\end{enumerate}
\end{proposition}

\begin{proof}
The first item combines the definition \eqref{eq:s7v52-Mk} with
\Cref{thm:s7v47-zero}.  The second is
\Cref{prop:s7v52-no-cutoff}.  The third is a domain statement: classical
coordinate differentiation is not defined on the operator compression.
\end{proof}

\begin{firewall}[Gevrey bounds do not silently repair the count]
\label{fire:s7v52-Gevrey}
A compact smooth cutoff may obey
\(M_k\lesssim R^k(k!)^s\) for some \(s>1\).  In a star tree, one vertex
has degree \(m-1\), so \eqref{eq:s7v52-degree} then contains
\(((m-1)!)^s\).  The single \(1/n!\) normalization does not absorb this.
Any Gevrey resummation would need a new combinatorial mechanism and is not
part of the present proof.
\end{firewall}

\section{The operator-Duhamel replacement}
\label{sec:s7v52-Duhamel}

Let
\begin{equation}
 H(\lambda)=H_0+\sum_X\lambda_XV_X
 \label{eq:s7v52-Hlambda}
\end{equation}
on a finite-volume Hilbert space, with bounded local
\(V_X=P_MV_XP_M\).  At inverse time \(\beta\), define
\[
 Z_\beta(\lambda)=\operatorname{Tr}e^{-\beta H(\lambda)}.
\]

\begin{lemma}[Bounded insertions require no cutoff derivatives]
\label{lem:s7v52-Duhamel}
Every mixed derivative of \(\log Z_\beta\) is an integrated connected
time-ordered Duhamel correlation of the bounded operators \(V_X\).
Each insertion enters through its operator norm; no derivative of
\(P_M\) occurs.
\end{lemma}

\begin{proof}
Duhamel's formula is
\[
 {d\over d\lambda}e^{-\beta(H+\lambda V)}
 =
 -\int_0^\beta
 e^{-(\beta-t)(H+\lambda V)}
 V e^{-t(H+\lambda V)}\,dt.
\]
Iterating gives time-ordered products with one \(V_X\) for each
derivative.  Applying the moment--cumulant Möbius relation to derivatives
of the logarithm retains their connected part.  The projection \(P_M\) is
already part of each bounded \(V_X\) and is never differentiated.
\end{proof}

\begin{firewall}[Finite-temperature norm bounds are not enough]
\label{fire:s7v52-beta}
A crude Duhamel estimate costs \(\beta^n\prod_X\|V_X\|\), which is useless
as \(\beta\to\infty\).  The operator route closes only if a reference
spectral gap and spatial propagation theorem yield an integrable
spacetime tree edge, uniformly in \(\beta\).  This is the precise V53
target.
\end{firewall}

\section{Forest-formula status}
\label{sec:s7v52-status}

\begin{theorem}[What V52 proves]
\label{thm:s7v52-result}
The Gaussian BKAR identity gives the exact normalized labelled-tree
combinatorics and proves the desired tree lift for functions satisfying
the derivative budget \eqref{eq:s7v52-analytic}.  The actual V47
occupation-window interaction is outside that classical domain, and no
nonzero compact smooth field cutoff can satisfy the required all-order
budget.  Operator Duhamel expansion is therefore the selected noncircular
replacement.
\end{theorem}

\begin{proof}
Use
\Cref{thm:s7v52-BKAR,cor:s7v52-analytic,
prop:s7v52-no-cutoff,prop:s7v52-classes,
lem:s7v52-Duhamel}.
\end{proof}

\begin{assessment}[Progress at seventh-series V52]
\label{ass:s7v52-status}
The missing forest combinatorics is now proved, and the absence of an extra
factorial is explicit.  The route fails at a clearly identified functional
domain boundary rather than at tree counting.  This prevents the
occupation cutoff from being treated as if it were a harmless smooth
Gaussian multiplier and redirects the programme to a representation in
which bounded operator insertions are native.
\end{assessment}

\begin{decision}[V53 acquisition order]
\label{dec:s7v52-V53}
V53 will formulate the spacetime operator tree card needed for the
Duhamel cumulants.  It will prove that an edge
\(e^{-\nu d(X,Y)}e^{-\gamma|\tau-\sigma|}\) integrates to the V51 spatial
edge with factor \(2/\gamma\), perform the normalized labelled-tree
summation, and isolate the genuinely dynamical input: a volume-uniform
reference gap plus a Lieb--Robinson or resolvent propagation estimate for
the gauge-reduced oscillator Hamiltonian.
\end{decision}

\begin{firewall}[Claim boundary after seventh-series V52]
\label{fire:s7v52-claim}
V52 does not prove the operator spacetime tree card, a uniform reference
many-body gap, slow-label sensitivity, weighted physical-source incidence,
the higher Wilson tail, interacting chart exit, patchwise coarse-action
cocycles, terminal strong-coupling matching, cutoff removal, or the full
Yang--Mills mass gap.
\end{firewall}

\paragraph{Parent-version record.}
V52 was formed from
\texttt{Renewal\_Yang\_Mills\_Mass\_Gap\_Research\_Notes\_V51.tex}.
The V51 parent had \(21\,638\) logical source lines, \(781\,190\) bytes,
and SHA--256
\[
 \texttt{F900BF6905254177A1B4FB5FE5DBB111E44B1E6626157F4C64F36A45FEA4710D}.
\]
The next compulsory companion audit is V55.

% =============================================================================
% END APPENDED SEVENTH-SERIES V52 MATERIAL
% =============================================================================

% =============================================================================
% BEGIN APPENDED SEVENTH-SERIES V53 MATERIAL
% =============================================================================

\chapter{Spacetime Duhamel trees and the quadratic propagation edge}
\label{ch:s7v53}

\section{Purpose}
\label{sec:s7v53-context}

V52 selected an operator expansion because the occupation projection is a
native bounded operator there and is never differentiated.  Two logically
different ingredients are needed:
\begin{enumerate}
\item a spacetime tree bound for connected Duhamel correlations of the
bounded local insertions;
\item summation of the imaginary-time variables and spatial tree
positions.
\end{enumerate}
This note proves the second step exactly.  It also proves the elementary
spacetime propagation edge for a gapped finite-range one-particle
Hamiltonian and its vacuum two-point function.  The remaining many-body
step is to lift that edge to connected Duhamel cumulants of the actual
gauge-reduced occupation-window operators.

\section{Weighted semigroup propagation}
\label{sec:s7v53-semigroup}

Let \(A=(A_{xy})\) be a self-adjoint block operator on
\(\ell^2(\Lambda;{\mathbb C}^d)\).  Assume
\begin{align}
 A&\geq\gamma_0 I,\qquad \gamma_0>0,
 \label{eq:s7v53-gap}\\
 A_{xy}&=0\quad\hbox{if }d(x,y)>R_0,
 \label{eq:s7v53-range}\\
 \sup_x\sum_{y\ne x}\|A_{xy}\|,
 \ \sup_y\sum_{x\ne y}\|A_{xy}\|
 &\leq J.
 \label{eq:s7v53-hop}
\end{align}
The diagonal blocks do not enter \(J\).

\begin{theorem}[Gapped finite-range imaginary-time propagation]
\label{thm:s7v53-semigroup}
If
\begin{equation}
 \epsilon_\mu:=J(e^{\mu R_0}-1)<\gamma_0,
 \label{eq:s7v53-epsilon}
\end{equation}
then, for \(t\geq0\),
\begin{equation}
 \|(e^{-tA})_{xy}\|
 \leq
 e^{-\mu d(x,y)}
 e^{-(\gamma_0-\epsilon_\mu)t}.
 \label{eq:s7v53-semigroup}
\end{equation}
The estimate is uniform in finite volume when
\(\gamma_0,R_0,J\) are.
\end{theorem}

\begin{proof}
Fix the target \(y\) and first use the bounded weight
\[
 (W_Nu)(z)=
 e^{\mu\min\{d(z,y),N\}}u(z).
\]
As in the V46 Combes--Thomas argument, the two-sided Schur test gives
\begin{equation}
 \|W_NAW_N^{-1}-A\|
 \leq J(e^{\mu R_0}-1)=\epsilon_\mu.
 \label{eq:s7v53-conjugate}
\end{equation}
Indeed the diagonal is unchanged, and on every nonzero off-diagonal block
the ratio of the two weights lies between
\(e^{-\mu R_0}\) and \(e^{\mu R_0}\).

Put \(A_N=W_NAW_N^{-1}\).  Although \(A_N\) need not be self-adjoint,
\eqref{eq:s7v53-gap} and \eqref{eq:s7v53-conjugate} imply
\[
 \operatorname{Re}\langle u,A_Nu\rangle
 \geq(\gamma_0-\epsilon_\mu)\|u\|^2.
\]
For \(v(t)=e^{-tA_N}v_0\),
\[
 {d\over dt}\|v(t)\|^2
 =-2\operatorname{Re}\langle v(t),A_Nv(t)\rangle
 \leq-2(\gamma_0-\epsilon_\mu)\|v(t)\|^2.
\]
Gronwall's inequality gives
\(\|e^{-tA_N}\|\leq e^{-(\gamma_0-\epsilon_\mu)t}\).
Since
\[
 W_Ne^{-tA}W_N^{-1}=e^{-tA_N},
\]
the \(xy\) block on the left is
\[
 e^{\mu\min\{d(x,y),N\}}(e^{-tA})_{xy},
\]
because the weight at \(y\) is one.  Bound it by the semigroup norm and
let \(N\to\infty\), obtaining \eqref{eq:s7v53-semigroup}.
\end{proof}

\begin{corollary}[Explicit half-gap edge]
\label{cor:s7v53-half}
With
\begin{equation}
 \mu_*={1\over R_0}
 \log\left(1+{\gamma_0\over2J}\right),
 \label{eq:s7v53-mustar}
\end{equation}
every \(0<\mu\leq\mu_*\) gives
\begin{equation}
 \|(e^{-tA})_{xy}\|
 \leq e^{-\mu d(x,y)}e^{-\gamma_0t/2}.
 \label{eq:s7v53-half}
\end{equation}
\end{corollary}

\begin{proof}
Then \(\epsilon_\mu\leq\gamma_0/2\) in
\Cref{thm:s7v53-semigroup}.
\end{proof}

\section{The oscillator two-point function}
\label{sec:s7v53-Fock}

Let \({\cal F}\) be bosonic Fock space over the one-particle space and
\[
 H_0=d\Gamma(A).
\]
The vacuum \(\Omega\) is the unique ground state and the excitation gap is
at least \(\gamma_0\).

\begin{lemma}[Second-quantized vacuum edge]
\label{lem:s7v53-vacuum}
For \(t\geq0\) and one-particle vectors \(f,g\),
\begin{equation}
 \langle\Omega,a(f)e^{-tH_0}a^*(g)\Omega\rangle
 =
 \langle f,e^{-tA}g\rangle.
 \label{eq:s7v53-vacuum}
\end{equation}
Consequently localized creation and annihilation operators have the
spacetime edge \eqref{eq:s7v53-semigroup}.
\end{lemma}

\begin{proof}
The vector \(a^*(g)\Omega\) lies in the one-particle sector, on which
\(d\Gamma(A)=A\).  Hence
\[
 e^{-tH_0}a^*(g)\Omega=a^*(e^{-tA}g)\Omega.
\]
The canonical commutation relation and \(a(f)\Omega=0\) give
\eqref{eq:s7v53-vacuum}.
\end{proof}

\begin{firewall}[Number-conserving reference]
\label{fire:s7v53-number}
\Cref{lem:s7v53-vacuum} is stated for \(d\Gamma(A)\).  A quadratic
Hamiltonian with pairing terms first requires a symplectic Bogoliubov
reduction whose kernels and inverse are uniformly local.  V46 supplies
some relevant functional calculus but, after the V50 projector correction,
does not by itself prove that full Bogoliubov locality card.
\end{firewall}

\section{Spacetime tree integration}
\label{sec:s7v53-time}

Let
\begin{equation}
 k_\gamma(\tau,\sigma):=e^{-\gamma|\tau-\sigma|}.
 \label{eq:s7v53-ktime}
\end{equation}
Its row integral is
\begin{equation}
 \sup_\sigma\int_{\mathbb R}k_\gamma(\tau,\sigma)\,d\tau
 ={2\over\gamma}.
 \label{eq:s7v53-time-row}
\end{equation}

\begin{lemma}[A tree with two equal-time anchors]
\label{lem:s7v53-time-tree}
Let \(T\) be a tree on
\(\{a,b,1,\ldots,n\}\), set \(\tau_a=\tau_b=0\), and integrate the other
times over \(\mathbb R\).  Then
\begin{equation}
 \int_{\mathbb R^n}
 \prod_{\{i,j\}\in E(T)}
 e^{-\gamma|\tau_i-\tau_j|}
 \,d\tau_1\cdots d\tau_n
 \leq\left({2\over\gamma}\right)^n.
 \label{eq:s7v53-time-tree}
\end{equation}
\end{lemma}

\begin{proof}
Remove and integrate every leaf other than \(a,b\).  Each such integration
costs \(2/\gamma\).  The remaining graph is the unique path from \(a\) to
\(b\), with \(k\) edges and \(k-1\) internal variables.  Its integral is
the convolution kernel \(k_\gamma^{*k}(0,0)\).  Since
\(\sup_{\tau,\sigma}k_\gamma(\tau,\sigma)\leq1\) and its row integral is
\(2/\gamma\), successive convolution gives
\[
 k_\gamma^{*k}(0,0)\leq(2/\gamma)^{k-1}.
\]
The number of removed off-path vertices plus \(k-1\) is exactly \(n\),
which proves the claim.
\end{proof}

\begin{hypothesis}[Operator spacetime tree card]
\label{hyp:s7v53-card}
For bounded local operators \(O_i\), centred connected ground-state
Duhamel densities obey
\begin{align}
 |\kappa_0^{\rm D}(O_1(\tau_1),\ldots,O_m(\tau_m))|
 \leq
 C_{\rm D}^{m}\prod_{i=1}^m\|O_i\|
 \sum_{T\in{\mathfrak T}_m}
 \prod_{\{i,j\}\in E(T)}
 \left[
 q_N(x_i,x_j)e^{-\gamma|\tau_i-\tau_j|}
 \right],
 \label{eq:s7v53-card}
\end{align}
where \(q_N=p_N+B/N\) is the corrected V50 edge.
\end{hypothesis}

\begin{theorem}[Spacetime card implies the V51 spatial series]
\label{thm:s7v53-integrated}
Let \(A_x,B_y\) be equal-time anchors and let the local bounded interaction
operators satisfy \(\|V_z\|\leq\rho\).  If
\begin{equation}
 \alpha_{\rm D}:=
 {2eC_{\rm D}\rho R\over\gamma}<1,
 \label{eq:s7v53-alpha}
\end{equation}
then the absolute ground-state Duhamel series is bounded by
\begin{align}
 e^2C_{\rm D}^2\|A_x\|\|B_y\|
 \left[
 {R\over1-\alpha_{\rm D}}e^{-\nu d(x,y)}
 +{B\over N(1-\alpha_{\rm D})^2}
 \right].
 \label{eq:s7v53-integrated}
\end{align}
\end{theorem}

\begin{proof}
For order \(n\), apply \eqref{eq:s7v53-card} with
\(\tau_a=\tau_b=0\), integrate the \(n\) insertion times, and use
\Cref{lem:s7v53-time-tree}.  This replaces the activity \(\rho\) in the
V51 labelled-tree sum by \(2\rho/\gamma\).  Apply
\Cref{thm:s7v51-series} with
\(C_*=C_{\rm D}\) and
\(\alpha=\alpha_{\rm D}\).
\end{proof}

\begin{corollary}[Conditional thermodynamic mass scale]
\label{cor:s7v53-scale}
Under the hypotheses of \Cref{thm:s7v53-integrated} and existence of the
thermodynamic limit, the limiting equal-time connected correlation decays
at least as \(e^{-\nu d(x,y)}\).  The spatial inverse correlation length is
\(\nu\); the reference imaginary-time decay entering the convergence
radius is \(\gamma\).
\end{corollary}

\begin{proof}
Take \(N\to\infty\) in \eqref{eq:s7v53-integrated}.
\end{proof}

\begin{firewall}[Correlation decay is not yet the physical mass gap]
\label{fire:s7v53-mass}
Even after \eqref{eq:s7v53-card} is proved, one must establish the
Osterwalder--Schrader reconstruction hypotheses and show that the
Euclidean time-decay exponent persists through regulator removal.  Only
then does the spectral representation identify a positive Hamiltonian
mass gap.  The spatial exponent \(\nu\) alone is not that conclusion.
\end{firewall}

\section{What remains in the operator lift}
\label{sec:s7v53-lift}

The two-point identity \eqref{eq:s7v53-vacuum} does not automatically imply
\eqref{eq:s7v53-card} for arbitrary bounded local operators.  A valid lift
must address:
\begin{enumerate}
\item Wick or resolvent organization of all connected contractions with
no extra factorial beyond the labelled-tree sum;
\item uniform constants for the local occupation-window algebra as
\(M=L^{1/8}\) grows;
\item propagation after inserting the gauge constraint and the
slow/fast Feshbach map;
\item the finite-rank toron tail and localized boundary labels;
\item stability of the reference gap under the nonlinear self-energy.
\end{enumerate}

\begin{proposition}[Polynomial vacuum subcase]
\label{prop:s7v53-polynomial}
For a fixed number of normal-ordered creation and annihilation fields of
uniformly bounded total degree, every vacuum moment is a finite sum of
products of two-point edges \eqref{eq:s7v53-vacuum}, and its connected
cumulant retains only connected contraction multigraphs.
\end{proposition}

\begin{proof}
Wick's theorem expresses the moment as the sum over pairings.  In the
moment--cumulant Möbius relation, a pairing graph disconnected across a
partition factorizes and cancels.  Conversely every surviving pairing
graph is connected on the observable vertices.
\end{proof}

\begin{firewall}[Why the polynomial subcase is insufficient]
\label{fire:s7v53-polynomial}
The number of legs in a general matrix element of the \(M\)-occupation
operator algebra grows with \(M\), and compression by local spectral
projections need not preserve a fixed normal-ordered degree.  Therefore
\Cref{prop:s7v53-polynomial} gives a check on the edge structure, not a
uniform proof of \eqref{eq:s7v53-card}.
\end{firewall}

\section{Spacetime-card status}
\label{sec:s7v53-status}

\begin{theorem}[What V53 proves]
\label{thm:s7v53-result}
A gapped finite-range one-particle Hamiltonian has an explicit
volume-uniform exponential spacetime propagator.  The same kernel is the
vacuum two-point edge of its number-conserving second quantization.
Assuming the operator spacetime tree card, imaginary-time integration,
Cayley counting, the finite-rank correction, and thermodynamic clustering
all close under the explicit smallness condition
\eqref{eq:s7v53-alpha}.
\end{theorem}

\begin{proof}
Use
\Cref{thm:s7v53-semigroup,lem:s7v53-vacuum,
lem:s7v53-time-tree,thm:s7v53-integrated}.
\end{proof}

\begin{assessment}[Progress at seventh-series V53]
\label{ass:s7v53-status}
The combinatorial and integration consequences of an operator connected
tree card are now complete, including explicit dependence on the
reference gap.  The dynamical bottleneck is no longer phrased as generic
clustering: it is the uniform lift from the proved quadratic two-point
propagator to connected Duhamel cumulants of the growing
occupation-window gauge algebra.
\end{assessment}

\begin{decision}[V54 acquisition order]
\label{dec:s7v53-V54}
V54 will quantify the local occupation-window algebra.  It will bound the
normal-ordered expansion of matrix units on the truncated oscillator
space, determine its growth with \(M\), and test that growth against the
small activity \(gM^{3/2}+g^2M^2\).  If the matrix-unit basis is
exponentially ill-conditioned in \(M\), the programme will instead seek a
completely bounded quasi-free correlation inequality.
\end{decision}

\begin{firewall}[Claim boundary after seventh-series V53]
\label{fire:s7v53-claim}
V53 does not prove the operator spacetime tree card, its uniform
occupation-window constants, the full Bogoliubov locality card,
slow-label sensitivity, weighted physical-source incidence, the higher
Wilson tail, interacting chart exit, patchwise coarse-action cocycles,
terminal strong-coupling matching, cutoff removal, OS reconstruction, or
the full Yang--Mills mass gap.
\end{firewall}

\paragraph{Parent-version record.}
V53 was formed from
\texttt{Renewal\_Yang\_Mills\_Mass\_Gap\_Research\_Notes\_V52.tex}.
The V52 parent had \(22\,063\) logical source lines, \(795\,776\) bytes,
and SHA--256
\[
 \texttt{7B492CC138DBF33A51989A807AB7075CCE5CCDC2A48B62DBEC62AE4CF973EA43}.
\]
The next compulsory companion audit is V55.

% =============================================================================
% END APPENDED SEVENTH-SERIES V53 MATERIAL
% =============================================================================

% =============================================================================
% BEGIN APPENDED SEVENTH-SERIES V54 MATERIAL
% =============================================================================

\chapter{Occupation-window algebra constants and schedule pressure}
\label{ch:s7v54}

\section{Purpose}
\label{sec:s7v54-context}

The operator spacetime card of V53 must be uniform enough that
\(C_{\rm D}(M)\rho_M\to0\).  Since the occupation cutoff
\(M=L^{1/8}\) grows, finite dimensionality alone supplies no such
uniformity.  This note computes two exact pieces of the local algebra:
its dimension and matrix-unit expansion cost, and the particle-number
cost of second-quantized propagation.

The resulting naive matrix-unit route loses only a polynomial in \(M\),
not an exponential.  Nevertheless that polynomial is already too large
for the current \(M=L^{1/8}\) schedule when a block contains even one
oscillator mode.  A slower window, such as a power of \(\log L\), restores
the perturbative inequality for every fixed polynomial loss, subject to
rechecking the high-occupation return estimates.

\section{Dimension of the truncated local Fock space}
\label{sec:s7v54-dimension}

Fix \(r\) oscillator modes in one local block and write
\[
 {\cal F}_{r,\leq M}
 :=
 \operatorname{span}\{
 |\boldsymbol n\rangle:
 \boldsymbol n\in{\mathbb N}_0^r,\ |\boldsymbol n|\leq M\}.
\]

\begin{lemma}[Exact occupation-window dimension]
\label{lem:s7v54-dimension}
The dimension is
\begin{equation}
 D_{r,M}:=\dim{\cal F}_{r,\leq M}
 ={M+r\choose r}
 \leq(M+1)^r.
 \label{eq:s7v54-dimension}
\end{equation}
\end{lemma}

\begin{proof}
Introduce one slack variable
\(n_{r+1}=M-\sum_{j=1}^rn_j\).  Stars and bars counts the nonnegative
solutions of \(\sum_{j=1}^{r+1}n_j=M\) as
\({M+r\choose r}\).  Dropping the total-occupation constraint embeds the
index set into \(\{0,\ldots,M\}^r\), proving the upper bound.
\end{proof}

\section{Matrix units and their exact normal ordering}
\label{sec:s7v54-matrix}

For multi-indices \(\boldsymbol p,\boldsymbol q\) with total occupation at
most \(M\), let
\[
 E_{\boldsymbol p\boldsymbol q}
 :=|\boldsymbol p\rangle\langle\boldsymbol q|.
\]

\begin{lemma}[Coefficient cost of a matrix-unit expansion]
\label{lem:s7v54-lone}
Every \(O\in{\cal B}({\cal F}_{r,\leq M})\) has the expansion
\[
 O=\sum_{\boldsymbol p,\boldsymbol q}
 O_{\boldsymbol p\boldsymbol q}
 E_{\boldsymbol p\boldsymbol q},
\]
and
\begin{equation}
 \sum_{\boldsymbol p,\boldsymbol q}
 |O_{\boldsymbol p\boldsymbol q}|
 \leq D_{r,M}^{3/2}\|O\|.
 \label{eq:s7v54-lone}
\end{equation}
\end{lemma}

\begin{proof}
There are \(D_{r,M}^2\) coefficients.  Cauchy--Schwarz and the
Hilbert--Schmidt bound give
\[
 \sum_{\boldsymbol p,\boldsymbol q}|O_{\boldsymbol p\boldsymbol q}|
 \leq D_{r,M}\|O\|_{\rm HS}
 \leq D_{r,M}^{3/2}\|O\|.
\]
\end{proof}

Let
\[
 {\cal N}=\sum_{j=1}^r a_j^*a_j.
\]
The normal-ordered falling factorial is denoted by
\(:{\cal N}^k:\); it has eigenvalue
\(n(n-1)\cdots(n-k+1)\) on total occupation \(n\).

\begin{lemma}[Vacuum projection on the finite window]
\label{lem:s7v54-vacuum}
On \({\cal F}_{r,\leq M}\),
\begin{equation}
 |\boldsymbol0\rangle\langle\boldsymbol0|
 =
 \sum_{k=0}^M{(-1)^k\over k!}:{\cal N}^k:.
 \label{eq:s7v54-vacuum}
\end{equation}
Consequently
\begin{equation}
 E_{\boldsymbol p\boldsymbol q}
 =
 {(\boldsymbol a^*)^{\boldsymbol p}\over
  \sqrt{\boldsymbol p!}}
 \left(\sum_{k=0}^M{(-1)^k\over k!}:{\cal N}^k:\right)
 {\boldsymbol a^{\boldsymbol q}\over
  \sqrt{\boldsymbol q!}}
 \quad\hbox{on }{\cal F}_{r,\leq M}.
 \label{eq:s7v54-matrix-normal}
\end{equation}
\end{lemma}

\begin{proof}
On total occupation \(n\leq M\), the right side of
\eqref{eq:s7v54-vacuum} is
\[
 \sum_{k=0}^n{(-1)^k\over k!}{n!\over(n-k)!}
 =\sum_{k=0}^n(-1)^k{n\choose k}
 =\mathbf 1_{\{n=0\}}.
\]
This proves the first identity.  The normalized annihilation monomial on
the right of \eqref{eq:s7v54-matrix-normal} maps
\(|\boldsymbol q\rangle\) to the vacuum and the normalized creation
monomial maps the vacuum to \(|\boldsymbol p\rangle\); all other basis
vectors are killed by the middle projection.
\end{proof}

\begin{corollary}[Normal-order degree grows linearly with \(M\)]
\label{cor:s7v54-degree}
The representation \eqref{eq:s7v54-matrix-normal} uses normal-ordered
monomials of total degree at most \(4M\).
\end{corollary}

\begin{proof}
The \(k\)-th middle term has degree \(2k\), and
\(|\boldsymbol p|,|\boldsymbol q|,k\leq M\).
\end{proof}

\begin{assessment}[What the matrix expansion does and does not cost]
\label{ass:s7v54-matrix}
The change from operator norm to matrix-unit coefficient norm costs at
most \(D_{r,M}^{3/2}=O(M^{3r/2})\), so it is polynomial for fixed local
mode number \(r\).  A subsequent absolute Wick expansion of
\eqref{eq:s7v54-matrix-normal}, however, sees as many as \(4M\) legs.
Finite dimensionality therefore does not justify a uniform fixed-degree
Wick estimate.
\end{assessment}

\section{Particle-number propagation costs}
\label{sec:s7v54-tensor}

\begin{lemma}[Tensor-power telescoping]
\label{lem:s7v54-tensor}
Let \(S,T\) be contractions on a Hilbert space.  For \(n\geq1\),
\begin{equation}
 \|T^{\otimes_s n}-S^{\otimes_s n}\|
 \leq n\|T-S\|.
 \label{eq:s7v54-tensor}
\end{equation}
\end{lemma}

\begin{proof}
On the full tensor power,
\[
 T^{\otimes n}-S^{\otimes n}
 =
 \sum_{j=1}^n
 T^{\otimes(j-1)}\otimes(T-S)
 \otimes S^{\otimes(n-j)}.
\]
Every other tensor factor has norm at most one, so the norm is at most
\(n\|T-S\|\).  Compression to the symmetric tensor subspace cannot
increase the norm.
\end{proof}

\begin{corollary}[At most linear occupation loss]
\label{cor:s7v54-linear}
For second-quantized contraction propagation restricted to total
occupation \(n\leq M\), a one-particle decoupling error is amplified by at
most \(M\).
\end{corollary}

\begin{proof}
The \(n\)-particle propagator is the symmetric tensor power of its
one-particle propagator.  Apply \Cref{lem:s7v54-tensor} and maximize over
\(n\leq M\).
\end{proof}

\begin{firewall}[This is not yet a many-body tree proof]
\label{fire:s7v54-linear}
The linear factor in \Cref{cor:s7v54-linear} controls one
second-quantized decoupling.  Iterating decouplings while taking connected
cumulants requires an exact forest or telescoping construction on the
operator state.  Without it, one may not declare the V53 spacetime card
proved.
\end{firewall}

\section{A quantified naive matrix-unit route}
\label{sec:s7v54-naive}

\begin{proposition}[Polynomial ceiling, conditional on unit matrix
elements]
\label{prop:s7v54-ceiling}
Suppose an operator forest construction has the following two properties:
\begin{enumerate}
\item each normalized matrix unit has a spacetime tree estimate with a
constant \(C_0\) independent of \(M\);
\item every tree edge uses at most one second-quantized decoupling loss of
\Cref{cor:s7v54-linear}.
\end{enumerate}
Then expansion in the matrix-unit basis yields a card constant no worse
than
\begin{equation}
 C_{\rm D}(M)
 \leq C\,M D_{r,M}^{3/2}
 \leq C\,M(M+1)^{3r/2}.
 \label{eq:s7v54-CD}
\end{equation}
\end{proposition}

\begin{proof}
\Cref{lem:s7v54-lone} costs \(D_{r,M}^{3/2}\) at every observable
vertex.  By the second assumption, the \(m-1\) edge losses are at most
\(M^{m-1}\), which can be bounded by \(M^m\) and assigned one factor \(M\)
per vertex.  Absorb fixed endpoint and \(C_0\) constants into \(C^m\).
This gives the per-vertex constant in \eqref{eq:s7v54-CD}; use
\Cref{lem:s7v54-dimension}.
\end{proof}

\begin{firewall}[The ceiling is conditional, not an achieved card]
\label{fire:s7v54-ceiling}
The two assumptions of \Cref{prop:s7v54-ceiling} have not been proved for
the gauge-reduced reference state.  The proposition answers only the
schedule question: even a favourable matrix-unit construction naturally
produces a polynomial \(M\)-loss that must be budgeted.
\end{firewall}

\section{General schedule theorem}
\label{sec:s7v54-schedule}

Suppose
\begin{equation}
 C_{\rm D}(M)\leq C M^p
 \label{eq:s7v54-poly}
\end{equation}
for some fixed \(p\geq0\), and use the V47 local activity estimate
\begin{equation}
 \rho_M\leq
 C'\bigl(gM^{3/2}+g^2M^2+\rho_{\geq5}(g,M)\bigr).
 \label{eq:s7v54-rho}
\end{equation}

\begin{theorem}[Power-window compatibility criterion]
\label{thm:s7v54-power}
Let \(g=L^{-1/2}\) and \(M=L^\theta\).  If
\begin{equation}
 0<\theta<{1\over2p+3}
 \label{eq:s7v54-theta}
\end{equation}
and
\begin{equation}
 M^p\rho_{\geq5}(L^{-1/2},M)\longrightarrow0,
 \label{eq:s7v54-tail}
\end{equation}
then
\begin{equation}
 C_{\rm D}(M)\rho_M\longrightarrow0.
 \label{eq:s7v54-product}
\end{equation}
\end{theorem}

\begin{proof}
The cubic and quartic contributions after multiplication by \(M^p\) are
\[
 L^{-1/2+\theta(p+3/2)},\qquad
 L^{-1+\theta(p+2)}.
\]
The first exponent is negative exactly under
\eqref{eq:s7v54-theta}.  For \(p\geq0\), that inequality also implies
\(\theta<1/(p+2)\), so the quartic exponent is negative.  The tail is
\eqref{eq:s7v54-tail}.
\end{proof}

\begin{corollary}[Pressure on the \(L^{1/8}\) window]
\label{cor:s7v54-eighth}
For the naive ceiling
\[
 p=1+{3r\over2},
\]
\Cref{thm:s7v54-power} requires
\begin{equation}
 \theta<{1\over3r+5}.
 \label{eq:s7v54-naive-theta}
\end{equation}
Thus \(\theta=1/8\) is already borderline for \(r=1\) and fails this
strict smallness test for every \(r\geq1\).
\end{corollary}

\begin{proof}
Substitute \(p=1+3r/2\) in \eqref{eq:s7v54-theta}.
\end{proof}

\begin{theorem}[Logarithmic windows absorb every fixed polynomial loss]
\label{thm:s7v54-log}
Let
\begin{equation}
 M=(\log L)^a,\qquad a>0.
 \label{eq:s7v54-log}
\end{equation}
For every fixed \(p\),
\[
 M^p\left(L^{-1/2}M^{3/2}+L^{-1}M^2\right)\longrightarrow0.
\]
Hence \eqref{eq:s7v54-product} follows if
\[
 M^p\rho_{\geq5}(L^{-1/2},M)\longrightarrow0.
\]
\end{theorem}

\begin{proof}
Every fixed power of \(\log L\) is \(o(L^\epsilon)\) for every
\(\epsilon>0\).  Apply this separately to the cubic and quartic terms.
\end{proof}

\begin{decision}[Provisional schedule correction]
\label{dec:s7v54-schedule}
The operator-Duhamel branch will provisionally use
\(M=(\log L)^a\) until a sharper completely bounded tree theorem removes
the matrix-dimension loss.  No earlier \(M=L^{1/8}\) estimate is
discarded: V55 must re-evaluate the occupation Feshbach floor, mixing,
chart, and tail rows under the slower window before the correction is
adopted globally.
\end{decision}

\begin{firewall}[A slower window must still diverge]
\label{fire:s7v54-log}
Keeping \(M\) fixed would make the local algebra constants harmless but
would not remove the occupation cutoff.  The logarithmic choice tends to
infinity, yet its slower rate may weaken large-occupation tail estimates.
Those tails, not perturbative smallness, are the possible obstruction to
\eqref{eq:s7v54-log}.
\end{firewall}

\section{Occupation-algebra status}
\label{sec:s7v54-status}

\begin{theorem}[What V54 proves]
\label{thm:s7v54-result}
The local occupation-window algebra has polynomial dimension for fixed
mode number, its matrix-unit coefficient expansion costs at most
\(D_{r,M}^{3/2}\), and one-particle propagation errors amplify at most
linearly with total occupation.  Any spacetime card with polynomial loss
\(M^p\) is compatible with a power window precisely under the sufficient
criterion \(\theta<1/(2p+3)\), and with every logarithmic window subject to
the declared higher-tail estimate.
\end{theorem}

\begin{proof}
Use
\Cref{lem:s7v54-dimension,lem:s7v54-lone,
lem:s7v54-tensor,thm:s7v54-power,thm:s7v54-log}.
\end{proof}

\begin{assessment}[Progress at seventh-series V54]
\label{ass:s7v54-status}
The finite-dimensional operator route is not exponentially
ill-conditioned at the level of matrix-unit coefficients, but the
existing \(L^{1/8}\) window has no margin for its natural polynomial
losses.  The programme therefore acquires a quantitative schedule
criterion and a viable logarithmic alternative.  The unproved part is now
the operator forest construction and the compatibility of the slower
window with high-occupation return.
\end{assessment}

\begin{decision}[V55 audit and acquisition order]
\label{dec:s7v54-V55}
V55 is a compulsory companion audit.  After checking the current SM and
NS notes, it will rebuild the V42--V45 occupation and chart rate ledger for
\(M=(\log L)^a\).  It will determine which returns need only \(M\to\infty\)
and which require a positive power of \(L\), before continuing the
operator forest card.
\end{decision}

\begin{firewall}[Claim boundary after seventh-series V54]
\label{fire:s7v54-claim}
V54 does not prove the matrix-unit hypotheses of
\Cref{prop:s7v54-ceiling}, the operator spacetime tree card, the higher
Wilson tail under the logarithmic window, high-occupation return under the
revised schedule, the full Bogoliubov locality card, slow-label
sensitivity, weighted physical-source incidence, interacting chart exit,
patchwise coarse-action cocycles, terminal strong-coupling matching,
cutoff removal, OS reconstruction, or the full Yang--Mills mass gap.
\end{firewall}

\paragraph{Parent-version record.}
V54 was formed from
\texttt{Renewal\_Yang\_Mills\_Mass\_Gap\_Research\_Notes\_V53.tex}.
The V53 parent had \(22\,448\) logical source lines, \(808\,717\) bytes,
and SHA--256
\[
 \texttt{9A77D4A8A27CE18718F0A65B666E91586E4D7ED61AAA47A5B7556E4B51E6A7DC}.
\]
The next compulsory companion audit is V55.

% =============================================================================
% END APPENDED SEVENTH-SERIES V54 MATERIAL
% =============================================================================

% =============================================================================
% BEGIN APPENDED SEVENTH-SERIES V55 MATERIAL
% =============================================================================

\chapter{Companion audit and the logarithmic occupation-window ledger}
\label{ch:s7v55}

\section{Compulsory companion audit}
\label{sec:s7v55-audit}

The current active companion files inspected for this five-version
checkpoint were
\begin{align}
 &\texttt{Renewal\_SM\_Einstein\_Continuum\_Research\_Notes\_V98.tex},
 \notag\\[-2pt]
 &\quad 39\,963\ {\rm logical\ source\ lines},\quad
 1\,419\,500\ {\rm bytes},\notag\\[-2pt]
 &\quad
 \texttt{4496729452F0131ECD9DA4C453F29CD6CBC9FB9CAADC303C32556FBC3A090AE2},
 \label{eq:s7v55-SM}\\
 &\texttt{Renewal\_NS\_Stability\_Research\_Notes\_V26.tex},
 \notag\\[-2pt]
 &\quad 5\,319\ {\rm logical\ source\ lines},\quad
 278\,283\ {\rm bytes},\notag\\[-2pt]
 &\quad
 \texttt{82C887F8C2C69E4F28FAD7C3DC8DE5B9099CB5A4BF431EBA1FFF3E91935978AD}.
 \label{eq:s7v55-NS}
\end{align}

\subsection{SM V93--V98}
\label{sec:s7v55-SM}

SM V95 independently audits Yang--Mills V48--V51 and confirms the same
finite-rank projector obstruction proved in V50: global slow projectors
are bookkeeping objects, not uniformly weighted-local kernels.  Its
subsequent V96--V98 work concerns local symbol coefficients and correlated
heat-regulator limits.  Those results neither control a Yang--Mills
occupation window nor prove an operator Duhamel tree card.  The
fast-quotient/slow-ledger rule has already been imported in V50, so this
audit creates no second copy of it.

\subsection{NS V26}
\label{sec:s7v55-NS}

NS V26 is unchanged.  Its rank-one Oseen-kernel optimization has a
different state space, time evolution, and input cone.  It supplies no
estimate for the oscillator-number return, Wilson Taylor tail, or
gauge-reduced Duhamel correlations.

\begin{theorem}[V55 companion-audit decision]
\label{thm:s7v55-audit}
No new companion theorem is imported.  The SM audit corroborates the V50
finite-rank correction, while the current SM and NS additions do not alter
the occupation-window acquisition order.
\end{theorem}

\begin{proof}
The corroborated statement is already proved internally by
\Cref{thm:s7v50-projector}.  All other inspected results concern
determinant heat limits or parabolic Oseen kernels, whose operators and
norms do not match the Yang--Mills window problem.
\end{proof}

\section{Decoupling the three growing scales}
\label{sec:s7v55-scales}

Let
\begin{equation}
 L=\log{1\over a\Lambda_0},\qquad L>e.
 \label{eq:s7v55-L}
\end{equation}
The V36 momentum-filter scale remains
\begin{equation}
 \kappa=L^{1/16},\qquad
 N_\kappa=O(\kappa^2)=O(L^{1/8}),
 \label{eq:s7v55-kappa}
\end{equation}
where \(N_\kappa\) is the literal polynomial-filter range.  Independently,
for a fixed \(\chi>0\), set
\begin{equation}
 M=(\log L)^\chi,\qquad
 R=M^{1/2}=(\log L)^{\chi/2}.
 \label{eq:s7v55-MR}
\end{equation}
Here \(M\) is the oscillator-number window and \(R\) is the local
weak-field chart radius.

\begin{lemma}[Basic scale compatibility]
\label{lem:s7v55-basic}
With \(g=O(L^{-1/2})\),
\begin{equation}
 M\to\infty,\quad R\to\infty,\quad
 gR=O\!\left(L^{-1/2}(\log L)^{\chi/2}\right)\to0.
 \label{eq:s7v55-basic}
\end{equation}
The momentum/filter rates of V36 remain unchanged.
\end{lemma}

\begin{proof}
Every fixed power of \(\log L\) is \(o(L^\epsilon)\) for every
\(\epsilon>0\).  The V36 estimates depend on \(\kappa\) and
\(N_\kappa\), not on the independently chosen \(M,R\).
\end{proof}

\begin{firewall}[An old numerical coincidence is not a scale identity]
\label{fire:s7v55-coincidence}
In the former schedule, both the literal filter range \(N_\kappa\) and the
occupation window \(M\) happened to scale as \(L^{1/8}\).  They arise from
different estimates.  V36 requires
\(N_\kappa=O(\kappa^2)\); V42--V43 require only a growing \(M\), a chart
radius with \(gR\to0\), and \(c_R=o(M)\).  No proved equation requires
\(N_\kappa=M\).
\end{firewall}

\section{Relative-form and occupation return}
\label{sec:s7v55-occupation}

The V43 fixed-block card is
\[
 \eta_{g,R}=O(gR),\qquad
 c_R=O(R^{-2}+g^2R^2).
\]

\begin{theorem}[Logarithmic-window relative-form rates]
\label{thm:s7v55-form}
Under \eqref{eq:s7v55-MR},
\begin{align}
 \eta_{g,R}
 &=O\!\left(L^{-1/2}(\log L)^{\chi/2}\right)=o(1),
 \label{eq:s7v55-eta}\\
 c_R
 &=O\!\left((\log L)^{-\chi}
 +L^{-1}(\log L)^\chi\right),
 \label{eq:s7v55-cR}\\
 {c_R\over M}
 &=O\!\left((\log L)^{-2\chi}+L^{-1}\right)=o(1).
 \label{eq:s7v55-cM}
\end{align}
Thus the V42 high-occupation floor obeys
\begin{equation}
 \Delta_M(a)\geq{\omega_*M\over2}
 \label{eq:s7v55-floor}
\end{equation}
for all sufficiently small \(a\).
\end{theorem}

\begin{proof}
Insert \(R=(\log L)^{\chi/2}\) and \(g=O(L^{-1/2})\) into the V43 card.
Then use the exact V42 formula
\[
 \Delta_M=(1-\eta_{g,R})\omega_*(M+1)-c_R.
\]
Equations \eqref{eq:s7v55-eta} and \eqref{eq:s7v55-cM} give
\eqref{eq:s7v55-floor}.
\end{proof}

Assume temporarily the V42 remainder condition
\begin{equation}
 \rho_a(M)=o(gM^{3/2}).
 \label{eq:s7v55-rho-assumed}
\end{equation}

\begin{theorem}[Logarithmic occupation Feshbach ledger]
\label{thm:s7v55-Feshbach}
For bounded physical spectral parameter \(z\), the V42 estimates become
\begin{align}
 \rho_{\rm occ,mix}
 &=
 O\!\left(L^{-1/2}(\log L)^{\chi/2}\right),
 \label{eq:s7v55-mix}\\
 \Sigma_{\rm occ}
 &=
 O\!\left({L^{-1}(\log L)^{2\chi}\over a}\right),
 \label{eq:s7v55-Sigma}\\
 \rho_{\rm occ,der}
 &=
 O\!\left(L^{-1}(\log L)^\chi\right).
 \label{eq:s7v55-der}
\end{align}
The first and third rows tend to zero.  The middle row remains a low-sector
self-energy rather than a residual tail.
\end{theorem}

\begin{proof}
By \Cref{thm:s7v55-form},
\(\delta_M(z)\asymp M\).  The leading window vertex is
\(k_M=O(gM^{3/2})\).  Insert these in the exact V42 bounds
\[
 {k_M\over\delta_M},\qquad
 {k_M^2\over a\delta_M},\qquad
 {k_M^2\over\delta_M^2}.
\]
They give \(gM^{1/2}\), \(g^2M^2/a\), and \(g^2M\), respectively.
Substitute \eqref{eq:s7v55-MR}.
\end{proof}

\section{The fixed-block higher Wilson tail}
\label{sec:s7v55-Wilson}

Let \(S_{\mathscr B}(q)\) be the fixed-block magnetic Wilson potential of
V43 and let \(T_4S_{\mathscr B}\) be its Taylor polynomial through degree
four at the vacuum.  Let \(\chi_R(y)\) be the V39 chart localization,
supported where \(|y|\leq R\), and define the dimensionless localized
remainder
\begin{equation}
 {\cal R}_{5,g,R}(y)
 :=
 \chi_R(y)g^{-2}
 \bigl[S_{\mathscr B}(gy)-T_4S_{\mathscr B}(gy)\bigr]
 \chi_R(y).
 \label{eq:s7v55-R5}
\end{equation}

\begin{lemma}[Oscillator coordinate moments on a number window]
\label{lem:s7v55-coordinate}
For every fixed integer \(k\geq0\) and fixed block oscillator family,
there is \(C_{k,\mathscr B}\) such that
\begin{equation}
 \||y|^kP_M\|
 \leq C_{k,\mathscr B}(M+k+1)^{k/2}.
 \label{eq:s7v55-coordinate}
\end{equation}
Moreover,
\begin{equation}
 \|P_M|y|^5P_M\|
 \leq C_{\mathscr B}(M+6)^{5/2}.
 \label{eq:s7v55-fifth-moment}
\end{equation}
\end{lemma}

\begin{proof}
Each coordinate is a fixed linear combination of a creation and an
annihilation operator.  Norm equivalence in the fixed coordinate
dimension reduces \(\||y|^k\psi\|\) to a fixed sum of
\(\||y_j|^k\psi\|=\|y_j^k\psi\|\).  Expanding each integer coordinate
monomial of degree \(k\) and using the V42 word estimate proves
\eqref{eq:s7v55-coordinate}.

For the second assertion and \(\psi\in\operatorname{Ran}P_M\),
Cauchy--Schwarz in coordinate space gives
\[
 \langle\psi,|y|^5\psi\rangle
 \leq\||y|^2\psi\|\,\||y|^3\psi\|
 \leq C_{\mathscr B}(M+6)^{5/2}\|\psi\|^2.
\]
Polarization, or the variational characterization of a positive
compression, gives \eqref{eq:s7v55-fifth-moment}.
\end{proof}

\begin{theorem}[Fifth-order magnetic remainder on the occupation window]
\label{thm:s7v55-Wilson-tail}
For \(gR\) inside the fixed injectivity radius,
\begin{equation}
 \|Q_M{\cal R}_{5,g,R}P_M\|
 +\|P_M{\cal R}_{5,g,R}P_M\|
 \leq C_{\mathscr B}g^3(M+6)^{5/2}.
 \label{eq:s7v55-Wilson-tail}
\end{equation}
Under \eqref{eq:s7v55-MR},
\begin{equation}
 \rho_{\geq5}^{\rm mag}
 =
 O\!\left(
 L^{-3/2}(\log L)^{5\chi/2}
 \right).
 \label{eq:s7v55-Wilson-rate}
\end{equation}
This is \(o(gM^{3/2})\), and
\(M^p\rho_{\geq5}^{\rm mag}\to0\) for every fixed \(p\).
\end{theorem}

\begin{proof}
Taylor's theorem on the fixed compact coordinate ball gives
\[
 |S_{\mathscr B}(q)-T_4S_{\mathscr B}(q)|
 \leq C_{\mathscr B}|q|^5.
\]
With \(q=gy\) and the prefactor \(g^{-2}\), the localized multiplication
remainder is bounded in absolute value by
\(C_{\mathscr B}g^3|y|^5\).  Therefore
\[
 \|{\cal R}_{5,g,R}P_M\|
 \leq C_{\mathscr B}g^3\||y|^5P_M\|
 \leq C_{\mathscr B}g^3(M+6)^{5/2}
\]
by \Cref{lem:s7v55-coordinate}.  Left multiplication by either
\(P_M\) or \(Q_M\) cannot increase the norm, proving
\eqref{eq:s7v55-Wilson-tail}.  Substitution of \(g,M\) gives
\eqref{eq:s7v55-Wilson-rate}.  Its ratio to \(gM^{3/2}\) is
\[
 g^2M=O\!\left(L^{-1}(\log L)^\chi\right)\to0.
\]
Multiplication by any fixed power of \(M\) adds only a power of
\(\log L\), still dominated by \(L^{-3/2}\).
\end{proof}

\begin{firewall}[Scope of the new tail theorem]
\label{fire:s7v55-Wilson}
\Cref{thm:s7v55-Wilson-tail} closes the chart-localized fixed-block
magnetic Taylor remainder.  It does not control chart exit, Haar kinetic
coefficients outside the V43 relative-form comparison, boundary gluing,
or the global quasi-local filter.  Thus
\eqref{eq:s7v55-rho-assumed} is proved for this magnetic component, not
for every term formerly hidden in \(\rho_a(M)\).
\end{firewall}

\section{Chart and filter rows}
\label{sec:s7v55-chart}

\begin{theorem}[Complete explicit local rate table]
\label{thm:s7v55-table}
With \eqref{eq:s7v55-kappa}--\eqref{eq:s7v55-MR}, the explicitly proved
rows are
\begin{align}
 \rho_{\rm filt}
 &\leq Ce^{-hL^{1/16}},
 \label{eq:s7v55-filter}\\
 \rho_{\rm qloc}
 &=O(L^{-5/16}),
 \qquad
 \rho_{\rm fr}=O(L^{-1/8}),
 \label{eq:s7v55-incidence}\\
 \rho_{\rm mom}
 &\leq e^{-\tau\sqrt\alpha/a},
 \label{eq:s7v55-momentum}\\
 \rho_3
 &=O\!\left(L^{-1/2}(\log L)^{3\chi/2}\right),
 \label{eq:s7v55-rho3}\\
 \rho_4
 &=O\!\left(L^{-1}(\log L)^{2\chi}\right),
 \label{eq:s7v55-rho4}\\
 \rho_{\rm occ,mix}
 &=O\!\left(L^{-1/2}(\log L)^{\chi/2}\right),
 \label{eq:s7v55-table-mix}\\
 \rho_{\rm occ,der}
 &=O\!\left(L^{-1}(\log L)^\chi\right),
 \label{eq:s7v55-table-der}\\
 \eta_{\rm chart}
 &=O\!\left(L^{-1/2}(\log L)^{\chi/2}\right),
 \label{eq:s7v55-chart}\\
 c_{\rm IMS}
 &=O\!\left((\log L)^{-\chi}\right),
 \label{eq:s7v55-IMS}\\
 \rho_{\rm Gexit}
 &\leq C\exp\!\left[-c(\log L)^\chi\right],
 \label{eq:s7v55-Gexit}\\
 \rho_{\geq5}^{\rm mag}
 &=O\!\left(L^{-3/2}(\log L)^{5\chi/2}\right).
 \label{eq:s7v55-table-tail}
\end{align}
Every displayed residual row tends to zero for every fixed \(\chi>0\).
\end{theorem}

\begin{proof}
The first three lines are the unchanged V34--V36 filter and momentum
estimates.  The cubic and quartic lines are their number-window operator
norms.  Use
\Cref{thm:s7v55-Feshbach,thm:s7v55-form,
thm:s7v55-Wilson-tail} for the remaining algebraic rows.  V38 gives
\(\rho_{\rm Gexit}\leq Ce^{-cR^2}\), and \(R^2=M\).
\end{proof}

\begin{corollary}[Compatibility with every fixed polynomial operator loss]
\label{cor:s7v55-polynomial}
For each fixed \(p\geq0\),
\begin{equation}
 M^p\bigl(
 \rho_3+\rho_4+\rho_{\rm occ,mix}
 +\rho_{\rm occ,der}+\eta_{\rm chart}
 +\rho_{\rm Gexit}+\rho_{\geq5}^{\rm mag}
 \bigr)\longrightarrow0.
 \label{eq:s7v55-polynomial}
\end{equation}
The IMS row also vanishes after multiplication by \(M^p\) only when
\(p<1\); for larger losses it must remain a separate form-localization
debit rather than an operator activity.
\end{corollary}

\begin{proof}
All rows except IMS contain either a negative power of \(L\) or an
exponential in \(M\), which dominates every fixed power of \(\log L\).
For IMS,
\[
 M^pc_{\rm IMS}
 =O((\log L)^{\chi(p-1)}),
\]
which tends to zero exactly for \(p<1\), is merely bounded at \(p=1\), and
grows for \(p>1\).
\end{proof}

\begin{firewall}[The logarithmic window does not close interacting exit]
\label{fire:s7v55-exit}
The Gaussian reference tail \eqref{eq:s7v55-Gexit} is proved.  A normalized
interacting chart-exit estimate must compare the exterior weight with the
interacting partition function or ground state uniformly in volume and
slow labels.  The fact that
\(\exp[-c(\log L)^\chi]\to0\) does not by itself pay an extensive
denominator.
\end{firewall}

\section{Schedule verdict}
\label{sec:s7v55-verdict}

\begin{theorem}[What V55 proves]
\label{thm:s7v55-result}
The logarithmic occupation schedule
\[
 \kappa=L^{1/16},\qquad
 M=(\log L)^\chi,\qquad R=M^{1/2}
\]
preserves the V36 filter rates, the V43 relative-form hypotheses, the V42
high-occupation floor and vanishing mixing/energy-dependence rows, and the
V38 Gaussian chart tail.  It also yields the new fixed-block magnetic
remainder bound \eqref{eq:s7v55-Wilson-tail}.  All operator-activity rows
survive every fixed polynomial \(M\)-loss; the IMS form debit must remain
outside that multiplication.
\end{theorem}

\begin{proof}
Use
\Cref{lem:s7v55-basic,thm:s7v55-form,
thm:s7v55-Feshbach,thm:s7v55-Wilson-tail,
thm:s7v55-table,cor:s7v55-polynomial}.
\end{proof}

\begin{assessment}[Progress at seventh-series V55]
\label{ass:s7v55-status}
The provisional V54 schedule change passes every explicit local
occupation, chart, and momentum test.  The formerly symbolic higher
magnetic Wilson tail is now proved on the localized fixed block.  Two
items remain outside this ledger: the operator forest/Duhamel card and the
normalized interacting chart-exit estimate.  The IMS term is correctly
kept as a relative-form floor debit rather than multiplied by a
matrix-algebra constant.
\end{assessment}

\begin{decision}[V56 acquisition order]
\label{dec:s7v55-V56}
V56 will return to the operator spacetime card.  It will construct a
decoupling interpolation for a finite-volume number-conserving quadratic
Fock Hamiltonian, differentiate its Gibbs or ground-state functional by
Duhamel's formula, and determine whether the connected forest weights have
a completely bounded estimate independent of the local matrix dimension.
The logarithmic schedule will be used if only a fixed polynomial
\(M\)-loss is obtained.
\end{decision}

\begin{firewall}[Claim boundary after seventh-series V55]
\label{fire:s7v55-claim}
V55 does not prove the operator spacetime tree card, a completely bounded
quasi-free forest inequality, the full remainder
\eqref{eq:s7v55-rho-assumed}, normalized interacting chart exit, the full
Bogoliubov locality card, slow-label sensitivity, weighted physical-source
incidence, patchwise coarse-action cocycles, terminal strong-coupling
matching, cutoff removal, OS reconstruction, or the full Yang--Mills mass
gap.
\end{firewall}

\paragraph{Parent-version record.}
V55 was formed from
\texttt{Renewal\_Yang\_Mills\_Mass\_Gap\_Research\_Notes\_V54.tex}.
The V54 parent had \(22\,867\) logical source lines, \(822\,329\) bytes,
and SHA--256
\[
 \texttt{2405E6305086309F2E5ABEC7BABB6B29A77FD3B7E441D213099D9C584E450AAA}.
\]
The next compulsory companion audit is V60.

% =============================================================================
% END APPENDED SEVENTH-SERIES V55 MATERIAL
% =============================================================================

% =============================================================================
% BEGIN APPENDED SEVENTH-SERIES V56 MATERIAL
% =============================================================================

\chapter{Quasi-free Weyl cumulants and a dimension-free tree card}
\label{ch:s7v56}

\section{Purpose}
\label{sec:s7v56-context}

V53 asks for a spacetime tree estimate for bounded operator insertions.
V54 shows that a generic matrix-unit expansion introduces polynomial
occupation losses.  The Weyl subalgebra has better structure: its
operators are unitary, its quasi-free Euclidean moments factor exactly
into pair terms, and its connected cumulants are Mayer connected-graph
sums.  A stable tree-graph inequality then gives precisely the normalized
labelled-tree card, with no local matrix-dimension factor.

This note proves that statement for the number-conserving quadratic
reference.  The remaining issue is representation: the compressed Wilson
vertices and occupation projections must be written as controlled
superpositions or derivatives of Weyl operators without losing the
dimension-free constant.

\section{A connected-graph identity}
\label{sec:s7v56-Mayer}

Let \(Y_1,\ldots,Y_m\) be formal variables whose subset moments have the
pair-product form
\begin{equation}
 {\cal M}(I)
 =
 \prod_{i\in I}m_i
 \prod_{\{i,j\}\subset I}(1+\zeta_{ij}),
 \qquad {\cal M}(\varnothing)=1.
 \label{eq:s7v56-moment}
\end{equation}
Write \({\mathfrak C}_m\) for the connected simple graphs on
\(\{1,\ldots,m\}\).

\begin{lemma}[Mayer moment--cumulant identity]
\label{lem:s7v56-Mayer}
The joint cumulant determined by \eqref{eq:s7v56-moment} is
\begin{equation}
 \kappa(Y_1,\ldots,Y_m)
 =
 \left(\prod_{i=1}^m m_i\right)
 \sum_{G\in{\mathfrak C}_m}
 \prod_{\{i,j\}\in E(G)}\zeta_{ij}.
 \label{eq:s7v56-Mayer}
\end{equation}
\end{lemma}

\begin{proof}
Expand the pair product in \eqref{eq:s7v56-moment}.  It is the sum over all
simple graphs on \(I\), with graph weight
\(\prod_{e\in E(G)}\zeta_e\).  Every graph factors uniquely over its
connected components.  Therefore the family of all-graph weights is the
set-partition product of the connected-graph weights.  Möbius inversion
between moments and cumulants selects the one-component graphs, proving
\eqref{eq:s7v56-Mayer}.
\end{proof}

\section{A stable tree-graph inequality}
\label{sec:s7v56-treegraph}

For a Hermitian family \(u_{ji}=\overline{u_{ij}}\), put
\begin{equation}
 \zeta_{ij}=e^{-u_{ij}}-1.
 \label{eq:s7v56-zeta}
\end{equation}

\begin{definition}[Vertex-weighted stability]
\label{def:s7v56-stable}
Given \(b_i\geq0\), the family is stable with vertex weights \(b_i\) if,
for every \(I\subset\{1,\ldots,m\}\),
\begin{equation}
 \sum_{\{i,j\}\subset I}\operatorname{Re}u_{ij}
 \geq-\sum_{i\in I}b_i.
 \label{eq:s7v56-stable}
\end{equation}
The uniform \(B\)-stable case is \(b_i=B\).
\end{definition}

\begin{lemma}[Stable connected graphs are bounded by trees]
\label{lem:s7v56-treegraph}
If \eqref{eq:s7v56-stable} holds, then
\begin{align}
 \sum_{G\in{\mathfrak C}_m}
 \prod_{e\in E(G)}(e^{-u_e}-1)
 =
 (-1)^{m-1}\sum_{T\in{\mathfrak T}_m}
 \left(\prod_{e\in E(T)}u_e\right)
 \int_{[0,1]^{E(T)}}
 \exp\!\left[-\sum_{i<j}s^T_{ij}(t)u_{ij}\right]dt .
 \label{eq:s7v56-BF}
\end{align}
Consequently
\begin{equation}
 \left|
 \sum_{G\in{\mathfrak C}_m}
 \prod_{e\in E(G)}(e^{-u_e}-1)
 \right|
 \leq
 \exp\!\left(\sum_{i=1}^m b_i\right)
 \sum_{T\in{\mathfrak T}_m}
 \prod_{e\in E(T)}|u_e|.
 \label{eq:s7v56-treegraph}
\end{equation}
\end{lemma}

\begin{proof}
Apply the scalar BKAR forest formula to
\[
 \Phi(s):=\exp\!\left[-\sum_{i<j}s_{ij}u_{ij}\right].
\]
Every forest edge derivative contributes \(-u_e\).  As in the connected
selection proof of \Cref{thm:s7v52-BKAR}, the forest term factorizes over
its components, so moment--cumulant Möbius inversion retains exactly the
trees.  This proves \eqref{eq:s7v56-BF}.

It remains to bound the interpolated exponential.  The V52 representation
of the forest matrix is
\[
 s^T_{ij}(t)=\int_0^1R_v(i,j)\,dv,
\]
where, at level \(v\), \(R_v(i,j)\) is the indicator that \(i,j\) belong
to the same component of the subforest containing edges with \(t_e\geq v\).
For the component partition \({\cal P}_v\), stability gives
\begin{align*}
 \sum_{i<j}R_v(i,j)\operatorname{Re}u_{ij}
 &=
 \sum_{C\in{\cal P}_v}
 \sum_{\{i,j\}\subset C}\operatorname{Re}u_{ij}\\
 &\geq-\sum_{C\in{\cal P}_v}\sum_{i\in C}b_i
 =-\sum_{i=1}^m b_i.
\end{align*}
Integrate over \(v\).  Therefore
\[
 \left|
 \exp\!\left[-\sum_{i<j}s^T_{ij}(t)u_{ij}\right]
 \right|\leq\exp\!\left(\sum_{i=1}^m b_i\right).
\]
Take absolute values in \eqref{eq:s7v56-BF} and sum over trees.
\end{proof}

\begin{firewall}[Role of stability]
\label{fire:s7v56-stability}
Without \eqref{eq:s7v56-stable}, the product of all non-tree Mayer bonds
can cost an exponential in \(m^2\).  Stability is what reduces this to a
constant per vertex and is essential for the normalized \(1/n!\)
connected expansion.
\end{firewall}

\section{Euclidean Weyl moments}
\label{sec:s7v56-Weyl}

Let \(A\geq\gamma_0I>0\) be the one-particle operator of V53,
\(H_0=d\Gamma(A)\), and \(\Omega\) the Fock vacuum.  Use the Weyl
convention
\begin{equation}
 W(f):=\exp(a^*(f)-a(f)).
 \label{eq:s7v56-W}
\end{equation}
For ordered times
\(\tau_1\geq\cdots\geq\tau_m\), define
\begin{align}
 {\cal M}(\{1,\ldots,m\})
 :=
 \langle\Omega,&
 W(f_1)e^{-(\tau_1-\tau_2)H_0}W(f_2)
 \cdots
 e^{-(\tau_{m-1}-\tau_m)H_0}W(f_m)
 \Omega\rangle .
 \label{eq:s7v56-Euclidean}
\end{align}
Subset moments retain the inherited time order.

\begin{theorem}[Exact quasi-free pair product]
\label{thm:s7v56-pair}
The moment \eqref{eq:s7v56-Euclidean} is
\begin{equation}
 {\cal M}(\{1,\ldots,m\})
 =
 \prod_{i=1}^m e^{-\|f_i\|^2/2}
 \prod_{1\leq i<j\leq m}
 \exp\!\left[
 -\langle f_i,e^{-(\tau_i-\tau_j)A}f_j\rangle
 \right].
 \label{eq:s7v56-pair}
\end{equation}
\end{theorem}

\begin{proof}
The normal-order identity is
\[
 W(f)=e^{-\|f\|^2/2}e^{a^*(f)}e^{-a(f)}.
\]
On coherent vectors,
\[
 e^{-tH_0}e^{a^*(g)}\Omega
 =e^{a^*(e^{-tA}g)}\Omega.
\]
Moving every annihilation exponential to the right uses
\[
 e^{-a(f)}e^{a^*(g)}
 =e^{-\langle f,g\rangle}
 e^{a^*(g)}e^{-a(f)}.
\]
Each ordered pair \(i<j\) therefore contributes
\(-\langle f_i,e^{-(\tau_i-\tau_j)A}f_j\rangle\);
the vacuum kills the remaining normal-ordered creation and annihilation
factors.  This proves \eqref{eq:s7v56-pair}.
\end{proof}

\section{Stability of the Euclidean pair kernel}
\label{sec:s7v56-kernel}

For arbitrary times, define the Hermitian kernel
\begin{equation}
 U_{ij}:=
 \begin{cases}
 \langle f_i,e^{-(\tau_i-\tau_j)A}f_j\rangle,
 &\tau_i\geq\tau_j,\\
 \overline{U_{ji}},&\tau_i<\tau_j.
 \end{cases}
 \label{eq:s7v56-U}
\end{equation}

\begin{lemma}[Ornstein--Uhlenbeck positivity]
\label{lem:s7v56-OU}
The matrix \(U=(U_{ij})\) is positive semidefinite.
\end{lemma}

\begin{proof}
By the spectral theorem it suffices to prove positivity of
\(e^{-\lambda|\tau-\sigma|}\) for \(\lambda>0\).  Its Fourier
representation is
\begin{equation}
 e^{-\lambda|t|}
 =
 {1\over\pi}\int_{\mathbb R}
 {\lambda\over\lambda^2+\omega^2}e^{i\omega t}\,d\omega,
 \label{eq:s7v56-Fourier}
\end{equation}
whose density is nonnegative.  Integrating this scalar positive kernel
against the positive spectral measure of \(A\) proves the operator-valued
claim.
\end{proof}

\begin{corollary}[Exact vertex stability]
\label{cor:s7v56-stable}
The off-diagonal family \(u_{ij}=U_{ij}\) is stable with
\begin{equation}
 b_i={\|f_i\|^2\over2}.
 \label{eq:s7v56-B}
\end{equation}
\end{corollary}

\begin{proof}
For every subset \(I\), positivity gives
\[
 0\leq\sum_{i,j\in I}U_{ij}
 =\sum_{i\in I}\|f_i\|^2
 +2\sum_{\{i,j\}\subset I}\operatorname{Re}U_{ij}.
\]
Rearrange.
\end{proof}

\section{The dimension-free Weyl tree card}
\label{sec:s7v56-card}

\begin{theorem}[Connected Euclidean Weyl tree bound]
\label{thm:s7v56-card}
Let \(\kappa_0^{\rm W}\) be the cumulant formed from the ordered subset
moments \eqref{eq:s7v56-Euclidean}.  Then
\begin{equation}
 |\kappa_0^{\rm W}(f_1,\tau_1;\ldots;f_m,\tau_m)|
 \leq
 \sum_{T\in{\mathfrak T}_m}
 \prod_{\{i,j\}\in E(T)}
 \left|
 \langle f_i,e^{-|\tau_i-\tau_j|A}f_j\rangle
 \right|.
 \label{eq:s7v56-card}
\end{equation}
\end{theorem}

\begin{proof}
\Cref{thm:s7v56-pair} has the form
\eqref{eq:s7v56-moment}, with
\[
 m_i=e^{-\|f_i\|^2/2},\qquad
 \zeta_{ij}=e^{-U_{ij}}-1.
\]
Apply
\Cref{lem:s7v56-Mayer,lem:s7v56-treegraph,
cor:s7v56-stable}.  The tree-graph stability factor is
\[
 \exp\!\left(\sum_i{\|f_i\|^2\over2}\right),
\]
which cancels exactly against
\(\prod_i|m_i|=\exp(-\sum_i\|f_i\|^2/2)\).
\end{proof}

\begin{corollary}[Spacetime-local Weyl card]
\label{cor:s7v56-spacetime}
Suppose \(f_i\) is supported in one coordinate block \(x_i\), with
\(\|f_i\|\leq F\), and choose the V53 half-gap decay parameters
\(\mu,\gamma=\gamma_0/2\).  Then
\begin{align}
 |\kappa_0^{\rm W}|
 \leq
 \sum_{T\in{\mathfrak T}_m}
 \prod_{\{i,j\}\in E(T)}
 \left[
 F^2e^{-\mu d(x_i,x_j)}
 e^{-\gamma|\tau_i-\tau_j|}
 \right].
 \label{eq:s7v56-spacetime}
\end{align}
\end{corollary}

\begin{proof}
Apply the V53 kernel estimate
\eqref{eq:s7v53-half} to every edge in
\eqref{eq:s7v56-card}.
\end{proof}

\begin{corollary}[Dimension-free connected-series closure for bounded
Weyl amplitudes]
\label{cor:s7v56-series}
For Weyl vertices with \(\|f_i\|\leq F\), the V53 operator spacetime card
holds with a per-vertex constant depending on \(F\) but not on the
occupation cutoff \(M\) or the dimension \(D_{r,M}\).  Consequently the
V53 time and tree summation applies whenever the associated Weyl
superposition activity satisfies its explicit smallness condition.
\end{corollary}

\begin{proof}
Absorb \(F^2\) per tree edge into the V53 card constant.  Then invoke
\Cref{thm:s7v53-integrated}.
\end{proof}

\begin{firewall}[No occupation cutoff was differentiated]
\label{fire:s7v56-cutoff}
The proof uses exact quasi-free Weyl moments and never replaces \(P_M\) by
a smooth classical field cutoff.  It therefore avoids the V52 derivative
obstruction.  It has not yet represented \(P_MV_{\rm nl}P_M\) in this
Weyl class.
\end{firewall}

\section{Finite-rank spatial correction}
\label{sec:s7v56-finite}

If the one-particle pair kernel has the V50 decomposition
\[
 |\langle f_i,e^{-|\tau_i-\tau_j|A}f_j\rangle|
 \leq F^2e^{-\gamma|\tau_i-\tau_j|}
 \left[p_N(x_i,x_j)+{B\over N}\right],
\]
then \eqref{eq:s7v56-card} has exactly the corrected V51 edge.

\begin{proposition}[Weyl card with toron tail]
\label{prop:s7v56-toron}
Under the preceding two-scale pair bound, the connected Weyl series has a
finite-volume estimate of the form
\begin{equation}
 C(F,\rho)\left[
 e^{-\nu d(x,y)}+{1\over N}
 \right]
 \label{eq:s7v56-toron}
\end{equation}
whenever the V53 smallness parameter is below one.  The \(N^{-1}\) term
vanishes in the thermodynamic limit.
\end{proposition}

\begin{proof}
Insert the two-scale edge into
\Cref{thm:s7v56-card}; then apply
\Cref{thm:s7v53-integrated}.
\end{proof}

\section{The remaining representation norm}
\label{sec:s7v56-representation}

For a local trace-class operator \(O\), define its Weyl characteristic
\[
 \widehat O(f):=\operatorname{Tr}(O\,W(-f)).
\]
Formally, Weyl inversion writes \(O\) as an integral of
\(\widehat O(f)W(f)\).  The stability exponential has already cancelled,
but every tree edge contributes one power of the Weyl label norm at each
endpoint.  The correct remaining quantities are therefore
degree-resolved moments.

\begin{definition}[Degree-resolved Weyl representation budget]
\label{def:s7v56-budget}
A family \(O_M\) has moments \({\cal W}_{M,k}\) if it admits
\begin{equation}
 O_M=\int W(f)\,d\nu_M(f),\qquad
 {\cal W}_{M,k}(O_M):=
 \int \|f\|^k\,d|\nu_M|(f)<\infty
 \quad(k\geq0).
 \label{eq:s7v56-budget}
\end{equation}
\end{definition}

\begin{proposition}[Sufficient representation interface]
\label{prop:s7v56-interface}
If every local compressed interaction has a Weyl representation with
\[
 {\cal W}_{M,k}(O_M)
 \leq C M^pR_0^k\|O_M\|
 \qquad(k\geq0)
\]
for fixed \(R_0,p,C\), then the Weyl tree theorem produces at most a fixed
polynomial occupation loss.  The logarithmic V55 schedule absorbs that
loss for every explicit operator-activity row.
\end{proposition}

\begin{proof}
Insert \eqref{eq:s7v56-budget} multilinearly into
\eqref{eq:s7v56-card}.  A vertex of tree degree \(k\) costs
\({\cal W}_{M,k}\); the factors \(R_0^k\) pair across edges and are
absorbed into the edge constant.  Total variation gives one polynomial
\(M^p\) loss per vertex.  Apply
\Cref{cor:s7v55-polynomial}.
\end{proof}

\begin{firewall}[The Weyl budget may be the new bottleneck]
\label{fire:s7v56-budget}
Finite rank implies a Weyl inversion formula but not the uniform
degree-resolved estimate in
\Cref{prop:s7v56-interface}.  Wigner functions of highly excited number
states oscillate, and their absolute moments may grow rapidly with both
\(M\) and \(k\).
The next note must compute this norm rather than infer it from
finite-dimensionality.
\end{firewall}

\section{Weyl-sector status}
\label{sec:s7v56-status}

\begin{theorem}[What V56 proves]
\label{thm:s7v56-result}
For the gapped number-conserving quadratic vacuum, bounded Weyl
insertions satisfy a normalized spacetime labelled-tree estimate with a
constant independent of the occupation cutoff and local matrix
dimension.  This follows from exact pair factorization, positivity and
stability of the Euclidean Ornstein--Uhlenbeck kernel, and the stable
tree-graph inequality.
\end{theorem}

\begin{proof}
Use
\Cref{thm:s7v56-pair,lem:s7v56-OU,
cor:s7v56-stable,thm:s7v56-card,
cor:s7v56-spacetime}.
\end{proof}

\begin{assessment}[Progress at seventh-series V56]
\label{ass:s7v56-status}
The operator spacetime card is now proved on a dimension-free bounded
subalgebra large enough to generate the local canonical algebra.  The
dynamical propagation and normalized tree combinatorics are no longer
hypotheses there.  The remaining bridge is quantitative rather than
formal: represent the compressed Wilson vertices in Weyl form with a
variation norm compatible with the logarithmic occupation schedule.
\end{assessment}

\begin{decision}[V57 acquisition order]
\label{dec:s7v56-V57}
V57 will compute the Weyl characteristic functions of truncated
one-mode matrix units using Laguerre functions, derive upper and lower
bounds for the weighted variation norm, and decide whether
\Cref{prop:s7v56-interface} can be polynomial in \(M\).  For several fixed
block modes it will tensorize only after the one-mode growth is known.
\end{decision}

\begin{firewall}[Claim boundary after seventh-series V56]
\label{fire:s7v56-claim}
V56 does not prove a polynomial Weyl representation budget for the
occupation projector or compressed Wilson vertices, the analogous
Bogoliubov/pairing theorem, normalized interacting chart exit,
slow-label sensitivity, weighted physical-source incidence, patchwise
coarse-action cocycles, terminal strong-coupling matching, cutoff removal,
OS reconstruction, or the full Yang--Mills mass gap.
\end{firewall}

\paragraph{Parent-version record.}
V56 was formed from
\texttt{Renewal\_Yang\_Mills\_Mass\_Gap\_Research\_Notes\_V55.tex}.
The V55 parent had \(23\,339\) logical source lines, \(837\,751\) bytes,
and SHA--256
\[
 \texttt{F6B2F2C52B74B45881FCEE3C9D752949AA4849B609EAA14E9E08B78A464DFD0D}.
\]
The next compulsory companion audit is V60.

% =============================================================================
% END APPENDED SEVENTH-SERIES V56 MATERIAL
% =============================================================================

% =============================================================================
% BEGIN APPENDED SEVENTH-SERIES V57 MATERIAL
% =============================================================================

\chapter{Weyl--Plancherel moments and degree-weighted tree summation}
\label{ch:s7v57}

\section{Purpose}
\label{sec:s7v57-context}

V56 proves the exact spacetime tree bound for Weyl vertices.  To transfer
it to a general operator on a finite occupation window, one must integrate
one power of the Weyl label for every incident tree edge.  A uniform
geometric moment bound, as proposed provisionally in V56, is impossible
for a Weyl characteristic density with unbounded phase-space support.

This note replaces that overstrong target by the actual moment growth.
Weyl--Plancherel converts phase-space powers into oscillator commutators.
The resulting degree weight grows like
\(\Gamma(k/2+c)\), but Prüfer's enumeration divides it by
\((k-1)!\) at each vertex.  The degree-weighted labelled-tree sum is
therefore only \(C^m(m-2)!\), exactly compatible with the normalized
\(1/(m-2)!\) interaction series.

\section{Local Weyl inversion}
\label{sec:s7v57-inversion}

Fix \(r\) oscillator modes in one spatial block.  For
\(z=(z_1,\ldots,z_r)\in{\mathbb C}^r\), put
\begin{equation}
 W(z):=\exp\left(\sum_{j=1}^r
 z_ja_j^*-\overline z_ja_j\right).
 \label{eq:s7v57-W}
\end{equation}
For a trace-class local operator \(O\), define
\begin{equation}
 \widehat O(z):=\operatorname{Tr}(O\,W(-z)).
 \label{eq:s7v57-characteristic}
\end{equation}

\begin{theorem}[Weyl inversion and Plancherel]
\label{thm:s7v57-Plancherel}
With \(d\lambda_r(z):=\pi^{-r}d^{2r}z\),
\begin{align}
 O&=\int_{{\mathbb C}^r}\widehat O(z)W(z)\,d\lambda_r(z),
 \label{eq:s7v57-inversion}\\
 \int_{{\mathbb C}^r}|\widehat O(z)|^2\,d\lambda_r(z)
 &=\|O\|_{\rm HS}^2.
 \label{eq:s7v57-Plancherel}
\end{align}
The first identity holds weakly, and in Hilbert--Schmidt norm for
Hilbert--Schmidt \(O\).
\end{theorem}

\begin{proof}
The canonical commutation relations give the distributional
orthogonality
\[
 \operatorname{Tr}(W(z)^*W(w))
 =\pi^r\delta^{(2r)}(w-z).
\]
Pairing the right side of \eqref{eq:s7v57-inversion} with an arbitrary
finite-rank operator recovers its Hilbert--Schmidt inner product with
\(O\).  This proves weak inversion on finite rank and then by density on
Hilbert--Schmidt operators.  Taking the Hilbert--Schmidt norm gives
\eqref{eq:s7v57-Plancherel}.
\end{proof}

\section{Phase-space powers as commutators}
\label{sec:s7v57-commutator}

\begin{lemma}[Weyl multiplier identities]
\label{lem:s7v57-multiplier}
Up to signs of modulus one,
\begin{align}
 z_j\widehat O(z)
 &=\widehat{[O,a_j]}(z),
 \label{eq:s7v57-z}\\
 \overline z_j\widehat O(z)
 &=\widehat{[O,a_j^*]}(z).
 \label{eq:s7v57-zbar}
\end{align}
\end{lemma}

\begin{proof}
The Weyl relations give
\[
 [a_j,W(-z)]=-z_jW(-z),\qquad
 [a_j^*,W(-z)]=-\overline z_jW(-z).
\]
Insert these identities in \eqref{eq:s7v57-characteristic} and use
cyclicity of the trace.  The convention-dependent signs disappear in all
norm estimates.
\end{proof}

Let \(E_{\boldsymbol p\boldsymbol q}\) be a matrix unit with
\(|\boldsymbol p|,|\boldsymbol q|\leq M\).

\begin{lemma}[Iterated commutator bound]
\label{lem:s7v57-commutator}
Every word \({\rm ad}_{b_1}\cdots{\rm ad}_{b_\ell}
E_{\boldsymbol p\boldsymbol q}\), where
\(b_j\) is a creation or annihilation operator in the fixed block, obeys
\begin{equation}
 \|
 {\rm ad}_{b_1}\cdots{\rm ad}_{b_\ell}
 E_{\boldsymbol p\boldsymbol q}
 \|_{\rm HS}
 \leq
 \bigl[2\sqrt{M+\ell+1}\bigr]^\ell.
 \label{eq:s7v57-commutator}
\end{equation}
\end{lemma}

\begin{proof}
A commutator is the difference of left and right multiplication.
At the \(s\)-th step, every matrix unit appearing has occupation indices
at most \(M+s\).  Left or right multiplication by a creation or
annihilation operator therefore multiplies its Hilbert--Schmidt norm by at
most \(\sqrt{M+s+1}\).  The triangle inequality supplies the factor two.
Multiply these bounds for \(s=1,\ldots,\ell\).
\end{proof}

\begin{theorem}[Weighted \(L^2\) characteristic moments]
\label{thm:s7v57-L2}
For every integer \(\ell\geq0\),
\begin{equation}
 \int_{{\mathbb C}^r}
 \|z\|^{2\ell}
 |\widehat E_{\boldsymbol p\boldsymbol q}(z)|^2
 \,d\lambda_r(z)
 \leq
 \bigl[C_r(M+\ell+1)\bigr]^\ell.
 \label{eq:s7v57-L2}
\end{equation}
\end{theorem}

\begin{proof}
Expand
\[
 \|z\|^{2\ell}
 =\left(\sum_{j=1}^r|z_j|^2\right)^\ell
\]
by the multinomial theorem.  Each summand is
\(|z^{\boldsymbol\alpha}|^2\) with
\(|\boldsymbol\alpha|=\ell\).  By
\Cref{lem:s7v57-multiplier,thm:s7v57-Plancherel}, its integral against
\(|\widehat E|^2\) is the squared Hilbert--Schmidt norm of an
\(\ell\)-fold commutator, up to an order choice.  Apply
\Cref{lem:s7v57-commutator}.  The multinomial coefficients sum to
\(r^\ell\), which is absorbed into \(C_r^\ell\).
\end{proof}

\section{Absolute characteristic moments}
\label{sec:s7v57-L1}

Define
\begin{equation}
 {\cal W}_k(O):=
 \int_{{\mathbb C}^r}
 \|z\|^k|\widehat O(z)|\,d\lambda_r(z).
 \label{eq:s7v57-Wk}
\end{equation}

\begin{theorem}[Matrix-unit absolute moment bound]
\label{thm:s7v57-matrix-moment}
There are constants \(C_r,c_r>0\), independent of \(M,k\), such that
\begin{align}
 {\cal W}_k(E_{\boldsymbol p\boldsymbol q})
 &\leq
 C_r^{k+1}(M+1)^{k/2+(r+1)/2}
 \Gamma\left({k\over2}+c_r\right),
 \qquad k\geq0.
 \label{eq:s7v57-matrix-moment}
\end{align}
One may take \(c_r=r+2\) after enlarging \(C_r\).
\end{theorem}

\begin{proof}
Since the real dimension of phase space is \(2r\),
\[
 c_r':=
 \int_{{\mathbb C}^r}
 (1+\|z\|^2)^{-(r+1)}d\lambda_r(z)<\infty.
\]
Cauchy--Schwarz gives
\begin{align}
 {\cal W}_k(E)^2
 \leq c_r'
 \int(1+\|z\|^2)^{r+1}\|z\|^{2k}
 |\widehat E(z)|^2\,d\lambda_r(z).
 \label{eq:s7v57-CS}
\end{align}
Expand the fixed integer power \(r+1\) and apply
\Cref{thm:s7v57-L2} with
\(\ell=k,k+1,\ldots,k+r+1\).  This yields
\begin{equation}
 {\cal W}_k(E)
 \leq
 C_r^{k+1}
 (M+k+r+2)^{(k+r+1)/2}.
 \label{eq:s7v57-raw}
\end{equation}
Now
\[
 M+k+r+2\leq(M+1)(k+r+2).
\]
Stirling's lower bound for the gamma function shows, after increasing the
fixed \(C_r\), that
\[
 (k+r+2)^{(k+r+1)/2}
 \leq C_r^{k+1}
 \Gamma\left({k\over2}+r+2\right).
\]
Insert these two estimates in \eqref{eq:s7v57-raw}.
\end{proof}

\begin{corollary}[General occupation-window operator]
\label{cor:s7v57-operator}
If \(O=P_MOP_M\) acts on the \(r\)-mode window, then
\begin{align}
 {\cal W}_k(O)
 \leq&
 C_r^{k+1}D_{r,M}^{3/2}
 (M+1)^{k/2+(r+1)/2}
 \Gamma\left({k\over2}+r+2\right)\|O\|.
 \label{eq:s7v57-operator}
\end{align}
\end{corollary}

\begin{proof}
Expand \(O\) in matrix units, apply
\Cref{lem:s7v54-lone}, and use the triangle inequality in
\eqref{eq:s7v57-Wk}.
\end{proof}

\begin{firewall}[The provisional geometric moment target is impossible]
\label{fire:s7v57-geometric}
The characteristic measure in Weyl inversion has unbounded support for
every nonzero finite-rank oscillator operator.  If
\({\cal W}_k\leq CR_0^k\) held for every \(k\), then the \(L^k\) norms of
\(\|z\|\) under its normalized total variation would remain bounded by
\(R_0\), forcing that measure to be supported in
\(\{\|z\|\leq R_0\}\).  Thus the sufficient condition proposed in V56 is
not the correct target.  The gamma growth in
\eqref{eq:s7v57-operator} must be summed with the tree degrees.
\end{firewall}

\section{Prüfer summation of the degree weights}
\label{sec:s7v57-Prufer}

Put
\begin{equation}
 w_d:=\Gamma\left({d\over2}+r+2\right),\qquad d\geq1.
 \label{eq:s7v57-wd}
\end{equation}

\begin{lemma}[Summability of the shifted degree sequence]
\label{lem:s7v57-generating}
The constant
\begin{equation}
 A_r:=\sum_{e=0}^\infty {w_{e+1}\over e!}
 \label{eq:s7v57-Ar}
\end{equation}
is finite.
\end{lemma}

\begin{proof}
Treat even and odd \(e\) separately.  The ratio of the term at \(e+2\)
to that at \(e\) is
\[
 {\,{(e+1)/2+r+2}\,\over(e+1)(e+2)}
 =O(e^{-1}).
\]
Both parity subseries converge by the ratio test.
\end{proof}

\begin{theorem}[Degree-weighted labelled-tree count]
\label{thm:s7v57-Prufer}
For \(m\geq2\),
\begin{equation}
 \sum_{T\in{\mathfrak T}_m}
 \prod_{i=1}^m w_{\deg_T(i)}
 \leq
 (m-2)!A_r^m.
 \label{eq:s7v57-Prufer}
\end{equation}
\end{theorem}

\begin{proof}
For positive integers \(d_i\) with
\(\sum_i d_i=2m-2\), the Prüfer code gives exactly
\[
 {(m-2)!\over\prod_i(d_i-1)!}
\]
labelled trees with degree sequence \((d_i)\).  Put \(e_i=d_i-1\);
then \(\sum_ie_i=m-2\).  Therefore the left side of
\eqref{eq:s7v57-Prufer} equals
\[
 (m-2)!
 \sum_{\substack{e_i\geq0\\\sum e_i=m-2}}
 \prod_{i=1}^m{w_{e_i+1}\over e_i!}.
\]
The constrained positive sum is at most the unconstrained product
\(A_r^m\).
\end{proof}

\section{Connected operator series}
\label{sec:s7v57-series}

Let \(q_N(x,y)e^{-\gamma|\tau-\sigma|}\) be the V53 two-scale spacetime
edge, with spatial row constant \(R\).  For a local operator
\(O_i=P_MO_iP_M\), insert its Weyl inversion into the exact V56 cumulant.
At a vertex of tree degree \(d_i\), use
\eqref{eq:s7v57-operator}.

\begin{theorem}[Occupation-window quasi-free tree bound after degree
summation]
\label{thm:s7v57-operator-tree}
For \(m\geq2\), after summing labelled tree topologies and the \(m-2\)
internal spacetime positions as in V51 and V53, the normalized factor
\(1/(m-2)!\) cancels the degree-weighted tree count.  The resulting
per-internal-vertex loss is bounded by
\begin{equation}
 {\cal C}_{r,M}
 \leq C_r
 D_{r,M}^{3/2}(M+1)^{(r+3)/2}
 \leq C_r(M+1)^{2r+3/2}.
 \label{eq:s7v57-CM}
\end{equation}
More precisely, the product of the degree-dependent powers in
\eqref{eq:s7v57-operator} contributes
\((M+1)^{m-1}\), one factor for every tree edge, while the remaining
per-vertex factor is
\(D_{r,M}^{3/2}(M+1)^{(r+1)/2}\).
\end{theorem}

\begin{proof}
The V56 edge is bilinear in the two Weyl labels.  Thus a vertex of degree
\(d_i\) costs \({\cal W}_{d_i}(O_i)\).  Multiplying
\eqref{eq:s7v57-operator} over vertices gives
\[
 \prod_i(M+1)^{d_i/2}
 =(M+1)^{m-1},
\]
because \(\sum_i d_i=2m-2\).  The gamma factors are summed by
\Cref{thm:s7v57-Prufer}, and its \((m-2)!\) is cancelled by the connected
series normalization.  The fixed exponential constants, time row
\(2/\gamma\), spatial row \(R\), and \(A_r\) are absorbed into \(C_r\).
Each internal vertex may be charged one of the \(m-1\) edge factors
\(M+1\), giving the first bound in \eqref{eq:s7v57-CM}.  Finally use
\(D_{r,M}\leq(M+1)^r\).
\end{proof}

\begin{corollary}[Convergence under the V55 logarithmic window]
\label{cor:s7v57-convergence}
Let every local interaction insertion have norm at most \(\rho_M\).  If
\begin{equation}
 C_r(M+1)^{2r+3/2}\rho_M<1
 \label{eq:s7v57-small}
\end{equation}
with the fixed time and spatial row constants included in \(C_r\), then
the absolute quasi-free connected operator series converges.  Under the
V55 schedule \(M=(\log L)^\chi\), every explicit activity row containing
a negative power of \(L\), including the fifth-order magnetic tail,
satisfies this condition for sufficiently large \(L\).
\end{corollary}

\begin{proof}
\Cref{thm:s7v57-operator-tree} makes each added interaction cost at most
the left side of \eqref{eq:s7v57-small}.  Sum the resulting geometric
series.  The asymptotic assertion is
\Cref{cor:s7v55-polynomial} with \(p=2r+3/2\).
\end{proof}

\begin{firewall}[Endpoint norms and the dense OS core]
\label{fire:s7v57-endpoint}
The generic bound \eqref{eq:s7v57-CM} also charges polynomial \(M\)-factors
to the two external operators.  A cutoff-uniform clustering theorem on the
dense OS core needs endpoint-specific Weyl moment bounds or a declared
renormalization whose norms remain uniform.  Interaction-series
convergence alone does not prove that endpoint statement.
\end{firewall}

\begin{firewall}[Reference state versus Feshbach-corrected dynamics]
\label{fire:s7v57-reference}
\Cref{thm:s7v57-operator-tree} is a theorem for the quadratic quasi-free
reference.  In the full programme, self-energies are inserted into the low
Hamiltonian and the occupation complement is returned by Feshbach.  One
must show that these operations preserve the spacetime tree constants or
are accounted for as small activities.  That stability step is not
contained in Weyl--Plancherel alone.
\end{firewall}

\section{Weyl-moment status}
\label{sec:s7v57-status}

\begin{theorem}[What V57 proves]
\label{thm:s7v57-result}
Every operator on a fixed-mode occupation window has explicit
degree-resolved Weyl characteristic moments.  Their gamma growth is
exactly compatible with labelled-tree normalization after Prüfer
summation.  The resulting quasi-free connected series has only the fixed
polynomial occupation loss \(O(M^{2r+3/2})\) per interaction vertex, which
is absorbed by the V55 logarithmic schedule.
\end{theorem}

\begin{proof}
Use
\Cref{thm:s7v57-Plancherel,thm:s7v57-L2,
thm:s7v57-matrix-moment,cor:s7v57-operator,
thm:s7v57-Prufer,thm:s7v57-operator-tree,
cor:s7v57-convergence}.
\end{proof}

\begin{assessment}[Progress at seventh-series V57]
\label{ass:s7v57-status}
The Weyl representation bottleneck is closed for fixed local mode number:
the characteristic measure is not polynomial in every degree separately,
but its actual gamma growth is summable against the exact tree-degree
multiplicity.  This proves quasi-free connected-series convergence for
general occupation-window interactions under the revised schedule.
The next problem is stability under the nonlinear Feshbach self-energy and
uniform control of physical endpoint observables.
\end{assessment}

\begin{decision}[V58 acquisition order]
\label{dec:s7v57-V58}
V58 will first derive endpoint-specific Weyl moments for fixed-degree
gauge-invariant polynomial observables and Wilson exponentials, avoiding
the generic matrix-algebra factor.  It will then formulate a resolvent
identity comparing the quadratic and self-energy-corrected low
propagators, with a weighted spacetime norm smallness condition that
preserves the V57 tree constant.
\end{decision}

\begin{firewall}[Claim boundary after seventh-series V57]
\label{fire:s7v57-claim}
V57 does not prove cutoff-uniform endpoint bounds on the dense OS core,
stability of the quasi-free tree theorem under the Feshbach self-energy,
the analogous pairing/Bogoliubov theorem, normalized interacting chart
exit, slow-label sensitivity, weighted physical-source incidence,
patchwise coarse-action cocycles, terminal strong-coupling matching,
cutoff removal, OS reconstruction, or the full Yang--Mills mass gap.
\end{firewall}

\paragraph{Parent-version record.}
V57 was formed from
\texttt{Renewal\_Yang\_Mills\_Mass\_Gap\_Research\_Notes\_V56.tex}.
The V56 parent had \(23\,843\) logical source lines, \(853\,375\) bytes,
and SHA--256
\[
 \texttt{C88EAF4D16DBB3691769EAF2ED68C90C089E26FC05AF3876B62910A2689E2F8D}.
\]
The next compulsory companion audit is V60.

% =============================================================================
% END APPENDED SEVENTH-SERIES V57 MATERIAL
% =============================================================================

% =============================================================================
% BEGIN APPENDED SEVENTH-SERIES V58 MATERIAL
% =============================================================================

\chapter{Uniform Weyl endpoints and quadratic self-energy stability}
\label{ch:s7v58}

\section{Purpose}
\label{sec:s7v58-context}

V57 controls arbitrary occupation-window interaction vertices, but its
generic matrix-unit bound also charges a polynomial \(M\)-factor to
external observables.  Such a factor is harmless for convergence at fixed
cutoff but not for a cutoff-uniform dense-core clustering theorem.  The
first part of this note selects a dense gauge-invariant Weyl-cylinder core
whose fixed endpoints carry no occupation loss.

The second part addresses the renormalized closed returns of V42.  A
self-energy is not discarded as a small activity.  When its quadratic
local part is small in a weighted kernel norm, it is inserted into the
one-particle reference.  A weighted semigroup perturbation theorem proves
that the spacetime edge and its gap persist.  Nonquadratic and delocalized
parts remain separate ledger entries.

\section{A dense Weyl-cylinder endpoint core}
\label{sec:s7v58-core}

On one oscillator coordinate space, let \(\Omega(q)>0\) be the Gaussian
vacuum wavefunction.  For \(p\in{\mathbb R}^d\), multiplication by
\(e^{ip\cdot q}\) is a Weyl operator with a purely configuration-space
label.

\begin{lemma}[Density of Gaussian exponential vectors]
\label{lem:s7v58-density}
The linear span
\begin{equation}
 {\cal E}:=
 \operatorname{span}\{e^{ip\cdot q}\Omega(q):
 p\in{\mathbb R}^d\}
 \label{eq:s7v58-E}
\end{equation}
is dense in \(L^2({\mathbb R}^d,dq)\).
\end{lemma}

\begin{proof}
If \(h\in L^2(dq)\) is orthogonal to \({\cal E}\), then
\[
 \int_{\mathbb R^d}\overline{h(q)}\Omega(q)e^{ip\cdot q}\,dq=0
 \qquad\hbox{for every }p.
\]
The product \(\overline h\,\Omega\) lies in \(L^1\) by
Cauchy--Schwarz.  Its Fourier transform vanishes identically, so Fourier
uniqueness gives \(h\Omega=0\) almost everywhere.  Since
\(\Omega>0\), \(h=0\).
\end{proof}

\begin{corollary}[Finite-cylinder density]
\label{cor:s7v58-cylinder}
Finite products of local configuration Weyl operators applied to the
quadratic vacuum are dense in the finite-cylinder Fock space.
\end{corollary}

\begin{proof}
Apply \Cref{lem:s7v58-density} on each member of a finite tensor product;
finite tensors of dense sets are dense.
\end{proof}

Let \({\cal G}\) be the compact residual gauge group acting unitarily on a
finite cylinder and let
\begin{equation}
 {\mathbb P}_{\cal G}v:=
 \int_{\cal G}U(g)v\,dg
 \label{eq:s7v58-gauge-average}
\end{equation}
be normalized Haar averaging.

\begin{lemma}[Gauge averaging preserves density in the invariant sector]
\label{lem:s7v58-gauge}
The gauge averages of the Weyl-cylinder vectors are dense in
\(\operatorname{Ran}{\mathbb P}_{\cal G}\).
\end{lemma}

\begin{proof}
\({\mathbb P}_{\cal G}\) is a bounded orthogonal projection.  The image
under a continuous map of a dense set is dense in the image of the whole
space: approximate \(v\) by Weyl-cylinder vectors \(v_n\), then
\({\mathbb P}_{\cal G}v_n\to{\mathbb P}_{\cal G}v\).
\end{proof}

\begin{definition}[Reference Weyl OS core]
\label{def:s7v58-core}
The reference core \({\cal D}_{\rm W}^{\cal G}\) is the linear span of
finite-time, finite-spatial-cylinder products of configuration Weyl
operators, averaged by \({\mathbb P}_{\cal G}\).
\end{definition}

\begin{firewall}[Reference density versus interacting OS density]
\label{fire:s7v58-core}
\Cref{lem:s7v58-gauge} proves density in the gauge-invariant quadratic
reference Hilbert space.  To use the V25 mass-gap criterion after cutoff
removal, one must identify this core with a dense core of the limiting
interacting OS Hilbert space.  Uniform integrability and null-space
compatibility are still required.
\end{firewall}

\section{Endpoint-specific tree summation}
\label{sec:s7v58-endpoint}

Fix two Weyl endpoint labels \(f_A,f_B\), independent of \(M\), and let
\[
 F_{AB}:=\max\{1,\|f_A\|,\|f_B\|\}.
\]
The \(n\) interaction vertices retain the V57 characteristic moment
bound.

\begin{theorem}[No occupation loss at fixed Weyl endpoints]
\label{thm:s7v58-endpoint}
In the V57 degree-weighted tree calculation with
\(m=n+2\) vertices, replacing the two generic endpoint operators by the
fixed Weyl operators \(W(f_A),W(f_B)\) removes their factors
\[
 D_{r,M}^{3/2}(M+1)^{(r+1)/2}.
\]
After summing tree degrees and dividing by \(n!\), the connected series has
an \(M\)-independent endpoint prefactor
\(C(f_A,f_B)\).  Every added interaction still costs at most
\begin{equation}
 C_r(M+1)^{2r+3/2}\rho_M.
 \label{eq:s7v58-cost}
\end{equation}
\end{theorem}

\begin{proof}
At an endpoint of degree \(d\), the exact V56 edge contributes only
\(\|f_A\|^d\) or \(\|f_B\|^d\), with no Weyl inversion integral.  Since
\[
 F_{AB}^d
 \leq C(F_{AB})^d
 \Gamma\left({d\over2}+r+2\right),
\]
the same Prüfer generating series used in
\Cref{thm:s7v57-Prufer} remains finite after changing its two endpoint
coefficients.  Its bound is
\(C(f_A,f_B)C_r^n n!\).  Division by \(n!\) cancels the factorial.
Only the \(n\) interaction vertices use
\Cref{cor:s7v57-operator}; assigning the edge factors as in
\Cref{thm:s7v57-operator-tree} gives
\eqref{eq:s7v58-cost}.
\end{proof}

\begin{corollary}[Uniform reference clustering on the Weyl core]
\label{cor:s7v58-clustering}
Under the V55 logarithmic schedule and the V57 smallness condition, the
quadratic-reference connected expansion between two fixed elements of
\({\cal D}_{\rm W}^{\cal G}\) has an \(M\)-uniform prefactor and the V51
spatial bound
\begin{equation}
 C(A,B)\left(e^{-\nu d(A,B)}+N^{-1}\right).
 \label{eq:s7v58-clustering}
\end{equation}
\end{corollary}

\begin{proof}
Expand the two fixed core elements into finitely many Weyl-cylinder terms,
apply \Cref{thm:s7v58-endpoint}, and then use the V51 path-resolved
finite-rank theorem.
\end{proof}

\section{Weighted perturbation of the one-particle reference}
\label{sec:s7v58-perturb}

Let \(A\) satisfy the V53 gap, range, and hopping assumptions.  For a
self-adjoint block kernel \(E\), define
\begin{align}
 \epsilon_0(E)&:=\|E\|_{{\cal S}_0},
 \label{eq:s7v58-e0}\\
 \epsilon_\mu(E)&:=
 \max\left\{
 \sup_x\sum_y(e^{\mu d(x,y)}-1)\|E_{xy}\|,
 \sup_y\sum_x(e^{\mu d(x,y)}-1)\|E_{xy}\|
 \right\}.
 \label{eq:s7v58-emu}
\end{align}
The diagonal part contributes to \(\epsilon_0\) but not
\(\epsilon_\mu\).

\begin{theorem}[Weighted quadratic self-energy stability]
\label{thm:s7v58-stability}
Let \(A_E=A+E\).  If
\begin{equation}
 \Gamma_E:=
 \gamma_0-\epsilon_0(E)
 -J(e^{\mu R_0}-1)-\epsilon_\mu(E)>0,
 \label{eq:s7v58-Gamma}
\end{equation}
then
\begin{equation}
 \|(e^{-tA_E})_{xy}\|
 \leq e^{-\mu d(x,y)}e^{-\Gamma_Et},
 \qquad t\geq0.
 \label{eq:s7v58-semigroup}
\end{equation}
\end{theorem}

\begin{proof}
The Schur test gives \(\|E\|\leq\epsilon_0(E)\), hence
\[
 A_E\geq(\gamma_0-\epsilon_0(E))I.
\]
Use the target-centred bounded exponential weight \(W_N\) of V53.  The
finite-range part obeys
\[
 \|W_NAW_N^{-1}-A\|
 \leq J(e^{\mu R_0}-1).
\]
For \(E\), the same row/column estimate gives
\[
 \|W_NEW_N^{-1}-E\|
 \leq\epsilon_\mu(E),
\]
because each kernel block is multiplied by a weight ratio whose deviation
from one is at most \(e^{\mu d(x,y)}-1\).
Thus the real part of the conjugated \(A_E\) is at least \(\Gamma_E\).
The energy estimate and target-weight orientation in the proof of
\Cref{thm:s7v53-semigroup} now give
\eqref{eq:s7v58-semigroup}.
\end{proof}

\begin{corollary}[Uniform persistence of the quasi-free tree edge]
\label{cor:s7v58-edge}
If
\begin{equation}
 \epsilon_0(E)+\epsilon_\mu(E)\longrightarrow0
 \label{eq:s7v58-smallE}
\end{equation}
uniformly in volume and slow labels, then any fixed
\(\mu\) satisfying the strict V53 half-gap condition gives
\(\Gamma_E\geq\gamma_0/4\) for all sufficiently small cutoff.  The V56
Weyl tree theorem and V57 operator-series theorem then hold with uniform
modified constants.
\end{corollary}

\begin{proof}
Choose \(\mu\) so that
\(J(e^{\mu R_0}-1)\leq\gamma_0/2\).  Eventually the two \(E\)-terms sum to
at most \(\gamma_0/4\).  Apply
\Cref{thm:s7v58-stability}; the V56 proof uses only positivity and the
resulting semigroup edge.
\end{proof}

\section{Integrated difference of the spacetime edge}
\label{sec:s7v58-difference}

Assume the graph has a uniform volume-growth sum at every exponent below
\(\mu\).  Fix \(0<\nu<\mu\).

\begin{lemma}[Rate-loss conversion to a weighted Schur semigroup]
\label{lem:s7v58-rate-loss}
If a kernel obeys
\[
 \|K_t(x,y)\|\leq e^{-\mu d(x,y)}e^{-\Gamma t},
\]
then
\begin{equation}
 \|K_t\|_{{\cal S}_\nu}
 \leq Q_{\mu-\nu}e^{-\Gamma t},
 \qquad
 Q_{\mu-\nu}:=
 \sup_x\sum_y e^{-(\mu-\nu)d(x,y)}.
 \label{eq:s7v58-rate-loss}
\end{equation}
\end{lemma}

\begin{proof}
Multiply the pointwise estimate by \(e^{\nu d(x,y)}\) and sum rows and
columns.
\end{proof}

\begin{theorem}[Integrated semigroup response]
\label{thm:s7v58-response}
Suppose both \(A\) and \(A_E\) have the preceding \({\cal S}_\nu\)
semigroup bounds with common decay \(\Gamma>0\).  Then
\begin{equation}
 \int_0^\infty
 \|e^{-tA_E}-e^{-tA}\|_{{\cal S}_\nu}\,dt
 \leq
 {Q_{\mu-\nu}^2\over\Gamma^2}
 \|E\|_{{\cal S}_\nu}.
 \label{eq:s7v58-response}
\end{equation}
\end{theorem}

\begin{proof}
Duhamel's identity gives
\[
 e^{-tA_E}-e^{-tA}
 =
 -\int_0^t e^{-(t-s)A_E}Ee^{-sA}\,ds.
\]
The weighted Schur algebra is submultiplicative.  Apply
\Cref{lem:s7v58-rate-loss} to both semigroups:
\[
 \|e^{-tA_E}-e^{-tA}\|_{{\cal S}_\nu}
 \leq Q_{\mu-\nu}^2
 t e^{-\Gamma t}\|E\|_{{\cal S}_\nu}.
\]
Integrate \(t e^{-\Gamma t}\), whose integral is \(\Gamma^{-2}\).
\end{proof}

\section{Application to the occupation self-energy}
\label{sec:s7v58-selfenergy}

Let the dimensionless occupation self-energy be decomposed as
\begin{equation}
 a\Sigma_{\rm occ}(0)
 =
 e_{\rm vac}I+d\Gamma(E_{\rm quad})+V_{\rm rem}.
 \label{eq:s7v58-decomposition}
\end{equation}
The scalar \(e_{\rm vac}\) is absorbed into the common vacuum
normalization.

\begin{definition}[Quadratic self-energy card]
\label{def:s7v58-QSEC}
QSEC consists of
\begin{align}
 \|E_{\rm quad}\|_{{\cal S}_0}
 +\|E_{\rm quad}\|_{{\cal S}_\mu}
 &\leq Cg^2M^2,
 \label{eq:s7v58-QSEC}\\
 \|V_{\rm rem}\|_{\rm activity}
 &=o(1)
 \label{eq:s7v58-rem}
\end{align}
in the V57 connected activity norm, uniformly in slow labels.
\end{definition}

\begin{proposition}[Conditional self-energy insertion]
\label{prop:s7v58-selfenergy}
Under QSEC and the V55 logarithmic schedule, the quadratic self-energy can
be inserted into the reference while preserving a uniform spacetime edge,
and its integrated change tends to zero in every lower weighted rate
\(\nu<\mu\).
\end{proposition}

\begin{proof}
The bound is
\[
 g^2M^2=O\!\left(L^{-1}(\log L)^{2\chi}\right)\to0.
\]
Apply
\Cref{cor:s7v58-edge,thm:s7v58-response}.
\end{proof}

\begin{firewall}[QSEC is a decomposition theorem, not notation]
\label{fire:s7v58-QSEC}
V42 bounds the total self-energy norm after a closed return.  It does not
prove that its quadratic part has the weighted locality in
\eqref{eq:s7v58-QSEC}, nor that the nonquadratic remainder is a small
connected activity.  Those properties require a support-resolved
Feshbach-kernel expansion.  Writing
\eqref{eq:s7v58-decomposition} does not establish them.
\end{firewall}

\begin{firewall}[Delocalized slow terms remain outside \(E_{\rm quad}\)]
\label{fire:s7v58-delocalized}
A toron contribution with a \(1/N\) kernel is not placed in
\(\epsilon_\mu(E_{\rm quad})\), whose exponential weight would diverge.
It remains in the V50 finite-rank ledger and is handled by the V51
path-resolved thermodynamic limit.
\end{firewall}

\section{Endpoint and self-energy status}
\label{sec:s7v58-status}

\begin{theorem}[What V58 proves]
\label{thm:s7v58-result}
Gauge-averaged finite Weyl cylinders form a dense invariant core for the
quadratic reference.  Fixed Weyl endpoints have a cutoff-uniform
connected-series prefactor, while arbitrary occupation-window interaction
vertices retain the convergent V57 polynomial loss.  A small local
quadratic self-energy in the declared weighted norm preserves both the
one-particle gap and the spacetime tree edge, with the integrated response
bound \eqref{eq:s7v58-response}.
\end{theorem}

\begin{proof}
Use
\Cref{lem:s7v58-density,lem:s7v58-gauge,
thm:s7v58-endpoint,cor:s7v58-clustering,
thm:s7v58-stability,thm:s7v58-response}.
\end{proof}

\begin{assessment}[Progress at seventh-series V58]
\label{ass:s7v58-status}
The quasi-free connected expansion now has cutoff-uniform endpoints on a
dense reference core, eliminating the endpoint defect left in V57.
Quadratic renormalization is stable under a precise weighted
self-energy card.  The next upstream task is to prove that card from the
support-resolved occupation Feshbach map rather than from its global norm.
\end{assessment}

\begin{decision}[V59 acquisition order]
\label{dec:s7v58-V59}
V59 will expand the occupation Feshbach self-energy
\(K_MR_M(0)K_M^*\) by local support labels, retain the resolvent
propagation between the two vertices, and derive an exponential
weighted-row estimate from the high-occupation gap.  It will separate the
vacuum, quadratic, higher local, and finite-rank slow contributions needed
by QSEC.  If the high-occupation resolvent lacks a Lieb--Robinson or
Combes--Thomas estimate in configuration space, that missing propagation
will be isolated explicitly.
\end{decision}

\begin{firewall}[Claim boundary after seventh-series V58]
\label{fire:s7v58-claim}
V58 does not prove QSEC, support-resolved high-occupation resolvent
locality, density of the Weyl core in the limiting interacting OS space,
the analogous pairing/Bogoliubov theorem, normalized interacting chart
exit, slow-label sensitivity, weighted physical-source incidence,
patchwise coarse-action cocycles, terminal strong-coupling matching,
cutoff removal, OS reconstruction, or the full Yang--Mills mass gap.
\end{firewall}

\paragraph{Parent-version record.}
V58 was formed from
\texttt{Renewal\_Yang\_Mills\_Mass\_Gap\_Research\_Notes\_V57.tex}.
The V57 parent had \(24\,303\) logical source lines, \(868\,452\) bytes,
and SHA--256
\[
 \texttt{53E2B4F4EF02FC1A6ADCA56374B801B4B88D1C575905565B9482322575387CC9}.
\]
The next compulsory companion audit is V60.

% =============================================================================
% END APPENDED SEVENTH-SERIES V58 MATERIAL
% =============================================================================

% =============================================================================
% BEGIN APPENDED SEVENTH-SERIES V59 MATERIAL
% =============================================================================

\chapter{Occupation-boundary support of the Feshbach self-energy}
\label{ch:s7v59}

\section{Purpose and correction of the preceding acquisition plan}
\label{sec:s7v59-purpose}

V58 proposed to decompose the occupation self-energy into a vacuum
constant, a quadratic one-particle correction, and a small nonlinear
remainder.  That is a legitimate abstract stability interface, but it is
not the first structure one should seek for the particular Feshbach map
which eliminates occupations above \(M\).  A polynomial of number degree
at most \(r\) can cross the boundary \(M\) only from the top \(r\) number
shells.  Consequently its closed Feshbach return is supported on those
same shells.  For \(M>r+1\) it vanishes identically on the vacuum and
one-particle spaces.

This elementary selection rule changes the direction of the programme.
The cubic--quartic occupation self-energy is not a bulk quadratic
renormalization.  It is an artificial occupation-boundary operator which
is small relative to the free number form.  The analytic fifth-order
chart remainder does not have finite number bandwidth and is therefore
kept as a separate norm-small error.  The present note proves these
claims exactly, corrects the use of the V58 quadratic self-energy card,
and states the genuinely missing locality input without pretending that
a spectral gap alone supplies it.

\section{Number bandwidth}
\label{sec:s7v59-bandwidth}

Let \({\cal N}\) denote the total oscillator number operator and put
\[
 P_{\leq n}:=\mathbf 1_{[0,n]}({\cal N}),\qquad
 P_M:=P_{\leq M},\qquad Q_M:=I-P_M.
\]
For an integer \(r\geq1\) and \(M\geq r\), define the upper boundary band
\begin{equation}
 B_{M,r}:=
 \mathbf 1_{(M-r,M]}({\cal N})
 =P_M-P_{\leq M-r}.
 \label{eq:s7v59-band}
\end{equation}

\begin{definition}[Number bandwidth]
\label{def:s7v59-bandwidth}
An operator \(V\), defined at least on the finite-particle core, has
number bandwidth at most \(r\) if
\begin{equation}
 \mathbf 1_{\{n\}}({\cal N})V\mathbf 1_{\{m\}}({\cal N})=0
 \quad\hbox{whenever}\quad |n-m|>r .
 \label{eq:s7v59-bandwidth}
\end{equation}
\end{definition}

\begin{lemma}[Polynomial number-band lemma]
\label{lem:s7v59-polynomial-band}
Every normally ordered monomial containing at most \(r\) creation and
annihilation operators has number bandwidth at most \(r\).  Any finite
linear combination of such monomials has the same property.
\end{lemma}

\begin{proof}
A monomial with \(p\) creation and \(q\) annihilation operators changes
the number eigenvalue by exactly \(p-q\) on every number eigenspace on
which it is nonzero.  Since \(p+q\leq r\), one has
\(|p-q|\leq r\).  Taking a finite linear combination cannot create a
matrix element outside the union of the individual bandwidths.
\end{proof}

\begin{lemma}[One-crossing boundary identity]
\label{lem:s7v59-crossing}
If \(V\) has number bandwidth at most \(r\), then
\begin{align}
 Q_MVP_M&=Q_MVB_{M,r},                                  \label{eq:s7v59-QVP}\\
 P_MVQ_M&=B_{M,r}VQ_M.                                  \label{eq:s7v59-PVQ}
\end{align}
The second identity follows from the first for \(V^*\), and in
particular both hold when \(V=V^*\).
\end{lemma}

\begin{proof}
If \(0\leq n\leq M-r\), the bandwidth condition gives
\[
 V\mathbf 1_{\{n\}}({\cal N})
 =
 \mathbf 1_{[0,n+r]}({\cal N})
 V\mathbf 1_{\{n\}}({\cal N}),
\]
and \(n+r\leq M\).  Thus
\(Q_MVP_{\leq M-r}=0\), which is
\eqref{eq:s7v59-QVP}.  Apply this identity to \(V^*\) and take adjoints
to obtain \eqref{eq:s7v59-PVQ}.
\end{proof}

\begin{remark}[Why the analytic remainder is different]
\label{rem:s7v59-analytic}
A nonpolynomial Wilson multiplication operator generally has matrix
elements between number states at arbitrarily large separation.  The
chart localization used in V55 makes its fifth Taylor remainder small in
the relevant \(P_M\)-to-\(Q_M\) block, but it does not give finite number
bandwidth.  Therefore
\Cref{lem:s7v59-crossing} may be applied to \(H_3+H_4\), not to the
unexpanded analytic remainder.
\end{remark}

\section{Exact shell support of the polynomial self-energy}
\label{sec:s7v59-shell}

Write the interaction as
\begin{equation}
 W_a=W_{{\rm pol},a}+W_{{\rm tail},a},\qquad
 W_{{\rm pol},a}:={1\over a}(gH_3+g^2H_4),
 \label{eq:s7v59-split}
\end{equation}
and set
\begin{align}
 K_{\rm pol}&:=P_MW_{{\rm pol},a}Q_M,                   \label{eq:s7v59-Kpol}\\
 K_{\rm tail}&:=P_MW_{{\rm tail},a}Q_M,                 \label{eq:s7v59-Ktail}\\
 R_M(z)&:=\bigl(Q_M(H_a-z)Q_M\bigr)^{-1}.               \label{eq:s7v59-R}
\end{align}
The V42 high-occupation floor ensures that \(R_M(z)\) is positive for
real \(z\) with
\(\delta_M(z):=\Delta_M(a)-az>0\).
Put
\begin{equation}
 \Sigma_{\alpha\beta}(z)
 :=K_\alpha R_M(z)K_\beta^*,
 \qquad \alpha,\beta\in\{{\rm pol},{\rm tail}\}.
 \label{eq:s7v59-Sab}
\end{equation}

\begin{theorem}[Occupation-boundary self-energy theorem]
\label{thm:s7v59-shell}
Assume that \(H_3\) and \(H_4\) have number bandwidth at most three and
four, respectively.  Then, with \(B_M:=B_{M,4}\),
\begin{align}
 K_{\rm pol}&=B_MK_{\rm pol},                            \label{eq:s7v59-BK}\\
 K_{\rm pol}^*&=K_{\rm pol}^*B_M,                       \label{eq:s7v59-KstarB}\\
 \Sigma_{{\rm pol},{\rm pol}}(z)
 &=B_M\Sigma_{{\rm pol},{\rm pol}}(z)B_M.               \label{eq:s7v59-BSigmaB}
\end{align}
In particular, if \(M>5\), then
\begin{equation}
 P_{\leq1}\Sigma_{{\rm pol},{\rm pol}}(z)
 =\Sigma_{{\rm pol},{\rm pol}}(z)P_{\leq1}=0.
 \label{eq:s7v59-lowzero}
\end{equation}
More generally, the self-energy vanishes on both sides of
\(P_{\leq M-4}\).
\end{theorem}

\begin{proof}
\Cref{lem:s7v59-polynomial-band} gives bandwidth at most four for
\(gH_3+g^2H_4\).  The free quadratic part commutes with \(P_M\), so only
the interaction occurs in the off-diagonal block.
\Cref{lem:s7v59-crossing} gives
\[
 P_MW_{{\rm pol},a}Q_M
 =B_MW_{{\rm pol},a}Q_M=B_MK_{\rm pol},
\]
and its adjoint gives \eqref{eq:s7v59-KstarB}.  Insert these identities
on the two sides of \(R_M(z)\) to prove
\eqref{eq:s7v59-BSigmaB}.  Since
\(P_{\leq1}B_M=0\) when \(M-4\geq1\), the stated low-sector vanishing
follows.  Replacing \(P_{\leq1}\) by \(P_{\leq M-4}\) proves the final
claim.
\end{proof}

\begin{corollary}[No polynomial occupation correction to the vacuum or
one-particle operator]
\label{cor:s7v59-noquadratic}
For \(M>5\), the vacuum and one-particle matrix elements of
\(\Sigma_{{\rm pol},{\rm pol}}(z)\) are exactly zero.  Hence this
particular closed return cannot supply either the vacuum constant or the
one-particle operator \(d\Gamma(E_{\rm quad})\) postulated in the V58
decomposition.
\end{corollary}

\begin{proof}
Both the vacuum and the one-particle subspaces lie in
\(\operatorname{Ran}P_{\leq1}\).  Apply
\eqref{eq:s7v59-lowzero}.
\end{proof}

\begin{firewall}[Scope of the correction to QSEC]
\label{fire:s7v59-QSEC}
\Cref{cor:s7v59-noquadratic} concerns only the Feshbach map across the
total occupation boundary \(M\), generated by the finite-degree cubic
and quartic vertices.  Other coarse-graining operations can and do
generate bulk quadratic self-energies.  The V58 weighted semigroup
theorem remains valid for those corrections.  What is withdrawn is the
unproved identification of the present polynomial occupation return
with such a correction.
\end{firewall}

\section{A relative-form estimate at the artificial boundary}
\label{sec:s7v59-relative}

Let \(H_{0,a}\) be the positive quadratic oscillator Hamiltonian on the
transverse sector.  The standing oscillator floor is
\begin{equation}
 aH_{0,a}\geq\omega_*{\cal N}.
 \label{eq:s7v59-H0floor}
\end{equation}
Define the polynomial block size
\begin{equation}
 k_{{\rm pol},M}
 :=C_3g(M+4)^{3/2}+C_4g^2(M+5)^2,
 \qquad
 \|K_{\rm pol}\|\leq{k_{{\rm pol},M}\over a}.
 \label{eq:s7v59-kpol}
\end{equation}

\begin{lemma}[Boundary projection controlled by number energy]
\label{lem:s7v59-Bform}
For \(M>4\), as quadratic forms on \(\operatorname{Ran}P_M\),
\begin{equation}
 0\leq B_M
 \leq{{\cal N}\over M-4}
 \leq{aH_{0,a}\over\omega_*(M-4)}.
 \label{eq:s7v59-Bform}
\end{equation}
\end{lemma}

\begin{proof}
On every number eigenspace with \(B_M=1\), the eigenvalue of
\({\cal N}\) is at least \(M-3\), and hence certainly exceeds \(M-4\).
On its orthogonal complement the left side vanishes.  This proves the
first inequality by spectral calculus.  The second is
\eqref{eq:s7v59-H0floor}.
\end{proof}

\begin{theorem}[Relative-form smallness of the polynomial return]
\label{thm:s7v59-relative}
Let \(M>4\) and \(\delta_M(z)>0\).  Then
\begin{equation}
 0\leq a\Sigma_{{\rm pol},{\rm pol}}(z)
 \leq
 \varepsilon_{{\rm occ},M}(z)\,aH_{0,a}
 \quad\hbox{on }\operatorname{Ran}P_M,
 \label{eq:s7v59-relative}
\end{equation}
where
\begin{equation}
 \varepsilon_{{\rm occ},M}(z)
 :=
 {k_{{\rm pol},M}^2
  \over
  \delta_M(z)\,\omega_*(M-4)}.
 \label{eq:s7v59-epsilon}
\end{equation}
\end{theorem}

\begin{proof}
Positivity of \(Q_M(H_a-z)Q_M\) gives \(R_M(z)\geq0\), so the
self-energy is positive.  The V42 resolvent estimate and
\eqref{eq:s7v59-kpol} give
\[
 \|a\Sigma_{{\rm pol},{\rm pol}}(z)\|
 \leq {k_{{\rm pol},M}^2\over\delta_M(z)}.
\]
By \eqref{eq:s7v59-BSigmaB}, this operator is supported in
\(\operatorname{Ran}B_M\).  A positive bounded operator \(A=B_MAB_M\)
satisfies \(0\leq A\leq\|A\|B_M\).  Apply this fact and then
\Cref{lem:s7v59-Bform}.
\end{proof}

\begin{corollary}[Logarithmic-window rate]
\label{cor:s7v59-lograte}
Use the V55 schedule
\[
 g=O(L^{-1/2}),\qquad M=(\log L)^\chi,
\]
and suppose, uniformly for the spectral interval under consideration,
\(\delta_M(z)\geq cM\) with \(c>0\).  Then
\begin{equation}
 \varepsilon_{{\rm occ},M}(z)
 =
 O\!\left(L^{-1}(\log L)^\chi\right)
 +O\!\left(L^{-3/2}(\log L)^{3\chi/2}\right)
 +O\!\left(L^{-2}(\log L)^{2\chi}\right)
 =o(1).
 \label{eq:s7v59-lograte}
\end{equation}
\end{corollary}

\begin{proof}
Squaring \eqref{eq:s7v59-kpol} produces terms of orders
\[
 g^2M^3,\qquad g^3M^{7/2},\qquad g^4M^4.
\]
The denominator in \eqref{eq:s7v59-epsilon} is bounded below by a
constant times \(M^2\).  Division and substitution of the schedule give
the three displayed rates.  Every fixed power of \(\log L\) is
dominated by every positive power of \(L\), so all three tend to zero.
\end{proof}

\begin{corollary}[Derivative ledger]
\label{cor:s7v59-derivative}
Under the same hypotheses,
\begin{align}
 \Sigma_{{\rm pol},{\rm pol}}'(z)
 &=K_{\rm pol}R_M(z)^2K_{\rm pol}^*\geq0,               \label{eq:s7v59-Sprime}\\
 \|\Sigma_{{\rm pol},{\rm pol}}'(z)\|
 &\leq {k_{{\rm pol},M}^2\over\delta_M(z)^2}
 =O\!\left(L^{-1}(\log L)^\chi\right).                 \label{eq:s7v59-Sprime-rate}
\end{align}
It has the same exact \(B_M\)-support as
\(\Sigma_{{\rm pol},{\rm pol}}(z)\).
\end{corollary}

\begin{proof}
Differentiate the resolvent and use
\(R_M'(z)=R_M(z)^2\).  The norm bound is the V42 derivative estimate
with \(k_M\) replaced by \(k_{{\rm pol},M}\).  The boundary factors in
\eqref{eq:s7v59-BK}--\eqref{eq:s7v59-KstarB} remain on the outside.
Finally \(k_{{\rm pol},M}^2/\delta_M(z)^2
=O(g^2M)+O(g^3M^{3/2})+O(g^4M^2)\), which yields the stated leading
rate.
\end{proof}

\section{The analytic tail is norm-small and remains separate}
\label{sec:s7v59-tail}

Let
\begin{equation}
 \|K_{\rm tail}\|\leq{\rho_{{\rm tail},M}\over a}.
 \label{eq:s7v59-rhotail}
\end{equation}
For the chart-localized magnetic fifth remainder proved in V55 one may
take
\begin{equation}
 \rho_{{\rm tail},M}
 \leq Cg^3(M+6)^{5/2}.
 \label{eq:s7v59-rho5}
\end{equation}

\begin{proposition}[Mixed and tail return]
\label{prop:s7v59-tail}
For \(\delta_M(z)>0\),
\begin{align}
 a\bigl\|
 \Sigma_{{\rm pol},{\rm tail}}(z)
 +\Sigma_{{\rm tail},{\rm pol}}(z)
 \bigr\|
 &\leq{2k_{{\rm pol},M}\rho_{{\rm tail},M}
        \over\delta_M(z)},                              \label{eq:s7v59-cross}\\
 a\|\Sigma_{{\rm tail},{\rm tail}}(z)\|
 &\leq{\rho_{{\rm tail},M}^2\over\delta_M(z)}.          \label{eq:s7v59-tailtail}
\end{align}
Under the V55 logarithmic schedule and \(\delta_M(z)\geq cM\), the
right sides are respectively
\begin{align}
 O\!\left(L^{-2}(\log L)^{3\chi}\right),
 \qquad
 O\!\left(L^{-3}(\log L)^{4\chi}\right).                \label{eq:s7v59-tailrates}
\end{align}
\end{proposition}

\begin{proof}
Use
\(\|R_M(z)\|\leq a/\delta_M(z)\),
\(\|K_{\rm pol}\|\leq k_{{\rm pol},M}/a\), and
\eqref{eq:s7v59-rhotail} in each product
\eqref{eq:s7v59-Sab}.  This proves
\eqref{eq:s7v59-cross}--\eqref{eq:s7v59-tailtail}.
The leading polynomial block is
\(O(gM^{3/2})\), while
\(\rho_{{\rm tail},M}=O(g^3M^{5/2})\).  After division by
\(\delta_M(z)\asymp M\), the mixed term is \(O(g^4M^3)\) and the
tail--tail term is \(O(g^6M^4)\).  Substitute
\(g=O(L^{-1/2})\) and \(M=(\log L)^\chi\).
\end{proof}

\begin{remark}[No false support claim for the tail]
\label{rem:s7v59-tail-support}
The estimates in \Cref{prop:s7v59-tail} are operator-norm estimates.
Neither mixed term is asserted to have \(B_M\)-support on the side
occupied by \(K_{\rm tail}\), and the tail--tail term need not be
boundary-supported at all.  Its admissibility in a connected
thermodynamic expansion still requires support-resolved incidence, not
merely a small global norm.
\end{remark}

\section{A support-resolved resolvent interface}
\label{sec:s7v59-locality}

The occupation-band theorem concerns number space.  It does not by
itself control the spatial separation of two local vertices.  To state
the missing input without conflating the two geometries, write
\[
 K_{\rm pol}=\sum_XK_X
\]
over finite physical supports \(X\), and define
\[
 \Sigma_{XY}(z):=K_XR_M(z)K_Y^*.
\]
Let \(d(X,Y)\) be the block-lattice separation and, for a support-labelled
matrix \(A=(A_{XY})\), set
\begin{equation}
 \|A\|_{{\cal S}_\mu}
 :=
 \sup_X\sum_Y e^{\mu d(X,Y)}\|A_{XY}\|.
 \label{eq:s7v59-Schur}
\end{equation}

\begin{definition}[High-occupation locality card]
\label{def:s7v59-HOLC}
A support resolution satisfies \({\rm HOLC}(\mu)\) if there are
dimensionless matrices
\[
 \widetilde K:=aK,\qquad \widetilde R:=a^{-1}R_M(z)
\]
on a common support/configuration index set such that
\begin{align}
 \|\widetilde K\|_{{\cal S}_\mu}
 +\|\widetilde K^*\|_{{\cal S}_\mu}
 &\leq Ck_{{\rm pol},M},                               \label{eq:s7v59-HOLC-K}\\
 \|\widetilde R\|_{{\cal S}_\mu}
 &\leq {C\over\delta_M(z)}.                            \label{eq:s7v59-HOLC-R}
\end{align}
The resolution must reconstruct the actual product:
\(a\Sigma_{{\rm pol},{\rm pol}}
=\widetilde K\widetilde R\widetilde K^*\).
\end{definition}

\begin{lemma}[Weighted Schur algebra]
\label{lem:s7v59-Schur-algebra}
For kernels on a discrete metric index set,
\begin{equation}
 \|AB\|_{{\cal S}_\mu}
 \leq\|A\|_{{\cal S}_\mu}\|B\|_{{\cal S}_\mu}.
 \label{eq:s7v59-Schur-algebra}
\end{equation}
\end{lemma}

\begin{proof}
The triangle inequality for the metric gives
\(e^{\mu d(X,Z)}
\leq e^{\mu d(X,Y)}e^{\mu d(Y,Z)}\).  Hence
\begin{align*}
 \sum_Ze^{\mu d(X,Z)}\|(AB)_{XZ}\|
 &\leq
 \sum_Ye^{\mu d(X,Y)}\|A_{XY}\|
 \sum_Ze^{\mu d(Y,Z)}\|B_{YZ}\|\\
 &\leq
 \|B\|_{{\cal S}_\mu}
 \sum_Ye^{\mu d(X,Y)}\|A_{XY}\|.
\end{align*}
Take the supremum over \(X\).
\end{proof}

\begin{proposition}[What HOLC would yield]
\label{prop:s7v59-HOLC}
Under \({\rm HOLC}(\mu)\),
\begin{equation}
 \|a\Sigma_{{\rm pol},{\rm pol}}(z)\|_{{\cal S}_\mu}
 \leq {C^3k_{{\rm pol},M}^2\over\delta_M(z)}.
 \label{eq:s7v59-HOLC-Sigma}
\end{equation}
If the analogous bound holds for the two resolvent factors, then
\begin{equation}
 \|\Sigma_{{\rm pol},{\rm pol}}'(z)\|_{{\cal S}_\mu}
 \leq {C^4k_{{\rm pol},M}^2\over\delta_M(z)^2}.
 \label{eq:s7v59-HOLC-der}
\end{equation}
\end{proposition}

\begin{proof}
Use
\[
 a\Sigma_{{\rm pol},{\rm pol}}
 =\widetilde K\widetilde R\widetilde K^*
\]
and apply \Cref{lem:s7v59-Schur-algebra} twice.  For the derivative,
\[
 \Sigma_{{\rm pol},{\rm pol}}'
 =\widetilde K\widetilde R^2\widetilde K^*,
\]
and a third application gives the claim.
\end{proof}

\begin{firewall}[A gap is not a spatial propagation estimate]
\label{fire:s7v59-gap-locality}
The operator inequality
\[
 Q_M(H_a-z)Q_M\geq{\delta_M(z)\over a}Q_M
\]
proves only
\(\|\widetilde R\|\leq\delta_M(z)^{-1}\).  It does not prove
\eqref{eq:s7v59-HOLC-R}.  A weighted resolvent estimate requires control
of commutators with spatial weights, a Lieb--Robinson argument, or a
Combes--Thomas theorem on a specified configuration graph.  None follows
from the spectral floor alone.
\end{firewall}

\begin{remark}[Configuration distance versus support distance]
\label{rem:s7v59-config}
At occupation \(M\), a local hopping Hamiltonian can have configuration
graph row sum of order \(M/a\), while the high-occupation floor is also
of order \(M/a\).  This makes a cutoff-uniform Combes--Thomas exponent
plausible at sufficiently small weight, but it does not finish the
argument.  A high-occupation configuration may already carry quanta near
both \(X\) and \(Y\), so configuration distance between intermediate
states need not equal the physical separation \(d(X,Y)\).  The
support-labelled estimate must exploit connected incidence or a
conditional expectation which removes spectator excitations.
\end{remark}

\section{Corrected low-sector handoff}
\label{sec:s7v59-handoff}

\begin{theorem}[What V59 proves]
\label{thm:s7v59-result}
For the finite-degree cubic--quartic interaction, the occupation
Feshbach self-energy has exact support on the top four shells of
\(\operatorname{Ran}P_M\).  It vanishes on the vacuum and one-particle
sectors and obeys the relative-form estimate
\eqref{eq:s7v59-relative}, whose coefficient is
\(O(L^{-1}(\log L)^\chi)\) under the V55 logarithmic schedule.  The
nonpolynomial fifth-order chart remainder contributes only the
norm-small rates \eqref{eq:s7v59-tailrates}.  A support-resolved weighted
Schur bound follows algebraically from HOLC, but HOLC itself is not
deduced from the high-occupation spectral floor.
\end{theorem}

\begin{proof}
Combine
\Cref{thm:s7v59-shell,thm:s7v59-relative,
cor:s7v59-lograte,prop:s7v59-tail,prop:s7v59-HOLC}.
\end{proof}

\begin{assessment}[Progress at seventh-series V59]
\label{ass:s7v59-status}
The misleading route of treating the polynomial occupation return as a
bulk quadratic correction has been removed.  The return is instead a
positive operator subtracted at an artificial number boundary, with a
vanishing relative-form coefficient and a vanishing energy derivative.
This is stronger than a global norm estimate for low-energy vectors and
requires no spatial self-energy decomposition.  What remains is to show
that such a boundary form perturbation preserves the normalized
connected clustering argument, while support locality is required only
for the non-band-limited analytic tail and for genuine coarse
self-energies.
\end{assessment}

\begin{decision}[V60 audit and acquisition order]
\label{dec:s7v59-V60}
V60 is the compulsory five-version companion audit.  It will inspect the
current SM and NS notes for a theorem controlling small positive
relative-form subtractions in a gapped Duhamel or resolvent expansion,
and for any spectator-removal mechanism relevant to HOLC.  It will then
either import a proved interface or derive the needed form-stability
lemma directly.  The analytic tail will remain a distinct
support-incidence debit.
\end{decision}

\begin{firewall}[Claim boundary after seventh-series V59]
\label{fire:s7v59-claim}
V59 does not prove HOLC, support-resolved locality of the analytic-tail
return, stability of the normalized connected expansion under the
occupation-boundary form subtraction, density of the Weyl core in the
limiting interacting OS space, the pairing/Bogoliubov analogue,
normalized interacting chart exit, slow-label sensitivity, weighted
physical-source incidence, patchwise coarse-action cocycles, terminal
strong-coupling matching, cutoff removal, OS reconstruction, or the full
Yang--Mills mass gap.
\end{firewall}

\paragraph{Parent-version record.}
V59 was formed from
\texttt{Renewal\_Yang\_Mills\_Mass\_Gap\_Research\_Notes\_V58.tex}.
The V58 parent had \(24\,741\) logical source lines, \(883\,036\) bytes,
and SHA--256
\[
 \texttt{A515FFEBF73F899269597201B576AF07D5920FBEEC45548AAD8D791407FDF4FC}.
\]

% =============================================================================
% END APPENDED SEVENTH-SERIES V59 MATERIAL
% =============================================================================

% =============================================================================
% BEGIN APPENDED SEVENTH-SERIES V60 MATERIAL
% =============================================================================

\chapter{Companion audit and occupation-boundary shielding}
\label{ch:s7v60}

\section{Compulsory five-version audit}
\label{sec:s7v60-audit}

\subsection{Files inspected}
\label{sec:s7v60-files}

The active companion files inspected for this audit were
\begin{center}
\small\ttfamily
Renewal\_SM\_Einstein\_Continuum\_Research\_Notes\_V13.tex,\\
Renewal\_NS\_Stability\_Research\_Notes\_V26.tex.
\end{center}
At the audit time, the SM file had \(4\,611\) logical source lines,
\(163\,016\) bytes, and SHA--256
\[
 \texttt{C33B5654C9D7B50A0083F1F34ABD00B43F7AAC2F917A2CD15D13099C40860EF1}.
\]
The NS file had \(5\,319\) logical source lines, \(278\,283\) bytes, and
SHA--256
\[
 \texttt{82C887F8C2C69E4F28FAD7C3DC8DE5B9099CB5A4BF431EBA1FFF3E91935978AD}.
\]

\subsection{Standard Model--Einstein transfer}
\label{sec:s7v60-SM}

The current SM note contains two results relevant to the logical order of
the Yang--Mills programme.  First, a full-density \(L^q\) bound transfers
higher reference moments by Holder's inequality.  Second, a relative
entropy bound transfers a reference exponential moment.  Together with
negative-Sobolev compact embedding, these give a conditional route from a
reference cylinder measure to interacting Radon tightness.

Neither result controls a Hamiltonian form subtraction, a real-time or
Euclidean semigroup generated by that subtraction, or spatial clustering.
Moreover, the density estimate must cover the complete normalized
interacting density; control of an isolated action factor is explicitly
insufficient.  Thus no SM estimate can be substituted for the present
occupation-boundary calculation.

\begin{proposition}[Type-correct SM consequence]
\label{prop:s7v60-SM}
The SM density-loading theorem supplies a later completion interface for
the V58 Weyl/cylinder core: once the full Yang--Mills density relative to
a tight reference is uniformly bounded in \(L^q\), sufficiently high
reference Sobolev moments transfer to the interacting measures.  It does
not imply the form, gap, locality, or clustering estimates needed below.
\end{proposition}

\begin{proof}
The transferred quantity in the SM theorem is a static moment
\(\int X^p\,d\mu\), obtained from
\[
 \int X^p\,d\mu
 \leq
 \left\|{d\mu\over d\nu}\right\|_{L^q(\nu)}
 \left(\int X^{pq'}\,d\nu\right)^{1/q'}.
\]
No generator or spatially weighted kernel occurs in this inequality.
Consequently it is applicable to tightness after a full density bound is
proved, but it has no premise from which the present generator estimates
could follow.
\end{proof}

\subsection{Navier--Stokes transfer}
\label{sec:s7v60-NS}

The NS file remains V26.  Its new terminal material computes pointwise
rank-one norms of derivatives of the Oseen kernel.  The surrounding
programme concerns an ancient-profile Duhamel formula, explicit
parabolic constants, and a neutral time-translation mode.  It contains
no occupation projection, Feshbach map, positive Hamiltonian form, Gibbs
density, or connected Euclidean field expansion.

\begin{proposition}[NS non-transfer at V60]
\label{prop:s7v60-NS}
No numerical or analytic estimate in the current NS note loads the
occupation-boundary form or spatial-locality rows.
\end{proposition}

\begin{proof}
The NS kernels act on deterministic velocity profiles and their
derivatives.  The Yang--Mills rows concern a positive operator on
oscillator Fock space and a support-resolved Euclidean expectation.
There is no map in either programme identifying these spaces, operators,
or norms.  Transfer would therefore be type-incorrect.
\end{proof}

\begin{audit}[V60 companion decision]
\label{aud:s7v60}
The audit makes two decisions.
\begin{enumerate}[label=\textup{(\roman*)},leftmargin=3.8em]
\item Prove form-gap and resolvent stability directly, since neither
      companion supplies it.
\item Retain the SM full-density moment loader as the later interface
      from the dense reference Weyl core to interacting tightness; do not
      use it as a substitute for clustering.
\end{enumerate}
The NS direction does not alter the acquisition order.
\end{audit}

\section{Gap stability under a positive relative-form subtraction}
\label{sec:s7v60-form}

We first isolate the operator-theoretic statement required by V59.  It is
important that a Feshbach self-energy is positive and is subtracted from
the low Hamiltonian.

\begin{theorem}[Relative-subtraction gap theorem]
\label{thm:s7v60-gap}
Let \(A\geq0\) be self-adjoint on a Hilbert space \({\cal H}\), let
\(\Omega\) be a unit vector, and set
\[
 P_0:=|\Omega\rangle\langle\Omega|,\qquad Q_0:=I-P_0.
\]
Assume
\begin{equation}
 A\Omega=0,\qquad A\geq\gamma Q_0,\qquad \gamma>0.
 \label{eq:s7v60-Agap}
\end{equation}
Let \(B\) be a nonnegative quadratic form on
\({\cal D}(A^{1/2})\) satisfying
\begin{equation}
 0\leq B\leq\varepsilon A,\qquad 0\leq\varepsilon<1.
 \label{eq:s7v60-Bform}
\end{equation}
Then the form \(H:=A-B\) is closed on
\({\cal D}(A^{1/2})\), defines a nonnegative self-adjoint operator,
\[
 H\Omega=0,\qquad
 H\geq(1-\varepsilon)\gamma Q_0,                       \label{eq:s7v60-Hgap}
\]
and its kernel is exactly \(\mathbb C\Omega\).
\end{theorem}

\begin{proof}
The form norm of \(A+I\) makes
\({\cal D}(A^{1/2})\) complete.  From
\eqref{eq:s7v60-Bform},
\[
 (1-\varepsilon)\langle\psi,A\psi\rangle
 \leq
 \langle\psi,H\psi\rangle
 \leq
 \langle\psi,A\psi\rangle .
\]
Thus the form norm of \(H+I\) is equivalent to that of \(A+I\), so \(H\)
is closed and nonnegative.  The representation theorem gives its
self-adjoint operator.

Putting \(\psi=\Omega\) in \eqref{eq:s7v60-Bform} gives
\(\|B^{1/2}\Omega\|^2=0\).  Hence \(B^{1/2}\Omega=0\), so the mixed
form entries with \(\Omega\) vanish and \(H\Omega=0\).
For \(\psi\perp\Omega\),
\[
 \langle\psi,H\psi\rangle
 \geq(1-\varepsilon)\langle\psi,A\psi\rangle
 \geq(1-\varepsilon)\gamma\|\psi\|^2.
\]
This proves \eqref{eq:s7v60-Hgap} and excludes every additional kernel
vector.
\end{proof}

\begin{theorem}[Form-resolvent comparison]
\label{thm:s7v60-resolvent}
Under the hypotheses of \Cref{thm:s7v60-gap}, both \(A\) and \(H\)
reduce \(P_0\oplus Q_0\).  For every \(\lambda\geq0\), with inverses
understood on \(Q_0{\cal H}\) when \(\lambda=0\),
\begin{equation}
 \left\|
 Q_0\bigl[(H+\lambda)^{-1}-(A+\lambda)^{-1}\bigr]Q_0
 \right\|
 \leq
 {\varepsilon\over1-\varepsilon}\,
 {1\over\gamma+\lambda}.
 \label{eq:s7v60-resolvent}
\end{equation}
\end{theorem}

\begin{proof}
The preceding proof shows that \(B^{1/2}\Omega=0\); hence \(B\), in the
form sense, has no mixed entries between \(P_0\) and \(Q_0\).
Therefore both forms reduce the decomposition.

On \(Q_0{\cal H}\), set
\[
 C_\lambda:=(A+\lambda)^{-1/2}
 B(A+\lambda)^{-1/2}.
\]
The form inequality gives \(0\leq C_\lambda\leq\varepsilon I\), because
\(A\leq A+\lambda\).  Form factorization yields
\[
 H+\lambda
 =(A+\lambda)^{1/2}(I-C_\lambda)(A+\lambda)^{1/2},
\]
and therefore
\[
 (H+\lambda)^{-1}-(A+\lambda)^{-1}
 =
 (A+\lambda)^{-1/2}
 \bigl[(I-C_\lambda)^{-1}-I\bigr]
 (A+\lambda)^{-1/2}.
\]
The middle norm is at most
\(\varepsilon/(1-\varepsilon)\), while
\(\|(A+\lambda)^{-1/2}Q_0\|
\leq(\gamma+\lambda)^{-1/2}\).  Multiply the three bounds.
\end{proof}

\begin{corollary}[Temporal clustering survives]
\label{cor:s7v60-time}
Under the same hypotheses,
\begin{equation}
 \|e^{-tH}Q_0\|
 \leq e^{-(1-\varepsilon)\gamma t},
 \qquad t\geq0.                                       \label{eq:s7v60-semigroup}
\end{equation}
For bounded \(X,Y\),
\begin{align}
 &\left|
 \langle\Omega,Xe^{-tH}Y\Omega\rangle
 -\langle\Omega,X\Omega\rangle
  \langle\Omega,Y\Omega\rangle
 \right| \notag\\
 &\hspace{3em}\leq
 \|Q_0X^*\Omega\|\,\|Q_0Y\Omega\|\,
 e^{-(1-\varepsilon)\gamma t}.                        \label{eq:s7v60-timecorr}
\end{align}
\end{corollary}

\begin{proof}
The spectral theorem and \eqref{eq:s7v60-Hgap} give
\eqref{eq:s7v60-semigroup}.  Since \(e^{-tH}\Omega=\Omega\), subtracting
the vacuum product inserts \(Q_0\) on both sides of the semigroup.
Cauchy--Schwarz then proves \eqref{eq:s7v60-timecorr}.
\end{proof}

\begin{firewall}[Temporal gap is not spatial clustering]
\label{fire:s7v60-spatial}
\Cref{cor:s7v60-time} controls Euclidean time separation.  A form
inequality, even with a uniform spectral gap, does not give decay in
spatial separation.  Spatial clustering still requires either the
connected activity expansion or a weighted propagation theorem.
\end{firewall}

\section{Application to the occupation boundary}
\label{sec:s7v60-application}

Let \(A_M\) denote the baseline low Hamiltonian after the interactions
already controlled by the V57 connected expansion have been assembled,
but before the polynomial occupation return is subtracted.  Let its
vacuum energy be normalized to zero.

\begin{definition}[Baseline coercivity card]
\label{def:s7v60-BCC}
The baseline coercivity card, BCC, consists of constants
\(c_0,\gamma_0>0\), independent of the cutoff, such that
\begin{equation}
 A_M\geq c_0H_{0,a},\qquad
 A_M\geq\gamma_0 Q_0,\qquad A_M\Omega=0.
 \label{eq:s7v60-BCC}
\end{equation}
\end{definition}

\begin{proposition}[Conditional occupation-gap insertion]
\label{prop:s7v60-insertion}
Assume BCC and let
\[
 B_M^{\rm F}:=\Sigma_{{\rm pol},{\rm pol}}(0).
\]
Then
\begin{equation}
 0\leq B_M^{\rm F}
 \leq{\varepsilon_{{\rm occ},M}(0)\over c_0}\,A_M.
 \label{eq:s7v60-BF}
\end{equation}
If \(\varepsilon_{{\rm occ},M}(0)<c_0\), the Feshbach-corrected low
Hamiltonian
\[
 H_M^{\rm eff}:=A_M-B_M^{\rm F}
\]
has the same vacuum and gap at least
\begin{equation}
 \gamma_M^{\rm eff}
 \geq
 \left(1-{\varepsilon_{{\rm occ},M}(0)\over c_0}\right)\gamma_0.
 \label{eq:s7v60-effgap}
\end{equation}
Under the V55 schedule, the relative loss tends to zero.
\end{proposition}

\begin{proof}
V59 gives
\(B_M^{\rm F}\leq\varepsilon_{{\rm occ},M}(0)H_{0,a}\).
BCC gives \(H_{0,a}\leq c_0^{-1}A_M\) as forms.  This proves
\eqref{eq:s7v60-BF}.  Apply \Cref{thm:s7v60-gap} with
\(\varepsilon=\varepsilon_{{\rm occ},M}(0)/c_0\).  The logarithmic rate
is \Cref{cor:s7v59-lograte}.
\end{proof}

\begin{firewall}[BCC remains a physical loader]
\label{fire:s7v60-BCC}
The V42 high-occupation floor is a lower bound on \(Q_MH_aQ_M\); it is
not BCC for the interacting low block \(P_MH_aP_M\).  The local connected
expansion and vacuum normalization must still prove a cutoff-uniform
low-sector gap or coercive comparison.  Thus
\Cref{prop:s7v60-insertion} closes the functional-analytic insertion
step, not the physical baseline gap.
\end{firewall}

\section{Exact number-distance shielding}
\label{sec:s7v60-shielding}

The shell support proved in V59 has a second consequence: a
finite-bandwidth word starting at low number cannot see the occupation
boundary until its length is proportional to \(M\).

\begin{lemma}[Number growth along a word]
\label{lem:s7v60-word}
Let \(V_1,\ldots,V_n\) have number bandwidth at most \(r\), and let
\(\psi=P_{\leq s}\psi\).  Then
\begin{equation}
 V_n\cdots V_1\psi
 =
 P_{\leq s+nr}V_n\cdots V_1\psi.
 \label{eq:s7v60-word}
\end{equation}
The same conclusion holds when arbitrary number-preserving operators are
inserted between the \(V_j\).
\end{lemma}

\begin{proof}
For one factor, the bandwidth definition gives
\(V_1P_{\leq s}=P_{\leq s+r}V_1P_{\leq s}\).  Iterate this identity.
A number-preserving operator leaves every spectral subspace of
\({\cal N}\) invariant and therefore does not change the induction.
\end{proof}

\begin{corollary}[Boundary invisibility below a linear order]
\label{cor:s7v60-invisible}
Let \(B_{M,r_0}=\mathbf1_{(M-r_0,M]}({\cal N})\).  Under the hypotheses
of \Cref{lem:s7v60-word},
\begin{equation}
 B_{M,r_0}V_n\cdots V_1P_{\leq s}=0
 \quad\hbox{if}\quad s+nr\leq M-r_0.                   \label{eq:s7v60-invisible}
\end{equation}
Thus the first possibly nonzero order is
\begin{equation}
 n_*(M;s,r,r_0)
 :=
 \left\lceil{M-r_0+1-s\over r}\right\rceil.
 \label{eq:s7v60-nstar}
\end{equation}
\end{corollary}

\begin{proof}
The range in \eqref{eq:s7v60-word} is orthogonal to
\(\operatorname{Ran}B_{M,r_0}\) under the displayed inequality.
\end{proof}

\begin{theorem}[Two-sided shielding of a boundary return]
\label{thm:s7v60-two-sided}
Let \(S=B_{M,r_0}SB_{M,r_0}\).  Let \(\psi_L,\psi_R\) have number
supports at most \(s_L,s_R\), and let
\({\cal W}_L,{\cal W}_R\) be words containing \(n_L,n_R\) vertices of
bandwidth at most \(r\), with arbitrary number-preserving factors between
them.  Then
\begin{equation}
 \langle\psi_L,{\cal W}_L^*S{\cal W}_R\psi_R\rangle=0
 \label{eq:s7v60-two-sided}
\end{equation}
unless simultaneously
\begin{equation}
 n_L\geq n_*(M;s_L,r,r_0),\qquad
 n_R\geq n_*(M;s_R,r,r_0).
 \label{eq:s7v60-two-thresholds}
\end{equation}
\end{theorem}

\begin{proof}
Insert \(B_{M,r_0}\) on each side of \(S\).  If the right inequality in
\eqref{eq:s7v60-two-thresholds} fails,
\Cref{cor:s7v60-invisible} annihilates the ket.  If the left inequality
fails, apply the same argument to the adjoint word and the bra.
\end{proof}

\begin{corollary}[Cubic--quartic return is perturbatively remote]
\label{cor:s7v60-remote}
For the V59 polynomial return, take \(r=r_0=4\).  From fixed-degree
endpoint vectors supported in \(P_{\leq s}\), every word containing that
return vanishes unless at least
\[
 2\left\lceil{M-3-s\over4}\right\rceil
\]
cubic--quartic vertices occur in total on its two sides.
\end{corollary}

\begin{proof}
Apply \Cref{thm:s7v60-two-sided} with \(s_L=s_R=s\).
\end{proof}

\section{Fixed Weyl endpoints have superexponentially small shell tails}
\label{sec:s7v60-Weyl}

The V58 endpoint core uses Weyl vectors rather than finite-number
vectors.  Their occupation-boundary component is nevertheless explicit.

\begin{lemma}[Poisson tail of a coherent vector]
\label{lem:s7v60-Poisson}
Let \(\Omega\) be the bosonic vacuum, let \(W(f)\) be a Weyl operator,
and put \(\lambda=\|f\|^2\).  For every integer \(N>\lambda\),
\begin{equation}
 \|\mathbf1_{[N,\infty)}({\cal N})W(f)\Omega\|^2
 =
 e^{-\lambda}\sum_{n=N}^{\infty}{\lambda^n\over n!}
 \leq
 \left({e\lambda\over N}\right)^N.
 \label{eq:s7v60-Poisson}
\end{equation}
\end{lemma}

\begin{proof}
The coherent-vector expansion gives the equality.  For \(\theta>0\),
Markov's inequality and the Poisson moment generating function give
\[
 {\mathbb P}(X\geq N)
 \leq
 \exp\{\lambda(e^\theta-1)-\theta N\}.
\]
Choose \(e^\theta=N/\lambda\).  The resulting bound is
\[
 \exp\{-N\log(N/\lambda)+N-\lambda\}
 \leq(e\lambda/N)^N.
\]
\end{proof}

\begin{corollary}[Boundary matrix element between Weyl vectors]
\label{cor:s7v60-Weyl-boundary}
Let \(S=B_{M,r_0}SB_{M,r_0}\), set \(N=M-r_0+1\), and suppose
\(N>\max\{\|f\|^2,\|h\|^2\}\).  Then
\begin{align}
 |\langle W(f)\Omega,SW(h)\Omega\rangle|
 &\leq
 \|S\|
 \left({e\|f\|^2\over N}\right)^{N/2}
 \left({e\|h\|^2\over N}\right)^{N/2}.                \label{eq:s7v60-Weyl-boundary}
\end{align}
For fixed \(f,h,r_0\), the factor multiplying \(\|S\|\) decays faster
than \(e^{-cM}\) for every fixed \(c>0\).
\end{corollary}

\begin{proof}
Insert \(B_{M,r_0}\) on both sides of \(S\), apply Cauchy--Schwarz and
\Cref{lem:s7v60-Poisson}.  The logarithm of either tail factor is
\[
 -{N\over2}\log N+O(N),
\]
which is eventually smaller than \(-cM\) for any fixed \(c\).
\end{proof}

\begin{remark}[Gauge averaging preserves the bound]
\label{rem:s7v60-gauge-average}
If the compact gauge action commutes with \({\cal N}\) and acts unitarily
on one-particle labels, then \(\|gf\|=\|f\|\) and
\(B_{M,r_0}\) commutes with gauge averaging.  The triangle inequality
therefore gives the same Poisson bound for every gauge average of a fixed
Weyl vector.
\end{remark}

\begin{corollary}[Direct fixed-endpoint estimate for the Feshbach return]
\label{cor:s7v60-Feshbach-Weyl}
With \(S=a\Sigma_{{\rm pol},{\rm pol}}(z)\), \(r_0=4\), and
\(\delta_M(z)>0\),
\begin{align}
 &|\langle W(f)\Omega,
 a\Sigma_{{\rm pol},{\rm pol}}(z)W(h)\Omega\rangle|
 \notag\\
 &\quad\leq
 {k_{{\rm pol},M}^2\over\delta_M(z)}
 \left({e\|f\|^2\over M-3}\right)^{(M-3)/2}
 \left({e\|h\|^2\over M-3}\right)^{(M-3)/2}.           \label{eq:s7v60-direct}
\end{align}
Under the logarithmic schedule, this tends to zero
superexponentially in \(M\) up to a polynomial prefactor.
\end{corollary}

\begin{proof}
Use the V42 norm estimate with the V59 polynomial block size and apply
\Cref{cor:s7v60-Weyl-boundary}.
\end{proof}

\section{A normalized-series shielding interface}
\label{sec:s7v60-series}

The preceding exact selection rule can replace a spatial locality bound
only if the high-order tail of the normalized Duhamel series is summable.
The next lemma records the precise required majorant.

\begin{definition}[Two-sided Duhamel majorant]
\label{def:s7v60-TDM}
For fixed endpoint number degrees \(s_L,s_R\), TDM with parameter
\(0<q_M<1\) means that the sum of absolute normalized connected weights
with a distinguished bounded insertion \(S\), \(n_L\) finite-bandwidth
vertices to its left, and \(n_R\) to its right is bounded by
\begin{equation}
 C_{\rm end}\|S\|_{\rm dim}
 q_M^{n_L+n_R},
 \label{eq:s7v60-TDM}
\end{equation}
after summing times, spatial labels, trees, and all other declared
indices.  Here \(\|\cdot\|_{\rm dim}\) is the dimensionless insertion
norm.
\end{definition}

\begin{proposition}[Shielded tail under TDM]
\label{prop:s7v60-TDM}
Let \(S=B_{M,r_0}SB_{M,r_0}\), and assume TDM for vertices of number
bandwidth at most \(r\).  Then the entire connected contribution
containing one distinguished \(S\) is bounded by
\begin{equation}
 C_{\rm end}\|S\|_{\rm dim}
 {q_M^{\,n_*(M;s_L,r,r_0)+n_*(M;s_R,r,r_0)}
  \over(1-q_M)^2}.
 \label{eq:s7v60-TDM-tail}
\end{equation}
\end{proposition}

\begin{proof}
\Cref{thm:s7v60-two-sided} makes every term outside the two threshold
ranges vanish.  Sum the two geometric series:
\[
 \sum_{n_L\geq n_{*,L}}\sum_{n_R\geq n_{*,R}}
 q_M^{n_L+n_R}
 =
 {q_M^{n_{*,L}+n_{*,R}}\over(1-q_M)^2}.
\]
\end{proof}

\begin{firewall}[TDM has not yet been extracted from the tree theorem]
\label{fire:s7v60-TDM}
V57 proves a normalized degree-weighted tree majorant for the reference
connected series.  TDM additionally singles out a possibly nonlocal
boundary insertion and separates the remaining vertices into its two
time sides.  This refinement is plausible but is not a formal consequence
of the stated V57 theorem.  It must be proved by reopening the ordered
Duhamel/tree interpolation and retaining the distinguished insertion.
\end{firewall}

\section{Status and next acquisition}
\label{sec:s7v60-status}

\begin{theorem}[What V60 proves]
\label{thm:s7v60-result}
A positive subtraction \(B\leq\varepsilon A\), \(\varepsilon<1\),
preserves the vacuum and at least a fraction \(1-\varepsilon\) of the
baseline gap, with the explicit resolvent comparison
\eqref{eq:s7v60-resolvent}.  Applied under BCC, the V59 occupation return
therefore has a vanishing gap cost.  Independently, exact number
bandwidth forces every finite-degree word seeing that return to have
length proportional to \(M\) on each side.  Fixed Weyl endpoints see the
same boundary only through the superexponentially small Poisson tail
\eqref{eq:s7v60-direct}.  Under the explicitly declared TDM, the whole
connected boundary contribution obeys
\eqref{eq:s7v60-TDM-tail}.
\end{theorem}

\begin{proof}
Combine
\Cref{thm:s7v60-gap,thm:s7v60-resolvent,
prop:s7v60-insertion,thm:s7v60-two-sided,
cor:s7v60-Weyl-boundary,prop:s7v60-TDM}.
\end{proof}

\begin{assessment}[Progress at seventh-series V60]
\label{ass:s7v60-status}
The functional-analytic treatment of the polynomial occupation return is
now closed conditional only on a baseline low-sector coercivity card.
The more serious spatial-locality demand can be bypassed on fixed
finite-degree endpoints if the V57 normalized tree proof is refined to
TDM: the boundary is perturbatively remote by an order proportional to
\(M\).  Fixed Weyl endpoints already have an independent
superexponential shell estimate.  This is a noncircular alternative to
trying to derive physical-space locality from the high-occupation gap.
\end{assessment}

\begin{decision}[V61 acquisition order]
\label{dec:s7v60-V61}
V61 will reopen the V56--V57 ordered connected expansion with one
distinguished bounded insertion.  It will track the number of vertices on
the two time sides and test whether the Prüfer degree sum and stable
tree-graph bound prove TDM with the same logarithmic activity.  If a
distinguished vertex incurs an uncontrolled spatial sum, the proof will
move upstream to an anchored, rather than global, boundary insertion
norm.  The analytic tail and BCC remain separate loaders.
\end{decision}

\begin{firewall}[Claim boundary after seventh-series V60]
\label{fire:s7v60-claim}
V60 does not prove BCC, TDM, HOLC, support-resolved analytic-tail
locality, interacting full-density \(L^q\) control, tightness of the
interacting field measures, density of the Weyl core in the limiting
interacting OS space, the pairing/Bogoliubov analogue, normalized
interacting chart exit, slow-label sensitivity, weighted physical-source
incidence, patchwise coarse-action cocycles, terminal strong-coupling
matching, cutoff removal, OS reconstruction, or the full Yang--Mills
mass gap.
\end{firewall}

\paragraph{Parent-version record.}
V60 was formed from
\texttt{Renewal\_Yang\_Mills\_Mass\_Gap\_Research\_Notes\_V59.tex}.
The V59 parent had \(25\,326\) logical source lines, \(903\,648\) bytes,
and SHA--256
\[
 \texttt{7845079A5A5D987B4777EDEE7CEC541B066DF861C4A2E29E1FB03CD006117B96}.
\]
The next compulsory companion audit is V65.

% =============================================================================
% END APPENDED SEVENTH-SERIES V60 MATERIAL
% =============================================================================

% =============================================================================
% BEGIN APPENDED SEVENTH-SERIES V61 MATERIAL
% =============================================================================

\chapter{Anchored boundary insertions and the locality no-go theorem}
\label{ch:s7v61}

\section{Purpose}
\label{sec:s7v61-purpose}

V60 formulated a two-sided Duhamel majorant for a distinguished
occupation-boundary insertion.  Reopening the V56--V57 proof shows that
the desired refinement is available for an insertion supported on a
fixed number of oscillator blocks.  The insertion costs one
Weyl--Plancherel polynomial and one harmless linear factor in the
perturbation order.  Exact number-distance shielding then starts the
series at order proportional to \(M\) on each side.

The actual Feshbach self-energy has not yet been proved to possess such an
anchored decomposition.  This is not a technicality: positivity,
top-shell support, and a small operator norm do not imply any spatial
decay.  The present note gives an explicit positive rank-one
counterexample.  It therefore replaces the overbroad TDM card by an
anchored theorem and isolates the exact additional locality datum needed
for the physical self-energy.

\section{Ordered cumulants retain number shielding}
\label{sec:s7v61-cumulant}

Let \(H_0=d\Gamma(A)\) be the number-conserving quadratic reference.  For
operators \(O_1,\ldots,O_m\) at ordered Euclidean times, write
\[
 {\cal M}_0(I)
\]
for the vacuum moment of the subword indexed by \(I\), in inherited time
order, and define the ordered cumulant by
\begin{equation}
 \kappa_0(O_1,\ldots,O_m)
 =
 \sum_{\pi\in\Pi_m}
 (|\pi|-1)!(-1)^{|\pi|-1}
 \prod_{C\in\pi}{\cal M}_0(C).
 \label{eq:s7v61-cumulant}
\end{equation}

\begin{lemma}[Moment block containing a boundary insertion]
\label{lem:s7v61-moment-zero}
Let \(S=B_{M,r_0}SB_{M,r_0}\).  In an ordered vacuum moment containing
\(S\), suppose all other internal vertices have number bandwidth at most
\(r\).  Let \(n_L,n_R\) be the numbers of such vertices occurring to the
left and right of \(S\) in that moment.  If the outer endpoint on the
left creates at most \(s_L\) particles from the vacuum and the outer
endpoint on the right creates at most \(s_R\), with \(s=0\) when the
corresponding endpoint is absent from the moment, then the moment
vanishes unless
\begin{align}
 n_L&\geq n_*(M;s_L,r,r_0),                            \label{eq:s7v61-left}\\
 n_R&\geq n_*(M;s_R,r,r_0),                            \label{eq:s7v61-right}
\end{align}
where \(n_*\) is defined in \eqref{eq:s7v60-nstar}.
\end{lemma}

\begin{proof}
Every reference semigroup between insertions commutes with
\({\cal N}\).  Starting from the right vacuum, the right endpoint and
the \(n_R\) finite-bandwidth vertices have number support at most
\(s_R+n_Rr\).  The right boundary projection in \(S\) therefore
annihilates this vector unless \eqref{eq:s7v61-right} holds, by
\Cref{cor:s7v60-invisible}.  Apply the same argument to the adjoint
left subword for \eqref{eq:s7v61-left}.
\end{proof}

\begin{theorem}[Cumulant shielding]
\label{thm:s7v61-cumulant-zero}
Consider an ordered cumulant with one distinguished insertion
\(S=B_{M,r_0}SB_{M,r_0}\), fixed finite-number endpoints of degrees at
most \(s_L,s_R\), and \(n_L,n_R\) finite-bandwidth interaction vertices
on the two time sides of \(S\).  The cumulant is zero unless
\begin{equation}
 n_L\geq n_*(M;s_L,r,r_0),\qquad
 n_R\geq n_*(M;s_R,r,r_0).
 \label{eq:s7v61-cumulant-threshold}
\end{equation}
\end{theorem}

\begin{proof}
Expand the cumulant by \eqref{eq:s7v61-cumulant}.  In every partition,
there is a unique block \(C_S\) containing \(S\).  If an outer endpoint
is absent from \(C_S\), the corresponding source degree for that block
is zero; if it is present, it is at most \(s_L\) or \(s_R\).
The numbers of internal vertices on the two sides within \(C_S\) are at
most the global \(n_L,n_R\).  If either inequality in
\eqref{eq:s7v61-cumulant-threshold} fails,
\Cref{lem:s7v61-moment-zero} makes
\({\cal M}_0(C_S)=0\).  Every partition summand therefore vanishes.
\end{proof}

\begin{remark}[Why the cancellation must precede absolute Weyl inversion]
\label{rem:s7v61-order}
An individual Weyl operator has infinite number bandwidth.  If every
finite-degree vertex is first expanded into Weyl operators and absolute
values are then taken, the exact zero in
\Cref{thm:s7v61-cumulant-zero} is lost.  The correct order is:
retain the polynomial operators, apply the number-selection rule to each
ordered coefficient, discard the zero coefficients, and only then use
Weyl inversion to bound the surviving high-order tail.
\end{remark}

\section{One distinguished local vertex in the tree sum}
\label{sec:s7v61-distinguished}

Fix two endpoint observables in the V58 cutoff-uniform Weyl core.  Let
\(V_X\) be local interaction vertices, supported on at most \(r\)
oscillator modes, with
\[
 \|V_X\|\leq\rho_M.
\]
Let \(S_Z=P_MS_ZP_M\) be one distinguished operator supported on a fixed
set \(Z\) containing at most \(r_S\) modes.  Its spacetime position is
integrated, but the physical anchor \(Z\) is not summed.

Denote by
\begin{equation}
 {\cal C}_{r,M}
 :=
 C_r(M+1)^{2r+3/2}                                    \label{eq:s7v61-CrM}
\end{equation}
the V57 per-interaction Weyl loss, including the fixed time and spatial
row constants, and set
\begin{equation}
 q_M:={\cal C}_{r,M}\rho_M.                            \label{eq:s7v61-qM}
\end{equation}
For the distinguished vertex put
\begin{equation}
 {\cal D}_{r_S,M}
 :=
 C_{r_S}(M+1)^{2r_S+3/2}.                              \label{eq:s7v61-DrM}
\end{equation}

\begin{lemma}[Distinguished Prüfer normalization]
\label{lem:s7v61-Prufer}
Let a connected cumulant contain two fixed endpoints, \(n\) ordinary
interaction vertices, and one distinguished local vertex.  After the
degree-weighted Prüfer sum, the ratio between the tree factorial and the
\(1/n!\) expansion coefficient is \(n+1\).
\end{lemma}

\begin{proof}
There are \(m=n+3\) total vertices.  The degree-weighted tree theorem
\eqref{eq:s7v57-Prufer} supplies \((m-2)!=(n+1)!\).
The expansion of a normalized expectation to first order in the
distinguished coupling and order \(n\) in the ordinary interaction has
coefficient \(1/n!\).  Equivalently, among the \(n+1\) internal time
labels, one is distinguished.  Their ratio is
\((n+1)!/n!=n+1\).
\end{proof}

\begin{theorem}[Anchored distinguished-insertion tree bound]
\label{thm:s7v61-anchored}
Assume \(q_M<1\).  At order \(n\) in the ordinary interaction, the
absolute normalized connected contribution containing one fixed
spatially anchored \(S_Z\) is bounded by
\begin{equation}
 C_{\rm end}(n+1)
 {\cal D}_{r_S,M}\|S_Z\|_{\rm dim}\,q_M^n.
 \label{eq:s7v61-anchored}
\end{equation}
Here \(\|S_Z\|_{\rm dim}\) is the dimensionless operator norm, and
\(C_{\rm end}\) depends on the two fixed endpoint Weyl labels and the
chosen lower spacetime decay rate, but not on \(M\), the volume, or
\(n\).
\end{theorem}

\begin{proof}
Use Weyl inversion for every surviving coefficient.  At each ordinary
vertex, the V57 absolute characteristic moment and spatial/time edge
summation cost at most
\({\cal C}_{r,M}\rho_M=q_M\).  At the distinguished vertex the same
argument with \(r_S\) modes costs
\({\cal D}_{r_S,M}\|S_Z\|_{\rm dim}\).  The V58 fixed Weyl endpoints
have no occupation-window loss.  Stability cancels the Gaussian
one-vertex factors exactly as in V56.  The degree-weighted Prüfer sum is
normalized as in \Cref{lem:s7v61-Prufer}.  All internal physical
positions are summed along the tree-edge row bound; \(Z\) remains fixed,
so no volume factor occurs.  Multiplication of these factors gives
\eqref{eq:s7v61-anchored}.
\end{proof}

\begin{corollary}[Absolute sum with one anchored insertion]
\label{cor:s7v61-sum}
Under the hypotheses of \Cref{thm:s7v61-anchored},
\begin{equation}
 \sum_{n=0}^\infty
 C_{\rm end}(n+1)
 {\cal D}_{r_S,M}\|S_Z\|_{\rm dim}q_M^n
 =
 {C_{\rm end}{\cal D}_{r_S,M}\|S_Z\|_{\rm dim}
  \over(1-q_M)^2}.
 \label{eq:s7v61-sum}
\end{equation}
\end{corollary}

\begin{proof}
Use \(\sum_{n\geq0}(n+1)q^n=(1-q)^{-2}\).
\end{proof}

\section{Anchored number-shielded tail}
\label{sec:s7v61-tail}

For fixed finite-degree endpoints it is useful to retain the two
perturbation orders \(n_L,n_R\).  The estimate
\eqref{eq:s7v61-anchored} may be distributed among them with
\(n=n_L+n_R\).  Summing over the location of the distinguished Euclidean
time costs at most the already present factor \(n+1\).

\begin{theorem}[Local TDM is proved]
\label{thm:s7v61-local-TDM}
Assume the distinguished local insertion also satisfies
\[
 S_Z=B_{M,r_0}S_ZB_{M,r_0}.
\]
Let the two finite-number endpoints have degrees at most \(s_L,s_R\),
and put
\[
 a_M:=n_*(M;s_L,r,r_0),\qquad
 b_M:=n_*(M;s_R,r,r_0).
\]
Then the absolute connected sum containing one \(S_Z\) is bounded by
\begin{align}
 &C_{\rm end}{\cal D}_{r_S,M}\|S_Z\|_{\rm dim}
 q_M^{a_M+b_M}
 \left[
 {a_M+b_M+1\over(1-q_M)^2}
 +{2q_M\over(1-q_M)^3}
 \right].
 \label{eq:s7v61-shielded}
\end{align}
\end{theorem}

\begin{proof}
\Cref{thm:s7v61-cumulant-zero} restricts the sum to
\(n_L\geq a_M\), \(n_R\geq b_M\).
\Cref{thm:s7v61-anchored} bounds a surviving pair of orders by
\[
 C_{\rm end}{\cal D}_{r_S,M}\|S_Z\|_{\rm dim}
 (n_L+n_R+1)q_M^{n_L+n_R}.
\]
Writing \(n_L=a_M+i\), \(n_R=b_M+j\), the remaining scalar sum is
\begin{align*}
 q_M^{a_M+b_M}
 \sum_{i,j\geq0}(a_M+b_M+i+j+1)q_M^{i+j}
 &=
 q_M^{a_M+b_M}
 \left[
 {a_M+b_M+1\over(1-q_M)^2}
 +{2q_M\over(1-q_M)^3}
 \right].
\end{align*}
\end{proof}

\begin{corollary}[Logarithmic-window disappearance]
\label{cor:s7v61-log}
Suppose \(r,r_0,r_S,s_L,s_R\) are fixed,
\(M=(\log L)^\chi\), and, for some fixed
\(\alpha>0\) and \(\beta\in\mathbb R\),
\[
 q_M\leq L^{-\alpha}(\log L)^\beta
\]
for all sufficiently large \(L\).  If
\({\cal D}_{r_S,M}\|S_Z\|_{\rm dim}\) grows at most as a fixed power of
\(M\), then the right side of \eqref{eq:s7v61-shielded} tends to zero
faster than \(L^{-p}\) for every fixed \(p>0\).
\end{corollary}

\begin{proof}
For all sufficiently large \(L\), \(q_M\leq1/2\), while
\(a_M+b_M\geq cM-C\) for fixed \(c>0,C\).  Hence
\[
 \log q_M^{a_M+b_M}
 \leq
 -c\alpha(\log L)^{\chi+1}
 +O((\log L)^\chi\log\log L).
\]
This dominates \(-p\log L\) for every fixed \(p\) when \(\chi>0\).
All remaining factors in \eqref{eq:s7v61-shielded} are polynomial in
\(M\).
\end{proof}

\begin{remark}[Fixed Weyl endpoints]
\label{rem:s7v61-Weyl}
The exact threshold theorem uses finite-number endpoints.  On the V58
Weyl core, split each coherent endpoint at a number level \(s_M<M\).
The low part obeys \eqref{eq:s7v61-shielded}; the high part is controlled
by the Poisson tail \eqref{eq:s7v60-Poisson}.  Taking, for example,
\(s_M=\lfloor M/2\rfloor\) leaves a threshold proportional to \(M\) and
a superexponentially small endpoint tail.  This extension requires only
the contraction of normalized expectations by endpoint vector norms.
\end{remark}

\section{Why shell support cannot replace spatial locality}
\label{sec:s7v61-nogo}

\begin{theorem}[Positive top-shell operators need not be spatially local]
\label{thm:s7v61-nogo}
Let \(x,y\) be two oscillator blocks at arbitrary separation.  On the
two-block Fock space, let
\[
 u_x:=|M\rangle_x\otimes|0\rangle_y,\qquad
 u_y:=|0\rangle_x\otimes|M\rangle_y,
\qquad
 \phi_{xy}:={u_x+u_y\over\sqrt2}.
\]
With total number cutoff \(P_M\), define
\begin{equation}
 S_{xy}:=|\phi_{xy}\rangle\langle\phi_{xy}|.
 \label{eq:s7v61-counter-S}
\end{equation}
Then
\begin{equation}
 S_{xy}\geq0,\qquad \|S_{xy}\|=1,\qquad
 S_{xy}=B_{M,r_0}S_{xy}B_{M,r_0}                     \label{eq:s7v61-counter-properties}
\end{equation}
for every fixed \(r_0\geq1\), but
\begin{equation}
 \langle u_x,S_{xy}u_y\rangle={1\over2}               \label{eq:s7v61-counter-matrix}
\end{equation}
independently of \(d(x,y)\).
\end{theorem}

\begin{proof}
Both \(u_x\) and \(u_y\) have total number exactly \(M\), so
\(\phi_{xy}\) lies in every upper band \(B_{M,r_0}\).
Equation \eqref{eq:s7v61-counter-S} is a rank-one orthogonal projection,
which proves positivity and norm one.  Finally
\[
 \langle u_x,S_{xy}u_y\rangle
 =
 \langle u_x,\phi_{xy}\rangle
 \langle\phi_{xy},u_y\rangle
 ={1\over2}.
\]
\end{proof}

\begin{corollary}[No spatial decay from the V59 data alone]
\label{cor:s7v61-nogo}
There is no function \(F(R)\to0\) such that every positive operator
\(S=B_{M,r_0}SB_{M,r_0}\) satisfies
\[
 |\langle u_x,Su_y\rangle|
 \leq F(d(x,y))\|S\|
\]
uniformly in the number of blocks.
\end{corollary}

\begin{proof}
Apply the proposed inequality to \(S_{xy}\) and use
\eqref{eq:s7v61-counter-matrix}; then let \(d(x,y)\to\infty\).
\end{proof}

\begin{firewall}[Interpretation of the no-go theorem]
\label{fire:s7v61-nogo}
The counterexample is not asserted to equal the Yang--Mills Feshbach
self-energy.  It proves that positivity, a small norm, and exact
occupation-shell support are logically insufficient to establish
spatial locality.  Additional information from the local Hamiltonian and
its high-sector resolvent is indispensable.
\end{firewall}

\section{The exact anchored locality card}
\label{sec:s7v61-card}

\begin{definition}[Anchored boundary decomposition card, ABDC]
\label{def:s7v61-ABDC}
ABDC at rate \(\mu>0\) is a decomposition
\begin{equation}
 a\Sigma_{{\rm pol},{\rm pol}}(z)
 =
 \sum_Z S_Z(z),\qquad
 S_Z=B_{M,4}S_ZB_{M,4},                               \label{eq:s7v61-ABDC-decomp}
\end{equation}
where every \(S_Z\) is supported on at most a fixed number \(r_S\) of
oscillator modes and
\begin{equation}
 \sup_x\sum_{Z\ni x}
 e^{\mu\,{\rm diam}(Z)}
 {\cal D}_{r_S,M}\|S_Z(z)\|
 \leq \sigma_M(z).                                    \label{eq:s7v61-ABDC}
\end{equation}
The target is \(\sigma_M(z)=o(1)\), uniformly in the declared low-energy
spectral interval and slow labels.
\end{definition}

\begin{proposition}[ABDC closes the boundary-activity row]
\label{prop:s7v61-ABDC}
Assume ABDC and \(q_M<1\).  Then the normalized connected contribution
with one polynomial occupation return is spatially summable.  For
finite-degree endpoints it additionally carries the shielding factor in
\eqref{eq:s7v61-shielded}; for fixed Weyl endpoints the low/high split of
\Cref{rem:s7v61-Weyl} gives the same vanishing conclusion.
\end{proposition}

\begin{proof}
Apply \Cref{thm:s7v61-local-TDM} to each \(S_Z\).
The anchored polymer/tree summation charges each support containing a
fixed root \(x\) by the left side of
\eqref{eq:s7v61-ABDC}; the exponential diameter weight pays the standard
loss when the tree is reanchored through \(Z\).  Taking the supremum over
the root gives \(\sigma_M(z)\) with no volume factor.  The number
shielding factors are uniform in \(Z\) and factor out of the support sum.
\end{proof}

\begin{firewall}[ABDC is stronger than HOLC notation]
\label{fire:s7v61-ABDC}
The product identity
\(\Sigma_{XY}=K_XR_MK_Y^*\) does not by itself give local operators
\(S_Z\) satisfying \eqref{eq:s7v61-ABDC}.  The high-sector resolvent can
depend on spectator excitations outside \(X\cup Y\).  One must either
prove a quasi-local resolvent decomposition by conditional expectations
or run the connected expansion before taking the Feshbach map.
\end{firewall}

\section{Status and next acquisition}
\label{sec:s7v61-status}

\begin{theorem}[What V61 proves]
\label{thm:s7v61-result}
Exact number-distance cancellation survives moment--cumulant inversion
provided it is applied before absolute Weyl inversion.  A distinguished
operator on fixed local support satisfies the anchored tree bound
\eqref{eq:s7v61-anchored}; if it is top-shell supported, the complete
connected contribution obeys the shielded estimate
\eqref{eq:s7v61-shielded}.  Conversely, the positive rank-one example
\eqref{eq:s7v61-counter-S} proves that top-shell support, positivity, and
operator norm control alone cannot imply spatial locality.  ABDC is a
sufficient, explicitly weighted interface for the actual self-energy.
\end{theorem}

\begin{proof}
Use
\Cref{thm:s7v61-cumulant-zero,thm:s7v61-anchored,
thm:s7v61-local-TDM,thm:s7v61-nogo,prop:s7v61-ABDC}.
\end{proof}

\begin{assessment}[Progress at seventh-series V61]
\label{ass:s7v61-status}
The TDM question is now settled at the correct level: it is a theorem for
anchored local boundary insertions and false as a consequence of shell
support alone.  Number shielding makes every anchored boundary return
negligible beyond all algebraic orders in the logarithmic schedule.
The remaining task is not more tree combinatorics; it is to derive ABDC
from the local high-sector dynamics, or to avoid the self-energy
localization problem by performing the connected expansion in the full
local Hamiltonian before Feshbach elimination.
\end{assessment}

\begin{decision}[V62 acquisition order]
\label{dec:s7v61-V62}
V62 will compare the two remaining routes.  It will first formulate a
conditional-expectation/Mobius decomposition of the high-sector
resolvent and identify the commutator estimate that would prove ABDC.
It will then test the alternative full-Hamiltonian route: expand the
local cutoff theory before Feshbach elimination and use number shielding
only for adjacent-cutoff differences.  The route with the weaker new
loader will be selected.
\end{decision}

\begin{firewall}[Claim boundary after seventh-series V61]
\label{fire:s7v61-claim}
V61 does not prove ABDC or HOLC for the physical self-energy, BCC,
support-resolved analytic-tail locality, a Cauchy theorem for occupation
cutoff removal, interacting full-density \(L^q\) control, tightness of the
interacting measures, the pairing/Bogoliubov analogue, normalized
interacting chart exit, slow-label sensitivity, weighted physical-source
incidence, patchwise coarse-action cocycles, terminal strong-coupling
matching, cutoff removal, OS reconstruction, or the full Yang--Mills
mass gap.
\end{firewall}

\paragraph{Parent-version record.}
V61 was formed from
\texttt{Renewal\_Yang\_Mills\_Mass\_Gap\_Research\_Notes\_V60.tex}.
The V60 parent had \(25\,955\) logical source lines, \(925\,562\) bytes,
and SHA--256
\[
 \texttt{965AB1F289BF9610A7527F923C4D529DE3E14369BA939BB47EB56CB8A129FC52}.
\]

% =============================================================================
% END APPENDED SEVENTH-SERIES V61 MATERIAL
% =============================================================================

% =============================================================================
% BEGIN APPENDED SEVENTH-SERIES V62 MATERIAL
% =============================================================================

\chapter{Coefficient stabilization and diagonal occupation-cutoff
independence}
\label{ch:s7v62}

\section{Purpose and route selection}
\label{sec:s7v62-purpose}

V61 leaves two possible treatments of the spatially nonlocal Feshbach
self-energy.  The first is to prove an anchored quasi-local decomposition
of the high-sector resolvent.  That route requires a new weighted
many-body propagation theorem and must solve the spectator-excitation
problem.  The second is to perform the connected expansion for the full
local cutoff Hamiltonian before applying any Feshbach map.

The second route has a weaker missing premise.  At every fixed
perturbative order, finite number bandwidth makes the coefficient exactly
independent of the occupation cutoff once the cutoff exceeds a linear
function of that order.  The already proved geometric majorant then
controls the remaining high orders.  This note proves the resulting
diagonal cutoff-independence theorem.  It does not construct the
continuum measure by itself; it removes dependence on the artificial
occupation regulator within the admissible weak-coupling schedules.

\section{Compressed ordered moments}
\label{sec:s7v62-moments}

Let \(H_0=d\Gamma(A)\) with \(A\geq\gamma_0I>0\), and let
\({\cal N}\) be the number operator.  Let
\[
 P_M=\mathbf1_{[0,M]}({\cal N}).
\]
For a local per-block cutoff, the same symbol will denote the product of
the corresponding local projections on the finite cylinder under
consideration.  Since a bound on total number bounds every local number,
the arguments below apply to either convention.

Let \(V\) have number bandwidth at most \(r\), and define
\[
 V^{(M)}:=P_MVP_M.
\]
Let \(O_1,\ldots,O_\ell\) be endpoint or source operators with number
bandwidths at most \(s_1,\ldots,s_\ell\), and put
\begin{equation}
 S_O:=\sum_{j=1}^{\ell}s_j.                             \label{eq:s7v62-SO}
\end{equation}
Every ordered reference moment coefficient is a vacuum matrix element of
a word made of the \(O_j\), \(n\) copies of \(V^{(M)}\), and
number-preserving reference semigroups.

\begin{lemma}[Cutoff projections are inert below the reachable number]
\label{lem:s7v62-inert}
Let \({\cal W}\) be any such word with \(n\) interaction vertices.
If
\begin{equation}
 M\geq S_O+nr,                                         \label{eq:s7v62-inert-condition}
\end{equation}
then its compressed vacuum matrix element equals the corresponding
uncompressed matrix element.  Consequently it is identical for every
cutoff \(M'\geq M\).
\end{lemma}

\begin{proof}
Read the word from the right vacuum to the left.  The vacuum starts in
number zero.  A number-preserving semigroup does not change its number
support.  Each occurrence of \(O_j\) can increase the largest possible
number by at most \(s_j\), and each occurrence of \(V\) can increase it
by at most \(r\).  Therefore every intermediate vector lies in
\(\operatorname{Ran}P_{S_O+nr}\).  Under
\eqref{eq:s7v62-inert-condition}, every inserted \(P_M\) acts as the
identity on every intermediate vector.  Removing them one at a time
gives the uncompressed word.  The same proof applies to \(M'\).
\end{proof}

\begin{theorem}[Connected-coefficient stabilization]
\label{thm:s7v62-coefficient}
Let \(c_n^{(M)}(O_1,\ldots,O_\ell)\) be any coefficient obtained as a
finite linear combination of products of ordered subset moments, with a
total of \(n\) interaction insertions; in particular it may be the
moment--cumulant coefficient of a normalized connected expectation.  If
\(M\geq S_O+nr\), then
\begin{equation}
 c_n^{(M)}(O_1,\ldots,O_\ell)
 =
 c_n^{(M')}(O_1,\ldots,O_\ell)
 \qquad(M'\geq M).
 \label{eq:s7v62-coefficient}
\end{equation}
\end{theorem}

\begin{proof}
In every product of subset moments, the number of endpoint occurrences
is at most \(\ell\), their total bandwidth is at most \(S_O\), and the
total number of interaction occurrences is \(n\).  Each factor therefore
satisfies the premise of \Cref{lem:s7v62-inert}.  Every factor, every
product, and hence the declared linear combination is independent of the
cutoff.
\end{proof}

\begin{corollary}[Gauge averages and finite sources]
\label{cor:s7v62-gauge}
If the compact gauge action commutes with \({\cal N}\), the stabilization
identity survives normalized Haar averaging of any finite collection of
endpoints or sources.
\end{corollary}

\begin{proof}
The bandwidth bound is invariant under the gauge action, and the equality
\eqref{eq:s7v62-coefficient} holds pointwise in the finitely many gauge
parameters.  Integrate it against normalized Haar measure.
\end{proof}

\section{A geometric tail gives cutoff independence}
\label{sec:s7v62-tail}

Let \(C_{A,B}^{(M)}(x,y)\) denote a normalized connected two-point
function for two fixed finite-bandwidth local endpoints \(A_x,B_y\), and
write its convergent expansion as
\begin{equation}
 C_{A,B}^{(M)}(x,y)
 =
 \sum_{n=0}^{\infty}c_n^{(M)}(x,y).
 \label{eq:s7v62-series}
\end{equation}

\begin{definition}[Uniform coefficient majorant]
\label{def:s7v62-UCM}
For a family \({\cal I}_L\) of admissible occupation cutoffs, UCM at
scale \(L\) means that there are
\(C_{A,B},\nu,\gamma>0\) and \(0<q_L<1\) such that
\begin{equation}
 |c_n^{(M)}(x,\tau;y,\sigma)|
 \leq
 C_{A,B}
 e^{-\nu d(x,y)}e^{-\gamma|\tau-\sigma|}
 q_L^n                                                   \label{eq:s7v62-UCM}
\end{equation}
for every \(M\in{\cal I}_L\), every \(n\geq0\), and all finite volumes.
\end{definition}

\begin{remark}[Loader already supplied by the reference tree theorem]
\label{rem:s7v62-UCM}
For the polynomial cubic--quartic interaction on a fixed local mode
space, V56 gives the spacetime edge, V57 gives the degree-weighted
interaction loss, and V58 removes the occupation loss at fixed Weyl
endpoints.  For fixed finite-rank endpoints the V57 endpoint factor is a
constant depending on their fixed number support.  Hence those theorems
supply UCM whenever the same physical interaction expansion is being
used and
\[
 \sup_{M\in{\cal I}_L}{\cal C}_{r,M}\rho_M\leq q_L<1.
\]
UCM is restated here because chart exit and slow-label uniformity are
not included in that conclusion.
\end{remark}

\begin{theorem}[Two-cutoff comparison]
\label{thm:s7v62-two-cutoff}
Assume UCM.  Let \(M_1,M_2\in{\cal I}_L\), put
\[
 M_\wedge:=\min\{M_1,M_2\},\qquad
 N_\wedge:=
 \left\lfloor{M_\wedge-S_O\over r}\right\rfloor,
\]
and suppose \(M_\wedge\geq S_O\).  Then
\begin{align}
 |C_{A,B}^{(M_1)}(x,\tau;y,\sigma)
  -C_{A,B}^{(M_2)}(x,\tau;y,\sigma)|
 &\leq
 {2C_{A,B}\over1-q_L}
 e^{-\nu d(x,y)}e^{-\gamma|\tau-\sigma|}
 q_L^{N_\wedge+1}.
 \label{eq:s7v62-two-cutoff}
\end{align}
\end{theorem}

\begin{proof}
\Cref{thm:s7v62-coefficient} gives
\[
 c_n^{(M_1)}=c_n^{(M_2)}
 \qquad(0\leq n\leq N_\wedge).
\]
Subtract the two series \eqref{eq:s7v62-series}, cancel these terms, use
\eqref{eq:s7v62-UCM} on both remaining tails, and sum the geometric
series.
\end{proof}

\begin{corollary}[Diagonal occupation-regulator independence]
\label{cor:s7v62-diagonal}
Suppose two cutoff schedules satisfy, for fixed
\(\chi,c_-,c_+>0\),
\begin{equation}
 c_-(\log L)^\chi
 \leq M_i(L)\leq c_+(\log L)^\chi,\qquad i=1,2,         \label{eq:s7v62-window-family}
\end{equation}
and UCM holds on the interval between them with
\begin{equation}
 q_L\leq L^{-\alpha}(\log L)^\beta                     \label{eq:s7v62-qrate}
\end{equation}
for fixed \(\alpha>0\), \(\beta\in\mathbb R\).  Then, for fixed
finite-bandwidth endpoints, the difference in
\eqref{eq:s7v62-two-cutoff} is
\begin{equation}
 O\!\left(
 e^{-\nu d(x,y)}e^{-\gamma|\tau-\sigma|}
 \exp[-c(\log L)^{\chi+1}]
 \right)                                               \label{eq:s7v62-super}
\end{equation}
for some \(c>0\), up to lower-order
\((\log L)^\chi\log\log L\) in the exponent.  In particular it is
smaller than \(L^{-p}\) for every fixed \(p>0\).
\end{corollary}

\begin{proof}
For large \(L\),
\[
 N_\wedge+1
 \geq {c_-\over2r}(\log L)^\chi.
\]
Taking logarithms in \eqref{eq:s7v62-qrate} gives
\[
 (N_\wedge+1)\log q_L
 \leq
 -{\alpha c_-\over2r}(\log L)^{\chi+1}
 +O((\log L)^\chi\log\log L).
\]
The factor \((1-q_L)^{-1}\) is bounded, and the displayed leading term
dominates \(p\log L\) for every fixed \(p\).
\end{proof}

\begin{corollary}[Cofinal equivalence class of occupation schedules]
\label{cor:s7v62-equivalence}
Under the hypotheses of
\Cref{cor:s7v62-diagonal}, if the connected correlations converge along
one admissible occupation schedule, they converge to the same values
along every other schedule satisfying
\eqref{eq:s7v62-window-family}--\eqref{eq:s7v62-qrate}.  The limiting
two-point function inherits the same spacetime exponential bound.
\end{corollary}

\begin{proof}
The two-schedule difference tends to zero by
\eqref{eq:s7v62-super}.  The UCM sum bounds each finite-cutoff
correlation by
\[
 {C_{A,B}\over1-q_L}
 e^{-\nu d(x,y)}e^{-\gamma|\tau-\sigma|}.
\]
Since \(q_L\to0\), pass to the limit at fixed spacetime points.
\end{proof}

\begin{firewall}[What diagonal independence does and does not remove]
\label{fire:s7v62-diagonal}
\Cref{cor:s7v62-equivalence} proves independence among admissible
cofinal occupation schedules.  It does not construct an interacting
theory at fixed lattice spacing with the occupation cutoff sent to
infinity while the bare coupling is held fixed.  The present continuum
programme uses the joint asymptotically-free schedule; a fixed-coupling
cutoff theorem would require a majorant uniform for arbitrarily large
\(M\), which V57 does not provide.
\end{firewall}

\section{The analytic Wilson tail can be restored locally}
\label{sec:s7v62-Wilson}

Let \(V_{\rm pol}=gH_3+g^2H_4\) denote the dimensionless polynomial
vertex and let \(T_M\) be the chart-localized fifth and higher Wilson
tail.  V55 proves on each fixed local block
\begin{equation}
 \|T_M\|\leq\rho_{{\rm tail},M}
 \leq Cg^3(M+6)^{5/2}.                                 \label{eq:s7v62-T}
\end{equation}
Introduce \(V(t)=V_{\rm pol}+tT_M\), \(0\leq t\leq1\).

\begin{theorem}[Local tail restoration]
\label{thm:s7v62-restoration}
Suppose the V57--V58 connected tree majorant holds uniformly for
\(V(t)\), with total smallness \(q_L<1\).  If the tail is supported on at
most \(r_T\) local modes, then
\begin{align}
 |C_{A,B}^{V(1)}-C_{A,B}^{V(0)}|
 &\leq
 {C_{A,B}{\cal D}_{r_T,M}\rho_{{\rm tail},M}
  \over(1-q_L)^2}
 e^{-\nu d(x,y)}e^{-\gamma|\tau-\sigma|}.              \label{eq:s7v62-restoration}
\end{align}
\end{theorem}

\begin{proof}
Differentiate the normalized connected expectation with respect to
\(t\).  The derivative is a connected expansion with one distinguished
local insertion \(T_M\) and arbitrary \(V(t)\)-vertices.  Apply
\Cref{cor:s7v61-sum}, including the endpoint spacetime edge, and use
\eqref{eq:s7v62-T}.  The bound is uniform in \(t\); integrate over
\([0,1]\).
\end{proof}

\begin{corollary}[Tail restoration vanishes on the V55 schedule]
\label{cor:s7v62-restoration-rate}
Under \(g=O(L^{-1/2})\), \(M=(\log L)^\chi\), and fixed \(r_T\),
\begin{equation}
 {\cal D}_{r_T,M}\rho_{{\rm tail},M}
 =
 O\!\left(
 L^{-3/2}
 (\log L)^{(2r_T+4)\chi}
 \right)
 =o(1).                                                \label{eq:s7v62-restoration-rate}
\end{equation}
\end{corollary}

\begin{proof}
By \eqref{eq:s7v61-DrM},
\({\cal D}_{r_T,M}=O(M^{2r_T+3/2})\).
Multiply by \(\rho_{{\rm tail},M}=O(g^3M^{5/2})\);
the total power of \(M\) is \(2r_T+4\).  Substitute the schedule.
\end{proof}

\begin{remark}[Why no number shielding is claimed for the tail]
\label{rem:s7v62-tail}
The analytic tail has unbounded number bandwidth, so the exact
coefficient stabilization theorem does not apply to insertions of
\(T_M\).  Instead \Cref{thm:s7v62-restoration} first removes the tail
with a vanishing local activity, applies the exact cutoff comparison to
the polynomial theory, and restores the tail at the end.  This avoids an
invalid finite-bandwidth assertion.
\end{remark}

\section{Dense finite-particle core and Weyl endpoints}
\label{sec:s7v62-core}

\begin{lemma}[Finite-particle cylinders form a dense reference core]
\label{lem:s7v62-core}
On every finite oscillator cylinder, the union
\[
 {\cal F}_{\rm fin}:=\bigcup_{s\geq0}P_{\leq s}{\cal H}
\]
is dense.  If a compact gauge group commutes with \({\cal N}\), the gauge
averages of these vectors are dense in the invariant subspace.
\end{lemma}

\begin{proof}
The spectral projections \(P_{\leq s}\) increase strongly to the
identity.  Hence \(P_{\leq s}\psi\to\psi\) for every \(\psi\).
Gauge averaging is a bounded projection commuting with every
\(P_{\leq s}\), so it preserves this convergence on the invariant
subspace.
\end{proof}

\begin{proposition}[Weyl endpoint approximation interface]
\label{prop:s7v62-Weyl}
Let \(W(f)\Omega,W(h)\Omega\) be fixed endpoint vectors and take
\[
 s_M=\lfloor M_\wedge/4\rfloor.
\]
Their reference high-number tails obey
\begin{equation}
 \|(I-P_{\leq s_M})W(f)\Omega\|
 +\|(I-P_{\leq s_M})W(h)\Omega\|
 \leq
 \left({e\|f\|^2\over s_M+1}\right)^{(s_M+1)/2}
 +
 \left({e\|h\|^2\over s_M+1}\right)^{(s_M+1)/2}.
 \label{eq:s7v62-Weyl-tail}
\end{equation}
The low endpoint pieces leave
\[
 N_\wedge\geq {M_\wedge\over8}-O(1)
\]
stabilized interaction orders when \(r=4\).
\end{proposition}

\begin{proof}
The tail estimate is \Cref{lem:s7v60-Poisson}.  The two low endpoint
bandwidths total at most \(2s_M\leq M_\wedge/2\).  Therefore
\[
 \left\lfloor{M_\wedge-2s_M\over4}\right\rfloor
 \geq {M_\wedge\over8}-O(1).
\]
\end{proof}

\begin{firewall}[Interacting endpoint continuity is still required]
\label{fire:s7v62-endpoint}
\Cref{prop:s7v62-Weyl} is a reference-vector approximation.  Extending
the cutoff-independent connected functional from the finite-particle
core to the full V58 Weyl core requires a uniform interacting endpoint
continuity bound.  Positivity of a normalized finite-volume state gives
operator-norm continuity, but the number truncation in
\eqref{eq:s7v62-Weyl-tail} is a vector approximation, not automatically
an operator-norm approximation of \(W(f)\).  This completion step is
therefore not claimed here.
\end{firewall}

\section{Status}
\label{sec:s7v62-status}

\begin{theorem}[What V62 proves]
\label{thm:s7v62-result}
For finite-number-bandwidth endpoints and the cubic--quartic local
interaction, every connected perturbative coefficient is exactly
independent of the occupation cutoff above
\(M=S_O+nr\).  Under the V57 geometric majorant this yields the
two-cutoff estimate \eqref{eq:s7v62-two-cutoff}; all logarithmic
occupation schedules in
\eqref{eq:s7v62-window-family} define the same cofinal correlations up
to an error beyond every algebraic order.  The local analytic Wilson tail
can be removed and restored with the vanishing error
\eqref{eq:s7v62-restoration-rate}.  Thus Euclidean occupation-regulator
independence can be proved before Feshbach elimination and does not
require ABDC.
\end{theorem}

\begin{proof}
Use
\Cref{thm:s7v62-coefficient,thm:s7v62-two-cutoff,
cor:s7v62-diagonal,cor:s7v62-equivalence,
thm:s7v62-restoration,cor:s7v62-restoration-rate}.
\end{proof}

\begin{assessment}[Progress at seventh-series V62]
\label{ass:s7v62-status}
The full-Hamiltonian route is strictly cheaper for Euclidean
correlations than localizing the Feshbach self-energy: it uses the
already available connected-series majorant and an exact algebraic
stabilization rule.  ABDC remains relevant only if one insists on a
spatially local low-sector Hamiltonian after elimination.  The next
upstream bottleneck is now visible: the reference connected expansion
must be loaded uniformly through the actual normalized chart-localized
interacting measure, including chart exit and slow labels, so that UCM
applies to the physical family rather than a formal fixed-chart series.
\end{assessment}

\begin{decision}[V63 acquisition order]
\label{dec:s7v62-V63}
V63 will attack normalized interacting chart exit.  It will distinguish
the Gaussian reference tail from the tilted probability, derive an exact
density-to-exit inequality, and identify whether the V56--V57 connected
activity bound supplies the required \(L^q\) norm of the complete local
density.  If it supplies only a partition-function expansion, the note
will state the missing exponential-moment card and go upstream to its
local action oscillation.
\end{decision}

\begin{firewall}[Claim boundary after seventh-series V62]
\label{fire:s7v62-claim}
V62 does not prove UCM uniformly through physical chart boundaries and
slow labels, interacting endpoint continuity on the full Weyl core, a
full-density \(L^q\) or entropy bound, normalized interacting chart exit,
BCC, ABDC, the pairing/Bogoliubov analogue, weighted physical-source
incidence, patchwise coarse-action cocycles, terminal strong-coupling
matching, removal of all regulators, OS reconstruction, or the full
Yang--Mills mass gap.
\end{firewall}

\paragraph{Parent-version record.}
V62 was formed from
\texttt{Renewal\_Yang\_Mills\_Mass\_Gap\_Research\_Notes\_V61.tex}.
The V61 parent had \(26\,446\) logical source lines, \(943\,962\) bytes,
and SHA--256
\[
 \texttt{5231B514EF5437D95AC80EAFA2DC59C8F37F8D3E85DE98DE3C8C3AB17536C7F5}.
\]

% =============================================================================
% END APPENDED SEVENTH-SERIES V62 MATERIAL
% =============================================================================

% =============================================================================
% BEGIN APPENDED SEVENTH-SERIES V63 MATERIAL
% =============================================================================

\chapter{Normalized chart exit as a rare bad-block polymer}
\label{ch:s7v63}

\section{Purpose and correction of the global target}
\label{sec:s7v63-purpose}

The Gaussian estimate of V38 controls the probability that one fixed
block leaves a weak-field chart.  Earlier notes correctly warned that
this does not automatically survive normalization by the interacting
density, but they continued to phrase the missing result as an
interacting chart-exit probability uniform over all blocks.  In an
infinite system that global no-exit event is the wrong target.  Even for
independent Gaussian blocks, some block exits with probability tending to
one.

The correct renormalization-group object is a gas of bad blocks.  One
must bound the probability that every block of a specified finite set is
bad by \(p_M^{|S|}\), with \(p_M\to0\).  Connected bad sets are then
summable polymers.  This note proves the Gaussian multi-block tail,
derives the exact density-to-polymer transfer, and identifies a
two-replica localized pressure defect as the missing normalized-density
loader.

\section{A global single-chart event is impossible}
\label{sec:s7v63-nogo}

\begin{proposition}[Independent-block no-go theorem]
\label{prop:s7v63-nogo}
Let \(E_x\) be the event that block \(x\) exits the chosen chart, and
suppose under a product reference probability that
\(\nu(E_x)=p\) with \(0<p<1\).  For a volume \(\Lambda\) of \(N\)
blocks,
\begin{equation}
 \nu\!\left(\bigcup_{x\in\Lambda}E_x\right)
 =1-(1-p)^N\longrightarrow1
 \quad(N\to\infty).                                    \label{eq:s7v63-nogo}
\end{equation}
Thus no positive one-block tail, however small, can make the event that
every block lies in one chart volume-uniform.
\end{proposition}

\begin{proof}
Independence gives
\(\nu(\cap_xE_x^c)=\prod_x(1-p)=(1-p)^N\).
Take complements and then the limit.
\end{proof}

\begin{firewall}[Meaning for the Yang--Mills construction]
\label{fire:s7v63-global}
The proposition does not say that weak-field coordinates are useless.
It says that they must be used locally, with chart-exit blocks retained
as polymers or covered by other compact-group charts.  Discarding every
configuration with one bad block cannot define the infinite-volume
theory.
\end{firewall}

\section{Correlated Gaussian multi-block tail}
\label{sec:s7v63-Gaussian}

Let \(q=(q_x)_{x\in\Lambda}\) be a centered real Gaussian field with
\(d_0\) coordinates per block and covariance \(C_\Lambda\).  Assume
\begin{equation}
 \|C_\Lambda\|_{\ell^2\to\ell^2}\leq\sigma^2            \label{eq:s7v63-cov}
\end{equation}
uniformly in the volume and in the declared slow-label patch.  For
\(R>0\), put
\[
 E_x(R):=\{|q_x|\geq R\},\qquad
 E_S(R):=\bigcap_{x\in S}E_x(R).
\]

\begin{lemma}[Quadratic exponential moment on a subset]
\label{lem:s7v63-mgf}
For every finite \(S\subset\Lambda\) and
\(\theta=1/(4\sigma^2)\),
\begin{equation}
 {\mathbb E}_\nu
 \exp\!\left(\theta\sum_{x\in S}|q_x|^2\right)
 \leq \exp\!\left({d_0\over2}|S|\right).
 \label{eq:s7v63-mgf}
\end{equation}
\end{lemma}

\begin{proof}
Let \(C_S\) be the covariance of the \(d_0|S|\)-dimensional marginal.
Compression cannot increase operator norm, so
\(\|C_S\|\leq\sigma^2\).  The Gaussian determinant identity gives
\[
 {\mathbb E}
 e^{\theta|q_S|^2}
 =
 \det(I-2\theta C_S)^{-1/2}.
\]
If \(\lambda_j\) are the eigenvalues of \(C_S\), then
\(0\leq2\theta\lambda_j\leq1/2\).  The elementary inequality
\(-\log(1-u)\leq2u\) on \(0\leq u\leq1/2\) yields
\begin{align*}
 \log{\mathbb E}e^{\theta|q_S|^2}
 &=
 {1\over2}\sum_j-\log(1-2\theta\lambda_j)\\
 &\leq
 2\theta\sum_j\lambda_j
 \leq
 2\theta d_0|S|\sigma^2
 ={d_0\over2}|S|.
\end{align*}
\end{proof}

\begin{theorem}[Gaussian bad-set bound]
\label{thm:s7v63-Gaussian}
Under \eqref{eq:s7v63-cov},
\begin{equation}
 \nu(E_S(R))
 \leq
 \exp\!\left[
 -\left({R^2\over4\sigma^2}-{d_0\over2}\right)|S|
 \right].
 \label{eq:s7v63-Gaussian}
\end{equation}
\end{theorem}

\begin{proof}
On \(E_S(R)\),
\(\sum_{x\in S}|q_x|^2\geq|S|R^2\).
Exponential Markov inequality and \Cref{lem:s7v63-mgf} give
\[
 \nu(E_S(R))
 \leq
 e^{-\theta|S|R^2}
 {\mathbb E}e^{\theta\sum_{x\in S}|q_x|^2}.
\]
Insert \(\theta=1/(4\sigma^2)\).
\end{proof}

\begin{remark}[No independence was used]
\label{rem:s7v63-correlated}
The multi-block exponent follows from a uniform covariance operator norm,
not from a product decomposition.  This is the form compatible with the
V46 patchwise quadratic covariance bound.  Exact gauge and toron modes
must be removed or retained as slow variables before
\eqref{eq:s7v63-cov} can hold.
\end{remark}

\section{Local marginal density transfer}
\label{sec:s7v63-density}

Let \(\mu_\Lambda\ll\nu_\Lambda\) be the complete normalized interacting
probability on a finite volume.  Denote its marginal on a finite set
\(S\) by \(\mu_S\), the Gaussian marginal by \(\nu_S\), and
\[
 r_S:={d\mu_S\over d\nu_S}.
\]

\begin{definition}[Local marginal density card, LMDC]
\label{def:s7v63-LMDC}
LMDC with parameters \(q>1\) and \(K_q\geq0\) is
\begin{equation}
 \|r_S\|_{L^q(\nu_S)}
 \leq e^{K_q|S|}
 \label{eq:s7v63-LMDC}
\end{equation}
for every finite \(S\), uniformly in the ambient volume, cutoff, and
declared slow-label patch.
\end{definition}

\begin{theorem}[Normalized interacting bad-set bound]
\label{thm:s7v63-interacting}
Assume \eqref{eq:s7v63-cov} and LMDC.  With
\(q'=q/(q-1)\),
\begin{equation}
 \mu_\Lambda(E_S(R))
 \leq p_R^{|S|},                                       \label{eq:s7v63-bad}
\end{equation}
where
\begin{equation}
 p_R:=
 \exp\!\left[
 K_q+{d_0\over2q'}
 -{R^2\over4\sigma^2q'}
 \right].
 \label{eq:s7v63-pR}
\end{equation}
\end{theorem}

\begin{proof}
The event \(E_S(R)\) is measurable with respect to the \(S\)-marginal.
Holder's inequality gives
\[
 \mu_\Lambda(E_S(R))
 =
 \int_{E_S(R)}r_S\,d\nu_S
 \leq
 \|r_S\|_{L^q(\nu_S)}
 \nu(E_S(R))^{1/q'}.
\]
Use \eqref{eq:s7v63-LMDC} and
\eqref{eq:s7v63-Gaussian}.
\end{proof}

\begin{corollary}[Logarithmic chart schedule]
\label{cor:s7v63-log}
With \(R^2=M=(\log L)^\chi\) and fixed LMDC constants,
\begin{equation}
 p_M
 \leq C\exp[-c(\log L)^\chi]\longrightarrow0           \label{eq:s7v63-log}
\end{equation}
for constants \(C,c>0\).  Every fixed polynomial in \(M\) times \(p_M\)
tends to zero.
\end{corollary}

\begin{proof}
Apply \eqref{eq:s7v63-pR}; the positive terms in its exponent are
constant, while the negative term is a fixed positive multiple of
\(M\).  Exponential decay in \(M\) dominates every fixed power of \(M\).
\end{proof}

\begin{firewall}[A global density norm is the wrong quantity]
\label{fire:s7v63-global-density}
For an extensive interaction,
\(\|d\mu_\Lambda/d\nu_\Lambda\|_{L^q}\) generally grows exponentially in
\(|\Lambda|\).  Substituting that global norm in Holder's inequality
would destroy \eqref{eq:s7v63-bad}.  LMDC permits precisely the expected
exponential cost in the size of the observed bad set and no cost in the
ambient volume.
\end{firewall}

\section{Bad-block polymers are summable}
\label{sec:s7v63-polymer}

Let the block incidence graph have maximum degree \(\Delta\).  A bad
polymer is a nonempty connected set of bad blocks.

\begin{lemma}[Coarse lattice-animal count]
\label{lem:s7v63-animals}
The number \(a_n(x)\) of connected vertex sets of size \(n\) containing a
fixed root \(x\) satisfies
\begin{equation}
 a_n(x)\leq\Delta^{2(n-1)}.                             \label{eq:s7v63-animals}
\end{equation}
\end{lemma}

\begin{proof}
Every finite connected set has a spanning tree rooted at \(x\).  A
depth-first traversal of that tree starts at \(x\), has length
\(2(n-1)\), and visits every vertex of the set.  Fix a deterministic
ordering of incident edges and choose, for each set, the lexicographically
first such traversal.  This injects the connected sets into walks of
length \(2(n-1)\) from \(x\).  There are at most
\(\Delta^{2(n-1)}\) such walks.
\end{proof}

\begin{theorem}[Bad-polymer summability]
\label{thm:s7v63-polymer}
Assume \(\mu(E_S(R))\leq p_R^{|S|}\).  For every \(\alpha\geq0\) such
that
\begin{equation}
 \Delta^2e^\alpha p_R<1,                               \label{eq:s7v63-KP}
\end{equation}
one has the rooted bound
\begin{equation}
 \sup_x
 \sum_{\substack{S\ni x\\S\ {\rm connected}}}
 e^{\alpha|S|}\mu(E_S(R))
 \leq
 {e^\alpha p_R
  \over1-\Delta^2e^\alpha p_R}.                        \label{eq:s7v63-polymer}
\end{equation}
\end{theorem}

\begin{proof}
Group the connected sets by their size and use
\Cref{lem:s7v63-animals}:
\begin{align*}
 \sum_{\substack{S\ni x\\S\ {\rm connected}}}
 e^{\alpha|S|}\mu(E_S(R))
 &\leq
 \sum_{n\geq1}
 \Delta^{2(n-1)}e^{\alpha n}p_R^n\\
 &=
 {e^\alpha p_R\over1-\Delta^2e^\alpha p_R}.
\end{align*}
\end{proof}

\begin{corollary}[Normalized chart exit becomes a vanishing activity]
\label{cor:s7v63-activity}
Under LMDC and the logarithmic chart schedule, the bad-block family
satisfies \eqref{eq:s7v63-KP} for every fixed \(\alpha\) and all
sufficiently large \(L\).  Its rooted activity is
\(O(e^{-cM})\) and survives every fixed polynomial occupation loss.
\end{corollary}

\begin{proof}
Use \Cref{cor:s7v63-log,thm:s7v63-polymer}.
\end{proof}

\section{Exact replica representation of the local density norm}
\label{sec:s7v63-replica}

It remains to load LMDC from a normalized connected expansion.  The
case \(q=2\) has an exact two-replica representation.

Write a configuration as \((x,y)\), where \(x\) is the field on \(S\)
and \(y\) the field on \(S^c\).  Disintegrate
\[
 d\nu_\Lambda(x,y)=d\nu_S(x)\,d\nu(dy\mid x).
\]
Let the complete density be
\[
 {d\mu_\Lambda\over d\nu_\Lambda}(x,y)
 ={e^{\ell(x,y)}\over Z},\qquad
 Z:=\int e^\ell\,d\nu_\Lambda.
\]
Define the two-replica partition function glued on \(S\) by
\begin{align}
 Z_S^{(2)}
 &:=
 \int d\nu_S(x)
 \prod_{a=1}^2
 \left[
 \int e^{\ell(x,y_a)}\,d\nu(dy_a\mid x)
 \right].
 \label{eq:s7v63-Z2}
\end{align}

\begin{theorem}[Replica identity for the marginal \(L^2\) norm]
\label{thm:s7v63-replica}
With the preceding notation,
\begin{equation}
 \|r_S\|_{L^2(\nu_S)}^2
 ={Z_S^{(2)}\over Z^2}.                                \label{eq:s7v63-replica}
\end{equation}
\end{theorem}

\begin{proof}
Disintegration gives
\[
 r_S(x)
 ={1\over Z}\int e^{\ell(x,y)}\,d\nu(dy\mid x).
\]
Square this expression, integrate against \(d\nu_S(x)\), and introduce
two independent conditional variables \(y_1,y_2\).  The result is
\eqref{eq:s7v63-Z2}/\(Z^2\).
\end{proof}

\begin{definition}[Replica pressure locality card, RPLC]
\label{def:s7v63-RPLC}
RPLC with constant \(K_{\rm rep}\) is the bound
\begin{equation}
 \log Z_S^{(2)}-2\log Z
 \leq K_{\rm rep}|S|                                  \label{eq:s7v63-RPLC}
\end{equation}
for every finite \(S\), uniformly in ambient volume, cutoffs, and slow
labels.
\end{definition}

\begin{corollary}[RPLC implies LMDC]
\label{cor:s7v63-RPLC}
RPLC implies LMDC at \(q=2\) with
\begin{equation}
 K_2={K_{\rm rep}\over2}.                              \label{eq:s7v63-K2}
\end{equation}
\end{corollary}

\begin{proof}
Take the logarithm of \eqref{eq:s7v63-replica}, apply
\eqref{eq:s7v63-RPLC}, and divide by two.
\end{proof}

\section{How a convergent polymer pressure would prove RPLC}
\label{sec:s7v63-pressure}

\begin{theorem}[Localized replica-defect theorem]
\label{thm:s7v63-defect}
Suppose the ordinary and \(S\)-glued two-replica logarithmic partition
functions admit absolutely convergent connected-polymer expansions
\begin{align}
 \log Z&=\sum_\Gamma\varphi(\Gamma),                    \label{eq:s7v63-logZ}\\
 \log Z_S^{(2)}
 &=\sum_{\Gamma^{(2)}}\varphi_S^{(2)}(\Gamma^{(2)}).
 \label{eq:s7v63-logZ2}
\end{align}
Assume:
\begin{enumerate}[label=\textup{(\roman*)},leftmargin=3.8em]
\item every glued-replica polymer disjoint from \(S\) factorizes into one
      of two ordinary copies, so its contribution cancels in
      \(\log Z_S^{(2)}-2\log Z\);
\item the rooted absolute cluster sums satisfy
\begin{align}
 \sup_x\sum_{\Gamma\ni x}|\varphi(\Gamma)|&\leq K,
 \label{eq:s7v63-root-original}\\
 \sup_x\sum_{\Gamma^{(2)}\ni x}
 |\varphi_S^{(2)}(\Gamma^{(2)})|&\leq K^{(2)}
 \label{eq:s7v63-root-replica}
\end{align}
uniformly in \(S\).
\end{enumerate}
Then RPLC holds with
\begin{equation}
 K_{\rm rep}=K^{(2)}+2K.                               \label{eq:s7v63-Krep}
\end{equation}
\end{theorem}

\begin{proof}
By assumption (i), only clusters meeting the glued set \(S\) remain in
the pressure difference.  Bound their absolute sum and overcount by a
chosen intersection point:
\begin{align*}
 |\log Z_S^{(2)}-2\log Z|
 &\leq
 \sum_{x\in S}
 \left[
 \sum_{\Gamma^{(2)}\ni x}|\varphi_S^{(2)}(\Gamma^{(2)})|
 +2\sum_{\Gamma\ni x}|\varphi(\Gamma)|
 \right]\\
 &\leq(K^{(2)}+2K)|S|.
\end{align*}
The desired one-sided inequality follows.
\end{proof}

\begin{firewall}[The ordinary tree theorem is not yet the replica theorem]
\label{fire:s7v63-replica}
V56--V58 control ordinary connected insertions in a fixed reference
sector.  RPLC additionally requires an absolutely convergent expansion
of the measure in which two complements share the same field on \(S\),
and exact cancellation of all clusters disjoint from \(S\).  This is a
localized defect expansion, not a consequence of a global free-energy
bound.  Its root constants must be proved uniformly in slow labels.
\end{firewall}

\section{Status and next acquisition}
\label{sec:s7v63-status}

\begin{theorem}[What V63 proves]
\label{thm:s7v63-result}
The global single-chart target is invalid in infinite volume.  Under a
uniform Gaussian covariance norm and the local marginal density card,
the normalized interacting probability that every block in \(S\) exits
the chart is at most \(p_M^{|S|}\), with
\(p_M=O(e^{-cM})\).  Connected bad sets then satisfy the rooted polymer
bound \eqref{eq:s7v63-polymer}.  The missing LMDC at \(q=2\) is exactly
equivalent to the two-replica pressure ratio
\eqref{eq:s7v63-replica}, and a localized connected-polymer expansion
with root bounds proves it by
\Cref{thm:s7v63-defect}.
\end{theorem}

\begin{proof}
Use
\Cref{prop:s7v63-nogo,thm:s7v63-Gaussian,
thm:s7v63-interacting,thm:s7v63-polymer,
thm:s7v63-replica,thm:s7v63-defect}.
\end{proof}

\begin{assessment}[Progress at seventh-series V63]
\label{ass:s7v63-status}
Normalized interacting chart exit has been reformulated at the correct
local level.  No extensive denominator must be paid: the needed object
is a replica pressure defect proportional to the observed bad set, not
the ambient volume.  Once RPLC is proved, the former chart-exit firewall
becomes a standard exponentially small polymer row.  The next work is to
construct the two-replica defect expansion from the same Weyl/tree
technology while checking that gluing on \(S\) preserves the reference
gap and the weighted spacetime edge.
\end{assessment}

\begin{decision}[V64 acquisition order]
\label{dec:s7v63-V64}
V64 will analyze the glued quadratic reference.  It will diagonalize the
two replicas into symmetric and antisymmetric fields away from \(S\),
identify the boundary condition imposed by equality on \(S\), and prove
or refute a uniform gap and weighted covariance bound.  If the glued
reference preserves the V56 stability kernel, the ordinary tree proof
will be lifted to the replica defect and RPLC; otherwise the missing
boundary covariance estimate will be isolated.
\end{decision}

\begin{firewall}[Claim boundary after seventh-series V63]
\label{fire:s7v63-claim}
V63 does not prove RPLC or LMDC for the physical Yang--Mills density,
uniform glued-reference propagation, the full bad-block polymer
expansion, interacting endpoint continuity, UCM through all slow-label
patches, BCC, the pairing/Bogoliubov analogue, weighted physical-source
incidence, patchwise coarse-action cocycles, terminal strong-coupling
matching, removal of all regulators, OS reconstruction, or the full
Yang--Mills mass gap.
\end{firewall}

\paragraph{Parent-version record.}
V63 was formed from
\texttt{Renewal\_Yang\_Mills\_Mass\_Gap\_Research\_Notes\_V62.tex}.
The V62 parent had \(26\,916\) logical source lines, \(961\,508\) bytes,
and SHA--256
\[
 \texttt{D559406640DA883813024610CA08DC1AE7328059280E12CE2444C13B44013261}.
\]

% =============================================================================
% END APPENDED SEVENTH-SERIES V63 MATERIAL
% =============================================================================

% =============================================================================
% BEGIN APPENDED SEVENTH-SERIES V64 MATERIAL
% =============================================================================

\chapter{The glued Gaussian reference and the replica-defect tree bound}
\label{ch:s7v64}

\section{Purpose}
\label{sec:s7v64-purpose}

V63 reduces normalized bad-block control to the pressure difference
\(\log Z_S^{(2)}-2\log Z\).  The new question is whether identifying the
two replica fields on \(S\) destroys the uniform Gaussian gap or creates
an extensive coupling between the replicas.  It does neither.

The conditional Gaussian variables diagonalize into symmetric and
antisymmetric fields.  Each single-replica marginal remains the original
Gaussian.  The cross-replica covariance outside \(S\) is the difference
between the full covariance and the Dirichlet covariance on \(S^c\);
the Schur formula factors this difference through \(S\).  Its total
weighted bond mass is therefore proportional to \(|S|\), not the ambient
volume.  Marking one cross-replica edge in every mixed tree converts this
factorization into the replica pressure locality card for every
interaction already controlled by a uniform Weyl/tree activity.

\section{Conditional Gaussian block algebra}
\label{sec:s7v64-block}

Let the real one-particle coordinate space split as
\({\cal K}_S\oplus{\cal K}_T\), where \(T=S^c\), and let the positive
precision operator be
\begin{equation}
 L=
 \begin{pmatrix}
 A&B\\ B^*&D
 \end{pmatrix},
 \qquad L\geq\gamma I,\quad\gamma>0.                   \label{eq:s7v64-L}
\end{equation}
Put
\begin{equation}
 K:=D^{-1},\qquad
 F:=A-BD^{-1}B^*.                                     \label{eq:s7v64-KF}
\end{equation}

\begin{lemma}[Schur covariance formulas]
\label{lem:s7v64-Schur}
The covariance \(C=L^{-1}\) has blocks
\begin{align}
 C_{SS}&=F^{-1},                                       \label{eq:s7v64-CSS}\\
 C_{ST}&=-F^{-1}BK,                                   \label{eq:s7v64-CST}\\
 C_{TS}&=-KB^*F^{-1},                                 \label{eq:s7v64-CTS}\\
 C_{TT}&=K+KB^*F^{-1}BK.                              \label{eq:s7v64-CTT}
\end{align}
Moreover \(D\geq\gamma I\), \(F\geq\gamma I\), and hence
\begin{equation}
 \|K\|\leq\gamma^{-1},\qquad
 \|F^{-1}\|\leq\gamma^{-1}.                            \label{eq:s7v64-gaps}
\end{equation}
\end{lemma}

\begin{proof}
The block inverse formula gives
\eqref{eq:s7v64-CSS}--\eqref{eq:s7v64-CTT}.
Compression of \eqref{eq:s7v64-L} to \({\cal K}_T\) gives
\(D\geq\gamma I\).  For \(x\in{\cal K}_S\), minimize
\[
 \langle(x,y),L(x,y)\rangle
\]
over \(y\).  The minimizer is \(y=-D^{-1}B^*x\), and the minimum is
\(\langle x,Fx\rangle\).  Since
\[
 \langle(x,y),L(x,y)\rangle
 \geq\gamma(\|x\|^2+\|y\|^2),
\]
the minimum is at least \(\gamma\|x\|^2\).  Thus
\(F\geq\gamma I\), proving \eqref{eq:s7v64-gaps}.
\end{proof}

Let \(X\) denote the field on \(S\).  Conditional on \(X=x\), one
replica field \(Y\) on \(T\) is Gaussian with
\begin{equation}
 {\mathbb E}(Y\mid X=x)=-KB^*x,\qquad
 {\rm Cov}(Y\mid X)=K.                                \label{eq:s7v64-conditional}
\end{equation}
The glued reference uses two conditionally independent copies
\(Y_1,Y_2\) with the same \(X\).

\begin{proposition}[Symmetric--antisymmetric diagonalization]
\label{prop:s7v64-diagonal}
Put
\[
 Z_+:={Y_1+Y_2\over\sqrt2},\qquad
 Z_-:={Y_1-Y_2\over\sqrt2}.
\]
Conditional on \(X=x\),
\begin{align}
 {\mathbb E}(Z_+\mid X=x)&=-\sqrt2KB^*x,&
 {\rm Cov}(Z_+\mid X)&=K,                              \label{eq:s7v64-plus}\\
 {\mathbb E}(Z_-\mid X=x)&=0,&
 {\rm Cov}(Z_-\mid X)&=K,                              \label{eq:s7v64-minus}
\end{align}
and \(Z_+,Z_-\) are conditionally independent.  In the centered
coordinates
\[
 X,\quad Z_++\sqrt2KB^*X,\quad Z_-,
\]
the glued Gaussian precision is the direct sum
\begin{equation}
 F\oplus D\oplus D,                                    \label{eq:s7v64-direct}
\end{equation}
so every component retains the lower floor \(\gamma\).
\end{proposition}

\begin{proof}
Write
\[
 Y_a=-KB^*X+\xi_a,
\]
where \(\xi_1,\xi_2\) are conditionally independent centered Gaussians
of covariance \(K\).  Their orthogonal sum and difference are independent
centered Gaussians with the same covariance \(K\).  The marginal
precision of \(X\) is \(F\), while the two centered conditional variables
have precision \(D\).  Apply \Cref{lem:s7v64-Schur}.
\end{proof}

\section{The cross-replica covariance factors through the glued set}
\label{sec:s7v64-cross}

Let \(G^{ab}\), \(a,b\in\{1,2\}\), denote the covariance between replica
\(a\) and replica \(b\), with the convention that the two replica
coordinates on \(S\) are the same variable \(X\).

\begin{theorem}[Exact glued covariance]
\label{thm:s7v64-covariance}
The one-replica covariances are unchanged:
\begin{equation}
 G^{11}=G^{22}=C.                                      \label{eq:s7v64-diagonal-cov}
\end{equation}
The cross-replica blocks are
\begin{align}
 G^{12}_{SS}&=C_{SS},&
 G^{12}_{ST}&=C_{ST},                                  \label{eq:s7v64-cross-S}\\
 G^{12}_{TS}&=C_{TS},&
 G^{12}_{TT}&=C_{TT}-K
 =KB^*F^{-1}BK.                                       \label{eq:s7v64-cross-T}
\end{align}
In particular, the cross covariance on \(T\times T\) factors through
\(S\).
\end{theorem}

\begin{proof}
Each pair \((X,Y_a)\) has exactly the original disintegration, proving
\eqref{eq:s7v64-diagonal-cov} and the three cross blocks with at least
one index in \(S\).  Conditional independence gives
\begin{align*}
 {\rm Cov}(Y_1,Y_2)
 &=
 {\rm Cov}\bigl({\mathbb E}(Y_1\mid X),
                 {\mathbb E}(Y_2\mid X)\bigr)\\
 &=KB^*{\rm Cov}(X)BK
 =KB^*F^{-1}BK.
\end{align*}
Equation \eqref{eq:s7v64-CTT} identifies this with \(C_{TT}-K\).
\end{proof}

\begin{remark}[Positivity]
\label{rem:s7v64-positive}
The full \(2\times2\) replica covariance \(G\) is positive semidefinite
because it is the covariance of the explicit conditional construction.
Thus the V56 stability cancellation applies to Weyl moments formed with
\(G\); no entrywise sign condition is needed.
\end{remark}

\section{Weighted bond mass anchored on \(S\)}
\label{sec:s7v64-weighted}

For a possibly rectangular block kernel \(T=(T_{xy})\), define the
two-sided weighted Schur norm
\begin{align}
 \|T\|_{{\cal B}_\mu}
 :=
 \max\left\{
 \sup_x\sum_y e^{\mu d(x,y)}\|T_{xy}\|,
 \sup_y\sum_x e^{\mu d(x,y)}\|T_{xy}\|
 \right\}.
 \label{eq:s7v64-Bmu}
\end{align}
The triangle inequality proves the algebra property
\begin{equation}
 \|T_1T_2\|_{{\cal B}_\mu}
 \leq
 \|T_1\|_{{\cal B}_\mu}\|T_2\|_{{\cal B}_\mu}.          \label{eq:s7v64-algebra}
\end{equation}

Assume the V46-type bounds
\begin{equation}
 \|C\|_{{\cal B}_\mu}\leq Q,\qquad
 \|K\|_{{\cal B}_\mu}\leq Q_D,\qquad
 \|B\|_{{\cal B}_\mu}\leq J_B                           \label{eq:s7v64-bounds}
\end{equation}
uniformly in \(S\), volume, and slow label.  The bound for \(K\) follows
from the same Combes--Thomas proof on the principal subdomain \(T\),
because compression preserves the lower floor and cannot enlarge the
finite-range hopping row.

\begin{lemma}[Boundary-to-bulk regression]
\label{lem:s7v64-regression}
For \(R_T:=KB^*:{\cal K}_S\to{\cal K}_T\),
\begin{equation}
 \|R_T\|_{{\cal B}_\mu}\leq Q_DJ_B.                    \label{eq:s7v64-regression}
\end{equation}
\end{lemma}

\begin{proof}
Apply \eqref{eq:s7v64-algebra} to \(K B^*\).
\end{proof}

\begin{theorem}[Anchored cross-bond estimate]
\label{thm:s7v64-cross-mass}
Under \eqref{eq:s7v64-bounds}, there is a constant
\begin{equation}
 Q_\times
 :=
 3Q+Q(Q_DJ_B)^2                                       \label{eq:s7v64-Qcross}
\end{equation}
such that
\begin{equation}
 \sum_{u,v}
 e^{\mu d(u,v)}\|G^{12}_{uv}\|
 \leq Q_\times |S|.                                   \label{eq:s7v64-cross-mass}
\end{equation}
The sums include the appropriate \(S\) and \(T\) coordinate blocks but
not the replica-color multiplicity.
\end{theorem}

\begin{proof}
For the blocks with at least one endpoint in \(S\), use
\eqref{eq:s7v64-cross-S}.  Summing a row or column of the corresponding
subblock of \(C\) costs at most \(Q\).  Allowing separately for
\(S\times S\), \(S\times T\), and \(T\times S\) gives the coarse bound
\(3Q|S|\).

For \(T\times T\), use
\[
 G^{12}_{TT}=R_TF^{-1}R_T^*.
\]
Write the kernel as a sum over \(s,t\in S\).  Since
\[
 d(x,y)\leq d(x,s)+d(s,t)+d(t,y),
\]
the exponential weight factors among the three kernels.  The two-sided
row/column bounds, \Cref{lem:s7v64-regression}, and
\(\|F^{-1}\|_{{\cal B}_\mu}\leq\|C\|_{{\cal B}_\mu}\leq Q\)
give
\begin{align*}
 \sum_{x,y\in T}e^{\mu d(x,y)}
 \|(G^{12}_{TT})_{xy}\|
 &\leq
 \sum_{s\in S}
 (Q_DJ_B)\,Q\,(Q_DJ_B)\\
 &=Q(Q_DJ_B)^2|S|.
\end{align*}
Add the two contributions.
\end{proof}

\begin{firewall}[Uniform row bound versus anchored total mass]
\label{fire:s7v64-mass}
A uniform weighted row bound on \(G^{12}\) alone would still allow a
total mass proportional to the ambient volume.  The factorization
\eqref{eq:s7v64-cross-T} is what improves the full double sum to
\eqref{eq:s7v64-cross-mass}, proportional only to the glued set.
\end{firewall}

\section{Mixed replica trees contain an anchored edge}
\label{sec:s7v64-tree}

Color every interaction vertex by its replica number.  A tree using
vertices of both colors must contain a cross-color edge.

\begin{lemma}[Marked cross-edge bound]
\label{lem:s7v64-marked}
Let a positive Gaussian Weyl cumulant have \(m\geq2\) local vertices
colored \(1,2\), and suppose both colors occur.  In the V56 tree bound,
mark the first cross-color edge according to a fixed ordering of tree
edges.  After summing all spatial positions, the marked edge costs at
most \(Q_\times|S|\), while every unmarked edge costs the ordinary
weighted row constant \(Q_G\), where
\begin{equation}
 Q_G:=2Q+Q_D.                                          \label{eq:s7v64-QG}
\end{equation}
There are at most \(m-1\) choices of the marked edge.
\end{lemma}

\begin{proof}
Every path in a tree joining a color-one vertex to a color-two vertex
contains a cross-color edge.  Sum a chosen cross edge by
\eqref{eq:s7v64-cross-mass}.  Removing that edge leaves two rooted
forests.  Sum their vertices outward from the two roots, using the
weighted row bounds for \(C\), \(K\), and their difference.  The coarse
constant \(Q_G=2Q+Q_D\) bounds every covariance block.  A tree has
\(m-1\) edges, which bounds the number of possible marks.
\end{proof}

\begin{theorem}[Replica-defect tree theorem]
\label{thm:s7v64-replica-tree}
Assume:
\begin{enumerate}[label=\textup{(\roman*)},leftmargin=3.8em]
\item the complete local interaction in each replica has the
      V57 degree-weighted Weyl representation;
\item the same representation and V56 stability argument hold for the
      positive glued covariance \(G\);
\item after all ordinary edge and degree sums, each added interaction
      vertex costs at most \(q_{\rm rep}<1\), uniformly in \(S\), volume,
      cutoff, and slow label.
\end{enumerate}
Then the sum of all connected mixed-replica pressure clusters satisfies
\begin{equation}
 |\log Z_S^{(2)}-2\log Z|
 \leq
 C_\times |S|
 {q_{\rm rep}\over(1-q_{\rm rep})^2},                  \label{eq:s7v64-replica-tree}
\end{equation}
where \(C_\times\) depends on \(Q_\times,Q_G\) and the fixed local mode
number, but not on \(S\) or the ambient volume.
\end{theorem}

\begin{proof}
Every connected cluster using only color one is an ordinary cluster in
the first replica; the same holds for color two.  These two pure-color
sums are exactly \(2\log Z\) and cancel from the difference.  Every
remaining cluster uses both colors.  Apply the V56 stable tree bound and
mark a cross edge as in \Cref{lem:s7v64-marked}.  Its spatial sum costs
\(Q_\times|S|\); all remaining branches are paid by ordinary row sums.
The V57 degree-weighted Prüfer normalization absorbs the local Weyl
moments.  At order \(m\), marking an edge leaves at most a factor \(m\),
so the normalized absolute order is bounded by
\[
 C_\times|S|\,m q_{\rm rep}^m.
\]
Sum
\(\sum_{m\geq1}m q_{\rm rep}^m
=q_{\rm rep}(1-q_{\rm rep})^{-2}\).
\end{proof}

\begin{corollary}[RPLC and LMDC for the replica-admissible density]
\label{cor:s7v64-RPLC}
Under the hypotheses of
\Cref{thm:s7v64-replica-tree}, RPLC holds with
\begin{equation}
 K_{\rm rep}
 =
 C_\times{q_{\rm rep}\over(1-q_{\rm rep})^2},           \label{eq:s7v64-Krep}
\end{equation}
and LMDC holds at \(q=2\) with
\begin{equation}
 K_2={C_\times\over2}
 {q_{\rm rep}\over(1-q_{\rm rep})^2}.                  \label{eq:s7v64-K2}
\end{equation}
\end{corollary}

\begin{proof}
Use \Cref{thm:s7v63-replica,cor:s7v63-RPLC} and
\eqref{eq:s7v64-replica-tree}.
\end{proof}

\section{Loading the replica smallness under the logarithmic schedule}
\label{sec:s7v64-schedule}

If one physical block contains \(r\) oscillator modes, the glued
two-replica block contains at most \(2r\).  Vertices on \(S\) may occur
twice because the two actions are evaluated on the common field.
Consequently a sufficient replica activity is
\begin{equation}
 q_{\rm rep}
 \leq
 2C_{2r}(M+1)^{4r+3/2}\rho_M,                          \label{eq:s7v64-qrep}
\end{equation}
with all covariance row constants absorbed into \(C_{2r}\).

\begin{proposition}[Replica activity vanishes]
\label{prop:s7v64-qrep}
If \(\rho_M\) is any explicit V55 local activity containing a negative
power of \(L\), and \(M=(\log L)^\chi\), then
\begin{equation}
 q_{\rm rep}\longrightarrow0.                          \label{eq:s7v64-qrep-zero}
\end{equation}
The same holds after adding the V55 fifth-order magnetic tail.
\end{proposition}

\begin{proof}
The prefactor in \eqref{eq:s7v64-qrep} is a fixed power of \(M\).
\Cref{cor:s7v55-polynomial} states that every such loss is absorbed by
each V55 operator-activity row.  The fifth-order tail is one of those
rows.
\end{proof}

\begin{corollary}[Good-field normalized exit row]
\label{cor:s7v64-exit}
For every density whose complete local interaction satisfies the three
hypotheses of \Cref{thm:s7v64-replica-tree}, the logarithmic chart
schedule obeys LMDC and hence
\begin{equation}
 \mu(E_S(R))\leq
 \exp[-(cM-C)|S|].                                     \label{eq:s7v64-exit}
\end{equation}
Its connected bad-block activity is \(O(e^{-cM})\).
\end{corollary}

\begin{proof}
\Cref{prop:s7v64-qrep} bounds \(K_2\) in
\eqref{eq:s7v64-K2}; indeed \(K_2\to0\).
Apply
\Cref{thm:s7v63-interacting,thm:s7v63-polymer} with \(R^2=M\).
\end{proof}

\section{The remaining circularity is global-chart completeness}
\label{sec:s7v64-firewall}

\begin{firewall}[Replica-admissible does not yet mean full compact-group
density]
\label{fire:s7v64-complete}
The V55--V57 interaction bounds are proved after localization to the
weak-field chart and occupation window.  The exact compact-link
Yang--Mills density also includes configurations outside that chart.
Using \Cref{cor:s7v64-exit} to justify discarding those configurations
would be circular if the density used to prove the replica bound had
already discarded them.

To close the physical chart-exit row, one must construct a simultaneous
large-field polymer expansion: good blocks use the replica-admissible
Weyl activity, bad blocks use the exact compact-group action and an atlas
or energy-barrier estimate, and mixed boundaries have a summable
transition activity.  V64 proves that normalization and Gaussian
replica gluing add no further obstruction once that local decomposition
is supplied.
\end{firewall}

\begin{remark}[Why the result is still upstream progress]
\label{rem:s7v64-progress}
The former normalized-denominator problem has been eliminated.  It is
now enough to prove local weights for the good/bad atlas decomposition.
No global partition-function lower bound and no all-block single-chart
event is required.  The unresolved input is geometric and local: a
compact-link large-field block card with transition compatibility.
\end{remark}

\section{Status}
\label{sec:s7v64-status}

\begin{theorem}[What V64 proves]
\label{thm:s7v64-result}
The two-replica Gaussian glued on \(S\) retains the original
single-replica covariance and a uniform conditional gap.  Its
cross-replica covariance factors through \(S\) and has total weighted
mass at most \(Q_\times|S|\).  Hence every mixed-replica connected tree
has an anchored cross edge, and the degree-weighted Weyl expansion proves
the pressure defect bound
\eqref{eq:s7v64-replica-tree}.  RPLC, LMDC, and a normalized
\(e^{-cM|S|}\) bad-set bound follow for every complete density belonging
to this replica-admissible activity class.
\end{theorem}

\begin{proof}
Combine
\Cref{prop:s7v64-diagonal,thm:s7v64-covariance,
thm:s7v64-cross-mass,lem:s7v64-marked,
thm:s7v64-replica-tree,cor:s7v64-RPLC,
cor:s7v64-exit}.
\end{proof}

\begin{assessment}[Progress at seventh-series V64]
\label{ass:s7v64-status}
The localized replica pressure card is no longer an independent
probabilistic hypothesis for the good-field theory; it follows from the
same Weyl/tree smallness after one cross edge is marked.  The physical
bottleneck has moved upstream to the exact compact geometry: construct
the simultaneous good-chart/bad-chart polymer decomposition without
assuming the bad sector away.  This is the appropriate point for the V65
companion audit.
\end{assessment}

\begin{decision}[V65 audit and acquisition order]
\label{dec:s7v64-V65}
V65 is the compulsory companion audit.  It will inspect the latest SM
and NS files for an atlas, coarea, large-field barrier, conditional
density, or exceptional-set polymer mechanism.  Any transfer will be
limited to an abstract local decomposition theorem.  The subsequent
attack will construct the compact-link good/bad block card, beginning
with a finite atlas and exact partition of unity rather than a global
logarithm.
\end{decision}

\begin{firewall}[Claim boundary after seventh-series V64]
\label{fire:s7v64-claim}
V64 does not prove the full compact-group good/bad polymer decomposition,
transition-chart activities, a large-field block barrier, physical LMDC
outside the replica-admissible chart sector, interacting endpoint
continuity, UCM through all slow-label patches, BCC, the
pairing/Bogoliubov analogue, weighted physical-source incidence,
patchwise coarse-action cocycles, terminal strong-coupling matching,
removal of all regulators, OS reconstruction, or the full Yang--Mills
mass gap.
\end{firewall}

\paragraph{Parent-version record.}
V64 was formed from
\texttt{Renewal\_Yang\_Mills\_Mass\_Gap\_Research\_Notes\_V63.tex}.
The V63 parent had \(27\,420\) logical source lines, \(978\,059\) bytes,
and SHA--256
\[
 \texttt{825489342970D8755A66349711D7BAF3B7CB52762C1087010D634EE2FCB78B40}.
\]
The next compulsory companion audit is V65.

% =============================================================================
% END APPENDED SEVENTH-SERIES V64 MATERIAL
% =============================================================================

% =============================================================================
% BEGIN APPENDED SEVENTH-SERIES V65 MATERIAL
% =============================================================================

\chapter{Companion audit and an exact compact-link atlas barrier}
\label{ch:s7v65}

\section{Compulsory five-version audit}
\label{sec:s7v65-audit}

\subsection{Files inspected}
\label{sec:s7v65-files}

The active companion files inspected were
\begin{center}
\small\ttfamily
Renewal\_SM\_Einstein\_Continuum\_Research\_Notes\_V22.tex,\\
Renewal\_NS\_Stability\_Research\_Notes\_V26.tex.
\end{center}
At the audit time, the SM file had \(7\,671\) logical source lines,
\(265\,160\) bytes, and SHA--256
\[
 \texttt{4097241C1D1F354B043E57CC7C6A653AA3448274A0F6CEC75CA3570444BAE948}.
\]
The NS file had \(5\,319\) logical source lines, \(278\,283\) bytes, and
SHA--256
\[
 \texttt{82C887F8C2C69E4F28FAD7C3DC8DE5B9099CB5A4BF431EBA1FFF3E91935978AD}.
\]

\subsection{Standard Model--Einstein content}
\label{sec:s7v65-SM}

Three SM results are structurally relevant.
\begin{enumerate}[label=\textup{(\roman*)},leftmargin=3.8em]
\item V9 treats a finite exceptional tube and a regular polar-chart
      complement separately; normalization is never inferred from an
      unnormalized numerator.
\item V18 proves that adding a chart or refinement label through a
      normalized conditional probability kernel leaves the original
      marginal exactly unchanged.  An unnormalized label kernel instead
      generates an effective action.
\item V21 proves the full-density Holder exit inequality and warns that a
      cutoff-dependent density norm can make a tail argument circular.
\end{enumerate}
The V63--V64 replica construction strengthens item (iii) in the direction
needed here: it seeks a marginal norm exponential in the observed set,
not a global norm exponential in the ambient volume.

\begin{proposition}[Type-correct SM transfer]
\label{prop:s7v65-SM}
The normalized-spectator theorem transfers directly to a finite
compact-link atlas: a partition-of-unity chart label may be added exactly
without changing the Yang--Mills measure.  The SM determinant--coarea
crossing powers and metric-support estimates do not transfer to the
Wilson barrier.
\end{proposition}

\begin{proof}
A partition of unity is a probability vector at each compact-link
configuration, so it is precisely a normalized conditional kernel.
The crossing and metric results concern different operators and
geometric variables and provide no inequality for the Wilson action.
\end{proof}

\subsection{Navier--Stokes content}
\label{sec:s7v65-NS}

NS V26 is unchanged.  Its rank-one Oseen derivative calculation sharpens
deterministic parabolic constants but contains no compact configuration
space, probability density, atlas, or large-field energy.

\begin{proposition}[NS non-transfer at V65]
\label{prop:s7v65-NS}
No current NS theorem loads the compact-link chart or bad-polymer rows.
\end{proposition}

\begin{proof}
There is no declared map from the ancient velocity-profile variables and
Oseen norms to compact gauge links, Haar measure, or the Wilson action.
\end{proof}

\begin{audit}[V65 companion decision]
\label{aud:s7v65}
The audit imports only the normalized-spectator principle.  The Yang--Mills
construction will:
\begin{enumerate}[label=\textup{(\roman*)},leftmargin=3.8em]
\item introduce a finite compact atlas by an exact partition of unity;
\item define bad blocks relative to the full compact action, not a
      pre-truncated Gaussian density; and
\item retain gauge, flat, toron, and boundary minima as a compact slow
      manifold before applying a transverse energy barrier.
\end{enumerate}
\end{audit}

\section{A finite atlas as an exact normalized spectator}
\label{sec:s7v65-atlas}

\begin{theorem}[Exact finite-atlas lift]
\label{thm:s7v65-atlas}
Let \(X\) be a compact smooth configuration manifold and let
\(\{U_\alpha\}_{\alpha\in{\cal A}}\) be a finite coordinate cover with a
smooth partition of unity
\begin{equation}
 \chi_\alpha\geq0,\qquad
 {\rm supp}\,\chi_\alpha\subset U_\alpha,\qquad
 \sum_{\alpha\in{\cal A}}\chi_\alpha(x)=1.             \label{eq:s7v65-POU}
\end{equation}
For every probability \(\mu\) on \(X\), define a probability on
\(X\times{\cal A}\) by
\begin{equation}
 \widetilde\mu(dx,\alpha)
 :=\mu(dx)\chi_\alpha(x).                              \label{eq:s7v65-lift}
\end{equation}
Then the \(X\)-marginal of \(\widetilde\mu\) is exactly \(\mu\), and for
every integrable \(F\),
\begin{equation}
 \int_XF\,d\mu
 =
 \sum_{\alpha\in{\cal A}}
 \int_{U_\alpha}F(x)\chi_\alpha(x)\,d\mu(x).           \label{eq:s7v65-atlas-identity}
\end{equation}
\end{theorem}

\begin{proof}
Summing \eqref{eq:s7v65-lift} over \(\alpha\) and using
\eqref{eq:s7v65-POU} gives \(\mu(dx)\).  Integrating \(F\) gives
\eqref{eq:s7v65-atlas-identity}.
\end{proof}

\begin{corollary}[Product-block atlas]
\label{cor:s7v65-product}
For a finite family of blocks \(b\in\Lambda\), the conditional label law
\begin{equation}
 {\mathbb P}(\alpha_\Lambda\mid x_\Lambda)
 =
 \prod_{b\in\Lambda}\chi_{\alpha_b}(x_b)               \label{eq:s7v65-label-law}
\end{equation}
is normalized for every \(x_\Lambda\).  Adding all chart labels therefore
changes neither the compact-link marginal nor its partition function.
\end{corollary}

\begin{proof}
The sum over all label configurations factorizes as
\[
 \sum_{\alpha_\Lambda}
 \prod_{b\in\Lambda}\chi_{\alpha_b}(x_b)
 =
 \prod_{b\in\Lambda}\sum_{\alpha_b}\chi_{\alpha_b}(x_b)
 =1.
\]
Apply \Cref{thm:s7v65-atlas}.
\end{proof}

\begin{firewall}[Overlaps are retained, not double counted]
\label{fire:s7v65-overlap}
The label lift does not choose one chart discontinuously.  A point in an
overlap carries several labels with total conditional mass one.
Transition Jacobians and slow-coordinate changes must be computed in the
labelled terms, but no overlap multiplicity enters the underlying
probability.
\end{firewall}

\section{Wilson action zero set and a compact quadratic barrier}
\label{sec:s7v65-Wilson}

Let \(G\) be a compact Lie group and
\(\rho:G\to U(d_\rho)\) a finite-dimensional unitary representation.
The one-plaquette Wilson function is
\begin{equation}
 s_\rho(U):=
 1-{1\over d_\rho}{\rm Re}\,\Tr\rho(U).                \label{eq:s7v65-Wilson}
\end{equation}

\begin{lemma}[Exact zero set]
\label{lem:s7v65-zero}
\(s_\rho\geq0\), and
\begin{equation}
 s_\rho(U)=0
 \quad\Longleftrightarrow\quad
 \rho(U)=I.                                            \label{eq:s7v65-zero}
\end{equation}
In particular, for a faithful representation the unique zero is the
identity.
\end{lemma}

\begin{proof}
Every eigenvalue of \(\rho(U)\) lies on the unit circle, so the real part
of each eigenvalue is at most one.  Their average is therefore at most
one.  Equality holds only if every eigenvalue has real part one, hence
every eigenvalue equals one.  A unitary matrix all of whose eigenvalues
are one is the identity.  Faithfulness identifies its inverse image with
the group identity.
\end{proof}

\begin{lemma}[Quadratic group-distance barrier]
\label{lem:s7v65-group-barrier}
If \(\rho\) is faithful, then for any fixed bi-invariant Riemannian
distance on \(G\) there are constants \(c,C>0\) such that
\begin{equation}
 c\,d_G(U,e)^2
 \leq s_\rho(U)
 \leq C\,d_G(U,e)^2
 \qquad(U\in G).                                       \label{eq:s7v65-group-barrier}
\end{equation}
\end{lemma}

\begin{proof}
In exponential coordinates \(U=e^X\), unitarity gives
\[
 s_\rho(e^X)
 =
 {1\over2d_\rho}\Tr\bigl(d\rho(X)^*d\rho(X)\bigr)
 +O(|X|^3).
\]
Faithfulness of the group representation makes \(d\rho\) injective on
the Lie algebra up to any discrete kernel, so the quadratic form is
positive definite.  Thus \eqref{eq:s7v65-group-barrier} holds in a
neighborhood of \(e\).  On the compact complement of a smaller
neighborhood, \(s_\rho\) has a positive minimum by
\Cref{lem:s7v65-zero}, while \(d_G(U,e)^2\) is bounded above and away
from zero.  Enlarging the constants joins the two regions.
\end{proof}

\begin{remark}[Block zero manifolds]
\label{rem:s7v65-flat}
For a gauge block, the sum of plaquette Wilson functions generally
vanishes on a compact manifold of flat connections rather than at one
coordinate point.  Gauge orbits, torons, and boundary holonomies belong
to that zero manifold and must be quotiented or retained as slow
variables.  The relevant statement is therefore a quadratic barrier in
the normal distance to a compact minimum manifold.
\end{remark}

\section{A compact Morse--Bott Gibbs tail}
\label{sec:s7v65-MB}

Let \(X\) be a compact Riemannian manifold with a probability reference
measure \(m\) having a smooth strictly positive density.  Let
\({\cal K}\subset X\) be a compact embedded submanifold of codimension
\(d_f\geq1\).  Let \(S:X\to[0,\infty)\) be continuous on \(X\), smooth
near \({\cal K}\), and suppose
\begin{equation}
 S^{-1}(0)={\cal K},\qquad
 c_0\,{\rm dist}(x,{\cal K})^2
 \leq S(x)
 \leq C_0\,{\rm dist}(x,{\cal K})^2                   \label{eq:s7v65-MB}
\end{equation}
in a fixed tubular neighborhood, with the lower inequality extended to
all \(X\) after reducing \(c_0>0\).

\begin{lemma}[Global extension of the lower barrier]
\label{lem:s7v65-global}
If the lower inequality in \eqref{eq:s7v65-MB} holds in a neighborhood
of \({\cal K}\) and \(S^{-1}(0)={\cal K}\), then it holds on all of \(X\)
after decreasing \(c_0\).
\end{lemma}

\begin{proof}
On the compact complement of a smaller tubular neighborhood,
\[
 {S(x)\over{\rm dist}(x,{\cal K})^2}
\]
is continuous and strictly positive.  Its minimum is positive.  Take
the smaller of that minimum and the local constant.
\end{proof}

For \(\beta\geq1\), define
\begin{equation}
 d\mu_\beta(x):=
 {e^{-\beta S(x)}\over Z_\beta}\,dm(x),\qquad
 Z_\beta:=\int_Xe^{-\beta S}\,dm.                      \label{eq:s7v65-Gibbs}
\end{equation}

\begin{theorem}[Compact Morse--Bott chart-exit bound]
\label{thm:s7v65-MB}
Under the preceding hypotheses, there are constants
\(C,c,r_0>0\) such that for
\begin{equation}
 1\leq R\leq r_0\sqrt\beta                              \label{eq:s7v65-Rrange}
\end{equation}
one has
\begin{equation}
 \mu_\beta\!\left\{
 {\rm dist}(x,{\cal K})\geq{R\over\sqrt\beta}
 \right\}
 \leq
 C(1+R)^{d_f}
 e^{-cR^2}.                                            \label{eq:s7v65-MB-tail}
\end{equation}
The constants are uniform for a compact family of slow labels whenever
the tubular radii, measure densities, and constants in
\eqref{eq:s7v65-MB} are uniform.
\end{theorem}

\begin{proof}
In a fixed tubular neighborhood write \(x=(z,v)\), with
\(z\in{\cal K}\) and \(v\) normal of dimension \(d_f\).  Compactness and
strict positivity of the density give uniform upper and lower bounds for
the tubular Jacobian.

For the denominator, integrate over \(|v|\leq\beta^{-1/2}\).  The upper
bound in \eqref{eq:s7v65-MB} gives
\[
 Z_\beta\geq c_1\beta^{-d_f/2}.
\]
For the numerator within the tube, the global lower bound and polar
coordinates give
\begin{align*}
 \int_{|v|\geq R/\sqrt\beta}
 e^{-\beta S}\,dm
 &\leq
 C_1\int_{R/\sqrt\beta}^{r_0}
 e^{-c_0\beta r^2}r^{d_f-1}\,dr\\
 &=
 C_1\beta^{-d_f/2}
 \int_R^{r_0\sqrt\beta}
 e^{-c_0t^2}t^{d_f-1}\,dt\\
 &\leq
 C_2\beta^{-d_f/2}(1+R)^{d_f}e^{-cR^2}.
\end{align*}
On the complement of the tube, compactness gives
\(S\geq s_0>0\), so that contribution is at most
\(C_3e^{-\beta s_0}\).  Under
\eqref{eq:s7v65-Rrange}, decrease \(r_0\) so it is bounded by the same
right side.  Divide by the denominator lower bound.
\end{proof}

\begin{corollary}[Perturbed local density]
\label{cor:s7v65-perturbed}
Let \(h_\beta\) be positive and suppose
\begin{equation}
 e^{-A}\leq h_\beta(x)\leq e^A                         \label{eq:s7v65-h}
\end{equation}
uniformly.  Replacing \(e^{-\beta S}dm\) by
\(h_\beta e^{-\beta S}dm\) multiplies the right side of
\eqref{eq:s7v65-MB-tail} by at most \(e^{2A}\).
\end{corollary}

\begin{proof}
The numerator is multiplied by at most \(e^A\), and the denominator by
at least \(e^{-A}\).
\end{proof}

\begin{corollary}[Yang--Mills logarithmic scaling]
\label{cor:s7v65-YM-tail}
Let \(\beta=g^{-2}\), use the rescaled transverse coordinate
\(y=g^{-1}v\), and set \(R^2=M=(\log L)^\chi\).  If
\(R\leq r_0/g\), then the exact compact-block Gibbs probability outside
the transverse weak-field tube obeys
\begin{equation}
 \mu_\beta\{|y|\geq R\}
 \leq
 C(1+\sqrt M)^{d_f}e^{-cM}.                            \label{eq:s7v65-YM-tail}
\end{equation}
This survives every fixed polynomial loss in \(M\).
\end{corollary}

\begin{proof}
The event \(|y|\geq R\) is
\({\rm dist}(x,{\cal K})\geq gR=R/\sqrt\beta\).
Apply \Cref{thm:s7v65-MB}; the range condition follows for sufficiently
small \(g\) because \(R\) is polylogarithmic in \(L\), whereas \(g^{-1}\)
grows like \(L^{1/2}\).  Exponential decay in \(M\) absorbs all fixed
powers.
\end{proof}

\section{Application to a gauge-reduced Wilson block}
\label{sec:s7v65-block}

\begin{definition}[Compact transverse block barrier, CTBB]
\label{def:s7v65-CTBB}
CTBB consists of:
\begin{enumerate}[label=\textup{(B\arabic*)},leftmargin=4em]
\item a compact gauge-reduced block configuration \(X\);
\item a compact slow minimum manifold \({\cal K}\) containing all flat,
      toron, boundary-holonomy, and retained transition directions;
\item a uniformly positive Hessian of the Wilson block action normal to
      \({\cal K}\);
\item a finite tubular atlas with uniform Haar/Faddeev--Popov density
      bounds and an exact normalized partition of unity;
\item uniformity over the finite buffered slow-label cover; and
\item bounded local multiplicative rows as in
      \eqref{eq:s7v65-h}, or a separately declared \(L^q\) replacement.
\end{enumerate}
\end{definition}

\begin{proposition}[CTBB gives a noncircular one-block bad weight]
\label{prop:s7v65-CTBB}
Under CTBB, the exact compact-link local Gibbs measure, before any
weak-field truncation, assigns the transverse bad event the weight
\begin{equation}
 p_M^{\rm block}
 \leq C M^{d_f/2}e^{-cM}.                              \label{eq:s7v65-block-weight}
\end{equation}
The atlas labels do not modify this probability.
\end{proposition}

\begin{proof}
The normal Hessian floor and compactness give
\eqref{eq:s7v65-MB} by the Morse--Bott lemma and
\Cref{lem:s7v65-global}.  Apply
\Cref{cor:s7v65-perturbed,cor:s7v65-YM-tail}.
The label marginal is exact by \Cref{thm:s7v65-atlas}.
\end{proof}

\begin{firewall}[CTBB has not yet been verified for the coupled block
family]
\label{fire:s7v65-CTBB}
Earlier transverse oscillator analysis proves a local Hessian floor on
declared patches, but it has not yet established every global item of
CTBB: Gribov/slice boundaries, all flat strata, Faddeev--Popov zeros, and
transition charts must be included in the compact slow atlas.  Moreover,
\Cref{prop:s7v65-CTBB} is a one-block Gibbs estimate.  Couplings across
block boundaries must be assembled without replacing them by a uniformly
bounded factor whose oscillation grows with the boundary field.
\end{firewall}

\section{What the audit changes}
\label{sec:s7v65-status}

\begin{theorem}[What V65 proves]
\label{thm:s7v65-result}
A finite compact atlas can be introduced as an exact normalized
spectator, so overlap multiplicity never changes the Yang--Mills
probability.  A compact nonnegative action with a uniformly nondegenerate
Morse--Bott minimum manifold has the exact normalized tail
\eqref{eq:s7v65-MB-tail}.  Under CTBB this gives the full compact-block
Wilson bad weight \(CM^{d_f/2}e^{-cM}\) before weak-field truncation.
Thus a local large-field weight can be obtained noncircularly; the
remaining problem is its coupled polymer assembly and the verification
of the global gauge-reduced atlas hypotheses.
\end{theorem}

\begin{proof}
Use
\Cref{thm:s7v65-atlas,lem:s7v65-group-barrier,
thm:s7v65-MB,cor:s7v65-YM-tail,
prop:s7v65-CTBB}.
\end{proof}

\begin{assessment}[Progress at seventh-series V65]
\label{ass:s7v65-status}
The companion audit prevents two errors: chart labels must be normalized
spectators, and an unnormalized compact numerator cannot be called a
probability.  The new compact Morse--Bott estimate resolves both issues
for one gauge-reduced block and reproduces the Gaussian
\(e^{-cM}\) scale with only a fixed polynomial prefactor.  V64 handles
normalization inside the good-field connected theory.  The next step is
to combine the exact compact bad weight and the good-field replica
activity across block boundaries.
\end{assessment}

\begin{decision}[V66 acquisition order]
\label{dec:s7v65-V66}
V66 will formulate a buffered block coloring so that interiors are
conditionally disjoint, prove a domination theorem for bad blocks in the
presence of finite-range boundary interactions, and identify the exact
boundary oscillation needed to retain the weight
\eqref{eq:s7v65-block-weight}.  If uniform conditional domination fails
because boundary fields can move the minimum manifold, those fields will
be enlarged into the slow label rather than estimated as perturbations.
\end{decision}

\begin{firewall}[Claim boundary after seventh-series V65]
\label{fire:s7v65-claim}
V65 does not verify CTBB for every physical gauge block, prove coupled
bad-polymer domination, transition-chart activities, physical LMDC
across good and bad sectors, interacting endpoint continuity, UCM
through all slow-label patches, BCC, the pairing/Bogoliubov analogue,
weighted physical-source incidence, patchwise coarse-action cocycles,
terminal strong-coupling matching, removal of all regulators, OS
reconstruction, or the full Yang--Mills mass gap.
\end{firewall}

\paragraph{Parent-version record.}
V65 was formed from
\texttt{Renewal\_Yang\_Mills\_Mass\_Gap\_Research\_Notes\_V64.tex}.
The V64 parent had \(27\,938\) logical source lines, \(997\,017\) bytes,
and SHA--256
\[
 \texttt{906E9F4F732FE3AB019B49A0C314A6A3C6820C823EF507703172936CD9AA7375}.
\]
The next compulsory companion audit is V70.

% =============================================================================
% END APPENDED SEVENTH-SERIES V65 MATERIAL
% =============================================================================

% =============================================================================
% BEGIN APPENDED SEVENTH-SERIES V66 MATERIAL
% =============================================================================

\chapter{Buffered conditional domination and coupled bad-polymer assembly}
\label{ch:s7v66}

\section{Purpose}
\label{sec:s7v66-purpose}

V65 proves an exact bad-field tail for one compact gauge-reduced block,
provided its transverse minimum manifold and Hessian are controlled
uniformly.  A lattice theory couples neighboring blocks, so a one-block
probability cannot simply be multiplied over an arbitrary set.  The
conditional energy depends on a boundary collar, and a neighboring bad
event may itself depend on the interior variable through that collar.

This note resolves the abstract assembly problem.  A finite coloring of
buffered interiors makes same-color conditional laws independent and
makes their bad events measurable outside the other same-color
interiors.  A uniform conditional compact-block barrier then yields
\(p_M^{|S|/K}\) for an arbitrary bad set.  Finally a marked-bad-site
animal sum combines this weight with the good-field connected activity.
The remaining physical input is stated exactly: a uniform conditional
transverse block barrier over every boundary slow label.

\section{Finite-range buffered coloring}
\label{sec:s7v66-color}

Let \({\cal B}\) be the graph of coarse blocks, of maximum degree
\(\Delta\).  Suppose every interaction meeting an interior block \(b\)
is contained in its collar \(b^+\), whose graph radius is at most
\(R_0\).  Define the conflict graph by joining \(b,b'\) whenever their
collars interact or either bad event can depend on the other's interior.
Let its maximum degree be \(\Delta_+\).

\begin{lemma}[Finite buffered coloring]
\label{lem:s7v66-color}
The conflict graph admits a proper coloring
\begin{equation}
 {\cal B}={\cal B}_1\sqcup\cdots\sqcup{\cal B}_K,
 \qquad K\leq\Delta_++1.                               \label{eq:s7v66-color}
\end{equation}
Within one color, the interiors have no common interaction term and no
bad event of one block depends on the interior variable of another.
\end{lemma}

\begin{proof}
The greedy coloring algorithm uses at most one more color than the
maximum degree.  Properness is exactly the declared absence of a
conflict edge, which has the two stated meanings by construction.
\end{proof}

\begin{remark}[The coloring cost is fixed]
\label{rem:s7v66-color}
For a finite-range action on a fixed-dimensional bounded-degree lattice,
\(\Delta_+\) and hence \(K\) depend only on the interaction range and
block geometry.  They do not grow with the volume, occupation cutoff, or
continuum scale.
\end{remark}

\section{Uniform conditional transverse barrier}
\label{sec:s7v66-conditional}

Fix a block \(b\) and condition on every variable outside its interior.
Denote the resulting boundary/collar label by \(\eta\).  Let
\({\cal K}_b(\eta)\) be the compact conditional minimum manifold, and
subtract the conditional minimum energy:
\begin{equation}
 \widehat S_b(x;\eta)
 :=
 S_b(x;\eta)-\min_zS_b(z;\eta)\geq0.                  \label{eq:s7v66-excess}
\end{equation}
The scalar minimum cancels from the normalized conditional law.

\begin{definition}[Uniform conditional transverse block barrier, UCTBB]
\label{def:s7v66-UCTBB}
UCTBB consists of the V65 CTBB data uniformly for every admissible
boundary label \(\eta\), applied to the excess action
\(\widehat S_b(\cdot;\eta)\), together with a residual logarithmic factor
\(u_b(x;\eta)\) satisfying
\begin{equation}
 \operatorname{osc}_{x}u_b(x;\eta)\leq D_M             \label{eq:s7v66-osc}
\end{equation}
uniformly in \(b,\eta\).  The conditional bad event is
\begin{equation}
 E_b(\eta)
 :=
 \left\{
 {\rm dist}(x,{\cal K}_b(\eta))
 \geq g\sqrt M
 \right\}.                                             \label{eq:s7v66-Eb}
\end{equation}
\end{definition}

\begin{lemma}[Oscillation comparison of normalized laws]
\label{lem:s7v66-osc}
Let \(\nu\) be a probability and
\[
 {d\mu\over d\nu}
 ={e^u\over\int e^u\,d\nu}.
\]
If \(\operatorname{osc}u\leq D\), then
\begin{equation}
 {d\mu\over d\nu}\leq e^D
 \quad \nu\hbox{-almost surely}.                       \label{eq:s7v66-density}
\end{equation}
\end{lemma}

\begin{proof}
Let \(m=\mathop{\rm ess\,inf}u\).  Then
\(\int e^u\,d\nu\geq e^m\), while \(e^u\leq e^{m+D}\).
Divide the two inequalities.
\end{proof}

\begin{theorem}[Uniform conditional bad-block probability]
\label{thm:s7v66-one-block}
Assume UCTBB and \(\beta=g^{-2}\).  There are constants
\(C,c>0\), independent of \(b,\eta,M\), such that
\begin{equation}
 \mu(E_b\mid{\cal F}_{b^c})
 \leq
 p_M:=
 CM^{d_f/2}\exp[-cM+D_M]                               \label{eq:s7v66-pM}
\end{equation}
almost surely.  If
\begin{equation}
 D_M\leq\delta M,\qquad 0\leq\delta<c,                 \label{eq:s7v66-margin}
\end{equation}
then \(p_M\leq CM^{d_f/2}e^{-(c-\delta)M}\to0\).
\end{theorem}

\begin{proof}
For fixed exterior data \(\eta\), let \(\nu_b^\eta\) be the normalized
conditional law with density proportional to
\(e^{-\beta\widehat S_b(x;\eta)}\).  UCTBB and
\Cref{thm:s7v65-MB} give
\[
 \nu_b^\eta(E_b(\eta))
 \leq CM^{d_f/2}e^{-cM}
\]
uniformly in \(\eta\).  The full conditional law is the normalized tilt
of \(\nu_b^\eta\) by \(e^{u_b}\).
\Cref{lem:s7v66-osc} multiplies the event probability by at most
\(e^{D_M}\), proving \eqref{eq:s7v66-pM}.  Insert
\eqref{eq:s7v66-margin} for the final statement.
\end{proof}

\begin{firewall}[The full conditional action must be used]
\label{fire:s7v66-conditional}
It is invalid to put only the interior plaquette action into
\(\widehat S_b\) and call every boundary interaction a bounded residual
without proving \eqref{eq:s7v66-osc}.  Boundary fields can shift the
minimum manifold or destroy a transverse floor.  Such fields must be
included in \(\eta\), and all terms depending substantially on the
interior variable must be included in the conditional excess action.
\end{firewall}

\section{Same-color factorization and arbitrary-set domination}
\label{sec:s7v66-domination}

\begin{lemma}[Conditional factorization on one color]
\label{lem:s7v66-factor}
Fix a color \(j\) and condition on all variables outside the interiors
of \({\cal B}_j\).  For every finite
\(S\subset{\cal B}_j\), the conditional law of the interiors in \(S\)
factorizes over \(b\in S\), and
\begin{equation}
 \mu\!\left(\bigcap_{b\in S}E_b\,
 \middle|\,{\cal F}_{{\cal B}_j^c}\right)
 \leq p_M^{|S|}.                                      \label{eq:s7v66-color-dom}
\end{equation}
\end{lemma}

\begin{proof}
By \Cref{lem:s7v66-color}, no interaction contains two same-color
interiors, so the conditional density factorizes.  Each bad event
\(E_b(\eta)\) depends only on its own interior and on conditioned collar
variables.  Apply \Cref{thm:s7v66-one-block} to every factor and
multiply.
\end{proof}

\begin{theorem}[Arbitrary bad-set domination]
\label{thm:s7v66-domination}
For every finite set \(S\subset{\cal B}\),
\begin{equation}
 \mu\!\left(\bigcap_{b\in S}E_b\right)
 \leq
 \overline p_M^{\,|S|},
 \qquad
 \overline p_M:=p_M^{1/K}.                             \label{eq:s7v66-domination}
\end{equation}
\end{theorem}

\begin{proof}
Among the \(K\) color classes, choose \(j\) with
\[
 |S\cap{\cal B}_j|\geq{|S|\over K}.
\]
The event that every block of \(S\) is bad is contained in the event that
every block of \(S\cap{\cal B}_j\) is bad.  Integrating
\eqref{eq:s7v66-color-dom} over the conditioned variables gives
\[
 \mu\!\left(\bigcap_{b\in S}E_b\right)
 \leq
 p_M^{|S\cap{\cal B}_j|}
 \leq p_M^{|S|/K}
 =\overline p_M^{|S|},
\]
because \(0<p_M<1\) for large \(M\).
\end{proof}

\begin{corollary}[Coupled bad blocks remain exponentially rare]
\label{cor:s7v66-rate}
Under \eqref{eq:s7v66-margin},
\begin{equation}
 \overline p_M
 \leq
 C^{1/K}M^{d_f/(2K)}
 \exp\!\left[-{c-\delta\over K}M\right].               \label{eq:s7v66-rate}
\end{equation}
Thus finite-range coupling and coloring reduce the exponent only by the
fixed factor \(K\).
\end{corollary}

\begin{proof}
Take the \(K\)-th root of \eqref{eq:s7v66-pM}.
\end{proof}

\section{Assembly with the good-field activity}
\label{sec:s7v66-assembly}

Let \(q_M\) be the rooted per-block majorant of the replica-admissible
good-field connected expansion, including all fixed polynomial
Weyl--Plancherel losses.  Let a mixed polymer support
\(\Gamma\) carry at least one bad block.  Suppose its absolute weight,
after internal labels are summed, obeys the product majorant
\begin{equation}
 |w(\Gamma,\mathfrak b)|
 \leq
 \overline p_M^{|\mathfrak b|}
 q_M^{|\Gamma\setminus\mathfrak b|},                  \label{eq:s7v66-mixed-weight}
\end{equation}
where \(\varnothing\ne\mathfrak b\subset\Gamma\) is the set of bad
blocks.

\begin{lemma}[Marked-bad color sum]
\label{lem:s7v66-marked}
For a fixed support \(\Gamma\) of size \(n\),
\begin{align}
 \sum_{\varnothing\ne\mathfrak b\subset\Gamma}
 \overline p_M^{|\mathfrak b|}
 q_M^{n-|\mathfrak b|}
 &=
 (q_M+\overline p_M)^n-q_M^n                         \label{eq:s7v66-binomial}\\
 &\leq
 n\overline p_M(q_M+\overline p_M)^{n-1}.             \label{eq:s7v66-marked}
\end{align}
\end{lemma}

\begin{proof}
The first identity is the binomial theorem with the empty bad set
removed.  For \(p,q\geq0\), the mean-value theorem applied to
\(t\mapsto(q+t)^n\) gives
\[
 (q+p)^n-q^n
 \leq np(q+p)^{n-1}.
\]
\end{proof}

\begin{theorem}[Mixed good/bad polymer summability]
\label{thm:s7v66-mixed}
Let the block graph have maximum degree \(\Delta\), and suppose
\eqref{eq:s7v66-mixed-weight} holds.  If
\begin{equation}
 \zeta_M:=
 \Delta^2e^\alpha(q_M+\overline p_M)<1,                \label{eq:s7v66-zeta}
\end{equation}
then
\begin{align}
 &\sup_x
 \sum_{\substack{\Gamma\ni x\\\Gamma\ {\rm connected}}}
 e^{\alpha|\Gamma|}
 \sum_{\varnothing\ne\mathfrak b\subset\Gamma}
 |w(\Gamma,\mathfrak b)|
 \leq
 {e^\alpha\overline p_M\over(1-\zeta_M)^2}.            \label{eq:s7v66-mixed}
\end{align}
\end{theorem}

\begin{proof}
Use the lattice-animal count
\eqref{eq:s7v63-animals} and
\Cref{lem:s7v66-marked}:
\begin{align*}
 \text{left side}
 &\leq
 \sum_{n\geq1}
 \Delta^{2(n-1)}e^{\alpha n}
 n\overline p_M(q_M+\overline p_M)^{n-1}\\
 &=
 e^\alpha\overline p_M
 \sum_{n\geq1}n\zeta_M^{n-1}
 =
 {e^\alpha\overline p_M\over(1-\zeta_M)^2}.
\end{align*}
\end{proof}

\begin{corollary}[Bad sectors are a vanishing correction to the good
expansion]
\label{cor:s7v66-vanishing}
Suppose \(q_M\to0\) by the V55--V57 logarithmic schedule and UCTBB holds
with \(\delta<c\).  Then \(\zeta_M\to0\), the mixed polymer expansion
converges, and its rooted total weight is
\begin{equation}
 O\!\left(
 M^{d_f/(2K)}
 e^{-(c-\delta)M/K}
 \right).                                              \label{eq:s7v66-vanishing}
\end{equation}
It survives every fixed polynomial occupation loss.
\end{corollary}

\begin{proof}
\Cref{cor:s7v66-rate} gives
\(\overline p_M\to0\), while \(q_M\to0\) is the V57 smallness condition
under V55.  Hence \eqref{eq:s7v66-zeta} holds for large \(M\).
Apply \Cref{thm:s7v66-mixed}; its denominator tends to one.
\end{proof}

\begin{firewall}[The product majorant is an interface, not automatic
factorization]
\label{fire:s7v66-product}
Equation \eqref{eq:s7v66-mixed-weight} requires the good-chart boundary
activities, atlas transition Jacobians, and exact compact bad-block
weights to be assigned once and estimated in one polymer representation.
The probabilistic domination
\eqref{eq:s7v66-domination} and the good-field tree bound do not by
themselves prove that their weights multiply without an additional
boundary factor.  That factor belongs in \(D_M\) or \(q_M\) and must
satisfy the same uniform margin.
\end{firewall}

\section{A precise closure card}
\label{sec:s7v66-card}

\begin{definition}[Buffered compact chart polymer card, BCCPC]
\label{def:s7v66-BCCPC}
BCCPC consists of:
\begin{enumerate}[label=\textup{(C\arabic*)},leftmargin=4em]
\item a fixed-range buffered block decomposition and coloring
      \eqref{eq:s7v66-color};
\item UCTBB for the full conditional excess action, uniformly over every
      boundary slow label;
\item the oscillation margin \(D_M\leq\delta M\) with \(\delta<c\);
\item the exact normalized atlas labels of V65;
\item the replica-admissible good activity \(q_M\to0\);
\item the mixed product majorant \eqref{eq:s7v66-mixed-weight}, including
      all transition-boundary rows; and
\item compatibility of the blockwise slow minima on collar overlaps.
\end{enumerate}
\end{definition}

\begin{theorem}[Conditional compact chart closure]
\label{thm:s7v66-closure}
BCCPC implies a volume-uniform, normalized, convergent good/bad chart
polymer expansion.  The sum of all polymers containing a bad block is
the vanishing rooted activity \eqref{eq:s7v66-vanishing}; the all-good
sector is controlled by the V64 replica-defect tree theorem.
\end{theorem}

\begin{proof}
UCTBB and the coloring give
\Cref{thm:s7v66-domination,cor:s7v66-rate}.
The product majorant and the smallness row give
\Cref{thm:s7v66-mixed,cor:s7v66-vanishing}.
The all-good sector is \Cref{thm:s7v64-replica-tree}; exact atlas
normalization is \Cref{thm:s7v65-atlas}.
\end{proof}

\section{Status}
\label{sec:s7v66-status}

\begin{theorem}[What V66 proves]
\label{thm:s7v66-result}
A uniform conditional compact-block barrier with logarithmic boundary
oscillation \(D_M< cM\) implies arbitrary bad-set domination
\[
 \mu(\cap_{b\in S}E_b)\leq p_M^{|S|/K}.
\]
The fixed coloring cost \(K\) is independent of volume and cutoff.
When the good and bad local weights satisfy the declared product
majorant, every mixed connected polymer is summable and the total
bad-sector activity obeys \eqref{eq:s7v66-vanishing}.  Thus the abstract
coupled chart-exit assembly is closed by BCCPC.
\end{theorem}

\begin{proof}
Use
\Cref{thm:s7v66-one-block,thm:s7v66-domination,
thm:s7v66-mixed,cor:s7v66-vanishing,
thm:s7v66-closure}.
\end{proof}

\begin{assessment}[Progress at seventh-series V66]
\label{ass:s7v66-status}
The former volume-normalization bottleneck has been reduced to two local
geometric loaders: UCTBB and the mixed boundary product majorant.  Once
they hold, bad compact-link charts are not discarded; they form a
convergent polymer family whose weight vanishes exponentially in
\(M=(\log L)^\chi\).  The next upstream question is whether the
gauge-reduced Wilson block has a uniform Morse--Bott normal Hessian for
all boundary labels after all flat and Gribov strata are retained.
\end{assessment}

\begin{decision}[V67 acquisition order]
\label{dec:s7v66-V67}
V67 will analyze the conditional Wilson Hessian geometrically.  It will
write the second variation as the covariant cochain Laplacian
\(d_A^*d_A\), identify its kernel under fixed boundary holonomy, and
separate regular flat strata from reducible or Gribov-singular strata.
The target is a finite stratified slow atlas with a uniform positive
normal eigenvalue on each buffered regular stratum; failure at a
singular stratum will be retained as a separate slow polymer label.
\end{decision}

\begin{firewall}[Claim boundary after seventh-series V66]
\label{fire:s7v66-claim}
V66 does not verify UCTBB, the mixed product majorant, or BCCPC for the
physical Wilson lattice family; construct the stratified gauge slice;
prove interacting endpoint continuity; load UCM through every slow
patch; prove BCC; handle the pairing/Bogoliubov sector, weighted physical
sources, patchwise coarse-action cocycles, or terminal strong-coupling
matching; remove all regulators; reconstruct the OS theory; or prove the
full Yang--Mills mass gap.
\end{firewall}

\paragraph{Parent-version record.}
V66 was formed from
\texttt{Renewal\_Yang\_Mills\_Mass\_Gap\_Research\_Notes\_V65.tex}.
The V65 parent had \(28\,427\) logical source lines, \(1\,015\,191\)
bytes, and SHA--256
\[
 \texttt{5B041AC8A42094B415D0C5DD1A9DA7DABCDF244D327D34DD3526A95D0F1F9C41}.
\]

% =============================================================================
% END APPENDED SEVENTH-SERIES V66 MATERIAL
% =============================================================================

% =============================================================================
% BEGIN APPENDED SEVENTH-SERIES V67 MATERIAL
% =============================================================================

\chapter{Conditional Wilson Hessians, cohomological kernels, and singular-label transfer}
\label{ch:s7v67}

\section{Purpose and correction of the acquisition target}
\label{sec:s7v67-purpose}

V66 asks for a conditional Wilson Hessian analysis.  Two results already in
the active notes must be kept distinct.  V1 proves a uniform Hessian floor
for the \emph{complete} incident action on the fully good collar
\(\mathcal G^\eta_{a_*,b_*}\).  V45 already proves the abstract finite
buffered spectral cover of a compact slow-label space.  Repeating either
argument would not advance the programme.

The missing point lies between them.  At a flat connection the Wilson
Hessian is exactly a covariant cochain square.  At a frustrated nonflat
conditional minimum it is not: a curvature-stress term remains.  Moreover,
a finite spectral cover can always retain eigenvalues approaching zero, but
that statement alone neither proves the Morse--Bott hypothesis nor makes
the retained directions rare.  This chapter proves those assertions
precisely and replaces the all-boundary UCTBB hypothesis by a
regular/singular-label domination theorem.

\section{Exact second variation of a plaquette}
\label{sec:s7v67-variation}

Equip \(\mathfrak{su}(3)\) with
\begin{equation}
 \langle X,Y\rangle:=-\operatorname{Tr}(XY),
 \qquad |X|^2=-\operatorname{Tr}(X^2).                 \label{eq:s7v67-metric}
\end{equation}
Let \(p\) be an oriented plaquette, and vary each positively oriented link
by a product-metric geodesic.  After replacing a negatively oriented link
by its inverse and transporting every tangent to the initial vertex of
\(p\), write
\begin{equation}
 Y_p(X):=\dot W_p(0)W_p(0)^{-1}.                       \label{eq:s7v67-Y}
\end{equation}
The product rule gives a second Lie-algebra expression \(Z_p(X)\) such
that
\begin{equation}
 \ddot W_p(0)W_p(0)^{-1}=Y_p(X)^2+Z_p(X),
 \qquad Z_p(X)\in\mathfrak{su}(3).                    \label{eq:s7v67-Z}
\end{equation}
Both \(Y_p\) and \(Z_p\) include the appropriate adjoint parallel
transports.

\begin{lemma}[Plaquette second variation]
\label{lem:s7v67-plaquette}
For \(\phi(W)=3-\operatorname{ReTr}W\),
\begin{equation}
 {d^2\over dt^2}\phi(W_p(t))\bigg|_{t=0}
 =
 -\operatorname{ReTr}
 \bigl((Y_p(X)^2+Z_p(X))W_p\bigr).                    \label{eq:s7v67-second}
\end{equation}
If \(W_p=I\), then
\begin{equation}
 {d^2\over dt^2}\phi(W_p(t))\bigg|_{t=0}
 =|Y_p(X)|^2.                                         \label{eq:s7v67-flat-second}
\end{equation}
For a fixed plaquette length there is a numerical constant \(C_p\) such
that
\begin{equation}
 \left|
 {d^2\over dt^2}\phi(W_p(t))\bigg|_{t=0}
 -|Y_p(X)|^2
 \right|
 \leq
 C_p\sqrt{\phi(W_p)}
 \sum_{e\in\partial p}|X_e|^2.                        \label{eq:s7v67-curvature-error}
\end{equation}
\end{lemma}

\begin{proof}
Differentiate \(\phi(W_p(t))\) twice and use
\eqref{eq:s7v67-Z}; this proves
\eqref{eq:s7v67-second}.  At \(W_p=I\),
\(\operatorname{Tr}Z_p(X)=0\), whereas
\(-\operatorname{Tr}Y_p(X)^2=|Y_p(X)|^2\), proving
\eqref{eq:s7v67-flat-second}.

For the last assertion, subtract the flat expression:
\[
 -\operatorname{ReTr}\bigl(Y_p^2(W_p-I)\bigr)
 -\operatorname{ReTr}\bigl(Z_p(W_p-I)\bigr).
\]
Adjoint transport is isometric for \eqref{eq:s7v67-metric}.  The product
rule on four factors therefore gives
\[
 |Y_p|^2+|Z_p|
 \leq C'_p\sum_{e\in\partial p}|X_e|^2.
\]
The exact unitary identity
\begin{equation}
 \|W-I\|_{\rm HS}^2
 =6-2\operatorname{ReTr}W
 =2\phi(W)                                             \label{eq:s7v67-defect-distance}
\end{equation}
and Cauchy--Schwarz for the trace give
\eqref{eq:s7v67-curvature-error}, after increasing the fixed constant.
\end{proof}

\begin{theorem}[Flat Wilson Hessian is a covariant cochain square]
\label{thm:s7v67-flat-Hessian}
Let \(K\) be a finite plaquette complex and let \(U\) be flat on every
plaquette of \(K\).  In left-trivialized link coordinates, let
\[
 d_U^1:C^1(K;\mathfrak{su}(3)_{\operatorname{Ad}U})
 \longrightarrow
 C^2(K;\mathfrak{su}(3)_{\operatorname{Ad}U})
\]
be the linearized plaquette-holonomy map.  Then
\begin{equation}
 \operatorname{Hess}_U
 \sum_{p\subset K}\phi(W_p)[X,X]
 =
 \|d_U^1X\|_2^2,                                      \label{eq:s7v67-cochain-square}
\end{equation}
and hence
\begin{equation}
 \operatorname{Hess}_U S_K=(d_U^1)^*d_U^1.            \label{eq:s7v67-laplacian}
\end{equation}
\end{theorem}

\begin{proof}
The \(p\)-component of \(d_U^1X\) is precisely
\(Y_p(X)\) in \eqref{eq:s7v67-Y}.  Apply
\eqref{eq:s7v67-flat-second} to every plaquette and sum.  Polarization
identifies the associated self-adjoint operator with
\((d_U^1)^*d_U^1\).
\end{proof}

\begin{firewall}[The cochain-square formula is not a formula at an
arbitrary conditional minimum]
\label{fire:s7v67-nonflat}
At a nonflat critical point, the sum of first variations vanishes, but
the individual \(Z_p(X)\)-terms in
\eqref{eq:s7v67-second} need not vanish.  Thus criticality does not turn
the Hessian into \((d_U^1)^*d_U^1\).  The latter is the leading weak-curvature
form, with the explicit error
\eqref{eq:s7v67-curvature-error}; using it as an exact identity at a
frustrated minimum would be false.
\end{firewall}

\begin{corollary}[Finite-complex weak-curvature comparison]
\label{cor:s7v67-weak}
For every fixed finite plaquette complex there is \(C_K<\infty\) such
that
\begin{equation}
 \left|
 \operatorname{Hess}_US_K[X,X]-\|d_U^1X\|_2^2
 \right|
 \leq
 C_K
 \max_{p\subset K}\sqrt{\phi(W_p)}
 \,\|X\|_2^2.                                         \label{eq:s7v67-weak}
\end{equation}
Consequently, if
\(\|d_U^1X\|_2^2\geq\lambda\|X\|_2^2\) on a declared
normal space and
\[
 C_K\max_p\sqrt{\phi(W_p)}\leq{\lambda\over2},
\]
then the Wilson Hessian on that space is at least \(\lambda/2\).
\end{corollary}

\begin{proof}
Sum \eqref{eq:s7v67-curvature-error}.  Each link occurs in only finitely
many plaquettes of the fixed complex, so the incidence multiplicity is
absorbed into \(C_K\).  The final assertion follows by subtraction.
\end{proof}

\section{The kernel is twisted cohomology, not a slice singularity}
\label{sec:s7v67-kernel}

Let \(d_U^0\) be the infinitesimal gauge coboundary.  If a boundary
subcomplex \(B\) is fixed, use relative cochains: link variations on
\(B\) and gauge parameters on its fixed vertices vanish.

\begin{theorem}[Flat kernel and physical infinitesimal moduli]
\label{thm:s7v67-cohomology}
At a flat connection,
\begin{equation}
 d_U^1d_U^0=0,\qquad
 \ker\operatorname{Hess}_US_K=\ker d_U^1.             \label{eq:s7v67-complex}
\end{equation}
After quotienting infinitesimal gauge transformations, the physical
flat Hessian kernel is canonically
\begin{equation}
 {\ker d_U^1\over\operatorname{im}d_U^0}
 \cong
 H^1(K,B;\mathfrak{su}(3)_{\operatorname{Ad}U}).       \label{eq:s7v67-H1}
\end{equation}
The infinitesimal stabilizer is
\(H^0(K,B;\mathfrak{su}(3)_{\operatorname{Ad}U})\).
\end{theorem}

\begin{proof}
Gauge transformations preserve every plaquette holonomy up to
conjugation.  Linearizing this statement at a flat connection gives
\(d_U^1d_U^0=0\).  The kernel identity follows from
\eqref{eq:s7v67-cochain-square}, since a squared norm vanishes exactly
when its argument does.  The quotient in
\eqref{eq:s7v67-H1} is the definition of first cohomology of the
flat adjoint local system; the degree-zero statement is the analogous
definition of its invariant sections.
\end{proof}

\begin{corollary}[Simply connected collar and the exact V1 floor]
\label{cor:s7v67-collar}
On a simply connected flat collar, parallel transport trivializes the
adjoint local system.  If the chosen full-tree gauge removes
\(\operatorname{im}d_U^0\), no physical flat kernel remains.  For the
specific V1 principal collar, the certified \(60\)-chord curvature matrix
therefore gives
\begin{equation}
 \operatorname{Hess}{\cal V}_+(U^{\rm fl})[X,X]
 \geq102^{-59}\|X\|^2,                                \label{eq:s7v67-v1-floor}
\end{equation}
and the complete fully good incident action has the already proved floor
\(\lambda_*=\frac14\,102^{-59}\).
\end{corollary}

\begin{proof}
A flat local system on a simply connected complex has path-independent
parallel transport and is gauge equivalent to the constant coefficient
system.  Its first cohomology vanishes when the collar has trivial first
cohomology relative to the fixed tree data.  The numerical statement is
exactly the unimodular-minor certificate and perturbation theorem
\Cref{thm:s7v1-complete-Hessian}.
\end{proof}

\begin{remark}[Torons, reducibles, and Gribov boundaries]
\label{rem:s7v67-three}
On a non-simply-connected block,
\eqref{eq:s7v67-H1} contains flat holonomy and toron directions.  At a
reducible holonomy, \(H^0\) and sometimes \(H^1\) can jump.  Those are
genuine slow directions and must be retained.  A zero of a particular
Faddeev--Popov determinant is instead a failure of that coordinate slice;
it does not, by itself, prove that the gauge-invariant quotient Hessian has
a new physical kernel.  These three phenomena cannot be merged into one
\emph{Gribov zero mode}.
\end{remark}

\section{What compact spectral buffering does and does not prove}
\label{sec:s7v67-buffer}

Let \(X_b\) be one compact finite-dimensional block carrier, with all
exact gauge directions either quotiented or explicitly retained, and let
\(\Xi_b\) be the compact boundary-label carrier.  For a continuous
conditional Wilson action \(S_b(x;\eta)\), put
\begin{equation}
 m_b(\eta):=\min_xS_b(x;\eta),\qquad
 {\mathfrak M}_b:=
 \{(x,\eta):S_b(x;\eta)=m_b(\eta)\}.                  \label{eq:s7v67-min-graph}
\end{equation}

\begin{lemma}[Compact graph of conditional minima]
\label{lem:s7v67-min-graph}
The function \(m_b\) is continuous and
\({\mathfrak M}_b\) is compact.
\end{lemma}

\begin{proof}
Uniform continuity of \(S_b\) on \(X_b\times\Xi_b\) gives
\[
 |m_b(\eta)-m_b(\eta')|
 \leq\sup_{x\in X_b}|S_b(x;\eta)-S_b(x;\eta')|,
\]
so \(m_b\) is continuous.  The equality defining
\({\mathfrak M}_b\) is therefore closed, and the ambient product is
compact.
\end{proof}

\begin{theorem}[Buffered conditional-minimum atlas]
\label{thm:s7v67-min-atlas}
Let \(H_{x,\eta}\geq0\) be the gauge-reduced Hessian on
\({\mathfrak M}_b\), expressed in any one of a finite family of
orthonormal tangent trivializations.  For every
\(z_0=(x_0,\eta_0)\in{\mathfrak M}_b\), choose a closed interval
\([a_{z_0},b_{z_0}]\subset(0,\infty)\) disjoint from
\(\operatorname{spec}H_{z_0}\).  There is a neighborhood \(U_{z_0}\)
on which:
\begin{enumerate}[label=\textup{(\roman*)},leftmargin=3.5em]
\item the interval remains in the resolvent set;
\item the spectral subspace below \(a_{z_0}\) has constant rank; and
\item the complementary Hessian has floor \(b_{z_0}\).
\end{enumerate}
A finite collection of these neighborhoods covers
\({\mathfrak M}_b\).  On an overlap, enlarging the slow space to the span
of the two patchwise slow spaces preserves each available complementary
floor on the smaller common complement.
\end{theorem}

\begin{proof}
In a tangent trivialization the Hessian is a continuous finite Hermitian
matrix.  The ordered eigenvalues are operator-norm Lipschitz, and the
resolvent identity preserves a spectral interval separated from the
spectrum.  The corresponding Riesz projection is continuous, hence has
locally constant integer rank.  Its complement has spectrum above the
upper edge \(b_{z_0}\).  Compactness from
\Cref{lem:s7v67-min-graph} gives a finite subcover; changes of
orthonormal trivialization are unitary and do not affect the spectrum.
If \(P_i,P_j\) are slow projections on an overlap, the orthogonal
complement of \(\operatorname{Ran}P_i+\operatorname{Ran}P_j\) is
contained in both individual complements, so the corresponding quadratic
form floors remain valid there.
\end{proof}

\begin{firewall}[A buffered atlas is not UCTBB]
\label{fire:s7v67-buffer}
\Cref{thm:s7v67-min-atlas} is finite-dimensional spectral bookkeeping.
At a point where every eigenvalue is small it may retain the entire tangent
space, leaving a zero-dimensional residual complement.  It also does not
show that the conditional minimum set is a manifold or that
\[
 \ker H_{x,\eta}=T_x{\cal K}_b(\eta).
\]
The elementary function \(t\mapsto t^4\) has a point minimum and zero
Hessian; retaining its Hessian kernel does not manufacture a quadratic
Morse--Bott tail.  Thus V45 buffering handles rank changes after a normal
floor is identified, but cannot prove that floor or the UCTBB probability
estimate.
\end{firewall}

\begin{proposition}[Exact scope of the V1 regular certificate]
\label{prop:s7v67-V1-scope}
Let \(x\) be a conditional minimizer of the \(60\)-chord incident action.
If
\begin{equation}
 {\cal V}_+(x)<a_*,
 \qquad
 \max_{p\in\partial P}\phi(W_p(x,\eta))<b_*,           \label{eq:s7v67-V1-certificate}
\end{equation}
then \(x\) is an isolated nondegenerate minimum in the chord fibre and
its Hessian is at least \(\lambda_*\).  Conversely, failure to certify a
conditional minimizer by \eqref{eq:s7v67-V1-certificate} has the explicit
witness
\begin{equation}
 {\cal V}_+(x)\geq a_*
 \quad\hbox{or}\quad
 \max_{p\in\partial P}\phi(W_p(x,\eta))\geq b_* .       \label{eq:s7v67-witness}
\end{equation}
Every crossing witness in the second alternative is assigned by V2 to
the principal defect event of its unique owner.
\end{proposition}

\begin{proof}
The first assertion is
\Cref{thm:s7v1-complete-Hessian}; a positive Hessian at a minimum makes
the minimum isolated in that fibre.  The displayed alternative is the
logical complement of \eqref{eq:s7v67-V1-certificate}.  Ownership is
\Cref{lem:s7v2-owner-principal}.
\end{proof}

\begin{firewall}[Conditional subtraction erases an absolute witness]
\label{fire:s7v67-subtraction}
The witness \eqref{eq:s7v67-witness} is an absolute Wilson defect.  In the
normalized conditional law, the scalar
\(\min_xS_b(x;\eta)\) cancels.  If the boundary forces every conditional
minimum to carry such a defect, that defect is not automatically a
conditional suppression factor.  Its cost must be charged in the
\emph{global} measure to an owner or an outward shell.  This is exactly
the distributional transfer row left open by V1--V5.
\end{firewall}

\section{Regular/singular conditional domination}
\label{sec:s7v67-domination}

The preceding analysis suggests a weaker hypothesis than UCTBB.  A
boundary label is called regular if every retained conditional well has
the Morse--Bott normal floor and uniform density data required by V65.
Let \(\Sigma_b\) be the complementary singular-label event.  It depends
only on the exterior/collar label and is therefore measurable outside the
interior of \(b\).

\begin{definition}[Regular/singular conditional barrier, RSCB]
\label{def:s7v67-RSCB}
For the buffered coloring of V66, RSCB consists of constants
\(p_M,s_M\in[0,1]\) such that:
\begin{enumerate}[label=\textup{(R\arabic*)},leftmargin=4em]
\item on \(\Sigma_b^c\),
\begin{equation}
 \mu(E_b\mid{\cal F}_{b^c})\leq p_M                  \label{eq:s7v67-regular-bound}
\end{equation}
almost surely, with
\(p_M=CM^{d_f/2}e^{-cM+D_M}\);
\item for every finite \(T\subset{\cal B}\),
\begin{equation}
 \mu\left(\bigcap_{b\in T}\Sigma_b\right)
 \leq s_M^{|T|};                                      \label{eq:s7v67-singular-dom}
\end{equation}
\item exact flat toron, boundary-holonomy, stabilizer, and buffered
transition modes have been retained as slow labels before
\(\Sigma_b\) is defined.
\end{enumerate}
\end{definition}

\begin{theorem}[Regular/singular replacement for uniform UCTBB]
\label{thm:s7v67-RSCB}
Assume the V66 conflict coloring and RSCB.  For every finite
\(S\subset{\cal B}\),
\begin{equation}
 \mu\left(\bigcap_{b\in S}E_b\right)
 \leq
 (p_M+s_M)^{|S|/K}.                                   \label{eq:s7v67-RSCB-dom}
\end{equation}
In particular, uniform control of every boundary label is unnecessary:
it is enough that the unretained singular labels themselves have
product domination.
\end{theorem}

\begin{proof}
First let \(S\) lie in one color.  Condition on the complement of the
same-color interiors.  The interiors factor as in
\Cref{lem:s7v66-factor}, and each \(\Sigma_b\) is measurable for that
conditioning.  On a regular label use
\eqref{eq:s7v67-regular-bound}; on a singular label use the trivial bound
one.  Therefore
\[
 \mu\left(\bigcap_{b\in S}E_b\right)
 \leq
 {\mathbb E}\prod_{b\in S}
 \bigl(p_M+{\bf1}_{\Sigma_b}\bigr).
\]
Expand the product and apply
\eqref{eq:s7v67-singular-dom}:
\[
 \sum_{T\subset S}
 p_M^{|S|-|T|}
 \mu\left(\bigcap_{b\in T}\Sigma_b\right)
 \leq
 \sum_{T\subset S}p_M^{|S|-|T|}s_M^{|T|}
 =(p_M+s_M)^{|S|}.
\]
For arbitrary \(S\), select a color containing at least \(|S|/K\)
members and discard the other bad-event requirements, exactly as in
\Cref{thm:s7v66-domination}.  Since \(p_M+s_M\leq1\) in the useful
regime, this proves \eqref{eq:s7v67-RSCB-dom}.  If the sum exceeds one,
the same statement remains true after replacing the right side by its
minimum with one.
\end{proof}

\begin{corollary}[Mixed-polymer assembly with singular labels]
\label{cor:s7v67-mixed}
Every conclusion of
\Cref{thm:s7v66-mixed,cor:s7v66-vanishing} remains valid after replacing
\(\overline p_M\) by
\begin{equation}
 \widetilde p_M:=(p_M+s_M)^{1/K},                     \label{eq:s7v67-ptilde}
\end{equation}
provided the mixed product majorant is proved with this replacement.
If \(p_M+s_M\to0\), the bad-or-singular rooted activity vanishes.
\end{corollary}

\begin{proof}
The proof of the marked-bad binomial sum uses only an arbitrary-set
product domination parameter.  Supply
\eqref{eq:s7v67-RSCB-dom} in place of
\eqref{eq:s7v66-domination}; every subsequent animal sum is unchanged.
\end{proof}

\begin{definition}[Singular-boundary transfer card, SBTC]
\label{def:s7v67-SBTC}
SBTC is the physical loader for
\eqref{eq:s7v67-singular-dom}.  It requires:
\begin{enumerate}[label=\textup{(S\arabic*)},leftmargin=4em]
\item cohomological flat modes from
\eqref{eq:s7v67-H1} and all buffered eigenvalue crossings are retained,
not declared rare;
\item loss of the V1 regular certificate is converted by
\eqref{eq:s7v67-witness} into a principal or uniquely owned crossing
defect;
\item the absolute defect is charged before conditional-minimum
subtraction, in a normalized global or shell measure;
\item the randomized plaquette transfer and martingale leakage rows of
V4--V5 produce a strict owner/shell contraction; and
\item iterated witnesses yield the multi-event domination
\eqref{eq:s7v67-singular-dom}, not merely a one-point probability.
\end{enumerate}
\end{definition}

\begin{theorem}[Geometric reduction of the conditional barrier]
\label{thm:s7v67-reduction}
For the Wilson collar, the conditional geometric problem splits into
three disjoint tasks:
\begin{enumerate}[label=\textup{(\roman*)},leftmargin=3.5em]
\item exact flat gauge and topology kernels are classified by
\eqref{eq:s7v67-H1} and loaded into the slow space;
\item weak-curvature regular minima are certified by
\Cref{cor:s7v67-weak,prop:s7v67-V1-scope};
\item all remaining unretained labels are handled by SBTC.
\end{enumerate}
If the regular labels satisfy the V65 local density bounds and SBTC gives
\(s_M\to0\), then RSCB replaces the uniform-all-\(\eta\) row of BCCPC.
\end{theorem}

\begin{proof}
The first two items are the preceding cohomology and Hessian theorems.
The logical complement of the regular certificate has the defect witness
\eqref{eq:s7v67-witness}; SBTC specifies the additional probabilistic
steps needed to turn those witnesses into
\eqref{eq:s7v67-singular-dom}.  The final conclusion is
\Cref{thm:s7v67-RSCB,cor:s7v67-mixed}.
\end{proof}

\section{Status and next acquisition}
\label{sec:s7v67-status}

\begin{theorem}[What V67 proves]
\label{thm:s7v67-result}
V67 proves the exact flat Wilson identity
\(\operatorname{Hess}S=(d_U^1)^*d_U^1\), its twisted-cohomology kernel, and
the explicit weak-curvature error at a nonflat configuration.  It shows
that the V1 collar already supplies the complete good-sector floor, while
V45 buffering supplies only a residual spectral atlas and not a
Morse--Bott theorem.  Finally it proves the regular/singular domination
\eqref{eq:s7v67-RSCB-dom}, which replaces V66's uniform boundary
hypothesis by a regular conditional tail plus product domination of the
unretained singular labels.
\end{theorem}

\begin{proof}
Use
\Cref{thm:s7v67-flat-Hessian,thm:s7v67-cohomology,
cor:s7v67-weak,prop:s7v67-V1-scope,
thm:s7v67-min-atlas,thm:s7v67-RSCB}.
\end{proof}

\begin{assessment}[Progress at seventh-series V67]
\label{ass:s7v67-status}
The attempted all-boundary Hessian route has been narrowed rather than
assumed.  Flat reducible and toron directions are slow; Gribov slice
failures are chart transitions; weak-curvature minima have the certified
V1 floor; and only the remaining frustrated labels require probability
transfer.  The live bottleneck is now the multi-event estimate
\eqref{eq:s7v67-singular-dom}.  A one-point unconditional bad probability
is still insufficient.
\end{assessment}

\begin{decision}[V68 acquisition order]
\label{dec:s7v67-V68}
V68 will attack SBTC upstream.  It will formulate the absolute global
Wilson-defect measure before conditional subtraction, use the V2 unique
owner and V4 randomized plaquette transfer to derive a multi-source
generating-function inequality, and test whether the V5 martingale
leakage debit leaves a strict radius of convergence.  If the transfer
still permits witness cycles with no paid root, the update blocks will be
ordered by outward shells so every singular chain terminates at a
principal paid defect.
\end{decision}

\begin{firewall}[Claim boundary after seventh-series V67]
\label{fire:s7v67-claim}
V67 does not prove SBTC, singular-label product domination, the physical
mixed product majorant, interacting endpoint continuity, UCM through all
slow patches, BCC, the pairing/Bogoliubov sector, weighted physical-source
incidence, patchwise coarse-action cocycles, terminal strong-coupling
matching, removal of all regulators, Osterwalder--Schrader reconstruction,
or the full Yang--Mills mass gap.
\end{firewall}

\paragraph{Parent-version record.}
V67 was formed from
\texttt{Renewal\_Yang\_Mills\_Mass\_Gap\_Research\_Notes\_V66.tex}.
The V66 parent had \(28\,877\) logical source lines, \(1\,031\,339\)
bytes, and SHA--256
\[
 \texttt{8DE9AF7186FFE4B293BD024C491D1EF72C5344445A90E644C03E215AA9C94DE9}.
\]
The next compulsory companion audit is V70.

% =============================================================================
% END APPENDED SEVENTH-SERIES V67 MATERIAL
% =============================================================================

% =============================================================================
% BEGIN APPENDED SEVENTH-SERIES V68 MATERIAL
% =============================================================================

\chapter{Multi-source owner domination and the absolute-energy repair criterion}
\label{ch:s7v68}

\section{Purpose}
\label{sec:s7v68-purpose}

V67 reduces the conditional chart problem to product domination of
unretained singular boundary labels.  The V2 owner map is deterministic,
but its earlier use was only one-source.  This chapter proves the exact
multi-source compiler.  A source has at most \(179\) possible principal
or crossing witnesses, while a fixed owner can receive at most \(239\)
of those source occurrences.  Thus product domination of principal
defect events with parameter \(r\) implies product domination of
\emph{actual-configuration} certificate failures with parameter
\(179r^{1/239}\).

The qualifier is essential.  A boundary-singular label is witnessed at a
conditional minimizer, not necessarily at the sampled configuration.
The chapter proves an absolute-energy repair lemma which can bridge that
gap, and shows why the Dirichlet-form transfers of V4--V5 do not by
themselves provide the needed probability estimate.

\section{An abstract bounded-incidence owner compiler}
\label{sec:s7v68-owner}

Let \({\cal B}\) be a finite source family and \({\cal O}\) an owner
family.  For each \(b\in{\cal B}\), let \({\cal W}(b)\) be a finite set
of witness occurrences and let
\[
 o:{\cal W}(b)\longrightarrow{\cal O}
\]
be the owner map.  Write \(\Sigma_b\) for a source event and \(A_o\) for
an owner event.

\begin{theorem}[Bounded-incidence product-domination compiler]
\label{thm:s7v68-owner}
Assume
\begin{align}
 \Sigma_b
 &\subset
 \bigcup_{w\in{\cal W}(b)} A_{o(w)},                  \label{eq:s7v68-cover}\\
 |{\cal W}(b)|
 &\leq m,                                             \label{eq:s7v68-row}\\
 \#\{(b,w):b\in{\cal B},\ w\in{\cal W}(b),\
             o(w)=o\}
 &\leq\ell                                           \label{eq:s7v68-column}
\end{align}
for every source and owner.  Suppose that, for every finite
\(T\subset{\cal O}\),
\begin{equation}
 \mu\left(\bigcap_{o\in T}A_o\right)\leq r^{|T|},
 \qquad 0\leq r\leq1.                                 \label{eq:s7v68-owner-dom}
\end{equation}
Then, for every finite \(S\subset{\cal B}\),
\begin{equation}
 \mu\left(\bigcap_{b\in S}\Sigma_b\right)
 \leq
 \left(mr^{1/\ell}\right)^{|S|}.                      \label{eq:s7v68-source-dom}
\end{equation}
The right side may be replaced by its minimum with one.
\end{theorem}

\begin{proof}
Use \eqref{eq:s7v68-cover} for every \(b\in S\) and distribute the
intersections over the unions.  This gives at most \(m^{|S|}\) witness
assignments
\[
 f:S\longrightarrow\bigcup_{b\in S}{\cal W}(b),
 \qquad f(b)\in{\cal W}(b).
\]
For one assignment let
\[
 T_f:=\{o(f(b)):b\in S\}.
\]
By \eqref{eq:s7v68-column}, one owner can be the image of at most
\(\ell\) selected source occurrences.  Hence
\[
 |T_f|\geq{|S|\over\ell}.
\]
The event belonging to \(f\) is contained in
\(\cap_{o\in T_f}A_o\), whose probability is at most
\[
 r^{|T_f|}\leq r^{|S|/\ell}
\]
because \(r\leq1\).  Sum over the at most \(m^{|S|}\) assignments.
\end{proof}

\begin{corollary}[Generating-function form]
\label{cor:s7v68-generating}
Under the assumptions of
\Cref{thm:s7v68-owner}, put \(s=mr^{1/\ell}\).  For every finite
\(S\subset{\cal B}\) and \(t\geq0\),
\begin{align}
 {\mathbb E}_\mu
 \exp\left(t\sum_{b\in S}{\bf1}_{\Sigma_b}\right)
 &\leq
 \bigl(1+s(e^t-1)\bigr)^{|S|}.                        \label{eq:s7v68-mgf}
\end{align}
\end{corollary}

\begin{proof}
Use
\[
 e^{t{\bf1}_{\Sigma_b}}
 =1+(e^t-1){\bf1}_{\Sigma_b},
\]
expand the product over \(b\in S\), and apply
\eqref{eq:s7v68-source-dom} to every selected subset.  The binomial
theorem gives \eqref{eq:s7v68-mgf}.
\end{proof}

\begin{remark}[Collisions cost an exponent, not a volume factor]
\label{rem:s7v68-collisions}
The column loss \(\ell\) is unavoidable for this information alone:
\(\ell\) distinct sources may all be witnessed by one owner event.  The
compiler therefore takes an \(\ell\)-th root of \(r\), but introduces no
factor depending on the lattice volume.
\end{remark}

\section{Exact Wilson-collar constants}
\label{sec:s7v68-Wilson}

For a translated chord block \(b\), define the actual-configuration
failure event
\begin{equation}
 \widehat\Sigma_b
 :=
 \left\{
 {\cal V}_{+,b}\geq a_*
 \ \hbox{or}\
 \max_{p\in\partial P_b}\phi_p\geq b_*
 \right\}.                                            \label{eq:s7v68-Sigmahat}
\end{equation}
Set \(c_*:=\min(a_*,b_*)\) and
\begin{equation}
 A_o^*:=
 \left\{
 \sum_{q\in P_o^+}\phi_q\geq c_*
 \right\}.                                            \label{eq:s7v68-Astar}
\end{equation}

\begin{lemma}[The \(179\)-by-\(239\) witness incidence]
\label{lem:s7v68-179-239}
The events \(\widehat\Sigma_b\) and \(A_o^*\) satisfy
\eqref{eq:s7v68-cover}--\eqref{eq:s7v68-column} with
\begin{equation}
 m=179,\qquad \ell=239.                               \label{eq:s7v68-constants}
\end{equation}
\end{lemma}

\begin{proof}
The principal alternative in \eqref{eq:s7v68-Sigmahat} is one witness,
owned by \(b\) itself.  Each of the \(178\) crossing plaquettes is one
additional witness.  If the principal alternative holds, then
\({\cal V}_{+,b}\geq a_*\geq c_*\), so \(A_b^*\) holds.  If a crossing
plaquette satisfies \(\phi_p\geq b_*\), the V2 owner rule places \(p\)
in \(P_{o(p)}^+\), and hence
\[
 \sum_{q\in P_{o(p)}^+}\phi_q
 \geq b_*\geq c_*.
\]
This proves the cover and the row count \(1+178=179\).

V2 proves that a fixed owner receives at most \(238\) crossing
occurrences.  Adding its one self-owned principal occurrence gives the
column bound \(239\).
\end{proof}

\begin{theorem}[Principal defects compile to actual certificate failures]
\label{thm:s7v68-Wilson-compiler}
If the principal defect family obeys
\begin{equation}
 \mu\left(\bigcap_{o\in T}A_o^*\right)
 \leq r_M^{|T|}                                      \label{eq:s7v68-principal-dom}
\end{equation}
for every finite owner set \(T\), then
\begin{equation}
 \mu\left(\bigcap_{b\in S}\widehat\Sigma_b\right)
 \leq
 \widehat s_M^{|S|},
 \qquad
 \widehat s_M:=179\,r_M^{1/239}.                      \label{eq:s7v68-Wilson-dom}
\end{equation}
In particular, if
\[
 r_M\leq C M^q e^{-cM},
\]
then
\begin{equation}
 \widehat s_M
 \leq
 179\,C^{1/239}M^{q/239}e^{-cM/239}\longrightarrow0.  \label{eq:s7v68-rate}
\end{equation}
\end{theorem}

\begin{proof}
Apply
\Cref{thm:s7v68-owner,lem:s7v68-179-239}.  The displayed rate is obtained
by taking the \(239\)-th root.
\end{proof}

\begin{assessment}[What the exact constants accomplish]
\label{ass:s7v68-constants}
The large incidence loss is not fatal at the logarithmic scale: an
\(e^{-cM}\) principal product weight remains exponentially small after
the \(239\)-th root.  Thus there is no need to demand the much stronger
one-step \(L^2\) contraction \(q<42364^{-1/2}\) merely to compile
multi-event probabilities.  The missing input is now
\eqref{eq:s7v68-principal-dom}.
\end{assessment}

\section{Actual failures are not boundary-singular labels}
\label{sec:s7v68-mismatch}

Let \(\Sigma_b^{\rm bdry}\) denote the V67 event that the fixed boundary
label \(\eta_b\) has an unretained conditional minimum which fails the
regular certificate.  For such a label, choose a conditional minimizer
\(x_b^*(\eta_b)\).  The V67 witness theorem applies to the completed
configuration
\[
 (x_b^*(\eta_b),\eta_b),
\]
whereas \(\widehat\Sigma_b\) in
\eqref{eq:s7v68-Sigmahat} is evaluated at the sampled interior variable
\(x_b\).

\begin{proposition}[No event inclusion without a completion comparison]
\label{prop:s7v68-no-inclusion}
The implication
\begin{equation}
 \Sigma_b^{\rm bdry}\Longrightarrow\widehat\Sigma_b   \label{eq:s7v68-false}
\end{equation}
does not follow from the definitions.  It can fail whenever the sampled
interior \(x_b\) and the selected conditional minimizer
\(x_b^*(\eta_b)\) lie on opposite sides of the certificate threshold.
Consequently
\Cref{thm:s7v68-Wilson-compiler} cannot be inserted directly as the
\(s_M\)-row of V67 RSCB.
\end{proposition}

\begin{proof}
\(\Sigma_b^{\rm bdry}\) is a statement about the value of the Wilson
defects at \(x_b^*(\eta_b)\); it imposes no pointwise equality between
that completion and the sampled \(x_b\).  Since the defect functions are
nonconstant continuous functions of \(x_b\), a sampled point may satisfy
the regular inequalities even when the selected minimizer does not, or
conversely.  Therefore the proposed implication is not a logical
consequence of either event.  An energy or measure comparison between
the two completions is required.
\end{proof}

\begin{firewall}[The hybrid-minimum problem]
\label{fire:s7v68-hybrid}
For several source blocks, independently selected conditional minimizers
need not form one global configuration when their collars overlap.
Even for one color, where interiors are conditionally disjoint, a later
repair of an owner boundary variable can invalidate the earlier
minimizer selections.  Thus a multi-source proof needs either a
simultaneous repair map, a normalized repair kernel, or a
reflection-positive estimate.  A list of incompatible virtual
completions cannot be charged as though it were an event of the original
field.
\end{firewall}

\section{An absolute-energy repair theorem}
\label{sec:s7v68-repair}

Let \((\Omega,m)\) be a compact reference probability space and
\begin{equation}
 d\mu_\beta(U)
 ={e^{-\beta S(U)}\over Z_\beta}\,dm(U),\qquad
 S\geq0.                                               \label{eq:s7v68-Gibbs}
\end{equation}
The following lemma is formulated so that a repair may be many-to-one.

\begin{definition}[Repair branch]
\label{def:s7v68-repair}
For a measurable cell \(C\subset\Omega\), a repair branch consists of a
measurable map \(F:C\to\Omega\), an integer \(n\geq1\), and constants
\(\delta,h>0\) such that
\begin{align}
 S(FU)&\leq S(U)-\delta n,                             \label{eq:s7v68-energy-drop}\\
 \int_C g(FU)\,dm(U)
 &\leq e^{hn}\int_\Omega g(V)\,dm(V)                  \label{eq:s7v68-push}
\end{align}
for every nonnegative measurable \(g\).
\end{definition}

\begin{theorem}[Multi-event repair compiler]
\label{thm:s7v68-repair}
Let \(A_o\), \(o\in{\cal O}\), be events.  Suppose that for every finite
\(T\subset{\cal O}\), the intersection
\(A_T:=\cap_{o\in T}A_o\) is covered by at most \(a_0^{|T|}\)
measurable cells.  On every cell there is a repair branch with
\begin{equation}
 n\geq{|T|\over\kappa}                                \label{eq:s7v68-n}
\end{equation}
and the same \(\delta,h,\kappa\).  If
\(\beta\delta>h\), then
\begin{equation}
 \mu_\beta(A_T)
 \leq r_\beta^{|T|},
 \qquad
 r_\beta:=
 a_0\exp\left[-{\beta\delta-h\over\kappa}\right].      \label{eq:s7v68-repair-dom}
\end{equation}
\end{theorem}

\begin{proof}
On one cell, \eqref{eq:s7v68-energy-drop} gives
\[
 e^{-\beta S(U)}
 \leq e^{-\beta\delta n}e^{-\beta S(FU)}.
\]
Integrate and use \eqref{eq:s7v68-push} with
\(g=e^{-\beta S}\):
\[
 \int_Ce^{-\beta S(U)}\,dm(U)
 \leq
 e^{-(\beta\delta-h)n}
 \int_\Omega e^{-\beta S(V)}\,dm(V)
 =
 e^{-(\beta\delta-h)n}Z_\beta.
\]
Divide by \(Z_\beta\), use
\eqref{eq:s7v68-n}, and sum over at most \(a_0^{|T|}\) cells.  This is
exactly \eqref{eq:s7v68-repair-dom}.
\end{proof}

\begin{lemma}[Independent repair selection]
\label{lem:s7v68-independent}
If candidate owner repairs have a conflict graph of maximum degree
\(\Delta_{\rm r}\), every finite candidate set \(T\) contains a
pairwise nonconflicting subset \(I\) with
\begin{equation}
 |I|\geq{|T|\over\Delta_{\rm r}+1}.                    \label{eq:s7v68-independent}
\end{equation}
\end{lemma}

\begin{proof}
Greedily select one vertex and remove it together with its at most
\(\Delta_{\rm r}\) neighbors.  Each selection removes at most
\(\Delta_{\rm r}+1\) candidates, which proves the bound.
\end{proof}

\begin{corollary}[Local additive repairs give principal product domination]
\label{cor:s7v68-local-repair}
Suppose a finite cell decomposition chooses one local repair for each
principal bad owner, nonconflicting repairs have additive action decrease
at least \(\delta\) per owner, and their combined reference-measure
pushforward loss is at most \(e^{h|I|}\).  Then
\Cref{thm:s7v68-repair} applies with
\(\kappa=\Delta_{\rm r}+1\).  Combining it with
\Cref{thm:s7v68-Wilson-compiler} gives
\begin{equation}
 \widehat s_\beta
 \leq
 179\,a_0^{1/239}
 \exp\left[
 -{\beta\delta-h\over239(\Delta_{\rm r}+1)}
 \right].                                             \label{eq:s7v68-combined}
\end{equation}
\end{corollary}

\begin{proof}
Choose \(I\) by
\Cref{lem:s7v68-independent}; the hypotheses give a repair branch with
\(n=|I|\).  Apply
\Cref{thm:s7v68-repair} and then take the \(239\)-th root in
\Cref{thm:s7v68-Wilson-compiler}.
\end{proof}

\begin{remark}[Why conditional-minimum subtraction is absent]
\label{rem:s7v68-absolute}
The proof uses the absolute action difference
\(S(U)-S(FU)\) before normalization.  No boundary-dependent scalar
minimum is subtracted.  This is precisely why the factor
\(e^{-\beta\delta n}\) survives division by \(Z_\beta\).
\end{remark}

\section{Why the V4--V5 form transfer is not the repair input}
\label{sec:s7v68-types}

\begin{proposition}[A uniform Dirichlet-form gap does not control event
rarity]
\label{prop:s7v68-form-no-event}
For every \(p\in(0,1)\), let \(\mu_p\) be the Bernoulli law on
\(\{0,1\}\) with \(\mu_p\{1\}=p\), and let \(Kf=\int f\,d\mu_p\).
Then
\begin{equation}
 \langle f,(I-K)f\rangle_{L^2(\mu_p)}
 =\operatorname{Var}_{\mu_p}(f)                       \label{eq:s7v68-Bernoulli-form}
\end{equation}
with constant one for every \(p\), while the event
\(A=\{1\}\) has probability \(p\), which can approach one.
\end{proposition}

\begin{proof}
\(K\) is the orthogonal projection onto constants, so
\(\langle f,(I-K)f\rangle=\|f-Kf\|_2^2\), equal to the variance.
The event probability is its definition and is independent of the form
identity.
\end{proof}

\begin{corollary}[Type separation for the existing transfer cards]
\label{cor:s7v68-type}
The V4 randomized plaquette form identity and the V5 martingale leakage
bound cannot, without an additional Gibbs-density or energy-repair
estimate, imply the principal-event product domination
\eqref{eq:s7v68-principal-dom}.  They remain relevant to conditional
Poincar\'e and response control after the probability sector is closed.
\end{corollary}

\begin{proof}
\Cref{prop:s7v68-form-no-event} gives a family with a fixed exact form
constant and arbitrarily non-rare events.  Hence no implication based
only on such form constants can yield a vanishing event parameter.
\end{proof}

\begin{assessment}[Upstream change after the V68 type check]
\label{ass:s7v68-upstream}
The generating-function calculation does not consume the V5 leakage
budget.  The probability bottleneck is an absolute Gibbs-measure
estimate: either construct the repair branches of
\Cref{thm:s7v68-repair} or use reflection positivity to obtain
\eqref{eq:s7v68-principal-dom}.  Only after that step should the
Dirichlet-form and response ledgers be reattached.
\end{assessment}

\section{Status and next acquisition}
\label{sec:s7v68-status}

\begin{theorem}[What V68 proves]
\label{thm:s7v68-result}
V68 proves:
\begin{enumerate}[label=\textup{(\roman*)},leftmargin=3.5em]
\item the bounded-incidence multi-source and moment-generating-function
      compilers;
\item the exact Wilson constants \(m=179\) and \(\ell=239\);
\item the resulting rate
      \(\widehat s_M=179r_M^{1/239}\);
\item the logical mismatch between actual certificate failures and
      boundary-singular conditional minima;
\item an absolute-energy repair theorem which yields principal product
      domination; and
\item the type obstruction preventing a pure \(L^2\) form estimate from
      supplying event rarity.
\end{enumerate}
\end{theorem}

\begin{proof}
Use
\Cref{thm:s7v68-owner,cor:s7v68-generating,
lem:s7v68-179-239,thm:s7v68-Wilson-compiler,
prop:s7v68-no-inclusion,thm:s7v68-repair,
prop:s7v68-form-no-event}.
\end{proof}

\begin{decision}[V69 acquisition order]
\label{dec:s7v68-V69}
V69 will test the reflection-positive route to
\eqref{eq:s7v68-principal-dom}.  It will prove the nonnegative character
expansion of the \(SU(3)\) Wilson plaquette weight, state the finite-torus
reflection form explicitly, and derive the chessboard product estimate
for reflection-generated block events.  It will then bound the
disseminated principal-defect event by an identity-neighborhood lower
bound on the partition function.  If the translated \(60\)-chord blocks
do not form reflection tiles, a fixed coloring and an enclosing
rectangular tile will be used, with every overlap loss shown.
\end{decision}

\begin{firewall}[Claim boundary after seventh-series V68]
\label{fire:s7v68-claim}
V68 does not prove the physical principal-defect product estimate,
construct a Wilson repair map, verify reflection positivity for the
chosen block event, bridge virtual conditional minima to actual fields,
prove SBTC or the physical mixed product majorant, establish interacting
endpoint continuity, UCM, BCC, the pairing/Bogoliubov sector, weighted
physical-source incidence, patchwise coarse-action cocycles, terminal
strong-coupling matching, regulator removal, Osterwalder--Schrader
reconstruction, or the full Yang--Mills mass gap.
\end{firewall}

\paragraph{Parent-version record.}
V68 was formed from
\texttt{Renewal\_Yang\_Mills\_Mass\_Gap\_Research\_Notes\_V67.tex}.
The V67 parent had \(29\,457\) logical source lines, \(1\,054\,229\)
bytes, and SHA--256
\[
 \texttt{12F2F24BAA107915453F7472152718C343F450C5722C52D38E03A5007DBE5D09}.
\]
The next compulsory companion audit is V70.

% =============================================================================
% END APPENDED SEVENTH-SERIES V68 MATERIAL
% =============================================================================

% =============================================================================
% BEGIN APPENDED SEVENTH-SERIES V69 MATERIAL
% =============================================================================

\chapter{Reflection-positive principal-defect domination}
\label{ch:s7v69}

\section{Purpose}
\label{sec:s7v69-purpose}

V68 isolates the principal-event product estimate
\eqref{eq:s7v68-principal-dom}.  This chapter proves that estimate on a
cofinal family of periodic Wilson regulators.  The proof has four parts:
the \(SU(3)\) Wilson plaquette weight has nonnegative character
coefficients; this gives reflection positivity; repeated reflection
Cauchy--Schwarz gives a chessboard estimate; and a disseminated
principal-defect event has an extensive absolute Wilson-action cost.
An explicit identity-neighborhood lower bound on the partition function
preserves that cost after normalization.

The resulting bound applies to defects of the sampled field.  It does
not yet solve the hybrid conditional-minimum problem isolated in V68.

\section{Positive-type Wilson plaquette kernel}
\label{sec:s7v69-character}

Let \(\widehat{SU(3)}\) denote the unitary dual and \(d_\lambda\) the
dimension of the irreducible representation \(\lambda\).  Put
\begin{equation}
 w_\beta(g):=\exp\{\beta\operatorname{ReTr}g\}.         \label{eq:s7v69-weight}
\end{equation}
The harmless factor \(e^{-3\beta}\) converts \(w_\beta\) into
\(e^{-\beta\phi(g)}\).

\begin{theorem}[Nonnegative character expansion]
\label{thm:s7v69-character}
For every \(\beta\geq0\),
\begin{equation}
 w_\beta(g)
 =
 \sum_{\lambda\in\widehat{SU(3)}}
 c_\lambda(\beta)\chi_\lambda(g),
 \qquad c_\lambda(\beta)\geq0,                         \label{eq:s7v69-character}
\end{equation}
with absolute and uniform convergence on \(SU(3)\).
\end{theorem}

\begin{proof}
Expand
\begin{equation}
 w_\beta(g)
 =
 \sum_{m,n\geq0}
 {(\beta/2)^{m+n}\over m!\,n!}
 (\operatorname{Tr}g)^m
 (\operatorname{Tr}g^{-1})^n.                         \label{eq:s7v69-expansion}
\end{equation}
The product
\((\operatorname{Tr}g)^m(\operatorname{Tr}g^{-1})^n\)
is the character of
\[
 {\bf3}^{\otimes m}\otimes\overline{\bf3}^{\otimes n}.
\]
Its irreducible decomposition has nonnegative integer multiplicities.
Thus every finite partial sum in
\eqref{eq:s7v69-expansion} is a nonnegative linear combination of
irreducible characters, and the coefficient of every fixed character
increases to a nonnegative limit.

Moreover, the absolute value of the term indexed by \(m,n\) is at most
\[
 {(\beta/2)^{m+n}\over m!\,n!}\,3^{m+n}.
\]
The double majorant sums to \(e^{3\beta}\).  The Weierstrass test gives
absolute uniform convergence of
\eqref{eq:s7v69-expansion}, and regrouping by irreducible character is
therefore legitimate.
\end{proof}

\begin{corollary}[Positive convolution kernel]
\label{cor:s7v69-kernel}
For every finite collection \(g_1,\ldots,g_N\in SU(3)\) and
\(z_1,\ldots,z_N\in\mathbb C\),
\begin{equation}
 \sum_{a,b=1}^N
 \overline z_a z_b\,w_\beta(g_ag_b^{-1})\geq0.         \label{eq:s7v69-positive-kernel}
\end{equation}
\end{corollary}

\begin{proof}
For a unitary representative \(D^\lambda\),
\[
 \chi_\lambda(g_ag_b^{-1})
 =
 \sum_{i,j}
 D^\lambda_{ij}(g_a)
 \overline{D^\lambda_{ij}(g_b)}.
\]
Insert \eqref{eq:s7v69-character}.  The quadratic form becomes a sum of
\[
 c_\lambda(\beta)
 \sum_{i,j}
 \left|\sum_a z_a\overline{D^\lambda_{ij}(g_a)}\right|^2,
\]
which is nonnegative.
\end{proof}

\section{Reflection positivity of the finite Wilson measure}
\label{sec:s7v69-RP}

Let \(\Lambda\) be an even four-dimensional periodic hypercubic lattice.
Its normalized Wilson measure is
\begin{equation}
 d\mu_{\Lambda,\beta}(U)
 =
 {1\over Z_{\Lambda,\beta}}
 \exp\left[-\beta\sum_{p\in{\cal P}_\Lambda}\phi(W_p)\right]
 \prod_{e\in E_\Lambda}dU_e.                          \label{eq:s7v69-measure}
\end{equation}
Choose a coordinate reflection \(\vartheta\) through a
reflection-adapted cut, including orientation reversal on reflected
links, and let \(\Theta F=\overline{F\circ\vartheta}\).

\begin{theorem}[Wilson reflection positivity]
\label{thm:s7v69-RP}
For every bounded function \(F\) depending only on the positive half of
the link variables,
\begin{equation}
 \int F\,\Theta F\,d\mu_{\Lambda,\beta}\geq0.          \label{eq:s7v69-RP}
\end{equation}
The assertion holds for each coordinate reflection used by an even
rectangular chessboard tiling.
\end{theorem}

\begin{proof}
Plaquettes strictly in the two open halves contribute a factor
\(B\,\Theta B\).  After the usual orientation reversal at the cut, each
crossing plaquette word splits into a positive-half word \(g\) and the
reflected inverse word, so its nonconstant Boltzmann factor has the form
\[
 w_\beta(g\,\vartheta(g)^{-1}).
\]
By \Cref{thm:s7v69-character},
\begin{equation}
 w_\beta(g\,\vartheta(g)^{-1})
 =
 \sum_{\lambda,i,j}
 c_\lambda(\beta)
 D^\lambda_{ij}(g)\,
 \Theta D^\lambda_{ij}(g).                            \label{eq:s7v69-cross-factor}
\end{equation}
The product over crossing plaquettes is again a sum of terms
\(C_\alpha\Theta C_\alpha\) with nonnegative coefficients, because
products of such sums have that form.  The factors \(e^{-3\beta}\) are
positive constants.

The product Haar measure is invariant under the reflection and
factorizes between paired half-lattice variables.  Therefore the
unnormalized left side of \eqref{eq:s7v69-RP} is a nonnegative sum of
squared absolute half-lattice integrals,
\[
 \sum_\alpha c_\alpha
 \left|\int FBC_\alpha\,dm_+\right|^2.
\]
Division by the positive partition function preserves the inequality.
The same factorization applies to every coordinate cut of an even
rectangular tiling.
\end{proof}

\begin{remark}[Gauge fixing is not used]
\label{rem:s7v69-no-gauge-fix}
The proof is on the compact link carrier with product Haar measure.
There is no Faddeev--Popov determinant, Gribov chart, or tree-gauge
Jacobian in \eqref{eq:s7v69-RP}.  Gauge reduction may be performed only
after the global probability estimate has been obtained.
\end{remark}

\section{Chessboard dissemination}
\label{sec:s7v69-chessboard}

Tile \(\Lambda\) by \(N\) congruent rectangular cells and assume that
successive reflections in their faces generate the tile positions.
If \(A\) is a tile event supported away from the tile boundary, denote
by \(A^\#\) the event obtained by imposing its successively reflected
copy on every tile.  Reflections may reverse the local pattern; the
plaquette defect \(\phi(W)=\phi(W^{-1})\) is unaffected.

\begin{theorem}[Chessboard estimate]
\label{thm:s7v69-chessboard}
Let \(A_j\) be events placed on distinct tiles, with reflection
orientation chosen according to their positions.  Then
\begin{equation}
 \mu_{\Lambda,\beta}\left(\bigcap_{j\in J}A_j\right)
 \leq
 \prod_{j\in J}
 \mu_{\Lambda,\beta}(A_j^\#)^{1/N}.                   \label{eq:s7v69-chessboard}
\end{equation}
\end{theorem}

\begin{proof}
For one reflection dividing the torus into two congruent halves,
reflection positivity and Cauchy--Schwarz give
\[
 |\langle F,\Theta G\rangle|
 \leq
 \langle F,\Theta F\rangle^{1/2}
 \langle G,\Theta G\rangle^{1/2}.
\]
Take \(F,G\) to be products of the event indicators lying in the two
halves.  The two factors on the right are the configurations obtained by
reflecting all events from one half into the other.  Repeat this operation
for every family of parallel tile faces and every coordinate direction.
After the full sequence, each original event has been reflected onto all
\(N\) tiles.  Each Cauchy--Schwarz step halves the exponent attached to
that original event; after the complete dissemination its exponent is
\(1/N\).  Multiplying the resulting factors gives
\eqref{eq:s7v69-chessboard}.  Indicator functions follow directly by
bounded approximation if a cut boundary has Haar-null ambiguities.
\end{proof}

\begin{corollary}[One-pattern product bound]
\label{cor:s7v69-one-pattern}
If every disseminated event in
\eqref{eq:s7v69-chessboard} has probability at most \(\rho^N\), then
\begin{equation}
 \mu_{\Lambda,\beta}\left(\bigcap_{j\in J}A_j\right)
 \leq\rho^{|J|}.                                      \label{eq:s7v69-one-pattern}
\end{equation}
\end{corollary}

\begin{proof}
Insert \(\mu(A_j^\#)^{1/N}\leq\rho\) into
\eqref{eq:s7v69-chessboard}.
\end{proof}

\section{Normalized cost of a disseminated defect event}
\label{sec:s7v69-cost}

Let \(Q\) be a finite plaquette pattern inside one tile and define
\begin{equation}
 D_Q(U):=\sum_{p\in Q}\phi(W_p(U)),\qquad
 A_Q(c):=\{D_Q\geq c\}.                               \label{eq:s7v69-AQ}
\end{equation}
Suppose the reflected copies \(Q_j\), \(1\leq j\leq N\), count any
plaquette at most \(\nu\) times.

\begin{lemma}[Dissemination forces extensive action]
\label{lem:s7v69-extensive}
On \(A_Q(c)^\#\),
\begin{equation}
 \sum_{p\in{\cal P}_\Lambda}\phi(W_p)
 \geq{Nc\over\nu}.                                    \label{eq:s7v69-extensive}
\end{equation}
\end{lemma}

\begin{proof}
Every tile event gives
\(\sum_{p\in Q_j}\phi(W_p)\geq c\).  Sum over \(j\).
The left side counts any global plaquette at most \(\nu\) times, so it is
at most \(\nu\sum_p\phi(W_p)\).  Rearrangement proves the claim.
\end{proof}

\begin{lemma}[Identity-neighborhood partition lower bound]
\label{lem:s7v69-Z}
There are group constants \(r_0,v_0>0\) such that, for
\(0<\rho\leq r_0\) and \(\beta\geq1\),
\begin{equation}
 Z_{\Lambda,\beta}
 \geq
 \left(v_0\rho^8\beta^{-4}\right)^{|E_\Lambda|}
 \exp\left[-8\rho^2|{\cal P}_\Lambda|\right].         \label{eq:s7v69-Z}
\end{equation}
\end{lemma}

\begin{proof}
For sufficiently small \(r\), the normalized Haar measure of the
geodesic ball \(B_r(I)\subset SU(3)\) is at least \(v_0r^8\), because
\(\dim SU(3)=8\) and exponential coordinates have a smooth positive
Jacobian at the origin.  Restrict every link integral to
\[
 U_e\in B_{\rho/\sqrt\beta}(I).
\]
The bi-invariant distance and the triangle inequality give
\[
 d(W_p,I)\leq{4\rho\over\sqrt\beta}.
\]
If \(W=\exp X\) with \(X\) a shortest logarithm, then
\[
 \phi(W)
 ={1\over2}\|W-I\|_{\rm HS}^2
 \leq{1\over2}|X|^2
 ={1\over2}d(W,I)^2.
\]
Thus \(\phi(W_p)\leq8\rho^2/\beta\), and the Boltzmann density on the
restricted set is at least
\(\exp[-8\rho^2|{\cal P}_\Lambda|]\).  Its product Haar measure is at
least
\((v_0\rho^8\beta^{-4})^{|E_\Lambda|}\), proving
\eqref{eq:s7v69-Z}.
\end{proof}

\begin{theorem}[Disseminated defect probability]
\label{thm:s7v69-disseminated}
Let
\[
 e_Q:={|E_\Lambda|\over N},\qquad
 p_Q:={|{\cal P}_\Lambda|\over N},
\]
which are fixed by the tile geometry.  Then
\begin{equation}
 \mu_{\Lambda,\beta}(A_Q(c)^\#)^{1/N}
 \leq
 C_Q\beta^{4e_Q}
 \exp\left[-{\beta c\over\nu}\right],                 \label{eq:s7v69-disseminated}
\end{equation}
where
\begin{equation}
 C_Q:=
 (v_0\rho^8)^{-e_Q}e^{8\rho^2p_Q}.                   \label{eq:s7v69-CQ}
\end{equation}
In particular, the right side tends to zero exponentially in \(\beta\).
\end{theorem}

\begin{proof}
By \Cref{lem:s7v69-extensive} and normalization,
\[
 \mu_{\Lambda,\beta}(A_Q(c)^\#)
 \leq
 {e^{-\beta Nc/\nu}\over Z_{\Lambda,\beta}}.
\]
Insert \eqref{eq:s7v69-Z}, take the \(N\)-th root, and use the definitions
of \(e_Q,p_Q,C_Q\).
\end{proof}

\begin{firewall}[Why the denominator cannot be discarded]
\label{fire:s7v69-denominator}
The numerator bound \(e^{-\beta Nc/\nu}\) alone is not a probability
estimate, because \(Z_{\Lambda,\beta}\) is volume dependent.  The local
identity-ball lower bound produces only the per-tile entropy loss
\(C_Q\beta^{4e_Q}\), which is dominated by the fixed positive action
margin as \(\beta\to\infty\).
\end{firewall}

\section{Application to every translated principal collar}
\label{sec:s7v69-collar}

The \(72\) principal plaquettes \(P_o^+\) of one oriented collar have
fixed finite diameter.  Choose an integer \(h\) so large that each
translated support fits strictly inside a cube of side \(h\).  Restrict
the cofinal periodic regulators to side lengths compatible with repeated
reflection tiling by these cubes.

\begin{lemma}[Finite reflection classes of owner blocks]
\label{lem:s7v69-classes}
The positive oriented links can be partitioned into at most
\begin{equation}
 K_{\rm tile}:=4(2h)^4                                \label{eq:s7v69-Ktile}
\end{equation}
classes such that, within one class:
\begin{enumerate}[label=\textup{(\roman*)},leftmargin=3.5em]
\item all principal supports lie in distinct reflection tiles;
\item all block centers differ by translations in \((2h\mathbb Z)^4\);
      and
\item their local events are copies with the same orientation after an
      even number of face reflections.
\end{enumerate}
\end{lemma}

\begin{proof}
Classify a positive link by its direction and by the four coordinates of
its base point modulo \(2h\).  This gives \(4(2h)^4\) classes.  The
choice of \(h\) places the finite support inside its assigned tile.
Translations by \(2h\) in one coordinate are products of two adjacent
face reflections, so they preserve the local orientation.  Distinct
centers in one residue class lie in distinct tiles.
\end{proof}

\begin{theorem}[Principal Wilson-defect product domination]
\label{thm:s7v69-principal}
On the reflection-compatible periodic regulators, the V68 principal
events
\[
 A_o^*=
 \left\{\sum_{q\in P_o^+}\phi(W_q)\geq c_*\right\}
\]
satisfy, for every finite owner set \(T\),
\begin{equation}
 \mu_{\Lambda,\beta}\left(\bigcap_{o\in T}A_o^*\right)
 \leq
 r_\beta^{|T|},                                       \label{eq:s7v69-principal}
\end{equation}
where
\begin{align}
 r_\beta
 &:=
 \left(
 C_Q\beta^{4e_Q}e^{-\beta c_*/\nu}
 \right)^{1/K_{\rm tile}}                             \label{eq:s7v69-rbeta}\\
 &\longrightarrow0.
\end{align}
With the disjoint interior placement of
\Cref{lem:s7v69-classes}, one may take \(\nu=1\); retaining \(\nu\)
makes the statement stable under a bounded-overlap enlargement.
\end{theorem}

\begin{proof}
Choose a class \(j\) containing at least
\(|T|/K_{\rm tile}\) members of \(T\), and call that subset \(T_j\).
Discarding event requirements increases probability:
\[
 \mu\left(\bigcap_{o\in T}A_o^*\right)
 \leq
 \mu\left(\bigcap_{o\in T_j}A_o^*\right).
\]
Apply the chessboard estimate
\Cref{thm:s7v69-chessboard} to this one local pattern.  Reflected copies
on the intervening tiles may reverse the shape, but each contains the
same number of plaquettes and the same defect threshold.  Therefore
\Cref{thm:s7v69-disseminated} bounds every chessboard norm by
\[
 \rho_\beta:=C_Q\beta^{4e_Q}e^{-\beta c_*/\nu}.
\]
It follows that
\[
 \mu\left(\bigcap_{o\in T_j}A_o^*\right)
 \leq
 \rho_\beta^{|T_j|}
 \leq
 \rho_\beta^{|T|/K_{\rm tile}}
 =r_\beta^{|T|},
\]
for all sufficiently large \(\beta\) such that \(\rho_\beta\leq1\).
For the remaining finite range replace \(r_\beta\) by its minimum-safe
upper bound one.  Exponential decay dominates the polynomial prefactor,
so \(r_\beta\to0\).
\end{proof}

\begin{corollary}[Actual fully-good failures are product dominated]
\label{cor:s7v69-actual}
The sampled-field certificate failures
\(\widehat\Sigma_b\) of V68 obey
\begin{equation}
 \mu_{\Lambda,\beta}
 \left(\bigcap_{b\in S}\widehat\Sigma_b\right)
 \leq
 \widehat s_\beta^{|S|},                              \label{eq:s7v69-actual}
\end{equation}
with
\begin{equation}
 \widehat s_\beta
 \leq
 179\,
 C_Q^{1/(239K_{\rm tile})}
 \beta^{4e_Q/(239K_{\rm tile})}
 \exp\left[
 -{\beta c_*\over239\nu K_{\rm tile}}
 \right]
 \longrightarrow0.                                   \label{eq:s7v69-actual-rate}
\end{equation}
\end{corollary}

\begin{proof}
Insert \eqref{eq:s7v69-rbeta} into
\Cref{thm:s7v68-Wilson-compiler}.
\end{proof}

\begin{assessment}[What the reflection route closes]
\label{ass:s7v69-progress}
The principal product estimate
\eqref{eq:s7v68-principal-dom}, left open in V68, is proved on a cofinal
reflection-compatible family of finite tori.  Its normalization, volume
uniformity, overlap cost, and entropy prefactor are explicit.  Therefore
actual sampled configurations outside the V1 fully good certificate form
an exponentially dilute family at weak coupling.
\end{assessment}

\begin{firewall}[The conditional-minimum bridge remains open]
\label{fire:s7v69-bridge}
\Cref{cor:s7v69-actual} does not imply product domination of
\(\Sigma_b^{\rm bdry}\), because the latter is evaluated at virtual
conditional minimizers.  Nor does it prove that arbitrary nonperiodic
regulators have the same chessboard bound, or that the cofinal periodic
subnet has a unique regulator-independent limit.  The V68 hybrid
completion/repair problem is unchanged and must not be hidden inside
reflection positivity.
\end{firewall}

\section{Status and next acquisition}
\label{sec:s7v69-status}

\begin{theorem}[What V69 proves]
\label{thm:s7v69-result}
V69 proves the nonnegative \(SU(3)\) Wilson character expansion,
reflection positivity of the finite periodic Wilson measure, the
chessboard product estimate, the normalized disseminated-defect bound,
and the principal product domination
\eqref{eq:s7v69-principal}.  Combined with V68, this gives the explicit
sampled-field certificate-failure rate
\eqref{eq:s7v69-actual-rate}.
\end{theorem}

\begin{proof}
Use
\Cref{thm:s7v69-character,thm:s7v69-RP,
thm:s7v69-chessboard,thm:s7v69-disseminated,
thm:s7v69-principal,cor:s7v69-actual}.
\end{proof}

\begin{decision}[V70 acquisition order]
\label{dec:s7v69-V70}
V70 is the compulsory companion audit.  It will inspect the current SM
and NS notes for a normalized completion kernel, simultaneous repair,
or reflection-positive conditioning argument.  It will then attack the
hybrid-minimum bridge: either compare a boundary-singular label to the
positive conditional probability of an actual certificate failure, or
construct an extended-space repair kernel with a controlled
Radon--Nikodym loss.  No companion result will be imported across field
types without reproducing its measure comparison.
\end{decision}

\begin{firewall}[Claim boundary after seventh-series V69]
\label{fire:s7v69-claim}
V69 does not bridge boundary-singular conditional minima to sampled
fields, prove SBTC, the physical mixed product majorant, interacting
endpoint continuity, UCM, BCC, the pairing/Bogoliubov sector, weighted
physical-source incidence, patchwise coarse-action cocycles, terminal
strong-coupling matching, regulator independence, removal of all
cutoffs, Osterwalder--Schrader reconstruction, or the full Yang--Mills
mass gap.
\end{firewall}

\paragraph{Parent-version record.}
V69 was formed from
\texttt{Renewal\_Yang\_Mills\_Mass\_Gap\_Research\_Notes\_V68.tex}.
The V68 parent had \(29\,958\) logical source lines, \(1\,072\,290\)
bytes, and SHA--256
\[
 \texttt{C05180F2C249C152BC49137AEA4A06E7FE18D59D5C4D68C880EB8C88948BE8FD}.
\]
The next compulsory companion audit is V70.

% =============================================================================
% END APPENDED SEVENTH-SERIES V69 MATERIAL
% =============================================================================

% =============================================================================
% BEGIN APPENDED SEVENTH-SERIES V70 MATERIAL
% =============================================================================

\chapter{V70 companion audit and the heat-bath bridge for robust singular labels}
\label{ch:s7v70}

\section{Compulsory SM/NS audit}
\label{sec:s7v70-audit}

\begin{audit}[Current companion files read at V70]
\label{audit:s7v70-files}
The active companion notes are:
\begin{enumerate}[label=\textup{(\roman*)},leftmargin=3.8em]
\item
\texttt{Renewal\_SM\_Einstein\_Continuum\_Research\_Notes\_V29.tex},
with \(10\,066\) logical source lines, \(360\,500\) bytes, and SHA--256
\[
 \texttt{CAB3DB39F6498AE274619391763FC5F64D47C7A8E3C67F5BDF5BED9FD9D87786};
\]
\item
\texttt{Renewal\_NS\_Stability\_Research\_Notes\_V26.tex},
with \(5\,319\) logical source lines, \(278\,283\) bytes, and SHA--256
\[
 \texttt{82C887F8C2C69E4F28FAD7C3DC8DE5B9099CB5A4BF431EBA1FFF3E91935978AD}.
\]
\end{enumerate}
The SM file has advanced from V22 at the preceding YM audit to V29.
The NS file is unchanged.
\end{audit}

\begin{theorem}[Transferable conclusions of the V70 audit]
\label{thm:s7v70-audit}
The companion audit gives the following type-correct guidance.
\begin{enumerate}[leftmargin=3em]
\item SM V18 proves that a normalized spectator probability kernel leaves
      the old marginal unchanged.  This measure-theoretic principle can
      be strengthened here by using the \emph{exact} Wilson conditional
      heat-bath kernel.
\item SM V25--V26 prove finite-color multiplication and a normalized
      coercive conditional tail only under a uniform boundary-label
      coercivity/oscillation hypothesis.  That hypothesis is the SM
      analogue of the YM row being repaired and cannot be imported.
\item SM V27--V29 concern stress measures, integration by parts, and
      closable metric gradients.  They provide no Wilson reflection or
      conditional-minimum estimate.
\item NS V26 supplies only pointwise Oseen-kernel derivative constants and
      contains no compact Gibbs kernel or product-event domination.
\end{enumerate}
\end{theorem}

\begin{proof}
These statements are the hypotheses, conclusions, and claim firewalls of
the cited active companion sections fixed by
\Cref{audit:s7v70-files}.  Only the normalized-kernel identity is
field-independent.  All quantitative conditional estimates in the SM and
NS notes act on different carriers and require different physical
loaders.
\end{proof}

\begin{assessment}[Audit-guided change]
\label{ass:s7v70-audit}
The correct bridge is not an artificial minimizer inserted with
unnormalized weight.  It is an independent draw from the exact
conditional Wilson law.  Same-color conditional factorization makes all
such draws compatible, and heat-bath invariance makes the resampled field
have exactly the original Wilson distribution.  The remaining question
is whether a boundary-singular label forces that conditional draw into an
actual certificate-failure event with high probability.
\end{assessment}

\section{Uniform energy-excess concentration without a Hessian hypothesis}
\label{sec:s7v70-Laplace}

Let \(X\) be a compact \(n\)-dimensional Riemannian manifold with
normalized smooth measure \(m\), and let \(\Xi\) be compact.  Suppose
\(S:X\times\Xi\to\mathbb R\) is continuous and uniformly
\(L\)-Lipschitz in \(x\).  Put
\begin{align}
 m_S(\eta)&:=\min_{x\in X}S(x;\eta),                  \label{eq:s7v70-min}\\
 {\cal M}(\eta)&:=\{x:S(x;\eta)=m_S(\eta)\},          \label{eq:s7v70-M}\\
 d\nu_{\beta,\eta}(x)
 &:=
 {e^{-\beta(S(x;\eta)-m_S(\eta))}\over Z_{\beta,\eta}}
 \,dm(x).                                             \label{eq:s7v70-nu}
\end{align}

\begin{theorem}[Uniform compact-family excess concentration]
\label{thm:s7v70-Laplace}
There are geometric constants \(r_0,v_0>0\) such that, for every
energy buffer \(\tau>0\), with
\begin{equation}
 r_\tau:=
 \min\left\{r_0,{\tau\over2\max(L,1)}\right\},\qquad
 C_\tau:=(v_0r_\tau^n)^{-1},                          \label{eq:s7v70-rtau}
\end{equation}
one has
\begin{equation}
 \sup_{\eta\in\Xi}
 \nu_{\beta,\eta}
 \left\{S(x;\eta)-m_S(\eta)\geq\tau\right\}
 \leq
 \varepsilon_{\beta,\tau}:=
 C_\tau e^{-\beta\tau/2}.                            \label{eq:s7v70-Laplace}
\end{equation}
No differentiability, Hessian floor, continuity of the argmin set, or
uniqueness of the conditional minimum is required.
\end{theorem}

\begin{proof}
Compact Riemannian geometry gives
\[
 m(B_r(x))\geq v_0r^n
\]
uniformly in \(x\) and \(0<r\leq r_0\).  For each \(\eta\), choose
\(x_\eta\in{\cal M}(\eta)\).  On \(B_{r_\tau}(x_\eta)\),
\[
 S(x;\eta)-m_S(\eta)\leq Lr_\tau\leq\tau/2.
\]
Hence
\[
 Z_{\beta,\eta}
 \geq v_0r_\tau^n e^{-\beta\tau/2}.
\]
On the excess-\(\tau\) set the unnormalized density is at most
\(e^{-\beta\tau}\), and \(m(X)=1\).  Divide by the denominator lower
bound to obtain \eqref{eq:s7v70-Laplace}.
\end{proof}

\begin{firewall}[Fixed-energy concentration is not the V65 Gaussian
tube]
\label{fire:s7v70-scale}
Theorem \ref{thm:s7v70-Laplace} controls a fixed positive
conditional-energy excess.  It gives no conversion from a shrinking
geometric distance \(g\sqrt M\) to an excess of order \(g^2M\).  That
conversion still requires a uniform quadratic normal floor.  V70 uses
the theorem only when a fixed energy buffer separates a robust singular
label from the actual regular certificate.
\end{firewall}

\section{Exact same-color heat-bath invariance}
\label{sec:s7v70-heatbath}

Return to the finite periodic Wilson measure
\(\mu_{\Lambda,\beta}\).  Fix one V66 conflict color
\({\cal B}_j\).  Conditional on all variables outside its interiors, no
interaction contains two of those interiors, so the conditional law is a
product
\begin{equation}
 \bigotimes_{b\in{\cal B}_j}
 \nu_{b,\beta,\eta_b}(dx_b).                          \label{eq:s7v70-factor}
\end{equation}
Let \(K_j(U,dU')\) retain the exterior and independently redraw every
same-color interior from \eqref{eq:s7v70-factor}.

\begin{lemma}[Heat-bath invariance and reversibility]
\label{lem:s7v70-invariance}
The kernel \(K_j\) satisfies
\begin{equation}
 \mu_{\Lambda,\beta}K_j=\mu_{\Lambda,\beta}           \label{eq:s7v70-invariance}
\end{equation}
and is reversible on \(L^2(\mu_{\Lambda,\beta})\).
The same is true if only a finite subset of one color is redrawn.
\end{lemma}

\begin{proof}
Disintegrate \(\mu_{\Lambda,\beta}\) into its exterior marginal and the
conditional product law \eqref{eq:s7v70-factor}.  Redrawing from that same
conditional probability leaves each conditional integral unchanged,
which proves \eqref{eq:s7v70-invariance}.  Given the exterior, the old and
new interiors are independent draws from the same product law; their joint
law is symmetric under interchange.  Integrating over the exterior proves
detailed balance and hence reversibility.  The argument for a subset is
identical.
\end{proof}

\begin{remark}[Normalized spectator interpretation]
\label{rem:s7v70-spectator}
One may first regard the redrawn interiors as normalized auxiliary
spectators, in the sense suggested by the SM audit.  Replacing the old
interiors by those spectators then yields
\Cref{lem:s7v70-invariance}.  The normalization is exact for every
exterior field; no field-dependent partition factor is dropped.
\end{remark}

\section{Robust boundary-singular labels}
\label{sec:s7v70-robust}

For block \(b\), let
\(\widehat\Sigma_b(x_b,\eta_b)\)
be the actual certificate-failure event
\eqref{eq:s7v68-Sigmahat}, and let
\({\cal M}_b(\eta_b)\) be the conditional minimum set.  Define the
\(\tau\)-energy-robust boundary-singular event
\begin{equation}
 \Sigma_{b,\tau}^{\rm rob}
 :=
 \left\{
 \inf_{x\in\widehat\Sigma_b(\,\cdot\,,\eta_b)^c}
 \left[
 S_b(x;\eta_b)-m_b(\eta_b)
 \right]\geq\tau
 \right\}.                                            \label{eq:s7v70-robust}
\end{equation}
The infimum over an empty complement is \(+\infty\).  Thus every
regular-certificate configuration lies at conditional energy at least
\(\tau\) above the minimum.  The event depends only on the exterior
label.

\begin{lemma}[A robust label forces a resampled actual failure]
\label{lem:s7v70-force}
For every exterior label in
\(\Sigma_{b,\tau}^{\rm rob}\),
\begin{equation}
 \nu_{b,\beta,\eta_b}
 \left(\widehat\Sigma_b\right)
 \geq1-\varepsilon_{\beta,\tau},                     \label{eq:s7v70-force}
\end{equation}
where \(\varepsilon_{\beta,\tau}\) is the uniform bound in
\eqref{eq:s7v70-Laplace} for the finite block family.
\end{lemma}

\begin{proof}
By \eqref{eq:s7v70-robust}, the complement of
\(\widehat\Sigma_b\) has conditional excess at least \(\tau\).
Apply \Cref{thm:s7v70-Laplace} to that complement.
\end{proof}

\begin{theorem}[Heat-bath bridge to the V69 product estimate]
\label{thm:s7v70-bridge}
Let \(S\) be a finite set of blocks in one conflict color.  If
\(\varepsilon_{\beta,\tau}<1\), then
\begin{equation}
 \mu_{\Lambda,\beta}
 \left(\bigcap_{b\in S}\Sigma_{b,\tau}^{\rm rob}\right)
 \leq
 \left(
 {\widehat s_\beta\over1-\varepsilon_{\beta,\tau}}
 \right)^{|S|},                                       \label{eq:s7v70-bridge}
\end{equation}
where \(\widehat s_\beta\) is the sampled-field rate
\eqref{eq:s7v69-actual-rate}.  Hence the robust boundary-singular labels
are product dominated on each color by a parameter tending to zero.
\end{theorem}

\begin{proof}
Condition on the complement of the same-color interiors.  On the event
\(\cap_{b\in S}\Sigma_{b,\tau}^{\rm rob}\), conditional factorization
and \Cref{lem:s7v70-force} give
\[
 K_j\left(U,
 \bigcap_{b\in S}
 \{\widehat\Sigma_b(U')\}\right)
 \geq
 (1-\varepsilon_{\beta,\tau})^{|S|}.
\]
Therefore, pointwise in the conditioned exterior,
\[
 {\bf1}_{\cap_{b\in S}\Sigma_{b,\tau}^{\rm rob}}
 \leq
 (1-\varepsilon_{\beta,\tau})^{-|S|}
 K_j{\bf1}_{\cap_{b\in S}\widehat\Sigma_b}.
\]
Integrate against \(\mu_{\Lambda,\beta}\).  Heat-bath invariance
\eqref{eq:s7v70-invariance} replaces the integral of the last term by
\(\mu_{\Lambda,\beta}(\cap_{b\in S}\widehat\Sigma_b)\), which is at most
\(\widehat s_\beta^{|S|}\) by
\Cref{cor:s7v69-actual}.  This proves
\eqref{eq:s7v70-bridge}.
\end{proof}

\begin{corollary}[Arbitrary-set robust domination]
\label{cor:s7v70-arbitrary}
For arbitrary finite \(S\subset{\cal B}\),
\begin{equation}
 \mu_{\Lambda,\beta}
 \left(\bigcap_{b\in S}\Sigma_{b,\tau}^{\rm rob}\right)
 \leq
 \left(
 {\widehat s_\beta\over1-\varepsilon_{\beta,\tau}}
 \right)^{|S|/K}.                                    \label{eq:s7v70-arbitrary}
\end{equation}
\end{corollary}

\begin{proof}
Select a color containing at least \(|S|/K\) blocks, discard the other
event requirements, and apply \Cref{thm:s7v70-bridge}.
\end{proof}

\section{Transition labels and the refined regular/singular compiler}
\label{sec:s7v70-transition}

A singular label which is not \(\tau\)-energy-robust lies within the fixed
transition band
\begin{equation}
 {\cal T}_{b,\tau}
 :=
 \Sigma_b^{\rm bdry}\cap
 \left\{
 0\leq
 \inf_{x\in\widehat\Sigma_b(\,\cdot\,,\eta_b)^c}
 \left[
 S_b(x;\eta_b)-m_b(\eta_b)
 \right]<\tau
 \right\}.                                            \label{eq:s7v70-transition}
\end{equation}
Such a label may have order-one probability.  It must be retained in the
slow transition atlas rather than included in the rare singular event.

\begin{theorem}[Same-color regular/robust compiler]
\label{thm:s7v70-compiler}
Assume:
\begin{enumerate}[label=\textup{(\roman*)},leftmargin=3.5em]
\item every transition label
      \eqref{eq:s7v70-transition} is retained in the low/slow chart
      family;
\item on the remaining regular labels,
\begin{equation}
 \mu(E_b\mid{\cal F}_{b^c})\leq p_M;                  \label{eq:s7v70-pM}
\end{equation}
\item the only unretained nonregular labels are
      \(\Sigma_{b,\tau}^{\rm rob}\).
\end{enumerate}
Put
\begin{equation}
 s_{\beta,\tau}^{(0)}
 :=
 {\widehat s_\beta\over1-\varepsilon_{\beta,\tau}}.   \label{eq:s7v70-s0}
\end{equation}
If \(p_M+s_{\beta,\tau}^{(0)}\leq1\), then for every finite
\(S\subset{\cal B}\),
\begin{equation}
 \mu\left(\bigcap_{b\in S}E_b\right)
 \leq
 \left(p_M+s_{\beta,\tau}^{(0)}\right)^{|S|/K}.        \label{eq:s7v70-final-dom}
\end{equation}
\end{theorem}

\begin{proof}
First suppose \(S\) lies in one color and condition on all complementary
variables.  On a regular label use \eqref{eq:s7v70-pM}; on an unretained
robust label use the trivial bound one.  As in V67,
\[
 \mu\left(\bigcap_{b\in S}E_b\right)
 \leq
 {\mathbb E}\prod_{b\in S}
 \left(p_M+
 {\bf1}_{\Sigma_{b,\tau}^{\rm rob}}\right).
\]
Expand over subsets \(T\subset S\) and apply the same-color domination
\eqref{eq:s7v70-bridge}:
\[
 \sum_{T\subset S}
 p_M^{|S|-|T|}
 \left(s_{\beta,\tau}^{(0)}\right)^{|T|}
 =
 \left(p_M+s_{\beta,\tau}^{(0)}\right)^{|S|}.
\]
For arbitrary \(S\), retain a color containing at least \(|S|/K\)
members and discard the other requirements.
\end{proof}

\begin{corollary}[Weak-coupling rate]
\label{cor:s7v70-rate}
For fixed \(\tau>0\),
\begin{equation}
 \varepsilon_{\beta,\tau}
 =O(e^{-\beta\tau/2}),\qquad
 \widehat s_\beta
 =
 O\left(
 \beta^{q_{\rm tile}}
 e^{-c_{\rm tile}\beta}
 \right).                                             \label{eq:s7v70-rates}
\end{equation}
Thus \(s_{\beta,\tau}^{(0)}\to0\).  If the regular-sector V65 rate
\(p_M=CM^{d_f/2}e^{-cM+D_M}\) also tends to zero, the full unretained
bad-or-robust-singular parameter in
\eqref{eq:s7v70-final-dom} tends to zero.
\end{corollary}

\begin{proof}
The first estimate is \eqref{eq:s7v70-Laplace}; the second is
\eqref{eq:s7v69-actual-rate}.  Their quotient in
\eqref{eq:s7v70-s0} has denominator tending to one.
\end{proof}

\begin{firewall}[Retaining a transition band is a low-sector obligation]
\label{fire:s7v70-transition}
V45 guarantees that a finite buffered spectral cover can retain
transition modes while preserving a floor on the complementary band.
It does not prove the interacting effective dynamics, compatibility, or
locality of the enlarged transition sector.  Thus
\Cref{thm:s7v70-compiler} closes the \emph{probability} bridge conditional
on retaining \({\cal T}_{b,\tau}\); it does not close UCM, BCC, or the
low-sector mass estimate.
\end{firewall}

\section{Status and next acquisition}
\label{sec:s7v70-status}

\begin{theorem}[What V70 proves]
\label{thm:s7v70-result}
V70 completes the compulsory companion audit and proves:
\begin{enumerate}[label=\textup{(\roman*)},leftmargin=3.5em]
\item uniform fixed-distance Laplace concentration for a compact family
      of conditional actions, without a Hessian assumption;
\item exact invariance and reversibility of the same-color Wilson
      heat-bath kernel;
\item the bridge
      \eqref{eq:s7v70-bridge} from robust boundary-singular labels to
      V69 sampled-field failures; and
\item the refined regular/robust arbitrary-set domination
      \eqref{eq:s7v70-final-dom}.
\end{enumerate}
Consequently the robust part of the V68 hybrid-minimum obstruction is
closed on the reflection-compatible periodic regulators.  The fixed
transition band is retained rather than falsely assigned a small
probability.
\end{theorem}

\begin{proof}
Use
\Cref{thm:s7v70-audit,thm:s7v70-Laplace,
lem:s7v70-invariance,thm:s7v70-bridge,
thm:s7v70-compiler}.
\end{proof}

\begin{assessment}[Progress at seventh-series V70]
\label{ass:s7v70-status}
The distributional boundary transfer requested in V1 is now substantially
sharper.  Actual bad collars are exponentially dilute by reflection
positivity; robustly singular conditional labels are converted to those
actual collars by an invariant heat-bath redraw; and near-threshold labels
are identified as slow transition data.  No essential supremum over every
exterior field is used.
\end{assessment}

\begin{decision}[V71 acquisition order]
\label{dec:s7v70-V71}
V71 will load the retained transition band into the low Hamiltonian
without destroying locality.  It will define a finite normalized
transition-label partition subordinate to the V45 buffered cover, prove
an overlap cocycle identity, and derive the blockwise residual floor on
the complement of the enlarged slow space.  The first test is whether
the transition projector has fixed spatial range; if it is delocalized,
the V50 no-go requires an explicit finite-rank ledger rather than a
weighted-locality claim.
\end{decision}

\begin{firewall}[Claim boundary after seventh-series V70]
\label{fire:s7v70-claim}
V70 does not prove the regular-sector physical UCTBB data, transition-band
effective dynamics, the physical mixed product majorant, interacting
endpoint continuity, UCM, BCC, the pairing/Bogoliubov sector, weighted
physical-source incidence, patchwise coarse-action cocycles, terminal
strong-coupling matching, regulator independence, cutoff removal,
Osterwalder--Schrader reconstruction, or the full Yang--Mills mass gap.
\end{firewall}

\paragraph{Parent-version record.}
V70 was formed from
\texttt{Renewal\_Yang\_Mills\_Mass\_Gap\_Research\_Notes\_V69.tex}.
The V69 parent had \(30\,503\) logical source lines, \(1\,091\,306\)
bytes, and SHA--256
\[
 \texttt{0736C0ED51E0C1262E217EB478591C7A141B7FF19E34F463113F041B6983045D}.
\]
The next compulsory companion audit is V75.

% =============================================================================
% END APPENDED SEVENTH-SERIES V70 MATERIAL
% =============================================================================

% =============================================================================
% BEGIN APPENDED SEVENTH-SERIES V71 MATERIAL
% =============================================================================

\chapter{Transition-atlas range constraints and compatible slow projections}
\label{ch:s7v71}

\section{Purpose}
\label{sec:s7v71-purpose}

V70 retains a fixed transition band of boundary labels.  V65 already
shows that a partition-of-unity label is a normalized probabilistic
spectator, and V45 gives local buffered spectral projections.  Those
facts do not by themselves produce one orthogonal slow projection on the
physical Hilbert space.  The labelled atlas is a redundant tight frame;
its image satisfies a nontrivial range constraint, and a diagonal family
of patch projections preserves that constraint only when the overlap
cocycles intertwine the projections.

This chapter proves the exact range projector, quantifies the
incompatibility, constructs the pointwise joined slow space on transition
overlaps, and proves that its residual curvature floor survives.  It also
derives the exact covariant IMS debit created when boundary labels become
dynamical.

\section{The labelled atlas is an isometry with a range constraint}
\label{sec:s7v71-frame}

Let \(\Xi\) be the compact boundary-label carrier, with measure \(d\nu\),
and let \({\cal H}_\xi\) be a finite-cutoff block Hilbert fibre.  Choose a
finite cover \(\{U_i\}_{i=1}^N\), a smooth partition of unity
\begin{equation}
 \chi_i\geq0,\qquad
 \operatorname{supp}\chi_i\subset U_i,\qquad
 \sum_{i=1}^N\chi_i=1,                                \label{eq:s7v71-POU}
\end{equation}
and unitary trivializations
\[
 V_i(\xi):{\cal K}_i\longrightarrow{\cal H}_\xi.
\]
On overlaps define
\begin{equation}
 G_{ij}(\xi):=V_i(\xi)^*V_j(\xi):
 {\cal K}_j\longrightarrow{\cal K}_i.                 \label{eq:s7v71-G}
\end{equation}
Then
\begin{equation}
 G_{ii}=I,\qquad G_{ij}G_{jk}=G_{ik},\qquad
 G_{ij}^*=G_{ji}.                                     \label{eq:s7v71-cocycle}
\end{equation}

Let
\[
 {\cal H}:=\int_\Xi^\oplus{\cal H}_\xi\,d\nu(\xi),
 \qquad
 {\cal H}_{\rm lab}:=
 \bigoplus_{i=1}^N
 L^2(U_i,d\nu;{\cal K}_i),
\]
where components are extended by zero outside \(U_i\).  Define
\begin{equation}
 ({\cal W}\psi)_i(\xi)
 :=
 \sqrt{\chi_i(\xi)}\,V_i(\xi)^*\psi(\xi).             \label{eq:s7v71-W}
\end{equation}

\begin{theorem}[Exact atlas range projector]
\label{thm:s7v71-range}
The analysis map \({\cal W}:{\cal H}\to{\cal H}_{\rm lab}\) is an
isometry.  Its range projection
\({\cal R}:={\cal W}{\cal W}^*\) has block kernel
\begin{equation}
 {\cal R}_{ij}(\xi)
 =
 \sqrt{\chi_i(\xi)\chi_j(\xi)}\,G_{ij}(\xi).           \label{eq:s7v71-R}
\end{equation}
In particular,
\begin{equation}
 {\cal R}^*={\cal R},\qquad {\cal R}^2={\cal R}.       \label{eq:s7v71-Rprojection}
\end{equation}
\end{theorem}

\begin{proof}
By unitarity and \eqref{eq:s7v71-POU},
\[
 \|{\cal W}\psi\|^2
 =
 \int_\Xi\sum_i\chi_i(\xi)\|\psi(\xi)\|^2\,d\nu(\xi)
 =
 \|\psi\|^2.
\]
Thus \({\cal W}^*{\cal W}=I\), and
\({\cal R}={\cal W}{\cal W}^*\) is automatically an orthogonal
projection.  Direct substitution gives \eqref{eq:s7v71-R}.  Alternatively,
its square has blocks
\begin{align*}
 ({\cal R}^2)_{ij}
 &=
 \sum_k
 \sqrt{\chi_i\chi_k}\,G_{ik}
 \sqrt{\chi_k\chi_j}\,G_{kj}\\
 &=
 \sqrt{\chi_i\chi_j}
 \left(\sum_k\chi_k\right)G_{ij}
 =
 {\cal R}_{ij},
\end{align*}
using \eqref{eq:s7v71-cocycle}.
\end{proof}

\begin{firewall}[Atlas labels are not independent Hilbert copies]
\label{fire:s7v71-range}
Replacing \({\cal H}\) by the full direct sum
\({\cal H}_{\rm lab}\) without imposing
\({\cal R}\Psi=\Psi\) creates spurious states.  Exact probabilistic
normalization in V65 proves \({\cal W}^*{\cal W}=I\); it does not make
\({\cal W}{\cal W}^*=I\).
\end{firewall}

\section{When patchwise slow projections glue}
\label{sec:s7v71-glue}

Let \(P_i(\xi)\) be an orthogonal slow projection on \({\cal K}_i\), and
put
\begin{equation}
 {\bf P}:=\operatorname{diag}(P_1,\ldots,P_N)
 \quad\hbox{on }{\cal H}_{\rm lab}.                   \label{eq:s7v71-boldP}
\end{equation}

\begin{theorem}[Cocycle compatibility criterion]
\label{thm:s7v71-compatible}
The labelled slow projection preserves the physical atlas range,
\begin{equation}
 [{\bf P},{\cal R}]=0,                                \label{eq:s7v71-commute}
\end{equation}
if and only if, almost everywhere on every active overlap,
\begin{equation}
 P_iG_{ij}=G_{ij}P_j
 \qquad\hbox{whenever }\chi_i\chi_j>0.                \label{eq:s7v71-intertwine}
\end{equation}
Under this condition,
\begin{equation}
 P_{\rm phys}:={\cal W}^*{\bf P}{\cal W}              \label{eq:s7v71-Pphys}
\end{equation}
is an orthogonal projection on \({\cal H}\), and
\begin{equation}
 {\bf P}{\cal W}={\cal W}P_{\rm phys}.                \label{eq:s7v71-intertwine-W}
\end{equation}
\end{theorem}

\begin{proof}
By \eqref{eq:s7v71-R},
\begin{equation}
 [{\bf P},{\cal R}]_{ij}
 =
 \sqrt{\chi_i\chi_j}
 \left(P_iG_{ij}-G_{ij}P_j\right).                   \label{eq:s7v71-comm-block}
\end{equation}
This proves the equivalence.  If the commutator vanishes, then
\begin{align*}
 P_{\rm phys}^2
 &=
 {\cal W}^*{\bf P}{\cal R}{\bf P}{\cal W}\\
 &=
 {\cal W}^*{\cal R}{\bf P}^2{\cal W}
 =
 {\cal W}^*{\bf P}{\cal W},
\end{align*}
because \({\cal R}{\cal W}={\cal W}\) and
\({\cal W}^*{\cal R}={\cal W}^*\).  Self-adjointness is immediate.
Finally,
\[
 {\cal W}P_{\rm phys}
 =
 {\cal R}{\bf P}{\cal W}
 =
 {\bf P}{\cal R}{\cal W}
 =
 {\bf P}{\cal W}.
\]
\end{proof}

\begin{corollary}[Quantified projection mismatch]
\label{cor:s7v71-mismatch}
Define
\begin{align}
 r_{\rm ov}
 &:=
 \sup_i\sum_j
 \left\|
 \sqrt{\chi_i\chi_j}
 (P_iG_{ij}-G_{ij}P_j)
 \right\|,                                            \label{eq:s7v71-row}\\
 c_{\rm ov}
 &:=
 \sup_j\sum_i
 \left\|
 \sqrt{\chi_i\chi_j}
 (P_iG_{ij}-G_{ij}P_j)
 \right\|.                                            \label{eq:s7v71-column}
\end{align}
Then
\begin{equation}
 \|[{\bf P},{\cal R}]\|
 \leq\sqrt{r_{\rm ov}c_{\rm ov}}.                     \label{eq:s7v71-Schur}
\end{equation}
For the physical positive contraction
\(F:={\cal W}^*{\bf P}{\cal W}\),
\begin{equation}
 \|F^2-F\|
 \leq r_{\rm ov}c_{\rm ov}.                           \label{eq:s7v71-idempotency}
\end{equation}
\end{corollary}

\begin{proof}
Equation \eqref{eq:s7v71-Schur} is the operator-valued Schur test applied
to \eqref{eq:s7v71-comm-block}.  Since
\[
 F^2-F
 =
 {\cal W}^*{\bf P}({\cal R}-I){\bf P}{\cal W}
 =
 -\bigl((I-{\cal R}){\bf P}{\cal W}\bigr)^*
   \bigl((I-{\cal R}){\bf P}{\cal W}\bigr),
\]
and
\[
 (I-{\cal R}){\bf P}{\cal W}
 =
 (I-{\cal R})[{\bf P},{\cal R}]{\cal W},
\]
its norm is at most
\(\|[{\bf P},{\cal R}]\|^2\).  Apply
\eqref{eq:s7v71-Schur}.
\end{proof}

\begin{assessment}[Why a diagonal patch projection was insufficient]
\label{ass:s7v71-mismatch}
Unless \eqref{eq:s7v71-intertwine} holds, the synthesized operator \(F\)
is not a projection.  It therefore cannot be inserted as the low
projection of a Feshbach map merely because every \(P_i\) is a projection.
This is the Hilbert-space analogue of the V65 warning that normalized
chart labels do not erase transition data.
\end{assessment}

\section{Joined transition slow spaces}
\label{sec:s7v71-join}

Let \(L(\xi)\geq0\) be the physical finite-dimensional curvature matrix
on \({\cal H}_\xi\).  On \(U_i\), suppose
\[
 \widetilde P_i(\xi):=
 V_iP_iV_i^*
\]
is a spectral projection of \(L(\xi)\), and
\begin{equation}
 L(\xi)\geq\beta_i\bigl(I-\widetilde P_i(\xi)\bigr),
 \qquad \beta_i>0.                                    \label{eq:s7v71-patch-floor}
\end{equation}
For each \(\xi\), let
\begin{equation}
 I(\xi):=\{i:\chi_i(\xi)>0\},\qquad
 P_\vee(\xi):=
 \bigvee_{i\in I(\xi)}\widetilde P_i(\xi),\qquad
 Q_\vee:=I-P_\vee.                                    \label{eq:s7v71-join}
\end{equation}
Because the \(\widetilde P_i\) are spectral projections of the same
matrix, they commute and their join is the projection onto the sum of
their ranges.

\begin{theorem}[Joined overlap floor and exact compatibility]
\label{thm:s7v71-join}
Put
\begin{equation}
 \beta_*:=\min_{1\leq i\leq N}\beta_i>0.              \label{eq:s7v71-beta}
\end{equation}
Then, fibrewise,
\begin{equation}
 Q_\vee LQ_\vee\geq\beta_*Q_\vee.                     \label{eq:s7v71-join-floor}
\end{equation}
Define on every active chart
\begin{equation}
 P_i^\vee:=V_i^*P_\vee V_i.                           \label{eq:s7v71-Pivee}
\end{equation}
The family \(P_i^\vee\) satisfies
\begin{equation}
 P_i^\vee G_{ij}=G_{ij}P_j^\vee                      \label{eq:s7v71-join-compatible}
\end{equation}
on every active overlap and therefore gives one physical orthogonal
projection through \Cref{thm:s7v71-compatible}.
\end{theorem}

\begin{proof}
For every active \(i\), \(P_\vee\geq\widetilde P_i\), so
\(\operatorname{Ran}Q_\vee\subset
\operatorname{Ran}(I-\widetilde P_i)\).  Since all these projections
commute with \(L\), \eqref{eq:s7v71-patch-floor} restricted to
\(\operatorname{Ran}Q_\vee\) gives
\[
 Q_\vee LQ_\vee\geq\beta_iQ_\vee\geq\beta_*Q_\vee.
\]
For compatibility, use \(G_{ij}=V_i^*V_j\):
\[
 P_i^\vee G_{ij}
 =
 V_i^*P_\vee V_j
 =
 G_{ij}P_j^\vee.
\]
Apply \Cref{thm:s7v71-compatible}.
\end{proof}

\begin{corollary}[No block-volume growth of the joined rank]
\label{cor:s7v71-rank}
If the finite block fibre has dimension \(d_b\), then
\begin{equation}
 \operatorname{rank}P_\vee(\xi)\leq d_b              \label{eq:s7v71-rank}
\end{equation}
uniformly in volume and cutoff labels which do not change the chosen
finite block fibre.  The join introduces no spatial support beyond that
block and its fixed collar.
\end{corollary}

\begin{proof}
The join is a subspace of one \(d_b\)-dimensional fibre.  Every constituent
projection acts on the same fixed block variables.
\end{proof}

\begin{firewall}[Rank and regularity are different]
\label{fire:s7v71-rank}
The rank bound is local and finite, but \(P_\vee(\xi)\) may jump when the
active label set \(I(\xi)\) changes.  Measurability is enough for a
fibrewise form decomposition; derivatives in a dynamical boundary label
require a smooth refinement or must pay the IMS debit below.  Also, if
\operatorname{rank}P_\vee=d_b\), the residual floor is vacuous and all
block dynamics has been moved into the slow sector.
\end{firewall}

\section{Exact covariant atlas IMS identity}
\label{sec:s7v71-IMS}

Assume now that \(\xi\) is a smooth dynamical variable.  Put
\begin{equation}
 {\cal A}_i:=V_i^*dV_i,\qquad
 D_i:=d+{\cal A}_i.                                   \label{eq:s7v71-connection}
\end{equation}
For \(\psi(\xi)\in{\cal H}_\xi\), write
\(\psi_i=V_i^*\psi\).  Then
\begin{equation}
 D_i\psi_i=V_i^*d\psi.                                \label{eq:s7v71-covariant}
\end{equation}

\begin{theorem}[Covariant atlas IMS formula]
\label{thm:s7v71-IMS}
For every smooth compactly supported fibre section,
\begin{equation}
 \sum_i
 \left\|
 D_i\bigl(\sqrt{\chi_i}\,\psi_i\bigr)
 \right\|_{L^2}^2
 =
 \|d\psi\|_{L^2}^2
 \int_\Xi
 \left(\sum_i|d\sqrt{\chi_i}|^2\right)
 \|\psi(\xi)\|^2\,d\nu(\xi).                          \label{eq:s7v71-IMS}
\end{equation}
The overlap coordinate changes contribute through the covariant
connections and create no additional bulk term.
\end{theorem}

\begin{proof}
Differentiate \(V_i^*\psi\).  Since
\(dV_i^*\,V_i=-V_i^*dV_i=-{\cal A}_i\), one obtains
\eqref{eq:s7v71-covariant}.  Therefore
\[
 D_i(\sqrt{\chi_i}\psi_i)
 =
 d\sqrt{\chi_i}\,\psi_i
 +\sqrt{\chi_i}\,V_i^*d\psi.
\]
Sum the squared norms.  The pure second terms give
\(\|d\psi\|^2\) by \(\sum_i\chi_i=1\).  The cross term is proportional
to
\[
 \sum_i\sqrt{\chi_i}\,d\sqrt{\chi_i}
 ={1\over2}d\sum_i\chi_i=0.
\]
The remaining pure first terms are exactly the last term of
\eqref{eq:s7v71-IMS}.
\end{proof}

\begin{corollary}[Finite atlas debit]
\label{cor:s7v71-IMS}
For a fixed smooth finite atlas,
\begin{equation}
 C_{\rm IMS}:=
 \sup_{\xi\in\Xi}\sum_i|d\sqrt{\chi_i(\xi)}|^2<\infty, \label{eq:s7v71-CIMS}
\end{equation}
and the localization debit in \eqref{eq:s7v71-IMS} is at most
\(C_{\rm IMS}\|\psi\|^2\).
\end{corollary}

\begin{proof}
The displayed sum is continuous on the compact carrier when the square
roots are chosen smooth on their supports, so it has a finite maximum.
\end{proof}

\begin{firewall}[A finite IMS constant need not be small]
\label{fire:s7v71-IMS}
The physical electric kinetic term multiplies
\eqref{eq:s7v71-IMS} by its scale-dependent coefficient.  Finiteness of
\(C_{\rm IMS}\) does not show that the resulting energy debit is small
relative to the transition or mass margin.  The atlas width and coupling
schedule must be inserted before the debit can enter UCM or BCC.
\end{firewall}

\section{Status and next acquisition}
\label{sec:s7v71-status}

\begin{theorem}[What V71 proves]
\label{thm:s7v71-result}
V71 proves the exact labelled-atlas range projector
\eqref{eq:s7v71-R}; the necessary and sufficient overlap-cocycle
condition for patchwise slow projections to define one physical
projection; the Schur bound on projection mismatch; the joined transition
projection with residual floor \eqref{eq:s7v71-join-floor}; and the exact
covariant IMS identity \eqref{eq:s7v71-IMS}.  Thus the retained V70
transition labels can be loaded into a compatible local slow space
without reducing the already available residual curvature floor.
\end{theorem}

\begin{proof}
Use
\Cref{thm:s7v71-range,thm:s7v71-compatible,
cor:s7v71-mismatch,thm:s7v71-join,
thm:s7v71-IMS}.
\end{proof}

\begin{assessment}[Progress at seventh-series V71]
\label{ass:s7v71-status}
The transition atlas no longer hides either duplicate Hilbert states or
incompatible Feshbach projections.  All overlap modes are joined before
the residual complement is taken, so the floor is monotone under
retention.  Because the join is confined to one fixed block, the V50
delocalized-global-projector obstruction is not triggered at this stage.
The live quantitative issue is the scale of the atlas IMS debit and the
interacting dynamics on the joined slow space.
\end{assessment}

\begin{decision}[V72 acquisition order]
\label{dec:s7v71-V72}
V72 will insert the Kogut--Susskind electric coefficient and the
weak-coupling block scaling into \eqref{eq:s7v71-IMS}.  It will compare a
fixed transition width, a coupling-dependent width, and a soft spectral
frame, and determine which choice keeps the IMS debit below the residual
oscillator margin.  If no shrinking width is compatible, the transition
coordinate will remain an unlocalized slow variable and its dynamics will
be assigned to the finite-rank low Hamiltonian.
\end{decision}

\begin{firewall}[Claim boundary after seventh-series V71]
\label{fire:s7v71-claim}
V71 does not make the atlas IMS debit small, prove the joined slow-sector
interacting gap, verify the remaining regular-sector UCTBB data, establish
the physical mixed product majorant, interacting endpoint continuity,
UCM, BCC, the pairing/Bogoliubov sector, weighted physical-source
incidence, patchwise coarse-action cocycles beyond the exact coordinate
identity, terminal strong-coupling matching, regulator independence,
cutoff removal, Osterwalder--Schrader reconstruction, or the full
Yang--Mills mass gap.
\end{firewall}

\paragraph{Parent-version record.}
V71 was formed from
\texttt{Renewal\_Yang\_Mills\_Mass\_Gap\_Research\_Notes\_V70.tex}.
The V70 parent had \(30\,977\) logical source lines, \(1\,108\,727\)
bytes, and SHA--256
\[
 \texttt{C34491AFB920926E9767B4D593951EBCC73C6B51FE45D9577991FC0717B780F2}.
\]
The next compulsory companion audit is V75.

% =============================================================================
% END APPENDED SEVENTH-SERIES V71 MATERIAL
% =============================================================================

% =============================================================================
% BEGIN APPENDED SEVENTH-SERIES V72 MATERIAL
% =============================================================================

\chapter{Kogut--Susskind scaling of transition-atlas and projector response}
\label{ch:s7v72}

\section{Purpose and normalization}
\label{sec:s7v72-purpose}

V71 isolates two derivative costs which were absent from the purely
fibrewise transition construction: localization derivatives of the atlas
partition and parameter derivatives of the buffered slow projection.
This chapter inserts their physical weak-coupling scale.

Use the dimensionless Kogut--Susskind convention
\begin{equation}
 H_{\Lambda,a,g}
 =
 {g^2\over2a}\sum_{e\in E_\Lambda}(-\Delta_e)
 +{1\over ag^2}\sum_{p\in{\cal P}_\Lambda}
 \bigl(3-\operatorname{ReTr}W_p\bigr).                \label{eq:s7v72-KS}
\end{equation}
Fixed representation-dependent constants may be absorbed into the
positive constants below.  In a weak-field chart \(U_e=\exp(gA_e)\), the
factor \(g^2\) in the electric term cancels the \(g^{-2}\) created by
the coordinate derivative, while the magnetic Taylor coefficient
similarly cancels the explicit \(g^{-2}\).  The local residual oscillator
therefore has an energy floor of order \(a^{-1}\).

\section{Atlas-width IMS debit}
\label{sec:s7v72-IMS}

Let \(\eta\) denote the finite collection of dynamical boundary-label
links.  Assume its electric form has coefficient bounded above by
\begin{equation}
 {\cal E}_{\eta}(\psi)
 \leq {c_Eg^2\over2a}\|\nabla_\eta\psi\|_2^2          \label{eq:s7v72-electric-upper}
\end{equation}
in the declared compact charts, with \(c_E\) independent of volume.
Let a smooth atlas partition of width \(w>0\) satisfy
\begin{equation}
 \sup_\eta\sum_i
 |\nabla_\eta\sqrt{\chi_i(\eta)}|^2
 \leq{C_\chi\over w^2}.                               \label{eq:s7v72-width}
\end{equation}

\begin{theorem}[Physical atlas IMS scale]
\label{thm:s7v72-IMS}
The exact V71 covariant atlas decomposition has localization debit
\begin{equation}
 0\leq V_{\rm IMS}
 \leq
 {c_EC_\chi g^2\over2aw^2}\,I.                        \label{eq:s7v72-IMS}
\end{equation}
If the joined residual sector obeys
\begin{equation}
 Q_\vee h_{\rm r}Q_\vee
 \geq {c_{\rm r}\over a}Q_\vee,\qquad c_{\rm r}>0,    \label{eq:s7v72-residual}
\end{equation}
then on that sector
\begin{equation}
 Q_\vee(h_{\rm r}-V_{\rm IMS})Q_\vee
 \geq
 {c_{\rm r}\over a}(1-\theta_{\rm IMS})Q_\vee,        \label{eq:s7v72-absorbed}
\end{equation}
where
\begin{equation}
 \theta_{\rm IMS}
 :=
 {c_EC_\chi\over2c_{\rm r}}
 \left({g\over w}\right)^2.                           \label{eq:s7v72-theta}
\end{equation}
\end{theorem}

\begin{proof}
Multiply the last term of the exact identity
\eqref{eq:s7v71-IMS} by the electric coefficient in
\eqref{eq:s7v72-electric-upper} and use
\eqref{eq:s7v72-width}.  This proves
\eqref{eq:s7v72-IMS}.  Restrict the scalar upper bound to
\(\operatorname{Ran}Q_\vee\), subtract it from
\eqref{eq:s7v72-residual}, and factor out \(c_{\rm r}/a\).
\end{proof}

\begin{corollary}[The three width regimes]
\label{cor:s7v72-widths}
The relative debit has the following behavior.
\begin{enumerate}[label=\textup{(\roman*)},leftmargin=3.5em]
\item For fixed \(w=w_0>0\),
\begin{equation}
 \theta_{\rm IMS}=O(g^2)\longrightarrow0.             \label{eq:s7v72-fixed}
\end{equation}
\item For a weak-field width \(w=gR\),
\begin{equation}
 \theta_{\rm IMS}=O(R^{-2}),                          \label{eq:s7v72-weak-width}
\end{equation}
so any \(R\to\infty\) closes the debit.
\item If \(w/g\to0\), then the bound diverges and supplies no residual
      floor.
\end{enumerate}
\end{corollary}

\begin{proof}
Insert the three width schedules into
\eqref{eq:s7v72-theta}.
\end{proof}

\begin{remark}[Relation to the earlier weak-field IMS estimate]
\label{rem:s7v72-V38}
The choice \(w=gR\) reproduces the \(R^{-2}\) scaling already obtained for
the V38 compact-link weak-field cutoff.  The new point is that the same
scaling governs the \emph{boundary-label transition atlas}; the two IMS
rows must be listed separately if both partitions are used.
\end{remark}

\section{Buffered Riesz-projector response}
\label{sec:s7v72-projector}

Let \(L(\eta)\) be a \(C^1\) finite-dimensional Hermitian curvature
matrix on one transition patch.  Suppose a contour \(\Gamma\) encloses
the retained slow spectrum and obeys
\begin{equation}
 \operatorname{dist}
 \bigl(\Gamma,\operatorname{spec}L(\eta)\bigr)
 \geq d_{\rm sp}>0                                    \label{eq:s7v72-distance}
\end{equation}
uniformly on the patch.  The buffered Riesz projection is
\begin{equation}
 P(\eta)
 =
 {1\over2\pi i}\oint_\Gamma
 (z-L(\eta))^{-1}\,dz.                                \label{eq:s7v72-Riesz}
\end{equation}

\begin{lemma}[Riesz derivative bound]
\label{lem:s7v72-Riesz}
For every tangent derivative \(\partial\),
\begin{equation}
 \|\partial P(\eta)\|
 \leq
 {|\Gamma|\over2\pi d_{\rm sp}^2}
 \|\partial L(\eta)\|.                                \label{eq:s7v72-Riesz-bound}
\end{equation}
\end{lemma}

\begin{proof}
Differentiate \eqref{eq:s7v72-Riesz} under the finite contour integral:
\[
 \partial P
 =
 {1\over2\pi i}\oint_\Gamma
 (z-L)^{-1}(\partial L)(z-L)^{-1}\,dz.
\]
Each resolvent norm is at most \(d_{\rm sp}^{-1}\).  Integrating its norm
over the contour gives \eqref{eq:s7v72-Riesz-bound}.
\end{proof}

\begin{lemma}[Exact vertical derivative of a slow section]
\label{lem:s7v72-vertical}
Let \(P(\eta)\) be a differentiable orthogonal projection,
\(Q=I-P\), and let \(P\psi=\psi\).  Then
\begin{equation}
 Q\,\partial\psi
 =
 Q(\partial P)\psi,                                   \label{eq:s7v72-vertical}
\end{equation}
and therefore
\begin{equation}
 \|Q\,\partial\psi\|
 \leq\|\partial P\|\,\|\psi\|.                        \label{eq:s7v72-vertical-bound}
\end{equation}
\end{lemma}

\begin{proof}
Differentiate \(P\psi=\psi\):
\[
 (\partial P)\psi+P\partial\psi=\partial\psi.
\]
Apply \(Q\) and use \(QP=0\).
\end{proof}

\begin{theorem}[Weak-coupling projector-response coefficient]
\label{thm:s7v72-response}
Assume
\begin{equation}
 \sup_\eta\sum_\alpha
 \|\partial_\alpha L(\eta)\|^2\leq C_L^2             \label{eq:s7v72-CL}
\end{equation}
over the finite boundary-label tangent frame.  Then the electric energy
of the high-component derivative generated by a normalized slow section
is bounded by
\begin{equation}
 {c_Eg^2\over2a}
 \sum_\alpha\|Q\partial_\alpha\psi\|^2
 \leq
 {c_Eg^2\over2a}
 \left({|\Gamma|C_L\over2\pi d_{\rm sp}^2}\right)^2
 \|\psi\|^2.                                          \label{eq:s7v72-response}
\end{equation}
Relative to the residual floor \(c_{\rm r}/a\), its coefficient is
\begin{equation}
 \theta_{\rm P}
 :=
 {c_Eg^2\over2c_{\rm r}}
 \left({|\Gamma|C_L\over2\pi d_{\rm sp}^2}\right)^2.  \label{eq:s7v72-thetaP}
\end{equation}
For fixed spectral buffer and \(C_L=O(1)\),
\(\theta_{\rm P}=O(g^2)\to0\).
\end{theorem}

\begin{proof}
Apply \Cref{lem:s7v72-vertical} in every tangent direction, then use
\Cref{lem:s7v72-Riesz} and
\eqref{eq:s7v72-CL}.  Divide the resulting coefficient by
\(c_{\rm r}/a\).
\end{proof}

\begin{corollary}[Constraint on a shrinking spectral buffer]
\label{cor:s7v72-shrinking}
If \(|\Gamma|\) and \(C_L\) remain bounded while
\(d_{\rm sp}=d_{\rm sp}(g)\downarrow0\), the displayed estimate closes
only if
\begin{equation}
 {g\over d_{\rm sp}(g)^2}\longrightarrow0.            \label{eq:s7v72-shrinking}
\end{equation}
In particular, a buffer shrinking like \(g^\alpha\) is allowed by this
bound only for \(\alpha<1/2\).
\end{corollary}

\begin{proof}
Equation \eqref{eq:s7v72-thetaP} is of order
\(g^2d_{\rm sp}^{-4}\).  This tends to zero precisely under the stated
sufficient condition; insert \(d_{\rm sp}=g^\alpha\).
\end{proof}

\begin{firewall}[Projector response is not an interacting gap theorem]
\label{fire:s7v72-response}
\Cref{thm:s7v72-response} controls the geometric high component created
when a slow basis moves with the boundary label.  It does not estimate
the nonlinear Wilson vertices inside the joined slow sector or prove
that its effective Hamiltonian has a mass gap.  A discontinuous
\(P_\vee\) from V71 is not differentiable; it must first be represented
by fixed buffered patch projections and a smooth labelled refinement.
\end{firewall}

\section{A scale-compatible transition card}
\label{sec:s7v72-card}

\begin{definition}[Smooth atlas transition card, SATC]
\label{def:s7v72-SATC}
SATC consists of:
\begin{enumerate}[label=\textup{(A\arabic*)},leftmargin=4em]
\item the exact labelled range projector and cocycle compatibility of
      V71;
\item a smooth partition satisfying
      \eqref{eq:s7v72-width};
\item a joined residual floor \(c_{\rm r}/a\);
\item buffered Riesz contours satisfying
      \eqref{eq:s7v72-distance,eq:s7v72-CL};
\item the strict total geometric budget
\begin{equation}
 \theta_{\rm atlas}:=
 \theta_{\rm IMS}+\theta_{\rm P}<1;                   \label{eq:s7v72-budget}
\end{equation}
\item uniformity of all constants in block position, regulator, and
      retained transition label.
\end{enumerate}
\end{definition}

\begin{theorem}[Fixed-width transition atlases are scale compatible]
\label{thm:s7v72-SATC}
If \(w\), \(d_{\rm sp}\), \(C_\chi\), \(C_L\), \(c_E\), and
\(c_{\rm r}\) have positive regulator-uniform bounds, then asymptotic
freedom \(g(a)\to0\) makes
\begin{equation}
 \theta_{\rm atlas}=O(g(a)^2)\longrightarrow0.         \label{eq:s7v72-SATC}
\end{equation}
Hence the atlas localization and projector response preserve a positive
fraction of the block residual floor for sufficiently small \(a\).
\end{theorem}

\begin{proof}
Combine
\eqref{eq:s7v72-theta,eq:s7v72-thetaP} with the fixed uniform
constants.  For small enough \(a\), their sum is below one, and
\eqref{eq:s7v72-absorbed} preserves the declared fraction of the floor.
\end{proof}

\begin{firewall}[Uniform finite-cutoff constants remain physical data]
\label{fire:s7v72-uniform}
Compactness gives finite constants at each fixed block cutoff.  It does
not prove that \(C_L\), the fibre dimension, the contour length, or the
spectral separation remain uniform when occupation and continuum
regulators are removed.  SATC records precisely those missing
uniformities.
\end{firewall}

\section{Status and next acquisition}
\label{sec:s7v72-status}

\begin{theorem}[What V72 proves]
\label{thm:s7v72-result}
V72 proves that the transition-atlas IMS debit is relatively
\(O((g/w)^2)\), and the buffered-projector response is relatively
\(O(g^2d_{\rm sp}^{-4})\).  A fixed-width atlas with fixed spectral
buffer therefore costs only \(O(g^2)\) against an \(a^{-1}\) residual
oscillator floor.  The weak-field width \(w=gR\) reproduces the
\(O(R^{-2})\) schedule, while widths \(w\ll g\) are rejected.
\end{theorem}

\begin{proof}
Use
\Cref{thm:s7v72-IMS,cor:s7v72-widths,
thm:s7v72-response,cor:s7v72-shrinking,
thm:s7v72-SATC}.
\end{proof}

\begin{assessment}[Progress at seventh-series V72]
\label{ass:s7v72-status}
The retained transition band is not, by itself, an ultraviolet
obstruction.  With a fixed local atlas and spectral buffer, both of its
geometric kinetic costs vanish relatively by asymptotic freedom.  The
live issue has moved to the interacting slow Hamiltonian and to uniform
control as its finite occupation cutoff is removed.
\end{assessment}

\begin{decision}[V73 acquisition order]
\label{dec:s7v72-V73}
V73 will return to the V66 mixed product majorant.  It will use the exact
normalized atlas range constraint, V70 bad-set domination, and the V64
replica activity to derive a labelled polymer convolution theorem.  The
target is to show that transition labels contribute through normalized
conditional sums and local cocycles, not an overlap-multiplicity factor.
Any unbounded Jacobian or projector-response insertion will be placed in
the rooted activity rather than hidden in the probability weight.
\end{decision}

\begin{firewall}[Claim boundary after seventh-series V72]
\label{fire:s7v72-claim}
V72 does not prove regulator-uniform SATC, the interacting joined
slow-sector gap, the remaining regular-sector UCTBB data, the physical
mixed product majorant, interacting endpoint continuity, UCM, BCC, the
pairing/Bogoliubov sector, weighted physical-source incidence, nonlinear
patchwise coarse-action cocycles, terminal strong-coupling matching,
regulator independence, cutoff removal, Osterwalder--Schrader
reconstruction, or the full Yang--Mills mass gap.
\end{firewall}

\paragraph{Parent-version record.}
V72 was formed from
\texttt{Renewal\_Yang\_Mills\_Mass\_Gap\_Research\_Notes\_V71.tex}.
The V71 parent had \(31\,463\) logical source lines, \(1\,124\,545\)
bytes, and SHA--256
\[
 \texttt{29C05C574508170CEF2C314FC004D6DE9EBEE91487EBA4B093BBFE5823FC2C8F}.
\]
The next compulsory companion audit is V75.

% =============================================================================
% END APPENDED SEVENTH-SERIES V72 MATERIAL
% =============================================================================

% =============================================================================
% BEGIN APPENDED SEVENTH-SERIES V73 MATERIAL
% =============================================================================

\chapter{Normalized labelled \(L^2\) assembly of mixed good/bad polymers}
\label{ch:s7v73}

\section{Purpose}
\label{sec:s7v73-purpose}

V66 leaves the mixed product majorant as a separate physical interface.
V70 now supplies arbitrary-set domination of every unretained bad or
robust-singular block, and V71 supplies an exact normalized transition
atlas.  Multiplying a good-field estimate by a bad-event probability
pointwise would still be unjustified.  The correct interface is a local
\(L^2\) activity: conditional Jensen removes the atlas-label
multiplicity, and Cauchy--Schwarz combines that activity with the
bad-set probability at the cost of one square root.

This chapter proves that compiler and the resulting connected-animal
sum.  It turns the V66 product majorant from a pointwise hypothesis into
a labelled local \(L^2\) hypothesis which is aligned with the V63--V64
replica machinery.

\section{Normalized label kernels and conditional Jensen}
\label{sec:s7v73-labels}

Let \((\Omega,\mu)\) be a probability space.  For a finite polymer
support \(\Gamma\), let \({\cal A}_\Gamma\) be its finite atlas-label
set and let
\begin{equation}
 \kappa_\Gamma(\alpha\mid U)\geq0,\qquad
 \sum_{\alpha\in{\cal A}_\Gamma}
 \kappa_\Gamma(\alpha\mid U)=1                       \label{eq:s7v73-kernel}
\end{equation}
for \(\mu\)-almost every \(U\).  The V65 product partition of unity is
the principal example.  For a labelled activity
\(G(U,\alpha)\), define
\begin{align}
 \overline G(U)
 &:=
 \sum_\alpha\kappa_\Gamma(\alpha\mid U)G(U,\alpha),
                                                               \label{eq:s7v73-Gbar}\\
 \|G\|_{2,\kappa}^2
 &:=
 \int_\Omega
 \sum_\alpha\kappa_\Gamma(\alpha\mid U)
 |G(U,\alpha)|^2\,d\mu(U).                            \label{eq:s7v73-L2}
\end{align}

\begin{lemma}[Conditional label contraction]
\label{lem:s7v73-Jensen}
\begin{equation}
 \|\overline G\|_{L^2(\mu)}
 \leq\|G\|_{2,\kappa}.                                \label{eq:s7v73-Jensen}
\end{equation}
The constant is one and is independent of the number of atlas labels or
overlap multiplicity.
\end{lemma}

\begin{proof}
For each fixed \(U\), Jensen's inequality for the probability vector
\(\kappa_\Gamma(\cdot\mid U)\) gives
\[
 |\overline G(U)|^2
 \leq
 \sum_\alpha\kappa_\Gamma(\alpha\mid U)|G(U,\alpha)|^2.
\]
Integrate over \(U\).
\end{proof}

\begin{corollary}[Unitary cocycles cost no label factor]
\label{cor:s7v73-unitary}
Suppose every chart transition appearing in \(G(U,\alpha)\) is a unitary
cocycle \(G_{ij}\) from \eqref{eq:s7v71-cocycle}.  Replacing any labelled
vector factor by its cocycle transport leaves its pointwise norm, and
hence \(\|G\|_{2,\kappa}\), unchanged.
\end{corollary}

\begin{proof}
Each \(G_{ij}\) is unitary.  Apply this inside every summand of
\eqref{eq:s7v73-L2}.
\end{proof}

\begin{firewall}[Coordinate densities are activities, not label
multiplicities]
\label{fire:s7v73-Jacobian}
If a chart calculation changes variables and introduces a nonunit
Haar/Faddeev--Popov density, that density belongs inside \(G\) and its
\(L^2\) norm.  Normalization of \(\kappa_\Gamma\) removes only the
combinatorial overlap sum; it does not bound an unbounded Jacobian.
\end{firewall}

\section{The mixed \(L^2\)-probability compiler}
\label{sec:s7v73-mixed}

For each block \(b\), let \(E_b\) be its unretained bad-or-robust-singular
event.  Assume the arbitrary-set domination
\begin{equation}
 \mu\left(\bigcap_{b\in B}E_b\right)
 \leq\pi_M^{|B|}
 \qquad(B\ \hbox{finite}),                            \label{eq:s7v73-event-dom}
\end{equation}
where V70 gives the admissible choice
\begin{equation}
 \pi_M
 :=
 \left(p_M+s_{\beta,\tau}^{(0)}\right)^{1/K}.         \label{eq:s7v73-pi}
\end{equation}

Let \(\varnothing\ne B\subset\Gamma\) be the marked bad set and let
\(G_{\Gamma,B}(U,\alpha)\) contain every good-field vertex, transition
cocycle, response insertion, and local counterterm assigned to that
marked polymer.  Define its physical weight by
\begin{equation}
 w(\Gamma,B)
 :=
 \int
 {\bf1}_{\cap_{b\in B}E_b}(U)
 \sum_\alpha
 \kappa_\Gamma(\alpha\mid U)
 G_{\Gamma,B}(U,\alpha)\,d\mu(U).                     \label{eq:s7v73-weight}
\end{equation}

\begin{theorem}[Labelled \(L^2\) mixed product theorem]
\label{thm:s7v73-mixed}
If
\begin{equation}
 \|G_{\Gamma,B}\|_{2,\kappa}
 \leq q_M^{|\Gamma\setminus B|}                      \label{eq:s7v73-activity}
\end{equation}
for every connected \(\Gamma\) and nonempty \(B\subset\Gamma\), then
\begin{equation}
 |w(\Gamma,B)|
 \leq
 q_M^{|\Gamma\setminus B|}
 b_M^{|B|},
 \qquad
 b_M:=\pi_M^{1/2}.                                    \label{eq:s7v73-product}
\end{equation}
\end{theorem}

\begin{proof}
Apply Cauchy--Schwarz in \(L^2(\mu)\) to
\eqref{eq:s7v73-weight}:
\[
 |w(\Gamma,B)|
 \leq
 \|\overline G_{\Gamma,B}\|_2
 \mu\left(\bigcap_{b\in B}E_b\right)^{1/2}.
\]
Use
\Cref{lem:s7v73-Jensen}, \eqref{eq:s7v73-activity}, and
\eqref{eq:s7v73-event-dom}.  This gives
\[
 |w(\Gamma,B)|
 \leq
 q_M^{|\Gamma\setminus B|}
 \pi_M^{|B|/2},
\]
which is \eqref{eq:s7v73-product}.
\end{proof}

\begin{remark}[The square-root loss is harmless]
\label{rem:s7v73-root}
If \(p_M=O(M^qe^{-cM})\) and
\(s_{\beta,\tau}^{(0)}=O(\beta^{q'}e^{-c'\beta})\), then
\(b_M\to0\) exponentially in the smaller of \(M\) and \(\beta\), after
only the fixed coloring and square-root losses.  Every fixed polynomial
occupation factor is still absorbed.
\end{remark}

\begin{corollary}[A pointwise sufficient activity row]
\label{cor:s7v73-pointwise}
The \(L^2\) condition \eqref{eq:s7v73-activity} follows if
\begin{equation}
 \sum_\alpha\kappa_\Gamma(\alpha\mid U)
 |G_{\Gamma,B}(U,\alpha)|^2
 \leq q_M^{2|\Gamma\setminus B|}                      \label{eq:s7v73-pointwise}
\end{equation}
for \(\mu\)-almost every \(U\).
\end{corollary}

\begin{proof}
Integrate \eqref{eq:s7v73-pointwise} and use
\(\mu(\Omega)=1\).
\end{proof}

\section{Connected mixed-polymer summation}
\label{sec:s7v73-sum}

\begin{theorem}[Rooted labelled mixed-polymer bound]
\label{thm:s7v73-sum}
Let the block graph have maximum degree \(\Delta\), and assume
\eqref{eq:s7v73-activity}.  If
\begin{equation}
 \zeta_M^{(2)}
 :=
 \Delta^2e^\alpha(q_M+b_M)<1,                         \label{eq:s7v73-zeta}
\end{equation}
then
\begin{equation}
 \sup_x
 \sum_{\substack{\Gamma\ni x\\\Gamma\ {\rm connected}}}
 e^{\alpha|\Gamma|}
 \sum_{\varnothing\ne B\subset\Gamma}
 |w(\Gamma,B)|
 \leq
 {e^\alpha b_M\over(1-\zeta_M^{(2)})^2}.              \label{eq:s7v73-sum}
\end{equation}
\end{theorem}

\begin{proof}
By \Cref{thm:s7v73-mixed}, for a support of size \(n\),
\begin{align*}
 \sum_{\varnothing\ne B\subset\Gamma}|w(\Gamma,B)|
 &\leq
 \sum_{\varnothing\ne B\subset\Gamma}
 b_M^{|B|}q_M^{n-|B|}\\
 &=
 (q_M+b_M)^n-q_M^n\\
 &\leq
 nb_M(q_M+b_M)^{n-1}.
\end{align*}
There are at most \(\Delta^{2(n-1)}\) connected \(n\)-block supports
through a fixed root.  Sum
\[
 e^\alpha b_M
 \sum_{n\geq1}
 n\left[
 \Delta^2e^\alpha(q_M+b_M)
 \right]^{n-1}
 =
 {e^\alpha b_M\over(1-\zeta_M^{(2)})^2}.
\]
\end{proof}

\begin{corollary}[Vanishing mixed sector]
\label{cor:s7v73-vanishing}
If the V64 replica-admissible good activity satisfies \(q_M\to0\), the
V70 probability parameters satisfy
\(p_M+s_{\beta,\tau}^{(0)}\to0\), and
\eqref{eq:s7v73-activity} holds, then the labelled mixed expansion
converges for all sufficiently large \(M,\beta\), and its rooted
bad-containing weight is \(O(b_M)\).
\end{corollary}

\begin{proof}
Equation \eqref{eq:s7v73-pi} gives \(b_M\to0\).  Thus
\(\zeta_M^{(2)}\to0\); apply \Cref{thm:s7v73-sum}.
\end{proof}

\section{Relation to the replica activity}
\label{sec:s7v73-replica}

\begin{definition}[Mixed labelled \(L^2\) activity card, ML2AC]
\label{def:s7v73-ML2AC}
ML2AC consists of:
\begin{enumerate}[label=\textup{(L\arabic*)},leftmargin=4em]
\item the exact normalized label kernel
      \eqref{eq:s7v73-kernel};
\item the physical event domination
      \eqref{eq:s7v73-event-dom};
\item the labelled local activity bound
      \eqref{eq:s7v73-activity}, uniformly in slow and transition labels;
\item every nonunit coordinate density, atlas derivative, and projector
      response insertion included in \(G_{\Gamma,B}\);
\item \(q_M+b_M\) satisfying \eqref{eq:s7v73-zeta}; and
\item the same assignment of boundary interactions in the good,
      bad, and transition representations, so no plaquette is omitted or
      counted twice.
\end{enumerate}
\end{definition}

\begin{theorem}[ML2AC closes the V66 mixed product interface]
\label{thm:s7v73-closure}
Under ML2AC, the V66 mixed product majorant holds with its bad activity
\(\overline p_M\) replaced by \(b_M=\pi_M^{1/2}\).  All mixed
bad-containing polymers obey \eqref{eq:s7v73-sum}; normalized atlas labels
and unitary overlap cocycles introduce no extra combinatorial factor.
\end{theorem}

\begin{proof}
Rows (L1)--(L4) give \Cref{thm:s7v73-mixed}; row (L5) gives
\Cref{thm:s7v73-sum}; row (L6) makes the formal polymer assignment an
exact decomposition of the physical interaction.
\end{proof}

\begin{assessment}[What V64 supplies and what it does not]
\label{ass:s7v73-V64}
The V64 glued-replica theorem is designed to prove \(L^2\) marginal
density bounds with constants proportional to observed support.  It
therefore supplies the correct analytic mechanism for the all-good part
of \eqref{eq:s7v73-activity}.  It has not yet been verified with every
compact transition density, atlas derivative, and mixed-boundary
assignment included.  ML2AC is strictly narrower than the old pointwise
product hypothesis, but its physical uniformity remains a real
calculation.
\end{assessment}

\begin{firewall}[A global \(L^2\) density norm is still forbidden]
\label{fire:s7v73-global}
Equation \eqref{eq:s7v73-activity} is a local polymer activity bound.
Replacing it by the \(L^2\) norm of the complete volume density can
produce an exponential volume factor and would undo the V63 replica
localization.  The support exponent must be
\(|\Gamma\setminus B|\), not the size of the ambient regulator.
\end{firewall}

\section{Status and next acquisition}
\label{sec:s7v73-status}

\begin{theorem}[What V73 proves]
\label{thm:s7v73-result}
V73 proves the normalized-label contraction with constant one, the
labelled \(L^2\) mixed product estimate
\eqref{eq:s7v73-product}, and the rooted connected sum
\eqref{eq:s7v73-sum}.  Consequently V66's pointwise mixed product
interface is replaced by ML2AC, a local \(L^2\) activity aligned with the
existing replica-pressure method.  The only probability loss is the
explicit square root \(b_M=\pi_M^{1/2}\).
\end{theorem}

\begin{proof}
Use
\Cref{lem:s7v73-Jensen,thm:s7v73-mixed,
thm:s7v73-sum,thm:s7v73-closure}.
\end{proof}

\begin{assessment}[Progress at seventh-series V73]
\label{ass:s7v73-status}
The compact good/bad assembly no longer requires an unexplained
pointwise multiplication of a good activity and a bad probability.
Normalized labels, unitary cocycles, and Cauchy--Schwarz give the exact
interface.  The live task is now to prove ML2AC for the complete
transition-labelled Wilson vertex family, including the V72 kinetic
insertions.
\end{assessment}

\begin{decision}[V74 acquisition order]
\label{dec:s7v73-V74}
V74 will build the two-replica representation of the ML2AC norm.  It
will glue two copies on the marked polymer support, identify which
transition-label and projector-response insertions cancel outside that
support, and derive a marked-cross-edge tree bound extending V64.  The
target is a concrete \(q_M\) containing the SATC debits rather than a new
abstract card.
\end{decision}

\begin{firewall}[Claim boundary after seventh-series V73]
\label{fire:s7v73-claim}
V73 does not verify ML2AC for the full physical Wilson vertex family,
prove regulator-uniform SATC, the interacting joined slow-sector gap,
the remaining regular-sector UCTBB data, interacting endpoint
continuity, UCM, BCC, the pairing/Bogoliubov sector, weighted
physical-source incidence, nonlinear coarse-action cocycles, terminal
strong-coupling matching, regulator independence, cutoff removal,
Osterwalder--Schrader reconstruction, or the full Yang--Mills mass gap.
\end{firewall}

\paragraph{Parent-version record.}
V73 was formed from
\texttt{Renewal\_Yang\_Mills\_Mass\_Gap\_Research\_Notes\_V72.tex}.
The V72 parent had \(31\,826\) logical source lines, \(1\,137\,603\)
bytes, and SHA--256
\[
 \texttt{EFA13B17F3E23A527C10A7969A124598727BFC5CC63AE45A7804EEBDB6CC37F2}.
\]
The next compulsory companion audit is V75.

% =============================================================================
% END APPENDED SEVENTH-SERIES V73 MATERIAL
% =============================================================================

% =============================================================================
% BEGIN APPENDED SEVENTH-SERIES V74 MATERIAL
% =============================================================================

\chapter{Replica-local density transfer to the mixed labelled \(L^2\) activity}
\label{ch:s7v74}

\section{Purpose}
\label{sec:s7v74-purpose}

V73 reduces the physical mixed-polymer interface to the local norm
\(\|G_{\Gamma,B}\|_{2,\kappa}\).  V63--V64 control a different but
compatible object: the \(L^2\) norm of the physical marginal density
relative to a gapped reference measure.  This chapter proves the exact
H\"older bridge.  A labelled fourth moment under the reference measure,
together with the replica marginal-density bound, yields the required
physical labelled second moment.

The density cost is exponential only in the polymer support and can be
distributed once per block.  Transition or bad-block insertions may carry
a bounded or polynomial factor \(t_M\); the V70--V73 bad activity remains
exponentially small after that factor.

\section{Local density and labelled reference moments}
\label{sec:s7v74-local}

Let \(S\) be a finite set of reference variables.  Denote the reference
and physical marginals by \(\nu_S,\mu_S\), and suppose
\begin{equation}
 d\mu_S=r_S\,d\nu_S,\qquad
 \|r_S\|_{L^2(\nu_S)}
 \leq e^{K_2|S|}.                                     \label{eq:s7v74-LMDC}
\end{equation}
For a normalized local label kernel
\(\kappa(\alpha\mid x)\), define
\begin{equation}
 \|G\|_{4,\kappa,\nu_S}
 :=
 \left[
 \int
 \sum_\alpha\kappa(\alpha\mid x)|G(x,\alpha)|^4
 \,d\nu_S(x)
 \right]^{1/4}.                                       \label{eq:s7v74-L4}
\end{equation}

\begin{theorem}[Replica-density to labelled-activity transfer]
\label{thm:s7v74-transfer}
Under \eqref{eq:s7v74-LMDC},
\begin{equation}
 \|G\|_{2,\kappa,\mu_S}
 \leq
 e^{K_2|S|/2}
 \|G\|_{4,\kappa,\nu_S}.                              \label{eq:s7v74-transfer}
\end{equation}
The label sum has constant one.
\end{theorem}

\begin{proof}
Set
\[
 A(x):=\sum_\alpha\kappa(\alpha\mid x)|G(x,\alpha)|^2.
\]
Then
\[
 \|G\|_{2,\kappa,\mu_S}^2
 =
 \int r_S(x)A(x)\,d\nu_S(x).
\]
Cauchy--Schwarz gives
\[
 \|G\|_{2,\kappa,\mu_S}^2
 \leq
 \|r_S\|_2\|A\|_2.
\]
Because \(\kappa(\cdot\mid x)\) is a probability vector, Jensen gives
\[
 A(x)^2
 \leq
 \sum_\alpha
 \kappa(\alpha\mid x)|G(x,\alpha)|^4.
\]
Thus \(\|A\|_2\leq\|G\|_{4,\kappa,\nu_S}^2\).
Insert \eqref{eq:s7v74-LMDC} and take a square root.
\end{proof}

\begin{remark}[Why two powers are lost]
\label{rem:s7v74-powers}
The physical \(L^2\) activity contains the product of one marginal
density and \(|G|^2\).  H\"older with exponents \(2,2\) therefore
requires a reference fourth moment.  No global density norm appears:
the density is first marginalized to the finite support \(S\).
\end{remark}

\begin{corollary}[RPLC loads the density factor]
\label{cor:s7v74-RPLC}
Under the V64 replica-defect hypotheses,
\eqref{eq:s7v74-LMDC} holds with
\begin{equation}
 K_2=
 {C_\times\over2}
 {q_{\rm rep}\over(1-q_{\rm rep})^2},                 \label{eq:s7v74-K2}
\end{equation}
which tends to zero under the V64 logarithmic activity schedule.
\end{corollary}

\begin{proof}
This is \Cref{cor:s7v64-RPLC,prop:s7v64-qrep}.
\end{proof}

\section{Per-block activity conversion}
\label{sec:s7v74-block}

For a connected support \(\Gamma\) and marked bad set
\(\varnothing\ne B\subset\Gamma\), let \(S(\Gamma)\) be the finite
reference-variable support of the complete labelled activity.  Assume
\begin{equation}
 |S(\Gamma)|\leq c_S|\Gamma|                          \label{eq:s7v74-support}
\end{equation}
for a fixed incidence constant \(c_S\), and suppose the reference moment
obeys
\begin{equation}
 \|G_{\Gamma,B}\|_{4,\kappa,\nu_{S(\Gamma)}}
 \leq
 q_{0,M}^{|\Gamma\setminus B|}
 t_{0,M}^{|B|}.                                       \label{eq:s7v74-reference}
\end{equation}

\begin{theorem}[Physical ML2AC from a reference fourth moment]
\label{thm:s7v74-ML2}
Define
\begin{equation}
 q_M:=e^{c_SK_2/2}q_{0,M},\qquad
 t_M:=e^{c_SK_2/2}t_{0,M}.                            \label{eq:s7v74-qt}
\end{equation}
Then
\begin{equation}
 \|G_{\Gamma,B}\|_{2,\kappa,\mu}
 \leq
 q_M^{|\Gamma\setminus B|}
 t_M^{|B|}.                                           \label{eq:s7v74-physical}
\end{equation}
\end{theorem}

\begin{proof}
Apply \Cref{thm:s7v74-transfer} on \(S(\Gamma)\), then use
\eqref{eq:s7v74-support,eq:s7v74-reference}:
\begin{align*}
 \|G_{\Gamma,B}\|_{2,\kappa,\mu}
 &\leq
 e^{K_2c_S|\Gamma|/2}
 q_{0,M}^{|\Gamma\setminus B|}
 t_{0,M}^{|B|}\\
 &=
 \left(e^{c_SK_2/2}q_{0,M}\right)^{|\Gamma\setminus B|}
 \left(e^{c_SK_2/2}t_{0,M}\right)^{|B|}.
\end{align*}
\end{proof}

\begin{corollary}[Mixed product with transition insertions]
\label{cor:s7v74-product}
Let \(b_M=\pi_M^{1/2}\) be the V73 probability activity.  Under
\eqref{eq:s7v74-physical}, the physical mixed weight satisfies
\begin{equation}
 |w(\Gamma,B)|
 \leq
 q_M^{|\Gamma\setminus B|}
 \bigl(t_Mb_M\bigr)^{|B|}.                            \label{eq:s7v74-product}
\end{equation}
\end{corollary}

\begin{proof}
Repeat the proof of \Cref{thm:s7v73-mixed}, using
\eqref{eq:s7v74-physical} in place of
\eqref{eq:s7v73-activity}.  The event factor is still
\(b_M^{|B|}\).
\end{proof}

\begin{theorem}[Rooted mixed sum with a bad insertion factor]
\label{thm:s7v74-sum}
If
\begin{equation}
 \zeta_M^{(4)}
 :=
 \Delta^2e^\alpha(q_M+t_Mb_M)<1,                     \label{eq:s7v74-zeta}
\end{equation}
then
\begin{equation}
 \sup_x
 \sum_{\substack{\Gamma\ni x\\\Gamma\ {\rm connected}}}
 e^{\alpha|\Gamma|}
 \sum_{\varnothing\ne B\subset\Gamma}
 |w(\Gamma,B)|
 \leq
 {e^\alpha t_Mb_M\over(1-\zeta_M^{(4)})^2}.           \label{eq:s7v74-sum}
\end{equation}
\end{theorem}

\begin{proof}
Use \eqref{eq:s7v74-product} in the binomial and lattice-animal proof of
\Cref{thm:s7v73-sum}, replacing \(b_M\) there by \(t_Mb_M\).
\end{proof}

\begin{corollary}[Polynomial transition losses are absorbed]
\label{cor:s7v74-polynomial}
Suppose
\begin{equation}
 q_{0,M}\to0,\qquad
 K_2\to0,\qquad
 t_{0,M}\leq C(1+M)^d,                               \label{eq:s7v74-poly}
\end{equation}
while \(b_M\leq C'(1+M)^{d'}e^{-cM}\) after the fixed V70 coloring and
square-root losses.  Then
\begin{equation}
 q_M\to0,\qquad t_Mb_M\to0,                           \label{eq:s7v74-small}
\end{equation}
and \eqref{eq:s7v74-zeta} holds for all sufficiently large \(M\).
\end{corollary}

\begin{proof}
The exponential \(e^{c_SK_2/2}\) tends to one.  A fixed polynomial is
absorbed by \(e^{-cM}\).  Thus both terms in
\eqref{eq:s7v74-zeta} tend to zero.
\end{proof}

\section{Loading transition and projector insertions}
\label{sec:s7v74-insertions}

\begin{lemma}[Bounded local insertions preserve the reference card]
\label{lem:s7v74-insertions}
Suppose an uninserted reference activity obeys
\eqref{eq:s7v74-reference} with \(t_{0,M}=1\).  If every marked bad or
transition block contributes a multiplier \(J_b\) satisfying
\begin{equation}
 \|J_b\|_{L^\infty(\nu)}\leq A_M,                     \label{eq:s7v74-J}
\end{equation}
then the inserted activity obeys
\eqref{eq:s7v74-reference} with \(t_{0,M}=A_M\).
\end{lemma}

\begin{proof}
The product of the \(|B|\) insertion norms is at most \(A_M^{|B|}\)
pointwise.  Pull it out of the labelled fourth moment.
\end{proof}

\begin{corollary}[Scale of SATC insertions]
\label{cor:s7v74-SATC}
If the V72 atlas and projector-response insertions enter the normalized
activity through bounded local multipliers of size
\begin{equation}
 A_M\leq
 C(1+M)^d
 \bigl(1+\theta_{\rm IMS}+\theta_{\rm P}\bigr),        \label{eq:s7v74-SATC}
\end{equation}
then fixed-width SATC makes \(A_M\) at most polynomial and
\Cref{cor:s7v74-polynomial} absorbs it.
\end{corollary}

\begin{proof}
V72 gives
\(\theta_{\rm IMS}+\theta_{\rm P}=O(g^2)\) under fixed uniform buffers.
Apply \Cref{lem:s7v74-insertions,cor:s7v74-polynomial}.
\end{proof}

\begin{firewall}[A form debit is not automatically a bounded multiplier]
\label{fire:s7v74-form}
The hypothesis \eqref{eq:s7v74-SATC} must be verified in the chosen
polymer representation.  V72 bounds quadratic-form debits; an unbounded
electric derivative insertion need not have an \(L^\infty\) multiplier
representation.  In that case its reference \(L^4\) moment must be proved
directly and entered in \eqref{eq:s7v74-reference}.
\end{firewall}

\section{A reference fourth-moment card}
\label{sec:s7v74-card}

\begin{definition}[Replica-labelled fourth-moment card, RL4C]
\label{def:s7v74-RL4C}
RL4C consists of:
\begin{enumerate}[label=\textup{(F\arabic*)},leftmargin=4em]
\item the V64 RPLC/LMDC bound
      \eqref{eq:s7v74-LMDC};
\item finite incidence \eqref{eq:s7v74-support};
\item the complete reference fourth moment
      \eqref{eq:s7v74-reference}, including transition labels, coordinate
      densities, and kinetic insertions;
\item the smallness and bad-insertion schedule
      \eqref{eq:s7v74-zeta};
\item uniformity over every retained slow and transition label; and
\item exact local ownership of every Wilson and electric term.
\end{enumerate}
\end{definition}

\begin{theorem}[RL4C loads ML2AC and closes mixed assembly]
\label{thm:s7v74-closure}
RL4C implies the physical labelled second-moment estimate
\eqref{eq:s7v74-physical}, the mixed product
\eqref{eq:s7v74-product}, and the rooted summability
\eqref{eq:s7v74-sum}.  Thus RL4C is a sufficient, reference-measure
loader for the V73 ML2AC card.
\end{theorem}

\begin{proof}
Rows (F1)--(F3) give
\Cref{thm:s7v74-transfer,thm:s7v74-ML2}; row (F4) gives
\Cref{thm:s7v74-sum}; rows (F5)--(F6) provide the physical uniformity and
exact decomposition required by ML2AC.
\end{proof}

\begin{firewall}[V64 does not yet prove the complete fourth moment]
\label{fire:s7v74-RL4C}
V64 proves the local marginal-density factor and the replica pressure
bound for the replica-admissible good-field interaction.  It does not
automatically prove \eqref{eq:s7v74-reference} after all compact
transition densities and electric insertions are included.  That
finite-support reference moment is now the sole analytic loader of the
mixed product interface.
\end{firewall}

\section{Status and next acquisition}
\label{sec:s7v74-status}

\begin{theorem}[What V74 proves]
\label{thm:s7v74-result}
V74 proves the exact transfer
\[
 \|G\|_{2,\kappa,\mu_S}
 \leq e^{K_2|S|/2}\|G\|_{4,\kappa,\nu_S},
\]
converts a per-block reference fourth moment into the physical ML2AC
bound, and proves the mixed rooted sum with bad insertion
\(t_Mb_M\).  The V64 replica factor tends to one, and every fixed
polynomial transition loss is absorbed by the exponentially small V70
bad activity.
\end{theorem}

\begin{proof}
Use
\Cref{thm:s7v74-transfer,cor:s7v74-RPLC,
thm:s7v74-ML2,thm:s7v74-sum,
cor:s7v74-polynomial,thm:s7v74-closure}.
\end{proof}

\begin{assessment}[Progress at seventh-series V74]
\label{ass:s7v74-status}
The mixed compact-chart assembly has been reduced to a finite-support
reference \(L^4\) estimate.  The normalization, label overlap, physical
density, bad probability, and connected-animal summation are no longer
independent gaps.  The remaining RL4C moment is the correct place to
insert the polynomial/Weyl bounds and the electric derivative moments.
\end{assessment}

\begin{decision}[V75 audit and acquisition order]
\label{dec:s7v74-V75}
V75 is the compulsory companion audit.  It will inspect the newest SM
and NS notes for a type-compatible reference fourth-moment or closable
score estimate.  It will then prove the local compact-group derivative
moments needed for RL4C, beginning with Peter--Weyl electric insertions
under the finite occupation cutoff and keeping every cutoff power
explicit.
\end{decision}

\begin{firewall}[Claim boundary after seventh-series V74]
\label{fire:s7v74-claim}
V74 does not prove RL4C for the complete Wilson/electric vertex family,
regulator-uniform SATC, the interacting joined slow-sector gap, the
remaining regular-sector UCTBB data, interacting endpoint continuity,
UCM, BCC, the pairing/Bogoliubov sector, weighted physical-source
incidence, nonlinear coarse-action cocycles, terminal strong-coupling
matching, regulator independence, cutoff removal, Osterwalder--Schrader
reconstruction, or the full Yang--Mills mass gap.
\end{firewall}

\paragraph{Parent-version record.}
V74 was formed from
\texttt{Renewal\_Yang\_Mills\_Mass\_Gap\_Research\_Notes\_V73.tex}.
The V73 parent had \(32\,200\) logical source lines, \(1\,150\,509\)
bytes, and SHA--256
\[
 \texttt{E5ACD30866839956E854C7A10D4E68A05258329975B8B68DF95B668BCABF82A7}.
\]
The next compulsory companion audit is V75.

% =============================================================================
% END APPENDED SEVENTH-SERIES V74 MATERIAL
% =============================================================================

% =============================================================================
% BEGIN APPENDED SEVENTH-SERIES V75 MATERIAL
% =============================================================================

\chapter{V75 companion audit and cutoff-explicit compact-group derivative moments}
\label{ch:s7v75}

\section{Compulsory SM/NS audit}
\label{sec:s7v75-audit}

\begin{audit}[Current companion files read at V75]
\label{audit:s7v75-files}
The active companion notes are:
\begin{enumerate}[label=\textup{(\roman*)},leftmargin=3.8em]
\item
\texttt{Renewal\_SM\_Einstein\_Continuum\_Research\_Notes\_V34.tex},
with \(11\,882\) logical source lines, \(430\,833\) bytes, and SHA--256
\[
 \texttt{D69E88B6D183562F04B47A3CC3A265FD8250EB8B3FD5EA74A9FA45A1EE6C3689};
\]
\item
\texttt{Renewal\_NS\_Stability\_Research\_Notes\_V26.tex},
with \(5\,319\) logical source lines, \(278\,283\) bytes, and SHA--256
\[
 \texttt{82C887F8C2C69E4F28FAD7C3DC8DE5B9099CB5A4BF431EBA1FFF3E91935978AD}.
\]
\end{enumerate}
The SM file has advanced from V29 at the V70 audit to V34.  The NS file
is unchanged.
\end{audit}

\begin{theorem}[Transferable conclusions of the V75 audit]
\label{thm:s7v75-audit}
The audit gives four precise conclusions.
\begin{enumerate}[leftmargin=3em]
\item SM V30 proves that event rarity does not imply an \(L^2\) score
      bound; a higher score moment is needed on the singular sector.  This
      agrees with the separate probability and RL4C ledgers in V73--V74.
\item SM V31 derives Fisher--Hessian identities for a complete logarithmic
      derivative.  The integration-by-parts algebra transfers to compact
      Haar carriers and is reproved below.
\item SM V32--V34 derive Gaussian trace and relative-Hessian criteria for
      metric scores.  Their trace ideals, constraint pullbacks, and
      conformal rows do not transfer to lattice gauge links.
\item NS V26 contains no probability score or occupation cutoff and gives
      no RL4C moment.
\end{enumerate}
\end{theorem}

\begin{proof}
These are the terminal result and firewall statements of the active
companion sections fixed by \Cref{audit:s7v75-files}.  The Fisher identity
uses only a normalized integration-by-parts law and therefore has a
type-compatible compact-group analogue.  The quantitative SM and NS
operators act on different spaces and supply no Yang--Mills constant.
\end{proof}

\begin{assessment}[Audit-guided separation]
\label{ass:s7v75-audit}
Two cutoffs must not be conflated.  A Peter--Weyl Casimir cutoff is a
finite compact-group representation regulator.  The \(P_M\) used in the
weak-field interaction estimates is an oscillator-number cutoff after
rescaling about a flat chart.  V75 records derivative powers for both;
matching their cofinal removal remains a separate row.
\end{assessment}

\section{Compact Haar Fisher identities}
\label{sec:s7v75-Fisher}

Let \(G^E\) carry normalized product Haar measure \(dm\), let
\begin{equation}
 d\mu={e^{-V}\over Z}\,dm,                             \label{eq:s7v75-mu}
\end{equation}
and let \(X\) be a real left- or right-invariant vector field on one link.
It is divergence free for \(dm\).  Define the complete logarithmic score
\begin{equation}
 b_X:=XV.                                              \label{eq:s7v75-score}
\end{equation}

\begin{theorem}[Compact-group Fisher--Hessian hierarchy]
\label{thm:s7v75-Fisher}
For every integer \(m\geq1\) for which the displayed quantities are
integrable,
\begin{equation}
 {\mathbb E}_\mu[b_X^{2m}]
 =
 (2m-1)
 {\mathbb E}_\mu
 \left[b_X^{2m-2}Xb_X\right].                         \label{eq:s7v75-Fisher}
\end{equation}
In particular,
\begin{equation}
 {\mathbb E}_\mu[b_X^2]
 =
 {\mathbb E}_\mu[X^2V].                               \label{eq:s7v75-Fisher2}
\end{equation}
There is no boundary-flux term.
\end{theorem}

\begin{proof}
Haar invariance gives
\(\int Xf\,dm=0\) for every smooth \(f\).  Apply this to
\[
 f=b_X^{2m-1}e^{-V}.
\]
The product and chain rules give
\[
 0=
 (2m-1)\int b_X^{2m-2}(Xb_X)e^{-V}\,dm
 -\int b_X^{2m}e^{-V}\,dm.
\]
Divide by \(Z\).  The case \(m=1\) is
\eqref{eq:s7v75-Fisher2}.
\end{proof}

\begin{corollary}[Higher score moments from a bounded differentiated
score]
\label{cor:s7v75-score-moments}
If
\begin{equation}
 \|Xb_X\|_{L^\infty(\mu)}\leq B,                      \label{eq:s7v75-B}
\end{equation}
then
\begin{equation}
 {\mathbb E}_\mu|b_X|^{2m}
 \leq
 (2m-1)!!\,B^m.                                       \label{eq:s7v75-score-moments}
\end{equation}
Consequently
\begin{equation}
 \|b_X\|_{L^{2m}(\mu)}
 \leq
 ((2m-1)!!)^{1/(2m)}B^{1/2}.                         \label{eq:s7v75-score-norm}
\end{equation}
\end{corollary}

\begin{proof}
Take absolute values on the right of
\eqref{eq:s7v75-Fisher} and use \eqref{eq:s7v75-B}:
\[
 {\mathbb E}|b_X|^{2m}
 \leq
 (2m-1)B\,{\mathbb E}|b_X|^{2m-2}.
\]
Iterate from \({\mathbb E}1=1\).
\end{proof}

\begin{corollary}[Wilson one-link score is polynomially controlled]
\label{cor:s7v75-Wilson-score}
For the finite Wilson measure with
\[
 V=\beta\sum_p\phi(W_p),
\]
there is a fixed incidence constant \(C_{\rm W}\) such that
\begin{equation}
 \|Xb_X\|_\infty
 =
 \left\|\beta X^2\sum_p\phi(W_p)\right\|_\infty
 \leq C_{\rm W}\beta.                                \label{eq:s7v75-Wilson-Hessian}
\end{equation}
Hence, for fixed \(m\),
\begin{equation}
 \|b_X\|_{L^{2m}(\mu)}
 \leq C_m\beta^{1/2}.                                 \label{eq:s7v75-Wilson-moment}
\end{equation}
\end{corollary}

\begin{proof}
Only the finitely many plaquettes incident on the differentiated link
contribute.  Every first and second invariant derivative of
\(\operatorname{ReTr}W_p\) is uniformly bounded on the compact group,
so their sum is at most a fixed incidence constant.  Apply
\Cref{cor:s7v75-score-moments} with \(B=C_{\rm W}\beta\).
\end{proof}

\begin{firewall}[The complete score must be differentiated]
\label{fire:s7v75-score}
If a gauge fixing, chart density, or counterterm changes the physical
density, its logarithmic derivative belongs in \(b_X\) before
\eqref{eq:s7v75-Fisher} is used.  Applying the identity to the Wilson
piece alone while differentiating a different normalized measure would
omit score terms and is invalid.
\end{firewall}

\section{Peter--Weyl Casimir cutoff}
\label{sec:s7v75-PW}

Let
\begin{equation}
 {\cal C}_E:=\sum_{e\in E}(-\Delta_e),\qquad
 \Pi_\Lambda:=\mathbf1_{[0,\Lambda]}({\cal C}_E).      \label{eq:s7v75-Casimir}
\end{equation}
Choose an orthonormal invariant tangent frame
\(\{X_{e,a}\}\), so
\begin{equation}
 -\Delta_e=\sum_aX_{e,a}^*X_{e,a}.                    \label{eq:s7v75-Laplacian}
\end{equation}

\begin{lemma}[One derivative under a Casimir cutoff]
\label{lem:s7v75-one}
For every \(e,a\),
\begin{equation}
 \|X_{e,a}\Pi_\Lambda\|\leq\Lambda^{1/2},\qquad
 \|(-\Delta_e)\Pi_\Lambda\|\leq\Lambda.               \label{eq:s7v75-one}
\end{equation}
\end{lemma}

\begin{proof}
The bi-invariant Casimir commutes with every invariant derivative and
with \(\Pi_\Lambda\).  Positivity and
\eqref{eq:s7v75-Laplacian} give
\[
 \|X_{e,a}\Pi_\Lambda\|^2
 =
 \|\Pi_\Lambda X_{e,a}^*X_{e,a}\Pi_\Lambda\|
 \leq
 \|\Pi_\Lambda{\cal C}_E\Pi_\Lambda\|
 \leq\Lambda.
\]
Also \(0\leq-\Delta_e\leq{\cal C}_E\), proving the second bound.
\end{proof}

\begin{theorem}[Derivative-word Casimir bound]
\label{thm:s7v75-word}
For every word of \(k\) invariant derivatives,
\begin{equation}
 \left\|
 X_{e_1,a_1}\cdots X_{e_k,a_k}\Pi_\Lambda
 \right\|
 \leq\Lambda^{k/2}.                                   \label{eq:s7v75-word}
\end{equation}
If
\begin{equation}
 D=\sum_{|w|\leq k}f_w(U)X_w                         \label{eq:s7v75-D}
\end{equation}
is written with all coefficient differentiations already expanded, then
\begin{equation}
 \|D\Pi_\Lambda\|
 \leq
 \sum_{|w|\leq k}\|f_w\|_\infty
 (1+\Lambda)^{|w|/2}.                                \label{eq:s7v75-D-bound}
\end{equation}
\end{theorem}

\begin{proof}
Every invariant derivative commutes with the central total Casimir and
its spectral projection.  Insert \(\Pi_\Lambda\) after each derivative
and apply \Cref{lem:s7v75-one} \(k\) times.  Multiplication by \(f_w\)
has operator norm \(\|f_w\|_\infty\); sum the word bounds.  The harmless
\(1+\Lambda\) includes the zeroth-order term.
\end{proof}

\begin{corollary}[Finite-cutoff fourth moments]
\label{cor:s7v75-PW-L4}
For any normalized positive state \(\omega\) supported in
\(\operatorname{Ran}\Pi_\Lambda\),
\begin{equation}
 \omega\bigl(|D\Pi_\Lambda|^4\bigr)^{1/4}
 \leq\|D\Pi_\Lambda\|
 \leq C_D(1+\Lambda)^{k/2}.                           \label{eq:s7v75-PW-L4}
\end{equation}
\end{corollary}

\begin{proof}
Positive normalized functionals are contractive on positive bounded
operators, so
\(\omega(|A|^4)\leq\|A\|^4\).  Apply
\eqref{eq:s7v75-D-bound}.
\end{proof}

\section{Weak-field oscillator-number cutoff}
\label{sec:s7v75-oscillator}

Let \(a_j,a_j^\dagger\), \(1\leq j\leq r\), be oscillator annihilation
and creation operators, let
\begin{equation}
 {\cal N}:=\sum_{j=1}^ra_j^\dagger a_j,\qquad
 P_M:=\mathbf1_{[0,M]}({\cal N}).                     \label{eq:s7v75-PM}
\end{equation}

\begin{lemma}[Creation and annihilation bounds]
\label{lem:s7v75-aa}
\begin{equation}
 \|a_jP_M\|\leq\sqrt M,\qquad
 \|a_j^\dagger P_M\|\leq\sqrt{M+1}.                  \label{eq:s7v75-aa}
\end{equation}
\end{lemma}

\begin{proof}
On a number basis,
\(a_j\) multiplies by \(\sqrt{n_j}\) and
\(a_j^\dagger\) by \(\sqrt{n_j+1}\).  On total occupation at most \(M\),
\(n_j\leq M\).
\end{proof}

\begin{theorem}[Oscillator-word cutoff bound]
\label{thm:s7v75-osc-word}
Let \(A_1\cdots A_k\) be a word in creation and annihilation operators.
Then
\begin{equation}
 \|A_1\cdots A_kP_M\|
 \leq(M+k)^{k/2}.                                     \label{eq:s7v75-osc-word}
\end{equation}
If the oscillator frequencies lie in
\(0<\omega_*\leq\omega_j\leq\omega^*<\infty\), every degree-\(k\)
word in \(y_j,\partial_{y_j}\) satisfies
\begin{equation}
 \|D_k(y,\partial_y)P_M\|
 \leq C_{k,\omega_*,\omega^*}(M+k)^{k/2}.             \label{eq:s7v75-y-word}
\end{equation}
\end{theorem}

\begin{proof}
After the rightmost operator acts, total occupation is at most \(M+1\);
after \(j\) factors it is at most \(M+j\).  Apply
\Cref{lem:s7v75-aa} successively and bound every factor by
\(\sqrt{M+k}\).  The standard identities
\[
 y_j={a_j+a_j^\dagger\over\sqrt{2\omega_j}},
 \qquad
 \partial_{y_j}=
 \sqrt{\omega_j\over2}(a_j-a_j^\dagger)
\]
write a degree-\(k\) coordinate/derivative word as at most \(2^k\)
creation/annihilation words with uniformly bounded frequency
coefficients.  Sum their bounds.
\end{proof}

\begin{corollary}[Oscillator fourth-moment bound]
\label{cor:s7v75-osc-L4}
For every normalized state supported in
\(\operatorname{Ran}P_M\),
\begin{equation}
 \omega\bigl(|D_kP_M|^4\bigr)^{1/4}
 \leq
 C_{k,\omega_*,\omega^*}(M+k)^{k/2}.                 \label{eq:s7v75-osc-L4}
\end{equation}
\end{corollary}

\begin{proof}
Use the positive-functional argument of
\Cref{cor:s7v75-PW-L4} and
\eqref{eq:s7v75-y-word}.
\end{proof}

\begin{firewall}[The two cutoffs are not the same projection]
\label{fire:s7v75-cutoffs}
\(\Pi_\Lambda\) is diagonal in compact-group irreducible
representations, whereas \(P_M\) is diagonal in the weak-field oscillator
number.  Neither commutes exactly with the nonlinear chart map to the
other basis.  The bounds above may be used in their respective
representations, but replacing one cutoff by the other requires a
separate cofinal comparison.
\end{firewall}

\section{Loading the RL4C polynomial row}
\label{sec:s7v75-RL4C}

\begin{theorem}[Finite derivative lists give polynomial RL4C losses]
\label{thm:s7v75-RL4C}
Suppose each local transition/electric activity contains at most
\(k_*\) invariant or oscillator derivatives and at most \(n_*\) such
factors per block, with bounded coefficient functions.  At finite
cutoffs,
\begin{equation}
 t_{0,M,\Lambda}
 \leq
 C(1+\Lambda)^{d_\Lambda}(1+M)^{d_M}                 \label{eq:s7v75-t}
\end{equation}
for fixed exponents depending only on \(k_*,n_*\) and the block
geometry.  Consequently, if
\begin{equation}
 \Lambda\leq C_\Lambda(1+M)^s                         \label{eq:s7v75-cofinal}
\end{equation}
along the declared cofinal regulator and the V70 bad activity is
exponential in \(M\), then
\begin{equation}
 t_{0,M,\Lambda}b_M\longrightarrow0.                 \label{eq:s7v75-absorb}
\end{equation}
\end{theorem}

\begin{proof}
Apply \Cref{thm:s7v75-word} to every compact-group derivative and
\Cref{thm:s7v75-osc-word} to every oscillator derivative.  A fixed
product of their bounds is a fixed polynomial, giving
\eqref{eq:s7v75-t}.  Under \eqref{eq:s7v75-cofinal} it becomes a fixed
power of \(1+M\), which is absorbed by the exponential factor in \(b_M\)
as in \Cref{cor:s7v74-polynomial}.
\end{proof}

\begin{corollary}[Good-vertex derivative condition]
\label{cor:s7v75-good}
If a degree-\(k\) good-field derivative vertex carries a coupling factor
\(g^\rho\), its cutoff norm is bounded by
\begin{equation}
 Cg^\rho(M+k)^{k/2}.                                  \label{eq:s7v75-good}
\end{equation}
It is a vanishing good activity whenever
\begin{equation}
 g^\rho(M+k)^{k/2}\longrightarrow0.                  \label{eq:s7v75-good-small}
\end{equation}
The logarithmic occupation schedule used in V55--V62 satisfies this for
every fixed \(k\) and \(\rho>0\).
\end{corollary}

\begin{proof}
Combine the coupling factor with
\eqref{eq:s7v75-y-word}.  Under the logarithmic schedule, \(M\) grows as
a fixed power of a logarithm while every positive power of the
asymptotically free coupling dominates that growth in the previously
declared V55 scaling.
\end{proof}

\begin{firewall}[Polynomial bounds do not enumerate the complete vertex
family]
\label{fire:s7v75-enumeration}
\Cref{thm:s7v75-RL4C} closes the cutoff-growth part of RL4C for any fixed
finite derivative list.  The physical proof must still show that the
complete transition-labelled Wilson Hamiltonian has such a list after
all commutators, Duhamel insertions, counterterms, and atlas derivatives
are expanded, and that their reference fourth moments use the same local
support.
\end{firewall}

\section{Status and next acquisition}
\label{sec:s7v75-status}

\begin{theorem}[What V75 proves]
\label{thm:s7v75-result}
V75 completes the mandatory companion audit; proves the compact Haar
Fisher--Hessian hierarchy and its Wilson score moments; proves
cutoff-explicit derivative-word bounds
\[
 \|X_1\cdots X_k\Pi_\Lambda\|\leq\Lambda^{k/2},
 \qquad
 \|A_1\cdots A_kP_M\|\leq(M+k)^{k/2};
\]
and shows that every fixed finite derivative list contributes only a
polynomial RL4C loss, absorbed by the V70 exponential bad activity.
\end{theorem}

\begin{proof}
Use
\Cref{thm:s7v75-audit,thm:s7v75-Fisher,
cor:s7v75-Wilson-score,thm:s7v75-word,
thm:s7v75-osc-word,thm:s7v75-RL4C}.
\end{proof}

\begin{assessment}[Progress at seventh-series V75]
\label{ass:s7v75-status}
Derivative insertions are no longer an unspecified source of cutoff
growth: at fixed local order their fourth moments are polynomial in the
declared representation or occupation cutoff.  The next physical task is
finite but nontrivial: enumerate the complete local vertex family and
verify that its coupling powers beat those polynomial losses in the
RL4C reference measure.
\end{assessment}

\begin{decision}[V76 acquisition order]
\label{dec:s7v75-V76}
V76 will construct the complete local transition vertex ledger.  It will
separate Wilson multiplication vertices, oscillator polynomial vertices,
atlas IMS multipliers, Riesz-projector derivatives, and electric
commutators; assign each its derivative order, coupling power, support,
and \(L^4\) cutoff exponent; and test all rows against the V55 logarithmic
schedule.  Any vertex with zero coupling power and growing cutoff norm
will be moved into the reference or slow Hamiltonian rather than treated
perturbatively.
\end{decision}

\begin{firewall}[Claim boundary after seventh-series V75]
\label{fire:s7v75-claim}
V75 does not enumerate or prove RL4C for the complete physical vertex
family, compare Peter--Weyl and oscillator cutoff removal, prove
regulator-uniform SATC, the interacting joined slow-sector gap, the
remaining regular-sector UCTBB data, interacting endpoint continuity,
UCM, BCC, the pairing/Bogoliubov sector, weighted physical-source
incidence, nonlinear coarse-action cocycles, terminal strong-coupling
matching, regulator independence, cutoff removal, Osterwalder--Schrader
reconstruction, or the full Yang--Mills mass gap.
\end{firewall}

\paragraph{Parent-version record.}
V75 was formed from
\texttt{Renewal\_Yang\_Mills\_Mass\_Gap\_Research\_Notes\_V74.tex}.
The V74 parent had \(32\,586\) logical source lines, \(1\,163\,359\)
bytes, and SHA--256
\[
 \texttt{847665ADE48B984160FBCBA4A98C6CD79D1E6EDF15D0A710F4E947773D3E1A2C}.
\]
The next compulsory companion audit is V80.

% =============================================================================
% END APPENDED SEVENTH-SERIES V75 MATERIAL
% =============================================================================

% =============================================================================
% BEGIN APPENDED SEVENTH-SERIES V76 MATERIAL
% =============================================================================

\chapter{A generator-complete transition-vertex ledger and the
inverse-coupling correction}
\label{ch:s7v76}

\section{Purpose and correction}
\label{sec:s7v76-purpose}

V75 proves polynomial cutoff bounds for every fixed derivative word.
That statement is necessary but not sufficient for RL4C.  Two different
questions had still been conflated:
\begin{enumerate}[label=\textup{(\roman*)},leftmargin=3.5em]
\item which elementary local terms can actually be generated from the
      Wilson Hamiltonian, its weak-field charts, and the transition
      atlas; and
\item whether an arbitrary number of occurrences of those elementary
      terms may be charged to one bad block.
\end{enumerate}
The first question has a finite answer at the level of generator types.
The second is a resummation question and is not answered by bounding one
derivative word.

This chapter gives the generator-complete ledger.  It also corrects the
schedule statement in V75.  Haar score moments can grow as a fixed power
of \(g^{-1}\).  Such a power is absorbed by an
\(\exp(-cM)\) bad-block activity when
\[
 M=(\log L)^\chi,\qquad \chi>1,
\]
but not as a consequence of \(\chi>0\) alone.  The correction does not
alter any of the V55 small-coupling estimates.  Finally, the ledger shows
that electric principal-part and atlas derivative terms should be
resummed as relative forms; expanding them as arbitrary marked
\(L^4\)-vertices would create an artificial multiplicity obstruction.

\section{The local differential algebra}
\label{sec:s7v76-algebra}

Fix a gauge-reduced finite block and a weak-field coordinate
\(q=gy\).  Let \({\cal P}\) denote either the oscillator cutoff \(P_M\)
or, in a compact boundary variable, the Peter--Weyl cutoff
\(\Pi_\Lambda\).  A local operator monomial has the form
\begin{equation}
 {\cal V}_{\rho,d,j}
 =
 g^\rho f_d(y,\eta)D^j,                              \label{eq:s7v76-monomial}
\end{equation}
where \(f_d\) is polynomial of \(y\)-degree at most \(d\), its
\(\eta\)-dependent coefficients and their declared derivatives are
bounded on the compact chart, and \(D^j\) is a word of \(j\) oscillator
or compact invariant derivatives.  The physical Hamiltonian coefficient
\(a^{-1}\) is suppressed in this dimensionless notation.

\begin{lemma}[Cutoff norm of one generated monomial]
\label{lem:s7v76-word}
For fixed \(d,j\), there is a block constant \(C_{d,j}\) such that
\begin{align}
 \|{\cal V}_{\rho,d,j}P_M\|
 &\leq
 C_{d,j}g^\rho(M+d+j+1)^{(d+j)/2},                  \label{eq:s7v76-osc-bound}\\
 \|{\cal V}_{\rho,0,j}\Pi_\Lambda\|
 &\leq
 C_jg^\rho(1+\Lambda)^{j/2}.                         \label{eq:s7v76-PW-bound}
\end{align}
For a mixed word the product of the two displayed polynomial factors is
valid.
\end{lemma}

\begin{proof}
Write every coordinate and oscillator derivative as a fixed linear
combination of creation and annihilation operators.  The total word
length is at most \(d+j\), so
\Cref{thm:s7v75-osc-word} gives
\eqref{eq:s7v76-osc-bound}.  Equation
\eqref{eq:s7v76-PW-bound} is
\Cref{thm:s7v75-word}.  Operators acting in the two tensor factors
commute, so their bounds multiply.
\end{proof}

The electric part is most naturally kept in quadratic-form normal form.
Let
\begin{equation}
 {\mathfrak e}_{A}(u,v)
 :=
 \sum_{\alpha,\beta}
 \langle X_\alpha u,A_{\alpha\beta}X_\beta v\rangle, \label{eq:s7v76-form}
\end{equation}
where \(X_\alpha\) are invariant or chart derivatives.

\begin{lemma}[First-derivative closure of form localization]
\label{lem:s7v76-form-Leibniz}
Let \(B\) be a bounded matrix-valued multiplier preserving the local
form domain.  Then
\begin{align}
 {\mathfrak e}_{A}(Bu,Bv)
 &=
 \sum_{\alpha,\beta}
 \langle
 B X_\alpha u+(X_\alpha B)u,\,
 A_{\alpha\beta}
 [B X_\beta v+(X_\beta B)v]
 \rangle .                                           \label{eq:s7v76-form-Leibniz}
\end{align}
Thus localization of a second-order electric operator at the form level
generates only the four types
\[
 (XB)^*A(XB),\quad
 (XB)^*ABX,\quad
 X^*B^*A(XB),\quad
 X^*B^*ABX,
\]
and requires only first derivatives of \(B\).
\end{lemma}

\begin{proof}
The invariant-vector-field product rule gives
\(X_\alpha(Bu)=B X_\alpha u+(X_\alpha B)u\).
Insert this identity on both sides of
\eqref{eq:s7v76-form} and expand the four terms.
\end{proof}

\begin{corollary}[No hidden higher derivative in the atlas form]
\label{cor:s7v76-nohigher}
Take
\[
 B_i(\eta)=\sqrt{\chi_i(\eta)}\,U_i(\eta)P_i(\eta),
\]
where \(U_i\) is a unitary transition frame and \(P_i\) a buffered slow
projection.  The exact localized electric form contains \(B_i\),
\(X B_i\), and no second derivative of \(B_i\).  Consequently the V71
IMS term and V72 projector response exhaust the geometric form types.
\end{corollary}

\begin{proof}
Apply \Cref{lem:s7v76-form-Leibniz}.  The product rule for \(XB_i\)
produces \(X\sqrt{\chi_i}\), \(XU_i\), and \(XP_i\).  The pure
\(X\sqrt{\chi_i}\)-square is the IMS row; the component
\(Q_i(XP_i)P_i\) is the projector-response row.  The remaining cross
terms are the corresponding connection forms.  No other term is
generated.
\end{proof}

\begin{firewall}[Operator normal form needs more smoothness]
\label{fire:s7v76-C2}
If the localized form is expanded into a pointwise second-order
operator, derivatives can strike \(XB_i\), and \(X^2B_i\) appears.
V72 assumes only the first-derivative data required by the form.
Therefore a claimed operator-vertex ledger with only the V72 hypotheses
would be false.  The present ledger is complete at the closed-form
level; an operator-normal-form version needs a regulator-uniform
\(C^2\) atlas card.
\end{firewall}

\section{Taylor families generated by the Wilson Hamiltonian}
\label{sec:s7v76-Taylor}

The magnetic Taylor expansion on the weak-field chart is
\begin{equation}
 g^{-2}S(gy)
 =
 H_2(y)+\sum_{r\geq3}g^{r-2}H_r(y),                  \label{eq:s7v76-magnetic}
\end{equation}
with \(H_r\) homogeneous of degree \(r\) after collecting its fixed
block coefficients.  The electric inverse metric has an analytic
expansion
\begin{equation}
 A(gy)=I+\sum_{m\geq1}g^mA_m(y),                     \label{eq:s7v76-metric}
\end{equation}
where \(A_m\) has coordinate degree at most \(m\).  Haar-density
conjugation produces lower differential order but no new nonanalytic
generator.

\begin{proposition}[Magnetic and electric raw cutoff rows]
\label{prop:s7v76-raw}
On the oscillator window,
\begin{align}
 \|g^{r-2}H_rP_M\|
 &\leq C_r g^{r-2}(M+r+1)^{r/2},                    \label{eq:s7v76-mag-row}\\
 \|P_M X^*g^mA_mX P_M\|
 &\leq C_m g^m(M+m+3)^{(m+2)/2}.                    \label{eq:s7v76-electric-row}
\end{align}
Terms created when a derivative strikes \(A_m\) have the same coupling
power and strictly smaller total oscillator word degree.
\end{proposition}

\begin{proof}
The first line is
\Cref{lem:s7v76-word} with
\((\rho,d,j)=(r-2,r,0)\).  In the second line,
\(X^*A_mX\) is a sum of words of total coordinate/derivative degree at
most \(m+2\); apply the same lemma.  Differentiating a degree-\(m\)
coefficient lowers its coordinate degree by one while consuming one
derivative, so it cannot increase the total word degree.
\end{proof}

\begin{lemma}[Analytic-family summation]
\label{lem:s7v76-analytic}
Suppose \(C_r\leq C_0C_1^r\).  If
\begin{equation}
 C_1g\sqrt{M+c}\leq\frac12,                           \label{eq:s7v76-radius}
\end{equation}
then
\begin{align}
 \sum_{r\geq5}
 C_rg^{r-2}(M+r+c)^{r/2}
 &\leq Cg^3(M+c')^{5/2},                             \label{eq:s7v76-mag-tail}\\
 \sum_{m\geq1}
 C_mg^m(M+m+c)^{(m+2)/2}
 &\leq Cg(M+c')^{3/2}.                               \label{eq:s7v76-elec-tail}
\end{align}
Here \(c,c'\) and \(C\) depend only on the fixed analytic block data.
\end{lemma}

\begin{proof}
Analyticity on a fixed complex neighborhood gives the geometric
coefficient bound.  After enlarging \(c'\), the ratio between consecutive
majorant terms is at most \(2C_1g\sqrt{M+c'}\).  Under
\eqref{eq:s7v76-radius} both series are bounded by twice their first
terms.  Those terms have the displayed powers.
\end{proof}

\begin{remark}[Raw operator size versus form size]
\label{rem:s7v76-raw-form}
The leading raw electric row is \(O(gM^{3/2})\), whereas the exact
V43 form comparison is \(O(gR)\) on \(|y|\leq R\).  With \(R=\sqrt M\)
the form coefficient is \(O(g\sqrt M)\).  The extra factor \(M\) in
\eqref{eq:s7v76-electric-row} is the energy already present in the
quadratic reference.  Charging it again as a bounded vertex is both
wasteful and the source of a false multiplicity problem.
\end{remark}

\section{The generator-complete ledger}
\label{sec:s7v76-ledger}

\begin{definition}[Declared finite-block generator class]
\label{def:s7v76-class}
The declared class consists of the Kogut--Susskind Wilson Hamiltonian
\eqref{eq:s7v72-KS}, fixed-block weak-field dilation and Haar-density
conjugation, the V39 gauge-invariant chart partition, the V71 labelled
transition isometry and range projector, buffered \(C^1\) Riesz
projections, and compression by \(P_M\) or \(\Pi_\Lambda\).  Terms are
kept as closed quadratic forms whenever an electric derivative is
present.
\end{definition}

\begin{theorem}[Exhaustive generator-type ledger]
\label{thm:s7v76-ledger}
Every local term produced by the class in
\Cref{def:s7v76-class} belongs to exactly one of the following rows,
up to a finite sum and adjoint:
\begin{enumerate}[label=\textup{(L\arabic*)},leftmargin=4em]
\item the global quadratic transverse Hamiltonian \(H_2/a\);
\item the exact compact slow Hamiltonian, including torons, boundary
      holonomies, and every retained transition mode;
\item a unit-norm atlas transition or range cocycle;
\item a normalized compact coordinate density or Jacobian;
\item a magnetic multiplier
      \(a^{-1}g^{r-2}H_r\), \(r\geq3\);
\item an electric metric form
      \(a^{-1}\langle Xu,g^mA_mXv\rangle\), \(m\geq1\),
      together with its Haar lower-order partners;
\item a chart or atlas connection form containing one derivative of
      \(\sqrt\chi\), a transition frame, or a slow projection;
\item the scalar IMS square, of relative size
      \(O((g/w)^2)\) for the transition atlas or \(O(R^{-2})\) for the
      weak-field chart;
\item the buffered-projector response, of relative size
      \(O(g^2d_{\rm sp}^{-4})\);
\item a compact invariant derivative or score attached to a marked
      observable; or
\item a cutoff boundary term generated by \(P_M\) or \(\Pi_\Lambda\).
\end{enumerate}
Rows (L1)--(L4) are exact representation or reference data.  Rows
(L5)--(L9) are the only dynamical perturbation/form rows.  Rows
(L10)--(L11) are insertions, not new Hamiltonian generators.
\end{theorem}

\begin{proof}
The magnetic Hamiltonian is multiplication, and analyticity gives
exactly row (L5) after its quadratic and slow restrictions are placed in
(L1)--(L2).  The electric Hamiltonian is second order in invariant
derivatives.  Its coefficient Taylor expansion gives (L6).
\Cref{lem:s7v76-form-Leibniz} proves that multiplication by a chart,
frame, or projection generates only the cross and square forms
(L7)--(L9).  V71 proves that the remaining atlas maps are unitary
cocycles acting on the exact labelled range, giving (L3); change of
coordinates gives (L4).  Differentiating a marked observable gives
(L10), and compression is either inert or produces the spectral-boundary
row (L11).  These operations exhaust
\Cref{def:s7v76-class}, so no further generator type can occur.
\end{proof}

\begin{corollary}[Support and ownership]
\label{cor:s7v76-support}
Every row in \Cref{thm:s7v76-ledger} is supported on the union of one
fixed plaquette star, one fixed chart collar, and one fixed transition
patch.  Hence there is a regulator-independent incidence constant
\(c_{\rm led}\) such that a connected family of \(n\) ledger entries
uses at most \(c_{\rm led}n\) reference-variable blocks.
\end{corollary}

\begin{proof}
Wilson terms are supported on one plaquette and electric terms on one
link.  The V39 cutoff, V71 atlas, and buffered conditional Hessian were
all defined on fixed-radius blocks.  Taking the union of finitely many
fixed supports enlarges the radius and cardinality by a constant only.
Assign each term to the lexicographically first block meeting its
distinguished link or plaquette.  This gives exact ownership and the
linear incidence bound.
\end{proof}

\section{Coupling verdict under the logarithmic schedule}
\label{sec:s7v76-schedule}

Retain
\begin{equation}
 g=O(L^{-1/2}),\qquad
 M=(\log L)^\chi,\qquad
 R=\sqrt M,\qquad
 \Lambda\leq C_\Lambda(1+M)^s.                       \label{eq:s7v76-schedule}
\end{equation}

\begin{theorem}[Every positive-coupling ledger row closes]
\label{thm:s7v76-positive}
For fixed Taylor/derivative order, every row carrying \(g^\rho\) with
\(\rho>0\) has vanishing cutoff norm under
\eqref{eq:s7v76-schedule}.  In particular,
\begin{align}
 gM^{3/2}&\longrightarrow0,&
 g^2M^2&\longrightarrow0,&
 g^3M^{5/2}&\longrightarrow0,                        \label{eq:s7v76-basic-rates}\\
 g^2\Lambda^{1/2}w^{-1}&\longrightarrow0,&
 g^2\Lambda w^{-2}&\longrightarrow0                 \label{eq:s7v76-atlas-rates}
\end{align}
for fixed \(w>0\), and the same conclusion holds after multiplication
by any fixed power of \(M\) or \(\Lambda\).
\end{theorem}

\begin{proof}
Every displayed expression is a negative fixed power of \(L\) times a
fixed power of \(\log L\).  This tends to zero.  The general statement
follows from \Cref{lem:s7v76-word} and
\(\Lambda\leq C(1+M)^s\).  The analytic families are summed by
\Cref{lem:s7v76-analytic}, because
\(g\sqrt M\to0\).
\end{proof}

\begin{theorem}[Zero-coupling placement rule]
\label{thm:s7v76-zero}
Suppose a generated row has no positive coupling power and its cutoff
norm is unbounded.  Then it cannot be part of a vanishing RL4C good
activity.  In the declared class it is one of:
\begin{enumerate}[label=\textup{(\alph*)},leftmargin=3.5em]
\item the global quadratic Hamiltonian;
\item the exact joined slow Hamiltonian;
\item a representation cutoff or transition map; or
\item a derivative of an observable.
\end{enumerate}
The first three must remain in the reference/slow representation.  The
fourth must be paid by an insertion norm and cannot be relabelled as a
small interaction.
\end{theorem}

\begin{proof}
Inspect the exhaustive list in
\Cref{thm:s7v76-ledger}.  Rows (L5)--(L9) either have positive coupling
power or a separately vanishing relative IMS coefficient.  The only
remaining zero-coupling entries are (L1)--(L4), (L10), and (L11).
Unitary cocycles and normalized densities are bounded representation
data.  This leaves precisely the four cases stated.
\end{proof}

\section{The inverse-coupling score row}
\label{sec:s7v76-score}

For \(d\mu_\beta=Z_\beta^{-1}e^{-V_\beta}dm\) with
\(V_\beta=\beta S\) and \(\beta=g^{-2}\), V75 gives, for a fixed
invariant derivative \(X\),
\begin{equation}
 \|XV_\beta\|_{L^{2m}(\mu_\beta)}
 \leq C_m\beta^{1/2}=C_mg^{-1}.                      \label{eq:s7v76-score}
\end{equation}

\begin{lemma}[Fixed score multiplicity]
\label{lem:s7v76-score-product}
Let \(N\) be fixed.  A product of at most \(N\) compact scores and a
fixed finite derivative list has a reference fourth moment bounded by
\begin{equation}
 t_{0,M,\Lambda}^{(N)}
 \leq
 C_Ng^{-N}
 (1+M)^{d_M}(1+\Lambda)^{d_\Lambda}                  \label{eq:s7v76-tN}
\end{equation}
for fixed exponents.
\end{lemma}

\begin{proof}
Apply generalized H\"older with exponent \(4N\) to the score factors.
The full even-moment hierarchy in
\Cref{thm:s7v75-Fisher} bounds each required score norm by a constant
times \(g^{-1}\).  Apply the V75 Peter--Weyl and oscillator word bounds
to the remaining finite derivative factors.
\end{proof}

\begin{theorem}[Correct bad-score absorption condition]
\label{thm:s7v76-absorb}
Suppose
\begin{equation}
 b_M\leq C(1+M)^d e^{-cM},\qquad c>0,                 \label{eq:s7v76-bad}
\end{equation}
and use \eqref{eq:s7v76-schedule} with \(\chi>1\).  Then for every fixed
\(N,d_M,d_\Lambda\),
\begin{equation}
 t_{0,M,\Lambda}^{(N)}b_M\longrightarrow0.           \label{eq:s7v76-score-absorb}
\end{equation}
If only \(0<\chi\leq1\) is assumed, this conclusion does not follow
from \eqref{eq:s7v76-bad} for arbitrary fixed \(N\).
\end{theorem}

\begin{proof}
After using \(\Lambda\leq C(1+M)^s\), the logarithm of the left side is
at most
\begin{equation}
 -c(\log L)^\chi+{N\over2}\log L+O(\log\log L).
 \label{eq:s7v76-log}
\end{equation}
For \(\chi>1\), the negative first term dominates and the expression
tends to \(-\infty\).  For \(\chi=1\), the sign depends on
\(c-N/2\), and for \(\chi<1\) the positive \(N\log L/2\) dominates.
Thus the weaker schedule alone cannot prove the general statement.
\end{proof}

\begin{corollary}[The schedule strengthening preserves all prior rows]
\label{cor:s7v76-compatible}
Replacing \(\chi>0\) by any fixed \(\chi>1\) in V55 does not spoil the
filter, weak-field, occupation, chart-exit, Feshbach, analytic-tail, or
diagonal cutoff-independence rates proved in V55--V62.
\end{corollary}

\begin{proof}
All those estimates use only that every fixed power of \(\log L\) is
dominated by every positive power of \(L\), that \(M,R\to\infty\), and
that \(gR\to0\).  These remain true for every fixed \(\chi>1\).
The V62 diagonal estimate improves as \(\chi\) increases.
\end{proof}

\section{Why the full RL4C still needs a normal form}
\label{sec:s7v76-multiplicity}

\begin{firewall}[One rarity factor cannot pay arbitrary score
multiplicity]
\label{fire:s7v76-multiplicity}
The event \(\{b\hbox{ is bad}\}\) occurs once per spatial block.  If an
expansion allows \(N_b\) score insertions on that same block, the bound
\eqref{eq:s7v76-tN} costs \(g^{-N_b}\), while the probability contributes
only one \(b_M\).  No fixed choice of \(\chi\) absorbs this uniformly as
\(N_b\to\infty\).  Therefore
\Cref{thm:s7v76-absorb} closes only bounded marked multiplicity; it is
not a proof of full RL4C.
\end{firewall}

\begin{definition}[Bounded marked-insertion normal form]
\label{def:s7v76-BMINF}
BMINF is a representation of each connected physical activity in which:
\begin{enumerate}[label=\textup{(B\arabic*)},leftmargin=4em]
\item every unbounded electric principal part is retained in a closed
      reference form or a form-normalized resolvent factor;
\item every bad spatial block carries at most \(N_*\) explicit compact
      score/derivative insertions, with \(N_*\) independent of
      perturbative order and regulator;
\item further local occurrences are resummed into a bounded semigroup,
      resolvent, or normalized conditional density; and
\item the resummation preserves exact local ownership and the weighted
      support incidence of \Cref{cor:s7v76-support}.
\end{enumerate}
\end{definition}

\begin{proposition}[Conditional closure of RL4C]
\label{prop:s7v76-RL4C}
Assume BMINF, regulator-uniform fixed-width SATC, the V64 replica
density card, and the V70 bad activity
\eqref{eq:s7v76-bad}.  Use \(\chi>1\).  Then every ledger row satisfies
the RL4C reference fourth-moment and smallness conditions.
\end{proposition}

\begin{proof}
Rows with positive coupling power vanish by
\Cref{thm:s7v76-positive}.  Exact reference/slow rows are not charged as
activities by \Cref{thm:s7v76-zero}.  BMINF bounds the score multiplicity
by \(N_*\), so
\Cref{thm:s7v76-absorb} pays every marked bad insertion.  Form-resummed
electric rows incur only the declared relative SATC and V43 factors.
Finite incidence is \Cref{cor:s7v76-support}, and the replica density
factor is V64.  These are exactly RL4C rows (F1)--(F6).
\end{proof}

\begin{firewall}[BMINF is not inferred from second order]
\label{fire:s7v76-BMINF}
The Hamiltonian being second order bounds the derivative order of one
elementary generator.  It does not bound how many times that generator
can occur in a Duhamel series.  BMINF therefore requires an actual
resummation theorem.  Treating it as automatic would repeat the
multiplicity error isolated above.
\end{firewall}

\section{Status and next acquisition}
\label{sec:s7v76-status}

\begin{theorem}[What V76 proves]
\label{thm:s7v76-result}
For the declared finite-block Wilson, chart, and transition operations,
V76 gives an exhaustive generator-type ledger and proves:
\begin{enumerate}[label=\textup{(\roman*)},leftmargin=3.5em]
\item no localized electric form requires more than first derivatives
      of the chart/frame/projector;
\item every positive-coupling finite-order row vanishes under the V55
      logarithmic window;
\item every zero-coupling growing row must remain in the exact
      reference/slow Hamiltonian or be treated as an observable
      insertion;
\item fixed compact-score multiplicity is paid by bad-block rarity when
      \(\chi>1\); and
\item arbitrary same-block multiplicity is not controlled by a
      one-block rarity factor.
\end{enumerate}
\end{theorem}

\begin{proof}
Use
\Cref{lem:s7v76-form-Leibniz,thm:s7v76-ledger,
thm:s7v76-positive,thm:s7v76-zero,
thm:s7v76-absorb,fire:s7v76-multiplicity}.
\end{proof}

\begin{assessment}[Progress at seventh-series V76]
\label{ass:s7v76-status}
The physical finite-block generator family is no longer an unspecified
list, and an incorrect schedule inference has been removed.  The
remaining RL4C obstruction is sharply localized: it is not derivative
order or cutoff growth, but repeated unbounded insertions charged to one
rare block.  The upstream repair is to resum relative electric forms
before the bad-block expansion.
\end{assessment}

\begin{decision}[V77 acquisition order]
\label{dec:s7v76-V77}
V77 will construct a symmetric form-resolvent normal form.  For a
relative perturbation \(v\), it will factor
\[
 H+\sigma
 =
 (H_0+\sigma)^{1/2}(I+T_\sigma)
 (H_0+\sigma)^{1/2}
\]
and resum arbitrary insertions through
\((I+T_\sigma)^{-1}\).  It will then prove a weighted block-Schur
version, so local relative electric and atlas forms cost a geometric
factor rather than one score moment per Duhamel occurrence.  This is the
candidate proof of BMINF row (B1); if weighted locality fails, the
failure will be stated as the next explicit loader.
\end{decision}

\begin{firewall}[Claim boundary after seventh-series V76]
\label{fire:s7v76-claim}
V76 does not prove BMINF, full RL4C, regulator-uniform SATC, the
interacting joined slow-sector gap, the remaining regular-sector UCTBB
data, interacting endpoint continuity, UCM, BCC, the
pairing/Bogoliubov sector, weighted physical-source incidence, nonlinear
coarse-action cocycles, terminal strong-coupling matching, regulator
independence, cutoff removal, Osterwalder--Schrader reconstruction, or
the full Yang--Mills mass gap.
\end{firewall}

\paragraph{Parent-version record.}
V76 was formed from
\texttt{Renewal\_Yang\_Mills\_Mass\_Gap\_Research\_Notes\_V75.tex}.
The V75 parent had \(33\,082\) logical source lines,
\(1\,180\,460\) bytes, and SHA--256
\[
 \texttt{FBFE5A350A6328F5D555ACAEC9C02288E4A29659521F5DE995FFA27E8656D2F1}.
\]
The next compulsory companion audit is V80.

% =============================================================================
% END APPENDED SEVENTH-SERIES V76 MATERIAL
% =============================================================================

% =============================================================================
% BEGIN APPENDED SEVENTH-SERIES V77 MATERIAL
% =============================================================================

\chapter{Symmetric form-resolvent resummation and weighted locality}
\label{ch:s7v77}

\section{Purpose}
\label{sec:s7v77-purpose}

V76 isolates an artificial obstruction: if a relative electric form is
expanded occurrence by occurrence, every occurrence can generate
unbounded derivative or score factors, although the complete form is
small relative to the Gaussian energy.  The correct operation is to
normalize by that energy and invert once.

This chapter proves the exact symmetric form-resolvent factorization,
its convergent Neumann resummation, and a weighted block-Schur version.
It also gives a sufficient local factorization which converts
energy-normalized Gaussian propagation into weighted locality of the
resummed perturbation.  These theorems close the resolvent part of BMINF
without imposing a bound on Duhamel multiplicity.  A separate
semigroup/normalized-correlation transfer is still required for the
Euclidean RL4C expansion.

\section{Energy normalization of a relative form}
\label{sec:s7v77-normalization}

Let \(H_0\geq0\) be self-adjoint on a Hilbert space \({\cal H}\), with
closed form \(h_0\) and form domain
\({\cal Q}=D(H_0^{1/2})\).  Fix \(\sigma>0\) and write
\begin{equation}
 h_{0,\sigma}[u,v]
 :=
 h_0[u,v]+\sigma\langle u,v\rangle.                  \label{eq:s7v77-hsigma}
\end{equation}
Let \(v\) be a symmetric form on \({\cal Q}\) satisfying
\begin{equation}
 |v[u,u]|
 \leq\theta h_{0,\sigma}[u,u],
 \qquad 0\leq\theta<1.                               \label{eq:s7v77-relative}
\end{equation}

\begin{lemma}[Polarized relative bound]
\label{lem:s7v77-polarization}
There is a bounded self-adjoint operator \(T_\sigma\) on \({\cal H}\)
such that
\begin{equation}
 v[u,v]
 =
 \left\langle
 (H_0+\sigma)^{1/2}u,\,
 T_\sigma(H_0+\sigma)^{1/2}v
 \right\rangle,                                      \label{eq:s7v77-T}
\end{equation}
and
\begin{equation}
 \|T_\sigma\|\leq\theta.                              \label{eq:s7v77-T-bound}
\end{equation}
\end{lemma}

\begin{proof}
Equip \({\cal Q}\) with the Hilbert norm
\(\|u\|_{0,\sigma}^2=h_{0,\sigma}[u,u]\).
The symmetric quadratic estimate
\eqref{eq:s7v77-relative}, applied to
\(u+t v\) and optimized in real and imaginary \(t\), gives
\[
 |v[u,v]|
 \leq\theta\|u\|_{0,\sigma}\|v\|_{0,\sigma}.
\]
Thus \(v\) is represented by a bounded self-adjoint operator of norm at
most \(\theta\) on this form Hilbert space.  The map
\((H_0+\sigma)^{1/2}:{\cal Q}\to{\cal H}\) is unitary when
\({\cal Q}\) carries this norm.  Transport the representing operator
through that unitary to obtain \(T_\sigma\).
\end{proof}

\begin{theorem}[Exact symmetric form-resolvent factorization]
\label{thm:s7v77-factor}
Let \(H\) be the self-adjoint operator associated with the closed form
\(h=h_0+v\).  Under \eqref{eq:s7v77-relative},
\begin{equation}
 (H+\sigma)^{-1}
 =
 (H_0+\sigma)^{-1/2}
 (I+T_\sigma)^{-1}
 (H_0+\sigma)^{-1/2}.                                \label{eq:s7v77-resolvent}
\end{equation}
Moreover,
\begin{align}
 (I+T_\sigma)^{-1}
 &=\sum_{n=0}^{\infty}(-T_\sigma)^n,                 \label{eq:s7v77-Neumann}\\
 \|(I+T_\sigma)^{-1}\|
 &\leq{1\over1-\theta},                              \label{eq:s7v77-inverse-bound}\\
 \left\|
 (I+T_\sigma)^{-1}
 -\sum_{n=0}^{N}(-T_\sigma)^n
 \right\|
 &\leq{\theta^{N+1}\over1-\theta}.                   \label{eq:s7v77-tail}
\end{align}
\end{theorem}

\begin{proof}
\Cref{lem:s7v77-polarization} gives the form identity
\[
 h[u,v]+\sigma\langle u,v\rangle
 =
 \left\langle
 (H_0+\sigma)^{1/2}u,\,
 (I+T_\sigma)(H_0+\sigma)^{1/2}v
 \right\rangle.
\]
Since \(T_\sigma=T_\sigma^*\) and \(\|T_\sigma\|<1\),
\(I+T_\sigma\geq(1-\theta)I\).  The representation theorem for closed
forms then gives \eqref{eq:s7v77-resolvent}.  The remaining statements
are the norm-convergent geometric series and its tail bound.
\end{proof}

\begin{corollary}[Arbitrary form multiplicity is resummed]
\label{cor:s7v77-multiplicity}
All orders in the relative form \(v\) contribute to the middle factor
in \eqref{eq:s7v77-resolvent} with total norm at most
\((1-\theta)^{-1}\).  No estimate of the form
\(C^ng^{-n}\) is incurred.
\end{corollary}

\begin{proof}
The \(n\)-th occurrence is the \(n\)-th Neumann term in
\eqref{eq:s7v77-Neumann}.  Sum its norm by
\eqref{eq:s7v77-inverse-bound}.
\end{proof}

\begin{firewall}[The normalization is essential]
\label{fire:s7v77-raw}
The series is in
\[
 T_\sigma
 =
 (H_0+\sigma)^{-1/2}v(H_0+\sigma)^{-1/2}
\]
in the form sense, not in the generally unbounded raw operator \(v\).
Replacing \(T_\sigma\) by a cutoff norm of \(vP_M\) loses the two energy
denominators and recreates the spurious factor \(M\) identified in
\Cref{rem:s7v76-raw-form}.
\end{firewall}

\section{Weighted block-Schur algebra}
\label{sec:s7v77-Schur}

Let
\({\cal H}=\bigoplus_{x\in{\mathbb X}}{\cal H}_x\),
where \(({\mathbb X},d)\) is a countable metric graph.  For an operator
matrix \(A=(A_{xy})\), define
\begin{equation}
 \|A\|_{{\mathfrak S}_\mu}
 :=
 \max\left\{
 \sup_x\sum_y e^{\mu d(x,y)}\|A_{xy}\|,
 \sup_y\sum_x e^{\mu d(x,y)}\|A_{xy}\|
 \right\}.                                           \label{eq:s7v77-Schur}
\end{equation}

\begin{lemma}[Weighted Schur Banach star algebra]
\label{lem:s7v77-Schur}
Whenever the displayed norms are finite,
\begin{align}
 \|AB\|_{{\mathfrak S}_\mu}
 &\leq
 \|A\|_{{\mathfrak S}_\mu}
 \|B\|_{{\mathfrak S}_\mu},                          \label{eq:s7v77-product}\\
 \|A^*\|_{{\mathfrak S}_\mu}
 &=\|A\|_{{\mathfrak S}_\mu},                        \label{eq:s7v77-adjoint}\\
 \|A\|_{{\cal B}({\cal H})}
 &\leq\|A\|_{{\mathfrak S}_\mu}.                     \label{eq:s7v77-operator}
\end{align}
\end{lemma}

\begin{proof}
The triangle inequality gives
\[
 e^{\mu d(x,z)}
 \leq
 e^{\mu d(x,y)}e^{\mu d(y,z)}.
\]
Insert this in the row sum of
\((AB)_{xz}=\sum_yA_{xy}B_{yz}\), then take the two suprema.
Adjunction interchanges the row and column suprema.  The last statement
is the ordinary two-sided Schur test, since the exponential weight is at
least one.
\end{proof}

\begin{theorem}[Weighted inverse and tail]
\label{thm:s7v77-weighted-inverse}
If
\begin{equation}
 \|T_\sigma\|_{{\mathfrak S}_\mu}
 \leq\vartheta_\mu<1,                                \label{eq:s7v77-vartheta}
\end{equation}
then
\begin{align}
 \|(I+T_\sigma)^{-1}\|_{{\mathfrak S}_\mu}
 &\leq{1\over1-\vartheta_\mu},                       \label{eq:s7v77-weighted-bound}\\
 \left\|
 (I+T_\sigma)^{-1}
 -\sum_{n=0}^{N}(-T_\sigma)^n
 \right\|_{{\mathfrak S}_\mu}
 &\leq{\vartheta_\mu^{N+1}\over1-\vartheta_\mu}.
                                                               \label{eq:s7v77-weighted-tail}
\end{align}
In particular, arbitrary chains of local form insertions preserve
exponential spatial summability.
\end{theorem}

\begin{proof}
Use the Neumann series in the Banach algebra of
\Cref{lem:s7v77-Schur} and sum the geometric majorant.
\end{proof}

\section{Loading weighted locality from Gaussian propagation}
\label{sec:s7v77-loading}

The energy normalizers in \(T_\sigma\) are not strictly local.  Their
weighted decay must be proved from the Gaussian reference rather than
assumed to follow from the locality of \(v\).

\begin{lemma}[Square-root resolvent from a weighted semigroup]
\label{lem:s7v77-square-root}
Suppose
\begin{equation}
 \|e^{-tH_0}\|_{{\mathfrak S}_\mu}
 \leq C_0e^{-\gamma_0t},
 \qquad t\geq0.                                      \label{eq:s7v77-semigroup}
\end{equation}
Then
\begin{equation}
 \|(H_0+\sigma)^{-1/2}\|_{{\mathfrak S}_\mu}
 \leq{C_0\over\sqrt{\gamma_0+\sigma}}.               \label{eq:s7v77-sqrt}
\end{equation}
\end{lemma}

\begin{proof}
Functional calculus gives the norm-convergent improper integral
\[
 (H_0+\sigma)^{-1/2}
 =
 {1\over\sqrt\pi}
 \int_0^\infty t^{-1/2}e^{-t\sigma}e^{-tH_0}\,dt.
\]
Use \eqref{eq:s7v77-semigroup}, integrate in the weighted Banach space,
and apply
\(\int_0^\infty t^{-1/2}e^{-ct}\,dt=\sqrt{\pi/c}\).
\end{proof}

Let the relative form admit a local energy-factor decomposition
\begin{equation}
 v[u,v]
 =
 \sum_{X\in{\cal X}}
 \langle C_Xu,J_XC_Xv\rangle,                        \label{eq:s7v77-factor-local}
\end{equation}
where \(J_X=J_X^*\) and \(C_X\) is supported in the fixed neighborhood
\(X\).  Put
\[
 B_X:=C_X(H_0+\sigma)^{-1/2}.
\]

\begin{definition}[Weighted energy-factor card, WEFC]
\label{def:s7v77-WEFC}
WEFC at weight \(\mu\) means that the block operators \(B_X\) reconstruct
the form-normalized operator
\begin{equation}
 T_\sigma=\sum_XB_X^*J_XB_X                          \label{eq:s7v77-T-local}
\end{equation}
and satisfy
\begin{equation}
 \left\|
 \sum_XB_X^*J_XB_X
 \right\|_{{\mathfrak S}_\mu}
 \leq\vartheta_\mu<1,                                \label{eq:s7v77-WEFC}
\end{equation}
uniformly in volume, cutoff, and retained slow labels.
\end{definition}

\begin{proposition}[A checkable sufficient bound for WEFC]
\label{prop:s7v77-WEFC}
Suppose the index \(X\) has overlap at most \(n_*\), the local factors
satisfy
\begin{equation}
 \sup_X\|J_X\|\leq j_*,
 \qquad
 \left\|\bigoplus_X B_X\right\|_{{\mathfrak S}_\mu}
 \leq b_\mu,                                         \label{eq:s7v77-B-bound}
\end{equation}
with the direct-sum incidence included in the norm.  Then
\begin{equation}
 \|T_\sigma\|_{{\mathfrak S}_\mu}
 \leq j_*b_\mu^2.                                    \label{eq:s7v77-WEFC-bound}
\end{equation}
Thus \(j_*b_\mu^2<1\) implies WEFC and the weighted resummation
\eqref{eq:s7v77-weighted-bound}.
\end{proposition}

\begin{proof}
Write
\[
 T_\sigma=B^*JB,
\qquad
 B=\bigoplus_XB_X,\qquad J=\bigoplus_XJ_X.
\]
The bounded overlap is part of the block incidence defining \(B\).
\Cref{lem:s7v77-Schur} gives
\[
 \|B^*JB\|_{{\mathfrak S}_\mu}
 \leq
 \|B^*\|_{{\mathfrak S}_\mu}
 \|J\|_{{\mathfrak S}_\mu}
 \|B\|_{{\mathfrak S}_\mu}
 \leq j_*b_\mu^2.
\]
\end{proof}

\begin{firewall}[An ordinary relative bound is not WEFC]
\label{fire:s7v77-WEFC}
Equation \eqref{eq:s7v77-relative} implies
\(\|T_\sigma\|\leq\theta\), but it does not imply the exponentially
weighted estimate \eqref{eq:s7v77-WEFC}.  WEFC additionally uses
weighted propagation of the Gaussian energy normalizers and exact local
incidence.  This is the same type distinction already exposed for HOLC
in V59.
\end{firewall}

\section{Wilson and atlas relative coefficient}
\label{sec:s7v77-Wilson}

On the V55 window \(R=\sqrt M\), collect the V43 weak-field relative
form, the V39 scalar IMS debit relative to the residual number floor,
and the fixed-width V72 atlas terms into
\begin{equation}
 \theta_{\rm form}
 \leq
 C\left(
 g\sqrt M+M^{-2}+g^2
 +g^2d_{\rm sp}^{-4}
 \right).                                            \label{eq:s7v77-theta-form}
\end{equation}
Here \(M^{-2}\) is the ratio of the V43 scalar debit
\(c_R=O(M^{-1})\) to the high-occupation floor \(O(M)\).

\begin{proposition}[Numerical closure of form resummation]
\label{prop:s7v77-rates}
Under
\[
 g=O(L^{-1/2}),\qquad M=(\log L)^\chi,
\]
with fixed \(\chi>1\) and fixed spectral buffer \(d_{\rm sp}>0\),
\begin{equation}
 \theta_{\rm form}
 =
 O\!\left(
 L^{-1/2}(\log L)^{\chi/2}
 +(\log L)^{-2\chi}
 +L^{-1}
 \right)
 \longrightarrow0.                                  \label{eq:s7v77-rates}
\end{equation}
Consequently the unweighted form inverse costs
\(1+o(1)\).
\end{proposition}

\begin{proof}
Insert the schedule into
\eqref{eq:s7v77-theta-form}.  Every term tends to zero.  Apply
\eqref{eq:s7v77-inverse-bound}.
\end{proof}

\begin{remark}[Two different residual comparisons]
\label{rem:s7v77-residuals}
Against the fixed transverse oscillator floor \(O(a^{-1})\), the V43
IMS coefficient is \(O(M^{-1})\), as in V55.  Against the
high-occupation floor \(O(M/a)\), it is \(O(M^{-2})\), as used in
\eqref{eq:s7v77-theta-form}.  Neither value should be substituted for
the other without specifying the residual sector.
\end{remark}

\begin{corollary}[Resolvent BMINF row]
\label{cor:s7v77-BMINF}
If WEFC holds with
\(\vartheta_\mu=O(\theta_{\rm form})\), then BMINF row (B1) is proved
for the resolvent representation: all electric metric, Haar, chart IMS,
atlas connection, and projector-response occurrences are resummed into
a weighted factor of norm \(1+o(1)\).
\end{corollary}

\begin{proof}
Use
\Cref{thm:s7v77-weighted-inverse,prop:s7v77-rates}.
The list of resummed form types is exhaustive by
\Cref{thm:s7v76-ledger}.
\end{proof}

\section{Spectral handoff}
\label{sec:s7v77-gap}

\begin{theorem}[Relative-form preservation of a residual gap]
\label{thm:s7v77-gap}
Let \(Q\) be a reducing projection for \(H_0\) such that
\begin{equation}
 QH_0Q\geq{\gamma_0\over a}Q.                        \label{eq:s7v77-gap0}
\end{equation}
Suppose on \(Q{\cal Q}\)
\begin{equation}
 |v[u,u]|
 \leq\eta h_0[u,u]+{c\over a}\|u\|^2,
 \qquad 0\leq\eta<1.                                 \label{eq:s7v77-gap-relative}
\end{equation}
Then
\begin{equation}
 QHQ
 \geq
 {[(1-\eta)\gamma_0-c]\over a}Q                      \label{eq:s7v77-gap-transfer}
\end{equation}
in form sense.  In particular the residual gap remains positive if
\(\eta+c/\gamma_0<1\).
\end{theorem}

\begin{proof}
For \(u=Qu\),
\[
 h[u,u]\geq(1-\eta)h_0[u,u]-{c\over a}\|u\|^2.
\]
Insert \eqref{eq:s7v77-gap0}.
\end{proof}

\begin{firewall}[Residual gap is not the physical slow gap]
\label{fire:s7v77-slow}
\Cref{thm:s7v77-gap} preserves a gap already present on the residual
complement \(Q\).  The joined slow sector contains the physical
long-distance degrees of freedom and has no assumed \(\gamma_0/a\)
floor.  Therefore this theorem does not prove the Yang--Mills mass gap;
it protects the elimination step while leaving the interacting slow
Hamiltonian as a separate problem.
\end{firewall}

\section{Status and next acquisition}
\label{sec:s7v77-status}

\begin{theorem}[What V77 proves]
\label{thm:s7v77-result}
V77 proves an exact energy-normalized resolvent factorization for every
strictly relatively bounded symmetric form, resums arbitrary form
multiplicity with cost \((1-\theta)^{-1}\), and proves the analogous
weighted block-Schur inverse.  It supplies WEFC as a precise sufficient
locality loader and shows that the Wilson/atlas relative coefficient
tends to zero under the corrected logarithmic schedule.
\end{theorem}

\begin{proof}
Use
\Cref{thm:s7v77-factor,thm:s7v77-weighted-inverse,
prop:s7v77-WEFC,prop:s7v77-rates,
cor:s7v77-BMINF,thm:s7v77-gap}.
\end{proof}

\begin{assessment}[Progress at seventh-series V77]
\label{ass:s7v77-status}
Repeated electric insertions no longer require repeated score moments
at the resolvent level.  The remaining locality datum is not a vague
derivative estimate but the weighted energy-factor card WEFC.  The
remaining Euclidean datum is a transfer from the resummed form resolvent
to the normalized semigroup/correlation expansion.
\end{assessment}

\begin{decision}[V78 acquisition order]
\label{dec:s7v77-V78}
V78 will attack WEFC upstream.  It will use the Gaussian spacetime
kernel already proved in V53--V56 to control the energy normalizers,
derive a finite-incidence weighted factor estimate for the V43 and V72
forms, and distinguish the fixed transverse sector from the retained
slow sector.  If the slow block lacks a uniform Gaussian floor, it will
be excluded from WEFC and kept inside the exact joined slow Hamiltonian.
\end{decision}

\begin{firewall}[Claim boundary after seventh-series V77]
\label{fire:s7v77-claim}
V77 does not verify WEFC for the full physical form, transfer the
weighted resolvent estimate to normalized Euclidean correlations, prove
BMINF beyond the resolvent representation, prove full RL4C,
regulator-uniform SATC, the interacting joined slow-sector gap, the
remaining regular-sector UCTBB data, interacting endpoint continuity,
UCM, BCC, the pairing/Bogoliubov sector, nonlinear coarse-action
cocycles, terminal matching, regulator independence, cutoff removal,
Osterwalder--Schrader reconstruction, or the full Yang--Mills mass gap.
\end{firewall}

\paragraph{Parent-version record.}
V77 was formed from
\texttt{Renewal\_Yang\_Mills\_Mass\_Gap\_Research\_Notes\_V76.tex}.
The V76 parent had \(33\,673\) logical source lines,
\(1\,203\,877\) bytes, and SHA--256
\[
 \texttt{848219BA5A8C8658AA62DFFB374DEFE996B8FEAE793F7146F82D30BA61082FBB}.
\]
The next compulsory companion audit is V80.

% =============================================================================
% END APPENDED SEVENTH-SERIES V77 MATERIAL
% =============================================================================

% =============================================================================
% BEGIN APPENDED SEVENTH-SERIES V78 MATERIAL
% =============================================================================

\chapter{The spectator no-go for many-body WEFC and a local form repair}
\label{ch:s7v78}

\section{Purpose and correction}
\label{sec:s7v78-purpose}

V77 proposes WEFC as a weighted operator-Schur estimate for the
energy-normalized form perturbation.  The V46 and V53 Gaussian
propagation theorems do prove that estimate for a one-particle quadratic
perturbation.  They do not prove it on the full bosonic Fock space.
Indeed, the full operator norm sees spectator particles simultaneously
present at two distant supports, and those particles remove any spatial
decay even when the oscillators are completely decoupled.

This chapter proves both statements.  It then replaces the false
many-body operator target by two type-correct mechanisms:
\begin{enumerate}[label=\textup{(\roman*)},leftmargin=3.5em]
\item good local forms are normalized by a local oscillator energy and
      expanded with their actual small activity; and
\item the internal Hamiltonian of a connected bad cluster is retained
      exactly, so its arbitrarily many internal occurrences cost a
      contraction rather than repeated compact scores.
\end{enumerate}
The remaining physical interface is the cross-boundary interaction of
the exact bad cluster with its good exterior.

\section{WEFC is valid in the one-particle quadratic sector}
\label{sec:s7v78-oneparticle}

Let \(A\geq\gamma_0I\) be the finite-range one-particle operator of V53,
and let \(E=E^*\) be a quadratic perturbation with
\(\|E\|_{{\mathfrak S}_\mu}<\infty\).  Set
\begin{equation}
 T^{(1)}_\sigma
 :=
 (A+\sigma)^{-1/2}E(A+\sigma)^{-1/2}.                \label{eq:s7v78-T1}
\end{equation}

\begin{lemma}[Weighted square root with rate loss]
\label{lem:s7v78-sqrt}
Choose \(0<\mu<\mu_*\), where \(\mu_*\) is a V53 half-gap rate.  Then
\begin{equation}
 \|(A+\sigma)^{-1/2}\|_{{\mathfrak S}_\mu}
 \leq
 {Q_{\mu_*-\mu}\over\sqrt{\sigma+\gamma_0/2}},
 \qquad
 Q_\delta:=
 \sup_x\sum_y e^{-\delta d(x,y)}.                    \label{eq:s7v78-sqrt}
\end{equation}
\end{lemma}

\begin{proof}
\Cref{cor:s7v53-half} gives the pointwise bound
\[
 \|(e^{-tA})_{xy}\|
 \leq e^{-\mu_*d(x,y)}e^{-\gamma_0t/2}.
\]
Multiply by \(e^{\mu d(x,y)}\), sum rows and columns, and obtain
\[
 \|e^{-tA}\|_{{\mathfrak S}_\mu}
 \leq
 Q_{\mu_*-\mu}e^{-\gamma_0t/2}.
\]
Insert this estimate in the square-root integral used in
\Cref{lem:s7v77-square-root}.
\end{proof}

\begin{theorem}[Quadratic one-particle WEFC]
\label{thm:s7v78-oneparticle}
Under the preceding hypotheses,
\begin{equation}
 \|T^{(1)}_\sigma\|_{{\mathfrak S}_\mu}
 \leq
 {Q_{\mu_*-\mu}^2\over\sigma+\gamma_0/2}
 \|E\|_{{\mathfrak S}_\mu}.                          \label{eq:s7v78-T1-bound}
\end{equation}
If the right side is below one, the weighted resolvent resummation of V77
applies to \(A+E\).
\end{theorem}

\begin{proof}
Apply the weighted Schur product theorem
\Cref{lem:s7v77-Schur} twice to
\eqref{eq:s7v78-T1}, then use
\eqref{eq:s7v78-sqrt}.  The last statement is
\Cref{thm:s7v77-weighted-inverse}.
\end{proof}

\begin{corollary}[Closed quadratic rows]
\label{cor:s7v78-quadratic}
Every genuinely quadratic local row whose kernel satisfies the V58
weighted perturbation estimate has WEFC.  In particular, a local
quadratic self-energy or a fixed-rank atlas correction can be inserted
into the one-particle reference whenever its
\({\mathfrak S}_\mu\) norm tends to zero.
\end{corollary}

\begin{proof}
Use \Cref{thm:s7v78-oneparticle}; the V58 hypotheses give the required
small weighted norm uniformly in volume and slow labels.
\end{proof}

\section{A spectator counterexample on Fock space}
\label{sec:s7v78-spectator}

Consider two independent bosonic modes at sites \(x\ne y\), with
\begin{equation}
 H_0=N_x+N_y,\qquad
 R_\sigma=(H_0+\sigma)^{-1},\qquad \sigma>0.          \label{eq:s7v78-two}
\end{equation}
The Hamiltonian has no hopping at all.  Let
\[
 C_x=a_x,\qquad C_y=a_y.
\]

\begin{theorem}[Spectator no-go]
\label{thm:s7v78-spectator}
The cross operator
\begin{equation}
 K_{yx}:=C_xR_\sigma C_y^*                           \label{eq:s7v78-K}
\end{equation}
obeys
\begin{equation}
 \|K_{yx}\|\geq{1\over2+\sigma}                      \label{eq:s7v78-lower}
\end{equation}
independently of the distance \(d(x,y)\).  Hence no estimate
\[
 \|K_{yx}\|\leq Ce^{-\mu d(x,y)}
\]
with fixed \(C,\mu>0\) can hold uniformly over arbitrarily distant pairs,
even though the one-particle hopping is identically zero.
\end{theorem}

\begin{proof}
Let
\[
 |1_x,0_y\rangle=a_x^*\Omega.
\]
Then
\begin{align*}
 C_y^*|1_x,0_y\rangle
 &=|1_x,1_y\rangle,\\
 R_\sigma|1_x,1_y\rangle
 &={1\over2+\sigma}|1_x,1_y\rangle,\\
 C_x|1_x,1_y\rangle
 &=|0_x,1_y\rangle.
\end{align*}
Both the input and output vectors are normalized, so
\eqref{eq:s7v78-lower} follows.
\end{proof}

\begin{corollary}[Full-Fock operator WEFC is false]
\label{cor:s7v78-WEFC-no}
The ordinary many-body operator blocks obtained by placing local form
factors at \(x\) and \(y\) need not satisfy
\eqref{eq:s7v77-WEFC}, even for a gapped, decoupled quadratic
Hamiltonian.  Therefore the V53 one-particle semigroup bound cannot load
full-Fock WEFC by a Schur-product argument.
\end{corollary}

\begin{proof}
The kernel \eqref{eq:s7v78-K} is exactly a cross term of the form
\(C_x(H_0+\sigma)^{-1}C_y^*\).  Its lower bound contradicts exponential
decay as \(d(x,y)\to\infty\).
\end{proof}

\begin{remark}[The particle at \(x\) is the spectator]
\label{rem:s7v78-spectator}
The resolvent does not propagate a quantum from \(y\) to \(x\).  The
input already contains a quantum at \(x\); \(a_y^*\) creates a separate
one at \(y\), and \(a_x\) removes the spectator.  Connected vacuum
correlations cancel precisely this mechanism.  Full operator norm does
not.
\end{remark}

\begin{firewall}[Correction to the V77 WEFC plan]
\label{fire:s7v78-correction}
\Cref{cor:s7v77-BMINF} remains a correct conditional implication, but
full-Fock WEFC is not an admissible physical loader.  WEFC may be used
for a one-particle quadratic correction or for a connected/vacuum-reduced
kernel.  It may not be inferred for nonlinear many-body form factors
from Gaussian one-particle locality.
\end{firewall}

\section{Strictly local energy normalization on good blocks}
\label{sec:s7v78-local}

Let the Gaussian form be written as a bounded-overlap sum
\begin{equation}
 h_0=\sum_{X}h_{0,X},\qquad h_{0,X}\geq0,             \label{eq:s7v78-hlocal}
\end{equation}
where \(X\) ranges over fixed-radius plaquette stars.  Let \(v_X\) be a
local symmetric form supported in \(X\) and suppose
\begin{equation}
 |v_X[u,u]|
 \leq\eta_X\bigl(h_{0,X}[u,u]+\sigma_X\|u\|^2\bigr). \label{eq:s7v78-local-relative}
\end{equation}

\begin{lemma}[Local bounded form representative]
\label{lem:s7v78-local-representative}
On the local Hilbert factor there is a self-adjoint contraction-scale
operator \(S_X\) such that
\begin{equation}
 v_X[u,v]
 =
 \left\langle
 (H_{0,X}+\sigma_X)^{1/2}u,\,
 S_X(H_{0,X}+\sigma_X)^{1/2}v
 \right\rangle,\qquad
 \|S_X\|\leq\eta_X.                                  \label{eq:s7v78-SX}
\end{equation}
Extended by the identity outside \(X\), \(S_X\) is strictly local.
\end{lemma}

\begin{proof}
Apply \Cref{lem:s7v77-polarization} on the finite local factor with
\(H_0\) replaced by \(H_{0,X}\).  Tensor extension does not change the
operator norm or support.
\end{proof}

Let \(P_{M,X}\) be the local number cutoff and assume the local oscillator
frequencies are bounded above by \(\omega^*\).  Then
\begin{equation}
 P_{M,X}H_{0,X}P_{M,X}
 \leq{C_X(M+1)\over a}P_{M,X}.                       \label{eq:s7v78-local-upper}
\end{equation}

\begin{corollary}[Compressed local form activity]
\label{cor:s7v78-local-activity}
The dimensionless compressed local perturbation satisfies
\begin{equation}
 a\|P_{M,X}V_XP_{M,X}\|
 \leq C_X\eta_X(M+1).                                \label{eq:s7v78-local-activity}
\end{equation}
For the V43 principal electric form,
\(\eta_X=O(gR)\), \(R=\sqrt M\), and hence
\begin{equation}
 a\|P_{M,X}V_XP_{M,X}\|
 =
 O(gM^{3/2}),                                        \label{eq:s7v78-local-rate}
\end{equation}
the same leading activity as the cubic Wilson vertex.
\end{corollary}

\begin{proof}
Use \eqref{eq:s7v78-SX} and the local energy upper bound on both sides.
The two square roots together contribute \(C_X(M+1)/a\).  Insert
\(\eta_X=O(g\sqrt M)\).
\end{proof}

\begin{theorem}[Repeated good-form occurrences are paid per occurrence]
\label{thm:s7v78-good-repeat}
Suppose the V57 local Weyl representation assigns a fixed polynomial
loss \(C(1+M)^p\) to each compressed local form occurrence.  Define
\begin{equation}
 q_{{\rm form},M}
 :=
 C(1+M)^p\,\eta_X(M+1).                              \label{eq:s7v78-qform}
\end{equation}
Under \(g=O(L^{-1/2})\), \(M=(\log L)^\chi\), and
\(\eta_X=O(g\sqrt M)\),
\begin{equation}
 q_{{\rm form},M}
 =
 O\!\left(
 L^{-1/2}(\log L)^{\chi(p+3/2)}
 \right)\longrightarrow0.                           \label{eq:s7v78-qform-rate}
\end{equation}
An \(n\)-occurrence connected coefficient is therefore bounded by the
same tree majorant with a factor \(q_{{\rm form},M}^n\); arbitrary
same-block multiplicity is summable once \(q_{{\rm form},M}<1\).
\end{theorem}

\begin{proof}
\Cref{cor:s7v78-local-activity} gives one factor
\(\eta_X(M+1)\); the assumed V57 representation gives the remaining
polynomial.  Substitute the schedule.  The V57 degree-weighted tree
theorem is multiplicative over interaction occurrences, including
occurrences with the same spatial owner, and its normalized factorial
summation closes whenever the per-occurrence activity is below its
fixed convergence threshold.
\end{proof}

\begin{firewall}[This is a good-field statement]
\label{fire:s7v78-good}
On a bad block there is no small coefficient
\(\eta_X=O(gR)\), because the weak-field chart hypothesis has failed.
\Cref{thm:s7v78-good-repeat} cannot be extended to that block by
replacing \(\eta_X\) with a compact supremum.  The bad internal
Hamiltonian must instead be retained exactly.
\end{firewall}

\section{Exact resummation inside a bad cluster}
\label{sec:s7v78-bad}

Let \(B\) be a finite connected set of bad blocks enlarged by the fixed
ownership collar, and let \(H_B\) contain every electric and Wilson term
whose declared owner and full support lie in that enlarged cluster.
Subtract its ground energy
\[
 E_B:=\inf\operatorname{spec}H_B,\qquad
 \widehat H_B:=H_B-E_B.
\]

\begin{lemma}[Bad-cluster contraction]
\label{lem:s7v78-bad-contraction}
For every \(t\geq0\),
\begin{equation}
 \|e^{-t\widehat H_B}\|\leq1.                         \label{eq:s7v78-bad-contraction}
\end{equation}
This bound is independent of the internal Duhamel multiplicity, local
representation cutoff, and value of \(g\).
\end{lemma}

\begin{proof}
\(\widehat H_B\geq0\) by the spectral theorem.  Hence
\(\|e^{-t\widehat H_B}\|
=\sup_{\lambda\in\operatorname{spec}\widehat H_B}e^{-t\lambda}\leq1\).
\end{proof}

\begin{definition}[Bad-cluster resummed form card, BCRF]
\label{def:s7v78-BCRF}
BCRF consists of:
\begin{enumerate}[label=\textup{(C\arabic*)},leftmargin=4em]
\item saturation of every bad component by a fixed ownership collar;
\item the exact internal contraction
      \eqref{eq:s7v78-bad-contraction};
\item a decomposition of all terms crossing the saturated boundary into
      owned local activities with a summable reference moment;
\item a ground-energy normalization whose cluster dependence cancels
      between numerator and partition function; and
\item uniformity over slow/transition labels and finite regulators.
\end{enumerate}
\end{definition}

\begin{proposition}[What BCRF would repair]
\label{prop:s7v78-BCRF}
Assume BCRF and the V70 product domination of bad components.  Then
arbitrarily many internal electric and Wilson occurrences in a bad
component contribute the contraction \eqref{eq:s7v78-bad-contraction},
not a factor \(g^{-N_B}\).  Only the finitely owned cross-boundary
activities enter the RL4C insertion factor.
\end{proposition}

\begin{proof}
Group the time-ordered expansion according to maximal connected
saturated bad components.  By (C1), every term either lies wholly inside
one component or is assigned to its boundary.  Resum the internal terms
into \(e^{-t\widehat H_B}\) and apply (C2).  The scalar
\(e^{-tE_B}\) cancels by (C4).  The remaining terms are exactly the
boundary activities of (C3), while component probabilities are supplied
by V70.
\end{proof}

\begin{firewall}[The boundary row is not yet proved]
\label{fire:s7v78-boundary}
The contraction theorem does not make a Wilson plaquette crossing the
bad-cluster boundary small.  Such a plaquette may have coefficient
\(g^{-2}\), and the exterior chart may share its links.  BCRF row (C3)
requires an exact collar/ownership construction and a conditional
boundary moment.  It is the remaining obstruction to using exact bad
cluster dynamics in RL4C.
\end{firewall}

\section{Revised normal-form interface}
\label{sec:s7v78-interface}

\begin{theorem}[Type-correct replacement for full-Fock WEFC]
\label{thm:s7v78-replacement}
The following three-way split is sufficient for the multiplicity part of
BMINF:
\begin{enumerate}[label=\textup{(\alph*)},leftmargin=3.5em]
\item quadratic one-particle corrections obey
      \Cref{thm:s7v78-oneparticle};
\item good nonlinear local forms obey
      \Cref{thm:s7v78-good-repeat}; and
\item bad internal forms obey BCRF.
\end{enumerate}
No full-Fock operator-Schur WEFC is required.
\end{theorem}

\begin{proof}
Part (a) is resummed into the quasi-free reference while preserving its
weighted edge.  Part (b) supplies one vanishing connected activity per
occurrence and hence sums arbitrary multiplicity.  Part (c) resums
arbitrary internal bad multiplicity into a contraction.  These are the
only form rows by the exhaustive V76 ledger.
\end{proof}

\section{Status and next acquisition}
\label{sec:s7v78-status}

\begin{theorem}[What V78 proves]
\label{thm:s7v78-result}
V78 proves quadratic one-particle WEFC from the V53 Gaussian edge,
disproves the corresponding full-Fock operator statement by an explicit
two-mode spectator counterexample, and constructs a type-correct local
repair.  Good local relative forms have the same
\(O(gM^{3/2})\) leading compressed activity as the cubic Wilson vertex
and are summable with every fixed polynomial loss.  Exact internal bad
cluster dynamics is a contraction after ground normalization.
\end{theorem}

\begin{proof}
Use
\Cref{thm:s7v78-oneparticle,thm:s7v78-spectator,
cor:s7v78-local-activity,thm:s7v78-good-repeat,
lem:s7v78-bad-contraction,thm:s7v78-replacement}.
\end{proof}

\begin{assessment}[Progress at seventh-series V78]
\label{ass:s7v78-status}
An overstrong weighted operator target has been removed before it could
become a circular hypothesis.  Quadratic, good nonlinear, and bad
internal forms now have distinct valid treatments.  The next bottleneck
is concrete: saturate bad components by the exact Wilson ownership
collar and prove that all cross-boundary terms can be charged either to a
good small activity or to one enlarged rare component.
\end{assessment}

\begin{decision}[V79 acquisition order]
\label{dec:s7v78-V79}
V79 will construct the saturated bad-component geometry from the exact
V15/V68 collar incidence.  It will prove component product domination
after finite-radius saturation, assign every crossing Wilson plaquette
to one saturated component, and derive the resulting boundary counting
factor.  It will then identify whether the boundary operator is bounded
after cluster resummation or still requires a conditional trace/ground
state moment.
\end{decision}

\begin{firewall}[Claim boundary after seventh-series V78]
\label{fire:s7v78-claim}
V78 does not prove BCRF row (C3) or (C4), full BMINF, full RL4C,
regulator-uniform SATC, the interacting joined slow-sector gap, the
remaining regular-sector UCTBB data, interacting endpoint continuity,
UCM, BCC, the pairing/Bogoliubov sector, nonlinear coarse-action
cocycles, terminal matching, regulator independence, cutoff removal,
Osterwalder--Schrader reconstruction, or the full Yang--Mills mass gap.
\end{firewall}

\paragraph{Parent-version record.}
V78 was formed from
\texttt{Renewal\_Yang\_Mills\_Mass\_Gap\_Research\_Notes\_V77.tex}.
The V77 parent had \(34\,182\) logical source lines,
\(1\,220\,556\) bytes, and SHA--256
\[
 \texttt{18E18847220951E2AC2FEE7F27823402D8D0E0AA8E6D35BB3A21A2B2DC77BB30}.
\]
The next compulsory companion audit is V80.

% =============================================================================
% END APPENDED SEVENTH-SERIES V78 MATERIAL
% =============================================================================

% =============================================================================
% BEGIN APPENDED SEVENTH-SERIES V79 MATERIAL
% =============================================================================

\chapter{Saturated bad components and one-insertion boundary forests}
\label{ch:s7v79}

\section{Purpose}
\label{sec:s7v79-purpose}

V78 replaces arbitrary expansions inside a bad region by its exact
semigroup.  Two points remain: enlarging bad blocks by the ownership
collar must preserve product domination, and interactions crossing the
enlarged boundary must not reintroduce arbitrary multiplicity.

The first point is another bounded-incidence owner theorem.  The second
is solved by interpolating each distinct boundary interaction once and
using a forest derivative: an edge of a forest is never differentiated
twice.  At finite compact or occupation cutoff, the resulting boundary
loss is a fixed power per saturated block and is absorbed by the
exponential bad activity under the corrected \(\chi>1\) schedule.

The construction is exact on the extended tensor product of link
Hilbert spaces.  The Gauss-law projection does not factor over a cut
boundary; it leaves boundary electric-flux labels.  Loading those labels
without an uncontrolled multiplicity is the remaining gluing card.

\section{Finite-radius saturation preserves product domination}
\label{sec:s7v79-saturation}

Let \(({\cal B},d)\) be the coarse block graph, with bounded degree.  For
an integer \(R\geq0\), put
\begin{equation}
 D_R:=\sup_{b\in{\cal B}}\#\{c:d(b,c)\leq R\}<\infty. \label{eq:s7v79-DR}
\end{equation}
For a random bad set
\({\cal F}=\{b:E_b\hbox{ occurs}\}\), define its saturation
\begin{equation}
 {\cal F}^{+R}:=\{x:d(x,{\cal F})\leq R\}.            \label{eq:s7v79-saturation}
\end{equation}

\begin{theorem}[Product domination after saturation]
\label{thm:s7v79-saturation}
Suppose that for every finite \(S\subset{\cal B}\),
\begin{equation}
 \mu\left(\bigcap_{b\in S}E_b\right)
 \leq\pi^{|S|},\qquad0\leq\pi\leq1.                  \label{eq:s7v79-original-dom}
\end{equation}
Then for every finite \(T\subset{\cal B}\),
\begin{equation}
 \mu\{T\subset{\cal F}^{+R}\}
 \leq
 \left(D_R\pi^{1/D_R}\right)^{|T|}.                  \label{eq:s7v79-saturated-dom}
\end{equation}
The right side may be replaced by its minimum with one.
\end{theorem}

\begin{proof}
For each \(x\in T\), the event
\(\{x\in{\cal F}^{+R}\}\) is covered by the at most \(D_R\) events
\(E_b\) with \(d(x,b)\leq R\).  A fixed bad witness \(b\) can cover at
most \(D_R\) members of \(T\).  Apply the bounded-incidence compiler
\Cref{thm:s7v68-owner} with row and column bounds
\(m=\ell=D_R\).
\end{proof}

\begin{corollary}[Wilson rate survives saturation]
\label{cor:s7v79-rate}
If
\[
 \pi_M\leq C M^q e^{-cM},
\]
then
\begin{equation}
 \pi_M^{+R}
 :=
 D_R\pi_M^{1/D_R}
 \leq
 D_RC^{1/D_R}M^{q/D_R}e^{-cM/D_R}
 \longrightarrow0.                                  \label{eq:s7v79-rate}
\end{equation}
The conclusion also holds with the stronger V69
\(\exp(-c\beta)\) rate.
\end{corollary}

\begin{proof}
Insert the assumed bound in
\eqref{eq:s7v79-saturated-dom}.  \(D_R\) is fixed.
\end{proof}

\section{Saturated component geometry}
\label{sec:s7v79-components}

Let every Hamiltonian interaction support have graph diameter at most
\(r_{\rm int}\).  Choose
\begin{equation}
 R\geq r_{\rm int}+R_{\rm own},                       \label{eq:s7v79-R}
\end{equation}
where \(R_{\rm own}\) contains the exact V68 Wilson ownership collar.
Let \(C\) be a connected component of
\({\cal F}^{+R}\).

\begin{lemma}[No interaction touching a bad core crosses its saturated
component]
\label{lem:s7v79-containment}
If an interaction support \(X\) meets
\({\cal F}\cap C\), then
\begin{equation}
 X\subset C.                                         \label{eq:s7v79-containment}
\end{equation}
\end{lemma}

\begin{proof}
Choose \(b\in X\cap{\cal F}\).  Every \(x\in X\) obeys
\(d(x,b)\leq r_{\rm int}\leq R\), so
\(x\in{\cal F}^{+R}\).  Since \(X\) is connected and meets \(C\), all
of \(X\) lies in the same connected component \(C\).
\end{proof}

\begin{corollary}[Every crossing support is fully good]
\label{cor:s7v79-cross-good}
If \(X\) meets both \(C\) and \(C^c\), then
\begin{equation}
 X\cap{\cal F}=\varnothing.                           \label{eq:s7v79-cross-good}
\end{equation}
Thus every chart-dependent nonlinear remainder on a crossing support has
the good-field bounds of V43 and V76.
\end{corollary}

\begin{proof}
Otherwise \(X\) meets a bad core, and
\Cref{lem:s7v79-containment} gives \(X\subset C\), a contradiction.
\end{proof}

\begin{lemma}[Boundary incidence]
\label{lem:s7v79-boundary-count}
There is a fixed constant \(d_\partial\) such that the number of owned
interaction supports meeting both \(C\) and \(C^c\) is at most
\begin{equation}
 |\partial_{\rm int}C|
 \leq d_\partial|C|.                                 \label{eq:s7v79-boundary-count}
\end{equation}
Each boundary support can be assigned to a block of \(C\) with at most
\(d_\partial\) assignments per block.
\end{lemma}

\begin{proof}
Every crossing support contains a block of \(C\).  Assign it to the
least such block in a fixed ordering.  Finite interaction range,
bounded graph degree, and finitely many Wilson/electric types imply a
uniform bound on the number of supports assignable to one block.  Sum
over \(C\).
\end{proof}

\begin{remark}[Relation to the exact \(179\)-by-\(239\) compiler]
\label{rem:s7v79-exact}
The exact V68 constants enter the probability of the original failure
set.  The additional \(D_R\) in
\eqref{eq:s7v79-saturated-dom} is the cost of geometric saturation.
They are distinct finite incidence losses and should not be silently
identified.
\end{remark}

\section{Boundary interpolation on the extended link Hilbert space}
\label{sec:s7v79-interpolation}

Before imposing the Gauss-law projection, partition the link variables
between \(C\) and \(C^c\).  The electric Laplacian is a sum over links.
Every coupling between the two tensor factors is therefore a Wilson
plaquette multiplier.  Group the finitely many crossing plaquettes into
nonnegative operators
\begin{equation}
 Q_e:={1\over ag^2}
 \sum_{p\in{\cal P}_e}
 \bigl(3-\operatorname{ReTr}W_p\bigr)\geq0,           \label{eq:s7v79-Qe}
\end{equation}
where each crossing plaquette occurs in exactly one \(Q_e\).

Let
\begin{equation}
 H(\boldsymbol\lambda)
 =
 \widehat H_C\otimes I
 +I\otimes\widehat H_{C^c}
 +\sum_{e\in\partial_{\rm int}C}\lambda_eQ_e,
 \qquad0\leq\lambda_e\leq1,                          \label{eq:s7v79-Hlambda}
\end{equation}
where the two decoupled Hamiltonians have their separate ground energies
subtracted.  Then \(H(\boldsymbol\lambda)\geq0\).

\begin{lemma}[One derivative of a positive boundary interaction]
\label{lem:s7v79-Duhamel}
At a finite representation or occupation cutoff,
\begin{equation}
 {\partial\over\partial\lambda_e}
 e^{-tH(\boldsymbol\lambda)}
 =
 -\int_0^t
 e^{-(t-s)H(\boldsymbol\lambda)}
 Q_e
 e^{-sH(\boldsymbol\lambda)}\,ds                    \label{eq:s7v79-Duhamel}
\end{equation}
and
\begin{equation}
 \left\|
 {\partial\over\partial\lambda_e}
 e^{-tH(\boldsymbol\lambda)}
 \right\|
 \leq t\|Q_e\|.                                      \label{eq:s7v79-one}
\end{equation}
\end{lemma}

\begin{proof}
The finite cutoff makes all operators bounded.  Differentiate the
exponential by Duhamel's formula.  Positivity of
\eqref{eq:s7v79-Hlambda} gives
\(\|e^{-sH(\boldsymbol\lambda)}\|\leq1\) for \(s\geq0\);
integrate the product norm.
\end{proof}

\begin{theorem}[Distinct-edge mixed derivative]
\label{thm:s7v79-mixed}
For a set \(F\) of \(n\) distinct boundary labels,
\begin{equation}
 \left\|
 \prod_{e\in F}{\partial\over\partial\lambda_e}
 e^{-tH(\boldsymbol\lambda)}
 \right\|
 \leq
 t^n\prod_{e\in F}\|Q_e\|.                           \label{eq:s7v79-mixed}
\end{equation}
Every \(Q_e\) occurs exactly once on the right.
\end{theorem}

\begin{proof}
Iterated Duhamel differentiation gives a sum over the \(n!\) orders of
the distinct insertions.  For each order, the \(n\) insertion times lie
in one ordered simplex of volume \(t^n/n!\).  Every intervening
semigroup is a contraction.  The norm of one ordered term is therefore
at most
\[
 {t^n\over n!}\prod_{e\in F}\|Q_e\|.
\]
Sum the \(n!\) orders.
\end{proof}

\begin{corollary}[Forest bounded multiplicity]
\label{cor:s7v79-forest}
In a BKAR forest interpolation in the variables
\(\{\lambda_e\}\), each chosen forest edge is differentiated once.
Consequently a saturated component is charged at most
\(d_\partial\) boundary insertion norms per one of its blocks, with no
dependence on perturbative order inside the component.
\end{corollary}

\begin{proof}
A forest is a simple graph and contains each edge at most once.  Apply
\Cref{thm:s7v79-mixed} and the ownership/counting rule
\Cref{lem:s7v79-boundary-count}.
\end{proof}

\section{Boundary norm and rarity}
\label{sec:s7v79-absorption}

On the full compact link Hilbert space,
\[
 0\leq3-\operatorname{ReTr}W_p\leq6.
\]
If each \(Q_e\) contains at most \(n_p\) plaquettes, then
\begin{equation}
 a\|Q_e\|
 \leq6n_pg^{-2}.                                     \label{eq:s7v79-Q-bound}
\end{equation}

\begin{theorem}[A fixed boundary power is absorbed]
\label{thm:s7v79-absorb}
Let the saturated bad-block activity obey
\[
 \pi_M^{+R}\leq C M^q e^{-cM}.
\]
For fixed \(t/a\) and fixed \(d_\partial,n_p\), define the worst
dimensionless boundary factor
\begin{equation}
 A_{\partial,M}
 :=
 C_\partial g^{-2d_\partial}(1+M)^{p_\partial}
 (1+\Lambda)^{q_\partial},                           \label{eq:s7v79-A}
\end{equation}
where the polynomial factors include the fixed endpoint and cutoff
losses.  Under
\[
 g=O(L^{-1/2}),\qquad
 M=(\log L)^\chi,\qquad
 \Lambda\leq C(1+M)^s,\qquad\chi>1,
\]
one has
\begin{equation}
 A_{\partial,M}\pi_M^{+R}\longrightarrow0.           \label{eq:s7v79-absorb}
\end{equation}
\end{theorem}

\begin{proof}
After eliminating \(\Lambda\), the logarithm of the product is at most
\[
 -c(\log L)^\chi
 +d_\partial\log L+O(\log\log L).
\]
The negative term dominates because \(\chi>1\).
\end{proof}

\begin{remark}[The V69 rate is stronger]
\label{rem:s7v79-beta}
If one uses the reflection-positive V69 rate
\(\exp(-c\beta)\) with \(\beta=g^{-2}\), any fixed power of
\(g^{-1}\), \(M\), and \(\Lambda\) is absorbed without the
\(\chi>1\) strengthening.  The present theorem uses only the weaker
\(\exp(-cM)\) activity needed to include chart exit.
\end{remark}

\begin{theorem}[Boundary row of BCRF on the extended space]
\label{thm:s7v79-BCRF}
At finite cutoff on the extended tensor-product link Hilbert space,
the saturated geometry, exact internal cluster semigroup, and distinct
boundary-edge forest prove BCRF rows (C1)--(C3).  The per-block marked
activity is bounded by
\begin{equation}
 \widetilde b_M
 :=
 A_{\partial,M}\pi_M^{+R}\longrightarrow0.           \label{eq:s7v79-btilde}
\end{equation}
\end{theorem}

\begin{proof}
Saturation and containment are
\Cref{thm:s7v79-saturation,lem:s7v79-containment}.
Internal evolution is the contraction
\Cref{lem:s7v78-bad-contraction}.  Boundary terms are each inserted once
by \Cref{cor:s7v79-forest}, their number per saturated block is bounded
by \Cref{lem:s7v79-boundary-count}, and their total fixed power is paid
by \Cref{thm:s7v79-absorb}.
\end{proof}

\section{Gauss-law gluing}
\label{sec:s7v79-Gauss}

Let \({\mathbb P}_{\cal G}\) be the normalized gauge projection.  Every
\(H(\boldsymbol\lambda)\) is gauge invariant and hence commutes with
\({\mathbb P}_{\cal G}\).  This does not make the physical Hilbert space
a tensor product across the cut.

\begin{lemma}[Commutation survives interpolation]
\label{lem:s7v79-gauge}
For all \(\boldsymbol\lambda\) and \(t\geq0\),
\begin{equation}
 {\mathbb P}_{\cal G}e^{-tH(\boldsymbol\lambda)}
 =
 e^{-tH(\boldsymbol\lambda)}{\mathbb P}_{\cal G}.     \label{eq:s7v79-gauge}
\end{equation}
The same identity holds after every mixed boundary derivative.
\end{lemma}

\begin{proof}
Each electric term and each Wilson trace in
\eqref{eq:s7v79-Hlambda} is gauge invariant.  Therefore its sum
commutes with the unitary gauge action and with the Haar average.
Differentiate the commuting identity in the interpolation parameters.
\end{proof}

\begin{firewall}[Commutation is not factorization]
\label{fire:s7v79-Gauss}
At \(\boldsymbol\lambda=0\), the extended link Hamiltonian factorizes,
but the Gauss projector contains vertex transformations acting on links
on both sides of the cut.  Thus a physical trace or matrix element does
not factor into an interior and exterior product merely because
\eqref{eq:s7v79-gauge} holds.  Decomposing the cut into boundary
representation/flux labels introduces a sum whose normalization and
cutoff growth must be controlled.
\end{firewall}

\begin{definition}[Boundary Gauss factorization card, BGFC]
\label{def:s7v79-BGFC}
BGFC consists of:
\begin{enumerate}[label=\textup{(G\arabic*)},leftmargin=4em]
\item an exact Peter--Weyl decomposition of the cut Gauss projector into
      boundary flux labels;
\item normalized conditional weights whose sum is one;
\item a polynomial bound in \(\Lambda\) per cut boundary block, not per
      internal Duhamel occurrence;
\item compatibility with the V71 transition cocycles and slow labels;
      and
\item preservation of the forest ownership in
      \Cref{cor:s7v79-forest}.
\end{enumerate}
\end{definition}

\begin{proposition}[Conditional physical BCRF]
\label{prop:s7v79-physical}
BGFC upgrades \Cref{thm:s7v79-BCRF} from the extended link Hilbert space
to the physical gauge-invariant sector.  Its fixed polynomial
\(\Lambda\)-loss is already included in
\eqref{eq:s7v79-A} and is absorbed by
\eqref{eq:s7v79-absorb}.
\end{proposition}

\begin{proof}
Use BGFC (G1)--(G2) to disintegrate the physical matrix element into
normalized boundary-label sectors.  Apply the extended-space forest
bound in each sector.  Rows (G3)--(G5) preserve bounded multiplicity and
contribute only the polynomial factor already allowed in
\eqref{eq:s7v79-A}.  Sum the normalized labels.
\end{proof}

\section{Status and next acquisition}
\label{sec:s7v79-status}

\begin{theorem}[What V79 proves]
\label{thm:s7v79-result}
V79 proves that finite-radius saturation preserves exponential product
domination, that no interaction touching an actual bad core crosses the
saturated component, and that every crossing support is fully good.  It
also proves an exact distinct-edge Duhamel derivative bound: a boundary
forest inserts each crossing Wilson interaction once.  The resulting
fixed inverse-coupling and cutoff power per block is absorbed by bad
rarity under \(\chi>1\).  Hence BCRF rows (C1)--(C3) close on the
extended link Hilbert space.
\end{theorem}

\begin{proof}
Use
\Cref{thm:s7v79-saturation,lem:s7v79-containment,
cor:s7v79-cross-good,thm:s7v79-mixed,
cor:s7v79-forest,thm:s7v79-absorb,
thm:s7v79-BCRF}.
\end{proof}

\begin{assessment}[Progress at seventh-series V79]
\label{ass:s7v79-status}
The bad-cluster boundary no longer carries arbitrary Duhamel
multiplicity.  Its remaining physical obstruction is not a Wilson
moment: it is factorization of the Gauss constraint across the cut with
only a polynomial boundary-label loss.  This is BGFC, a finite compact
representation problem which can be attacked directly.
\end{assessment}

\begin{decision}[V80 audit and acquisition order]
\label{dec:s7v79-V80}
V80 is the compulsory companion audit.  It will inspect the current SM
and NS notes for a boundary-constraint disintegration, normalized sector
sum, or finite-interface trace estimate.  It will then derive BGFC for a
finite cut using Peter--Weyl orthogonality, keeping the representation
dimension and cutoff powers explicit.  If the sector sum is exponential
in boundary area, it will be charged per saturated boundary block and
tested against \eqref{eq:s7v79-rate}.
\end{decision}

\begin{firewall}[Claim boundary after seventh-series V79]
\label{fire:s7v79-claim}
V79 does not prove BGFC, physical-sector BCRF, full BMINF or RL4C,
regulator-uniform SATC, the interacting joined slow-sector gap, the
remaining regular-sector UCTBB data, interacting endpoint continuity,
UCM, BCC, the pairing/Bogoliubov sector, nonlinear coarse-action
cocycles, terminal matching, regulator independence, cutoff removal,
Osterwalder--Schrader reconstruction, or the full Yang--Mills mass gap.
\end{firewall}

\paragraph{Parent-version record.}
V79 was formed from
\texttt{Renewal\_Yang\_Mills\_Mass\_Gap\_Research\_Notes\_V78.tex}.
The V78 parent had \(34\,641\) logical source lines,
\(1\,237\,424\) bytes, and SHA--256
\[
 \texttt{FE8920C9F245A1523485707384D79DB2C83C6D6B52CFBA4315575A9558EEF192}.
\]
The next compulsory companion audit is V80.

% =============================================================================
% END APPENDED SEVENTH-SERIES V79 MATERIAL
% =============================================================================

% =============================================================================
% BEGIN APPENDED SEVENTH-SERIES V80 MATERIAL
% =============================================================================

\chapter{Companion audit and exact boundary Gauss-sector
disintegration}
\label{ch:s7v80}

\section{Compulsory companion audit}
\label{sec:s7v80-audit}

This is the five-version audit announced in V75 and V79.  At the time of
inspection the current companion files were
\begin{align*}
 &\texttt{Renewal\_SM\_Einstein\_Continuum\_Research\_Notes\_V41.tex},\\
 &\texttt{Renewal\_NS\_Stability\_Research\_Notes\_V26.tex}.
\end{align*}
The SM file had \(14\,586\) logical source lines, \(529\,943\) bytes,
and SHA--256
\[
 \texttt{B0915676F931373B8A9F27B577F1DFD849A96F2ADE282384351C98961EEAA887}.
\]
The NS file had \(5\,319\) lines, \(278\,283\) bytes, and SHA--256
\[
 \texttt{82C887F8C2C69E4F28FAD7C3DC8DE5B9099CB5A4BF431EBA1FFF3E91935978AD}.
\]

The SM V40--V41 increment independently proves two distinctions directly
relevant here: all fixed polynomial moments need not sum over arbitrary
insertion multiplicity, and rarity of a set does not reduce the operator
norm of a multiplication form on vectors supported on that set.  It then
uses the same energy-normalized form-resolvent mechanism as V77, while
explicitly leaving semigroup and weighted locality open.  This
corroborates the V76--V78 decision to separate fixed RL4C moments,
resolvent resummation, and connected locality.  It supplies no
Yang--Mills Gauss-sector theorem.  The NS file contains no
type-compatible compact boundary representation or normalized sector
sum.

\begin{theorem}[Audit verdict]
\label{thm:s7v80-audit}
No physical estimate is imported from the companion programmes.  The SM
audit confirms that neither finite score moments nor rare support
controls arbitrary operator multiplicity.  The Yang--Mills route
therefore continues with the exact compact boundary representation
problem BGFC, rather than reverting to a moment-only argument.
\end{theorem}

\begin{proof}
The SM V40 lognormal example separates all finite moments from
exponential summability, and its V41 direct-integral test function
separates rarity from operator norm.  Both statements concern abstract
functional analysis, not Yang--Mills dynamics.  The file searches found
no compact Gauss-boundary factorization in either companion note.
\end{proof}

\section{Two-sided compact boundary representation}
\label{sec:s7v80-two}

Let \(K\) be a compact group; in the lattice application it is the
product of one copy of \(\mathrm{SU}(3)\) at each cut vertex.  Let
\({\cal H}_{L},{\cal H}_{R}\) carry unitary representations
\(U_L,U_R\) of \(K\).  At a finite Peter--Weyl cutoff, complete
reducibility gives
\begin{align}
 {\cal H}_L
 &=
 \bigoplus_{\lambda\in\widehat K_\Lambda}
 V_\lambda\otimes{\cal M}_{L,\lambda},               \label{eq:s7v80-HL}\\
 {\cal H}_R
 &=
 \bigoplus_{\mu\in\widehat K_\Lambda}
 V_\mu\otimes{\cal M}_{R,\mu},                       \label{eq:s7v80-HR}
\end{align}
where the group acts only on \(V_\lambda,V_\mu\).

Let
\begin{equation}
 {\mathbb P}_K
 :=
 \int_K U_L(k)\otimes U_R(k)\,dk                     \label{eq:s7v80-P}
\end{equation}
be normalized diagonal gauge averaging.

\begin{lemma}[Invariant carrier pairing]
\label{lem:s7v80-pair}
For irreducibles \(V_\lambda,V_\mu\),
\begin{equation}
 \dim(V_\lambda\otimes V_\mu)^K
 =
 \begin{cases}
 1,&\mu\simeq\lambda^*,\\
 0,&\mu\not\simeq\lambda^*.
 \end{cases}                                         \label{eq:s7v80-pair}
\end{equation}
When \(\mu=\lambda^*\), a normalized invariant vector is
\begin{equation}
 \Omega_\lambda
 =
 {1\over\sqrt{d_\lambda}}
 \sum_{j=1}^{d_\lambda}e_j\otimes e_j^*.             \label{eq:s7v80-Omega}
\end{equation}
\end{lemma}

\begin{proof}
There are canonical isomorphisms
\[
 (V_\lambda\otimes V_\mu)^K
 \simeq
 \operatorname{Hom}_K(V_\lambda^*,V_\mu).
\]
Schur's lemma gives \eqref{eq:s7v80-pair}.  Direct application of
\(U_\lambda(k)\otimes U_{\lambda^*}(k)\) shows that
\eqref{eq:s7v80-Omega} is invariant, and its squared norm is
\(d_\lambda^{-1}\sum_j1=1\).
\end{proof}

\begin{theorem}[Exact cut Gauss projector]
\label{thm:s7v80-projector}
Under \eqref{eq:s7v80-HL}--\eqref{eq:s7v80-HR},
\begin{equation}
 {\mathbb P}_K
 =
 \bigoplus_{\lambda\in\widehat K_\Lambda}
 |\Omega_\lambda\rangle\langle\Omega_\lambda|
 \otimes I_{{\cal M}_{L,\lambda}}
 \otimes I_{{\cal M}_{R,\lambda^*}},                 \label{eq:s7v80-projector}
\end{equation}
with zero action on unmatched isotypic pairs.  The sum is orthogonal.
\end{theorem}

\begin{proof}
Normalized Haar averaging is the orthogonal projection onto the
invariant subspace.  Decompose the tensor product using
\eqref{eq:s7v80-HL}--\eqref{eq:s7v80-HR} and apply
\Cref{lem:s7v80-pair} on every irreducible carrier pair.  Multiplicity
spaces are fixed by the group.
\end{proof}

\section{Normalized heat-trace disintegration}
\label{sec:s7v80-trace}

Suppose \(H_L,H_R\) commute with their respective \(K\)-actions.  By
Schur's lemma,
\begin{align}
 H_L&=\bigoplus_\lambda
 I_{V_\lambda}\otimes h_{L,\lambda},                 \label{eq:s7v80-hL}\\
 H_R&=\bigoplus_\mu
 I_{V_\mu}\otimes h_{R,\mu}.                         \label{eq:s7v80-hR}
\end{align}
Let \(A_L,A_R\) be bounded separately gauge-invariant observables, with
analogous multiplicity-space blocks \(a_{L,\lambda},a_{R,\mu}\).

\begin{theorem}[Exact normalized sector mixture]
\label{thm:s7v80-mixture}
For \(t>0\), define
\begin{align}
 Z_{L,\lambda}
 &:=
 \operatorname{Tr}_{{\cal M}_{L,\lambda}}
 e^{-t h_{L,\lambda}},                               \label{eq:s7v80-ZL}\\
 Z_{R,\lambda^*}
 &:=
 \operatorname{Tr}_{{\cal M}_{R,\lambda^*}}
 e^{-t h_{R,\lambda^*}},                             \label{eq:s7v80-ZR}\\
 Z_K&:=\sum_\lambda Z_{L,\lambda}Z_{R,\lambda^*}.
                                                               \label{eq:s7v80-Z}
\end{align}
Then
\begin{align}
 &{\operatorname{Tr}\!\left[
 {\mathbb P}_K(A_L\otimes A_R)
 e^{-t(H_L\otimes I+I\otimes H_R)}
 \right]\over Z_K}                                   \notag\\
 &\qquad=
 \sum_\lambda w_\lambda\,
 \omega_{L,\lambda}(A_L)
 \omega_{R,\lambda^*}(A_R),                          \label{eq:s7v80-mixture}
\end{align}
where
\begin{align}
 w_\lambda
 &:={Z_{L,\lambda}Z_{R,\lambda^*}\over Z_K}\geq0,
 &\sum_\lambda w_\lambda&=1,                         \label{eq:s7v80-weights}\\
 \omega_{L,\lambda}(A_L)
 &:=
 {\operatorname{Tr}(a_{L,\lambda}e^{-t h_{L,\lambda}})
  \over Z_{L,\lambda}},
 \quad
 \omega_{R,\lambda^*}(A_R)
 :=
 {\operatorname{Tr}(a_{R,\lambda^*}e^{-t h_{R,\lambda^*}})
  \over Z_{R,\lambda^*}}.                            \label{eq:s7v80-sector-states}
\end{align}
Terms with a zero sector partition function are omitted.
\end{theorem}

\begin{proof}
Insert the projector
\eqref{eq:s7v80-projector}.  On a matched carrier,
\[
 \operatorname{Tr}_{V_\lambda\otimes V_{\lambda^*}}
 |\Omega_\lambda\rangle\langle\Omega_\lambda|=1.
\]
Thus no factor \(d_\lambda\) remains.  The multiplicity-space traces
factor and give
\[
 \sum_\lambda
 \operatorname{Tr}(a_{L,\lambda}e^{-t h_{L,\lambda}})
 \operatorname{Tr}(a_{R,\lambda^*}e^{-t h_{R,\lambda^*}}).
\]
Taking \(A_L=A_R=I\) gives \eqref{eq:s7v80-Z}; divide by it and use the
definitions above.
\end{proof}

\begin{corollary}[No boundary-label cardinality loss]
\label{cor:s7v80-no-cardinality}
If
\[
 |\omega_{L,\lambda}(A_L)|\leq C_L,\qquad
 |\omega_{R,\lambda^*}(A_R)|\leq C_R
\]
uniformly in \(\lambda\), then the physical expectation in
\eqref{eq:s7v80-mixture} is bounded by \(C_LC_R\), independently of
\(\#\widehat K_\Lambda\), the representation dimensions, and the number
of retained labels.
\end{corollary}

\begin{proof}
Use positivity and normalization of the weights
\eqref{eq:s7v80-weights}.
\end{proof}

\begin{corollary}[Labelled \(L^p\) Jensen bound]
\label{cor:s7v80-Jensen}
For \(p\geq1\) and any sector family \(F_\lambda\),
\begin{equation}
 \left|
 \sum_\lambda w_\lambda F_\lambda
 \right|^p
 \leq
 \sum_\lambda w_\lambda|F_\lambda|^p.                \label{eq:s7v80-Jensen}
\end{equation}
\end{corollary}

\begin{proof}
This is Jensen's inequality for the probability weights
\(w_\lambda\).
\end{proof}

\section{Boundary insertions need not be resolved into labels}
\label{sec:s7v80-insertions}

The crossing Wilson operator \(Q_e\) commutes with the diagonal boundary
action but need not commute with the separate left and right actions.
It can therefore change the intermediate flux label.  Resolving every
such change into matrix elements is unnecessary.

\begin{lemma}[State bound before label resolution]
\label{lem:s7v80-state}
Let \(\omega\) be any positive normalized functional, let
\(S_j\) be contractions, and let \(Q_1,\ldots,Q_n\) be bounded.
Then
\begin{equation}
 |\omega(S_0Q_1S_1\cdots Q_nS_n)|
 \leq\prod_{j=1}^n\|Q_j\|.                           \label{eq:s7v80-state}
\end{equation}
\end{lemma}

\begin{proof}
A positive normalized functional has norm one.  Apply
\(|\omega(A)|\leq\|A\|\), submultiplicativity, and
\(\|S_j\|\leq1\).
\end{proof}

\begin{corollary}[Flux-changing forest edges have no sector-sum loss]
\label{cor:s7v80-flux}
Apply \Cref{lem:s7v80-state} to the mixed Duhamel word in
\Cref{thm:s7v79-mixed}, after the physical Gauss projection.  The same
bound \eqref{eq:s7v79-mixed} holds without inserting intermediate
Peter--Weyl resolutions.  Hence a flux-changing \(Q_e\) contributes its
operator norm once and no representation-label cardinality.
\end{corollary}

\begin{proof}
\Cref{lem:s7v79-gauge} permits the Gauss projection to be placed in the
normalized state.  Every interpolated semigroup is a contraction and
each distinct \(Q_e\) occurs once.  Use
\Cref{lem:s7v80-state} before resolving labels.
\end{proof}

\begin{firewall}[Derivative observables still see Casimir growth]
\label{fire:s7v80-derivative}
The no-cardinality theorem concerns separately invariant observables and
unresolved bounded forest edges.  An explicit electric derivative
acting on a boundary carrier has norm
\(O(\Lambda^{1/2})\) per derivative by V75.  Such a marked derivative
retains its polynomial Casimir loss.  It does not acquire an additional
factor equal to the number of sectors.
\end{firewall}

\section{Completion of the boundary Gauss card}
\label{sec:s7v80-BGFC}

\begin{theorem}[Finite-cutoff BGFC]
\label{thm:s7v80-BGFC}
For a finite cut of the compact lattice gauge theory, BGFC holds:
\begin{enumerate}[label=\textup{(\roman*)},leftmargin=3.5em]
\item \eqref{eq:s7v80-projector} is the exact Peter--Weyl cut
      decomposition;
\item \eqref{eq:s7v80-mixture} gives normalized positive sector weights;
\item invariant component observables have no sector-cardinality or
      representation-dimension loss;
\item unresolved flux-changing forest insertions obey
      \eqref{eq:s7v80-state}; and
\item any explicit boundary derivative has only the V75 polynomial
      Casimir loss.
\end{enumerate}
All statements are compatible with gauge-equivariant unitary transition
cocycles because their conjugation preserves the isotypic decomposition
and all operator norms.
\end{theorem}

\begin{proof}
Items (i)--(ii) are
\Cref{thm:s7v80-projector,thm:s7v80-mixture}.
Item (iii) is \Cref{cor:s7v80-no-cardinality}; item (iv) is
\Cref{cor:s7v80-flux}; item (v) is
\Cref{fire:s7v80-derivative} together with V75.  A gauge-equivariant
change of boundary frame conjugates the representation and multiplicity
blocks unitarily, so it changes none of the bounds.
\end{proof}

\begin{corollary}[Physical BCRF at finite regulator]
\label{cor:s7v80-BCRF}
The extended-space BCRF theorem
\Cref{thm:s7v79-BCRF} holds on the physical gauge-invariant Hilbert
space.  Its effective saturated bad activity remains
\begin{equation}
 \widetilde b_M
 =
 A_{\partial,M}\pi_M^{+R}\longrightarrow0            \label{eq:s7v80-bad}
\end{equation}
under the corrected schedule, with no additional factor from the number
of boundary flux sectors.
\end{corollary}

\begin{proof}
Apply \Cref{thm:s7v80-BGFC} to
\Cref{prop:s7v79-physical}.  The only possible extra marked derivative
loss is polynomial in \(\Lambda\) and was already allowed in
\(A_{\partial,M}\).
\end{proof}

\begin{corollary}[Ground-energy normalization is harmless in the sector
mixture]
\label{cor:s7v80-ground}
Sector-dependent component ground energies enter only through the
positive factors \(Z_{L,\lambda}Z_{R,\lambda^*}\) and hence through the
normalized weights \(w_\lambda\).  No uniform comparison of individual
sector partition functions is needed for
\Cref{cor:s7v80-no-cardinality}.
\end{corollary}

\begin{proof}
The derivation of \eqref{eq:s7v80-mixture} uses only positivity and
finiteness of the sector heat traces.  Multiplying a sector partition
function by its ground-energy exponential changes \(w_\lambda\), but
the weights remain nonnegative and sum to one.
\end{proof}

\begin{firewall}[Uniform sector estimates remain necessary]
\label{fire:s7v80-uniform}
Convex normalization removes the number of labels, not label dependence
of the conditional estimate.  If a sector state has a boundary
insertion norm or a residual gap diverging with \(\lambda\), the supremum
still diverges.  V75 supplies polynomial derivative bounds at finite
cutoff; regulator-uniform SATC and the interacting joined slow-sector
gap remain separate.
\end{firewall}

\section{Status and next acquisition}
\label{sec:s7v80-status}

\begin{theorem}[What V80 proves]
\label{thm:s7v80-result}
V80 completes the mandatory companion audit and proves BGFC at finite
cutoff.  The cut Gauss projector pairs dual boundary irreducibles through
a normalized invariant vector; normalized physical heat traces are
convex mixtures of sector states; and neither representation dimensions
nor the number of boundary labels appears in a uniform invariant
estimate.  Flux-changing boundary forest edges may be bounded before
label resolution.  Consequently the physical BCRF bad-cluster boundary
activity still tends to zero.
\end{theorem}

\begin{proof}
Use
\Cref{thm:s7v80-audit,thm:s7v80-projector,
thm:s7v80-mixture,cor:s7v80-no-cardinality,
cor:s7v80-flux,thm:s7v80-BGFC,
cor:s7v80-BCRF}.
\end{proof}

\begin{assessment}[Progress at seventh-series V80]
\label{ass:s7v80-status}
The Gauss constraint no longer obstructs finite-cut bad-cluster
factorization: its boundary sectors form a normalized convex mixture,
not an unweighted representation sum.  The mixed good/bad multiplicity
problem is therefore reduced to the already explicit local good activity
and bad rarity.  The leading unresolved physical issue is now the
interacting Hamiltonian on the joined slow sector and uniformity as its
cutoffs are removed.
\end{assessment}

\begin{decision}[V81 acquisition order]
\label{dec:s7v80-V81}
V81 will turn to the joined slow Hamiltonian.  It will first isolate its
exact degrees of freedom after Gauss reduction and transition-band
joining, then prove a finite-volume spectral comparison separating
compact torons, boundary flux sectors, and propagating transverse slow
modes.  It will test whether the quadratic slow floor survives the
nonlinear form by the V77 relative theorem.  Any true zero mode will be
kept as a vacuum/collective coordinate rather than assigned a false
mass.
\end{decision}

\begin{firewall}[Claim boundary after seventh-series V80]
\label{fire:s7v80-claim}
V80 does not prove regulator-uniform SATC, the interacting joined
slow-sector gap, full RL4C uniformly in all removed cutoffs, the
remaining regular-sector UCTBB data, interacting endpoint continuity,
UCM, BCC, the pairing/Bogoliubov sector, nonlinear coarse-action
cocycles, terminal matching, regulator independence, cutoff removal,
Osterwalder--Schrader reconstruction, or the full Yang--Mills mass gap.
\end{firewall}

\paragraph{Parent-version record.}
V80 was formed from
\texttt{Renewal\_Yang\_Mills\_Mass\_Gap\_Research\_Notes\_V79.tex}.
The V79 parent had \(35\,120\) logical source lines,
\(1\,254\,402\) bytes, and SHA--256
\[
 \texttt{D212C789BD47EA0F92F63BF2AAA37C5A6B0198CACEDE9226CA026329CF0B72CC}.
\]
The next compulsory companion audit is V85.

% =============================================================================
% END APPENDED SEVENTH-SERIES V80 MATERIAL
% =============================================================================

% =============================================================================
% BEGIN APPENDED SEVENTH-SERIES V81 MATERIAL
% =============================================================================

\chapter{The joined slow sector and the infrared mass-generation
dichotomy}
\label{ch:s7v81}

\section{Purpose}
\label{sec:s7v81-purpose}

V71--V80 protect the residual complement, organize transition labels,
and close the finite-cut bad-cluster boundary.  None of those results
creates a mass in the retained infrared sector.  V29 already proves that
exact Gaussian blocking preserves the dispersion
\[
 \omega_a(k)
 =
 {2\over a}
 \left(\sum_i\sin^2{ak_i\over2}\right)^{1/2}
\]
and hence preserves modes with energy tending to zero as the physical
volume grows.

This chapter makes the remaining logical alternative precise.  A
perturbation which is only small relative to the massless slow form and
has vanishing expectation on infrared physical trial states cannot
create a uniform gap.  Therefore a proof of the Yang--Mills gap must
establish at least one genuinely nonperturbative effect: infrared trial
states are removed by the exact Gauss projection/dressing, or their
energy receives a positive physical-scale contribution, or the
interacting Hilbert norm becomes singular relative to the Gaussian
carrier.  Ultraviolet residual estimates alone prove none of these.

\section{An abstract variational obstruction}
\label{sec:s7v81-variational}

Let \({\cal H}_L\) be a finite-volume Hilbert space with normalized
vacuum \(\Omega_L\), and let \(H_L\geq0\) be self-adjoint with
\(H_L\Omega_L=0\).  Define its nonvacuum gap by
\begin{equation}
 \operatorname{gap}(H_L)
 :=
 \inf_{\substack{\psi\in D(H_L^{1/2})\\
                  \psi\perp\Omega_L,\ \|\psi\|=1}}
 \langle\psi,H_L\psi\rangle.                          \label{eq:s7v81-gap}
\end{equation}

\begin{lemma}[One infrared trial state bounds the gap]
\label{lem:s7v81-Rayleigh}
For every normalized
\(\psi_L\perp\Omega_L\) in the form domain,
\begin{equation}
 \operatorname{gap}(H_L)
 \leq\langle\psi_L,H_L\psi_L\rangle.                 \label{eq:s7v81-Rayleigh}
\end{equation}
\end{lemma}

\begin{proof}
The right side is one member of the infimum
\eqref{eq:s7v81-gap}.
\end{proof}

Let \(H_L=H_{0,L}+V_L\) in form sense, with common vacuum normalization.

\begin{theorem}[Relative control cannot raise a soft physical trial
family]
\label{thm:s7v81-relative-no-gap}
Suppose there are normalized physical vectors
\(\psi_L\perp\Omega_L\) and numbers
\(\varepsilon_L,\delta_L\geq0\) such that
\begin{align}
 \langle\psi_L,H_{0,L}\psi_L\rangle
 &\leq\varepsilon_L,                                 \label{eq:s7v81-soft}\\
 |\langle\psi_L,V_L\psi_L\rangle|
 &\leq
 \eta_L\langle\psi_L,H_{0,L}\psi_L\rangle+\delta_L,
 \qquad \eta_L\geq0.                                 \label{eq:s7v81-relative}
\end{align}
Then
\begin{equation}
 \operatorname{gap}(H_L)
 \leq(1+\eta_L)\varepsilon_L+\delta_L.               \label{eq:s7v81-gap-upper}
\end{equation}
Consequently, if
\[
 (1+\eta_L)\varepsilon_L+\delta_L\longrightarrow0,
\]
the family has no positive volume-uniform gap.
\end{theorem}

\begin{proof}
Add \eqref{eq:s7v81-soft} and the upper side of
\eqref{eq:s7v81-relative}, then apply
\Cref{lem:s7v81-Rayleigh}.
\end{proof}

\begin{corollary}[Pure relative smallness preserves gaplessness]
\label{cor:s7v81-pure-relative}
If \(\delta_L=0\), \(\sup_L\eta_L<\infty\), and the physical reference
has a normalized trial family with \(\varepsilon_L\to0\), then
\(\inf_L\operatorname{gap}(H_L)=0\).
\end{corollary}

\begin{proof}
Use \eqref{eq:s7v81-gap-upper}.
\end{proof}

\begin{firewall}[An additive infrared self-energy is not excluded]
\label{fire:s7v81-additive}
The theorem does not say that a small ultraviolet coupling cannot
generate a mass.  Such a mass is precisely a contribution for which
\(\delta_L\) stays positive on the normalized infrared physical trial
state as \(L\to\infty\).  The theorem says that a proof recording only a
multiple of the massless reference energy has not established that
contribution.
\end{firewall}

\section{The explicit Gaussian singlet test}
\label{sec:s7v81-singlet}

On the V29 transverse Fock space, fix a nonzero momentum \(k\) with
\(k\ne-k\) modulo the periodic lattice and a transverse polarization.
Let \(a_{k,c}^*\) create a quantum of adjoint colour
\(c=1,\ldots,d_G\), where \(d_G=8\) for \(\mathrm{SU}(3)\).  Put
\begin{equation}
 \Phi_k
 :=
 {1\over\sqrt{d_G}}
 \sum_{c=1}^{d_G}
 a_{k,c}^*a_{-k,c}^*\Omega.                           \label{eq:s7v81-Phi}
\end{equation}

\begin{lemma}[Normalized global-colour singlet pair]
\label{lem:s7v81-singlet}
\(\|\Phi_k\|=1\), \(\Phi_k\perp\Omega\), and
\(\Phi_k\) is invariant under the residual global colour action.
Moreover,
\begin{equation}
 \langle\Phi_k,H_a^{(0)}\Phi_k\rangle
 =2\omega_a(k).                                      \label{eq:s7v81-singlet-energy}
\end{equation}
\end{lemma}

\begin{proof}
The modes \(k\) and \(-k\) are distinct, so the \(d_G\) two-particle
vectors in \eqref{eq:s7v81-Phi} are orthonormal.  This proves
normalization and vacuum orthogonality.  The tensor
\(\sum_c e_c\otimes e_c\) is the invariant metric on the real adjoint
representation, so the colour contraction is globally invariant.  The
quadratic Hamiltonian counts one quantum of energy \(\omega_a(k)\) in
each of the two modes.
\end{proof}

\begin{corollary}[Massless singlet energy]
\label{cor:s7v81-singlet-soft}
For the least nonzero lattice momentum
\(k_L=(2\pi/L,0,0)\),
\begin{equation}
 \langle\Phi_{k_L},H_a^{(0)}\Phi_{k_L}\rangle
 =
 {4\over a}\sin{\pi a\over L}
 \longrightarrow0
 \qquad(L\to\infty).                                 \label{eq:s7v81-singlet-soft}
\end{equation}
\end{corollary}

\begin{proof}
Insert \(k_L\) in \eqref{eq:s7v29-omega} and use
\eqref{eq:s7v81-singlet-energy}.
\end{proof}

\begin{firewall}[Global colour invariance is not the nonlinear Gauss
law]
\label{fire:s7v81-local-Gauss}
\(\Phi_k\) is a legitimate trial state in the linearly gauge-reduced
Gaussian carrier and is invariant under residual global
\(\mathrm{SU}(3)\).  It has not been proved to lie in the exact
nonlinear local-Gauss physical Hilbert space.  Projecting or dressing it
may suppress its norm or change its energy.  That operation is one of
the nonperturbative alternatives isolated below.
\end{firewall}

\section{Projection and dressing dichotomy}
\label{sec:s7v81-projection}

Let \({\mathbb P}_{G,L}\) denote the exact finite-volume Gauss projector
on a common kinematic carrier, and define
\begin{equation}
 \rho_L:=\|{\mathbb P}_{G,L}\Phi_{k_L}\|.             \label{eq:s7v81-rho}
\end{equation}
When \(\rho_L>0\), subtract the vacuum component and normalize:
\begin{equation}
 \Psi_L
 :=
 { (I-|\Omega_L\rangle\langle\Omega_L|)
    {\mathbb P}_{G,L}\Phi_{k_L}
  \over
  \|
  (I-|\Omega_L\rangle\langle\Omega_L|)
    {\mathbb P}_{G,L}\Phi_{k_L}
  \|}.                                                \label{eq:s7v81-Psi}
\end{equation}

\begin{theorem}[Infrared physical-carrier alternative]
\label{thm:s7v81-dichotomy}
Assume \(\Psi_L\) is defined along a cofinal volume sequence.  If
\begin{align}
 \langle\Psi_L,H_{0,L}\Psi_L\rangle&\longrightarrow0,
                                                               \label{eq:s7v81-projected-soft}\\
 \langle\Psi_L,V_L\Psi_L\rangle&\longrightarrow0,     \label{eq:s7v81-projected-V}
\end{align}
then the interacting physical gap cannot have a positive
volume-uniform lower bound.  Therefore a positive mass gap requires at
least one of the following:
\begin{enumerate}[label=\textup{(D\arabic*)},leftmargin=4em]
\item the projected/dressed infrared family disappears or becomes
      vacuum, so \eqref{eq:s7v81-Psi} has no stable nonzero limit;
\item its interacting energy has a positive lower limit;
\item the correct interacting physical vectors are not controlled in
      the Gaussian Hilbert norm used in
      \eqref{eq:s7v81-projected-soft}; or
\item a different physical infrared family must be tested and receives
      the same nonperturbative treatment.
\end{enumerate}
\end{theorem}

\begin{proof}
Under
\eqref{eq:s7v81-projected-soft}--\eqref{eq:s7v81-projected-V},
\(\langle\Psi_L,H_L\Psi_L\rangle\to0\).  Apply
\Cref{lem:s7v81-Rayleigh}.  The listed alternatives are the logical
negations of using this stable soft family to contradict a positive
gap; item (D4) records that one trial family alone is not a dense-core
statement.
\end{proof}

\begin{remark}[Why bad-set rarity does not decide the alternative]
\label{rem:s7v81-rarity}
V70 and V80 show that bad clusters are dilute and their normalized
boundary sums are controlled.  A normalized infrared state is extended
over the whole volume.  Small probability of a local bad event neither
forces \eqref{eq:s7v81-projected-V} nor proves a positive energy in
item (D2).  The SM V41 essential-supremum warning is the analogous
operator-norm distinction.
\end{remark}

\section{Self-energy formulation}
\label{sec:s7v81-selfenergy}

Suppose, only for diagnostic purposes, that the joined physical slow
Hamiltonian has a translation-invariant quadratic one-particle block
\begin{equation}
 E_{\rm eff}(k)
 =
 \omega_a(k)I+\Sigma_{\rm phys}(k).                  \label{eq:s7v81-Eeff}
\end{equation}

\begin{lemma}[Infrared-relative symbols remain massless]
\label{lem:s7v81-symbol}
If
\begin{equation}
 \|\Sigma_{\rm phys}(k)\|
 \leq\eta\,\omega_a(k),
 \qquad |k|\leq k_0,                                 \label{eq:s7v81-symbol-relative}
\end{equation}
with fixed \(\eta<\infty\), then
\begin{equation}
 \inf_{0<|k|\leq k_0}
 \lambda_{\min}E_{\rm eff}(k)=0.                     \label{eq:s7v81-symbol-gapless}
\end{equation}
\end{lemma}

\begin{proof}
The upper variational bound on any normalized colour-polarization vector
is at most
\((1+\eta)\omega_a(k)\), which tends to zero as \(k\to0\).
\end{proof}

\begin{corollary}[A quadratic mass needs a nonrelative infrared row]
\label{cor:s7v81-mass-row}
Within the diagnostic quadratic description, a mass \(m_*>0\) requires
a positive infrared contribution not bounded by a finite multiple of
\(\omega_a(k)\); for example,
\[
 \liminf_{k\to0}
 \lambda_{\min}\Sigma_{\rm phys}(k)\geq m_*.
\]
In the full gauge theory the alternative is a nonquadratic confinement
mechanism which removes the coloured quasiparticle description while
still giving the dense gauge-invariant decay criterion of V25.
\end{corollary}

\begin{proof}
The first statement is the contrapositive of
\Cref{lem:s7v81-symbol}.  The second records that the mass-gap criterion
\Cref{thm:s7v25-OS-gap} concerns gauge-invariant local states and does
not require a coloured one-particle pole.
\end{proof}

\begin{firewall}[A gluon self-energy is not the mass-gap theorem]
\label{fire:s7v81-gluon}
Even a positive \(\Sigma_{\rm phys}(0)\) in a gauge-fixed propagator
would not by itself prove the Yang--Mills mass gap.  The final theorem
must give one common exponential time-decay rate on a dense
gauge-invariant OS core, as required by V25.
\end{firewall}

\section{The terminal nonperturbative card}
\label{sec:s7v81-terminal}

The V28 NTRGC already gives the exact operator recursion which transports
a terminal transfer gap through controlled coarse matching.  V81 adds
the missing infrared content of its terminal row.

\begin{definition}[Infrared singlet mass card, ISMC]
\label{def:s7v81-ISMC}
At a terminal physical scale \(a_*\), ISMC consists of either:
\begin{enumerate}[label=\textup{(I\arabic*)},leftmargin=4em]
\item a reflection-positive terminal Euclidean state whose connected
      correlations obey
      \[
      |\omega(A\,\tau_tB)-\omega(A)\omega(B)|
      \leq C_{A,B}e^{-m_*t}
      \]
      for a dense gauge-invariant local OS core and one \(m_*>0\); or
\item a terminal physical transfer operator with a unique vacuum and
      nonvacuum norm at most \(e^{-m_*a_*}\).
\end{enumerate}
The estimate must be uniform in terminal spatial volume and every
retained boundary/transition sector.  It must also resolve the
alternative in \Cref{thm:s7v81-dichotomy}.
\end{definition}

\begin{theorem}[ISMC is the genuine mass-generating row]
\label{thm:s7v81-ISMC}
NTRGC ultraviolet matching and residual contraction, together with
ISMC and summable operator errors below the terminal gap margin, imply
the regulated gauge-invariant exponential decay required by V25.
Conversely, the V29 Gaussian terminal theory fails ISMC in the
thermodynamic limit.
\end{theorem}

\begin{proof}
The forward implication is the V28 transfer recursion followed by
\Cref{thm:s7v25-OS-gap}; ISMC states its previously abstract terminal
gap in the correct physical sector.  The Gaussian failure is
\Cref{thm:s7v29-gapless} and the explicit singlet diagnostic
\Cref{cor:s7v81-singlet-soft}, subject to the local-Gauss firewall.
\end{proof}

\section{Status and next acquisition}
\label{sec:s7v81-status}

\begin{theorem}[What V81 proves]
\label{thm:s7v81-result}
V81 proves a variational no-gap theorem for any stable physical
infrared trial family whose reference and interaction energies both
vanish.  It constructs the explicit globally colour-singlet Gaussian
two-mode diagnostic with energy tending to zero, isolates the nonlinear
local-Gauss dressing qualification, and proves that an
infrared-relative quadratic self-energy remains massless.  ISMC states
the nonperturbative terminal datum in the gauge-invariant sector.
\end{theorem}

\begin{proof}
Use
\Cref{thm:s7v81-relative-no-gap,lem:s7v81-singlet,
cor:s7v81-singlet-soft,thm:s7v81-dichotomy,
lem:s7v81-symbol,thm:s7v81-ISMC}.
\end{proof}

\begin{assessment}[Progress at seventh-series V81]
\label{ass:s7v81-status}
The programme has reached the correct conceptual bottleneck.  Residual
coercivity, rare bad clusters, and polynomial cutoff control cannot be
relabelled as mass generation.  A terminal proof must now produce
uniform gauge-invariant exponential decay at a physical coarse scale.
The most concrete available route is the compact strong-coupling regime,
where conditional influences are genuinely contractive.
\end{assessment}

\begin{decision}[V82 acquisition order]
\label{dec:s7v81-V82}
V82 will prove a terminal strong-coupling ISMC for the compact Wilson
measure at sufficiently small plaquette coupling.  It will compute a
finite-range Dobrushin influence matrix for link conditionals, derive
exponential spacetime covariance decay for bounded local
gauge-invariant observables, and insert reflection positivity into the
V25 OS criterion.  The theorem will be a terminal card; matching the
asymptotically free ultraviolet flow to its small-coupling hypothesis
will remain a separate RG problem.
\end{decision}

\begin{firewall}[Claim boundary after seventh-series V81]
\label{fire:s7v81-claim}
V81 does not prove ISMC, the UV-to-terminal RG matching,
regulator-uniform SATC, the interacting joined slow-sector gap, full
RL4C through cutoff removal, interacting endpoint continuity, UCM, BCC,
the pairing/Bogoliubov sector, nonlinear coarse-action cocycles,
regulator independence, Osterwalder--Schrader reconstruction, or the
full Yang--Mills mass gap.
\end{firewall}

\paragraph{Parent-version record.}
V81 was formed from
\texttt{Renewal\_Yang\_Mills\_Mass\_Gap\_Research\_Notes\_V80.tex}.
The V80 parent had \(35\,571\) logical source lines,
\(1\,270\,763\) bytes, and SHA--256
\[
 \texttt{C14E143490EEC3EE796B5D4F76B194D9547DA1B4FB7E456C8B67686056BAFAA8}.
\]
The next compulsory companion audit is V85.

% =============================================================================
% END APPENDED SEVENTH-SERIES V81 MATERIAL
% =============================================================================

% =============================================================================
% BEGIN APPENDED SEVENTH-SERIES V82 MATERIAL
% =============================================================================

\chapter{A terminal strong-coupling mass gap from link-conditional
contraction}
\label{ch:s7v82}

\section{Purpose and scope}
\label{sec:s7v82-purpose}

V81 isolates ISMC, a genuinely nonperturbative terminal gap.  This
chapter proves ISMC for the four-dimensional compact Wilson measure when
the terminal plaquette coupling is sufficiently small.  The proof is
elementary and quantitative: changing one exterior link changes only
finitely many plaquettes in another link's conditional law; a
Dobrushin row sum below one yields exponential covariance decay; and
the reflection positivity proved in V69 converts that decay into an
Osterwalder--Schrader Hamiltonian gap.

This is a terminal theorem, not the full continuum theorem.  The
asymptotically free ultraviolet theory begins at large
\(\beta=g^{-2}\), whereas the theorem below assumes small terminal
\(\beta\).  One must still construct a gauge-covariant RG trajectory
whose effective terminal action lies in the contraction domain while
preserving the operator matching errors of NTRGC.

\section{Wilson conditionals and their oscillation}
\label{sec:s7v82-conditionals}

Let \(\Lambda\) be a finite four-dimensional periodic lattice with link
set \(E\) and plaquette set \(P\).  Put \(G=\mathrm{SU}(3)\) and
\begin{equation}
 \phi(W):=3-\operatorname{ReTr}W,\qquad
 0\leq\phi(W)\leq B_\phi,\qquad B_\phi:=6.            \label{eq:s7v82-phi}
\end{equation}
The safe value \(6\) uses only
\(-3\leq\operatorname{ReTr}W\leq3\).
Define
\begin{equation}
 d\mu_{\Lambda,\beta}(U)
 :=
 {1\over Z_{\Lambda,\beta}}
 \exp\left[-\beta\sum_{p\in P}\phi(W_p(U))\right]
 \prod_{e\in E}dU_e.                                 \label{eq:s7v82-measure}
\end{equation}

For links \(e\ne f\), let
\begin{equation}
 n_{ef}:=\#\{p:p\ni e,\ p\ni f\},                    \label{eq:s7v82-nef}
\end{equation}
and set
\begin{equation}
 N_{\rm nb}:=
 \sup_e\#\{f\ne e:n_{ef}>0\},\qquad
 n_{\rm sh}:=\sup_{e\ne f}n_{ef}.                    \label{eq:s7v82-incidence}
\end{equation}
Both are finite constants independent of volume; the crude
four-dimensional bound \(N_{\rm nb}\leq18\) follows by assigning the
three other links of each of the six plaquettes incident to \(e\).

\begin{lemma}[Normalized-tilt total-variation bound]
\label{lem:s7v82-tilt}
Let \(\nu_0,\nu_1\) be probability measures with positive densities and
suppose
\begin{equation}
 \operatorname{osc}\log{d\nu_1\over d\nu_0}
 \leq D.                                              \label{eq:s7v82-logosc}
\end{equation}
Then, with total variation normalized as
\(\|\nu_1-\nu_0\|_{\rm TV}
=\frac12\int|d\nu_1-d\nu_0|\),
\begin{equation}
 \|\nu_1-\nu_0\|_{\rm TV}
 \leq
 \min\left\{1,{e^D-1\over2}\right\}.                 \label{eq:s7v82-TV}
\end{equation}
\end{lemma}

\begin{proof}
Let \(r=d\nu_1/d\nu_0\).  The oscillation bound and
\(\int r\,d\nu_0=1\) imply
\(e^{-D}\leq r\leq e^D\): the essential infimum is at most one and the
essential supremum at least one, while their ratio is at most \(e^D\).
Thus \(|r-1|\leq e^D-1\), and
\[
 \|\nu_1-\nu_0\|_{\rm TV}
 ={1\over2}\int|r-1|\,d\nu_0
 \leq{e^D-1\over2}.
\]
Total variation is always at most one.
\end{proof}

\begin{lemma}[One-link conditional influence]
\label{lem:s7v82-influence}
Let \(\mu_e(\cdot\mid U_{E\setminus\{e\}})\) be the one-link
conditional law from \eqref{eq:s7v82-measure}.  If two exterior
configurations differ only at \(f\ne e\), then
\begin{equation}
 \left\|
 \mu_e(\cdot\mid U)-\mu_e(\cdot\mid U')
 \right\|_{\rm TV}
 \leq
 c_{ef}:=
 \min\left\{
 1,{e^{2\beta B_\phi n_{ef}}-1\over2}
 \right\}.                                           \label{eq:s7v82-cef}
\end{equation}
If \(n_{ef}=0\), then \(c_{ef}=0\).
\end{lemma}

\begin{proof}
Only plaquettes containing both \(e\) and \(f\) change.  As a function
of the conditioned variable \(U_e\), the difference between the two
conditional log densities is
\[
 h(U_e)
 =
 -\beta\sum_{\substack{p\ni e\\p\ni f}}
 [\phi(W_p(U_e,U))-\phi(W_p(U_e,U'))].
\]
Each bracket takes values in an interval of length at most
\(2B_\phi\), so
\(\operatorname{osc}h\leq2\beta B_\phi n_{ef}\).
Apply \Cref{lem:s7v82-tilt}.  If no plaquette is shared, the two
conditionals coincide.
\end{proof}

\begin{corollary}[Explicit Dobrushin row bound]
\label{cor:s7v82-alpha}
The influence matrix \(C=(c_{ef})\) obeys
\begin{equation}
 \alpha_\beta:=
 \sup_e\sum_{f\ne e}c_{ef}
 \leq
 {N_{\rm nb}\over2}
 \left(e^{2\beta B_\phi n_{\rm sh}}-1\right).         \label{eq:s7v82-alpha}
\end{equation}
Hence \(\alpha_\beta<1\) whenever
\begin{equation}
 \beta<
 \beta_{\rm D}
 :=
 {1\over2B_\phi n_{\rm sh}}
 \log\left(1+{2\over N_{\rm nb}}\right).             \label{eq:s7v82-betaD}
\end{equation}
\end{corollary}

\begin{proof}
There are at most \(N_{\rm nb}\) nonzero entries in each row, and each
is bounded by the untruncated right side with
\(n_{ef}\leq n_{\rm sh}\).  Solve the resulting strict inequality for
\(\beta\).
\end{proof}

\section{Dobrushin comparison with finite propagation}
\label{sec:s7v82-Dobrushin}

For a bounded function \(F\) on \(G^E\), define its one-link
oscillation
\begin{equation}
 \delta_eF
 :=
 \sup\{|F(U)-F(U')|:
 U_{E\setminus\{e\}}=U'_{E\setminus\{e\}}\}.          \label{eq:s7v82-delta}
\end{equation}
Let \(d_E\) be the interaction-graph distance on links:
\(d_E(e,f)=1\) when \(n_{ef}>0\).

\begin{theorem}[Dobrushin covariance comparison]
\label{thm:s7v82-Dobrushin}
Suppose \(\alpha_\beta<1\) and define
\begin{equation}
 D:=(I-C)^{-1}=\sum_{n=0}^{\infty}C^n.               \label{eq:s7v82-D}
\end{equation}
Then the finite-volume Gibbs measure is unique for each boundary
condition and, for bounded cylinder functions \(F,G\),
\begin{equation}
 |\operatorname{Cov}_{\mu_{\Lambda,\beta}}(F,G)|
 \leq
 \sum_{e,f\in E}\delta_eF\,D_{ef}\,\delta_fG.         \label{eq:s7v82-cov}
\end{equation}
The constant one is not sharp but is uniform in volume.
\end{theorem}

\begin{proof}
Couple two single-link conditional updates maximally.  If their exterior
configurations disagree on a set with indicator vector \(b\), the
probability that the updated link \(e\) disagrees is at most
\(\sum_fc_{ef}b_f\): change the disagreeing exterior links one at a time
and use the definition of \(c_{ef}\) and the triangle inequality for
total variation.  Iterating complete conditional sweeps gives the
disagreement majorant
\[
 b+Cb+C^2b+\cdots=Db.
\]
Because \(\|C\|_{\ell^\infty\to\ell^\infty}\leq\alpha_\beta<1\),
the coupled boundary influence tends to zero, proving uniqueness.

To compare conditional expectations of \(F\) under two exteriors,
change \(F\) along the coupled disagreements:
\[
 |\mathbb E F(U)-\mathbb E F(U')|
 \leq\sum_e\delta_eF\,\mathbb P\{U_e\ne U'_e\}.
\]
Use the \(D\)-majorant and then expose the variables on which \(G\)
depends one at a time in the martingale covariance identity.  Each
exposure costs at most \(\delta_fG\).  Summing the resulting bound gives
\eqref{eq:s7v82-cov}.  This proof deliberately retains the safe constant
one; the usual symmetric normalization improves it but is unnecessary.
\end{proof}

\begin{lemma}[Exponential tail of the influence resolvent]
\label{lem:s7v82-D-tail}
If \(d_E(e,f)>n\), then
\((C^j)_{ef}=0\) for \(j\leq n\).  Consequently
\begin{equation}
 \sup_e\sum_{\{f:d_E(e,f)\geq r\}}D_{ef}
 \leq
 {\alpha_\beta^{\,\lceil r\rceil}\over1-\alpha_\beta}.
                                                               \label{eq:s7v82-D-tail}
\end{equation}
\end{lemma}

\begin{proof}
One multiplication by \(C\) moves at most one edge in the interaction
graph, so a product of \(j\) entries cannot connect graph distance
greater than \(j\).  The row sum of \(C^j\) is at most
\(\alpha_\beta^j\).  Sum the geometric tail from
\(\lceil r\rceil\).
\end{proof}

\begin{theorem}[Uniform exponential clustering]
\label{thm:s7v82-clustering}
Let \(F,G\) be bounded cylinder functions with link supports
\(S_F,S_G\), and let
\[
 r:=d_E(S_F,S_G).
\]
If \(\alpha_\beta<1\), then
\begin{equation}
 |\operatorname{Cov}_{\mu_{\Lambda,\beta}}(F,G)|
 \leq
 {4|S_F|\|F\|_\infty\|G\|_\infty\over1-\alpha_\beta}
 \alpha_\beta^{\,r}.                                 \label{eq:s7v82-clustering}
\end{equation}
The bound is uniform in periodic volume and in arbitrary fixed boundary
conditions.
\end{theorem}

\begin{proof}
\(\delta_eF=0\) off \(S_F\) and is at most
\(2\|F\|_\infty\) on it; similarly for \(G\).
For each \(e\in S_F\), sum \(D_{ef}\) over
\(f\in S_G\), all of which have distance at least \(r\), and apply
\Cref{lem:s7v82-D-tail}.  Insert the two oscillation factors in
\eqref{eq:s7v82-cov}.
\end{proof}

\begin{corollary}[Infinite-volume strong-coupling state]
\label{cor:s7v82-infinite}
For \(\beta<\beta_{\rm D}\), the periodic finite-volume measures have a
unique infinite-volume Gibbs limit on cylinder functions, and the limit
obeys \eqref{eq:s7v82-clustering}.
\end{corollary}

\begin{proof}
Compactness gives subsequential limits.  The boundary-influence
coupling in the proof of \Cref{thm:s7v82-Dobrushin}, together with
\eqref{eq:s7v82-D-tail}, shows that two subsequential limits agree on
every cylinder.  Pass the volume-uniform covariance bound to the limit.
\end{proof}

\section{Physical Euclidean time rate}
\label{sec:s7v82-time}

Let \(a_*>0\) be the terminal lattice spacing.  If two local supports are
separated by \(n\) terminal time slices, their link-interaction distance
is at least \(n-c_{F,G}\), where the fixed constant only accounts for
the temporal thicknesses of the observables.  Define
\begin{equation}
 m_{\rm D}
 :=
 {-\log\alpha_\beta\over a_*}>0.                     \label{eq:s7v82-mD}
\end{equation}

\begin{corollary}[Uniform terminal time decay]
\label{cor:s7v82-time}
For every fixed bounded local \(F,G\), there is
\(C_{F,G}<\infty\) such that
\begin{equation}
 |\operatorname{Cov}(F,\tau_tG)|
 \leq C_{F,G}e^{-m_{\rm D}t}                         \label{eq:s7v82-time}
\end{equation}
whenever \(t=na_*\) is larger than their fixed temporal thickness.
\end{corollary}

\begin{proof}
Use \eqref{eq:s7v82-clustering} with
\(r\geq n-c_{F,G}\), and absorb
\(\alpha_\beta^{-c_{F,G}}\) into the prefactor.
\end{proof}

\begin{remark}[Nonoptimal rate]
\label{rem:s7v82-rate}
The safe influence and graph-distance estimates are deliberately crude.
They prove a strictly positive mass \(m_{\rm D}\) on a nonempty interval
\((0,\beta_{\rm D})\); optimization is irrelevant to the existence
claim.
\end{remark}

\section{Reflection positivity and the terminal OS gap}
\label{sec:s7v82-OS}

V69 proves reflection positivity of the finite periodic Wilson measure
from its nonnegative character expansion.  The property is an inequality
between finitely supported expectations and therefore passes to the
unique infinite-volume limit of
\Cref{cor:s7v82-infinite}.

\begin{lemma}[Density of bounded gauge-invariant cylinders]
\label{lem:s7v82-density}
The linear span of bounded continuous gauge-invariant cylinder
functions is dense in the Osterwalder--Schrader Hilbert space generated
by the Wilson state.  Their centered classes are dense in the orthogonal
complement of the vacuum class.
\end{lemma}

\begin{proof}
The OS Hilbert space is, by construction, the completion of
positive-time cylinder functions modulo the null space.  Haar averaging
is a contractive projection onto gauge invariants.  Peter--Weyl
polynomials are uniformly dense in continuous functions on each finite
product of compact link groups, and averaging preserves convergence.
Subtracting the vacuum expectation is the orthogonal projection away
from the vacuum class.
\end{proof}

\begin{theorem}[Terminal Wilson mass-gap theorem]
\label{thm:s7v82-gap}
Let \(0<\beta<\beta_{\rm D}\).  On a reflection-compatible cofinal
family of periodic regulators, the infinite-volume Wilson state is
reflection positive, has a unique OS vacuum, and its reconstructed
Hamiltonian satisfies
\begin{equation}
 \operatorname{spec}(H_{\rm OS}|_{\Omega^\perp})
 \subset[m_{\rm D},\infty),                           \label{eq:s7v82-gap}
\end{equation}
with \(m_{\rm D}\) from \eqref{eq:s7v82-mD}.
\end{theorem}

\begin{proof}
Reflection positivity is V69 and passes to the unique limit as stated
above.  For every bounded gauge-invariant local \(A\), apply
\Cref{cor:s7v82-time} to
\((\Theta A^\circ,\tau_tA^\circ)\).  Reflection positivity makes this
correlator nonnegative, while \eqref{eq:s7v82-time} bounds it by
\(C_Ae^{-m_{\rm D}t}\).  The centered local classes are dense by
\Cref{lem:s7v82-density}.  Apply the dense-core spectral theorem
\Cref{thm:s7v25-dense-core}; it simultaneously excludes any second
zero-energy vector and proves \eqref{eq:s7v82-gap}.
\end{proof}

\begin{corollary}[Strong-coupling ISMC]
\label{cor:s7v82-ISMC}
The compact Wilson terminal theory satisfies ISMC for every
\(\beta<\beta_{\rm D}\), with physical mass lower bound
\[
 m_*={-\log\alpha_\beta\over a_*}.
\]
\end{corollary}

\begin{proof}
\Cref{thm:s7v82-gap} is ISMC item (I1).
\end{proof}

\begin{firewall}[Gauge-invariant gap, not a coloured propagator pole]
\label{fire:s7v82-colour}
The theorem concerns the OS Hilbert space generated by gauge-invariant
local cylinders.  It neither constructs nor requires a coloured
one-gluon asymptotic state.  This is the physical alternative allowed by
\Cref{cor:s7v81-mass-row}.
\end{firewall}

\section{The missing ultraviolet-to-terminal bridge}
\label{sec:s7v82-bridge}

\begin{definition}[Terminal Dobrushin matching card, TDMC]
\label{def:s7v82-TDMC}
TDMC consists of a finite RG scale \(J(a)\) and a terminal effective
action such that:
\begin{enumerate}[label=\textup{(T\arabic*)},leftmargin=4em]
\item its local link conditionals have an influence matrix with a
      regulator-uniform row sum \(\alpha_*<1\);
\item it is reflection positive or has a transfer-operator substitute;
\item its terminal spacing \(a_J\) stays comparable to a positive
      physical scale;
\item the V28 matching and off-diagonal errors accumulated through
      scale \(J\) are strictly below the terminal gap margin; and
\item all induced nonplaquette interactions, boundary labels, torons,
      and bad polymers are included in the conditional influence bound.
\end{enumerate}
\end{definition}

\begin{theorem}[Conditional continuum handoff]
\label{thm:s7v82-handoff}
NTRGC through the terminal scale, TDMC, regulator-independent
convergence of the gauge-invariant local state, and the V25 OS
hypotheses imply a continuum Yang--Mills mass gap bounded below by the
terminal decay rate minus the accumulated matching debit.
\end{theorem}

\begin{proof}
TDMC (T1)--(T3) repeat the Dobrushin and OS argument above for the
complete terminal effective action.  Item (T4) permits the V28 operator
recursion to transport the strict terminal contraction back through the
ultraviolet scales.  Item (T5) prevents omission of generated
interactions.  Regulator convergence passes the common decay rate to the
limiting state, and
\Cref{cor:s7v25-regulator} gives the mass gap.
\end{proof}

\begin{firewall}[Bare small beta is not an RG trajectory]
\label{fire:s7v82-trajectory}
\Cref{thm:s7v82-gap} proves a mass gap for the lattice Wilson theory in
its strong-coupling domain.  It does not prove that the continuum
asymptotically free trajectory reaches that domain after finitely many
controlled blockings.  In particular, replacing the terminal effective
action by a bare Wilson action and merely changing its coupling would
discard generated interactions and violate TDMC (T5).
\end{firewall}

\section{Status and next acquisition}
\label{sec:s7v82-status}

\begin{theorem}[What V82 proves]
\label{thm:s7v82-result}
V82 proves an explicit nonempty strong-coupling domain
\(\beta<\beta_{\rm D}\) in which the compact Wilson measure has
volume-uniform exponential spacetime clustering.  Reflection positivity
and density of bounded gauge-invariant cylinders then give an OS
Hamiltonian mass gap at least
\(-a_*^{-1}\log\alpha_\beta\).  Thus the terminal ISMC itself is no
longer an abstract missing theorem for the bare Wilson action.
\end{theorem}

\begin{proof}
Use
\Cref{lem:s7v82-influence,cor:s7v82-alpha,
thm:s7v82-Dobrushin,thm:s7v82-clustering,
cor:s7v82-time,thm:s7v82-gap,
cor:s7v82-ISMC}.
\end{proof}

\begin{assessment}[Progress at seventh-series V82]
\label{ass:s7v82-status}
A genuine gauge-invariant mass-generating regime has been proved with an
explicit contraction constant.  The central bottleneck has moved from
\emph{proving some terminal gap} to the sharper RG statement TDMC:
show that the complete effective action produced from the ultraviolet
continuum trajectory enters a conditional-influence ball with all
generated interactions retained.
\end{assessment}

\begin{decision}[V83 acquisition order]
\label{dec:s7v82-V83}
V83 will extend the Dobrushin calculation from the single Wilson
plaquette action to a general finite-range gauge-invariant effective
interaction.  It will express the terminal criterion in a weighted sum
of interaction oscillations, prove stability under a summable polymer
tail, and determine exactly which RG norm must be driven below one.
This removes the false requirement that the terminal action be a bare
one-parameter Wilson action.
\end{decision}

\begin{firewall}[Claim boundary after seventh-series V82]
\label{fire:s7v82-claim}
V82 does not prove TDMC for the ultraviolet trajectory, construct the
nonlinear coarse-action cocycle, establish regulator-uniform SATC or
full RL4C through cutoff removal, prove interacting endpoint continuity,
UCM, BCC, the pairing/Bogoliubov sector, regulator independence,
continuum OS reconstruction, or the full Yang--Mills mass gap.
\end{firewall}

\paragraph{Parent-version record.}
V82 was formed from
\texttt{Renewal\_Yang\_Mills\_Mass\_Gap\_Research\_Notes\_V81.tex}.
The V81 parent had \(35\,989\) logical source lines,
\(1\,286\,457\) bytes, and SHA--256
\[
 \texttt{155C2D032C0A8E698DCEC3324BCE3D3E7171EC70A75563A4C0BA7C44C2356C14}.
\]
The next compulsory companion audit is V85.

% =============================================================================
% END APPENDED SEVENTH-SERIES V82 MATERIAL
% =============================================================================

% =============================================================================
% BEGIN APPENDED SEVENTH-SERIES V83 MATERIAL
% =============================================================================

\chapter{A Dobrushin terminal ball for complete quasi-local effective
actions}
\label{ch:s7v83}

\section{Purpose}
\label{sec:s7v83-purpose}

V82 proves a terminal mass gap for a bare Wilson plaquette action at
small coupling.  An exact renormalization step does not preserve the
one-plaquette functional form: it produces larger Wilson loops,
multi-loop interactions, boundary-label terms, and polymer corrections.
The correct terminal target is therefore a normed open set of complete
effective actions, not a one-dimensional interval of bare couplings.

This chapter defines that open set directly through mixed conditional
oscillations.  A weighted Dobrushin norm below one gives exponential
clustering even for exponentially decaying infinite-range interactions.
The norm is stable under a summable generated tail.  Finally, an exact
reflection-equivariant coarse pushforward is shown to preserve reflection
positivity, so individual generated couplings need not separately admit
the bare Wilson character proof.

\section{Mixed oscillation of an interaction}
\label{sec:s7v83-oscillation}

Let \(E\) be a finite or countable link set with metric \(d_E\), and let
\(U_e\in G\) for a compact group \(G\).  An interaction is a family of
bounded real functions
\begin{equation}
 \Phi=\{\Phi_X:X\Subset E\},\qquad
 \Phi_X=\Phi_X(U_X).                                  \label{eq:s7v83-interaction}
\end{equation}
For distinct \(e,f\in X\), define the rectangular mixed oscillation
\begin{align}
 \partial_{ef}\Phi_X
 &:=
 \sup
 \bigl|
 \Phi_X(u_e,u_f,\zeta)
 -\Phi_X(u'_e,u_f,\zeta)                             \notag\\
 &\hspace{5em}
 -\Phi_X(u_e,u'_f,\zeta)
 +\Phi_X(u'_e,u'_f,\zeta)
 \bigr|,                                             \label{eq:s7v83-mixed}
\end{align}
where the supremum is over the displayed variables and all remaining
\(\zeta\in G^{X\setminus\{e,f\}}\).  Put
\begin{equation}
 \kappa_{ef}(\Phi)
 :=
 \sum_{X\ni e,f}\partial_{ef}\Phi_X.                 \label{eq:s7v83-kappa}
\end{equation}

\begin{lemma}[Elementary mixed-oscillation bounds]
\label{lem:s7v83-mixed}
\begin{align}
 \partial_{ef}(\Phi_X+\Psi_X)
 &\leq\partial_{ef}\Phi_X+\partial_{ef}\Psi_X,
                                                               \label{eq:s7v83-subadd}\\
 \partial_{ef}\Phi_X
 &\leq2\operatorname{osc}\Phi_X.                    \label{eq:s7v83-osc-bound}
\end{align}
\end{lemma}

\begin{proof}
The first statement is the triangle inequality for the rectangular
increment.  For the second, group the first two and last two terms; each
difference is bounded by \(\operatorname{osc}\Phi_X\).
\end{proof}

\section{The weighted conditional-influence norm}
\label{sec:s7v83-norm}

Define
\begin{equation}
 {\mathfrak D}_\mu(\Phi)
 :=
 \sup_{e\in E}
 \sum_{f\ne e}
 e^{\mu d_E(e,f)}
 {e^{\kappa_{ef}(\Phi)}-1\over2},
 \qquad\mu\geq0.                                     \label{eq:s7v83-Dnorm}
\end{equation}

\begin{theorem}[Effective-action conditional influence]
\label{thm:s7v83-influence}
For the Gibbs specification
\begin{equation}
 d\mu_\Phi(U)
 \propto
 \exp\left[-\sum_X\Phi_X(U_X)\right]\prod_edU_e,      \label{eq:s7v83-Gibbs}
\end{equation}
the one-link Dobrushin coefficient satisfies
\begin{equation}
 c_{ef}\leq
 {e^{\kappa_{ef}(\Phi)}-1\over2}.                    \label{eq:s7v83-cef}
\end{equation}
Consequently
\begin{equation}
 \sup_e\sum_{f\ne e}e^{\mu d_E(e,f)}c_{ef}
 \leq{\mathfrak D}_\mu(\Phi).                        \label{eq:s7v83-weighted-row}
\end{equation}
\end{theorem}

\begin{proof}
Change the conditioned exterior link \(f\) from \(u_f\) to \(u'_f\).
The difference of the two conditional log densities as a function of
\(u_e\) is
\[
 -\sum_{X\ni e,f}
 [\Phi_X(u_e,u_f,\zeta_X)
  -\Phi_X(u_e,u'_f,\zeta_X)].
\]
Its oscillation in \(u_e\) is at most
\(\kappa_{ef}(\Phi)\) by the definition of the rectangular increment.
Apply \Cref{lem:s7v82-tilt} and then sum with the exponential weights.
\end{proof}

\begin{theorem}[Weighted Dobrushin terminal ball]
\label{thm:s7v83-ball}
If, for some \(\mu>0\),
\begin{equation}
 {\mathfrak D}_\mu(\Phi)=:\alpha_\mu<1,               \label{eq:s7v83-ball}
\end{equation}
then the Gibbs specification has a unique infinite-volume state and,
for bounded cylinder functions \(F,G\) at support distance \(r\),
\begin{equation}
 |\operatorname{Cov}_{\mu_\Phi}(F,G)|
 \leq
 {4|S_F|\|F\|_\infty\|G\|_\infty\over1-\alpha_\mu}
 e^{-\mu r}.                                         \label{eq:s7v83-clustering}
\end{equation}
\end{theorem}

\begin{proof}
Let \(C=(c_{ef})\).  The weighted row norm of \(C\) is at most
\(\alpha_\mu\) by
\Cref{thm:s7v83-influence}.  Submultiplicativity of the exponentially
weighted row norm gives
\[
 \sup_e\sum_f e^{\mu d_E(e,f)}(C^n)_{ef}
 \leq\alpha_\mu^n.
\]
Thus
\[
 \sup_e\sum_f e^{\mu d_E(e,f)}D_{ef}
 \leq{1\over1-\alpha_\mu},
 \qquad D=\sum_{n\geq0}C^n.
\]
Use the Dobrushin coupling/covariance theorem
\Cref{thm:s7v82-Dobrushin}.  For
\(e\in S_F,f\in S_G\), extract
\(e^{-\mu r}\) from the weighted sum and use
\(\delta_eF\leq2\|F\|_\infty\),
\(\delta_fG\leq2\|G\|_\infty\).
The same weighted coupling makes boundary influence vanish and proves
uniqueness.
\end{proof}

\begin{corollary}[Terminal mass for a complete effective action]
\label{cor:s7v83-gap}
Suppose the effective state in
\Cref{thm:s7v83-ball} is reflection positive at terminal spacing
\(a_*\), and bounded gauge-invariant cylinders generate the physical OS
space.  Then its Hamiltonian has mass gap at least
\begin{equation}
 m_\Phi={\mu\over a_*}.                               \label{eq:s7v83-gap}
\end{equation}
\end{corollary}

\begin{proof}
Time separation by \(n\) terminal slices gives
\(r\geq n-O(1)\).  Apply
\eqref{eq:s7v83-clustering}, absorb the fixed thickness into the
prefactor, and use the V25 dense-core argument exactly as in
\Cref{thm:s7v82-gap}.
\end{proof}

\section{A linear sufficient norm}
\label{sec:s7v83-linear}

For estimates produced by an RG expansion, define
\begin{equation}
 {\mathfrak O}_\mu(\Phi)
 :=
 \sup_e
 \sum_{X\ni e}
 (|X|-1)e^{\mu\operatorname{diam}X}
 \operatorname{osc}\Phi_X.                           \label{eq:s7v83-Onorm}
\end{equation}
Also put
\begin{equation}
 \kappa_*(\Phi):=\sup_{e\ne f}\kappa_{ef}(\Phi).
                                                               \label{eq:s7v83-kappastar}
\end{equation}

\begin{lemma}[Linear-to-Dobrushin conversion]
\label{lem:s7v83-linear}
If \(\kappa_*(\Phi)\leq K<\infty\), then
\begin{equation}
 {\mathfrak D}_\mu(\Phi)
 \leq e^K{\mathfrak O}_\mu(\Phi).                    \label{eq:s7v83-linear}
\end{equation}
\end{lemma}

\begin{proof}
For \(0\leq x\leq K\),
\((e^x-1)/2\leq e^Kx/2\).  Therefore
\begin{align*}
 {\mathfrak D}_\mu(\Phi)
 &\leq
 {e^K\over2}
 \sup_e\sum_{f\ne e}e^{\mu d_E(e,f)}
 \sum_{X\ni e,f}\partial_{ef}\Phi_X\\
 &\leq
 e^K
 \sup_e\sum_{X\ni e}
 (|X|-1)e^{\mu\operatorname{diam}X}
 \operatorname{osc}\Phi_X,
\end{align*}
where \Cref{lem:s7v83-mixed} was used in the second line.
\end{proof}

\begin{corollary}[Checkable terminal criterion]
\label{cor:s7v83-linear}
The inequality
\begin{equation}
 e^{\kappa_*(\Phi)}{\mathfrak O}_\mu(\Phi)<1          \label{eq:s7v83-linear-criterion}
\end{equation}
is sufficient for the terminal clustering and mass-gap conclusions.
\end{corollary}

\begin{proof}
Use
\Cref{lem:s7v83-linear,thm:s7v83-ball,
cor:s7v83-gap}.
\end{proof}

\begin{remark}[Constants and one-body terms]
\label{rem:s7v83-onebody}
Adding a constant or a one-link term changes neither
\(\partial_{ef}\Phi_X\) nor the Dobrushin influence.  This is correct:
such terms can change a one-site marginal but cannot transmit a boundary
perturbation between distinct links.
\end{remark}

\section{Stability under a generated interaction tail}
\label{sec:s7v83-tail}

Let \(\Phi^0\) be a terminal reference interaction and let
\(\Psi\) be the complete generated tail.  Put
\begin{equation}
 K_0:=\kappa_*(\Phi^0),\qquad
 K_\Psi:=\kappa_*(\Psi).                              \label{eq:s7v83-Ks}
\end{equation}

\begin{theorem}[Open-ball stability]
\label{thm:s7v83-stability}
The combined action obeys
\begin{equation}
 {\mathfrak D}_\mu(\Phi^0+\Psi)
 \leq
 {\mathfrak D}_\mu(\Phi^0)
 +e^{K_0+K_\Psi}{\mathfrak O}_\mu(\Psi).              \label{eq:s7v83-stability}
\end{equation}
Consequently, if
\begin{equation}
 {\mathfrak D}_\mu(\Phi^0)
 +e^{K_0+K_\Psi}{\mathfrak O}_\mu(\Psi)<1,            \label{eq:s7v83-stable-ball}
\end{equation}
the full generated action belongs to the terminal ball.
\end{theorem}

\begin{proof}
By \Cref{lem:s7v83-mixed},
\[
 \kappa_{ef}(\Phi^0+\Psi)
 \leq
 \kappa_{ef}(\Phi^0)+\kappa_{ef}(\Psi).
\]
For \(x,y\geq0\),
\[
 e^{x+y}-1
 =
 (e^x-1)+e^x(e^y-1)
 \leq
 (e^x-1)+e^{K_0+K_\Psi}y
\]
when \(x\leq K_0,y\leq K_\Psi\).  Multiply by
\(\frac12e^{\mu d_E(e,f)}\), sum, and convert the \(y\)-sum by the proof
of \Cref{lem:s7v83-linear}.
\end{proof}

\begin{corollary}[Strong-coupling Wilson action plus a tail]
\label{cor:s7v83-Wilson-tail}
Let \(\Phi^0\) be a Wilson plaquette action at a coupling strictly inside
the V82 Dobrushin domain.  There is
\(\epsilon_{\rm tail}>0\), depending only on its strict margin and local
incidence constants, such that every gauge-invariant generated
\(\Psi\) with
\[
 K_\Psi\leq1,\qquad
 {\mathfrak O}_\mu(\Psi)<\epsilon_{\rm tail}
\]
preserves exponential clustering.
\end{corollary}

\begin{proof}
At \(\mu=0\), V82 gives a strict margin
\({\mathfrak D}_0(\Phi^0)<1\).  Finite range makes
\({\mathfrak D}_\mu(\Phi^0)\) continuous at \(\mu=0\), so it remains
below one for sufficiently small \(\mu>0\).  Choose
\(\epsilon_{\rm tail}\) so that the second term of
\eqref{eq:s7v83-stability} is smaller than the remaining margin.
\end{proof}

\begin{firewall}[Activity norm must control mixed oscillation]
\label{fire:s7v83-activity}
A convergent partition-function or polymer activity norm does not
automatically bound \({\mathfrak O}_\mu(\Psi)\).  Constants and
oscillatory cancellations may make the pressure small while a
conditional derivative is large.  The RG output must include a
support-resolved oscillation estimate for the complete logarithmic
effective density.
\end{firewall}

\section{Reflection positivity under exact coarse pushforward}
\label{sec:s7v83-RP}

Let \((\Omega,\mu,\Theta)\) be reflection positive:
\[
 \int\overline{F(\Theta\omega)}F(\omega)\,d\mu(\omega)
 \geq0
\]
for positive-side observables \(F\).  Let
\({\cal B}:\Omega\to\Omega'\) be a measurable block map and
\(\Theta'\) a coarse reflection satisfying
\begin{equation}
 {\cal B}\Theta=\Theta'{\cal B}.                      \label{eq:s7v83-equivariant}
\end{equation}

\begin{theorem}[Reflection positivity survives exact equivariant
pushforward]
\label{thm:s7v83-RP}
The pushforward measure
\(\mu':={\cal B}_*\mu\) is reflection positive with respect to
\(\Theta'\).
\end{theorem}

\begin{proof}
For every coarse positive-side observable \(f\), the pullback
\(F=f\circ{\cal B}\) is a fine positive-side observable, and
\[
 \int_{\Omega'}
 \overline{f(\Theta'\omega')}f(\omega')\,d\mu'(\omega')
 =
 \int_\Omega
 \overline{F(\Theta\omega)}F(\omega)\,d\mu(\omega)
 \geq0
\]
by \eqref{eq:s7v83-equivariant} and fine reflection positivity.
\end{proof}

\begin{corollary}[Generated couplings need no termwise character test]
\label{cor:s7v83-RP}
If the terminal effective measure is the exact pushforward of the
reflection-positive Wilson measure under a reflection-equivariant
gauge-covariant block map, it is reflection positive even when its
logarithmic effective action contains multi-loop terms with no known
termwise nonnegative character expansion.
\end{corollary}

\begin{proof}
Apply \Cref{thm:s7v83-RP}.  Reflection positivity is a property of the
measure, not of a chosen decomposition of its logarithm.
\end{proof}

\begin{firewall}[Approximate RG maps need an operator error theorem]
\label{fire:s7v83-approximate}
An approximate local action fitted to an exact coarse marginal need not
be reflection positive.  Small coefficient error in
\({\mathfrak O}_\mu\) does not by itself preserve a quadratic
reflection-positivity inequality with zero lower margin.  One must
either retain the exact pushforward state or use the V28
transfer-operator error route.
\end{firewall}

\section{The exact TDMC norm}
\label{sec:s7v83-TDMC}

\begin{definition}[Effective Dobrushin RG card, EDRC]
\label{def:s7v83-EDRC}
At a terminal scale, EDRC consists of:
\begin{enumerate}[label=\textup{(E\arabic*)},leftmargin=4em]
\item the complete logarithmic effective interaction
      \(\Phi^{(J)}=\{\Phi_X^{(J)}\}\), with no generated support omitted;
\item a weight \(\mu_*>0\) and the strict estimate
      \[
      {\mathfrak D}_{\mu_*}(\Phi^{(J)})\leq\alpha_*<1;
      \]
\item uniformity in volume, transition labels, boundary flux sectors,
      torons, and removed ultraviolet cutoffs;
\item an exact reflection-equivariant pushforward representation or the
      NTRGC transfer substitute; and
\item a terminal spacing bounded above and below by positive physical
      constants.
\end{enumerate}
\end{definition}

\begin{theorem}[EDRC loads TDMC and ISMC]
\label{thm:s7v83-EDRC}
EDRC implies terminal exponential clustering with mass
\begin{equation}
 m_*={\mu_*\over a_J}>0,                              \label{eq:s7v83-mass}
\end{equation}
uniformly in the declared regulators.  Hence it loads TDMC (T1)--(T3)
and ISMC.  With the NTRGC accumulated-error row it loads the conditional
continuum handoff of V82.
\end{theorem}

\begin{proof}
Rows (E1)--(E3) permit
\Cref{thm:s7v83-ball} uniformly.  Rows (E4)--(E5) and
\Cref{thm:s7v83-RP,cor:s7v83-gap} give the terminal physical gap.
The remaining ultraviolet transport is exactly NTRGC/TDMC (T4).
\end{proof}

\section{Status and next acquisition}
\label{sec:s7v83-status}

\begin{theorem}[What V83 proves]
\label{thm:s7v83-result}
V83 replaces the bare-Wilson terminal target by the open weighted
Dobrushin ball
\({\mathfrak D}_\mu(\Phi)<1\) for a complete quasi-local effective
action.  It proves exponential clustering and a terminal OS mass in
that ball, gives the checkable linear norm
\({\mathfrak O}_\mu\), proves stability under a summable generated tail,
and proves preservation of reflection positivity under an exact
reflection-equivariant coarse pushforward.
\end{theorem}

\begin{proof}
Use
\Cref{thm:s7v83-influence,thm:s7v83-ball,
cor:s7v83-gap,lem:s7v83-linear,
thm:s7v83-stability,thm:s7v83-RP,
thm:s7v83-EDRC}.
\end{proof}

\begin{assessment}[Progress at seventh-series V83]
\label{ass:s7v83-status}
The terminal target is now stable under the interactions an exact RG
actually generates.  The remaining central problem is dynamical:
control the complete coarse marginal from the ultraviolet good/bad
expansion until its conditional-influence norm enters the EDRC ball.
This is a sharper and weaker target than matching a bare Wilson action.
\end{assessment}

\begin{decision}[V84 acquisition order]
\label{dec:s7v83-V84}
V84 will derive a marginalization theorem for conditional influences.
For a joint retained/eliminated specification whose eliminated
conditional block is Dobrushin-contracting, it will prove the Schur-path
bound
\[
 C_{\rm eff}
 \leq
 C_{RR}+C_{RE}(I-C_{EE})^{-1}C_{ER}.
\]
The result will identify the exact influence transmitted through one
controlled elimination step and provide an EDRC-compatible RG recursion.
\end{decision}

\begin{firewall}[Claim boundary after seventh-series V83]
\label{fire:s7v83-claim}
V83 does not prove EDRC for the ultraviolet trajectory or the
marginalization recursion proposed for V84, construct the nonlinear
coarse-action cocycle, establish regulator-uniform SATC or full RL4C
through cutoff removal, prove interacting endpoint continuity, UCM,
BCC, the pairing/Bogoliubov sector, regulator independence, continuum
OS reconstruction, or the full Yang--Mills mass gap.
\end{firewall}

\paragraph{Parent-version record.}
V83 was formed from
\texttt{Renewal\_Yang\_Mills\_Mass\_Gap\_Research\_Notes\_V82.tex}.
The V82 parent had \(36\,487\) logical source lines,
\(1\,304\,833\) bytes, and SHA--256
\[
 \texttt{8FCC680D72115D7546114B38329599497C138DDBEB4530DFD144954DAAA2D7AF}.
\]
The next compulsory companion audit is V85.

% =============================================================================
% END APPENDED SEVENTH-SERIES V83 MATERIAL
% =============================================================================

% =============================================================================
% BEGIN APPENDED SEVENTH-SERIES V84 MATERIAL
% =============================================================================

\chapter{Marginal conditional influences and the corrected Schur-path
recursion}
\label{ch:s7v84}

\section{Purpose and correction of the proposed formula}
\label{sec:s7v84-purpose}

V83 proposed the schematic marginal influence
\[
 C_{RR}+C_{RE}(I-C_{EE})^{-1}C_{ER}.
\]
That expression accounts for a boundary perturbation which enters the
eliminated variables and returns to a retained target.  It omits a
self-feedback loop: while the retained target is sampled, it can change
the eliminated conditional law, which can in turn change the target.
The safe bound contains the denominator
\[
 1-\bigl[C_{RE}(I-C_{EE})^{-1}C_{ER}\bigr]_{ii}.
\]

This chapter proves that corrected formula by a block maximal coupling.
It derives a weighted norm recursion suitable for EDRC and shows exactly
what one controlled elimination step must verify.  It also records a
crucial scope boundary: bad-set probability does not control a
Dobrushin coefficient, which is a worst-exterior supremum.  A polymer
RG must produce a uniform mixed-oscillation norm, not only an average
rare-event estimate.

\section{Retained and eliminated variables}
\label{sec:s7v84-split}

Let a Gibbs specification have sites
\[
 {\cal I}=R\sqcup E,
\]
where \(R\) is retained and \(E\) is integrated out.  Let
\(C=(c_{ij})\) be a Dobrushin interdependence matrix for the joint
specification, written in blocks
\begin{equation}
 C=
 \begin{pmatrix}
 C_{RR}&C_{RE}\\
 C_{ER}&C_{EE}
 \end{pmatrix}.                                      \label{eq:s7v84-C}
\end{equation}
The row index is the updated variable and the column index is the
exterior variable which is changed.

Assume
\begin{equation}
 \|C_{EE}\|_{\infty}
 :=
 \sup_{e\in E}\sum_{f\in E}c_{ef}
 =:\epsilon_E<1,                                     \label{eq:s7v84-E-contract}
\end{equation}
and define the nonnegative matrix
\begin{equation}
 D_E:=(I-C_{EE})^{-1}
 =\sum_{n=0}^{\infty}C_{EE}^n.                       \label{eq:s7v84-DE}
\end{equation}

\begin{lemma}[Eliminated disagreement response]
\label{lem:s7v84-E-response}
Fix a retained target \(i\in R\).  Couple two conditional systems on
\(\{i\}\cup E\), with exterior retained data differing only at
\(j\in R\setminus\{i\}\).  If \(p_i\) and
\(p_E=(p_e)_{e\in E}\) denote the resulting disagreement
probabilities, they may be chosen to satisfy
\begin{align}
 p_i
 &\leq c_{ij}+\sum_{e\in E}c_{ie}p_e,                \label{eq:s7v84-pi}\\
 p_E
 &\leq C_{Ei}p_i+C_{Ej}+C_{EE}p_E                   \label{eq:s7v84-pE}
\end{align}
entrywise.
\end{lemma}

\begin{proof}
Use maximal coupling at each single-site conditional update.  The
conditional law at \(i\) can differ directly because the boundary site
\(j\) differs, and indirectly through every disagreeing eliminated
site, giving \eqref{eq:s7v84-pi}.  An eliminated conditional can differ
because of the target \(i\), the boundary \(j\), or another eliminated
site, giving \eqref{eq:s7v84-pE}.  Change disagreeing exterior
coordinates one at a time and sum their Dobrushin coefficients.  Iterate
the coupled updates from zero internal disagreement and take the
monotone fixed-point majorant.
\end{proof}

\begin{corollary}[Eliminated resolvent response]
\label{cor:s7v84-E-response}
\begin{equation}
 p_E
 \leq
 D_E(C_{Ei}p_i+C_{Ej}).                              \label{eq:s7v84-pE-solved}
\end{equation}
\end{corollary}

\begin{proof}
Iterate \eqref{eq:s7v84-pE} and use the convergent nonnegative series
\eqref{eq:s7v84-DE}.
\end{proof}

\section{The corrected marginal influence}
\label{sec:s7v84-marginal}

Put
\begin{equation}
 A:=C_{RE}D_EC_{ER}.                                 \label{eq:s7v84-A}
\end{equation}

\begin{theorem}[Single-site marginal influence]
\label{thm:s7v84-marginal}
Let \(c^{\rm eff}_{ij}\) be the Dobrushin coefficient of retained site
\(i\) with respect to retained exterior site \(j\) after marginalizing
all variables in \(E\).  If
\begin{equation}
 A_{ii}<1,                                            \label{eq:s7v84-loop}
\end{equation}
then
\begin{equation}
 c^{\rm eff}_{ij}
 \leq
 {c_{ij}+A_{ij}\over1-A_{ii}},
 \qquad i\ne j.                                      \label{eq:s7v84-effective}
\end{equation}
\end{theorem}

\begin{proof}
Insert \eqref{eq:s7v84-pE-solved} into
\eqref{eq:s7v84-pi}:
\begin{align*}
 p_i
 &\leq
 c_{ij}
 +C_{iE}D_EC_{Ei}\,p_i
 +C_{iE}D_EC_{Ej}\\
 &=c_{ij}+A_{ii}p_i+A_{ij}.
\end{align*}
Move the loop term to the left and use
\eqref{eq:s7v84-loop}.  The marginal disagreement probability at \(i\)
under this coupling bounds the total variation between the two
effective one-site conditionals, which is \(c^{\rm eff}_{ij}\).
\end{proof}

\begin{firewall}[Why the denominator cannot be omitted]
\label{fire:s7v84-loop}
Even when \(C_{EE}=0\), a target can influence an eliminated variable
and be influenced back.  The loop strength is
\(C_{iE}C_{Ei}\).  If it approaches one, an arbitrarily small change at
a different retained boundary site can be amplified.  The schematic
formula in the V83 decision is therefore valid only when this feedback
has separately been absorbed.
\end{firewall}

\begin{remark}[No sign or cancellation is used]
\label{rem:s7v84-positive}
All matrices in \eqref{eq:s7v84-effective} are nonnegative influence
majorants.  The theorem cannot exploit cancellations between different
paths, which is appropriate for a uniform conditional criterion.
\end{remark}

\section{Weighted norm recursion}
\label{sec:s7v84-weighted}

Use one common metric on \(R\sqcup E\) and the weighted row norm
\begin{equation}
 \|K\|_\mu
 :=
 \sup_x\sum_y e^{\mu d(x,y)}K_{xy}                  \label{eq:s7v84-weight}
\end{equation}
for nonnegative rectangular matrices, with the source and target sets
understood from context.  The triangle inequality makes this norm
submultiplicative.

Define
\begin{align}
 \rho&:=\|C_{RR}\|_\mu,&
 r&:=\|C_{RE}\|_\mu,&
 s&:=\|C_{ER}\|_\mu,&
 \epsilon&:=\|C_{EE}\|_\mu.                          \label{eq:s7v84-blocknorms}
\end{align}

\begin{lemma}[Weighted eliminated inverse]
\label{lem:s7v84-weighted-DE}
If \(\epsilon<1\), then
\begin{equation}
 \|D_E\|_\mu\leq{1\over1-\epsilon},\qquad
 \|A\|_\mu\leq
 \tau:={rs\over1-\epsilon}.                          \label{eq:s7v84-tau}
\end{equation}
\end{lemma}

\begin{proof}
Sum the Neumann series for \(D_E\) in the weighted row algebra, then use
submultiplicativity in
\(A=C_{RE}D_EC_{ER}\).
\end{proof}

\begin{theorem}[One-step marginal influence recursion]
\label{thm:s7v84-recursion}
If
\begin{equation}
 \epsilon<1,\qquad \tau<1,                           \label{eq:s7v84-recursion-hyp}
\end{equation}
then the retained marginal has a Dobrushin matrix satisfying
\begin{equation}
 \|C^{\rm eff}\|_\mu
 \leq
 {\rho+\tau\over1-\tau}.                             \label{eq:s7v84-recursion}
\end{equation}
In particular, the effective specification lies in the EDRC terminal
ball whenever
\begin{equation}
 \rho+2\tau<1.                                       \label{eq:s7v84-entry}
\end{equation}
\end{theorem}

\begin{proof}
\Cref{lem:s7v84-weighted-DE} gives
\(A_{ii}\leq\|A\|_\mu\leq\tau\).  Apply
\eqref{eq:s7v84-effective}, multiply by
\(e^{\mu d(i,j)}\), and sum over \(j\ne i\).  The numerator row is
bounded by \(\rho+\tau\), and
\((1-A_{ii})^{-1}\leq(1-\tau)^{-1}\).  This proves
\eqref{eq:s7v84-recursion}.  Its right side is below one exactly under
\eqref{eq:s7v84-entry}.
\end{proof}

\begin{corollary}[Scale-by-scale influence ledger]
\label{cor:s7v84-scale}
For an RG step \(j\), let
\(\rho_j,r_j,s_j,\epsilon_j\) be the four exact conditional-influence
block norms before eliminating the \(E_j\) variables, and put
\[
 \tau_j={r_js_j\over1-\epsilon_j}.
\]
Then
\begin{equation}
 \alpha_{j+1}
 \leq
 {\rho_j+\tau_j\over1-\tau_j}                        \label{eq:s7v84-scale}
\end{equation}
whenever \(\epsilon_j,\tau_j<1\).  This is an
EDRC-compatible influence recursion with all return paths through the
eliminated block included.
\end{corollary}

\begin{proof}
Apply \Cref{thm:s7v84-recursion} at each scale.
\end{proof}

\section{Relation to Gaussian residual elimination}
\label{sec:s7v84-Gaussian}

The V29 fine complement has a spectral/semigroup contraction.  A
Dobrushin coefficient is different: it controls sensitivity of a
conditional probability kernel under a worst-case exterior change.

\begin{firewall}[A spectral gap is not \(\epsilon<1\)]
\label{fire:s7v84-gap}
The fine-sector energy floor and transfer contraction do not imply
\(\|C_{EE}\|_\mu<1\).  The latter requires a normalized conditional-law
comparison, including nonlinear interactions and the complete exterior
label.  V42 occupation resolvents and V53 Gaussian propagators can help
bound it only after a measure/conditional kernel has been specified.
\end{firewall}

\begin{proposition}[Quadratic Gaussian diagnostic]
\label{prop:s7v84-Gaussian}
For a finite-dimensional centered Gaussian with precision matrix
\[
 Q=
 \begin{pmatrix}
 Q_{RR}&Q_{RE}\\
 Q_{ER}&Q_{EE}
 \end{pmatrix},
\qquad Q_{EE}>0,
\]
exact marginalization gives the retained precision
\begin{equation}
 Q_{\rm eff}
 =
 Q_{RR}-Q_{RE}Q_{EE}^{-1}Q_{ER}.                     \label{eq:s7v84-Gaussian-Schur}
\end{equation}
This algebraic Schur complement is the signed quadratic analogue of the
positive influence-path return \(A\), but it does not by itself bound
the nonlinear Dobrushin norm.
\end{proposition}

\begin{proof}
Complete the square in the eliminated Gaussian variables and perform
their normalized integral.  The resulting quadratic form in \(R\) is
\eqref{eq:s7v84-Gaussian-Schur}.  Dobrushin influence takes absolute
conditional changes and therefore replaces cancellation by positive
path sums.
\end{proof}

\section{Rare bad blocks and worst-case influence}
\label{sec:s7v84-rarity}

\begin{theorem}[Rarity does not imply a Dobrushin bound]
\label{thm:s7v84-rarity-no}
For every \(0<p<1\) and every \(0<c<1\), there is a probability measure
on three binary variables \((B,X,Y)\) such that
\[
 \mathbb P(B=1)=p,
\]
but the conditional influence of \(X\) on \(Y\) given \(B=1\) is one.
Consequently no bound depending only on \(p\) can force a
worst-exterior Dobrushin coefficient below \(c\).
\end{theorem}

\begin{proof}
Let \(B\) be Bernoulli with parameter \(p\).  Conditional on \(B=0\),
take \(X,Y\) independent fair signs.  Conditional on \(B=1\), take
\(X\) fair and set \(Y=X\).  Then changing the conditioned value of
\(X\) when \(B=1\) changes the conditional law of \(Y\) between two
point masses, whose total variation is one.  This holds for arbitrarily
small \(p\).
\end{proof}

\begin{corollary}[V70 rarity is not EDRC]
\label{cor:s7v84-rarity-no}
The V70 product domination of bad blocks, even after V79 saturation,
does not by itself imply any row of
\eqref{eq:s7v84-blocknorms}.  To enter the V83 terminal ball, the
coarse marginal must be represented by a complete effective interaction
whose mixed oscillation norm is uniformly small, or one must replace
worst-case Dobrushin by an explicitly averaged disagreement-percolation
theorem.
\end{corollary}

\begin{proof}
Apply \Cref{thm:s7v84-rarity-no} with the bad indicator as \(B\).
Product rarity strengthens multi-event probabilities but leaves the
conditional law on the rare event unrestricted.
\end{proof}

\section{A marginal-influence RG card}
\label{sec:s7v84-card}

\begin{definition}[Marginal influence recursion card, MIRC]
\label{def:s7v84-MIRC}
At each coarse step, MIRC consists of:
\begin{enumerate}[label=\textup{(M\arabic*)},leftmargin=4em]
\item an exact retained/eliminated disintegration of the normalized
      physical measure, compatible with Gauss and transition labels;
\item joint conditional coefficients assembled into the four blocks of
      \eqref{eq:s7v84-C};
\item weighted bounds
      \(\epsilon_j<1,\tau_j<1\), uniform in regulator and admissible
      exterior data;
\item the recursion \eqref{eq:s7v84-scale} for the exact coarse
      marginal;
\item a scale \(J\) at which
      \(\rho_J+2\tau_J<1\); and
\item exact reflection-equivariant pushforward or the NTRGC transfer
      substitute.
\end{enumerate}
\end{definition}

\begin{theorem}[MIRC loads the terminal EDRC entry]
\label{thm:s7v84-MIRC}
MIRC implies that the scale-\(J+1\) exact marginal lies in the weighted
Dobrushin terminal ball, with
\begin{equation}
 \alpha_{J+1}
 \leq
 {\rho_J+\tau_J\over1-\tau_J}<1.                     \label{eq:s7v84-MIRC}
\end{equation}
Together with its reflection and physical-spacing rows, it proves the
terminal ISMC.
\end{theorem}

\begin{proof}
Items (M1)--(M4) permit
\Cref{thm:s7v84-recursion}.  Item (M5) makes the right side strictly
less than one.  Apply
\Cref{thm:s7v83-ball,cor:s7v83-gap}; item (M6) supplies the required
reflection-positive or transfer representation.
\end{proof}

\begin{firewall}[MIRC is not yet loaded by the ultraviolet expansion]
\label{fire:s7v84-MIRC}
V64--V80 control normalized densities, good activities, rare bad
clusters, and boundary labels in integrated \(L^2/L^4\) and polymer
norms.  MIRC asks for worst-exterior conditional total-variation
coefficients of the exact coarse marginal.  A new theorem must convert
the support-resolved polymer expansion into the mixed-oscillation norm
of V83, or develop an averaged terminal criterion which retains the bad
probability.
\end{firewall}

\section{Status and next acquisition}
\label{sec:s7v84-status}

\begin{theorem}[What V84 proves]
\label{thm:s7v84-result}
V84 proves the corrected marginal influence bound
\[
 c^{\rm eff}_{ij}
 \leq
 {c_{ij}+
 [C_{RE}(I-C_{EE})^{-1}C_{ER}]_{ij}
 \over
 1-[C_{RE}(I-C_{EE})^{-1}C_{ER}]_{ii}},
\]
and the weighted recursion
\[
 \|C^{\rm eff}\|_\mu
 \leq{\rho+\tau\over1-\tau},
 \qquad
 \tau={rs\over1-\epsilon}.
\]
It also proves that rare bad-block probability cannot control a
worst-case Dobrushin coefficient.
\end{theorem}

\begin{proof}
Use
\Cref{lem:s7v84-E-response,thm:s7v84-marginal,
lem:s7v84-weighted-DE,thm:s7v84-recursion,
thm:s7v84-rarity-no,thm:s7v84-MIRC}.
\end{proof}

\begin{assessment}[Progress at seventh-series V84]
\label{ass:s7v84-status}
The terminal entry problem now has an exact conditional-influence
recursion and a necessary feedback denominator.  The next upstream
choice is clear.  Either prove uniform mixed oscillations of the complete
coarse logarithmic density, or weaken the terminal criterion to an
averaged disagreement process in which V70 rarity can participate.
\end{assessment}

\begin{decision}[V85 audit and acquisition order]
\label{dec:s7v84-V85}
V85 is the compulsory companion audit.  It will inspect the current SM
and NS notes for a type-compatible conversion from an integrated
good/bad polymer norm to conditional sensitivity or averaged
disagreement.  It will then choose between:
\begin{enumerate}[label=\textup{(\alph*)},leftmargin=3.5em]
\item a uniform logarithmic-density oscillation theorem loading MIRC; or
\item a rare-defect disagreement-percolation theorem which replaces the
worst-case terminal row while retaining exponential OS clustering.
\end{enumerate}
\end{decision}

\begin{firewall}[Claim boundary after seventh-series V84]
\label{fire:s7v84-claim}
V84 does not prove MIRC for the ultraviolet trajectory, EDRC entry, an
averaged rare-defect terminal theorem, the nonlinear coarse-action
cocycle, regulator-uniform SATC or full RL4C through cutoff removal,
interacting endpoint continuity, UCM, BCC, the pairing/Bogoliubov
sector, regulator independence, continuum OS reconstruction, or the
full Yang--Mills mass gap.
\end{firewall}

\paragraph{Parent-version record.}
V84 was formed from
\texttt{Renewal\_Yang\_Mills\_Mass\_Gap\_Research\_Notes\_V83.tex}.
The V83 parent had \(36\,989\) logical source lines,
\(1\,321\,442\) bytes, and SHA--256
\[
 \texttt{7432FE7CAFE2DB6DD3D971B0BA8219CEE1A0991C29DE37D4CD767C2BDAF08CC0}.
\]
The next compulsory companion audit is V85.

% =============================================================================
% END APPENDED SEVENTH-SERIES V84 MATERIAL
% =============================================================================

% =============================================================================
% BEGIN APPENDED SEVENTH-SERIES V85 MATERIAL
% =============================================================================

\chapter{Companion audit and an averaged rare-defect disagreement
criterion}
\label{ch:s7v85}

\section{Purpose}
\label{sec:s7v85-purpose}

V84 proved that a rare bad set cannot improve a worst-exterior
Dobrushin coefficient.  The present compulsory audit confirms that this
is not merely a feature of the Yang--Mills estimates: the companion
Standard Model analysis reaches the same conclusion for a quartic-tail
decomposition.  We therefore take the second branch of the V84
decision.  The target is an averaged, two-replica path estimate in which
the V70 product domination can legitimately enter.

The result below is not a relabelling of Dobrushin uniqueness.  It is a
different sufficient terminal criterion.  Its regular transmission
number is averaged inside an exact replica-connected expansion, while
bad blocks are paid by a product probability.  The theorem proves
exponential clustering and an OS mass once that expansion is loaded.
It also identifies the remaining quantitative obstruction: the
regular-block transmission must itself be subcritical after path
entropy.

\section{Compulsory companion audit}
\label{sec:s7v85-audit}

\subsection{Files inspected}

The active companion files inspected read-only were
\begin{align}
 &\texttt{Renewal\_SM\_Einstein\_Continuum\_Research\_Notes\_V49.tex},
 \label{eq:s7v85-SM-file}\\
 &\texttt{Renewal\_NS\_Stability\_Research\_Notes\_V26.tex}.
 \label{eq:s7v85-NS-file}
\end{align}
The SM file had \(17\,762\) counted source lines, \(643\,054\) bytes,
and SHA--256
\[
 \texttt{13353EDC7A50BA52C1A26E0013354522E31F4D43E022825E1D6A3260719AEC87}.
\]
The NS file had \(5\,320\) counted source lines, \(278\,283\) bytes,
and SHA--256
\[
 \texttt{82C887F8C2C69E4F28FAD7C3DC8DE5B9099CB5A4BF431EBA1FFF3E91935978AD}.
\]
Here the line convention counts the final empty line when a file ends
with a newline.

\subsection{Mathematical comparison}

SM V45 proves that normalized sector weights remove a label-cardinality
factor but do not repair a closing Schur margin.  It also proves that an
absolute one-particle estimate can amplify with spectator number,
whereas an energy-relative form estimate lifts uniformly to Fock
space.  These facts agree with the Yang--Mills V78 and V80 scope
firewalls; they provide no new conditional-mixing coefficient.

SM V49 proves a uniform conditional \(e^{-cR^4}\) Higgs tail under a
quartic coercive lower bound and a uniform local partition floor.  It
then splits a boundary insertion into a deterministic good part and an
integrated bad expectation.  Its explicit firewall states that rarity
pays only the bad expectation and cannot improve the good-sector
operator norm.  This is type-compatible with the V84 conclusion:
probability belongs in an averaged correlation expansion, not in a
worst-exterior supremum.  The SM theorem itself cannot be imported,
because its noncompact Higgs coercivity and exterior-source assumptions
are not statements about compact Yang--Mills links.

NS V26 sharpens pointwise norms for derivatives of a rank-one pressure
kernel in a hypothetical singularity programme.  It contains no Gibbs
specification, replica product domination, conditional-influence
estimate, or reflection-positive transfer theorem.  No NS estimate is
imported.

\begin{theorem}[V85 audit decision]
\label{thm:s7v85-audit}
The companion audit supports the averaged rare-defect branch and
supplies no missing Yang--Mills numerical estimate.  In particular:
\begin{enumerate}[label=\textup{(\roman*)},leftmargin=3.5em]
\item MIRC remains a valid sufficient route if a uniform mixed
      oscillation estimate is later proved;
\item V70 rarity must enter through a normalized expectation or replica
      expansion rather than a worst-case Dobrushin row; and
\item the present acquisition is an averaged replica disagreement
      theorem with an explicit regular-transmission hypothesis.
\end{enumerate}
\end{theorem}

\begin{proof}
The first two conclusions are the common logical content of SM V45,
SM V49, and Yang--Mills V78--V84.  The absence of a type-compatible NS
input follows from the subject and hypotheses of NS V26.  Since V64
already writes Yang--Mills covariances on two replicas and V70--V79
control bad-block products and saturation, the third route uses
available objects without changing the claim type.
\end{proof}

\section{Block graph and bad replicas}
\label{sec:s7v85-graph}

Let \({\cal G}=({\cal V},{\cal E})\) be the fixed-scale block adjacency
graph.  We assume that, from any starting block, the number of oriented
self-avoiding paths of length \(n\) is at most
\begin{equation}
 \Delta^n,\qquad n\geq0,                             \label{eq:s7v85-path-count}
\end{equation}
for a finite geometric constant \(\Delta\geq1\).  This deliberately
uses a crude path-entropy constant; replacing \(\Delta\) by the
connective constant is allowed but is not needed for the theorem.

Let \(\mu\) be one finite-volume normalized physical measure and let
\(B_v\in\{0,1\}\) be its saturated bad-block indicator.  Assume the
product domination
\begin{equation}
 \mu(B_v=1\text{ for every }v\in S)
 \leq p^{|S|}                                       \label{eq:s7v85-product}
\end{equation}
for every finite set \(S\) of distinct blocks, uniformly in the
declared regulators and physical boundary labels.

Take two independent replicas \(U^{(1)},U^{(2)}\) with law
\(\mu\otimes\mu\), and define
\begin{equation}
 D_v:=
 \mathbf 1\{B_v(U^{(1)})=1\text{ or }B_v(U^{(2)})=1\}.
                                                               \label{eq:s7v85-D}
\end{equation}

\begin{lemma}[Replica-union product domination]
\label{lem:s7v85-replica-product}
Under \eqref{eq:s7v85-product},
\begin{equation}
 (\mu\otimes\mu)(D_v=1\text{ for every }v\in S)
 \leq p_2^{|S|},
 \qquad p_2:=\min\{1,2p\}.                           \label{eq:s7v85-p2}
\end{equation}
\end{lemma}

\begin{proof}
For every \(v\in S\), choose a replica in which that block is bad.
The union over all maps \(\sigma:S\to\{1,2\}\) covers the event on the
left.  For a fixed map, independence and
\eqref{eq:s7v85-product} give
\[
 \mu(B_v=1,\ v\in\sigma^{-1}(1))\,
 \mu(B_v=1,\ v\in\sigma^{-1}(2))
 \leq p^{|S|}.
\]
There are \(2^{|S|}\) maps.  This gives \((2p)^{|S|}\); the trivial
probability bound gives the minimum with one.
\end{proof}

\begin{corollary}[Saturated V70 rate on two replicas]
\label{cor:s7v85-saturated}
If the unsaturated indicators obey V70 product domination with rate
\(\pi\), and saturation has radius \(R\), then V79 gives
\begin{equation}
 p\leq D_R\pi^{1/D_R},\qquad
 p_2\leq
 \min\{1,2D_R\pi^{1/D_R}\}.                         \label{eq:s7v85-saturated}
\end{equation}
\end{corollary}

\begin{proof}
The first inequality is the V79 finite-radius saturation theorem.
Apply \Cref{lem:s7v85-replica-product} to its saturated indicators.
\end{proof}

\section{Regular and bad path weights}
\label{sec:s7v85-weights}

Fix a number \(q\in[0,1]\).  It represents the largest transmission
factor through a regular block in the support-resolved two-replica
expansion.  A bad block is assigned the trivial factor one.  Thus
\begin{equation}
 w_v:=q+(1-q)D_v
 =
 \begin{cases}
 q,&D_v=0,\\
 1,&D_v=1.
 \end{cases}                                        \label{eq:s7v85-weight}
\end{equation}
The number \(q\) is not asserted to be a single-site Dobrushin
coefficient.

\begin{lemma}[Annealed weight of one path]
\label{lem:s7v85-one-path}
If \(\gamma\) is self-avoiding and contains \(n\) transmitted blocks,
then
\begin{equation}
 \mathbb E_{\mu\otimes\mu}
 \prod_{v\in\gamma}w_v
 \leq
 \bar q^{\,n},
 \qquad
 \bar q:=q+(1-q)p_2.                                \label{eq:s7v85-qbar}
\end{equation}
\end{lemma}

\begin{proof}
Expand the product:
\[
 \prod_{v\in\gamma}\bigl(q+(1-q)D_v\bigr)
 =
 \sum_{S\subseteq\gamma}
 q^{n-|S|}(1-q)^{|S|}
 \prod_{v\in S}D_v.
\]
The vertices of a self-avoiding path are distinct, so
\Cref{lem:s7v85-replica-product} bounds the expectation of the last
product by \(p_2^{|S|}\).  The binomial theorem gives
\eqref{eq:s7v85-qbar}.
\end{proof}

\begin{remark}[No independence of blocks is assumed]
\label{rem:s7v85-no-independence}
Only the two replicas are independent.  Bad indicators within either
replica may be strongly dependent; the product domination
\eqref{eq:s7v85-product} is precisely what replaces block
independence.
\end{remark}

\section{The replica path-transmission hypothesis}
\label{sec:s7v85-RPT}

For a bounded observable \(F\), write
\[
 \operatorname{osc}F:=\sup F-\inf F
\]
for real \(F\), and use the diameter of its range for complex \(F\).
Let \({\cal P}_n(A,C)\) be the self-avoiding block paths of length \(n\)
which start in \(A\) and reach \(C\).

\begin{definition}[Replica path-transmission card, RPTC]
\label{def:s7v85-RPTC}
At one fixed physical scale, RPTC consists of the following uniform
statement.  There is \(K_{\rm pt}<\infty\) and \(q\in[0,1]\) such that
for every pair of bounded gauge-invariant cylinder observables \(F,G\)
with block supports \(A,C\),
\begin{equation}
 |\operatorname{Cov}_\mu(F,G)|
 \leq
 K_{\rm pt}\operatorname{osc}F\operatorname{osc}G
 \sum_{n\geq d(A,C)}
 \ \sum_{\gamma\in{\cal P}_n(A,C)}
 \mathbb E_{\mu\otimes\mu}
 \prod_{v\in\gamma^\circ}
 \bigl(q+(1-q)D_v\bigr).                            \label{eq:s7v85-RPTC}
\end{equation}
Here \(\gamma^\circ\) contains one transmitted arrival block per path
step, so \(|\gamma^\circ|=n\).  The estimate must hold for the complete
normalized physical measure, not for an unnormalized numerator.
\end{definition}

\begin{firewall}[RPTC is an expansion theorem, not a definition of
smallness]
\label{fire:s7v85-RPTC}
The right side of \eqref{eq:s7v85-RPTC} is useful only after a proof
that every connected replica contribution is assigned to a path, every
regular transmitted block costs \(q\), every bad block is exposed
through \(D_v\), and the remaining decorations sum into
\(K_{\rm pt}\).  V64 gives the two-replica covariance identity, but
that identity alone does not prove RPTC.
\end{firewall}

\begin{proposition}[A sufficient connected-expansion loader]
\label{prop:s7v85-loader}
Suppose an absolutely convergent two-replica covariance expansion has
terms indexed by finite connected decorated block graphs
\(\Gamma\) meeting both \(A\) and \(C\).  Suppose there is a
deterministic rule selecting from each \(\Gamma\) a self-avoiding path
\(\gamma(\Gamma)\) from \(A\) to \(C\), and the absolute sum of all
decorations with one selected path \(\gamma\) is at most
\begin{equation}
 K_{\rm pt}\operatorname{osc}F\operatorname{osc}G
 \mathbb E_{\mu\otimes\mu}
 \prod_{v\in\gamma^\circ}w_v.                       \label{eq:s7v85-loader}
\end{equation}
Then RPTC holds.
\end{proposition}

\begin{proof}
Every connected \(\Gamma\) contains a path between the two supports.
Select the first path in any fixed ordering.  Absolute convergence
permits regrouping by the selected path.  Sum
\eqref{eq:s7v85-loader} over all path lengths and paths.
\end{proof}

\begin{remark}[Why selection matters]
\label{rem:s7v85-selection}
Summing over every path contained in one connected graph can introduce
an uncontrolled multiplicity.  A deterministic selected-path rule
counts each decorated graph exactly once.  This is the replica analogue
of the V79 distinct-edge ownership rule.
\end{remark}

\section{Averaged rare-defect clustering}
\label{sec:s7v85-clustering}

\begin{theorem}[Rare-defect replica path theorem]
\label{thm:s7v85-clustering}
Assume \eqref{eq:s7v85-path-count},
\eqref{eq:s7v85-product}, and RPTC.  Put
\begin{equation}
 \vartheta
 :=
 \Delta\bigl[q+(1-q)p_2\bigr].                      \label{eq:s7v85-theta}
\end{equation}
If
\begin{equation}
 \vartheta<1,                                       \label{eq:s7v85-subcritical}
\end{equation}
then for \(r=d(A,C)\),
\begin{equation}
 |\operatorname{Cov}_\mu(F,G)|
 \leq
 {K_{\rm pt}|A|\over1-\vartheta}\,
 \operatorname{osc}F\operatorname{osc}G\,
 \vartheta^r.                                       \label{eq:s7v85-covariance}
\end{equation}
The estimate is uniform whenever the hypotheses are uniform.
\end{theorem}

\begin{proof}
By \Cref{lem:s7v85-one-path}, every path of length \(n\) has expected
weight at most \(\bar q^n\).  From each of the \(|A|\) starting blocks,
\eqref{eq:s7v85-path-count} bounds the number of candidate paths by
\(\Delta^n\); retaining paths which actually reach \(C\) only lowers
the sum.  Insert these bounds in \eqref{eq:s7v85-RPTC}:
\[
 |\operatorname{Cov}_\mu(F,G)|
 \leq
 K_{\rm pt}|A|\operatorname{osc}F\operatorname{osc}G
 \sum_{n\geq r}(\Delta\bar q)^n.
\]
The geometric series is
\(\vartheta^r/(1-\vartheta)\).
\end{proof}

\begin{corollary}[Explicit rarity margin]
\label{cor:s7v85-margin}
If \(q<1/\Delta\), then
\eqref{eq:s7v85-subcritical} holds whenever
\begin{equation}
 p_2<
 {\,\Delta^{-1}-q\,\over1-q}.                       \label{eq:s7v85-p-margin}
\end{equation}
In particular, if
\[
 q\leq q_*<1/\Delta,\qquad
 \pi\longrightarrow0,
\]
then the V79 saturated rate
\eqref{eq:s7v85-saturated} eventually satisfies the terminal
criterion for every fixed saturation radius.
\end{corollary}

\begin{proof}
Solve
\(\Delta[q+(1-q)p_2]<1\) for \(p_2\).  For fixed \(R\),
\(D_R\pi^{1/D_R}\to0\), so use
\Cref{cor:s7v85-saturated}.
\end{proof}

\begin{firewall}[Rarity cannot repair a supercritical regular path]
\label{fire:s7v85-regular}
As \(p_2\downarrow0\), the criterion tends to
\(\Delta q<1\).  Thus arbitrarily rare bad blocks do not close RPTC
when the regular transmission already exceeds path entropy.  The
missing regular estimate is not a cosmetic constant; it is a necessary
loader for this sufficient route.
\end{firewall}

\begin{remark}[Comparison with worst-case Dobrushin]
\label{rem:s7v85-comparison}
A rare block may have conditional influence one, so V84 correctly
rejects it from a uniform Dobrushin row.  In
\Cref{lem:s7v85-one-path}, the same block costs one only on the event
that it occurs, and product domination pays that event after expansion.
This is the exact location at which averaging changes the conclusion.
\end{remark}

\section{Reflection positivity and the terminal mass}
\label{sec:s7v85-OS}

\begin{theorem}[OS gap from uniform RPTC]
\label{thm:s7v85-OS}
Assume the hypotheses of \Cref{thm:s7v85-clustering} uniformly in
volume and removed auxiliary cutoffs.  Assume in addition:
\begin{enumerate}[label=\textup{(\roman*)},leftmargin=3.5em]
\item the limiting physical state is translation invariant and
      reflection positive;
\item the positive-time gauge-invariant cylinder classes form a dense
      OS Hilbert-space core;
\item one block time step has physical size \(a_*>0\); and
\item the OS time-translation operator \(T\) is a positive
      self-adjoint contraction.
\end{enumerate}
Then the vacuum is unique on the cylinder sector and the reconstructed
Hamiltonian \(H=-a_*^{-1}\log T\) obeys
\begin{equation}
 \operatorname{spec}(H)\setminus\{0\}
 \subseteq
 [m_{\rm rpt},\infty),
 \qquad
 m_{\rm rpt}:=-{1\over a_*}\log\vartheta>0.          \label{eq:s7v85-mass}
\end{equation}
\end{theorem}

\begin{proof}
Exponential clustering from
\eqref{eq:s7v85-covariance} implies uniqueness of the translation
invariant vacuum vector on the cylinder sector by the standard
clustering argument.  Let \([F]\) be a mean-zero vector from the dense
OS cylinder core.  Reflection positivity and time translation identify
\[
 \langle[F],T^n[F]\rangle
 \]
with a reflected Euclidean two-point function whose two supports are
separated by \(n+O_F(1)\) block steps.  Hence
\[
 0\leq\langle[F],T^n[F]\rangle\leq C_F\vartheta^n
\]
after absorbing the fixed support offset into \(C_F\).

Let \(\nu_F\) be the spectral measure of \(T\) for \([F]\).  If it gave
positive mass to \([\vartheta+\epsilon,1]\) for some \(\epsilon>0\),
then
\[
 \int t^n\,d\nu_F(t)
 \geq
 (\vartheta+\epsilon)^n
 \nu_F([\vartheta+\epsilon,1]),
\]
contradicting the preceding bound as \(n\to\infty\).  Thus
\(\nu_F\) is supported in \([0,\vartheta]\).  Density of the core and
boundedness of spectral projections extend this support statement to
the vacuum-orthogonal OS space.  Spectral calculus for
\(H=-a_*^{-1}\log T\) gives \eqref{eq:s7v85-mass}.
\end{proof}

\begin{firewall}[Terminal theorem versus continuum handoff]
\label{fire:s7v85-handoff}
\Cref{thm:s7v85-OS} is a terminal-scale theorem.  It proves a continuum
mass gap only after the exact RG/transfer construction supplies the
uniform limiting state, physical time spacing, dense OS core, and
accumulated-error handoff.  RPTC cannot replace those NTRGC/TDMC rows.
\end{firewall}

\section{The Yang--Mills averaged terminal card}
\label{sec:s7v85-card}

\begin{definition}[Averaged rare-defect path card, ARDPC]
\label{def:s7v85-ARDPC}
ARDPC consists of:
\begin{enumerate}[label=\textup{(A\arabic*)},leftmargin=4em]
\item the exact normalized two-replica covariance expansion of the
      physical Gauss-reduced measure;
\item a selected-path regrouping satisfying
      \eqref{eq:s7v85-loader}, uniformly over volumes, transition
      charts, flux sectors, torons, and removed auxiliary cutoffs;
\item saturated bad-block product domination with rate \(p\);
\item a regular transmission number \(q\) and path entropy \(\Delta\)
      satisfying
      \[
       \Delta[q+(1-q)\min\{1,2p\}]<1;
      \]
\item reflection positivity, a dense gauge-invariant cylinder core,
      and a fixed positive physical block spacing; and
\item the NTRGC accumulated-error and regulator-independence rows
      required to compare the terminal state with the ultraviolet
      Yang--Mills trajectory.
\end{enumerate}
\end{definition}

\begin{theorem}[ARDPC loads the averaged terminal mass card]
\label{thm:s7v85-ARDPC}
ARDPC implies regulator-uniform exponential clustering of bounded
gauge-invariant cylinder observables at the terminal scale and a
positive OS spectral gap bounded below by
\[
 -a_*^{-1}
 \log\!\left(\Delta[q+(1-q)\min\{1,2p\}]\right).
\]
With item (A6), it loads TDMC and the ISMC handoff.
\end{theorem}

\begin{proof}
Items (A1)--(A4) are RPTC and the hypotheses of
\Cref{thm:s7v85-clustering}.  Item (A5) permits
\Cref{thm:s7v85-OS}.  Item (A6) is precisely the still-separate
ultraviolet-to-terminal comparison.
\end{proof}

\begin{proposition}[What the existing good/bad machinery supplies]
\label{prop:s7v85-existing}
The existing Yang--Mills chain supplies the following proper subsets of
ARDPC:
\begin{enumerate}[label=\textup{(\alph*)},leftmargin=3.5em]
\item V64 supplies the normalized two-replica covariance identity;
\item V70 and V79 supply the one-replica product rate and its saturated
      version;
\item V79 supplies distinct-edge forest ownership and finite boundary
      incidence;
\item V80 removes a spurious boundary representation-cardinality
      factor after Gauss projection; and
\item V83 supplies reflection positivity for an exact
      reflection-equivariant coarse pushforward.
\end{enumerate}
They do not yet supply the selected-path absolute expansion or a
uniform \(q<1/\Delta\).
\end{proposition}

\begin{proof}
Each positive statement is the corresponding cited result.  The
missing statements are not consequences of a replica identity, product
domination, boundary ownership, sector normalization, or reflection
positivity separately; they are exactly ARDPC (A2) and the regular part
of (A4).
\end{proof}

\section{Status and next upstream acquisition}
\label{sec:s7v85-status}

\begin{theorem}[What V85 proves]
\label{thm:s7v85-result}
V85 proves that one-replica product domination with rate \(p\) becomes
two-replica bad-union domination with rate \(\min\{1,2p\}\).  For an
exact selected-path covariance expansion with regular transmission
\(q\), it proves exponential clustering under
\[
 \Delta[q+(1-q)\min\{1,2p\}]<1,
\]
and, under the stated reflection-positive OS hypotheses, a spectral
gap of size at least \(-a_*^{-1}\log\vartheta\).
\end{theorem}

\begin{proof}
Use
\Cref{lem:s7v85-replica-product,lem:s7v85-one-path,
prop:s7v85-loader,thm:s7v85-clustering,
thm:s7v85-OS,thm:s7v85-ARDPC}.
\end{proof}

\begin{assessment}[Progress at seventh-series V85]
\label{ass:s7v85-status}
The rare-block probability now participates in a rigorous terminal
criterion without being mistaken for a worst-case influence bound.
The criterion is strictly weaker in that respect than MIRC, but it
exposes a new non-negotiable threshold: regular replica transmission
must beat block-path entropy.  The current ultraviolet machinery has
not yet proved that threshold.
\end{assessment}

\begin{decision}[V86 acquisition order]
\label{dec:s7v85-V86}
V86 will go upstream to the missing regular loader.  It will start from
the V64 normalized two-replica covariance identity, the V76 complete
operation ledger, the V78 good-form activity, and the V79 selected-edge
ownership.  It will derive an abstract selected-path expansion with all
off-path decorations resummed, or prove by counterexample that the
available \(L^2/L^4\) activity estimates are insufficient.  The target
number is the regular transmission \(q\), not another bad-set tail.
\end{decision}

\begin{firewall}[Claim boundary after seventh-series V85]
\label{fire:s7v85-claim}
V85 does not prove RPTC or ARDPC for the Yang--Mills trajectory, the
regular bound \(q<1/\Delta\), MIRC, EDRC entry, the nonlinear
coarse-action cocycle, regulator-uniform SATC or full RL4C through
cutoff removal, interacting endpoint continuity, UCM, BCC, the
pairing/Bogoliubov sector, regulator independence, continuum OS
reconstruction, or the full Yang--Mills mass gap.
\end{firewall}

\paragraph{Parent-version record.}
V85 was formed from
\texttt{Renewal\_Yang\_Mills\_Mass\_Gap\_Research\_Notes\_V84.tex}.
The V84 parent had \(37\,465\) counted source lines,
\(1\,337\,700\) bytes, and SHA--256
\[
 \texttt{6FB9582D4AAB62A939596684A9680A9C541A441E2CC2DD6F8BB5C0B6C8990C1F}.
\]
The next compulsory companion audit is V90.

% =============================================================================
% END APPENDED SEVENTH-SERIES V85 MATERIAL
% =============================================================================

% =============================================================================
% BEGIN APPENDED SEVENTH-SERIES V86 MATERIAL
% =============================================================================

\chapter{Regular transmission, bridge resolvents, and the slow-carrier
obstruction}
\label{ch:s7v86}

\section{Purpose}
\label{sec:s7v86-purpose}

V85 reduced its averaged terminal theorem to an exact selected-path
expansion and a regular transmission below path entropy.  Several
earlier estimates look superficially sufficient: V56 proves a connected
Weyl tree bound, V64 has a replica activity tending to zero, V78 gives a
vanishing good-form occurrence, and V79 pays bad boundaries once.
This chapter checks the types rather than combining their constants
formally.

There are two conclusions.  First, an exact bridge-resolvent theorem
shows how rare polymers preserve exponential decay of a regular
background once a selected bridge expansion exists.  Second, the
existing gapped Weyl/tree background is a conditional fast-sector
theorem.  The law of total covariance leaves a separate slow-label
covariance which can carry all long-range correlation.  Thus the V64 and
V78 small activities do not by themselves load the V85 regular number
for physical Wilson observables.

\section{A weighted kernel algebra}
\label{sec:s7v86-algebra}

Let \(K=(K_{xy})_{x,y\in{\cal V}}\) be a nonnegative kernel on a block
graph.  For \(\mu\geq0\), define
\begin{equation}
 \|K\|_\mu
 :=
 \sup_x\sum_y e^{\mu d(x,y)}K_{xy}.                 \label{eq:s7v86-norm}
\end{equation}

\begin{lemma}[Weighted convolution algebra]
\label{lem:s7v86-algebra}
For nonnegative kernels \(K,L\),
\begin{equation}
 \|KL\|_\mu\leq\|K\|_\mu\|L\|_\mu.                  \label{eq:s7v86-product}
\end{equation}
If \(\|K\|_\mu<1\), then
\begin{equation}
 (I-K)^{-1}=\sum_{n=0}^{\infty}K^n,\qquad
 \|(I-K)^{-1}\|_\mu\leq{1\over1-\|K\|_\mu}.         \label{eq:s7v86-inverse}
\end{equation}
\end{lemma}

\begin{proof}
The triangle inequality for the metric gives
\[
 e^{\mu d(x,z)}
 \leq e^{\mu d(x,y)}e^{\mu d(y,z)}.
\]
Insert this inequality into
\(\sum_z e^{\mu d(x,z)}\sum_yK_{xy}L_{yz}\), exchange the nonnegative
sums, and take the supremum in \(x\).  This proves
\eqref{eq:s7v86-product}.  The Neumann series and its geometric norm
bound then give \eqref{eq:s7v86-inverse}.
\end{proof}

\begin{corollary}[Pointwise decay from a weighted row]
\label{cor:s7v86-pointwise}
For every nonnegative \(K\),
\begin{equation}
 K_{xy}\leq\|K\|_\mu e^{-\mu d(x,y)}.                \label{eq:s7v86-pointwise}
\end{equation}
\end{corollary}

\begin{proof}
The \(y\)-summand in the \(x\)-row of
\eqref{eq:s7v86-norm} is bounded by the whole row.
\end{proof}

\section{Alternating regular and exceptional bridges}
\label{sec:s7v86-bridges}

Let \(G\) be a nonnegative kernel which majorizes a fully resummed
connected correlation through the regular background.  Let \(B\)
majorize one marked exceptional bridge, including its saturated
component, boundary insertions, and all internal decorations.  A chain
with \(n\) exceptional bridges has kernel
\[
 G(BG)^n.
\]

\begin{theorem}[Bridge-resolvent theorem]
\label{thm:s7v86-bridge}
Assume
\begin{equation}
 g_\mu:=\|G\|_\mu<\infty,\qquad
 b_\mu:=\|B\|_\mu<\infty,\qquad
 b_\mu g_\mu<1.                                     \label{eq:s7v86-bridge-hyp}
\end{equation}
Then the complete alternating-chain kernel
\begin{equation}
 {\cal R}:=
 \sum_{n=0}^{\infty}G(BG)^n
 =
 G(I-BG)^{-1}                                       \label{eq:s7v86-R}
\end{equation}
converges in the weighted row algebra and satisfies
\begin{align}
 \|{\cal R}\|_\mu
 &\leq{g_\mu\over1-b_\mu g_\mu},                    \label{eq:s7v86-Rnorm}\\
 {\cal R}_{xy}
 &\leq
 {g_\mu\over1-b_\mu g_\mu}
 e^{-\mu d(x,y)}.                                   \label{eq:s7v86-Rdecay}
\end{align}
\end{theorem}

\begin{proof}
\Cref{lem:s7v86-algebra} gives
\(\|BG\|_\mu\leq b_\mu g_\mu<1\).  Sum its Neumann
series, multiply once by \(G\), and use
\eqref{eq:s7v86-product}.  This proves
\eqref{eq:s7v86-Rnorm}; then apply
\Cref{cor:s7v86-pointwise}.
\end{proof}

\begin{remark}[No smallness of the regular row is required]
\label{rem:s7v86-regular-large}
The regular kernel may have \(g_\mu>1\).  Only the product of one
exceptional bridge and the regular propagation between successive
bridges must be below one.  This is sharper than treating every regular
block as one step of the crude V85 \(\Delta q\) path count.
\end{remark}

\begin{definition}[Selected exceptional-bridge expansion, SEBE]
\label{def:s7v86-SEBE}
SEBE means that the absolute connected covariance of two bounded local
observables is bounded by
\begin{equation}
 C_{F,G}\sum_{n=0}^{\infty}[G(BG)^n]_{AC},           \label{eq:s7v86-SEBE}
\end{equation}
where each connected decorated graph is assigned exactly once to an
ordered chain of saturated exceptional components, all regular
decorations between consecutive components are included in \(G\), and
all decorations owned by one exceptional component are included in
\(B\).
\end{definition}

\begin{corollary}[SEBE clustering]
\label{cor:s7v86-SEBE}
Under SEBE and \eqref{eq:s7v86-bridge-hyp}, the complete covariance
decays exponentially at rate \(\mu\), with prefactor bounded by
\(C_{F,G}g_\mu/(1-b_\mu g_\mu)\).
\end{corollary}

\begin{proof}
Use \Cref{thm:s7v86-bridge} in
\eqref{eq:s7v86-SEBE}.
\end{proof}

\begin{firewall}[A norm bound is not a bridge expansion]
\label{fire:s7v86-SEBE}
V78 controls each good form occurrence and V79 controls each selected
bad boundary edge.  Those bounds do not prove that the normalized
covariance equals, or is absolutely dominated by, the alternating
series \eqref{eq:s7v86-SEBE}.  A numerator expansion with a separately
estimated partition function is insufficient; normalization and
connected selection must occur before absolute values.
\end{firewall}

\section{Rooted bad polymers produce a bridge row}
\label{sec:s7v86-polymer}

The following elementary estimate records the quantitative content
needed from a saturated bad-polymer theorem.  Let \(N_n(x)\) be the
number of connected block sets of cardinality \(n\) containing \(x\).
On a bounded-degree graph there is a geometric constant
\(K_{\rm conn}<\infty\) such that
\begin{equation}
 N_n(x)\leq K_{\rm conn}^{\,n-1}.                   \label{eq:s7v86-connected-count}
\end{equation}

\begin{lemma}[Rooted exceptional-component row]
\label{lem:s7v86-bad-row}
Suppose a connected saturated component \(C\) has absolute marked
weight
\begin{equation}
 |z(C)|\leq z_*^{|C|},                              \label{eq:s7v86-z}
\end{equation}
and its bridge kernel is
\begin{equation}
 B_{xy}:=
 \sum_{\substack{C\ {\rm connected}\\x,y\in C}}|z(C)|.
                                                               \label{eq:s7v86-B}
\end{equation}
Then
\begin{equation}
 \|B\|_\mu
 \leq
 \sum_{n\geq1}
 nK_{\rm conn}^{\,n-1}
 e^{\mu(n-1)}z_*^n.                                 \label{eq:s7v86-Bseries}
\end{equation}
In particular, if
\begin{equation}
 \kappa:=K_{\rm conn}e^\mu z_*<1,                   \label{eq:s7v86-kappa}
\end{equation}
then
\begin{equation}
 \|B\|_\mu
 \leq {z_*\over(1-\kappa)^2}.                       \label{eq:s7v86-Bbound}
\end{equation}
\end{lemma}

\begin{proof}
For a connected \(n\)-block set containing \(x\), every \(y\in C\)
satisfies \(d(x,y)\leq n-1\), and there are at most \(n\) choices of
\(y\).  Use \eqref{eq:s7v86-connected-count} and
\eqref{eq:s7v86-z}, then sum over \(n\).  Since
\(\sum_{n\geq1}n\kappa^{n-1}=(1-\kappa)^{-2}\),
\eqref{eq:s7v86-Bbound} follows.
\end{proof}

\begin{corollary}[Vanishing bad bridge]
\label{cor:s7v86-bad-vanish}
If the V79 marked saturated activity provides
\[
 z_*=O(\widetilde b_M),\qquad
 \widetilde b_M\longrightarrow0,
\]
with fixed block geometry, then for every fixed sufficiently small
\(\mu>0\),
\begin{equation}
 b_\mu=\|B\|_\mu\longrightarrow0.                   \label{eq:s7v86-b-vanish}
\end{equation}
\end{corollary}

\begin{proof}
For large \(M\), \eqref{eq:s7v86-kappa} holds and
\eqref{eq:s7v86-Bbound} tends to zero.
\end{proof}

\begin{firewall}[Probability versus signed polymer weight]
\label{fire:s7v86-probability}
V70 product domination and V79 boundary absorption give the correct
size for \(z_*\), but \eqref{eq:s7v86-z} is an absolute connected
polymer-weight statement.  It must be proved for the same normalized
replica expansion used in SEBE.  A probability bound on the union of
bad blocks does not by itself construct the signed connected weights.
\end{firewall}

\section{What the all-good Weyl expansion actually proves}
\label{sec:s7v86-good}

V53 assumes a gapped number-conserving one-particle operator
\(A\geq\gamma_0I\), and V56 proves its exact Weyl tree card.  For
replica-admissible good interactions, V64 proves that the per-vertex
activity \(q_{\rm rep}\) tends to zero under the logarithmic schedule.
Applying the two-anchor V53 summation therefore gives a regular
fast-sector correlation kernel of the form
\begin{equation}
 G^{\rm fast}_{xy}
 \leq
 C_{\rm fast}e^{-\nu d(x,y)}
 \quad\text{after the thermodynamic limit},         \label{eq:s7v86-fast-decay}
\end{equation}
for the declared compressed Weyl algebra and for slow labels on which
all constants are uniform.

\begin{proposition}[Weighted row of the fast kernel]
\label{prop:s7v86-fast-row}
If the block graph has sphere growth
\begin{equation}
 \#\{y:d(x,y)=r\}\leq C_{\rm sph}e^{\eta r},         \label{eq:s7v86-spheres}
\end{equation}
then for every \(0<\mu<\nu-\eta\),
\begin{equation}
 g_\mu^{\rm fast}
 :=
 \|G^{\rm fast}\|_\mu
 \leq
 {C_{\rm fast}C_{\rm sph}\over
  1-e^{-(\nu-\eta-\mu)}}<\infty.                    \label{eq:s7v86-fast-row}
\end{equation}
\end{proposition}

\begin{proof}
Sum \eqref{eq:s7v86-fast-decay} over distance spheres:
\[
 \sum_y e^{\mu d(x,y)}G^{\rm fast}_{xy}
 \leq
 C_{\rm fast}C_{\rm sph}
 \sum_{r\geq0}e^{-(\nu-\eta-\mu)r}.
\]
The series is geometric under the stated inequality.
\end{proof}

\begin{corollary}[Conditional rare-bridge closure in the fast algebra]
\label{cor:s7v86-fast-bridge}
If SEBE is proved for the same fast Weyl algebra and its exceptional
kernel satisfies \eqref{eq:s7v86-b-vanish}, then
\[
 b_\mu g_\mu^{\rm fast}<1
\]
for all sufficiently large regulators, and the fast covariance retains
exponential decay through the rare saturated components.
\end{corollary}

\begin{proof}
\Cref{prop:s7v86-fast-row} gives a fixed finite row constant and
\Cref{cor:s7v86-bad-vanish} makes the bridge row vanish.  Apply
\Cref{cor:s7v86-SEBE}.
\end{proof}

\begin{remark}[Role of the V78 good-form activity]
\label{rem:s7v86-form}
The V78 quantity \(q_{{\rm form},M}\to0\) is a contribution to the
regular interaction activity which enters \(C_{\rm fast}\) and the
convergence margin of the V53/V56 series.  It is not the whole kernel
\(G^{\rm fast}\): even at zero interaction activity the reference
propagator remains.
\end{remark}

\section{The exact slow--fast covariance split}
\label{sec:s7v86-slow}

Let \({\cal S}\) be the sigma algebra generated by the slow field,
coarse Wilson labels, torons, transition data, and any other variables
held fixed in the fast conditional construction.

\begin{theorem}[Law of total covariance]
\label{thm:s7v86-total-covariance}
For square-integrable observables \(F,G\),
\begin{align}
 \operatorname{Cov}_\mu(F,G)
 &=
 \mathbb E_\mu\!
 \left[
  \operatorname{Cov}_\mu(F,G\mid{\cal S})
 \right]
 \notag\\
 &\quad+
 \operatorname{Cov}_\mu\!
 \left(
  \mathbb E_\mu[F\mid{\cal S}],
  \mathbb E_\mu[G\mid{\cal S}]
 \right).
                                                               \label{eq:s7v86-total-covariance}
\end{align}
\end{theorem}

\begin{proof}
Write \(F=F_0+\mathbb E[F\mid{\cal S}]\) and
\(G=G_0+\mathbb E[G\mid{\cal S}]\), where
\(\mathbb E[F_0\mid{\cal S}]
=\mathbb E[G_0\mid{\cal S}]=0\).
The two cross terms have zero expectation by the tower property.  The
expectation of \(F_0G_0\) is the first term in
\eqref{eq:s7v86-total-covariance}; the product of the conditional means
gives the second after subtracting the global means.
\end{proof}

\begin{corollary}[Fast decay leaves a slow carrier]
\label{cor:s7v86-slow-carrier}
Uniform exponential decay of
\(\operatorname{Cov}(F,G\mid{\cal S})\) controls only the first term of
\eqref{eq:s7v86-total-covariance}.  If \(F\) and \(G\) are
\({\cal S}\)-measurable, that term is identically zero and the complete
covariance is the slow covariance.
\end{corollary}

\begin{proof}
For \({\cal S}\)-measurable \(F,G\), their conditional values are fixed,
so their conditional covariance is zero and their conditional means
are \(F,G\).
\end{proof}

\begin{firewall}[The fast oscillator gap is not the physical mass gap]
\label{fire:s7v86-fast-gap}
The V53 lower bound \(A\geq\gamma_0I\) belongs to the conditional
number-conserving fast reference.  Physical Wilson observables have
nontrivial conditional means as functions of the slow gauge field.
Unless the second term of
\eqref{eq:s7v86-total-covariance} is shown to decay, the fast gap and
the conditional bridge closure give no full Yang--Mills mass gap.
\end{firewall}

\section{Local integrated smallness is insufficient}
\label{sec:s7v86-no-go}

The next example prevents replacing the missing connected selection by
fixed-support \(L^2\) or \(L^4\) closeness.

\begin{theorem}[Small local density defects can hide a global carrier]
\label{thm:s7v86-lp-no-go}
Let \(\nu_N\) be the product law of \(N\) independent fair signs, and
let
\begin{equation}
 \mu_{N,\epsilon}
 :=
 (1-\epsilon)\nu_N
 +{\epsilon\over2}
  \bigl(\delta_{(+1,\ldots,+1)}
       +\delta_{(-1,\ldots,-1)}\bigr),
 \qquad0<\epsilon<1.                                \label{eq:s7v86-mixture}
\end{equation}
For every fixed set \(A\) of \(k\) coordinates, the marginal density
\(f_A=d\mu_A/d\nu_A\) satisfies
\begin{equation}
 \|f_A-1\|_{L^2(\nu_A)}
 =
 \epsilon\sqrt{2^{k-1}-1}.                          \label{eq:s7v86-local-L2}
\end{equation}
Nevertheless, for every \(i\ne j\),
\begin{equation}
 \operatorname{Cov}_{\mu_{N,\epsilon}}(X_i,X_j)
 =\epsilon,                                         \label{eq:s7v86-long}
\end{equation}
independently of their separation.  Hence arbitrarily small
fixed-support normalized \(L^2\) defects do not imply any
volume-uniform exponential covariance bound.
\end{theorem}

\begin{proof}
On the two constant configurations of \(A\),
\[
 f_A=1-\epsilon+\epsilon2^{k-1};
\]
on the other \(2^k-2\) configurations, \(f_A=1-\epsilon\).
Therefore
\begin{align*}
 \|f_A-1\|_2^2
 &=
 2^{1-k}\epsilon^2(2^{k-1}-1)^2
 +(1-2^{1-k})\epsilon^2\\
 &=\epsilon^2(2^{k-1}-1),
\end{align*}
which proves \eqref{eq:s7v86-local-L2}.  Symmetry gives
\(\mathbb E X_i=0\).  The product component has
\(\mathbb E X_iX_j=0\), whereas either constant configuration has
\(X_iX_j=1\), proving \eqref{eq:s7v86-long}.  For fixed
\(\epsilon>0\), a bound \(Ce^{-m d(i,j)}\) with \(m>0\) and
volume-independent \(C\) fails as the separation grows.
\end{proof}

\begin{corollary}[What the V64 normalized defect does not say]
\label{cor:s7v86-V64-scope}
A normalized local replica-pressure or \(L^2\) density bound, even with
a small constant on every fixed support, cannot replace SEBE or RPTC.
One also needs connected locality which excludes a latent slow carrier.
\end{corollary}

\begin{proof}
\Cref{thm:s7v86-lp-no-go} has normalized local densities arbitrarily
close to the product reference while its global latent sign produces
nondecaying covariance.
\end{proof}

\section{Reference transport cannot be omitted}
\label{sec:s7v86-reference}

\begin{theorem}[Uniform exponential decay is closed under local limits]
\label{thm:s7v86-limit}
Let \(\mu_n\) converge locally to \(\mu_0\).  Suppose for every fixed
pair of bounded local observables \(F,G\) and every translate
\(\tau_xG\),
\begin{equation}
 |\operatorname{Cov}_{\mu_n}(F,\tau_xG)|
 \leq C_{F,G}e^{-m|x|}                              \label{eq:s7v86-uniform-decay}
\end{equation}
with \(m>0\) and \(C_{F,G}\) independent of \(n,x\).  Then the same
bound holds for \(\mu_0\).
\end{theorem}

\begin{proof}
For fixed \(x\), local convergence passes the three expectations in the
covariance to the limit.  Take \(n\to\infty\) in
\eqref{eq:s7v86-uniform-decay}.
\end{proof}

\begin{corollary}[A massless limiting reference forbids a uniform
perturbative gap bound]
\label{cor:s7v86-massless}
If a limiting reference state has a local two-point function which does
not obey any positive exponential rate, then no approximating family
converging locally to it can satisfy
\eqref{eq:s7v86-uniform-decay} with a common positive \(m\) and common
local prefactor.  Vanishing interaction activity alone therefore cannot
be used as a proof of a regulator-uniform physical mass.
\end{corollary}

\begin{proof}
The contrapositive is \Cref{thm:s7v86-limit}.
\end{proof}

\begin{remark}[Compatibility with V81]
\label{rem:s7v86-V81}
V81 gave the spectral version of the same warning: a globally
colour-singlet two-gluon trial state of a massless quadratic reference
has energy tending to zero unless nonlinear local Gauss projection or
confining dressing removes it.  The present covariance argument shows
where that carrier sits in the slow--fast disintegration.
\end{remark}

\section{Revised regular-transmission card}
\label{sec:s7v86-card}

\begin{definition}[Slow-aware selected bridge card, SASBC]
\label{def:s7v86-SASBC}
SASBC consists of:
\begin{enumerate}[label=\textup{(S\arabic*)},leftmargin=4em]
\item a common normalized physical measure disintegrated over the slow
      sigma algebra \({\cal S}\);
\item the uniform conditional fast Weyl/tree estimate
      \eqref{eq:s7v86-fast-decay};
\item a normalized connected SEBE for rare saturated components, with
      \(b_\mu g_\mu^{\rm fast}<1\);
\item a separate selected-path or effective-action theorem for the
      slow covariance in
      \eqref{eq:s7v86-total-covariance};
\item compatibility of the two expansions with Gauss projection,
      transition labels, torons, reflection, and cutoff removal; and
\item the TDMC/NTRGC continuum handoff.
\end{enumerate}
\end{definition}

\begin{theorem}[SASBC identifies the complete covariance load]
\label{thm:s7v86-SASBC}
Rows (S2)--(S3) give exponential decay of the conditional fast term in
\eqref{eq:s7v86-total-covariance}.  Row (S4) gives exponential decay of
the slow term.  If both rates are positive and the prefactors are
uniform, their sum decays with the smaller rate.  Rows (S5)--(S6) then
permit the V85 OS and continuum handoffs.
\end{theorem}

\begin{proof}
Apply \Cref{cor:s7v86-fast-bridge} conditionally and integrate its
uniform bound.  Apply (S4) to the second term of
\eqref{eq:s7v86-total-covariance}.  Add the two bounds and use
\(\min\{\mu_{\rm fast},\mu_{\rm slow}\}\).  The remaining conclusions
are \Cref{thm:s7v85-OS,thm:s7v85-ARDPC}.
\end{proof}

\begin{proposition}[Present status of SASBC]
\label{prop:s7v86-status-card}
The existing V53, V56, V64, V78, and V79 estimates provide candidate
inputs for (S2) and the numerical smallness in (S3), subject to the
normalization and selected-expansion firewall.  They provide no proof
of (S4) for the full slow physical measure.
\end{proposition}

\begin{proof}
The first assertion is
\Cref{cor:s7v53-half,thm:s7v53-integrated,
thm:s7v56-card,prop:s7v64-qrep,
thm:s7v78-good-repeat,thm:s7v79-BCRF}.
The second is \Cref{cor:s7v86-slow-carrier,
fire:s7v86-fast-gap}.
\end{proof}

\section{Status and next acquisition}
\label{sec:s7v86-status}

\begin{theorem}[What V86 proves]
\label{thm:s7v86-result}
V86 proves the bridge-resolvent bound
\[
 \|G(I-BG)^{-1}\|_\mu
 \leq{\|G\|_\mu\over1-\|B\|_\mu\|G\|_\mu},
\]
derives a vanishing weighted bridge row from an absolute rooted
bad-polymer activity, and shows how these facts preserve an already
exponentially clustering regular background.  It also proves the exact
slow--fast covariance split and gives a finite-spin counterexample
showing that arbitrarily small fixed-support normalized \(L^2\) defects
can coexist with nondecaying correlations.
\end{theorem}

\begin{proof}
Use
\Cref{thm:s7v86-bridge,lem:s7v86-bad-row,
cor:s7v86-fast-bridge,thm:s7v86-total-covariance,
thm:s7v86-lp-no-go,thm:s7v86-SASBC}.
\end{proof}

\begin{assessment}[Progress at seventh-series V86]
\label{ass:s7v86-status}
The regular-transmission question has split cleanly.  The gapped fast
Weyl sector can absorb rare exceptional bridges once a normalized
selected expansion is proved.  The physical obstruction is the slow
conditional mean, not another fast occurrence estimate.  Attempting to
manufacture a full gap solely from the vanishing V64/V78 activity would
be circular or false.
\end{assessment}

\begin{decision}[V87 acquisition order]
\label{dec:s7v86-V87}
V87 will attack the slow term directly.  It will derive the exact
effective slow marginal and its covariance response under one coarse
integration, distinguish the gauge-orbit zero directions from physical
curvature variables, and test whether the V72--V75 atlas/source cards
give a support-resolved slow interaction norm.  If the slow Gaussian
carrier is massless, the note will identify the first genuinely
nonperturbative Wilson-loop or flux term required to enter either MIRC
or ARDPC.
\end{decision}

\begin{firewall}[Claim boundary after seventh-series V86]
\label{fire:s7v86-claim}
V86 does not prove SEBE, SASBC (S3)--(S6), slow-sector exponential
clustering, RPTC or ARDPC for the physical Yang--Mills measure, MIRC,
EDRC entry, the nonlinear coarse-action cocycle, regulator-uniform SATC
or full RL4C through cutoff removal, interacting endpoint continuity,
UCM, BCC, the pairing/Bogoliubov sector, regulator independence,
continuum OS reconstruction, or the full Yang--Mills mass gap.
\end{firewall}

\paragraph{Parent-version record.}
V86 was formed from
\texttt{Renewal\_Yang\_Mills\_Mass\_Gap\_Research\_Notes\_V85.tex}.
The V85 parent had \(38\,059\) counted source lines,
\(1\,360\,031\) bytes, and SHA--256
\[
 \texttt{1061F2D25F14BCF468078D13A89D8469C077A90BDCC29093E9D21DBC9A0003A2}.
\]
The next compulsory companion audit is V90.

% =============================================================================
% END APPENDED SEVENTH-SERIES V86 MATERIAL
% =============================================================================

% =============================================================================
% BEGIN APPENDED SEVENTH-SERIES V87 MATERIAL
% =============================================================================

\chapter{The exact slow marginal, score covariance, and the local
mass-term obstruction}
\label{ch:s7v87}

\section{Purpose}
\label{sec:s7v87-purpose}

V86 located the uncontrolled part of the physical covariance in the
conditional means over the slow gauge field.  We now compute the slow
marginal exactly.  Differentiating its normalized fast partition
function produces an expected Hessian minus a score covariance.  Under
a uniform conditional fast-clustering hypothesis this makes the slow
Hessian quasi-local.

That fact does not produce a gauge-boson mass term.  Local gauge
invariance forces the quadratic symbol of every analytic
contractible-loop action to vanish at zero momentum.  Thus a
strong-convexity attack on the slow connection would either violate
gauge invariance or assume the desired infrared conclusion.  The
remaining route must control gauge-invariant Wilson/flux propagation
nonperturbatively.

\section{Exact disintegration}
\label{sec:s7v87-disintegration}

At a finite regulator, write the variables as a slow field \(s\) and a
fast field \(f\).  Let \(d\lambda(s)\) and \(d\rho_s(f)\) be positive
reference measures, allowing the fast reference to depend covariantly
on \(s\).  Put
\begin{equation}
 d\mu(s,f)
 =
 Z^{-1}
 e^{-S_0(s)}
 e^{-V(s,f)}
 d\rho_s(f)\,d\lambda(s).                            \label{eq:s7v87-joint}
\end{equation}
For the differentiation formulas below, choose a common local
trivialization in which
\begin{equation}
 d\rho_s(f)=e^{-R(s,f)}\,d\rho(f),\qquad
 W(s,f):=V(s,f)+R(s,f).                              \label{eq:s7v87-W}
\end{equation}
Thus all reference-density derivatives are included in \(W\).

Define
\begin{align}
 Z_{\rm f}(s)
 &:=
 \int e^{-W(s,f)}\,d\rho(f),                         \label{eq:s7v87-Zf}\\
 S_{\rm eff}(s)
 &:=
 S_0(s)-\log Z_{\rm f}(s),                           \label{eq:s7v87-Seff}\\
 d\mu_s(f)
 &:=
 Z_{\rm f}(s)^{-1}e^{-W(s,f)}\,d\rho(f).             \label{eq:s7v87-mus}
\end{align}

\begin{lemma}[Exact slow marginal]
\label{lem:s7v87-marginal}
The slow marginal of \(\mu\) is
\begin{equation}
 d\bar\mu(s)
 =
 \bar Z^{-1}e^{-S_{\rm eff}(s)}\,d\lambda(s),         \label{eq:s7v87-marginal}
\end{equation}
and \(\mu_s\) is the normalized conditional fast law.
\end{lemma}

\begin{proof}
Integrate \eqref{eq:s7v87-joint} in \(f\), using
\eqref{eq:s7v87-W}.  The factor multiplying \(d\lambda(s)\) is
\(e^{-S_0(s)}Z_{\rm f}(s)=e^{-S_{\rm eff}(s)}\).
Dividing the fast integrand by \(Z_{\rm f}(s)\) gives
\eqref{eq:s7v87-mus}.
\end{proof}

\section{First and second slow derivatives}
\label{sec:s7v87-derivatives}

Let \(D_h\) denote a directional derivative along a slow tangent field
\(h\) in the chosen chart.  All following formulas hold under dominated
differentiation; at finite cutoff this follows from the V75
fixed-word bounds once the displayed integrable majorants are verified.

\begin{theorem}[Effective score and Hessian]
\label{thm:s7v87-Hessian}
For slow tangent fields \(h,k\),
\begin{align}
 D_hS_{\rm eff}
 &=
 D_hS_0+\mathbb E_s[D_hW],                           \label{eq:s7v87-first}\\
 D_kD_hS_{\rm eff}
 &=
 D_kD_hS_0
 +\mathbb E_s[D_kD_hW]
 -\operatorname{Cov}_s(D_hW,D_kW).                  \label{eq:s7v87-second}
\end{align}
\end{theorem}

\begin{proof}
Differentiate \eqref{eq:s7v87-Zf}:
\[
 D_hZ_{\rm f}
 =-\int(D_hW)e^{-W}\,d\rho.
\]
Hence
\(-D_h\log Z_{\rm f}=\mathbb E_s[D_hW]\), proving
\eqref{eq:s7v87-first}.  For any differentiable \(A(s,f)\),
\begin{equation}
 D_k\mathbb E_s[A]
 =
 \mathbb E_s[D_kA]
 -\operatorname{Cov}_s(A,D_kW).                     \label{eq:s7v87-response}
\end{equation}
Indeed, differentiate both the numerator and the normalizing
denominator in \eqref{eq:s7v87-mus}.  Apply
\eqref{eq:s7v87-response} to \(A=D_hW\).
\end{proof}

\begin{corollary}[Conditional-mean response]
\label{cor:s7v87-mean-response}
For a slow-dependent observable \(F(s,f)\),
\begin{equation}
 D_h\mathbb E_s[F]
 =
 \mathbb E_s[D_hF]
 -\operatorname{Cov}_s(F,D_hW).                     \label{eq:s7v87-mean-response}
\end{equation}
\end{corollary}

\begin{proof}
This is \eqref{eq:s7v87-response}.
\end{proof}

\begin{firewall}[The covariance term has a fixed sign]
\label{fire:s7v87-sign}
For \(h=k\), the final term in \eqref{eq:s7v87-second} is
\(-\operatorname{Var}_s(D_hW)\leq0\).  Integrating out fast variables
can reduce a slow convexity floor.  Dropping the score covariance gives
an upper, not a lower, Hessian estimate.
\end{firewall}

\begin{remark}[Relation to V75 Fisher moments]
\label{rem:s7v87-Fisher}
V75 controls moments of compact Haar/Wilson scores at finite cutoff.
Those bounds justify and estimate individual terms in
\eqref{eq:s7v87-second}; they do not determine the sign or spatial
cancellation of a cross covariance.
\end{remark}

\section{Quasi-locality inherited from the fast conditional law}
\label{sec:s7v87-quasilocal}

Index slow tangent directions by blocks.  Let \(D_x\) act only in a
fixed collar of block \(x\).  Assume
\begin{equation}
 D_xD_yS_0=0,\qquad D_xD_yW=0
 \quad\text{when }d(x,y)>R_{\rm dir}                \label{eq:s7v87-direct-range}
\end{equation}
for a fixed range \(R_{\rm dir}\).  Suppose the conditional fast law
obeys, uniformly in \(s\),
\begin{equation}
 |\operatorname{Cov}_s(A_X,B_Y)|
 \leq
 C_{\rm f}\|A_X\|_{\diamond}\|B_Y\|_{\diamond}
 e^{-\mu_{\rm f}d(X,Y)}.                            \label{eq:s7v87-fast-cov}
\end{equation}
Here \(\|\cdot\|_\diamond\) is a declared local norm compatible with
the V56/V78 expansion.  Assume
\begin{equation}
 \|D_xW\|_\diamond\leq J_x\leq J_*                 \label{eq:s7v87-score-row}
\end{equation}
uniformly.

\begin{theorem}[Slow effective Hessian is quasi-local]
\label{thm:s7v87-quasilocal}
For \(d(x,y)>R_{\rm dir}\),
\begin{equation}
 |D_yD_xS_{\rm eff}|
 \leq
 C_{\rm f}J_*^2
 e^{-\mu_{\rm f}(d(x,y)-2R_{\rm dir})}.             \label{eq:s7v87-Hessian-decay}
\end{equation}
Consequently, on a block graph with exponential sphere growth exponent
\(\eta<\mu_{\rm f}\), every weighted row
\begin{equation}
 \sup_x\sum_y e^{\mu d(x,y)}
 |D_yD_xS_{\rm eff}|                                \label{eq:s7v87-Hessian-row}
\end{equation}
is finite for \(0<\mu<\mu_{\rm f}-\eta\), provided the finitely many
near-diagonal rows are uniformly bounded.
\end{theorem}

\begin{proof}
Outside the direct range, \eqref{eq:s7v87-second} reduces to
\[
 D_yD_xS_{\rm eff}
 =-\operatorname{Cov}_s(D_xW,D_yW).
\]
The two score supports lie in fixed \(R_{\rm dir}\)-collars, whose
distance is at least \(d(x,y)-2R_{\rm dir}\).  Apply
\eqref{eq:s7v87-fast-cov} and
\eqref{eq:s7v87-score-row}.  The weighted row is then a geometric
sphere sum, as in \Cref{prop:s7v86-fast-row}; the finite direct collar
is controlled by the stated near-diagonal bound.
\end{proof}

\begin{corollary}[Quasi-local response of conditional means]
\label{cor:s7v87-response-local}
If \(F\) is supported in \(A\), has \(D_xF=0\) away from a fixed collar
of \(A\), and \(\|F\|_\diamond<\infty\), then outside that collar
\begin{equation}
 |D_x\mathbb E_s[F]|
 \leq
 C_{\rm f}\|F\|_\diamond J_*
 e^{-\mu_{\rm f}(d(A,x)-R_{\rm dir})}.               \label{eq:s7v87-response-local}
\end{equation}
\end{corollary}

\begin{proof}
Use \eqref{eq:s7v87-mean-response}; the direct derivative vanishes, and
the covariance is bounded by \eqref{eq:s7v87-fast-cov}.
\end{proof}

\begin{firewall}[Finite row is not a Dobrushin margin]
\label{fire:s7v87-row}
The estimate \eqref{eq:s7v87-Hessian-row} proves summability, not
smallness.  The near-diagonal bare Wilson Hessian is of order
\(\beta=g^{-2}\), and taking absolute values in
\eqref{eq:s7v87-second} cannot use cancellation against the score
covariance.  Thus this theorem does not load the V83 inequality
\({\mathfrak D}_\mu<1\).
\end{firewall}

\section{Gauge invariance of the slow marginal}
\label{sec:s7v87-gauge}

Let \({\cal G}\) be the regulator gauge group acting on \(s\) and \(f\).
Assume \(d\lambda\), \(d\rho\), \(S_0\), and \(W\) are invariant under
the simultaneous action, and the fast change of variables induced by a
gauge transformation preserves \(d\rho\).

\begin{lemma}[Fast integration preserves gauge invariance]
\label{lem:s7v87-gauge}
For every \(g\in{\cal G}\),
\begin{equation}
 Z_{\rm f}(g\!\cdot\!s)=Z_{\rm f}(s),\qquad
 S_{\rm eff}(g\!\cdot\!s)=S_{\rm eff}(s).            \label{eq:s7v87-gauge}
\end{equation}
\end{lemma}

\begin{proof}
In \eqref{eq:s7v87-Zf}, change variables
\(f\mapsto g^{-1}\!\cdot f\).  Invariance of \(W\) under the
simultaneous action and invariance of \(d\rho\) give the first identity.
Combine it with invariance of \(S_0\) in
\eqref{eq:s7v87-Seff}.
\end{proof}

\begin{corollary}[Gauge-orbit Ward zero]
\label{cor:s7v87-Ward}
Let \(X_\omega\) be the infinitesimal slow gauge-orbit vector generated
by \(\omega\).  Then
\begin{equation}
 D_{X_\omega}S_{\rm eff}=0.                          \label{eq:s7v87-Ward-first}
\end{equation}
At a stationary slow field,
\begin{equation}
 D_hD_{X_\omega}S_{\rm eff}=0                       \label{eq:s7v87-Ward-second}
\end{equation}
for every tangent \(h\).
\end{corollary}

\begin{proof}
Differentiate \eqref{eq:s7v87-gauge} along the gauge orbit.  Differentiate
once more.  At a stationary point, the derivative of the orbit vector
does not contribute because it is paired with the first derivative of
\(S_{\rm eff}\).
\end{proof}

\begin{firewall}[Convexity must be on a physical quotient]
\label{fire:s7v87-quotient}
No gauge-invariant effective action can have a strictly positive Hessian
on the unreduced connection space, because
\eqref{eq:s7v87-Ward-second} supplies gauge zero directions.  Gauge
fixing may remove those directions, but it cannot turn a bound on gauge
coordinates into a gauge-invariant mass statement without the Gauss and
ghost/Jacobian rows.
\end{firewall}

\section{Contractible Wilson loops have no local Proca row}
\label{sec:s7v87-loop}

We now examine physical transverse directions.  Work near the identity
and write a lattice connection as
\[
 U_e(t)=\exp(tA_e),\qquad A_e\in\mathfrak{su}(3).
\]
For an oriented loop \(\ell\), let
\(\epsilon_{\ell e}\in\mathbb Z\) be its signed edge incidence and put
\begin{equation}
 q_\ell(A):=\sum_e\epsilon_{\ell e}A_e.              \label{eq:s7v87-incidence}
\end{equation}
The linearized holonomy is
\begin{equation}
 W_\ell(U(t))
 =
 I+tq_\ell(A)+O(t^2).                               \label{eq:s7v87-linear-holonomy}
\end{equation}

\begin{lemma}[Closed-loop incidence vanishes at zero momentum]
\label{lem:s7v87-incidence}
For every finite contractible loop \(\ell\), the Fourier incidence
vector \(\widehat q_{\ell,\mu}(p)\) satisfies
\begin{equation}
 \widehat q_{\ell,\mu}(0)=0,\qquad
 |\widehat q_{\ell,\mu}(p)|\leq C_\ell|p|           \label{eq:s7v87-q-zero}
\end{equation}
near \(p=0\).
\end{lemma}

\begin{proof}
A constant Abelianized connection assigns the same Lie-algebra element
to all parallel edges.  A closed contractible lattice path has zero net
oriented displacement in every coordinate direction, so
\(\sum_{e\parallel\mu}\epsilon_{\ell e}=0\).  This is precisely
\(\widehat q_{\ell,\mu}(0)=0\).  Since the loop has finite length, its
Fourier incidence is a finite trigonometric polynomial.  The mean value
theorem then gives the linear bound near zero.
\end{proof}

\begin{theorem}[Local analytic loop actions have a massless quadratic
symbol]
\label{thm:s7v87-no-Proca}
Let a translation-invariant slow action near the identity be an
absolutely summable analytic combination of finite contractible Wilson
loops,
\begin{equation}
 S_{\rm loop}(U)
 =
 \sum_{\ell,x}c_\ell\,
 \phi_\ell(W_{\ell+x}(U)),                           \label{eq:s7v87-loop-action}
\end{equation}
where every class function \(\phi_\ell\) is stationary at the identity
and the coefficient-weighted second derivatives are summable with two
loop-length moments.  Its quadratic Fourier symbol \(K(p)\) obeys
\begin{equation}
 \|K(p)\|\leq C|p|^2,\qquad K(0)=0.                 \label{eq:s7v87-K-zero}
\end{equation}
In particular, it contains no positive local Proca term
\(m^2\sum_\mu\|A_\mu\|^2\).
\end{theorem}

\begin{proof}
By \eqref{eq:s7v87-linear-holonomy} and stationarity of each class
function, its quadratic contribution is a bounded bilinear form applied
to \(q_\ell(A)\) twice.  In momentum space it therefore has the form
\[
 \widehat A_\mu(p)^*
 \widehat q_{\ell,\mu}(p)^*
 H_\ell
 \widehat q_{\ell,\nu}(p)
 \widehat A_\nu(p).
\]
\Cref{lem:s7v87-incidence} bounds the two incidence factors by
\(C_\ell^2|p|^2\).  The assumed coefficient and length-moment
summability permits summation over loops and gives
\eqref{eq:s7v87-K-zero}.  A Proca row has nonzero symbol
\(m^2I\) at \(p=0\), so it is absent.
\end{proof}

\begin{corollary}[Fast integration cannot perturbatively generate a
gauge-field mass row]
\label{cor:s7v87-no-Proca}
If the exact slow effective action remains in the quasi-local analytic
contractible-loop class of \Cref{thm:s7v87-no-Proca}, then its quadratic
connection symbol still vanishes at zero momentum, including the
score-covariance contribution in \eqref{eq:s7v87-second}.
\end{corollary}

\begin{proof}
\Cref{lem:s7v87-gauge} preserves gauge invariance, and the stated
quasi-local analytic representation places the complete effective
action under \Cref{thm:s7v87-no-Proca}.
\end{proof}

\begin{firewall}[Confinement is not a gluon Proca mass]
\label{fire:s7v87-confinement}
The Yang--Mills mass gap concerns the gauge-invariant OS spectrum.  It
need not arise from a quadratic mass term for the gauge potential, and
gauge invariance forbids such a local term in the class above.
Therefore failure of a Proca row is not evidence against a mass gap; it
is evidence that the proof must use non-Gaussian Wilson/flux disorder,
not strict connection-space convexity.
\end{firewall}

\section{The nonperturbative slow card}
\label{sec:s7v87-card}

\begin{definition}[Confining slow-marginal card, CSMC]
\label{def:s7v87-CSMC}
CSMC consists of:
\begin{enumerate}[label=\textup{(C\arabic*)},leftmargin=4em]
\item the exact normalized marginal
      \eqref{eq:s7v87-marginal}, including transition-chart and
      reference-density terms in \(W\);
\item gauge invariance, Gauss compatibility, and reflection positivity
      of the slow marginal or of its exact pushforward state;
\item a complete quasi-local effective interaction with a
      support-resolved norm uniform in removed fast cutoffs;
\item a nonperturbative correlation mechanism for gauge-invariant
      Wilson/flux observables, expressed as MIRC, ARDPC, a transfer
      contraction, or an equivalent terminal card;
\item uniform physical time spacing and a dense OS cylinder core; and
\item the NTRGC/TDMC comparison with the ultraviolet physical state.
\end{enumerate}
\end{definition}

\begin{theorem}[CSMC closes the V86 slow carrier]
\label{thm:s7v87-CSMC}
If CSMC (C4) gives exponential slow covariance with rate
\(\mu_{\rm slow}>0\), and the conditional fast term satisfies SASBC
(S2)--(S3) with rate \(\mu_{\rm fast}>0\), then the full physical
covariance decays with rate
\[
 \min\{\mu_{\rm slow},\mu_{\rm fast}\}.
\]
Rows (C2), (C5), and (C6) then give the corresponding OS and continuum
handoffs.
\end{theorem}

\begin{proof}
Apply the total-covariance identity
\eqref{eq:s7v86-total-covariance}.  The first term is controlled by the
conditional SASBC rows and the second by CSMC (C4).  Add the bounds.
The spectral and handoff conclusions are
\Cref{thm:s7v85-OS,thm:s7v86-SASBC}.
\end{proof}

\begin{proposition}[What V72--V75 contribute]
\label{prop:s7v87-V72-V75}
V72--V75 contribute finite-cutoff atlas form debits, normalized label
contractions, mixed \(L^2/L^4\) polymer summability, and polynomial
score/derivative bounds.  They help justify the differentiations and
near-diagonal constants in
\eqref{eq:s7v87-second}--\eqref{eq:s7v87-Hessian-row}.  They do not
prove CSMC (C3)--(C6) or a nonperturbative slow correlation mechanism.
\end{proposition}

\begin{proof}
Use
\Cref{thm:s7v72-SATC,thm:s7v73-closure,
thm:s7v74-closure,thm:s7v75-RL4C}.
Their claim firewalls retain regulator-uniform SATC/RL4C, global
physical connected normalization, and the terminal mass handoff.
\end{proof}

\section{Status and next acquisition}
\label{sec:s7v87-status}

\begin{theorem}[What V87 proves]
\label{thm:s7v87-result}
V87 proves the exact slow effective action and its derivative identities
\[
 D^2S_{\rm eff}
 =
 D^2S_0+\mathbb E_sD^2W
 -\operatorname{Cov}_s(DW,DW).
\]
Uniform conditional fast clustering makes its off-diagonal Hessian
quasi-local.  Gauge invariance supplies orbit zero modes, and every
analytic summable contractible-loop action has a quadratic connection
symbol \(K(p)=O(|p|^2)\), so no local Proca row can close the slow
physical carrier.
\end{theorem}

\begin{proof}
Use
\Cref{lem:s7v87-marginal,thm:s7v87-Hessian,
thm:s7v87-quasilocal,lem:s7v87-gauge,
cor:s7v87-Ward,thm:s7v87-no-Proca,
thm:s7v87-CSMC}.
\end{proof}

\begin{assessment}[Progress at seventh-series V87]
\label{ass:s7v87-status}
The slow marginal is no longer a schematic remainder: its normalization,
score response, quasi-locality interface, and gauge Ward structure are
explicit.  The analysis rules out the wrong acquisition target of
strict local connection convexity.  The remaining mass mechanism must
be formulated directly on gauge-invariant transfer or Wilson/flux
observables.
\end{assessment}

\begin{decision}[V88 acquisition order]
\label{dec:s7v87-V88}
V88 will formulate and prove a terminal transfer criterion directly in
compact-group representation channels.  It will separate the vacuum
character from nontrivial electric-flux characters, include normalized
Gauss boundary pairing from V80, and determine which uniform character
ratio or transfer-kernel minorization would give a positive
gauge-invariant gap without a Proca term.  It will then test whether the
small-\(\beta\) V82 theorem is the first instance of that criterion.
\end{decision}

\begin{firewall}[Claim boundary after seventh-series V87]
\label{fire:s7v87-claim}
V87 does not prove CSMC, a nonperturbative slow correlation bound,
MIRC or ARDPC for the physical slow marginal, terminal entry of the
ultraviolet trajectory, regulator-uniform SATC or full RL4C, the
nonlinear coarse-action cocycle, interacting endpoint continuity, UCM,
BCC, the pairing/Bogoliubov sector, regulator independence, continuum
OS reconstruction, or the full Yang--Mills mass gap.
\end{firewall}

\paragraph{Parent-version record.}
V87 was formed from
\texttt{Renewal\_Yang\_Mills\_Mass\_Gap\_Research\_Notes\_V86.tex}.
The V86 parent had \(38\,689\) counted source lines,
\(1\,382\,380\) bytes, and SHA--256
\[
 \texttt{B5D4170727DBCB0AF75A047303373EC0A19D1F7F8A0CC2CC39C884AC1B8D3A70}.
\]
The next compulsory companion audit is V90.

% =============================================================================
% END APPENDED SEVENTH-SERIES V87 MATERIAL
% =============================================================================

% =============================================================================
% BEGIN APPENDED SEVENTH-SERIES V88 MATERIAL
% =============================================================================

\chapter{Compact-group character channels and a gauge-invariant
transfer gap}
\label{ch:s7v88}

\section{Purpose}
\label{sec:s7v88-purpose}

V87 showed why a local quadratic mass for the connection is the wrong
target.  This chapter works directly with the reflection-positive
transfer operator.  A central compact-group convolution diagonalizes in
Peter--Weyl character channels.  If every nontrivial channel has a
uniform eigenvalue ratio below one, tensorization and Gauss projection
preserve that contraction without a representation-cardinality loss.
The result gives a genuine gauge-invariant spectral gap for an
ultralocal electric transfer.

We then formulate the corresponding block-channel Schur test for an
interacting transfer.  This separates the proved compact-group theorem
from the still-missing physical estimate: spatial plaquette interactions
mix cut-flux channels, and their total off-diagonal row must be controlled
by a connected local expansion rather than a volume-extensive operator
norm.

\section{Central convolution on a compact group}
\label{sec:s7v88-convolution}

Let \(G\) be a compact group with normalized Haar measure.  Let
\(k\in L^1(G)\) be central, real, nonnegative, and normalized by
\begin{equation}
 \int_Gk(g)\,dg=1.                                   \label{eq:s7v88-normalized}
\end{equation}
Define
\begin{equation}
 (K_kf)(g)
 :=
 \int_G k(gh^{-1})f(h)\,dh.                         \label{eq:s7v88-convolution}
\end{equation}
For an irreducible unitary representation \(\lambda\), write
\(d_\lambda=\dim V_\lambda\) and
\(\chi_\lambda=\operatorname{Tr}D^\lambda\).  Put
\begin{equation}
 a_\lambda
 :=
 {1\over d_\lambda}
 \int_G k(g)\overline{\chi_\lambda(g)}\,dg.          \label{eq:s7v88-alambda}
\end{equation}

\begin{lemma}[Character diagonalization]
\label{lem:s7v88-diagonal}
Every matrix coefficient of \(D^\lambda\) is an eigenfunction of
\(K_k\) with eigenvalue \(a_\lambda\).  The trivial coefficient is
\begin{equation}
 a_{\mathbf1}=1.                                    \label{eq:s7v88-trivial}
\end{equation}
If \(k\) is positive definite, then
\begin{equation}
 a_\lambda\geq0                                     \label{eq:s7v88-positive}
\end{equation}
for every \(\lambda\).
\end{lemma}

\begin{proof}
The Fourier transform
\[
 \widehat k(\lambda)
 :=
 \int_G k(g)D^\lambda(g^{-1})\,dg
\]
commutes with every \(D^\lambda(u)\), because \(k\) is central.
Schur's lemma makes it a scalar \(a_\lambda I\).  Taking traces gives
\eqref{eq:s7v88-alambda}.  The convolution formula then gives the
eigenvalue statement.  Normalization yields
\eqref{eq:s7v88-trivial}.  Positive definiteness of \(k\) is equivalent,
by the compact-group Bochner theorem, to positive semidefiniteness of
all Fourier transforms; hence their scalar eigenvalues are nonnegative.
\end{proof}

\begin{theorem}[One-link character gap]
\label{thm:s7v88-one-link}
Assume \(k\) is positive definite and
\begin{equation}
 r_k:=\sup_{\lambda\ne\mathbf1}a_\lambda<1.          \label{eq:s7v88-rk}
\end{equation}
Then \(K_k\) is a positive self-adjoint contraction, its constant
eigenfunction is the unique vacuum whenever no nontrivial
\(a_\lambda\) equals one, and
\begin{equation}
 \|K_k\|_{\mathbf1^\perp}=r_k.                       \label{eq:s7v88-one-gap}
\end{equation}
For a time step \(a_t>0\), the Hamiltonian
\(H_k=-a_t^{-1}\log K_k\) has gap
\begin{equation}
 m_k=-a_t^{-1}\log r_k>0.                            \label{eq:s7v88-one-mass}
\end{equation}
\end{theorem}

\begin{proof}
Peter--Weyl matrix coefficients form a complete orthogonal basis of
\(L^2(G)\).  By \Cref{lem:s7v88-diagonal}, the spectrum consists of the
numbers \(a_\lambda\), with the usual multiplicities.  Positivity gives
\(a_\lambda\geq0\), while the Markov normalization gives operator norm
one.  The constant sector is \(\lambda=\mathbf1\); on its orthogonal
complement the norm is the supremum
\eqref{eq:s7v88-rk}.  Spectral calculus gives
\eqref{eq:s7v88-one-mass}.
\end{proof}

\section{An explicit small-action character bound}
\label{sec:s7v88-small}

Let \(\phi:G\to\mathbb R\) be central with
\begin{equation}
 |\phi(g)|\leq B,                                    \label{eq:s7v88-B}
\end{equation}
and define
\begin{equation}
 k_\beta(g)
 :=
 {e^{\beta\phi(g)}\over Z_\beta},
 \qquad
 Z_\beta:=\int_Ge^{\beta\phi(g)}\,dg.                \label{eq:s7v88-kbeta}
\end{equation}

\begin{theorem}[Uniform nontrivial character ratio]
\label{thm:s7v88-ratio}
For every nontrivial irreducible \(\lambda\),
\begin{equation}
 |a_\lambda(\beta)|
 \leq
 {e^{\beta B}(e^{\beta B}-1)\over d_\lambda}
 \leq
 r_\beta:=
 e^{\beta B}(e^{\beta B}-1).                        \label{eq:s7v88-rbeta}
\end{equation}
Consequently \(r_\beta<1\) whenever
\begin{equation}
 \beta B<
 \log{1+\sqrt5\over2}.                              \label{eq:s7v88-beta-domain}
\end{equation}
If \(k_\beta\) is positive definite, this gives the one-link gap
\(-a_t^{-1}\log r_\beta\).
\end{theorem}

\begin{proof}
For nontrivial \(\lambda\), Haar orthogonality gives
\(\int_G\chi_\lambda\,dg=0\).  Hence
\[
 a_\lambda(\beta)
 =
 {1\over d_\lambda Z_\beta}
 \int_G(e^{\beta\phi(g)}-1)
 \overline{\chi_\lambda(g)}\,dg.
\]
The partition function obeys \(Z_\beta\geq e^{-\beta B}\), and
\[
 |e^{\beta\phi}-1|\leq e^{\beta B}-1.
\]
Character orthogonality and Cauchy--Schwarz give
\[
 \int_G|\chi_\lambda(g)|\,dg
 \leq
 \left(\int_G|\chi_\lambda(g)|^2\,dg\right)^{1/2}
 =1.
\]
These three bounds prove \eqref{eq:s7v88-rbeta}.  Writing
\(x=e^{\beta B}\), the inequality \(x(x-1)<1\) is equivalent to
\(x<(1+\sqrt5)/2\), which is
\eqref{eq:s7v88-beta-domain}.  Apply
\Cref{thm:s7v88-one-link} under positive definiteness.
\end{proof}

\begin{lemma}[Wilson one-link density is positive definite]
\label{lem:s7v88-Wilson-positive}
For \(G=SU(N)\) and
\begin{equation}
 \phi(g)=\operatorname{ReTr}g
 ={1\over2}\bigl(\chi_{\rm f}(g)+
                  \chi_{\overline{\rm f}}(g)\bigr), \label{eq:s7v88-Wilson-phi}
\end{equation}
\(e^{\beta\phi}\), and hence its normalized density, is positive
definite for every \(\beta\geq0\).
\end{lemma}

\begin{proof}
Expand
\[
 e^{\beta\phi}
 =
 \sum_{m,n\geq0}
 {(\beta/2)^{m+n}\over m!\,n!}
 \chi_{\rm f}^{\,m}\chi_{\overline{\rm f}}^{\,n}.
\]
Each product is the character of a tensor product representation, and
its irreducible decomposition has nonnegative integer multiplicities.
Thus the full character expansion has nonnegative coefficients.
Compact-group Bochner positivity follows.  Division by the positive
constant \(Z_\beta\) preserves it.
\end{proof}

\begin{corollary}[Explicit ultralocal Wilson character gap]
\label{cor:s7v88-Wilson-gap}
For the normalized \(SU(3)\) density
\[
 k_\beta(g)\propto
 e^{\beta\operatorname{ReTr}g},
\]
one may take \(B=3\) in
\eqref{eq:s7v88-beta-domain}.  Therefore the convolution transfer has a
strict character gap whenever
\begin{equation}
 \beta<
 {1\over3}\log{1+\sqrt5\over2}.                     \label{eq:s7v88-SU3-domain}
\end{equation}
\end{corollary}

\begin{proof}
For a unitary \(3\times3\) matrix,
\(|\operatorname{ReTr}g|\leq3\).  Apply
\Cref{thm:s7v88-ratio,lem:s7v88-Wilson-positive}.
\end{proof}

\begin{remark}[The numerical domain is deliberately crude]
\label{rem:s7v88-crude}
Exact character coefficients give a larger domain.  The value in
\eqref{eq:s7v88-SU3-domain} is useful because it is uniform over all
nontrivial representations and follows from only Haar orthogonality and
boundedness.
\end{remark}

\section{Tensorization and Gauss projection}
\label{sec:s7v88-Gauss}

For a finite spatial slice with link set \(E\), define
\begin{equation}
 T_E:=\bigotimes_{e\in E}K_k
 \quad\text{on}\quad
 {\cal H}_E:=L^2(G^E).                               \label{eq:s7v88-product-transfer}
\end{equation}
Let \(\Omega_E=1\) be the constant vector and
\(Q_E=I-|\Omega_E\rangle\langle\Omega_E|\).

\begin{theorem}[Volume-uniform tensor character gap]
\label{thm:s7v88-tensor}
Under \eqref{eq:s7v88-rk},
\begin{equation}
 \|Q_ET_EQ_E\|\leq r_k                              \label{eq:s7v88-tensor-gap}
\end{equation}
for every finite \(E\), independently of \(|E|\).
\end{theorem}

\begin{proof}
The tensor Peter--Weyl basis is indexed by one irreducible
\(\lambda_e\) on each link.  Its transfer eigenvalue is
\(\prod_ea_{\lambda_e}\).  The vacuum has every
\(\lambda_e=\mathbf1\).  Every orthogonal basis vector has at least one
nontrivial label, whose factor is at most \(r_k\); all remaining factors
are at most one.  Taking the supremum proves
\eqref{eq:s7v88-tensor-gap}.
\end{proof}

Let \(P_{\cal G}\) be the orthogonal projector onto the local
gauge-invariant slice Hilbert space.  Central convolution commutes with
left and right translations and hence with \(P_{\cal G}\).

\begin{corollary}[Gauss projection preserves the character gap]
\label{cor:s7v88-Gauss}
On
\({\cal H}_{\rm phys}:=P_{\cal G}{\cal H}_E\),
\begin{equation}
 \left\|
 (P_{\cal G}Q_E)
 T_E
 (Q_EP_{\cal G})
 \right\|
 \leq r_k.                                          \label{eq:s7v88-Gauss-gap}
\end{equation}
Thus Gauss reduction introduces neither a representation-cardinality
factor nor a volume loss.
\end{corollary}

\begin{proof}
The constant vacuum is gauge invariant, so \(P_{\cal G}\) commutes with
\(Q_E\).  Compression by an orthogonal projection cannot increase
operator norm.  Apply \Cref{thm:s7v88-tensor}.
\end{proof}

\begin{firewall}[Nontrivial physical states are flux networks]
\label{fire:s7v88-flux}
Gauss projection removes isolated charged link states, but it retains
closed spin networks and paired boundary fluxes.  Their tensor labels
still contain nontrivial irreducibles, so
\eqref{eq:s7v88-Gauss-gap} applies.  No gauge-potential mass term is
used.
\end{firewall}

\section{Interacting flux-channel Schur test}
\label{sec:s7v88-channel}

The physical interacting transfer need not preserve each cut-flux
label.  We therefore use a block operator bound.  Let
\begin{equation}
 {\cal H}
 =
 \mathbb C\Omega
 \oplus
 \bigoplus_{\lambda\in\Lambda_*}{\cal M}_\lambda,    \label{eq:s7v88-decomposition}
\end{equation}
where \(\Lambda_*\) includes the nonvacuum part of the trivial boundary
sector as well as every nontrivial paired flux sector.  Let \(T\) be a
positive self-adjoint contraction with \(T\Omega=\Omega\), and write
\[
 T_{\lambda\mu}:=
 P_\lambda TP_\mu
\]
on the vacuum-orthogonal sum.

\begin{lemma}[Vacuum cross blocks vanish]
\label{lem:s7v88-vacuum-block}
Self-adjointness and \(T\Omega=\Omega\) imply
\begin{equation}
 QT\Omega=0,\qquad
 \langle\Omega,TQ\,\cdot\,\rangle=0,                \label{eq:s7v88-vacuum-zero}
\end{equation}
where \(Q=I-|\Omega\rangle\langle\Omega|\).
\end{lemma}

\begin{proof}
The first identity is \(Q\Omega=0\).  For \(\psi\perp\Omega\),
\(\langle\Omega,T\psi\rangle
=\langle T\Omega,\psi\rangle=0\).
\end{proof}

\begin{theorem}[Flux-channel Schur contraction]
\label{thm:s7v88-Schur}
Assume
\begin{equation}
 s_*:=
 \sup_{\lambda\in\Lambda_*}
 \sum_{\mu\in\Lambda_*}
 \|T_{\lambda\mu}\|<1.                              \label{eq:s7v88-Schur-row}
\end{equation}
Then
\begin{equation}
 \|QTQ\|\leq s_*                                   \label{eq:s7v88-QTQ}
\end{equation}
and the reconstructed Hamiltonian has gap at least
\begin{equation}
 -a_t^{-1}\log s_*.                                 \label{eq:s7v88-Schur-mass}
\end{equation}
The conclusion is independent of the number of flux sectors.
\end{theorem}

\begin{proof}
Put \(b_{\lambda\mu}:=\|T_{\lambda\mu}\|\).  Self-adjointness gives
\(b_{\lambda\mu}=b_{\mu\lambda}\), so both row and column sums of the
nonnegative scalar matrix \(b\) are at most \(s_*\).  The scalar Schur
test gives \(\|b\|_{\ell^2\to\ell^2}\leq s_*\).
For \(\psi=(\psi_\mu)\),
\[
 \|(T\psi)_\lambda\|
 \leq\sum_\mu b_{\lambda\mu}\|\psi_\mu\|.
\]
Taking the \(\ell^2\) norm proves
\(\|QTQ\|\leq\|b\|\leq s_*\).
\Cref{lem:s7v88-vacuum-block} separates the vacuum, and spectral
calculus gives \eqref{eq:s7v88-Schur-mass}.
\end{proof}

\begin{corollary}[Diagonal margin plus flux-changing row]
\label{cor:s7v88-diagonal-off}
If
\begin{align}
 \sup_\lambda\|T_{\lambda\lambda}\|&\leq r_{\rm d},
                                                               \label{eq:s7v88-rd}\\
 \sup_\lambda\sum_{\mu\ne\lambda}
 \|T_{\lambda\mu}\|&\leq\epsilon_{\rm flux},         \label{eq:s7v88-eps}
\end{align}
and
\begin{equation}
 r_{\rm d}+\epsilon_{\rm flux}<1,                   \label{eq:s7v88-margin}
\end{equation}
then the gap is at least
\(-a_t^{-1}\log(r_{\rm d}+\epsilon_{\rm flux})\).
\end{corollary}

\begin{proof}
The full row in \eqref{eq:s7v88-Schur-row} is at most
\(r_{\rm d}+\epsilon_{\rm flux}\).  Apply
\Cref{thm:s7v88-Schur}.
\end{proof}

\begin{remark}[Relation to normalized V80 sectors]
\label{rem:s7v88-V80}
V80 proves that normalized Gauss boundary pairing introduces no
\(d_\lambda\) factor in traces or convex sector mixtures.  The operator
criterion \eqref{eq:s7v88-Schur-row} is stronger: when plaquettes change
flux, it asks for an absolutely summable transition row.  Normalized
weights prevent artificial cardinality growth but do not prove that
row.
\end{remark}

\section{Why a global perturbation norm is the wrong loader}
\label{sec:s7v88-volume}

An interacting reflection-positive transfer often has the schematic
form
\begin{equation}
 T_V=M_{e^{-V/2}}T_EM_{e^{-V/2}},                   \label{eq:s7v88-TV}
\end{equation}
where \(V\) is a sum of spatial plaquette interactions on a slice.

\begin{proposition}[Volume-extensive norm obstruction]
\label{prop:s7v88-volume}
If \(V=\sum_{p\subset{\rm slice}}V_p\) with
\(\|V_p\|\leq v_*\), then the direct estimate
\begin{equation}
 \|V\|\leq v_*\,\#\{p\subset{\rm slice}\}            \label{eq:s7v88-volume-norm}
\end{equation}
grows with spatial volume.  Consequently a perturbation argument based
only on
\(\|T_V-T_E\|\) cannot yield a thermodynamic uniform channel margin
from local smallness.
\end{proposition}

\begin{proof}
The displayed inequality is the triangle bound.  Even if it is not
sharp, it is the only consequence of the individual norm bounds without
connected cancellation or locality.  Its right side diverges with
volume, so it cannot imply a fixed
\(\epsilon_{\rm flux}\) in \eqref{eq:s7v88-eps}.
\end{proof}

\begin{firewall}[Required interacting estimate]
\label{fire:s7v88-interacting}
The spatial interaction must be organized as a normalized connected
cluster or forest attached to a tested flux channel.  V79 ensures each
selected boundary edge occurs once and V80 removes sector-cardinality
loss, but neither theorem sums the complete off-diagonal transfer row
\eqref{eq:s7v88-eps}.  That is the next loader.
\end{firewall}

\section{Relation to the V82 terminal theorem}
\label{sec:s7v88-V82}

\begin{proposition}[V82 is a nonperturbative instance of channel
contraction]
\label{prop:s7v88-V82}
Under the full hypotheses of the V82 small-\(\beta\) terminal theorem,
the OS transfer operator satisfies
\begin{equation}
 \|QTQ\|\leq\alpha_\beta<1,                          \label{eq:s7v88-V82}
\end{equation}
where \(\alpha_\beta\) is its Dobrushin contraction constant.  Hence the
aggregate one-block decomposition
\(\Lambda_*=\{\mathrm{all\ nonvacuum\ channels}\}\) satisfies
\Cref{thm:s7v88-Schur} with \(s_*=\alpha_\beta\).
\end{proposition}

\begin{proof}
V82 proves reflected time-correlation decay
\(O(\alpha_\beta^n)\) on a dense gauge-invariant cylinder core.
The spectral-measure argument in \Cref{thm:s7v85-OS} places the
vacuum-orthogonal transfer spectrum in \([0,\alpha_\beta]\), giving
\eqref{eq:s7v88-V82}.  With one aggregate nonvacuum block, its operator
norm is the Schur row.
\end{proof}

\begin{remark}[Character and Dobrushin proofs are complementary]
\label{rem:s7v88-complementary}
The character theorem diagonalizes the ultralocal electric transfer and
makes Gauss compatibility explicit.  V82 controls the interacting
Euclidean measure by conditional mixing.  Neither derivation is needed
inside the other, but both land on the same gauge-invariant transfer
contraction.
\end{remark}

\begin{firewall}[Small terminal \(\beta\) is not the ultraviolet
trajectory]
\label{fire:s7v88-UV}
The continuum ultraviolet Yang--Mills coupling corresponds to large
Wilson inverse coupling, whereas
\eqref{eq:s7v88-SU3-domain} and V82 are strong-coupling terminal
statements.  A complete proof must show that an exact coarse marginal
enters such a channel-contraction domain, or establish the general
Schur margin directly.  Bare-coupling substitution would reverse the
renormalization direction.
\end{firewall}

\section{Flux-channel contraction card}
\label{sec:s7v88-card}

\begin{definition}[Flux-channel contraction card, FCCC]
\label{def:s7v88-FCCC}
FCCC consists of:
\begin{enumerate}[label=\textup{(F\arabic*)},leftmargin=4em]
\item a normalized positive self-adjoint transfer operator \(T\) on the
      physical Gauss-reduced slice Hilbert space, with
      \(T\Omega=\Omega\);
\item a complete paired boundary-flux decomposition of
      \(\Omega^\perp\), including the nonvacuum trivial-flux sector;
\item uniform diagonal and flux-changing bounds
      \eqref{eq:s7v88-rd}--\eqref{eq:s7v88-margin};
\item no hidden representation multiplicity, transition-chart, toron,
      or removed-cutoff loss in those rows;
\item a physical time step bounded above and below away from zero; and
\item the NTRGC/TDMC comparison to the ultraviolet state.
\end{enumerate}
\end{definition}

\begin{theorem}[FCCC gives a gauge-invariant mass gap]
\label{thm:s7v88-FCCC}
FCCC (F1)--(F5) gives a terminal gauge-invariant mass
\begin{equation}
 m_{\rm FCCC}
 \geq
 -a_t^{-1}\log(r_{\rm d}+\epsilon_{\rm flux})>0.     \label{eq:s7v88-FCCC-mass}
\end{equation}
With (F6), it loads the continuum ISMC handoff.
\end{theorem}

\begin{proof}
Rows (F1)--(F4) permit
\Cref{cor:s7v88-diagonal-off}.  Row (F5) converts the dimensionless
transfer ratio to the physical mass
\eqref{eq:s7v88-FCCC-mass}.  Row (F6) is the remaining state comparison.
\end{proof}

\section{Status and next acquisition}
\label{sec:s7v88-status}

\begin{theorem}[What V88 proves]
\label{thm:s7v88-result}
V88 proves a compact-group character-channel gap for every normalized
positive-definite central convolution with
\(\sup_{\lambda\ne\mathbf1}a_\lambda<1\), an explicit uniform
small-action ratio, volume-independent tensorization, and preservation
under Gauss projection.  It also proves the interacting flux-channel
Schur criterion
\[
 \sup_\lambda\sum_\mu\|T_{\lambda\mu}\|<1
\quad\Longrightarrow\quad
 m\geq-a_t^{-1}\log s_*.
\]
\end{theorem}

\begin{proof}
Use
\Cref{lem:s7v88-diagonal,thm:s7v88-one-link,
thm:s7v88-ratio,lem:s7v88-Wilson-positive,
thm:s7v88-tensor,cor:s7v88-Gauss,
thm:s7v88-Schur,thm:s7v88-FCCC}.
\end{proof}

\begin{assessment}[Progress at seventh-series V88]
\label{ass:s7v88-status}
The terminal mass mechanism is now stated on physical compact-group
flux channels and no longer resembles a forbidden connection mass.
The ultralocal theorem is complete.  For the interacting slow marginal,
the single missing quantitative row is a volume-uniform connected bound
on flux-changing transfer blocks.
\end{assessment}

\begin{decision}[V89 acquisition order]
\label{dec:s7v88-V89}
V89 will derive a local linked-cluster stability theorem for FCCC.
Starting with a diagonal character transfer having margin \(1-r_{\rm d}\),
it will organize spatial plaquette insertions by connected supports,
prove an operator-valued Schur row bound from rooted polymer activities,
and state the exact inequality which keeps
\(r_{\rm d}+\epsilon_{\rm flux}<1\).  It will explicitly exclude a
volume-global perturbation norm.
\end{decision}

\begin{firewall}[Claim boundary after seventh-series V88]
\label{fire:s7v88-claim}
V88 does not prove the interacting flux-changing row, FCCC for the
coarse Yang--Mills marginal, RG entry into the terminal character or
Dobrushin domain, CSMC, SEBE, MIRC, ARDPC, regulator-uniform SATC or
full RL4C, the nonlinear coarse-action cocycle, interacting endpoint
continuity, UCM, BCC, pairing/Bogoliubov control, regulator
independence, continuum OS reconstruction, or the full Yang--Mills mass
gap.
\end{firewall}

\paragraph{Parent-version record.}
V88 was formed from
\texttt{Renewal\_Yang\_Mills\_Mass\_Gap\_Research\_Notes\_V87.tex}.
The V87 parent had \(39\,223\) counted source lines,
\(1\,401\,589\) bytes, and SHA--256
\[
 \texttt{6DA0FA93F6FAA4013AEA002B1C084230695AD2447A735005FC3A3E60C34E4287}.
\]
The next compulsory companion audit is V90.

% =============================================================================
% END APPENDED SEVENTH-SERIES V88 MATERIAL
% =============================================================================

% =============================================================================
% BEGIN APPENDED SEVENTH-SERIES V89 MATERIAL
% =============================================================================

\chapter{Defect-anchored transfer polymers and a volume-free flux row}
\label{ch:s7v89}

\section{Purpose}
\label{sec:s7v89-purpose}

V88 reduced interacting terminal stability to the flux-channel Schur
row.  A direct norm of the spatial potential is volume extensive.  The
repair is to normalize by the exact vacuum and expand only connected
transfer decorations attached to an existing nonvacuum defect.  This
chapter proves the abstract combinatorial theorem behind that repair.

The theorem has a simple quantitative output.  If \(r_0<1\) is the
reference contraction per elementary flux defect and \(a\) is the
rooted channel-summed polymer activity, then the full nonvacuum transfer
row is at most
\[
 r_0e^a.
\]
Thus the terminal gap persists when \(r_0e^a<1\).  The proof counts
each connected decoration once, anchors it to a source defect, and
drops compatibility only after the anchor has removed the spatial
volume.

\section{Localized channel configurations}
\label{sec:s7v89-configurations}

Let \({\cal I}\) be a finite set of slice cells.  A channel
configuration \(\lambda\) has a finite defect support
\begin{equation}
 D(\lambda)\subseteq{\cal I},\qquad
 n(\lambda):=|D(\lambda)|.                           \label{eq:s7v89-defect}
\end{equation}
The vacuum is the unique configuration \(\boldsymbol0\) with
\(D(\boldsymbol0)=\varnothing\).  The channel Hilbert space is
\begin{equation}
 {\cal H}
 =
 \mathbb C\Omega
 \oplus
 \bigoplus_{\lambda\ne\boldsymbol0}{\cal M}_\lambda. \label{eq:s7v89-H}
\end{equation}
The multiplicity spaces include Gauss-paired intertwiners and the
nonvacuum part of the trivial cut-flux sector.

Let \(T_0\) be a diagonal positive transfer satisfying
\begin{equation}
 \|(T_0)_{\lambda\lambda}\|
 \leq r_0^{\,n(\lambda)},
 \qquad0<r_0<1.                                     \label{eq:s7v89-reference}
\end{equation}
For the ultralocal character transfer of V88, a defect cell is a link
carrying a nontrivial Peter--Weyl label and
\eqref{eq:s7v89-reference} follows by tensorization.

\begin{remark}[Closed networks still have a nonempty defect support]
\label{rem:s7v89-networks}
A Gauss-invariant spin network may have trivial total boundary charge,
but every nonconstant network uses nontrivial link representations.
Its local support \(D(\lambda)\) is therefore nonempty.  The vacuum is
not identified merely by a total representation label.
\end{remark}

\section{Channel-summed connected decorations}
\label{sec:s7v89-decorations}

A transfer polymer is a finite connected set
\(X\subseteq{\cal I}\).  For each incoming local channel label
\(\alpha\) on a collar of \(X\), let
\({\cal K}_X^{\alpha\to\beta}\) be its operator decoration into an
outgoing local label \(\beta\).  Define the channel-summed activity
\begin{equation}
 z(X):=
 \sup_\alpha
 \sum_\beta
 \|{\cal K}_X^{\alpha\to\beta}\|.                   \label{eq:s7v89-z}
\end{equation}
All representation multiplicities, transition labels, derivative
words, and internal bad-component contractions owned by \(X\) must be
included before \(z(X)\) is declared.

Fix a total ordering of \({\cal I}\).  If
\(X\cap D(\lambda)\ne\varnothing\), define its source anchor
\begin{equation}
 a_\lambda(X):=\min(X\cap D(\lambda)).               \label{eq:s7v89-anchor}
\end{equation}

\begin{definition}[Defect-anchored normalized transfer expansion,
DANTE]
\label{def:s7v89-DANTE}
DANTE means that, after exact ground normalization and removal of the
vacuum block, the channel blocks of \(T\) admit an absolutely convergent
expansion by compatible finite families \({\cal F}\) of connected
polymers such that:
\begin{enumerate}[label=\textup{(D\arabic*)},leftmargin=4em]
\item every \(X\in{\cal F}\) meets the incoming defect set
      \(D(\lambda)\);
\item every decorated contribution is assigned exactly once to its
      polymer family and local outgoing labels;
\item after summing all outgoing channel labels, its norm is at most
      \[
       r_0^{n(\lambda)}
       \prod_{X\in{\cal F}}z(X);
      \]
\item empty families give the diagonal reference term; and
\item the exact vacuum obeys \(T\Omega=\Omega\), so no unanchored
      vacuum polymer remains in \(QTQ\).
\end{enumerate}
\end{definition}

\begin{firewall}[Ground normalization is structural]
\label{fire:s7v89-normalization}
Condition (D1) is false for an unnormalized transfer numerator:
vacuum bubbles may be placed anywhere in the slice.  Dividing by a
separately bounded partition function does not establish their
cancellation inside an operator block.  DANTE requires a connected
logarithm, an exact vacuum intertwiner, or another construction which
removes unanchored components before norms are taken.
\end{firewall}

\section{The rooted activity}
\label{sec:s7v89-rooted}

Define
\begin{equation}
 a_*:=
 \sup_{x\in{\cal I}}
 \sum_{\substack{X\ {\rm connected}\\x\in X}}z(X).
                                                               \label{eq:s7v89-rooted}
\end{equation}
This unweighted form is enough for the channel row.  Spatially weighted
versions may be retained for simultaneous covariance estimates.

\begin{lemma}[Families anchored to a fixed defect set]
\label{lem:s7v89-families}
For every nonempty finite \(D\subseteq{\cal I}\),
\begin{equation}
 \sum_{\substack{{\cal F}\ {\rm finite,\ compatible}\\
                  X\cap D\ne\varnothing\ \forall X\in{\cal F}}}
 \prod_{X\in{\cal F}}z(X)
 \leq e^{a_*|D|}.                                   \label{eq:s7v89-family-bound}
\end{equation}
\end{lemma}

\begin{proof}
Assign every polymer \(X\) to the first point of \(X\cap D\) in the
fixed ordering.  Dropping compatibility and allowing repeated choices
only enlarges the nonnegative sum.  For one anchor \(x\), the sum over
an arbitrary finite set of assigned polymers is bounded by
\[
 \prod_{\substack{X\ {\rm connected}\\a_D(X)=x}}
 (1+z(X))
 \leq
 \exp\!\left(
  \sum_{\substack{X\ {\rm connected}\\a_D(X)=x}}z(X)
 \right)
 \leq e^{a_*}.
\]
Multiply this estimate over the \(|D|\) anchors.  Every original
compatible family has one unique assignment and is counted once in the
enlarged product.
\end{proof}

\begin{remark}[No spatial volume factor]
\label{rem:s7v89-no-volume}
The last product is over source defects, not over every slice cell.
This is exactly where DANTE (D1) removes the volume divergence found in
V88.
\end{remark}

\section{The channel-row theorem}
\label{sec:s7v89-row}

\begin{theorem}[Defect-anchored flux-row bound]
\label{thm:s7v89-row}
Assume \eqref{eq:s7v89-reference}, DANTE, and
\(a_*<\infty\).  Then for every nonvacuum incoming channel
\(\lambda\),
\begin{equation}
 \sum_{\mu\ne\boldsymbol0}\|T_{\lambda\mu}\|
 \leq
 \bigl(r_0e^{a_*}\bigr)^{n(\lambda)}.                \label{eq:s7v89-row}
\end{equation}
If
\begin{equation}
 s_{\rm DA}:=r_0e^{a_*}<1,                           \label{eq:s7v89-margin}
\end{equation}
then
\begin{equation}
 \sup_{\lambda\ne\boldsymbol0}
 \sum_{\mu\ne\boldsymbol0}\|T_{\lambda\mu}\|
 \leq s_{\rm DA}<1.                                 \label{eq:s7v89-uniform-row}
\end{equation}
\end{theorem}

\begin{proof}
Fix \(\lambda\).  DANTE (D2)--(D3), absolute convergence, and the
triangle inequality permit the outgoing channel sum to be performed
inside the polymer majorant.  DANTE (D1) and
\Cref{lem:s7v89-families} then give
\[
 \sum_\mu\|T_{\lambda\mu}\|
 \leq
 r_0^{n(\lambda)}e^{a_*n(\lambda)},
\]
which is \eqref{eq:s7v89-row}.  Every nonvacuum channel has
\(n(\lambda)\geq1\).  If \(s_{\rm DA}<1\), then
\(s_{\rm DA}^{n(\lambda)}\leq s_{\rm DA}\), proving
\eqref{eq:s7v89-uniform-row}.
\end{proof}

\begin{corollary}[Interacting transfer gap from rooted activity]
\label{cor:s7v89-gap}
If \(T\) is positive and self-adjoint, DANTE holds, and
\eqref{eq:s7v89-margin} holds, then
\begin{equation}
 \|QTQ\|\leq r_0e^{a_*}                             \label{eq:s7v89-QTQ}
\end{equation}
and the physical transfer Hamiltonian has gap at least
\begin{equation}
 m_{\rm DA}
 :=
 -a_t^{-1}\log(r_0e^{a_*})>0.                       \label{eq:s7v89-mass}
\end{equation}
\end{corollary}

\begin{proof}
Use \eqref{eq:s7v89-uniform-row} in the flux-channel Schur theorem
\Cref{thm:s7v88-Schur}.
\end{proof}

\begin{corollary}[Explicit stability budget]
\label{cor:s7v89-budget}
Let \(\delta_0:=-\log r_0>0\).  DANTE preserves a positive gap whenever
\begin{equation}
 a_*<\delta_0,                                      \label{eq:s7v89-budget}
\end{equation}
and then
\begin{equation}
 m_{\rm DA}\geq a_t^{-1}(\delta_0-a_*).              \label{eq:s7v89-linear-gap}
\end{equation}
\end{corollary}

\begin{proof}
The margin \(r_0e^{a_*}<1\) is equivalent to
\(a_*<-\log r_0\).  Take the logarithm in
\eqref{eq:s7v89-mass}.
\end{proof}

\section{Representation branching}
\label{sec:s7v89-branching}

The definition \eqref{eq:s7v89-z} requires a channel sum.  For
\(SU(3)\), fixed fundamental Wilson words have uniformly finite
branching.

\begin{lemma}[\(SU(3)\) fundamental branching]
\label{lem:s7v89-SU3}
Label irreducibles by highest weights \((p,q)\), \(p,q\geq0\).  Tensoring
once with the fundamental or antifundamental representation produces
at most three irreducible summands:
\begin{align}
 (p,q)\otimes(1,0)
 &=
 (p+1,q)\oplus(p-1,q+1)\oplus(p,q-1),               \label{eq:s7v89-fund}\\
 (p,q)\otimes(0,1)
 &=
 (p,q+1)\oplus(p+1,q-1)\oplus(p-1,q),               \label{eq:s7v89-antifund}
\end{align}
with terms having a negative label omitted.  Hence a word of \(k\)
fundamental or antifundamental multipliers has at most \(3^k\) outgoing
highest-weight branches from any incoming \((p,q)\).
\end{lemma}

\begin{proof}
The two displayed decompositions are the \(SU(3)\) Pieri rules.
At each multiplication there are at most three choices.  Iteration
gives at most \(3^k\) paths; merging equal final irreducibles only
reduces their number.
\end{proof}

\begin{corollary}[Fixed Wilson words fit the local channel sum]
\label{cor:s7v89-fixed-word}
If a polymer \(X\) contains \(k(X)\) fundamental/antifundamental
matrix multipliers and every resolved branch operator has norm at most
\(u(X)\), then
\begin{equation}
 z(X)\leq3^{k(X)}u(X).                               \label{eq:s7v89-branch-bound}
\end{equation}
If \(k(X)\leq c_{\rm w}|X|\), the branching loss is only exponential per
polymer block and may be absorbed into its activity.
\end{corollary}

\begin{proof}
Use \Cref{lem:s7v89-SU3} and sum the branch norms.  The second statement
uses \(3^{k(X)}\leq(3^{c_{\rm w}})^{|X|}\).
\end{proof}

\begin{remark}[Unresolved V79 bounds are stronger]
\label{rem:s7v89-unresolved}
When a flux-changing boundary insertion is bounded before resolving
intermediate representation labels, as in V79--V80, its operator norm
already controls the complete local channel sum and no explicit
\(3^{k(X)}\) loss is needed.  The branching lemma is a fallback for a
resolved construction.
\end{remark}

\section{A rooted activity from per-block polymer decay}
\label{sec:s7v89-activity}

Assume the connected-set count
\eqref{eq:s7v86-connected-count}.  Suppose
\begin{equation}
 z(X)\leq \zeta^{|X|}.                               \label{eq:s7v89-zeta}
\end{equation}

\begin{lemma}[Explicit rooted sum]
\label{lem:s7v89-root-sum}
If \(K_{\rm conn}\zeta<1\), then
\begin{equation}
 a_*
 \leq
 {\zeta\over1-K_{\rm conn}\zeta}.                   \label{eq:s7v89-a-bound}
\end{equation}
\end{lemma}

\begin{proof}
For a fixed root \(x\), sum
\eqref{eq:s7v89-zeta} over connected sets of size \(n\) and use
\eqref{eq:s7v86-connected-count}:
\[
 \sum_{X\ni x}z(X)
 \leq
 \sum_{n\geq1}K_{\rm conn}^{n-1}\zeta^n
 =
 {\zeta\over1-K_{\rm conn}\zeta}.
\]
Take the supremum over \(x\).
\end{proof}

\begin{corollary}[Checkable linked-cluster margin]
\label{cor:s7v89-checkable}
Under DANTE and \eqref{eq:s7v89-zeta}, a sufficient gap condition is
\begin{equation}
 {\zeta\over1-K_{\rm conn}\zeta}<-\log r_0.          \label{eq:s7v89-checkable}
\end{equation}
\end{corollary}

\begin{proof}
Combine \Cref{lem:s7v89-root-sum,cor:s7v89-budget}.
\end{proof}

\section{Good, bad, and transition contributions}
\label{sec:s7v89-ledger}

Suppose the complete connected polymer activity has been partitioned by
ownership into
\begin{equation}
 z(X)
 \leq
 z_{\rm good}(X)+z_{\rm bad}(X)+z_{\rm tr}(X),       \label{eq:s7v89-split}
\end{equation}
and define their rooted sums
\(a_{\rm good},a_{\rm bad},a_{\rm tr}\).

\begin{proposition}[Additive rooted ledger]
\label{prop:s7v89-ledger}
\begin{equation}
 a_*
 \leq
 a_{\rm good}+a_{\rm bad}+a_{\rm tr}.               \label{eq:s7v89-ledger}
\end{equation}
If the right side tends to zero while a terminal reference has a fixed
\(r_0<1\), then \eqref{eq:s7v89-budget} holds eventually.
\end{proposition}

\begin{proof}
Sum \eqref{eq:s7v89-split} over polymers containing a fixed root and
take the supremum.  A quantity tending to zero is eventually smaller
than the fixed number \(-\log r_0\).
\end{proof}

\begin{proposition}[Conditional loading by the existing activities]
\label{prop:s7v89-existing}
Assume a complete normalized DANTE construction identifies:
\begin{enumerate}[label=\textup{(\roman*)},leftmargin=3.5em]
\item \(a_{\rm good}\) with a fixed polynomial multiple of the V64
      replica activity and V78 good-form activity;
\item \(a_{\rm bad}\) with a fixed geometric multiple of the V79 marked
      saturated activity \(\widetilde b_M\); and
\item \(a_{\rm tr}\) with the fully enumerated V72--V76 transition and
      chart rows under the V74/V75 mixed norm.
\end{enumerate}
If each identified row tends to zero, then
\begin{equation}
 a_*\longrightarrow0,                               \label{eq:s7v89-a-zero}
\end{equation}
and every fixed terminal character margin \(r_0<1\) is stable.
\end{proposition}

\begin{proof}
Under the stated identifications, each of the three rooted terms in
\eqref{eq:s7v89-ledger} tends to zero by the cited activity estimates.
Apply \Cref{prop:s7v89-ledger,cor:s7v89-budget}.
\end{proof}

\begin{firewall}[The conditional clause is the remaining theorem]
\label{fire:s7v89-existing}
The earlier notes prove the numerical bounds for declared local rows.
They do not construct the exact ground-normalized operator expansion
DANTE or prove that its complete polymer list is exhausted by
\eqref{eq:s7v89-split}.  Therefore
\Cref{prop:s7v89-existing} is a precise loader, not a proof that the
physical Yang--Mills transfer already has the gap.
\end{firewall}

\section{Spectator stability and the defect exponent}
\label{sec:s7v89-spectator}

\begin{theorem}[Why the defect exponent prevents spectator growth]
\label{thm:s7v89-spectator}
Under \eqref{eq:s7v89-row}, a channel with \(n\) defects has row at most
\(s_{\rm DA}^n\).  Hence adding spectator defects decreases the row
when \(s_{\rm DA}<1\), and the worst nonvacuum row occurs at
\(n=1\).
\end{theorem}

\begin{proof}
This is the monotonicity
\[
 s_{\rm DA}^n\leq s_{\rm DA},\qquad n\geq1,
\]
when \(0<s_{\rm DA}<1\).
\end{proof}

\begin{firewall}[An additive occurrence bound would fail]
\label{fire:s7v89-additive}
If the polymer estimate produced \(na_*\) as an additive operator error
to a fixed \(r_0\), its supremum over unrestricted defect number would
diverge, exactly as in the V78 spectator counterexample.  DANTE
factorizes the correction per defect and combines it with the
reference factor \(r_0^n\); this multiplicative structure is essential.
\end{firewall}

\section{The defect-anchored flux card}
\label{sec:s7v89-card}

\begin{definition}[Defect-anchored flux stability card, DAFSC]
\label{def:s7v89-DAFSC}
DAFSC consists of:
\begin{enumerate}[label=\textup{(A\arabic*)},leftmargin=4em]
\item a physical positive self-adjoint transfer with exact normalized
      vacuum;
\item localized complete channel configurations with
      \(n(\lambda)\geq1\) off the vacuum;
\item the reference defect contraction
      \eqref{eq:s7v89-reference};
\item a complete absolutely convergent DANTE expansion;
\item a regulator-uniform rooted channel-summed activity satisfying
      \(a_*<-\log r_0\);
\item Gauss, transition, toron, and reflection compatibility; and
\item the NTRGC/TDMC ultraviolet comparison.
\end{enumerate}
\end{definition}

\begin{theorem}[DAFSC loads FCCC]
\label{thm:s7v89-DAFSC}
DAFSC (A1)--(A6) implies FCCC with
\[
 r_{\rm d}+\epsilon_{\rm flux}
 \leq r_0e^{a_*}<1
\]
and terminal mass at least
\[
 a_t^{-1}(-\log r_0-a_*).
\]
With (A7), it loads the continuum ISMC handoff.
\end{theorem}

\begin{proof}
Use
\Cref{thm:s7v89-row,cor:s7v89-gap,
cor:s7v89-budget,thm:s7v88-FCCC}.
\end{proof}

\section{Status and next acquisition}
\label{sec:s7v89-status}

\begin{theorem}[What V89 proves]
\label{thm:s7v89-result}
V89 proves that an exact defect-anchored normalized transfer expansion
with rooted channel-summed activity \(a_*\) changes a reference
per-defect contraction \(r_0\) to at most \(r_0e^{a_*}\).  The bound is
uniform in volume, representation count, and spectator defect number.
It also gives the checkable polymer condition
\[
 {\zeta\over1-K_{\rm conn}\zeta}<-\log r_0.
\]
\end{theorem}

\begin{proof}
Use
\Cref{lem:s7v89-families,thm:s7v89-row,
cor:s7v89-gap,lem:s7v89-SU3,
lem:s7v89-root-sum,thm:s7v89-spectator,
thm:s7v89-DAFSC}.
\end{proof}

\begin{assessment}[Progress at seventh-series V89]
\label{ass:s7v89-status}
The volume-free operator combinatorics required by V88 is now proved
abstractly.  The remaining issue is constructive rather than algebraic:
build DANTE for the exact ground-normalized coarse Yang--Mills transfer
and verify that every connected decoration is included with the rooted
activities already estimated.
\end{assessment}

\begin{decision}[V90 audit and acquisition order]
\label{dec:s7v89-V90}
V90 is the compulsory companion audit.  It will inspect the current SM
and NS files for an exact normalized linked-transfer, vacuum-bubble
cancellation, or defect-anchoring mechanism.  It will then attack the
DANTE construction at finite cutoff, beginning with the algebraic
connected logarithm and the distinction between a normalized state
polymer and an operator transfer polymer.
\end{decision}

\begin{firewall}[Claim boundary after seventh-series V89]
\label{fire:s7v89-claim}
V89 does not construct DANTE, verify DAFSC for the physical coarse
transfer, prove RG entry into the terminal character/Dobrushin domain,
close CSMC, SEBE, MIRC, or ARDPC, establish regulator-uniform SATC or
full RL4C, construct the nonlinear coarse-action cocycle, prove
interacting endpoint continuity, UCM, BCC, pairing/Bogoliubov control,
regulator independence, continuum OS reconstruction, or the full
Yang--Mills mass gap.
\end{firewall}

\paragraph{Parent-version record.}
V89 was formed from
\texttt{Renewal\_Yang\_Mills\_Mass\_Gap\_Research\_Notes\_V88.tex}.
The V88 parent had \(39\,825\) counted source lines,
\(1\,422\,522\) bytes, and SHA--256
\[
 \texttt{199735C32FC1B8D91352B4C13C7556F087F3B276A76B080ACFEC02E4AFB17B73}.
\]
The next compulsory companion audit is V90.

% =============================================================================
% END APPENDED SEVENTH-SERIES V89 MATERIAL
% =============================================================================

% =============================================================================
% BEGIN APPENDED SEVENTH-SERIES V90 MATERIAL
% =============================================================================

\chapter{V90 companion audit and boundary-source connected transfer}
\label{ch:s7v90}

\section{Purpose}
\label{sec:s7v90-purpose}

V89 proved the volume-free summation of a defect-anchored operator
expansion but left its construction open.  The compulsory audit finds a
weaker route in the companion Standard Model notes: an analytic
connected logarithm with sources cancels all clusters disjoint from the
tested observables.  In a reflection-positive Euclidean theory,
exponential decay of those source derivatives already determines the
vacuum-orthogonal transfer spectrum.  An operator-norm polymer
expansion is therefore optional rather than compulsory.

This chapter records the audit, proves the exact ground-state transform
for the operator route, proves the boundary-source connected theorem,
and gives a checkable rooted cluster norm.  The physical missing step is
then sharply stated: construct this source expansion for the exact
coarse Yang--Mills marginal with all good, bad, transition, and Gauss
rows included.

\section{Compulsory companion audit}
\label{sec:s7v90-audit}

\subsection{Files inspected}

The active companion files inspected read-only were
\begin{align}
 &\texttt{Renewal\_SM\_Einstein\_Continuum\_Research\_Notes\_V56.tex},
 \label{eq:s7v90-SM-file}\\
 &\texttt{Renewal\_NS\_Stability\_Research\_Notes\_V26.tex}.
 \label{eq:s7v90-NS-file}
\end{align}
The SM file had \(20\,742\) counted source lines, \(745\,315\) bytes,
and SHA--256
\[
 \texttt{82123934C5E9DF581B5C6F6629249673E28852825A69DFA8AB640BEAE7627B6B}.
\]
The NS file had \(5\,320\) counted source lines, \(278\,283\) bytes,
and SHA--256
\[
 \texttt{82C887F8C2C69E4F28FAD7C3DC8DE5B9099CB5A4BF431EBA1FFF3E91935978AD}.
\]
The line convention includes the final empty line when present.

\subsection{Transfer decision}

SM V50--V52 develops an averaged regular/defect disagreement theorem
for a positive conditional Higgs law.  Its logic agrees with
Yang--Mills V84--V86, but its quartic coercivity and exterior-amplitude
conditions do not provide a compact-gauge estimate.

SM V54 proves that, when an analytic connected source logarithm
converges absolutely, two source derivatives retain only clusters
meeting both observable supports.  This is an algebraic normalization
mechanism of exactly the type needed after V89.  It is scalar rather
than operator-valued, but reflection positivity makes scalar
time-correlation decay sufficient for a spectral gap.

SM V55 gives the same slow--fast covariance firewall as V86, and SM V56
gives the same expected-Hessian-minus-covariance formula as V87.  These
are independent confirmations, not new Yang--Mills estimates.  NS V26
contains no normalized Gibbs, connected-source, compact-character, or
transfer input.

\begin{theorem}[V90 audit decision]
\label{thm:s7v90-audit}
The programme retains DANTE as a sufficient operator route but no
longer treats it as necessary.  The preferred next loader is a
boundary-source connected logarithm for the exact reflection-positive
coarse Yang--Mills measure.  No numerical activity or physical
Yang--Mills bound is imported from the companion files.
\end{theorem}

\begin{proof}
The source-logarithm cancellation in SM V54 depends only on analytic
differentiation and connected support, so its abstract algebra is
type-compatible.  Its complex fermion phase estimates and all SM
V50--V56 field-specific constants are not.  The spectral sufficiency of
the scalar route is proved below in
\Cref{thm:s7v90-source-gap}.
\end{proof}

\section{Exact ground-state normalization}
\label{sec:s7v90-Doob}

The operator route has an exact normalization even when the interacting
vacuum is not the constant function.  Let \((\Omega_0,\nu)\) be a
configuration space and let \(T\) be a positivity-preserving
self-adjoint integral operator on \(L^2(\nu)\) with kernel
\(K(x,y)=K(y,x)\geq0\).  Assume
\begin{equation}
 T\psi_0=\lambda_0\psi_0,\qquad
 \lambda_0>0,\qquad \psi_0(x)>0\quad\nu\text{-a.e.} \label{eq:s7v90-ground}
\end{equation}
and normalize \(\|\psi_0\|_{L^2(\nu)}=1\).

Define
\begin{align}
 d\pi(x)&:=\psi_0(x)^2\,d\nu(x),                    \label{eq:s7v90-pi}\\
 P(x,dy)&:=
 {K(x,y)\psi_0(y)\over\lambda_0\psi_0(x)}\,d\nu(y).
                                                               \label{eq:s7v90-P}
\end{align}

\begin{theorem}[Ground-state Doob transform]
\label{thm:s7v90-Doob}
The kernel \(P\) is Markov and reversible with invariant probability
\(\pi\).  The map
\begin{equation}
 U:L^2(\pi)\to L^2(\nu),\qquad Uf=\psi_0f           \label{eq:s7v90-U}
\end{equation}
is unitary and
\begin{equation}
 UPU^{-1}={T\over\lambda_0}.                         \label{eq:s7v90-equivalence}
\end{equation}
Consequently the nontrivial spectral contraction of \(P\) is exactly
the transfer ratio \(\lambda_1/\lambda_0\).
\end{theorem}

\begin{proof}
By \eqref{eq:s7v90-ground},
\[
 \int P(x,dy)
 ={(T\psi_0)(x)\over\lambda_0\psi_0(x)}=1.
\]
Furthermore,
\[
 \pi(dx)P(x,dy)
 ={1\over\lambda_0}
 \psi_0(x)K(x,y)\psi_0(y)\,d\nu(x)d\nu(y),
\]
which is symmetric in \(x,y\); this proves reversibility and invariance.
The normalization of \(\psi_0\) makes \(\pi\) a probability, and
\(\|Uf\|_{L^2(\nu)}=\|f\|_{L^2(\pi)}\).
Finally,
\[
 (UPf)(x)
 ={1\over\lambda_0}
 \int K(x,y)\psi_0(y)f(y)\,d\nu(y)
 ={1\over\lambda_0}(TUf)(x).
\]
This proves \eqref{eq:s7v90-equivalence}; unitary equivalence gives the
spectral statement.
\end{proof}

\begin{firewall}[The ground-state ratio may be nonlocal]
\label{fire:s7v90-ground-locality}
Equation \eqref{eq:s7v90-P} cancels vacuum normalization exactly, but a
local transfer kernel does not by itself prove that
\(\psi_0(y)/\psi_0(x)\) has a quasi-local polymer expansion.  Assuming
such locality in order to prove the gap can be circular.  The
boundary-source route below avoids constructing \(\psi_0\) explicitly.
\end{firewall}

\section{A scalar pressure does not determine the gap}
\label{sec:s7v90-pressure}

\begin{proposition}[Equal vacuum data, different subleading spectrum]
\label{prop:s7v90-pressure-no}
There are positive self-adjoint contractions with the same top
eigenvalue and the same trace but arbitrarily different
vacuum-orthogonal norms.
\end{proposition}

\begin{proof}
On \(\mathbb C^3\), fix \(0<c<1\) and compare
\[
 T_r=\operatorname{diag}(1,r,c-r),
 \qquad
 0\leq r\leq c,
\]
in the subrange where both nonvacuum entries are nonnegative.  Every
\(T_r\) has top eigenvalue one and trace \(1+c\), while its
vacuum-orthogonal norm is
\(\max\{r,c-r\}\), which varies with \(r\).  Taking a higher-dimensional
version with fixed sum permits one subleading eigenvalue to approach
one while redistributing the remaining trace.
\end{proof}

\begin{firewall}[Source-free pressure is insufficient]
\label{fire:s7v90-pressure}
Cancellation of vacuum bubbles in \(\log Z(0)\) controls a pressure or
top eigenvalue.  A mass gap requires source derivatives separated in
time, an operator channel bound, or equivalent information about the
subleading spectrum.  A source-free cluster expansion alone does not
load FCCC.
\end{firewall}

\section{Boundary sources on a Euclidean slab}
\label{sec:s7v90-sources}

Let \(\mu_{\Lambda,n}\) be a normalized finite-volume Euclidean
Yang--Mills measure on a slab of \(n\) block-time steps, with a
reflection plane and physical boundary conditions compatible with the
transfer construction.  Let \(F\) and \(G\) be bounded gauge-invariant
cylinder observables supported in fixed collars \(A_-\) and \(A_+\) of
the lower and upper slab boundaries.  Put
\begin{equation}
 Z_{\Lambda,n}(J,K)
 :=
 \mathbb E_{\mu_{\Lambda,n}}
 \exp(JF+KG),                                       \label{eq:s7v90-ZJK}
\end{equation}
and choose
\begin{equation}
 W_{\Lambda,n}(J,K):=\log Z_{\Lambda,n}(J,K)         \label{eq:s7v90-WJK}
\end{equation}
on a zero-free polydisc about \((0,0)\).

\begin{lemma}[Boundary-source derivatives]
\label{lem:s7v90-source-derivatives}
At \(J=K=0\),
\begin{align}
 \partial_JW_{\Lambda,n}&=\mathbb E[F],              \label{eq:s7v90-first}\\
 \partial_KW_{\Lambda,n}&=\mathbb E[G],              \label{eq:s7v90-first-K}\\
 \partial_J\partial_KW_{\Lambda,n}
 &=\operatorname{Cov}_{\mu_{\Lambda,n}}(F,G).        \label{eq:s7v90-second}
\end{align}
\end{lemma}

\begin{proof}
Differentiate \(\log Z\) once and twice.  Since
\(Z_{\Lambda,n}(0,0)=1\), the mixed second derivative is
\(\mathbb E[FG]-\mathbb E[F]\mathbb E[G]\).
\end{proof}

\section{Connected source expansion}
\label{sec:s7v90-connected}

Let \({\cal X}_{\Lambda,n}\) be the connected finite subsets of the
slab block graph.  Assume
\begin{equation}
 W_{\Lambda,n}(J,K)
 =
 \sum_{X\in{\cal X}_{\Lambda,n}}w_X(J,K)             \label{eq:s7v90-cluster}
\end{equation}
absolutely and locally uniformly in a source polydisc, with the support
rules
\begin{align}
 \partial_Jw_X&=0\quad\text{if }X\cap A_-=\varnothing,
                                                               \label{eq:s7v90-support-J}\\
 \partial_Kw_X&=0\quad\text{if }X\cap A_+=\varnothing.
                                                               \label{eq:s7v90-support-K}
\end{align}
Suppose
\begin{equation}
 |\partial_J\partial_Kw_X(0,0)|
 \leq
 c_X\operatorname{osc}F\operatorname{osc}G          \label{eq:s7v90-cX}
\end{equation}
and, for some \(\mu>0\),
\begin{equation}
 C_\mu:=
 \sup_z\sum_{X\ni z}
 e^{\mu\operatorname{diam}X}c_X<\infty              \label{eq:s7v90-Cmu}
\end{equation}
uniformly in slab size, spatial volume, slow/transition labels, and
removed auxiliary cutoffs.

\begin{theorem}[Boundary-source connected decay]
\label{thm:s7v90-connected}
Under \eqref{eq:s7v90-cluster}--\eqref{eq:s7v90-Cmu},
\begin{equation}
 |\operatorname{Cov}_{\mu_{\Lambda,n}}(F,G)|
 \leq
 C_\mu|A_-|
 \operatorname{osc}F\operatorname{osc}G\,
 e^{-\mu d(A_-,A_+)}.                               \label{eq:s7v90-connected-decay}
\end{equation}
Every source-free vacuum cluster cancels before the absolute bound.
\end{theorem}

\begin{proof}
Local uniform absolute convergence permits two source derivatives of
\eqref{eq:s7v90-cluster}.  By
\eqref{eq:s7v90-support-J}--\eqref{eq:s7v90-support-K}, a surviving
connected \(X\) meets both boundary collars, so
\(\operatorname{diam}X\geq d(A_-,A_+)\).  Therefore
\begin{align*}
 |\operatorname{Cov}(F,G)|
 &\leq
 \operatorname{osc}F\operatorname{osc}G
 \sum_{\substack{X:X\cap A_-\ne\varnothing\\
                   X\cap A_+\ne\varnothing}}c_X\\
 &\leq
 e^{-\mu d(A_-,A_+)}
 \operatorname{osc}F\operatorname{osc}G
 \sum_{z\in A_-}
 \sum_{X\ni z}e^{\mu\operatorname{diam}X}c_X,
\end{align*}
which is \eqref{eq:s7v90-connected-decay}.
Terms \(w_X(0,0)\) vanish after either source derivative.
\end{proof}

\begin{remark}[Why this is weaker than DANTE]
\label{rem:s7v90-weaker}
DANTE bounds the full vacuum-orthogonal transfer operator by a channel
row.  \Cref{thm:s7v90-connected} bounds matrix elements generated by a
dense local cylinder algebra.  Reflection positivity and the spectral
theorem are needed to pass from the latter to the full
vacuum-orthogonal spectrum.
\end{remark}

\section{From boundary clusters to a spectral gap}
\label{sec:s7v90-gap}

\begin{theorem}[Boundary-source transfer gap]
\label{thm:s7v90-source-gap}
Assume:
\begin{enumerate}[label=\textup{(\roman*)},leftmargin=3.5em]
\item the connected bound
      \eqref{eq:s7v90-connected-decay} holds uniformly and passes to the
      infinite-volume/cofinal state;
\item the state is reflection positive and time translation gives a
      positive self-adjoint contraction \(T\);
\item gauge-invariant positive-time cylinder classes are dense in the
      OS Hilbert space; and
\item one block-time step has physical size \(a_t>0\).
\end{enumerate}
Then the OS Hamiltonian has vacuum-orthogonal gap at least
\begin{equation}
 m_{\rm BS}
 \geq{\mu\over a_t}.                                \label{eq:s7v90-mass}
\end{equation}
\end{theorem}

\begin{proof}
For a mean-zero cylinder vector \([F]\), place its reflected copy at the
lower boundary and its translate after \(n\) steps at the upper
boundary.  Their collar separation is \(n-O_F(1)\).  Reflection
positivity identifies the covariance with
\(\langle[F],T^n[F]\rangle\), so
\[
 0\leq\langle[F],T^n[F]\rangle
 \leq C_Fe^{-\mu n}.
\]
The spectral-measure argument of
\Cref{thm:s7v85-OS} shows that the spectral measure of every such vector
is supported in \([0,e^{-\mu}]\).  Density extends this to the full
vacuum-orthogonal OS space.  Thus
\(\|QTQ\|\leq e^{-\mu}\), and
\(-a_t^{-1}\log T\) has gap at least \(\mu/a_t\).
\end{proof}

\begin{corollary}[Scalar connected expansion loads FCCC]
\label{cor:s7v90-FCCC}
Under \Cref{thm:s7v90-source-gap}, FCCC holds in the aggregate
nonvacuum channel with
\begin{equation}
 s_*\leq e^{-\mu}.                                  \label{eq:s7v90-sstar}
\end{equation}
\end{corollary}

\begin{proof}
The proof of \Cref{thm:s7v90-source-gap} gives
\(\|QTQ\|\leq e^{-\mu}\).  Use one block for the whole nonvacuum space
in \Cref{thm:s7v88-Schur}.
\end{proof}

\section{A checkable rooted coefficient norm}
\label{sec:s7v90-rooted}

Assume the slab graph has the connected-set bound
\begin{equation}
 \#\{X\ni z:X\text{ connected},\ |X|=m\}
 \leq K_{\rm slab}^{m-1}.                           \label{eq:s7v90-count}
\end{equation}
Every connected \(m\)-block set has
\(\operatorname{diam}X\leq m-1\).

\begin{lemma}[Per-block activity implies a connected source norm]
\label{lem:s7v90-rooted}
If
\begin{equation}
 c_X\leq c_0\zeta^{|X|},                            \label{eq:s7v90-zeta}
\end{equation}
then
\begin{equation}
 C_\mu
 \leq
 {c_0\zeta\over
  1-K_{\rm slab}e^\mu\zeta}                         \label{eq:s7v90-Cmu-bound}
\end{equation}
whenever
\begin{equation}
 K_{\rm slab}e^\mu\zeta<1.                          \label{eq:s7v90-root-margin}
\end{equation}
\end{lemma}

\begin{proof}
For fixed \(z\), sum over cardinality:
\begin{align*}
 \sum_{X\ni z}e^{\mu\operatorname{diam}X}c_X
 &\leq
 c_0\sum_{m\geq1}
 K_{\rm slab}^{m-1}e^{\mu(m-1)}\zeta^m\\
 &=
 {c_0\zeta\over1-K_{\rm slab}e^\mu\zeta}.
\end{align*}
Take the supremum in \(z\).
\end{proof}

\begin{corollary}[Explicit decay exponent]
\label{cor:s7v90-exponent}
If \(K_{\rm slab}\zeta<1\), every
\begin{equation}
 0<\mu<-\log(K_{\rm slab}\zeta)                     \label{eq:s7v90-mu-domain}
\end{equation}
is an admissible exponential rate in
\Cref{lem:s7v90-rooted}, subject to the construction of
\eqref{eq:s7v90-cluster} and its source derivative bound.
\end{corollary}

\begin{proof}
The displayed interval is exactly
\(K_{\rm slab}e^\mu\zeta<1\).
\end{proof}

\section{Mixed good/bad source activity}
\label{sec:s7v90-mixed}

Let a selected connected source cluster mark its regular blocks, its
saturated bad components, and its transition/atlas collars.  Suppose
the complete absolute coefficient satisfies
\begin{equation}
 c_X
 \leq
 c_0
 \bigl(\zeta_{\rm good}
       +\zeta_{\rm bad}
       +\zeta_{\rm tr}\bigr)^{|X|}.                 \label{eq:s7v90-mixed-zeta}
\end{equation}
Put
\begin{equation}
 \zeta_{\rm tot}
 :=
 \zeta_{\rm good}+\zeta_{\rm bad}+\zeta_{\rm tr}.
                                                               \label{eq:s7v90-total-zeta}
\end{equation}

\begin{proposition}[Source-cluster terminal criterion]
\label{prop:s7v90-terminal}
If the exact connected expansion obeys
\eqref{eq:s7v90-mixed-zeta} and
\begin{equation}
 K_{\rm slab}\zeta_{\rm tot}<1,                     \label{eq:s7v90-terminal}
\end{equation}
then every
\[
 0<\mu<-\log(K_{\rm slab}\zeta_{\rm tot})
\]
gives exponential boundary-source covariance decay.  Under the OS
hypotheses, it gives mass at least \(\mu/a_t\).
\end{proposition}

\begin{proof}
Apply
\Cref{lem:s7v90-rooted,cor:s7v90-exponent,
thm:s7v90-connected,thm:s7v90-source-gap}
with \(\zeta=\zeta_{\rm tot}\).
\end{proof}

\begin{firewall}[A binomial majorant must follow the exact expansion]
\label{fire:s7v90-mixed}
Equation \eqref{eq:s7v90-mixed-zeta} is not obtained by merely adding
three previously proved numbers.  Every physical connected source
coefficient must first be decomposed with unique ownership, correct
normalization, and no omitted operation type.  V76 supplies the
operation ledger, while V73--V75 and V78--V80 supply candidate local
bounds; the global source expansion joining them remains to be built.
\end{firewall}

\section{Comparison of the two terminal routes}
\label{sec:s7v90-routes}

\begin{theorem}[DANTE and boundary sources are alternative sufficient
loaders]
\label{thm:s7v90-alternatives}
At a reflection-positive terminal scale, either of the following proves
a mass gap:
\begin{enumerate}[label=\textup{(\alph*)},leftmargin=3.5em]
\item DANTE with \(a_*<-\log r_0\), giving
      \(m\geq a_t^{-1}(-\log r_0-a_*)\);
\item the connected boundary-source expansion with
      \(K_{\rm slab}\zeta_{\rm tot}<1\), giving every
      \(m< a_t^{-1}[-\log(K_{\rm slab}\zeta_{\rm tot})]\).
\end{enumerate}
Neither implication requires the other expansion.
\end{theorem}

\begin{proof}
Part (a) is \Cref{cor:s7v89-gap,cor:s7v89-budget}.
Part (b) is \Cref{prop:s7v90-terminal}.  The first controls an operator
row directly; the second controls a dense family of reflected matrix
elements and uses spectral reconstruction.
\end{proof}

\begin{decision}[Preferred route after the V90 audit]
\label{dec:s7v90-route}
The boundary-source route is preferred for the next iteration because
the existing Yang--Mills machinery is Euclidean, replica-normalized,
and organized by connected forests.  DANTE remains available if an
exact quasi-local ground-state transform is later constructed.
\end{decision}

\section{Boundary-source connected transfer card}
\label{sec:s7v90-card}

\begin{definition}[Boundary-source connected transfer card, BSCTC]
\label{def:s7v90-BSCTC}
BSCTC consists of:
\begin{enumerate}[label=\textup{(B\arabic*)},leftmargin=4em]
\item exact normalized reflection-positive slab measures for the
      coarse physical Yang--Mills state;
\item a zero-free source polydisc and the locally uniformly absolutely
      convergent connected logarithm
      \eqref{eq:s7v90-cluster};
\item exact source-support rules
      \eqref{eq:s7v90-support-J}--\eqref{eq:s7v90-support-K};
\item a regulator-uniform connected derivative norm
      \eqref{eq:s7v90-Cmu} for some \(\mu>0\);
\item Gauss, boundary-flux, toron, transition-chart, and removed-cutoff
      compatibility;
\item a dense gauge-invariant OS cylinder core and physical block-time
      spacing; and
\item the NTRGC/TDMC comparison with the ultraviolet state.
\end{enumerate}
\end{definition}

\begin{theorem}[BSCTC gives the terminal and continuum gap handoff]
\label{thm:s7v90-BSCTC}
BSCTC (B1)--(B6) gives a terminal gap at least \(\mu/a_t\).  With (B7),
it loads the continuum ISMC handoff.
\end{theorem}

\begin{proof}
Rows (B1)--(B4) permit
\Cref{thm:s7v90-connected}; rows (B5)--(B6) permit
\Cref{thm:s7v90-source-gap}.  Row (B7) is the ultraviolet comparison.
\end{proof}

\section{Status and next acquisition}
\label{sec:s7v90-status}

\begin{theorem}[What V90 proves]
\label{thm:s7v90-result}
V90 proves the exact reversible ground-state transform of a positive
transfer and a separate boundary-source connected theorem.  A
locally uniformly convergent source logarithm with rooted derivative
activity \(\zeta\) has exponential rate
\[
 0<\mu<-\log(K_{\rm slab}\zeta),
\]
and reflection positivity converts that decay into an OS mass
\(\mu/a_t\).  Thus exact source-connected normalization is a sufficient
alternative to the stronger operator expansion DANTE.
\end{theorem}

\begin{proof}
Use
\Cref{thm:s7v90-Doob,thm:s7v90-connected,
thm:s7v90-source-gap,lem:s7v90-rooted,
prop:s7v90-terminal,thm:s7v90-alternatives,
thm:s7v90-BSCTC}.
\end{proof}

\begin{assessment}[Progress at seventh-series V90]
\label{ass:s7v90-status}
The vacuum-normalization bottleneck has been split into two rigorous
options.  The source route matches the existing Euclidean forest
architecture and avoids explicit construction of the interacting
ground vector.  The remaining task is to derive the exact connected
source expansion for the complete coarse action and load its rooted
coefficient with the already separated good, bad, and transition rows.
\end{assessment}

\begin{decision}[V91 acquisition order]
\label{dec:s7v90-V91}
V91 will build the finite-cutoff boundary-source expansion by coupling
interpolation.  It will differentiate plaquette, electric/form, atlas,
and source couplings, use a connected forest selection to cancel
source-disjoint components, and determine whether noncommuting
Hamiltonian form rows admit the same scalar Euclidean interpolation as
bounded multiplier rows.  If not, it will isolate the exact
Trotter/Duhamel loader rather than merging the two types.
\end{decision}

\begin{firewall}[Claim boundary after seventh-series V90]
\label{fire:s7v90-claim}
V90 does not construct BSCTC (B2)--(B7) for the physical coarse
Yang--Mills measure, prove DANTE, verify the complete mixed coefficient
\eqref{eq:s7v90-mixed-zeta}, establish RG entry into a terminal domain,
close CSMC, MIRC, or ARDPC, prove regulator-uniform SATC or full RL4C,
construct the nonlinear coarse-action cocycle, prove interacting
endpoint continuity, UCM, BCC, pairing/Bogoliubov control, regulator
independence, continuum OS reconstruction, or the full Yang--Mills mass
gap.
\end{firewall}

\paragraph{Parent-version record.}
V90 was formed from
\texttt{Renewal\_Yang\_Mills\_Mass\_Gap\_Research\_Notes\_V89.tex}.
The V89 parent had \(40\,381\) counted source lines,
\(1\,441\,870\) bytes, and SHA--256
\[
 \texttt{EC8550D0098222AD3EB20BE2D62114BC366DBFCE4BA79C69081C8F7029FAC44E}.
\]
The next compulsory companion audit is V95.

% =============================================================================
% END APPENDED SEVENTH-SERIES V90 MATERIAL
% =============================================================================

% =============================================================================
% BEGIN APPENDED SEVENTH-SERIES V91 MATERIAL
% =============================================================================

\chapter{Compact-link Euclideanization and the bounded-multiplier
source expansion}
\label{ch:s7v91}

\section{Purpose}
\label{sec:s7v91-purpose}

V90 chose the boundary-source route and asked whether the electric/form
rows require the same scalar coupling interpolation as Wilson
multipliers.  In the exact compact-link Hamiltonian they do not.  The
electric Laplacian generates the positive heat-kernel reference
measure; the spatial Wilson potential appears as a bounded scalar
Feynman--Kac weight.  Chart metrics, IMS terms, and projector responses
are introduced only when this exact object is represented in local
oscillator coordinates.

This chapter proves the finite-regulator Trotter/Feynman--Kac
scalarization, derives the exact coupling-cumulant identities, and
constructs the bounded-multiplier polymer expansion.  It gives a
complete BSCTC construction in a small-activity heat-kernel terminal
domain.  The result is a second terminal proof, not an ultraviolet RG
entry theorem.

\section{The exact compact-link transfer}
\label{sec:s7v91-transfer}

Let \(E\) be the finite spatial-link set and
\[
 {\cal H}_E=L^2(G^E,dU),
 \qquad G=SU(3),
\]
with normalized product Haar measure.  Define
\begin{align}
 H_{\rm el}
 &:=
 {c_Eg^2\over2a_s}
 \sum_{e\in E}(-\Delta_e),                           \label{eq:s7v91-Hel}\\
 V_{\rm W}(U)
 &:=
 {c_B\over a_sg^2}
 \sum_{p\subset{\rm slice}}
 \bigl(3-\operatorname{ReTr}W_p(U)\bigr).           \label{eq:s7v91-VW}
\end{align}
At finite volume, \(H_{\rm el}\geq0\) is self-adjoint and
\(V_{\rm W}\) is a bounded nonnegative multiplier.

Let \(p_t^E(U,V)\) be the product heat kernel of
\(H_{\rm el}\).  It is positive, symmetric, and satisfies
\begin{equation}
 \int p_t^E(U,V)\,dV=1.                              \label{eq:s7v91-Markov}
\end{equation}

\begin{theorem}[Trotter compact-link path formula]
\label{thm:s7v91-Trotter}
For \(t>0\),
\begin{equation}
 e^{-t(H_{\rm el}+V_{\rm W})}
 =
 \operatorname*{s-lim}_{N\to\infty}
 \left(
  e^{-tH_{\rm el}/N}e^{-tV_{\rm W}/N}
 \right)^N.                                         \label{eq:s7v91-Trotter}
\end{equation}
For bounded \(\Psi_-,\Psi_+\), its matrix element is
\begin{align}
 \langle\Psi_-,
 e^{-t(H_{\rm el}+V_{\rm W})}\Psi_+\rangle
 &=
 \lim_{N\to\infty}
 \int
 \overline{\Psi_-(U_0)}\Psi_+(U_N)
 \notag\\
 &\quad\times
 \prod_{j=1}^N
 \left[
  p_{t/N}^E(U_{j-1},U_j)
  e^{-(t/N)V_{\rm W}(U_j)}
 \right]
 \prod_{j=0}^NdU_j.                                 \label{eq:s7v91-path}
\end{align}
\end{theorem}

\begin{proof}
The bounded perturbation \(V_{\rm W}\) and the self-adjoint
semibounded \(H_{\rm el}\) satisfy the Lie--Trotter product theorem,
giving \eqref{eq:s7v91-Trotter}.  Insert the heat kernel of each
electric semigroup and the diagonal kernel of each multiplier into the
finite product.  Fubini's theorem gives the \(N\)-slice integral.
Strong convergence and bounded endpoint multipliers give
\eqref{eq:s7v91-path}.
\end{proof}

\begin{corollary}[Physical rows become scalar weights]
\label{cor:s7v91-scalar}
In \eqref{eq:s7v91-path}, the electric operator is entirely contained
in the positive Markov reference
\(\prod_jp_{t/N}^E\), while every spatial Wilson interaction and
boundary observable is a scalar multiplication factor.
\end{corollary}

\begin{proof}
This is the displayed factorization in
\eqref{eq:s7v91-path}.
\end{proof}

\begin{firewall}[Heat-kernel and Wilson temporal actions]
\label{fire:s7v91-regularization}
The exact time-sliced transfer naturally produces a heat-kernel
temporal action.  Replacing it by the standard Wilson temporal
plaquette action is a change of regulator, not an algebraic identity.
A final continuum theorem must prove regulator independence or carry
one chosen reflection-positive discretization through the complete
construction.
\end{firewall}

\section{Coordinate rows are not new physical interactions}
\label{sec:s7v91-coordinate}

Let \(\{\chi_\alpha\}\) be an exact smooth partition of unity on the
compact link configuration space, and let each chart represent the
electric Dirichlet form in local coordinates.  The variable metric,
coordinate density, IMS debit, and projector response of V72--V76 arise
from this representation.

\begin{proposition}[Exact atlas recombination]
\label{prop:s7v91-atlas}
If the chart decomposition is exact on the form domain, then summing all
chart contributions reconstructs the single physical form of
\(H_{\rm el}\).  Consequently those chart rows must be used to justify
the local oscillator expansion and its cutoff removal, but they are not
additional scalar interactions in the compact-link path measure
\eqref{eq:s7v91-path}.
\end{proposition}

\begin{proof}
The IMS identity is an equality:
\[
 H_{\rm el}
 =
 \sum_\alpha\chi_\alpha H_{\rm el}\chi_\alpha
 -\sum_\alpha|\nabla\chi_\alpha|^2
\]
at the quadratic-form level, with the coordinate metric and density
being the pullbacks of the same Haar Laplacian.  Projector response
records differentiation of a chosen moving coordinate subspace.
Adding these terms according to the exact identity returns
\(H_{\rm el}\); counting them again in
\eqref{eq:s7v91-path} would double-count a coordinate
representation.
\end{proof}

\begin{firewall}[The coordinate proof is still needed]
\label{fire:s7v91-coordinate}
Removing coordinate rows from the physical scalar interaction list does
not prove their regulator estimates.  They remain essential to compare
the exact compact transfer with the truncated oscillator construction,
to control chart overlaps, and to justify cofinal limits.  The
distinction is ownership, not dispensability.
\end{firewall}

\section{Exact coupling derivatives}
\label{sec:s7v91-couplings}

On a finite time-sliced compact path space, let \(\nu_0\) be the
normalized product heat-kernel reference.  Let
\(\{\Phi_b\}_{b\in{\cal B}}\) be bounded local scalar interactions,
including time-sliced spatial Wilson plaquettes and any already-proved
bounded coarse loop term.  Put
\begin{equation}
 Z(\boldsymbol\lambda;J,K)
 :=
 \mathbb E_{\nu_0}
 \exp\!\left[
  -\sum_{b\in{\cal B}}\lambda_b\Phi_b
  +JF+KG
 \right],                                           \label{eq:s7v91-Z}
\end{equation}
and \(W=\log Z\) on a zero-free polydisc.

\begin{theorem}[Coupling derivatives are connected cumulants]
\label{thm:s7v91-cumulants}
For distinct \(b_1,\ldots,b_m\),
\begin{equation}
 \partial_{\lambda_{b_1}}\cdots
 \partial_{\lambda_{b_m}}
 \partial_J\partial_KW
 =
 (-1)^m
 \kappa_{\boldsymbol\lambda,J,K}
 (\Phi_{b_1},\ldots,\Phi_{b_m},F,G),                 \label{eq:s7v91-cumulant}
\end{equation}
where \(\kappa\) is the joint cumulant in the normalized tilted law.
\end{theorem}

\begin{proof}
For any parameter \(\theta\) coupled linearly to an observable \(A\),
\(\partial_\theta\log Z=\mathbb E_\theta A\).
Repeated differentiation subtracts every partition product of lower
moments, which is the moment--cumulant recursion.  Each
\(\lambda_b\)-derivative inserts \(-\Phi_b\), producing the sign.
\end{proof}

\begin{corollary}[Source-disconnected terms vanish algebraically]
\label{cor:s7v91-source-connected}
If the tilted reference factorizes into two independent variable
families, one containing \(F\) and the other containing \(G\), then the
cumulant in \eqref{eq:s7v91-cumulant} is zero.
\end{corollary}

\begin{proof}
A joint cumulant vanishes when its arguments split into two nonempty
independent subfamilies.  This follows directly by factorizing the
moment generating function; its logarithm is the sum of the two
independent logarithms and has no mixed derivative.
\end{proof}

\section{Bounded interaction factors}
\label{sec:s7v91-factors}

For each local interaction put
\begin{equation}
 R_b:=e^{-\Phi_b}-1,\qquad
 \rho_b:=\|R_b\|_\infty.                             \label{eq:s7v91-Rb}
\end{equation}
Then
\begin{equation}
 \rho_b\leq e^{\|\Phi_b\|_\infty}-1.                \label{eq:s7v91-rho}
\end{equation}
The interacting density expands exactly as
\begin{equation}
 \prod_{b\in{\cal B}}e^{-\Phi_b}
 =
 \sum_{A\subseteq{\cal B}}\prod_{b\in A}R_b.         \label{eq:s7v91-subsets}
\end{equation}

Join two interaction labels when their variable supports overlap in the
heat-kernel reference factorization.  Every finite \(A\) decomposes
uniquely into connected components
\(\gamma_1,\ldots,\gamma_k\).  Define the connected polymer activity
\begin{equation}
 z(\gamma)
 :=
 \mathbb E_{\nu_0}\prod_{b\in\gamma}R_b.             \label{eq:s7v91-polymer}
\end{equation}

\begin{lemma}[Polymer activity bound]
\label{lem:s7v91-polymer}
If disjoint support components are independent under \(\nu_0\), then
\begin{align}
 \mathbb E_{\nu_0}\prod_{b\in A}R_b
 &=
 \prod_{j=1}^kz(\gamma_j),                           \label{eq:s7v91-factorization}\\
 |z(\gamma)|
 &\leq\prod_{b\in\gamma}\rho_b.                     \label{eq:s7v91-z-bound}
\end{align}
\end{lemma}

\begin{proof}
Independence gives \eqref{eq:s7v91-factorization}.  The expectation of
an absolute product is at most its supremum norm, giving
\eqref{eq:s7v91-z-bound}.
\end{proof}

\begin{remark}[Heat-kernel dependency supports are columns]
\label{rem:s7v91-columns}
The electric reference couples a fixed spatial link through Euclidean
time, so its independent factors are link histories rather than
individual spacetime edges.  The support graph used above must join
plaquettes sharing one such history.  Incorrectly treating distinct
times on the same link as independent would undercount polymers.
\end{remark}

\section{A convergent hard-core polymer logarithm}
\label{sec:s7v91-log}

Let polymers be incompatible when their dependency supports overlap.
The partition ratio is the hard-core gas
\begin{equation}
 \Xi
 =
 \sum_{\substack{\Gamma\ {\rm finite}\\
                  \text{pairwise compatible}}}
 \prod_{\gamma\in\Gamma}z(\gamma).                  \label{eq:s7v91-Xi}
\end{equation}

\begin{theorem}[Kotecky--Preiss source domain]
\label{thm:s7v91-KP}
Suppose there is a nonnegative polymer size \(a(\gamma)\) such that
\begin{equation}
 \sum_{\gamma'\not\sim\gamma}
 |z(\gamma')|e^{a(\gamma')}
 \leq a(\gamma)                                     \label{eq:s7v91-KP}
\end{equation}
for every polymer \(\gamma\).  Then \(\Xi\ne0\), its logarithm has the
absolutely convergent Ursell expansion
\begin{equation}
 \log\Xi
 =
 \sum_{n\geq1}{1\over n!}
 \sum_{\gamma_1,\ldots,\gamma_n}
 \varphi^T(\gamma_1,\ldots,\gamma_n)
 \prod_{i=1}^nz(\gamma_i),                          \label{eq:s7v91-Ursell}
\end{equation}
and the rooted bound
\begin{equation}
 \sum_{n\geq1}{1\over(n-1)!}
 \sum_{\gamma_2,\ldots,\gamma_n}
 |\varphi^T(\gamma,\gamma_2,\ldots,\gamma_n)|
 \prod_{i=2}^n|z(\gamma_i)|
 \leq e^{a(\gamma)}.                                \label{eq:s7v91-rooted}
\end{equation}
\end{theorem}

\begin{proof}
Expand the hard-core compatibility factors
\[
 \prod_{i<j}\bigl(1+\zeta_{ij}\bigr),
 \qquad
 \zeta_{ij}=
 \begin{cases}
 -1,&\gamma_i\not\sim\gamma_j,\\
 0,&\gamma_i\sim\gamma_j.
 \end{cases}
\]
Moment--cumulant inversion selects connected incompatibility graphs and
gives \eqref{eq:s7v91-Ursell}.  The tree-graph bound replaces their
absolute sum by spanning trees.  Root a tree at \(\gamma\), sum its
leaves inward, and apply \eqref{eq:s7v91-KP} at each parent.  The
exponential weights pay the unordered children through
\(\sum_{k\geq0}a^k/k!=e^a\), yielding
\eqref{eq:s7v91-rooted}.  The same rooted majorant converges uniformly
in volume and in a sufficiently small activity polydisc; analytic
continuation from zero shows that \(\Xi\) has no zero there.
\end{proof}

\begin{firewall}[The KP theorem needs the correct reference
factorization]
\label{fire:s7v91-KP}
The activity \eqref{eq:s7v91-polymer} and incompatibility relation must
be built from the exact heat-kernel path measure, including Gauss
projection or an exact gauge-reduced representation.  A polymer gas
formed after deleting shared link histories does not represent the
physical partition function.
\end{firewall}

\section{Boundary sources inside the polymer gas}
\label{sec:s7v91-source}

Introduce two marked source polymers \(s_-\) and \(s_+\), with factors
\[
 e^{JF}-1,\qquad e^{KG}-1,
\]
and dependency supports in the two boundary collars.  Apply
\Cref{thm:s7v91-KP} in a small source polydisc.

\begin{theorem}[Connected boundary-source construction]
\label{thm:s7v91-source}
Under the KP condition uniformly in slab size and physical labels,
\(\log Z(J,K)\) has a locally uniformly absolutely convergent cluster
expansion.  Its mixed derivative at the origin is a sum only over
Ursell clusters containing both marked source polymers.  If the rooted
two-source coefficients satisfy
\begin{equation}
 \sup_z\sum_{X\ni z}
 e^{\mu\operatorname{diam}X}c_X<\infty,             \label{eq:s7v91-source-norm}
\end{equation}
then BSCTC (B2)--(B4) holds.
\end{theorem}

\begin{proof}
The source factors are bounded and analytic in \(J,K\), so the KP
margin persists in a sufficiently small polydisc.  Differentiate the
locally uniform Ursell series.  A term without \(s_-\) is independent of
\(J\), and a term without \(s_+\) is independent of \(K\).  Hence a
mixed derivative contains both.  Group the connected union of their
dependency supports into \(X\); the assumed weighted rooted bound is
exactly \eqref{eq:s7v90-Cmu}, and source support is automatic.
\end{proof}

\begin{corollary}[Uniform bounded plaquette criterion]
\label{cor:s7v91-plaquette}
Suppose every elementary interaction has
\(\rho_b\leq\rho\), each dependency cell meets at most \(D_{\rm int}\)
interaction labels, and the number of connected \(m\)-label polymers
containing a fixed label is at most \(K_{\rm int}^{m-1}\).  Then a
sufficient small-activity regime for KP and a positive source decay
rate has the form
\begin{equation}
 C_{\rm KP}K_{\rm int}D_{\rm int}\rho<1,             \label{eq:s7v91-small}
\end{equation}
where \(C_{\rm KP}\) is a fixed tree/size-weight constant.  For a
Wilson multiplier \(\Phi_b=\beta\phi_b\) with
\(\|\phi_b\|_\infty\leq B_\phi\),
\begin{equation}
 \rho\leq e^{\beta B_\phi}-1=O(\beta).               \label{eq:s7v91-Wilson-rho}
\end{equation}
\end{corollary}

\begin{proof}
\Cref{lem:s7v91-polymer} gives
\(|z(\gamma)|\leq\rho^{|\gamma|}\).  Sum incompatible rooted polymers
by their first overlapping label and the connected-set count.  Choosing
a linear size weight in \eqref{eq:s7v91-KP} gives a convergent geometric
series under a fixed inequality of the displayed form.  The final bound
is \eqref{eq:s7v91-rho}.
\end{proof}

\begin{remark}[Agreement with V82]
\label{rem:s7v91-V82}
Both \eqref{eq:s7v91-small} and the V82 Dobrushin criterion hold for
sufficiently small terminal Wilson coupling.  Their constants differ
because one sums connected polymers and the other sums conditional
influences.  Either yields a legitimate reflection-positive terminal
mass theorem.
\end{remark}

\section{Bounded multipliers versus form rows}
\label{sec:s7v91-types}

\begin{theorem}[Type separation for BSCTC]
\label{thm:s7v91-types}
For the exact compact Kogut--Susskind transfer:
\begin{enumerate}[label=\textup{(\roman*)},leftmargin=3.5em]
\item \(H_{\rm el}\) belongs to the heat-kernel reference and is not
      differentiated as a polymer activity;
\item spatial Wilson and bounded generated loop terms enter the scalar
      factors \eqref{eq:s7v91-Rb};
\item endpoint gauge-invariant observables enter as marked source
      factors;
\item chart, IMS, moving-projector, and oscillator form rows belong to
      the comparison proof between local coordinates and the exact
      compact transfer; and
\item a genuine modification of the electric generator requires a
      separate Duhamel/form interpolation and cannot be inserted into
      \eqref{eq:s7v91-subsets} as a scalar multiplier.
\end{enumerate}
\end{theorem}

\begin{proof}
Items (i)--(iii) are
\Cref{thm:s7v91-Trotter,cor:s7v91-scalar,
thm:s7v91-source}.  Item (iv) is
\Cref{prop:s7v91-atlas}.  A differential or quadratic-form operator is
not pointwise multiplication on path configurations, proving (v).
\end{proof}

\begin{firewall}[Generated electric operators]
\label{fire:s7v91-electric}
If the RG map changes the electric generator rather than merely
integrating the original compact-link measure, the heat-kernel
reference becomes scale dependent.  BSCTC then needs uniform heat-kernel
comparison or the V77 relative-form resolvent/Duhamel route.  V91 removes
coordinate artifacts from the physical vertex list; it does not forbid
genuine generated kinetic terms.
\end{firewall}

\section{A finite-cutoff BSCTC theorem}
\label{sec:s7v91-card}

\begin{theorem}[Small-activity heat-kernel BSCTC]
\label{thm:s7v91-BSCTC}
Assume at a fixed terminal scale:
\begin{enumerate}[label=\textup{(\alph*)},leftmargin=3.5em]
\item the exact compact transfer has the heat-kernel path representation
      \eqref{eq:s7v91-path};
\item all nonreference physical interactions are bounded scalar local
      factors with a uniform KP margin;
\item the source-marked rooted norm
      \eqref{eq:s7v91-source-norm} holds for some \(\mu>0\);
\item Gauss reduction and transition labels preserve the reference
      factorization or have an exact normalized sector
      disintegration; and
\item reflection positivity and the cofinal limiting rows of BSCTC
      hold.
\end{enumerate}
Then the terminal OS Hamiltonian has mass at least \(\mu/a_t\).
\end{theorem}

\begin{proof}
\Cref{thm:s7v91-Trotter,thm:s7v91-KP,
thm:s7v91-source} construct BSCTC (B1)--(B4).
Item (d) supplies its physical label row and item (e) its remaining
OS/cofinal rows.  Apply \Cref{thm:s7v90-BSCTC}.
\end{proof}

\begin{firewall}[Terminal closure is not trajectory entry]
\label{fire:s7v91-entry}
At the ultraviolet weak-coupling regulator,
\(\beta\) is large and
\eqref{eq:s7v91-Wilson-rho} is not small.  The theorem becomes useful
only after an exact coarse marginal has entered a small complete
interaction norm, or if a different nonperturbative reference absorbs
the large local Wilson term.  V91 does not prove that RG statement.
\end{firewall}

\section{Status and next acquisition}
\label{sec:s7v91-status}

\begin{theorem}[What V91 proves]
\label{thm:s7v91-result}
V91 proves the exact compact-link Trotter path representation, places
the electric Laplacian in the positive heat-kernel reference, and
constructs the connected boundary-source logarithm for bounded local
multiplier interactions under a KP margin.  With uniform Gauss,
reflection, and cofinal rows this gives a terminal OS mass
\(\mu/a_t\).  It also proves that chart/electric form rows used in the
oscillator comparison must not be double-counted as physical scalar
vertices.
\end{theorem}

\begin{proof}
Use
\Cref{thm:s7v91-Trotter,prop:s7v91-atlas,
thm:s7v91-cumulants,lem:s7v91-polymer,
thm:s7v91-KP,thm:s7v91-source,
thm:s7v91-types,thm:s7v91-BSCTC}.
\end{proof}

\begin{assessment}[Progress at seventh-series V91]
\label{ass:s7v91-status}
The finite-cutoff scalar source expansion is now constructed for the
correct compact-link types.  This removes an unnecessary attempt to
place coordinate form debits inside the physical polymer gas.  The
principal unresolved step remains the same hard one: produce the exact
coarse marginal and show that its complete bounded interaction tail
enters the KP/BSCTC domain uniformly.
\end{assessment}

\begin{decision}[V92 acquisition order]
\label{dec:s7v91-V92}
V92 will attack one exact coarse step.  For a reflection-equivariant
block map it will disintegrate the fine compact-link path measure,
write the coarse density as a conditional partition function, and
derive its connected logarithm by a source-dependent polymer
pushforward.  The target is a quantitative recursion for the complete
coarse KP norm, including the feedback denominator identified in V84.
\end{decision}

\begin{firewall}[Claim boundary after seventh-series V91]
\label{fire:s7v91-claim}
V91 does not prove ultraviolet entry into the small-activity terminal
domain, a regulator-uniform coarse KP recursion, the complete nonlinear
coarse-action cocycle, BSCTC through all cutoff removals, CSMC, MIRC, or
ARDPC for the physical trajectory, regulator-uniform SATC or full
RL4C, interacting endpoint continuity, UCM, BCC,
pairing/Bogoliubov control, regulator independence, continuum OS
reconstruction, or the full Yang--Mills mass gap.
\end{firewall}

\paragraph{Parent-version record.}
V91 was formed from
\texttt{Renewal\_Yang\_Mills\_Mass\_Gap\_Research\_Notes\_V90.tex}.
The V90 parent had \(41\,000\) counted source lines,
\(1\,463\,967\) bytes, and SHA--256
\[
 \texttt{37C9816724025B300F2517DB348DA2095BE0C49C7169E92BD41F1EF5D7511F29}.
\]
The next compulsory companion audit is V95.

% =============================================================================
% END APPENDED SEVENTH-SERIES V91 MATERIAL
% =============================================================================

% =============================================================================
% BEGIN APPENDED SEVENTH-SERIES V92 MATERIAL
% =============================================================================

\chapter{One exact coarse pushforward and its connected norm recursion}
\label{ch:s7v92}

\section{Purpose}
\label{sec:s7v92-purpose}

V91 constructed the terminal connected-source expansion for bounded
compact-link interactions.  The missing dynamical step is to show that
an exact coarse marginal inherits a controlled complete interaction.
This chapter carries out the measure-theoretic and combinatorial part of
one such step.

The exact coarse density is a conditional partition function.  Under a
locally uniformly absolutely convergent conditional cluster expansion,
clusters may be regrouped by their retained footprint to give the
complete coarse logarithmic interaction.  A typed retained/eliminated
tree majorant gives the Schur-path term
\[
 \rho+{rs\over1-\epsilon}.
\]
The V84 feedback denominator does not belong to the logarithmic
interaction itself; it appears only when that interaction is converted
to worst-exterior conditional influence.  Keeping these two recursions
separate removes a possible double count.

\section{Exact disintegration under a block map}
\label{sec:s7v92-disintegration}

Let \((\Omega,\nu_0)\) be a finite-regulator compact-link path space and
let
\begin{equation}
 {\cal B}:\Omega\to\overline\Omega                 \label{eq:s7v92-B}
\end{equation}
be a measurable gauge-covariant block map.  Disintegrate
\begin{equation}
 \nu_0(d\omega)
 =
 \nu_0^y(d\omega)\,\overline\nu_0(dy),
 \qquad y={\cal B}\omega.                            \label{eq:s7v92-disintegration}
\end{equation}
Let the fine physical measure be
\begin{equation}
 \mu(d\omega)
 =
 Z^{-1}e^{-\Phi(\omega)}\,\nu_0(d\omega),            \label{eq:s7v92-fine}
\end{equation}
where \(\Phi\) denotes the complete fine interaction.

Define the conditional partition function
\begin{equation}
 {\cal Z}(y)
 :=
 \int e^{-\Phi(\omega)}\,\nu_0^y(d\omega).           \label{eq:s7v92-Zy}
\end{equation}

\begin{theorem}[Exact coarse density]
\label{thm:s7v92-density}
The pushforward \(\overline\mu={\cal B}_*\mu\) satisfies
\begin{equation}
 {d\overline\mu\over d\overline\nu_0}(y)
 =
 {{\cal Z}(y)\over Z}.                               \label{eq:s7v92-density}
\end{equation}
Equivalently, with
\begin{equation}
 \overline\Phi(y):=-\log{\cal Z}(y),                 \label{eq:s7v92-Phi-bar}
\end{equation}
the coarse measure is proportional to
\(e^{-\overline\Phi}\overline\nu_0\).  Additive constants in
\(\overline\Phi\) are immaterial.
\end{theorem}

\begin{proof}
For bounded measurable \(F\) on \(\overline\Omega\),
\begin{align*}
 \int F(y)\,\overline\mu(dy)
 &=
 {1\over Z}\int F({\cal B}\omega)
 e^{-\Phi(\omega)}\,\nu_0(d\omega)\\
 &=
 {1\over Z}\int F(y)
 \left[
  \int e^{-\Phi(\omega)}\,\nu_0^y(d\omega)
 \right]\overline\nu_0(dy).
\end{align*}
This is \eqref{eq:s7v92-density}.  Taking the negative logarithm gives
\eqref{eq:s7v92-Phi-bar}.
\end{proof}

\begin{corollary}[Reflection positivity survives the exact step]
\label{cor:s7v92-RP}
If the fine measure is reflection positive and
\({\cal B}\Theta=\overline\Theta{\cal B}\), with positive-side coarse
observables pulling back to positive-side fine observables, then
\(\overline\mu\) is reflection positive.
\end{corollary}

\begin{proof}
This is the exact equivariant pushforward theorem
\Cref{thm:s7v83-RP} applied to \({\cal B}\).
\end{proof}

\section{Conditional connected logarithm}
\label{sec:s7v92-log}

Let \({\cal C}\) denote finite connected fine block clusters.  Assume
that on a common zero-free domain the conditional logarithm has the
locally uniformly absolutely convergent expansion
\begin{equation}
 \log{\cal Z}(y)
 =
 c_\varnothing+
 \sum_{C\in{\cal C}}w_C(y).                         \label{eq:s7v92-log}
\end{equation}
Let \(R(C)\) be the nonempty set of retained coarse blocks on which
\(w_C\) depends.  Terms with empty retained footprint are included in
the constant after conditional integration.

For every finite nonempty coarse set \(Y\), define
\begin{equation}
 \Psi_Y(y_Y)
 :=
 -\sum_{\substack{C\in{\cal C}\\R(C)=Y}}w_C(y_Y).
                                                               \label{eq:s7v92-PsiY}
\end{equation}

\begin{theorem}[Exact footprint regrouping]
\label{thm:s7v92-regroup}
If \eqref{eq:s7v92-log} converges absolutely and locally uniformly, then
\begin{equation}
 \overline\Phi(y)
 =
 \overline c+\sum_{\varnothing\ne Y}\Psi_Y(y_Y)      \label{eq:s7v92-coarse-action}
\end{equation}
with locally uniform absolute convergence.  No generated support is
omitted.
\end{theorem}

\begin{proof}
Use \eqref{eq:s7v92-Phi-bar} and
\eqref{eq:s7v92-log}.  Absolute convergence permits regrouping the
terms by the unique retained footprint \(R(C)\).  The empty-footprint
sum changes only the additive constant.  Every nonempty cluster occurs
in exactly one \(\Psi_Y\), proving completeness.
\end{proof}

\begin{firewall}[A formal conditional logarithm is insufficient]
\label{fire:s7v92-log}
The pointwise identity
\(\overline\Phi=-\log{\cal Z}\) does not imply
\eqref{eq:s7v92-log}.  One needs a zero-free conditional partition
function and a cluster expansion uniform in the retained field,
boundary flux, transition chart, volume, and removed fast cutoffs.
Otherwise differentiation and footprint regrouping may fail.
\end{firewall}

\section{The exact coarse oscillation row}
\label{sec:s7v92-oscillation}

For a function \(h(y_Y)\), define its mixed two-block oscillation
\begin{equation}
 \partial_{ij}h
 :=
 \sup_{\substack{y,y':\,y_{(ij)^c}=y'_{(ij)^c}}}
 |h(y_i,y_j)-h(y_i',y_j)
  -h(y_i,y_j')+h(y_i',y_j')|.                       \label{eq:s7v92-mixed}
\end{equation}
The precise supremum may equivalently be taken over four configurations
agreeing off \(i,j\).

Define the coarse linear interaction row
\begin{equation}
 {\cal O}_\mu(\overline\Phi)
 :=
 \sup_i\sum_{Y\ni i}
 (|Y|-1)e^{\mu\operatorname{diam}Y}
 \operatorname{osc}\Psi_Y.                          \label{eq:s7v92-O}
\end{equation}

\begin{lemma}[Cluster-to-coarse norm]
\label{lem:s7v92-cluster-norm}
Under \eqref{eq:s7v92-PsiY},
\begin{equation}
 {\cal O}_\mu(\overline\Phi)
 \leq
 \sup_i
 \sum_{\substack{C\in{\cal C}\\i\in R(C)}}
 (|R(C)|-1)e^{\mu\operatorname{diam}R(C)}
 \operatorname{osc}w_C.                             \label{eq:s7v92-cluster-norm}
\end{equation}
\end{lemma}

\begin{proof}
Oscillation is subadditive.  Insert
\eqref{eq:s7v92-PsiY} into \eqref{eq:s7v92-O}, bound the oscillation of
each footprint sum by the sum of cluster oscillations, and undo the
regrouping.  Absolute convergence justifies the exchange of sums.
\end{proof}

\begin{corollary}[Complete mixed-oscillation interaction]
\label{cor:s7v92-complete}
If the right side of \eqref{eq:s7v92-cluster-norm} is finite, the exact
coarse logarithm lies in the V83 linear mixed-oscillation class.  In
particular,
\begin{equation}
 {\mathfrak D}_\mu(\overline\Phi)
 \leq
 e^{\kappa_*(\overline\Phi)}
 {\cal O}_\mu(\overline\Phi).                        \label{eq:s7v92-D-linear}
\end{equation}
\end{corollary}

\begin{proof}
Apply the V83 linear comparison
\Cref{lem:s7v83-linear} to the complete interaction
\eqref{eq:s7v92-coarse-action}.
\end{proof}

\section{Typed retained/eliminated path summation}
\label{sec:s7v92-typed}

We now give a sufficient majorant for the right side of
\eqref{eq:s7v92-cluster-norm}.  Split fine variables into retained type
\(R\) and eliminated type \(E\).  Let four nonnegative weighted
connection kernels have row norms
\begin{align}
 \|K_{RR}\|_\mu&\leq\rho,&
 \|K_{RE}\|_\mu&\leq r,&
 \|K_{ER}\|_\mu&\leq s,&
 \|K_{EE}\|_\mu&\leq\epsilon.                       \label{eq:s7v92-four}
\end{align}
The first index is the arrival type and the second the departure type,
as in V84.

\begin{lemma}[Eliminated path resolvent]
\label{lem:s7v92-E-path}
If \(\epsilon<1\), the weighted norm of all paths which leave a
retained block, use only eliminated internal vertices, and return to a
retained block is at most
\begin{equation}
 \tau:={rs\over1-\epsilon}.                          \label{eq:s7v92-tau}
\end{equation}
Including a direct retained connection gives total two-retained
connection row at most
\begin{equation}
 \rho+\tau.                                         \label{eq:s7v92-two-retained}
\end{equation}
\end{lemma}

\begin{proof}
The eliminated internal path sum is
\[
 K_{RE}\sum_{n\geq0}K_{EE}^nK_{ER}.
\]
Use the weighted convolution algebra
\Cref{lem:s7v86-algebra}; its norm is bounded by
\(rs\sum_{n\geq0}\epsilon^n=\tau\).
Add the direct kernel \(K_{RR}\).
\end{proof}

Off-path connected decorations must also be summed.  Let
\(D_{\rm dec}\geq1\) be a proved uniform factor which includes every
rooted branch attached to a selected retained-to-retained spine, with
each cluster assigned to one spine by a deterministic ordering.

\begin{definition}[Coarse selected-tree majorant, CSTM]
\label{def:s7v92-CSTM}
CSTM means that the complete conditional clusters in
\eqref{eq:s7v92-log}, after multiplication by
\((|R(C)|-1)\operatorname{osc}w_C\), admit a deterministic selected
retained spanning tree such that:
\begin{enumerate}[label=\textup{(T\arabic*)},leftmargin=4em]
\item each selected tree edge is either one \(K_{RR}\) connection or
      one \(R\)-to-\(E\), \(E\)-path, \(E\)-to-\(R\) connection;
\item all branches not on the selected retained tree have total rooted
      factor at most \(D_{\rm dec}\);
\item every conditional cluster is counted once; and
\item the weighted metric of a projected coarse edge is dominated by
      the fine path metric used in \eqref{eq:s7v92-four}.
\end{enumerate}
\end{definition}

\begin{theorem}[One-step coarse logarithmic norm recursion]
\label{thm:s7v92-log-recursion}
Under CSTM and \(\epsilon<1\),
\begin{equation}
 {\cal O}_\mu(\overline\Phi)
 \leq
 D_{\rm dec}
 \left(\rho+{rs\over1-\epsilon}\right).              \label{eq:s7v92-log-recursion}
\end{equation}
\end{theorem}

\begin{proof}
Root the selected retained spanning tree at the tested retained block
\(i\).  The factor \(|R(C)|-1\) may be assigned by marking the first
selected tree edge leading toward one of the other retained vertices.
For a direct marked edge, the weighted row sum is at most \(\rho\).
For an edge whose internal vertices are eliminated, sum its path by
\Cref{lem:s7v92-E-path}, costing at most \(\tau\).  Sum all unmarked
rooted branches using CSTM (T2).  Unique selection (T3) prevents
multiplicity, and (T4) preserves the exponential weight.  Taking the
supremum over \(i\) gives
\eqref{eq:s7v92-log-recursion}.
\end{proof}

\begin{firewall}[CSTM is the constructive core]
\label{fire:s7v92-CSTM}
The four row constants alone do not prove
\eqref{eq:s7v92-log-recursion}.  CSTM asserts absolute connected
normalization, a unique selected retained tree, and summability of every
off-spine decoration.  These are precisely the properties to be
constructed from the V91 conditional polymer gas.
\end{firewall}

\section{Where the feedback denominator belongs}
\label{sec:s7v92-feedback}

There are two different outputs of one marginalization step.

\begin{theorem}[Logarithmic norm versus conditional influence]
\label{thm:s7v92-two-recursions}
Let
\[
 \tau={rs\over1-\epsilon}<1.
\]
Then:
\begin{enumerate}[label=\textup{(\roman*)},leftmargin=3.5em]
\item the selected connected logarithmic interaction has the sufficient
      bound
      \[
       {\cal O}_\mu(\overline\Phi)
       \leq D_{\rm dec}(\rho+\tau);
      \]
\item a direct Dobrushin marginalization of a joint specification has
      the V84 bound
      \[
       \|C^{\rm eff}\|_\mu
       \leq{\rho+\tau\over1-\tau};
      \]
\item converting the complete coarse logarithm to influence instead
      gives
      \[
       {\mathfrak D}_\mu(\overline\Phi)
       \leq
       e^{\kappa_*(\overline\Phi)}
       D_{\rm dec}(\rho+\tau).
      \]
\end{enumerate}
The denominators or exponential conversion factors are alternatives
belonging to different proofs and must not be multiplied without a
separate argument.
\end{theorem}

\begin{proof}
Part (i) is \Cref{thm:s7v92-log-recursion}.  Part (ii) is the V84
conditional feedback theorem
\Cref{thm:s7v84-recursion}.  Part (iii) combines
\eqref{eq:s7v92-D-linear} with part (i).  Since (ii) couples conditional
single-site updates while (iii) estimates oscillations of the
already-integrated logarithm, they are two routes to a coarse influence
bound, not successive mandatory factors.
\end{proof}

\begin{corollary}[Two sufficient terminal entry tests]
\label{cor:s7v92-entry}
Either of the following is sufficient for a terminal Dobrushin entry:
\begin{align}
 \rho+2\tau&<1
 &&\text{by direct marginal influence},             \label{eq:s7v92-direct-entry}\\
 e^{\kappa_*(\overline\Phi)}
 D_{\rm dec}(\rho+\tau)&<1
 &&\text{by the complete logarithm}.                 \label{eq:s7v92-log-entry}
\end{align}
\end{corollary}

\begin{proof}
The first is V84.  The second and
\Cref{thm:s7v92-two-recursions} give
\({\mathfrak D}_\mu(\overline\Phi)<1\).
\end{proof}

\section{A scale recursion}
\label{sec:s7v92-scale}

At coarse step \(j\), attach subscripts \(j\) to
\(\rho,r,s,\epsilon,D_{\rm dec}\), and let
\[
 \zeta_{j+1}:=
 {\cal O}_{\mu_{j+1}}(\Phi^{(j+1)}).
\]

\begin{corollary}[Exact-step KP/oscillation recursion]
\label{cor:s7v92-scale}
If the exact step has CSTM, then
\begin{equation}
 \zeta_{j+1}
 \leq
 D_{{\rm dec},j}
 \left(
  \rho_j+
  {r_js_j\over1-\epsilon_j}
 \right).                                           \label{eq:s7v92-scale}
\end{equation}
If the right side is below the V91 KP threshold after conversion to
bounded factors, the next-scale measure has BSCTC.  If
\eqref{eq:s7v92-log-entry} holds, it has EDRC directly.
\end{corollary}

\begin{proof}
Apply \Cref{thm:s7v92-log-recursion} at scale \(j\).
The two terminal conclusions are
\Cref{thm:s7v91-BSCTC,cor:s7v92-entry}.
\end{proof}

\begin{firewall}[Norm conversion must be shown]
\label{fire:s7v92-conversion}
The oscillation norm \(\zeta_{j+1}\) does not automatically equal the
bounded-factor activity
\(\|e^{-\Psi_Y}-1\|_\infty\) used by V91.  One may use
\[
 \|e^{-\Psi_Y}-1\|_\infty
 \leq e^{\|\Psi_Y\|_\infty}-1,
\]
but this also needs a normalized representative controlling the
one-body and constant parts of \(\Psi_Y\).  A pressure constant may be
discarded; an unbounded one-body row may not.
\end{firewall}

\section{Gauge, boundary, and source compatibility}
\label{sec:s7v92-physical}

\begin{proposition}[Exact symmetries of the coarse logarithm]
\label{prop:s7v92-symmetry}
If \({\cal B}\), the fine interaction, and the reference disintegration
are gauge covariant, then \({\cal Z}(y)\) and
\(\overline\Phi(y)\) are gauge invariant.  If boundary representation
labels are normalized as in V80, footprint regrouping introduces no
unweighted representation-cardinality factor.
\end{proposition}

\begin{proof}
For a gauge transform of \(y\), change variables by the corresponding
fine gauge transform in \eqref{eq:s7v92-Zy}.  Conditional reference
covariance and invariance of \(\Phi\) preserve the integral.  The V80
Gauss disintegration is a convex normalized sector mixture, so summing
uniform cluster bounds over its labels has total weight one.
\end{proof}

\begin{proposition}[Coarse boundary sources pull back exactly]
\label{prop:s7v92-sources}
For a coarse observable \(F(y)\), insertion of
\(e^{JF(y)}\) in the coarse conditional partition function equals
insertion of \(e^{JF({\cal B}\omega)}\) in the fine measure.
Consequently a connected fine source expansion pushes forward to the
coarse source logarithm without an observable approximation error.
\end{proposition}

\begin{proof}
The source factor is constant on every conditional fibre
\({\cal B}^{-1}(y)\), so it factors through
\eqref{eq:s7v92-Zy}.  Apply the disintegration identity in the proof of
\Cref{thm:s7v92-density}.
\end{proof}

\section{One-step coarse pushforward card}
\label{sec:s7v92-card}

\begin{definition}[Coarse conditional logarithm card, CCLC]
\label{def:s7v92-CCLC}
For one RG step, CCLC consists of:
\begin{enumerate}[label=\textup{(C\arabic*)},leftmargin=4em]
\item an exact gauge- and reflection-equivariant block map and
      disintegration \eqref{eq:s7v92-disintegration};
\item a zero-free, locally uniformly absolutely convergent conditional
      logarithm \eqref{eq:s7v92-log};
\item complete retained-footprint dependence and the regrouping
      \eqref{eq:s7v92-PsiY};
\item a CSTM selected-tree majorant with
      \(\epsilon<1\) and finite \(D_{\rm dec}\);
\item uniformity in volume, flux sectors, torons, charts, and removed
      fast cutoffs;
\item a proved conversion to the next-scale KP or EDRC norm; and
\item exact source pullback and reflection-positive pushforward.
\end{enumerate}
\end{definition}

\begin{theorem}[CCLC advances the terminal recursion]
\label{thm:s7v92-CCLC}
CCLC produces the exact complete coarse interaction
\eqref{eq:s7v92-coarse-action}, preserves reflection positivity, and
obeys the norm recursion \eqref{eq:s7v92-scale}.  If either terminal
entry test in \Cref{cor:s7v92-entry} or the V91 KP threshold is reached,
the coarse state has the corresponding terminal mass-gap theorem.
\end{theorem}

\begin{proof}
Use
\Cref{thm:s7v92-density,cor:s7v92-RP,
thm:s7v92-regroup,thm:s7v92-log-recursion,
cor:s7v92-scale,prop:s7v92-sources}.
\end{proof}

\section{Status and next acquisition}
\label{sec:s7v92-status}

\begin{theorem}[What V92 proves]
\label{thm:s7v92-result}
V92 proves that one exact coarse pushforward has density equal to a
conditional partition function and that an absolutely convergent
conditional cluster logarithm regroups exactly into the complete coarse
interaction.  Under the selected-tree majorant it obeys
\[
 {\cal O}_\mu(\overline\Phi)
 \leq
 D_{\rm dec}
 \left(\rho+{rs\over1-\epsilon}\right).
\]
It also proves that the V84 feedback denominator belongs to the direct
conditional-influence route and must not be automatically multiplied
into the logarithmic-interaction route.
\end{theorem}

\begin{proof}
Use
\Cref{thm:s7v92-density,thm:s7v92-regroup,
lem:s7v92-cluster-norm,lem:s7v92-E-path,
thm:s7v92-log-recursion,thm:s7v92-two-recursions,
thm:s7v92-CCLC}.
\end{proof}

\begin{assessment}[Progress at seventh-series V92]
\label{ass:s7v92-status}
The algebra of an exact RG step and the correct coarse norm recursion
are now explicit.  The remaining constructive burden has narrowed to
CSTM: prove a conditional selected-tree expansion and bound its four
typed kernels and decoration factor uniformly along the physical
trajectory.
\end{assessment}

\begin{decision}[V93 acquisition order]
\label{dec:s7v92-V93}
V93 will prove a conditional polymer-to-CSTM theorem.  Starting from a
KP expansion on each fibre, it will select the minimal retained
spanning tree, resum eliminated-only branches, and derive explicit
\(\rho,r,s,\epsilon,D_{\rm dec}\) from support-resolved activities.
It will test carefully whether the conditional fibre measure actually
factorizes at the block scale; any failure will be recorded as a
reference-covariance kernel rather than assumed away.
\end{decision}

\begin{firewall}[Claim boundary after seventh-series V92]
\label{fire:s7v92-claim}
V92 does not prove CCLC or CSTM for the physical Yang--Mills coarse
step, bound the four typed kernels along the ultraviolet trajectory,
prove eventual terminal entry, construct the complete nonlinear
coarse-action cocycle, close BSCTC through all cutoff removals, CSMC,
MIRC, or ARDPC for the physical trajectory, establish
regulator-uniform SATC or full RL4C, prove interacting endpoint
continuity, UCM, BCC, pairing/Bogoliubov control, regulator
independence, continuum OS reconstruction, or the full Yang--Mills mass
gap.
\end{firewall}

\paragraph{Parent-version record.}
V92 was formed from
\texttt{Renewal\_Yang\_Mills\_Mass\_Gap\_Research\_Notes\_V91.tex}.
The V91 parent had \(41\,567\) counted source lines,
\(1\,485\,066\) bytes, and SHA--256
\[
 \texttt{38EC98F9390A9A917273599B3D575AC5EB45D606325FB301E410CBAB21E027A6}.
\]
The next compulsory companion audit is V95.

% =============================================================================
% END APPENDED SEVENTH-SERIES V92 MATERIAL
% =============================================================================

% =============================================================================
% BEGIN APPENDED SEVENTH-SERIES V93 MATERIAL
% =============================================================================

\chapter{Conditional cumulant trees and a constructive CSTM loader}
\label{ch:s7v93}

\section{Purpose}
\label{sec:s7v93-purpose}

V92 reduced one coarse step to CSTM but did not derive that card from a
conditional fibre measure.  The fibre generally does not factorize:
conditioning on a block variable can correlate eliminated links which
were independent before conditioning.  We therefore retain those
correlations as a reference cumulant kernel and prove a tree loader
which includes them.

\section{A uniform conditional cumulant card}
\label{sec:s7v93-cumulant}

For a retained field \(y\), let \(\nu^y\) be the normalized eliminated
fibre law.  Let \(A_i\) be bounded observables supported on fine block
sets \(X_i\).  Assume a nonnegative support kernel \(Q_\mu(X,Y)\) whose
weighted one-block row is
\begin{equation}
 {\cal Q}_\mu:=
 \sup_x\sum_Y e^{\mu d(x,Y)}Q_\mu(\{x\},Y)<\infty.  \label{eq:s7v93-Q}
\end{equation}

\begin{definition}[Uniform conditional cumulant tree card, UCCTC]
\label{def:s7v93-UCCTC}
UCCTC means that, uniformly in \(y\) and all physical labels,
\begin{equation}
 |\kappa_{\nu^y}(A_1,\ldots,A_m)|
 \leq
 C_{\rm c}^{m}
 \prod_{i=1}^m\|A_i\|_\infty
 \sum_{T\in{\mathfrak T}_m}
 \prod_{\{i,j\}\in E(T)}Q_\mu(X_i,X_j)              \label{eq:s7v93-tree}
\end{equation}
for \(m\geq2\), with the evident one-point bound for \(m=1\).
\end{definition}

\begin{remark}[Product fibres are a special case]
\label{rem:s7v93-product}
For a product fibre, cumulants vanish when the overlap graph is
disconnected.  UCCTC then holds with \(Q\) equal to a finite-range
overlap kernel.  A correlated fibre has an additional reference
connection in \(Q\).
\end{remark}

\section{Interaction vertices and the tree parameter}
\label{sec:s7v93-vertices}

Let local bounded factors have a degree-resolved representation whose
total activity at one fine block is at most \(u\).  All fixed cutoff,
chart, label, and Weyl moments are included in \(u\).  Put
\begin{equation}
 q_{\rm br}:=eC_{\rm c}u{\cal Q}_\mu.                \label{eq:s7v93-q}
\end{equation}

\begin{lemma}[Rooted conditional tree sum]
\label{lem:s7v93-rooted}
Under UCCTC, if \(q_{\rm br}<1\), the absolute sum of all interaction
trees rooted at one declared vertex is bounded by
\begin{equation}
 D_{\rm root}\leq{1\over1-q_{\rm br}}.              \label{eq:s7v93-root}
\end{equation}
The same sum with one marked tree edge is bounded by
\begin{equation}
 D_{\rm mark}\leq{1\over(1-q_{\rm br})^2}.          \label{eq:s7v93-mark}
\end{equation}
\end{lemma}

\begin{proof}
At order \(n\), sum positions outward from a fixed root using
\eqref{eq:s7v93-Q}; each added vertex costs
\(C_{\rm c}u{\cal Q}_\mu\).  Cayley's formula and
\[
 {n^{n-2}\over(n-1)!}\leq e^{n-1}
\]
bound the normalized labelled-tree multiplicity by one factor \(e\)
per added vertex.  Thus the rooted order is at most
\(q_{\rm br}^{n-1}\), whose sum is
\eqref{eq:s7v93-root}.  Marking one of at most \(n\) edges replaces the
geometric sum by
\(\sum_{k\geq0}(k+1)q_{\rm br}^k\), giving
\eqref{eq:s7v93-mark}.
\end{proof}

\section{Typed marked connections}
\label{sec:s7v93-typed}

Color every support vertex retained \(R\) or eliminated \(E\).  In the
UCCTC tree sum, mark the first edge on the path from a fixed retained
root toward another retained vertex.  After all unmarked branches are
summed, assume the four possible marked edge rows obey
\begin{align}
 \rho_0&:=\|L_{RR}\|_\mu,&
 r_0&:=\|L_{RE}\|_\mu,&
 s_0&:=\|L_{ER}\|_\mu,&
 \epsilon_0&:=\|L_{EE}\|_\mu.                       \label{eq:s7v93-four}
\end{align}
These kernels include both overlap interactions and conditional
reference covariance.

\begin{theorem}[Conditional tree loader for CSTM]
\label{thm:s7v93-CSTM}
Assume UCCTC, \(q_{\rm br}<1\), and
\begin{equation}
 \epsilon:=D_{\rm root}\epsilon_0<1.                \label{eq:s7v93-epsilon}
\end{equation}
Then the conditional logarithm has CSTM with
\begin{align}
 \rho&=D_{\rm root}\rho_0,&
 r&=D_{\rm root}r_0,&
 s&=D_{\rm root}s_0,&
 \epsilon&=D_{\rm root}\epsilon_0,                  \label{eq:s7v93-dressed}\\
 D_{\rm dec}&\leq D_{\rm mark}.                     \label{eq:s7v93-Ddec}
\end{align}
Consequently
\begin{equation}
 {\cal O}_\mu(\overline\Phi)
 \leq
 {1\over(1-q_{\rm br})^2}
 \left[
  {\,\rho_0\over1-q_{\rm br}}
  +
  {r_0s_0\over
   (1-q_{\rm br})^2
   \left(1-{\epsilon_0\over1-q_{\rm br}}\right)}
 \right].                                           \label{eq:s7v93-recursion}
\end{equation}
\end{theorem}

\begin{proof}
Expand the conditional logarithm in cumulants of local interaction
vertices.  UCCTC bounds every cumulant by labelled trees.  For each
tree with at least two retained vertices, root it at the first retained
vertex and select the lexicographically first retained spanning subtree
obtained by suppressing eliminated degree-two chains.  Mark its first
edge toward a second retained vertex.  This deterministic rule counts
each tree once.

Every branch not on a selected connection is a rooted tree and costs at
most \(D_{\rm root}\) by
\Cref{lem:s7v93-rooted}.  Absorb this factor into the four marked
connection kernels, giving \eqref{eq:s7v93-dressed}.  A chain of
eliminated internal connections sums geometrically because
\eqref{eq:s7v93-epsilon} holds.  The remaining marked-tree
multiplicity costs \(D_{\rm mark}\).  Thus CSTM holds with
\eqref{eq:s7v93-Ddec}.  Substitute these quantities into the V92
recursion \eqref{eq:s7v92-log-recursion} and simplify.
\end{proof}

\begin{corollary}[A fully numerical sufficient entry test]
\label{cor:s7v93-entry}
If the right side of \eqref{eq:s7v93-recursion}, multiplied by
\(e^{\kappa_*(\overline\Phi)}\), is strictly below one, then the exact
coarse marginal belongs to the V83 Dobrushin ball.
\end{corollary}

\begin{proof}
Use \Cref{cor:s7v92-complete,thm:s7v93-CSTM}.
\end{proof}

\section{Why two-point fibre decay is insufficient}
\label{sec:s7v93-high}

\begin{theorem}[A global conditional constraint can hide in high
cumulants]
\label{thm:s7v93-parity}
Let \(X_1,\ldots,X_N\) be independent fair signs conditioned on
\begin{equation}
 \prod_{i=1}^NX_i=1.                                \label{eq:s7v93-parity}
\end{equation}
Every marginal of fewer than \(N\) variables is the independent fair
law.  Hence all distinct two-point covariances vanish, but
\begin{equation}
 \mathbb E\prod_{i=1}^NX_i=1.                       \label{eq:s7v93-Npoint}
\end{equation}
Thus pair decay alone does not imply UCCTC uniformly in cumulant order.
\end{theorem}

\begin{proof}
After fixing any \(k<N\) signs, exactly half of the remaining
configurations satisfy \eqref{eq:s7v93-parity}; the \(k\)-variable
marginal is therefore uniform.  The full product is one by the
conditioning event.
\end{proof}

\begin{firewall}[Conditional factorization must be proved]
\label{fire:s7v93-factorization}
A local block map may impose shared holonomy, Gauss, or transition
constraints on an entire fibre.  Even if its conditional two-point
kernel decays, \Cref{thm:s7v93-parity} shows that a high-order global
constraint can invalidate a tree cumulant estimate.  UCCTC must be
proved from an explicit fibre parametrization and normalized constraint
disintegration.
\end{firewall}

\section{Loading UCCTC from a conditional KP gas}
\label{sec:s7v93-KP}

\begin{theorem}[Uniform conditional KP implies UCCTC]
\label{thm:s7v93-KP}
Suppose that for every retained \(y\), the normalized fibre law is a
hard-core polymer perturbation of a product fibre reference, with a KP
margin uniform in \(y\).  Suppose also that insertion of each local
observable changes only polymer activities meeting its support and that
the source-dependent KP margin is uniform in a fixed polydisc.  Then
the conditional cumulants admit UCCTC with a kernel \(Q\) equal to the
sum of the finite-range overlap kernel and the exponentially weighted
connected-polymer response kernel.
\end{theorem}

\begin{proof}
Add independent sources \(J_iA_i\) to the fibre partition function.
The uniform KP theorem gives a locally uniformly absolutely convergent
connected source logarithm.  Differentiate \(m\) times.  Only connected
source-polymer clusters meeting every \(X_i\) survive.  Apply the
tree-graph inequality to their incompatibility graph and select a
spanning tree on the marked sources after contracting unmarked polymer
branches.  Each selected edge is either a direct overlap or a connected
polymer chain; their absolute sums define \(Q\).  The uniform rooted KP
bound pays all contracted branches by one constant per source.  This is
\eqref{eq:s7v93-tree}.
\end{proof}

\begin{corollary}[Product-fibre assumption is avoidable]
\label{cor:s7v93-reference}
It is enough that the conditional fibre itself has a uniform
source-dependent KP representation.  Exact product factorization after
conditioning is not required.
\end{corollary}

\begin{proof}
Apply \Cref{thm:s7v93-KP}; its resulting response kernel records the
conditional reference correlations.
\end{proof}

\section{The conditional selected-tree card}
\label{sec:s7v93-card}

\begin{definition}[Conditional fibre tree card, CFTC]
\label{def:s7v93-CFTC}
CFTC consists of:
\begin{enumerate}[label=\textup{(F\arabic*)},leftmargin=4em]
\item an exact normalized fibre disintegration for the block map;
\item a source-dependent conditional KP expansion uniform in retained
      fields and all physical labels;
\item the UCCTC bound \eqref{eq:s7v93-tree};
\item explicit \(u,{\cal Q}_\mu\) with \(q_{\rm br}<1\);
\item typed marked rows \eqref{eq:s7v93-four} satisfying
      \(\epsilon_0<1-q_{\rm br}\);
\item the complete operation/ownership ledger in every activity; and
\item compatibility with coarse source pullback and reflection.
\end{enumerate}
\end{definition}

\begin{theorem}[CFTC loads CCLC]
\label{thm:s7v93-CFTC}
CFTC gives CSTM with the explicit recursion
\eqref{eq:s7v93-recursion}; together with the V92 exact disintegration
and norm conversion, it loads CCLC for one coarse step.
\end{theorem}

\begin{proof}
Use
\Cref{thm:s7v93-KP,thm:s7v93-CSTM,
thm:s7v92-CCLC}.
\end{proof}

\section{Status and next acquisition}
\label{sec:s7v93-status}

\begin{theorem}[What V93 proves]
\label{thm:s7v93-result}
V93 proves that a uniform source-dependent conditional KP expansion
implies a full cumulant tree card and hence the CSTM selected-tree
majorant.  It derives the explicit coarse norm recursion
\eqref{eq:s7v93-recursion}, including reference fibre correlations and
all eliminated return paths.  It also proves that conditional
two-point decay alone cannot replace the required all-order card.
\end{theorem}

\begin{proof}
Use
\Cref{lem:s7v93-rooted,thm:s7v93-CSTM,
thm:s7v93-parity,thm:s7v93-KP,
thm:s7v93-CFTC}.
\end{proof}

\begin{assessment}[Progress at seventh-series V93]
\label{ass:s7v93-status}
The abstract one-step recursion is now constructively loaded by a
conditional KP gas without assuming product fibres.  The remaining
physical task is quantitative: build that uniform conditional KP
representation for an actual Yang--Mills block map and compute its
typed rows along the weak-to-terminal trajectory.
\end{assessment}

\begin{decision}[V94 acquisition order]
\label{dec:s7v93-V94}
V94 will examine an explicit compact-group block map based on ordered
path products and fluctuation links.  It will compute the fibre Haar
Jacobian, Gauss constraint, and conditional reference dependencies,
then determine whether they have finite-range product structure or a
uniform KP response kernel.  This goes upstream of the typed constants
rather than assuming CFTC.
\end{decision}

\begin{firewall}[Claim boundary after seventh-series V93]
\label{fire:s7v93-claim}
V93 does not prove CFTC for a physical Yang--Mills block map, compute
the typed rows on the ultraviolet trajectory, prove terminal entry,
construct the complete nonlinear coarse-action cocycle, close BSCTC
through all cutoff removals, CSMC, MIRC, or ARDPC, establish
regulator-uniform SATC or full RL4C, prove interacting endpoint
continuity, UCM, BCC, pairing/Bogoliubov control, regulator
independence, continuum OS reconstruction, or the full Yang--Mills mass
gap.
\end{firewall}

\paragraph{Parent-version record.}
V93 was formed from
\texttt{Renewal\_Yang\_Mills\_Mass\_Gap\_Research\_Notes\_V92.tex}.
The V92 parent had \(42\,149\) counted source lines,
\(1\,505\,470\) bytes, and SHA--256
\[
 \texttt{90B22DECFE116DCC4DE8CEF19DDCF47927D49F7954632E13480B9DFD47DAD62A}.
\]
The next compulsory companion audit is V95.

% =============================================================================
% END APPENDED SEVENTH-SERIES V93 MATERIAL
% =============================================================================

% =============================================================================
% BEGIN APPENDED SEVENTH-SERIES V94 MATERIAL
% =============================================================================

\chapter{Tree coordinates for compact blocks and exact Haar fibres}
\label{ch:s7v94}

\section{Purpose}
\label{sec:s7v94-purpose}

V93 made the conditional fibre law the next upstream object.  We now
compute its reference geometry for one finite connected gauge block.
A rooted spanning tree separates vertex-frame variables from chord
holonomies with unit Haar Jacobian.  Normalized Gauss integration removes
the frames, leaving product Haar measure on the chords modulo one root
conjugation.  Thus a block map built from disjoint chord coordinates has
an exact product reference fibre; all remaining dependence comes from
the action and explicitly shared boundary data.

\section{Rooted tree coordinates}
\label{sec:s7v94-tree}

Let \(B=(V,E)\) be a finite connected oriented graph, choose a root
\(o\), and choose a spanning tree \(T\subset E\).  Write
\(C=E\setminus T\) for the chords.  Set \(h_o=1\).  For each
\(v\ne o\), let \(h_v\) be the ordered parallel transport from \(o\) to
\(v\) along the unique tree path, using inverses when an edge orientation
is traversed backwards.  For a chord \(e:u\to v\), define
\begin{equation}
 z_e:=h_uU_eh_v^{-1}.                               \label{eq:s7v94-z}
\end{equation}

\begin{theorem}[Tree--chord Haar factorization]
\label{thm:s7v94-Haar}
The map
\begin{equation}
 \{U_e:e\in E\}
 \longleftrightarrow
 \bigl(\{h_v:v\ne o\},\{z_e:e\in C\}\bigr)           \label{eq:s7v94-map}
\end{equation}
is a measurable bijection up to Haar-null coordinate choices, and
\begin{equation}
 \prod_{e\in E}dU_e
 =
 \prod_{v\ne o}dh_v
 \prod_{e\in C}dz_e.                                \label{eq:s7v94-Haar}
\end{equation}
\end{theorem}

\begin{proof}
Order tree vertices so every parent precedes its children.  If the tree
edge from parent \(p\) to child \(v\) has the chosen orientation, then
\(U_{pv}=h_p^{-1}h_v\); with the opposite orientation use its inverse.
Thus the tree links are reconstructed successively from the \(h_v\).
Every chord is reconstructed from
\(U_e=h_u^{-1}z_eh_v\), proving bijectivity.

For fixed earlier variables, the change from one tree link to the next
\(h_v\) is left or right translation, possibly followed by inversion.
Normalized Haar measure is invariant under all three operations.
Conditioned on the frames, the change \(U_e\mapsto z_e\) for each chord
is again a left-right translation.  Applying Fubini in the parent-first
order proves \eqref{eq:s7v94-Haar}.
\end{proof}

\section{Gauge action and the normalized quotient}
\label{sec:s7v94-Gauss}

Use the convention
\begin{equation}
 U_{uv}\mapsto k_uU_{uv}k_v^{-1}.                  \label{eq:s7v94-gauge}
\end{equation}
Then tree transport transforms as
\begin{equation}
 h_v\mapsto k_oh_vk_v^{-1},                         \label{eq:s7v94-h-transform}
\end{equation}
and every chord coordinate transforms by the single root conjugation
\begin{equation}
 z_e\mapsto k_oz_ek_o^{-1}.                         \label{eq:s7v94-z-transform}
\end{equation}

\begin{corollary}[Local Gauss variables separate]
\label{cor:s7v94-Gauss}
Normalized integration over gauge transformations at \(v\ne o\)
removes the frame variables \(h_v\) with total factor one.  The
gauge-invariant block reference is product Haar on \(G^C\), followed
only by normalized simultaneous conjugation at the root.
\end{corollary}

\begin{proof}
Equations \eqref{eq:s7v94-h-transform} and
\eqref{eq:s7v94-z-transform} show that nonroot gauge transformations act
only on the frame coordinates after the tree change of variables.
Their normalized Haar integrals equal one by
\eqref{eq:s7v94-Haar}.  The residual root transformation conjugates all
chords simultaneously.
\end{proof}

\begin{remark}[No representation dimension appears]
\label{rem:s7v94-dimension}
Resolving the final root conjugation into irreducibles gives the
normalized invariant pairing of V80.  Staying in Haar variables makes
the absence of an unweighted \(d_\lambda\) factor immediate.
\end{remark}

\section{A product-fibre block map}
\label{sec:s7v94-block-map}

Partition the chord set as
\begin{equation}
 C=C_{\rm ret}\sqcup C_{\rm elim}.                  \label{eq:s7v94-split}
\end{equation}
Let the coarse block variable be
\begin{equation}
 y=(z_e)_{e\in C_{\rm ret}}                         \label{eq:s7v94-y}
\end{equation}
before the final root-conjugation quotient.

\begin{theorem}[Exact reference disintegration]
\label{thm:s7v94-fibre}
For the product Haar link reference, the conditional law given \(y\) is
\begin{equation}
 \nu_0^y
 =
 \left(\prod_{v\ne o}dh_v\right)
 \left(\prod_{e\in C_{\rm elim}}dz_e\right),         \label{eq:s7v94-fibre}
\end{equation}
independent of \(y\).  After normalized nonroot Gauss integration it is
product Haar on the eliminated chords.  Hence the reference part of
UCCTC is a finite-range overlap kernel with no additional conditional
correlation.
\end{theorem}

\begin{proof}
Factor \eqref{eq:s7v94-Haar} according to
\eqref{eq:s7v94-split}.  Conditioning fixes only the retained chord
coordinates; all remaining factors stay normalized Haar and contain no
\(y\)-dependent density.  Apply
\Cref{cor:s7v94-Gauss}.
\end{proof}

\begin{corollary}[Independent disjoint blocks]
\label{cor:s7v94-blocks}
If spatial blocks have disjoint owned link sets and all shared boundary
links are retained, the product Haar reference fibres of distinct
blocks are independent.  Fine interactions crossing blocks are then
the only reference to both fibres and belong in the conditional polymer
activity.
\end{corollary}

\begin{proof}
Apply \Cref{thm:s7v94-fibre} in each disjoint owned link set and use
product Haar factorization.  Retained boundary links are conditioned
data rather than integrated variables.
\end{proof}

\section{Path-product coarse variables}
\label{sec:s7v94-products}

A coarse link is often an ordered product of several chord variables.
Haar measure still permits a triangular product coordinate.

\begin{lemma}[Haar product coordinate]
\label{lem:s7v94-product}
If \(z_1,\ldots,z_m\) are independent Haar variables and
\begin{equation}
 y=z_1\cdots z_m,                                    \label{eq:s7v94-product}
\end{equation}
then the change
\((z_1,\ldots,z_m)\leftrightarrow(y,z_2,\ldots,z_m)\)
preserves product Haar measure.  Conditional on \(y\), the remaining
\(m-1\) coordinates have product Haar law.
\end{lemma}

\begin{proof}
For fixed \(z_2,\ldots,z_m\),
\(z_1=y(z_2\cdots z_m)^{-1}\) is a right translation of \(y\).
Haar invariance and Fubini prove the measure identity and the
conditional statement.
\end{proof}

\begin{corollary}[Disjoint ordered-path block map]
\label{cor:s7v94-path-map}
If each retained coarse variable is an ordered product of a disjoint
family of chord coordinates, then repeated application of
\Cref{lem:s7v94-product} gives an exact product Haar fibre.
\end{corollary}

\begin{proof}
Apply the triangular change separately to each disjoint family.
\end{proof}

\begin{firewall}[Overlapping coarse paths create constraints]
\label{fire:s7v94-overlap}
If two retained path products share an eliminated chord, conditioning
on both imposes a joint group constraint and the fibre need not
factorize.  The parity phenomenon of V93 has a compact-group analogue.
One must redesign the ownership paths, retain the shared chord, or
include the resulting constraint in the UCCTC reference kernel.
\end{firewall}

\section{Interaction loading}
\label{sec:s7v94-interaction}

\begin{proposition}[Where conditional dependence now resides]
\label{prop:s7v94-dependence}
For a disjoint-path tree block map, the conditional physical fibre is
\begin{equation}
 d\mu^y_{\rm f}
 =
 {\exp[-\Phi(y,h,z_{\rm elim})]\over{\cal Z}(y)}
 \prod_{v\ne o}dh_v
 \prod_{e\in C_{\rm elim}}dz_e.                    \label{eq:s7v94-physical-fibre}
\end{equation}
Thus all nonproduct dependence is generated by the explicitly known
fine interaction \(\Phi\), including crossing plaquettes.  If its
bounded factors satisfy the uniform conditional KP criterion of V93,
CFTC follows with no extra reference-covariance row.
\end{proposition}

\begin{proof}
Insert the Haar disintegration
\eqref{eq:s7v94-fibre} into the exact conditional density
\eqref{eq:s7v92-Zy}.  The final statement is
\Cref{thm:s7v93-KP}.
\end{proof}

\begin{firewall}[Weak coupling remains outside the elementary KP ball]
\label{fire:s7v94-weak}
The exact product fibre removes a geometric obstruction, but a Wilson
factor of size \(\beta\gg1\) is not a small bounded perturbation of Haar.
At the ultraviolet scale the reference must absorb the dominant
near-flat action, using the gapped conditional Gaussian/heat-kernel
construction and the good/bad atlas.  V94 does not turn large
\(\beta\) into a Haar-KP activity.
\end{firewall}

\section{Status and next acquisition}
\label{sec:s7v94-status}

\begin{theorem}[What V94 proves]
\label{thm:s7v94-result}
V94 constructs rooted tree--chord coordinates with exact unit Haar
Jacobian, separates normalized local Gauss frames from physical chord
holonomies, and proves product reference fibres for disjoint retained
chords or disjoint ordered path products.  For such block maps, every
conditional fibre correlation is assigned to the physical interaction,
not to an unknown reference constraint.
\end{theorem}

\begin{proof}
Use
\Cref{thm:s7v94-Haar,cor:s7v94-Gauss,
thm:s7v94-fibre,lem:s7v94-product,
cor:s7v94-path-map,prop:s7v94-dependence}.
\end{proof}

\begin{assessment}[Progress at seventh-series V94]
\label{ass:s7v94-status}
The fibre factorization row of CFTC is now proved for a concrete,
ownership-compatible compact-group block map.  The remaining difficulty
is analytic rather than geometric: at weak coupling one needs a
near-flat conditional reference which absorbs the large Wilson
quadratic action while preserving the exact Haar and bad-chart
decomposition.
\end{assessment}

\begin{decision}[V95 audit and acquisition order]
\label{dec:s7v94-V95}
V95 is the compulsory companion audit.  It will then combine the exact
tree--chord Haar fibres with a conditional near-flat Gaussian reference,
derive the Radon--Nikodym remainder on good blocks, and test its KP
activity together with saturated bad components.  The calculation will
retain the compact Haar density and crossing plaquettes explicitly.
\end{decision}

\begin{firewall}[Claim boundary after seventh-series V94]
\label{fire:s7v94-claim}
V94 does not prove the weak-coupling conditional KP margin, CFTC for the
full physical interaction, terminal entry, the complete nonlinear
coarse-action cocycle, BSCTC through cutoff removal, CSMC, MIRC, or
ARDPC, regulator-uniform SATC or full RL4C, interacting endpoint
continuity, UCM, BCC, pairing/Bogoliubov control, regulator
independence, continuum OS reconstruction, or the full Yang--Mills mass
gap.
\end{firewall}

\paragraph{Parent-version record.}
V94 was formed from
\texttt{Renewal\_Yang\_Mills\_Mass\_Gap\_Research\_Notes\_V93.tex}.
The V93 parent had \(42\,487\) counted source lines,
\(1\,518\,364\) bytes, and SHA--256
\[
 \texttt{4362CAA89AACD6302634C46561E0A96C40C89EBF4A387FA3D6141EA762EC5A9E}.
\]
The next compulsory companion audit is V95.

% =============================================================================
% END APPENDED SEVENTH-SERIES V94 MATERIAL
% =============================================================================

% =============================================================================
% BEGIN APPENDED SEVENTH-SERIES V95 MATERIAL
% =============================================================================

\chapter{V95 audit and the conditional near-flat Gaussian remainder}
\label{ch:s7v95}

\section{Purpose and compulsory audit}
\label{sec:s7v95-purpose}

V94 proved exact product Haar fibres for an ownership-compatible
tree--chord block map.  At ultraviolet weak coupling, however, Haar is
the wrong perturbative reference: the large Wilson quadratic term must
be absorbed into a near-flat Gaussian.  This note performs the V95
companion audit and proves the finite-dimensional rescaled remainder
estimate needed for that replacement.

The active companion files inspected were
\begin{align}
 &\texttt{Renewal\_SM\_Einstein\_Continuum\_Research\_Notes\_V62.tex},
 \label{eq:s7v95-SM}\\
 &\texttt{Renewal\_NS\_Stability\_Research\_Notes\_V26.tex}.
 \label{eq:s7v95-NS}
\end{align}
SM V62 had \(23\,187\) counted source lines, \(836\,565\) bytes, and
SHA--256
\[
 \texttt{D313C9CB0BE50DF3382B161D67F7BE824DCEB557CB3932350C9F0BE64797B8BB}.
\]
NS V26 had \(5\,320\) counted source lines, \(278\,283\) bytes, and
SHA--256
\[
 \texttt{82C887F8C2C69E4F28FAD7C3DC8DE5B9099CB5A4BF431EBA1FFF3E91935978AD}.
\]

SM V61 independently derives a marked conditional coarse pushforward.
SM V62 proves an exact Duhamel formula for principal heat-kernel marks
and computes Gaussian bridge Fisher information
\(D(m-1)/2\).  The transferable lesson is a type restriction:
coarse-dependent Gaussian covariance and its determinant belong to the
reference/coarse action; they cannot be assigned the small scalar
activity of a bounded nonlinear remainder.  NS V26 contains no
type-compatible input.

\begin{theorem}[V95 audit decision]
\label{thm:s7v95-audit}
The near-flat conditional reference will absorb the full quadratic
Wilson Hessian.  Its normalization, determinant, and any principal
variation remain explicit owned coarse rows.  Only cubic-and-higher
Wilson terms and the Haar-coordinate density are tested in the
conditional KP norm.
\end{theorem}

\begin{proof}
The Gaussian bridge computation in the audited SM note rules out a
uniform small scalar mark for a varying principal covariance.  The
Yang--Mills V76 type ledger and V91 scalarization already separate
physical multipliers from principal form rows.  Combining those two
scope statements gives the decision.
\end{proof}

\section{A stationary conditional chart}
\label{sec:s7v95-chart}

Fix one good tree--chord fibre and a retained field \(y\).  Let
\(\xi\in\mathbb R^d\) be eliminated exponential coordinates centered at
a conditional background.  Write the Wilson action as
\begin{equation}
 S_y(\xi)
 =
 S_y(0)+{1\over2}\langle\xi,L_y\xi\rangle+R_y(\xi).
                                                               \label{eq:s7v95-Taylor}
\end{equation}
Assume
\begin{align}
 D_\xi S_y(0)&=0,                                    \label{eq:s7v95-stationary}\\
 \lambda I\leq L_y&\leq\Lambda I,\qquad\lambda>0,    \label{eq:s7v95-Hessian}\\
 |R_y(\xi)|
 &\leq C_3|\xi|^3+C_4|\xi|^4                       \label{eq:s7v95-remainder}
\end{align}
on the good chart, uniformly in admissible \(y\).

\begin{firewall}[The stationary center is essential]
\label{fire:s7v95-linear}
If \(D_\xi S_y(0)\ne0\), the rescaled action contains a term of order
\(\sqrt\beta\,D_\xi S_y(0)\cdot\eta\).  It is not a small KP
remainder.  The background must be the conditional minimizer/harmonic
extension, or the linear source must be absorbed exactly by completing
the Gaussian square.
\end{firewall}

\section{Gaussian rescaling}
\label{sec:s7v95-rescaling}

Let \(\beta=g^{-2}\) and define
\begin{equation}
 \eta:=\sqrt\beta\,L_y^{1/2}\xi.                    \label{eq:s7v95-eta}
\end{equation}
The Gaussian reference is
\begin{equation}
 d\gamma_y(\eta)
 =
 (2\pi)^{-d/2}e^{-|\eta|^2/2}\,d\eta.               \label{eq:s7v95-gamma}
\end{equation}
Let the good occupation region be
\begin{equation}
 G_M:=\{|\eta|\leq\sqrt M\}.                         \label{eq:s7v95-GM}
\end{equation}

\begin{lemma}[Rescaled Wilson remainder]
\label{lem:s7v95-Wilson}
On \(G_M\),
\begin{equation}
 |\beta R_y(\xi)|
 \leq
 \delta_{{\rm W},M}
 :=
 C_3\lambda^{-3/2}\beta^{-1/2}M^{3/2}
 +C_4\lambda^{-2}\beta^{-1}M^2.                    \label{eq:s7v95-deltaW}
\end{equation}
\end{lemma}

\begin{proof}
From \eqref{eq:s7v95-eta} and
\eqref{eq:s7v95-Hessian},
\[
 |\xi|\leq\beta^{-1/2}\lambda^{-1/2}|\eta|.
\]
Insert this into \eqref{eq:s7v95-remainder}, multiply by \(\beta\), and
use \(|\eta|\leq\sqrt M\).
\end{proof}

Let \(J_y(\xi)\) be the Haar-coordinate density relative to Lebesgue
measure.  Assume, after extracting its value at the origin,
\begin{equation}
 |\log J_y(\xi)-\log J_y(0)|
 \leq C_{J,2}|\xi|^2+C_{J,3}|\xi|^3                \label{eq:s7v95-Haar-log}
\end{equation}
on the good chart.

\begin{lemma}[Rescaled Haar remainder]
\label{lem:s7v95-Haar}
On \(G_M\),
\begin{equation}
 |\log J_y(\xi)-\log J_y(0)|
 \leq
 \delta_{{\rm H},M}
 :=
 C_{J,2}\lambda^{-1}\beta^{-1}M
 +C_{J,3}\lambda^{-3/2}\beta^{-3/2}M^{3/2}.
                                                               \label{eq:s7v95-deltaH}
\end{equation}
\end{lemma}

\begin{proof}
Use the same bound on \(|\xi|\) as in
\Cref{lem:s7v95-Wilson} and insert it into
\eqref{eq:s7v95-Haar-log}.
\end{proof}

\section{The normalized good-fibre density}
\label{sec:s7v95-density}

Define the scalar remainder
\begin{equation}
 A_y(\eta)
 :=
 \beta R_y(\xi(\eta))
 -\log{J_y(\xi(\eta))\over J_y(0)}.                 \label{eq:s7v95-A}
\end{equation}
On \(G_M\), put
\begin{equation}
 d\mu_{y,M}^{\rm good}
 :=
 {e^{-A_y(\eta)}\mathbf1_{G_M}\over
  Z_{y,M}^{\rm good}}\,d\gamma_y(\eta),
 \qquad
 Z_{y,M}^{\rm good}
 :=
 \mathbb E_{\gamma_y}
 [e^{-A_y}\mathbf1_{G_M}].                          \label{eq:s7v95-good-law}
\end{equation}
Set
\begin{equation}
 \delta_M:=\delta_{{\rm W},M}+\delta_{{\rm H},M}.    \label{eq:s7v95-delta}
\end{equation}

\begin{theorem}[Uniform normalized remainder bound]
\label{thm:s7v95-normalized}
On \(G_M\),
\begin{equation}
 |A_y|\leq\delta_M.                                 \label{eq:s7v95-A-bound}
\end{equation}
Let \(\gamma_{y,M}\) be \(\gamma_y\) conditioned on \(G_M\).  Then
\begin{align}
 e^{-2\delta_M}
 \leq
 {d\mu_{y,M}^{\rm good}\over d\gamma_{y,M}}
 &\leq e^{2\delta_M},                               \label{eq:s7v95-RN}\\
 \operatorname{osc}_{G_M}
 \log{d\mu_{y,M}^{\rm good}\over d\gamma_{y,M}}
 &\leq2\delta_M,                                    \label{eq:s7v95-log-osc}\\
 \left\|
 {d\mu_{y,M}^{\rm good}\over d\gamma_{y,M}}-1
 \right\|_\infty
 &\leq e^{2\delta_M}-1.                             \label{eq:s7v95-factor}
\end{align}
\end{theorem}

\begin{proof}
\Cref{lem:s7v95-Wilson,lem:s7v95-Haar} give
\eqref{eq:s7v95-A-bound}.  Relative to the conditioned Gaussian,
\[
 {d\mu_{y,M}^{\rm good}\over d\gamma_{y,M}}
 =
 {e^{-A_y}\over
  \mathbb E_{\gamma_{y,M}}e^{-A_y}}.
\]
Both numerator and denominator lie between
\(e^{-\delta_M}\) and \(e^{\delta_M}\), proving
\eqref{eq:s7v95-RN}.  The normalizing denominator is constant in
\(\eta\), so the logarithmic oscillation equals
\(\operatorname{osc}A_y\leq2\delta_M\).
The last bound follows from \eqref{eq:s7v95-RN}.
\end{proof}

\begin{corollary}[Logarithmic schedule]
\label{cor:s7v95-schedule}
If
\begin{equation}
 M=(\log\beta)^\chi
 \quad\text{for fixed }\chi>0,                      \label{eq:s7v95-schedule}
\end{equation}
and the constants in
\eqref{eq:s7v95-Hessian}--\eqref{eq:s7v95-Haar-log} are uniform, then
\begin{equation}
 \delta_M\longrightarrow0,\qquad
 e^{2\delta_M}-1=O(\delta_M).                       \label{eq:s7v95-delta-zero}
\end{equation}
\end{corollary}

\begin{proof}
Every positive power of \(\log\beta\) is
\(o(\beta^\varepsilon)\) for every \(\varepsilon>0\).  Apply this to
the four terms in \eqref{eq:s7v95-deltaW} and
\eqref{eq:s7v95-deltaH}.  The exponential expansion gives the second
claim.
\end{proof}

\section{Conditional KP activity on good blocks}
\label{sec:s7v95-KP}

Assume the action remainder decomposes by owned fixed-radius supports,
\begin{equation}
 A_y(\eta)=\sum_XA_{y,X}(\eta_X),                   \label{eq:s7v95-local-sum}
\end{equation}
and each support satisfies the rescaled bound
\begin{equation}
 \operatorname{osc}A_{y,X}\leq2\delta_{M,X}.         \label{eq:s7v95-local-osc}
\end{equation}
Put
\begin{equation}
 \zeta_{{\rm good},M}
 :=
 \sup_x\sum_{X\ni x}
 \left(e^{2\delta_{M,X}}-1\right)
 e^{a|X|}.                                          \label{eq:s7v95-zeta-good}
\end{equation}

\begin{theorem}[Good conditional KP loader]
\label{thm:s7v95-good-KP}
If the support incidence is uniformly finite and
\(\sup_X\delta_{M,X}\to0\), then
\begin{equation}
 \zeta_{{\rm good},M}\longrightarrow0.              \label{eq:s7v95-good-zero}
\end{equation}
For all sufficiently large \(\beta\), the normalized good-fibre
remainder satisfies a uniform conditional KP criterion.
\end{theorem}

\begin{proof}
Normalize each local factor by subtracting any chosen constant; its
factor activity is at most
\[
 \|e^{-(A_{y,X}-c_X)}-1\|_\infty
 \leq e^{\operatorname{osc}A_{y,X}}-1
 \leq e^{2\delta_{M,X}}-1.
\]
Only a fixed number of supports of fixed cardinality meet a given
block.  Hence the weighted row
\eqref{eq:s7v95-zeta-good} tends to zero.  The standard KP inequality
then holds eventually with a strict margin.
\end{proof}

\begin{firewall}[Uniform Taylor data are not yet proved]
\label{fire:s7v95-uniform}
The theorem assumes a stationary fibre center, a regulator-uniform
positive Hessian on eliminated physical directions, fixed support
incidence, and uniform third/fourth derivative and Haar-density
constants.  V72--V76 provide finite-cutoff components of this ledger,
but regulator-uniform SATC and the complete nonlinear conditional
Taylor card remain open.
\end{firewall}

\section{Bad blocks and the common schedule}
\label{sec:s7v95-bad}

Let \(\pi_M^{+R}\) be the V79 saturated bad-block rate and let
\(A_{\partial,M}\) be its boundary insertion factor.  Define
\begin{equation}
 \zeta_{{\rm bad},M}
 :=
 A_{\partial,M}\pi_M^{+R}.                           \label{eq:s7v95-zeta-bad}
\end{equation}

\begin{proposition}[Good/bad joint activity]
\label{prop:s7v95-joint}
Under the V79 schedule
\[
 \beta=g^{-2},\qquad
 M=(\log\beta)^\chi,\qquad\chi>1,
\]
and the uniform good Taylor hypotheses,
\begin{equation}
 \zeta_{{\rm good},M}
 +\zeta_{{\rm bad},M}
 \longrightarrow0.                                 \label{eq:s7v95-joint-zero}
\end{equation}
Every fixed finite transition/atlas polynomial can be absorbed into the
bad term, while every positive power of \(\beta^{-1/2}\) times a fixed
power of \(M\) vanishes in the good term.
\end{proposition}

\begin{proof}
\Cref{thm:s7v95-good-KP} gives the good limit.  V79 proves
\(A_{\partial,M}\pi_M^{+R}\to0\) for \(\chi>1\).
Its proof also absorbs fixed cutoff and transition polynomials.  The
last good statement is the logarithm-versus-power estimate used in
\Cref{cor:s7v95-schedule}.
\end{proof}

\begin{firewall}[The mixture still needs exact conditional ownership]
\label{fire:s7v95-mixture}
Equation \eqref{eq:s7v95-joint-zero} compares numerical activities.
To load CFTC, one must construct a single normalized conditional
expansion in which maximal bad components are retained exactly, every
crossing support is good by V79 saturation, and every remaining good
factor is measured relative to the Gaussian reference used above.
\end{firewall}

\section{Reference determinant and coarse ownership}
\label{sec:s7v95-determinant}

Changing variables in the Gaussian quadratic term produces the local
normalization
\begin{equation}
 Z_{\rm G}(y)
 =
 (2\pi/\beta)^{d/2}(\det L_y)^{-1/2}                \label{eq:s7v95-ZG}
\end{equation}
before cutoff tails.

\begin{proposition}[Gaussian normalization is a coarse action row]
\label{prop:s7v95-determinant}
The contribution of the conditional Gaussian normalization to the
coarse effective action is
\begin{equation}
 -\log Z_{\rm G}(y)
 =
 {1\over2}\log\det L_y
 +{d\over2}\log{\beta\over2\pi}.                    \label{eq:s7v95-logdet}
\end{equation}
The second term is a pressure constant; the first is a retained-field
interaction and must be included in the complete coarse logarithm.  It
is not part of \(\zeta_{{\rm good},M}\).
\end{proposition}

\begin{proof}
Take the negative logarithm of
\eqref{eq:s7v95-ZG}.  The \(\beta\)-term is independent of \(y\),
whereas \(L_y\) depends on the retained background.
\end{proof}

\begin{firewall}[A determinant can carry the slow interaction]
\label{fire:s7v95-determinant}
Small nonlinear remainder does not imply a small complete coarse
action.  The mixed oscillation and locality of
\(\frac12\log\det L_y\), together with the coarse reference
pushforward, must satisfy the V92 recursion.  Omitting this row would
discard the principal Gaussian response highlighted by the V95 audit.
\end{firewall}

\section{Near-flat conditional Gaussian card}
\label{sec:s7v95-card}

\begin{definition}[Near-flat conditional Gaussian card, NCGC]
\label{def:s7v95-NCGC}
NCGC consists of:
\begin{enumerate}[label=\textup{(N\arabic*)},leftmargin=4em]
\item the exact tree--chord Haar fibre of V94;
\item a stationary conditional center and uniform eliminated Hessian
      floor \eqref{eq:s7v95-Hessian};
\item the complete uniform Taylor and Haar estimates
      \eqref{eq:s7v95-remainder} and
      \eqref{eq:s7v95-Haar-log};
\item fixed-radius ownership of all good remainders;
\item exact maximal bad-component retention and saturated boundary
      ownership;
\item explicit ownership of the Gaussian determinant, cutoff tail, and
      principal reference variations in the coarse action;
\item a single normalized conditional good/bad KP expansion; and
\item Gauss, transition, reflection, and cofinal uniformity.
\end{enumerate}
\end{definition}

\begin{theorem}[NCGC loads the conditional KP part of CFTC]
\label{thm:s7v95-NCGC}
Under NCGC and the logarithmic schedule with \(\chi>1\), the scalar
good and saturated bad activities tend to zero as in
\eqref{eq:s7v95-joint-zero}.  Hence the bounded-remainder part of the
conditional fibre satisfies a strict KP margin for sufficiently large
\(\beta\).  The determinant and principal rows remain in the typed
coarse/reference kernels required by CFTC.
\end{theorem}

\begin{proof}
Use
\Cref{thm:s7v95-normalized,cor:s7v95-schedule,
thm:s7v95-good-KP,prop:s7v95-joint,
prop:s7v95-determinant}.
NCGC (N7)--(N8) supplies the exact common expansion and physical
uniformity needed to apply the conditional KP theorem
\Cref{thm:s7v93-KP}.
\end{proof}

\section{Status and next acquisition}
\label{sec:s7v95-status}

\begin{theorem}[What V95 proves]
\label{thm:s7v95-result}
V95 proves that, after centering at a stationary conditional background
and absorbing the full Hessian, the good-fibre Wilson/Haar remainder has
normalized factor activity
\[
 O\!\left(
 \beta^{-1/2}M^{3/2}+\beta^{-1}M^2
 \right)
\]
up to fixed uniform constants.  With the V79 saturated bad activity,
both parts vanish under \(M=(\log\beta)^\chi\), \(\chi>1\).
It also proves that the Gaussian log determinant is a separate complete
coarse-action row and cannot be hidden in that small activity.
\end{theorem}

\begin{proof}
Use
\Cref{lem:s7v95-Wilson,lem:s7v95-Haar,
thm:s7v95-normalized,thm:s7v95-good-KP,
prop:s7v95-joint,prop:s7v95-determinant,
thm:s7v95-NCGC}.
\end{proof}

\begin{assessment}[Progress at seventh-series V95]
\label{ass:s7v95-status}
The weak-coupling conditional activity now has the correct scaling and
type separation.  The bottleneck has moved to two concrete uniform
rows: prove the stationary eliminated Hessian/Taylor card for the
physical block map, and control the retained mixed oscillation of the
Gaussian determinant and coarse reference.
\end{assessment}

\begin{decision}[V96 acquisition order]
\label{dec:s7v95-V96}
V96 will compute derivatives of
\(\frac12\log\det L_y\):
\[
 D\log\det L=\operatorname{Tr}(L^{-1}DL),\qquad
 D^2\log\det L=
 \operatorname{Tr}(L^{-1}D^2L
 -L^{-1}DL\,L^{-1}DL).
\]
It will derive a support-resolved weighted row bound using the V53
exponential inverse kernel and identify any zero-mode or volume trace
which prevents regulator-uniform coarse locality.
\end{decision}

\begin{firewall}[Claim boundary after seventh-series V95]
\label{fire:s7v95-claim}
V95 does not prove NCGC for the full physical block map, the
regulator-uniform Hessian/Taylor data, determinant mixed-oscillation
row, CFTC along the ultraviolet trajectory, terminal entry, the
complete nonlinear coarse-action cocycle, BSCTC through cutoff removal,
CSMC, MIRC, or ARDPC, regulator-uniform SATC or full RL4C, interacting
endpoint continuity, UCM, BCC, pairing/Bogoliubov control, regulator
independence, continuum OS reconstruction, or the full Yang--Mills mass
gap.
\end{firewall}

\paragraph{Parent-version record.}
V95 was formed from
\texttt{Renewal\_Yang\_Mills\_Mass\_Gap\_Research\_Notes\_V94.tex}.
The V94 parent had \(42\,788\) counted source lines,
\(1\,529\,909\) bytes, and SHA--256
\[
 \texttt{183043E50D2360566EC9DDBA206510069EC38ED88ABD996FC29E069F7B084495}.
\]
The next compulsory companion audit is V100.

% =============================================================================
% END APPENDED SEVENTH-SERIES V95 MATERIAL
% =============================================================================

% =============================================================================
% BEGIN APPENDED SEVENTH-SERIES V96 MATERIAL
% =============================================================================

\chapter{The Gaussian log determinant and its weighted retained-field
row}
\label{ch:s7v96}

\section{Purpose}
\label{sec:s7v96-purpose}

V95 isolated \(\frac12\log\det L_y\) from the small nonlinear
remainder.  We now prove its exact derivative formulas and a
support-resolved locality bound.  The determinant is quasi-local when
the eliminated Hessian has a uniform gap and exponentially localized
inverse.  If its first and second retained-field derivatives carry the
expected weak-field factors \(O(g)\) and \(O(g^2)\), its mixed row
vanishes under the logarithmic schedule.

\section{Finite-regulator differentiation}
\label{sec:s7v96-derivatives}

Let \(L(y)\) be a \(C^2\) family of positive matrices on the
finite-regulator eliminated physical space, with
\begin{equation}
 L(y)\geq\lambda I,\qquad\lambda>0.                 \label{eq:s7v96-gap}
\end{equation}
Put \(G=L^{-1}\), \(B_x=D_xL\), and
\(C_{xy}=D_yD_xL\).

\begin{theorem}[Log-determinant derivatives]
\label{thm:s7v96-derivatives}
\begin{align}
 D_x\log\det L
 &=\operatorname{Tr}(GB_x),                         \label{eq:s7v96-first}\\
 D_yD_x\log\det L
 &=\operatorname{Tr}(GC_{xy})
 -\operatorname{Tr}(GB_yGB_x).                     \label{eq:s7v96-second}
\end{align}
\end{theorem}

\begin{proof}
For an invertible matrix,
\[
 D\det L=(\det L)\operatorname{Tr}(L^{-1}DL),
\]
which proves \eqref{eq:s7v96-first}.  Differentiate once more and use
\(D_yG=-GB_yG\).
\end{proof}

\begin{remark}[Sign and ownership]
\label{rem:s7v96-sign}
The second term in \eqref{eq:s7v96-second} is the Gaussian response
covariance.  It is nonpositive when \(x=y\) and \(B_x\) is self-adjoint.
It is part of the complete retained Hessian, not a scalar normalization
constant.
\end{remark}

\section{Local derivative supports}
\label{sec:s7v96-support}

Let the eliminated mode space be decomposed into fixed-dimensional
block fibres.  For each retained block \(x\), let \(P_x\) project onto a
collar \(S_x\) of at most \(r_*\) eliminated modes, and assume
\begin{align}
 B_x&=P_xB_xP_x,& \|B_x\|&\leq b_1,                 \label{eq:s7v96-B}\\
 C_{xy}&=P_{xy}C_{xy}P_{xy},&
 \|C_{xy}\|&\leq b_2\,\mathbf1_{\{d(x,y)\leq R_2\}},
                                                               \label{eq:s7v96-C}
\end{align}
where \(P_{xy}\) has rank at most \(r_{**}\).
Assume the inverse has the off-diagonal block bound
\begin{equation}
 \|P_xGP_y\|
 \leq C_Ge^{-m_G(d(x,y)-2R_1)_+}.                  \label{eq:s7v96-G}
\end{equation}

\begin{theorem}[Mixed determinant locality]
\label{thm:s7v96-locality}
There is a fixed constant \(C_*\), depending only on
\(r_*,r_{**},C_G,R_1,R_2\), such that
\begin{equation}
 |D_yD_x\tfrac12\log\det L|
 \leq
 C_*
 \left[
  b_2\mathbf1_{\{d(x,y)\leq R_2\}}
  +b_1^2e^{-2m_G(d(x,y)-2R_1)_+}
 \right].                                           \label{eq:s7v96-locality}
\end{equation}
\end{theorem}

\begin{proof}
For the direct term,
\[
 |\operatorname{Tr}(GC_{xy})|
 \leq\operatorname{rank}(P_{xy})\|G\|\|C_{xy}\|
 \leq r_{**}\lambda^{-1}b_2
 \mathbf1_{\{d(x,y)\leq R_2\}}.
\]
For the response term, cyclically restrict the trace:
\[
 \operatorname{Tr}(GB_yGB_x)
 =
 \operatorname{Tr}
 (P_xGP_yB_yP_yGP_xB_x).
\]
The trace norm is at most rank \(P_x\) times the product of operator
norms.  Use \eqref{eq:s7v96-B} and two copies of
\eqref{eq:s7v96-G}.  Divide the sum of the two estimates by two.
\end{proof}

\begin{corollary}[Weighted retained row]
\label{cor:s7v96-row}
On a polynomial-growth block lattice, for every
\(0<\mu<2m_G\),
\begin{equation}
 \sup_x\sum_y e^{\mu d(x,y)}
 |D_yD_x\tfrac12\log\det L|
 \leq
 C_\mu(b_2+b_1^2),                                  \label{eq:s7v96-row}
\end{equation}
with \(C_\mu\) independent of total volume.
\end{corollary}

\begin{proof}
The direct term has fixed range.  The second term is summable because
the number of blocks at distance \(n\) grows polynomially while
\(e^{-(2m_G-\mu)n}\) decays exponentially.
\end{proof}

\section{Inverse locality from the eliminated gap}
\label{sec:s7v96-inverse}

\begin{lemma}[Gapped finite-range inverse]
\label{lem:s7v96-inverse}
Suppose \(L\geq\lambda I\) has interaction range \(R_0\) and
off-diagonal row norm at most \(J\).  Then for every
\begin{equation}
 0<\nu\leq {1\over R_0}
 \log\left(1+{\lambda\over2J}\right),                \label{eq:s7v96-nu}
\end{equation}
there is \(C_G\leq2/\lambda\) such that
\begin{equation}
 \|G_{uv}\|\leq C_Ge^{-\nu d(u,v)}.                 \label{eq:s7v96-inverse-decay}
\end{equation}
\end{lemma}

\begin{proof}
The V53 weighted-conjugation argument gives
\(\|(e^{-tL})_{uv}\|
\leq e^{-\nu d(u,v)}e^{-\lambda t/2}\).
Since
\[
 L^{-1}=\int_0^\infty e^{-tL}\,dt,
\]
integration gives \(2\lambda^{-1}e^{-\nu d(u,v)}\).
\end{proof}

\begin{corollary}[Uniform eliminated gap loads the inverse hypothesis]
\label{cor:s7v96-inverse}
Uniform bounds on \(\lambda,R_0,J\) imply
\eqref{eq:s7v96-G} with regulator-independent constants.
\end{corollary}

\begin{proof}
Apply \Cref{lem:s7v96-inverse} to the fixed collars \(P_x,P_y\), losing
only their fixed radii in the distance.
\end{proof}

\section{From Hessian derivatives to mixed oscillation}
\label{sec:s7v96-oscillation}

Assume every retained block coordinate lies in a geodesically connected
good set of diameter at most \(D_{\rm ret}\).

\begin{lemma}[Two-variable fundamental theorem]
\label{lem:s7v96-FTC}
For a \(C^2\) scalar function \(F\), its mixed block oscillation obeys
\begin{equation}
 \partial_{xy}F
 \leq
 D_{\rm ret}^2
 \sup_y\|D_yD_xF\|.                                 \label{eq:s7v96-FTC}
\end{equation}
\end{lemma}

\begin{proof}
Join the two choices of the \(x\)-coordinate and the two choices of the
\(y\)-coordinate by constant-speed minimizing paths.  Apply the
fundamental theorem of calculus once along each path.  The mixed
difference is a double integral of \(D_yD_xF\), and the product of path
lengths is at most \(D_{\rm ret}^2\).
\end{proof}

\begin{corollary}[Determinant mixed-oscillation row]
\label{cor:s7v96-oscillation}
Under the hypotheses of \Cref{cor:s7v96-row},
\begin{equation}
 \sup_x\sum_y e^{\mu d(x,y)}
 \partial_{xy}\!\left(\tfrac12\log\det L\right)
 \leq
 D_{\rm ret}^2C_\mu(b_2+b_1^2).                    \label{eq:s7v96-oscillation}
\end{equation}
\end{corollary}

\begin{proof}
Combine \Cref{lem:s7v96-FTC,cor:s7v96-row}.
\end{proof}

\section{Weak-field scaling}
\label{sec:s7v96-scaling}

Suppose the rescaled conditional Hessian has
\begin{equation}
 L_y=L_0+E_y,\qquad
 \|D_xE_y\|\leq C_1gP_1(M),\qquad
 \|D_yD_xE_y\|\leq C_2g^2P_2(M),                   \label{eq:s7v96-weak}
\end{equation}
where \(P_1,P_2\) are fixed polynomials and the derivative supports are
as above.

\begin{theorem}[Vanishing determinant response row]
\label{thm:s7v96-vanishing}
Under the uniform gap and inverse-locality hypotheses,
\begin{equation}
 {\cal L}_{{\rm det},M}
 :=
 \sup_x\sum_y e^{\mu d(x,y)}
 \partial_{xy}\!\left(\tfrac12\log\det L_y\right)
 \leq
 C
 g^2\bigl(P_2(M)+P_1(M)^2\bigr).                   \label{eq:s7v96-det-rate}
\end{equation}
If \(g=\beta^{-1/2}\) and
\(M=(\log\beta)^\chi\), then
\begin{equation}
 {\cal L}_{{\rm det},M}\longrightarrow0.            \label{eq:s7v96-det-zero}
\end{equation}
\end{theorem}

\begin{proof}
In \eqref{eq:s7v96-oscillation}, take
\(b_1=C_1gP_1(M)\) and
\(b_2=C_2g^2P_2(M)\).  This gives
\eqref{eq:s7v96-det-rate}.  A fixed logarithmic power is dominated by
\(\beta\), proving \eqref{eq:s7v96-det-zero}.
\end{proof}

\begin{firewall}[The derivative scaling is a physical hypothesis]
\label{fire:s7v96-scaling}
The \(O(g)\) and \(O(g^2)\) rows in
\eqref{eq:s7v96-weak} must be derived from the actual conditional
Wilson Hessian after gauge reduction and moving-projector correction.
They do not follow merely from \(L_y-L_0\to0\) in operator norm.
\end{firewall}

\section{Zero modes and the trace obstruction}
\label{sec:s7v96-zero}

\begin{theorem}[Uniform determinant locality fails without an
eliminated gap]
\label{thm:s7v96-zero}
For a translation-invariant scalar model with
\(L_m=-\Delta+m^2\), the inverse kernel loses every fixed exponential
decay rate as \(m\downarrow0\).  In four dimensions its massless
response product behaves schematically as
\begin{equation}
 G_0(x,y)^2\asymp |x-y|^{-4},                       \label{eq:s7v96-massless-square}
\end{equation}
whose absolute row has a logarithmic infrared divergence.
\end{theorem}

\begin{proof}
The Fourier multiplier is
\((\widehat\Delta(p)+m^2)^{-1}\); its nearest complex singularity
approaches the real axis with \(m\), so no uniform exponential rate is
possible.  In four continuum dimensions
\(G_0(x,y)\asymp|x-y|^{-2}\).  Squaring and summing over a shell of
radius \(r\), which has size \(O(r^3)\), gives \(O(r^{-1})\); its radial
sum diverges logarithmically.
\end{proof}

\begin{firewall}[Only the eliminated physical subspace may enter]
\label{fire:s7v96-zero}
Gauge directions, torons, and retained slow curvature modes must be
projected out before \(\det L_y\) is formed.  Their determinant belongs
neither to the gapped eliminated Gaussian nor to
\Cref{thm:s7v96-locality}.  A regulator-dependent pseudo-determinant
with a closing gap would invalidate the claimed uniform row.
\end{firewall}

\section{Determinant locality card}
\label{sec:s7v96-card}

\begin{definition}[Eliminated determinant locality card, EDLC]
\label{def:s7v96-EDLC}
EDLC requires:
\begin{enumerate}[label=\textup{(D\arabic*)},leftmargin=4em]
\item a fixed-rank local decomposition of the eliminated physical
      Hessian, with all zero/slow modes removed;
\item a regulator-uniform spectral gap, range, and off-diagonal row;
\item local first and second retained derivatives
      \eqref{eq:s7v96-B}--\eqref{eq:s7v96-C};
\item the weak-field scaling \eqref{eq:s7v96-weak};
\item moving-projector and transition-chart derivatives included once;
\item a fixed retained good-set diameter; and
\item compatibility with the complete coarse logarithm and cutoff
      removal.
\end{enumerate}
\end{definition}

\begin{theorem}[EDLC closes the V95 determinant row]
\label{thm:s7v96-EDLC}
Under EDLC, the Gaussian determinant has a regulator-uniform
quasi-local mixed-oscillation row and that row tends to zero under the
logarithmic weak-coupling schedule.  It may therefore be included in
the retained part of the V92/CFTC recursion without spoiling a strict
small-activity margin.
\end{theorem}

\begin{proof}
Use
\Cref{thm:s7v96-derivatives,thm:s7v96-locality,
lem:s7v96-inverse,cor:s7v96-oscillation,
thm:s7v96-vanishing}.
\end{proof}

\section{Status and next acquisition}
\label{sec:s7v96-status}

\begin{theorem}[What V96 proves]
\label{thm:s7v96-result}
V96 proves the exact first and second derivatives of the Gaussian log
determinant, exponential locality of its response under a gapped
finite-range eliminated Hessian, and the row bound
\[
 {\cal L}_{{\rm det},M}
 \leq
 Cg^2(P_2(M)+P_1(M)^2).
\]
It also proves that an unremoved massless mode destroys the required
uniform exponential determinant locality.
\end{theorem}

\begin{proof}
Use
\Cref{thm:s7v96-derivatives,thm:s7v96-locality,
cor:s7v96-row,thm:s7v96-vanishing,
thm:s7v96-zero,thm:s7v96-EDLC}.
\end{proof}

\begin{assessment}[Progress at seventh-series V96]
\label{ass:s7v96-status}
The principal Gaussian normalization is now quantitatively compatible
with the coarse interaction norm, conditional on the physical
weak-field derivative scaling.  The remaining upstream task is to
prove that scaling and the uniform eliminated Hessian floor for the
actual tree--chord Wilson block.
\end{assessment}

\begin{decision}[V97 acquisition order]
\label{dec:s7v96-V97}
V97 will compute the Wilson plaquette Hessian in tree--chord exponential
coordinates around a stationary flat background.  It will identify its
gauge kernel, prove the transverse quadratic form, and estimate the
first two retained-background derivatives after the
\(\eta=\sqrt\beta L^{1/2}\xi\) rescaling.  The target is EDLC (D1)--(D5).
\end{decision}

\begin{firewall}[Claim boundary after seventh-series V96]
\label{fire:s7v96-claim}
V96 does not prove EDLC or the weak-field derivative scaling for the
physical Wilson block, the regulator-uniform eliminated Hessian floor,
NCGC/CFTC along the ultraviolet trajectory, terminal entry, the
complete nonlinear coarse-action cocycle, BSCTC through cutoff removal,
CSMC, MIRC, or ARDPC, regulator-uniform SATC or full RL4C, interacting
endpoint continuity, UCM, BCC, pairing/Bogoliubov control, regulator
independence, continuum OS reconstruction, or the full Yang--Mills mass
gap.
\end{firewall}

\paragraph{Parent-version record.}
V96 was formed from
\texttt{Renewal\_Yang\_Mills\_Mass\_Gap\_Research\_Notes\_V95.tex}.
The V95 parent had \(43\,300\) counted source lines,
\(1\,547\,404\) bytes, and SHA--256
\[
 \texttt{9E329756E55FFD92870F5D64781C2142AA0069A4AE3C03B155AF0778241C20DB}.
\]
The next compulsory companion audit is V100.

% =============================================================================
% END APPENDED SEVENTH-SERIES V96 MATERIAL
% =============================================================================

\clearpage
% =============================================================================
% BEGIN APPENDED SEVENTH-SERIES V97 MATERIAL
% =============================================================================

\part{Seventh-series V97: Wilson transverse Hessian and physical-field
derivative powers}
\label{part:s7v97}

\section{Purpose and dependency}
\label{sec:s7v97-purpose}

V96 reduced the Gaussian-normalization problem to a concrete local
question.  If \(L_y\) is the eliminated Wilson Hessian in a retained
background \(y\), one must prove
\[
 L_y\geq\lambda I,\qquad
 \|\partial_{y_z}L_y\|=O(g),\qquad
 \|\partial_{y_w}\partial_{y_z}L_y\|=O(g^2),
\]
with local supports and constants uniform in the regulator.  These
powers cannot be inserted as a generic smoothness hypothesis.  This
note obtains them in a fixed tree--chord weak-field chart by solving
the stationary equation in the physical retained field
\(u=gy\).  It also identifies exactly which zero modes have to be
removed before the determinant is formed.

\begin{scope}[What is and is not proved in V97]
\label{scope:s7v97}
All assertions below concern a fixed finite block, or a uniformly
finite family of block types, inside one injectivity chart of
\(\mathrm{SU}(N)\).  The constants are independent of the ultraviolet
regulator when the combinatorial block types and coordinate
normalizations are fixed.  V97 does not assert a uniform Poincare
constant for blocks whose graph diameter tends to infinity, nor does
it cover non-flat stationary branches.
\end{scope}

\section{Linearized Wilson form on a finite two-complex}
\label{sec:s7v97-linear}

Let \(B=(V,E,P)\) be a finite oriented two-complex.  Choose one
orientation of every edge and plaquette.  Write
\[
 C^k_B=C^k(B;\mathfrak{su}(N)),\qquad k=0,1,2,
\]
with the counting inner product induced by
\(\langle X,Y\rangle_{\mathfrak g}=-\operatorname{Re}\operatorname{Tr}(XY)\).
The cellular coboundaries are denoted by
\[
 d_0:C^0_B\longrightarrow C^1_B,\qquad
 d_1:C^1_B\longrightarrow C^2_B.
\]
Thus \(d_1d_0=0\).

For \(A\in C^1_B\), let \(U_e(t)=\exp(tA_e)\), with the inverse link
used when the boundary orientation is opposite to the chosen edge
orientation.  Let \(\Omega_p(tA)\) be the ordered plaquette holonomy
and put
\[
 w(A)=\sum_{p\in P}\left(
 N-\operatorname{Re}\operatorname{Tr}\Omega_p(A)\right).
                                                        \tag{97.1}\label{eq:s7v97-action}
\]

\begin{lemma}[Second variation of one unitary loop]
\label{lem:s7v97-loop}
If \(Q:(-\epsilon,\epsilon)\to\mathrm{SU}(N)\) is \(C^2\),
\(Q(0)=I\), and \(Q'(0)=X\in\mathfrak{su}(N)\), then
\[
 \left.\frac{d}{dt}\right|_{t=0}
 \bigl(N-\operatorname{Re}\operatorname{Tr}Q(t)\bigr)=0,
\qquad
 \left.\frac{d^2}{dt^2}\right|_{t=0}
 \bigl(N-\operatorname{Re}\operatorname{Tr}Q(t)\bigr)
 =\|X\|_{\mathfrak g}^2.                              \tag{97.2}
\]
\end{lemma}

\begin{proof}
Differentiate \(Q(t)^*Q(t)=I\) twice at zero.  The first derivative
gives \(X^*=-X\).  The second gives
\[
 Q''(0)^*+Q''(0)+2X^*X=0.
\]
Taking real traces yields
\(
 -\operatorname{Re}\operatorname{Tr}Q''(0)
 =\operatorname{Tr}(X^*X)=\|X\|_{\mathfrak g}^2.
\)
The first identity follows from
\(\operatorname{Re}\operatorname{Tr}X=0\).
\end{proof}

\begin{lemma}[Linearized plaquette holonomy]
\label{lem:s7v97-plaquette}
For every \(p\in P\),
\[
 \left.\frac{d}{dt}\right|_{t=0}\Omega_p(tA)
 =(d_1A)_p.                                           \tag{97.3}
\]
Consequently,
\[
 w(tA)=\frac{t^2}{2}\|d_1A\|_{C^2_B}^2+O_A(t^3).
                                                        \tag{97.4}\label{eq:s7v97-expansion}
\]
\end{lemma}

\begin{proof}
At \(t=0\) every factor in the ordered product is \(I\).
The product rule therefore gives the signed sum of the edge
generators around \(p\), which is precisely \((d_1A)_p\).
Apply \Cref{lem:s7v97-loop} to each plaquette and sum.  Analyticity of
matrix multiplication and the exponential supplies the remainder.
\end{proof}

\begin{theorem}[Exact flat-background Wilson Hessian]
\label{thm:s7v97-hessian}
The Hessian of \eqref{eq:s7v97-action} at the trivial connection is
\[
 D^2w(0)=d_1^*d_1.                                   \tag{97.5}
\]
In particular,
\[
 D^2w(0)[A,A]=\|d_1A\|^2,\qquad
 \operatorname{im}d_0\subseteq\ker D^2w(0).           \tag{97.6}\label{eq:s7v97-gauge-kernel}
\]
\end{theorem}

\begin{proof}
Equation \eqref{eq:s7v97-expansion} identifies the diagonal quadratic form.
Polarization identifies its symmetric bilinear form with
\(\langle d_1A,d_1B\rangle\), hence with \(d_1^*d_1\).
If \(A=d_0\phi\), then \(d_1A=d_1d_0\phi=0\).
\end{proof}

\begin{firewall}[Gauge fixing alone need not remove every zero mode]
\label{fire:s7v97-cohomology}
The reverse inclusion in \eqref{eq:s7v97-gauge-kernel} is generally false.  Flat
one-cochains not in \(\operatorname{im}d_0\) represent discrete
cohomology, boundary transport, or toron directions.  They must be
retained as slow variables or otherwise fixed by an independently
justified boundary condition.  Putting them into the eliminated
determinant gives a zero eigenvalue.
\end{firewall}

\section{The eliminated transverse space}
\label{sec:s7v97-transverse}

Define the harmonic representative space and the transverse
eliminated space by
\[
 {\cal H}^1_B:=\ker d_1\cap(\operatorname{im}d_0)^\perp,
\qquad
 {\cal E}_B:=(\ker d_1)^\perp.
                                                        \tag{97.7}
\]
Thus the orthogonal decomposition is
\[
 C^1_B=\overline{\operatorname{im}d_0}
 \oplus{\cal H}^1_B\oplus{\cal E}_B,                  \tag{97.8}\label{eq:s7v97-decomp}
\]
where the overline is harmless in finite dimension.

\begin{lemma}[Finite-block transverse Poincare floor]
\label{lem:s7v97-floor}
There is a number \(\lambda_B>0\) such that
\[
 \|d_1A\|^2\geq\lambda_B\|A\|^2
 \quad\hbox{for every }A\in{\cal E}_B.                \tag{97.9}
\]
Equivalently,
\[
 P_{\cal E}d_1^*d_1P_{\cal E}\geq
 \lambda_BP_{\cal E}.                                \tag{97.10}
\]
\end{lemma}

\begin{proof}
The restriction of \(d_1\) to \({\cal E}_B=(\ker d_1)^\perp\) is
injective.  The continuous function
\(A\mapsto\|d_1A\|^2\) is therefore strictly positive on the compact
unit sphere of \({\cal E}_B\).  Its minimum is \(\lambda_B>0\).
\end{proof}

\begin{corollary}[Uniformity for a finite block dictionary]
\label{cor:s7v97-dictionary}
If every ultraviolet cell is assigned one of finitely many
two-complexes \(B_1,\ldots,B_J\), with the same normalized cochain
inner products, then
\[
 \lambda_*:=\min_{1\leq j\leq J}\lambda_{B_j}>0       \tag{97.11}
\]
is uniform over the regulator and the cell location.
\end{corollary}

\begin{proof}
Apply \Cref{lem:s7v97-floor} to a finite set.
\end{proof}

\begin{remark}[Tree--chord realization]
\label{rem:s7v97-tree}
The V94 spanning-tree change of variables removes the
\(\operatorname{im}d_0\) factor with unit product-Haar Jacobian.
Chord coordinates spanning \({\cal H}^1_B\), or boundary chord
coordinates that carry the corresponding transport, are retained.
The remaining chord combinations may be linearly identified with
\({\cal E}_B\).  Thus \eqref{eq:s7v97-decomp} is the infinitesimal form of the
actual V94 product-fibre chart, not an additional gauge quotient.
\end{remark}

\begin{firewall}[Growing blocks]
\label{fire:s7v97-growing}
Compactness proves \(\lambda_B>0\) for each fixed block but says
nothing uniform when \(B=B_\ell\) grows.  For Laplace-type complexes
one normally expects the first transverse eigenvalue to be of order
\(\ell^{-2}\).  Any cofinal construction with growing block diameter
must expose that loss in the determinant and polymer rows.
\end{firewall}

\section{Stationary elimination in physical coordinates}
\label{sec:s7v97-stationary}

Let a tree--chord chart split the non-gauge link coordinates as
\[
 a=x+y,\qquad x\in{\cal E},\quad y\in{\cal R},         \tag{97.12}\label{eq:s7v97-split}
\]
where \({\cal R}\) contains all retained chord, harmonic, and boundary
coordinates.  Let \(\iota_{\cal E}\) and \(\iota_{\cal R}\) denote the
corresponding linear injections into \(C^1_B\).  In physical
coordinates \(b_x\in{\cal E}\), \(u\in{\cal R}\), define
\[
 {\mathfrak w}(b_x,u)
 :=
 \sum_{p\in P}\left[
 N-\operatorname{Re}\operatorname{Tr}
 \Omega_p(\iota_{\cal E}b_x+\iota_{\cal R}u)
 \right].                                             \tag{97.13}
\]
Here each link is the exponential of its displayed coordinate; any
fixed ordering convention in a multi-coordinate link chart is
allowed.  The function is real analytic in a neighbourhood of zero.

The flat quadratic form in the eliminated variables is
\[
 L_0:=D^2_{b_xb_x}{\mathfrak w}(0,0)
 =\iota_{\cal E}^*d_1^*d_1\iota_{\cal E}.             \tag{97.14}
\]
By construction and \Cref{lem:s7v97-floor}, \(L_0\geq\lambda_BI\).

\begin{theorem}[Physical stationary-centre map]
\label{thm:s7v97-centre}
There are neighbourhoods \(0\in{\cal U}\subset{\cal R}\) and
\(0\in{\cal V}\subset{\cal E}\), and a unique real-analytic map
\[
 \widehat x:{\cal U}\longrightarrow{\cal V}           \tag{97.15}\label{eq:s7v97-xhat}
\]
such that
\[
 \widehat x(0)=0,\qquad
 D_{b_x}{\mathfrak w}(\widehat x(u),u)=0.              \tag{97.16}\label{eq:s7v97-stationary-eq}
\]
Its first two derivatives are bounded on every relatively compact
ball \(\overline{B_{\cal R}(0,r)}\subset{\cal U}\).  At zero,
\[
 D\widehat x(0)
 =-L_0^{-1}D^2_{ub_x}{\mathfrak w}(0,0).              \tag{97.17}\label{eq:s7v97-xhat-der}
\]
\end{theorem}

\begin{proof}
Set \(F(b_x,u)=D_{b_x}{\mathfrak w}(b_x,u)\).
The flat connection is stationary, so \(F(0,0)=0\), while
\(D_{b_x}F(0,0)=L_0\) is invertible.  The finite-dimensional analytic
implicit-function theorem gives
\eqref{eq:s7v97-xhat}--\eqref{eq:s7v97-stationary-eq} and
uniqueness after shrinking the neighbourhoods.  Differentiate
\(F(\widehat x(u),u)=0\) at zero to get
\eqref{eq:s7v97-xhat-der}.  Analyticity
gives bounded derivatives on compact subsets.
\end{proof}

\begin{definition}[Rescaled centre and centred Hessian]
\label{def:s7v97-rescaled}
For \(g>0\) and \(gy\in{\cal U}\), set
\[
 x_*(y,g):=g^{-1}\widehat x(gy),                       \tag{97.18}
\]
and define the weak-coupling action and its centred eliminated Hessian
by
\[
 {\mathfrak S}_g(x,y):=g^{-2}{\mathfrak w}(gx,gy),
\qquad
 L_{g,y}:=D^2_{xx}{\mathfrak S}_g(x_*(y,g),y).
                                                        \tag{97.19}\label{eq:s7v97-scaled-action}
\]
\end{definition}

\begin{lemma}[Cancellation of the apparent singular scaling]
\label{lem:s7v97-cancel}
Let
\[
 {\cal L}(u):=
 D^2_{b_xb_x}{\mathfrak w}(\widehat x(u),u).           \tag{97.20}\label{eq:s7v97-Lcal}
\]
Then
\[
 L_{g,y}={\cal L}(gy).                                \tag{97.21}\label{eq:s7v97-composition}
\]
Moreover,
\[
 D_yx_*(y,g)=D\widehat x(gy),\qquad
 D_y^2x_*(y,g)=gD^2\widehat x(gy).                    \tag{97.22}\label{eq:s7v97-center-der}
\]
\end{lemma}

\begin{proof}
Each \(x\)-derivative of
\(g^{-2}{\mathfrak w}(gx,gy)\) contributes one factor \(g\);
the two factors cancel \(g^{-2}\).  Substitution of
\(gx_*(y,g)=\widehat x(gy)\) gives
\eqref{eq:s7v97-composition}.
The chain rule gives \eqref{eq:s7v97-center-der}.
\end{proof}

\begin{theorem}[Wilson weak-field derivative powers]
\label{thm:s7v97-powers}
Choose \(r>0\) with
\(\overline{B_{\cal R}(0,r)}\subset{\cal U}\).
There are finite constants \(C_0,C_1,C_2\), depending only on the
chosen block and chart, such that whenever \(\|gy\|\leq r\),
\begin{align}
 \|L_{g,y}-L_0\|&\leq C_0\|gy\|,                      \tag{97.23}\label{eq:s7v97-Ldiff}\\
 \|D_yL_{g,y}\|&\leq C_1g,                            \tag{97.24}\label{eq:s7v97-DL}\\
 \|D_y^2L_{g,y}\|&\leq C_2g^2.                        \tag{97.25}\label{eq:s7v97-DDL}
\end{align}
If \(C_0\|gy\|\leq\lambda_B/2\), then
\[
 L_{g,y}\geq\frac{\lambda_B}{2}I.                    \tag{97.26}
\]
\end{theorem}

\begin{proof}
The map \({\cal L}\) in \eqref{eq:s7v97-Lcal} is analytic on
\({\cal U}\) and
\({\cal L}(0)=L_0\).  The mean-value formula on the physical ball
gives \eqref{eq:s7v97-Ldiff}.  Differentiating the exact identity
\eqref{eq:s7v97-composition} once and twice gives
\[
 D_yL_{g,y}=g(D{\cal L})(gy),\qquad
 D_y^2L_{g,y}=g^2(D^2{\cal L})(gy).
                                                        \tag{97.27}
\]
The first two derivatives of \({\cal L}\) are bounded on the compact
physical ball, proving
\eqref{eq:s7v97-DL}--\eqref{eq:s7v97-DDL}.  Finally,
the min--max principle and \eqref{eq:s7v97-Ldiff} give
\[
 \lambda_{\min}(L_{g,y})
 \geq\lambda_B-\|L_{g,y}-L_0\|
 \geq\lambda_B/2.
\]
\end{proof}

\begin{corollary}[Logarithmic good-set schedule]
\label{cor:s7v97-log}
Let \(g=\beta^{-1/2}\) and let
\(\|y\|\leq cM^{1/2}\), where
\(M=(\log\beta)^\chi\) with fixed \(\chi>0\).  Then, for all
sufficiently large \(\beta\),
\[
 \|gy\|\leq r,\qquad L_{g,y}\geq\lambda_B/2,
\qquad
 \|D_yL_{g,y}\|\leq C_1\beta^{-1/2},\quad
 \|D_y^2L_{g,y}\|\leq C_2\beta^{-1}.                  \tag{97.28}
\]
\end{corollary}

\begin{proof}
Every fixed power of \(\log\beta\) is \(o(\beta)\), so
\(\beta^{-1/2}M^{1/2}\to0\).  Apply
\Cref{thm:s7v97-powers}.
\end{proof}

\begin{remark}[Why centering does not lose a power of \(g\)]
\label{rem:s7v97-centering}
The derivative \(D_yx_*\) can be of order one; indeed
\eqref{eq:s7v97-xhat-der} need not vanish.  The powers in
\eqref{eq:s7v97-DL}--\eqref{eq:s7v97-DDL} survive because the Hessian is an
analytic function of the \emph{physical} background \(u=gy\).
Furthermore \(D_y^2x_*=O(g)\) by
\eqref{eq:s7v97-center-der}.  Differentiating a
stationary centre directly in the rescaled variables without this
physical-coordinate identity obscures the cancellation and can lead
to a circular \(O(1)\) estimate.
\end{remark}

\section{Local supports and the fixed-block row}
\label{sec:s7v97-local}

Index retained coordinates by \(z\) and write
\[
 B_z(g,y):=\partial_{y_z}L_{g,y},\qquad
 C_{zw}(g,y):=\partial_{y_w}\partial_{y_z}L_{g,y}.
                                                        \tag{97.29}
\]

\begin{lemma}[Conditional block support]
\label{lem:s7v97-support}
Suppose the fine lattice is partitioned into conditional interiors
whose Wilson actions share only retained boundary variables.  Solving
\eqref{eq:s7v97-stationary-eq} independently in each interior has the following
properties.
\begin{enumerate}[label=\textup{(\roman*)}]
\item \(B_z\) vanishes on every eliminated block not incident to
      the retained coordinate \(z\).
\item \(C_{zw}\) vanishes unless some eliminated block is incident
      to both \(z\) and \(w\).
\item If the block dictionary and boundary incidence are uniformly
      finite, the support radii and support multiplicities in V96
      are regulator-uniform.
\end{enumerate}
\end{lemma}

\begin{proof}
The full conditional action is a sum of block actions, and its
eliminated variables are disjoint between interiors.  The stationary
equation for one block therefore depends only on the retained
coordinates incident to that block.  The analytic solution
\(\widehat x_B\), its Hessian \({\cal L}_B\), and all of their
derivatives have the same dependence.  This proves (i) and (ii).
The maxima over a finite dictionary and finite incidence number prove
(iii).
\end{proof}

\begin{theorem}[Physical realization of the V96 EDLC derivative rows]
\label{thm:s7v97-EDLCrows}
For a uniformly finite tree--chord block dictionary, on the
logarithmic good set and for sufficiently large \(\beta\), the actual
centred Wilson Hessian satisfies the V96 assumptions
\[
 \lambda I\leq L_{g,y},\qquad
 \|B_z(g,y)\|\leq Cg,\qquad
 \|C_{zw}(g,y)\|\leq Cg^2,                            \tag{97.30}\label{eq:s7v97-EDLC-row}
\]
with fixed local support radius and multiplicity.  Here one may take
\(\lambda=\lambda_*/2\).  Therefore the V96 determinant-response
estimate specializes to
\[
 {\cal L}_{{\rm det},M}\leq C_{\rm det}g^2
 =C_{\rm det}\beta^{-1}.                              \tag{97.31}
\]
\end{theorem}

\begin{proof}
Use \Cref{cor:s7v97-dictionary,cor:s7v97-log,
lem:s7v97-support}.  Equations
\eqref{eq:s7v97-DL}--\eqref{eq:s7v97-DDL} imply the coordinatewise
bounds in \eqref{eq:s7v97-EDLC-row}.  Apply
\Cref{thm:s7v96-locality,cor:s7v96-row,
thm:s7v96-vanishing} with constant polynomials \(P_1,P_2\).
\end{proof}

\begin{firewall}[What locality of the stationary solve means]
\label{fire:s7v97-locality}
Inside one block, \(D\widehat x\) can be dense: an incident retained
boundary coordinate may move every eliminated link in that block.
The support statement is block-local, not plaquette-local.  It is
uniform only while block diameter and the block dictionary stay
uniformly bounded.  For growing blocks one must replace this argument
by an inverse-decay estimate and retain its scale dependence.
\end{firewall}

\section{Gauge covariance and chart ownership}
\label{sec:s7v97-gauge}

\begin{lemma}[Equivariance of a unique stationary branch]
\label{lem:s7v97-equivariance}
Let a residual compact group \(K\) act linearly and isometrically on
\({\cal E}\oplus{\cal R}\), preserve the splitting, and leave
\({\mathfrak w}\) invariant.  If the stationary solution in
\Cref{thm:s7v97-centre} is unique in a \(K\)-invariant neighbourhood,
then
\[
 \widehat x(k u)=k\widehat x(u)\qquad(k\in K).         \tag{97.32}\label{eq:s7v97-equivariance}
\]
Consequently \(L_{g,ky}=kL_{g,y}k^{-1}\), so
\(\det L_{g,y}\) is residual-gauge invariant.
\end{lemma}

\begin{proof}
Invariance and preservation of the splitting imply that
\(k\widehat x(u)\) solves the eliminated stationary equation at
\(ku\).  Uniqueness gives \eqref{eq:s7v97-equivariance}.  Differentiate the invariant
action twice at the corresponding centres to obtain conjugacy of the
Hessians.  Determinants are invariant under conjugacy.
\end{proof}

\begin{proposition}[No moving projector inside one tree chart]
\label{prop:s7v97-fixed-slice}
In the V94 tree--chord chart, the gauge variables are integrated out
with product Haar measure and the vector spaces
\({\cal E}\) and \({\cal R}\) in \eqref{eq:s7v97-split} are fixed coordinate
spaces.  Thus the Hessian in \eqref{eq:s7v97-scaled-action} is already the physical
eliminated Hessian on a fixed slice; no derivative of a
background-dependent orthogonal projector is omitted in
\eqref{eq:s7v97-DL}--\eqref{eq:s7v97-DDL}.
\end{proposition}

\begin{proof}
V94 gives an exact coordinate bijection between link variables and
root gauge frames plus chord variables, with product Haar Jacobian.
Choosing a fixed subset of chord coordinates as retained makes the
remaining coordinate subspace fixed.  Differentiation in \(y\) is
therefore performed after, not before, this exact coordinate
reduction.  There is no \(P_y\) in
\eqref{eq:s7v97-scaled-action}.
\end{proof}

\begin{firewall}[Chart transitions remain an owned row]
\label{fire:s7v97-transition}
\Cref{prop:s7v97-fixed-slice} removes moving-projector terms within
the flat chart.  It does not prove that a cofinal good set is covered
by this one chart, or that transition Jacobians between several
tree--chord/exponential charts have the same small mixed-oscillation
row.  Those transition terms must be included exactly once in the
complete coarse logarithm.
\end{firewall}

\section{Wilson transverse Hessian card}
\label{sec:s7v97-card}

\begin{definition}[Wilson transverse Hessian card, WTHC]
\label{def:s7v97-WTHC}
A conditional block passes WTHC on a retained set \(G\) if:
\begin{enumerate}[label=\textup{(W\arabic*)},leftmargin=4em]
\item an exact tree--chord product-Haar chart removes gauge frames;
\item every flat cohomology, toron, and boundary transport direction
      is retained rather than put in the determinant;
\item the remaining flat quadratic form has a stated floor
      \(\lambda_B>0\);
\item the physical stationary equation has a unique branch
      \(\widehat x(u)\) for all \(u=gy\), \(y\in G\);
\item the centred Hessian is exactly
      \(L_{g,y}={\cal L}(gy)\);
\item first and second retained derivatives have the \(g,g^2\)
      powers and a stated support geometry;
\item any growing-block deterioration of the floor, support radius,
      or derivative constants is exposed; and
\item transition-chart and cutoff-boundary rows have explicit owners.
\end{enumerate}
\end{definition}

\begin{theorem}[Flat fixed-block WTHC]
\label{thm:s7v97-WTHC}
For a fixed finite block in the trivial-connection injectivity chart,
items \textup{(W1)--(W6)} of WTHC hold on
\[
 G_{\beta,M}:=\{y:\|y\|\leq cM^{1/2}\},
 \qquad g=\beta^{-1/2},\quad M=(\log\beta)^\chi,
                                                        \tag{97.33}
\]
for all sufficiently large \(\beta\).  For a finite block dictionary
the constants are uniform.  Items \textup{(W7)--(W8)} are firewalls,
not automatic consequences.
\end{theorem}

\begin{proof}
Item (W1) is V94 and
\Cref{prop:s7v97-fixed-slice}.  Item (W2) is the decomposition
\eqref{eq:s7v97-decomp}.  Item (W3) is
\Cref{lem:s7v97-floor}.  Item (W4) is
\Cref{thm:s7v97-centre} and the fact that
\(gM^{1/2}\to0\).  Items (W5)--(W6) are
\Cref{lem:s7v97-cancel,thm:s7v97-powers,
lem:s7v97-support}.  Uniformity follows from taking maxima and minima
over the finite dictionary.
\end{proof}

\section{Status and next acquisition}
\label{sec:s7v97-status}

\begin{theorem}[What V97 proves]
\label{thm:s7v97-result}
In an actual finite tree--chord Wilson block, V97 proves:
\begin{enumerate}[label=\textup{(\alph*)}]
\item the flat Hessian is \(d_1^*d_1\);
\item its kernel consists of gauge gradients together with flat
      cohomological directions, all of which are excluded from the
      eliminated determinant;
\item the transverse eliminated form has a positive finite-block
      floor;
\item the physical stationary centre exists uniquely near the flat
      branch;
\item the centred Hessian has the exact composition form
      \(L_{g,y}={\cal L}(gy)\); and
\item its first two retained derivatives are \(O(g)\) and \(O(g^2)\),
      giving an \(O(\beta^{-1})\) Gaussian determinant response row
      for a uniformly finite block dictionary.
\end{enumerate}
\end{theorem}

\begin{proof}
Combine
\Cref{thm:s7v97-hessian,lem:s7v97-floor,
thm:s7v97-centre,lem:s7v97-cancel,
thm:s7v97-powers,thm:s7v97-EDLCrows}.
\end{proof}

\begin{assessment}[Progress at seventh-series V97]
\label{ass:s7v97-status}
The V96 weak-field derivative hypothesis is no longer an assumption
on a fixed flat Wilson block: it follows from the physical-coordinate
stationary map.  The unresolved scale question is now sharply
separated.  A fixed block dictionary gives a uniform transverse floor
and finite-range response, whereas a growing block can lose both.
The next note must decide which geometry the cofinal ultraviolet
trajectory actually requires and propagate the resulting scale
factors through the determinant and cluster margins.
\end{assessment}

\begin{decision}[V98 acquisition order]
\label{dec:s7v97-V98}
V98 will compare fixed-physical-block and growing-lattice-block
schemes.  It will compute the scale dependence of the discrete
transverse Poincare floor, inverse range, retained derivative row, and
Gaussian determinant response.  The purpose is to determine whether
the present terminal schedule closes with a fixed block dictionary or
whether a multiscale Schur complement is compulsory.
\end{decision}

\begin{firewall}[Claim boundary after seventh-series V97]
\label{fire:s7v97-claim}
V97 does not prove a regulator-uniform WTHC for growing blocks, a
global stationary-branch cover, transition-chart control, NCGC/CFTC
along the complete ultraviolet trajectory, terminal entry, the
complete nonlinear coarse-action cocycle, BSCTC through cutoff
removal, CSMC, MIRC, or ARDPC, regulator-uniform SATC or full RL4C,
interacting endpoint continuity, UCM, BCC, pairing/Bogoliubov
control, regulator independence, continuum OS reconstruction, or the
full Yang--Mills mass gap.
\end{firewall}

\paragraph{Parent-version record.}
V97 was formed from
\texttt{Renewal\_Yang\_Mills\_Mass\_Gap\_Research\_Notes\_V96.tex}.
The V96 parent had \(43\,696\) counted source lines,
\(1\,560\,740\) bytes, and SHA--256
\[
 \texttt{C14B8DBD6CF49735BC2B62D12994C3673FEBBE76D47B8F82D81CBF61D4EEB8DA}.
\]
The next compulsory companion audit is V100.

% =============================================================================
% END APPENDED SEVENTH-SERIES V97 MATERIAL
% =============================================================================

\clearpage
% =============================================================================
% BEGIN APPENDED SEVENTH-SERIES V98 MATERIAL
% =============================================================================

\part{Seventh-series V98: fixed-factor cofinal blocking versus one-shot
growing fibres}
\label{part:s7v98}

\section{Question left by V97}
\label{sec:s7v98-question}

V97 proved the Wilson transverse-Hessian card on a fixed finite block
and warned that the Poincare floor can close on a growing block.  The
cofinal continuum trajectory does require a physical distance of
order one to contain more and more ultraviolet links.  It does
\emph{not}, however, require those links to be eliminated in one
conditional integral.  This note separates two constructions:
\begin{enumerate}[label=\textup{(\roman*)}]
\item a one-shot fibre whose lattice diameter
      \(\ell(a)\asymp a^{-1}\) grows; and
\item an iteration of a fixed integer blocking factor \(b\), whose
      number of steps grows while every elementary fibre belongs to a
      finite dictionary.
\end{enumerate}
The first construction loses the transverse floor.  The second
preserves the V97 floor exactly, but moves the remaining burden to a
uniform scale-to-scale contraction and a projectively compatible
exact cocycle.

\section{The growing-torus transverse eigenvalue}
\label{sec:s7v98-fourier}

The following calculation is an exact model of the infrared loss in a
large conditional fibre.  Let
\(\mathbb T_\ell^d=(\mathbb Z/\ell\mathbb Z)^d\), \(d\geq2\), with
unit lattice spacing and the counting inner product.  Work one colour
component at a time; the colour multiplicity does not change the
spectrum.  For momentum
\[
 k\in\{0,\ldots,\ell-1\}^d,\qquad
 q_\mu(k):=e^{2\pi i k_\mu/\ell}-1,
                                                        \label{eq:s7v98-q}
\]
the symbols of the coboundaries are
\[
 (\widehat d_0\phi)_\mu=q_\mu\phi,\qquad
 (\widehat d_1A)_{\mu\nu}=q_\mu A_\nu-q_\nu A_\mu.
                                                        \label{eq:s7v98-symbol}
\]

\begin{lemma}[Fourier transverse identity]
\label{lem:s7v98-fourier}
For every \(q,A\in\mathbb C^d\),
\[
 \sum_{\mu<\nu}|q_\mu A_\nu-q_\nu A_\mu|^2
 =
 |q|^2|A|^2-|\langle q,A\rangle|^2.                 \label{eq:s7v98-lagrange}
\]
Hence on the momentumwise transverse subspace
\(\langle q,A\rangle=0\), the symbol of \(d_1^*d_1\) is
\(|q|^2I\).
\end{lemma}

\begin{proof}
Expand the left side.  The diagonal terms contribute
\(\sum_{\mu\neq\nu}|q_\mu|^2|A_\nu|^2\), while the two cross terms
give the negative off-diagonal part of
\(|\langle q,A\rangle|^2\).  Adding its diagonal part yields the
displayed identity.  The last assertion follows from
\(\|\widehat d_1A\|^2=\langle A,\widehat d_1^*\widehat d_1A\rangle\).
\end{proof}

\begin{theorem}[Exact periodic transverse floor]
\label{thm:s7v98-floor}
Remove the \(k=0\) constant one-forms as torons and remove the
longitudinal gauge gradients at every \(k\neq0\).  On the remaining
transverse space,
\[
 \lambda_\ell
 :=
 \lambda_{\min}(d_1^*d_1)
 =
 4\sin^2\!\left(\frac{\pi}{\ell}\right).             \label{eq:s7v98-floor}
\]
In particular, for \(\ell\geq2\),
\[
 \frac{16}{\ell^2}
 \leq\lambda_\ell
 \leq\frac{4\pi^2}{\ell^2}.                          \label{eq:s7v98-floor-bounds}
\]
\end{theorem}

\begin{proof}
By \Cref{lem:s7v98-fourier}, every nonzero momentum contributes the
eigenvalue
\[
 |q(k)|^2
 =4\sum_{\mu=1}^d\sin^2\!\left(\frac{\pi k_\mu}{\ell}\right).
\]
The smallest nonzero value is obtained when one component of \(k\) is
\(1\) or \(\ell-1\) and the rest vanish, proving
\eqref{eq:s7v98-floor}.  For \(0\leq x\leq\pi/2\),
\(2x/\pi\leq\sin x\leq x\).  Substitute \(x=\pi/\ell\)
to obtain \eqref{eq:s7v98-floor-bounds}.
\end{proof}

\begin{corollary}[No uniform inverse row for a one-shot growing fibre]
\label{cor:s7v98-inverse}
Let \(G_\ell\) be the inverse of \(d_1^*d_1\) on the preceding
transverse space.  Then
\[
 \|G_\ell\|_{\rm op}
 =\lambda_\ell^{-1}
 \asymp\ell^2.                                      \label{eq:s7v98-inverse}
\]
Thus neither the V96 uniform inverse hypothesis nor the V97 fixed
floor survives a one-shot limit \(\ell\to\infty\).
\end{corollary}

\begin{proof}
Spectral calculus gives the equality; use
\eqref{eq:s7v98-floor-bounds}.
\end{proof}

\begin{remark}[Boundary conditions]
\label{rem:s7v98-boundary}
Dirichlet and mixed-boundary boxes have the same
\(\ell^{-2}\) scale, with different constants, by the discrete sine
transform or the usual path Poincare inequality.  The periodic
calculation is used because it makes the gauge-gradient and toron
subtractions transparent.
\end{remark}

\section{What the V96 determinant estimate pays on a large fibre}
\label{sec:s7v98-determinant}

\begin{lemma}[Finite-dimensional local-rank trace bound]
\label{lem:s7v98-trace}
Let \(L\geq\lambda I>0\), \(G=L^{-1}\), and let \(B_z,B_w,C_{zw}\)
be matrices whose ranges lie in a common subspace of dimension at
most \(r\).  Then
\[
 \left|
 \frac12\operatorname{Tr}(GC_{zw})
 -\frac12\operatorname{Tr}(GB_wGB_z)
 \right|
 \leq
 \frac r2
 \left(
 \lambda^{-1}\|C_{zw}\|
 +\lambda^{-2}\|B_w\|\|B_z\|
 \right).                                           \label{eq:s7v98-trace}
\]
\end{lemma}

\begin{proof}
If \(T\) has range of dimension at most \(r\), then
\(|\operatorname{Tr}T|\leq r\|T\|\).
Use \(\|G\|\leq\lambda^{-1}\) on each trace and submultiplicativity.
\end{proof}

\begin{corollary}[Optimistic one-shot scaling]
\label{cor:s7v98-one-shot}
Suppose, optimistically, that a growing fibre still has uniformly
bounded derivative support rank and
\[
 \|B_z\|\leq Cg,\qquad \|C_{zw}\|\leq Cg^2.
                                                        \label{eq:s7v98-optimistic}
\]
The direct absolute V96 trace estimate on the periodic
\(\ell\)-block is then only
\[
 \left|\partial_{y_w}\partial_{y_z}
 \frac12\log\det L\right|
 \leq Cg^2(\ell^2+\ell^4).                           \label{eq:s7v98-one-shot}
\]
Its sufficient vanishing condition is \(g\ell^2\to0\).
\end{corollary}

\begin{proof}
Insert \(\lambda_\ell^{-1}=O(\ell^2)\) from
\Cref{thm:s7v98-floor} and
\eqref{eq:s7v98-optimistic} into
\eqref{eq:s7v98-trace}.
\end{proof}

\begin{firewall}[Failure of a majorant is not a divergence theorem]
\label{fire:s7v98-majorant}
For asymptotic freedom,
\(g(a)^2\) is only inverse-logarithmic while a one-shot physical block
has \(\ell(a)\asymp a^{-1}\); therefore the sufficient condition
\(g\ell^2\to0\) fails violently.  This proves that the present
absolute EDLC argument cannot close on such a fibre.  It does not
prove that the exact determinant itself diverges: cancellations or a
different renormalized representation would require a separate
theorem.  Moreover, the stationary response can make \(B_z\) dense,
which worsens rather than improves the displayed absolute estimate.
\end{firewall}

\section{Fixed-factor cofinal factorization}
\label{sec:s7v98-fixed}

Fix an integer \(b\geq2\).  Let the initial lattice spacing be
\(a_0\), and put
\[
 a_j=b^ja_0,\qquad 0\leq j\leq N(a_0),               \label{eq:s7v98-spacings}
\]
where \(N(a_0)\) is chosen so that
\[
 a_*\leq a_{N(a_0)}<ba_*                             \label{eq:s7v98-terminal-window}
\]
for a fixed physical matching scale \(a_*>0\).  Each elementary map
\[
 \pi_j:X_j\longrightarrow X_{j+1}                   \label{eq:s7v98-pi}
\]
eliminates only a \(b\)-by-\(\cdots\)-by-\(b\) cell dictionary in
level-\(j\) lattice units.

\begin{lemma}[Cofinal depth and bounded elementary geometry]
\label{lem:s7v98-depth}
As \(a_0\downarrow0\),
\[
 N(a_0)
 =
 \frac{\log(a_*/a_0)}{\log b}+O(1)
 \longrightarrow\infty.                             \label{eq:s7v98-depth}
\]
Nevertheless every elementary fibre has bounded combinatorial
diameter, rank, boundary incidence, and V97 transverse floor
\(\lambda_*>0\), independently of \(j\) and \(a_0\), provided the
same finite block dictionary is used.
\end{lemma}

\begin{proof}
The first assertion follows by taking logarithms in
\eqref{eq:s7v98-terminal-window}.  The second is
\Cref{cor:s7v97-dictionary}: changing the physical lattice spacing
does not change the normalized finite cochain matrices of one
blocking step.
\end{proof}

\begin{theorem}[Exact pushforward associativity]
\label{thm:s7v98-associativity}
Let \(\mu_0\) be a finite positive lattice gauge measure, and suppose
each \(\pi_j\) is measurable.  With
\[
 \Pi_{0,N}:=\pi_{N-1}\circ\cdots\circ\pi_0,
                                                        \label{eq:s7v98-composite}
\]
one has the exact identity
\[
 (\Pi_{0,N})_\#\mu_0
 =
 (\pi_{N-1})_\#\cdots(\pi_0)_\#\mu_0.                \label{eq:s7v98-pushforward}
\]
If the conditional densities exist, iterated conditional integration
and one-shot integration give the same terminal density almost
everywhere.
\end{theorem}

\begin{proof}
For every bounded measurable \(F\) on \(X_N\), both sides integrate
\(F\) to
\[
 \int_{X_0}
 F(\pi_{N-1}\circ\cdots\circ\pi_0(x))\,d\mu_0(x).
\]
Equality of finite measures follows.  The density statement is the
uniqueness of Radon--Nikodym derivatives, or equivalently repeated
Fubini disintegration.
\end{proof}

\begin{corollary}[A growing one-shot fibre is unnecessary]
\label{cor:s7v98-unnecessary}
The exact terminal marginal at physical scale \(a_*\) can be
constructed through \(N(a_0)\) fixed-factor conditional
pushforwards.  Hence the cofinal growth
\(\ell(a_0)\asymp a_*/a_0\) is carried by the \emph{number} of exact
maps and need not appear in the dimension or diameter of any
elementary eliminated Hessian.
\end{corollary}

\begin{proof}
Use \Cref{lem:s7v98-depth,thm:s7v98-associativity}.
\end{proof}

\begin{proposition}[One-step WTHC and determinant row]
\label{prop:s7v98-one-step}
Assume every elementary map uses the V94 tree--chord disintegration
and all its retained physical backgrounds lie in the V97 flat
weak-field ball.  Then at every weak-coupling step,
\[
 L_j\geq\frac{\lambda_*}{2}I,\qquad
 \|\partial L_j\|\leq Cg_j,\qquad
 \|\partial^2L_j\|\leq Cg_j^2,                       \label{eq:s7v98-step-Hessian}
\]
and
\[
 {\cal L}_{{\rm det},j}\leq C_{\rm det}g_j^2,        \label{eq:s7v98-step-det}
\]
with constants independent of \(j\) and \(a_0\).
\end{proposition}

\begin{proof}
The finite dictionary gives the common floor by
\Cref{cor:s7v97-dictionary}.  Apply
\Cref{thm:s7v97-powers,lem:s7v97-support,
thm:s7v97-EDLCrows} at each scale.
\end{proof}

\begin{firewall}[The determinant is a principal row, not discarded
error]
\label{fire:s7v98-principal}
Equation \eqref{eq:s7v98-step-det} does not authorize deletion of
\(\frac12\log\det L_j\) at each step.  By V95 it is an owned part of
the exact coarse action.  Summing \(Cg_j^2\) over
\(N(a_0)\) steps is neither the exact action nor a valid cutoff-error
estimate.  The correct object is the exact cocycle of
\Cref{thm:s7v98-associativity}, followed by a propagated comparison
of two cofinal trajectories.
\end{firewall}

\section{Propagation rather than raw summation}
\label{sec:s7v98-propagation}

The elementary algebra needed for a cofinal comparison is recorded
next.  Let \(\|\cdot\|_j\) be the complete interaction norm at scale
\(j\).  Suppose a linearized difference \(h_j\) between two
compatible coarse trajectories obeys
\[
 \|h_{j+1}\|_{j+1}
 \leq\kappa_j\|h_j\|_j+e_j.                          \label{eq:s7v98-recursion}
\]

\begin{lemma}[Nonautonomous discrete Duhamel formula]
\label{lem:s7v98-duhamel}
For \(m<n\),
\[
 \|h_n\|_n
 \leq
 \left(\prod_{k=m}^{n-1}\kappa_k\right)\|h_m\|_m
 +
 \sum_{j=m}^{n-1}
 \left(\prod_{k=j+1}^{n-1}\kappa_k\right)e_j,        \label{eq:s7v98-duhamel}
\]
where an empty product equals one.
\end{lemma}

\begin{proof}
Induct on \(n-m\).  The case one is
\eqref{eq:s7v98-recursion}.  Insert the induction hypothesis into the
last application of \eqref{eq:s7v98-recursion}; multiplication by
\(\kappa_{n-1}\) extends every product by one factor and the new
\(e_{n-1}\) has the empty product.
\end{proof}

\begin{corollary}[Uniformly contractive tail]
\label{cor:s7v98-contractive}
If \(0\leq\kappa_j\leq\kappa<1\), then
\[
 \|h_n\|_n
 \leq
 \kappa^{n-m}\|h_m\|_m+
 \sum_{j=m}^{n-1}\kappa^{n-1-j}e_j.                 \label{eq:s7v98-contractive}
\]
If \(e_j\leq e_*\), the second term is at most
\(e_*/(1-\kappa)\).  If the defects occur only at ultraviolet depths
whose distance to the terminal scale tends to infinity, their
terminal contribution tends to zero even when they are merely
bounded.
\end{corollary}

\begin{proof}
Bound every product in \eqref{eq:s7v98-duhamel} by the corresponding
power of \(\kappa\).  For a defect at \(j\), its terminal weight is
\(\kappa^{n-1-j}\); dominated convergence against the geometric
series proves the last assertion.
\end{proof}

\begin{remark}[Common head and ultraviolet tail]
\label{rem:s7v98-head-tail}
When two regulators are aligned by physical scale, their final
finitely many coarse steps form a common head.  Only the deeper
ultraviolet tail changes.  Principal determinant rows in the common
head are not cutoff defects.  What must vanish is the effect of
additional ultraviolet steps after transport through the intervening
exact maps.
\end{remark}

\begin{theorem}[Fixed-factor scale alternative]
\label{thm:s7v98-alternative}
For the V97 Hessian mechanism, the cofinal geometry has the following
rigorous alternative.
\begin{enumerate}[label=\textup{(\alph*)}]
\item A one-shot growing fibre has
      \(\lambda_\ell\asymp\ell^{-2}\) already at the flat Gaussian
      level and does not satisfy uniform EDLC by the V96 absolute
      argument.
\item A fixed-factor exact iteration keeps the elementary transverse
      floor, derivative support, and \(O(g_j^2)\) determinant row
      uniform, regardless of the number of steps.
\item In case \textup{(b)}, cofinal closure is reduced to proving
      projective compatibility and a propagated scale recursion such
      as \eqref{eq:s7v98-recursion}; raw summability of the principal
      determinant norms is neither required nor sufficient.
\end{enumerate}
\end{theorem}

\begin{proof}
Part (a) is
\Cref{thm:s7v98-floor,cor:s7v98-inverse,
cor:s7v98-one-shot}.  Part (b) is
\Cref{lem:s7v98-depth,thm:s7v98-associativity,
prop:s7v98-one-step}.  Part (c) follows from the exact ownership
firewall \Cref{fire:s7v98-principal} and the propagated comparison
\Cref{lem:s7v98-duhamel}.
\end{proof}

\section{Fixed-factor cofinal Hessian card}
\label{sec:s7v98-card}

\begin{definition}[Fixed-factor cofinal Hessian card, FCHC]
\label{def:s7v98-FCHC}
FCHC requires:
\begin{enumerate}[label=\textup{(F\arabic*)},leftmargin=4em]
\item one fixed blocking factor and a finite block dictionary;
\item exact gauge- and reflection-equivariant elementary
      disintegrations;
\item scalewise retention of gauge, toron, cohomological, and other
      slow directions;
\item a V97 physical stationary branch and transverse floor on every
      weak-field step;
\item exact ownership of Gaussian covariance, determinant, Haar, and
      nonlinear Wilson rows;
\item projective alignment of coarse coordinates across regulators;
\item a uniform derivative or difference bound for the exact coarse
      map in the complete interaction norm; and
\item a terminally vanishing propagated ultraviolet-tail defect.
\end{enumerate}
\end{definition}

\begin{proposition}[Rows now proved and rows still open]
\label{prop:s7v98-FCHC}
V94 and V97, together with
\Cref{thm:s7v98-associativity}, prove FCHC
\textup{(F1), (F3)--(F5)} locally in the flat weak-field chart and
the measure-theoretic associativity part of \textup{(F2)}.
FCHC \textup{(F6)--(F8)}, reflection compatibility in every chart,
and continuation beyond the weak-field branch remain open.
\end{proposition}

\begin{proof}
This is a direct dependency audit:
\Cref{rem:s7v97-tree,thm:s7v97-WTHC,
thm:s7v98-associativity,prop:s7v98-one-step,
fire:s7v97-transition,fire:s7v98-principal}.
\end{proof}

\section{Status and next acquisition}
\label{sec:s7v98-status}

\begin{theorem}[What V98 proves]
\label{thm:s7v98-result}
V98 proves the exact \(\ell^{-2}\) collapse of the flat transverse
Wilson floor on a growing periodic fibre and quantifies the consequent
loss in the direct determinant majorant.  It then proves that an exact
fixed-factor cofinal factorization avoids that loss: the growing
physical separation is represented by a growing number of
fixed-dictionary pushforwards, each retaining the uniform V97 Hessian
floor and derivative powers.  The remaining problem is a
scale-propagated cocycle estimate, not a growing-fibre Hessian
estimate.
\end{theorem}

\begin{proof}
Use
\Cref{thm:s7v98-floor,cor:s7v98-one-shot,
cor:s7v98-unnecessary,prop:s7v98-one-step,
thm:s7v98-alternative}.
\end{proof}

\begin{assessment}[Progress at seventh-series V98]
\label{ass:s7v98-status}
The growing-block branch is now rejected for the present absolute
determinant method, while the fixed-factor branch is mathematically
compatible with cofinal refinement.  This removes one apparent
uniformity obstruction but exposes the actual bottleneck: no theorem
yet bounds the derivative of the complete exact Yang--Mills coarse map
by a scale-uniform contraction after the principal Gaussian rows are
included.
\end{assessment}

\begin{decision}[V99 acquisition order]
\label{dec:s7v98-V99}
V99 will formulate the exact nonlinear coarse-map derivative from the
conditional score identity, separate its principal Gaussian
linearization from the Wilson/Haar KP remainder, and derive a
checkable contraction criterion in the complete oscillation norm.
If a marginal or gauge slow direction forces eigenvalue one, it will
be retained explicitly instead of being hidden in a scalar
contraction.
\end{decision}

\begin{firewall}[Claim boundary after seventh-series V98]
\label{fire:s7v98-claim}
V98 does not prove FCHC \textup{(F6)--(F8)}, a uniform contraction of
the physical coarse map, global branch/chart control, NCGC/CFTC along
the full ultraviolet trajectory, terminal entry, the complete
nonlinear coarse-action cocycle, BSCTC through cutoff removal, CSMC,
MIRC, or ARDPC, regulator-uniform SATC or full RL4C, interacting
endpoint continuity, UCM, BCC, pairing/Bogoliubov control, regulator
independence, continuum OS reconstruction, or the full Yang--Mills
mass gap.
\end{firewall}

\paragraph{Parent-version record.}
V98 was formed from
\texttt{Renewal\_Yang\_Mills\_Mass\_Gap\_Research\_Notes\_V97.tex}.
The V97 parent had \(44\,352\) counted source lines,
\(1\,585\,198\) bytes, and SHA--256
\[
 \texttt{F7650DD4FE9380480861AB1AFB8A1E311B39B8396134FBD982A0D9E0662D0C9D}.
\]
The next compulsory companion audit is V100.

% =============================================================================
% END APPENDED SEVENTH-SERIES V98 MATERIAL
% =============================================================================

\clearpage
% =============================================================================
% BEGIN APPENDED SEVENTH-SERIES V99 MATERIAL
% =============================================================================

\part{Seventh-series V99: exact coarse differential, unit directions,
and calibrated contraction}
\label{part:s7v99}

\section{Why a full contraction cannot be assumed}
\label{sec:s7v99-purpose}

V98 replaced a growing one-shot fibre by an exact sequence of
fixed-factor maps.  Its remaining desired estimate was a
scale-uniform contraction of the complete coarse map.  There is an
immediate obstruction: any interaction already measurable in the
retained variable passes through conditional integration exactly.
The derivative of the RG map therefore has a right inverse and an
eigenvalue-one sector.  This note derives the exact differential,
identifies that sector, and states the corrected target: contraction
only after the marginal/relevant retained directions have been
calibrated and removed.

\section{The exact conditional-score differential}
\label{sec:s7v99-score}

Let \(\pi:X\to Y\) be one finite compact block map and let
\(\nu_y\) be a measurable reference disintegration on its fibres.
For a bounded real action \(S\), define
\[
 {\cal R}_\pi(S)(y)
 :=
 -\log Z_S(y),\qquad
 Z_S(y):=\int_{\pi^{-1}(y)}e^{-S(x)}\,d\nu_y(x).
                                                        \label{eq:s7v99-R}
\]
Let \(\mu_{S,y}\) be the normalized conditional probability
\[
 d\mu_{S,y}(x)
 :=
 Z_S(y)^{-1}e^{-S(x)}\,d\nu_y(x).                    \label{eq:s7v99-mu}
\]

\begin{theorem}[Exact first and second coarse differentials]
\label{thm:s7v99-differential}
For bounded directions \(h,k\),
\begin{align}
 D{\cal R}_{\pi,S}[h](y)
 &=
 \mathbb E_{S,y}h,                                   \label{eq:s7v99-first}\\
 D^2{\cal R}_{\pi,S}[h,k](y)
 &=
 -\operatorname{Cov}_{S,y}(h,k).                    \label{eq:s7v99-second}
\end{align}
More generally, the \(n\)-th differential is
\[
 D^n{\cal R}_{\pi,S}[h_1,\ldots,h_n](y)
 =
 (-1)^{n+1}
 \kappa_{S,y}(h_1,\ldots,h_n),                       \label{eq:s7v99-cumulant}
\]
where \(\kappa\) is the joint conditional cumulant.
\end{theorem}

\begin{proof}
Boundedness permits differentiation under the finite integral.
For one parameter,
\[
 {\cal R}_\pi(S+th)(y)-{\cal R}_\pi(S)(y)
 =
 -\log\mathbb E_{S,y}e^{-th}.
                                                        \label{eq:s7v99-mgf}
\]
The logarithm of the moment-generating function is the cumulant
generating function.  Its \(n\)-th derivative at zero is
\((-1)^n\kappa_{S,y}(h,\ldots,h)\); the outer minus sign gives
\eqref{eq:s7v99-cumulant}.  Polarization gives distinct directions.
The cases \(n=1,2\) are
\eqref{eq:s7v99-first}--\eqref{eq:s7v99-second}.
\end{proof}

\begin{corollary}[Conditional score identity]
\label{cor:s7v99-score}
For a bounded observable \(f\),
\[
 D_S\!\left(\mathbb E_{S,y}f\right)[h]
 =
 -\operatorname{Cov}_{S,y}(f,h).                    \label{eq:s7v99-score}
\]
\end{corollary}

\begin{proof}
Differentiate the normalized numerator and denominator in
\eqref{eq:s7v99-mu}; the two denominator terms subtract the product
of the conditional means.
\end{proof}

\begin{remark}[Normalization modulo constants]
\label{rem:s7v99-normalization}
If actions are normalized by
\(\widehat{\cal R}_\pi(S)(y)
={\cal R}_\pi(S)(y)-{\cal R}_\pi(S)(y_0)\), then
\eqref{eq:s7v99-first} becomes
\(\mathbb E_{S,y}h-\mathbb E_{S,y_0}h\).
All conclusions below hold modulo constants.
\end{remark}

\section{The exact unit sector}
\label{sec:s7v99-unit}

Let \(J\) be retained pullback:
\[
 (Jq)(x):=q(\pi(x)).                                  \label{eq:s7v99-J}
\]

\begin{theorem}[Exact retained-pullback identity]
\label{thm:s7v99-pullback}
For every bounded retained interaction \(q\),
\[
 {\cal R}_\pi(S+Jq)
 =
 {\cal R}_\pi(S)+q.                                  \label{eq:s7v99-pullback}
\]
Consequently,
\[
 D{\cal R}_{\pi,S}\,J=I.                             \label{eq:s7v99-right-inverse}
\]
\end{theorem}

\begin{proof}
On the fibre \(\pi^{-1}(y)\), \(Jq=q(y)\) is constant.  Thus
\[
 Z_{S+Jq}(y)=e^{-q(y)}Z_S(y).
\]
Taking minus logarithms proves
\eqref{eq:s7v99-pullback}; differentiation gives
\eqref{eq:s7v99-right-inverse}.
\end{proof}

\begin{corollary}[No strict contraction on the complete action space]
\label{cor:s7v99-no-full}
Suppose the fine and coarse norms satisfy
\(\|Jq\|_{\rm fine}=\|q\|_{\rm coarse}\).  Then
\[
 \|D{\cal R}_{\pi,S}\|_{\rm fine\to coarse}\geq1.
                                                        \label{eq:s7v99-nocontraction}
\]
In particular, no estimate
\(\|D{\cal R}_{\pi,S}\|\leq\kappa<1\) can hold on the
complete interaction space.
\end{corollary}

\begin{proof}
For nonzero \(q\), apply \eqref{eq:s7v99-right-inverse} to \(Jq\):
\[
 \frac{\|D{\cal R}_{\pi,S}Jq\|_{\rm coarse}}
 {\|Jq\|_{\rm fine}}=1.
\]
\end{proof}

\begin{firewall}[The unit sector is physical, not a proof defect]
\label{fire:s7v99-unit}
The identity direction contains the plaquette coupling and every
other interaction written solely in retained links.  Gauge symmetry
does not remove it.  Calling the complete RG map contractive would
silently erase the running coupling.  These directions require a
beta-function or renormalization-condition trajectory, while
contraction is sought only in the complementary irrelevant sector.
\end{firewall}

\section{Exact conditional centering}
\label{sec:s7v99-centering}

Define the conditional expectation operator
\[
 K_Sh(y):=\mathbb E_{S,y}h
                                                        \label{eq:s7v99-K}
\]
and the action-space operators
\[
 P_Sh:=JK_Sh,\qquad Q_Sh:=h-P_Sh.                    \label{eq:s7v99-PQ}
\]

\begin{lemma}[Conditional projection algebra]
\label{lem:s7v99-projection}
One has
\[
 K_SJ=I,\qquad P_S^2=P_S,\qquad K_SQ_S=0.            \label{eq:s7v99-projection}
\]
Thus every perturbation has the unique conditional split
\[
 h=Jm+v,\qquad m=K_Sh,\quad K_Sv=0.                  \label{eq:s7v99-split}
\]
\end{lemma}

\begin{proof}
The first identity holds because \(Jq\) is constant on each fibre.
Then \(P_S^2=JK_SJK_S=JK_S=P_S\) and
\(K_SQ_S=K_S-K_SJK_S=0\).  The displayed split follows, and
uniqueness follows by applying \(K_S\).
\end{proof}

\begin{theorem}[Exact marginal plus centred nonlinear remainder]
\label{thm:s7v99-exact-split}
For the decomposition \eqref{eq:s7v99-split},
\[
 {\cal R}_\pi(S+h)-{\cal R}_\pi(S)
 =
 m-\log\mathbb E_{S,y}e^{-v}.                        \label{eq:s7v99-exact-split}
\]
If
\[
 \delta(y):=
 \operatorname*{ess\,sup}_{\mu_{S,y}}v
 -
 \operatorname*{ess\,inf}_{\mu_{S,y}}v,
                                                        \label{eq:s7v99-delta}
\]
then pointwise in \(y\),
\[
 -\frac{\delta(y)^2}{8}
 \leq
 {\cal R}_\pi(S+h)-{\cal R}_\pi(S)-m
 \leq0.                                               \label{eq:s7v99-hoeffding}
\]
\end{theorem}

\begin{proof}
Since \(h=Jm+v\), factor \(e^{-m(y)}\) out of the conditional
integral to obtain \eqref{eq:s7v99-exact-split}.
The conditional mean of \(v\) is zero.  Jensen gives
\(\log\mathbb E e^{-v}\geq-\mathbb Ev=0\).
Hoeffding's lemma for a real variable of range \(\delta(y)\) gives
\[
 \log\mathbb E e^{-v}\leq\delta(y)^2/8.
\]
Negating proves \eqref{eq:s7v99-hoeffding}.
\end{proof}

\begin{corollary}[Fibre-centred directions are quadratically
irrelevant pointwise]
\label{cor:s7v99-quadratic}
If \(K_Sh=0\) and \(\delta(y)\leq\delta_*\) uniformly, then
\[
 \left|
 {\cal R}_\pi(S+h)(y)-{\cal R}_\pi(S)(y)
 \right|
 \leq\delta_*^2/8.                                   \label{eq:s7v99-quadratic}
\]
\end{corollary}

\begin{proof}
Set \(m=0\) in \eqref{eq:s7v99-hoeffding}.
\end{proof}

\begin{firewall}[Pointwise quadratic size is not yet a KP row]
\label{fire:s7v99-pointwise}
The bound \eqref{eq:s7v99-quadratic} does not control the support
decomposition or mixed oscillation of the output as \(y\) changes.
To enter the V92 recursion, the centred logarithm must still be
expanded by the V93 conditional cumulant tree and regrouped by
retained footprint.  Hoeffding's lemma cannot replace CFTC.
\end{firewall}

\section{The principal Gaussian differential}
\label{sec:s7v99-gaussian}

Consider a positive quadratic action with retained/eliminated block
matrix
\[
 H=
 \begin{pmatrix}
 A&B\\ B^*&C
 \end{pmatrix},
 \qquad C>0.                                         \label{eq:s7v99-H}
\]
Gaussian elimination gives the retained quadratic Hessian and
normalization
\[
 {\mathscr S}(H)=A-BC^{-1}B^*,\qquad
 {\mathscr D}(C)=\frac12\log\det C.                  \label{eq:s7v99-Schur}
\]

\begin{theorem}[Differential of the Gaussian principal rows]
\label{thm:s7v99-Schur}
For
\(\dot H=\left(\begin{smallmatrix}\dot A&\dot B\\
\dot B^*&\dot C\end{smallmatrix}\right)\),
\begin{align}
 D{\mathscr S}_H[\dot H]
 &=
 \dot A-\dot B C^{-1}B^*
 -BC^{-1}\dot B^*
 +BC^{-1}\dot C C^{-1}B^*,                           \label{eq:s7v99-DS}\\
 D{\mathscr D}_C[\dot C]
 &=
 \frac12\operatorname{Tr}(C^{-1}\dot C).             \label{eq:s7v99-DD}
\end{align}
If
\[
 \theta:=\max\{\|BC^{-1}\|,\|C^{-1}B^*\|\},
                                                        \label{eq:s7v99-theta}
\]
then
\[
 \|D{\mathscr S}_H[\dot H]\|
 \leq
 \|\dot A\|+2\theta\|\dot B\|
 +\theta^2\|\dot C\|.                                \label{eq:s7v99-DS-bound}
\]
\end{theorem}

\begin{proof}
Differentiate \(C^{-1}\) using
\(D(C^{-1})[\dot C]=-C^{-1}\dot C C^{-1}\), then apply the product
rule to \eqref{eq:s7v99-Schur}.  The determinant formula is V96's
first derivative.  Submultiplicativity gives
\eqref{eq:s7v99-DS-bound}.
\end{proof}

\begin{corollary}[Gaussian unit direction]
\label{cor:s7v99-gaussian-unit}
If \(\dot B=\dot C=0\), then
\[
 D{\mathscr S}_H[(\dot A,0,0)]=\dot A.               \label{eq:s7v99-gaussian-unit}
\]
Thus the principal Gaussian map exhibits the same retained
eigenvalue-one sector as \Cref{thm:s7v99-pullback}.
\end{corollary}

\begin{proof}
Set \(\dot B=\dot C=0\) in \eqref{eq:s7v99-DS}.
\end{proof}

\begin{remark}[Typed stable candidates]
\label{rem:s7v99-typed}
After \(\dot A\) is assigned to the calibrated retained sector, the
remaining Schur response is bounded by
\(2\theta\|\dot B\|+\theta^2\|\dot C\|\).  This is the quadratic
counterpart of the V92/V93 retained--eliminated path types.  It is
small only if the actual principal Gaussian cross-transfer
\(\theta\) is below the geometry and norm threshold.
\end{remark}

\section{Comparison with the nonlinear Wilson--Haar conditional law}
\label{sec:s7v99-comparison}

\begin{lemma}[Conditional tilt comparison]
\label{lem:s7v99-tilt}
Let \(\gamma_y\) be a reference conditional probability and let
\[
 d\mu_y=
 \frac{e^{-V}}{\mathbb E_{\gamma_y}e^{-V}}\,d\gamma_y.
                                                        \label{eq:s7v99-tilt}
\]
If
\(\operatorname{osc}_{\gamma_y}V\leq\delta\), then every bounded
real \(f\) satisfies
\[
 \left|\mathbb E_{\mu_y}f-\mathbb E_{\gamma_y}f\right|
 \leq
 (e^\delta-1)\operatorname{osc}_{\gamma_y}f.         \label{eq:s7v99-tilt-bound}
\]
\end{lemma}

\begin{proof}
Subtract the essential infimum of \(V\), which leaves both measures
unchanged, so \(0\leq V\leq\delta\).  The normalized density
\(r=e^{-V}/\mathbb E e^{-V}\) obeys
\(e^{-\delta}\leq r\leq e^\delta\), hence
\(|r-1|\leq e^\delta-1\).  For any constant \(c\),
\[
 |\mathbb E_\gamma[(r-1)f]|
 =
 |\mathbb E_\gamma[(r-1)(f-c)]|
 \leq(e^\delta-1)\|f-c\|_\infty.
\]
Choose \(c\) in the essential range of \(f\) and bound
\(\|f-c\|_\infty\) by its oscillation.
\end{proof}

\begin{corollary}[Principal plus nonlinear differential]
\label{cor:s7v99-principal}
Let \(K_{\rm G}\) be conditional expectation under the centred
Gaussian reference and \(K_{\rm WH}\) conditional expectation under
the normalized Wilson--Haar law on the same good fibre.  Under the V95
Radon--Nikodym oscillation bound \(\delta_M\),
\[
 \|(K_{\rm WH}-K_{\rm G})f\|_\infty
 \leq
 (e^{\delta_M}-1)
 \sup_y\operatorname{osc}_{\gamma_y}f.               \label{eq:s7v99-K-comparison}
\]
Thus the exact coarse differential is the owned Gaussian conditional
expectation plus a perturbation tending to zero on the logarithmic
good set.
\end{corollary}

\begin{proof}
Apply \Cref{lem:s7v99-tilt} fibrewise.  V95 gives
\(\delta_M\to0\); changing between its symmetric density-ratio
convention and the oscillation convention only changes the harmless
constant in the exponential.
\end{proof}

\begin{firewall}[Supremum comparison versus complete interaction norm]
\label{fire:s7v99-norm}
Equation \eqref{eq:s7v99-K-comparison} proves a conditional
supremum-norm comparison.  A contraction in the V92 complete
interaction norm additionally requires support-resolved Gaussian
transfer, CFTC regrouping of the nonlinear correction, and uniform
chart/bad-block rows.  Those are separate hypotheses in the theorem
below.
\end{firewall}

\section{The corrected contraction theorem}
\label{sec:s7v99-contraction}

Let \({\mathfrak B}_j\) be the Banach space of complete interactions
at scale \(j\), modulo constants.  Choose bounded projections
\[
 M_j:{\mathfrak B}_j\to{\mathfrak M}_j,\qquad
 Q_j:=I-M_j,                                         \label{eq:s7v99-MQ}
\]
where \({\mathfrak M}_j\) contains the calibrated gauge coupling and
every other relevant or marginal retained row.

\begin{definition}[Calibrated coarse slice]
\label{def:s7v99-slice}
A set \(\Sigma_j(m_j)\subset{\mathfrak B}_j\) is a calibrated slice
if \(M_jS=m_j\) for every \(S\in\Sigma_j(m_j)\).  The exact coarse map
\({\cal T}_j\) respects the calibration if
\[
 {\cal T}_j\bigl(\Sigma_j(m_j)\bigr)
 \subseteq\Sigma_{j+1}(m_{j+1})                     \label{eq:s7v99-slice}
\]
for the prescribed running marginal coordinate \(m_{j+1}\).
\end{definition}

\begin{theorem}[Calibrated stable-sector contraction]
\label{thm:s7v99-contraction}
Let \({\cal D}_j\subset\Sigma_j(m_j)\) be convex.  Assume:
\begin{enumerate}[label=\textup{(\roman*)},leftmargin=3.5em]
\item \({\cal T}_j\) respects the calibration;
\item its principal Gaussian differential \(K_{{\rm G},j}\) obeys
      \[
       \|Q_{j+1}K_{{\rm G},j}Q_j\|
       \leq\kappa_{\rm G}<1;                         \label{eq:s7v99-kappaG}
      \]
\item the complete support-resolved Wilson--Haar correction obeys
      \[
       \sup_{S\in{\cal D}_j}
       \|Q_{j+1}(D{\cal T}_{j,S}-K_{{\rm G},j})Q_j\|
       \leq\eta_j;                                   \label{eq:s7v99-eta}
      \]
\item \(\kappa_{\rm G}+\eta_j\leq\kappa<1\).
\end{enumerate}
Then for \(S,S'\in{\cal D}_j\),
\[
 \|Q_{j+1}({\cal T}_jS-{\cal T}_jS')\|_{j+1}
 \leq
 \kappa\|Q_j(S-S')\|_j.                              \label{eq:s7v99-stable-contraction}
\]
\end{theorem}

\begin{proof}
Because \(S,S'\) lie in the same calibrated slice,
\(M_j(S-S')=0\), so \(S-S'=Q_j(S-S')\).
Integrate the Frechet derivative along
\(S_t=S'+t(S-S')\), which remains in the convex set:
\[
 Q_{j+1}({\cal T}_jS-{\cal T}_jS')
 =
 \int_0^1Q_{j+1}D{\cal T}_{j,S_t}Q_j(S-S')\,dt.
\]
Use \eqref{eq:s7v99-kappaG}--\eqref{eq:s7v99-eta} and
the triangle inequality.
\end{proof}

\begin{corollary}[Cofinal stable tail]
\label{cor:s7v99-cofinal}
If the hypotheses of \Cref{thm:s7v99-contraction} hold uniformly in
the fixed-factor trajectory and the marginal path \(m_j\) is common,
then the stable difference of two aligned regulators obeys the V98
discrete Duhamel bound with \(\kappa_j\leq\kappa<1\).
Any bounded defect inserted at a scale whose distance from the
terminal scale tends to infinity has vanishing terminal stable
effect.
\end{corollary}

\begin{proof}
Apply \eqref{eq:s7v99-stable-contraction} successively, inserting any
additive compatibility defects, and then use
\Cref{lem:s7v98-duhamel,cor:s7v98-contractive}.
\end{proof}

\begin{theorem}[Necessary three-row split for the YM coarse map]
\label{thm:s7v99-three-row}
A valid fixed-factor contraction proof must separately establish:
\begin{enumerate}[label=\textup{(\alph*)}]
\item a \emph{marginal trajectory row}, specifying and controlling
      the running gauge coupling and any other \(M_j\)-coordinates;
\item a \emph{principal Gaussian stable row}, proving
      \eqref{eq:s7v99-kappaG} after all slow Gaussian directions are
      moved into \(M_j\); and
\item a \emph{nonlinear conditional row}, proving
      \eqref{eq:s7v99-eta} from Wilson/Haar good activities, Gaussian
      determinant ownership, bad-block saturation, and CFTC.
\end{enumerate}
No combination of V95--V98 alone proves item \textup{(a)} or the
strict inequality in item \textup{(b)}.
\end{theorem}

\begin{proof}
The necessity of (a) follows from the exact right inverse
\eqref{eq:s7v99-right-inverse}.  The Gaussian decomposition
\Cref{thm:s7v99-Schur,cor:s7v99-gaussian-unit} gives (b) as the
remaining principal linear problem.  The comparison
\Cref{cor:s7v99-principal} and the norm firewall
\Cref{fire:s7v99-norm} identify (c).  The contraction theorem requires
all three.
\end{proof}

\section{Calibrated coarse differential card}
\label{sec:s7v99-card}

\begin{definition}[Calibrated coarse differential card, CCDC]
\label{def:s7v99-CCDC}
CCDC requires:
\begin{enumerate}[label=\textup{(C\arabic*)},leftmargin=4em]
\item the exact conditional-score formulas
      \eqref{eq:s7v99-first}--\eqref{eq:s7v99-second};
\item explicit marginal/relevant projections \(M_j\) containing the
      full pullback unit sector relevant to the physical trajectory;
\item a prescribed and regulator-compatible marginal path \(m_j\);
\item a Gaussian stable complement with
      \(\kappa_{\rm G}<1\);
\item a support-resolved nonlinear comparison
      \(\eta_j\to0\) or \(\sup_j\eta_j<1-\kappa_{\rm G}\);
\item exact ownership of determinant, covariance, Haar, chart, and
      bad-block terms;
\item calibrated-slice invariance; and
\item the propagated cofinal-tail estimate.
\end{enumerate}
\end{definition}

\begin{proposition}[Current CCDC ledger]
\label{prop:s7v99-CCDC}
V99 proves CCDC \textup{(C1)} and the necessity and algebraic form of
\textup{(C2)--(C4), (C7)--(C8)}.  V95--V97 give the local good-chart
smallness mechanism for part of \textup{(C5)--(C6)}.  The actual
Yang--Mills marginal path, a strict principal Gaussian stable
constant, global chart rows, and the normalized bad-block/CFTC
comparison are not yet proved.
\end{proposition}

\begin{proof}
Use
\Cref{thm:s7v99-differential,thm:s7v99-pullback,
thm:s7v99-Schur,thm:s7v99-contraction,
cor:s7v99-principal,thm:s7v99-three-row}.
\end{proof}

\section{Status and next acquisition}
\label{sec:s7v99-status}

\begin{theorem}[What V99 proves]
\label{thm:s7v99-result}
V99 proves the exact first, second, and all-order differentials of one
coarse action as conditional cumulants.  It proves an exact retained
pullback identity and therefore rules out strict contraction of the
complete RG map.  It isolates fibre-centred perturbations, whose
coarse effect is quadratic pointwise, computes the principal Gaussian
Schur/determinant differential, compares it to the Wilson--Haar
conditional differential, and proves the correct calibrated
stable-sector contraction theorem.
\end{theorem}

\begin{proof}
Combine
\Cref{thm:s7v99-differential,thm:s7v99-pullback,
thm:s7v99-exact-split,thm:s7v99-Schur,
cor:s7v99-principal,thm:s7v99-contraction}.
\end{proof}

\begin{assessment}[Progress at seventh-series V99]
\label{ass:s7v99-status}
The V98 request for a uniform complete-map contraction was too strong:
the exact coarse differential contains an identity sector.  The
programme now has the correct noncircular target.  It must construct
the running marginal Yang--Mills trajectory, prove a strict
contraction of the principal Gaussian map only on its calibrated
stable complement, and show that the complete nonlinear conditional
row fits inside the remaining margin.
\end{assessment}

\begin{decision}[V100 acquisition order]
\label{dec:s7v99-V100}
V100 will first perform the compulsory SM/NS companion audit.  It will
then attack the principal Gaussian stable constant in Fourier space:
classify the fixed-factor block symbol into pullback/marginal and
detail subspaces, and determine whether the detail transfer has norm
strictly below one in a scale-compatible energy norm.  Any eigenvalue
one found outside the declared marginal sector will be moved upstream
into the calibration ledger.
\end{decision}

\begin{firewall}[Claim boundary after seventh-series V99]
\label{fire:s7v99-claim}
V99 does not construct the Yang--Mills beta-function trajectory,
prove a strict principal Gaussian detail contraction, prove the
support-resolved nonlinear bound \eqref{eq:s7v99-eta}, close FCHC or
CCDC, establish global branch/chart control, NCGC/CFTC along the full
ultraviolet trajectory, terminal entry, BSCTC through cutoff removal,
CSMC, MIRC, or ARDPC, regulator-uniform SATC or full RL4C, interacting
endpoint continuity, UCM, BCC, pairing/Bogoliubov control, regulator
independence, continuum OS reconstruction, or the full Yang--Mills
mass gap.
\end{firewall}

\paragraph{Parent-version record.}
V99 was formed from
\texttt{Renewal\_Yang\_Mills\_Mass\_Gap\_Research\_Notes\_V98.tex}.
The V98 parent had \(44\,882\) counted source lines,
\(1\,604\,066\) bytes, and SHA--256
\[
 \texttt{72092C7F7697E498DA58DB9970F3BB6C6DFD75B405629A035CBB80ADE06F2582}.
\]
The next compulsory companion audit is V100.

% =============================================================================
% END APPENDED SEVENTH-SERIES V99 MATERIAL
% =============================================================================

\clearpage
% =============================================================================
% BEGIN APPENDED SEVENTH-SERIES V100 MATERIAL
% =============================================================================

\part{Seventh-series V100: companion audit and the Gaussian calibrated
detail split}
\label{part:s7v100}

\section{Compulsory V100 companion audit}
\label{sec:s7v100-audit}

The current SM and Navier--Stokes files in
\texttt{latex\_notes} were inspected read-only.  The SM file advanced
from V68 to V69 during the audit, so the checkpoint was refreshed
before the decision was made:
\begin{align}
 &\texttt{Renewal\_SM\_Einstein\_Continuum\_Research\_Notes\_V69.tex},
 \notag\\[-2pt]
 &\qquad 26\,102\ \hbox{counted source lines},\quad
 940\,539\ \hbox{bytes},                              \label{eq:s7v100-SM-size}\\
 &\qquad
 \texttt{95349B1F2214BEDF6A603CDDE2FE7371CF878DCDE9DEBCC929A021C0C80C5B12},
 \notag\\[4pt]
 &\texttt{Renewal\_NS\_Stability\_Research\_Notes\_V26.tex},
 \notag\\[-2pt]
 &\qquad 5\,319\ \hbox{counted source lines},\quad
 278\,283\ \hbox{bytes},                              \label{eq:s7v100-NS-size}\\
 &\qquad
 \texttt{82C887F8C2C69E4F28FAD7C3DC8DE5B9099CB5A4BF431EBA1FFF3E91935978AD}.
 \notag
\end{align}

\paragraph{SM V63--V69.}
V63 identifies the massless source norm as \(H^{-1}\), so normalized
averaging cannot remove a zero mode.  V64 keeps the infrared block
explicit and integrates only gapped annuli, paying the exact
Feshbach debit.  V65 reproduces the determinant derivative and
inverse-locality ledger used independently in YM V96.  V66 gives a
positive resolvent telescope with the sharp energy-normalized slice
bound
\[
 \left\|A^{1/2}
 \bigl[(A+m_j^2)^{-1}-(A+m_{j-1}^2)^{-1}\bigr]
 A^{1/2}\right\|
 \leq\frac{1-q}{1+q}
                                                        \label{eq:s7v100-SM-shell}
\]
for \(m_j=qm_{j-1}\), but also proves that per-shell smallness can
accumulate logarithmically.  V67 proves that these time-resolvent
differences have a negative spectral residue and therefore fail
Osterwalder--Schrader positivity.  It repairs the free decomposition
by spatial sandwiching with an exact square partition.  V68 proves
that an analytic nonzero filter cannot have an exact spectral gap;
smooth square filters give arbitrary polynomial locality and exact
divided-difference marks instead.  V69 proves that one interacting
shell marginal preserves reflection positivity only from a jointly
reflection-positive fine density; positivity at each fixed value of
a dynamical slow metric is insufficient.

\paragraph{Navier--Stokes V26.}
The NS file is unchanged since the V95 audit.  Its rank-one Oseen
kernel refinement improves a second-derivative constant on an
intermediate radial range but leaves the first-derivative norm
unchanged and does not alter the global-regularity conclusion.  It
has no transferable gauge-block or Gaussian marginal theorem.

\begin{decision}[Transfer from the V100 audit]
\label{dec:s7v100-transfer}
The Yang--Mills programme adopts the following audited rules.
\begin{enumerate}[label=\textup{(\arabic*)}]
\item The SM energy-shell number
      \((1-q)/(1+q)\) is a norm guide, not an admissible
      reflection-positive covariance for YM integration.
\item The actual YM coarse measure continues to be defined by the
      V83 reflection-equivariant exact pushforward; no signed
      time-resolvent slice is substituted for it.
\item Every Gaussian detail split must retain its zero/slow sector
      before inversion and must record accumulation over the number
      of scales.
\item Reflection positivity is checked for the joint fine measure and
      block map.  Conditional positivity at fixed retained background
      is not used as a substitute.
\item The next principal contraction is sought in an
      energy-orthogonal conditional split, followed by canonical
      rescaling; marginal curvature densities remain explicit.
\end{enumerate}
\end{decision}

\begin{proof}
Items (1)--(4) are the direct type-correct consequences of the audited
SM V63--V69 results.  Item (5) addresses the unit-sector obstruction
of \Cref{thm:s7v99-pullback} without changing the exact coarse
measure.  The NS result concerns a different kernel and supplies no
additional row.
\end{proof}

\section{Minimum-energy interpolation}
\label{sec:s7v100-energy}

Let \({\cal H}_f,{\cal H}_c\) be finite-dimensional real Hilbert
spaces, let \(L:{\cal H}_f\to{\cal H}_f\) be positive definite, and
let \(Q:{\cal H}_f\to{\cal H}_c\) be surjective.  The retained
variable is \(y=Qa\).  Put
\[
 C=L^{-1},\qquad
 C_c=QCQ^*,\qquad
 L_c=C_c^{-1},\qquad
 J=CQ^*C_c^{-1}.                                    \label{eq:s7v100-J}
\]

\begin{lemma}[Right inverse and energy orthogonality]
\label{lem:s7v100-right}
The interpolation \(J\) satisfies
\[
 QJ=I_{{\cal H}_c}.                                  \label{eq:s7v100-QJ}
\]
For every \(e\in\ker Q\) and \(y\in{\cal H}_c\),
\[
 \langle e,LJy\rangle_f=0.                           \label{eq:s7v100-orthogonal}
\]
\end{lemma}

\begin{proof}
By \eqref{eq:s7v100-J},
\[
 QJ=QCQ^*(QCQ^*)^{-1}=I.
\]
Also \(LJ=Q^*C_c^{-1}\), so
\[
 \langle e,LJy\rangle_f
 =
 \langle Qe,C_c^{-1}y\rangle_c=0.
\]
\end{proof}

\begin{theorem}[Exact energy decomposition]
\label{thm:s7v100-energy}
Define
\[
 \Pi:=JQ,\qquad D:=I-\Pi.                            \label{eq:s7v100-Pi}
\]
Then \(\Pi^2=\Pi\), \(\ker\Pi=\ker Q\),
\(\operatorname{ran}D=\ker Q\), and
\[
 a=JQa+Da                                             \label{eq:s7v100-field-split}
\]
is \(L\)-orthogonal.  In particular,
\[
 \langle a,La\rangle_f
 =
 \langle Qa,L_cQa\rangle_c+
 \langle Da,LDa\rangle_f.                            \label{eq:s7v100-energy-split}
\]
\end{theorem}

\begin{proof}
\Cref{lem:s7v100-right} gives
\(\Pi^2=JQJQ=JQ=\Pi\), while
\(\Pi a=0\) iff \(Qa=0\) because \(J\) is injective.
Furthermore \(QD=Q-QJQ=0\), and \(De=e\) for \(e\in\ker Q\).
Thus the two terms in \eqref{eq:s7v100-field-split} lie in
\(\operatorname{ran}J\) and \(\ker Q\), respectively, and they are
energy-orthogonal by \eqref{eq:s7v100-orthogonal}.  Finally,
\[
 \langle Jy,LJy\rangle_f
 =
 \langle y,C_c^{-1}y\rangle_c,
\]
which proves \eqref{eq:s7v100-energy-split}.
\end{proof}

\begin{corollary}[Variational characterization]
\label{cor:s7v100-minimum}
For fixed \(y\), \(Jy\) is the unique minimizer of
\(\frac12\langle a,La\rangle\) subject to \(Qa=y\), and the minimum
is \(\frac12\langle y,L_cy\rangle\).
\end{corollary}

\begin{proof}
Every admissible \(a\) is \(Jy+e\) with \(e\in\ker Q\).
Use \eqref{eq:s7v100-energy-split}; strict positivity of \(L\) gives
uniqueness.
\end{proof}

\section{Exact Gaussian disintegration}
\label{sec:s7v100-Gaussian}

Let \(\gamma_L\) be the centred Gaussian probability on
\({\cal H}_f\) with precision \(L\).  Let \(L_E\) be the restriction
of the quadratic form \(L\) to \(E:=\ker Q\), and let
\(\gamma_E\) be the centred Gaussian on \(E\) with precision \(L_E\).

\begin{theorem}[Independent retained and detail Gaussian fields]
\label{thm:s7v100-disintegration}
Under \(\gamma_L\), the random variables
\[
 y=Qa,\qquad e=Da                                   \label{eq:s7v100-ye}
\]
are independent centred Gaussians with precisions \(L_c\) and \(L_E\),
and
\[
 a=Jy+e.                                             \label{eq:s7v100-a}
\]
Equivalently, the conditional law of \(a\) given \(Qa=y\) is the
translate
\[
 a=Jy+e,\qquad e\sim\gamma_E,                        \label{eq:s7v100-conditional}
\]
whose covariance and normalization are independent of \(y\).
\end{theorem}

\begin{proof}
The linear bijection
\({\cal H}_c\oplus E\to{\cal H}_f\),
\((y,e)\mapsto Jy+e\), has a constant nonzero Jacobian.
The exponent splits by \eqref{eq:s7v100-energy-split}:
\[
 \frac12\langle Jy+e,L(Jy+e)\rangle
 =
 \frac12\langle y,L_cy\rangle+
 \frac12\langle e,L_Ee\rangle.
\]
The transformed density therefore factorizes.  Normalization gives
the two Gaussian probabilities and proves independence and the
conditional statement.
\end{proof}

\begin{corollary}[Flat determinant response vanishes]
\label{cor:s7v100-flat-det}
For the exact flat Gaussian split, the eliminated normalization
\(\frac12\log\det L_E\) is independent of \(y\).  Hence its first and
second retained derivatives vanish.
\end{corollary}

\begin{proof}
The detail precision \(L_E\) in
\Cref{thm:s7v100-disintegration} does not depend on the conditioned
value \(y\).
\end{proof}

Let \({\cal J}q=q\circ Q\) be action pullback and define the Gaussian
conditional expectation
\[
 K_{\rm G}h(y)
 :=
 \int_E h(Jy+e)\,d\gamma_E(e).                       \label{eq:s7v100-KG}
\]

\begin{theorem}[Zero transfer on the Gaussian-centred detail space]
\label{thm:s7v100-zero}
Set
\[
 {\cal Q}_{\rm G}:=I-{\cal J}K_{\rm G}.              \label{eq:s7v100-QG}
\]
Then
\[
 K_{\rm G}{\cal J}=I,\qquad
 K_{\rm G}{\cal Q}_{\rm G}=0.                       \label{eq:s7v100-zero}
\]
Thus the principal Gaussian differential has stable norm
\[
 \left\|K_{\rm G}\big|_{\operatorname{ran}{\cal Q}_{\rm G}}\right\|
 =0                                                   \label{eq:s7v100-kappa-zero}
\]
in every norm for which the displayed operators are defined.
\end{theorem}

\begin{proof}
If \(h={\cal J}q\), then
\(h(Jy+e)=q(QJy+Qe)=q(y)\) by
\eqref{eq:s7v100-QJ}.  Hence \(K_{\rm G}{\cal J}=I\), and
\[
 K_{\rm G}{\cal Q}_{\rm G}
 =
 K_{\rm G}-K_{\rm G}{\cal J}K_{\rm G}=0.
\]
The norm assertion is immediate.
\end{proof}

\begin{firewall}[A dynamically centred complement can be nonlocal]
\label{fire:s7v100-centering}
The algebraic constant zero in \eqref{eq:s7v100-kappa-zero} is not by
itself a local RG theorem.  The interpolation
\(J=L^{-1}Q^*(QL^{-1}Q^*)^{-1}\) can be nonlocal for a massless global
operator, and \(K_{\rm G}h\) can have a large retained footprint.  The
V92/V93 norm requires a blockwise conditional construction in which
the eliminated Dirichlet Hessian has the V97 floor and all retained
boundary links are explicit.  Otherwise the definition of the
stable complement merely hides the missing long-range row.
\end{firewall}

\begin{proposition}[Block-local realization]
\label{prop:s7v100-local}
Suppose \(E=\bigoplus_B E_B\), the eliminated quadratic form is
block diagonal on \(E\), each \(E_B\) has diameter at most \(R\) and
floor at least \(\lambda_*\), and every coupling between blocks is a
retained boundary coordinate.  Then the conditional Gaussian in
\Cref{thm:s7v100-disintegration} factorizes over \(B\);
\({\cal Q}_{\rm G}\) is block-local with radius \(R\); and its
determinant normalization is the sum of block-owned constants.
\end{proposition}

\begin{proof}
Under the hypotheses,
\[
 \langle e,L_Ee\rangle=\sum_B\langle e_B,L_{E_B}e_B\rangle.
\]
Therefore \(\gamma_E=\bigotimes_B\gamma_{E_B}\).
Conditional expectations and centering of a block-supported
interaction use only its incident retained boundary variables.
The determinant of a block-diagonal matrix is the product of its
block determinants, so logarithms add.
\end{proof}

\section{Reflection equivariance of the energy split}
\label{sec:s7v100-reflection}

\begin{lemma}[Equivariant interpolation]
\label{lem:s7v100-reflection}
Let \(\Theta_f,\Theta_c\) be orthogonal involutions satisfying
\[
 \Theta_fL=L\Theta_f,\qquad
 Q\Theta_f=\Theta_cQ.                                \label{eq:s7v100-reflection}
\]
Then
\[
 \Theta_fJ=J\Theta_c,\qquad
 \Theta_f\Pi=\Pi\Theta_f,\qquad
 \Theta_fD=D\Theta_f.                                \label{eq:s7v100-reflection-commute}
\]
\end{lemma}

\begin{proof}
From \(Q\Theta_f=\Theta_cQ\), take adjoints to get
\(\Theta_fQ^*=Q^*\Theta_c\).  The first hypothesis implies that
\(\Theta_f\) commutes with \(C=L^{-1}\), and consequently
\(\Theta_c\) commutes with \(C_c=QCQ^*\).  Substitute these
commutations into \eqref{eq:s7v100-J}.  The identities for
\(\Pi=JQ\) and \(D=I-\Pi\) follow.
\end{proof}

\begin{corollary}[Compatibility with the exact RP pushforward]
\label{cor:s7v100-RP}
If the fine Gaussian measure is reflection positive and the retained
map \(Q\) is reflection-equivariant and maps the positive-time
algebra into itself, then its exact retained pushforward is reflection
positive.  The energy split introduces no additional signed
covariance.
\end{corollary}

\begin{proof}
The commutation is \Cref{lem:s7v100-reflection}.  Reflection positivity
of the pushforward is the V83 pullback proof: for a positive-time
retained observable \(F\), the fine observable \(F\circ Q\) is
positive-time and
\[
 \langle\Theta F,F\rangle_{Q_\#\gamma_L}
 =
 \langle\Theta(F\circ Q),F\circ Q\rangle_{\gamma_L}\geq0.
\]
\end{proof}

\begin{remark}[Audited distinction]
\label{rem:s7v100-audited}
The corollary uses one exact pushforward of a positive measure.  It
does not use the difference of two massive time resolvents, which SM
V67 correctly rejects as an OS-positive shell covariance.
\end{remark}

\section{Canonical dilation of retained curvature jets}
\label{sec:s7v100-dilation}

The conditional split removes fibre-centred detail perturbations at
the flat Gaussian differential, but retained interactions still pass
with eigenvalue one before rescaling.  The missing part of the RG map
is the canonical dilation.  In four dimensions set
\[
 x=bx',\qquad
 A'_\mu(x')=bA_\mu(bx'),\qquad b>1.                  \label{eq:s7v100-dilation}
\]
Then
\[
 F'_{\mu\nu}(x')=b^2F_{\mu\nu}(bx'),\qquad
 D'_\mu=bD_\mu.                                      \label{eq:s7v100-Fscale}
\]

\begin{lemma}[Homogeneous coupling multiplier]
\label{lem:s7v100-scaling}
Let \({\cal O}_\Delta(A,D)\) be a homogeneous local scalar of canonical
dimension \(\Delta\), with \([A]=[D]=1\) and \([F]=2\).  Then
\[
 c\int_{\mathbb R^4}{\cal O}_\Delta(A,D)\,d^4x
 =
 c\,b^{4-\Delta}
 \int_{\mathbb R^4}{\cal O}_\Delta(A',D')\,d^4x'.
                                                        \label{eq:s7v100-coupling-scale}
\]
\end{lemma}

\begin{proof}
Equation \eqref{eq:s7v100-dilation} gives
\(A(bx')=b^{-1}A'(x')\) and
\(\partial_x=b^{-1}\partial_{x'}\), so a homogeneous dimension
\(\Delta\) scalar contributes \(b^{-\Delta}\).  The Jacobian is
\(d^4x=b^4d^4x'\).
\end{proof}

\begin{lemma}[Low-dimensional pure-gauge curvature scalars]
\label{lem:s7v100-classification}
Consider the local polynomial algebra generated by traces of
covariant derivatives of \(F\), with all Euclidean indices contracted
using \(\delta\) and \(\epsilon\), modulo total derivatives.  For
\(\mathfrak{su}(N)\), its scalar elements of dimension at most four
are spanned by
\[
 1,\qquad
 {\cal O}_{\rm YM}=\operatorname{Tr}F_{\mu\nu}F_{\mu\nu},
 \qquad
 {\cal O}_\theta=
 \operatorname{Tr}F_{\mu\nu}\widetilde F_{\mu\nu}.   \label{eq:s7v100-low-ops}
\]
There is no nonconstant scalar of dimension one, two, three, or five.
Consequently every other scalar in this algebra has
\(\Delta\geq6\).
\end{lemma}

\begin{proof}
A nonconstant gauge-invariant curvature jet contains \(F\) or a
covariant derivative of \(F\).  A term linear in one such factor has
vanishing colour trace because it is \(\mathfrak{su}(N)\)-valued;
after integration, contracted covariant derivatives are also total
derivatives or have zero trace.  The first nonzero colour trace
therefore contains at least two curvature factors, of total dimension
at least four.  At dimension four, two antisymmetric two-tensors have
only the metric contraction and the epsilon contraction, giving the
two displayed densities.  A dimension-five monomial would contain
two \(F\)'s and one covariant derivative.  It has five Euclidean
indices.  Contractions by tensors of even rank \(\delta\) and
\(\epsilon\) cannot produce a scalar from an odd number of indices.
Terms with more factors have dimension at least six.
\end{proof}

\begin{corollary}[Strict canonical factor off the marginal sector]
\label{cor:s7v100-canonical}
Retain the coefficients of \(1\),
\({\cal O}_{\rm YM}\), and \({\cal O}_\theta\) in the calibrated
sector.  On the coefficient \(\ell^1\)-norm of homogeneous curvature
jets with \(\Delta\geq6\), canonical dilation has operator norm at
most
\[
 \kappa_{\rm can}=b^{-2}<1.                          \label{eq:s7v100-kappa-can}
\]
\end{corollary}

\begin{proof}
By \Cref{lem:s7v100-scaling}, each coefficient is multiplied by
\(b^{4-\Delta}\).  Its absolute value is at most \(b^{-2}\) for
\(\Delta\geq6\).  Sum in the coefficient \(\ell^1\)-norm.
\end{proof}

\begin{firewall}[Curvature jets are not the complete lattice
interaction space]
\label{fire:s7v100-jets}
\Cref{cor:s7v100-canonical} is exact on the declared local polynomial
jet algebra.  It does not prove that every compact-link Wilson
interaction has a uniformly summable curvature-jet expansion, control
extended loops under rescaling, or show that the V92 complete
oscillation norm is equivalent to the graded coefficient norm.  Those
statements are the nonlinear localization bridge still required
before \(b^{-2}\) can be installed as the full
\(\kappa_{\rm G}\) in CCDC.
\end{firewall}

\section{The calibrated free detail theorem}
\label{sec:s7v100-detail}

\begin{theorem}[Free fixed-factor calibrated detail constant]
\label{thm:s7v100-detail}
Consider one reflection-equivariant fixed-factor Gaussian block step
with the following declared decomposition:
\begin{enumerate}[label=\textup{(\roman*)},leftmargin=3.5em]
\item fibre-centred interactions are projected by
      \({\cal Q}_{\rm G}\);
\item retained local interactions are expanded in the curvature-jet
      algebra of \Cref{lem:s7v100-classification};
\item \(1,{\cal O}_{\rm YM},{\cal O}_\theta\) are calibrated and
      removed from the stable norm; and
\item the remaining interaction is canonically dilated by \(b\).
\end{enumerate}
Then the principal free Gaussian derivative has stable constant
\[
 \kappa_{\rm G}\leq b^{-2}<1                         \label{eq:s7v100-kappaG}
\]
on the direct-sum norm consisting of the Gaussian-centred fibre norm
and the \(\ell^1\) norm of retained homogeneous jets of dimension at
least six.
\end{theorem}

\begin{proof}
On the fibre-centred summand, the conditional differential is zero by
\eqref{eq:s7v100-kappa-zero}.  On the retained stable jet summand,
conditional integration is the identity by
\Cref{thm:s7v99-pullback}, after which canonical dilation has norm at
most \(b^{-2}\) by \Cref{cor:s7v100-canonical}.  The direct-sum
operator norm is the maximum of \(0\) and \(b^{-2}\).
\end{proof}

\begin{corollary}[Available nonlinear margin]
\label{cor:s7v100-margin}
Within the declared free detail space, the V99 calibrated contraction
closes whenever the complete support-resolved nonlinear correction
satisfies
\[
 \eta_j<1-b^{-2}.                                    \label{eq:s7v100-margin}
\]
Under a uniform strict version of this inequality, the cofinal stable
tail contracts.
\end{corollary}

\begin{proof}
Insert \eqref{eq:s7v100-kappaG} into
\Cref{thm:s7v99-contraction,cor:s7v99-cofinal}.
\end{proof}

\begin{firewall}[The running coupling remains outside the contraction]
\label{fire:s7v100-coupling}
The coefficient of
\(\operatorname{Tr}F^2\) has canonical multiplier one in four
dimensions.  Its logarithmic quantum running, the topological-sector
coefficient, and wave-function normalization must be fixed by a
regulator-compatible marginal trajectory.  The free detail theorem
does not construct the Yang--Mills beta function or prove asymptotic
freedom nonperturbatively.
\end{firewall}

\section{Gaussian calibrated-detail card}
\label{sec:s7v100-card}

\begin{definition}[Gaussian calibrated-detail card, GCDC]
\label{def:s7v100-GCDC}
GCDC requires:
\begin{enumerate}[label=\textup{(G\arabic*)},leftmargin=4em]
\item a reflection-equivariant retained map \(Q\);
\item prior removal or retention of every zero/slow mode;
\item a block-local minimum-energy interpolation and uniformly gapped
      detail Hessian;
\item exact Gaussian conditional centering;
\item a complete marginal list containing at least
      \(1,\operatorname{Tr}F^2,\operatorname{Tr}F\widetilde F\);
\item a proved localization map from the compact-link complete
      interaction norm to the graded retained stable norm;
\item canonical stable multiplier at most \(b^{-2}\);
\item a support-resolved Wilson/Haar/chart/bad correction
      \(\eta_j<1-b^{-2}\); and
\item a regulator-compatible running marginal path and accumulated
      scale ledger.
\end{enumerate}
\end{definition}

\begin{proposition}[GCDC ledger after V100]
\label{prop:s7v100-GCDC}
V94, V97, V99, and V100 prove the local flat-chart forms of GCDC
\textup{(G1)--(G5), (G7)} and identify the strict margin in
\textup{(G8)}.  GCDC \textup{(G6), (G8), (G9)} remain open for the
complete physical Yang--Mills trajectory; so do global chart and
bad-block compatibility.
\end{proposition}

\begin{proof}
Use
\Cref{prop:s7v100-local,cor:s7v100-RP,
thm:s7v100-zero,lem:s7v100-classification,
thm:s7v100-detail,cor:s7v100-margin,
fire:s7v100-centering,fire:s7v100-jets,
fire:s7v100-coupling}.
\end{proof}

\section{Status at the V100 archive boundary}
\label{sec:s7v100-status}

\begin{theorem}[What V100 proves]
\label{thm:s7v100-result}
After the compulsory companion audit, V100 proves the exact
minimum-energy Gaussian block interpolation, its orthogonal
retained/detail decomposition, the independent conditional Gaussian
law, and zero linear transfer on the conditionally centred detail
space.  It proves reflection equivariance of that split.  After
retaining the dimension-four Yang--Mills and topological densities,
canonical dilation contracts the remaining local curvature jets by
at least \(b^{-2}\).  Hence the declared free calibrated detail space
has the strict principal constant
\(\kappa_{\rm G}\leq b^{-2}\).
\end{theorem}

\begin{proof}
Combine
\Cref{dec:s7v100-transfer,thm:s7v100-energy,
thm:s7v100-disintegration,thm:s7v100-zero,
cor:s7v100-RP,lem:s7v100-classification,
thm:s7v100-detail}.
\end{proof}

\begin{assessment}[Progress at seventh-series V100]
\label{ass:s7v100-status}
The principal Gaussian contraction bottleneck is closed on a precise
local calibrated subspace, not on the complete lattice interaction
space.  Two upstream bridges now dominate: prove that compact-link
coarse interactions split into the calibrated curvature-jet sector
plus a complete norm-controlled remainder, and construct the
regulator-compatible marginal gauge-coupling trajectory.  Only after
those bridges are in place can the V95 nonlinear smallness be compared
with the genuine margin \(1-b^{-2}\).
\end{assessment}

\begin{decision}[Post-archive acquisition order]
\label{dec:s7v100-next}
After freezing V100 and restarting at V1 as required, the next note
will summarize the established chain and attack GCDC (G6): a
gauge-invariant localization/Taylor map from small compact-link
Wilson interactions to curvature jets, with an integral remainder
bounded in the V92 complete oscillation norm.  It will keep
\(\operatorname{Tr}F^2\) and
\(\operatorname{Tr}F\widetilde F\) as explicit marginal owners.
\end{decision}

\begin{firewall}[Claim boundary after seventh-series V100]
\label{fire:s7v100-claim}
V100 does not prove GCDC (G6), (G8), or (G9), construct the
nonperturbative running coupling, prove a global curvature-jet
localization of compact-link interactions, close FCHC or CCDC,
establish global branch/chart control, NCGC/CFTC along the full
ultraviolet trajectory, terminal entry, BSCTC through cutoff removal,
CSMC, MIRC, or ARDPC, regulator-uniform SATC or full RL4C, interacting
endpoint continuity, UCM, BCC, pairing/Bogoliubov control, regulator
independence, continuum OS reconstruction, or the full Yang--Mills
mass gap.
\end{firewall}

\paragraph{Parent-version record.}
V100 was formed from
\texttt{Renewal\_Yang\_Mills\_Mass\_Gap\_Research\_Notes\_V99.tex}.
The V99 parent had \(45\,516\) counted source lines,
\(1\,625\,638\) bytes, and SHA--256
\[
 \texttt{7CA9D687D343116056B357B4B3C8245E4B0FC6CB0F01BFBCA0FFA4C21B0A6421}.
\]
This V100 file is the seventh-series archive boundary.

% =============================================================================
% END APPENDED SEVENTH-SERIES V100 MATERIAL
% =============================================================================

\end{document}
